% Created 2026-09-24 Thu 17:29
% Intended LaTeX compiler: pdflatex
\documentclass[11pt]{amsart}
\usepackage[utf8]{inputenc}
\usepackage[T1]{fontenc}
\usepackage{graphicx}
\usepackage{longtable}
\usepackage{wrapfig}
\usepackage{rotating}
\usepackage[normalem]{ulem}
\usepackage{amsmath}
\usepackage{amssymb}
\usepackage{capt-of}

\usepackage[numbered,framed]{matlab-prettifier}
\usepackage[T1]{fontenc}
\usepackage[utf8]{inputenc}
\usepackage[margin=1.3in]{geometry}
\usepackage{amssymb,amsmath,amsfonts}
\usepackage{amsthm}
\usepackage{microtype}
\usepackage{nicefrac}
\usepackage[colorlinks=true,linkcolor=blue,citecolor=blue,urlcolor=blue]{hyperref}
\usepackage{enumitem}
\usepackage{mathtools}
\usepackage{seqsplit}

\setlength{\emergencystretch}{3em}

%%% Notation
\newcommand{\R}{\mathbb{R}}
\newcommand{\N}{\mathbb{N}}
\newcommand{\fint}{\mathop{\mathchoice%
{\,\rlap{$\displaystyle-$}\!\int}%
{\,\rlap{$\textstyle-$}\!\int}%
{\,\rlap{$\scriptstyle-$}\!\int}%
{\,\rlap{$\scriptscriptstyle-$}\!\int}}\nolimits}
\DeclareMathOperator*{\esssup}{ess\,sup}
\DeclareMathOperator*{\essinf}{ess\,inf}
\DeclareMathOperator{\dv}{div}
\DeclareMathOperator{\supp}{supp}
\DeclareMathOperator{\osc}{osc}
\DeclareMathOperator{\dist}{dist}
\newcommand{\eLpNormText}{\mathrm{eLpNorm}}
% Macros invented by translation agents — collected here so the document compiles
\newcommand{\Q}{\mathbb{Q}}
\newcommand{\Mst}{M_{\mathrm{st}}}
\newcommand{\Eext}{E_{\mathrm{ext}}}
\newcommand{\BMO}{\mathrm{BMO}}
\newcommand{\vol}{\mathrm{vol}}
\newcommand{\lam}{\lambda}
\newcommand{\Cpoinc}{C_{\mathrm{poinc}}}
\newcommand{\CJN}{C_{\mathrm{JN}}}
\newcommand{\loc}{\mathrm{loc}}
\newcommand{\tsupp}{\mathrm{tsupp}}
\newcommand{\unitBallDilate}{\mathrm{unitBallDilate}}
\newcommand{\unitBallShellFormula}{\mathrm{unitBallShellFormula}}
\newcommand{\smoothUnitBallExtensionApprox}{\mathrm{smoothUnitBallExtensionApprox}}
\newcommand{\smoothUnitBallExtensionGrad}{\mathrm{smoothUnitBallExtensionGrad}}
\newcommand{\tsupport}{\mathrm{tsupport}}
\newcommand{\sSup}{\mathrm{sSup}}
\newcommand{\sInf}{\mathrm{sInf}}
\newcommand{\IsSolution}{\mathrm{IsSolution}}
\newcommand{\IsSupersolution}{\mathrm{IsSupersolution}}
\newcommand{\CunitBallExtensionGrad}{C_{\mathrm{ext},\nabla}}
\newcommand{\CunitBallExtensionFun}{C_{\mathrm{ext}}}
\newcommand{\ENNReal}{\mathrm{ENNReal}}
\newcommand{\ENNRealofReal}{\mathrm{ENNReal.ofReal}}
\newcommand{\Fin}{\mathrm{Fin}}
\newcommand{\IsHomogeneousWeakSolution}{\mathrm{IsHomogeneousWeakSolution}}
\newcommand{\fderiv}{\mathrm{fderiv}}
\newcommand{\hw}{\mathrm{hw}}
\newcommand{\hwExt}{\mathrm{hwExt}}
\newcommand{\Lp}[1]{L^{#1}}
\newcommand{\Wkp}[2]{W^{#1,#2}}
\newcommand{\Wop}[1]{W_0^{1,#1}}
\newcommand{\Ball}[2]{B_{#1}(#2)}
\newcommand{\Bz}[1]{B_{#1}}
% Allow line breaks in lean names at underscores; use \small monospace
\newcommand{\leanname}[1]{%
\begingroup
\textnormal{\small\texttt{%
\let\oldunderscore\_%
\def\_{\oldunderscore\discretionary{}{}{}}%
\seqsplit{#1}}}%
\endgroup}
\emergencystretch=3em

%%% Theorem environments
\newtheorem{theorem}{Theorem}[section]
\newtheorem{proposition}[theorem]{Proposition}
\newtheorem{lemma}[theorem]{Lemma}
\newtheorem{corollary}[theorem]{Corollary}

\theoremstyle{definition}
\newtheorem{definition}[theorem]{Definition}

\theoremstyle{remark}
\newtheorem{remark}[theorem]{Remark}
\newtheorem{notation}[theorem]{Notation}



%%% Title
\title{Natural Language Translation of the De~Giorgi--Nash--Moser Lean Formalization}

\author{Machine-generated from Lean 4 source code}

\date{}
%{{{ lean stuff
\usepackage{fontspec}
\setmonofont{JuliaMono}   % broad coverage: blackboard bold, ⊤, ⊆, ∞, ∫, etc.
\usepackage{minted}
\setminted{fontsize=\small, breaklines}

\usepackage[T1]{fontenc}
\usepackage[utf8]{inputenc}
\usepackage{listings}
\usepackage{amssymb}
\usepackage{color}

\definecolor{keywordcolor}{rgb}{0.7, 0.1, 0.1} % red
\definecolor{tacticcolor}{rgb}{0.0, 0.1, 0.6}  % blue
\definecolor{commentcolor}{rgb}{0.4, 0.4, 0.4} % grey
\definecolor{symbolcolor}{rgb}{0.0, 0.1, 0.6}  % blue
\definecolor{sortcolor}{rgb}{0.1, 0.5, 0.1}    % green
\definecolor{attributecolor}{rgb}{0.7, 0.1, 0.1} % red

\def\lstlanguagefiles{lstlean.tex}
\lstset{language=lean}
\usepackage[numbered,framed]{matlab-prettifier}


%}}}

\begin{document}



\tableofcontents

%% ========================================================================
\section{Introduction}\label{sec:intro}
%% ========================================================================

This document is a complete natural-language translation of the Lean~4 formalization of the De~Giorgi--Nash--Moser regularity theory for uniformly elliptic divergence-form equations with bounded measurable coefficients. The formalization is completely sorry-free and built on top of Mathlib. The Lean source is available at \url{https://github.com/scottnarmstrong/DeGiorgi}.

The main results formalized are:
\begin{enumerate}[label=(\roman*)]
  \item \textbf{Local boundedness} of weak subsolutions via De~Giorgi iteration;
  \item \textbf{Weak Harnack inequality} for nonnegative weak supersolutions via Moser iteration;
  \item \textbf{Harnack inequality} for weak solutions;
  \item \textbf{Interior H\"older regularity} for weak solutions via Moser's method.
\end{enumerate}

Throughout, we work in dimension $d \ge 3$ with functions $u \colon \R^d \to \R$. We write $\Bz{r} = \{x \in \R^d : |x| < r\}$ for the open ball of radius~$r$ centered at the origin, and $\Ball{r}{x_0}$ for the ball centered at~$x_0$.

Each definition and theorem includes a reference to its Lean name in \texttt{typewriter font}, enabling cross-referencing with the source code.

\bigskip
\noindent\textbf{Notation.} We adopt the following conventions:
\begin{itemize}
  \item $\nabla u$ denotes the weak gradient of~$u$;
  \item $(u - k)_+ = \max(u - k, 0)$ denotes the positive part;
  \item $\|f\|_{\Lp{p}(\Omega)}$ denotes the $\Lp{p}$ norm over~$\Omega$;
  \item $\lambda, \Lambda > 0$ denote the ellipticity constants (lower and upper bounds);
  \item $|E|$ denotes the Lebesgue measure of a measurable set~$E$;
  \item $\fint_\Omega f\, dx = |\Omega|^{-1}\int_\Omega f\, dx$ denotes the average integral.
\end{itemize}



%% ========================================================================
\section{Sobolev Spaces and Weak Derivatives}\label{sec:sobolev}
%% ========================================================================

% Section 2: Sobolev Spaces and Weak Derivatives
% Translation of:
%   DeGiorgi/SobolevSpace.lean
%   DeGiorgi/SobolevSpace/WeakDerivatives.lean
%   DeGiorgi/SobolevSpace/Witnesses.lean
%   DeGiorgi/SobolevSpace/Approximation.lean
%   DeGiorgi/SobolevSpace/PositivePartPrelude.lean
%   DeGiorgi/WholeSpaceSobolev.lean

Throughout this section we fix a positive integer $d$ and write $E = \R^d$ (in
the Lean formalization, $E$ is the Euclidean space of dimension $d$, with the
canonical orthonormal basis $(e_i)_{i=1}^d$). All open sets $\Omega \subseteq E$
are equipped with Lebesgue measure. For a smooth scalar test function $\varphi
\in C_c^\infty(E)$, $i \in \{1,\ldots,d\}$ and $x \in E$ we abbreviate the
$i$-th partial derivative by $\partial_i \varphi(x) = (D\varphi(x))(e_i)$.

\subsection{Weak partial derivatives and weak gradients}

\begin{definition}[\leanname{DeGiorgi.HasWeakPartialDeriv}]\label{def:HasWeakPartialDeriv}
Let $\Omega \subseteq E$ be a set, $i \in \{1,\ldots,d\}$, and $f, g : E \to \R$.
We say that $g$ is the \emph{weak $i$-th partial derivative of $f$ on $\Omega$},
written $\partial_i^w f = g$ on $\Omega$, if for every test function
$\varphi \in C^\infty(E,\R)$ with compact support and $\supp \varphi \subseteq
\Omega$ we have
\begin{equation*}
  \int_\Omega f(x) \, \partial_i \varphi(x) \, dx
    = - \int_\Omega g(x) \, \varphi(x) \, dx.
\end{equation*}
\end{definition}

\begin{lstlisting}[style=Matlab-editor,basicstyle=\mlttfamily,escapechar=",mlshowsectionrules=true,mathescape=true,morecomment={[s]{\%\{}{\%\}}},language=lean,numbers=none]
/-- `HasWeakPartialDeriv i g f Ω` means that `g` is the weak partial derivative
of `f` with respect to coordinate `i` on `Ω`. -/
def HasWeakPartialDeriv (i : Fin d) (g f : E → ℝ) (Ω : Set E) : Prop :=
  ∀ φ : E → ℝ,
    ContDiff ℝ (⊤ : ℕ∞) φ →
    HasCompactSupport φ →
    tsupport φ ⊆ Ω →
    ∫ x in Ω, f x * (fderiv ℝ φ x) (EuclideanSpace.single i 1) =
      -∫ x in Ω, g x * φ x

\end{lstlisting}
\begin{definition}[\leanname{DeGiorgi.HasWeakGrad}, \leanname{DeGiorgi.HasWeakDiv}]\label{def:HasWeakGrad}
A vector field $G : E \to E$ is a \emph{weak gradient} of $f$ on $\Omega$,
written $\nabla^w f = G$, if each component $G_i$ is the weak $i$-th partial
derivative of $f$ on $\Omega$.

A scalar function $g : E \to \R$ is a \emph{weak divergence} of a vector field
$F : E \to E$ on $\Omega$ if for every $\varphi \in C_c^\infty$ with $\supp
\varphi \subseteq \Omega$
\begin{equation*}
  \int_\Omega \sum_{i=1}^d F_i(x) \, \partial_i \varphi(x) \, dx
    = - \int_\Omega g(x)\, \varphi(x)\, dx.
\end{equation*}
\end{definition}

\begin{lstlisting}[style=Matlab-editor,basicstyle=\mlttfamily,escapechar=",mlshowsectionrules=true,mathescape=true,morecomment={[s]{\%\{}{\%\}}},language=lean,numbers=none]
/-- `HasWeakGrad G f Ω` means that `G` is the weak gradient of `f` on `Ω`. -/
def HasWeakGrad (G : E → E) (f : E → ℝ) (Ω : Set E) : Prop :=
  ∀ i : Fin d, HasWeakPartialDeriv i (fun x => G x i) f Ω

/-- `HasWeakDiv g F Ω` means that `g` is the weak divergence of `F` on `Ω`. -/
def HasWeakDiv (g : E → ℝ) (F : E → E) (Ω : Set E) : Prop :=
  ∀ φ : E → ℝ,
    ContDiff ℝ (⊤ : ℕ∞) φ →
    HasCompactSupport φ →
    tsupport φ ⊆ Ω →
    ∫ x in Ω, (∑ i, F x i * (fderiv ℝ φ x) (EuclideanSpace.single i 1)) =
      -∫ x in Ω, g x * φ x

\end{lstlisting}
\begin{lemma}[Uniqueness; \leanname{DeGiorgi.HasWeakPartialDeriv.ae\_eq}]\label{lem:weak_partial_unique}
Let $\Omega \subseteq E$ be open, $i \in \{1,\ldots,d\}$ and $f, g_1, g_2 :
E \to \R$ with $g_1$ and $g_2$ locally integrable on $\Omega$. If both $g_1$
and $g_2$ are weak $i$-th partial derivatives of $f$ on $\Omega$, then
$g_1 = g_2$ almost everywhere on $\Omega$.
\end{lemma}

\begin{lstlisting}[style=Matlab-editor,basicstyle=\mlttfamily,escapechar=",mlshowsectionrules=true,mathescape=true,morecomment={[s]{\%\{}{\%\}}},language=lean,numbers=none]
omit [NeZero d] in
/-- A weak partial derivative on an open set is unique up to a.e. equality. -/
theorem HasWeakPartialDeriv.ae_eq {Ω : Set E} (hΩ : IsOpen Ω)
    {i : Fin d} {g₁ g₂ f : E → ℝ}
    (h1 : HasWeakPartialDeriv i g₁ f Ω) (h2 : HasWeakPartialDeriv i g₂ f Ω)
    (hg₁ : LocallyIntegrable g₁ (volume.restrict Ω))
    (hg₂ : LocallyIntegrable g₂ (volume.restrict Ω)) :
    g₁ =ᵐ[volume.restrict Ω] g₂ := by
  suffices h : ∀ᵐ x ∂(volume.restrict Ω), (g₁ - g₂) x = 0 by
    filter_upwards [h] with x hx
    simp [sub_eq_zero] at hx
    exact hx
  rw [ae_restrict_iff' hΩ.measurableSet]
  apply IsOpen.ae_eq_zero_of_integral_contDiff_smul_eq_zero hΩ
  · exact locallyIntegrableOn_of_locallyIntegrable_restrict (hg₁.sub hg₂)
  · intro φ hφ hφ_supp hφ_tsupport
    have eq1 := h1 φ hφ hφ_supp hφ_tsupport
    have eq2 := h2 φ hφ hφ_supp hφ_tsupport
    have h_eq : ∫ x in Ω, g₁ x * φ x = ∫ x in Ω, g₂ x * φ x := by
      linarith
    have h_zero : ∀ x, x ∉ Ω → φ x • (g₁ - g₂) x = 0 := by
      intro x hx
      have hφx : φ x = 0 := by
        by_contra h
        exact hx (hφ_tsupport (subset_tsupport _ h))
      simp [hφx]
    rw [← setIntegral_eq_integral_of_forall_compl_eq_zero h_zero]
    simp_rw [Pi.sub_apply, smul_eq_mul, mul_sub]
    rw [integral_sub]
    · simp_rw [mul_comm (φ _)]
      linarith
    ·
      have : Integrable (fun x => φ x • g₁ x) (volume.restrict Ω) :=
        hg₁.integrable_smul_left_of_hasCompactSupport hφ.continuous hφ_supp
      simpa [smul_eq_mul] using this
    ·
      have : Integrable (fun x => φ x • g₂ x) (volume.restrict Ω) :=
        hg₂.integrable_smul_left_of_hasCompactSupport hφ.continuous hφ_supp
      simpa [smul_eq_mul] using this

\end{lstlisting}
\begin{proof}
It suffices to show that $g_1 - g_2 = 0$ a.e.\ on $\Omega$. The function $g_1 -
g_2$ is locally integrable on $\Omega$, and for any $\varphi \in C_c^\infty(E)$
with $\supp \varphi \subseteq \Omega$ we have, subtracting the two test
identities,
\begin{equation*}
  \int_\Omega g_1(x)\,\varphi(x)\,dx = \int_\Omega g_2(x)\,\varphi(x)\,dx,
\end{equation*}
so $\int_\Omega (g_1 - g_2)(x)\,\varphi(x)\,dx = 0$. Since $\varphi$ vanishes
outside its support which lies in $\Omega$, this is the same as integrating over
the whole space $E$. By the standard fact that a locally integrable function
whose integral against every smooth compactly supported test function vanishes
is zero almost everywhere (\leanname{IsOpen.ae\_eq\_zero\_of\_integral\_contDiff\_smul\_eq\_zero}),
we conclude $g_1 = g_2$ a.e.\ on $\Omega$.
\end{proof}

\begin{lemma}[Classical $\Rightarrow$ weak derivative; \leanname{DeGiorgi.HasWeakPartialDeriv.of\_contDiff}]\label{lem:weak_of_contDiff}
Let $\Omega \subseteq E$ be open and $f \in C^1(E,\R)$. Then $\partial_i f$ is
the weak $i$-th partial derivative of $f$ on $\Omega$.
\end{lemma}

\begin{lstlisting}[style=Matlab-editor,basicstyle=\mlttfamily,escapechar=",mlshowsectionrules=true,mathescape=true,morecomment={[s]{\%\{}{\%\}}},language=lean,numbers=none]
omit [NeZero d] in
/-- Classical `C¹` derivatives give weak derivatives on open sets. -/
theorem HasWeakPartialDeriv.of_contDiff {Ω : Set E} (hΩ : IsOpen Ω)
    {i : Fin d} {f : E → ℝ} (hf : ContDiff ℝ 1 f) :
    HasWeakPartialDeriv i (fun x => (fderiv ℝ f x) (EuclideanSpace.single i 1)) f Ω := by
  let _ := hΩ
  intro φ hφ hφ_supp hφ_sub
  let v := EuclideanSpace.single i (1 : ℝ)
  have h_fderiv_supp : tsupport (fun x => (fderiv ℝ φ x) v) ⊆ Ω :=
    (tsupport_fderiv_apply_subset ℝ v).trans hφ_sub
  have hf_diff : Differentiable ℝ f := hf.differentiable one_ne_zero
  have hφ_diff : Differentiable ℝ φ := hφ.differentiable (by simp)
  have hf_cont : Continuous f := hf_diff.continuous
  have hφ_cont : Continuous φ := hφ_diff.continuous
  have hfderiv_φ_cont : Continuous (fun x => (fderiv ℝ φ x) v) :=
    (hφ.continuous_fderiv (by simp)).clm_apply continuous_const
  have hfderiv_f_cont : Continuous (fun x => (fderiv ℝ f x) v) :=
    (hf.continuous_fderiv one_ne_zero).clm_apply continuous_const
  have hφ_fderiv_supp : HasCompactSupport (fun x => (fderiv ℝ φ x) v) :=
    hφ_supp.fderiv_apply (𝕜 := ℝ) v
  rw [setIntegral_eq_integral_of_forall_compl_eq_zero, setIntegral_eq_integral_of_forall_compl_eq_zero]
  ·
    exact integral_mul_fderiv_eq_neg_fderiv_mul_of_integrable
      ((hfderiv_f_cont.mul hφ_cont).integrable_of_hasCompactSupport hφ_supp.mul_left)
      ((hf_cont.mul hfderiv_φ_cont).integrable_of_hasCompactSupport hφ_fderiv_supp.mul_left)
      ((hf_cont.mul hφ_cont).integrable_of_hasCompactSupport hφ_supp.mul_left)
      (fun x _ => hf_diff x) (fun x _ => hφ_diff x)
  ·
    intro x hx
    have : x ∉ tsupport φ := fun h => hx (hφ_sub h)
    simp only [mul_eq_zero]
    right
    by_contra hf'
    exact this (subset_tsupport φ (mem_support.mpr hf'))
  ·
    intro x hx
    have : x ∉ tsupport (fun x => (fderiv ℝ φ x) v) := fun h => hx (h_fderiv_supp h)
    simp only [mul_eq_zero]
    right
    by_contra hf'
    exact this (subset_tsupport _ (mem_support.mpr hf'))

\end{lstlisting}
\begin{proof}
Take any $\varphi \in C_c^\infty(E)$ with $\supp\varphi \subseteq \Omega$. Both
$f \cdot \partial_i \varphi$ and $\partial_i f \cdot \varphi$ vanish outside
$\supp\varphi \subseteq \Omega$, so the corresponding integrals on $\Omega$ and
on $E$ coincide. The classical integration-by-parts formula
(\leanname{integral\_mul\_fderiv\_eq\_neg\_fderiv\_mul\_of\_integrable}) for
differentiable functions with compactly supported product gives
\begin{equation*}
  \int_E f(x)\,\partial_i \varphi(x)\,dx
    = - \int_E (\partial_i f)(x)\,\varphi(x)\,dx,
\end{equation*}
which is the weak-derivative identity on $\Omega$.
\end{proof}

\begin{lemma}[Product rule against a smooth factor; \leanname{DeGiorgi.HasWeakPartialDeriv.mul\_smooth}]\label{lem:weak_product_rule}
Let $\Omega \subseteq E$ be open, $i \in \{1,\ldots,d\}$, and let $f, g, \eta :
E \to \R$ with $\eta \in C^\infty$, $f$ and $g$ locally integrable on $\Omega$,
and $g$ a weak $i$-th partial derivative of $f$ on $\Omega$. Then
\begin{equation*}
  x \mapsto \eta(x)\,g(x) + (\partial_i \eta(x))\,f(x)
\end{equation*}
is the weak $i$-th partial derivative of $\eta f$ on $\Omega$.
\end{lemma}

\begin{lstlisting}[style=Matlab-editor,basicstyle=\mlttfamily,escapechar=",mlshowsectionrules=true,mathescape=true,morecomment={[s]{\%\{}{\%\}}},language=lean,numbers=none]
omit [NeZero d] in
/-- Product rule for weak derivatives against a smooth scalar factor. -/
theorem HasWeakPartialDeriv.mul_smooth {Ω : Set E} (hΩ : IsOpen Ω)
    {i : Fin d} {g f η : E → ℝ}
    (hf : HasWeakPartialDeriv i g f Ω)
    (hη : ContDiff ℝ (⊤ : ℕ∞) η)
    (hf_int : LocallyIntegrable f (volume.restrict Ω))
    (hg_int : LocallyIntegrable g (volume.restrict Ω)) :
    HasWeakPartialDeriv i
      (fun x => η x * g x + (fderiv ℝ η x) (EuclideanSpace.single i 1) * f x)
      (fun x => η x * f x) Ω := by
  let _ := hΩ
  intro φ hφ_smooth hφ_supp hφ_sub
  have hηφ_smooth : ContDiff ℝ (⊤ : ℕ∞) (fun x => η x * φ x) := hη.mul hφ_smooth
  have hηφ_supp : HasCompactSupport (fun x => η x * φ x) := hφ_supp.mul_left
  have hηφ_tsub : tsupport (fun x => η x * φ x) ⊆ Ω :=
    (tsupport_smul_subset_right η φ).trans hφ_sub
  have hη_diff : Differentiable ℝ η := hη.differentiable (by simp)
  have hφ_diff : Differentiable ℝ φ := hφ_smooth.differentiable (by simp)
  let ei : E := EuclideanSpace.single i (1 : ℝ)
  have fderiv_prod : ∀ x : E, (fderiv ℝ (fun y => η y * φ y) x) ei =
      η x * (fderiv ℝ φ x) ei + φ x * (fderiv ℝ η x) ei := by
    intro x
    rw [fderiv_fun_mul hη_diff.differentiableAt hφ_diff.differentiableAt]
    simp [ContinuousLinearMap.add_apply, ContinuousLinearMap.smul_apply, smul_eq_mul]
  have key :
      ∫ x in Ω, f x * (η x * (fderiv ℝ φ x) ei + φ x * (fderiv ℝ η x) ei) =
        -∫ x in Ω, g x * (η x * φ x) := by
    simpa [fderiv_prod, ei] using hf (fun x => η x * φ x) hηφ_smooth hηφ_supp hηφ_tsub
  have hη_deriv_cont : Continuous (fun x => (fderiv ℝ η x) ei) :=
    (hη.continuous_fderiv (by simp)).clm_apply continuous_const
  have hφ_deriv_cont : Continuous (fun x => (fderiv ℝ φ x) ei) :=
    (hφ_smooth.continuous_fderiv (by simp)).clm_apply continuous_const
  have hφ_deriv_supp : HasCompactSupport (fun x => (fderiv ℝ φ x) ei) :=
    hφ_supp.fderiv_apply (𝕜 := ℝ) ei
  have int_f_φ_dη :
      Integrable (fun x => f x * (φ x * (fderiv ℝ η x) ei)) (volume.restrict Ω) := by
    have h := hf_int.integrable_smul_right_of_hasCompactSupport
      (hφ_smooth.continuous.mul hη_deriv_cont) hφ_supp.smul_right
    simpa [smul_eq_mul] using h
  have int_f_η_dφ :
      Integrable (fun x => f x * (η x * (fderiv ℝ φ x) ei)) (volume.restrict Ω) := by
    have h := hf_int.integrable_smul_right_of_hasCompactSupport
      (hη.continuous.mul hφ_deriv_cont) hφ_deriv_supp.mul_left
    simpa [smul_eq_mul] using h
  have int_g_ηφ : Integrable (fun x => g x * (η x * φ x)) (volume.restrict Ω) := by
    have h := hg_int.integrable_smul_right_of_hasCompactSupport
      (hη.continuous.mul hφ_smooth.continuous) hηφ_supp
    simpa [smul_eq_mul] using h
  have key2 :
      ∫ x in Ω, f x * (η x * (fderiv ℝ φ x) ei) + f x * (φ x * (fderiv ℝ η x) ei) =
        -∫ x in Ω, g x * (η x * φ x) := by
    simp_rw [← mul_add] at key ⊢
    exact key
  rw [integral_add int_f_η_dφ int_f_φ_dη] at key2
  show ∫ x in Ω, η x * f x * (fderiv ℝ φ x) ei =
    -∫ x in Ω, (η x * g x + (fderiv ℝ η x) ei * f x) * φ x
  have goal_rhs_conv :
      (fun x => (η x * g x + (fderiv ℝ η x) ei * f x) * φ x) =
        (fun x => g x * (η x * φ x) + f x * (φ x * (fderiv ℝ η x) ei)) := by
    ext x
    ring
  simp_rw [goal_rhs_conv]
  rw [integral_add int_g_ηφ int_f_φ_dη, neg_add]
  have goal_lhs_conv :
      (fun x => η x * f x * (fderiv ℝ φ x) ei) =
        (fun x => f x * (η x * (fderiv ℝ φ x) ei)) := by
    ext x
    ring
  simp_rw [goal_lhs_conv]
  linarith

\end{lstlisting}
\begin{proof}
Fix a test function $\varphi \in C_c^\infty(E)$ with $\supp\varphi \subseteq
\Omega$. Then $\eta \varphi$ is also smooth, compactly supported, with support
contained in $\Omega$. Applying the weak-derivative identity for $f$ to the test
function $\eta \varphi$ and using the smooth product rule
$\partial_i(\eta\varphi) = \eta\,\partial_i\varphi + \varphi\,\partial_i\eta$
gives
\begin{equation*}
  \int_\Omega f(x)\bigl(\eta(x)\,\partial_i\varphi(x) + \varphi(x)\,\partial_i\eta(x)\bigr)\,dx
    = -\int_\Omega g(x)\,\eta(x)\,\varphi(x)\,dx.
\end{equation*}
The local integrability of $f$ and $g$, together with continuity of $\eta$,
$\partial_i\eta$, $\varphi$, $\partial_i\varphi$ and compactness of supports,
ensures that all integrals are absolutely convergent and may be split. Splitting
the left-hand side and rearranging yields
\begin{align*}
  \int_\Omega \eta(x)\,f(x)\,\partial_i\varphi(x)\,dx
  &= -\int_\Omega g(x)\,\eta(x)\,\varphi(x)\,dx
     - \int_\Omega f(x)\,\varphi(x)\,\partial_i\eta(x)\,dx \\
  &= -\int_\Omega \bigl(\eta(x)g(x) + (\partial_i\eta(x))f(x)\bigr)\,\varphi(x)\,dx,
\end{align*}
which is the desired weak-derivative identity for $\eta f$.
\end{proof}

\begin{lemma}[Restriction to a sub-open set; \leanname{DeGiorgi.HasWeakPartialDeriv.restrict}]\label{lem:weak_restrict}
Let $\Omega' \subseteq \Omega \subseteq E$ with $\Omega'$ open. If $g$ is the
weak $i$-th partial derivative of $f$ on $\Omega$, then it is also the weak
$i$-th partial derivative of $f$ on $\Omega'$.
\end{lemma}

\begin{lstlisting}[style=Matlab-editor,basicstyle=\mlttfamily,escapechar=",mlshowsectionrules=true,mathescape=true,morecomment={[s]{\%\{}{\%\}}},language=lean,numbers=none]
omit [NeZero d] in
theorem HasWeakPartialDeriv.restrict {Ω Ω' : Set E}
    (hΩ' : IsOpen Ω') (h_sub : Ω' ⊆ Ω)
    {i : Fin d} {g f : E → ℝ}
    (hf : HasWeakPartialDeriv i g f Ω) :
    HasWeakPartialDeriv i g f Ω' := by
  let _ := hΩ'
  intro φ hφ_smooth hφ_compact hφ_supp
  have hφ_supp_Ω : tsupport φ ⊆ Ω := hφ_supp.trans h_sub
  have key := hf φ hφ_smooth hφ_compact hφ_supp_Ω
  have h1 : ∀ x, x ∉ Ω' → f x * (fderiv ℝ φ x) (EuclideanSpace.single i 1) = 0 := by
    intro x hx
    have hx_notin : x ∉ tsupport φ := fun h => hx (hφ_supp h)
    have hφ_eq : φ =ᶠ[𝓝 x] 0 :=
      (isClosed_tsupport (f := φ)).isOpen_compl.eventually_mem hx_notin |>.mono
        (fun y hy => image_eq_zero_of_notMem_tsupport hy)
    rw [Filter.EventuallyEq.fderiv_eq hφ_eq]
    simp
  have h2 : ∀ x, x ∉ Ω' → g x * φ x = 0 := by
    intro x hx
    simp [image_eq_zero_of_notMem_tsupport (fun h => hx (hφ_supp h))]
  have h3 : ∀ x, x ∉ Ω → f x * (fderiv ℝ φ x) (EuclideanSpace.single i 1) = 0 :=
    fun x hx => h1 x (fun h => hx (h_sub h))
  have h4 : ∀ x, x ∉ Ω → g x * φ x = 0 :=
    fun x hx => h2 x (fun h => hx (h_sub h))
  rw [setIntegral_eq_integral_of_forall_compl_eq_zero h1,
    setIntegral_eq_integral_of_forall_compl_eq_zero h2,
    ← setIntegral_eq_integral_of_forall_compl_eq_zero h3,
    ← setIntegral_eq_integral_of_forall_compl_eq_zero h4,
    key]

\end{lstlisting}
\begin{proof}
Let $\varphi \in C_c^\infty(E)$ with $\supp\varphi \subseteq \Omega'$. Then a
fortiori $\supp\varphi \subseteq \Omega$, so applying the hypothesis on
$\Omega$ gives the integral identity over $\Omega$. Outside $\Omega'$ the test
function $\varphi$ and the field $\partial_i\varphi$ both vanish (the latter
because $\varphi$ is locally constant equal to zero), so the integrals over
$\Omega$ and $\Omega'$ of $f\,\partial_i\varphi$ and $g\,\varphi$ agree. The
result follows.
\end{proof}

\begin{lemma}[$L^p$-closure of weak partial derivatives; \leanname{DeGiorgi.HasWeakPartialDeriv.of\_eLpNormApprox\_p}]\label{lem:weak_lp_closure}
Let $\Omega \subseteq E$ be open and $1 < p < \infty$. Suppose $f, g \in
\Lp{p}(\Omega)$ and there are sequences $(\psi_n)$ and $(g_{\psi_n})$ such that:
\begin{itemize}
  \item $g_{\psi_n}$ is the weak $i$-th partial derivative of $\psi_n$ on
        $\Omega$ for every $n$;
  \item $\psi_n - f$ and $g_{\psi_n} - g$ lie in $\Lp{p}(\Omega)$ and tend to
        $0$ in $\Lp{p}(\Omega)$ as $n \to \infty$.
\end{itemize}
Then $g$ is the weak $i$-th partial derivative of $f$ on $\Omega$.
\end{lemma}

\begin{lstlisting}[style=Matlab-editor,basicstyle=\mlttfamily,escapechar=",mlshowsectionrules=true,mathescape=true,morecomment={[s]{\%\{}{\%\}}},language=lean,numbers=none]
omit [NeZero d] in
/-- `L^p`-closure of weak partial derivatives on an open set, for `1 < p < ∞`. -/
theorem HasWeakPartialDeriv.of_eLpNormApprox_p
    {Ω : Set E} (hΩ : IsOpen Ω)
    {p : ℝ} (hp : 1 < p)
    {i : Fin d} {f g : E → ℝ} {ψ : ℕ → E → ℝ} {gψ : ℕ → E → ℝ}
    (hf_memLp : MemLp f (ENNReal.ofReal p) (volume.restrict Ω))
    (hg_memLp : MemLp g (ENNReal.ofReal p) (volume.restrict Ω))
    (hψ_wd : ∀ n, HasWeakPartialDeriv i (gψ n) (ψ n) Ω)
    (hψ_fun_memLp : ∀ n, MemLp (fun x => ψ n x - f x) (ENNReal.ofReal p) (volume.restrict Ω))
    (hψ_fun :
      Tendsto (fun n => eLpNorm (fun x => ψ n x - f x) (ENNReal.ofReal p) (volume.restrict Ω))
        atTop (nhds 0))
    (hψ_grad_memLp : ∀ n, MemLp (fun x => gψ n x - g x) (ENNReal.ofReal p) (volume.restrict Ω))
    (hψ_grad :
      Tendsto (fun n => eLpNorm (fun x => gψ n x - g x) (ENNReal.ofReal p) (volume.restrict Ω))
        atTop (nhds 0)) :
    HasWeakPartialDeriv i g f Ω := by
  let _ := hΩ
  intro φ hφ hφ_supp hφ_sub
  let μ : Measure E := volume.restrict Ω
  let dφ : E → ℝ := fun x => (fderiv ℝ φ x) (EuclideanSpace.single i 1)
  let q : ℝ := Real.conjExponent p
  have hpqR : p.HolderConjugate q := Real.HolderConjugate.conjExponent hp
  have hpq : p⁻¹ + q⁻¹ = 1 := (Real.holderConjugate_iff.mp hpqR).2
  have hq : 1 < q := (Real.holderConjugate_iff.mp hpqR.symm).1
  letI : (ENNReal.ofReal p).HolderTriple (ENNReal.ofReal q) 1 :=
    ENNReal.HolderConjugate.of_toReal <| by
      simpa [ENNReal.toReal_ofReal (by linarith : 0 ≤ p),
        ENNReal.toReal_ofReal (by linarith : 0 ≤ q)] using hpqR
  have hdφ_memLp : MemLp dφ (ENNReal.ofReal q) μ := by
    have hcont : Continuous dφ := by
      simpa [dφ] using
        ((hφ.continuous_fderiv (by simp : ((⊤ : ℕ∞) : WithTop ℕ∞) ≠ 0)).clm_apply
          continuous_const)
    have hcpt : HasCompactSupport dφ := by
      simpa [dφ] using hφ_supp.fderiv_apply (𝕜 := ℝ) (EuclideanSpace.single i 1)
    exact (hcont.memLp_of_hasCompactSupport hcpt).restrict _
  have hφ_memLp : MemLp φ (ENNReal.ofReal q) μ :=
    (hφ.continuous.memLp_of_hasCompactSupport hφ_supp).restrict _
  have h_fun_int :
      Tendsto (fun n => ∫ x, dφ x * (ψ n x - f x) ∂μ) atTop (nhds 0) :=
    tendsto_integral_mul_of_eLpNorm_tendsto_zero_p hpq hp hq hdφ_memLp hψ_fun_memLp hψ_fun
  have h_grad_int :
      Tendsto (fun n => ∫ x, φ x * (gψ n x - g x) ∂μ) atTop (nhds 0) :=
    tendsto_integral_mul_of_eLpNorm_tendsto_zero_p hpq hp hq hφ_memLp hψ_grad_memLp hψ_grad
  have h_lhs_tendsto :
      Tendsto (fun n => ∫ x in Ω, ψ n x * dφ x)
        atTop (nhds (∫ x in Ω, f x * dφ x)) := by
    have h_eq :
        (fun n => ∫ x in Ω, ψ n x * dφ x) =
          fun n => (∫ x in Ω, f x * dφ x) + ∫ x, dφ x * (ψ n x - f x) ∂μ := by
      funext n
      have hfi : Integrable (fun x => f x * dφ x) μ := by
        simpa [mul_comm] using hdφ_memLp.integrable_mul hf_memLp
      have hdi : Integrable (fun x => dφ x * (ψ n x - f x)) μ :=
        hdφ_memLp.integrable_mul (hψ_fun_memLp n)
      calc
        ∫ x in Ω, ψ n x * dφ x
            = ∫ x, (f x * dφ x) + dφ x * (ψ n x - f x) ∂μ := by
                congr with x
                ring
        _ = (∫ x, f x * dφ x ∂μ) + ∫ x, dφ x * (ψ n x - f x) ∂μ :=
              integral_add hfi hdi
    rw [h_eq]
    simpa [μ] using Tendsto.const_add _ h_fun_int
  have h_rhs_tendsto :
      Tendsto (fun n => -∫ x in Ω, gψ n x * φ x)
        atTop (nhds (-∫ x in Ω, g x * φ x)) := by
    have h_eq :
        (fun n => -∫ x in Ω, gψ n x * φ x) =
          fun n => -(∫ x in Ω, g x * φ x) - ∫ x, φ x * (gψ n x - g x) ∂μ := by
      funext n
      have hgi : Integrable (fun x => g x * φ x) μ := by
        simpa [mul_comm] using hφ_memLp.integrable_mul hg_memLp
      have hdi : Integrable (fun x => φ x * (gψ n x - g x)) μ :=
        hφ_memLp.integrable_mul (hψ_grad_memLp n)
      have : ∫ x, gψ n x * φ x ∂μ =
          (∫ x, g x * φ x ∂μ) + ∫ x, φ x * (gψ n x - g x) ∂μ := by
        calc
          ∫ x, gψ n x * φ x ∂μ
              = ∫ x, (g x * φ x) + φ x * (gψ n x - g x) ∂μ := by
                  congr with x
                  ring
          _ = (∫ x, g x * φ x ∂μ) + ∫ x, φ x * (gψ n x - g x) ∂μ :=
                integral_add hgi hdi
      linarith
    have h_aux :
        Tendsto
          (fun n => -(∫ x, g x * φ x ∂μ) - ∫ x, φ x * (gψ n x - g x) ∂μ)
          atTop (nhds (-(∫ x, g x * φ x ∂μ) - 0)) :=
      Tendsto.const_sub _ h_grad_int
    simpa [μ, h_eq] using h_aux
  have h_eq_n :
      ∀ n,
        ∫ x in Ω, ψ n x * dφ x = -∫ x in Ω, gψ n x * φ x := by
    intro n
    exact hψ_wd n φ hφ hφ_supp hφ_sub
  have h_eq_tendsto :
      Tendsto (fun n => ∫ x in Ω, ψ n x * dφ x)
        atTop (nhds (-∫ x in Ω, g x * φ x)) := by
    simpa [h_eq_n] using h_rhs_tendsto
  exact tendsto_nhds_unique h_lhs_tendsto h_eq_tendsto


\end{lstlisting}
\begin{proof}
Fix $\varphi \in C_c^\infty(E)$ with $\supp\varphi \subseteq \Omega$. Let $q$
denote the H{\"o}lder conjugate of $p$, so that $1 < q < \infty$ and $1/p + 1/q
= 1$. Both $\varphi$ and $\partial_i\varphi$ are continuous, compactly
supported, and hence belong to $\Lp{q}(\Omega)$.

By H{\"o}lder's inequality on $(\Omega, dx)$,
\begin{align*}
  \left| \int_\Omega \partial_i\varphi(x) (\psi_n(x) - f(x))\,dx \right|
  &\le \|\partial_i\varphi\|_{\Lp{q}(\Omega)}\,\|\psi_n - f\|_{\Lp{p}(\Omega)},\\
  \left| \int_\Omega \varphi(x) (g_{\psi_n}(x) - g(x))\,dx \right|
  &\le \|\varphi\|_{\Lp{q}(\Omega)}\,\|g_{\psi_n} - g\|_{\Lp{p}(\Omega)},
\end{align*}
both of which tend to zero by hypothesis. Hence
\begin{align*}
  \int_\Omega \psi_n(x)\,\partial_i\varphi(x)\,dx
    &\;\to\; \int_\Omega f(x)\,\partial_i\varphi(x)\,dx, \\
  \int_\Omega g_{\psi_n}(x)\,\varphi(x)\,dx
    &\;\to\; \int_\Omega g(x)\,\varphi(x)\,dx.
\end{align*}
By assumption, for every $n$
\begin{equation*}
  \int_\Omega \psi_n(x)\,\partial_i\varphi(x)\,dx
    = -\int_\Omega g_{\psi_n}(x)\,\varphi(x)\,dx,
\end{equation*}
and passing to the limit yields
\begin{equation*}
  \int_\Omega f(x)\,\partial_i\varphi(x)\,dx
    = -\int_\Omega g(x)\,\varphi(x)\,dx,
\end{equation*}
which is the weak-derivative identity for $f$ and $g$.
\end{proof}

\subsection{Sobolev membership and weak-gradient witnesses}

\begin{definition}[\leanname{DeGiorgi.MemW1p}]\label{def:MemW1p}
Let $p \in [1,\infty]$, $\Omega \subseteq E$ a set, and $f : E \to \R$. We say
that $f \in \Wkp{1}{p}(\Omega)$ if $f \in \Lp{p}(\Omega)$ and for every $i \in
\{1,\ldots,d\}$ there exists $g_i \in \Lp{p}(\Omega)$ such that $g_i$ is the
weak $i$-th partial derivative of $f$ on $\Omega$.
\end{definition}

\begin{lstlisting}[style=Matlab-editor,basicstyle=\mlttfamily,escapechar=",mlshowsectionrules=true,mathescape=true,morecomment={[s]{\%\{}{\%\}}},language=lean,numbers=none]
/-- `MemW1p p f Ω μ` means `f ∈ W^{1,p}(Ω)` with respect to `μ`. -/
def MemW1p (p : ℝ≥0∞) (f : E → ℝ) (Ω : Set E)
    (μ : Measure E := volume) : Prop :=
  MemLp f p (μ.restrict Ω) ∧
  ∀ i : Fin d, ∃ g : E → ℝ,
    MemLp g p (μ.restrict Ω) ∧ HasWeakPartialDeriv i g f Ω

\end{lstlisting}
\begin{definition}[\leanname{DeGiorgi.MemW1pWitness}]\label{def:MemW1pWitness}
A \emph{$\Wkp{1}{p}$-witness} for $f$ on $\Omega$ is a tuple consisting of a
proof that $f \in \Lp{p}(\Omega)$, a vector field $G : E \to E$ (the weak
gradient field) such that each component $x \mapsto G_i(x)$ lies in
$\Lp{p}(\Omega)$, together with a proof that $G$ is a weak gradient of $f$ on
$\Omega$ (Definition~\ref{def:HasWeakGrad}).
\end{definition}

\begin{lstlisting}[style=Matlab-editor,basicstyle=\mlttfamily,escapechar=",mlshowsectionrules=true,mathescape=true,morecomment={[s]{\%\{}{\%\}}},language=lean,numbers=none]
/-- An explicit `W^{1,p}` witness carrying the weak gradient. -/
structure MemW1pWitness (p : ℝ≥0∞) (f : E → ℝ) (Ω : Set E)
    (μ : Measure E := volume) where
\end{lstlisting}
The Hilbert-space case $p = 2$ is also denoted $H^1(\Omega) := W^{1,2}(\Omega)$
in the formalization (\leanname{DeGiorgi.MemH1}).

\begin{definition}[\leanname{DeGiorgi.MemW01p}]\label{def:MemW01p}
We say that $f \in \Wop{p}(\Omega)$ if $f \in \Wkp{1}{p}(\Omega)$ and there
exist a $\Wkp{1}{p}$-witness $(G)$ for $f$ on $\Omega$ and a sequence
$(\varphi_n)_{n\in\N}$ of smooth compactly supported functions with
$\supp\varphi_n \subseteq \Omega$ such that
\begin{equation*}
  \|\varphi_n - f\|_{\Lp{p}(\Omega)} \to 0
  \quad\text{and}\quad
  \|\partial_i \varphi_n - G_i\|_{\Lp{p}(\Omega)} \to 0
  \;\text{ for every } i.
\end{equation*}
We write $H_0^1(\Omega) := W_0^{1,2}(\Omega)$ (\leanname{DeGiorgi.MemH01}).
\end{definition}

\begin{lstlisting}[style=Matlab-editor,basicstyle=\mlttfamily,escapechar=",mlshowsectionrules=true,mathescape=true,morecomment={[s]{\%\{}{\%\}}},language=lean,numbers=none]
/-- `W₀^{1,p}(Ω)` defined via approximation by smooth compactly supported
functions with convergence in both the function and gradient `L^p` norms. -/
def MemW01p (p : ℝ≥0∞) (f : E → ℝ) (Ω : Set E)
    (μ : Measure E := volume) : Prop :=
  MemW1p p f Ω μ ∧
  ∃ (hw : MemW1pWitness p f Ω μ) (φ : ℕ → E → ℝ),
    (∀ n, ContDiff ℝ (⊤ : ℕ∞) (φ n)) ∧
    (∀ n, HasCompactSupport (φ n)) ∧
    (∀ n, tsupport (φ n) ⊆ Ω) ∧
    Tendsto (fun n => eLpNorm (fun x => φ n x - f x) p (μ.restrict Ω))
      atTop (nhds 0) ∧
    ∀ i : Fin d,
      Tendsto (fun n => eLpNorm
        (fun x => (fderiv ℝ (φ n) x) (EuclideanSpace.single i 1) - hw.weakGrad x i)
        p (μ.restrict Ω))
        atTop (nhds 0)

/-- `H₀¹(Ω) = W₀^{1,2}(Ω)`. -/
abbrev MemH01 (f : E → ℝ) (Ω : Set E) (μ : Measure E := volume) :=
  MemW01p 2 f Ω μ

\end{lstlisting}
\begin{remark}[Existence of an explicit witness; \leanname{DeGiorgi.MemW1p.someWitness}]
Given $f \in W^{1,p}(\Omega)$, axiom of choice produces a witness as in
Definition~\ref{def:MemW1pWitness}: choose for each coordinate $i$ a function
$g_i$ as in Definition~\ref{def:MemW1p}, and assemble them into the vector
field $G(x) := (g_1(x),\ldots,g_d(x))$.
\end{remark}

\begin{lstlisting}[style=Matlab-editor,basicstyle=\mlttfamily,escapechar=",mlshowsectionrules=true,mathescape=true,morecomment={[s]{\%\{}{\%\}}},language=lean,numbers=none]
/-- Choose an explicit weak-gradient witness from `W^{1,p}` membership. -/
noncomputable def MemW1p.someWitness
    {p : ℝ≥0∞} {Ω : Set E} {f : E → ℝ} {μ : Measure E}
    (hf : MemW1p p f Ω μ) :
    MemW1pWitness p f Ω μ := by
  rcases hf with ⟨hf_memLp, hf_grad⟩
  choose g hg_memLp hg_weak using hf_grad
  let G : E → E := fun x => WithLp.toLp 2 fun i => g i x
  exact
    { memLp := hf_memLp
      weakGrad := G
      weakGrad_component_memLp := hg_memLp
      isWeakGrad := hg_weak }

\end{lstlisting}
\begin{lemma}[Forgetting the witness]
\label{lem:forgetting-witness}
(\leanname{DeGiorgi.MemW1pWitness.memW1p},
\leanname{DeGiorgi.MemW01p.memW1p})\\
Every witness gives rise to membership: a $\Wkp{1}{p}$-witness yields
$f \in W^{1,p}(\Omega)$, and an element of $\Wop{p}(\Omega)$ is in particular
in $\Wkp{1}{p}(\Omega)$.
\end{lemma}

\begin{lstlisting}[style=Matlab-editor,basicstyle=\mlttfamily,escapechar=",mlshowsectionrules=true,mathescape=true,morecomment={[s]{\%\{}{\%\}}},language=lean,numbers=none]
omit [NeZero d] in
/-- Forget the explicit gradient from a `W^{1,p}` witness. -/
theorem MemW1pWitness.memW1p
    {p : ℝ≥0∞} {Ω : Set E} {f : E → ℝ} {μ : Measure E}
    (hw : MemW1pWitness p f Ω μ) :
    MemW1p p f Ω μ := by
  refine ⟨hw.memLp, ?_⟩
  intro i
  exact ⟨fun x => hw.weakGrad x i, hw.weakGrad_component_memLp i, hw.isWeakGrad i⟩

omit [NeZero d] in
/-- Forget the zero-trace approximation data from a `W₀^{1,p}` witness. -/
theorem MemW01p.memW1p
    {p : ℝ≥0∞} {Ω : Set E} {f : E → ℝ} {μ : Measure E}
    (hf : MemW01p p f Ω μ) :
    MemW1p p f Ω μ :=
  hf.1

\end{lstlisting}
\begin{proof}
Direct from the definitions: take $g_i := G_i$.
\end{proof}

\begin{lemma}[Sum of witnesses; \leanname{DeGiorgi.MemW1pWitness.add}]\label{lem:witness_add}
For $u, v \in H^1(\Omega)$ with witnesses $G_u$ and $G_v$, the function $u + v$
admits the witness $G_u + G_v$.
\end{lemma}

\begin{lstlisting}[style=Matlab-editor,basicstyle=\mlttfamily,escapechar=",mlshowsectionrules=true,mathescape=true,morecomment={[s]{\%\{}{\%\}}},language=lean,numbers=none]
/-- Add two `W^{1,2}` witnesses on the same domain. -/
noncomputable def MemW1pWitness.add
    {Ω : Set E} {u v : E → ℝ}
    (hu : MemW1pWitness 2 u Ω)
    (hv : MemW1pWitness 2 v Ω) :
    MemW1pWitness 2 (fun x => u x + v x) Ω where
  memLp := hu.memLp.add hv.memLp
  weakGrad := fun x => hu.weakGrad x + hv.weakGrad x
  weakGrad_component_memLp := by
    intro i
    simpa using (hu.weakGrad_component_memLp i).add (hv.weakGrad_component_memLp i)
  isWeakGrad := by
    intro i φ hφ hφ_supp hφ_sub
    let ei : E := EuclideanSpace.single i (1 : ℝ)
    have hu_eq := hu.isWeakGrad i φ hφ hφ_supp hφ_sub
    have hv_eq := hv.isWeakGrad i φ hφ hφ_supp hφ_sub
    have hderiv_cont : Continuous (fun x => (fderiv ℝ φ x) ei) :=
      (hφ.continuous_fderiv (by simp : ((⊤ : ℕ∞) : WithTop ℕ∞) ≠ 0)).clm_apply
        continuous_const
    have hderiv_supp : HasCompactSupport (fun x => (fderiv ℝ φ x) ei) :=
      hφ_supp.fderiv_apply (𝕜 := ℝ) ei
    have hu_int : Integrable (fun x => u x * (fderiv ℝ φ x) ei) (volume.restrict Ω) := by
      have hu_loc : LocallyIntegrable u (volume.restrict Ω) :=
        hu.memLp.locallyIntegrable (by norm_num)
      simpa [smul_eq_mul] using
        hu_loc.integrable_smul_right_of_hasCompactSupport hderiv_cont hderiv_supp
    have hv_int : Integrable (fun x => v x * (fderiv ℝ φ x) ei) (volume.restrict Ω) := by
      have hv_loc : LocallyIntegrable v (volume.restrict Ω) :=
        hv.memLp.locallyIntegrable (by norm_num)
      simpa [smul_eq_mul] using
        hv_loc.integrable_smul_right_of_hasCompactSupport hderiv_cont hderiv_supp
    have hgu_int : Integrable (fun x => hu.weakGrad x i * φ x) (volume.restrict Ω) := by
      have hgu_loc : LocallyIntegrable (fun x => hu.weakGrad x i) (volume.restrict Ω) :=
        (hu.weakGrad_component_memLp i).locallyIntegrable (by norm_num)
      simpa [smul_eq_mul] using
        hgu_loc.integrable_smul_right_of_hasCompactSupport hφ.continuous hφ_supp
    have hgv_int : Integrable (fun x => hv.weakGrad x i * φ x) (volume.restrict Ω) := by
      have hgv_loc : LocallyIntegrable (fun x => hv.weakGrad x i) (volume.restrict Ω) :=
        (hv.weakGrad_component_memLp i).locallyIntegrable (by norm_num)
      simpa [smul_eq_mul] using
        hgv_loc.integrable_smul_right_of_hasCompactSupport hφ.continuous hφ_supp
    have hfun :
        (fun x => (u x + v x) * (fderiv ℝ φ x) ei) =
          (fun x => u x * (fderiv ℝ φ x) ei + v x * (fderiv ℝ φ x) ei) := by
      ext x
      ring
    have hgrad :
        (fun x => (hu.weakGrad x i + hv.weakGrad x i) * φ x) =
          (fun x => hu.weakGrad x i * φ x + hv.weakGrad x i * φ x) := by
      ext x
      ring
    calc
      ∫ x in Ω, (u x + v x) * (fderiv ℝ φ x) ei
          = (-∫ x in Ω, hu.weakGrad x i * φ x) + (-∫ x in Ω, hv.weakGrad x i * φ x) := by
            rw [hfun, integral_add hu_int hv_int, hu_eq, hv_eq]
      _ = -((∫ x in Ω, hu.weakGrad x i * φ x) + (∫ x in Ω, hv.weakGrad x i * φ x)) := by
            ring
      _ = -∫ x in Ω, (hu.weakGrad x i * φ x + hv.weakGrad x i * φ x) := by
            rw [integral_add hgu_int hgv_int]
      _ = -∫ x in Ω, ((fun x => hu.weakGrad x + hv.weakGrad x) x i * φ x) := by
            congr 2
            ext x
            simpa [Pi.add_apply] using (congrFun hgrad x).symm

\end{lstlisting}
\begin{proof}
Membership of $u+v$ in $\Lp{2}(\Omega)$ and of $G_u + G_v$ in $\Lp{2}(\Omega)$
in each coordinate follows from $\Lp{2}$ being a vector space. To verify the
weak-gradient identity, take $\varphi \in C_c^\infty(E)$ with $\supp\varphi
\subseteq \Omega$. The witnesses $G_u$ and $G_v$ give
\begin{align*}
  \int_\Omega u(x)\,\partial_i\varphi(x)\,dx
    &= -\int_\Omega (G_u)_i(x)\,\varphi(x)\,dx, \\
  \int_\Omega v(x)\,\partial_i\varphi(x)\,dx
    &= -\int_\Omega (G_v)_i(x)\,\varphi(x)\,dx.
\end{align*}
Adding these identities, using local integrability and the linearity of the
integral, yields the weak-gradient identity for $u + v$ with weak gradient
$G_u + G_v$.
\end{proof}

\begin{lemma}[Restriction of witnesses; \leanname{DeGiorgi.MemW1pWitness.restrict}]
Let $\Omega' \subseteq \Omega$ with $\Omega'$ open. Any
$\Wkp{1}{p}$-witness on $\Omega$ restricts to a $\Wkp{1}{p}$-witness on
$\Omega'$, with the same gradient field.
\end{lemma}

\begin{lstlisting}[style=Matlab-editor,basicstyle=\mlttfamily,escapechar=",mlshowsectionrules=true,mathescape=true,morecomment={[s]{\%\{}{\%\}}},language=lean,numbers=none]
noncomputable def MemW1pWitness.restrict
    {p : ℝ≥0∞} {Ω Ω' : Set E} {f : E → ℝ}
    (hΩ' : IsOpen Ω')
    (hsub : Ω' ⊆ Ω)
    (hw : MemW1pWitness p f Ω) :
    MemW1pWitness p f Ω' where
  memLp := hw.memLp.mono_measure (Measure.restrict_mono_set volume hsub)
  weakGrad := hw.weakGrad
  weakGrad_component_memLp := by
    intro i
    exact (hw.weakGrad_component_memLp i).mono_measure (Measure.restrict_mono_set volume hsub)
  isWeakGrad := by
    intro i
    exact HasWeakPartialDeriv.restrict hΩ' hsub (hw.isWeakGrad i)

\end{lstlisting}
\begin{proof}
$\Lp{p}$ membership is monotone in the measure (here the volume measure
restricted to a smaller set), and Lemma~\ref{lem:weak_restrict} gives the weak
gradient identity on $\Omega'$.
\end{proof}

\begin{lemma}[Scalar multiple of a witness; \leanname{DeGiorgi.MemW1pWitness.smul}]
For $u \in H^1(\Omega)$ with witness $G$ and $c \in \R$, the function $cu$
admits the witness $cG$.
\end{lemma}

\begin{lstlisting}[style=Matlab-editor,basicstyle=\mlttfamily,escapechar=",mlshowsectionrules=true,mathescape=true,morecomment={[s]{\%\{}{\%\}}},language=lean,numbers=none]
/-- Multiply a `W^{1,2}` witness by a scalar. -/
noncomputable def MemW1pWitness.smul
    {Ω : Set E} {u : E → ℝ}
    (hu : MemW1pWitness 2 u Ω) (c : ℝ) :
    MemW1pWitness 2 (fun x => c * u x) Ω where
  memLp := hu.memLp.const_mul c
  weakGrad := fun x => c • hu.weakGrad x
  weakGrad_component_memLp := by
    intro i
    simpa [Pi.smul_apply, smul_eq_mul] using (hu.weakGrad_component_memLp i).const_mul c
  isWeakGrad := by
    intro i φ hφ hφ_supp hφ_sub
    let ei : E := EuclideanSpace.single i (1 : ℝ)
    have hu_eq := hu.isWeakGrad i φ hφ hφ_supp hφ_sub
    have hderiv_cont : Continuous (fun x => (fderiv ℝ φ x) ei) :=
      (hφ.continuous_fderiv (by simp : ((⊤ : ℕ∞) : WithTop ℕ∞) ≠ 0)).clm_apply
        continuous_const
    have hderiv_supp : HasCompactSupport (fun x => (fderiv ℝ φ x) ei) :=
      hφ_supp.fderiv_apply (𝕜 := ℝ) ei
    have hu_int : Integrable (fun x => u x * (fderiv ℝ φ x) ei) (volume.restrict Ω) := by
      have hu_loc : LocallyIntegrable u (volume.restrict Ω) :=
        hu.memLp.locallyIntegrable (by norm_num)
      simpa [smul_eq_mul] using
        hu_loc.integrable_smul_right_of_hasCompactSupport hderiv_cont hderiv_supp
    have hgu_int : Integrable (fun x => hu.weakGrad x i * φ x) (volume.restrict Ω) := by
      have hgu_loc : LocallyIntegrable (fun x => hu.weakGrad x i) (volume.restrict Ω) :=
        (hu.weakGrad_component_memLp i).locallyIntegrable (by norm_num)
      simpa [smul_eq_mul] using
        hgu_loc.integrable_smul_right_of_hasCompactSupport hφ.continuous hφ_supp
    calc
      ∫ x in Ω, (c * u x) * (fderiv ℝ φ x) ei
          = ∫ x in Ω, c * (u x * (fderiv ℝ φ x) ei) := by
              congr with x
              ring
      _ = c * ∫ x in Ω, u x * (fderiv ℝ φ x) ei := by
            rw [integral_const_mul]
      _ = c * (-∫ x in Ω, hu.weakGrad x i * φ x) := by rw [hu_eq]
      _ = -(c * ∫ x in Ω, hu.weakGrad x i * φ x) := by ring
      _ = -∫ x in Ω, c * (hu.weakGrad x i * φ x) := by
            rw [integral_const_mul]
      _ = -∫ x in Ω, (c * hu.weakGrad x i) * φ x := by
            congr 2
            ext x
            ring
      _ = -∫ x in Ω, ((fun x => c • hu.weakGrad x) x i * φ x) := by
            simp [smul_eq_mul]

\end{lstlisting}
\begin{proof}
The $\Lp{2}$ memberships are immediate. For the test-function identity, fix
$\varphi$ as before; using linearity and the witness for $u$,
\begin{equation*}
  \int_\Omega (cu)(x)\partial_i\varphi(x)\,dx
    = c\int_\Omega u(x)\partial_i\varphi(x)\,dx
    = -c\int_\Omega G_i(x)\varphi(x)\,dx
    = -\int_\Omega (cG)_i(x)\varphi(x)\,dx. \qedhere
\end{equation*}
\end{proof}

\begin{proposition}[Multiplying by a bounded smooth scalar; \leanname{DeGiorgi.MemW1pWitness.mul\_smooth\_bounded\_p}]\label{prop:witness_mul_smooth_bounded}
Let $1 \le p \le \infty$, $\Omega$ open, $u$ a function with $\Wkp{1}{p}$-witness
$G$ on $\Omega$, and $\eta \in C^\infty(E)$ satisfying
$|\eta(x)| \le C_0$ and $\|D\eta(x)\| \le C_1$ for all $x$ and constants
$C_0, C_1 \ge 0$. Then $\eta u$ admits the $\Wkp{1}{p}$-witness whose gradient
field is
\begin{equation*}
  x \mapsto \eta(x)\,G(x) + \bigl(\partial_1\eta(x)\,u(x),\ldots,\partial_d\eta(x)\,u(x)\bigr).
\end{equation*}
\end{proposition}

\begin{lstlisting}[style=Matlab-editor,basicstyle=\mlttfamily,escapechar=",mlshowsectionrules=true,mathescape=true,morecomment={[s]{\%\{}{\%\}}},language=lean,numbers=none]
/-- Multiply a `W^{1,p}` witness by a bounded smooth scalar factor. -/
noncomputable def MemW1pWitness.mul_smooth_bounded_p
    {p : ℝ≥0∞} (hp : 1 ≤ p)
    {Ω : Set E} (hΩ : IsOpen Ω)
    {u η : E → ℝ} (hw : MemW1pWitness p u Ω)
    (hη : ContDiff ℝ (⊤ : ℕ∞) η)
    {C₀ C₁ : ℝ}
    (hC₀ : 0 ≤ C₀) (hC₁ : 0 ≤ C₁)
    (hη_bound : ∀ x, |η x| ≤ C₀)
    (hη_grad_bound : ∀ x, ‖fderiv ℝ η x‖ ≤ C₁) :
    MemW1pWitness p (fun x => η x * u x) Ω where
  memLp := by
    let _ := hC₀
    let _ := hC₁
    refine MemLp.of_le_mul (g := u) (c := C₀) hw.memLp ?_ ?_
    · exact hη.continuous.aestronglyMeasurable.mul hw.memLp.aestronglyMeasurable
    · exact Filter.Eventually.of_forall fun x => by
        calc
          ‖η x * u x‖ = |η x| * ‖u x‖ := by
            rw [norm_mul, Real.norm_eq_abs]
          _ ≤ C₀ * ‖u x‖ := by
            gcongr
            exact hη_bound x
  weakGrad := fun x =>
    η x • hw.weakGrad x +
      (show E from WithLp.toLp 2 fun i => (fderiv ℝ η x) (EuclideanSpace.single i 1) * u x)
  weakGrad_component_memLp := by
    intro i
    have hfirst : MemLp (fun x => η x * hw.weakGrad x i) p (volume.restrict Ω) := by
      refine MemLp.of_le_mul (g := fun x => hw.weakGrad x i) (c := C₀)
        (hw.weakGrad_component_memLp i) ?_ ?_
      · exact hη.continuous.aestronglyMeasurable.mul
          (hw.weakGrad_component_memLp i).aestronglyMeasurable
      · exact Filter.Eventually.of_forall fun x => by
          calc
            ‖η x * hw.weakGrad x i‖ = |η x| * ‖hw.weakGrad x i‖ := by
              rw [norm_mul, Real.norm_eq_abs]
            _ ≤ C₀ * ‖hw.weakGrad x i‖ := by
              gcongr
              exact hη_bound x
    have hderiv_cont : Continuous (fun x => (fderiv ℝ η x) (EuclideanSpace.single i 1)) :=
      (hη.continuous_fderiv (by simp : ((⊤ : ℕ∞) : WithTop ℕ∞) ≠ 0)).clm_apply continuous_const
    have hsecond :
        MemLp (fun x => (fderiv ℝ η x) (EuclideanSpace.single i 1) * u x) p
          (volume.restrict Ω) := by
      refine MemLp.of_le_mul (g := u) (c := C₁) hw.memLp ?_ ?_
      · exact hderiv_cont.aestronglyMeasurable.mul hw.memLp.aestronglyMeasurable
      · exact Filter.Eventually.of_forall fun x => by
          calc
            ‖(fderiv ℝ η x) (EuclideanSpace.single i 1) * u x‖
                = ‖(fderiv ℝ η x) (EuclideanSpace.single i 1)‖ * ‖u x‖ := by
                    rw [norm_mul]
            _ ≤ ‖fderiv ℝ η x‖ * ‖u x‖ := by
                  gcongr
                  simpa using (ContinuousLinearMap.le_opNorm (fderiv ℝ η x)
                    (EuclideanSpace.single i (1 : ℝ)))
            _ ≤ C₁ * ‖u x‖ := by
                  gcongr
                  exact hη_grad_bound x
    simpa using hfirst.add hsecond
  isWeakGrad := by
    intro i
    simpa using HasWeakPartialDeriv.mul_smooth hΩ
      (hf := hw.isWeakGrad i) hη
      (hw.memLp.locallyIntegrable hp)
      ((hw.weakGrad_component_memLp i).locallyIntegrable hp)

\end{lstlisting}
\begin{proof}
The $\Lp{p}$ bound on $\eta u$ follows from the pointwise inequality
$|\eta(x)u(x)| \le C_0 |u(x)|$ and $u \in \Lp{p}(\Omega)$. For each component
of the displayed gradient field:
\begin{itemize}
  \item $|\eta(x)G_i(x)| \le C_0 |G_i(x)|$, and $G_i \in \Lp{p}(\Omega)$;
  \item $|\partial_i\eta(x)\,u(x)| \le \|D\eta(x)\|\,|u(x)| \le C_1 |u(x)|$,
        and $u \in \Lp{p}(\Omega)$.
\end{itemize}
By Minkowski, the sum lies in $\Lp{p}(\Omega)$. The weak-gradient identity is
the product rule (Lemma~\ref{lem:weak_product_rule}) applied componentwise to
the weak partial derivative $G_i$ of $u$ and the smooth factor $\eta$.
\end{proof}

\begin{lemma}[Smooth compactly supported functions are $\Wkp{1}{p}$]
\label{lem:witness_smooth_cpt}
(\leanname{DeGiorgi.MemW1pWitness.of\_contDiff\_hasCompactSupport})\\
Let $f \in C^\infty_c(E,\R)$. Then $f$ admits a canonical $\Wkp{1}{p}$-witness
on the whole space $E$, with weak gradient $\nabla^w f = \nabla f$ given by the
classical gradient field.
\end{lemma}

\begin{lstlisting}[style=Matlab-editor,basicstyle=\mlttfamily,escapechar=",mlshowsectionrules=true,mathescape=true,morecomment={[s]{\%\{}{\%\}}},language=lean,numbers=none]
/-- A smooth compactly supported function on `ℝ^d` carries a canonical
`W^{1,p}` witness on the whole space. -/
noncomputable def MemW1pWitness.of_contDiff_hasCompactSupport
    {p : ℝ≥0∞} {f : E → ℝ}
    (hf : ContDiff ℝ (⊤ : ℕ∞) f) (hf_supp : HasCompactSupport f) :
    MemW1pWitness p f Set.univ where
  memLp := hf.continuous.memLp_of_hasCompactSupport hf_supp
  weakGrad := fun x =>
    WithLp.toLp 2 fun i => (fderiv ℝ f x) (EuclideanSpace.single i 1)
  weakGrad_component_memLp := by
    intro i
    have hderiv_smooth : ContDiff ℝ (⊤ : ℕ∞)
        (fun x => (fderiv ℝ f x) (EuclideanSpace.single i 1)) :=
      (hf.fderiv_right (m := (⊤ : ℕ∞)) (by norm_cast)).clm_apply contDiff_const
    exact hderiv_smooth.continuous.memLp_of_hasCompactSupport (hf_supp.fderiv_apply (𝕜 := ℝ) _)
  isWeakGrad := by
    intro i
    simpa [PiLp.toLp_apply] using
      (HasWeakPartialDeriv.of_contDiff (Ω := Set.univ) (i := i)
        (f := f) isOpen_univ (hf.of_le (by norm_cast)))

\end{lstlisting}
\begin{proof}
Smooth compactly supported functions are integrable to any power, so $f \in
\Lp{p}(E)$ and likewise each classical partial derivative. By
Lemma~\ref{lem:weak_of_contDiff} the classical partial derivatives are weak
partial derivatives.
\end{proof}

\begin{lemma}[Norm of weak gradient lies in $\Lp{p}$;
\leanname{DeGiorgi.MemW1pWitness.weakGrad\_memLp},
\leanname{DeGiorgi.MemW1pWitness.weakGrad\_norm\_memLp}]\label{lem:weak_grad_norm_memLp}
Let $G$ be the weak gradient field of a $\Wkp{1}{p}$-witness on $\Omega$. Then
$G \in \Lp{p}(\Omega; E)$, and in particular $\|G(\cdot)\| \in \Lp{p}(\Omega)$.
\end{lemma}

\begin{lstlisting}[style=Matlab-editor,basicstyle=\mlttfamily,escapechar=",mlshowsectionrules=true,mathescape=true,morecomment={[s]{\%\{}{\%\}}},language=lean,numbers=none]
omit [NeZero d] in
/-- The weak gradient field of a witness lies in `L^p(Ω)` as a Euclidean-space
valued map. -/
theorem MemW1pWitness.weakGrad_memLp
    {p : ℝ≥0∞} {Ω : Set E} {f : E → ℝ} {μ : Measure E}
    (hw : MemW1pWitness p f Ω μ) :
    MemLp hw.weakGrad p (μ.restrict Ω) := by
  refine MemLp.of_eval_piLp ?_
  intro i
  simpa using hw.weakGrad_component_memLp i

/-- The Euclidean norm of the weak gradient lies in `L^p(Ω)`. -/
theorem MemW1pWitness.weakGrad_norm_memLp
    {p : ℝ≥0∞} {Ω : Set E} {f : E → ℝ} {μ : Measure E}
    (hw : MemW1pWitness p f Ω μ) :
    MemLp (fun x => ‖hw.weakGrad x‖) p (μ.restrict Ω) := by
  let _ := (inferInstance : NeZero d)
  exact hw.weakGrad_memLp.norm

\end{lstlisting}
\begin{proof}
Each component $G_i$ lies in $\Lp{p}(\Omega)$ by definition of a witness. For
finite-dimensional Euclidean targets the $\Lp{p}$ membership of a vector field
is equivalent to the $\Lp{p}$ membership of each component
(\leanname{MemLp.of\_eval\_piLp}), so $G \in \Lp{p}(\Omega; E)$. Taking norms
preserves $\Lp{p}$ membership.
\end{proof}

\begin{proposition}[$H_0^1$ is a vector space; \leanname{DeGiorgi.MemW01p.add}, \leanname{DeGiorgi.MemW01p.smul}, \leanname{DeGiorgi.MemW01p.sub}]\label{prop:H01_vector_space}
The space $H_0^1(\Omega)$ is closed under addition, scalar multiplication, and
subtraction.
\end{proposition}

\begin{lstlisting}[style=Matlab-editor,basicstyle=\mlttfamily,escapechar=",mlshowsectionrules=true,mathescape=true,morecomment={[s]{\%\{}{\%\}}},language=lean,numbers=none]
/-- `H₀¹(Ω)` is closed under addition. -/
theorem MemW01p.add
    {Ω : Set E} {u v : E → ℝ}
    (hu : MemW01p 2 u Ω) (hv : MemW01p 2 v Ω) :
    MemW01p 2 (fun x => u x + v x) Ω := by
  let _ := (inferInstance : NeZero d)
  rcases hu with ⟨_, hwu, φu, hφu_smooth, hφu_compact, hφu_sub, hφu_fun, hφu_grad⟩
  rcases hv with ⟨_, hwv, φv, hφv_smooth, hφv_compact, hφv_sub, hφv_fun, hφv_grad⟩
  refine ⟨(hwu.add hwv).memW1p, hwu.add hwv, fun n x => φu n x + φv n x, ?_, ?_, ?_, ?_, ?_⟩
  · intro n
    exact (hφu_smooth n).add (hφv_smooth n)
  · intro n
    exact (hφu_compact n).add (hφv_compact n)
  · intro n
    exact (tsupport_add (φu n) (φv n)).trans <| union_subset (hφu_sub n) (hφv_sub n)
  ·
    have hupper :
        ∀ n,
          eLpNorm (fun x => (φu n x + φv n x) - (u x + v x)) 2 (volume.restrict Ω) ≤
            eLpNorm (fun x => φu n x - u x) 2 (volume.restrict Ω) +
              eLpNorm (fun x => φv n x - v x) 2 (volume.restrict Ω) := by
      intro n
      have hφu_mem : MemLp (φu n) 2 (volume.restrict Ω) :=
        ((hφu_smooth n).continuous.memLp_of_hasCompactSupport (hφu_compact n)).restrict Ω
      have hφv_mem : MemLp (φv n) 2 (volume.restrict Ω) :=
        ((hφv_smooth n).continuous.memLp_of_hasCompactSupport (hφv_compact n)).restrict Ω
      have hdu_mem : MemLp (fun x => φu n x - u x) 2 (volume.restrict Ω) :=
        hφu_mem.sub hwu.memLp
      have hdv_mem : MemLp (fun x => φv n x - v x) 2 (volume.restrict Ω) :=
        hφv_mem.sub hwv.memLp
      have hEq :
          (fun x => (φu n x + φv n x) - (u x + v x)) =
            (fun x => (φu n x - u x) + (φv n x - v x)) := by
        ext x
        ring
      rw [hEq]
      exact eLpNorm_add_le hdu_mem.aestronglyMeasurable hdv_mem.aestronglyMeasurable (by norm_num)
    have hsum :
        Tendsto
          (fun n =>
            eLpNorm (fun x => φu n x - u x) 2 (volume.restrict Ω) +
              eLpNorm (fun x => φv n x - v x) 2 (volume.restrict Ω))
          atTop (nhds (0 + 0)) :=
      hφu_fun.add hφv_fun
    refine tendsto_of_tendsto_of_tendsto_of_le_of_le tendsto_const_nhds ?_ (fun n => zero_le _) hupper
    simpa using hsum
  · intro i
    have hupper :
        ∀ n,
          eLpNorm
              (fun x =>
                (fderiv ℝ (fun y => φu n y + φv n y) x) (EuclideanSpace.single i 1) -
                  (hwu.add hwv).weakGrad x i)
              2 (volume.restrict Ω) ≤
            eLpNorm
                (fun x => (fderiv ℝ (φu n) x) (EuclideanSpace.single i 1) - hwu.weakGrad x i)
                2 (volume.restrict Ω) +
              eLpNorm
                (fun x => (fderiv ℝ (φv n) x) (EuclideanSpace.single i 1) - hwv.weakGrad x i)
                2 (volume.restrict Ω) := by
      intro n
      let ei : E := EuclideanSpace.single i (1 : ℝ)
      have hderiv_u_smooth : ContDiff ℝ (⊤ : ℕ∞)
          (fun x => (fderiv ℝ (φu n) x) ei) :=
        ((hφu_smooth n).fderiv_right (m := (⊤ : ℕ∞)) (by norm_cast)).clm_apply contDiff_const
      have hderiv_v_smooth : ContDiff ℝ (⊤ : ℕ∞)
          (fun x => (fderiv ℝ (φv n) x) ei) :=
        ((hφv_smooth n).fderiv_right (m := (⊤ : ℕ∞)) (by norm_cast)).clm_apply contDiff_const
      have hderiv_u_mem : MemLp (fun x => (fderiv ℝ (φu n) x) ei) 2 (volume.restrict Ω) :=
        (hderiv_u_smooth.continuous.memLp_of_hasCompactSupport ((hφu_compact n).fderiv_apply (𝕜 := ℝ) ei)).restrict Ω
      have hderiv_v_mem : MemLp (fun x => (fderiv ℝ (φv n) x) ei) 2 (volume.restrict Ω) :=
        (hderiv_v_smooth.continuous.memLp_of_hasCompactSupport ((hφv_compact n).fderiv_apply (𝕜 := ℝ) ei)).restrict Ω
      have hdu_mem :
          MemLp (fun x => (fderiv ℝ (φu n) x) ei - hwu.weakGrad x i) 2 (volume.restrict Ω) :=
        hderiv_u_mem.sub (hwu.weakGrad_component_memLp i)
      have hdv_mem :
          MemLp (fun x => (fderiv ℝ (φv n) x) ei - hwv.weakGrad x i) 2 (volume.restrict Ω) :=
        hderiv_v_mem.sub (hwv.weakGrad_component_memLp i)
      have hEq :
          (fun x =>
            (fderiv ℝ (fun y => φu n y + φv n y) x) ei - (hwu.add hwv).weakGrad x i) =
            (fun x =>
              ((fderiv ℝ (φu n) x) ei - hwu.weakGrad x i) +
                ((fderiv ℝ (φv n) x) ei - hwv.weakGrad x i)) := by
        ext x
        have hfd :
            fderiv ℝ (fun y => φu n y + φv n y) x =
              fderiv ℝ (φu n) x + fderiv ℝ (φv n) x := by
          have hu_diff : DifferentiableAt ℝ (φu n) x :=
            ((hφu_smooth n).contDiffAt).differentiableAt (by norm_num)
          have hv_diff : DifferentiableAt ℝ (φv n) x :=
            ((hφv_smooth n).contDiffAt).differentiableAt (by norm_num)
          exact fderiv_add hu_diff hv_diff
        simp [ei, MemW1pWitness.add, hfd]
        ring
      rw [hEq]
      exact eLpNorm_add_le hdu_mem.aestronglyMeasurable hdv_mem.aestronglyMeasurable (by norm_num)
    have hsum :
        Tendsto
          (fun n =>
            eLpNorm
                (fun x => (fderiv ℝ (φu n) x) (EuclideanSpace.single i 1) - hwu.weakGrad x i)
                2 (volume.restrict Ω) +
              eLpNorm
                (fun x => (fderiv ℝ (φv n) x) (EuclideanSpace.single i 1) - hwv.weakGrad x i)
                2 (volume.restrict Ω))
          atTop (nhds (0 + 0)) :=
      (hφu_grad i).add (hφv_grad i)
    refine tendsto_of_tendsto_of_tendsto_of_le_of_le tendsto_const_nhds ?_ (fun n => zero_le _) hupper
    simpa using hsum

/-- `H₀¹(Ω)` is closed under scalar multiplication. -/
theorem MemW01p.smul
    {Ω : Set E} {u : E → ℝ} (c : ℝ)
    (hu : MemW01p 2 u Ω) :
    MemW01p 2 (fun x => c * u x) Ω := by
  let _ := (inferInstance : NeZero d)
  rcases hu with ⟨_, hwu, φ, hφ_smooth, hφ_compact, hφ_sub, hφ_fun, hφ_grad⟩
  refine ⟨(hwu.smul c).memW1p, hwu.smul c, fun n x => c * φ n x, ?_, ?_, ?_, ?_, ?_⟩
  · intro n
    simpa [smul_eq_mul] using (contDiff_const.mul (hφ_smooth n))
  · intro n
    simpa [Pi.smul_apply, smul_eq_mul] using (hφ_compact n).smul_left (f := fun _ : E => c)
  · intro n
    simpa [Pi.smul_apply, smul_eq_mul] using
      (tsupport_smul_subset_right (fun _ : E => c) (φ n)).trans (hφ_sub n)
  ·
    have hscaled :
        Tendsto
          (fun n =>
            ‖c‖ₑ * eLpNorm (fun x => φ n x - u x) 2 (volume.restrict Ω))
          atTop (nhds (‖c‖ₑ * 0)) :=
      ENNReal.Tendsto.const_mul hφ_fun (Or.inr ENNReal.coe_ne_top)
    have hEq :
        (fun n => eLpNorm (fun x => c * φ n x - c * u x) 2 (volume.restrict Ω)) =
          (fun n => ‖c‖ₑ * eLpNorm (fun x => φ n x - u x) 2 (volume.restrict Ω)) := by
      funext n
      have hfun :
          (fun x => c * φ n x - c * u x) = c • (fun x => φ n x - u x) := by
        ext x
        simp [Pi.smul_apply, smul_eq_mul]
        ring
      rw [hfun, eLpNorm_const_smul]
    rw [hEq]
    simpa using hscaled
  · intro i
    have hscaled :
        Tendsto
          (fun n =>
            ‖c‖ₑ *
              eLpNorm
                (fun x => (fderiv ℝ (φ n) x) (EuclideanSpace.single i 1) - hwu.weakGrad x i)
                2 (volume.restrict Ω))
          atTop (nhds (‖c‖ₑ * 0)) :=
      ENNReal.Tendsto.const_mul (hφ_grad i) (Or.inr ENNReal.coe_ne_top)
    have hEq :
        (fun n =>
          eLpNorm
            (fun x =>
              (fderiv ℝ (fun y => c * φ n y) x) (EuclideanSpace.single i 1) -
                (hwu.smul c).weakGrad x i)
            2 (volume.restrict Ω)) =
          (fun n =>
            ‖c‖ₑ *
              eLpNorm
                (fun x => (fderiv ℝ (φ n) x) (EuclideanSpace.single i 1) - hwu.weakGrad x i)
                2 (volume.restrict Ω)) := by
      funext n
      have hfun :
          (fun x =>
            (fderiv ℝ (fun y => c * φ n y) x) (EuclideanSpace.single i 1) -
              (hwu.smul c).weakGrad x i) =
            c • (fun x =>
              (fderiv ℝ (φ n) x) (EuclideanSpace.single i 1) - hwu.weakGrad x i) := by
        ext x
        have hfd : fderiv ℝ (fun y => c * φ n y) x = c • fderiv ℝ (φ n) x := by
          simpa [smul_eq_mul] using congrFun (fderiv_const_smul_field (𝕜 := ℝ) (f := φ n) c) x
        simp [MemW1pWitness.smul, Pi.smul_apply, smul_eq_mul, hfd]
        ring
      rw [hfun, eLpNorm_const_smul]
    rw [hEq]
    simpa using hscaled

/-- `H₀¹(Ω)` is closed under subtraction. -/
theorem MemW01p.sub
    {Ω : Set E} {u v : E → ℝ}
    (hu : MemW01p 2 u Ω) (hv : MemW01p 2 v Ω) :
    MemW01p 2 (fun x => u x - v x) Ω := by
  simpa [sub_eq_add_neg, Pi.smul_apply, smul_eq_mul] using hu.add (hv.smul (-1))

\end{lstlisting}
\begin{proof}
For addition: Lemma~\ref{lem:witness_add} provides a witness for the sum.
Approximating sequences add componentwise; by the triangle inequality for
$\Lp{2}$ on differences,
\begin{equation*}
  \|\varphi_n^u + \varphi_n^v - (u+v)\|_{\Lp{2}(\Omega)}
  \le \|\varphi_n^u - u\|_{\Lp{2}(\Omega)} + \|\varphi_n^v - v\|_{\Lp{2}(\Omega)}
  \;\to\; 0,
\end{equation*}
and similarly for the gradients (using the linearity of the classical
derivative on smooth functions).

For scalar multiplication by $c$: rescale the approximating sequence by $c$ and
use the homogeneity of the $\Lp{2}$ norm,
$\|c\varphi_n - c u\|_{\Lp{2}} = |c|\,\|\varphi_n - u\|_{\Lp{2}}$.

Subtraction follows by writing $u - v = u + (-1)v$.
\end{proof}

\begin{lemma}[Smooth compactly supported $\Rightarrow$ $\Wop{p}$;
\leanname{DeGiorgi.memW01p\_of\_contDiff\_hasCompactSupport},
\leanname{DeGiorgi.memW01p\_of\_contDiff\_hasCompactSupport\_subset}]\label{lem:memW01p_of_smooth_cpt}
If $f \in C_c^\infty(E)$ then $f \in \Wop{p}(E)$. More generally, if $\Omega
\subseteq E$ is open and $f \in C_c^\infty(E)$ with $\supp f \subseteq \Omega$,
then $f \in \Wop{p}(\Omega)$.
\end{lemma}

\begin{lstlisting}[style=Matlab-editor,basicstyle=\mlttfamily,escapechar=",mlshowsectionrules=true,mathescape=true,morecomment={[s]{\%\{}{\%\}}},language=lean,numbers=none]
/-- A smooth compactly supported function belongs to `W₀^{1,p}(ℝ^d)`. -/
theorem memW01p_of_contDiff_hasCompactSupport
    {p : ℝ≥0∞} {f : E → ℝ}
    (hf : ContDiff ℝ (⊤ : ℕ∞) f) (hf_supp : HasCompactSupport f) :
    MemW01p p f Set.univ := by
  let _ := (inferInstance : NeZero d)
  let hw : MemW1pWitness p f Set.univ :=
    MemW1pWitness.of_contDiff_hasCompactSupport (p := p) hf hf_supp
  refine ⟨hw.memW1p, hw, fun _ => f, ?_, ?_, ?_, ?_, ?_⟩
  · intro n
    exact hf
  · intro n
    exact hf_supp
  · intro n
    simp
  · simp
  · intro i
    simp [hw, MemW1pWitness.of_contDiff_hasCompactSupport]

/-- A smooth compactly supported function whose support is contained in an open
set belongs to `W₀^{1,p}` on that set. -/
theorem memW01p_of_contDiff_hasCompactSupport_subset
    {p : ℝ≥0∞} {Ω : Set E} (hΩ : IsOpen Ω) {f : E → ℝ}
    (hf : ContDiff ℝ (⊤ : ℕ∞) f) (hf_supp : HasCompactSupport f)
    (hf_sub : tsupport f ⊆ Ω) :
    MemW01p p f Ω := by
  let _ := (inferInstance : NeZero d)
  let hw : MemW1pWitness p f Ω :=
    { memLp := (hf.continuous.memLp_of_hasCompactSupport hf_supp).restrict Ω
      weakGrad := fun x =>
        WithLp.toLp 2 fun i => (fderiv ℝ f x) (EuclideanSpace.single i 1)
      weakGrad_component_memLp := by
        intro i
        have hderiv_smooth : ContDiff ℝ (⊤ : ℕ∞)
            (fun x => (fderiv ℝ f x) (EuclideanSpace.single i 1)) :=
          (hf.fderiv_right (m := (⊤ : ℕ∞)) (by norm_cast)).clm_apply contDiff_const
        exact (hderiv_smooth.continuous.memLp_of_hasCompactSupport
          (hf_supp.fderiv_apply (𝕜 := ℝ) _)).restrict Ω
      isWeakGrad := by
        intro i
        simpa [PiLp.toLp_apply] using
          (HasWeakPartialDeriv.of_contDiff (Ω := Ω) (i := i)
            (f := f) hΩ (hf.of_le (by norm_cast))) }
  refine ⟨hw.memW1p, hw, fun _ => f, ?_, ?_, ?_, ?_, ?_⟩
  · intro n
    exact hf
  · intro n
    exact hf_supp
  · intro n
    exact hf_sub
  · simp
  · intro i
    simp [hw]

\end{lstlisting}
\begin{proof}
Take the constant approximating sequence $\varphi_n := f$. The $\Lp{p}$
differences are identically zero, so trivially converge to $0$, and
Lemma~\ref{lem:witness_smooth_cpt} provides the witness.
\end{proof}

\subsection{Cutoffs and zero-extension}

\begin{lemma}[Smooth cutoff between a compact set and an open neighborhood;
\leanname{DeGiorgi.exists\_smooth\_cutoff} (private)]\label{lem:smooth_cutoff}
Let $K \subseteq \Omega \subseteq E$ with $K$ compact and $\Omega$ open. Then
there exists $\eta \in C^\infty(E)$ with compact support such that
$0 \le \eta \le 1$, $\eta = 1$ on $K$, and $\supp\eta \subseteq \Omega$.
\end{lemma}

\begin{lstlisting}[style=Matlab-editor,basicstyle=\mlttfamily,escapechar=",mlshowsectionrules=true,mathescape=true,morecomment={[s]{\%\{}{\%\}}},language=lean,numbers=none]
omit [NeZero d] in
private lemma exists_smooth_cutoff
    {K Ω : Set E}
    (hK : IsCompact K) (hΩ : IsOpen Ω) (hKΩ : K ⊆ Ω) :
    ∃ η : E → ℝ,
      ContDiff ℝ (⊤ : ℕ∞) η ∧
      HasCompactSupport η ∧
      Set.range η ⊆ Set.Icc (0 : ℝ) 1 ∧
      (∀ x ∈ K, η x = 1) ∧
      tsupport η ⊆ Ω := by
  obtain ⟨δ, hδ_pos, hδΩ⟩ := hK.exists_cthickening_subset_open hΩ hKΩ
  rcases exists_contMDiff_support_eq_eq_one_iff
      (I := modelWithCornersSelf ℝ E) (s := Metric.thickening δ K) (t := K)
      isOpen_thickening hK.isClosed (self_subset_thickening hδ_pos K) with
    ⟨η, hη_smooth, hη_range, hη_support, hη_one_iff⟩
  refine ⟨η, contMDiff_iff_contDiff.mp hη_smooth, ?_, hη_range, ?_, ?_⟩
  · apply HasCompactSupport.intro' (K := Metric.cthickening δ K)
    · exact hK.cthickening (r := δ)
    · simpa using (isClosed_cthickening : IsClosed (Metric.cthickening δ K))
    · intro x hx
      have hxt : x ∉ tsupport η := by
        intro hxt
        have hx_closure : x ∈ closure (Metric.thickening δ K) := by
          rw [tsupport, hη_support] at hxt
          exact hxt
        exact hx ((Metric.closure_thickening_subset_cthickening δ K) hx_closure)
      exact image_eq_zero_of_notMem_tsupport hxt
  · intro x hx
    exact (hη_one_iff x).1 hx
  · rw [tsupport, hη_support]
    exact Metric.closure_thickening_subset_cthickening δ K |>.trans hδΩ

\end{lstlisting}
\begin{proof}
Since $K$ is compact and contained in the open set $\Omega$, there exists
$\delta > 0$ such that the closed $\delta$-thickening of $K$ is contained in
$\Omega$. The classical existence of a smooth bump (in Lean,
\leanname{exists\_contMDiff\_support\_eq\_eq\_one\_iff}) yields a smooth $\eta$
with $\eta(x) = 1$ on $K$, support equal to the open $\delta$-thickening of $K$,
and range in $[0,1]$. Its support is then contained in the closed
$\delta$-thickening of $K$ and hence in $\Omega$, and it is automatically
compactly supported.
\end{proof}

\begin{lemma}[Vanishing outside the support; \leanname{DeGiorgi.zero\_outside\_of\_tsupport\_subset}, \leanname{DeGiorgi.fderiv\_apply\_zero\_outside\_of\_tsupport\_subset}]\label{lem:zero_outside_tsupport}
If $\supp f \subseteq \Omega$ and $x \notin \Omega$, then $f(x) = 0$. If
furthermore $f \in C^\infty$ then $\partial_i f(x) = 0$ as well.
\end{lemma}

\begin{lstlisting}[style=Matlab-editor,basicstyle=\mlttfamily,escapechar=",mlshowsectionrules=true,mathescape=true,morecomment={[s]{\%\{}{\%\}}},language=lean,numbers=none]
omit [NeZero d] in
theorem zero_outside_of_tsupport_subset
    {Ω : Set E} {f : E → ℝ}
    (hsub : tsupport f ⊆ Ω) {x : E} (hx : x ∉ Ω) :
    f x = 0 := by
  by_cases hfx : f x = 0
  · exact hfx
  · exfalso
    exact hx (hsub (subset_tsupport f (mem_support.mpr hfx)))

theorem fderiv_apply_zero_outside_of_tsupport_subset
    {Ω : Set E} {f : E → ℝ}
    (hf : ContDiff ℝ (⊤ : ℕ∞) f)
    (hsub : tsupport f ⊆ Ω) {x : E} (hx : x ∉ Ω) (i : Fin d) :
    (fderiv ℝ f x) (EuclideanSpace.single i 1) = 0 := by
  let _ := (inferInstance : NeZero d)
  let _ := hf
  apply zero_outside_of_tsupport_subset
    (Ω := Ω)
    (f := fun y => (fderiv ℝ f y) (EuclideanSpace.single i 1))
  · exact (tsupport_fderiv_apply_subset ℝ (EuclideanSpace.single i (1 : ℝ))).trans hsub
  · exact hx

\end{lstlisting}
\begin{proof}
If $f(x) \ne 0$ then $x$ lies in the support of $f$, hence in $\Omega$,
contradicting $x \notin \Omega$. The same argument applied to $\partial_i f$
(whose support is contained in that of $f$) handles the second statement.
\end{proof}

\begin{theorem}[Zero-extension of a $\Wop{p}$ function; \leanname{DeGiorgi.zeroExtend\_memW1pWitness\_p}]\label{thm:zero_extend}
Let $\Omega \subseteq E$ be open, $1 < p < \infty$, and $v \in \Wop{p}(\Omega)$
with $\Wkp{1}{p}$-witness $G_v$; assume in addition that $v \in W_0^{1,p}(\Omega)$. Define $\tilde v := \mathbf{1}_\Omega \cdot v$
on $E$ and let $\tilde G$ be the vector field whose components are
$\mathbf{1}_\Omega \cdot (G_v)_i$. Then $\tilde G$ is a $\Wkp{1}{p}$-witness
for $\tilde v$ on the whole space $E$.
\end{theorem}

\begin{lstlisting}[style=Matlab-editor,basicstyle=\mlttfamily,escapechar=",mlshowsectionrules=true,mathescape=true,morecomment={[s]{\%\{}{\%\}}},language=lean,numbers=none]
/-- Zero-extension of a compactly supported local `W^{1,p}` witness to all of
`ℝ^d`, using the `MemW01p` approximation sequence on the source open set. -/
noncomputable def zeroExtend_memW1pWitness_p
    {Ω : Set E} (hΩ : IsOpen Ω)
    {p : ℝ} (hp : 1 < p)
    {v : E → ℝ}
    (hv : MemW01p (ENNReal.ofReal p) v Ω)
    (hw : MemW1pWitness (ENNReal.ofReal p) v Ω) :
    MemW1pWitness (ENNReal.ofReal p) (Ω.indicator v) Set.univ := by
  classical
  have hΩ_meas : MeasurableSet Ω := hΩ.measurableSet
  let hw₀ : MemW1pWitness (ENNReal.ofReal p) v Ω := Classical.choose hv.2
  let hExφ := Classical.choose_spec hv.2
  let φ : ℕ → E → ℝ := Classical.choose hExφ
  have hφspec := Classical.choose_spec hExφ
  have hφ_smooth : ∀ n, ContDiff ℝ (⊤ : ℕ∞) (φ n) := hφspec.1
  have hφ_compact : ∀ n, HasCompactSupport (φ n) := hφspec.2.1
  have hφ_sub : ∀ n, tsupport (φ n) ⊆ Ω := hφspec.2.2.1
  have hφ_fun :
      Tendsto
        (fun n => eLpNorm (fun x => φ n x - v x) (ENNReal.ofReal p) (volume.restrict Ω))
        atTop (nhds 0) := hφspec.2.2.2.1
  have hφ_grad :
      ∀ i : Fin d,
        Tendsto
          (fun n =>
            eLpNorm
              (fun x => (fderiv ℝ (φ n) x) (EuclideanSpace.single i 1) - hw₀.weakGrad x i)
              (ENNReal.ofReal p) (volume.restrict Ω))
          atTop (nhds 0) := hφspec.2.2.2.2
  let G : E → E := fun x =>
    WithLp.toLp 2 fun i => Ω.indicator (fun y => hw.weakGrad y i) x
  refine
    { memLp := ?_
      weakGrad := G
      weakGrad_component_memLp := ?_
      isWeakGrad := ?_ }
  · simpa using
      (MeasureTheory.memLp_indicator_iff_restrict
        (μ := volume) (s := Ω) (f := v) (p := ENNReal.ofReal p) hΩ_meas).2 hw.memLp
  · intro i
    simpa [G, PiLp.toLp_apply] using
      (MeasureTheory.memLp_indicator_iff_restrict
        (μ := volume) (s := Ω) (f := fun y => hw.weakGrad y i) (p := ENNReal.ofReal p) hΩ_meas).2
        (hw.weakGrad_component_memLp i)
  · intro i
    have hp_le : 1 ≤ p := le_of_lt hp
    have hcompΩ :
        (fun x => hw₀.weakGrad x i) =ᵐ[volume.restrict Ω] (fun x => hw.weakGrad x i) := by
      exact HasWeakPartialDeriv.ae_eq hΩ (hw₀.isWeakGrad i) (hw.isWeakGrad i)
        ((hw₀.weakGrad_component_memLp i).locallyIntegrable (by
          simpa using (ENNReal.ofReal_le_ofReal hp_le : ENNReal.ofReal (1 : ℝ) ≤ ENNReal.ofReal p)))
        ((hw.weakGrad_component_memLp i).locallyIntegrable (by
          simpa using (ENNReal.ofReal_le_ofReal hp_le : ENNReal.ofReal (1 : ℝ) ≤ ENNReal.ofReal p)))
    refine HasWeakPartialDeriv.of_eLpNormApprox_p
      (d := d) (Ω := Set.univ) (p := p) (hΩ := isOpen_univ) (hp := hp)
      (i := i) (f := Ω.indicator v) (g := fun x => G x i)
      (ψ := φ)
      (gψ := fun n x => (fderiv ℝ (φ n) x) (EuclideanSpace.single i 1))
      ?_ ?_ ?_ ?_ ?_ ?_ ?_
    · rw [Measure.restrict_univ]
      simpa using
        (MeasureTheory.memLp_indicator_iff_restrict
          (μ := volume) (s := Ω) (f := v) (p := ENNReal.ofReal p) hΩ_meas).2 hw.memLp
    · rw [Measure.restrict_univ]
      simpa [G, PiLp.toLp_apply] using
        (MeasureTheory.memLp_indicator_iff_restrict
          (μ := volume) (s := Ω) (f := fun y => hw.weakGrad y i) (p := ENNReal.ofReal p) hΩ_meas).2
          (hw.weakGrad_component_memLp i)
    · intro n
      simpa using
        HasWeakPartialDeriv.of_contDiff (Ω := Set.univ) (i := i) (f := φ n)
          isOpen_univ ((hφ_smooth n).of_le (by norm_cast))
    · intro n
      have hφ_memLp :
          MemLp (φ n) (ENNReal.ofReal p) (volume.restrict Ω) :=
        ((hφ_smooth n).continuous.memLp_of_hasCompactSupport (hφ_compact n)).restrict Ω
      have hEq :
          (fun x => φ n x - Ω.indicator v x) =
            Ω.indicator (fun x => φ n x - v x) := by
        ext x
        by_cases hx : x ∈ Ω
        · simp [hx]
        · have hφx : φ n x = 0 := zero_outside_of_tsupport_subset (Ω := Ω) (hφ_sub n) hx
          simp [hx, hφx]
      rw [Measure.restrict_univ, hEq,
        MeasureTheory.memLp_indicator_iff_restrict (μ := volume) (s := Ω)
          (p := ENNReal.ofReal p) hΩ_meas]
      exact hφ_memLp.sub hw.memLp
    · have hEq :
          (fun n => eLpNorm (fun x => φ n x - Ω.indicator v x) (ENNReal.ofReal p) volume) =
            fun n => eLpNorm (fun x => φ n x - v x) (ENNReal.ofReal p) (volume.restrict Ω) := by
        funext n
        have hFn :
            (fun x => φ n x - Ω.indicator v x) =
              Ω.indicator (fun x => φ n x - v x) := by
          ext x
          by_cases hx : x ∈ Ω
          · simp [hx]
          · have hφx : φ n x = 0 := zero_outside_of_tsupport_subset (Ω := Ω) (hφ_sub n) hx
            simp [hx, hφx]
        rw [hFn, MeasureTheory.eLpNorm_indicator_eq_eLpNorm_restrict
          (μ := volume) (s := Ω) (p := ENNReal.ofReal p) hΩ_meas]
      rw [Measure.restrict_univ, hEq]
      simpa using hφ_fun
    · intro n
      let ei : E := EuclideanSpace.single i (1 : ℝ)
      have hderiv_smooth :
          ContDiff ℝ (⊤ : ℕ∞) (fun x => (fderiv ℝ (φ n) x) ei) :=
        (hφ_smooth n).fderiv_right (m := (⊤ : ℕ∞)) (by norm_cast)
          |>.clm_apply contDiff_const
      have hderiv_memLp :
          MemLp (fun x => (fderiv ℝ (φ n) x) ei) (ENNReal.ofReal p) (volume.restrict Ω) :=
        (hderiv_smooth.continuous.memLp_of_hasCompactSupport
          ((hφ_compact n).fderiv_apply (𝕜 := ℝ) _)).restrict Ω
      have hEq :
          (fun x => (fderiv ℝ (φ n) x) ei - G x i) =
            Ω.indicator (fun x => (fderiv ℝ (φ n) x) ei - hw.weakGrad x i) := by
        ext x
        by_cases hx : x ∈ Ω
        · simp [G, hx]
        · have hdx :
              (fderiv ℝ (φ n) x) ei = 0 := by
            simpa [ei] using
              fderiv_apply_zero_outside_of_tsupport_subset (Ω := Ω) (hf := hφ_smooth n)
                (hsub := hφ_sub n) hx i
          simp [G, hx, hdx]
      rw [Measure.restrict_univ, hEq,
        MeasureTheory.memLp_indicator_iff_restrict (μ := volume) (s := Ω)
          (p := ENNReal.ofReal p) hΩ_meas]
      exact hderiv_memLp.sub (hw.weakGrad_component_memLp i)
    · have hEq :
          (fun n =>
            eLpNorm
              (fun x => (fderiv ℝ (φ n) x) (EuclideanSpace.single i 1) - G x i)
              (ENNReal.ofReal p) volume) =
            fun n =>
              eLpNorm
                (fun x => (fderiv ℝ (φ n) x) (EuclideanSpace.single i 1) - hw.weakGrad x i)
                (ENNReal.ofReal p) (volume.restrict Ω) := by
        funext n
        have hFn :
            (fun x => (fderiv ℝ (φ n) x) (EuclideanSpace.single i 1) - G x i) =
              Ω.indicator
                (fun x => (fderiv ℝ (φ n) x) (EuclideanSpace.single i 1) - hw.weakGrad x i) := by
          ext x
          by_cases hx : x ∈ Ω
          · simp [G, hx]
          · have hdx :
                (fderiv ℝ (φ n) x) (EuclideanSpace.single i 1) = 0 := by
              exact fderiv_apply_zero_outside_of_tsupport_subset (Ω := Ω) (hf := hφ_smooth n)
                (hsub := hφ_sub n) hx i
            simp [G, hx, hdx]
        rw [hFn, MeasureTheory.eLpNorm_indicator_eq_eLpNorm_restrict
          (μ := volume) (s := Ω) (p := ENNReal.ofReal p) hΩ_meas]
      rw [Measure.restrict_univ, hEq]
      have hEqSeq :
          (fun n =>
            eLpNorm
              (fun x => (fderiv ℝ (φ n) x) (EuclideanSpace.single i 1) - hw.weakGrad x i)
              (ENNReal.ofReal p) (volume.restrict Ω)) =
            (fun n =>
              eLpNorm
                (fun x => (fderiv ℝ (φ n) x) (EuclideanSpace.single i 1) - hw₀.weakGrad x i)
                (ENNReal.ofReal p) (volume.restrict Ω)) := by
        funext n
        have hDiffAe :
            (fun x => (fderiv ℝ (φ n) x) (EuclideanSpace.single i 1) - hw.weakGrad x i)
              =ᵐ[volume.restrict Ω]
            (fun x => (fderiv ℝ (φ n) x) (EuclideanSpace.single i 1) - hw₀.weakGrad x i) := by
          filter_upwards [hcompΩ] with x hx
          simp [hx]
        exact eLpNorm_congr_ae hDiffAe
      rw [hEqSeq]
      simpa using hφ_grad i

\end{lstlisting}
\begin{proof}
The function $v$ admits an approximating sequence $(\varphi_n)$ of smooth
compactly supported test functions with $\supp\varphi_n \subseteq \Omega$, and
$\varphi_n \to v$, $\partial_i\varphi_n \to (G_v)_i$ in $\Lp{p}(\Omega)$. By
Lemma~\ref{lem:weak_partial_unique} the gradient field $G_v$ obtained from the
arbitrary witness coincides almost everywhere on $\Omega$ with the one from the
$\Wop{p}$ data.

The $\Lp{p}$ membership of $\tilde v$ on $E$ is equivalent to the $\Lp{p}$
membership of $v$ on $\Omega$, via
\leanname{MeasureTheory.memLp\_indicator\_iff\_restrict}, and similarly for each
component of $\tilde G$.

To check that $\tilde G$ is a weak gradient on $E$, we apply
Lemma~\ref{lem:weak_lp_closure} (closure of weak partial derivatives in
$\Lp{p}$, on $\Omega = E$) with the same approximating sequence
$(\varphi_n)$ and the classical gradients $\partial_i\varphi_n$. The
key observations are:
\begin{itemize}
  \item Outside $\Omega$ both $\varphi_n$ and $\partial_i\varphi_n$ vanish (by
        Lemma~\ref{lem:zero_outside_tsupport}); hence
        $\varphi_n - \tilde v = \mathbf{1}_\Omega(\varphi_n - v)$ and
        $\partial_i\varphi_n - \tilde G_i =
         \mathbf{1}_\Omega(\partial_i\varphi_n - (G_v)_i)$ as functions on $E$.
  \item The $\Lp{p}$-norm of an indicator of a measurable set equals the
        $\Lp{p}$-norm computed on the restricted measure
        (\leanname{MeasureTheory.eLpNorm\_indicator\_eq\_eLpNorm\_restrict}),
        so the convergence on $\Omega$ implies convergence on $E$.
\end{itemize}
The hypotheses of Lemma~\ref{lem:weak_lp_closure} are then satisfied, and
$\tilde G_i$ is the weak $i$-th partial derivative of $\tilde v$ on $E$.
\end{proof}

\subsection{Convolution smoothing}

For $n \in \N$, define a $C^\infty$ bump function on $E$ centered at $0$ with
inner radius $r_{\text{in}}^{(n)} := \frac{1}{2(n+1)}$ and outer radius
$r_{\text{out}}^{(n)} := \frac{1}{n+1}$ (\leanname{DeGiorgi.shrinkingBump} (private)).
We denote the corresponding mass-one normalized bump by $\rho_n :=
(\text{shrinkingBump}_n).\text{normed}$, so $\int_E \rho_n = 1$, $\rho_n \ge 0$,
and $\supp \rho_n \subseteq \Bz{r_{\text{out}}^{(n)}}$. The radii satisfy
\begin{equation*}
  r_{\text{out}}^{(n)} = \frac{1}{n+1} \;\to\; 0
  \quad\text{as } n \to \infty
\end{equation*}
(\leanname{DeGiorgi.tendsto\_shrinkingBump\_rOut} (private)).

\begin{lemma}[Support of a convolution; \leanname{DeGiorgi.tsupport\_convolution\_subset\_cthickening} (private), \leanname{DeGiorgi.tsupport\_normed\_convolution\_subset\_cthickening} (private)]\label{lem:tsupport_conv}
If $\supp \kappa \subseteq \overline{\Bz{r}}$ and $u$ has compact support, then
\begin{equation*}
  \supp(\kappa * u) \;\subseteq\; \{x \in E : \dist(x, \supp u) \le r\}.
\end{equation*}
In particular, $\supp(\rho_n * u) \subseteq \{x : \dist(x,\supp u) \le
r_{\text{out}}^{(n)}\}$.
\end{lemma}

\begin{lstlisting}[style=Matlab-editor,basicstyle=\mlttfamily,escapechar=",mlshowsectionrules=true,mathescape=true,morecomment={[s]{\%\{}{\%\}}},language=lean,numbers=none]
omit [NeZero d] in
private lemma tsupport_convolution_subset_cthickening
    {κ u : E → ℝ} {r : ℝ}
    (hκ : Function.support κ ⊆ Metric.closedBall (0 : E) r)
    (hu : HasCompactSupport u) :
    tsupport (κ ⋆[ContinuousLinearMap.lsmul ℝ ℝ, volume] u)
      ⊆ Metric.cthickening r (tsupport u) := by
  let _ := hu
  have hsupport :
      Function.support (κ ⋆[ContinuousLinearMap.lsmul ℝ ℝ, volume] u)
        ⊆ Function.support κ + Function.support u :=
    support_convolution_subset (L := ContinuousLinearMap.lsmul ℝ ℝ) (μ := volume)
  have hsum :
      Function.support κ + Function.support u ⊆ Metric.cthickening r (tsupport u) := by
    intro x hx
    rcases Set.mem_add.mp hx with ⟨a, ha, b, hb, rfl⟩
    have ha_ball : a ∈ Metric.closedBall (0 : E) r := hκ ha
    have hb_tsupport : b ∈ tsupport u := subset_tsupport u hb
    refine Metric.mem_cthickening_of_dist_le (x := a + b) (y := b) (δ := r)
      (tsupport u) hb_tsupport ?_
    simpa [Metric.mem_closedBall, dist_eq_norm, sub_eq_add_neg, add_comm, add_left_comm, add_assoc]
      using ha_ball
  simpa [tsupport] using
    closure_minimal (hsupport.trans hsum) isClosed_cthickening

omit [NeZero d] in
private lemma tsupport_normed_convolution_subset_cthickening
    (φ : ContDiffBump (0 : E)) {u : E → ℝ}
    (hu : HasCompactSupport u) :
    tsupport ((φ.normed volume) ⋆[ContinuousLinearMap.lsmul ℝ ℝ, volume] u)
      ⊆ Metric.cthickening φ.rOut (tsupport u) :=
  tsupport_convolution_subset_cthickening
    (r := φ.rOut)
    (fun x hx => (Metric.mem_ball.1 ((by simpa [φ.support_normed_eq] using hx))).le)
    hu

\end{lstlisting}
\begin{proof}
Pointwise, $(\kappa * u)(x) = \int \kappa(y) u(x - y)\,dy$ vanishes unless
there exists $y \in \supp\kappa$ and $z \in \supp u$ with $x = y + z$, hence
$\dist(x, z) = \|y\| \le r$. Closing under the topology, the closure is
contained in the closed $r$-thickening, which is closed.
\end{proof}

\begin{proposition}[$\Lp{p}$-contraction of normed convolution;
\leanname{DeGiorgi.eLpNorm\_normedConvolution\_le} (private)]\label{prop:Lp_contraction_conv}
Let $\rho \in C_c^\infty(E)$ with $\rho \ge 0$ and $\int_E \rho = 1$, and let
$1 \le p < \infty$. Then for every $f \in \Lp{p}(E)$,
\begin{equation*}
  \|\rho * f\|_{\Lp{p}(E)} \le \|f\|_{\Lp{p}(E)}.
\end{equation*}
\end{proposition}

\begin{lstlisting}[style=Matlab-editor,basicstyle=\mlttfamily,escapechar=",mlshowsectionrules=true,mathescape=true,morecomment={[s]{\%\{}{\%\}}},language=lean,numbers=none]
omit [NeZero d] in
private theorem eLpNorm_normedConvolution_le
    (φ : ContDiffBump (0 : E))
    {p : ℝ} (hp : 1 ≤ p) {f : E → ℝ}
    (hf : MemLp f (ENNReal.ofReal p) volume) :
    eLpNorm ((φ.normed volume) ⋆[ContinuousLinearMap.lsmul ℝ ℝ, volume] f)
      (ENNReal.ofReal p) volume ≤
    eLpNorm f (ENNReal.ofReal p) volume := by
  let ψ : E → ℝ := φ.normed volume
  let ρ : E → ℝ≥0∞ := fun y => ENNReal.ofReal (ψ y)
  let ν : Measure E := volume.withDensity ρ
  have hp_pos : 0 < p := lt_of_lt_of_le zero_lt_one hp
  have hp0 : ENNReal.ofReal p ≠ 0 := by
    exact ne_of_gt (ENNReal.ofReal_pos.mpr hp_pos)
  have hp_toReal : (ENNReal.ofReal p).toReal = p := by
    simp [hp_pos.le]
  have hρ_meas : Measurable ρ := by
    exact φ.continuous_normed.measurable.ennreal_ofReal
  have hρ_lt_top : ∀ᵐ y ∂volume, ρ y < ∞ := by
    filter_upwards with y
    simp [ρ]
  have hρ_ne_top : ∀ᵐ y ∂volume, ρ y ≠ ∞ := by
    filter_upwards [hρ_lt_top] with y hy
    exact hy.ne
  have hν_prob : IsProbabilityMeasure ν := by
    refine ⟨?_⟩
    rw [show ν = volume.withDensity ρ by rfl, withDensity_apply _ MeasurableSet.univ]
    rw [Measure.restrict_univ]
    rw [← MeasureTheory.ofReal_integral_eq_lintegral_ofReal
      φ.integrable_normed (Filter.Eventually.of_forall φ.nonneg_normed)]
    simpa [ρ, ψ] using φ.integral_normed (μ := volume)
  have hf_loc : LocallyIntegrable f volume := by
    exact hf.locallyIntegrable (by simpa using (ENNReal.ofReal_le_ofReal hp : ENNReal.ofReal (1 : ℝ) ≤ ENNReal.ofReal p))
  have hh_int : Integrable (fun x => ‖f x‖ ^ p) volume := by
    simpa [hp_toReal] using hf.integrable_norm_rpow hp0 ENNReal.ofReal_ne_top
  have hψ_int : Integrable ψ volume := by
    simpa [ψ] using φ.integrable_normed
  have hψ_cont : Continuous ψ := by
    simpa [ψ] using φ.continuous_normed
  have hψ_cpt : HasCompactSupport ψ := by
    simpa [ψ] using φ.hasCompactSupport_normed (μ := volume)
  have hconv_cont :
      Continuous (((φ.normed volume) ⋆[ContinuousLinearMap.lsmul ℝ ℝ, volume] f) : E → ℝ) := by
    exact hψ_cpt.continuous_convolution_left (L := ContinuousLinearMap.lsmul ℝ ℝ) hψ_cont hf_loc
  have hpointwise :
      ∀ x : E,
        ‖((ψ ⋆[ContinuousLinearMap.lsmul ℝ ℝ, volume] f) x)‖ ^ p ≤
          (ψ ⋆[ContinuousLinearMap.lsmul ℝ ℝ, volume] (fun y => ‖f y‖ ^ p)) x := by
    intro x
    have hce_f :
        ConvolutionExistsAt ψ f x (ContinuousLinearMap.lsmul ℝ ℝ) volume := by
      exact (hψ_cpt.convolutionExists_left (L := ContinuousLinearMap.lsmul ℝ ℝ) hψ_cont hf_loc) x
    have hce_h :
        ConvolutionExistsAt ψ (fun y => ‖f y‖ ^ p) x (ContinuousLinearMap.lsmul ℝ ℝ) volume := by
      exact (hψ_cpt.convolutionExists_left
        (L := ContinuousLinearMap.lsmul ℝ ℝ) hψ_cont hh_int.locallyIntegrable) x
    have hconv_ν :
        ((ψ ⋆[ContinuousLinearMap.lsmul ℝ ℝ, volume] f) x) = ∫ y, f (x - y) ∂ν := by
      rw [show ν = volume.withDensity ρ by rfl]
      rw [integral_withDensity_eq_integral_toReal_smul hρ_meas hρ_lt_top]
      rw [MeasureTheory.convolution_def]
      refine integral_congr_ae ?_
      filter_upwards with y
      simp [ρ, ψ, ContinuousLinearMap.lsmul_apply, smul_eq_mul, φ.nonneg_normed y, mul_comm]
    have h_int_norm_vol :
        Integrable (fun y => (ρ y).toReal * ‖f (x - y)‖) volume := by
      refine hce_f.integrable.norm.congr ?_
      filter_upwards with y
      simp [ρ, ψ, ContinuousLinearMap.lsmul_apply, smul_eq_mul, φ.nonneg_normed y]
    have h_int_norm :
        Integrable (fun y => ‖f (x - y)‖) ν := by
      exact (MeasureTheory.integrable_withDensity_iff_integrable_smul' hρ_meas hρ_lt_top).2 <| by
        simpa [ρ, ψ, φ.nonneg_normed] using h_int_norm_vol
    have h_int_rpow_vol :
        Integrable (fun y => (ρ y).toReal * ‖f (x - y)‖ ^ p) volume := by
      refine hce_h.integrable.congr ?_
      filter_upwards with y
      simp [ρ, ψ, ContinuousLinearMap.lsmul_apply, smul_eq_mul, φ.nonneg_normed y, mul_comm]
    have h_int_rpow :
        Integrable (fun y => ‖f (x - y)‖ ^ p) ν := by
      exact (MeasureTheory.integrable_withDensity_iff_integrable_smul' hρ_meas hρ_lt_top).2 <| by
        simpa [ρ, ψ, φ.nonneg_normed] using h_int_rpow_vol
    have hnorm_le : ‖((ψ ⋆[ContinuousLinearMap.lsmul ℝ ℝ, volume] f) x)‖ ≤ ∫ y, ‖f (x - y)‖ ∂ν := by
      rw [hconv_ν]
      exact MeasureTheory.norm_integral_le_integral_norm _
    have hJensen :
        (∫ y, ‖f (x - y)‖ ∂ν) ^ p ≤ ∫ y, ‖f (x - y)‖ ^ p ∂ν := by
      have hconv : ConvexOn ℝ (Set.Ici 0) (fun t : ℝ => t ^ p) := convexOn_rpow hp
      have hcont : ContinuousOn (fun t : ℝ => t ^ p) (Set.Ici (0 : ℝ)) := by
        exact continuousOn_id.rpow_const (by
          intro t ht
          exact Or.inr hp_pos.le)
      have hmem : ∀ᵐ y ∂ν, ‖f (x - y)‖ ∈ Set.Ici (0 : ℝ) := by
        filter_upwards with y
        exact norm_nonneg _
      exact hconv.map_integral_le hcont isClosed_Ici hmem h_int_norm h_int_rpow
    calc
      ‖((ψ ⋆[ContinuousLinearMap.lsmul ℝ ℝ, volume] f) x)‖ ^ p
          ≤ (∫ y, ‖f (x - y)‖ ∂ν) ^ p := by
            exact Real.rpow_le_rpow (norm_nonneg _) hnorm_le hp_pos.le
      _ ≤ ∫ y, ‖f (x - y)‖ ^ p ∂ν := hJensen
      _ = ∫ y, ψ y * ‖f (x - y)‖ ^ p ∂volume := by
            rw [show ν = volume.withDensity ρ by rfl]
            rw [integral_withDensity_eq_integral_toReal_smul hρ_meas hρ_lt_top]
            refine integral_congr_ae ?_
            filter_upwards with y
            simp [ρ, ψ, φ.nonneg_normed y, mul_comm]
      _ = (ψ ⋆[ContinuousLinearMap.lsmul ℝ ℝ, volume] (fun y => ‖f y‖ ^ p)) x := by
            simp [MeasureTheory.convolution_def, ContinuousLinearMap.lsmul_apply, mul_comm]
  have hconv_h_int :
      Integrable (ψ ⋆[ContinuousLinearMap.lsmul ℝ ℝ, volume] (fun y => ‖f y‖ ^ p)) volume := by
    exact hψ_int.integrable_convolution (L := ContinuousLinearMap.lsmul ℝ ℝ) hh_int
  have hpow_int :
      Integrable (fun x =>
        ‖((ψ ⋆[ContinuousLinearMap.lsmul ℝ ℝ, volume] f) x)‖ ^ p) volume := by
    refine Integrable.mono' hconv_h_int ?_ ?_
    · exact (hconv_cont.norm.rpow_const (by intro x; exact Or.inr hp_pos.le)).aestronglyMeasurable
    · filter_upwards with x
      have hnonneg : 0 ≤ |((ψ ⋆[ContinuousLinearMap.lsmul ℝ ℝ, volume] f) x)| ^ p := by
        positivity
      simpa [Real.norm_eq_abs, abs_of_nonneg hnonneg] using hpointwise x
  have h_int_le :
      ∫ x, ‖((ψ ⋆[ContinuousLinearMap.lsmul ℝ ℝ, volume] f) x)‖ ^ p ∂volume ≤
        ∫ x, (ψ ⋆[ContinuousLinearMap.lsmul ℝ ℝ, volume] (fun y => ‖f y‖ ^ p)) x ∂volume := by
    refine integral_mono_ae hpow_int hconv_h_int ?_
    filter_upwards with x
    exact hpointwise x
  have h_conv_eq :
      ∫ x, (ψ ⋆[ContinuousLinearMap.lsmul ℝ ℝ, volume] (fun y => ‖f y‖ ^ p)) x ∂volume =
        ∫ x, ‖f x‖ ^ p ∂volume := by
    calc
      ∫ x, (ψ ⋆[ContinuousLinearMap.lsmul ℝ ℝ, volume] (fun y => ‖f y‖ ^ p)) x ∂volume
          = (∫ x, ψ x ∂volume) * (∫ x, ‖f x‖ ^ p ∂volume) := by
              simpa [ContinuousLinearMap.lsmul_apply, mul_comm, mul_left_comm, mul_assoc] using
                (MeasureTheory.integral_convolution
                  (L := ContinuousLinearMap.lsmul ℝ ℝ)
                  (μ := volume) (ν := volume) hψ_int hh_int)
      _ = ∫ x, ‖f x‖ ^ p ∂volume := by
            simp [ψ, φ.integral_normed (μ := volume)]
  have hconv_memLp :
      MemLp ((ψ ⋆[ContinuousLinearMap.lsmul ℝ ℝ, volume] f) : E → ℝ)
        (ENNReal.ofReal p) volume := by
    exact (MeasureTheory.integrable_norm_rpow_iff
      hconv_cont.aestronglyMeasurable hp0 ENNReal.ofReal_ne_top).1 <| by
        simpa [hp_toReal] using hpow_int
  calc
    eLpNorm ((ψ ⋆[ContinuousLinearMap.lsmul ℝ ℝ, volume] f) : E → ℝ)
        (ENNReal.ofReal p) volume
      = ENNReal.ofReal ((∫ x, ‖((ψ ⋆[ContinuousLinearMap.lsmul ℝ ℝ, volume] f) x)‖ ^ p ∂volume) ^ p⁻¹) := by
          simpa [hp_toReal] using
            hconv_memLp.eLpNorm_eq_integral_rpow_norm hp0 ENNReal.ofReal_ne_top
    _ ≤ ENNReal.ofReal ((∫ x, ‖f x‖ ^ p ∂volume) ^ p⁻¹) := by
          apply ENNReal.ofReal_le_ofReal
          refine Real.rpow_le_rpow ?_ (h_int_le.trans_eq h_conv_eq) ?_
          · exact integral_nonneg fun _ => by positivity
          · exact inv_nonneg.mpr hp_pos.le
    _ = eLpNorm f (ENNReal.ofReal p) volume := by
          symm
          simpa [hp_toReal] using hf.eLpNorm_eq_integral_rpow_norm hp0 ENNReal.ofReal_ne_top

\end{lstlisting}
\begin{proof}
Define the probability measure $\nu$ on $E$ by $d\nu = \rho\,d\,\mathrm{vol}$,
which is a probability measure since $\rho \ge 0$ and $\int \rho = 1$. For each
$x \in E$,
\begin{equation*}
  (\rho * f)(x) = \int_E f(x - y)\,d\nu(y).
\end{equation*}
By the integral form of Jensen's inequality applied to the convex function
$t \mapsto t^p$ on $[0,\infty)$ (\leanname{ConvexOn.map\_integral\_le}),
\begin{equation*}
  |(\rho * f)(x)|^p
    \le \left(\int_E |f(x-y)|\,d\nu(y)\right)^p
    \le \int_E |f(x-y)|^p \,d\nu(y)
    = (\rho * |f|^p)(x).
\end{equation*}
Integrating in $x$ and using Fubini together with $\int \rho = 1$,
\begin{equation*}
  \int_E |(\rho * f)(x)|^p\,dx \le \int_E (\rho * |f|^p)(x)\,dx
    = \left(\int_E \rho\right)\left(\int_E |f|^p\right)
    = \|f\|_{\Lp{p}}^p.
\end{equation*}
Taking $p$-th roots gives the claim.
\end{proof}

\begin{corollary}[Convolution is linear on differences;
\leanname{DeGiorgi.normedConvolution\_sub\_eq} (private),
\leanname{DeGiorgi.eLpNorm\_normedConvolution\_sub\_le} (private)]\label{cor:conv_sub}
Under the assumptions of Proposition~\ref{prop:Lp_contraction_conv} we also
have, for $f, g \in \Lp{p}(E)$,
\begin{equation*}
  \|\rho * f - \rho * g\|_{\Lp{p}(E)} = \|\rho * (f - g)\|_{\Lp{p}(E)}
   \le \|f - g\|_{\Lp{p}(E)}.
\end{equation*}
\end{corollary}

\begin{lstlisting}[style=Matlab-editor,basicstyle=\mlttfamily,escapechar=",mlshowsectionrules=true,mathescape=true,morecomment={[s]{\%\{}{\%\}}},language=lean,numbers=none]
omit [NeZero d] in
private theorem normedConvolution_sub_eq
    (φ : ContDiffBump (0 : E)) {f g : E → ℝ}
    (hf_loc : LocallyIntegrable f volume) (hg_loc : LocallyIntegrable g volume) :
    (fun x =>
      ((φ.normed volume) ⋆[ContinuousLinearMap.lsmul ℝ ℝ, volume] f) x -
        ((φ.normed volume) ⋆[ContinuousLinearMap.lsmul ℝ ℝ, volume] g) x) =
    (fun x =>
      ((φ.normed volume) ⋆[ContinuousLinearMap.lsmul ℝ ℝ, volume] fun y => f y - g y) x) := by
  ext x
  have hce_f := (φ.hasCompactSupport_normed (μ := volume)).convolutionExists_left
    (L := ContinuousLinearMap.lsmul ℝ ℝ) φ.continuous_normed hf_loc
  have hce_g := (φ.hasCompactSupport_normed (μ := volume)).convolutionExists_left
    (L := ContinuousLinearMap.lsmul ℝ ℝ) φ.continuous_normed hg_loc
  have hf_int :
      Integrable (fun t => φ.normed volume t * f (x - t)) volume := by
    exact (hce_f x).congr <| Eventually.of_forall fun t => by
      simp [ContinuousLinearMap.lsmul_apply, smul_eq_mul]
  have hg_int :
      Integrable (fun t => φ.normed volume t * g (x - t)) volume := by
    exact (hce_g x).congr <| Eventually.of_forall fun t => by
      simp [ContinuousLinearMap.lsmul_apply, smul_eq_mul]
  simp only [MeasureTheory.convolution_def, ContinuousLinearMap.lsmul_apply, smul_eq_mul]
  rw [← integral_sub hf_int hg_int]
  congr 1
  ext t
  ring

omit [NeZero d] in
private theorem eLpNorm_normedConvolution_sub_le
    (φ : ContDiffBump (0 : E))
    {p : ℝ} (hp : 1 ≤ p) {f g : E → ℝ}
    (hf : MemLp f (ENNReal.ofReal p) volume)
    (hg : MemLp g (ENNReal.ofReal p) volume) :
    eLpNorm
      (fun x =>
        ((φ.normed volume) ⋆[ContinuousLinearMap.lsmul ℝ ℝ, volume] f) x -
          ((φ.normed volume) ⋆[ContinuousLinearMap.lsmul ℝ ℝ, volume] g) x)
      (ENNReal.ofReal p) volume ≤
      eLpNorm (fun x => f x - g x) (ENNReal.ofReal p) volume := by
  rw [normedConvolution_sub_eq (d := d) φ
    (hf.locallyIntegrable (by
      simpa using (ENNReal.ofReal_le_ofReal hp : ENNReal.ofReal (1 : ℝ) ≤ ENNReal.ofReal p)))
    (hg.locallyIntegrable (by
      simpa using (ENNReal.ofReal_le_ofReal hp : ENNReal.ofReal (1 : ℝ) ≤ ENNReal.ofReal p)))]
  exact eLpNorm_normedConvolution_le (d := d) φ hp (hf.sub hg)

\end{lstlisting}
\begin{proof}
Linearity of $f \mapsto \rho * f$ on locally integrable functions yields the
equality, after which Proposition~\ref{prop:Lp_contraction_conv} applies to
$f - g$.
\end{proof}

\begin{lemma}[Convolution by a shrinking bump approximates a continuous function;
\leanname{DeGiorgi.tendsto\_eLpNorm\_normedConvolution\_sub\_of\_continuous} (private)]\label{lem:conv_approx_cont}
Let $g \in C_c(E)$ and $1 \le p < \infty$. Then
\begin{equation*}
  \|\rho_n * g - g\|_{\Lp{p}(E)} \;\to\; 0 \quad\text{as } n \to \infty.
\end{equation*}
\end{lemma}

\begin{lstlisting}[style=Matlab-editor,basicstyle=\mlttfamily,escapechar=",mlshowsectionrules=true,mathescape=true,morecomment={[s]{\%\{}{\%\}}},language=lean,numbers=none]
omit [NeZero d] in
private theorem tendsto_eLpNorm_normedConvolution_sub_of_continuous
    {p : ℝ} (hp : 1 ≤ p) {g : E → ℝ}
    (hg_cont : Continuous g) (hg_comp : HasCompactSupport g) :
    Tendsto
      (fun n =>
        eLpNorm
          (fun x =>
            ((shrinkingBump (d := d) n).normed volume ⋆[ContinuousLinearMap.lsmul ℝ ℝ, volume] g) x -
              g x)
          (ENNReal.ofReal p) volume)
      atTop (nhds 0) := by
  let _ := hp
  let S : Set E := Metric.cthickening 1 (tsupport g)
  let T : Set E := Metric.cthickening (1 + 1) (tsupport g)
  have hS_comp : IsCompact S := hg_comp.isCompact.cthickening (r := 1)
  have hT_comp : IsCompact T := hg_comp.isCompact.cthickening (r := 1 + 1)
  have hS_meas : MeasurableSet S := hS_comp.measurableSet
  have hS_finite : volume S < ∞ := hS_comp.measure_lt_top
  have hS_sub_T : S ⊆ T := by
    simpa [S, T] using (cthickening_mono (show (1 : ℝ) ≤ 1 + 1 by norm_num) (tsupport g))
  have hg_support_sub : Function.support g ⊆ S := by
    intro x hx
    exact Metric.mem_cthickening_of_dist_le (x := x) (y := x) (δ := 1)
      (tsupport g) (subset_tsupport g hx) (by simp)
  have huc : UniformContinuousOn g T := hT_comp.uniformContinuousOn_of_continuous hg_cont.continuousOn
  refine ENNReal.tendsto_nhds_zero.2 ?_
  intro ε hε
  by_cases hε_top : ε = ∞
  · filter_upwards with n
    simp [hε_top]
  let εr : ℝ := ε.toReal
  have hεr_pos : 0 < εr := ENNReal.toReal_pos hε.ne' hε_top
  let M : ℝ≥0∞ := volume S ^ ((1 : ℝ) / (ENNReal.ofReal p).toReal)
  have hM_top : M ≠ ∞ := by
    simp [M, hS_finite.ne]
  let C : ℝ := M.toReal + 1
  have hC_pos : 0 < C := by
    have hM_nonneg : 0 ≤ M.toReal := ENNReal.toReal_nonneg
    dsimp [C]
    linarith
  let ε0 : ℝ := εr / C
  have hε0_pos : 0 < ε0 := by
    dsimp [ε0]
    positivity
  rcases Metric.uniformContinuousOn_iff.mp huc ε0 hε0_pos with ⟨δ0, hδ0_pos, hδ0⟩
  let δ : ℝ := min δ0 1
  have hδ_pos : 0 < δ := by
    dsimp [δ]
    exact lt_min hδ0_pos zero_lt_one
  have hsmall :
      ∀ᶠ n in atTop, (shrinkingBump (d := d) n).rOut < δ := by
    exact (tendsto_order.1 tendsto_shrinkingBump_rOut).2 _ hδ_pos
  filter_upwards [hsmall] with n hn
  have hrOut_le_one : (shrinkingBump (d := d) n).rOut ≤ 1 := by
    exact le_trans hn.le (min_le_right _ _)
  have hconv_support_sub :
      Function.support
          (((shrinkingBump (d := d) n).normed volume) ⋆[ContinuousLinearMap.lsmul ℝ ℝ, volume] g)
        ⊆ S := by
    intro x hx
    have hts :
        tsupport
            (((shrinkingBump (d := d) n).normed volume ⋆[ContinuousLinearMap.lsmul ℝ ℝ, volume] g))
          ⊆ S := by
      refine (tsupport_normed_convolution_subset_cthickening (d := d) (shrinkingBump (d := d) n) hg_comp).trans ?_
      simpa [S] using cthickening_mono hrOut_le_one (tsupport g)
    exact hts (subset_tsupport _ hx)
  have hdist :
      ∀ x : E,
        dist
            (((shrinkingBump (d := d) n).normed volume ⋆[ContinuousLinearMap.lsmul ℝ ℝ, volume] g) x)
            (g x) ≤ ε0 := by
    intro x
    by_cases hx : x ∈ S
    · have hxT : x ∈ T := hS_sub_T hx
      refine (shrinkingBump (d := d) n).dist_normed_convolution_le hg_cont.aestronglyMeasurable ?_
      intro y hy
      have hyT : y ∈ T := by
        have hyS :
            y ∈ Metric.cthickening 1 S := by
          refine Metric.mem_cthickening_of_dist_le (x := y) (y := x) (δ := 1) S hx ?_
          exact le_trans hy.le hrOut_le_one
        have hyT' : y ∈ Metric.cthickening 1 (Metric.cthickening 1 (tsupport g)) := by
          simpa [S] using hyS
        have hyT'' : y ∈ Metric.cthickening (1 + 1) (tsupport g) := by
          simpa [S, cthickening_cthickening (show (0 : ℝ) ≤ 1 by positivity)
            (show (0 : ℝ) ≤ 1 by positivity) (tsupport g)] using hyT'
        simpa [T] using hyT''
      exact (hδ0 y hyT x hxT (lt_of_lt_of_le hy (le_trans hn.le (min_le_left _ _)))).le
    · have hconvx :
          (((shrinkingBump (d := d) n).normed volume ⋆[ContinuousLinearMap.lsmul ℝ ℝ, volume] g) x) = 0 := by
        exact zero_outside_of_tsupport_subset
          (Ω := S)
          ((tsupport_normed_convolution_subset_cthickening (d := d) (shrinkingBump (d := d) n) hg_comp).trans
            (by simpa [S] using cthickening_mono hrOut_le_one (tsupport g)))
          hx
      have hgx : g x = 0 := zero_outside_of_tsupport_subset (Ω := S)
        (show tsupport g ⊆ S by
          intro y hy
          exact Metric.mem_cthickening_of_dist_le (x := y) (y := y) (δ := 1)
            (tsupport g) hy (by simp))
        hx
      simp [hconvx, hgx, hε0_pos.le]
  have hbound :
      eLpNorm
        (fun x =>
          ((shrinkingBump (d := d) n).normed volume ⋆[ContinuousLinearMap.lsmul ℝ ℝ, volume] g) x -
            g x)
        (ENNReal.ofReal p) volume ≤ ENNReal.ofReal ε0 * M := by
    refine eLpNorm_sub_le_of_dist_bdd
      (μ := volume) (p := ENNReal.ofReal p) (s := S) (hp := by simp) hS_meas hε0_pos.le hdist ?_ ?_
    · intro x hx
      exact hconv_support_sub hx
    · exact hg_support_sub
  have hfinal :
      ENNReal.ofReal ε0 * M < ε := by
    have hM_eq : M = ENNReal.ofReal (M.toReal) := by
      exact (ENNReal.ofReal_toReal hM_top).symm
    have hlt_real : ε0 * M.toReal < εr := by
      dsimp [ε0, C]
      have hfrac_lt_one : M.toReal / (M.toReal + 1) < 1 := by
        have hM_nonneg : 0 ≤ M.toReal := ENNReal.toReal_nonneg
        have hden_pos : 0 < M.toReal + 1 := by linarith
        have hnum_lt : M.toReal < M.toReal + 1 := by linarith
        exact (div_lt_one hden_pos).2 hnum_lt
      have hmul_lt : εr * (M.toReal / (M.toReal + 1)) < εr := by
        simpa [mul_comm, mul_left_comm, mul_assoc] using
          (mul_lt_mul_of_pos_left hfrac_lt_one hεr_pos)
      have hcalc :
          εr / (M.toReal + 1) * M.toReal = εr * (M.toReal / (M.toReal + 1)) := by
        field_simp [show M.toReal + 1 ≠ 0 by linarith]
      simpa [hcalc] using hmul_lt
    have hlt_ofReal :
        ENNReal.ofReal ε0 * M < ENNReal.ofReal εr := by
      rw [hM_eq, ← ENNReal.ofReal_mul (show 0 ≤ ε0 by positivity),
        ENNReal.ofReal_lt_ofReal_iff hεr_pos]
      exact hlt_real
    simpa [εr, hε_top] using hlt_ofReal
  exact hbound.trans hfinal.le

\end{lstlisting}
\begin{proof}
Choose $S := \{x : \dist(x, \supp g) \le 1\}$ and $T := \{x : \dist(x, \supp g)
\le 2\}$, both compact, with $S \subseteq T$. The function $g$ is uniformly
continuous on the compact set $T$. Fix $\varepsilon > 0$ and let $\varepsilon_0
:= \varepsilon / (\mathrm{vol}(S)^{1/p} + 1)$. By uniform continuity, there
exists $\delta_0 > 0$ such that $|g(y) - g(x)| < \varepsilon_0$ whenever
$x, y \in T$ with $\dist(x, y) < \delta_0$. Set $\delta := \min(\delta_0, 1)$.

For $n$ large enough that $r_{\text{out}}^{(n)} < \delta$, we have
$\supp(\rho_n * g) \subseteq S$ (Lemma~\ref{lem:tsupport_conv}, since
$r_{\text{out}}^{(n)} \le 1$). Pointwise, for $x \in S$,
\begin{equation*}
  |(\rho_n * g)(x) - g(x)|
    = \left|\int_E \rho_n(y) (g(x - y) - g(x))\,dy\right|
    \le \varepsilon_0,
\end{equation*}
since $\rho_n$ is supported in $\Bz{r_{\text{out}}^{(n)}} \subseteq \Bz{\delta}$
and both $x, x - y \in T$ for $y \in \supp\rho_n$. For $x \notin S$ both
$(\rho_n * g)(x)$ and $g(x)$ vanish.

The bound $\|\rho_n * g - g\|_{\Lp{p}(E)} \le \varepsilon_0
\,\mathrm{vol}(S)^{1/p}$ then follows from
\leanname{eLpNorm\_sub\_le\_of\_dist\_bdd}, and the choice of $\varepsilon_0$
gives this is less than $\varepsilon$.
\end{proof}

\begin{theorem}[Convolution by a shrinking bump approximates an $\Lp{p}$ function;
\leanname{DeGiorgi.tendsto\_eLpNorm\_normedConvolution\_sub} (private)]\label{thm:conv_approx_Lp}
Let $1 \le p < \infty$ and $f \in \Lp{p}(E)$. Then
\begin{equation*}
  \|\rho_n * f - f\|_{\Lp{p}(E)} \;\to\; 0 \quad\text{as } n \to \infty.
\end{equation*}
\end{theorem}

\begin{proof}
By density of smooth compactly supported functions in $\Lp{p}$ (Mathlib's
\leanname{MemLp.exist\_eLpNorm\_sub\_le}), for any $\varepsilon > 0$ there
exists $g \in C_c^\infty(E)$ with $\|f - g\|_{\Lp{p}} < \varepsilon/3$. Then by
the triangle inequality
\begin{equation*}
  \|\rho_n * f - f\|_{\Lp{p}} \le
  \|\rho_n * f - \rho_n * g\|_{\Lp{p}} + \|\rho_n * g - g\|_{\Lp{p}} +
  \|g - f\|_{\Lp{p}}.
\end{equation*}
By Corollary~\ref{cor:conv_sub} the first term is at most $\|f - g\|_{\Lp{p}} <
\varepsilon/3$. By Lemma~\ref{lem:conv_approx_cont} the middle term tends to
$0$, so for $n$ large it is less than $\varepsilon/3$. The last term is also
less than $\varepsilon/3$. Hence $\|\rho_n * f - f\|_{\Lp{p}} < \varepsilon$
for all sufficiently large $n$, proving convergence to zero.
\end{proof}

\subsection{Uniqueness of the weak gradient and global smooth approximation}

\begin{theorem}[Uniqueness of the weak gradient;
\leanname{DeGiorgi.MemW1pWitness.ae\_eq\_p}, \leanname{DeGiorgi.MemW1pWitness.ae\_eq}]\label{thm:weak_grad_unique}
Let $\Omega \subseteq E$ be open, $1 \le p < \infty$, and $f$ a function with
two $\Wkp{1}{p}$-witnesses on $\Omega$ giving rise to weak gradient fields
$G_1, G_2$. Then $G_1 = G_2$ a.e.\ on $\Omega$. The same holds for $H^1$
witnesses ($p = 2$).
\end{theorem}

\begin{lstlisting}[style=Matlab-editor,basicstyle=\mlttfamily,escapechar=",mlshowsectionrules=true,mathescape=true,morecomment={[s]{\%\{}{\%\}}},language=lean,numbers=none]
/-- Weak gradients of the same `W^{1,p}` function agree a.e. on the source open
set. -/
theorem MemW1pWitness.ae_eq_p
    {Ω : Set E} (hΩ : IsOpen Ω) {p : ℝ} (hp : 1 ≤ p) {f : E → ℝ}
    (hw₁ hw₂ : MemW1pWitness (ENNReal.ofReal p) f Ω) :
    hw₁.weakGrad =ᵐ[volume.restrict Ω] hw₂.weakGrad := by
  let _ := (inferInstance : NeZero d)
  have hp_enn : (1 : ℝ≥0∞) ≤ ENNReal.ofReal p := by
    simpa using (ENNReal.ofReal_le_ofReal hp : ENNReal.ofReal (1 : ℝ) ≤ ENNReal.ofReal p)
  have hcomp :
      ∀ i : Fin d,
        (fun x => hw₁.weakGrad x i) =ᵐ[volume.restrict Ω] (fun x => hw₂.weakGrad x i) := by
    intro i
    exact HasWeakPartialDeriv.ae_eq hΩ (hw₁.isWeakGrad i) (hw₂.isWeakGrad i)
      ((hw₁.weakGrad_component_memLp i).locallyIntegrable hp_enn)
      ((hw₂.weakGrad_component_memLp i).locallyIntegrable hp_enn)
  filter_upwards [ae_all_iff.2 hcomp] with x hx
  ext i
  exact hx i


\end{lstlisting}
\begin{proof}
Apply Lemma~\ref{lem:weak_partial_unique} componentwise to $(G_1)_i$ and
$(G_2)_i$, then assemble: a.e.\ all components agree, so the vectors
themselves agree a.e.
\end{proof}

\begin{lemma}[Convolution differentiates against the weak partial derivative;
\leanname{DeGiorgi.convolution\_fderiv\_eq\_convolution\_weakPartial\_univ} (private)]\label{lem:conv_fderiv_weak_partial}
Let $u \in L^1_{loc}(E)$ and let $g$ be the weak $i$-th partial derivative of
$u$ on $E$. Then for every $\varphi \in C_c^\infty(E)$ and every $x \in E$,
\begin{equation*}
  \bigl((\partial_i\varphi) * u\bigr)(x) = (\varphi * g)(x).
\end{equation*}
\end{lemma}

\begin{lstlisting}[style=Matlab-editor,basicstyle=\mlttfamily,escapechar=",mlshowsectionrules=true,mathescape=true,morecomment={[s]{\%\{}{\%\}}},language=lean,numbers=none]
omit [NeZero d] in
private theorem convolution_fderiv_eq_convolution_weakPartial_univ
    {u g φ : E → ℝ} {i : Fin d}
    (hweak : HasWeakPartialDeriv i g u Set.univ)
    (hφ_smooth : ContDiff ℝ (⊤ : ℕ∞) φ)
    (hφ_compact : HasCompactSupport φ) :
    ∀ x,
      ((fun y => (fderiv ℝ φ y) (EuclideanSpace.single i 1))
          ⋆[ContinuousLinearMap.lsmul ℝ ℝ, volume] u) x =
        (φ ⋆[ContinuousLinearMap.lsmul ℝ ℝ, volume] g) x := by
  intro x
  let T : Homeomorph E E := (Homeomorph.neg E).trans (Homeomorph.addLeft x)
  let ψ : E → ℝ := φ ∘ T
  have hψ_smooth : ContDiff ℝ (⊤ : ℕ∞) ψ := by
    simpa [ψ, T, Function.comp, sub_eq_add_neg, add_comm, add_left_comm, add_assoc] using
      hφ_smooth.comp (contDiff_const.add contDiff_id.neg)
  have hψ_compact : HasCompactSupport ψ := by
    simpa [ψ, T, Function.comp] using hφ_compact.comp_homeomorph T
  have key := hweak ψ hψ_smooth hψ_compact (by simp)
  have hderiv :
      ∀ t,
        (fderiv ℝ ψ t) (EuclideanSpace.single i 1) =
          - (fderiv ℝ φ (x + -t)) (EuclideanSpace.single i 1) := by
    intro t
    have hfd :
        HasFDerivAt ψ
          (((fderiv ℝ φ (x + -t)).comp ((0 : E →L[ℝ] E) - 1)) : E →L[ℝ] ℝ) t := by
      simpa [ψ, T, Function.comp, sub_eq_add_neg, add_comm, add_left_comm, add_assoc] using
        ((hφ_smooth.differentiable
            (show (((⊤ : ℕ∞) : WithTop ℕ∞)) ≠ 0 by simp) (x + -t)).hasFDerivAt.comp t
          ((hasFDerivAt_const x t).add (hasFDerivAt_id t).neg))
    calc
      (fderiv ℝ ψ t) (EuclideanSpace.single i 1)
        = (((fderiv ℝ φ (x + -t)).comp ((0 : E →L[ℝ] E) - 1))
            (EuclideanSpace.single i 1)) := by
              rw [hfd.fderiv]
      _ = (fderiv ℝ φ (x + -t))
            (((0 : E →L[ℝ] E) - 1) (EuclideanSpace.single i 1)) := by
              rfl
      _ = (fderiv ℝ φ (x + -t)) (- EuclideanSpace.single i 1) := by
            simp
      _ = - (fderiv ℝ φ (x + -t)) (EuclideanSpace.single i 1) := by
            simp
  have hkey :
      ∫ t, g t * ψ t ∂volume =
        ∫ t, u t * (fderiv ℝ φ (x + -t)) (EuclideanSpace.single i 1) ∂volume := by
    have hkey' :
        ∫ t, u t * (fderiv ℝ ψ t) (EuclideanSpace.single i 1) ∂volume =
          -∫ t, g t * ψ t ∂volume := by
      simpa [Measure.restrict_univ] using key
    have hderiv_int :
        ∫ t, u t * (fderiv ℝ ψ t) (EuclideanSpace.single i 1) ∂volume =
          -∫ t, u t * (fderiv ℝ φ (x + -t)) (EuclideanSpace.single i 1) ∂volume := by
      calc
        ∫ t, u t * (fderiv ℝ ψ t) (EuclideanSpace.single i 1) ∂volume
          = ∫ t, -(u t * (fderiv ℝ φ (x + -t)) (EuclideanSpace.single i 1)) ∂volume := by
              refine integral_congr_ae ?_
              filter_upwards with t
              rw [hderiv t]
              ring
        _ = -∫ t, u t * (fderiv ℝ φ (x + -t)) (EuclideanSpace.single i 1) ∂volume := by
              rw [integral_neg]
    linarith
  calc
      ((fun y => (fderiv ℝ φ y) (EuclideanSpace.single i 1))
        ⋆[ContinuousLinearMap.lsmul ℝ ℝ, volume] u) x
      = ∫ t, u t * (fderiv ℝ φ (x - t)) (EuclideanSpace.single i 1) ∂volume := by
          simpa [smul_eq_mul, mul_comm] using
            (MeasureTheory.convolution_lsmul_swap
              (f := fun y => (fderiv ℝ φ y) (EuclideanSpace.single i 1)) (g := u) (x := x)
              (μ := volume))
    _ = ∫ t, g t * ψ t ∂volume := hkey.symm
    _ = (φ ⋆[ContinuousLinearMap.lsmul ℝ ℝ, volume] g) x := by
          simpa [ψ, T, Function.comp, sub_eq_add_neg, add_comm, add_left_comm, add_assoc,
            smul_eq_mul, mul_comm] using
            (MeasureTheory.convolution_lsmul_swap (f := φ) (g := g) (x := x) (μ := volume)).symm

\end{lstlisting}
\begin{proof}
For fixed $x$, define the test function $\psi(t) := \varphi(x - t)$, which is
smooth and compactly supported. A direct computation shows
$(\partial_i \psi)(t) = -(\partial_i\varphi)(x - t)$ (chain rule with the
sign-flipping translation). The weak-derivative identity for $u$ tested against
$\psi$ then reads
\begin{equation*}
  \int_E u(t)\,(\partial_i\psi)(t)\,dt = -\int_E g(t)\,\psi(t)\,dt,
\end{equation*}
i.e.,
\begin{equation*}
  -\int_E u(t)\,(\partial_i\varphi)(x - t)\,dt = -\int_E g(t)\,\varphi(x - t)\,dt.
\end{equation*}
Cancelling the minus signs and rewriting both sides as convolutions yields the
claim.
\end{proof}

\begin{theorem}[Smooth compactly supported approximation on $E$ for compactly
supported $\Wkp{1}{p}$ functions;
\leanname{DeGiorgi.exists\_smooth\_compactSupport\_W1p\_approx\_univ}]\label{thm:smooth_approx_univ}
Let $1 < p < \infty$, and let $u : E \to \R$ be compactly supported with a
$\Wkp{1}{p}$-witness on $E$ with weak gradient $G$. Then there exists a sequence
$(\varphi_n)$ of smooth compactly supported functions on $E$ with
\begin{align*}
  \supp \varphi_n &\subseteq \{x : \dist(x, \supp u) \le r_{\text{out}}^{(n)}\}, \\
  \|\varphi_n - u\|_{\Lp{p}(E)} &\to 0, \\
  \|\partial_i\varphi_n - G_i\|_{\Lp{p}(E)} &\to 0 \quad\text{for every } i.
\end{align*}
\end{theorem}

\begin{lstlisting}[style=Matlab-editor,basicstyle=\mlttfamily,escapechar=",mlshowsectionrules=true,mathescape=true,morecomment={[s]{\%\{}{\%\}}},language=lean,numbers=none]
/-- Global smooth compactly supported approximation of a compactly supported
`W^{1,p}` function on `ℝ^d`, in the finite-`p` regime. -/
theorem exists_smooth_compactSupport_W1p_approx_univ
    {p : ℝ} (hp : 1 < p)
    {u : E → ℝ}
    (hw : MemW1pWitness (ENNReal.ofReal p) u Set.univ)
    (hu_compact : HasCompactSupport u) :
    ∃ φ : ℕ → E → ℝ,
      (∀ n, ContDiff ℝ (⊤ : ℕ∞) (φ n)) ∧
      (∀ n, HasCompactSupport (φ n)) ∧
      (∀ n,
        tsupport (φ n) ⊆
          Metric.cthickening (shrinkingBump (d := d) n).rOut (tsupport u)) ∧
      Tendsto
        (fun n => eLpNorm (fun x => φ n x - u x) (ENNReal.ofReal p) volume)
        atTop (nhds 0) ∧
      (∀ i : Fin d,
        Tendsto
          (fun n =>
            eLpNorm
              (fun x => (fderiv ℝ (φ n) x) (EuclideanSpace.single i 1) - hw.weakGrad x i)
              (ENNReal.ofReal p) volume)
          atTop (nhds 0)) := by
  let _ := (inferInstance : NeZero d)
  have hp_le : 1 ≤ p := le_of_lt hp
  have hp_enn : (1 : ℝ≥0∞) ≤ ENNReal.ofReal p := by
    simpa using (ENNReal.ofReal_le_ofReal hp_le : ENNReal.ofReal (1 : ℝ) ≤ ENNReal.ofReal p)
  have hu_memLp : MemLp u (ENNReal.ofReal p) volume := by
    simpa [Measure.restrict_univ] using hw.memLp
  let φ : ℕ → E → ℝ := fun n =>
    ((shrinkingBump (d := d) n).normed volume
      ⋆[ContinuousLinearMap.lsmul ℝ ℝ, volume] u)
  refine ⟨φ, ?_, ?_, ?_, ?_, ?_⟩
  · intro n
    exact (shrinkingBump (d := d) n).hasCompactSupport_normed.contDiff_convolution_left
      (L := ContinuousLinearMap.lsmul ℝ ℝ)
      ((shrinkingBump (d := d) n).contDiff_normed)
      (hu_memLp.locallyIntegrable hp_enn)
  · intro n
    apply HasCompactSupport.intro'
      (K := Metric.cthickening (shrinkingBump (d := d) n).rOut (tsupport u))
      (hu_compact.isCompact.cthickening (r := (shrinkingBump (d := d) n).rOut))
      (hu_compact.isCompact.cthickening (r := (shrinkingBump (d := d) n).rOut)).isClosed
    intro x hx
    exact zero_outside_of_tsupport_subset
      (Ω := Metric.cthickening (shrinkingBump (d := d) n).rOut (tsupport u))
      (tsupport_normed_convolution_subset_cthickening
        (d := d) (shrinkingBump (d := d) n) hu_compact)
      hx
  · intro n
    exact tsupport_normed_convolution_subset_cthickening
      (d := d) (shrinkingBump (d := d) n) hu_compact
  · simpa [φ] using
      tendsto_eLpNorm_normedConvolution_sub (d := d) hp_le hu_memLp
  · intro i
    have hGi_memLp : MemLp (fun y => hw.weakGrad y i) (ENNReal.ofReal p) volume := by
      simpa [Measure.restrict_univ] using hw.weakGrad_component_memLp i
    have hEq :
        (fun n =>
          eLpNorm
            (fun x => (fderiv ℝ (φ n) x) (EuclideanSpace.single i 1) - hw.weakGrad x i)
            (ENNReal.ofReal p) volume) =
          fun n =>
            eLpNorm
              (fun x =>
                ((shrinkingBump (d := d) n).normed volume
                  ⋆[ContinuousLinearMap.lsmul ℝ ℝ, volume] (fun y => hw.weakGrad y i)) x -
                  hw.weakGrad x i)
              (ENNReal.ofReal p) volume := by
      funext n
      refine eLpNorm_congr_ae ?_
      filter_upwards with x
      let bump : E → ℝ := (shrinkingBump (d := d) n).normed volume
      let Dbump : E → E →L[ℝ] ℝ := fderiv ℝ bump
      let ei : E := EuclideanSpace.single i 1
      have hbump1 : ContDiff ℝ 1 bump := by
        simpa [bump] using
          ((shrinkingBump (d := d) n).contDiff_normed : ContDiff ℝ 1 ((shrinkingBump (d := d) n).normed volume))
      have hfd :=
        (shrinkingBump (d := d) n).hasCompactSupport_normed.hasFDerivAt_convolution_left
          (L := ContinuousLinearMap.lsmul ℝ ℝ)
          hbump1
          (hu_memLp.locallyIntegrable hp_enn) x
      have hDbump_cont : Continuous Dbump := by
        simpa [Dbump, bump] using hbump1.continuous_fderiv one_ne_zero
      have hDbump_compact : HasCompactSupport Dbump := by
        simpa [Dbump, bump] using ((shrinkingBump (d := d) n).hasCompactSupport_normed.fderiv (𝕜 := ℝ))
      have hfd' :
          (fderiv ℝ (φ n) x) ei =
            (((Dbump ⋆[
                ContinuousLinearMap.precompL E (ContinuousLinearMap.lsmul ℝ ℝ), volume] u) x)
                ei) := by
        simpa [φ, bump, Dbump, ei] using
          congrArg (fun A : E →L[ℝ] ℝ => A ei) hfd.fderiv
      have hconv_apply :
          (((Dbump ⋆[
              ContinuousLinearMap.precompL E (ContinuousLinearMap.lsmul ℝ ℝ), volume] u) x) ei) =
            (((fun y => (fderiv ℝ bump y) ei)
              ⋆[ContinuousLinearMap.lsmul ℝ ℝ, volume] u) x) := by
        calc
          (((Dbump ⋆[
              ContinuousLinearMap.precompL E (ContinuousLinearMap.lsmul ℝ ℝ), volume] u) x) ei)
              =
            ((u ⋆[
              ContinuousLinearMap.precompR E (ContinuousLinearMap.lsmul ℝ ℝ).flip, volume] Dbump) x) ei := by
                simpa [ContinuousLinearMap.precompL, Dbump] using
                  congrArg (fun F => F x ei)
                    (MeasureTheory.convolution_flip
                      (L := ContinuousLinearMap.precompR E (ContinuousLinearMap.lsmul ℝ ℝ).flip)
                      (μ := volume) (f := u) (g := Dbump))
          _ = (u ⋆[(ContinuousLinearMap.lsmul ℝ ℝ).flip, volume] (fun y => Dbump y ei)) x := by
                simpa [Dbump, bump] using
                  (MeasureTheory.convolution_precompR_apply
                    (L := (ContinuousLinearMap.lsmul ℝ ℝ).flip)
                    (μ := volume) (f := u) (g := Dbump)
                    (hf := hu_memLp.locallyIntegrable hp_enn)
                    (hcg := hDbump_compact) (hg := hDbump_cont) (x₀ := x) (x := ei))
          _ = (((fun y => Dbump y ei) ⋆[ContinuousLinearMap.lsmul ℝ ℝ, volume] u) x) := by
                simpa [smul_eq_mul, mul_comm] using
                  congrArg (fun F => F x)
                    ((MeasureTheory.convolution_flip
                      (L := (ContinuousLinearMap.lsmul ℝ ℝ).flip)
                      (μ := volume) (f := u) (g := fun y => Dbump y ei)).symm)
      calc
        (fderiv ℝ (φ n) x) ei - hw.weakGrad x i
          = (((Dbump ⋆[
                ContinuousLinearMap.precompL E (ContinuousLinearMap.lsmul ℝ ℝ), volume] u) x)
                ei) - hw.weakGrad x i := by
                  rw [hfd']
        _ = (((fun y =>
              (fderiv ℝ bump y) ei)
                ⋆[ContinuousLinearMap.lsmul ℝ ℝ, volume] u) x) - hw.weakGrad x i := by
                  rw [hconv_apply]
        _ = ((bump)
              ⋆[ContinuousLinearMap.lsmul ℝ ℝ, volume] (fun y => hw.weakGrad y i)) x -
                hw.weakGrad x i := by
                  simpa [ei] using
                    congrArg (fun z => z - hw.weakGrad x i)
                      (convolution_fderiv_eq_convolution_weakPartial_univ
                        (d := d) (i := i) (u := u) (g := fun y => hw.weakGrad y i)
                        (φ := bump)
                        (hw.isWeakGrad i)
                        (by
                          simpa [bump] using
                            ((shrinkingBump (d := d) n).contDiff_normed :
                              ContDiff ℝ (⊤ : ℕ∞) ((shrinkingBump (d := d) n).normed volume)))
                        ((shrinkingBump (d := d) n).hasCompactSupport_normed) x)
    exact hEq ▸ tendsto_eLpNorm_normedConvolution_sub (d := d) hp_le hGi_memLp

\end{lstlisting}
\begin{proof}
Define $\varphi_n := \rho_n * u$. Since $u$ is compactly supported and $\rho_n
\in C_c^\infty(E)$, $\varphi_n$ is smooth and has support contained in the
$r_{\text{out}}^{(n)}$-thickening of $\supp u$ (Lemma~\ref{lem:tsupport_conv}).
The $\Lp{p}$ approximation $\varphi_n \to u$ is exactly
Theorem~\ref{thm:conv_approx_Lp} applied to $u$.

For the gradients, the classical fact
\[
  \partial_i (\rho_n * u) = (\partial_i \rho_n) * u
\]
(Mathlib's
\leanname{HasCompactSupport.hasFDerivAt\_convolution\_left})
combined with Lemma~\ref{lem:conv_fderiv_weak_partial} (with $u$ on the right
and the smooth $\rho_n$ on the left) yields
\begin{equation*}
  \partial_i \varphi_n = (\partial_i\rho_n) * u = \rho_n * G_i,
\end{equation*}
since $G_i$ is the weak partial derivative of $u$ on $E$. Applying
Theorem~\ref{thm:conv_approx_Lp} to $G_i \in \Lp{p}(E)$ then gives
$\partial_i\varphi_n - G_i = \rho_n * G_i - G_i \to 0$ in $\Lp{p}(E)$.
\end{proof}

\begin{theorem}[Smooth approximation for a witness supported in a compact subset;
\leanname{DeGiorgi.exists\_smooth\_W1p\_approx\_of\_supportedWitness}]\label{thm:smooth_approx_local}
Let $\Omega$ be open, $1 < p < \infty$, and let $u$ have a $\Wkp{1}{p}$-witness
on $\Omega$ with gradient field $G$. Suppose $K \subseteq \Omega$ is compact
with $\supp u \subseteq K$ and $\supp G_i \subseteq K$ for every $i$. Then
there exists a sequence $(\varphi_n)$ of smooth compactly supported functions
with $\supp\varphi_n \subseteq \Omega$ such that
\begin{align*}
  \|\varphi_n - u\|_{\Lp{p}(\Omega)} &\to 0, \\
  \|\partial_i\varphi_n - G_i\|_{\Lp{p}(\Omega)} &\to 0 \quad\text{for every } i.
\end{align*}
\end{theorem}

\begin{lstlisting}[style=Matlab-editor,basicstyle=\mlttfamily,escapechar=",mlshowsectionrules=true,mathescape=true,morecomment={[s]{\%\{}{\%\}}},language=lean,numbers=none]
theorem exists_smooth_W1p_approx_of_supportedWitness
    {Ω K : Set E} (hΩ : IsOpen Ω)
    {p : ℝ} (hp : 1 < p)
    {u : E → ℝ}
    (hw : MemW1pWitness (ENNReal.ofReal p) u Ω)
    (hK_compact : IsCompact K) (hKΩ : K ⊆ Ω)
    (hu_sub : tsupport u ⊆ K)
    (hgrad_sub : ∀ i : Fin d, tsupport (fun x => hw.weakGrad x i) ⊆ K) :
    ∃ φ : ℕ → E → ℝ,
      (∀ n, ContDiff ℝ (⊤ : ℕ∞) (φ n)) ∧
      (∀ n, HasCompactSupport (φ n)) ∧
      (∀ n, tsupport (φ n) ⊆ Ω) ∧
      Tendsto
        (fun n => eLpNorm (fun x => φ n x - u x) (ENNReal.ofReal p) (volume.restrict Ω))
        atTop (nhds 0) ∧
      (∀ i : Fin d,
        Tendsto
          (fun n =>
            eLpNorm
              (fun x => (fderiv ℝ (φ n) x) (EuclideanSpace.single i 1) - hw.weakGrad x i)
              (ENNReal.ofReal p) (volume.restrict Ω))
          atTop (nhds 0)) := by
  have hp_le : 1 ≤ p := le_of_lt hp
  have hK_compact_u : HasCompactSupport u := by
    apply HasCompactSupport.intro' (K := K) hK_compact hK_compact.isClosed
    intro x hx
    exact zero_outside_of_tsupport_subset (Ω := K) hu_sub hx
  obtain ⟨δ, hδ_pos, hδΩ⟩ := hK_compact.exists_cthickening_subset_open hΩ hKΩ
  let K' : Set E := Metric.cthickening δ K
  have hK'_compact : IsCompact K' := hK_compact.cthickening (r := δ)
  have hK'_Ω : K' ⊆ Ω := hδΩ
  obtain ⟨χ, hχ_smooth, hχ_compact, _hχ_range, hχ_one, hχ_sub⟩ :=
    exists_smooth_cutoff (d := d) hK'_compact hΩ hK'_Ω
  have hu_memLp_univ : MemLp u (ENNReal.ofReal p) volume := by
    have hu_eq_ind : Ω.indicator u = u := by
      ext x
      by_cases hx : x ∈ Ω
      · simp [hx]
      · simp [hx, zero_outside_of_tsupport_subset (Ω := Ω) (hu_sub.trans hKΩ) hx]
    rw [← hu_eq_ind]
    exact (MeasureTheory.memLp_indicator_iff_restrict
      (μ := volume) (s := Ω) (f := u) (p := ENNReal.ofReal p) hΩ.measurableSet).2 hw.memLp
  have hgrad_memLp_univ :
      ∀ i : Fin d, MemLp (fun x => hw.weakGrad x i) (ENNReal.ofReal p) volume := by
    intro i
    have hgi_eq_ind : Ω.indicator (fun x => hw.weakGrad x i) = fun x => hw.weakGrad x i := by
      ext x
      by_cases hx : x ∈ Ω
      · simp [hx]
      · simp [hx, zero_outside_of_tsupport_subset (Ω := Ω) (hgrad_sub i |>.trans hKΩ) hx]
    rw [← hgi_eq_ind]
    exact (MeasureTheory.memLp_indicator_iff_restrict
      (μ := volume) (s := Ω) (f := fun x => hw.weakGrad x i) (p := ENNReal.ofReal p)
      hΩ.measurableSet).2 (hw.weakGrad_component_memLp i)
  let hwUniv : MemW1pWitness (ENNReal.ofReal p) u Set.univ := by
    refine
      { memLp := by simpa [Measure.restrict_univ] using hu_memLp_univ
        weakGrad := hw.weakGrad
        weakGrad_component_memLp := by
          intro i
          simpa [Measure.restrict_univ] using hgrad_memLp_univ i
        isWeakGrad := ?_ }
    ·
      intro i φ hφ_smooth hφ_compact _hφ_sub
      let ψ : E → ℝ := fun x => χ x * φ x
      have hψ_smooth : ContDiff ℝ (⊤ : ℕ∞) ψ := hχ_smooth.mul hφ_smooth
      have hψ_compact : HasCompactSupport ψ := by
        simpa [ψ] using hχ_compact.mul_right (f' := φ)
      have hψ_sub : tsupport ψ ⊆ Ω := (tsupport_smul_subset_left χ φ).trans hχ_sub
      have key := hw.isWeakGrad i ψ hψ_smooth hψ_compact hψ_sub
      let ei : E := EuclideanSpace.single i 1
      have hχ_deriv_zero :
          ∀ x ∈ K, (fderiv ℝ χ x) ei = 0 := by
        intro x hx
        exact fderiv_apply_eq_zero_of_eq_one_on_cthickening
          (d := d) hδ_pos hχ_one hx i
      have hleft_pt :
          ∀ x,
            u x * (fderiv ℝ ψ x) ei =
              u x * (fderiv ℝ φ x) ei := by
        intro x
        by_cases hx : x ∈ K
        · have hχx : χ x = 1 := hχ_one x (Metric.self_subset_cthickening K hx)
          have hdχx : (fderiv ℝ χ x) ei = 0 := hχ_deriv_zero x hx
          have hχ_diff : Differentiable ℝ χ := hχ_smooth.differentiable (by simp)
          have hφ_diff : Differentiable ℝ φ := hφ_smooth.differentiable (by simp)
          rw [fderiv_fun_mul hχ_diff.differentiableAt hφ_diff.differentiableAt]
          simp [ei, ContinuousLinearMap.add_apply, ContinuousLinearMap.smul_apply,
            smul_eq_mul, hχx, hdχx]
        · have hux : u x = 0 := zero_outside_of_tsupport_subset (Ω := K) hu_sub hx
          simp [hux]
      have hright_pt :
          ∀ x,
            hw.weakGrad x i * ψ x =
              hw.weakGrad x i * φ x := by
        intro x
        by_cases hx : x ∈ K
        · have hχx : χ x = 1 := hχ_one x (Metric.self_subset_cthickening K hx)
          simp [ψ, hχx]
        · have hgx : hw.weakGrad x i = 0 :=
            zero_outside_of_tsupport_subset (Ω := K) (hgrad_sub i) hx
          simp [ψ, hgx]
      have hψ_deriv_sub : tsupport (fun x => (fderiv ℝ ψ x) ei) ⊆ Ω :=
        (tsupport_fderiv_apply_subset ℝ ei).trans hψ_sub
      have hleft_zero : ∀ x, x ∉ Ω → u x * (fderiv ℝ ψ x) ei = 0 := by
        intro x hx
        have hderivx : (fderiv ℝ ψ x) ei = 0 :=
          image_eq_zero_of_notMem_tsupport
            (f := fun y => (fderiv ℝ ψ y) ei) (fun h => hx (hψ_deriv_sub h))
        simp [hderivx]
      have hright_zero : ∀ x, x ∉ Ω → hw.weakGrad x i * ψ x = 0 := by
        intro x hx
        have hψx : ψ x = 0 := image_eq_zero_of_notMem_tsupport (fun h => hx (hψ_sub h))
        simp [hψx]
      have key_global :
          ∫ x, u x * (fderiv ℝ ψ x) ei ∂volume =
            -∫ x, hw.weakGrad x i * ψ x ∂volume := by
        rw [← setIntegral_eq_integral_of_forall_compl_eq_zero hleft_zero,
          ← setIntegral_eq_integral_of_forall_compl_eq_zero hright_zero]
        simpa [Measure.restrict_univ] using key
      have hleft_global :
          ∫ x, u x * (fderiv ℝ ψ x) ei ∂volume =
            ∫ x, u x * (fderiv ℝ φ x) ei ∂volume := by
        refine integral_congr_ae ?_
        filter_upwards with x
        exact hleft_pt x
      have hright_global :
          ∫ x, hw.weakGrad x i * ψ x ∂volume =
            ∫ x, hw.weakGrad x i * φ x ∂volume := by
        refine integral_congr_ae ?_
        filter_upwards with x
        exact hright_pt x
      rw [hleft_global, hright_global] at key_global
      simpa [Measure.restrict_univ] using key_global
  rcases exists_smooth_compactSupport_W1p_approx_univ (d := d) hp hwUniv hK_compact_u with
    ⟨φ, hφ_smooth, hφ_compact, hφ_sup, hφ_fun, hφ_grad⟩
  have hsmall_lt :
      ∀ᶠ n in atTop, (shrinkingBump (d := d) n).rOut < δ := by
    exact (tendsto_order.1 (tendsto_shrinkingBump_rOut (d := d))).2 _ hδ_pos
  have hsmall :
      ∀ᶠ n in atTop, (shrinkingBump (d := d) n).rOut ≤ δ := by
    exact hsmall_lt.mono fun _ hn => hn.le
  rcases Filter.eventually_atTop.1 hsmall with ⟨N, hN⟩
  let ψ : ℕ → E → ℝ := fun n => φ (n + N)
  refine ⟨ψ, ?_, ?_, ?_, ?_, ?_⟩
  · intro n
    exact hφ_smooth (n + N)
  · intro n
    exact hφ_compact (n + N)
  · intro n
    have hle : (shrinkingBump (d := d) (n + N)).rOut ≤ δ := hN (n + N) (Nat.le_add_left _ _)
    exact (hφ_sup (n + N)).trans <|
      (Metric.cthickening_subset_of_subset (shrinkingBump (d := d) (n + N)).rOut hu_sub).trans <|
        (Metric.cthickening_mono hle K).trans hδΩ
  ·
    have hψ_fun_global :
        Tendsto
          (fun n => eLpNorm (fun x => ψ n x - u x) (ENNReal.ofReal p) volume)
          atTop (nhds 0) := by
      simpa [ψ] using hφ_fun.comp (tendsto_add_atTop_nat N)
    have hψ_fun_nonneg :
        ∀ n,
          (0 : ℝ≥0∞) ≤ eLpNorm (fun x => ψ n x - u x) (ENNReal.ofReal p) (volume.restrict Ω) :=
      fun n => zero_le _
    have hψ_fun_bound :
        ∀ n,
          eLpNorm (fun x => ψ n x - u x) (ENNReal.ofReal p) (volume.restrict Ω) ≤
            eLpNorm (fun x => ψ n x - u x) (ENNReal.ofReal p) volume :=
      fun n =>
        eLpNorm_mono_measure (f := fun x => ψ n x - u x) Measure.restrict_le_self
    exact tendsto_of_tendsto_of_tendsto_of_le_of_le
      tendsto_const_nhds hψ_fun_global hψ_fun_nonneg hψ_fun_bound
  · intro i
    have hψ_grad_global :
        Tendsto
          (fun n =>
            eLpNorm
              (fun x => (fderiv ℝ (ψ n) x) (EuclideanSpace.single i 1) - hw.weakGrad x i)
              (ENNReal.ofReal p) volume)
          atTop (nhds 0) := by
      simpa [ψ] using (hφ_grad i).comp (tendsto_add_atTop_nat N)
    have hψ_grad_nonneg :
        ∀ n,
          (0 : ℝ≥0∞) ≤
            eLpNorm
              (fun x => (fderiv ℝ (ψ n) x) (EuclideanSpace.single i 1) - hw.weakGrad x i)
              (ENNReal.ofReal p) (volume.restrict Ω) :=
      fun n => zero_le _
    have hψ_grad_bound :
        ∀ n,
          eLpNorm
              (fun x => (fderiv ℝ (ψ n) x) (EuclideanSpace.single i 1) - hw.weakGrad x i)
              (ENNReal.ofReal p) (volume.restrict Ω) ≤
            eLpNorm
              (fun x => (fderiv ℝ (ψ n) x) (EuclideanSpace.single i 1) - hw.weakGrad x i)
              (ENNReal.ofReal p) volume :=
      fun n =>
        eLpNorm_mono_measure
          (f := fun x => (fderiv ℝ (ψ n) x) (EuclideanSpace.single i 1) - hw.weakGrad x i)
          Measure.restrict_le_self
    exact tendsto_of_tendsto_of_tendsto_of_le_of_le
      tendsto_const_nhds hψ_grad_global hψ_grad_nonneg hψ_grad_bound

\end{lstlisting}
\begin{proof}
Choose $\delta > 0$ so that the closed $\delta$-thickening $K' := \{x :
\dist(x,K) \le \delta\}$ is contained in $\Omega$
(\leanname{IsCompact.exists\_cthickening\_subset\_open}), and a smooth cutoff
$\chi$ as in Lemma~\ref{lem:smooth_cutoff} satisfying $\chi = 1$ on $K'$ and
$\supp \chi \subseteq \Omega$.

We promote the local witness on $\Omega$ to a witness on the whole space $E$
with the same gradient field. The $\Lp{p}$-membership extends to all of $E$
because $u$ and each $G_i$ vanish outside $K \subseteq \Omega$
(Lemma~\ref{lem:zero_outside_tsupport}). To verify the weak-gradient identity
on $E$, take any test function $\varphi \in C_c^\infty(E)$ and consider $\psi
:= \chi \varphi$. Then $\psi \in C_c^\infty$ and $\supp \psi \subseteq \Omega$,
so the witness on $\Omega$ gives
\begin{equation*}
  \int_\Omega u(x)\,\partial_i\psi(x)\,dx = -\int_\Omega G_i(x)\,\psi(x)\,dx.
\end{equation*}
On $K$ we have $\chi = 1$ and $\partial_i\chi = 0$ (because $\chi$ is constant
equal to $1$ on the larger neighborhood $K'$ of $K$); since $u$ and $G_i$ are
both supported in $K$, this gives pointwise on $E$
\begin{equation*}
  u(x)\,\partial_i\psi(x) = u(x)\,\partial_i\varphi(x), \qquad
  G_i(x)\,\psi(x) = G_i(x)\,\varphi(x),
\end{equation*}
and the integral identity transfers from $\psi$ to $\varphi$.

Now apply Theorem~\ref{thm:smooth_approx_univ} on the whole space to get a
sequence $(\widetilde\varphi_m)$ with
\begin{equation*}
  \supp \widetilde\varphi_m \subseteq \{x : \dist(x, \supp u) \le r_{\text{out}}^{(m)}\}.
\end{equation*}
For $m$ large enough that $r_{\text{out}}^{(m)} \le \delta$, this support is
contained in $K' \subseteq \Omega$. Setting $\varphi_n :=
\widetilde\varphi_{n+N}$ for an appropriate $N$ furnishes the desired sequence:
the support condition holds for every $n$, and $\Lp{p}(\Omega)$-convergence
follows from $\Lp{p}(E)$-convergence by monotonicity of the $\Lp{p}$-norm in
the underlying measure (\leanname{eLpNorm\_mono\_measure}).
\end{proof}

\begin{theorem}[Localization: compactly supported $\Wkp{1}{p}$ implies $\Wop{p}$;
\leanname{DeGiorgi.memW01p\_of\_memW1p\_of\_tsupport\_subset}]\label{thm:memW01p_of_compactly_supported}
Let $\Omega \subseteq E$ be open, $1 < p < \infty$, and $u \in W^{1,p}(\Omega)$
with compact support and $\supp u \subseteq \Omega$. Then $u \in \Wop{p}(\Omega)$.
\end{theorem}

\begin{lstlisting}[style=Matlab-editor,basicstyle=\mlttfamily,escapechar=",mlshowsectionrules=true,mathescape=true,morecomment={[s]{\%\{}{\%\}}},language=lean,numbers=none]
/-- Localization by compact support: a finite-`p` Sobolev function whose
support is compactly contained in an open set belongs to `W₀^{1,p}`. -/
theorem memW01p_of_memW1p_of_tsupport_subset
    {Ω : Set E} (hΩ : IsOpen Ω)
    {p : ℝ} (hp : 1 < p) {u : E → ℝ}
    (hu : MemW1p (ENNReal.ofReal p) u Ω)
    (hu_compact : HasCompactSupport u)
    (hu_sub : tsupport u ⊆ Ω) :
    MemW01p (ENNReal.ofReal p) u Ω := by
  have hp_enn : (1 : ℝ≥0∞) ≤ ENNReal.ofReal p := by
    simpa using (ENNReal.ofReal_le_ofReal (le_of_lt hp) : ENNReal.ofReal (1 : ℝ) ≤ ENNReal.ofReal p)
  rcases hu with ⟨hu_memLp, hu_grad⟩
  choose g hg_memLp hg_weak using hu_grad
  let G : E → E := fun x => WithLp.toLp 2 fun i => g i x
  let hw : MemW1pWitness (ENNReal.ofReal p) u Ω :=
    { memLp := hu_memLp
      weakGrad := G
      weakGrad_component_memLp := hg_memLp
      isWeakGrad := hg_weak }
  obtain ⟨δ, hδ_pos, hδΩ⟩ := hu_compact.isCompact.exists_cthickening_subset_open hΩ hu_sub
  let K : Set E := Metric.cthickening δ (tsupport u)
  have hK_compact : IsCompact K := hu_compact.isCompact.cthickening (r := δ)
  have hKΩ : K ⊆ Ω := hδΩ
  obtain ⟨η, hη_smooth, hη_compact, hη_range, hη_one, hη_sub⟩ :=
    exists_smooth_cutoff (d := d) hK_compact hΩ hKΩ
  have hη_bound : ∀ x, |η x| ≤ 1 := by
    intro x
    rcases hη_range ⟨x, rfl⟩ with ⟨hx0, hx1⟩
    rw [abs_of_nonneg hx0]
    exact hx1
  have hη_fderiv_compact : HasCompactSupport (fderiv ℝ η) := hη_compact.fderiv (𝕜 := ℝ)
  obtain ⟨Cη, hCη⟩ := hη_fderiv_compact.isCompact.exists_bound_of_continuousOn
    ((hη_smooth.continuous_fderiv (by simp)).continuousOn)
  let C₁ : ℝ := max Cη 0
  have hC₁_nonneg : 0 ≤ C₁ := le_max_right _ _
  have hη_grad_bound : ∀ x, ‖fderiv ℝ η x‖ ≤ C₁ := by
    intro x
    by_cases hx : x ∈ tsupport (fderiv ℝ η)
    · exact (hCη x hx).trans (le_max_left _ _)
    · have hzero : fderiv ℝ η x = 0 := image_eq_zero_of_notMem_tsupport hx
      simp [C₁, hzero]
  let v : E → ℝ := fun x => η x * u x
  let hwCut :=
    MemW1pWitness.mul_smooth_bounded_p
      (d := d) (p := ENNReal.ofReal p) hp_enn hΩ hw hη_smooth
      zero_le_one hC₁_nonneg hη_bound hη_grad_bound
  have hv_eq_u : v = u := by
    funext x
    by_cases hx : x ∈ tsupport u
    · have hηx : η x = 1 := hη_one x (Metric.self_subset_cthickening (tsupport u) hx)
      simp [v, hηx]
    · have hux : u x = 0 := image_eq_zero_of_notMem_tsupport hx
      simp [v, hux]
  have hv_sub : tsupport v ⊆ tsupport η := tsupport_smul_subset_left η u
  have hgrad_sub :
      ∀ i : Fin d, tsupport (fun x => hwCut.weakGrad x i) ⊆ tsupport η := by
    intro i
    exact tsupport_mul_smooth_bounded_p_weakGrad_component_subset
      (d := d) (p := ENNReal.ofReal p) hp_enn hΩ hw hη_smooth zero_le_one hC₁_nonneg hη_bound
      hη_grad_bound i
  rcases exists_smooth_W1p_approx_of_supportedWitness
      (d := d) (Ω := Ω) (K := tsupport η) hΩ hp hwCut
      hη_compact.isCompact hη_sub hv_sub hgrad_sub with
    ⟨φ, hφ_smooth, hφ_compact, hφ_sub, hφ_fun, hφ_grad⟩
  simpa [v, hv_eq_u] using
    memW01p_of_global_approx_supported hwCut φ hφ_smooth hφ_compact hφ_sub hφ_fun hφ_grad

\end{lstlisting}
\begin{proof}
Choose any $\Wkp{1}{p}$-witness for $u$ on $\Omega$, with gradient field $G$.
Pick $\delta > 0$ such that the closed $\delta$-thickening of $\supp u$ is
contained in $\Omega$, and let $K$ be that thickening. Apply
Lemma~\ref{lem:smooth_cutoff} to obtain a smooth cutoff $\eta$ equal to $1$ on
$K$ with $\supp \eta \subseteq \Omega$. Since $\eta$ is smooth and compactly
supported, both $|\eta|$ and $\|D\eta\|$ are bounded by some constants $C_0 :=
1$ and $C_1 \ge 0$ (the latter using continuity of $D\eta$ on its compact
support, \leanname{HasCompactSupport.exists\_bound\_of\_continuousOn}).

By Proposition~\ref{prop:witness_mul_smooth_bounded} the function $v := \eta u$
admits a $\Wkp{1}{p}$-witness. Since $\eta = 1$ on a neighborhood of $\supp u$
and $u = 0$ outside $\supp u$, we have $v = u$ pointwise. Moreover,
$\supp v \subseteq \supp \eta$ (a compact subset of $\Omega$), and the same is
true for each component of the new gradient, by
\leanname{DeGiorgi.tsupport\_mul\_smooth\_bounded\_p\_weakGrad\_component\_subset}
(private).

We are now in the situation of Theorem~\ref{thm:smooth_approx_local} with $K :=
\supp\eta$. The resulting smooth approximating sequence shows $u = v$ belongs
to $\Wop{p}(\Omega)$.
\end{proof}

\subsection{Whole-space Sobolev inequality}

\begin{definition}[Sobolev constant; \leanname{DeGiorgi.C\_gns}]
For $d \in \N_{\ge 1}$ and $p \in \R$, define
\begin{equation*}
  C_{\mathrm{gns}}(d, p) :=
  \mathrm{SNormLESNormFDerivOfEqConst_{\R}}\bigl(\mathrm{vol}_{\R^d}, p\bigr),
\end{equation*}
the constant supplied by Mathlib's general $W^{1,p} \hookrightarrow L^{p^*}$
embedding. By construction, $C_{\mathrm{gns}}(d,p) \ge 0$
(\leanname{DeGiorgi.C\_gns\_nonneg}).
\end{definition}

\begin{lstlisting}[style=Matlab-editor,basicstyle=\mlttfamily,escapechar=",mlshowsectionrules=true,mathescape=true,morecomment={[s]{\%\{}{\%\}}},language=lean,numbers=none]
/-! ## Section 1: Sobolev Constant -/

/-- The Sobolev constant for the general `W^{1,p} -> L^{p*}` embedding,
extracted from Mathlib. -/
noncomputable def C_gns (d : ℕ) [NeZero d] (p : ℝ) : ℝ :=
  (MeasureTheory.SNormLESNormFDerivOfEqConst (F := ℝ)
    (μ := (volume : Measure (EuclideanSpace ℝ (Fin d)))) p : ℝ≥0)

theorem C_gns_nonneg (d : ℕ) [NeZero d] (p : ℝ) : 0 ≤ C_gns d p :=
  NNReal.coe_nonneg _

\end{lstlisting}
\begin{lemma}[Sobolev exponent;
\leanname{DeGiorgi.sobolev\_exp\_pos} (private),
\leanname{DeGiorgi.sobolev\_exponent\_relation} (private)]\label{lem:sobolev_exp}
Let $1 \le p < d$. Then $\frac{dp}{d - p} > 0$, and
\begin{equation*}
  \left(\frac{dp}{d - p}\right)^{-1} = \frac{1}{p} - \frac{1}{d}.
\end{equation*}
\end{lemma}

\begin{lstlisting}[style=Matlab-editor,basicstyle=\mlttfamily,escapechar=",mlshowsectionrules=true,mathescape=true,morecomment={[s]{\%\{}{\%\}}},language=lean,numbers=none]
/-! ## Section 2: Exponent Arithmetic -/

private lemma sobolev_exp_pos {p : ℝ} (hp : 1 ≤ p) (hpd : p < (d : ℝ)) :
    0 < (d : ℝ) * p / ((d : ℝ) - p) := by
  apply div_pos
  · exact mul_pos (Nat.cast_pos.mpr (NeZero.pos d)) (by linarith)
  · linarith

/-- The key exponent relation: `(d*p/(d-p))⁻¹ = p⁻¹ - d⁻¹`. -/
private lemma sobolev_exponent_relation
    {p : ℝ} (hp : 0 < p) (hpd : p < (d : ℝ)) :
    ((d : ℝ) * p / ((d : ℝ) - p))⁻¹ = p⁻¹ - (d : ℝ)⁻¹ := by
  let _ := (inferInstance : NeZero d)
  have hd : (0 : ℝ) < d := lt_trans hp hpd
  rw [inv_div, inv_sub_inv (ne_of_gt hp) (ne_of_gt hd)]
  field_simp

\end{lstlisting}
\begin{proof}
Both numerator $dp$ and denominator $d - p$ are positive (since $d \ge 1$ and
$p < d$). The exponent identity follows by direct algebraic manipulation
$\frac{d - p}{dp} = \frac{1}{p} - \frac{1}{d}$.
\end{proof}

\begin{theorem}[Sobolev inequality for smooth compactly supported functions;
\leanname{DeGiorgi.sobolev\_smooth}]\label{thm:sobolev_smooth}
Let $1 \le p < d$, $u \in C_c^1(E)$. Then
\begin{equation*}
  \|u\|_{L^{p^*}(E)} \le C_{\mathrm{gns}}(d, p)\,\|Du\|_{L^p(E)},
  \qquad p^* = \frac{dp}{d-p}.
\end{equation*}
\end{theorem}

\begin{lstlisting}[style=Matlab-editor,basicstyle=\mlttfamily,escapechar=",mlshowsectionrules=true,mathescape=true,morecomment={[s]{\%\{}{\%\}}},language=lean,numbers=none]
/-! ## Section 3: GNS for Smooth Compactly Supported Functions -/

/-- Sobolev inequality for smooth compactly supported functions at general `p`.
Direct wrapper around Mathlib's `eLpNorm_le_eLpNorm_fderiv_of_eq`. -/
theorem sobolev_smooth {p : ℝ} (hp : 1 ≤ p) (hpd : p < (d : ℝ))
    {u : E → ℝ} (hu : ContDiff ℝ 1 u) (hu_cpt : HasCompactSupport u) :
    eLpNorm u (ENNReal.ofReal ((d : ℝ) * p / ((d : ℝ) - p))) volume ≤
    ENNReal.ofReal (C_gns d p) *
      eLpNorm (fderiv ℝ u) (ENNReal.ofReal p) volume := by
  set p_nn : ℝ≥0 := Real.toNNReal p with hp_nn_def
  set p'_nn : ℝ≥0 := Real.toNNReal ((d : ℝ) * p / ((d : ℝ) - p)) with hp'_nn_def
  have hp_nn_ge : (1 : ℝ≥0) ≤ p_nn := by
    rwa [← Real.toNNReal_one, Real.toNNReal_le_toNNReal_iff (by linarith)]
  have hd_pos : 0 < Module.finrank ℝ E := by
    rw [finrank_euclideanSpace_fin]
    exact NeZero.pos d
  have hrel : (p'_nn : ℝ)⁻¹ = (p_nn : ℝ)⁻¹ - (Module.finrank ℝ E : ℝ)⁻¹ :=
    sobolev_exponent_nnreal_relation hp hpd
  have hmain := MeasureTheory.eLpNorm_le_eLpNorm_fderiv_of_eq
    (F := ℝ) (μ := (volume : Measure E)) hu hu_cpt hp_nn_ge hd_pos hrel
  have hp_real : (p_nn : ℝ) = p := Real.coe_toNNReal p (by linarith)
  have hp_conv : (p_nn : ℝ≥0∞) = ENNReal.ofReal p := by
    simp [p_nn, ENNReal.ofReal, Real.toNNReal]
  have hp'_conv : (p'_nn : ℝ≥0∞) = ENNReal.ofReal ((d : ℝ) * p / ((d : ℝ) - p)) := by
    simp [p'_nn, ENNReal.ofReal, Real.toNNReal]
  rw [hp_conv, hp'_conv, hp_real] at hmain
  calc
    eLpNorm u (ENNReal.ofReal ((d : ℝ) * p / ((d : ℝ) - p))) volume
      ≤ ↑(MeasureTheory.SNormLESNormFDerivOfEqConst (F := ℝ) (μ := (volume : Measure E)) p) *
          eLpNorm (fderiv ℝ u) (ENNReal.ofReal p) volume := hmain
    _ = ENNReal.ofReal (C_gns d p) *
          eLpNorm (fderiv ℝ u) (ENNReal.ofReal p) volume := by
        congr 1
        simp only [C_gns]
        exact (ENNReal.ofReal_coe_nnreal).symm

\end{lstlisting}
\begin{proof}
This is a direct application of Mathlib's general result
\leanname{MeasureTheory.eLpNorm\_le\_eLpNorm\_fderiv\_of\_eq} with the exponent
relation supplied by Lemma~\ref{lem:sobolev_exp} (rewritten in terms of
$\mathrm{Real.toNNReal}$ to match the Mathlib API).
\end{proof}

\begin{lemma}[Operator norm of the differential equals partials norm;
\leanname{DeGiorgi.norm\_fderiv\_eq\_norm\_partials} (private)]\label{lem:norm_fderiv}
For any $f : E \to \R$ with derivative $Df(x)$ at $x$,
\begin{equation*}
  \|Df(x)\|
   = \left(\sum_{i=1}^d |\partial_i f(x)|^2\right)^{1/2}.
\end{equation*}
\end{lemma}

\begin{lstlisting}[style=Matlab-editor,basicstyle=\mlttfamily,escapechar=",mlshowsectionrules=true,mathescape=true,morecomment={[s]{\%\{}{\%\}}},language=lean,numbers=none]
/-! ## Section 4: Gradient norm comparison -/

omit [NeZero d] in
private theorem norm_fderiv_eq_norm_partials
    {f : E → ℝ} {x : E} :
    ‖fderiv ℝ f x‖ =
    ‖WithLp.toLp 2
      (fun i => (fderiv ℝ f x) (EuclideanSpace.single i 1))‖ := by
  let v : E := (InnerProductSpace.toDual ℝ E).symm (fderiv ℝ f x)
  have hv :
      v = WithLp.toLp 2
        (fun i => (fderiv ℝ f x) (EuclideanSpace.single i 1)) := by
    have hv_map : (InnerProductSpace.toDual ℝ E) v = fderiv ℝ f x := by
      simp [v]
    ext i
    calc
      v i = inner ℝ v (EuclideanSpace.single i (1 : ℝ)) := by
        simpa using
          (EuclideanSpace.inner_single_right (i := i) (a := (1 : ℝ)) v).symm
      _ = ((InnerProductSpace.toDual ℝ E) v) (EuclideanSpace.single i (1 : ℝ)) := by
        rw [InnerProductSpace.toDual_apply_apply]
      _ = (fderiv ℝ f x) (EuclideanSpace.single i (1 : ℝ)) := by
        rw [hv_map]
      _ = (WithLp.toLp 2
            (fun j => (fderiv ℝ f x) (EuclideanSpace.single j 1))) i := by
        simp
  calc
    ‖fderiv ℝ f x‖ = ‖v‖ := by
      simp [v]
    _ = ‖WithLp.toLp 2
        (fun i => (fderiv ℝ f x) (EuclideanSpace.single i 1))‖ := by
          rw [hv]

\end{lstlisting}
\begin{proof}
By Riesz representation in the Hilbert space $E$, $Df(x)$ corresponds to a
unique $v \in E$ with $Df(x) = \langle v, \cdot\rangle$, and $\|Df(x)\| = \|v\|$.
Evaluating at the standard basis vector $e_i$ gives $v_i = Df(x)(e_i) =
\partial_i f(x)$, so $\|v\|^2 = \sum_i |\partial_i f(x)|^2$.
\end{proof}

\begin{theorem}[Sobolev inequality from approximating sequences;
\leanname{DeGiorgi.sobolev\_of\_approx}]\label{thm:sobolev_of_approx}
Let $1 \le p < d$, $u : E \to \R$ a.e.\ strongly measurable, and $G : E \to E$
a vector field whose components are a.e.\ strongly measurable. Suppose there
is a sequence $(\varphi_n)$ of smooth compactly supported functions on $E$
with
\begin{equation*}
  \|\varphi_n - u\|_{\Lp{p}(E)} \to 0, \qquad
  \|\partial_i\varphi_n - G_i\|_{\Lp{p}(E)} \to 0 \;\text{ for every } i.
\end{equation*}
Then
\begin{equation*}
  \|u\|_{L^{p^*}(E)} \le C_{\mathrm{gns}}(d,p)\,\|\,\|G(\cdot)\|\,\|_{\Lp{p}(E)},
  \qquad p^* = \frac{dp}{d-p}.
\end{equation*}
\end{theorem}

\begin{lstlisting}[style=Matlab-editor,basicstyle=\mlttfamily,escapechar=",mlshowsectionrules=true,mathescape=true,morecomment={[s]{\%\{}{\%\}}},language=lean,numbers=none]
/-- The main whole-space Sobolev inequality for functions approximated by
smooth compactly supported functions. -/
theorem sobolev_of_approx {p : ℝ} (hp : 1 ≤ p) (hpd : p < (d : ℝ))
    {u : E → ℝ} {G : E → E}
    (hu_aesm : AEStronglyMeasurable u volume)
    (hG_comp_aesm : ∀ i : Fin d, AEStronglyMeasurable (fun x => G x i) volume)
    (φ : ℕ → E → ℝ)
    (hφ_smooth : ∀ n, ContDiff ℝ (⊤ : ℕ∞) (φ n))
    (hφ_cpt : ∀ n, HasCompactSupport (φ n))
    (hφ_fun : Tendsto (fun n => eLpNorm (fun x => φ n x - u x) (ENNReal.ofReal p) volume)
      atTop (nhds 0))
    (hφ_grad : ∀ i : Fin d, Tendsto (fun n => eLpNorm
      (fun x => (fderiv ℝ (φ n) x) (EuclideanSpace.single i 1) - G x i)
      (ENNReal.ofReal p) volume) atTop (nhds 0)) :
    eLpNorm u (ENNReal.ofReal ((d : ℝ) * p / ((d : ℝ) - p))) volume ≤
    ENNReal.ofReal (C_gns d p) *
      eLpNorm (fun x => ‖G x‖) (ENNReal.ofReal p) volume := by
  set p_star := ENNReal.ofReal ((d : ℝ) * p / ((d : ℝ) - p)) with hp_star_def
  set p_enn := ENNReal.ofReal p with hp_enn_def
  set C := ENNReal.ofReal (C_gns d p) with hC_def
  have hGNS : ∀ n,
      eLpNorm (φ n) p_star volume ≤
      C * eLpNorm (fderiv ℝ (φ n)) p_enn volume :=
    fun n => sobolev_smooth hp hpd ((hφ_smooth n).of_le (by norm_cast)) (hφ_cpt n)
  have hφ_aesm : ∀ n, AEStronglyMeasurable (φ n) volume :=
    fun n => (hφ_smooth n).continuous.aestronglyMeasurable
  have hp_ne : p_enn ≠ 0 := by
    simp [p_enn]
    linarith
  have hp_one_enn : (1 : ℝ≥0∞) ≤ p_enn := by
    rw [hp_enn_def, ← ENNReal.ofReal_one]
    exact ENNReal.ofReal_le_ofReal hp
  have hTIM : TendstoInMeasure volume φ atTop u :=
    tendstoInMeasure_of_eLpNorm_tendsto hp_ne hφ_aesm hu_aesm hφ_fun
  obtain ⟨σ, hσ_mono, hσ_ae⟩ := hTIM.exists_seq_tendsto_ae
  have hFatou : eLpNorm u p_star volume ≤
      atTop.liminf (fun n => eLpNorm (φ (σ n)) p_star volume) :=
    Lp.eLpNorm_lim_le_liminf_eLpNorm (fun n => hφ_aesm (σ n)) u hσ_ae
  let gradVec : ℕ → E → E := fun n x =>
    WithLp.toLp 2 (fun i => (fderiv ℝ (φ n) x) (EuclideanSpace.single i 1))
  have hBridge : ∀ n, (fun x => ‖fderiv ℝ (φ n) x‖) =ᵐ[volume]
      (fun x => ‖gradVec n x‖) := by
    intro n
    filter_upwards with x
    exact norm_fderiv_eq_norm_partials
  have hGNS' : ∀ n,
      eLpNorm (φ n) p_star volume ≤
      C * eLpNorm (fun x => ‖gradVec n x‖) p_enn volume := by
    intro n
    calc
      eLpNorm (φ n) p_star volume
        ≤ C * eLpNorm (fderiv ℝ (φ n)) p_enn volume := hGNS n
      _ = C * eLpNorm (fun x => ‖gradVec n x‖) p_enn volume := by
          congr 1
          rw [← eLpNorm_norm (fderiv ℝ (φ n))]
          exact eLpNorm_congr_ae (hBridge n)
  have hGradDiffTendstoZero : Tendsto
      (fun n => eLpNorm (fun x => gradVec n x - G x) p_enn volume) atTop (nhds 0) := by
    have hNormCompTendsto :=
      partialDiff_sum_tendsto_zero (p := p) (φ := φ) (G := G) hφ_grad
    have hBound :
        ∀ᶠ n in atTop,
          eLpNorm (fun x => gradVec n x - G x) p_enn volume ≤
            ∑ i : Fin d,
              eLpNorm
                (fun x => (fderiv ℝ (φ n) x) (EuclideanSpace.single i 1) - G x i)
                p_enn volume := by
      apply Eventually.of_forall
      intro n
      simpa [gradVec, hp_enn_def] using
        gradVec_eLpNorm_le_sum (p := p) (φ := φ) (G := G) hp hφ_smooth hG_comp_aesm n
    have hZeroLE :
        ∀ᶠ n in atTop, (0 : ℝ≥0∞) ≤ eLpNorm (fun x => gradVec n x - G x) p_enn volume :=
      Eventually.of_forall fun n => zero_le _
    exact tendsto_of_tendsto_of_tendsto_of_le_of_le'
      (show Tendsto (fun _ : ℕ => (0 : ℝ≥0∞)) atTop (nhds 0) from tendsto_const_nhds)
      hNormCompTendsto
      hZeroLE
      hBound
  have hBound : atTop.liminf (fun n => eLpNorm (φ (σ n)) p_star volume) ≤
      C * eLpNorm (fun x => ‖G x‖) p_enn volume := by
    have hC_ne_top : C ≠ ⊤ := ENNReal.ofReal_ne_top
    have hGnorm_aesm : AEStronglyMeasurable (fun x => ‖G x‖) volume := by
      exact (aestronglyMeasurable_euclidean_of_components hG_comp_aesm).norm
    have hGradNormUpper :
        ∀ᶠ n in atTop,
          eLpNorm (fun x => ‖gradVec (σ n) x‖) p_enn volume ≤
            eLpNorm (fun x => ‖G x‖) p_enn volume +
              eLpNorm (fun x => gradVec (σ n) x - G x) p_enn volume := by
      apply Eventually.of_forall
      intro n
      have hGradNormAesm :
          AEStronglyMeasurable (fun x => ‖gradVec (σ n) x‖ - ‖G x‖) volume := by
        have hGradVecCont : Continuous (gradVec (σ n)) := by
          have h_toLp_cont : Continuous (fun x : Fin d → ℝ => WithLp.toLp 2 x) := by
            simpa using (PiLp.continuous_toLp 2 (fun _ : Fin d => ℝ))
          refine h_toLp_cont.comp ?_
          rw [continuous_pi_iff]
          intro i
          exact ((hφ_smooth (σ n)).continuous_fderiv (by simp)).clm_apply continuous_const
        have hGradVecAesm : AEStronglyMeasurable (gradVec (σ n)) volume :=
          hGradVecCont.aestronglyMeasurable
        exact hGradVecAesm.norm.sub hGnorm_aesm
      have hNormDiff :
          eLpNorm (fun x => ‖gradVec (σ n) x‖ - ‖G x‖) p_enn volume ≤
            eLpNorm (fun x => gradVec (σ n) x - G x) p_enn volume := by
        calc
          eLpNorm (fun x => ‖gradVec (σ n) x‖ - ‖G x‖) p_enn volume
              = eLpNorm (fun x => |‖gradVec (σ n) x‖ - ‖G x‖|) p_enn volume := by
                  simpa [Real.norm_eq_abs] using
                    (eLpNorm_norm
                      (f := fun x => ‖gradVec (σ n) x‖ - ‖G x‖)
                      (p := p_enn) (μ := volume)).symm
          _ ≤ eLpNorm (fun x => gradVec (σ n) x - G x) p_enn volume := by
                refine eLpNorm_mono_ae ?_
                filter_upwards with x
                simpa [Real.norm_eq_abs] using
                  (abs_norm_sub_norm_le (gradVec (σ n) x) (G x))
      calc
        eLpNorm (fun x => ‖gradVec (σ n) x‖) p_enn volume
            = eLpNorm (fun x => (‖gradVec (σ n) x‖ - ‖G x‖) + ‖G x‖) p_enn volume := by
                congr 1
                ext x
                ring
        _ ≤ eLpNorm (fun x => ‖gradVec (σ n) x‖ - ‖G x‖) p_enn volume +
              eLpNorm (fun x => ‖G x‖) p_enn volume := by
                exact eLpNorm_add_le hGradNormAesm hGnorm_aesm hp_one_enn
        _ ≤ eLpNorm (fun x => gradVec (σ n) x - G x) p_enn volume +
              eLpNorm (fun x => ‖G x‖) p_enn volume := by
                exact add_le_add hNormDiff le_rfl
        _ = eLpNorm (fun x => ‖G x‖) p_enn volume +
              eLpNorm (fun x => gradVec (σ n) x - G x) p_enn volume := by
                rw [add_comm]
    have hGradBound :
        atTop.liminf (fun n => C * eLpNorm (fun x => ‖gradVec (σ n) x‖) p_enn volume) ≤
          C * eLpNorm (fun x => ‖G x‖) p_enn volume := by
      have hDiffZeroSubseq :
          Tendsto
            (fun n => C * eLpNorm (fun x => gradVec (σ n) x - G x) p_enn volume)
            atTop (nhds 0) := by
          simpa using
            (ENNReal.Tendsto.const_mul
              (hGradDiffTendstoZero.comp hσ_mono.tendsto_atTop)
              (Or.inr hC_ne_top))
      have hu_grad :
          atTop.IsBoundedUnder (· ≥ ·)
            (fun n => C * eLpNorm (fun x => ‖gradVec (σ n) x‖) p_enn volume) :=
        isBoundedUnder_of_eventually_ge (a := 0) <| Eventually.of_forall fun _ => by simp
      have hv_rhs :
          atTop.IsCoboundedUnder (· ≥ ·)
            (fun n =>
              C * eLpNorm (fun x => ‖G x‖) p_enn volume +
                C * eLpNorm (fun x => gradVec (σ n) x - G x) p_enn volume) :=
        isCoboundedUnder_ge_of_eventually_le atTop <|
          Eventually.of_forall fun _ => le_top
      calc
        atTop.liminf (fun n => C * eLpNorm (fun x => ‖gradVec (σ n) x‖) p_enn volume)
            ≤ atTop.liminf
                (fun n =>
                  C * eLpNorm (fun x => ‖G x‖) p_enn volume +
                    C * eLpNorm (fun x => gradVec (σ n) x - G x) p_enn volume) := by
                  exact Filter.liminf_le_liminf
                    (hGradNormUpper.mono fun n hn => by
                      simpa [mul_add, add_comm, add_left_comm, add_assoc] using
                        mul_le_mul_right hn C)
                    hu_grad hv_rhs
        _ = atTop.liminf (fun _ : ℕ => C * eLpNorm (fun x => ‖G x‖) p_enn volume) := by
              simpa using
                ENNReal.liminf_add_of_right_tendsto_zero hDiffZeroSubseq
                  (fun _ : ℕ => C * eLpNorm (fun x => ‖G x‖) p_enn volume)
        _ = C * eLpNorm (fun x => ‖G x‖) p_enn volume := by
              simp
    have hu_phi :
        atTop.IsBoundedUnder (· ≥ ·) (fun n => eLpNorm (φ (σ n)) p_star volume) :=
      isBoundedUnder_of_eventually_ge (a := 0) <| Eventually.of_forall fun _ => by simp
    have hv_grad :
        atTop.IsCoboundedUnder (· ≥ ·)
          (fun n => C * eLpNorm (fun x => ‖gradVec (σ n) x‖) p_enn volume) :=
      isCoboundedUnder_ge_of_eventually_le atTop <|
        Eventually.of_forall fun _ => le_top
    calc
      atTop.liminf (fun n => eLpNorm (φ (σ n)) p_star volume)
          ≤ atTop.liminf (fun n => C * eLpNorm (fun x => ‖gradVec (σ n) x‖) p_enn volume) := by
            exact Filter.liminf_le_liminf
              (Eventually.of_forall fun n => hGNS' (σ n))
              hu_phi hv_grad
      _ ≤ C * eLpNorm (fun x => ‖G x‖) p_enn volume := hGradBound
  exact le_trans hFatou hBound

\end{lstlisting}
\begin{proof}
Apply Theorem~\ref{thm:sobolev_smooth} to each $\varphi_n$:
\begin{equation*}
  \|\varphi_n\|_{L^{p^*}} \le C_{\mathrm{gns}}(d,p)\,\|D\varphi_n\|_{L^p}.
\end{equation*}
Define $\Phi_n(x) := (\partial_1\varphi_n(x),\ldots,\partial_d\varphi_n(x))$,
the vector of classical partial derivatives. By Lemma~\ref{lem:norm_fderiv},
$\|D\varphi_n(x)\| = \|\Phi_n(x)\|$, hence $\|D\varphi_n\|_{L^p} = \|\,\|\Phi_n\|\,\|_{L^p}$.

By the triangle inequality applied componentwise (and bounding the Euclidean
norm of a vector by the sum of the norms of its components, see
\leanname{DeGiorgi.gradVec\_eLpNorm\_le\_sum} (private)),
\begin{equation*}
  \|\Phi_n - G\|_{L^p(E; E)}
   \le \sum_{i=1}^d \|\partial_i\varphi_n - G_i\|_{\Lp{p}(E)} \;\to\; 0.
\end{equation*}
Hence $\|\,\|\Phi_n\|\,\|_{L^p} \to \|\,\|G\|\,\|_{L^p}$, by reverse triangle
inequality $|\,\|\Phi_n(x)\| - \|G(x)\|\,| \le \|\Phi_n(x) - G(x)\|$.

Convergence in $L^p$ implies convergence in measure
(\leanname{tendstoInMeasure\_of\_tendsto\_eLpNorm}), so by passing to a
subsequence we may assume $\varphi_{\sigma(n)} \to u$ a.e. By Fatou's lemma
applied to $L^{p^*}$ (\leanname{Lp.eLpNorm\_lim\_le\_liminf\_eLpNorm}),
\begin{equation*}
  \|u\|_{L^{p^*}}
   \le \liminf_{n\to\infty} \|\varphi_{\sigma(n)}\|_{L^{p^*}}
   \le \liminf_{n\to\infty} C_{\mathrm{gns}}(d,p)\,\|\,\|\Phi_{\sigma(n)}\|\,\|_{L^p}.
\end{equation*}
The right-hand side $\liminf$ is bounded above by $C_{\mathrm{gns}}(d,p)\,\|
\,\|G\|\,\|_{L^p}$, since $\|\,\|\Phi_{\sigma(n)}\|\,\|_{L^p} \le \|\,\|G\|\,\|_{L^p}
+ \|\Phi_{\sigma(n)} - G\|_{L^p}$ and the second term tends to zero, using
\leanname{ENNReal.liminf\_add\_of\_right\_tendsto\_zero}. Combining yields the
claim.
\end{proof}

\begin{corollary}[Sobolev inequality on $\Wop{p}(E)$;
\leanname{DeGiorgi.sobolev\_of\_memW01p\_univ}]\label{cor:sobolev_W01p}
Let $1 \le p < d$ and $u \in \Wop{p}(E)$. Then there exists a
$\Wkp{1}{p}$-witness for $u$ on $E$ such that
\begin{equation*}
  \|u\|_{L^{p^*}(E)} \le C_{\mathrm{gns}}(d,p)\,\|\,\|G(\cdot)\|\,\|_{\Lp{p}(E)},
  \qquad p^* = \frac{dp}{d - p},
\end{equation*}
where $G$ is the weak gradient field carried by the witness.
\end{corollary}

\begin{lstlisting}[style=Matlab-editor,basicstyle=\mlttfamily,escapechar=",mlshowsectionrules=true,mathescape=true,morecomment={[s]{\%\{}{\%\}}},language=lean,numbers=none]
/-- Whole-space Sobolev inequality specialized to the `MemW₀^{1,p}` data
already stored in the `MemW₀^{1,p}` witness predicate. -/
theorem sobolev_of_memW01p_univ
    {p : ℝ} (hp : 1 ≤ p) (hpd : p < (d : ℝ))
    {u : E → ℝ}
    (hu : MemW01p (ENNReal.ofReal p) u Set.univ volume) :
    ∃ hw : MemW1pWitness (ENNReal.ofReal p) u Set.univ volume,
      eLpNorm u (ENNReal.ofReal ((d : ℝ) * p / ((d : ℝ) - p))) volume ≤
        ENNReal.ofReal (C_gns d p) *
          eLpNorm (fun x => ‖hw.weakGrad x‖) (ENNReal.ofReal p) volume := by
  rcases hu with ⟨_, hw, φ, hφ_smooth, hφ_cpt, _hφ_sub, hφ_fun, hφ_grad⟩
  refine ⟨hw, ?_⟩
  have hu_aesm : AEStronglyMeasurable u volume := by
    simpa using hw.memLp.aestronglyMeasurable
  have hG_comp_aesm : ∀ i : Fin d,
      AEStronglyMeasurable (fun x => hw.weakGrad x i) volume := by
    intro i
    simpa using (hw.weakGrad_component_memLp i).aestronglyMeasurable
  have hφ_fun_univ :
      Tendsto (fun n => eLpNorm (fun x => φ n x - u x) (ENNReal.ofReal p) volume)
        atTop (nhds 0) := by
    simpa using hφ_fun
  have hφ_grad_univ : ∀ i : Fin d, Tendsto
      (fun n => eLpNorm
        (fun x => (fderiv ℝ (φ n) x) (EuclideanSpace.single i 1) - hw.weakGrad x i)
        (ENNReal.ofReal p) volume)
      atTop (nhds 0) := by
    intro i
    simpa using hφ_grad i
  simpa [eLpNorm_norm] using
    sobolev_of_approx hp hpd hu_aesm hG_comp_aesm φ
      hφ_smooth hφ_cpt hφ_fun_univ hφ_grad_univ


\end{lstlisting}
\begin{proof}
By Definition~\ref{def:MemW01p}, $u \in \Wop{p}(E)$ comes equipped with a
$\Wkp{1}{p}$-witness on $E$ and an approximating sequence $(\varphi_n)$ of
smooth compactly supported functions converging to $u$ and whose classical
gradients converge to the weak gradient $G$ in $\Lp{p}$. The hypotheses of
Theorem~\ref{thm:sobolev_of_approx} are satisfied, and its conclusion is
exactly the desired inequality.
\end{proof}

\subsection{Positive-part prelude}

\begin{lemma}[Indicator truncation preserves $L^2$;
\leanname{DeGiorgi.indicator\_component\_memLp}]\label{lem:indicator_component_memLp}
Let $\Omega \subseteq E$, $\sigma : E \to \R$ be a.e.\ measurable on $\Omega$
and $g \in \Lp{2}(\Omega)$. Then the function
\begin{equation*}
  x \mapsto \mathbf{1}_{\{\sigma > 0\}}(x)\,g(x)
\end{equation*}
also lies in $\Lp{2}(\Omega)$.
\end{lemma}

\begin{lstlisting}[style=Matlab-editor,basicstyle=\mlttfamily,escapechar=",mlshowsectionrules=true,mathescape=true,morecomment={[s]{\%\{}{\%\}}},language=lean,numbers=none]
/-! ## Initial Positive-Part Infrastructure -/

omit [NeZero d] in
/-- Truncating an `L²` function by the positivity set of an a.e.-measurable
scalar function preserves `L²`. -/
theorem indicator_component_memLp
    {Ω : Set E} {σ g : E → ℝ}
    (hσ_aemeas : AEMeasurable σ (volume.restrict Ω))
    (hg_memLp : MemLp g 2 (volume.restrict Ω)) :
    MemLp (fun x => if 0 < σ x then g x else 0) 2 (volume.restrict Ω) := by
  let h : ℝ × ℝ → ℝ := fun yz => if 0 < yz.1 then yz.2 else 0
  have hh_meas : Measurable h := by
    refine measurable_snd.piecewise ?_ measurable_const
    exact measurableSet_lt measurable_const measurable_fst
  have hpair_aemeas :
      AEMeasurable (fun x => (σ x, g x)) (volume.restrict Ω) :=
    hσ_aemeas.prodMk hg_memLp.aemeasurable
  have htrunc_aestrong :
      AEStronglyMeasurable (fun x => if 0 < σ x then g x else 0) (volume.restrict Ω) := by
    refine (hh_meas.comp_aemeasurable hpair_aemeas).aestronglyMeasurable.congr ?_
    filter_upwards [] with x
    by_cases hx : 0 < σ x
    · simp [h, hx]
    · simp [h, hx]
  refine hg_memLp.norm.mono' htrunc_aestrong ?_
  filter_upwards [] with x
  by_cases hx : 0 < σ x
  · simp [hx]
  · simp [hx]

\end{lstlisting}
\begin{proof}
The set $\{\sigma > 0\}$ is measurable since $\sigma$ is a.e.\ measurable, so
the indicator multiplication preserves a.e.\ strong measurability. Pointwise,
$|\mathbf{1}_{\{\sigma > 0\}}(x)\,g(x)| \le |g(x)|$, and squaring, integrating
and applying the $\Lp{2}$ membership of $g$ gives the conclusion.
\end{proof}

\begin{lemma}[Truncated weak-gradient component in $L^2$;
\leanname{DeGiorgi.posPartGrad\_component\_memLp}]\label{lem:posPartGrad}
Let $u \in H^1(\Omega)$ with $H^1$-witness gradient field $G$. For each
$i \in \{1,\ldots,d\}$, the function
\begin{equation*}
  x \mapsto \mathbf{1}_{\{u > 0\}}(x)\,G_i(x)
\end{equation*}
belongs to $\Lp{2}(\Omega)$.
\end{lemma}

\begin{lstlisting}[style=Matlab-editor,basicstyle=\mlttfamily,escapechar=",mlshowsectionrules=true,mathescape=true,morecomment={[s]{\%\{}{\%\}}},language=lean,numbers=none]
/-- Each component of the truncated positive-part gradient lies in `L²`. -/
theorem posPartGrad_component_memLp
    {Ω : Set E} {u : E → ℝ}
    (hw : MemW1pWitness 2 u Ω) (i : Fin d) :
    MemLp (fun x => (if 0 < u x then hw.weakGrad x else 0) i) 2
      (volume.restrict Ω) := by
  let _ := (inferInstance : NeZero d)
  convert
    indicator_component_memLp hw.memLp.aemeasurable (hw.weakGrad_component_memLp i) using 1
  ext x
  by_cases hx : 0 < u x
  · simp [hx]
  · simp [hx]

\end{lstlisting}
\begin{proof}
Apply Lemma~\ref{lem:indicator_component_memLp} with $\sigma := u$ and $g :=
G_i$, both of which lie in the relevant $\Lp{2}$ spaces (the second by the
witness assumption).
\end{proof}


%% ========================================================================
\section{Functional Inequalities}\label{sec:functional}
%% ========================================================================

% Section 3: Functional Inequalities
% Translation of DeGiorgi/{Poincare,SobolevPoincare,StampacchiaTruncation,Oscillation/*}.lean

This section gathers the functional analytic inequalities used throughout the De Giorgi
machinery: a representation--formula proof of the Poincar\'e inequality on the unit ball,
the Sobolev--Poincar\'e refinement obtained via extension and Gagliardo--Nirenberg--Sobolev,
the Stampacchia truncation lemma, and the BMO / John--Nirenberg / Campanato package
underlying the oscillation theory. Throughout the section, $E = \R^d$ with $d \ge 1$
(in most statements $d \ge 2$, supplied via a \texttt{NeZero} hypothesis), and $\Bz{r}$
denotes the open ball of radius $r$ centred at the origin.

\subsection{Poincar\'e inequality on the unit ball}

\begin{definition}[\leanname{C\_poinc\_val}]
For $d \in \N$, define the Poincar\'e constant
\[
C_{\mathrm{poinc}}(d) := 2^{d+1}\, d.
\]
The lemma \leanname{C\_poinc\_val\_pos} records that $C_{\mathrm{poinc}}(d) > 0$ for
$d > 0$.
\end{definition}

\begin{lstlisting}[style=Matlab-editor,basicstyle=\mlttfamily,escapechar=",mlshowsectionrules=true,mathescape=true,morecomment={[s]{\%\{}{\%\}}},language=lean,numbers=none]
noncomputable def C_poinc_val (d : ℕ) : ℝ :=
  2 ^ (d + 1) * d

omit [NeZero d] in
theorem C_poinc_val_pos (hd : 0 < d) : 0 < C_poinc_val d := by
  unfold C_poinc_val
  positivity

\end{lstlisting}
\begin{definition}[\leanname{smoothGradNorm}]
For a smooth scalar function $u : E \to \R$, the pointwise Euclidean gradient norm is
\[
\lvert \nabla u (x)\rvert
:= \Bigl(\sum_{i=1}^d \bigl(\partial_i u(x)\bigr)^{2}\Bigr)^{1/2}.
\]
\end{definition}

\begin{lstlisting}[style=Matlab-editor,basicstyle=\mlttfamily,escapechar=",mlshowsectionrules=true,mathescape=true,morecomment={[s]{\%\{}{\%\}}},language=lean,numbers=none]
/-- Pointwise Euclidean norm of the classical gradient of a smooth scalar
function, written in coordinates. This matches the weak-gradient norm used in
the Sobolev witness layer for smooth functions. -/
noncomputable def smoothGradNorm (u : E → ℝ) : E → ℝ :=
  fun x => ‖WithLp.toLp 2 (fun i => (fderiv ℝ u x) (EuclideanSpace.single i 1))‖

\end{lstlisting}
The lemma \leanname{norm\_fderiv\_eq\_smoothGradNorm} identifies this expression with the
operator norm of the Fr\'echet derivative $\lVert \mathrm{D}u(x)\rVert$, via the
Riesz representation $\mathrm{D}u(x)\mapsto \nabla u (x)$ in
$\R^d$. Auxiliary support lemmas
\leanname{smoothGradNorm\_eq\_zero\_of\_notMem\_tsupport} and
\leanname{support\_smoothGradNorm\_subset\_tsupport} record that the gradient vanishes
outside the closed support of $u$.

\begin{lemma}[\leanname{smoothCompactSupport\_L2\_bound\_on\_bounded\_ge\_two}]
Let $d \ge 2$ and let $\Omega \subseteq E$ be a bounded set. Then there exists
$C < \infty$ such that for every smooth, compactly supported $u$ with $\supp u
\subseteq \Omega$,
\[
\lVert u\rVert_{L^{2}(\Omega)}
\le C\, \lVert \lvert\nabla u\rvert\rVert_{L^{2}(\Omega)}.
\]
\end{lemma}

\begin{lstlisting}[style=Matlab-editor,basicstyle=\mlttfamily,escapechar=",mlshowsectionrules=true,mathescape=true,morecomment={[s]{\%\{}{\%\}}},language=lean,numbers=none]
omit [NeZero d] in
theorem smoothCompactSupport_L2_bound_on_bounded_ge_two
    {Ω : Set E} (hd : 2 ≤ d) (hΩ_bdd : Bornology.IsBounded Ω) :
    ∃ C : ℝ≥0∞, C < ∞ ∧
      ∀ {u : E → ℝ}, ContDiff ℝ (⊤ : ℕ∞) u → HasCompactSupport u → tsupport u ⊆ Ω →
        eLpNorm u 2 (volume.restrict Ω) ≤
          C * eLpNorm (smoothGradNorm u) 2 (volume.restrict Ω) := by
  obtain ⟨R, hR_pos, hΩ_sub⟩ := hΩ_bdd.subset_closedBall_lt (0 : ℝ) (0 : E)
  let s : Set E := Metric.closedBall (0 : E) R
  have hs_bdd : Bornology.IsBounded s := Metric.isBounded_closedBall
  have hs_finite : volume s < ∞ := measure_closedBall_lt_top
  by_cases hd_eq_two : d = 2
  · let Csob : ℝ≥0 := MeasureTheory.eLpNormLESNormFDerivOfLeConst ℝ volume s 1 2
    let C : ℝ≥0∞ := (Csob : ℝ≥0∞) * (volume s) ^ (1 / (1 : ℝ) - 1 / (2 : ℝ))
    refine ⟨C, ?_, ?_⟩
    · dsimp [C]
      exact ENNReal.mul_lt_top ENNReal.coe_lt_top <|
        ENNReal.rpow_lt_top_of_nonneg (by positivity) hs_finite.ne
    · intro u hu hu_cpt hu_sub
      have hu1 : ContDiff ℝ 1 u := hu.of_le (by norm_cast)
      have hu_support_sub : Function.support u ⊆ Ω := subset_closure.trans hu_sub
      have hu_support_sub_s : Function.support u ⊆ s := hu_support_sub.trans hΩ_sub
      have hgrad_support_sub : Function.support (smoothGradNorm u) ⊆ Ω :=
        support_smoothGradNorm_subset_tsupport.trans hu_sub
      have hgrad_support_sub_s : Function.support (smoothGradNorm u) ⊆ s :=
        hgrad_support_sub.trans hΩ_sub
      have hgrad_cont : Continuous (smoothGradNorm u) := by
        have hcont_fderiv : Continuous (fderiv ℝ u) :=
          hu.continuous_fderiv (by simp)
        have hsmooth_eq : smoothGradNorm u = fun x => ‖fderiv ℝ u x‖ := by
          ext x
          exact (norm_fderiv_eq_smoothGradNorm (u := u) (x := x)).symm
        rw [hsmooth_eq]
        exact continuous_norm.comp hcont_fderiv
      have hu_eq_full :
          eLpNorm u 2 (volume.restrict Ω) = eLpNorm u 2 volume := by
        rw [eLpNorm_restrict_eq_of_support_subset hu_support_sub]
      have hgrad_eq_full :
          eLpNorm (smoothGradNorm u) 2 (volume.restrict Ω) =
            eLpNorm (smoothGradNorm u) 2 volume := by
        rw [eLpNorm_restrict_eq_of_support_subset hgrad_support_sub]
      have hgrad_eq_s_one :
          eLpNorm (smoothGradNorm u) (ENNReal.ofReal 1) (volume.restrict s) =
            eLpNorm (smoothGradNorm u) (ENNReal.ofReal 1) volume := by
        rw [eLpNorm_restrict_eq_of_support_subset hgrad_support_sub_s]
      have hgrad_eq_s_two :
          eLpNorm (smoothGradNorm u) 2 (volume.restrict s) =
            eLpNorm (smoothGradNorm u) 2 volume := by
        rw [eLpNorm_restrict_eq_of_support_subset hgrad_support_sub_s]
      have hgrad_aesm_s :
          AEStronglyMeasurable (smoothGradNorm u) (volume.restrict s) :=
        hgrad_cont.aestronglyMeasurable
      have hsobolev :
        eLpNorm u 2 volume ≤
          (Csob : ℝ≥0∞) * eLpNorm (fderiv ℝ u) 1 volume := by
        have h2p : (1 : ℝ≥0) < Module.finrank ℝ E := by
          rw [finrank_euclideanSpace_fin, hd_eq_two]
          norm_num
        have hpq :
            ((1 : ℝ≥0) : ℝ)⁻¹ - (Module.finrank ℝ E : ℝ)⁻¹ ≤
              ((2 : ℝ≥0) : ℝ)⁻¹ := by
          rw [finrank_euclideanSpace_fin, hd_eq_two]
          norm_num
        simpa [Csob] using
          (MeasureTheory.eLpNorm_le_eLpNorm_fderiv_of_le
            (μ := volume) (F := ℝ) (u := u) (s := s) hu1 hu_support_sub_s
            (p := (1 : ℝ≥0)) (q := (2 : ℝ≥0)) (by norm_num) h2p hpq hs_bdd)
      have hgrad_eq_one :
          eLpNorm (fderiv ℝ u) 1 volume =
            eLpNorm (smoothGradNorm u) 1 volume := by
        calc
          eLpNorm (fderiv ℝ u) 1 volume
              = eLpNorm (fun x => ‖fderiv ℝ u x‖) 1 volume := by
                  exact (eLpNorm_norm (f := fderiv ℝ u) (p := (1 : ℝ≥0∞)) (μ := volume)).symm
          _ = eLpNorm (smoothGradNorm u) 1 volume := by
                exact eLpNorm_congr_ae (Eventually.of_forall fun x =>
                  norm_fderiv_eq_smoothGradNorm (u := u) (x := x))
      have hsobolev' :
        eLpNorm u 2 volume ≤
          (Csob : ℝ≥0∞) * eLpNorm (smoothGradNorm u) 1 volume := by
        rwa [hgrad_eq_one] at hsobolev
      have hcompare :
        eLpNorm (smoothGradNorm u) 1 (volume.restrict s) ≤
          eLpNorm (smoothGradNorm u) 2 (volume.restrict s) *
            (volume s) ^ (1 / (1 : ℝ) - 1 / (2 : ℝ)) := by
        simpa [s, Measure.restrict_apply_univ] using
          (eLpNorm_le_eLpNorm_mul_rpow_measure_univ
            (μ := volume.restrict s) (f := smoothGradNorm u)
            (p := (1 : ℝ≥0∞)) (q := (2 : ℝ≥0∞)) (by norm_num) hgrad_aesm_s)
      rw [hu_eq_full, hgrad_eq_full]
      calc
        eLpNorm u 2 volume
            ≤ (Csob : ℝ≥0∞) * eLpNorm (smoothGradNorm u) 1 volume := hsobolev'
        _ = (Csob : ℝ≥0∞) * eLpNorm (smoothGradNorm u) 1 (volume.restrict s) := by
              congr 1
              simpa using hgrad_eq_s_one.symm
        _ ≤ (Csob : ℝ≥0∞) *
              (eLpNorm (smoothGradNorm u) 2 (volume.restrict s) *
                (volume s) ^ (1 / (1 : ℝ) - 1 / (2 : ℝ))) := by
              gcongr
        _ = C * eLpNorm (smoothGradNorm u) 2 volume := by
              rw [hgrad_eq_s_two]
              dsimp [C]
              rw [mul_assoc, mul_comm (eLpNorm (smoothGradNorm u) 2 volume) ((volume s) ^ _)]
  · have hne : 2 ≠ d := by
      intro h
      exact hd_eq_two h.symm
    have hd_gt_two : 2 < d := lt_of_le_of_ne hd hne
    let Csob : ℝ≥0 := MeasureTheory.eLpNormLESNormFDerivOfLeConst ℝ volume s 2 2
    let C : ℝ≥0∞ := Csob
    refine ⟨C, ?_, ?_⟩
    · dsimp [C]
      exact ENNReal.coe_lt_top
    · intro u hu hu_cpt hu_sub
      have hu1 : ContDiff ℝ 1 u := hu.of_le (by norm_cast)
      have hu_support_sub : Function.support u ⊆ Ω := subset_closure.trans hu_sub
      have hu_support_sub_s : Function.support u ⊆ s := hu_support_sub.trans hΩ_sub
      have hgrad_support_sub : Function.support (smoothGradNorm u) ⊆ Ω :=
        support_smoothGradNorm_subset_tsupport.trans hu_sub
      have hu_eq_full :
          eLpNorm u 2 (volume.restrict Ω) = eLpNorm u 2 volume := by
        rw [eLpNorm_restrict_eq_of_support_subset hu_support_sub]
      have hgrad_eq_full :
          eLpNorm (smoothGradNorm u) 2 (volume.restrict Ω) =
            eLpNorm (smoothGradNorm u) 2 volume := by
        rw [eLpNorm_restrict_eq_of_support_subset hgrad_support_sub]
      have hsobolev :
        eLpNorm u 2 volume ≤
          (Csob : ℝ≥0∞) * eLpNorm (fderiv ℝ u) 2 volume := by
        have h2p : (2 : ℝ≥0) < Module.finrank ℝ E := by
          rw [finrank_euclideanSpace_fin]
          exact_mod_cast hd_gt_two
        simpa [Csob] using
          (MeasureTheory.eLpNorm_le_eLpNorm_fderiv
            (μ := volume) (F := ℝ) (u := u) (s := s) hu1 hu_support_sub_s
            (p := (2 : ℝ≥0)) (by norm_num) h2p hs_bdd)
      have hgrad_eq_two :
          eLpNorm (fderiv ℝ u) 2 volume =
            eLpNorm (smoothGradNorm u) 2 volume := by
        calc
          eLpNorm (fderiv ℝ u) 2 volume
              = eLpNorm (fun x => ‖fderiv ℝ u x‖) 2 volume := by
                  exact (eLpNorm_norm (f := fderiv ℝ u) (p := (2 : ℝ≥0∞)) (μ := volume)).symm
          _ = eLpNorm (smoothGradNorm u) 2 volume := by
                exact eLpNorm_congr_ae (Eventually.of_forall fun x =>
                  norm_fderiv_eq_smoothGradNorm (u := u) (x := x))
      have hsobolev' :
        eLpNorm u 2 volume ≤
          (Csob : ℝ≥0∞) * eLpNorm (smoothGradNorm u) 2 volume := by
        rwa [hgrad_eq_two] at hsobolev
      rw [hu_eq_full, hgrad_eq_full]
      calc
        eLpNorm u 2 volume
            ≤ (Csob : ℝ≥0∞) * eLpNorm (smoothGradNorm u) 2 volume := hsobolev'
        _ = C * eLpNorm (smoothGradNorm u) 2 volume := by
              rfl

\end{lstlisting}
\begin{proof}
Let $R$ be such that $\Omega \subseteq \overline{B}(0,R)$ and $s = \overline{B}(0,R)$.
When $d = 2$, the Mathlib Sobolev embedding
\texttt{eLpNormLESNormFDerivOfLeConst} provides the constant
$C_{\mathrm{sob}} = $ \texttt{eLpNormLESNormFDerivOfLeConst}\,$\R\,\mathrm{vol}\,s\,1\,2$
giving
$\lVert u\rVert_{L^{2}(s)} \le C_{\mathrm{sob}}\,\mathrm{vol}(s)^{1/1 - 1/2}
\lVert \mathrm{D}u\rVert_{L^{1}(s)}^{}$ after raising the exponent. For $d \ge 3$ one
applies the same Mathlib bound at the conjugate exponent $2$. In both cases we set
$C := C_{\mathrm{sob}}\,\mathrm{vol}(s)^{1/1 - 1/2}$ and conclude using the inclusion
$\supp u \subseteq \Omega \subseteq s$.
\end{proof}

The proof of the Poincar\'e inequality is via the classical representation formula
through spherical means. We collect the auxiliary lemmas before the main statement.

\begin{lemma}[\leanname{sphere\_radial\_integral\_ball}]
For $R > 0$ and a measurable, integrable, radial--bounded function $f$ on $\Ball{R}{x}$,
\[
\int_{\Ball{R}{x}} f(y)\, dy
= \int_{0}^{R} \int_{\partial \Ball{r}{x}} f \, d\sigma_{r}\, dr.
\]
\end{lemma}

\begin{lstlisting}[style=Matlab-editor,basicstyle=\mlttfamily,escapechar=",mlshowsectionrules=true,mathescape=true,morecomment={[s]{\%\{}{\%\}}},language=lean,numbers=none]
theorem sphere_radial_integral_ball
    {R : ℝ} (_hR : 0 < R) {F : E → ℝ}
    (hF : IntegrableOn F (Metric.ball (0 : E) R) volume) :
    ∫ z in Metric.ball (0 : E) R, F z ∂volume =
    ∫ ω in Set.univ,
      (∫ r in Set.Ioo (0 : ℝ) R, r ^ (d - 1 : ℕ) • F (r • (ω : E)))
      ∂(volume : Measure E).toSphere := by
  set G : E → ℝ := (Metric.ball (0 : E) R).indicator F
  have hstep1 : ∫ z in Metric.ball (0 : E) R, F z ∂volume = ∫ z, G z ∂volume := by
    simp [G, integral_indicator measurableSet_ball]
  have hstep1' : ∫ z, G z ∂volume =
      ∫ z : ({(0 : E)}ᶜ : Set E), G z ∂((volume : Measure E).comap (↑)) := by
    rw [integral_subtype_comap (measurableSet_singleton _).compl,
      restrict_compl_singleton]
  have hmp := (volume : Measure E).measurePreserving_homeomorphUnitSphereProd
  have hemb := (homeomorphUnitSphereProd E).measurableEmbedding
  have hstep2 :
      ∫ z : ({(0 : E)}ᶜ : Set E), G z ∂((volume : Measure E).comap (↑)) =
      ∫ p : Metric.sphere (0 : E) 1 × Set.Ioi (0 : ℝ),
        G (p.2.1 • (p.1 : E))
        ∂(volume : Measure E).toSphere.prod (Measure.volumeIoiPow (Module.finrank ℝ E - 1)) := by
    convert hmp.integral_comp hemb (fun p => G (p.2.1 • (p.1 : E))) using 1
    congr 1
    ext ⟨z, hz⟩
    simp [homeomorphUnitSphereProd, homeomorphSphereProd, Set.mem_compl_iff,
      Set.mem_singleton_iff] at hz ⊢
    congr 1
    have hne : ‖z‖ ≠ 0 := norm_ne_zero_iff.mpr hz
    rw [smul_smul, mul_inv_cancel₀ hne, one_smul]
  have hG_simp : ∀ (ω : Metric.sphere (0 : E) 1) (r : Set.Ioi (0 : ℝ)),
      G (r.1 • (ω : E)) = if r.1 < R then F (r.1 • (ω : E)) else 0 := by
    intro ⟨ω, hω⟩ ⟨r, hr⟩
    simp only [G, Set.indicator, Metric.mem_ball, dist_zero_right]
    have hω_norm : ‖ω‖ = 1 := by simpa using hω
    rw [norm_smul, Real.norm_of_nonneg (le_of_lt hr), hω_norm, mul_one]
  rw [hstep1, hstep1', hstep2]
  set σ := (volume : Measure E).toSphere
  set n := Module.finrank ℝ E - 1
  set ρ := Measure.volumeIoiPow n
  have hInteg : Integrable (fun p : Metric.sphere (0 : E) 1 × Set.Ioi (0 : ℝ) =>
      G (p.2.1 • (p.1 : E))) (σ.prod ρ) :=
    indicator_integrable_on_sphere_prod hF
  have hFubini :
      ∫ p : Metric.sphere (0 : E) 1 × Set.Ioi (0 : ℝ),
        G (p.2.1 • (p.1 : E)) ∂σ.prod ρ =
      ∫ ω : Metric.sphere (0 : E) 1,
        (∫ r : Set.Ioi (0 : ℝ), G (r.1 • (ω : E)) ∂ρ) ∂σ :=
    integral_prod
      (f := fun (p : Metric.sphere (0 : E) 1 × Set.Ioi (0 : ℝ)) =>
        G (p.2.1 • (p.1 : E)))
      hInteg
  rw [hFubini]
  rw [setIntegral_univ]
  congr 1
  ext ω
  simp only [ρ, Measure.volumeIoiPow]
  rw [integral_withDensity_eq_integral_toReal_smul
    ((measurable_subtype_coe.pow_const _).ennreal_ofReal)
    (by filter_upwards with ⟨r, hr⟩; exact ENNReal.ofReal_lt_top)]
  have h5b : ∫ r : Set.Ioi (0 : ℝ),
      (ENNReal.ofReal (r.1 ^ n)).toReal • G (r.1 • (ω : E))
      ∂(Measure.comap Subtype.val volume) =
      ∫ r in Set.Ioi (0 : ℝ), r ^ n • G (r • (ω : E)) ∂volume := by
    have : (fun (r : Set.Ioi (0 : ℝ)) =>
        (ENNReal.ofReal (r.1 ^ n)).toReal • G (r.1 • (ω : E))) =
        (fun (r : Set.Ioi (0 : ℝ)) => r.1 ^ n • G (r.1 • (ω : E))) := by
      ext ⟨r, hr⟩
      rw [ENNReal.toReal_ofReal (pow_nonneg hr.le _)]
    rw [this]
    exact integral_subtype_comap (α := ℝ) measurableSet_Ioi
      (fun r => r ^ n • G (r • (ω : E)))
  rw [h5b]
  have hfin : n = d - 1 := by simp [n]
  rw [show (∫ r in Set.Ioo (0 : ℝ) R, r ^ (d - 1 : ℕ) • F (r • (ω : E))) =
    ∫ r in Set.Ioi (0 : ℝ),
      (Set.Ioo (0 : ℝ) R).indicator (fun r => r ^ (d - 1 : ℕ) • F (r • (ω : E))) r from by
    rw [setIntegral_indicator measurableSet_Ioo,
      show Set.Ioi (0 : ℝ) ∩ Set.Ioo 0 R = Set.Ioo 0 R from
        Set.inter_eq_right.mpr Set.Ioo_subset_Ioi_self]]
  apply setIntegral_congr_fun measurableSet_Ioi
  intro r hr
  rw [hfin]
  have hr_pos : 0 < r := hr
  simp only [Set.indicator, Set.mem_Ioo, hr_pos, true_and, smul_eq_mul]
  rw [hG_simp ω ⟨r, hr⟩]
  by_cases hlt : r < R
  · simp [hlt]
  · simp [hlt, mul_zero]

\end{lstlisting}
\begin{lemma}[\leanname{polar\_identity\_ball}, \leanname{star\_shaped\_polar\_identity}]
Polar coordinates expansion in $\Ball{R}{x}$: any $y \in \Ball{R}{x}$ can be written
$y = x + r\omega$ with $r \in (0,R)$, $\omega \in S^{d-1}$, and the volume element
factors as $dy = r^{d-1}\, dr\, d\sigma(\omega)$.
\end{lemma}

\begin{lstlisting}[style=Matlab-editor,basicstyle=\mlttfamily,escapechar=",mlshowsectionrules=true,mathescape=true,morecomment={[s]{\%\{}{\%\}}},language=lean,numbers=none]
theorem polar_identity_ball
    {R : ℝ} (hR : 0 < R) {f : E → ℝ}
    (hf : IntegrableOn (fun z => f z * ‖z‖ ^ (1 - (d : ℝ))) (Metric.ball (0 : E) R) volume) :
    ∫ ω in Set.univ,
      (∫ s in Set.Ioo (0 : ℝ) R, f (s • (ω : E)))
      ∂(volume : Measure E).toSphere =
    ∫ z in Metric.ball (0 : E) R, f z * ‖z‖ ^ (1 - (d : ℝ)) ∂volume := by
  set g : E → ℝ := fun z => f z * ‖z‖ ^ (1 - (d : ℝ))
  have hpolar := sphere_radial_integral_ball hR hf
  rw [hpolar]
  congr 1
  ext ω
  apply setIntegral_congr_fun measurableSet_Ioo
  intro r ⟨hr_pos, hr_lt⟩
  simp only [g, smul_eq_mul]
  have hω_norm : ‖(ω : E)‖ = 1 := by
    have hm := ω.2
    simp only [Metric.mem_sphere, dist_zero_right] at hm
    exact hm
  rw [norm_smul, Real.norm_of_nonneg hr_pos.le, hω_norm, mul_one]
  have hd_ge : 1 ≤ d := Nat.one_le_iff_ne_zero.mpr (NeZero.ne d)
  have hrpow_nat : (r : ℝ) ^ (d - 1 : ℕ) * r ^ (1 - (d : ℝ)) = 1 := by
    rw [← Real.rpow_natCast r (d - 1), Nat.cast_sub hd_ge, ← Real.rpow_add hr_pos]
    norm_num
  have : r ^ (d - 1 : ℕ) * (f (r • (ω : E)) * r ^ (1 - (d : ℝ))) = f (r • (ω : E)) := by
    rw [mul_comm (f _) _, ← mul_assoc, hrpow_nat, one_mul]
  linarith

theorem star_shaped_polar_identity
    {R : ℝ} (hR : 0 < R) {x : E} (hx : x ∈ Metric.ball (0 : E) R)
    {f : E → ℝ} (hf : IntegrableOn (fun y => f y * ‖y - x‖ ^ (1 - (d : ℝ)))
      (Metric.ball (0 : E) R) volume) :
    ∫ ω in Set.univ,
      (∫ s in Set.Ioo (0 : ℝ) (2 * R),
        (Metric.ball (0 : E) R).indicator (fun y => f y) (x + s • (ω : E)))
      ∂(volume : Measure E).toSphere =
    ∫ y in Metric.ball (0 : E) R, f y * ‖y - x‖ ^ (1 - (d : ℝ)) ∂volume := by
  set B := Metric.ball (0 : E) R
  set g : E → ℝ := fun z => B.indicator (fun y => f y) (x + z)
  have hpolar := polar_identity_ball (show (0 : ℝ) < 2 * R by linarith)
    (f := g) (hf := by
      have heq : (fun z => g z * ‖z‖ ^ (1 - (d : ℝ))) =ᵐ[volume.restrict (Metric.ball 0 (2*R))]
          (fun z => B.indicator (fun y => f y * ‖y - x‖ ^ (1 - (d : ℝ))) (x + z)) := by
        filter_upwards with z
        simp only [g, Set.indicator]
        split_ifs with h
        · congr 1
          congr 1
          show ‖z‖ = ‖(x + z) - x‖
          abel_nf
        · simp
      rw [IntegrableOn, integrable_congr heq]
      apply Integrable.integrableOn
      set F := B.indicator (fun y => f y * ‖y - x‖ ^ (1 - (d : ℝ)))
      have hF_int : Integrable F volume := (integrable_indicator_iff measurableSet_ball).mpr hf
      rw [show (fun z => F (x + z)) = F ∘ (· + x) from by ext z; rw [add_comm]; rfl]
      exact (measurePreserving_add_right volume x).integrable_comp_of_integrable hF_int)
  have hRHS : ∫ z in Metric.ball (0 : E) (2 * R), g z * ‖z‖ ^ (1 - (d : ℝ)) ∂volume =
      ∫ y in B, f y * ‖y - x‖ ^ (1 - (d : ℝ)) ∂volume := by
    have hx_norm : ‖x‖ < R := by rwa [Metric.mem_ball, dist_zero_right] at hx
    have hfun_eq : ∀ z, g z * ‖z‖ ^ (1 - (d : ℝ)) =
        B.indicator (fun y => f y * ‖y - x‖ ^ (1 - (d : ℝ))) (x + z) := by
      intro z
      simp only [g, Set.indicator]
      split_ifs with h
      · congr 1
        show ‖z‖ ^ (1 - (d : ℝ)) = ‖x + z - x‖ ^ (1 - (d : ℝ))
        congr 1
        abel_nf
      · ring
    conv_lhs =>
      arg 2
      ext z
      rw [hfun_eq]
    set h : E → ℝ := B.indicator (fun y => f y * ‖y - x‖ ^ (1 - (d : ℝ)))
    calc ∫ z in Metric.ball (0 : E) (2 * R), h (x + z) ∂volume
        = ∫ z, h (x + z) ∂volume := by
          rw [← integral_indicator measurableSet_ball]
          apply integral_congr_ae
          filter_upwards with z
          simp only [Set.indicator]
          split_ifs with hz
          · rfl
          · simp only [h, Set.indicator]
            split_ifs with hb
            · exfalso
              apply hz
              simp only [B, Metric.mem_ball, dist_zero_right] at hb ⊢
              calc ‖z‖ ≤ ‖x + z‖ + ‖x‖ := by
                    calc ‖z‖ = ‖(x + z) + (-x)‖ := by abel_nf
                      _ ≤ ‖x + z‖ + ‖-x‖ := norm_add_le _ _
                      _ = ‖x + z‖ + ‖x‖ := by rw [norm_neg]
                _ < R + R := add_lt_add hb hx_norm
                _ = 2 * R := by ring
            · rfl
      _ = ∫ y, h y ∂volume := by
          rw [show (fun z => h (x + z)) = (fun z => h (z + x)) from by ext z; rw [add_comm]]
          exact integral_add_right_eq_self _ x
      _ = ∫ y in B, f y * ‖y - x‖ ^ (1 - (d : ℝ)) ∂volume := by
          exact integral_indicator measurableSet_ball
  calc ∫ ω in Set.univ,
        (∫ s in Set.Ioo (0 : ℝ) (2 * R), B.indicator (fun y => f y) (x + s • (ω : E)))
        ∂(volume : Measure E).toSphere
      = ∫ ω in Set.univ,
        (∫ s in Set.Ioo (0 : ℝ) (2 * R), g (s • (ω : E)))
        ∂(volume : Measure E).toSphere := by rfl
    _ = ∫ z in Metric.ball (0 : E) (2 * R), g z * ‖z‖ ^ (1 - (d : ℝ)) ∂volume :=
        hpolar
    _ = ∫ y in B, f y * ‖y - x‖ ^ (1 - (d : ℝ)) ∂volume := hRHS

\end{lstlisting}
\begin{lemma}[\leanname{norm\_sub\_le\_integral\_fderiv\_along\_segment}]
For $u \in C^{1}$ and any segment $[x,y] \subseteq E$,
\[
\lvert u(y) - u(x)\rvert
\le \int_{0}^{1} \lvert \nabla u(x + t(y-x))\rvert\, \lvert y-x\rvert\, dt.
\]
\end{lemma}

\begin{lstlisting}[style=Matlab-editor,basicstyle=\mlttfamily,escapechar=",mlshowsectionrules=true,mathescape=true,morecomment={[s]{\%\{}{\%\}}},language=lean,numbers=none]
omit [NeZero d] in
theorem norm_sub_le_integral_fderiv_along_segment
    {u : E → ℝ} (hu : ContDiff ℝ (⊤ : ℕ∞) u) (x y : E) :
    ‖u x - u y‖ ≤ ∫ t in (0 : ℝ)..‖x - y‖,
      ‖fderiv ℝ u (x + t • (‖x - y‖⁻¹ • (y - x)))‖ := by
  by_cases hxy : x = y
  · simp [hxy]
  set ω : E := ‖x - y‖⁻¹ • (y - x)
  set γ : ℝ → E := fun t => x + t • ω
  have hγ : ∀ t, HasDerivAt γ ω t := by
    intro t
    have : HasDerivAt (fun t => t • ω) ω t := by
      simpa using (hasDerivAt_id t).smul_const ω
    exact this.const_add x
  have huγ : ∀ t, HasDerivAt (u ∘ γ) ((fderiv ℝ u (γ t)) ω) t := by
    intro t
    exact (hu.differentiable (by norm_cast)).differentiableAt.hasFDerivAt.comp_hasDerivAt t
      (hγ t)
  have hγ0 : γ 0 = x := by simp [γ]
  have hγ1 : γ ‖x - y‖ = y := by
    simp only [γ, ω, smul_smul]
    rw [mul_inv_cancel₀ (norm_ne_zero_iff.mpr (sub_ne_zero.mpr hxy)), one_smul]
    abel
  have hfderiv_cont : Continuous (fderiv ℝ u) := by
    have h1 : ContDiff ℝ 1 u := hu.of_le (by norm_cast)
    exact h1.continuous_fderiv (by norm_num)
  have hγ_cont : Continuous γ := by fun_prop
  have hcont_comp : Continuous (fun t => (fderiv ℝ u (γ t)) ω) :=
    (hfderiv_cont.comp hγ_cont).clm_apply continuous_const
  have hint1 : IntervalIntegrable (fun t => (fderiv ℝ u (γ t)) ω) volume 0 ‖x - y‖ :=
    hcont_comp.intervalIntegrable _ _
  have hint2 : IntervalIntegrable (fun t => ‖fderiv ℝ u (γ t)‖) volume 0 ‖x - y‖ :=
    (continuous_norm.comp (hfderiv_cont.comp hγ_cont)).intervalIntegrable _ _
  have hftc : u y - u x = ∫ t in (0 : ℝ)..‖x - y‖, (fderiv ℝ u (γ t)) ω := by
    have := intervalIntegral.integral_eq_sub_of_hasDerivAt (fun t _ => huγ t) hint1
    simp only [Function.comp_apply, hγ0, hγ1] at this
    linarith
  have hω_norm : ‖ω‖ = 1 := by
    show ‖‖x - y‖⁻¹ • (y - x)‖ = 1
    rw [norm_smul, norm_inv, norm_norm, norm_sub_rev y x,
      inv_mul_cancel₀ (norm_ne_zero_iff.mpr (sub_ne_zero.mpr hxy))]
  rw [norm_sub_rev]
  calc ‖u y - u x‖ = ‖∫ t in (0 : ℝ)..‖x - y‖, (fderiv ℝ u (γ t)) ω‖ := by rw [hftc]
    _ ≤ ∫ t in (0 : ℝ)..‖x - y‖, ‖(fderiv ℝ u (γ t)) ω‖ :=
        intervalIntegral.norm_integral_le_integral_norm (norm_nonneg _)
    _ ≤ ∫ t in (0 : ℝ)..‖x - y‖, ‖fderiv ℝ u (γ t)‖ := by
        apply intervalIntegral.integral_mono_on (norm_nonneg _) hint1.norm hint2
        intro t ht
        calc ‖(fderiv ℝ u (γ t)) ω‖ ≤ ‖fderiv ℝ u (γ t)‖ * ‖ω‖ :=
              ContinuousLinearMap.le_opNorm _ _
          _ = ‖fderiv ℝ u (γ t)‖ := by rw [hω_norm, mul_one]

\end{lstlisting}
\begin{lemma}[\leanname{integral\_norm\_rpow\_neg\_ball},
\leanname{riesz\_kernel\_integrable}, \leanname{riesz\_kernel\_L1\_bound}]
For $R > 0$ and $0 < \alpha < d$ the function $z\mapsto \lvert z\rvert^{\alpha-d}$ is
integrable on $\Ball{R}{0}$, and for $x \in \Ball{R}{0}$ the Riesz kernel satisfies
\[
\int_{\Ball{R}{0}} \lvert x - y\rvert^{1-d}\, dy
\le d\, \mathrm{vol}(\Bz{1})\, (2R).
\]
\end{lemma}

\begin{lstlisting}[style=Matlab-editor,basicstyle=\mlttfamily,escapechar=",mlshowsectionrules=true,mathescape=true,morecomment={[s]{\%\{}{\%\}}},language=lean,numbers=none]
theorem integral_norm_rpow_neg_ball {R : ℝ} (hR : 0 < R) :
    ∫ x in Metric.ball (0 : E) R, ‖x‖ ^ (1 - (d : ℝ)) ∂volume =
    ↑d * (volume (Metric.ball (0 : E) 1)).toReal * R := by
  let f : ℝ → ℝ := fun r => if 0 < r ∧ r < R then r ^ (1 - (d : ℝ)) else 0
  have hconv : ∫ x in Metric.ball (0 : E) R, ‖x‖ ^ (1 - (d : ℝ)) ∂volume =
      ∫ x : E, f (‖x‖) ∂volume := by
    rw [← integral_indicator measurableSet_ball]
    refine integral_congr_ae ?_
    have hmeas0 : (volume : Measure E) {(0 : E)} = 0 := measure_singleton _
    have hae : ∀ᵐ x ∂(volume : Measure E), x ≠ (0 : E) :=
      compl_mem_ae_iff.mpr hmeas0
    filter_upwards [hae] with x hx
    simp only [f, indicator, Metric.mem_ball, dist_zero_right]
    have hpos : 0 < ‖x‖ := norm_pos_iff.mpr hx
    simp only [and_iff_right hpos]
  have hrad : ∫ x : E, f (‖x‖) ∂volume =
      ↑(Module.finrank ℝ E) • (volume : Measure E).real (Metric.ball 0 1) •
        ∫ y in Set.Ioi (0 : ℝ), y ^ (Module.finrank ℝ E - 1) • f y :=
    MeasureTheory.integral_fun_norm_addHaar (μ := (volume : Measure E)) f
  have hfin : Module.finrank ℝ E = d := finrank_euclideanSpace_fin
  have h1d : ∫ y in Set.Ioi (0 : ℝ), y ^ (Module.finrank ℝ E - 1) • f y = R := by
    rw [hfin]
    have hsupp : ∀ y ∈ Set.Ioi (0 : ℝ), y ^ (d - 1) • f y =
        Set.indicator (Set.Ioo 0 R) (fun _ => (1 : ℝ)) y := by
      intro y hy
      have hy_pos : 0 < y := hy
      by_cases hlt : y < R
      · simp only [f, smul_eq_mul, Set.indicator, Set.mem_Ioo, hy_pos, hlt, true_and,
          if_true]
        rw [← Real.rpow_natCast y (d - 1),
          Nat.cast_sub (Nat.one_le_iff_ne_zero.mpr (NeZero.ne d)),
          ← Real.rpow_add hy_pos]
        norm_num
      · simp only [f, smul_eq_mul, Set.indicator, Set.mem_Ioo, hy_pos, hlt, true_and,
          if_false, mul_zero]
    rw [setIntegral_congr_fun measurableSet_Ioi hsupp]
    rw [setIntegral_indicator measurableSet_Ioo]
    rw [show Set.Ioi (0 : ℝ) ∩ Set.Ioo 0 R = Set.Ioo 0 R from
      Set.inter_eq_right.mpr Set.Ioo_subset_Ioi_self]
    simp only [setIntegral_const, smul_eq_mul, mul_one, Measure.real,
      Real.volume_Ioo, ENNReal.toReal_ofReal (by linarith : (0 : ℝ) ≤ R - 0)]
    ring
  rw [hconv, hrad, h1d, hfin]
  simp [nsmul_eq_mul, Measure.real]
  ring

theorem riesz_kernel_integrable
    {R : ℝ} (_hR : 0 < R) :
    IntegrableOn (fun x : E => ‖x‖ ^ (1 - (d : ℝ)))
      (Metric.ball (0 : E) R) volume := by
  let g : ℝ → ℝ := fun r => if r < R then r ^ (1 - (d : ℝ)) else 0
  have hag : (fun x : E => ‖x‖ ^ (1 - (d : ℝ))) =ᵐ[volume.restrict (Metric.ball (0 : E) R)]
      (g ∘ (‖·‖)) := by
    filter_upwards [ae_restrict_mem measurableSet_ball] with x hx
    simp only [Function.comp_apply, g, Metric.mem_ball, dist_zero_right] at hx ⊢
    rw [if_pos hx]
  rw [IntegrableOn, integrable_congr hag]
  suffices h : Integrable (fun x : E => g ‖x‖) volume from h.integrableOn
  have h1d : IntegrableOn (fun y : ℝ => y ^ (Module.finrank ℝ E - 1) • g y) (Set.Ioi 0) := by
    have hfin : Module.finrank ℝ E = d := finrank_euclideanSpace_fin
    set h_ind : ℝ → ℝ := (Set.Ioo (0 : ℝ) R).indicator (fun _ => (1 : ℝ))
    have heq : Set.EqOn (fun y : ℝ => y ^ (Module.finrank ℝ E - 1) • g y)
        h_ind (Set.Ioi 0) := by
      intro r hr
      simp only [Set.mem_Ioi] at hr
      simp only [g, smul_eq_mul, h_ind, Set.indicator, Set.mem_Ioo]
      split_ifs with h1 h2
      · rw [hfin, ← Real.rpow_natCast r (d - 1),
          Nat.cast_sub (Nat.one_le_iff_ne_zero.mpr (NeZero.ne d)),
          ← Real.rpow_add hr]
        norm_num
      · next h2 => exact absurd ⟨hr, h1⟩ h2
      · next h1 h2 => exact absurd h2.2 h1
      · simp [mul_zero]
    have h_ind_int : IntegrableOn h_ind (Set.Ioi 0) := by
      apply Integrable.integrableOn
      exact integrable_indicator_iff measurableSet_Ioo |>.mpr (integrableOn_const (hs := measure_Ioo_lt_top.ne))
    exact h_ind_int.congr_fun heq.symm measurableSet_Ioi
  exact (MeasureTheory.integrable_fun_norm_addHaar (μ := volume) (f := g)).mpr h1d

theorem riesz_kernel_L1_bound
    {R : ℝ} (hR : 0 < R) (x : E) (hx : x ∈ Metric.ball (0 : E) R) :
    ∫ y in Metric.ball (0 : E) R,
      ‖x - y‖ ^ (1 - (d : ℝ)) ∂volume ≤
    ↑d * (volume (Metric.ball (0 : E) 1)).toReal * (2 * R) := by
  have hx_dist : dist x 0 < R := Metric.mem_ball.mp hx
  have hBsub : Metric.ball (0 : E) R ⊆ Metric.ball x (2 * R) := by
    intro y hy
    rw [Metric.mem_ball] at hy ⊢
    calc dist y x ≤ dist y 0 + dist 0 x := dist_triangle y 0 x
      _ < R + R := by rw [dist_comm] at hx_dist; linarith [Metric.mem_ball.mp hy]
      _ = 2 * R := by ring
  calc ∫ y in Metric.ball (0 : E) R, ‖x - y‖ ^ (1 - (d : ℝ)) ∂volume
      ≤ ∫ z in Metric.ball (0 : E) (2 * R), ‖z‖ ^ (1 - (d : ℝ)) ∂volume := by
        have hmp := measurePreserving_add_right (volume : Measure E) x
        have hemb := (MeasurableEquiv.addRight x : E ≃ᵐ E).measurableEmbedding
        have hpre : (· + x) ⁻¹' Metric.ball x (2 * R) = Metric.ball (0 : E) (2 * R) := by
          ext z; simp [Metric.mem_ball]
        have htrans : ∫ y in Metric.ball x (2 * R), ‖x - y‖ ^ (1 - (d : ℝ)) ∂volume =
            ∫ z in Metric.ball (0 : E) (2 * R), ‖z‖ ^ (1 - (d : ℝ)) ∂volume := by
          rw [← hmp.setIntegral_preimage_emb hemb, hpre]
          congr 1 with z
          show ‖x - (z + x)‖ ^ (1 - (d : ℝ)) = ‖z‖ ^ (1 - (d : ℝ))
          simp [norm_neg]
        have hinteg : IntegrableOn (fun y => ‖x - y‖ ^ (1 - (d : ℝ)))
            (Metric.ball x (2 * R)) volume := by
          have h2R : (0 : ℝ) < 2 * R := by linarith
          rw [← hmp.integrableOn_comp_preimage hemb, hpre]
          exact (riesz_kernel_integrable h2R).congr
            (ae_of_all _ (fun z => by
              show ‖z‖ ^ (1 - (d : ℝ)) = ‖x - (z + x)‖ ^ (1 - (d : ℝ))
              simp [norm_neg]))
        calc ∫ y in Metric.ball (0 : E) R, ‖x - y‖ ^ (1 - (d : ℝ)) ∂volume
            ≤ ∫ y in Metric.ball x (2 * R), ‖x - y‖ ^ (1 - (d : ℝ)) ∂volume := by
              apply setIntegral_mono_set hinteg
              · exact ae_of_all _ (fun y => Real.rpow_nonneg (norm_nonneg _) _)
              · exact hBsub.eventuallyLE
          _ = ∫ z in Metric.ball (0 : E) (2 * R), ‖z‖ ^ (1 - (d : ℝ)) ∂volume := htrans
    _ = ↑d * (volume (Metric.ball (0 : E) 1)).toReal * (2 * R) := by
        exact integral_norm_rpow_neg_ball (by linarith)

\end{lstlisting}
\begin{lemma}[Representation formula, \leanname{representation\_formula\_smooth}]
Let $u \in C^{\infty}$ and let $\bar u = \fint_{\Bz{1}} u$. For every $x \in \Bz{1}$,
\[
\lvert u(x) - \bar u\rvert
\le \frac{2^{d}}{d\,\mathrm{vol}(\Bz{1})}\,
\int_{\Bz{1}} \lvert \nabla u(y)\rvert\, \lvert x - y\rvert^{1-d}\, dy.
\]
\end{lemma}

\begin{lstlisting}[style=Matlab-editor,basicstyle=\mlttfamily,escapechar=",mlshowsectionrules=true,mathescape=true,morecomment={[s]{\%\{}{\%\}}},language=lean,numbers=none]
theorem representation_formula_smooth
    {R : ℝ} (hR : 0 < R)
    {u : E → ℝ} (hu : ContDiff ℝ (⊤ : ℕ∞) u) (x : E) (hx : x ∈ Metric.ball (0 : E) R) :
    ‖u x - ⨍ y in Metric.ball (0 : E) R, u y ∂volume‖ ≤
    (2 : ℝ) ^ d / ((d : ℝ) * (volume (Metric.ball (0 : E) 1)).toReal) *
      ∫ y in Metric.ball (0 : E) R,
        ‖fderiv ℝ u y‖ * ‖x - y‖ ^ (1 - (d : ℝ)) ∂volume := by
  set B := Metric.ball (0 : E) R
  set ubar := ⨍ y in B, u y ∂volume
  have h1 : ‖u x - ubar‖ ≤ ⨍ y in B, ‖u x - u y‖ ∂volume := by
    have hB_ne : volume B ≠ 0 := by
      simp [B]
      exact (measure_ball_pos volume 0 hR).ne'
    have hB_fin : volume B ≠ ⊤ := measure_ball_lt_top.ne
    have hu_int : IntegrableOn u B volume :=
      (hu.continuous.continuousOn.integrableOn_compact (isCompact_closedBall 0 R)).mono_set
        Metric.ball_subset_closedBall
    set μ := volume.restrict B
    haveI : IsFiniteMeasure μ := ⟨by rw [Measure.restrict_apply_univ]; exact hB_fin.lt_top⟩
    haveI : NeZero μ := ⟨by rwa [ne_eq, Measure.restrict_eq_zero]⟩
    have havg_eq : u x - ubar = ⨍ y, (u x - u y) ∂μ := by
      show u x - ⨍ y, u y ∂μ = ⨍ y, (u x - u y) ∂μ
      rw [average_eq, average_eq,
        integral_sub (integrable_const (u x)) hu_int.integrable,
        integral_const, smul_sub]
      have hne : (μ.real Set.univ) ≠ 0 := by
        rw [measureReal_def, Measure.restrict_apply_univ]
        exact ENNReal.toReal_ne_zero.mpr ⟨hB_ne, hB_fin⟩
      rw [inv_smul_smul₀ hne]
    rw [havg_eq]
    show ‖⨍ y, (u x - u y) ∂μ‖ ≤ ⨍ y, ‖u x - u y‖ ∂μ
    simp only [average_eq]
    calc ‖(μ.real Set.univ)⁻¹ • ∫ y, (u x - u y) ∂μ‖
        ≤ ‖(μ.real Set.univ)⁻¹‖ * ‖∫ y, (u x - u y) ∂μ‖ := norm_smul_le _ _
      _ ≤ ‖(μ.real Set.univ)⁻¹‖ * ∫ y, ‖u x - u y‖ ∂μ :=
          mul_le_mul_of_nonneg_left (norm_integral_le_integral_norm _) (norm_nonneg _)
      _ = (μ.real Set.univ)⁻¹ • ∫ y, ‖u x - u y‖ ∂μ := by
          rw [Real.norm_of_nonneg (inv_nonneg.mpr measureReal_nonneg), smul_eq_mul]
  set C_rep := (2 : ℝ) ^ d / ((d : ℝ) * (volume (Metric.ball (0 : E) 1)).toReal)
  suffices h2 : ⨍ y in B, ‖u x - u y‖ ∂volume ≤
      C_rep * ∫ y in B, ‖fderiv ℝ u y‖ * ‖x - y‖ ^ (1 - (d : ℝ)) ∂volume from
    le_trans h1 h2
  have hB_ne : volume B ≠ 0 := by
    simp [B]; exact (measure_ball_pos volume 0 hR).ne'
  have hB_fin : volume B ≠ ⊤ := measure_ball_lt_top.ne
  have hvol_pos : (0 : ℝ) < (volume B).toReal :=
    ENNReal.toReal_pos hB_ne hB_fin
  have hd_pos : (0 : ℝ) < d := Nat.cast_pos.mpr (NeZero.pos d)
  have hvol1_pos : (0 : ℝ) < (volume (Metric.ball (0 : E) 1)).toReal :=
    ENNReal.toReal_pos (measure_ball_pos volume 0 one_pos).ne' measure_ball_lt_top.ne
  set volB := (volume B).toReal
  have hH_int : IntegrableOn (fun y : E => ‖u x - u y‖) B volume := by
    have hcont : Continuous (fun y : E => ‖u x - u y‖) :=
      continuous_norm.comp (continuous_const.sub hu.continuous)
    exact (hcont.continuousOn.integrableOn_compact (isCompact_closedBall 0 R)).mono_set
      Metric.ball_subset_closedBall
  have hfderiv_cont : Continuous (fderiv ℝ u) := by
    have h1u : ContDiff ℝ 1 u := hu.of_le (by norm_cast)
    exact h1u.continuous_fderiv (by norm_num)
  have hgrad_cont : Continuous (fun y : E => ‖fderiv ℝ u y‖) :=
    continuous_norm.comp hfderiv_cont
  obtain ⟨C0, hC0⟩ :=
    IsCompact.exists_bound_of_continuousOn (isCompact_closedBall (0 : E) R) hgrad_cont.continuousOn
  set M : ℝ := max C0 0
  have hM_nonneg : 0 ≤ M := le_max_right _ _
  have hM_bound : ∀ y ∈ B, ‖fderiv ℝ u y‖ ≤ M := by
    intro y hy
    have hy' : y ∈ Metric.closedBall (0 : E) R := Metric.ball_subset_closedBall hy
    have hbound := hC0 y hy'
    rw [Real.norm_of_nonneg (norm_nonneg _)] at hbound
    exact hbound.trans (le_max_left _ _)
  have hBsub : B ⊆ Metric.ball x (2 * R) := by
    intro y hy
    rw [Metric.mem_ball] at hy ⊢
    have hx_dist : dist x 0 < R := Metric.mem_ball.mp hx
    calc dist y x ≤ dist y 0 + dist 0 x := dist_triangle y 0 x
      _ < R + R := by rw [dist_comm] at hx_dist; linarith [Metric.mem_ball.mp hy]
      _ = 2 * R := by ring
  have hk_big : IntegrableOn (fun y : E => ‖y - x‖ ^ (1 - (d : ℝ)))
      (Metric.ball x (2 * R)) volume := by
    have hmp := measurePreserving_add_right (volume : Measure E) x
    have hemb := (MeasurableEquiv.addRight x : E ≃ᵐ E).measurableEmbedding
    have hpre : (· + x) ⁻¹' Metric.ball x (2 * R) = Metric.ball (0 : E) (2 * R) := by
      ext z
      simp [Metric.mem_ball]
    rw [← hmp.integrableOn_comp_preimage hemb, hpre]
    have h2R : (0 : ℝ) < 2 * R := by linarith
    exact (riesz_kernel_integrable (d := d) h2R).congr
      (ae_of_all _ (fun z => by
        show ‖z‖ ^ (1 - (d : ℝ)) = ‖(z + x) - x‖ ^ (1 - (d : ℝ))
        abel_nf))
  have hk : IntegrableOn (fun y : E => ‖y - x‖ ^ (1 - (d : ℝ))) B volume :=
    hk_big.mono_set hBsub
  have hfG : IntegrableOn (fun y : E => ‖fderiv ℝ u y‖ * ‖y - x‖ ^ (1 - (d : ℝ))) B volume := by
    rw [IntegrableOn] at hk ⊢
    have hmajor : Integrable (fun y : E => M * ‖y - x‖ ^ (1 - (d : ℝ))) (volume.restrict B) :=
      hk.const_mul M
    apply Integrable.mono' hmajor
    · have hmeas :
          Measurable (fun y : E => ‖fderiv ℝ u y‖ * ‖y - x‖ ^ (1 - (d : ℝ))) := by
          measurability
      exact hmeas.aestronglyMeasurable.restrict
    · filter_upwards [ae_restrict_mem measurableSet_ball] with y hy
      have hk_nonneg : 0 ≤ ‖y - x‖ ^ (1 - (d : ℝ)) := Real.rpow_nonneg (norm_nonneg _) _
      have hleft_nonneg : 0 ≤ ‖fderiv ℝ u y‖ * ‖y - x‖ ^ (1 - (d : ℝ)) :=
        mul_nonneg (norm_nonneg _) hk_nonneg
      have hright_nonneg : 0 ≤ M * ‖y - x‖ ^ (1 - (d : ℝ)) :=
        mul_nonneg hM_nonneg hk_nonneg
      simpa [Real.norm_of_nonneg hleft_nonneg, Real.norm_of_nonneg hright_nonneg] using
        (mul_le_mul_of_nonneg_right (hM_bound y hy) hk_nonneg)
  have houterG_int := integrable_star_shaped_radial
    (d := d) (R := R) hR (x := x) hx (f := fun y => ‖fderiv ℝ u y‖) hfG
  have hkernel_symm :
      ∫ y in B, ‖fderiv ℝ u y‖ * ‖y - x‖ ^ (1 - (d : ℝ)) ∂volume =
      ∫ y in B, ‖fderiv ℝ u y‖ * ‖x - y‖ ^ (1 - (d : ℝ)) ∂volume := by
    apply setIntegral_congr_fun measurableSet_ball
    intro y hy
    simp [norm_sub_rev]
  have hrad_u :
      ∫ ω in Set.univ,
        (∫ s in Set.Ioo (0 : ℝ) (2 * R),
          s ^ (d - 1 : ℕ) • B.indicator (fun y => ‖u x - u y‖) (x + s • (ω : E)) ∂volume)
        ∂(volume : Measure E).toSphere
      = ∫ y in B, ‖u x - u y‖ ∂volume := by
    simpa [B] using translated_ball_sphere_radial_identity
      (d := d) (R := R) hR (x := x) hx (f := fun y => ‖u x - u y‖) hH_int
  have hpointwise :
      ∀ ω : Metric.sphere (0 : E) 1,
        ∫ s in Set.Ioo (0 : ℝ) (2 * R),
          s ^ (d - 1 : ℕ) • B.indicator (fun y => ‖u x - u y‖) (x + s • (ω : E)) ∂volume
        ≤ ((2 * R) ^ d / ↑d) *
          ∫ s in Set.Ioo (0 : ℝ) (2 * R),
            B.indicator (fun y => ‖fderiv ℝ u y‖) (x + s • (ω : E)) ∂volume := by
    intro ω
    set Cω := ∫ s in Set.Ioo (0 : ℝ) (2 * R),
      B.indicator (fun y => ‖fderiv ℝ u y‖) (x + s • (ω : E)) ∂volume
    have hCω_nonneg : 0 ≤ Cω := by
      apply integral_nonneg
      intro s
      by_cases hs_ball : x + s • (ω : E) ∈ B <;> simp [B, hs_ball]
    have hpow_int :
        IntegrableOn (fun s : ℝ => Cω * s ^ (d - 1 : ℕ))
          (Set.Ioo (0 : ℝ) (2 * R)) volume := by
      have hcont : Continuous (fun s : ℝ => Cω * s ^ (d - 1 : ℕ)) :=
        continuous_const.mul (continuous_id.pow _)
      have hsubset : Set.Ioo (0 : ℝ) (2 * R) ⊆ Set.Icc (0 : ℝ) (2 * R) := by
        intro s hs
        exact ⟨le_of_lt hs.1, le_of_lt hs.2⟩
      exact (hcont.continuousOn.integrableOn_compact
        (isCompact_Icc : IsCompact (Set.Icc (0 : ℝ) (2 * R)))).mono_set hsubset
    have hnonneg_lhs :
        0 ≤ᵐ[volume.restrict (Set.Ioo (0 : ℝ) (2 * R))]
          (fun s : ℝ => s ^ (d - 1 : ℕ) • B.indicator (fun y => ‖u x - u y‖) (x + s • (ω : E))) := by
      filter_upwards [ae_restrict_mem measurableSet_Ioo] with s hs
      by_cases hs_ball : x + s • (ω : E) ∈ B
      · simpa [B, hs_ball, smul_eq_mul] using
          mul_nonneg (pow_nonneg hs.1.le _) (norm_nonneg (u x - u (x + s • (ω : E))))
      · simp [B, hs_ball, smul_eq_mul]
    have hrad_le :
        ∫ s in Set.Ioo (0 : ℝ) (2 * R),
          s ^ (d - 1 : ℕ) • B.indicator (fun y => ‖u x - u y‖) (x + s • (ω : E)) ∂volume
        ≤ ∫ s in Set.Ioo (0 : ℝ) (2 * R), Cω * s ^ (d - 1 : ℕ) ∂volume := by
      apply integral_mono_of_nonneg hnonneg_lhs hpow_int.integrable
      filter_upwards [ae_restrict_mem measurableSet_Ioo] with s hs
      have hs_pos : 0 < s := hs.1
      by_cases hs_ball : x + s • (ω : E) ∈ B
      · have hω_norm : ‖(ω : E)‖ = 1 := by
          simpa only [Metric.mem_sphere, dist_zero_right] using ω.2
        have hnorm : ‖x - (x + s • (ω : E))‖ = s := by
          calc
            ‖x - (x + s • (ω : E))‖ = ‖-(s • (ω : E))‖ := by congr 1; abel_nf
            _ = ‖s • (ω : E)‖ := norm_neg _
            _ = s := by rw [norm_smul, Real.norm_of_nonneg hs_pos.le, hω_norm, mul_one]
        have hdir : s⁻¹ • ((x + s • (ω : E)) - x) = (ω : E) := by
          calc
            s⁻¹ • ((x + s • (ω : E)) - x) = s⁻¹ • (s • (ω : E)) := by congr 2; abel_nf
            _ = (s⁻¹ * s) • (ω : E) := by rw [smul_smul]
            _ = (ω : E) := by rw [inv_mul_cancel₀ hs_pos.ne', one_smul]
        have hseg0 := norm_sub_le_integral_fderiv_along_segment
          (d := d) (u := u) hu x (x + s • (ω : E))
        have hseg :
            ‖u x - u (x + s • (ω : E))‖ ≤
            ∫ t in (0 : ℝ)..s, ‖fderiv ℝ u (x + t • (ω : E))‖ := by
          calc
            ‖u x - u (x + s • (ω : E))‖
                ≤ ∫ t in (0 : ℝ)..‖x - (x + s • (ω : E))‖,
                    ‖fderiv ℝ u (x + t • (‖x - (x + s • (ω : E))‖⁻¹ • ((x + s • (ω : E)) - x)))‖ := hseg0
            _ = ∫ t in (0 : ℝ)..s, ‖fderiv ℝ u (x + t • (ω : E))‖ := by
                  rw [hnorm]
                  congr with t
                  simp [smul_smul, hs_pos.ne']
        have hrad_eq :
            ∫ t in (0 : ℝ)..s, ‖fderiv ℝ u (x + t • (ω : E))‖ =
            ∫ t in (0 : ℝ)..s,
              B.indicator (fun y => ‖fderiv ℝ u y‖) (x + t • (ω : E)) ∂volume := by
          rw [intervalIntegral.integral_of_le hs_pos.le, intervalIntegral.integral_of_le hs_pos.le]
          apply setIntegral_congr_fun measurableSet_Ioc
          intro t ht
          have htIcc : t ∈ Set.Icc (0 : ℝ) s := ⟨le_of_lt ht.1, ht.2⟩
          have hmem := ray_mem_ball_of_mem_ball (d := d) hx hs_ball hs_pos htIcc
          simp [B, hmem]
        have hsubset_s : Set.Ioo (0 : ℝ) s ⊆ Set.Ioo (0 : ℝ) (2 * R) := by
          intro t ht
          exact ⟨ht.1, lt_of_lt_of_le ht.2 hs.2.le⟩
        have hmono :
            ∫ t in (0 : ℝ)..s,
              B.indicator (fun y => ‖fderiv ℝ u y‖) (x + t • (ω : E)) ∂volume
            ≤ Cω := by
          rw [intervalIntegral.integral_of_le hs_pos.le, integral_Ioc_eq_integral_Ioo]
          apply setIntegral_mono_set
            (integrableOn_radial_indicator_fderiv (d := d) hR hu x ω)
          · filter_upwards [ae_restrict_mem measurableSet_Ioo] with t ht
            by_cases ht_ball : x + t • (ω : E) ∈ B <;> simp [B, ht_ball]
          · exact hsubset_s.eventuallyLE
        have hk_nonneg : 0 ≤ s ^ (d - 1 : ℕ) := pow_nonneg hs_pos.le _
        have hseg' : ‖u x - u (x + s • (ω : E))‖ ≤ Cω := by
          exact hseg.trans (by rw [hrad_eq]; exact hmono)
        have hmul : ‖u x - u (x + s • (ω : E))‖ * s ^ (d - 1 : ℕ) ≤ Cω * s ^ (d - 1 : ℕ) :=
          mul_le_mul_of_nonneg_right hseg' hk_nonneg
        simpa [B, hs_ball, smul_eq_mul, mul_comm] using hmul
      · have : 0 ≤ Cω * s ^ (d - 1 : ℕ) := mul_nonneg hCω_nonneg (pow_nonneg hs_pos.le _)
        simpa [B, hs_ball, smul_eq_mul] using this
    have hpow_eval :
        ∫ s in Set.Ioo (0 : ℝ) (2 * R), s ^ (d - 1 : ℕ) ∂volume =
          (2 * R) ^ d / ↑d := by
      have h2R : (0 : ℝ) < 2 * R := by linarith
      calc
        ∫ s in Set.Ioo (0 : ℝ) (2 * R), s ^ (d - 1 : ℕ) ∂volume
            = ∫ s in Set.Ioc (0 : ℝ) (2 * R), s ^ (d - 1 : ℕ) ∂volume := by
                rw [integral_Ioc_eq_integral_Ioo]
        _ = ∫ s in (0 : ℝ)..(2 * R), s ^ (d - 1 : ℕ) := by
              rw [← intervalIntegral.integral_of_le h2R.le]
        _ = ((2 * R) ^ d - 0 ^ d) / ↑d := by
              rw [integral_pow (a := (0 : ℝ)) (b := 2 * R) (n := d - 1)]
              simp [Nat.sub_add_cancel (Nat.one_le_iff_ne_zero.mpr (NeZero.ne d)),
                Nat.cast_sub (Nat.one_le_iff_ne_zero.mpr (NeZero.ne d))]
        _ = (2 * R) ^ d / ↑d := by
              simp [NeZero.ne d]
    calc
      ∫ s in Set.Ioo (0 : ℝ) (2 * R),
          s ^ (d - 1 : ℕ) • B.indicator (fun y => ‖u x - u y‖) (x + s • (ω : E)) ∂volume
          ≤ ∫ s in Set.Ioo (0 : ℝ) (2 * R), Cω * s ^ (d - 1 : ℕ) ∂volume := hrad_le
      _ = Cω * ((2 * R) ^ d / ↑d) := by
            rw [integral_const_mul, hpow_eval]
      _ = ((2 * R) ^ d / ↑d) * Cω := by ring
  have houter_nonneg :
      0 ≤ᵐ[(volume : Measure E).toSphere]
        (fun ω : Metric.sphere (0 : E) 1 =>
          ∫ s in Set.Ioo (0 : ℝ) (2 * R),
            s ^ (d - 1 : ℕ) • B.indicator (fun y => ‖u x - u y‖) (x + s • (ω : E)) ∂volume) := by
    filter_upwards with ω
    rw [← integral_indicator measurableSet_Ioo]
    apply integral_nonneg
    intro s
    by_cases hs : s ∈ Set.Ioo (0 : ℝ) (2 * R)
    · by_cases hs_ball : x + s • (ω : E) ∈ B
      · simpa [hs, B, hs_ball, smul_eq_mul] using
          mul_nonneg (pow_nonneg hs.1.le _) (norm_nonneg (u x - u (x + s • (ω : E))))
      · simp [hs, B, hs_ball, smul_eq_mul]
    · simp [hs]
  have houter_le :
      ∫ ω in Set.univ,
        (∫ s in Set.Ioo (0 : ℝ) (2 * R),
          s ^ (d - 1 : ℕ) • B.indicator (fun y => ‖u x - u y‖) (x + s • (ω : E)) ∂volume)
        ∂(volume : Measure E).toSphere
      ≤
      ∫ ω in Set.univ,
        (((2 * R) ^ d / ↑d) *
          ∫ s in Set.Ioo (0 : ℝ) (2 * R),
            B.indicator (fun y => ‖fderiv ℝ u y‖) (x + s • (ω : E)) ∂volume)
        ∂(volume : Measure E).toSphere := by
    simpa [setIntegral_univ] using integral_mono_of_nonneg houter_nonneg
      (houterG_int.const_mul ((2 * R) ^ d / ↑d))
      (ae_of_all _ hpointwise)
  have hrad_g :
      ∫ ω in Set.univ,
        (∫ s in Set.Ioo (0 : ℝ) (2 * R),
          B.indicator (fun y => ‖fderiv ℝ u y‖) (x + s • (ω : E)) ∂volume)
        ∂(volume : Measure E).toSphere
      = ∫ y in B, ‖fderiv ℝ u y‖ * ‖y - x‖ ^ (1 - (d : ℝ)) ∂volume := by
    simpa [B] using star_shaped_polar_identity
      (d := d) (R := R) hR (x := x) hx (f := fun y => ‖fderiv ℝ u y‖) hfG
  have hkey_integral : ∫ y in B, ‖u x - u y‖ ∂volume ≤
      ((2 * R) ^ d / ↑d) *
        ∫ y in B, ‖fderiv ℝ u y‖ * ‖x - y‖ ^ (1 - (d : ℝ)) ∂volume := by
    calc
      ∫ y in B, ‖u x - u y‖ ∂volume
          = ∫ ω in Set.univ,
              (∫ s in Set.Ioo (0 : ℝ) (2 * R),
                s ^ (d - 1 : ℕ) • B.indicator (fun y => ‖u x - u y‖) (x + s • (ω : E)) ∂volume)
              ∂(volume : Measure E).toSphere := by
                symm
                exact hrad_u
      _ ≤ ∫ ω in Set.univ,
            (((2 * R) ^ d / ↑d) *
              ∫ s in Set.Ioo (0 : ℝ) (2 * R),
                B.indicator (fun y => ‖fderiv ℝ u y‖) (x + s • (ω : E)) ∂volume)
            ∂(volume : Measure E).toSphere := houter_le
      _ = ((2 * R) ^ d / ↑d) *
            ∫ ω in Set.univ,
              (∫ s in Set.Ioo (0 : ℝ) (2 * R),
                B.indicator (fun y => ‖fderiv ℝ u y‖) (x + s • (ω : E)) ∂volume)
              ∂(volume : Measure E).toSphere := by
                rw [integral_const_mul]
      _ = ((2 * R) ^ d / ↑d) *
            ∫ y in B, ‖fderiv ℝ u y‖ * ‖y - x‖ ^ (1 - (d : ℝ)) ∂volume := by
                rw [hrad_g]
      _ = ((2 * R) ^ d / ↑d) *
            ∫ y in B, ‖fderiv ℝ u y‖ * ‖x - y‖ ^ (1 - (d : ℝ)) ∂volume := by
                rw [hkernel_symm]
  have hfin : Module.finrank ℝ E = d := finrank_euclideanSpace_fin
  have hvolB_eq : volB = (volume (Metric.ball (0 : E) 1)).toReal * R ^ d := by
    simp only [volB, B]
    rw [MeasureTheory.Measure.addHaar_ball_of_pos _ _ hR]
    rw [ENNReal.toReal_mul, hfin, ENNReal.toReal_ofReal (pow_nonneg hR.le d)]
    ring
  have havg_unfold : ⨍ y in B, ‖u x - u y‖ ∂volume =
      volB⁻¹ * ∫ y in B, ‖u x - u y‖ ∂volume := by
    show _ = (volume B).toReal⁻¹ * _
    rw [average_eq, smul_eq_mul, measureReal_restrict_apply_univ, measureReal_def]
  rw [havg_unfold]
  calc volB⁻¹ * ∫ y in B, ‖u x - u y‖ ∂volume
      ≤ volB⁻¹ * (((2 * R) ^ d / ↑d) *
          ∫ y in B, ‖fderiv ℝ u y‖ * ‖x - y‖ ^ (1 - (d : ℝ)) ∂volume) := by
        apply mul_le_mul_of_nonneg_left _ (inv_nonneg.mpr hvol_pos.le)
        exact hkey_integral
    _ = (volB⁻¹ * ((2 * R) ^ d / ↑d)) *
          ∫ y in B, ‖fderiv ℝ u y‖ * ‖x - y‖ ^ (1 - (d : ℝ)) ∂volume := by ring
    _ = C_rep * ∫ y in B, ‖fderiv ℝ u y‖ * ‖x - y‖ ^ (1 - (d : ℝ)) ∂volume := by
        congr 1
        simp only [C_rep, hvolB_eq, mul_pow]
        field_simp

-- Local copy of lintegral_rpow_norm_eq_eLpNorm_pow (defined in BallExtension, not imported here)
\end{lstlisting}
\begin{proof}
For $\omega \in S^{d-1}$ and $r$ such that $x + r\omega \in \Bz{1}$ one writes
$u(x+r\omega) - u(x) = \int_{0}^{r} \partial_{\omega} u(x + s\omega)\, ds$ and integrates
in $r$ along radial segments inside $\Bz{1}$. Averaging this identity over $\omega$
weighted by the chord length, and using
\leanname{star\_shaped\_polar\_identity} to convert the spherical integral into a
weighted volume integral with weight $\lvert x - y\rvert^{1-d}$, gives the stated
pointwise estimate. The constant $2^{d}/(d\,\mathrm{vol}(\Bz{1}))$ comes from estimating
the maximal chord length of $\Bz{1}$ from any point $x \in \Bz{1}$ by $2$.
\end{proof}

\begin{lemma}[Weighted power mean for set integrals,
\leanname{weighted\_power\_mean\_setIntegral}]
Let $p \ge 1$, let $w \ge 0$ be a weight on a measurable set $S$ with $W := \int_S w\, dy$
finite, and let $f \ge 0$. Then
\[
\Bigl(\int_S f\,w\, dy\Bigr)^{p}
\le W^{p-1}\, \int_S f^{p}\,w\, dy.
\]
\end{lemma}

\begin{lstlisting}[style=Matlab-editor,basicstyle=\mlttfamily,escapechar=",mlshowsectionrules=true,mathescape=true,morecomment={[s]{\%\{}{\%\}}},language=lean,numbers=none]
theorem weighted_power_mean_setIntegral
    {α : Type*} [MeasurableSpace α] {μ : Measure α}
    {s : Set α} (hs : MeasurableSet s)
    {p : ℝ} (hp : 1 < p)
    {f w : α → ℝ} (hf : ∀ x, 0 ≤ f x) (hw : ∀ x, 0 ≤ w x)
    (hf_meas : AEMeasurable f (μ.restrict s))
    (hwi : IntegrableOn w s μ)
    (hfpwi : IntegrableOn (fun x => (f x) ^ p * w x) s μ) :
    (∫ x in s, f x * w x ∂μ) ^ p ≤
      (∫ x in s, w x ∂μ) ^ (p - 1) * ∫ x in s, (f x) ^ p * w x ∂μ := by
  set μs : Measure α := μ.restrict s
  set ρ : α → ℝ≥0∞ := fun x => ENNReal.ofReal (w x)
  set ν : Measure α := μs.withDensity ρ
  let q : ℝ := p / (p - 1)
  have hp_pos : 0 < p := lt_trans zero_lt_one hp
  have hp_nonneg : 0 ≤ p := le_of_lt hp_pos
  have hpq : p.HolderConjugate q := by
    dsimp [q]
    exact Real.HolderConjugate.conjExponent hp
  have hq_pos : 0 < q := hpq.symm.pos
  have hρ_aemeas : AEMeasurable ρ μs := by
    simpa [ρ, μs] using hwi.aestronglyMeasurable.aemeasurable.ennreal_ofReal
  have hρ_lt_top : ∀ᵐ x ∂μs, ρ x < ∞ := by
    filter_upwards with x
    simp [ρ]
  have hρ_lint_ne_top : ∫⁻ x, ρ x ∂μs ≠ ∞ := by
    rw [← ofReal_integral_eq_lintegral_ofReal hwi (ae_of_all _ fun x => hw x)]
    simp
  haveI : IsFiniteMeasure ν := isFiniteMeasure_withDensity hρ_lint_ne_top
  have hf_meas_ν : AEMeasurable f ν := by
    exact hf_meas.mono_ac (withDensity_absolutelyContinuous _ _)
  have hpow_int_base : Integrable (fun x => (ρ x).toReal • (‖f x‖ ^ p)) μs := by
    refine hfpwi.congr ?_
    filter_upwards with x
    rw [smul_eq_mul, ENNReal.toReal_ofReal (hw x), Real.norm_of_nonneg (hf x)]
    ring
  have hpow_int : Integrable (fun x => ‖f x‖ ^ p) ν := by
    rw [show ν = μs.withDensity ρ by rfl]
    exact
      (MeasureTheory.integrable_withDensity_iff_integrable_smul₀'
        (μ := μs) hρ_aemeas hρ_lt_top).2 hpow_int_base
  have hf_mem : MemLp f (ENNReal.ofReal p) ν := by
    exact
      (MeasureTheory.integrable_norm_rpow_iff
        (μ := ν) hf_meas_ν.aestronglyMeasurable
        (by simp [hp_pos]) ENNReal.ofReal_ne_top).1 <| by
          simpa [ENNReal.toReal_ofReal hp_nonneg] using hpow_int
  have h_one_mem : MemLp (fun _ : α => (1 : ℝ)) (ENNReal.ofReal q) ν := by
    simpa [q] using (MeasureTheory.memLp_const (μ := ν) (p := ENNReal.ofReal q) (1 : ℝ))
  have hHolder :
      ∫ x, f x * (1 : ℝ) ∂ν ≤
        (∫ x, (f x) ^ p ∂ν) ^ (1 / p : ℝ) *
          (∫ x, (1 : ℝ) ^ q ∂ν) ^ (1 / q : ℝ) := by
    exact MeasureTheory.integral_mul_le_Lp_mul_Lq_of_nonneg
      (μ := ν) hpq
      (ae_of_all _ fun x => hf x)
      (ae_of_all _ fun _ => by positivity)
      hf_mem h_one_mem
  have hleft_eq : ∫ x, f x ∂ν = ∫ x in s, f x * w x ∂μ := by
    rw [show ν = μs.withDensity ρ by rfl]
    rw [integral_withDensity_eq_integral_toReal_smul₀ (μ := μs) hρ_aemeas hρ_lt_top]
    simp [μs, ρ, smul_eq_mul, ENNReal.toReal_ofReal, hw, mul_comm]
  have hpow_eq : ∫ x, (f x) ^ p ∂ν = ∫ x in s, (f x) ^ p * w x ∂μ := by
    rw [show ν = μs.withDensity ρ by rfl]
    rw [integral_withDensity_eq_integral_toReal_smul₀ (μ := μs) hρ_aemeas hρ_lt_top]
    simp [μs, ρ, smul_eq_mul, ENNReal.toReal_ofReal, hw, mul_comm]
  have hone_eq : ∫ x, (1 : ℝ) ^ q ∂ν = ∫ x in s, w x ∂μ := by
    rw [show ν = μs.withDensity ρ by rfl]
    rw [integral_withDensity_eq_integral_toReal_smul₀ (μ := μs) hρ_aemeas hρ_lt_top]
    simp [μs, ρ, q, smul_eq_mul, ENNReal.toReal_ofReal, hw]
  set A := ∫ x in s, f x * w x ∂μ
  set I := ∫ x in s, (f x) ^ p * w x ∂μ
  set W := ∫ x in s, w x ∂μ
  have hA_nonneg : 0 ≤ A := by
    exact setIntegral_nonneg hs (fun x _ => mul_nonneg (hf x) (hw x))
  have hI_nonneg : 0 ≤ I := by
    exact setIntegral_nonneg hs (fun x _ =>
      mul_nonneg (Real.rpow_nonneg (hf x) _) (hw x))
  have hW_nonneg : 0 ≤ W := by
    exact setIntegral_nonneg hs (fun x _ => hw x)
  -- If ∫ w = 0, both sides are 0
  by_cases hW_zero : ∫ x in s, w x ∂μ = 0
  · -- Here `w = 0` a.e., so the weighted integral also vanishes.
    have hfw_nonneg : 0 ≤ ∫ x in s, f x * w x ∂μ :=
      setIntegral_nonneg hs (fun x _ => mul_nonneg (hf x) (hw x))
    have hfw_zero : ∫ x in s, f x * w x ∂μ ≤ 0 := by
      by_cases hfwi : IntegrableOn (fun x => f x * w x) s μ
      · -- w ≥ 0 and ∫ w = 0 → w =ᵃᵉ 0 → f·w =ᵃᵉ 0 → ∫ f·w = 0
        have hwi := hwi
        have hw_ae : w =ᵐ[μ.restrict s] 0 := by
          rwa [setIntegral_eq_zero_iff_of_nonneg_ae
            (ae_restrict_of_ae (ae_of_all _ (fun x => hw x))) hwi] at hW_zero
        have : (fun x => f x * w x) =ᵐ[μ.restrict s] 0 :=
          hw_ae.mono (fun x hx => by simp [hx])
        rw [integral_congr_ae this]; simp
      · simp [integral_undef hfwi]
    have h0 := le_antisymm hfw_zero hfw_nonneg
    -- h0 : ∫ f·w = 0, hW_zero : ∫ w = 0. Goal: 0^p ≤ 0^{p-1} * ∫ f^p·w
    simpa [A, I, W, h0, hW_zero] using
      (show A ^ p ≤ W ^ (p - 1) * I by
        simp [A, I, W, h0, hW_zero,
          mul_nonneg (Real.rpow_nonneg (le_refl _) _)
            (setIntegral_nonneg hs (fun x _ => mul_nonneg (Real.rpow_nonneg (hf x) _) (hw x))),
          Real.zero_rpow (by linarith : p ≠ 0)])
  · -- ∫ w > 0
    have hW_pos : 0 < ∫ x in s, w x ∂μ :=
      lt_of_le_of_ne (setIntegral_nonneg hs (fun x _ => hw x)) (Ne.symm hW_zero)
    have hνreal_eq : ν.real Set.univ = W := by
      calc
        ν.real Set.univ = ∫ x, (1 : ℝ) ∂ν := by
          rw [integral_const]
          simp [Measure.real]
        _ = ∫ x, (1 : ℝ) ^ q ∂ν := by simp [q]
        _ = W := hone_eq
    have hHolder' : A ≤ W ^ (1 / q : ℝ) * I ^ (1 / p : ℝ) := by
      simpa [A, I, W, hleft_eq, hpow_eq, hνreal_eq, one_div, mul_comm, mul_left_comm, mul_assoc]
        using hHolder
    have hpow :=
      Real.rpow_le_rpow hA_nonneg hHolder' (le_of_lt hp_pos)
    have hWroot_nonneg : 0 ≤ W ^ (1 / q : ℝ) := Real.rpow_nonneg hW_nonneg _
    have hIroot_nonneg : 0 ≤ I ^ (1 / p : ℝ) := Real.rpow_nonneg hI_nonneg _
    have hq_ne_zero : q ≠ 0 := by linarith
    have hp_ne_zero : p ≠ 0 := by linarith
    have hrhs :
        (W ^ (1 / q : ℝ) * I ^ (1 / p : ℝ)) ^ p = W ^ (p - 1) * I := by
      rw [Real.mul_rpow hWroot_nonneg hIroot_nonneg]
      rw [← Real.rpow_mul hW_nonneg, ← Real.rpow_mul hI_nonneg]
      have hWq : (1 / q : ℝ) * p = p - 1 := by
        dsimp [q]
        field_simp [hp_ne_zero, show p - 1 ≠ 0 by linarith]
      have hIp : (1 / p : ℝ) * p = 1 := by
        field_simp [hp_ne_zero]
      rw [hWq, hIp, Real.rpow_one]
    simpa [A, I, W] using hpow.trans_eq hrhs

\end{lstlisting}
\begin{proof}
Normalise $w$ to a probability measure $d\nu = w/W\, dy$ and apply Jensen's inequality
to the convex function $t \mapsto t^{p}$, then rescale.
\end{proof}

\begin{theorem}[Poincar\'e on the unit ball for smooth functions,
\leanname{poincare\_smooth\_unitBall}]
Let $p \ge 1$ and $u \in C^{\infty}(E,\R)$. Then
\[
\Bigl\lVert u - \fint_{\Bz{1}} u \Bigr\rVert_{\Lp{p}(\Bz{1})}
\le C_{\mathrm{poinc}}(d)\, \bigl\lVert \lvert\nabla u\rvert \bigr\rVert_{\Lp{p}(\Bz{1})}.
\]
\end{theorem}

\begin{lstlisting}[style=Matlab-editor,basicstyle=\mlttfamily,escapechar=",mlshowsectionrules=true,mathescape=true,morecomment={[s]{\%\{}{\%\}}},language=lean,numbers=none]
theorem poincare_smooth_unitBall
    {p : ℝ} (hp : 1 ≤ p)
    {u : E → ℝ} (hu : ContDiff ℝ (⊤ : ℕ∞) u) :
    eLpNorm (fun x => u x - ⨍ y in Metric.ball (0 : E) 1, u y ∂volume)
      (ENNReal.ofReal p) (volume.restrict (Metric.ball (0 : E) 1)) ≤
    ENNReal.ofReal (C_poinc_val d) *
      eLpNorm (fun x => ‖fderiv ℝ u x‖) (ENNReal.ofReal p)
        (volume.restrict (Metric.ball (0 : E) 1)) := by
  set B := Metric.ball (0 : E) 1
  set ubar := ⨍ y in B, u y ∂volume
  set f : E → ℝ := fun x => u x - ubar
  set g : E → ℝ := fun x => ‖fderiv ℝ u x‖
  have hpw : ∀ x ∈ B, ‖f x‖ ≤
      (2 : ℝ) ^ d / ((d : ℝ) * (volume (Metric.ball (0 : E) 1)).toReal) *
        ∫ y in B, g y * ‖x - y‖ ^ (1 - (d : ℝ)) ∂volume :=
    fun x hx => representation_formula_smooth one_pos hu x hx
  set μ := volume.restrict B
  set C_rep := (2 : ℝ) ^ d / ((d : ℝ) * (volume (Metric.ball (0 : E) 1)).toReal)
  set K : E → ℝ := fun z => ‖z‖ ^ (1 - (d : ℝ))
  set h : E → ℝ := fun x => ∫ y in B, g y * K (x - y) ∂volume
  have hpw_ae : ∀ᵐ x ∂μ, ‖f x‖ ≤ C_rep * h x := by
    filter_upwards [ae_restrict_mem measurableSet_ball] with x hx_mem
    exact hpw x hx_mem
  have h_step2 : eLpNorm f (ENNReal.ofReal p) μ ≤
      eLpNorm (fun x => C_rep * h x) (ENNReal.ofReal p) μ :=
    eLpNorm_mono_ae_real hpw_ae
  set M := (d : ℝ) * (volume (Metric.ball (0 : E) 1)).toReal * (2 * 1)
  -- Young-type bound: ‖h‖_p ≤ M · ‖g‖_p
  -- Proof outline:
  --   (a) pointwise Jensen: |h(x)|^p ≤ M^{p-1} · ∫_B |g(y)|^p K(x-y) dy
  --   (b) integrate in x, Fubini swap, kernel L¹ bound → ∫_B |h|^p ≤ M^p ∫_B |g|^p
  --   (c) take p-th root: ‖h‖_p ≤ M · ‖g‖_p
  have hYoung : eLpNorm h (ENNReal.ofReal p) μ ≤
      ENNReal.ofReal M * eLpNorm g (ENNReal.ofReal p) μ := by
    have hp_pos : 0 < p := lt_of_lt_of_le one_pos hp
    have hp_nonneg : 0 ≤ p := le_of_lt hp_pos
    have hd_pos : (0 : ℝ) < d := Nat.cast_pos.mpr (NeZero.pos d)
    have hvol_pos : (0 : ℝ) < (volume (Metric.ball (0 : E) 1)).toReal :=
      ENNReal.toReal_pos (measure_ball_pos volume 0 one_pos).ne' measure_ball_lt_top.ne
    have hM_pos : 0 < M := by positivity
    have hM_nonneg : 0 ≤ M := le_of_lt hM_pos
    have h1 : ContDiff ℝ 1 u := hu.of_le (by norm_num)
    have hcont_fderiv : Continuous (fderiv ℝ u) := h1.continuous_fderiv (by norm_num)
    have hg_cont : Continuous g := by
      simpa [g] using (continuous_norm.comp hcont_fderiv)
    obtain ⟨z₀, hz₀_mem, hz₀_max⟩ :=
      (isCompact_closedBall (0 : E) 1).exists_isMaxOn
        (Metric.nonempty_closedBall.mpr (le_of_lt one_pos)) hg_cont.continuousOn
    set Cg := g z₀
    have hg_le : ∀ y ∈ Metric.closedBall (0 : E) 1, g y ≤ Cg := by
      intro y hy
      exact hz₀_max hy
    -- K ≥ 0
    have hK_nonneg : ∀ z : E, 0 ≤ K z := fun z =>
      Real.rpow_nonneg (norm_nonneg z) _
    have hK_meas : Measurable K := by
      fun_prop
    -- Kernel L¹ bound: for x ∈ B, ∫_B K(x-y) dy ≤ M
    have hK_L1 : ∀ x ∈ B, ∫ y in B, K (x - y) ∂volume ≤ M := by
      intro x hx
      have := riesz_kernel_L1_bound one_pos x hx
      simp only [mul_one] at this
      convert this using 1
      simp [M]
    have hK_int : ∀ x ∈ B, IntegrableOn (fun y => K (x - y)) B volume := by
      intro x hx
      have hball_int : IntegrableOn (fun y => K (x - y)) (Metric.ball x (2 : ℝ)) volume := by
        have hmp := measurePreserving_add_right (volume : Measure E) x
        have hemb := (MeasurableEquiv.addRight x : E ≃ᵐ E).measurableEmbedding
        have hpre : (fun z : E => z + x) ⁻¹' Metric.ball x (2 : ℝ) = Metric.ball (0 : E) (2 : ℝ) := by
          ext z
          simp [Metric.mem_ball]
        rw [← hmp.integrableOn_comp_preimage hemb, hpre]
        exact (riesz_kernel_integrable (by norm_num : (0 : ℝ) < 2)).congr
          (ae_of_all _ fun z => by
            show ‖z‖ ^ (1 - (d : ℝ)) = K (x - (z + x))
            simp [K, norm_neg])
      refine hball_int.mono_set ?_
      intro y hy
      have hx_ball : ‖x‖ < 1 := by
        simpa [B, Metric.mem_ball, dist_zero_right] using hx
      have hy_ball : ‖y‖ < 1 := by
        simpa [B, Metric.mem_ball, dist_zero_right] using hy
      rw [Metric.mem_ball, dist_eq_norm]
      calc
        ‖y - x‖ ≤ ‖y‖ + ‖x‖ := norm_sub_le _ _
        _ < 1 + 1 := add_lt_add hy_ball hx_ball
        _ = (2 : ℝ) := by ring
    -- g is nonneg
    have hg_nonneg : ∀ y, 0 ≤ g y := fun y => norm_nonneg _
    -- h is nonneg (integral of nonneg * nonneg)
    have hh_nonneg : ∀ x, 0 ≤ h x := by
      intro x; apply setIntegral_nonneg measurableSet_ball
      intro y _; exact mul_nonneg (hg_nonneg y) (hK_nonneg (x - y))
    -- |g y| = g y since g ≥ 0
    have habs_g : ∀ y, |g y| = g y := fun y => abs_of_nonneg (hg_nonneg y)
    -- |h x| = h x since h ≥ 0
    have habs_h : ∀ x, |h x| = h x := fun x => abs_of_nonneg (hh_nonneg x)
    -- (a) Pointwise Jensen: |h(x)|^p ≤ M^{p-1} · ∫_B |g(y)|^p · K(x-y) dy
    -- This is the weighted power-mean inequality for integrals:
    --   (∫ f·w)^p ≤ (∫ w)^{p-1} · ∫ f^p·w
    -- Proof: normalize w to probability measure ν = w/(∫w), apply Jensen
    --   (∫ f dν)^p ≤ ∫ f^p dν  (convexity of t^p)
    --   then scale back by (∫ w)^p on left, (∫ w) on right.
    have hJensen : ∀ x ∈ B,
        |h x| ^ p ≤ M ^ (p - 1) * ∫ y in B, |g y| ^ p * K (x - y) ∂volume := by
      intro x hx
      rw [habs_h]
      -- Let W = ∫_B K(x-y) dy, the kernel integral
      set W := ∫ y in B, K (x - y) ∂volume with hW_def
      have hW_nonneg : 0 ≤ W :=
        setIntegral_nonneg measurableSet_ball (fun y _ => hK_nonneg (x - y))
      have hW_le_M : W ≤ M := hK_L1 x hx
      -- Rewrite |g y|^p = (g y)^p since g ≥ 0
      have : ∫ y in B, |g y| ^ p * K (x - y) ∂volume =
          ∫ y in B, (g y) ^ p * K (x - y) ∂volume := by
        congr 1; ext y; rw [habs_g]
      rw [this]
      -- Core weighted power-mean inequality:
      -- (∫ g·K)^p ≤ W^{p-1} · ∫ g^p·K, then use W ≤ M
      -- Weighted power-mean inequality for integrals:
      -- (∫ f·w)^p ≤ (∫ w)^{p-1} · ∫ f^p·w, for nonneg f, w, and p ≥ 1
      -- Proof sketch: normalize w to a probability measure, apply Jensen
      -- (convexOn_rpow hp) to get (∫ f dν)^p ≤ ∫ f^p dν, then rescale.
      -- We bound (∫ K)^{p-1} ≤ M^{p-1} since ∫ K ≤ M.
      --         by Hölder with measure K·vol, exponents p and q = p/(p-1).
      calc (h x) ^ p
          ≤ W ^ (p - 1) * ∫ y in B, (g y) ^ p * K (x - y) ∂volume := by
            -- Weighted power-mean: (∫ g·K)^p ≤ (∫ K)^{p-1} · ∫ g^p·K
            -- via Hölder with f₁ = g · K^{1/p}, f₂ = K^{1/q}, exponents p, q
            -- For p = 1: trivial (LHS = RHS with W^0 = 1)
            by_cases hp1 : p = 1
            · -- p = 1: (h x)^1 = h x = ∫ g·K ≤ 1 · ∫ g·K = W^0 · ∫ g·K
              simp [hp1, h]
            · -- p > 1: use Hölder's inequality on lintegrals
              -- Set up: q = p/(p-1) is the conjugate exponent
              have hp_gt1 : 1 < p := lt_of_le_of_ne hp (Ne.symm hp1)
              have hpq : p.HolderConjugate (p / (p - 1)) := Real.HolderConjugate.conjExponent hp_gt1
              -- Convert Bochner integrals to lintegrals (nonneg functions)
              -- h x = ∫_B g·K = toReal(∫⁻_B ofReal(g·K))
              set μB := volume.restrict B
              -- Hölder: ∫⁻ (f₁ * f₂) ≤ (∫⁻ f₁^p)^{1/p} · (∫⁻ f₂^q)^{1/q}
              -- with f₁ y = ofReal(g y · K(x-y)^{1/p}), f₂ y = ofReal(K(x-y)^{1/q})
              -- f₁ * f₂ = ofReal(g · K^{1/p} · K^{1/q}) = ofReal(g · K)
              -- f₁^p = ofReal(g^p · K), f₂^q = ofReal(K)
              exact weighted_power_mean_setIntegral measurableSet_ball hp_gt1
                hg_nonneg (fun y => hK_nonneg (x - y))
                (hg_cont.aemeasurable.mono_measure Measure.restrict_le_self)
                (hK_int x hx)
                (by
                  have hgpow_cont : Continuous fun y => (g y) ^ p :=
                    hg_cont.rpow_const (fun _ => Or.inr hp_pos.le)
                  simpa [mul_comm] using
                    (hK_int x hx).mul_continuousOn_of_subset hgpow_cont.continuousOn
                      measurableSet_ball (isCompact_closedBall (0 : E) 1)
                      Metric.ball_subset_closedBall)
        _ ≤ M ^ (p - 1) * ∫ y in B, (g y) ^ p * K (x - y) ∂volume := by
            apply mul_le_mul_of_nonneg_right _ (setIntegral_nonneg measurableSet_ball
              (fun y _ => mul_nonneg (Real.rpow_nonneg (hg_nonneg y) _) (hK_nonneg _)))
            exact Real.rpow_le_rpow hW_nonneg hW_le_M (by linarith)
    -- (b) Integrate Jensen bound over B, Fubini, kernel bound → ∫_B |h|^p ≤ M^p · ∫_B |g|^p
    -- Kernel L¹ bound in swapped variables: ∀ y ∈ B, ∫_B K(x-y) dx ≤ M
    have hK_L1_swap : ∀ y ∈ B, ∫ x in B, K (x - y) ∂volume ≤ M := by
      intro y hy
      have hsym : ∫ x in B, K (x - y) ∂volume = ∫ x in B, K (y - x) ∂volume := by
        congr 1; ext x; simp only [K]; congr 1; rw [norm_sub_rev]
      rw [hsym]; exact hK_L1 y hy
    have hgpow_cont : Continuous fun y => |g y| ^ p :=
      (hg_cont.abs.rpow_const (fun _ => Or.inr hp_pos.le))
    have hF_aesm :
        AEStronglyMeasurable
          (fun xy : E × E => |g xy.2| ^ p * K (xy.1 - xy.2))
          ((volume.restrict B).prod (volume.restrict B)) := by
      exact (((hgpow_cont.measurable.comp_aemeasurable measurable_snd.aemeasurable).mul
        (hK_meas.comp_aemeasurable
          (measurable_fst.sub measurable_snd).aemeasurable))).aestronglyMeasurable
    have hprod_int :
        Integrable
          (fun xy : E × E => |g xy.2| ^ p * K (xy.1 - xy.2))
          ((volume.restrict B).prod (volume.restrict B)) := by
      rw [integrable_prod_iff hF_aesm]
      constructor
      · filter_upwards [ae_restrict_mem measurableSet_ball] with x hx
        simpa [mul_comm] using
          (hK_int x hx).mul_continuousOn_of_subset hgpow_cont.continuousOn
            measurableSet_ball (isCompact_closedBall (0 : E) 1)
            Metric.ball_subset_closedBall
      · have houter_bdd :
            ∀ᵐ x ∂(volume.restrict B).restrict Set.univ,
              ‖∫ y in B, ‖|g y| ^ p * K (x - y)‖ ∂volume‖ ≤ Cg ^ p * M := by
          filter_upwards [show ∀ᵐ x ∂(volume.restrict B).restrict Set.univ, x ∈ B by
            simpa [Measure.restrict_univ] using
              (ae_restrict_mem (μ := volume) measurableSet_ball)] with x hx
          have hslice_int : IntegrableOn (fun y => |g y| ^ p * K (x - y)) B volume := by
            simpa [mul_comm] using
              (hK_int x hx).mul_continuousOn_of_subset hgpow_cont.continuousOn
                measurableSet_ball (isCompact_closedBall (0 : E) 1)
                Metric.ball_subset_closedBall
          have hslice_const_int : IntegrableOn (fun y => Cg ^ p * K (x - y)) B volume := by
            exact (hK_int x hx).const_mul (Cg ^ p)
          rw [Real.norm_of_nonneg
            (setIntegral_nonneg measurableSet_ball (fun y _ => norm_nonneg _))]
          calc
            ∫ y in B, ‖|g y| ^ p * K (x - y)‖ ∂volume
                = ∫ y in B, |g y| ^ p * K (x - y) ∂volume := by
                    congr 1
                    ext y
                    rw [Real.norm_of_nonneg
                      (mul_nonneg (Real.rpow_nonneg (abs_nonneg _) _) (hK_nonneg _))]
            _ ≤ ∫ y in B, Cg ^ p * K (x - y) ∂volume := by
                exact setIntegral_mono_on hslice_int hslice_const_int measurableSet_ball
                  (fun y hy =>
                    mul_le_mul_of_nonneg_right
                      (Real.rpow_le_rpow (abs_nonneg _)
                        (by
                          rw [habs_g y]
                          simpa [Cg] using hg_le y (ball_subset_closedBall hy))
                        hp_nonneg)
                      (hK_nonneg _))
            _ = Cg ^ p * ∫ y in B, K (x - y) ∂volume := by
                rw [integral_const_mul]
            _ ≤ Cg ^ p * M := by
                apply mul_le_mul_of_nonneg_left (hK_L1 x hx)
                exact Real.rpow_nonneg (le_trans (hg_nonneg _) (hg_le z₀ hz₀_mem)) _
        simpa [Measure.restrict_univ] using
          (Measure.integrableOn_of_bounded
            (μ := volume.restrict B) (s := Set.univ)
            (by simpa [B] using measure_ball_lt_top.ne)
            (hF_aesm.norm.integral_prod_right') houter_bdd).integrable
    have hFh_aesm :
        AEStronglyMeasurable
          (fun xy : E × E => g xy.2 * K (xy.1 - xy.2))
          ((volume.restrict B).prod (volume.restrict B)) := by
      exact ((hg_cont.measurable.comp_aemeasurable measurable_snd.aemeasurable).mul
        (hK_meas.comp_aemeasurable
          (measurable_fst.sub measurable_snd).aemeasurable)).aestronglyMeasurable
    have hh_aesm : AEStronglyMeasurable h (volume.restrict B) := by
      change AEStronglyMeasurable (fun x => ∫ y, g y * K (x - y) ∂(volume.restrict B))
        (volume.restrict B)
      exact hFh_aesm.integral_prod_right'
    have hh_bound : ∀ x ∈ B, h x ≤ Cg * M := by
      intro x hx
      have hslice_int : IntegrableOn (fun y => g y * K (x - y)) B volume := by
        simpa [mul_comm] using
          (hK_int x hx).mul_continuousOn_of_subset hg_cont.continuousOn
            measurableSet_ball (isCompact_closedBall (0 : E) 1)
            Metric.ball_subset_closedBall
      have hslice_const_int : IntegrableOn (fun y => Cg * K (x - y)) B volume := by
        exact (hK_int x hx).const_mul Cg
      calc
        h x = ∫ y in B, g y * K (x - y) ∂volume := rfl
        _ ≤ ∫ y in B, Cg * K (x - y) ∂volume := by
            exact setIntegral_mono_on hslice_int hslice_const_int measurableSet_ball
              (fun y hy =>
                mul_le_mul_of_nonneg_right
                  (hg_le y (ball_subset_closedBall hy))
                  (hK_nonneg _))
        _ = Cg * ∫ y in B, K (x - y) ∂volume := by
            rw [integral_const_mul]
        _ ≤ Cg * M := by
            apply mul_le_mul_of_nonneg_left (hK_L1 x hx)
            exact le_trans (hg_nonneg _) (hg_le z₀ hz₀_mem)
    have hhpow_int : IntegrableOn (fun x => |h x| ^ p) B volume := by
      change Integrable (fun x => |h x| ^ p) (volume.restrict B)
      have hhpow_aesm : AEStronglyMeasurable (fun x => |h x| ^ p) (volume.restrict B) := by
        exact ((Real.continuous_rpow_const hp_pos.le).measurable.comp_aemeasurable
          hh_aesm.aemeasurable.abs).aestronglyMeasurable
      have hhpow_bdd :
          ∀ᵐ x ∂(volume.restrict B).restrict Set.univ, ‖|h x| ^ p‖ ≤ (Cg * M) ^ p := by
        filter_upwards [show ∀ᵐ x ∂(volume.restrict B).restrict Set.univ, x ∈ B by
          simpa [Measure.restrict_univ] using
            (ae_restrict_mem (μ := volume) measurableSet_ball)] with x hx
        rw [Real.norm_eq_abs, abs_of_nonneg (Real.rpow_nonneg (abs_nonneg _) _),
          abs_of_nonneg (hh_nonneg x)]
        exact Real.rpow_le_rpow (hh_nonneg x) (hh_bound x hx) hp_nonneg
      simpa [Measure.restrict_univ] using
        (Measure.integrableOn_of_bounded
          (μ := volume.restrict B) (s := Set.univ)
          (by simpa [B] using measure_ball_lt_top.ne)
          hhpow_aesm hhpow_bdd).integrable
    have hgpow_int : IntegrableOn (fun y => |g y| ^ p) B volume := by
      exact
        ((hgpow_cont.continuousOn.integrableOn_compact
          (isCompact_closedBall (0 : E) 1)).mono_set Metric.ball_subset_closedBall)
    have hIntBound : ∫ x in B, |h x| ^ p ∂volume ≤
        M ^ p * ∫ y in B, |g y| ^ p ∂volume := by
      have hstep1 : ∫ x in B, |h x| ^ p ∂volume ≤
          ∫ x in B, (M ^ (p - 1) * ∫ y in B, |g y| ^ p * K (x - y) ∂volume) ∂volume := by
        exact setIntegral_mono_on
          (hf := hhpow_int)
          (hg := by
            have houter_int :
                Integrable
                  (fun x => ∫ y, |g y| ^ p * K (x - y) ∂(volume.restrict B))
                  (volume.restrict B) := by
              refine hprod_int.integral_norm_prod_left.congr ?_
              filter_upwards with x
              congr 1
              ext y
              rw [Real.norm_of_nonneg
                (mul_nonneg (Real.rpow_nonneg (abs_nonneg _) _) (hK_nonneg _))]
            exact houter_int.const_mul (M ^ (p - 1))
          )
          measurableSet_ball hJensen
      have hFubini : ∫ x in B, (∫ y in B, |g y| ^ p * K (x - y) ∂volume) ∂volume =
          ∫ y in B, (∫ x in B, |g y| ^ p * K (x - y) ∂volume) ∂volume := by
        let μB := volume.restrict B
        change ∫ x, (∫ y, |g y| ^ p * K (x - y) ∂μB) ∂μB =
            ∫ y, (∫ x, |g y| ^ p * K (x - y) ∂μB) ∂μB
        exact integral_integral_swap hprod_int
      have hpull : ∀ y, ∫ x in B, |g y| ^ p * K (x - y) ∂volume =
          |g y| ^ p * ∫ x in B, K (x - y) ∂volume := by
        intro y; exact integral_const_mul _ _
      have hstep4 : ∫ y in B, |g y| ^ p * (∫ x in B, K (x - y) ∂volume) ∂volume ≤
          ∫ y in B, |g y| ^ p * M ∂volume := by
        apply setIntegral_mono_on
        · have hinner_int :
              Integrable
                (fun y => ∫ x, |g y| ^ p * K (x - y) ∂(volume.restrict B))
                (volume.restrict B) := by
            refine hprod_int.integral_norm_prod_right.congr ?_
            filter_upwards with y
            congr 1
            ext x
            rw [Real.norm_of_nonneg
              (mul_nonneg (Real.rpow_nonneg (abs_nonneg _) _) (hK_nonneg _))]
          refine hinner_int.congr ?_
          filter_upwards with y
          exact hpull y
        · simpa [mul_comm] using hgpow_int.mul_const M
        · exact measurableSet_ball
        · intro y hy
          apply mul_le_mul_of_nonneg_left (hK_L1_swap y hy)
          exact Real.rpow_nonneg (abs_nonneg _) _
      -- M^{p-1} * M = M^p
      have hMpow : M ^ (p - 1) * M = M ^ p := by
        rw [mul_comm]
        -- Goal: M * M ^ (p - 1) = M ^ p
        -- Convert the leading M to M ^ (1:ℝ) without touching M inside rpow
        change M * M ^ (p - 1) = M ^ p
        conv_lhs => lhs; rw [show (M : ℝ) = M ^ (1 : ℝ) from (Real.rpow_one M).symm]
        rw [← Real.rpow_add hM_pos]
        congr 1; ring
      calc ∫ x in B, |h x| ^ p ∂volume
          ≤ ∫ x in B, (M ^ (p - 1) * ∫ y in B, |g y| ^ p * K (x - y) ∂volume) ∂volume :=
            hstep1
        _ = M ^ (p - 1) * ∫ x in B, (∫ y in B, |g y| ^ p * K (x - y) ∂volume) ∂volume := by
            rw [integral_const_mul]
        _ = M ^ (p - 1) * ∫ y in B, (∫ x in B, |g y| ^ p * K (x - y) ∂volume) ∂volume := by
            rw [hFubini]
        _ = M ^ (p - 1) * ∫ y in B, |g y| ^ p * (∫ x in B, K (x - y) ∂volume) ∂volume := by
            congr 1; congr 1; ext y; exact hpull y
        _ ≤ M ^ (p - 1) * ∫ y in B, |g y| ^ p * M ∂volume := by
            apply mul_le_mul_of_nonneg_left hstep4
            exact Real.rpow_nonneg (le_of_lt hM_pos) _
        _ = M ^ (p - 1) * (M * ∫ y in B, |g y| ^ p ∂volume) := by
            congr 1; simp_rw [mul_comm _ M]; rw [integral_const_mul]
        _ = M ^ p * ∫ y in B, |g y| ^ p ∂volume := by
            rw [← mul_assoc, hMpow]
    -- (c) Convert the integral bound to an `L^p` bound.
    have hIntBoundμ : ∫ x, |h x| ^ p ∂μ ≤ M ^ p * ∫ y, |g y| ^ p ∂μ := by
      simpa [μ] using hIntBound
    have h_pow : eLpNorm h (ENNReal.ofReal p) μ ^ p ≤
        (ENNReal.ofReal M * eLpNorm g (ENNReal.ofReal p) μ) ^ p := by
      rw [← lintegral_rpow_norm_eq_eLpNorm_pow' hp_pos (f := h)]
      rw [ENNReal.mul_rpow_of_nonneg _ _ (le_of_lt hp_pos)]
      rw [← lintegral_rpow_norm_eq_eLpNorm_pow' hp_pos (f := g)]
      simp_rw [ENNReal.ofReal_rpow_of_nonneg (norm_nonneg _) (le_of_lt hp_pos),
        Real.norm_eq_abs]
      rw [← ofReal_integral_eq_lintegral_ofReal hhpow_int.integrable
          (ae_of_all _ (fun x => Real.rpow_nonneg (abs_nonneg _) _)),
        ← ofReal_integral_eq_lintegral_ofReal hgpow_int.integrable
          (ae_of_all _ (fun x => Real.rpow_nonneg (abs_nonneg _) _))]
      rw [ENNReal.ofReal_rpow_of_nonneg (le_of_lt hM_pos) (le_of_lt hp_pos),
        ← ENNReal.ofReal_mul (Real.rpow_nonneg (le_of_lt hM_pos) _)]
      exact ENNReal.ofReal_le_ofReal hIntBoundμ
    have hp_inv : 0 < (1 / p : ℝ) := div_pos one_pos hp_pos
    have hroot := ENNReal.rpow_le_rpow h_pow hp_inv.le
    rwa [← ENNReal.rpow_mul, ← ENNReal.rpow_mul,
      show p * (1 / p) = 1 from by field_simp [hp_pos.ne'],
      ENNReal.rpow_one, ENNReal.rpow_one] at hroot
  calc eLpNorm f (ENNReal.ofReal p) μ
      ≤ eLpNorm (fun x => C_rep * h x) (ENNReal.ofReal p) μ := h_step2
    _ ≤ ENNReal.ofReal C_rep * eLpNorm h (ENNReal.ofReal p) μ := by
        have : (fun x => C_rep * h x) = C_rep • h := by ext x; simp [Pi.smul_apply, smul_eq_mul]
        rw [this]
        calc eLpNorm (C_rep • h) (ENNReal.ofReal p) μ
            ≤ ‖C_rep‖ₑ * eLpNorm h (ENNReal.ofReal p) μ := eLpNorm_const_smul_le
          _ = ENNReal.ofReal C_rep * eLpNorm h (ENNReal.ofReal p) μ := by
              congr 1
              rw [Real.enorm_eq_ofReal_abs, abs_of_nonneg (by positivity)]
    _ ≤ ENNReal.ofReal C_rep * (ENNReal.ofReal M * eLpNorm g (ENNReal.ofReal p) μ) := by
        gcongr
    _ = ENNReal.ofReal (C_rep * M) * eLpNorm g (ENNReal.ofReal p) μ := by
        rw [← mul_assoc, ← ENNReal.ofReal_mul (by positivity)]
    _ ≤ ENNReal.ofReal (C_poinc_val d) * eLpNorm g (ENNReal.ofReal p) μ := by
        gcongr
        simp only [C_rep, M, C_poinc_val]
        have hd_pos : (0 : ℝ) < d := Nat.cast_pos.mpr (NeZero.pos d)
        have hvol_pos : (0 : ℝ) < (volume (Metric.ball (0 : E) 1)).toReal :=
          ENNReal.toReal_pos (measure_ball_pos volume 0 one_pos).ne' measure_ball_lt_top.ne
        have hdv_ne : (d : ℝ) * (volume (Metric.ball (0 : E) 1)).toReal ≠ 0 :=
          mul_ne_zero hd_pos.ne' hvol_pos.ne'
        rw [mul_one]
        rw [div_mul_comm, mul_comm ((d : ℝ) * _) 2, mul_div_cancel_of_imp (fun h => absurd h hdv_ne)]
        calc (2 : ℝ) ^ d * 2 * ↑d = 2 * 2 ^ d * ↑d := by ring
          _ ≥ 2 * 2 ^ d * 1 := by
              nlinarith [pow_pos (show (0 : ℝ) < 2 by norm_num) d,
                show (1 : ℝ) ≤ (d : ℝ) from by
                  exact_mod_cast Nat.one_le_iff_ne_zero.mpr (NeZero.ne d)]
          _ = 2 * 2 ^ d := by ring

\end{lstlisting}
\begin{proof}
Set $B = \Bz{1}$, $\bar u = \fint_B u$, $f = u - \bar u$, $g = \lvert \nabla u\rvert$,
$K(z) = \lvert z\rvert^{1-d}$ and
\[
h(x) = \int_B g(y)\, K(x-y)\, dy,\qquad
M = d\,\mathrm{vol}(B)\cdot 2,\qquad
C_{\mathrm{rep}} = \frac{2^{d}}{d\,\mathrm{vol}(B)}.
\]
The representation formula yields $\lvert f(x)\rvert \le C_{\mathrm{rep}}\, h(x)$ pointwise on
$B$. Therefore $\lVert f\rVert_{\Lp{p}(B)} \le C_{\mathrm{rep}}\,\lVert h\rVert_{\Lp{p}(B)}$.

For the Young-type estimate $\lVert h\rVert_{\Lp{p}(B)} \le M\,\lVert g\rVert_{\Lp{p}(B)}$
we apply the weighted power mean lemma to the kernel weight $w(y) = K(x-y)$, using the
$L^{1}$ bound $\int_B K(x-y)\, dy \le M$ provided by \leanname{riesz\_kernel\_L1\_bound}.
This gives
\[
h(x)^{p}
\le M^{p-1}\, \int_B g(y)^{p}\, K(x-y)\, dy.
\]
Integrating in $x \in B$, using Fubini and the same kernel bound in the $x$ variable,
\[
\int_B h(x)^{p}\, dx
\le M^{p-1}\, \int_B g(y)^{p}\Bigl(\int_B K(x-y)\, dx\Bigr) dy
\le M^{p}\, \int_B g(y)^{p}\, dy.
\]
Taking $p$-th roots gives $\lVert h\rVert_{\Lp{p}(B)} \le M\,\lVert g\rVert_{\Lp{p}(B)}$.
Combining,
\[
\lVert f\rVert_{\Lp{p}(B)}
\le C_{\mathrm{rep}}\, M\, \lVert g\rVert_{\Lp{p}(B)}
= 2^{d+1}\, d\, \lVert g\rVert_{\Lp{p}(B)}
= C_{\mathrm{poinc}}(d)\,\lVert g\rVert_{\Lp{p}(B)}. \qedhere
\]
\end{proof}

A density argument lifts this to general $\Wkp{1}{p}(\Bz{1})$ functions.

\begin{theorem}[Poincar\'e on the unit ball for $W^{1,p}$,
\leanname{poincare\_unitBall\_W1p\_public}]
Let $1 < p < \infty$ and let $u \in \Wkp{1}{p}(\Bz{1})$ with weak gradient $G$. Then
\[
\Bigl\lVert u - \fint_{\Bz{1}} u \Bigr\rVert_{\Lp{p}(\Bz{1})}
\le C_{\mathrm{poinc}}(d)\, \lVert \lvert G\rvert\rVert_{\Lp{p}(\Bz{1})}.
\]
\end{theorem}

\begin{lstlisting}[style=Matlab-editor,basicstyle=\mlttfamily,escapechar=",mlshowsectionrules=true,mathescape=true,morecomment={[s]{\%\{}{\%\}}},language=lean,numbers=none]
/-- Unit-ball Poincare estimate packaged for direct reuse. -/
theorem poincare_unitBall_W1p_public
    {p : ℝ} (hp : 1 < p)
    {u : E → ℝ}
    (hw : MemW1pWitness (ENNReal.ofReal p) u (Metric.ball (0 : E) 1)) :
    eLpNorm (fun x => u x - ⨍ y in Metric.ball (0 : E) 1, u y ∂volume)
      (ENNReal.ofReal p) (volume.restrict (Metric.ball (0 : E) 1)) ≤
    ENNReal.ofReal (C_poinc_val d) *
      eLpNorm (fun x => ‖hw.weakGrad x‖) (ENNReal.ofReal p)
        (volume.restrict (Metric.ball (0 : E) 1)) :=
  poincare_unitBall_W1p (d := d) hp hw

\end{lstlisting}
\begin{proof}
Approximate $u$ by smooth $\psi_n \to u$ in $\Wkp{1}{p}(\Bz{1})$. Apply the smooth
Poincar\'e inequality to each $\psi_n$, then pass to the limit using strong $L^{p}$
convergence of $\psi_n - \fint \psi_n$ to $u - \fint u$ and of
$\nabla \psi_n$ to $G$.
\end{proof}

\subsection{Sobolev--Poincar\'e on the unit ball}

\begin{definition}[\leanname{C\_extensionGrad}, \leanname{C\_sobolevPoincare}]
For $d \ge 1$ and $p > 1$,
\begin{align*}
C_{\mathrm{ext\!-\!grad}}(d,p)
  &:= 1 + C_{\mathrm{ubExt\,grad}}(d,p)\bigl(C_{\mathrm{poinc}}(d) + 1\bigr),\\
C_{\mathrm{sobPoinc}}(d,p)
  &:= C_{\mathrm{gns}}(d,p)\, C_{\mathrm{ext\!-\!grad}}(d,p),
\end{align*}
where $C_{\mathrm{gns}}$ is the whole-space Gagliardo--Nirenberg--Sobolev constant and
$C_{\mathrm{ubExt\,grad}}$ comes from the unit-ball extension.
\end{definition}

\begin{lstlisting}[style=Matlab-editor,basicstyle=\mlttfamily,escapechar=",mlshowsectionrules=true,mathescape=true,morecomment={[s]{\%\{}{\%\}}},language=lean,numbers=none]
/-! ## Constants -/

/-- Intermediate constant for the extension gradient bound. -/
private noncomputable def C_extensionGrad (d : ℕ) [NeZero d] (p : ℝ) : ℝ≥0∞ :=
  1 + C_unitBallExtensionGrad d p * (ENNReal.ofReal (C_poinc_val d) + 1)

/-- The Sobolev-Poincare constant on the unit ball. -/
noncomputable def C_sobolevPoincare (d : ℕ) [NeZero d] (p : ℝ) : ℝ≥0∞ :=
  ENNReal.ofReal (C_gns d p) * C_extensionGrad d p

\end{lstlisting}
\begin{lemma}[Subtracting a constant preserves $W^{1,p}$,
\leanname{memW1pWitness\_sub\_const\_unitBall},
\leanname{memW1pWitness\_sub\_const\_unitBall\_weakGrad}]
If $u \in \Wkp{1}{p}(\Bz{1})$ has weak gradient $G$, then for any $c \in \R$ the function
$u - c$ is also in $\Wkp{1}{p}(\Bz{1})$ with the same weak gradient $G$.
\end{lemma}

\begin{lstlisting}[style=Matlab-editor,basicstyle=\mlttfamily,escapechar=",mlshowsectionrules=true,mathescape=true,morecomment={[s]{\%\{}{\%\}}},language=lean,numbers=none]
/-! ## Subtracting constants -/

/-- Subtracting a constant from a `W^{1,p}` witness on the unit ball preserves
membership and leaves the weak gradient unchanged. -/
noncomputable def memW1pWitness_sub_const_unitBall
    {p : ℝ} (hp : 1 < p)
    {u : E → ℝ}
    (hw : MemW1pWitness (ENNReal.ofReal p) u (Metric.ball (0 : E) 1))
    (c : ℝ) :
    MemW1pWitness (ENNReal.ofReal p) (fun x => u x - c) (Metric.ball (0 : E) 1) := by
  haveI : IsFiniteMeasure (volume.restrict (Metric.ball (0 : E) 1)) :=
    ⟨by rw [Measure.restrict_apply_univ]; exact measure_ball_lt_top.lt_top⟩
  refine
    { memLp := hw.memLp.sub (memLp_const c)
      weakGrad := hw.weakGrad
      weakGrad_component_memLp := hw.weakGrad_component_memLp
      isWeakGrad := ?_ }
  intro i φ hφ hφ_supp hφ_sub
  let ei : E := EuclideanSpace.single i (1 : ℝ)
  have hkey := hw.isWeakGrad i φ hφ hφ_supp hφ_sub
  have hconst :
      HasWeakPartialDeriv i (fun _ : E => (0 : ℝ)) (fun _ : E => c)
        (Metric.ball (0 : E) 1) := by
    simpa [ei] using
      (HasWeakPartialDeriv.of_contDiff (Ω := Metric.ball (0 : E) 1) (i := i)
        (f := fun _ : E => c) isOpen_ball contDiff_const)
  have hconst_zero :
      ∫ x in Metric.ball (0 : E) 1, c * (fderiv ℝ φ x) ei = 0 := by
    simpa using hconst φ hφ hφ_supp hφ_sub
  have hderiv_cont : Continuous (fun x => (fderiv ℝ φ x) ei) :=
    (hφ.continuous_fderiv (by simp : ((⊤ : ℕ∞) : WithTop ℕ∞) ≠ 0)).clm_apply
      continuous_const
  have hderiv_supp : HasCompactSupport (fun x => (fderiv ℝ φ x) ei) :=
    hφ_supp.fderiv_apply (𝕜 := ℝ) ei
  have hu_int : Integrable (fun x => u x * (fderiv ℝ φ x) ei)
      (volume.restrict (Metric.ball (0 : E) 1)) := by
    have hu_loc : LocallyIntegrable u (volume.restrict (Metric.ball (0 : E) 1)) :=
      hw.memLp.locallyIntegrable (by
        rw [show (1 : ℝ≥0∞) = ENNReal.ofReal 1 from by simp]
        exact ENNReal.ofReal_le_ofReal (le_of_lt hp))
    simpa [smul_eq_mul] using
      hu_loc.integrable_smul_right_of_hasCompactSupport hderiv_cont hderiv_supp
  have hc_int : Integrable (fun x => c * (fderiv ℝ φ x) ei)
      (volume.restrict (Metric.ball (0 : E) 1)) :=
    (hderiv_cont.integrable_of_hasCompactSupport hderiv_supp).const_mul c
  have hsub_fun :
      (fun x => (u x - c) * (fderiv ℝ φ x) ei) =
        (fun x => u x * (fderiv ℝ φ x) ei - c * (fderiv ℝ φ x) ei) := by
    ext x
    ring
  rw [hsub_fun, integral_sub hu_int hc_int, hkey, hconst_zero]
  simp

/-- The weak gradient of the subtracted witness equals the original weak
gradient. -/
theorem memW1pWitness_sub_const_unitBall_weakGrad
    {p : ℝ} (hp : 1 < p) {u : E → ℝ}
    (hw : MemW1pWitness (ENNReal.ofReal p) u (Metric.ball (0 : E) 1)) (c : ℝ) :
    (memW1pWitness_sub_const_unitBall hp hw c).weakGrad = hw.weakGrad := by
  let _ := (inferInstance : NeZero d)
  simp [memW1pWitness_sub_const_unitBall]

\end{lstlisting}
\begin{proof}
$u - c \in L^{p}$ since $L^{p}$ is a vector space. Constants have zero weak gradient by
$\int_{\Bz{1}} c\, \partial_i \varphi = 0$ for compactly supported $\varphi$, hence
subtracting a constant from the test integral relation
$\int u\,\partial_i \varphi = -\int G_i \varphi$ leaves it unchanged.
\end{proof}

\begin{lemma}[Jensen for averages, \leanname{eLpNorm\_const\_average\_le}]
If $\mu$ is a finite measure, $1 \le p < \infty$ and $f \in L^{1}(\mu)$, then
\[
\Bigl\lVert x \mapsto \fint f\, d\mu \Bigr\rVert_{\Lp{p}(\mu)}
\le \lVert f \rVert_{\Lp{p}(\mu)}.
\]
\end{lemma}

\begin{lstlisting}[style=Matlab-editor,basicstyle=\mlttfamily,escapechar=",mlshowsectionrules=true,mathescape=true,morecomment={[s]{\%\{}{\%\}}},language=lean,numbers=none]
/-! ## Jensen for `eLpNorm` of averages -/

/-- **Jensen's inequality for `eLpNorm`**: the `L^p` norm of a constant equal to
the average is at most the `L^p` norm of the function.

For `1 ≤ p < ⊤`, an integrable `f`, and a finite measure `μ`:
  `eLpNorm (fun _ => ⨍ x, f x ∂μ) p μ ≤ eLpNorm f p μ`

This is a consequence of Jensen's inequality for the convex function `t ↦ |t|^p`,
proved here via Hölder's inequality (`eLpNorm_le_eLpNorm_mul_rpow_measure_univ`)
combined with the norm bound for the Bochner integral average.

The proof avoids `EuclideanSpace` and works for any real-valued integrable function. -/
theorem eLpNorm_const_average_le
    {α : Type*} [MeasurableSpace α] {μ : Measure α} [IsFiniteMeasure μ]
    {p : ℝ≥0∞} (hp : 1 ≤ p) (hp_top : p ≠ ⊤)
    {f : α → ℝ} (hf : Integrable f μ) :
    eLpNorm (fun _ => ⨍ x, f x ∂μ) p μ ≤ eLpNorm f p μ := by
  let _ := hp_top
  -- Trivial case: μ = 0
  rcases eq_or_ne μ 0 with rfl | hμ
  · simp [eLpNorm_measure_zero]
  have hμ_ne : μ Set.univ ≠ 0 := by rwa [ne_eq, Measure.measure_univ_eq_zero]
  have hμ_top : μ Set.univ ≠ ⊤ := measure_ne_top μ _
  have hp_ne : p ≠ 0 := (lt_of_lt_of_le one_pos hp).ne'
  -- Proof: ‖⨍ f‖ ≤ (μ.real univ)⁻¹ * ‖∫ f‖ ≤ (μ.real univ)⁻¹ * (eLpNorm f 1 μ).toReal
  -- Then convert to ENNReal.
  have h_avg_enorm_le :
      ‖⨍ x, f x ∂μ‖ₑ ≤ (μ Set.univ)⁻¹ * eLpNorm f 1 μ := by
    have h_avg_bound : ‖⨍ x, f x ∂μ‖ ≤
        (μ.real Set.univ)⁻¹ * (eLpNorm f 1 μ).toReal := by
      simp only [average_eq]
      calc ‖(μ.real Set.univ)⁻¹ • ∫ x, f x ∂μ‖
          ≤ ‖(μ.real Set.univ)⁻¹‖ * ‖∫ x, f x ∂μ‖ := norm_smul_le _ _
        _ = (μ.real Set.univ)⁻¹ * ‖∫ x, f x ∂μ‖ := by
            rw [Real.norm_of_nonneg (inv_nonneg.mpr measureReal_nonneg)]
        _ ≤ (μ.real Set.univ)⁻¹ * (eLpNorm f 1 μ).toReal := by
            gcongr
            calc ‖∫ x, f x ∂μ‖
                ≤ ∫ x, ‖f x‖ ∂μ := norm_integral_le_integral_norm _
              _ = (∫⁻ x, ‖f x‖ₑ ∂μ).toReal :=
                  integral_norm_eq_lintegral_enorm hf.aestronglyMeasurable
              _ = (eLpNorm f 1 μ).toReal := by rw [eLpNorm_one_eq_lintegral_enorm]
    -- Convert ℝ bound to ENNReal
    have h_inv_eq : ENNReal.ofReal (μ.real Set.univ)⁻¹ = (μ Set.univ)⁻¹ := by
      rw [measureReal_def, ENNReal.ofReal_inv_of_pos (ENNReal.toReal_pos hμ_ne hμ_top),
        ENNReal.ofReal_toReal hμ_top]
    calc ‖⨍ x, f x ∂μ‖ₑ
        ≤ ENNReal.ofReal ((μ.real Set.univ)⁻¹ * (eLpNorm f 1 μ).toReal) := by
          simp only [Real.enorm_eq_ofReal_abs, ← Real.norm_eq_abs]
          exact ENNReal.ofReal_le_ofReal h_avg_bound
      _ = ENNReal.ofReal (μ.real Set.univ)⁻¹ * ENNReal.ofReal (eLpNorm f 1 μ).toReal :=
          ENNReal.ofReal_mul (inv_nonneg.mpr measureReal_nonneg)
      _ ≤ (μ Set.univ)⁻¹ * eLpNorm f 1 μ := by
          rw [h_inv_eq]; gcongr; exact ENNReal.ofReal_toReal_le
  set exp := 1 / (1 : ℝ≥0∞).toReal - 1 / p.toReal with hexp_def
  have h_holder : eLpNorm f 1 μ ≤
      eLpNorm f p μ * μ Set.univ ^ exp :=
    eLpNorm_le_eLpNorm_mul_rpow_measure_univ hp hf.aestronglyMeasurable
  -- eLpNorm(const c) = ‖c‖ₑ * μ(univ)^{1/p}  [eLpNorm_const]
  -- ≤ (μ univ)⁻¹ * eLpNorm f p μ * μ(univ)^exp * μ(univ)^{1/p}
  -- = eLpNorm f p μ * ((μ univ)⁻¹ * μ(univ)^{exp + 1/p})
  -- Since exp + 1/p = 1/1 - 1/p + 1/p = 1: = eLpNorm f p μ * (μ univ)⁻¹ * μ univ = eLpNorm f p μ
  have hexp_simp : exp + 1 / p.toReal = 1 := by
    simp [hexp_def, show (1 : ℝ≥0∞).toReal = 1 from ENNReal.toReal_one]
  calc eLpNorm (fun _ => ⨍ x, f x ∂μ) p μ
      = ‖⨍ x, f x ∂μ‖ₑ * μ Set.univ ^ (1 / p.toReal) := eLpNorm_const _ hp_ne hμ
    _ ≤ ((μ Set.univ)⁻¹ * eLpNorm f 1 μ) * μ Set.univ ^ (1 / p.toReal) := by
        gcongr
    _ ≤ ((μ Set.univ)⁻¹ * (eLpNorm f p μ * μ Set.univ ^ exp)) *
          μ Set.univ ^ (1 / p.toReal) := by gcongr
    _ = eLpNorm f p μ * ((μ Set.univ)⁻¹ * (μ Set.univ ^ exp *
          μ Set.univ ^ (1 / p.toReal))) := by ring
    _ = eLpNorm f p μ * ((μ Set.univ)⁻¹ * μ Set.univ ^ (exp + 1 / p.toReal)) := by
        congr 2; exact (ENNReal.rpow_add exp (1 / p.toReal) hμ_ne hμ_top).symm
    _ = eLpNorm f p μ * ((μ Set.univ)⁻¹ * μ Set.univ ^ (1 : ℝ)) := by
        rw [hexp_simp]
    _ = eLpNorm f p μ * ((μ Set.univ)⁻¹ * μ Set.univ) := by
        rw [ENNReal.rpow_one]
    _ = eLpNorm f p μ := by rw [ENNReal.inv_mul_cancel hμ_ne hμ_top, mul_one]

\end{lstlisting}
\begin{proof}
Apply H\"older's inequality to bound the average $\lvert\fint f\rvert$ by the
$L^{p}$ norm of $f$ times $\mu(X)^{1-1/p}$, then integrate.
\end{proof}

\begin{theorem}[Sobolev--Poincar\'e on $\Bz{1}$, \leanname{sobolev\_poincare\_unitBall}]
Let $1 < p < d$ and $u \in \Wkp{1}{p}(\Bz{1})$ with weak gradient $G$, and write
$p^{*} = dp/(d-p)$. Then
\[
\Bigl\lVert u - \fint_{\Bz{1}} u \Bigr\rVert_{\Lp{p^{*}}(\Bz{1})}
\le C_{\mathrm{sobPoinc}}(d,p)\, \lVert\lvert G\rvert\rVert_{\Lp{p}(\Bz{1})}.
\]
\end{theorem}

\begin{lstlisting}[style=Matlab-editor,basicstyle=\mlttfamily,escapechar=",mlshowsectionrules=true,mathescape=true,morecomment={[s]{\%\{}{\%\}}},language=lean,numbers=none]
/-! ## Main theorem -/

/-- Sobolev-Poincare inequality on the unit ball. -/
theorem sobolev_poincare_unitBall
    {p : ℝ} (hp : 1 < p) (hpd : p < (d : ℝ))
    {u : E → ℝ}
    (hw : MemW1pWitness (ENNReal.ofReal p) u (Metric.ball (0 : E) 1)) :
    eLpNorm (fun x => u x - ⨍ y in Metric.ball (0 : E) 1, u y ∂volume)
      (ENNReal.ofReal ((d : ℝ) * p / ((d : ℝ) - p)))
      (volume.restrict (Metric.ball (0 : E) 1)) ≤
    C_sobolevPoincare d p *
      eLpNorm (fun x => ‖hw.weakGrad x‖) (ENNReal.ofReal p)
        (volume.restrict (Metric.ball (0 : E) 1)) := by
  let B := Metric.ball (0 : E) 1
  let ubar := ⨍ y in B, u y ∂volume
  let v : E → ℝ := fun x => u x - ubar
  let pstar := ENNReal.ofReal ((d : ℝ) * p / ((d : ℝ) - p))
  let hwv := memW1pWitness_sub_const_unitBall (d := d) hp hw ubar
  have hwv_grad : hwv.weakGrad = hw.weakGrad :=
    memW1pWitness_sub_const_unitBall_weakGrad hp hw ubar
  obtain ⟨hwExt, hcpt, hagree, _hfun_bound, hgrad_bound⟩ :=
    exists_unitBall_W1p_extension (d := d) hp hwv
  have hW01 : MemW01p (ENNReal.ofReal p) (unitBallExtension (d := d) v) Set.univ :=
    memW01p_of_memW1p_of_tsupport_subset isOpen_univ hp hwExt.memW1p hcpt
      (Set.subset_univ _)
  obtain ⟨hwSob, hSob⟩ := sobolev_of_memW01p_univ (d := d) (le_of_lt hp) hpd hW01
  have hae_grad : hwSob.weakGrad =ᵐ[volume] hwExt.weakGrad := by
    have h := MemW1pWitness.ae_eq_p (d := d) isOpen_univ (show 1 ≤ p from le_of_lt hp)
      hwSob hwExt
    simpa [Measure.restrict_univ] using h
  have hSob' :
      eLpNorm (unitBallExtension (d := d) v) pstar volume ≤
        ENNReal.ofReal (C_gns d p) *
          eLpNorm (fun x => ‖hwExt.weakGrad x‖) (ENNReal.ofReal p) volume := by
    calc
      eLpNorm (unitBallExtension (d := d) v) pstar volume
          ≤ ENNReal.ofReal (C_gns d p) *
              eLpNorm (fun x => ‖hwSob.weakGrad x‖) (ENNReal.ofReal p) volume := hSob
      _ = ENNReal.ofReal (C_gns d p) *
            eLpNorm (fun x => ‖hwExt.weakGrad x‖) (ENNReal.ofReal p) volume := by
          congr 1
          exact eLpNorm_congr_ae (hae_grad.fun_comp (‖·‖))
  have hrestrict :
      eLpNorm v pstar (volume.restrict B) ≤
        eLpNorm (unitBallExtension (d := d) v) pstar volume := by
    calc
      eLpNorm v pstar (volume.restrict B)
          = eLpNorm (unitBallExtension (d := d) v) pstar (volume.restrict B) := by
            apply eLpNorm_congr_ae
            have : ∀ᵐ x ∂(volume.restrict B), v x = unitBallExtension (d := d) v x := by
              filter_upwards [ae_restrict_mem measurableSet_ball] with x hx
              exact (hagree x hx).symm
            exact this
      _ ≤ eLpNorm (unitBallExtension (d := d) v) pstar volume :=
          eLpNorm_mono_measure _ Measure.restrict_le_self
  have hv_poinc_raw :
      eLpNorm (fun x => u x - ⨍ y in B, u y ∂volume) (ENNReal.ofReal p)
        (volume.restrict B) ≤
      ENNReal.ofReal (C_poinc_val d) *
        eLpNorm (fun x => ‖hw.weakGrad x‖) (ENNReal.ofReal p)
          (volume.restrict B) :=
    poincare_unitBall_W1p (d := d) hp hw
  have hv_poinc :
      eLpNorm v (ENNReal.ofReal p) (volume.restrict B) ≤
        ENNReal.ofReal (C_poinc_val d) *
          eLpNorm (fun x => ‖hwv.weakGrad x‖) (ENNReal.ofReal p)
            (volume.restrict B) := by
    rw [hwv_grad]
    exact hv_poinc_raw
  have hgrad_absorb :
      eLpNorm (fun x => ‖hwExt.weakGrad x‖) (ENNReal.ofReal p) volume ≤
        C_extensionGrad d p *
          eLpNorm (fun x => ‖hwv.weakGrad x‖) (ENNReal.ofReal p)
            (volume.restrict B) :=
    extension_gradient_eLpNorm_bound (d := d) hp hwv hgrad_bound hv_poinc
  have hwv_grad_eq : (fun x => ‖hwv.weakGrad x‖) = (fun x => ‖hw.weakGrad x‖) := by
    ext x
    rw [hwv_grad]
  calc
    eLpNorm v pstar (volume.restrict B)
        ≤ eLpNorm (unitBallExtension (d := d) v) pstar volume := hrestrict
    _ ≤ ENNReal.ofReal (C_gns d p) *
          eLpNorm (fun x => ‖hwExt.weakGrad x‖) (ENNReal.ofReal p) volume := hSob'
    _ ≤ ENNReal.ofReal (C_gns d p) *
          (C_extensionGrad d p *
            eLpNorm (fun x => ‖hwv.weakGrad x‖) (ENNReal.ofReal p)
              (volume.restrict B)) := by
          gcongr
    _ = ENNReal.ofReal (C_gns d p) * C_extensionGrad d p *
          eLpNorm (fun x => ‖hwv.weakGrad x‖) (ENNReal.ofReal p)
            (volume.restrict B) := by
          rw [mul_assoc]
    _ = C_sobolevPoincare d p *
          eLpNorm (fun x => ‖hwv.weakGrad x‖) (ENNReal.ofReal p)
            (volume.restrict B) := by
          rfl
    _ = C_sobolevPoincare d p *
          eLpNorm (fun x => ‖hw.weakGrad x‖) (ENNReal.ofReal p)
            (volume.restrict B) := by
          rw [hwv_grad_eq]

\end{lstlisting}
\begin{proof}
Set $\bar u = \fint_{\Bz{1}} u$ and $v = u - \bar u$. By the previous lemma
$v \in \Wkp{1}{p}(\Bz{1})$ with weak gradient $G$. Let
$\widetilde v = \mathrm{ext}(v) \in \Wop{p}$ be the unit-ball extension, which has
compact support, agrees with $v$ on $\Bz{1}$ and satisfies the lintegral gradient bound
\[
\int \lvert \nabla \widetilde v\rvert^{p}
\le \int_{\Bz{1}} \lvert G\rvert^{p}
+ C_{\mathrm{ubExt\,grad}}(d,p)\,
\Bigl(\int_{\Bz{1}} \lvert v\rvert^{p} + \int_{\Bz{1}} \lvert G\rvert^{p}\Bigr).
\]
Combining this with Poincar\'e
$\lVert v\rVert_{\Lp{p}(\Bz{1})} \le C_{\mathrm{poinc}}(d)\,\lVert\lvert G\rvert\rVert_{\Lp{p}(\Bz{1})}$
and raising to the $p$-th power gives
\[
\lVert \lvert \nabla \widetilde v\rvert\rVert_{\Lp{p}(\R^{d})}
\le C_{\mathrm{ext\!-\!grad}}(d,p)\,\lVert\lvert G\rvert\rVert_{\Lp{p}(\Bz{1})}
\]
(this is the auxiliary lemma \leanname{extension\_gradient\_eLpNorm\_bound}, which uses
the algebraic inequalities $a^{p}+b^{p} \le (a+b)^{p}$ for $p\ge 1$ to absorb the
constants). The whole-space Gagliardo--Nirenberg--Sobolev inequality applied to
$\widetilde v$ gives
\[
\lVert \widetilde v\rVert_{\Lp{p^{*}}(\R^{d})}
\le C_{\mathrm{gns}}(d,p)\,
\lVert \lvert\nabla \widetilde v\rvert\rVert_{\Lp{p}(\R^{d})}.
\]
Restricting the left side to $\Bz{1}$, where $\widetilde v = v$, and chaining the two
displayed estimates with the Poincar\'e bound yields the conclusion with constant
$C_{\mathrm{sobPoinc}}(d,p) = C_{\mathrm{gns}}(d,p)\,C_{\mathrm{ext\!-\!grad}}(d,p)$.
\end{proof}

\begin{corollary}[Existential form, \leanname{sobolev\_poincare\_unitBall'}]
For $1 < p < d$ and $u$ in the predicate $\Wkp{1}{p}(\Bz{1})$ (existential form), there
exist a weak gradient $G$ and a constant $C = C_{\mathrm{sobPoinc}}(d,p)$ realising the
above Sobolev--Poincar\'e estimate.
\end{corollary}

\begin{lstlisting}[style=Matlab-editor,basicstyle=\mlttfamily,escapechar=",mlshowsectionrules=true,mathescape=true,morecomment={[s]{\%\{}{\%\}}},language=lean,numbers=none]
/-! ## Existential version -/

/-- Sobolev-Poincare inequality for `MemW1p` in existential form. -/
theorem sobolev_poincare_unitBall'
    {p : ℝ} (hp : 1 < p) (hpd : p < (d : ℝ))
    {u : E → ℝ}
    (hu : MemW1p (ENNReal.ofReal p) u (Metric.ball (0 : E) 1)) :
    ∃ C : ℝ≥0∞, ∃ G : E → E,
      HasWeakGrad G u (Metric.ball (0 : E) 1) ∧
      eLpNorm (fun x => u x - ⨍ y in Metric.ball (0 : E) 1, u y ∂volume)
        (ENNReal.ofReal ((d : ℝ) * p / ((d : ℝ) - p)))
        (volume.restrict (Metric.ball (0 : E) 1)) ≤
      C * eLpNorm (fun x => ‖G x‖) (ENNReal.ofReal p)
        (volume.restrict (Metric.ball (0 : E) 1)) := by
  have hfLp := hu.1
  choose gi hgi using hu.2
  let G : E → E := fun x => WithLp.toLp 2 (fun i => gi i x)
  let hw : MemW1pWitness (ENNReal.ofReal p) u (Metric.ball (0 : E) 1) :=
    { memLp := hfLp
      weakGrad := G
      weakGrad_component_memLp := by
        intro i
        simpa [G, PiLp.toLp_apply] using (hgi i).1
      isWeakGrad := by
        intro i
        simpa [G, PiLp.toLp_apply] using (hgi i).2 }
  exact ⟨C_sobolevPoincare d p, G, hw.isWeakGrad,
    sobolev_poincare_unitBall hp hpd hw⟩

\end{lstlisting}
\begin{corollary}[Smooth version, \leanname{sobolev\_poincare\_smooth\_unitBall}]
For $1 < p < d$ and $u \in C^{\infty}(E,\R)$,
\[
\Bigl\lVert u - \fint_{\Bz{1}} u \Bigr\rVert_{\Lp{p^{*}}(\Bz{1})}
\le C_{\mathrm{sobPoinc}}(d,p)\,\lVert\lvert\nabla u\rvert\rVert_{\Lp{p}(\Bz{1})}.
\]
\end{corollary}

\begin{lstlisting}[style=Matlab-editor,basicstyle=\mlttfamily,escapechar=",mlshowsectionrules=true,mathescape=true,morecomment={[s]{\%\{}{\%\}}},language=lean,numbers=none]
/-- Sobolev-Poincare for smooth functions on the unit ball. -/
theorem sobolev_poincare_smooth_unitBall
    {p : ℝ} (hp : 1 < p) (hpd : p < (d : ℝ))
    {u : E → ℝ} (hu : ContDiff ℝ (⊤ : ℕ∞) u) :
    eLpNorm (fun x => u x - ⨍ y in Metric.ball (0 : E) 1, u y ∂volume)
      (ENNReal.ofReal ((d : ℝ) * p / ((d : ℝ) - p)))
      (volume.restrict (Metric.ball (0 : E) 1)) ≤
    C_sobolevPoincare d p *
      eLpNorm (fun x => ‖fderiv ℝ u x‖) (ENNReal.ofReal p)
        (volume.restrict (Metric.ball (0 : E) 1)) := by
  let hw := smooth_memW1pWitness_unitBall (d := d) hp hpd hu
  have hmain := sobolev_poincare_unitBall (d := d) hp hpd hw
  suffices hgrad_eq :
      eLpNorm (fun x => ‖hw.weakGrad x‖) (ENNReal.ofReal p)
        (volume.restrict (Metric.ball (0 : E) 1)) =
      eLpNorm (fun x => ‖fderiv ℝ u x‖) (ENNReal.ofReal p)
        (volume.restrict (Metric.ball (0 : E) 1)) by
    rw [← hgrad_eq]
    exact hmain
  apply eLpNorm_congr_ae
  filter_upwards with x
  show ‖WithLp.toLp 2 (fun i => (fderiv ℝ u x) (EuclideanSpace.single i 1))‖ =
    ‖fderiv ℝ u x‖
  exact norm_fderiv_eq_norm_partials_local (d := d).symm

\end{lstlisting}
\begin{proof}
Build the canonical $W^{1,p}$ witness for the smooth $u$ via
\leanname{smooth\_memW1pWitness\_unitBall} and apply the previous theorem; the weak
gradient agrees a.e.\ with the classical gradient on $\Bz{1}$.
\end{proof}

\subsection{Stampacchia truncation}

\begin{lemma}[\leanname{deriv\_ne\_zero\_implies\_isolated\_zero}]
Let $f : \R \to \R$ be differentiable at $x$ with $f(x)=0$ and $f'(x) = c \neq 0$. Then
there exists $\delta > 0$ such that $f(x+h) \neq 0$ for all $h$ with $0 < \lvert h\rvert < \delta$.
\end{lemma}

\begin{lstlisting}[style=Matlab-editor,basicstyle=\mlttfamily,escapechar=",mlshowsectionrules=true,mathescape=true,morecomment={[s]{\%\{}{\%\}}},language=lean,numbers=none]
/-! ## Section 1: Pointwise 1D Stampacchia -/

/-- **Core lemma**: if `f` is differentiable at `x` with `f(x) = 0` and
`f'(x) ≠ 0`, then `f(x+h) ≠ 0` for all sufficiently small `h ≠ 0`.

Proof: the little-o condition gives `|f(x+h) - c·h| ≤ (|c|/2)·|h|` for small `h`.
Since `f(x) = 0`, triangle inequality gives `|f(x+h)| ≥ |c|·|h|/2 > 0`. -/
theorem deriv_ne_zero_implies_isolated_zero {f : ℝ → ℝ} {x : ℝ} {c : ℝ}
    (hf : HasDerivAt f c x) (hfx : f x = 0) (hc : c ≠ 0) :
    ∃ δ > 0, ∀ h, 0 < |h| → |h| < δ → f (x + h) ≠ 0 := by
  rw [hasDerivAt_iff_isLittleO_nhds_zero] at hf
  have hc_pos : 0 < |c| := abs_pos.mpr hc
  rw [Asymptotics.isLittleO_iff] at hf
  have hfilt := hf (half_pos hc_pos)
  rw [Filter.eventually_iff_exists_mem] at hfilt
  obtain ⟨S, hS_mem, hS⟩ := hfilt
  obtain ⟨δ, hδ_pos, hδ_sub⟩ := Metric.mem_nhds_iff.mp hS_mem
  refine ⟨δ, hδ_pos, fun h hh_pos hh_lt hf_eq => ?_⟩
  have hh_in : h ∈ S := hδ_sub (by simp; exact hh_lt)
  have h1 := hS h hh_in
  simp only [hfx, sub_zero, Real.norm_eq_abs, smul_eq_mul] at h1
  rw [hf_eq, zero_sub, abs_neg, abs_mul] at h1
  nlinarith

\end{lstlisting}
\begin{proof}
The little-$o$ definition of differentiability gives, for sufficiently small $h$,
$\lvert f(x+h) - c h\rvert \le (\lvert c\rvert/2)\, \lvert h\rvert$. With $f(x)=0$
the reverse triangle inequality yields $\lvert f(x+h)\rvert \ge (\lvert c\rvert/2)\,
\lvert h\rvert > 0$.
\end{proof}

\begin{corollary}[\leanname{HasDerivAt.deriv\_eq\_zero\_of\_zero\_and\_not\_isolated}]
If $f$ is differentiable at $x$ with $f(x) = 0$ and $x$ is a non-isolated zero, then
$f'(x) = 0$.
\end{corollary}

\begin{lstlisting}[style=Matlab-editor,basicstyle=\mlttfamily,escapechar=",mlshowsectionrules=true,mathescape=true,morecomment={[s]{\%\{}{\%\}}},language=lean,numbers=none]
/-- Contrapositive: at a non-isolated zero where `f` is differentiable, `f'(x) = 0`. -/
theorem HasDerivAt.deriv_eq_zero_of_zero_and_not_isolated {f : ℝ → ℝ} {x c : ℝ}
    (hf : HasDerivAt f c x) (hfx : f x = 0)
    (hni : ∀ δ > 0, ∃ y, y ≠ x ∧ |y - x| < δ ∧ f y = 0) :
    c = 0 := by
  by_contra hc
  obtain ⟨δ, hδ_pos, hδ⟩ := deriv_ne_zero_implies_isolated_zero hf hfx hc
  obtain ⟨y, hy_ne, hy_lt, hy_zero⟩ := hni δ hδ_pos
  exact hδ (y - x) (by rwa [abs_pos, sub_ne_zero]) hy_lt (by rwa [add_sub_cancel])

\end{lstlisting}
\begin{lemma}[Isolated points are countable, \leanname{countable\_isolated}]
For any $S \subseteq \R$, the set of isolated points of $S$ is countable.
\end{lemma}

\begin{lstlisting}[style=Matlab-editor,basicstyle=\mlttfamily,escapechar=",mlshowsectionrules=true,mathescape=true,morecomment={[s]{\%\{}{\%\}}},language=lean,numbers=none]
/-! ## Section 2: Isolated points are countable → a.e. version -/

/-- In ℝ, the set of isolated points of any subset S is countable.

Proof: cover by ⋃_{p,q ∈ ℚ} {unique element of S in (p,q)}.
Each fiber has at most one element (uniqueness), and ℚ × ℚ is countable. -/
theorem countable_isolated (S : Set ℝ) :
    Set.Countable {x ∈ S | ∃ ε > 0, ∀ y ∈ S, dist y x < ε → y = x} := by
  set T := {x ∈ S | ∃ ε > 0, ∀ y ∈ S, dist y x < ε → y = x}
  apply Set.Countable.mono (show T ⊆
    ⋃ p : ℚ, ⋃ q : ℚ, {x ∈ S | (p : ℝ) < x ∧ x < q ∧ ∀ y ∈ S, (p : ℝ) < y → y < q → y = x}
    from ?_)
  · apply Set.countable_iUnion; intro p
    apply Set.countable_iUnion; intro q
    rcases Set.eq_empty_or_nonempty
      {x ∈ S | (p : ℝ) < x ∧ x < q ∧ ∀ y ∈ S, (p : ℝ) < y → y < q → y = x} with h | ⟨z, hz⟩
    · rw [h]; exact Set.countable_empty
    · apply (Set.countable_singleton z).mono
      intro y hy
      simp only [Set.mem_sep_iff] at hy hz
      exact hz.2.2.2 y hy.1 hy.2.1 hy.2.2.1
  · intro x ⟨hxS, ε, hε, hiso⟩
    simp only [Set.mem_iUnion, Set.mem_sep_iff]
    obtain ⟨p, hp1, hp2⟩ := exists_rat_btwn (show x - ε < x by linarith)
    obtain ⟨q, hq1, hq2⟩ := exists_rat_btwn (show x < x + ε by linarith)
    exact ⟨p, q, hxS, hp2, hq1, fun y hyS hpy hyq => hiso y hyS (by
      rw [Real.dist_eq, abs_lt]; constructor <;> linarith)⟩

\end{lstlisting}
\begin{proof}
Each isolated point $x$ admits a rational interval $(p,q) \ni x$ inside which $x$ is the
only point of $S$. Map $x \mapsto (p,q) \in \Q^{2}$; the map is well-defined and
injective on isolated points.
\end{proof}

\begin{lemma}[1D Stampacchia, a.e.\ form, \leanname{deriv\_eq\_zero\_ae\_on\_zeroSet}]
If $f : \R \to \R$ is differentiable a.e., then $f' = 0$ a.e.\ on $\{f = 0\}$.
\end{lemma}

\begin{lstlisting}[style=Matlab-editor,basicstyle=\mlttfamily,escapechar=",mlshowsectionrules=true,mathescape=true,morecomment={[s]{\%\{}{\%\}}},language=lean,numbers=none]
/-- **1D Stampacchia (a.e.)**: for `f` differentiable a.e., `f' = 0` a.e. on `{f = 0}`.

Proof: {f=0 ∧ f'≠0 ∧ f differentiable} ⊆ {isolated points of {f=0}}, which is
countable (hence measure 0). The remaining set {f not differentiable} has measure 0
by hypothesis. -/
theorem deriv_eq_zero_ae_on_zeroSet {f : ℝ → ℝ}
    (hf_diff : ∀ᵐ x, HasDerivAt f (deriv f x) x) :
    ∀ᵐ x, f x = 0 → deriv f x = 0 := by
  rw [ae_iff]
  have heq : {x | ¬(f x = 0 → deriv f x = 0)} = {x | f x = 0 ∧ deriv f x ≠ 0} := by
    ext x; simp only [mem_setOf_eq, Classical.not_imp, ne_eq]
  rw [heq]
  rw [ae_iff] at hf_diff
  apply le_antisymm _ (zero_le _)
  calc volume {x | f x = 0 ∧ deriv f x ≠ 0}
      ≤ volume ({x | f x = 0 ∧ deriv f x ≠ 0 ∧ HasDerivAt f (deriv f x) x} ∪
                {x | ¬HasDerivAt f (deriv f x) x}) := by
        apply measure_mono; intro x ⟨hfx, hdx⟩
        by_cases h : HasDerivAt f (deriv f x) x
        · left; exact ⟨hfx, hdx, h⟩
        · right; exact h
    _ ≤ volume {x | f x = 0 ∧ deriv f x ≠ 0 ∧ HasDerivAt f (deriv f x) x} +
        volume {x | ¬HasDerivAt f (deriv f x) x} := measure_union_le _ _
    _ ≤ volume {x | f x = 0 ∧ deriv f x ≠ 0 ∧ HasDerivAt f (deriv f x) x} + 0 := by
        gcongr; exact hf_diff.le
    _ = volume {x | f x = 0 ∧ deriv f x ≠ 0 ∧ HasDerivAt f (deriv f x) x} := add_zero _
    _ ≤ 0 := by
        -- Every such point is isolated in {f = 0}
        have hsub : {x | f x = 0 ∧ deriv f x ≠ 0 ∧ HasDerivAt f (deriv f x) x} ⊆
            {x ∈ {x | f x = 0} | ∃ ε > 0, ∀ y ∈ {x | f x = 0}, dist y x < ε → y = x} := by
          intro x ⟨hfx, hdx, hd⟩
          refine ⟨hfx, ?_⟩
          obtain ⟨δ, hδ, hiso⟩ := deriv_ne_zero_implies_isolated_zero hd hfx hdx
          exact ⟨δ, hδ, fun y hy hyd => by
            by_contra hne
            exact hiso (y - x) (by rwa [abs_pos, sub_ne_zero])
              (by rwa [Real.dist_eq] at hyd) (by rwa [add_sub_cancel])⟩
        calc volume _ ≤ volume _ := measure_mono hsub
          _ = 0 := (countable_isolated _).measure_zero _

\end{lstlisting}
\begin{proof}
The bad set $\{f = 0,\ f' \neq 0\}$ is contained in the union of the points where $f$
is not differentiable (a null set by hypothesis) and the set of isolated zeros of
$\{f = 0\}$, which is countable by the previous lemma. Both have measure zero.
\end{proof}

\begin{lemma}[Du Bois-Reymond fundamental lemma, \leanname{du\_bois\_reymond}]
Let $g \in L^{1}(a,b)$ satisfy $\int_a^b g\,\varphi' = 0$ for every test function
$\varphi \in C_c^{\infty}(a,b)$. Then $g$ is constant a.e.
\end{lemma}

\begin{lstlisting}[style=Matlab-editor,basicstyle=\mlttfamily,escapechar=",mlshowsectionrules=true,mathescape=true,morecomment={[s]{\%\{}{\%\}}},language=lean,numbers=none]
/-- **Du Bois-Reymond lemma**: if `g ∈ L^1(a,b)` and `∫ g · φ' = 0` for all
smooth compactly supported `φ` in `(a,b)`, then `g` is constant a.e.

Proof (Evans, Appendix C): fix `ψ₀ ∈ C_c^∞(a,b)` with `∫ψ₀ = 1`, set `c = ∫g·ψ₀`.
For any test `η`, decompose `η = [η - (∫η)·ψ₀] + (∫η)·ψ₀`. The first part has zero
integral, hence is `φ'` for some test `φ` (antiderivative lemma). The hypothesis gives
`∫g·φ' = 0`, so `∫g·η = c·∫η`. Then `∫(g-c)·η = 0` for all test `η`, and the
fundamental lemma (`ae_eq_zero_of_integral_contDiff_smul_eq_zero`) gives `g = c` a.e. -/
theorem du_bois_reymond
    {a b : ℝ} (hab : a < b)
    {g : ℝ → ℝ} (hg : IntegrableOn g (Ioo a b) volume)
    (htest : ∀ φ : ℝ → ℝ, ContDiff ℝ (⊤ : ℕ∞) φ → HasCompactSupport φ →
      tsupport φ ⊆ Ioo a b → ∫ x in Ioo a b, g x * deriv φ x = 0) :
    ∃ C : ℝ, g =ᵐ[volume.restrict (Ioo a b)] fun _ => C := by
  have ⟨ψ₀, hψ₀_smooth, hψ₀_cs, hψ₀_supp, hψ₀_int⟩ :
      ∃ ψ₀ : ℝ → ℝ, ContDiff ℝ (⊤ : ℕ∞) ψ₀ ∧ HasCompactSupport ψ₀ ∧
        tsupport ψ₀ ⊆ Ioo a b ∧ ∫ x, ψ₀ x = 1 := by
    set f : ContDiffBump ((a + b) / 2) :=
      ⟨(b - a) / 8, (b - a) / 4, by linarith, by linarith⟩
    refine ⟨f.normed volume, f.contDiff_normed, f.hasCompactSupport_normed, ?_,
      f.integral_normed⟩
    rw [f.tsupport_normed_eq]; intro x hx
    simp only [mem_closedBall, Real.dist_eq] at hx; rw [mem_Ioo]
    have := abs_le.mp (show |x - (a + b) / 2| ≤ (b - a) / 4 from hx)
    constructor <;> linarith
  set c := ∫ x in Ioo a b, g x * ψ₀ x
  have hmain : ∀ η : ℝ → ℝ, ContDiff ℝ (⊤ : ℕ∞) η → HasCompactSupport η →
      tsupport η ⊆ Ioo a b → ∫ x in Ioo a b, (g x - c) * η x = 0 := by
    intro η hη hη_cs hη_supp
    set α := ∫ x, η x
    set η₀ : ℝ → ℝ := fun x => η x - α * ψ₀ x with hη₀_def
    have hαψ₀_cs : HasCompactSupport (fun x => α * ψ₀ x) := by
      rw [hasCompactSupport_def]
      apply hψ₀_cs.isCompact.of_isClosed_subset isClosed_closure
      exact closure_mono (fun x hx => by
        simp only [mem_support] at hx ⊢; intro h; exact hx (by simp [h]))
    have hαψ₀_tsup : tsupport (fun x => α * ψ₀ x) ⊆ tsupport ψ₀ :=
      closure_mono (fun x hx => by simp only [mem_support] at hx ⊢; intro h; exact hx (by simp [h]))
    have hunion_closed : IsClosed (tsupport η ∪ tsupport (fun x => α * ψ₀ x)) :=
      isClosed_closure.union isClosed_closure
    have hη₀_smooth : ContDiff ℝ (⊤ : ℕ∞) η₀ := hη.sub (contDiff_const.mul hψ₀_smooth)
    have hη₀_cs : HasCompactSupport η₀ := by
      rw [hasCompactSupport_def]; exact (hη_cs.isCompact.union hαψ₀_cs.isCompact).of_isClosed_subset
        isClosed_closure (closure_minimal (show support η₀ ⊆ _ from hη₀_def ▸ support_sub_tsupport) hunion_closed)
    have hη₀_supp : tsupport η₀ ⊆ Ioo a b :=
      (closure_minimal (hη₀_def ▸ support_sub_tsupport) hunion_closed).trans
        (union_subset hη_supp (hαψ₀_tsup.trans hψ₀_supp))
    have hη₀_int : ∫ x, η₀ x = 0 := by
      have : ∫ x, η₀ x = (∫ x, η x) - ∫ x, α * ψ₀ x :=
        integral_sub (integrable_of_smooth_cs hη hη_cs)
          ((integrable_of_smooth_cs hψ₀_smooth hψ₀_cs).const_mul α)
      rw [this, MeasureTheory.integral_const_mul, hψ₀_int, mul_one, sub_self]
    obtain ⟨φ, hφ_smooth, hφ_cs, hφ_supp, hφ_deriv⟩ :=
      exists_antideriv_of_zero_integral hab hη₀_smooth hη₀_cs hη₀_supp hη₀_int
    have key : ∫ x in Ioo a b, g x * η₀ x = 0 := by
      rw [show (fun x => g x * η₀ x) = (fun x => g x * deriv φ x) from by ext x; rw [hφ_deriv]]
      exact htest φ hφ_smooth hφ_cs hφ_supp
    have hgη := integrableOn_mul_of_smooth_cs hg hη.continuous hη_cs
    have hgψ₀ := integrableOn_mul_of_smooth_cs hg hψ₀_smooth.continuous hψ₀_cs
    have gη_eq : ∫ x in Ioo a b, g x * η x = α * ∫ x in Ioo a b, g x * ψ₀ x := by
      have h_pw : ∀ x, g x * η₀ x = g x * η x - α * (g x * ψ₀ x) := by
        intro x; simp [hη₀_def, mul_sub, mul_comm α, mul_assoc]
      have h_sub : ∫ x in Ioo a b, g x * η₀ x =
          (∫ x in Ioo a b, g x * η x) - ∫ x in Ioo a b, α * (g x * ψ₀ x) := by
        simp_rw [h_pw]; exact integral_sub hgη (hgψ₀.const_mul α)
      have h_smul : ∫ x in Ioo a b, α * (g x * ψ₀ x) =
          α * ∫ x in Ioo a b, g x * ψ₀ x := by
        simp_rw [show ∀ x, α * (g x * ψ₀ x) = α • (g x * ψ₀ x) from
          fun x => (smul_eq_mul _ _).symm]; exact integral_smul α _
      linarith [key, h_sub, h_smul]
    have hcη : IntegrableOn (fun x => c * η x) (Ioo a b) :=
      ((integrable_of_smooth_cs hη hη_cs).integrableOn).const_mul c
    have hη_Ioo : ∫ x in Ioo a b, η x = α :=
      setIntegral_eq_integral_of_forall_compl_eq_zero (fun x hx =>
        image_eq_zero_of_notMem_tsupport (fun hmem => hx (hη_supp hmem)))
    have hcη_eq : ∫ x in Ioo a b, c * η x = c * α := by
      simp_rw [show ∀ x, c * η x = c • η x from fun x => (smul_eq_mul _ _).symm]
      rw [MeasureTheory.integral_smul, hη_Ioo, smul_eq_mul]
    simp_rw [fun x => show (g x - c) * η x = g x * η x - c * η x from by ring]
    rw [integral_sub hgη hcη, gη_eq, hcη_eq]; ring
  refine ⟨c, ?_⟩
  set h := (Ioo a b).indicator (fun x => g x - c)
  have hh_li : LocallyIntegrable h volume := by
    exact ((integrable_indicator_iff measurableSet_Ioo).mpr
      (hg.sub (integrableOn_const (hs := measure_Ioo_lt_top.ne)))).locallyIntegrable
  have hh_zero : ∀ ψ : ℝ → ℝ, ContDiff ℝ (⊤ : ℕ∞) ψ → HasCompactSupport ψ →
      ∫ x, ψ x • h x = 0 := by
    intro ψ hψ hψ_cs
    simp only [smul_eq_mul]
    rw [show (fun x => ψ x * h x) = (Ioo a b).indicator (fun x => ψ x * (g x - c)) from by
      ext x; by_cases hx : x ∈ Ioo a b <;> simp [indicator, hx, h]]
    rw [MeasureTheory.integral_indicator measurableSet_Ioo,
        show (fun x => ψ x * (g x - c)) = (fun x => (g x - c) * ψ x) from by ext; ring]
    exact smooth_cutoff_approximation hab
      (hg.sub (integrableOn_const (hs := measure_Ioo_lt_top.ne))) hψ hψ_cs hmain
  have hae := ae_eq_zero_of_integral_contDiff_smul_eq_zero hh_li hh_zero
  rw [Filter.EventuallyEq, ae_restrict_iff' measurableSet_Ioo]
  filter_upwards [hae] with x hx hx_mem
  simp only [h, indicator_of_mem hx_mem] at hx
  linarith

\end{lstlisting}
\begin{proof}
Fix $\psi_0 \in C_c^{\infty}(a,b)$ with $\int \psi_0 = 1$ and let
$c = \int g\,\psi_0$. For any test $\eta$ write $\eta = \eta_0 + \alpha \psi_0$ where
$\alpha = \int\eta$ and $\int \eta_0 = 0$. Since $\eta_0$ has zero integral, it is the
derivative of a test function $\varphi$ (antiderivative lemma), so the hypothesis gives
$\int g \eta_0 = 0$, hence $\int g\,\eta = c\,\int\eta$. Equivalently
$\int (g - c)\eta = 0$ for all test $\eta$, so by the fundamental lemma of the calculus
of variations $g = c$ a.e.
\end{proof}

\begin{lemma}[$W^{1,1}$ AC representative, \leanname{w11\_ae\_eq\_ac\_representative}]
Let $u, g \in L^{1}(a,b)$ with the weak-derivative relation
$\int u\,\varphi' = -\int g\,\varphi$ for all $\varphi \in C_c^{\infty}(a,b)$. Then there
is a constant $C$ such that
\[
u(x) = C + \int_a^{x} g(t)\, dt \qquad \text{for a.e. } x \in (a,b).
\]
\end{lemma}

\begin{lstlisting}[style=Matlab-editor,basicstyle=\mlttfamily,escapechar=",mlshowsectionrules=true,mathescape=true,morecomment={[s]{\%\{}{\%\}}},language=lean,numbers=none]
/-- For `u ∈ W^{1,1}(a,b)` with weak derivative `g`, `u` agrees a.e. with
`x ↦ C + ∫_a^x g(t) dt` for some constant `C`.

Proof: define `F(x) = ∫_a^x g`. By AC-IBP, `F` also has weak derivative `g` on `(a,b)`.
Then `u - F` has zero weak derivative. By `du_bois_reymond`, `u - F = D` a.e. -/
theorem w11_ae_eq_ac_representative
    {a b : ℝ} (hab : a < b) {u g : ℝ → ℝ}
    (hu : IntegrableOn u (Ioo a b) volume)
    (hg : IntegrableOn g (Ioo a b) volume)
    (hweak : ∀ φ : ℝ → ℝ, ContDiff ℝ (⊤ : ℕ∞) φ → HasCompactSupport φ →
      tsupport φ ⊆ Ioo a b →
      ∫ x in Ioo a b, u x * deriv φ x = -∫ x in Ioo a b, g x * φ x) :
    ∃ C : ℝ, u =ᵐ[volume.restrict (Ioo a b)]
      fun x => C + ∫ t in a..x, g t := by
  set F := fun x => ∫ t in a..x, g t
  have h_uF_test : ∀ φ : ℝ → ℝ, ContDiff ℝ (⊤ : ℕ∞) φ → HasCompactSupport φ →
      tsupport φ ⊆ Ioo a b →
      ∫ x in Ioo a b, (u x - F x) * deriv φ x = 0 := by
    intro φ hφ hφ_cs hφ_supp
    -- Convert set integral on Ioo to interval integral (they agree for volume)
    have conv := setIntegral_Ioo_eq_interval hab.le
    rw [conv]; simp_rw [sub_mul]
    -- Split ∫(u*φ' - F*φ') = ∫u*φ' - ∫F*φ'
    have hu_ii : IntervalIntegrable (fun x => u x * deriv φ x) volume a b :=
      integrableOn_Ioo_intervalIntegrable hab.le
        ((hu.bdd_mul (hφ.continuous_deriv (by norm_cast)).aestronglyMeasurable.restrict
          (ae_of_all _ fun x => (hφ_cs.deriv.exists_bound_of_continuous
            (hφ.continuous_deriv (by norm_cast))).choose_spec x)).congr
          (ae_of_all _ fun x => by ring))
    have hF_ii : IntervalIntegrable (fun x => (∫ t in a..x, g t) * deriv φ x) volume a b := by
      apply integrableOn_Ioo_intervalIntegrable hab.le
      have hcont := (intervalIntegral.continuousOn_primitive (μ := volume)
        (integrableOn_Icc_of_Ioo hg)).congr
        (fun x hx => by show ∫ t in a..x, g t = _; rw [intervalIntegral.integral_of_le hx.1])
      exact (hcont.mul (hφ.continuous_deriv (by norm_cast)).continuousOn).integrableOn_compact
        isCompact_Icc |>.mono_set Ioo_subset_Icc_self
    rw [intervalIntegral.integral_sub hu_ii hF_ii,
        ← conv, hweak φ hφ hφ_cs hφ_supp, conv, ac_ibp_step hab hg hφ hφ_supp]; ring
  obtain ⟨D, hD⟩ := du_bois_reymond hab
    (hu.sub (integrable_primitive hab hg))
    h_uF_test
  exact ⟨D, hD.mono (fun x hx => by simp at hx; linarith)⟩

\end{lstlisting}
\begin{proof}
Let $F(x) = \int_a^{x} g$. By integration by parts for absolutely continuous functions,
$F$ has weak derivative $g$. Hence $u - F$ has zero weak derivative, and du Bois-Reymond
gives $u - F = C$ a.e.
\end{proof}

\begin{theorem}[1D Stampacchia for $W^{1,1}$, \leanname{stampacchia\_1d}]
Under the hypotheses of the previous lemma, $g(x) = 0$ for a.e.\ $x \in (a,b)$
with $u(x) = 0$.
\end{theorem}

\begin{lstlisting}[style=Matlab-editor,basicstyle=\mlttfamily,escapechar=",mlshowsectionrules=true,mathescape=true,morecomment={[s]{\%\{}{\%\}}},language=lean,numbers=none]
/-- **1D Stampacchia for W^{1,1}**: the weak derivative `g` vanishes a.e. on `{u = 0}`.

Proof: replace `u` by its AC representative `F` (from `w11_ae_eq_ac_representative`).
Then `F' = g` a.e. by the FTC. Apply `deriv_eq_zero_ae_on_zeroSet` to `F`. -/
theorem stampacchia_1d {a b : ℝ} (hab : a < b) {u g : ℝ → ℝ}
    (hu : IntegrableOn u (Ioo a b) volume)
    (hg : IntegrableOn g (Ioo a b) volume)
    (hweak : ∀ φ : ℝ → ℝ, ContDiff ℝ (⊤ : ℕ∞) φ → HasCompactSupport φ →
      tsupport φ ⊆ Ioo a b →
      ∫ x in Ioo a b, u x * deriv φ x = -∫ x in Ioo a b, g x * φ x) :
    ∀ᵐ x ∂(volume.restrict (Ioo a b)), u x = 0 → g x = 0 := by
  obtain ⟨C, hC⟩ := w11_ae_eq_ac_representative hab hu hg hweak
  set F := fun x => C + ∫ t in a..x, g t
  -- We need global a.e. differentiability. Use g̃ = indicator (Ioo a b) g (globally integrable).
  -- F̃(x) = C + ∫_a^x g̃ agrees with F on (a,b) and is globally a.e. differentiable.
  set g' := (Ioo a b).indicator g
  have hg'_li : LocallyIntegrable g' volume :=
    ((integrable_indicator_iff measurableSet_Ioo).mpr hg).locallyIntegrable
  -- F̃ has derivative g̃ a.e. globally (by FTC for locally integrable functions)
  set F' := fun x => C + ∫ t in a..x, g' t
  have hF'_ae : ∀ᵐ x, HasDerivAt F' (g' x) x := by
    filter_upwards [LocallyIntegrable.ae_hasDerivAt_integral hg'_li] with x hx
    have := (hasDerivAt_const x C).add (hx a)
    simp only [zero_add] at this; exact this
  -- Apply deriv_eq_zero_ae_on_zeroSet to F'
  have hF'_deriv_ae : ∀ᵐ x, HasDerivAt F' (deriv F' x) x :=
    hF'_ae.mono fun x hx => by rwa [hx.deriv]
  have hF'_zero := deriv_eq_zero_ae_on_zeroSet hF'_deriv_ae
  -- On (a,b): g' = g, F' = F, so combine with hC and hF'_zero.
  have hint_eq : ∀ x ∈ Ioo a b,
      (∫ t in a..x, g' t) = ∫ t in a..x, g t := by
    intro x hx; apply intervalIntegral.integral_congr_ae; apply ae_of_all
    intro t ht; rw [uIoc_of_le hx.1.le] at ht
    exact indicator_of_mem (show t ∈ Ioo a b from ⟨ht.1, lt_of_le_of_lt ht.2 hx.2⟩) g
  filter_upwards [ae_restrict_of_ae (s := Ioo a b) hF'_ae,
                   ae_restrict_of_ae (s := Ioo a b) hF'_zero,
                   hC, ae_restrict_mem measurableSet_Ioo] with x hF'_d hF'_z hu_eq hx_mem
  intro hu_zero
  have hg'_eq : g' x = g x := indicator_of_mem hx_mem g
  have hF'_eq : F' x = C + ∫ t in a..x, g t := by
    change C + ∫ t in a..x, g' t = _; congr 1; exact hint_eq x hx_mem
  rw [← hg'_eq, ← hF'_d.deriv]
  exact hF'_z (by rw [hF'_eq]; show C + ∫ t in a..x, g t = 0; linarith [hu_eq])

\end{lstlisting}
\begin{proof}
Replace $u$ by its absolutely continuous representative $F(x) = C + \int_a^x g$. Then
$F' = g$ a.e.\ by the Lebesgue differentiation theorem for locally integrable
functions, and \leanname{deriv\_eq\_zero\_ae\_on\_zeroSet} applied to $F$ yields
$F' = 0$ a.e.\ on $\{F = 0\}$, i.e.\ $g = 0$ a.e.\ on $\{u = 0\}$.
\end{proof}

\begin{theorem}[$d$-dimensional Stampacchia, \leanname{weakGrad\_ae\_eq\_zero\_on\_zeroSet}]
Let $\Omega \subseteq E$ be open and $u \in \Wkp{1}{2}(\Omega)$ with weak gradient
$G : \Omega \to \R^{d}$. Then for each $i \in \{1,\dots,d\}$, $G_i(x) = 0$ for a.e.\
$x \in \Omega$ with $u(x) = 0$.
\end{theorem}

\begin{lstlisting}[style=Matlab-editor,basicstyle=\mlttfamily,escapechar=",mlshowsectionrules=true,mathescape=true,morecomment={[s]{\%\{}{\%\}}},language=lean,numbers=none]
theorem weakGrad_ae_eq_zero_on_zeroSet
    {Ω : Set E} (hΩ : IsOpen Ω) {u : E → ℝ}
    (hu : MemW1p 2 u Ω)
    {G : E → E} (hG : ∀ i, HasWeakPartialDeriv' i (fun x => G x i) u Ω) (i : Fin d)
    (hGi_mem : MemLp (fun x => G x i) 2 (volume.restrict Ω))
    (hu_meas : Measurable u) :
    ∀ᵐ x ∂(volume.restrict Ω), u x = 0 → G x i = 0 := by
  -- Define h(x) = indicator({u=0}) * G_i(x). We show h = 0 a.e. on Ω.
  set h : E → ℝ := fun x => if u x = 0 then G x i else 0 with hh_def
  -- It suffices to show h = 0 a.e. on Ω
  suffices hgoal : ∀ᵐ x ∂volume, x ∈ Ω → h x = 0 by
    rw [ae_restrict_iff' hΩ.measurableSet]
    filter_upwards [hgoal] with x hx hxΩ hux
    have := hx hxΩ
    simp only [hh_def, hux, ite_true] at this
    exact this
  -- Apply the fundamental lemma
  apply hΩ.ae_eq_zero_of_integral_contDiff_smul_eq_zero
  -- Sub-goal 1: h is locally integrable on Ω
  · exact weakGrad_locallyIntegrableOn hΩ hu hG i hGi_mem hu_meas
  -- Sub-goal 2: ∫ ψ • h = 0 for all smooth compactly supported ψ with supp ⊆ Ω
  · intro ψ hψ hψ_cs hψ_supp
    exact weakGrad_integral_test_eq_zero hΩ hu hu_meas hG i hGi_mem ψ hψ hψ_cs hψ_supp

\end{lstlisting}
\begin{proof}
Set $h_i(x) = \mathbf{1}_{\{u = 0\}}(x)\, G_i(x)$. We show $h_i = 0$ a.e.\ on $\Omega$
by reducing to the variational fundamental lemma: for any test function
$\psi \in C_c^{\infty}(\Omega)$ the integral $\int \psi\, h_i$ vanishes. By Fubini's
theorem, the line restrictions of $u$ along the $i$-th coordinate axis lie in
$W^{1,1}$ for a.e.\ line, with weak derivative the line restriction of $G_i$. The 1D
Stampacchia theorem applied on each line gives $G_i = 0$ a.e.\ on
$\{u = 0\}$ along almost every line. Recombining via Fubini gives the result.
\end{proof}

\subsection{BMO and the John--Nirenberg geometry}

\begin{definition}[BMO seminorm, \leanname{bmoNorm}]
For $u : E \to \R$ and $U \subseteq E$,
\[
\lVert u\rVert_{\mathrm{BMO}(U)}
:= \sup\Bigl\{\, \fint_{\Ball{r}{x_0}} \bigl|\, u - \tfrac{1}{|B|}\!\int_{\Ball{r}{x_0}} u\,\bigr|
\;:\; \overline{\Ball{r}{x_0}} \subseteq U,\ r > 0\,\Bigr\}.
\]
\end{definition}

\begin{lstlisting}[style=Matlab-editor,basicstyle=\mlttfamily,escapechar=",mlshowsectionrules=true,mathescape=true,morecomment={[s]{\%\{}{\%\}}},language=lean,numbers=none]
/-! ## BMO -/

/-- The BMO seminorm on a set `U`. -/
noncomputable def bmoNorm (u : E → ℝ) (U : Set E) : ℝ :=
  sSup {t : ℝ | ∃ (x₀ : E) (r : ℝ), 0 < r ∧
    Metric.closedBall x₀ r ⊆ U ∧
    t = (⨍ x in Metric.ball x₀ r, ‖u x - ⨍ y in Metric.ball x₀ r, u y ∂volume‖ ∂volume)}

\end{lstlisting}
\begin{definition}[John--Nirenberg constants, \leanname{C\_JN}, \leanname{C\_iter}]
\[
C_{\mathrm{JN}}(d) := 16 \cdot 36^{d}, \qquad C_{\mathrm{iter}} := 1024.
\]
Both are strictly positive (\leanname{C\_JN\_pos}, \leanname{C\_iter\_pos}).
\end{definition}

\begin{lstlisting}[style=Matlab-editor,basicstyle=\mlttfamily,escapechar=",mlshowsectionrules=true,mathescape=true,morecomment={[s]{\%\{}{\%\}}},language=lean,numbers=none]
/-! ## Constants -/

/-- The John-Nirenberg constant. -/
noncomputable def C_JN (d : ℕ) : ℝ :=
  16 * 36 ^ d

/-- The constant in the simple iteration lemma with exponent ξ = 2.
  Using δ = 1/4 gives C_iter = 4 · 144 = 576; we round up to 1024 for margin. -/
noncomputable def C_iter : ℝ := 1024

theorem C_JN_pos (d : ℕ) : 0 < C_JN d := by
  unfold C_JN
  positivity

theorem C_iter_pos : 0 < C_iter := by
  unfold C_iter; positivity

\end{lstlisting}
\begin{definition}[\leanname{JNBall}]
For $x_0 \in E$, $R > 0$, a \emph{JN-ball} is a record consisting of a centre
$c \in \Ball{R}{x_0}$ and a radius $0 < \rho \le R$. Its \emph{carrier} is
$\Ball{\rho}{c}$ and its \emph{five-fold enlargement} is $\Ball{5\rho}{c}$.
\end{definition}

\begin{lstlisting}[style=Matlab-editor,basicstyle=\mlttfamily,escapechar=",mlshowsectionrules=true,mathescape=true,morecomment={[s]{\%\{}{\%\}}},language=lean,numbers=none]
/-! ## John-Nirenberg Geometry -/

/-- Balls inside a fixed ambient ball, used in the Calderón-Zygmund selection behind the
John-Nirenberg argument. -/
structure JNBall (x₀ : E) (R : ℝ) where
  center : E
  radius : ℝ
  radius_pos : 0 < radius
  radius_le : radius ≤ R
  center_mem : center ∈ Metric.ball x₀ R

\end{lstlisting}
\begin{lemma}[Vitali $5r$-covering, \leanname{vitali\_covering\_lemma}]
Given a family of balls $(\Ball{r_i}{x_i})_{i \in \iota}$ with $r_i > 0$ uniformly bounded
above, there is a countable subfamily $S \subseteq \iota$ such that the carriers
$(\Ball{r_i}{x_i})_{i \in S}$ are pairwise disjoint, every original ball meets some
selected ball whose radius is at least half that of the original, and
$\bigcup_{i \in \iota} \Ball{r_i}{x_i} \subseteq \bigcup_{i \in S} \Ball{5 r_i}{x_i}$.
\end{lemma}

\begin{lstlisting}[style=Matlab-editor,basicstyle=\mlttfamily,escapechar=",mlshowsectionrules=true,mathescape=true,morecomment={[s]{\%\{}{\%\}}},language=lean,numbers=none]
/-! ## Vitali Covering Lemma -/

/-- Vitali covering lemma (5r-covering): from any family of balls, one can extract a
disjoint subfamily such that the 5× enlargements cover the union, and every original
ball meets a selected ball of at least half its radius. -/
theorem vitali_covering_lemma
    {ι : Type*} {x : ι → E} {r : ι → ℝ}
    (hr_pos : ∀ i, 0 < r i)
    (hr_bdd : BddAbove (Set.range r))
    (hfin : (⋃ i, Metric.ball (x i) (r i)).Nonempty) :
    ∃ S : Set ι, S.Countable ∧
      (∀ i ∈ S, ∀ j ∈ S, i ≠ j →
        Disjoint (Metric.ball (x i) (r i)) (Metric.ball (x j) (r j))) ∧
      (∀ i, ∃ j ∈ S,
        (Metric.ball (x i) (r i) ∩ Metric.ball (x j) (r j)).Nonempty ∧ r i ≤ 2 * r j) ∧
      (⋃ i, Metric.ball (x i) (r i)) ⊆ ⋃ i ∈ S, Metric.ball (x i) (5 * r i) := by
  classical
  let _ := hfin
  let t : Set ι := Set.univ
  obtain ⟨R, hR⟩ := hr_bdd
  rcases
      Vitali.exists_disjoint_subfamily_covering_enlargement
        (fun i => Metric.ball (x i) (r i)) t r 2 (by norm_num)
        (fun i _ => (hr_pos i).le) R (by
          intro i hi
          exact hR (Set.mem_range_self i)) (fun i _ =>
            nonempty_ball.2 (hr_pos i))
    with ⟨S, _, hS_disj, hS_hit⟩
  have hS_ball_disj : S.PairwiseDisjoint (fun i => Metric.ball (x i) (r i)) := by
    exact hS_disj
  have hS_count : S.Countable := by
    exact hS_ball_disj.countable_of_isOpen
      (fun i _ => isOpen_ball)
      (fun i _ => nonempty_ball.2 (hr_pos i))
  refine ⟨S, hS_count, ?_, ?_, ?_⟩
  · intro i hi j hj hij
    exact hS_ball_disj hi hj hij
  · intro i
    rcases hS_hit i (by simp [t]) with ⟨j, hjS, hij, hrad⟩
    exact ⟨j, hjS, hij, hrad⟩
  · intro y hy
    rcases Set.mem_iUnion.1 hy with ⟨i, hyi⟩
    rcases hS_hit i (by simp [t]) with ⟨j, hjS, hij, hrad⟩
    refine Set.mem_iUnion.2 ⟨j, Set.mem_iUnion.2 ⟨hjS, ?_⟩⟩
    rcases hij with ⟨z, hzi, hzj⟩
    have hyj : dist y (x j) < 5 * r j := by
      calc
        dist y (x j) ≤ dist y (x i) + dist (x i) z + dist z (x j) :=
          dist_triangle4 y (x i) z (x j)
        _ < r i + r i + r j := by
          have h1 : dist y (x i) < r i := by simpa [Metric.mem_ball] using hyi
          have h2 : dist (x i) z < r i := by rw [dist_comm]; simpa [Metric.mem_ball] using hzi
          have h3 : dist z (x j) < r j := by simpa [Metric.mem_ball] using hzj
          linarith
        _ ≤ 5 * r j := by linarith
    exact hyj


\end{lstlisting}
\begin{lemma}[Geometric absorption, \leanname{ball\_subset\_fivefold\_of\_inter\_nonempty}]
If $\Ball{\rho}{a} \subseteq \Ball{r}{c}$, $\Ball{r}{c} \cap \Ball{s}{d} \neq \emptyset$
and $r \le 2 s$, then $\Ball{\rho}{a} \subseteq \Ball{5 s}{d}$.
\end{lemma}

\begin{lstlisting}[style=Matlab-editor,basicstyle=\mlttfamily,escapechar=",mlshowsectionrules=true,mathescape=true,morecomment={[s]{\%\{}{\%\}}},language=lean,numbers=none]
lemma ball_subset_fivefold_of_inter_nonempty
    {a c d : E} {ρ r s : ℝ}
    (hsub : Metric.ball a ρ ⊆ Metric.ball c r)
    (hint : (Metric.ball c r ∩ Metric.ball d s).Nonempty)
    (hrad : r ≤ 2 * s) :
    Metric.ball a ρ ⊆ Metric.ball d (5 * s) := by
  intro y hy
  rcases hint with ⟨z, hzcr, hzds⟩
  have hycr : y ∈ Metric.ball c r := hsub hy
  have hzc : dist c z < r := by simpa [dist_comm] using hzcr
  have hyd : dist y d < 5 * s := by
    calc
      dist y d ≤ dist y c + dist c z + dist z d := by
        exact dist_triangle4 y c z d
      _ < r + r + s := by
        have hyc : dist y c < r := by simpa [Metric.mem_ball] using hycr
        have hzd : dist z d < s := by simpa [Metric.mem_ball] using hzds
        linarith
      _ ≤ 5 * s := by linarith
  exact hyd

\end{lstlisting}
\begin{lemma}[Iteration of one-step decay,
\leanname{john\_nirenberg\_iteration\_decay\_iterate},
\leanname{john\_nirenberg\_iteration}]
Let $E_{\ell} \subseteq B$ be an antitone family of measurable subsets of a measurable
set $B$ with $\mathrm{vol}(B) < \infty$. Suppose there exist $A > 0$ and
$0 < \theta < 1$ with
\[
\mathrm{vol}(E_{\ell + A}) \le \theta\, \mathrm{vol}(E_{\ell})
\quad\text{for all } \ell > 0.
\]
Then for every $\ell > 0$,
\[
\mathrm{vol}(E_{\ell})
\le \frac{1}{\theta}\, \mathrm{vol}(B)\,
\exp\!\Bigl(-\ell\, \frac{-\log \theta}{A}\Bigr).
\]
\end{lemma}

\begin{lstlisting}[style=Matlab-editor,basicstyle=\mlttfamily,escapechar=",mlshowsectionrules=true,mathescape=true,morecomment={[s]{\%\{}{\%\}}},language=lean,numbers=none]
/-- Iterating one-step level-set decay gives a geometric bound. -/
lemma john_nirenberg_iteration_decay_iterate
    {E_lam : ℝ → Set E} {A : ℝ} (hA : 0 < A) {θ : ℝ}
    (h_decay : ∀ lam : ℝ, 0 < lam → volume (E_lam (lam + A)) ≤ ENNReal.ofReal θ * volume (E_lam lam))
    (lam : ℝ) (hlam : 0 < lam) (n : ℕ) :
    volume (E_lam (lam + n * A)) ≤ ENNReal.ofReal θ ^ n * volume (E_lam lam) := by
  induction n with
  | zero =>
      simp
  | succ n ih =>
      have hlam_n_pos : 0 < lam + n * A := by positivity
      have h_step : volume (E_lam (lam + ↑(n + 1) * A)) ≤
          ENNReal.ofReal θ * volume (E_lam (lam + n * A)) := by
        have : lam + ↑(n + 1) * A = (lam + ↑n * A) + A := by
          push_cast
          ring
        rw [this]
        exact h_decay _ hlam_n_pos
      calc
        volume (E_lam (lam + ↑(n + 1) * A))
            ≤ ENNReal.ofReal θ * volume (E_lam (lam + ↑n * A)) := h_step
        _ ≤ ENNReal.ofReal θ * (ENNReal.ofReal θ ^ n * volume (E_lam lam)) := by
            gcongr
        _ = ENNReal.ofReal θ ^ (n + 1) * volume (E_lam lam) := by
            ring

/-- Iteration of level-set decay to exponential decay. -/
theorem john_nirenberg_iteration
    {B : Set E} {E_lam : ℝ → Set E}
    (_hB : MeasurableSet B) (_hB_fin : volume B ≠ ⊤)
    (hE_sub : ∀ lam, E_lam lam ⊆ B)
    (_hE_meas : ∀ lam, MeasurableSet (E_lam lam))
    (_hE_anti : ∀ lam₁ lam₂, lam₁ ≤ lam₂ → E_lam lam₂ ⊆ E_lam lam₁)
    {A : ℝ} (hA : 0 < A)
    {θ : ℝ} (hθ_pos : 0 < θ) (hθ_lt : θ < 1)
    (h_decay : ∀ lam : ℝ, 0 < lam → volume (E_lam (lam + A)) ≤ ENNReal.ofReal θ * volume (E_lam lam)) :
    ∀ lam : ℝ, 0 < lam →
      volume (E_lam lam) ≤ ENNReal.ofReal (1 / θ) * volume B *
        ENNReal.ofReal (Real.exp (-lam * (-Real.log θ / A))) := by
  intro lam hlam
  have hlamA_pos : 0 < lam / A := div_pos hlam hA
  have hceil_ge1 : 1 ≤ ⌈lam / A⌉₊ := Nat.ceil_pos.mpr hlamA_pos
  set n := ⌈lam / A⌉₊ - 1 with hn_def
  set lam₀ := lam - ↑n * A with hlam0_def
  have hlam0_pos : 0 < lam₀ := by
    have hn_lt : (↑n : ℝ) < lam / A := by
      rw [hn_def, Nat.cast_sub hceil_ge1]
      push_cast
      linarith [Nat.ceil_lt_add_one (le_of_lt hlamA_pos)]
    have := mul_lt_mul_of_pos_right hn_lt hA
    rw [div_mul_cancel₀] at this
    · linarith
    · exact ne_of_gt hA
  have hlam_eq : lam₀ + ↑n * A = lam := by
    ring
  have h_iter := john_nirenberg_iteration_decay_iterate hA h_decay lam₀ hlam0_pos n
  rw [hlam_eq] at h_iter
  have h_sub : volume (E_lam lam₀) ≤ volume B := measure_mono (hE_sub lam₀)
  have hn_ge : (↑n : ℝ) ≥ lam / A - 1 := by
    rw [hn_def, Nat.cast_sub hceil_ge1]
    push_cast
    linarith [Nat.le_ceil (lam / A)]
  calc
    volume (E_lam lam)
        ≤ ENNReal.ofReal θ ^ n * volume (E_lam lam₀) := h_iter
    _ ≤ ENNReal.ofReal θ ^ n * volume B := by
        gcongr
    _ = ENNReal.ofReal (θ ^ n) * volume B := by
        rw [john_nirenberg_iteration_ofReal_pow θ (le_of_lt hθ_pos) n]
    _ ≤ (ENNReal.ofReal (1 / θ) * ENNReal.ofReal (Real.exp (-lam * (-Real.log θ / A)))) * volume B := by
        gcongr
        exact john_nirenberg_iteration_rpow_bound θ lam A hθ_pos hθ_lt n hn_ge
    _ = ENNReal.ofReal (1 / θ) * volume B * ENNReal.ofReal (Real.exp (-lam * (-Real.log θ / A))) := by
        ring

\end{lstlisting}
\begin{proof}
Iterating one-step decay $n$ times gives
$\mathrm{vol}(E_{\ell_0 + n A}) \le \theta^{n}\,\mathrm{vol}(E_{\ell_0})$. Choose
$n = \lceil \ell/A\rceil - 1$, so $\ell_0 = \ell - n A \in (0, A]$, and bound
$\theta^{n} \le (1/\theta)\, \exp(-\ell\,(-\log\theta)/A)$ using
$n \ge \ell/A - 1$ and $\theta^{x} = \exp(x\log\theta)$.
\end{proof}

\begin{theorem}[John--Nirenberg from one-step decay, \leanname{john\_nirenberg\_iteration}]
Let $u : E \to \R$ be measurable and let $x_0, r$ with $r > 0$. If for some $A,\theta$
with $A > 0$, $0 < \theta < 1$ the sublevel sets of $|u - \fint_{\Ball{r}{x_0}} u|$
satisfy
\[
\mathrm{vol}\bigl\{ x \in \Ball{r}{x_0} : \lvert u - \tfrac{1}{|B|}\!\int_{\Ball{r}{x_0}} u\rvert > \ell + A\bigr\}
\le \theta\, \mathrm{vol}\bigl\{ x \in \Ball{r}{x_0} : \lvert u - \tfrac{1}{|B|}\!\int_{\Ball{r}{x_0}} u\rvert > \ell\bigr\}
\]
for all $\ell > 0$, then for every $t > 0$,
\[
\mathrm{vol}\bigl\{ x \in \Ball{r}{x_0} : \lvert u - \tfrac{1}{|B|}\!\int_{\Ball{r}{x_0}} u\rvert > t\bigr\}
\le \frac{1}{\theta}\, \mathrm{vol}(\Ball{r}{x_0})\,
\exp\!\Bigl(-t\,\frac{-\log\theta}{A}\Bigr).
\]
\end{theorem}

\begin{lstlisting}[style=Matlab-editor,basicstyle=\mlttfamily,escapechar=",mlshowsectionrules=true,mathescape=true,morecomment={[s]{\%\{}{\%\}}},language=lean,numbers=none]
/-- Iteration of level-set decay to exponential decay. -/
theorem john_nirenberg_iteration
    {B : Set E} {E_lam : ℝ → Set E}
    (_hB : MeasurableSet B) (_hB_fin : volume B ≠ ⊤)
    (hE_sub : ∀ lam, E_lam lam ⊆ B)
    (_hE_meas : ∀ lam, MeasurableSet (E_lam lam))
    (_hE_anti : ∀ lam₁ lam₂, lam₁ ≤ lam₂ → E_lam lam₂ ⊆ E_lam lam₁)
    {A : ℝ} (hA : 0 < A)
    {θ : ℝ} (hθ_pos : 0 < θ) (hθ_lt : θ < 1)
    (h_decay : ∀ lam : ℝ, 0 < lam → volume (E_lam (lam + A)) ≤ ENNReal.ofReal θ * volume (E_lam lam)) :
    ∀ lam : ℝ, 0 < lam →
      volume (E_lam lam) ≤ ENNReal.ofReal (1 / θ) * volume B *
        ENNReal.ofReal (Real.exp (-lam * (-Real.log θ / A))) := by
  intro lam hlam
  have hlamA_pos : 0 < lam / A := div_pos hlam hA
  have hceil_ge1 : 1 ≤ ⌈lam / A⌉₊ := Nat.ceil_pos.mpr hlamA_pos
  set n := ⌈lam / A⌉₊ - 1 with hn_def
  set lam₀ := lam - ↑n * A with hlam0_def
  have hlam0_pos : 0 < lam₀ := by
    have hn_lt : (↑n : ℝ) < lam / A := by
      rw [hn_def, Nat.cast_sub hceil_ge1]
      push_cast
      linarith [Nat.ceil_lt_add_one (le_of_lt hlamA_pos)]
    have := mul_lt_mul_of_pos_right hn_lt hA
    rw [div_mul_cancel₀] at this
    · linarith
    · exact ne_of_gt hA
  have hlam_eq : lam₀ + ↑n * A = lam := by
    ring
  have h_iter := john_nirenberg_iteration_decay_iterate hA h_decay lam₀ hlam0_pos n
  rw [hlam_eq] at h_iter
  have h_sub : volume (E_lam lam₀) ≤ volume B := measure_mono (hE_sub lam₀)
  have hn_ge : (↑n : ℝ) ≥ lam / A - 1 := by
    rw [hn_def, Nat.cast_sub hceil_ge1]
    push_cast
    linarith [Nat.le_ceil (lam / A)]
  calc
    volume (E_lam lam)
        ≤ ENNReal.ofReal θ ^ n * volume (E_lam lam₀) := h_iter
    _ ≤ ENNReal.ofReal θ ^ n * volume B := by
        gcongr
    _ = ENNReal.ofReal (θ ^ n) * volume B := by
        rw [john_nirenberg_iteration_ofReal_pow θ (le_of_lt hθ_pos) n]
    _ ≤ (ENNReal.ofReal (1 / θ) * ENNReal.ofReal (Real.exp (-lam * (-Real.log θ / A)))) * volume B := by
        gcongr
        exact john_nirenberg_iteration_rpow_bound θ lam A hθ_pos hθ_lt n hn_ge
    _ = ENNReal.ofReal (1 / θ) * volume B * ENNReal.ofReal (Real.exp (-lam * (-Real.log θ / A))) := by
        ring

\end{lstlisting}
\begin{lemma}[Indicator BMO, \leanname{indicator\_bmo\_bound}]
If $u$ has BMO seminorm at most $M$ on every closed sub-ball of $\Ball{r}{x_0}$, then
the same bound holds for the function $\mathbf{1}_{\Ball{r}{x_0}}\, u$ on those sub-balls,
since the average over $\Ball{s}{z} \subseteq \Ball{r}{x_0}$ coincides with that of $u$.
\end{lemma}

\begin{lstlisting}[style=Matlab-editor,basicstyle=\mlttfamily,escapechar=",mlshowsectionrules=true,mathescape=true,morecomment={[s]{\%\{}{\%\}}},language=lean,numbers=none]
/-- On balls staying inside `Metric.ball x₀ r`, the zero extension agrees with `u`. -/
theorem indicator_bmo_bound
    {u : E → ℝ} {x₀ : E} {r M : ℝ}
    (hu_bmo : ∀ (z : E) (s : ℝ), 0 < s → Metric.closedBall z s ⊆ Metric.ball x₀ r →
      (⨍ x in Metric.ball z s, ‖u x - ⨍ y in Metric.ball z s, u y ∂volume‖ ∂volume) ≤ M)
    (z : E) (s : ℝ) (hs : 0 < s) (hzs : Metric.closedBall z s ⊆ Metric.ball x₀ r) :
    (⨍ x in Metric.ball z s,
      ‖(Metric.ball x₀ r).indicator u x -
        ⨍ y in Metric.ball z s, (Metric.ball x₀ r).indicator u y ∂volume‖ ∂volume) ≤
      M := by
  have hball : Metric.ball z s ⊆ Metric.ball x₀ r := by
    exact Set.Subset.trans ball_subset_closedBall hzs
  have havg :
      (⨍ y in Metric.ball z s, (Metric.ball x₀ r).indicator u y ∂volume) =
        ⨍ y in Metric.ball z s, u y ∂volume := by
    apply MeasureTheory.setAverage_congr_fun isOpen_ball.measurableSet
    filter_upwards with y
    intro hy
    simp [hball hy]
  have hosc :
      (⨍ x in Metric.ball z s,
        ‖(Metric.ball x₀ r).indicator u x -
          ⨍ y in Metric.ball z s, (Metric.ball x₀ r).indicator u y ∂volume‖ ∂volume) =
        (⨍ x in Metric.ball z s, ‖u x - ⨍ y in Metric.ball z s, u y ∂volume‖ ∂volume) := by
    apply MeasureTheory.setAverage_congr_fun isOpen_ball.measurableSet
    filter_upwards with x
    intro hx
    simp [hball hx, havg]
  rw [hosc]
  exact hu_bmo z s hs hzs

\end{lstlisting}
\subsection{Simple iteration lemma}

\begin{lemma}[\leanname{simple\_iteration\_lemma}]
Let $A_{\mathrm{iter}} \ge 0$ and let $\rho : \R \to \R$ satisfy on $[1/2, 1)$:
\begin{enumerate}
\item $\rho \ge 0$;
\item $\sup_{1/2 \le t < 1}(1-t)^{2}\,\rho(t) < \infty$;
\item for all $1/2 \le s < t < 1$,
$\rho(s) \le \tfrac{1}{2}\rho(t) + A_{\mathrm{iter}}(t-s)^{-2}$.
\end{enumerate}
Then $\rho(1/2) \le C_{\mathrm{iter}}\, A_{\mathrm{iter}}$.
\end{lemma}

\begin{lstlisting}[style=Matlab-editor,basicstyle=\mlttfamily,escapechar=",mlshowsectionrules=true,mathescape=true,morecomment={[s]{\%\{}{\%\}}},language=lean,numbers=none]
/-! ## Simple Iteration Lemma (AKM, Lemma C.6) -/

/-- **Simple Iteration Lemma** (AKM, Appendix C, Lemma C.6, specialized to ξ = 2).

Suppose `ρ : ℝ → ℝ` satisfies:
1. `ρ ≥ 0` on `[1/2, 1)`,
2. `sup_{t ∈ [1/2,1)} (1-t)² ρ(t) < ∞` (finiteness of weighted supremum),
3. For all `1/2 ≤ s < t < 1`: `ρ(s) ≤ (1/2) ρ(t) + A_iter · (t - s)⁻²`

Then `ρ(1/2) ≤ C_iter · A_iter`.

**Proof sketch** (following AKM):
- Let `M = sup_{t ∈ [1/2,1)} (1-t)² ρ(t)`.
- Choose `δ = 1 - (2/3)^{1/2}` and set `t = 1 - (1-δ)(1-s)` in hypothesis (3).
- Multiply by `(1-s)²` to get: `(1-s)² ρ(s) ≤ (3/4) · (1-t)² ρ(t) + δ⁻² · A_iter`.
- Take supremum: `M ≤ (3/4) M + δ⁻² · A_iter`, hence `M ≤ 4 δ⁻² A_iter`.
- Evaluate at `s = 1/2`: `ρ(1/2) ≤ 2² · M ≤ 4 · 4 · δ⁻² · A_iter ≤ C_iter · A_iter`.
-/
theorem simple_iteration_lemma
    {ρ : ℝ → ℝ} {A_iter : ℝ}
    (hA_iter : 0 ≤ A_iter)
    (hρ_nonneg : ∀ t : ℝ, 1 / 2 ≤ t → t < 1 → 0 ≤ ρ t)
    (hρ_bdd : ∃ M : ℝ, ∀ t : ℝ, 1 / 2 ≤ t → t < 1 → (1 - t) ^ 2 * ρ t ≤ M)
    (hρ_contract : ∀ s t : ℝ, 1 / 2 ≤ s → s < t → t < 1 →
      ρ s ≤ 1 / 2 * ρ t + A_iter * (t - s) ⁻¹ ^ 2) :
    ρ (1 / 2) ≤ C_iter * A_iter := by
  /- Proof using δ = 1/4: for s ∈ [1/2,1), set t = s + (1-s)/4 = (3s+1)/4.
     Then t-s = (1-s)/4, 1-t = 3(1-s)/4, (1-s)/(1-t) = 4/3.
     Multiplying the contraction by (1-s)² gives
       (1-s)² ρ(s) ≤ (8/9)(1-t)² ρ(t) + 16 A_iter.
     Iterating: if M bounds (1-t)² ρ(t), then (8/9)M + 16 A_iter also does.
     After n iterations the bound is (8/9)^n M + 144(1-(8/9)^n) A_iter.
     As n → ∞ this tends to 144 A_iter. So (1/2)² ρ(1/2) ≤ 144 A_iter,
     hence ρ(1/2) ≤ 576 A_iter ≤ 1024 A_iter = C_iter A_iter. -/
  -- Extract M from hρ_bdd and make it nonneg
  obtain ⟨M₀, hM₀⟩ := hρ_bdd
  set M : ℝ := max M₀ 0 with hM_def
  have hM_nn : 0 ≤ M := le_max_right _ _
  have hM_bound : ∀ t : ℝ, 1 / 2 ≤ t → t < 1 → (1 - t) ^ 2 * ρ t ≤ M := by
    intro t ht1 ht2
    exact le_trans (hM₀ t ht1 ht2) (le_max_left _ _)
  -- Key one-step improvement: for any upper bound B ≥ 0 on (1-t)²ρ(t),
  -- (8/9)B + 16 A_iter is also an upper bound.
  have improvement : ∀ B : ℝ, 0 ≤ B →
      (∀ t : ℝ, 1 / 2 ≤ t → t < 1 → (1 - t) ^ 2 * ρ t ≤ B) →
      (∀ s : ℝ, 1 / 2 ≤ s → s < 1 → (1 - s) ^ 2 * ρ s ≤ 8 / 9 * B + 16 * A_iter) := by
    intro B hB_nn hB_bound s hs1 hs2
    -- Set t = (3s + 1) / 4
    set t := (3 * s + 1) / 4 with ht_def
    have ht_gt_s : s < t := by linarith
    have ht_lt_1 : t < 1 := by linarith
    have ht_ge : 1 / 2 ≤ t := by linarith
    have h1ms_pos : 0 < 1 - s := by linarith
    have h1mt_pos : 0 < 1 - t := by linarith
    have hts_pos : 0 < t - s := by linarith
    -- Compute key ratios
    have h1mt : 1 - t = 3 / 4 * (1 - s) := by rw [ht_def]; ring
    have hts : t - s = 1 / 4 * (1 - s) := by rw [ht_def]; ring
    -- Apply the contraction hypothesis
    have hcontr := hρ_contract s t hs1 ht_gt_s ht_lt_1
    -- Multiply both sides by (1-s)²
    have hρ_s_nn : 0 ≤ ρ s := hρ_nonneg s hs1 hs2
    have hρ_t_nn : 0 ≤ ρ t := hρ_nonneg t ht_ge ht_lt_1
    -- (1-s)² ρ(s) ≤ (1-s)² ((1/2) ρ(t) + A_iter (t-s)⁻²)
    have step1 : (1 - s) ^ 2 * ρ s ≤
        (1 - s) ^ 2 * (1 / 2 * ρ t + A_iter * (t - s)⁻¹ ^ 2) := by
      exact mul_le_mul_of_nonneg_left hcontr (sq_nonneg _)
    -- (1-s)² · (t-s)⁻¹² = 16
    have hprod1 : (1 - s) ^ 2 * ((t - s)⁻¹ ^ 2) = 16 := by
      rw [hts]
      have : 1 / 4 * (1 - s) ≠ 0 := by positivity
      field_simp
      ring
    -- Also (1-s)² · (1/2) ρ(t) = (1/2) ((1-s)/(1-t))² (1-t)² ρ(t)
    -- = (1/2)(4/3)²(1-t)² ρ(t) = (8/9)(1-t)² ρ(t)
    have hratio : (1 - s) ^ 2 * (1 / 2 * ρ t) = 8 / 9 * ((1 - t) ^ 2 * ρ t) := by
      rw [h1mt]; ring
    -- Combine
    calc (1 - s) ^ 2 * ρ s
        ≤ (1 - s) ^ 2 * (1 / 2 * ρ t + A_iter * (t - s)⁻¹ ^ 2) := step1
      _ = (1 - s) ^ 2 * (1 / 2 * ρ t) + A_iter * ((1 - s) ^ 2 * (t - s)⁻¹ ^ 2) := by ring
      _ = 8 / 9 * ((1 - t) ^ 2 * ρ t) + A_iter * 16 := by rw [hratio, hprod1]
      _ ≤ 8 / 9 * B + 16 * A_iter := by
          have := hB_bound t ht_ge ht_lt_1
          nlinarith
  -- Iterate the improvement n times: after n steps the bound is
  -- (8/9)^n M + 144(1 - (8/9)^n) A_iter
  -- Define the iterated bound
  let bound : ℕ → ℝ := fun n => (8 / 9) ^ n * M + (1 - (8 / 9) ^ n) * (144 * A_iter)
  have hbound_nn : ∀ n, 0 ≤ bound n := by
    intro n
    have h89 : (0 : ℝ) ≤ (8 / 9) ^ n := by positivity
    have h89_le1 : (8 / 9 : ℝ) ^ n ≤ 1 := pow_le_one₀ (by positivity) (by norm_num)
    nlinarith [mul_nonneg h89 hM_nn]
  -- Show that bound n is an upper bound for all n
  have hbound_is_bound : ∀ n, ∀ t : ℝ, 1 / 2 ≤ t → t < 1 →
      (1 - t) ^ 2 * ρ t ≤ bound n := by
    intro n
    induction n with
    | zero =>
      intro t ht1 ht2
      simp only [bound, pow_zero, one_mul, sub_self, zero_mul, add_zero]
      exact hM_bound t ht1 ht2
    | succ n ih =>
      intro s hs1 hs2
      have him := improvement (bound n) (hbound_nn n) ih s hs1 hs2
      -- Need: 8/9 * bound n + 16 * A_iter ≤ bound (n+1)
      -- bound (n+1) = (8/9)^(n+1) M + (1 - (8/9)^(n+1)) * 144 * A_iter
      -- 8/9 * bound n + 16 * A_iter
      --   = 8/9 * ((8/9)^n M + (1-(8/9)^n) * 144 A_iter) + 16 A_iter
      --   = (8/9)^(n+1) M + (8/9)(1-(8/9)^n) * 144 A_iter + 16 A_iter
      --   = (8/9)^(n+1) M + (8/9 - (8/9)^(n+1)) * 144 A_iter + 16 A_iter
      --   = (8/9)^(n+1) M + (8/9)*144 A_iter - (8/9)^(n+1)*144 A_iter + 16 A_iter
      --   = (8/9)^(n+1) M + 128 A_iter + 16 A_iter - (8/9)^(n+1) * 144 A_iter
      --   = (8/9)^(n+1) M + 144 A_iter - (8/9)^(n+1) * 144 A_iter
      --   = (8/9)^(n+1) M + (1 - (8/9)^(n+1)) * 144 A_iter = bound (n+1) ✓
      suffices 8 / 9 * bound n + 16 * A_iter = bound (n + 1) by linarith
      show 8 / 9 * ((8 / 9) ^ n * M + (1 - (8 / 9) ^ n) * (144 * A_iter)) + 16 * A_iter =
        (8 / 9) ^ (n + 1) * M + (1 - (8 / 9) ^ (n + 1)) * (144 * A_iter)
      rw [pow_succ]
      ring
  -- ρ(1/2) ≤ 4 * bound n for all n
  have hρ_le_4bound : ∀ n, ρ (1 / 2) ≤ 4 * bound n := by
    intro n
    have h12 : (1 : ℝ) / 2 ≤ 1 / 2 := le_refl _
    have h12_lt : (1 : ℝ) / 2 < 1 := by norm_num
    have := hbound_is_bound n (1 / 2) h12 h12_lt
    -- (1 - 1/2)² ρ(1/2) ≤ bound n, i.e., (1/4) ρ(1/2) ≤ bound n
    have h14 : (1 - 1 / 2 : ℝ) ^ 2 = 1 / 4 := by norm_num
    rw [h14] at this
    -- ρ(1/2) ≤ 4 * bound n
    have hρ12_nn : 0 ≤ ρ (1 / 2) := hρ_nonneg (1 / 2) h12 h12_lt
    linarith
  -- Now: bound n → 144 * A_iter as n → ∞ (since (8/9)^n → 0)
  -- We need: ρ(1/2) ≤ 576 * A_iter
  -- Use: for all ε > 0, exists n such that bound n ≤ 144 * A_iter + ε
  -- Then ρ(1/2) ≤ 4 * (144 * A_iter + ε) = 576 * A_iter + 4ε
  -- Since ε arbitrary, ρ(1/2) ≤ 576 * A_iter
  suffices h576 : ρ (1 / 2) ≤ 576 * A_iter by
    calc ρ (1 / 2) ≤ 576 * A_iter := h576
      _ ≤ C_iter * A_iter := by unfold C_iter; nlinarith [hA_iter]
  -- Prove via le_of_forall_pos_lt_add
  rw [show (576 : ℝ) * A_iter = 4 * (144 * A_iter) from by ring]
  apply le_of_forall_pos_lt_add
  intro ε hε
  -- Find n such that (8/9)^n * M < ε / 4
  have h89_lt : (8 : ℝ) / 9 < 1 := by norm_num
  have h89_nn : (0 : ℝ) ≤ 8 / 9 := by norm_num
  have hε4 : (0 : ℝ) < ε / 4 := by linarith
  have htend : Filter.Tendsto (fun n => (8 / 9 : ℝ) ^ n * M) Filter.atTop (nhds 0) := by
    have := (tendsto_pow_atTop_nhds_zero_of_lt_one h89_nn h89_lt).mul_const M
    rwa [zero_mul] at this
  rw [tendsto_atTop_nhds] at htend
  obtain ⟨N, hN⟩ := htend (Set.Iio (ε / 4)) (show (0 : ℝ) ∈ Set.Iio (ε / 4) from hε4)
    isOpen_Iio
  have hN' : (8 / 9) ^ N * M < ε / 4 := hN N (le_refl _)
  -- bound N ≤ (8/9)^N * M + 144 * A_iter
  have hboundN : bound N ≤ (8 / 9) ^ N * M + 144 * A_iter := by
    show (8 / 9) ^ N * M + (1 - (8 / 9) ^ N) * (144 * A_iter) ≤
      (8 / 9) ^ N * M + 144 * A_iter
    nlinarith [pow_nonneg h89_nn N, hA_iter]
  calc ρ (1 / 2)
      ≤ 4 * bound N := hρ_le_4bound N
    _ ≤ 4 * ((8 / 9) ^ N * M + 144 * A_iter) := by nlinarith
    _ < 4 * (ε / 4 + 144 * A_iter) := by nlinarith
    _ = 4 * (144 * A_iter) + ε := by ring

\end{lstlisting}
\begin{proof}
Let $M$ be a non-negative upper bound for $(1-t)^{2}\rho(t)$. With the choice
$t = (3s+1)/4$, one has $t-s = (1-s)/4$ and $1-t = 3(1-s)/4$, so multiplying the
contraction by $(1-s)^{2}$ and using $(1-s)^{2}\,(t-s)^{-2} = 16$ gives
\[
(1-s)^{2}\,\rho(s) \le \tfrac{8}{9}\,(1-t)^{2}\rho(t) + 16\, A_{\mathrm{iter}}
\le \tfrac{8}{9}\, M + 16\, A_{\mathrm{iter}}.
\]
Taking the supremum over $s$ shows $M \mapsto \tfrac{8}{9}M + 16 A_{\mathrm{iter}}$ also
bounds $(1-t)^{2}\rho(t)$. Iterating $n$ times gives the bound
$M_n = (8/9)^{n} M + (1-(8/9)^{n})\cdot 144\, A_{\mathrm{iter}}$, which tends to
$144\,A_{\mathrm{iter}}$. Evaluating at $s = 1/2$ where $(1-s)^{2} = 1/4$, we get
$\rho(1/2) \le 4 M_n \to 576\, A_{\mathrm{iter}} \le C_{\mathrm{iter}}\,A_{\mathrm{iter}}$.
\end{proof}

\subsection{Local John--Nirenberg theorem}

\begin{lemma}[Geometric decay step, \leanname{john\_nirenberg\_level\_set\_decay}]
Let $E_{\ell+A} \subseteq E_{\ell}$ be measurable sets with $\mathrm{vol}(E_{\ell}) < \infty$
and let $(B_i)$ be a countable family of pairwise disjoint JN-balls with carriers
$\bigcup_i B_i^{\mathrm{car}} \subseteq E_{\ell}$ and 5-fold enlargements covering
$E_{\ell}$. If
\[
\mathrm{vol}(E_{\ell+A} \cap B_i^{\mathrm{car}}) \le \tfrac{1}{2}\, \mathrm{vol}(B_i^{\mathrm{car}})
\quad \text{for every } i,
\]
then
\[
\mathrm{vol}(E_{\ell+A}) \le \Bigl(1 - \frac{1}{2 \cdot 5^{d}}\Bigr)\, \mathrm{vol}(E_{\ell}).
\]
\end{lemma}

\begin{lstlisting}[style=Matlab-editor,basicstyle=\mlttfamily,escapechar=",mlshowsectionrules=true,mathescape=true,morecomment={[s]{\%\{}{\%\}}},language=lean,numbers=none]
/-- Geometric decay of selected John-Nirenberg bad-ball unions from the local half-measure step. -/
theorem john_nirenberg_level_set_decay
    {x₀ : E} {r : ℝ} {E_lam E_lam_A : Set E}
    {ι : Type*} [Countable ι] (B : ι → JNBall x₀ r)
    (_hE_lam_meas : MeasurableSet E_lam)
    (hE_lam_A_meas : MeasurableSet E_lam_A)
    (hE_lam_fin : volume E_lam ≠ ⊤)
    (hE_lam_A_sub : E_lam_A ⊆ E_lam)
    (hcover : E_lam ⊆ ⋃ i, (B i).fivefold)
    (hU_sub_E_lam : (⋃ i, (B i).carrier) ⊆ E_lam)
    (hdisj : ∀ i j, i ≠ j → Disjoint ((B i).carrier) ((B j).carrier))
    (hmain : ∀ i,
      volume (E_lam_A ∩ (B i).carrier) ≤ ENNReal.ofReal (1 / 2) * volume ((B i).carrier)) :
    volume.real E_lam_A ≤ (1 - 1 / (2 * 5 ^ d)) * volume.real E_lam := by
  -- Finiteness of derived sets
  have hE_lam_A_fin : volume E_lam_A ≠ ⊤ := measure_ne_top_of_subset hE_lam_A_sub hE_lam_fin
  -- Set up U = ⋃ i, (B i).carrier
  set U := ⋃ i, (B i).carrier with hU_def
  have hU_meas : MeasurableSet U :=
    MeasurableSet.iUnion (fun i => (B i).measurableSet_carrier)
  have hU_fin : volume U ≠ ⊤ := measure_ne_top_of_subset hU_sub_E_lam hE_lam_fin
  -- Carrier finiteness
  have hcarrier_fin : ∀ i, volume ((B i).carrier) ≠ ⊤ :=
    fun i => measure_ne_top_of_subset (subset_iUnion (fun i => (B i).carrier) i) hU_fin
  -- Pairwise disjointness in Pairwise form
  have hdisj_pw : Pairwise (Function.onFun Disjoint fun i => (B i).carrier) :=
    fun i j hij => hdisj i j hij
  have hvolU : volume U = ∑' i, volume ((B i).carrier) :=
    measure_iUnion hdisj_pw (fun i => (B i).measurableSet_carrier)
  have hE_lam_le : volume.real E_lam ≤ (5 : ℝ) ^ d * volume.real U := by
    have h1 : volume E_lam ≤ volume (⋃ i, (B i).fivefold) := measure_mono hcover
    have h2 : volume (⋃ i, (B i).fivefold) ≤ ∑' i, volume ((B i).fivefold) :=
      measure_iUnion_le _
    have h3 : ∀ i, volume ((B i).fivefold) = ENNReal.ofReal ((5 : ℝ) ^ d) * volume ((B i).carrier) :=
      fun i => (B i).volume_fivefold
    have h4 : ∑' i, volume ((B i).fivefold) = ∑' i, (ENNReal.ofReal ((5 : ℝ) ^ d) * volume ((B i).carrier)) :=
      tsum_congr (fun i => h3 i)
    have h5 : ∑' i, (ENNReal.ofReal ((5 : ℝ) ^ d) * volume ((B i).carrier)) =
        ENNReal.ofReal ((5 : ℝ) ^ d) * ∑' i, volume ((B i).carrier) :=
      ENNReal.tsum_mul_left
    rw [h4, h5, ← hvolU] at h2
    have hfive_ne_top : ENNReal.ofReal ((5 : ℝ) ^ d) * volume U ≠ ⊤ :=
      ENNReal.mul_ne_top ENNReal.ofReal_ne_top hU_fin
    have hle := ENNReal.toReal_mono hfive_ne_top (h1.trans h2)
    rwa [ENNReal.toReal_mul, ENNReal.toReal_ofReal (by positivity : (0:ℝ) ≤ (5:ℝ) ^ d)] at hle
  have hA_inter_U : volume.real (E_lam_A ∩ U) ≤ 1 / 2 * volume.real U := by
    -- Work in ENNReal: volume (E_lam_A ∩ U) ≤ ENNReal.ofReal (1/2) * volume U
    have hinter_eq : E_lam_A ∩ U = ⋃ i, (E_lam_A ∩ (B i).carrier) := by
      simp [hU_def, inter_iUnion]
    have hdisj_inter : Pairwise (Function.onFun Disjoint fun i => E_lam_A ∩ (B i).carrier) := by
      intro i j hij
      exact Disjoint.mono inter_subset_right inter_subset_right (hdisj_pw hij)
    have hmeas_inter : ∀ i, MeasurableSet (E_lam_A ∩ (B i).carrier) :=
      fun i => hE_lam_A_meas.inter (B i).measurableSet_carrier
    have hvol_inter : volume (E_lam_A ∩ U) = ∑' i, volume (E_lam_A ∩ (B i).carrier) := by
      rw [hinter_eq, measure_iUnion hdisj_inter hmeas_inter]
    have henn_le : volume (E_lam_A ∩ U) ≤ ENNReal.ofReal (1 / 2) * volume U := by
      rw [hvol_inter, hvolU]
      calc ∑' i, volume (E_lam_A ∩ (B i).carrier)
          ≤ ∑' i, (ENNReal.ofReal (1 / 2) * volume ((B i).carrier)) :=
            ENNReal.tsum_le_tsum (fun i => hmain i)
        _ = ENNReal.ofReal (1 / 2) * ∑' i, volume ((B i).carrier) :=
            ENNReal.tsum_mul_left
    have hfin : ENNReal.ofReal (1 / 2) * volume U ≠ ⊤ :=
      ENNReal.mul_ne_top ENNReal.ofReal_ne_top hU_fin
    have := ENNReal.toReal_mono hfin henn_le
    rwa [ENNReal.toReal_mul, ENNReal.toReal_ofReal (by norm_num : (0:ℝ) ≤ 1/2)] at this
  have hAdiffU_sub : E_lam_A \ U ⊆ E_lam \ U :=
    diff_subset_diff_left hE_lam_A_sub
  -- volume.real E_lam_A = volume.real (E_lam_A ∩ U) + volume.real (E_lam_A \ U)
  have hE_lam_A_split : volume.real E_lam_A =
      volume.real (E_lam_A ∩ U) + volume.real (E_lam_A \ U) := by
    rw [← measureReal_inter_add_diff hU_meas hE_lam_A_fin]
  -- volume.real (E_lam \ U) = volume.real E_lam - volume.real U
  have hE_lam_diff : volume.real (E_lam \ U) = volume.real E_lam - volume.real U :=
    measureReal_diff hU_sub_E_lam hU_meas hE_lam_fin
  -- volume.real (E_lam_A \ U) ≤ volume.real (E_lam \ U)
  have hAdiffU_le : volume.real (E_lam_A \ U) ≤ volume.real (E_lam \ U) := by
    exact measureReal_mono hAdiffU_sub
      (ne_top_of_le_ne_top hE_lam_fin (measure_mono diff_subset))
  have hU_lower : 1 / (5 : ℝ) ^ d * volume.real E_lam ≤ volume.real U := by
    by_cases h5d : (5 : ℝ) ^ d = 0
    · simp [h5d]
    · rw [div_mul_eq_mul_div, one_mul]
      exact div_le_of_le_mul₀ (by positivity) measureReal_nonneg (by linarith [hE_lam_le])
  -- Now combine: volume.real E_lam_A ≤ 1/2 * volume.real U + (volume.real E_lam - volume.real U)
  --                                    = volume.real E_lam - 1/2 * volume.real U
  --                                    ≤ volume.real E_lam - 1/2 * 1/5^d * volume.real E_lam
  --                                    = (1 - 1/(2*5^d)) * volume.real E_lam
  have hU_real_le : volume.real U ≤ volume.real E_lam :=
    measureReal_mono hU_sub_E_lam hE_lam_fin
  calc volume.real E_lam_A
      = volume.real (E_lam_A ∩ U) + volume.real (E_lam_A \ U) := hE_lam_A_split
    _ ≤ 1 / 2 * volume.real U + volume.real (E_lam \ U) := by
        linarith [hA_inter_U, hAdiffU_le]
    _ = 1 / 2 * volume.real U + (volume.real E_lam - volume.real U) := by
        rw [hE_lam_diff]
    _ = volume.real E_lam - 1 / 2 * volume.real U := by ring
    _ ≤ volume.real E_lam - 1 / 2 * (1 / (5 : ℝ) ^ d * volume.real E_lam) := by
        linarith [hU_lower]
    _ = (1 - 1 / (2 * 5 ^ d)) * volume.real E_lam := by ring

\end{lstlisting}
\begin{proof}
Let $U = \bigcup_i B_i^{\mathrm{car}}$. By disjointness and the 5-fold cover,
$\mathrm{vol}(U) \ge 5^{-d}\, \mathrm{vol}(E_{\ell})$. The hypothesis on each carrier
sums to $\mathrm{vol}(E_{\ell+A}\cap U) \le \tfrac{1}{2}\,\mathrm{vol}(U)$. Splitting
$E_{\ell+A} = (E_{\ell+A}\cap U) \cup (E_{\ell+A}\setminus U)$ and using
$E_{\ell+A}\setminus U \subseteq E_{\ell}\setminus U$ and the lower bound on
$\mathrm{vol}(U)$ gives the stated estimate.
\end{proof}

\begin{theorem}[Iterated decay from a base level, \leanname{john\_nirenberg\_from\_base},
\leanname{level\_set\_family\_from\_base}]
If a measurable antitone family $E_{\ell} \subseteq B$ satisfies the one-step decay
$\mathrm{vol}(E_{\ell+A}) \le \theta\,\mathrm{vol}(E_{\ell})$ only for $\ell \ge A$, then
for every $t > 0$,
\[
\mathrm{vol}(E_{t}) \le \frac{1}{\theta^{2}}\, \mathrm{vol}(B)\,
\exp\!\Bigl(-t\,\frac{-\log\theta}{A}\Bigr).
\]
\end{theorem}

\begin{lstlisting}[style=Matlab-editor,basicstyle=\mlttfamily,escapechar=",mlshowsectionrules=true,mathescape=true,morecomment={[s]{\%\{}{\%\}}},language=lean,numbers=none]
/-- **John-Nirenberg iteration from a base level.**

This is a variant of `john_nirenberg` (and `john_nirenberg_iteration`) where the
one-step decay hypothesis `h_decay` is only assumed for `lam ≥ A` (instead of for
all `lam > 0`). The exponential decay conclusion holds for all `t > 0`, with a
slightly larger constant prefactor `1/θ²` to absorb the base case `t ≤ A`.

**Proof sketch (pure math)**:
- For `t > A`: decompose `t = A + k·A + r` with `0 < r ≤ A`. Iterate h_decay
  `k` times from level `A + r` to get `vol(E_t) ≤ θ^k · vol(E_{A+r}) ≤ θ^k · vol(B)`.
  The exponential bound follows from `θ^k ≤ (1/θ²) · exp(-t · (-log θ / A))`.
- For `t ≤ A`: `vol(E_t) ≤ vol(B)` trivially, and the RHS is at least
  `(1/θ²) · vol(B) · exp(-A · (-log θ / A)) = (1/θ²) · vol(B) · θ ≥ vol(B)/θ ≥ vol(B)`. -/
theorem john_nirenberg_from_base
    {u : E → ℝ} {x₀ : E} {r : ℝ} {A θ : ℝ}
    (_hr : 0 < r) (hu_meas : Measurable u)
    (hA : 0 < A) (hθ_pos : 0 < θ) (hθ_lt : θ < 1)
    (h_decay : ∀ lam : ℝ, A ≤ lam →
      volume ({x ∈ Metric.ball x₀ r |
        ‖u x - ⨍ y in Metric.ball x₀ r, u y ∂volume‖ > lam + A}) ≤
        ENNReal.ofReal θ *
          volume ({x ∈ Metric.ball x₀ r |
            ‖u x - ⨍ y in Metric.ball x₀ r, u y ∂volume‖ > lam}))
    (t : ℝ) (_ht : 0 < t) :
    volume ({x ∈ Metric.ball x₀ r |
      ‖u x - ⨍ y in Metric.ball x₀ r, u y ∂volume‖ > t}) ≤
    ENNReal.ofReal (1 / θ ^ 2) * volume (Metric.ball x₀ r) *
      ENNReal.ofReal (Real.exp (-t * (-Real.log θ / A))) := by
  -- This is a purely mathematical iteration argument (no CZ covering involved).
  -- The proof decomposes t into base + iterated steps, applies h_decay iteratively,
  -- and bounds the resulting geometric series by the exponential.
  set avg := ⨍ y in Metric.ball x₀ r, u y ∂volume
  set B := Metric.ball x₀ r with hB_def
  set F : ℝ → Set E := fun s => {x ∈ B | ‖u x - avg‖ > s + A} with hF_def
  have hB_meas : MeasurableSet B := by
    simpa [hB_def] using (measurableSet_ball : MeasurableSet (Metric.ball x₀ r))
  have hF_sub : ∀ s, F s ⊆ B := by
    intro s x hx
    exact hx.1
  have hF_meas : ∀ s, MeasurableSet (F s) := by
    intro s
    refine hB_meas.inter ?_
    exact measurableSet_lt measurable_const ((hu_meas.sub_const avg).norm)
  have hF_anti : ∀ s₁ s₂, s₁ ≤ s₂ → F s₂ ⊆ F s₁ := by
    intro s₁ s₂ hs x hx
    refine ⟨hx.1, ?_⟩
    have hs' : s₁ + A ≤ s₂ + A := by linarith
    exact lt_of_le_of_lt hs' hx.2
  have h_decayF : ∀ s : ℝ, 0 < s →
      volume (F (s + A)) ≤ ENNReal.ofReal θ * volume (F s) := by
    intro s hs
    have hsA : A ≤ s + A := by linarith
    simpa [F, add_assoc, avg, B] using h_decay (s + A) hsA
  have h_iter :=
    john_nirenberg_iteration hB_meas
      (measure_ball_lt_top (μ := volume) (x := x₀) (r := r)).ne
      hF_sub hF_meas hF_anti hA hθ_pos hθ_lt h_decayF
  have hAdiv : A * (-Real.log θ / A) = -Real.log θ := by
    field_simp [hA.ne']
  have hexpA : Real.exp (-A * (-Real.log θ / A)) = θ := by
    have hexp_arg : -A * (-Real.log θ / A) = Real.log θ := by
      calc
        -A * (-Real.log θ / A) = -(A * (-Real.log θ / A)) := by ring
        _ = -(-Real.log θ) := by rw [hAdiv]
        _ = Real.log θ := by ring
    rw [hexp_arg, Real.exp_log hθ_pos]
  by_cases htA : A < t
  · have hs_pos : 0 < t - A := by linarith
    have hmain :
        volume ({x ∈ Metric.ball x₀ r | ‖u x - avg‖ > t}) ≤
          ENNReal.ofReal (1 / θ) * volume (Metric.ball x₀ r) *
            ENNReal.ofReal (Real.exp (-(t - A) * (-Real.log θ / A))) := by
      simpa [F, hB_def, avg, sub_eq_add_neg, add_assoc, add_left_comm, add_comm] using
        h_iter (t - A) hs_pos
    have hexp_shift :
        Real.exp (-(t - A) * (-Real.log θ / A)) =
          (1 / θ) * Real.exp (-t * (-Real.log θ / A)) := by
      have h_inv : (1 / θ : ℝ) = Real.exp (-Real.log θ) := by
        rw [Real.exp_neg, Real.exp_log hθ_pos, one_div]
      rw [h_inv, ← Real.exp_add]
      congr 1
      calc
        -(t - A) * (-Real.log θ / A)
            = -t * (-Real.log θ / A) + A * (-Real.log θ / A) := by ring
        _ = -t * (-Real.log θ / A) + (-Real.log θ) := by rw [hAdiv]
        _ = -Real.log θ + -t * (-Real.log θ / A) := by ring
    calc
      volume ({x ∈ Metric.ball x₀ r | ‖u x - avg‖ > t})
          ≤ ENNReal.ofReal (1 / θ) * volume (Metric.ball x₀ r) *
              ENNReal.ofReal (Real.exp (-(t - A) * (-Real.log θ / A))) := hmain
      _ = ENNReal.ofReal (1 / θ ^ 2) * volume (Metric.ball x₀ r) *
            ENNReal.ofReal (Real.exp (-t * (-Real.log θ / A))) := by
        rw [hexp_shift, ENNReal.ofReal_mul (by positivity)]
        calc
          ENNReal.ofReal (1 / θ) * volume (Metric.ball x₀ r) *
              (ENNReal.ofReal (1 / θ) * ENNReal.ofReal (Real.exp (-t * (-Real.log θ / A))))
              =
              (ENNReal.ofReal (1 / θ) * ENNReal.ofReal (1 / θ)) * volume (Metric.ball x₀ r) *
                ENNReal.ofReal (Real.exp (-t * (-Real.log θ / A))) := by
                ring
          _ = ENNReal.ofReal ((1 / θ) * (1 / θ)) * volume (Metric.ball x₀ r) *
                ENNReal.ofReal (Real.exp (-t * (-Real.log θ / A))) := by
                rw [← ENNReal.ofReal_mul (by positivity)]
          _ = ENNReal.ofReal (1 / θ ^ 2) * volume (Metric.ball x₀ r) *
                ENNReal.ofReal (Real.exp (-t * (-Real.log θ / A))) := by
                congr 1
                field_simp [pow_two, hθ_pos.ne']
  · have ht_le_A : t ≤ A := le_of_not_gt htA
    have hsub :
        {x ∈ Metric.ball x₀ r | ‖u x - avg‖ > t} ⊆ Metric.ball x₀ r := by
      intro x hx
      exact hx.1
    have hrate_pos : 0 < -Real.log θ / A := by
      have hlog_neg : Real.log θ < 0 := Real.log_neg hθ_pos hθ_lt
      exact div_pos (by linarith) hA
    have hexp_lower : θ ≤ Real.exp (-t * (-Real.log θ / A)) := by
      have hArg : -A * (-Real.log θ / A) ≤ -t * (-Real.log θ / A) := by
        nlinarith [ht_le_A, le_of_lt hrate_pos]
      calc
        θ = Real.exp (-A * (-Real.log θ / A)) := by rw [hexpA]
        _ ≤ Real.exp (-t * (-Real.log θ / A)) := Real.exp_le_exp.2 hArg
    have hmul : 1 ≤ (1 / θ ^ 2) * θ := by
      have hθ_ne : θ ≠ 0 := hθ_pos.ne'
      rw [pow_two]
      field_simp [hθ_ne]
      linarith
    have hcoeff : 1 ≤ (1 / θ ^ 2) * Real.exp (-t * (-Real.log θ / A)) := by
      calc
        1 ≤ (1 / θ ^ 2) * θ := hmul
        _ ≤ (1 / θ ^ 2) * Real.exp (-t * (-Real.log θ / A)) := by
            gcongr
    calc
      volume ({x ∈ Metric.ball x₀ r | ‖u x - avg‖ > t})
          ≤ volume (Metric.ball x₀ r) := measure_mono hsub
      _ = ENNReal.ofReal 1 * volume (Metric.ball x₀ r) := by simp
      _ ≤ ENNReal.ofReal ((1 / θ ^ 2) * Real.exp (-t * (-Real.log θ / A))) *
            volume (Metric.ball x₀ r) := by
            have hcoeff' :
                ENNReal.ofReal (1 : ℝ) ≤
                  ENNReal.ofReal ((1 / θ ^ 2) * Real.exp (-t * (-Real.log θ / A))) :=
              ENNReal.ofReal_le_ofReal hcoeff
            exact mul_le_mul_of_nonneg_right hcoeff' (by positivity)
      _ = ENNReal.ofReal (1 / θ ^ 2) * volume (Metric.ball x₀ r) *
            ENNReal.ofReal (Real.exp (-t * (-Real.log θ / A))) := by
            rw [ENNReal.ofReal_mul (by positivity)]
            ring

/-- A set-family John-Nirenberg tail decay theorem from a base-level one-step
decay hypothesis.

This is the reusable set-family analogue of `john_nirenberg_from_base`. It
iterates one-step decay starting from the base level `A` for an arbitrary
antitone measurable family `E_lam`, and is useful when the Calderon-Zygmund /
stopping-time work has already been packaged into such a family and one wants the final
exponential tail bound without rebuilding the iteration argument. -/
theorem level_set_family_from_base
    {B : Set E} {E_lam : ℝ → Set E} {A θ : ℝ}
    (hB_meas : MeasurableSet B) (hB_fin : volume B ≠ ⊤)
    (hE_sub : ∀ lam, E_lam lam ⊆ B)
    (hE_meas : ∀ lam, MeasurableSet (E_lam lam))
    (hE_anti : ∀ lam₁ lam₂, lam₁ ≤ lam₂ → E_lam lam₂ ⊆ E_lam lam₁)
    (hA : 0 < A) (hθ_pos : 0 < θ) (hθ_lt : θ < 1)
    (h_decay : ∀ lam : ℝ, A ≤ lam →
      volume (E_lam (lam + A)) ≤ ENNReal.ofReal θ * volume (E_lam lam))
    (t : ℝ) (_ht : 0 < t) :
    volume (E_lam t) ≤
      ENNReal.ofReal (1 / θ ^ 2) * volume B *
        ENNReal.ofReal (Real.exp (-t * (-Real.log θ / A))) := by
  have h_decayF : ∀ s : ℝ, 0 < s →
      volume (E_lam (s + A + A)) ≤ ENNReal.ofReal θ * volume (E_lam (s + A)) := by
    intro s hs
    have hsA : A ≤ s + A := by linarith
    simpa [add_assoc] using h_decay (s + A) hsA
  set F : ℝ → Set E := fun s => E_lam (s + A) with hF_def
  have hF_sub : ∀ s, F s ⊆ B := by
    intro s
    simpa [F] using hE_sub (s + A)
  have hF_meas : ∀ s, MeasurableSet (F s) := by
    intro s
    simpa [F] using hE_meas (s + A)
  have hF_anti : ∀ s₁ s₂, s₁ ≤ s₂ → F s₂ ⊆ F s₁ := by
    intro s₁ s₂ hs
    simpa [F] using hE_anti (s₁ + A) (s₂ + A) (by linarith)
  have h_iter :=
    john_nirenberg_iteration hB_meas hB_fin hF_sub hF_meas hF_anti hA hθ_pos hθ_lt
      (by
        intro s hs
        simpa [F, add_assoc] using h_decayF s hs)
  have hAdiv : A * (-Real.log θ / A) = -Real.log θ := by
    field_simp [hA.ne']
  have hexpA : Real.exp (-A * (-Real.log θ / A)) = θ := by
    have hexp_arg : -A * (-Real.log θ / A) = Real.log θ := by
      calc
        -A * (-Real.log θ / A) = -(A * (-Real.log θ / A)) := by ring
        _ = -(-Real.log θ) := by rw [hAdiv]
        _ = Real.log θ := by ring
    rw [hexp_arg, Real.exp_log hθ_pos]
  by_cases htA : A < t
  · have hs_pos : 0 < t - A := by linarith
    have hmain :
        volume (E_lam t) ≤
          ENNReal.ofReal (1 / θ) * volume B *
            ENNReal.ofReal (Real.exp (-(t - A) * (-Real.log θ / A))) := by
      simpa [F, sub_eq_add_neg, add_assoc, add_left_comm, add_comm] using
        h_iter (t - A) hs_pos
    have hexp_shift :
        Real.exp (-(t - A) * (-Real.log θ / A)) =
          (1 / θ) * Real.exp (-t * (-Real.log θ / A)) := by
      have h_inv : (1 / θ : ℝ) = Real.exp (-Real.log θ) := by
        rw [Real.exp_neg, Real.exp_log hθ_pos, one_div]
      rw [h_inv, ← Real.exp_add]
      congr 1
      calc
        -(t - A) * (-Real.log θ / A)
            = -t * (-Real.log θ / A) + A * (-Real.log θ / A) := by ring
        _ = -t * (-Real.log θ / A) + (-Real.log θ) := by rw [hAdiv]
        _ = -Real.log θ + -t * (-Real.log θ / A) := by ring
    calc
      volume (E_lam t)
          ≤ ENNReal.ofReal (1 / θ) * volume B *
              ENNReal.ofReal (Real.exp (-(t - A) * (-Real.log θ / A))) := hmain
      _ = ENNReal.ofReal (1 / θ ^ 2) * volume B *
            ENNReal.ofReal (Real.exp (-t * (-Real.log θ / A))) := by
        rw [hexp_shift, ENNReal.ofReal_mul (by positivity)]
        calc
          ENNReal.ofReal (1 / θ) * volume B *
              (ENNReal.ofReal (1 / θ) * ENNReal.ofReal (Real.exp (-t * (-Real.log θ / A))))
              =
              (ENNReal.ofReal (1 / θ) * ENNReal.ofReal (1 / θ)) * volume B *
                ENNReal.ofReal (Real.exp (-t * (-Real.log θ / A))) := by
                ring
          _ = ENNReal.ofReal ((1 / θ) * (1 / θ)) * volume B *
                ENNReal.ofReal (Real.exp (-t * (-Real.log θ / A))) := by
                rw [← ENNReal.ofReal_mul (by positivity)]
          _ = ENNReal.ofReal (1 / θ ^ 2) * volume B *
                ENNReal.ofReal (Real.exp (-t * (-Real.log θ / A))) := by
                congr 1
                field_simp [pow_two, hθ_pos.ne']
  · have ht_le_A : t ≤ A := le_of_not_gt htA
    have hsub : E_lam t ⊆ B := hE_sub t
    have hrate_pos : 0 < -Real.log θ / A := by
      have hlog_neg : Real.log θ < 0 := Real.log_neg hθ_pos hθ_lt
      exact div_pos (by linarith) hA
    have hexp_lower : θ ≤ Real.exp (-t * (-Real.log θ / A)) := by
      have hArg : -A * (-Real.log θ / A) ≤ -t * (-Real.log θ / A) := by
        nlinarith [ht_le_A, le_of_lt hrate_pos]
      calc
        θ = Real.exp (-A * (-Real.log θ / A)) := by rw [hexpA]
        _ ≤ Real.exp (-t * (-Real.log θ / A)) := Real.exp_le_exp.2 hArg
    have hmul : 1 ≤ (1 / θ ^ 2) * θ := by
      have hθ_ne : θ ≠ 0 := hθ_pos.ne'
      rw [pow_two]
      field_simp [hθ_ne]
      linarith
    have hcoeff : 1 ≤ (1 / θ ^ 2) * Real.exp (-t * (-Real.log θ / A)) := by
      calc
        1 ≤ (1 / θ ^ 2) * θ := hmul
        _ ≤ (1 / θ ^ 2) * Real.exp (-t * (-Real.log θ / A)) := by
            gcongr
    calc
      volume (E_lam t) ≤ volume B := measure_mono hsub
      _ = ENNReal.ofReal 1 * volume B := by simp
      _ ≤ ENNReal.ofReal ((1 / θ ^ 2) * Real.exp (-t * (-Real.log θ / A))) * volume B := by
          have hcoeff' :
              ENNReal.ofReal (1 : ℝ) ≤
                ENNReal.ofReal ((1 / θ ^ 2) * Real.exp (-t * (-Real.log θ / A))) :=
            ENNReal.ofReal_le_ofReal hcoeff
          exact mul_le_mul_of_nonneg_right hcoeff' (by positivity)
      _ = ENNReal.ofReal (1 / θ ^ 2) * volume B *
            ENNReal.ofReal (Real.exp (-t * (-Real.log θ / A))) := by
          rw [ENNReal.ofReal_mul (by positivity)]
          ring

\end{lstlisting}
\begin{proof}
For $t > A$ apply \leanname{john\_nirenberg\_iteration} to the shifted family
$F_s = E_{s+A}$, picking up an extra factor $1/\theta$ from the shift; the bound
$\theta^{n} \le (1/\theta) \exp(\dots)$ contributes another $1/\theta$, totalling
$1/\theta^{2}$. For $t \le A$ the trivial inclusion
$E_{t} \subseteq B$ already lies below the right-hand side because the exponential factor
is bounded below by $\theta$.
\end{proof}

\begin{theorem}[Local John--Nirenberg, \leanname{john\_nirenberg\_local}]
Let $R > 0$, $M > 0$, and $u$ measurable on $\Ball{6R}{x_0}$ with the BMO bound
\[
\fint_{\Ball{s}{z}} \Bigl|\, u - \fint_{\Ball{s}{z}} u\, \Bigr|\, dy \le M
\]
for every closed sub-ball $\overline{\Ball{s}{z}} \subseteq \Ball{6R}{x_0}$. Then for every
$t > 0$,
\[
\mathrm{vol}\Bigl\{ x \in \Ball{R}{x_0} : \Bigl|\,u(x) - \fint_{\Ball{R}{x_0}} u\,\Bigr| > t\Bigr\}
\le C_{\mathrm{JN}}(d)\, \mathrm{vol}(\Ball{R}{x_0})\,
\exp\!\Bigl(-\frac{t}{C_{\mathrm{JN}}(d)\,M}\Bigr).
\]
\end{theorem}

\begin{lstlisting}[style=Matlab-editor,basicstyle=\mlttfamily,escapechar=",mlshowsectionrules=true,mathescape=true,morecomment={[s]{\%\{}{\%\}}},language=lean,numbers=none]
/-- The full local John-Nirenberg inequality: if `u` has BMO seminorm at most `M`
on all sub-balls of `ball(x₀, 6R)`, then the level sets of `|u - (u)_B|` on `B = ball(x₀, R)`
decay exponentially. -/
theorem john_nirenberg_local
    {u : E → ℝ} {x₀ : E} {R M : ℝ}
    (hR : 0 < R) (hM : 0 < M)
    (_hu_meas : Measurable u)
    (hu_int : IntegrableOn u (Metric.ball x₀ (6 * R)) volume)
    (hu_bmo : ∀ (z : E) (s : ℝ), 0 < s →
      Metric.closedBall z s ⊆ Metric.ball x₀ (6 * R) →
      (⨍ x in Metric.ball z s, ‖u x - ⨍ y in Metric.ball z s, u y ∂volume‖ ∂volume) ≤ M)
    {t : ℝ} (ht : 0 < t) :
    volume ({x ∈ Metric.ball x₀ R |
      ‖u x - ⨍ y in Metric.ball x₀ R, u y ∂volume‖ > t}) ≤
    ENNReal.ofReal (C_JN d) * volume (Metric.ball x₀ R) *
      ENNReal.ofReal (Real.exp (-t / (C_JN d * M))) := by
  classical
  set B := Metric.ball x₀ R with hB_def
  set sixB := Metric.ball x₀ (6 * R) with hsixB_def
  set avgR := ⨍ y in B, u y ∂volume with havgR_def
  set w : E → ℝ := fun x => ‖u x - avgR‖ with hw_def
  set θ : ℝ := 1 / 2 with hθ_def
  set A : ℝ := 8 * 25 ^ d * M with hA_def
  set C5 : ℝ := 2 * 5 ^ d * M with hC5_def
  set F : ℝ → Type := fun lam => JNBadWitness w x₀ R lam with hF_def
  have hB_sub_six : B ⊆ sixB := by
    intro x hx
    rw [hB_def, Metric.mem_ball] at hx
    rw [hsixB_def, Metric.mem_ball]
    linarith
  have hw_int : IntegrableOn w sixB volume := by
    simpa [w, hw_def, sixB, hsixB_def] using (hu_int.sub integrableOn_const).norm
  have hθ_pos : 0 < θ := by simp [hθ_def]
  have hθ_lt : θ < 1 := by
    rw [hθ_def]
    norm_num
  have hA_pos : 0 < A := by
    rw [hA_def]
    positivity
  have hC5_nonneg : 0 ≤ C5 := by
    rw [hC5_def]
    positivity
  have hC5_le : C5 ≤ 4 * 25 ^ d * M := by
    rw [hC5_def]
    have hpow : (5 : ℝ) ^ d ≤ 25 ^ d :=
      pow_le_pow_left₀ (by positivity : (0 : ℝ) ≤ 5)
        (by norm_num : (5 : ℝ) ≤ 25) d
    nlinarith [hpow, hM]
  have hA_sub_C5_pos : 0 < A - C5 := by
    rw [hA_def]
    linarith [hC5_le, hM]
  have hA_ge_base : 4 * 6 ^ (d : ℕ) * M ≤ A := by
    rw [hA_def]
    have hpow : (6 : ℝ) ^ d ≤ 25 ^ d :=
      pow_le_pow_left₀ (by norm_num : (0 : ℝ) ≤ 6) (by norm_num : (6 : ℝ) ≤ 25) d
    nlinarith [hpow, hM]
  have hlocal_coeff :
      2 * M * 5 ^ d / (A - C5) ≤ 1 / (2 * 5 ^ d) := by
    have hden_lower : 4 * 25 ^ d * M ≤ A - C5 := by
      rw [hA_def]
      linarith [hC5_le]
    calc
      2 * M * 5 ^ d / (A - C5) ≤ 2 * M * 5 ^ d / (4 * 25 ^ d * M) := by
        apply div_le_div_of_nonneg_left
        · positivity
        · positivity
        · exact hden_lower
      _ = 1 / (2 * 5 ^ d) := by
        rw [show (25 : ℝ) ^ d = (5 : ℝ) ^ d * (5 : ℝ) ^ d by
          rw [show (25 : ℝ) = 5 * 5 by norm_num, mul_pow]]
        field_simp [hM.ne', show (5 : ℝ) ^ d ≠ 0 by positivity]
        ring
  have havg_base : ∀ lam : ℝ, A ≤ lam →
      ∀ x ∈ Metric.ball x₀ R, ⨍ y in Metric.ball x R, w y ∂volume ≤ lam := by
    intro lam hlamA
    simpa [w, hw_def, avgR, hB_def] using
      (base_ball_average_le (u := u) (x₀ := x₀) (R := R) (M := M) (lam := lam)
        hR hM (le_trans hA_ge_base hlamA) hu_int hu_bmo)
  let stopR : ∀ lam : ℝ, A ≤ lam → F lam → ℝ := fun lam hlamA =>
    stoppingRadiusBad hR (havg_base lam hlamA)
  let stopParentR : ∀ lam : ℝ, A ≤ lam → F lam → ℝ := fun lam hlamA =>
    stoppingParentRadiusBad hR (havg_base lam hlamA)
  let Ebad : ℝ → Set E := fun lam =>
    if hlamA : A ≤ lam then
      ⋃ p : F lam, Metric.ball p.center (stopR lam hlamA p)
    else
      sixB
  have hEbad_sub : ∀ lam : ℝ, Ebad lam ⊆ sixB := by
    intro lam
    by_cases hlamA : A ≤ lam
    · intro x hx
      dsimp [Ebad] at hx
      rw [dif_pos hlamA] at hx
      rcases Set.mem_iUnion.1 hx with ⟨p, hp⟩
      have hcenter : p.center ∈ B := p.center_mem
      have hrad_le : stopR lam hlamA p ≤ R := by
        dsimp [stopR]
        exact stoppingRadiusBad_le_R hR (havg_base lam hlamA) p
      have hball : dist x p.center < stopR lam hlamA p := Metric.mem_ball.mp hp
      have hcenter_ball : dist p.center x₀ < R := by
        simpa [hB_def] using hcenter
      rw [hsixB_def, Metric.mem_ball]
      calc
        dist x x₀ ≤ dist x p.center + dist p.center x₀ := dist_triangle _ _ _
        _ < stopR lam hlamA p + R := by linarith
        _ ≤ 6 * R := by linarith [hrad_le, hR]
    · simp [Ebad, hlamA]
  have hEbad_meas : ∀ lam : ℝ, MeasurableSet (Ebad lam) := by
    intro lam
    by_cases hlamA : A ≤ lam
    · dsimp [Ebad]
      rw [dif_pos hlamA]
      exact ((isOpen_iUnion fun _ : F lam => isOpen_ball).measurableSet)
    · dsimp [Ebad]
      rw [dif_neg hlamA]
      simpa [hsixB_def] using
        (measurableSet_ball : MeasurableSet (Metric.ball x₀ (6 * R)))
  have hEbad_anti : ∀ lam₁ lam₂ : ℝ, lam₁ ≤ lam₂ → Ebad lam₂ ⊆ Ebad lam₁ := by
    intro lam₁ lam₂ h12
    by_cases h₁ : A ≤ lam₁
    · have h₂ : A ≤ lam₂ := le_trans h₁ h12
      intro x hx
      dsimp [Ebad] at hx
      rw [dif_pos h₂] at hx
      rcases Set.mem_iUnion.1 hx with ⟨p₂, hp₂⟩
      let p₁ : F lam₁ := weakenBadWitness h12 p₂
      have hrad :
          stopR lam₂ h₂ p₂ ≤ stopR lam₁ h₁ p₁ := by
        dsimp [stopR, p₁]
        exact stoppingRadiusBad_mono hR
          (havg_base lam₁ h₁) (havg_base lam₂ h₂) h12 p₂
      dsimp [Ebad]
      rw [dif_pos h₁]
      exact Set.mem_iUnion.2 ⟨p₁, Metric.ball_subset_ball hrad hp₂⟩
    · intro x hx
      dsimp [Ebad]
      rw [dif_neg h₁]
      exact hEbad_sub lam₂ hx
  have h_point_subset_ae : ∀ s : ℝ, {x ∈ B | w x > s} ≤ᵐ[volume] Ebad s := by
    intro s
    by_cases hsA : A ≤ s
    · have hbad_ae :
          {x ∈ B | w x > s} ≤ᵐ[volume] JNBadCenterSet w x₀ R s :=
        ae_pointwise_gt_subset_badCenter (w := w) (x₀ := x₀) (R := R) hR hw_int
      filter_upwards [hbad_ae] with y hy
      intro hy_point
      rcases hy hy_point with ⟨hyB, r, hr_pos, hr_le, hbad⟩
      let p : F s :=
        { center := y
          witnessRadius := r
          center_mem := hyB
          radius_pos := hr_pos
          radius_le := hr_le
          bad := hbad }
      dsimp [Ebad]
      rw [dif_pos hsA]
      exact Set.mem_iUnion.2 ⟨p, by
        simpa [p] using
          (Metric.mem_ball_self (stoppingRadiusBad_pos hR (havg_base s hsA) p) :
            p.center ∈ Metric.ball p.center (stopR s hsA p))⟩
    · exact Filter.Eventually.of_forall fun y hy => by
        dsimp [Ebad]
        rw [dif_neg hsA]
        exact hB_sub_six hy.1
  have h_decay_bad : ∀ lam : ℝ, A ≤ lam →
      volume (Ebad (lam + A)) ≤ ENNReal.ofReal θ * volume (Ebad lam) := by
    intro lam hlamA
    let μ : ℝ := lam + A
    have hμA : A ≤ μ := by
      dsimp [μ]
      linarith
    have hlamμ : lam ≤ μ := by
      dsimp [μ]
      linarith
    by_cases hFμ_nonempty : Nonempty (F μ)
    · have hFlam_nonempty : Nonempty (F lam) := by
        rcases hFμ_nonempty with ⟨p⟩
        exact ⟨weakenBadWitness hlamμ p⟩
      let rLam : F lam → ℝ := stopR lam hlamA
      let rMu : F μ → ℝ := stopR μ hμA
      let parentLam : F lam → ℝ := stopParentR lam hlamA
      have hrLam_pos : ∀ p : F lam, 0 < rLam p := by
        intro p
        dsimp [rLam, stopR]
        exact stoppingRadiusBad_pos hR (havg_base lam hlamA) p
      have hrMu_pos : ∀ p : F μ, 0 < rMu p := by
        intro p
        dsimp [rMu, stopR]
        exact stoppingRadiusBad_pos hR (havg_base μ hμA) p
      have hrLam_bdd : BddAbove (Set.range rLam) := by
        refine ⟨R, ?_⟩
        rintro r ⟨p, rfl⟩
        dsimp [rLam, stopR]
        exact stoppingRadiusBad_le_R hR (havg_base lam hlamA) p
      have hrMu_bdd : BddAbove (Set.range rMu) := by
        refine ⟨R, ?_⟩
        rintro r ⟨p, rfl⟩
        dsimp [rMu, stopR]
        exact stoppingRadiusBad_le_R hR (havg_base μ hμA) p
      have hrawLam_nonempty : (⋃ p : F lam, Metric.ball p.center (rLam p)).Nonempty := by
        rcases hFlam_nonempty with ⟨p⟩
        refine ⟨p.center, Set.mem_iUnion.2 ?_⟩
        exact ⟨p, by simpa using (Metric.mem_ball_self (hrLam_pos p) : p.center ∈ Metric.ball p.center (rLam p))⟩
      have hrawMu_nonempty : (⋃ p : F μ, Metric.ball p.center (rMu p)).Nonempty := by
        rcases hFμ_nonempty with ⟨p⟩
        refine ⟨p.center, Set.mem_iUnion.2 ?_⟩
        exact ⟨p, by simpa using (Metric.mem_ball_self (hrMu_pos p) : p.center ∈ Metric.ball p.center (rMu p))⟩
      obtain ⟨SLam, hSLam_count, hSLam_disj, hSLam_hit, hSLam_cover⟩ :=
        vitali_covering_lemma (ι := F lam) (x := fun p => p.center) (r := rLam)
          hrLam_pos hrLam_bdd hrawLam_nonempty
      obtain ⟨SMu, hSMu_count, hSMu_disj, hSMu_hit, hSMu_cover⟩ :=
        vitali_covering_lemma (ι := F μ) (x := fun p => p.center) (r := rMu)
          hrMu_pos hrMu_bdd hrawMu_nonempty
      haveI : Countable SLam := hSLam_count.to_subtype
      haveI : Countable SMu := hSMu_count.to_subtype
      let BLam : SLam → JNBall x₀ R := fun q =>
        { center := q.1.center
          radius := rLam q.1
          radius_pos := hrLam_pos q.1
          radius_le := by
            dsimp [rLam, stopR]
            exact stoppingRadiusBad_le_R hR (havg_base lam hlamA) q.1
          center_mem := q.1.center_mem }
      let BMu : SMu → JNBall x₀ R := fun q =>
        { center := q.1.center
          radius := rMu q.1
          radius_pos := hrMu_pos q.1
          radius_le := by
            dsimp [rMu, stopR]
            exact stoppingRadiusBad_le_R hR (havg_base μ hμA) q.1
          center_mem := q.1.center_mem }
      have hBLam_disj : ∀ i j : SLam, i ≠ j → Disjoint ((BLam i).carrier) ((BLam j).carrier) := by
        intro i j hij
        exact hSLam_disj i.1 i.2 j.1 j.2 (fun h => hij (Subtype.ext h))
      have hBMu_disj : ∀ i j : SMu, i ≠ j → Disjoint ((BMu i).carrier) ((BMu j).carrier) := by
        intro i j hij
        exact hSMu_disj i.1 i.2 j.1 j.2 (fun h => hij (Subtype.ext h))
      have hBMu_cover : Ebad μ ⊆ ⋃ q : SMu, (BMu q).fivefold := by
        simpa [Ebad, hμA, BMu, rMu] using hSMu_cover
      have hULam_sub : (⋃ q : SLam, (BLam q).carrier) ⊆ Ebad lam := by
        intro x hx
        dsimp [Ebad]
        rw [dif_pos hlamA]
        rcases Set.mem_iUnion.1 hx with ⟨q, hq⟩
        exact Set.mem_iUnion.2 ⟨q.1, by simpa [BLam, rLam] using hq⟩
      have hUMu_sub : (⋃ q : SMu, (BMu q).carrier) ⊆ Ebad μ := by
        intro x hx
        dsimp [Ebad]
        rw [dif_pos hμA]
        rcases Set.mem_iUnion.1 hx with ⟨q, hq⟩
        exact Set.mem_iUnion.2 ⟨q.1, by simpa [BMu, rMu] using hq⟩
      let toLow : SMu → F lam := fun p => weakenBadWitness hlamμ p.1
      let assign : SMu → SLam := fun p =>
        ⟨Classical.choose (hSLam_hit (toLow p)), (Classical.choose_spec (hSLam_hit (toLow p))).1⟩
      have hassign_hit : ∀ p : SMu,
          (Metric.ball (toLow p).center (rLam (toLow p)) ∩ (BLam (assign p)).carrier).Nonempty := by
        intro p
        exact (Classical.choose_spec (hSLam_hit (toLow p))).2.1
      have hassign_rad : ∀ p : SMu, rLam (toLow p) ≤ 2 * (BLam (assign p)).radius := by
        intro p
        exact (Classical.choose_spec (hSLam_hit (toLow p))).2.2
      have hrMu_le_low : ∀ p : SMu, rMu p.1 ≤ rLam (toLow p) := by
        intro p
        dsimp [rMu, rLam, toLow, stopR]
        exact stoppingRadiusBad_mono hR
          (havg_base lam hlamA) (havg_base μ hμA) hlamμ p.1
      have hassign_sub : ∀ p : SMu, (BMu p).carrier ⊆ (BLam (assign p)).fivefold := by
        intro p
        apply ball_subset_fivefold_of_inter_nonempty
        · exact Metric.ball_subset_ball (hrMu_le_low p)
        · exact hassign_hit p
        · exact hassign_rad p
      let V : SLam → Set E := fun q => ⋃ p : {p : SMu // assign p = q}, (BMu p.1).carrier
      have hV_sub : ∀ q : SLam, V q ⊆ (BLam q).fivefold := by
        intro q x hx
        rcases Set.mem_iUnion.1 hx with ⟨p, hp⟩
        have : assign p.1 = q := p.2
        simpa [V, this] using hassign_sub p.1 hp
      have hV_meas : ∀ q : SLam, MeasurableSet (V q) := by
        intro q
        dsimp [V]
        exact (isOpen_iUnion fun _ => isOpen_ball).measurableSet
      have havg5_le : ∀ q : SLam,
          ⨍ x in (BLam q).fivefold, w x ∂volume ≤ lam + C5 := by
        intro q
        let b := BLam q
        let parent := parentLam q.1
        let avg5 := ⨍ x in b.fivefold, w x ∂volume
        let avgParent := ⨍ x in Metric.ball b.center parent, w x ∂volume
        have h5r_pos : 0 < 5 * b.radius := by
          nlinarith [b.radius_pos]
        have hparent_pos : 0 < parent := by
          dsimp [parent, parentLam, stopParentR]
          exact stoppingParentRadiusBad_pos hR (havg_base lam hlamA) q.1
        have hparent_le : parent ≤ 2 * b.radius := by
          dsimp [parent, b, BLam, parentLam, stopParentR, rLam, stopR]
          exact stoppingParentRadiusBad_le_double hR (havg_base lam hlamA) q.1
        have hrad_le_parent : b.radius ≤ parent := by
          dsimp [parent, b, BLam, parentLam, stopParentR, rLam, stopR]
          exact stoppingRadiusBad_le_parent hR (havg_base lam hlamA) q.1
        have hw_int_five : IntegrableOn w b.fivefold volume :=
          hw_int.mono_set b.fivefold_subset_sixBall
        have hw_bmo_five :
            (⨍ x in b.fivefold, ‖w x - avg5‖ ∂volume) ≤ 2 * M := by
          have hu_bmo_five :=
            hu_bmo b.center (5 * b.radius) h5r_pos b.closedBall_fivefold_subset_sixBall
          simpa [w, hw_def, avg5] using
            (abs_sub_const_bmo_le_two (M := M) (u := u) (c := avgR) (z := b.center)
              (s := 5 * b.radius) h5r_pos
              (hu_int := hu_int.mono_set b.fivefold_subset_sixBall)
              (hu_bmo := hu_bmo_five))
        have hsub_parent_5 :
            Metric.ball b.center parent ⊆ b.fivefold :=
          Metric.ball_subset_ball (by linarith [hparent_le, b.radius_pos])
        have hratio_le :
            ((5 * b.radius) / parent) ^ d ≤ (5 : ℝ) ^ d := by
          have hfrac_le : (5 * b.radius) / parent ≤ (5 : ℝ) := by
            exact (div_le_iff₀ hparent_pos).2 (by nlinarith [hrad_le_parent])
          exact pow_le_pow_left₀ (by positivity : (0 : ℝ) ≤ (5 * b.radius) / parent)
            hfrac_le d
        have hdiff :
            |avgParent - avg5| ≤
              ((5 * b.radius) / parent) ^ d *
                (⨍ z in b.fivefold, ‖w z - avg5‖ ∂volume) := by
          simpa [avgParent, avg5, JNBall.fivefold] using
            (abs_subballAverage_sub_ballAverage_le (u := w)
              (x := b.center) (c := b.center) h5r_pos hparent_pos hsub_parent_5 hw_int_five)
        have havgParent_le : avgParent ≤ lam := by
          dsimp [avgParent, parent, parentLam, stopParentR]
          exact stoppingParentRadiusBad_le hR (havg_base lam hlamA) q.1
        have hnorm_nonneg : 0 ≤ ⨍ z in b.fivefold, ‖w z - avg5‖ ∂volume := by
          rw [MeasureTheory.setAverage_eq]
          exact smul_nonneg (by positivity) (MeasureTheory.integral_nonneg fun _ => norm_nonneg _)
        have hdiff' : |avgParent - avg5| ≤ (5 : ℝ) ^ d * (2 * M) := by
          calc
            |avgParent - avg5| ≤
                ((5 * b.radius) / parent) ^ d *
                  (⨍ z in b.fivefold, ‖w z - avg5‖ ∂volume) := hdiff
            _ ≤ (5 : ℝ) ^ d * (⨍ z in b.fivefold, ‖w z - avg5‖ ∂volume) := by
                exact mul_le_mul_of_nonneg_right hratio_le hnorm_nonneg
            _ ≤ (5 : ℝ) ^ d * (2 * M) := by
                exact mul_le_mul_of_nonneg_left hw_bmo_five (by positivity)
        have hdiff_upper : avg5 - avgParent ≤ C5 := by
          rw [hC5_def]
          have hpair := abs_le.mp hdiff'
          linarith
        linarith [havgParent_le, hdiff_upper]
      have hlarge : ∀ p : SMu, μ < ⨍ x in (BMu p).carrier, w x ∂volume := by
        intro p
        dsimp [BMu, rMu, stopR]
        exact stoppingRadiusBad_bad hR (havg_base μ hμA) p.1
      let UMu : Set E := ⋃ p : SMu, (BMu p).carrier
      let ULam : Set E := ⋃ p : SLam, (BLam p).carrier
      have hUMu_le_ULam :
          volume.real UMu ≤ (1 / (2 * 5 ^ d)) * volume.real ULam := by
        simpa [UMu, ULam] using
          assigned_union_half_measure_real
            (d := d) (_hR := hR) (_hM := hM) (hμ_eq := by rfl)
            (hA_sub_C5_pos := hA_sub_C5_pos) (hlocal_coeff := hlocal_coeff)
            (hu_int := hu_int) (hu_bmo := hu_bmo) (w := w) (hw_def := hw_def)
            (BLam := BLam) (BMu := BMu) (hBLam_disj := hBLam_disj)
            (hBMu_disj := hBMu_disj) (assign := assign) (hassign_sub := hassign_sub)
            (havg5_le := havg5_le) (hlarge := hlarge)
      have hEbadMu_real :
          volume.real (Ebad μ) ≤ (5 : ℝ) ^ d * volume.real UMu := by
        simpa [UMu] using
          fivefold_cover_real_le
            (d := d) (x₀ := x₀) (R := R) (S := Ebad μ)
            (hS_sub := by simpa [sixB, hsixB_def] using hEbad_sub μ)
            (B := BMu) (hcover := hBMu_cover) (hdisj := hBMu_disj)
      have hEbadLam_fin : volume (Ebad lam) ≠ ⊤ :=
        measure_ne_top_of_subset (hEbad_sub lam)
          (measure_ball_lt_top (μ := volume) (x := x₀) (r := 6 * R)).ne
      have hULam_le_Ebad : volume.real ULam ≤ volume.real (Ebad lam) := by
        exact measureReal_mono hULam_sub hEbadLam_fin
      have hreal_decay :
          volume.real (Ebad μ) ≤ (1 / 2) * volume.real (Ebad lam) := by
        calc
          volume.real (Ebad μ) ≤ (5 : ℝ) ^ d * volume.real UMu := hEbadMu_real
          _ ≤ (5 : ℝ) ^ d * ((1 / (2 * 5 ^ d)) * volume.real ULam) := by
              exact mul_le_mul_of_nonneg_left hUMu_le_ULam (by positivity)
          _ = (1 / 2) * volume.real ULam := by
              rw [← mul_assoc, five_pow_half_factor]
          _ ≤ (1 / 2) * volume.real (Ebad lam) := by
              exact mul_le_mul_of_nonneg_left hULam_le_Ebad (by norm_num)
      have hEbadMu_fin : volume (Ebad μ) ≠ ⊤ :=
        measure_ne_top_of_subset (hEbad_sub μ)
          (measure_ball_lt_top (μ := volume) (x := x₀) (r := 6 * R)).ne
      have henn_decay :
          volume (Ebad μ) ≤ ENNReal.ofReal (1 / 2) * volume (Ebad lam) :=
        volume_real_le_to_ennreal hEbadMu_fin hEbadLam_fin (by norm_num) hreal_decay
      simpa [Ebad, hμA, θ, hθ_def, μ] using henn_decay
    · haveI : IsEmpty (F μ) := not_nonempty_iff.mp hFμ_nonempty
      have hEbadμ_empty : Ebad μ = ∅ := by
        simp [Ebad, hμA]
      simp [hEbadμ_empty, θ, μ]
  have h_exp_bad :=
    level_set_family_from_base
      (by simpa [hsixB_def] using (measurableSet_ball : MeasurableSet (Metric.ball x₀ (6 * R))))
      (measure_ball_lt_top (μ := volume) (x := x₀) (r := 6 * R)).ne
      hEbad_sub hEbad_meas hEbad_anti hA_pos hθ_pos hθ_lt h_decay_bad t ht
  have h_exp :
      volume ({x ∈ B | w x > t}) ≤
        ENNReal.ofReal 4 * volume sixB *
          ENNReal.ofReal (Real.exp (-t * (-Real.log θ / A))) := by
    have hmain :
        volume ({x ∈ B | w x > t}) ≤
          ENNReal.ofReal (1 / θ ^ 2) * volume sixB *
            ENNReal.ofReal (Real.exp (-t * (-Real.log θ / A))) := by
      calc
        volume ({x ∈ B | w x > t}) ≤ volume (Ebad t) := measure_mono_ae (h_point_subset_ae t)
        _ ≤ ENNReal.ofReal (1 / θ ^ 2) * volume sixB *
              ENNReal.ofReal (Real.exp (-t * (-Real.log θ / A))) := h_exp_bad
    have hθsq : (1 / θ ^ 2 : ℝ) = 4 := by
      rw [hθ_def]
      norm_num
    simpa [hθsq] using hmain
  have hsixB_vol :
      volume sixB = ENNReal.ofReal ((6 : ℝ) ^ d) * volume B := by
    have hleft : volume sixB ≠ ⊤ := measure_ball_lt_top.ne
    have hright : ENNReal.ofReal ((6 : ℝ) ^ d) * volume B ≠ ⊤ :=
      ENNReal.mul_ne_top ENNReal.ofReal_ne_top measure_ball_lt_top.ne
    rw [← ENNReal.toReal_eq_toReal_iff' hleft hright]
    rw [hsixB_def, hB_def, ENNReal.toReal_mul,
      ENNReal.toReal_ofReal (by positivity : (0 : ℝ) ≤ (6 : ℝ) ^ d)]
    change volume.real (Metric.ball x₀ (6 * R)) = (6 : ℝ) ^ d * volume.real (Metric.ball x₀ R)
    rw [volumeReal_ball_eq x₀ (show 0 < 6 * R by nlinarith), volumeReal_ball_eq x₀ hR]
    rw [mul_pow]
    ring
  have h_coeff_ball :
      ENNReal.ofReal 4 * volume sixB ≤ ENNReal.ofReal (C_JN d) * volume B := by
    rw [hsixB_vol]
    have hcoeff : 4 * (6 : ℝ) ^ d ≤ C_JN d := by
      rw [C_JN]
      have h36 : (6 : ℝ) ^ d ≤ 36 ^ d :=
        pow_le_pow_left₀ (by positivity : (0 : ℝ) ≤ 6) (by norm_num : (6 : ℝ) ≤ 36) d
      calc
        4 * (6 : ℝ) ^ d ≤ 4 * 36 ^ d := by
          gcongr
        _ ≤ 16 * 36 ^ d := by
          have h36_nonneg : 0 ≤ (36 : ℝ) ^ d := by positivity
          nlinarith
    calc
      ENNReal.ofReal 4 * (ENNReal.ofReal ((6 : ℝ) ^ d) * volume B)
          = ENNReal.ofReal (4 * (6 : ℝ) ^ d) * volume B := by
              rw [← mul_assoc, ← ENNReal.ofReal_mul (by positivity)]
      _ ≤ ENNReal.ofReal (C_JN d) * volume B := by
          exact mul_le_mul_of_nonneg_right (ENNReal.ofReal_le_ofReal hcoeff) (by positivity)
  have h_const_exp : ∀ s : ℝ, 0 < s →
      -s * (-Real.log θ / A) ≤ -s / (C_JN d * M) := by
    intro s hs
    have hCJN_pos : 0 < C_JN d := C_JN_pos d
    have hCM_pos : 0 < C_JN d * M := mul_pos hCJN_pos hM
    suffices hkey : A ≤ C_JN d * M * (-Real.log θ) by
      have h_rate : 1 / (C_JN d * M) ≤ -Real.log θ / A := by
        rw [div_le_div_iff₀ hCM_pos hA_pos, one_mul]
        linarith [mul_comm (C_JN d * M) (-Real.log θ)]
      have := mul_le_mul_of_nonpos_left h_rate (neg_nonpos.mpr (le_of_lt hs))
      simp only [mul_div_assoc'] at this
      convert this using 1 <;> ring
    have hlog : Real.log θ ≤ θ - 1 := Real.log_le_sub_one_of_pos hθ_pos
    have hneg_log : (1 / 2 : ℝ) ≤ -Real.log θ := by
      rw [hθ_def] at hlog ⊢
      linarith
    calc
      A = 8 * 25 ^ d * M := hA_def
      _ ≤ 8 * 36 ^ d * M := by
          have hpow : (25 : ℝ) ^ d ≤ 36 ^ d :=
            pow_le_pow_left₀ (by positivity : (0 : ℝ) ≤ 25) (by norm_num : (25 : ℝ) ≤ 36) d
          nlinarith [hpow, hM]
      _ = C_JN d * M * (1 / 2) := by
          rw [C_JN]
          ring
      _ ≤ C_JN d * M * (-Real.log θ) := by
          exact mul_le_mul_of_nonneg_left hneg_log (by positivity : 0 ≤ C_JN d * M)
  have hexp_bound :
      ENNReal.ofReal (Real.exp (-t * (-Real.log θ / A))) ≤
        ENNReal.ofReal (Real.exp (-t / (C_JN d * M))) := by
    exact ENNReal.ofReal_le_ofReal (Real.exp_le_exp.2 (h_const_exp t ht))
  calc
    volume ({x ∈ Metric.ball x₀ R |
        ‖u x - ⨍ y in Metric.ball x₀ R, u y ∂volume‖ > t})
        = volume ({x ∈ B | w x > t}) := by
            simp [hB_def, w, avgR]
    _ ≤ ENNReal.ofReal 4 * volume sixB *
          ENNReal.ofReal (Real.exp (-t * (-Real.log θ / A))) := h_exp
    _ = (ENNReal.ofReal 4 * volume sixB) *
          ENNReal.ofReal (Real.exp (-t * (-Real.log θ / A))) := by
            rfl
    _ ≤ (ENNReal.ofReal (C_JN d) * volume B) *
          ENNReal.ofReal (Real.exp (-t * (-Real.log θ / A))) := by
            exact mul_le_mul_of_nonneg_right h_coeff_ball (by positivity)
    _ = ENNReal.ofReal (C_JN d) * volume B *
          ENNReal.ofReal (Real.exp (-t * (-Real.log θ / A))) := by
            rfl
    _ ≤ ENNReal.ofReal (C_JN d) * volume B *
          ENNReal.ofReal (Real.exp (-t / (C_JN d * M))) := by
            exact mul_le_mul_of_nonneg_left hexp_bound
              (by positivity : 0 ≤ ENNReal.ofReal (C_JN d) * volume B)


\end{lstlisting}
\begin{proof}
Set $w(x) = |u(x) - \fint_{\Ball{R}{x_0}} u|$, $\theta = 1/2$, $A = 8\cdot 25^{d}\, M$.
For each $\ell \ge A$, apply a Calder\'on--Zygmund stopping-time / Vitali selection
inside $\Ball{R}{x_0}$ on the level set $\{w > \ell\}$ to extract a disjoint family of
JN-balls $(B_i)$ whose 5-fold enlargements cover the level set (up to a null set), with
the property that on each parent ball the average of $w$ does not exceed $\ell$. The
BMO hypothesis applied on each $B_i$ then yields
\[
\mathrm{vol}(\{w > \ell + A\}\cap B_i^{\mathrm{car}})
\le \tfrac{1}{2}\, \mathrm{vol}(B_i^{\mathrm{car}}),
\]
because raising $w$ by $A$ over the parent average exceeds the BMO control of the
oscillation. Lemma \leanname{john\_nirenberg\_level\_set\_decay} converts this into the
geometric decay $\mathrm{vol}(\{w > \ell + A\}) \le \theta\, \mathrm{vol}(\{w > \ell\})$ for
all $\ell \ge A$. Theorem \leanname{john\_nirenberg\_from\_base} then yields the
exponential bound, where $-\log\theta/A = \log 2 / (8\cdot 25^{d}\, M)$ and the
constants combine into $C_{\mathrm{JN}}(d) = 16\cdot 36^{d}$.
\end{proof}

\subsection{Campanato characterisation of H\"older spaces}

\begin{definition}[\leanname{CampanatoBall}, \leanname{campanatoBallValue},
\leanname{campanatoSeminorm}]
For $x_0 \in E$ and $R > 0$, a \emph{Campanato ball} is a pair $(c,\rho)$ with
$0 < \rho \le R$ and $\Ball{\rho}{c} \subseteq \Ball{R}{x_0}$. For $\alpha > 0$ define
\[
N_{\alpha}(u; c,\rho)
:= \rho^{-\alpha}\, \fint_{\Ball{\rho}{c}}
\Bigl|\, u(x) - \fint_{\Ball{\rho}{c}} u\, \Bigr|\, dx,
\]
and the Campanato seminorm
\[
[u]_{\mathcal{L}^{\alpha}(\Ball{R}{x_0})}
:= \sup_{(c,\rho)\,\text{Campanato ball}} N_{\alpha}(u;c,\rho).
\]
\end{definition}

\begin{lstlisting}[style=Matlab-editor,basicstyle=\mlttfamily,escapechar=",mlshowsectionrules=true,mathescape=true,morecomment={[s]{\%\{}{\%\}}},language=lean,numbers=none]
/-! ## Campanato Characterization of Hölder Spaces -/

/-- Admissible sub-balls of `Metric.ball x₀ R` for the Campanato seminorm. -/
def CampanatoBall (x₀ : E) (R : ℝ) :=
  {p : E × ℝ // 0 < p.2 ∧ p.2 ≤ R ∧ Metric.ball p.1 p.2 ⊆ Metric.ball x₀ R}

/-- The Campanato seminorm on a ball B_R(x₀):
sup_{B ⊆ B_R} ⨍_B |u - (u)_B|. -/
noncomputable def campanatoBallValue
    (u : E → ℝ) (α : ℝ) {x₀ : E} {R : ℝ} (b : CampanatoBall x₀ R) : ℝ :=
  (CampanatoBall.radius b)⁻¹ ^ α *
    ⨍ x in Metric.ball (CampanatoBall.center b) (CampanatoBall.radius b),
      |u x - ⨍ z in Metric.ball (CampanatoBall.center b) (CampanatoBall.radius b), u z| ∂volume

noncomputable def campanatoSeminorm (u : E → ℝ) (x₀ : E) (R : ℝ) (α : ℝ) : ℝ :=
  sSup (Set.range (campanatoBallValue u α (x₀ := x₀) (R := R)))

\end{lstlisting}
\begin{definition}[\leanname{HasCampanatoBound}]
We write $\mathrm{HasCampanatoBound}(u; x_0, R, \alpha, C)$ if $u$ is locally integrable on every
admissible sub-ball and $N_{\alpha}(u; c, \rho) \le C$ for every Campanato ball
$(c,\rho) \subseteq \Ball{R}{x_0}$.
\end{definition}

\begin{lstlisting}[style=Matlab-editor,basicstyle=\mlttfamily,escapechar=",mlshowsectionrules=true,mathescape=true,morecomment={[s]{\%\{}{\%\}}},language=lean,numbers=none]
/-- Expanded local form of a Campanato bound on `Metric.ball x₀ R`. -/
def HasCampanatoBound (u : E → ℝ) (x₀ : E) (R α C : ℝ) : Prop :=
  ∀ b : CampanatoBall x₀ R,
    IntegrableOn u (Metric.ball (CampanatoBall.center b) (CampanatoBall.radius b)) volume ∧
      campanatoBallValue u α b ≤ C

\end{lstlisting}
\begin{definition}[Dyadic ball averages and limits,
\leanname{dyadicBallAverage}, \leanname{dyadicBallAverageLimit}]
For $\rho > 0$ set
\begin{gather*}
a_n(x) := \fint_{\Ball{\rho/2^{n}}{x}} u\, dz, \\
\widetilde u(x) := \lim_{n\to\infty} a_n(x)
\end{gather*}
when the limit exists; $\widetilde u$ is the candidate H\"older representative.
\end{definition}

\begin{lstlisting}[style=Matlab-editor,basicstyle=\mlttfamily,escapechar=",mlshowsectionrules=true,mathescape=true,morecomment={[s]{\%\{}{\%\}}},language=lean,numbers=none]
noncomputable def dyadicBallAverage (u : E → ℝ) (x : E) (ρ : ℝ) (n : ℕ) : ℝ :=
  ⨍ z in Metric.ball x (ρ / (2 : ℝ) ^ n), u z ∂volume

noncomputable def dyadicBallAverageLimit
    (u : E → ℝ) (x : E) (ρ : ℝ) : ℝ :=
  limUnder atTop (dyadicBallAverage u x ρ)

\end{lstlisting}
\begin{definition}[\leanname{C\_campanato\_holder}]
\[
C_{\mathrm{Camp\to H\ddot ol}}(d,\alpha) := \frac{2^{d+3}}{1 - 2^{-\alpha}}.
\]
\end{definition}

\begin{lstlisting}[style=Matlab-editor,basicstyle=\mlttfamily,escapechar=",mlshowsectionrules=true,mathescape=true,morecomment={[s]{\%\{}{\%\}}},language=lean,numbers=none]
/-- A conservative Hölder constant extracted from a Campanato bound. -/
noncomputable def C_campanato_holder (d : ℕ) (α : ℝ) : ℝ :=
  (2 : ℝ) ^ (d + 3) / (1 - (2 : ℝ) ^ (-α))

\end{lstlisting}
The chain of lemmas
\leanname{abs\_ballAverage\_half\_sub\_ballAverage\_le},
\leanname{ballAverage\_sub\_ballAverage\_sameRadius\_le\_campanato},
\leanname{dyadicBallAverage\_step\_le\_campanato},
\leanname{dyadicBallAverage\_step\_le\_geometric},
\leanname{dyadicBallAverage\_cauchy},
\leanname{tendsto\_dyadicBallAverageLimit},
\leanname{dyadicBallAverage\_tail\_le}
establishes that the dyadic averages form a
Cauchy sequence with geometric ratio $2^{-\alpha}$, and converge to a limit
$\widetilde u$ with quantitative tail control. Lemmas
\leanname{dyadicBallAverageLimit\_holder\_small} and
\leanname{dyadicBallAverageLimit\_holder\_large}
give pointwise H\"older estimates on the
limit in the regimes $|x-y| \le \rho$ and $|x-y| > \rho$ respectively.

\begin{theorem}[Campanato implies H\"older, \leanname{campanato\_implies\_holder}]
Let $0 < \alpha \le 1$, $R > 0$ and $\mathrm{HasCampanatoBound}(u; x_0, R, \alpha, C_{\mathrm{camp}})$.
Then there is a function $v : E \to \R$ that agrees with $u$ a.e.\ on $\Ball{R}{x_0}$ and
satisfies, for all $x, y \in \Ball{R/2}{x_0}$,
\[
\bigl|v(x) - v(y)\bigr|
\le C_{\mathrm{Camp\to H\ddot ol}}(d,\alpha)\, C_{\mathrm{camp}}\, \lvert x - y\rvert^{\alpha}.
\]
\end{theorem}

\begin{lstlisting}[style=Matlab-editor,basicstyle=\mlttfamily,escapechar=",mlshowsectionrules=true,mathescape=true,morecomment={[s]{\%\{}{\%\}}},language=lean,numbers=none]
/-- A Campanato bound determines a Hölder representative up to a.e. equality. -/
theorem campanato_implies_holder
    {u : E → ℝ} {x₀ : E} {R α C_camp : ℝ}
    (hα : 0 < α) (hα_le : α ≤ 1) (hR : 0 < R)
    (hcamp : HasCampanatoBound u x₀ R α C_camp) :
    ∃ v : E → ℝ,
      (∀ᵐ x ∂volume.restrict (Metric.ball x₀ R), v x = u x) ∧
      ∀ x ∈ Metric.ball x₀ (R / 2), ∀ y ∈ Metric.ball x₀ (R / 2),
        |v x - v y| ≤ C_campanato_holder d α * C_camp * ‖x - y‖ ^ α := by
  classical
  let outer : Set E := Metric.ball x₀ R
  let inner : Set E := Metric.ball x₀ (R / 2)
  let ρ : ℝ := R / 2
  let q : ℝ := (2 : ℝ) ^ (-α)
  let invDen : ℝ := (1 - q)⁻¹
  let v : E → ℝ := fun x =>
    if x ∈ inner then dyadicBallAverageLimit u x ρ else u x
  have hρ : 0 < ρ := by
    dsimp [ρ]
    positivity
  have hρR : ρ ≤ R := by
    dsimp [ρ]
    linarith
  have hC_nonneg : 0 ≤ C_camp := hcamp.nonneg hα hR
  have hq_lt : q < 1 := by
    dsimp [q]
    rw [Real.rpow_neg (by positivity)]
    have h2a : 1 < (2 : ℝ) ^ α := Real.one_lt_rpow (by norm_num) hα
    exact inv_lt_one_of_one_lt₀ h2a
  have hden_pos : 0 < 1 - q := sub_pos.mpr hq_lt
  have hinvDen_nonneg : 0 ≤ invDen := by
    dsimp [invDen]
    positivity
  have h2a_le : (2 : ℝ) ^ α ≤ 2 := by
    simpa using
      (Real.rpow_le_rpow_of_exponent_le (by norm_num : 1 ≤ (2 : ℝ)) hα_le)
  have hq_ge_half : (1 / 2 : ℝ) ≤ q := by
    dsimp [q]
    rw [Real.rpow_neg (by positivity)]
    have h2a_pos : 0 < (2 : ℝ) ^ α := Real.rpow_pos_of_pos (by norm_num) α
    simpa [one_div] using (one_div_le_one_div_of_le h2a_pos h2a_le)
  have hinvDen_ge_two : (2 : ℝ) ≤ invDen := by
    dsimp [invDen]
    have hden_half : 1 - q ≤ (1 / 2 : ℝ) := by linarith
    have hinv := one_div_le_one_div_of_le hden_pos hden_half
    simpa using hinv
  have hinner_eq_global :
      ∀ᵐ x ∂volume, x ∈ inner → dyadicBallAverageLimit u x ρ = u x := by
    have hinner_eq_ae :
        ∀ᵐ x ∂volume.restrict inner, dyadicBallAverageLimit u x ρ = u x := by
      simpa [inner, ρ] using ae_eq_dyadicBallAverageLimit_on_halfBall hα hR hcamp
    rw [ae_restrict_iff' measurableSet_ball] at hinner_eq_ae
    exact hinner_eq_ae
  refine ⟨v, ?_, ?_⟩
  · rw [ae_restrict_iff' measurableSet_ball]
    filter_upwards [hinner_eq_global] with x hx hx_outer
    by_cases hx_inner : x ∈ inner
    · dsimp [v]
      rw [if_pos hx_inner]
      exact hx hx_inner
    · dsimp [v]
      rw [if_neg hx_inner]
  · intro x hx y hy
    have hx_inner : x ∈ inner := by simpa [inner] using hx
    have hy_inner : y ∈ inner := by simpa [inner] using hy
    have hvx : v x = dyadicBallAverageLimit u x ρ := by
      dsimp [v]
      rw [if_pos hx_inner]
    have hvy : v y = dyadicBallAverageLimit u y ρ := by
      dsimp [v]
      rw [if_pos hy_inner]
    by_cases hsmall : ‖x - y‖ ≤ ρ
    · rw [hvx, hvy]
      simpa [ρ] using
        (dyadicBallAverageLimit_holder_small hα hα_le hR hcamp hx hy
          (by simpa [ρ] using hsmall))
    · have hlarge : ρ ≤ ‖x - y‖ := le_of_not_ge hsmall
      rw [hvx, hvy]
      simpa [ρ] using
        (dyadicBallAverageLimit_holder_large hα hα_le hR hcamp hx hy
          (by simpa [ρ] using hlarge))

\end{lstlisting}
\begin{proof}
Take $\rho = R/2$ and define $v(x) := \widetilde u(x)$ for $x \in \Ball{R/2}{x_0}$ and
$v(x) := u(x)$ otherwise. By
\leanname{ae\_eq\_dyadicBallAverageLimit\_on\_halfBall}, $\widetilde u = u$ a.e.\ on
$\Ball{R/2}{x_0}$, so $v = u$ a.e.\ on $\Ball{R}{x_0}$. For $x, y \in \Ball{R/2}{x_0}$,
split into the cases $|x - y| \le \rho$ and $|x - y| > \rho$, applying the corresponding
H\"older estimate for the dyadic limit. Both cases yield the same bound with constant
$C_{\mathrm{Camp\to H\ddot ol}}(d,\alpha)\, C_{\mathrm{camp}}$.
\end{proof}

\begin{lemma}[H\"older implies Campanato bound, \leanname{holder\_hasCampanatoBound}]
If $u$ is $\alpha$-H\"older on $\Ball{R}{x_0}$ with constant $C_{\mathrm{H\ddot ol}}$, then
$\mathrm{HasCampanatoBound}(u; x_0, R, \alpha, C_{\mathrm{H\ddot ol}}\cdot 2^{\alpha})$.
\end{lemma}

\begin{lstlisting}[style=Matlab-editor,basicstyle=\mlttfamily,escapechar=",mlshowsectionrules=true,mathescape=true,morecomment={[s]{\%\{}{\%\}}},language=lean,numbers=none]
lemma holder_hasCampanatoBound
    {u : E → ℝ} {x₀ : E} {R α C_hold : ℝ}
    (hα : 0 < α) (hC_hold : 0 ≤ C_hold)
    (hhold : ∀ x ∈ Metric.ball x₀ R, ∀ y ∈ Metric.ball x₀ R,
      |u x - u y| ≤ C_hold * ‖x - y‖ ^ α) :
    HasCampanatoBound u x₀ R α (C_hold * (2 : ℝ) ^ α) := by
  intro b
  rcases b with ⟨⟨y, r⟩, hbr⟩
  rcases hbr with ⟨hr, hrR, hsub⟩
  let αnn : ℝ≥0 := ⟨α, hα.le⟩
  let Chold : ℝ≥0 := ⟨C_hold, hC_hold⟩
  let B : Set E := Metric.ball y r
  let avg : ℝ := ⨍ z in B, u z ∂volume
  have hBmeas : MeasurableSet B := by
    simpa [B] using (Metric.isOpen_ball.measurableSet : MeasurableSet (Metric.ball y r))
  have hyB : y ∈ B := by
    show y ∈ Metric.ball y r
    exact Metric.mem_ball_self hr
  have hyR : y ∈ Metric.ball x₀ R := hsub hyB
  have hBfin : volume B ≠ ∞ := by
    simpa [B] using (measure_ball_lt_top (μ := volume) (x := y) (r := r)).ne
  have hB0 : volume B ≠ 0 := by
    exact (measure_ball_pos (μ := volume) y hr).ne'
  have hholderB : HolderOnWith Chold αnn u B := by
    intro x hx z hz
    have hChold : ENNReal.ofReal C_hold = (Chold : ℝ≥0∞) := by
      simpa [Chold] using (ENNReal.ofReal_eq_coe_nnreal hC_hold)
    have hdist0 : ENNReal.ofReal (dist (u x) (u z)) ≤
        ENNReal.ofReal C_hold * ENNReal.ofReal (dist x z) ^ α := by
      rw [ENNReal.ofReal_rpow_of_nonneg dist_nonneg hα.le,
        ← ENNReal.ofReal_mul hC_hold]
      exact ENNReal.ofReal_le_ofReal <| by
        simpa [Real.dist_eq, dist_eq_norm] using hhold x (hsub hx) z (hsub hz)
    have hdist : ENNReal.ofReal (dist (u x) (u z)) ≤
        (Chold : ℝ≥0∞) * ENNReal.ofReal (dist x z) ^ α := by
      calc
        ENNReal.ofReal (dist (u x) (u z))
            ≤ ENNReal.ofReal C_hold * ENNReal.ofReal (dist x z) ^ α := hdist0
        _ = (Chold : ℝ≥0∞) * ENNReal.ofReal (dist x z) ^ α := by rw [hChold]
    simpa [αnn, edist_dist] using hdist
  have hcontB : ContinuousOn u B := hholderB.continuousOn (show 0 < αnn by exact hα)
  have hu_int : IntegrableOn u B volume := by
    have hBfin_lt : volume B < ∞ := lt_top_iff_ne_top.2 hBfin
    refine IntegrableOn.of_bound hBfin_lt (hcontB.aestronglyMeasurable hBmeas)
      (‖u y‖ + C_hold * r ^ α) ?_
    filter_upwards [self_mem_ae_restrict hBmeas] with z hz
    have hzR : z ∈ Metric.ball x₀ R := hsub hz
    have hdist_le : dist z y ≤ r := le_of_lt (by simpa [B] using hz)
    have hrpow_le : dist z y ^ α ≤ r ^ α :=
      Real.rpow_le_rpow dist_nonneg hdist_le hα.le
    have hdiff : |u z - u y| ≤ C_hold * r ^ α := by
      calc
        |u z - u y| ≤ C_hold * dist z y ^ α := by
          simpa [Real.dist_eq] using hhold z hzR y hyR
        _ ≤ C_hold * r ^ α := by
          exact mul_le_mul_of_nonneg_left hrpow_le hC_hold
    calc
      ‖u z‖ ≤ |u z - u y| + ‖u y‖ := by
        simpa [Real.dist_eq] using (dist_triangle (u z) (u y) 0)
      _ ≤ ‖u y‖ + C_hold * r ^ α := by
          linarith
  refine ⟨hu_int, ?_⟩
  have hpoint : ∀ x ∈ B, |u x - avg| ≤ C_hold * (2 * r) ^ α := by
    obtain ⟨c, hcB, havg⟩ := MeasureTheory.exists_eq_setAverage
      (μ := volume) (s := B) (f := u) (isConnected_ball hr) hcontB hu_int hBfin hB0
    intro x hx
    have hxR : x ∈ Metric.ball x₀ R := hsub hx
    have hcR : c ∈ Metric.ball x₀ R := hsub hcB
    have hdist_le : dist x c ≤ 2 * r := by
      have hxy : dist x y < r := by simpa [B] using hx
      have hyc : dist c y < r := by simpa [B] using hcB
      have hyc' : dist y c < r := by simpa [dist_comm] using hyc
      have htri : dist x c ≤ dist x y + dist y c := dist_triangle _ _ _
      have hsum : dist x y + dist y c < 2 * r := by
        linarith
      exact htri.trans hsum.le
    calc
      |u x - avg| = |u x - u c| := by
        change |u x - (⨍ z in B, u z ∂volume)| = |u x - u c|
        rw [← havg]
      _ ≤ C_hold * dist x c ^ α := by
          simpa [Real.dist_eq] using hhold x hxR c hcR
      _ ≤ C_hold * (2 * r) ^ α := by
          exact mul_le_mul_of_nonneg_left
            (Real.rpow_le_rpow dist_nonneg hdist_le hα.le) hC_hold
  have hAvgBound : (⨍ x in B, |u x - avg| ∂volume) ≤ C_hold * (2 * r) ^ α := by
    have hg_int : IntegrableOn (fun x => |u x - avg|) B volume := by
      simpa [Real.norm_eq_abs, avg] using (hu_int.sub (integrableOn_const hBfin)).norm
    obtain ⟨x', hx'B, hAvgLe⟩ := MeasureTheory.exists_setAverage_le
        (μ := volume) (s := B) (f := fun x => |u x - avg|) hB0 hBfin hg_int
    exact hAvgLe.trans (hpoint x' hx'B)
  have hrpow :
      r⁻¹ ^ α * (2 * r) ^ α = (r⁻¹ * (2 * r)) ^ α := by
    simpa using
      (Real.mul_rpow (x := r⁻¹) (y := 2 * r) (z := α)
        (inv_nonneg.mpr hr.le) (by positivity)).symm
  have hmul : r⁻¹ * (2 * r) = (2 : ℝ) := by
    calc
      r⁻¹ * (2 * r) = 2 * (r⁻¹ * r) := by ring
      _ = 2 := by rw [inv_mul_cancel₀ (ne_of_gt hr), mul_one]
  calc
    r⁻¹ ^ α * (⨍ x in B, |u x - avg| ∂volume)
        ≤ r⁻¹ ^ α * (C_hold * (2 * r) ^ α) := by
        exact mul_le_mul_of_nonneg_left hAvgBound (Real.rpow_nonneg (inv_nonneg.mpr hr.le) _)
    _ = C_hold * (r⁻¹ ^ α * (2 * r) ^ α) := by ring
    _ = C_hold * ((r⁻¹ * (2 * r)) ^ α) := by rw [hrpow]
    _ = C_hold * (2 : ℝ) ^ α := by rw [hmul]

\end{lstlisting}
\begin{proof}
On any Campanato ball $\Ball{\rho}{y}$, for any $x, z \in \Ball{\rho}{y}$ we have
$|x - z| \le 2\rho$, so the H\"older bound gives
\[
|u(x) - u(z)| \le C_{\mathrm{H\ddot ol}}\,\lvert x - z\rvert^{\alpha}
\le C_{\mathrm{H\ddot ol}}\,(2\rho)^{\alpha}.
\]
Averaging over $z \in \Ball{\rho}{y}$ via $\overline u = \fint_{\Ball{\rho}{y}} u(z)\,dz$
yields
\[
|u(x) - \overline u| \le \fint_{\Ball{\rho}{y}} |u(x) - u(z)|\,dz
\le C_{\mathrm{H\ddot ol}}\,(2\rho)^{\alpha}.
\]
Multiplying by $\rho^{-\alpha}$ gives
$N_{\alpha}(u; y, \rho) \le C_{\mathrm{H\ddot ol}}\, 2^{\alpha}$.
\end{proof}

\begin{theorem}[H\"older implies Campanato seminorm bound,
\leanname{holder\_implies\_campanato}]
If $u$ is $\alpha$-H\"older on $\Ball{R}{x_0}$ with constant $C_{\mathrm{H\ddot ol}}$, then
\[
[u]_{\mathcal{L}^{\alpha}(\Ball{R}{x_0})} \le C_{\mathrm{H\ddot ol}}\, 2^{\alpha}.
\]
\end{theorem}

\begin{lstlisting}[style=Matlab-editor,basicstyle=\mlttfamily,escapechar=",mlshowsectionrules=true,mathescape=true,morecomment={[s]{\%\{}{\%\}}},language=lean,numbers=none]
/-- A Hölder continuous function has controlled Campanato seminorm. -/
theorem holder_implies_campanato
    {u : E → ℝ} {x₀ : E} {R α C_hold : ℝ}
    (hα : 0 < α) (hR : 0 < R) (hC_hold : 0 ≤ C_hold)
    (hhold : ∀ x ∈ Metric.ball x₀ R, ∀ y ∈ Metric.ball x₀ R,
      |u x - u y| ≤ C_hold * ‖x - y‖ ^ α) :
    campanatoSeminorm u x₀ R α ≤ C_hold * (2 : ℝ) ^ α := by
  exact (holder_hasCampanatoBound hα hC_hold hhold).campanatoSeminorm_le hR

\end{lstlisting}
\begin{proof}
Combine the previous lemma with
\leanname{HasCampanatoBound.campanatoSeminorm\_le}.
\end{proof}


%% ========================================================================
\section{Ball Extension, Approximation, and Supporting Tools}\label{sec:extension}
%% ========================================================================

% Section 4: Ball Extension, Approximation, and Supporting Tools.
% Faithful translation of the De Giorgi formalization.

This section collects the analytic infrastructure underlying the De Giorgi
development: a Sobolev chain rule on open sets, smooth approximation of
$W^{1,p}$ functions on the unit ball, an explicit unit-ball extension operator
together with its $L^p$ and $W^{1,p}$ estimates, the positive-part machinery,
and a small toolkit for $L^p$ limits, finite covers, and pointwise component
analysis. Throughout we work in the Euclidean space $E := \R^d$ with $d \ge 1$.

\subsection{Sobolev chain rule}

The first goal is the chain rule for compositions $\Phi \circ u$ when $u$ is a
$W^{1,2}$ function and $\Phi$ is a smooth scalar function with bounded
derivative.

\begin{theorem}[\leanname{sobolev\_chain\_rule\_unitBall}]
\label{thm:chain-rule-unitBall}
Let $u \in W^{1,2}(\Bz{1})$ with weak gradient $\nabla u$ packaged as a
\leanname{MemW1pWitness}, let $i \in \{1,\dots,d\}$, and let $\Phi \in
C^{\infty}(\R)$ satisfy $\Phi(0)=0$ and $|\Phi'(t)| \le M$ for all $t \in \R$.
Then $\partial_i \Phi(u) = \Phi'(u)\,\partial_i u$ in the weak sense on
$\Bz{1}$.
\end{theorem}

\begin{lstlisting}[style=Matlab-editor,basicstyle=\mlttfamily,escapechar=",mlshowsectionrules=true,mathescape=true,morecomment={[s]{\%\{}{\%\}}},language=lean,numbers=none]
theorem sobolev_chain_rule_unitBall
    {u : E → ℝ}
    (hw : MemW1pWitness 2 u (Metric.ball (0 : E) 1))
    (i : Fin d)
    (Φ : ℝ → ℝ) (hΦ : ContDiff ℝ (⊤ : ℕ∞) Φ) (hΦ0 : Φ 0 = 0)
    (hΦ'_bdd : ∃ M, ∀ t, |deriv Φ t| ≤ M) :
    HasWeakPartialDeriv' (d := d) i
      (fun x => deriv Φ (u x) * hw.weakGrad x i)
      (fun x => Φ (u x))
      (Metric.ball (0 : E) 1) := by
  obtain ⟨M, hM⟩ := hΦ'_bdd
  have hM_pos : 0 ≤ M := le_trans (abs_nonneg _) (hM 0)
  have hB_fin : volume (Metric.ball (0 : E) 1) < ⊤ :=
    lt_of_le_of_lt (measure_mono ball_subset_closedBall)
      (isCompact_closedBall (0 : E) 1).measure_lt_top
  obtain ⟨ψ, hψ_smooth, hψ_cs, hψ_L2, hψ_grad⟩ :=
    exists_smooth_W12_approx_on_unitBall (d := d) hw
  obtain ⟨ns, hns_mono, hns_ae⟩ := exists_ae_tendsto_of_eLpNorm_tendsto
    measurableSet_ball
    (fun n => (hψ_smooth n).continuous.aestronglyMeasurable.restrict)
    hw.memLp.aestronglyMeasurable
    hψ_L2
  intro φ hφ hφ_cs hφ_supp
  have hIBP : ∀ n,
      ∫ x in Metric.ball (0 : E) 1, Φ (ψ (ns n) x) *
        (fderiv ℝ φ x) (EuclideanSpace.single i 1) =
      -∫ x in Metric.ball (0 : E) 1, (deriv Φ (ψ (ns n) x) *
        (fderiv ℝ (ψ (ns n)) x) (EuclideanSpace.single i 1)) * φ x := by
    intro n
    exact chain_rule_smooth_ibp isOpen_ball i (hψ_smooth (ns n)) (hψ_cs (ns n)) hΦ hΦ0
      φ hφ hφ_cs hφ_supp
  -- === hLHS ===
  have hLHS : Tendsto
      (fun n => ∫ x in Metric.ball (0 : E) 1, Φ (ψ (ns n) x) *
        (fderiv ℝ φ x) (EuclideanSpace.single i 1))
      atTop
      (nhds (∫ x in Metric.ball (0 : E) 1, Φ (u x) *
        (fderiv ℝ φ x) (EuclideanSpace.single i 1))) := by
    have hΦ_lip : LipschitzWith ⟨M, hM_pos⟩ Φ :=
      lipschitzWith_of_nnnorm_deriv_le (hΦ.differentiable (by simp)) (fun t => by
        simp only [← NNReal.coe_le_coe, NNReal.coe_mk, coe_nnnorm]
        exact (Real.norm_eq_abs _).symm ▸ hM t)
    have hΦ_ptwise : ∀ n x, ‖Φ (ψ (ns n) x) - Φ (u x)‖ ≤ M * ‖ψ (ns n) x - u x‖ := by
      intro n x
      have := hΦ_lip.dist_le_mul (ψ (ns n) x) (u x)
      rwa [Real.dist_eq, Real.dist_eq, NNReal.coe_mk] at this
    set μ := volume.restrict (Metric.ball (0 : E) 1)
    have hΦψ_L2 : Tendsto (fun n => eLpNorm (fun x => Φ (ψ (ns n) x) - Φ (u x)) 2 μ)
        atTop (nhds 0) := by
      have hψ_sub : Tendsto (fun n => eLpNorm (fun x => ψ (ns n) x - u x) 2 μ)
          atTop (nhds 0) := hψ_L2.comp hns_mono.tendsto_atTop
      apply tendsto_of_tendsto_of_tendsto_of_le_of_le' tendsto_const_nhds
        (show Tendsto (fun n => ‖(M : ℝ)‖ₑ * eLpNorm (fun x => ψ (ns n) x - u x) 2 μ)
          atTop (nhds 0) from ?_)
        (Eventually.of_forall (fun _ => zero_le _))
        (Eventually.of_forall (fun n => ?_))
      · rw [show (0 : ℝ≥0∞) = ‖(M : ℝ)‖ₑ * 0 from by simp]
        exact ENNReal.Tendsto.const_mul hψ_sub (Or.inr enorm_ne_top)
      · calc eLpNorm (fun x => Φ (ψ (ns n) x) - Φ (u x)) 2 μ
            ≤ eLpNorm (fun x => M * (ψ (ns n) x - u x)) 2 μ := by
              apply eLpNorm_mono (fun x => ?_)
              calc ‖Φ (ψ (ns n) x) - Φ (u x)‖
                  ≤ M * ‖ψ (ns n) x - u x‖ := hΦ_ptwise n x
                _ = ‖M * (ψ (ns n) x - u x)‖ := by
                    simp only [Real.norm_eq_abs, abs_mul, abs_of_nonneg hM_pos]
          _ = ‖(M : ℝ)‖ₑ * eLpNorm (fun x => ψ (ns n) x - u x) 2 μ := by
              rw [show (fun x => M * (ψ (ns n) x - u x)) = M • (fun x => ψ (ns n) x - u x)
                from by ext; simp [smul_eq_mul]]
              exact eLpNorm_const_smul M _ _ _
    have hφ_deriv_cont : Continuous (fun x => (fderiv ℝ φ x) (EuclideanSpace.single i 1)) :=
      (hφ.continuous_fderiv (by simp)).clm_apply continuous_const
    have hφ_bdd : ∃ C, ∀ x, |(fderiv ℝ φ x) (EuclideanSpace.single i 1)| ≤ C := by
      obtain ⟨C, hC⟩ := (hφ_cs.fderiv_apply (𝕜 := ℝ) (EuclideanSpace.single i 1)).exists_bound_of_continuous
        hφ_deriv_cont
      exact ⟨C, fun x => by rw [← Real.norm_eq_abs]; exact hC x⟩
    exact integral_mul_tendsto_of_eLpNorm_tendsto measurableSet_ball hB_fin
      (fun n => hΦ.continuous.comp_aestronglyMeasurable
        ((hψ_smooth (ns n)).continuous.aestronglyMeasurable.restrict))
      (hΦ.continuous.comp_aestronglyMeasurable hw.memLp.aestronglyMeasurable)
      hΦψ_L2 hφ_bdd hφ_deriv_cont.aestronglyMeasurable.restrict
  -- === hRHS ===
  have hRHS : Tendsto
      (fun n => ∫ x in Metric.ball (0 : E) 1, (deriv Φ (ψ (ns n) x) *
        (fderiv ℝ (ψ (ns n)) x) (EuclideanSpace.single i 1)) * φ x)
      atTop
      (nhds (∫ x in Metric.ball (0 : E) 1,
        (deriv Φ (u x) * hw.weakGrad x i) * φ x)) := by
    obtain ⟨C_φ, hC_φ⟩ : ∃ C, ∀ x, |φ x| ≤ C := by
      obtain ⟨C, hC⟩ := hφ_cs.exists_bound_of_continuous hφ.continuous
      exact ⟨C, fun x => by rw [← Real.norm_eq_abs]; exact hC x⟩
    have hgi_int : IntegrableOn (fun x => hw.weakGrad x i)
        (Metric.ball (0 : E) 1) volume := by
      haveI : IsFiniteMeasure (volume.restrict (Metric.ball (0 : E) 1)) :=
        isFiniteMeasure_restrict.mpr (ne_top_of_lt hB_fin)
      exact (hw.weakGrad_component_memLp i).integrable one_le_two
    -- Term 1: ∫ Φ'(ψ_n)·(∂ᵢψ_n - gi)·φ → 0
    have hT1 : Tendsto
        (fun n => ∫ x in Metric.ball (0 : E) 1, deriv Φ (ψ (ns n) x) *
          ((fderiv ℝ (ψ (ns n)) x) (EuclideanSpace.single i 1) - hw.weakGrad x i) * φ x)
        atTop (nhds 0) := by
      have hgrad_L1 : Tendsto
          (fun n => ∫ x in Metric.ball (0 : E) 1,
            ‖(fderiv ℝ (ψ (ns n)) x) (EuclideanSpace.single i 1) - hw.weakGrad x i‖)
          atTop (nhds 0) :=
        eLpNorm_one_tendsto_of_eLpNorm_two_tendsto hB_fin
          (fun n => ((hψ_smooth (ns n)).continuous_fderiv (by norm_cast)
            |>.clm_apply continuous_const).aestronglyMeasurable.restrict)
          (hw.weakGrad_component_memLp i).aestronglyMeasurable
          ((hψ_grad i).comp hns_mono.tendsto_atTop)
      exact integral_bdd_mul_L1_tendsto measurableSet_ball hB_fin
        hM_pos (fun n x => hM _)
        (fun n =>
          Integrable.sub
            (((hψ_smooth (ns n)).continuous_fderiv (by norm_cast)
              |>.clm_apply continuous_const).continuousOn.integrableOn_compact
                (isCompact_closedBall (0 : E) 1) |>.mono_set ball_subset_closedBall)
            hgi_int)
        hgrad_L1 hC_φ
    -- Term 2: ∫ (Φ'(ψ_n) - Φ'(u))·gi·φ → 0
    have hT2 : Tendsto
        (fun n => ∫ x in Metric.ball (0 : E) 1,
          (deriv Φ (ψ (ns n) x) - deriv Φ (u x)) * (hw.weakGrad x i * φ x))
        atTop (nhds 0) := by
      apply integral_dominated_convergence_term measurableSet_ball
      · intro n x
        have h1 : |deriv Φ (ψ (ns n) x) - deriv Φ (u x)| ≤
            |deriv Φ (ψ (ns n) x)| + |deriv Φ (u x)| := by
          rw [← Real.norm_eq_abs, ← Real.norm_eq_abs (deriv Φ (ψ _ _)),
              ← Real.norm_eq_abs (deriv Φ (u _))]
          exact norm_sub_le _ _
        exact h1.trans (add_le_add (hM _) (hM _))
      · intro n
        have hΦ'_cont : Continuous (deriv Φ) := hΦ.continuous_deriv (by simp)
        exact ((hΦ'_cont.comp (hψ_smooth (ns n)).continuous).aestronglyMeasurable.restrict.sub
          (hΦ'_cont.comp_aestronglyMeasurable hw.memLp.aestronglyMeasurable))
      · filter_upwards [hns_ae] with x hx
        have hΦ'_cont : Continuous (deriv Φ) := hΦ.continuous_deriv (by simp)
        -- Φ'(ψ_n(x)) → Φ'(u(x)) since ψ_n(x) → u(x) and Φ' is continuous
        have h1 : Tendsto (fun k => deriv Φ (ψ (ns k) x)) atTop
            (nhds (deriv Φ (u x))) :=
          (hΦ'_cont.tendsto (u x)).comp hx
        rw [show (0 : ℝ) = deriv Φ (u x) - deriv Φ (u x) from by ring]
        exact h1.sub tendsto_const_nhds
      · exact hgi_int.mul_bdd hφ.continuous.aestronglyMeasurable.restrict
          (Eventually.of_forall (fun x => by rw [Real.norm_eq_abs]; exact hC_φ x))
    -- Algebraic rearrangement + combine T1, T2
    -- Key identity: RHS_n(x) = target(x) + T1_integrand(x) + T2_integrand(x)
    have h_split : ∀ n,
        ∫ x in Metric.ball (0 : E) 1, (deriv Φ (ψ (ns n) x) *
          (fderiv ℝ (ψ (ns n)) x) (EuclideanSpace.single i 1)) * φ x =
        (∫ x in Metric.ball (0 : E) 1, (deriv Φ (u x) * hw.weakGrad x i) * φ x) +
        (∫ x in Metric.ball (0 : E) 1, deriv Φ (ψ (ns n) x) *
          ((fderiv ℝ (ψ (ns n)) x) (EuclideanSpace.single i 1) - hw.weakGrad x i) * φ x) +
        (∫ x in Metric.ball (0 : E) 1, (deriv Φ (ψ (ns n) x) - deriv Φ (u x)) * (hw.weakGrad x i * φ x)) := by
      intro n
      -- Integrability of the three summands (each is bounded × L¹ × bounded on finite measure)
      have hΦ'_cont : Continuous (deriv Φ) := hΦ.continuous_deriv (by simp)
      have hgiφ_int := hgi_int.mul_bdd hφ.continuous.aestronglyMeasurable.restrict
          (Eventually.of_forall (fun x => by rw [Real.norm_eq_abs]; exact hC_φ x))
      have hT2_int : IntegrableOn (fun x =>
          (deriv Φ (ψ (ns n) x) - deriv Φ (u x)) * (hw.weakGrad x i * φ x))
          (Metric.ball (0 : E) 1) volume :=
        hgiφ_int.bdd_mul
          (((hΦ'_cont.comp (hψ_smooth (ns n)).continuous).aestronglyMeasurable.restrict.sub
            (hΦ'_cont.comp_aestronglyMeasurable hw.memLp.aestronglyMeasurable)))
          (Eventually.of_forall (fun x => by
            rw [Real.norm_eq_abs]
            calc |deriv Φ (ψ (ns n) x) - deriv Φ (u x)|
                = ‖deriv Φ (ψ (ns n) x) - deriv Φ (u x)‖ := (Real.norm_eq_abs _).symm
              _ ≤ ‖deriv Φ (ψ (ns n) x)‖ + ‖deriv Φ (u x)‖ := norm_sub_le _ _
              _ ≤ M + M := by
                  rw [Real.norm_eq_abs, Real.norm_eq_abs]
                  exact add_le_add (hM _) (hM _)))
      have hTarget_int : IntegrableOn (fun x =>
          (deriv Φ (u x) * hw.weakGrad x i) * φ x) (Metric.ball (0 : E) 1) volume :=
        (hgi_int.bdd_mul
          (hΦ'_cont.comp_aestronglyMeasurable hw.memLp.aestronglyMeasurable)
          (Eventually.of_forall (fun x => by rw [Real.norm_eq_abs]; exact hM _))).mul_bdd
          hφ.continuous.aestronglyMeasurable.restrict
          (Eventually.of_forall (fun x => by rw [Real.norm_eq_abs]; exact hC_φ x))
      have hT1_int : IntegrableOn (fun x =>
          deriv Φ (ψ (ns n) x) *
            ((fderiv ℝ (ψ (ns n)) x) (EuclideanSpace.single i 1) - hw.weakGrad x i) * φ x)
          (Metric.ball (0 : E) 1) volume :=
        ((Integrable.sub
            (((hψ_smooth (ns n)).continuous_fderiv (by norm_cast)
              |>.clm_apply continuous_const).continuousOn.integrableOn_compact
                (isCompact_closedBall (0 : E) 1) |>.mono_set ball_subset_closedBall)
            hgi_int).bdd_mul
          ((hΦ'_cont.comp (hψ_smooth (ns n)).continuous).aestronglyMeasurable.restrict)
          (Eventually.of_forall (fun x => by rw [Real.norm_eq_abs]; exact hM _))).mul_bdd
          hφ.continuous.aestronglyMeasurable.restrict
          (Eventually.of_forall (fun x => by rw [Real.norm_eq_abs]; exact hC_φ x))
      -- Use linearity of integral + pointwise identity (as function-level equality)
      have h_eq : (fun x => (deriv Φ (ψ (ns n) x) *
          (fderiv ℝ (ψ (ns n)) x) (EuclideanSpace.single i 1)) * φ x) =
          ((fun x => (deriv Φ (u x) * hw.weakGrad x i) * φ x) +
          (fun x => deriv Φ (ψ (ns n) x) *
            ((fderiv ℝ (ψ (ns n)) x) (EuclideanSpace.single i 1) - hw.weakGrad x i) * φ x)) +
          (fun x => (deriv Φ (ψ (ns n) x) - deriv Φ (u x)) * (hw.weakGrad x i * φ x)) := by
        ext x; simp only [Pi.add_apply]; ring
      rw [h_eq, integral_add' (hTarget_int.add hT1_int) hT2_int,
          integral_add' hTarget_int hT1_int]
    -- RHS_n = target + T1_n + T2_n, and T1_n + T2_n → 0
    -- Use sub_tendsto: show ∫ RHS_n - ∫ target = T1_n + T2_n → 0
    apply tendsto_sub_nhds_zero_iff.mp
    have h_sub_eq : ∀ n,
        (∫ x in Metric.ball (0 : E) 1, (deriv Φ (ψ (ns n) x) *
          (fderiv ℝ (ψ (ns n)) x) (EuclideanSpace.single i 1)) * φ x) -
        (∫ x in Metric.ball (0 : E) 1, (deriv Φ (u x) * hw.weakGrad x i) * φ x) =
        (∫ x in Metric.ball (0 : E) 1, deriv Φ (ψ (ns n) x) *
          ((fderiv ℝ (ψ (ns n)) x) (EuclideanSpace.single i 1) - hw.weakGrad x i) * φ x) +
        (∫ x in Metric.ball (0 : E) 1, (deriv Φ (ψ (ns n) x) - deriv Φ (u x)) * (hw.weakGrad x i * φ x)) := by
      intro n; simp only [h_split n]; ring
    refine (Filter.Tendsto.congr (fun n => (h_sub_eq n).symm) ?_)
    rw [show (0 : ℝ) = 0 + 0 from by ring]
    exact hT1.add hT2
  -- Uniqueness of limits: hIBP says LHS_n = -RHS_n, so limit of LHS = -(limit of RHS)
  have hLHS2 : Tendsto
      (fun n => ∫ x in Metric.ball (0 : E) 1, Φ (ψ (ns n) x) *
        (fderiv ℝ φ x) (EuclideanSpace.single i 1))
      atTop (nhds (-∫ x in Metric.ball (0 : E) 1,
        (deriv Φ (u x) * hw.weakGrad x i) * φ x)) := by
    have hIBP' : (fun n => ∫ x in Metric.ball (0 : E) 1, Φ (ψ (ns n) x) *
        (fderiv ℝ φ x) (EuclideanSpace.single i 1)) =
      (fun n => -∫ x in Metric.ball (0 : E) 1, (deriv Φ (ψ (ns n) x) *
        (fderiv ℝ (ψ (ns n)) x) (EuclideanSpace.single i 1)) * φ x) := by
      ext n; exact hIBP n
    rw [hIBP']
    exact hRHS.neg
  exact tendsto_nhds_unique hLHS hLHS2

/-- An explicit `W^{1,p}` witness carrying the weak gradient. -/
structure MemW1pWitness (p : ℝ≥0∞) (f : E → ℝ) (Ω : Set E)
    (μ : Measure E := volume) where
\end{lstlisting}
\begin{proof}
Let $M$ be the prescribed bound on $|\Phi'|$, so $M \ge 0$. Since $\Bz{1}$ has
finite Lebesgue measure, by \leanname{exists\_smooth\_W12\_approx\_on\_unitBall}
there is a sequence $\psi_n \in C^{\infty}_c(\R^d)$ with
\begin{align*}
\eLpNormText{\psi_n - u}_{L^2(\Bz{1})} &\to 0,\\
\eLpNormText{\partial_i \psi_n - \partial_i u}_{L^2(\Bz{1})} &\to 0.
\end{align*}
Passing to a subsequence (still denoted $\psi_n$) we may also assume $\psi_n
\to u$ a.e.\ on $\Bz{1}$ via
\leanname{exists\_ae\_tendsto\_of\_eLpNorm\_tendsto}.

Since each $\psi_n$ is smooth and compactly supported, the classical chain
rule and integration by parts yield, for every test function
$\varphi \in C^{\infty}_c(\Bz{1})$,
\begin{multline*}
\int_{\Bz{1}} \Phi(\psi_n)\, \partial_i \varphi \,dx
= -\int_{\Bz{1}} \Phi'(\psi_n)\, \partial_i \psi_n\, \varphi \,dx.
\end{multline*}
We pass to the limit on each side.

\emph{Left-hand side.} The function $\Phi$ is $M$-Lipschitz because $\Phi'$
is bounded and $\Phi$ is differentiable. Hence
$|\Phi(\psi_n(x)) - \Phi(u(x))| \le M\, |\psi_n(x)-u(x)|$ pointwise, so
\[
\eLpNormText{\Phi(\psi_n)-\Phi(u)}_{L^2(\Bz{1})}
\le M\, \eLpNormText{\psi_n - u}_{L^2(\Bz{1})}
\to 0.
\]
The function $\partial_i \varphi$ lies in $L^2(\Bz{1})$ (it is bounded and
compactly supported), so by H\"older's inequality
$\int_{\Bz{1}} \Phi(\psi_n)\,\partial_i \varphi \,dx \to
\int_{\Bz{1}} \Phi(u)\,\partial_i\varphi \,dx$.

\emph{Right-hand side.} Decompose
\[
\Phi'(\psi_n)\,\partial_i \psi_n - \Phi'(u)\,\partial_i u
= \Phi'(\psi_n)\,(\partial_i \psi_n - \partial_i u)
+ (\Phi'(\psi_n) - \Phi'(u))\,\partial_i u.
\]
The first term is bounded in $L^2$-norm by
$M\,\eLpNormText{\partial_i \psi_n - \partial_i u}_{L^2(\Bz{1})}\to 0$.
For the second term, $\Phi'$ is continuous and $\psi_n \to u$ a.e., so
$\Phi'(\psi_n) \to \Phi'(u)$ a.e.\ with uniform bound $M$. The dominated
convergence theorem applied to $|\Phi'(\psi_n) - \Phi'(u)|^2 |\partial_i u|^2$
(dominated by $4M^2 |\partial_i u|^2 \in L^1$) yields convergence of the
second term in $L^2$. Multiplying by $\varphi \in L^\infty_c$ and integrating
gives convergence of the right-hand side to
$-\int_{\Bz{1}} \Phi'(u)\,\partial_i u\, \varphi\,dx$. The chain-rule
identity follows.
\end{proof}

\begin{theorem}[\leanname{sobolev\_chain\_rule}]
\label{thm:chain-rule-open}
Let $\Omega \subseteq \R^d$ be open, $u \in W^{1,2}(\Omega)$, and let $g \in
L^2(\Omega)$ be a representative of $\partial_i u$ in the weak sense. Let
$\Phi \in C^\infty(\R)$ with $\Phi(0)=0$ and $|\Phi'| \le M$. Then
$\Phi'(u)\,g$ is a weak $i$-th partial derivative of $\Phi(u)$ on $\Omega$.
\end{theorem}

\begin{lstlisting}[style=Matlab-editor,basicstyle=\mlttfamily,escapechar=",mlshowsectionrules=true,mathescape=true,morecomment={[s]{\%\{}{\%\}}},language=lean,numbers=none]
theorem sobolev_chain_rule
    {Ω : Set E} (hΩ : IsOpen Ω)
    {u : E → ℝ} {g : E → ℝ} {i : Fin d}
    (hwd : HasWeakPartialDeriv' (d := d) i g u Ω)
    (hu : MemW1p 2 u Ω)
    (hg_Lp : MemLp g 2 (volume.restrict Ω))
    (Φ : ℝ → ℝ) (hΦ : ContDiff ℝ (⊤ : ℕ∞) Φ) (hΦ0 : Φ 0 = 0)
    (hΦ'_bdd : ∃ M, ∀ t, |deriv Φ t| ≤ M) :
    HasWeakPartialDeriv' (d := d) i (fun x => deriv Φ (u x) * g x) (fun x => Φ (u x)) Ω := by
  set hw := hu.toWitness with hw_def
  set gi := fun x => hw.weakGrad x i with gi_def
  have hgi_Lp : MemLp gi 2 (volume.restrict Ω) := hw.weakGrad_component_memLp i
  have hgi_weak : HasWeakPartialDeriv (d := d) i gi u Ω := hw.isWeakGrad i
  -- Both are weak i-derivatives of u, both in L²(Ω) ⊆ L^1_loc(Ω).
  have hg_ae : g =ᵐ[volume.restrict Ω] gi := by
    exact HasWeakPartialDeriv.ae_eq hΩ
      (show HasWeakPartialDeriv (d := d) i g u Ω by
        simpa [HasWeakPartialDeriv, HasWeakPartialDeriv'] using hwd)
      hgi_weak
      (hg_Lp.locallyIntegrable (by norm_num : (1 : ℝ≥0∞) ≤ 2))
      (hgi_Lp.locallyIntegrable (by norm_num : (1 : ℝ≥0∞) ≤ 2))
  suffices h : HasWeakPartialDeriv' (d := d) i
      (fun x => deriv Φ (u x) * gi x) (fun x => Φ (u x)) Ω by
    intro φ hφ hφ_cs hφ_supp
    have key := h φ hφ hφ_cs hφ_supp
    -- Replace gi by g in the RHS using g =ᵐ gi.
    convert key using 2
    -- Goal: ∫_Ω (Φ'(u) * g) * φ = ∫_Ω (Φ'(u) * gi) * φ
    -- Use integral_congr_ae on the restricted measure.
    apply integral_congr_ae
    -- Goal: ∀ᵐ x ∂(volume.restrict Ω), (Φ'(u x) * g x) * φ x = (Φ'(u x) * gi x) * φ x
    filter_upwards [hg_ae] with x hx
    rw [hx]
  obtain ⟨M, hM⟩ := hΦ'_bdd
  have hM_pos : 0 ≤ M := le_trans (abs_nonneg _) (hM 0)
  -- Use HasWeakPartialDeriv'_of_local: prove on every ball, then extend to Ω.
  apply HasWeakPartialDeriv'_of_local hΩ
  · -- Φ ∘ u is locally integrable on Ω
    have hu_loc : LocallyIntegrable u (volume.restrict Ω) :=
      hw.memLp.locallyIntegrable (by norm_num : (1 : ℝ≥0∞) ≤ 2)
    have hMu_loc : LocallyIntegrable (fun x => M * u x) (volume.restrict Ω) := by
      simpa [smul_eq_mul] using hu_loc.smul M
    have hΦ_lip : LipschitzWith ⟨M, hM_pos⟩ Φ :=
      lipschitzWith_of_nnnorm_deriv_le (hΦ.differentiable (by simp)) (fun t => by
        simp only [← NNReal.coe_le_coe, NNReal.coe_mk, coe_nnnorm]
        exact (Real.norm_eq_abs _).symm ▸ hM t)
    have hΦ_abs_le : ∀ t, |Φ t| ≤ M * |t| := by
      intro t
      have ht := hΦ_lip.dist_le_mul t 0
      simp [hΦ0, Real.norm_eq_abs] at ht
      exact ht
    exact hMu_loc.mono
      (hΦ.continuous.comp_aestronglyMeasurable hw.memLp.aestronglyMeasurable)
      (Eventually.of_forall fun x => by
        rw [Real.norm_eq_abs]
        simpa [Real.norm_eq_abs, abs_mul, abs_of_nonneg hM_pos] using hΦ_abs_le (u x))
  · -- Φ'(u) · gi is locally integrable on Ω
    have hgi_loc : LocallyIntegrable gi (volume.restrict Ω) :=
      hgi_Lp.locallyIntegrable (by norm_num : (1 : ℝ≥0∞) ≤ 2)
    have hMgi_loc : LocallyIntegrable (fun x => M * gi x) (volume.restrict Ω) := by
      simpa [smul_eq_mul] using hgi_loc.smul M
    exact hMgi_loc.mono
      (((hΦ.continuous_deriv (by simp)).comp_aestronglyMeasurable
        hw.memLp.aestronglyMeasurable).mul hgi_Lp.aestronglyMeasurable)
      (Eventually.of_forall fun x => by
        calc
          ‖deriv Φ (u x) * gi x‖ = |deriv Φ (u x)| * |gi x| := by
            rw [Real.norm_eq_abs, abs_mul]
          _ ≤ M * |gi x| := mul_le_mul_of_nonneg_right (hM (u x)) (abs_nonneg _)
          _ = ‖M * gi x‖ := by
            rw [Real.norm_eq_abs, abs_mul, abs_of_nonneg hM_pos])
  intro x₀ _hx₀ r hr hball
  -- On ball x₀ r ⊆ Ω, restrict the witness and apply sobolev_chain_rule_ball.
  let hw_ball := MemW1pWitness.restrict (d := d) isOpen_ball hball hw
  -- hw_ball.weakGrad = hw.weakGrad (same function, restricted domain)
  -- So hw_ball.weakGrad · i = gi
  -- Apply chain rule on the ball
  -- Apply chain rule on the ball.
  -- sobolev_chain_rule_ball gives:
  --   HasWeakPartialDeriv' i (fun x => deriv Φ (u x) * hw_ball.weakGrad x i) (Φ ∘ u) (ball x₀ r)
  -- Since hw_ball = MemW1pWitness.restrict isOpen_ball hball hw,
  -- hw_ball.weakGrad = hw.weakGrad, so hw_ball.weakGrad x i = gi x.
  -- We unfold to show the two functions are the same.
  exact sobolev_chain_rule_ball hr hw_ball (i := i) Φ hΦ hΦ0 ⟨M, hM⟩

\end{lstlisting}
\begin{proof}
Let $\nabla u$ be the canonical weak gradient extracted from the witness
$u \in W^{1,2}(\Omega)$, and write $g_i := (\nabla u)_i$. By uniqueness of
weak partial derivatives in $L^1_{\mathrm{loc}}$, we have $g = g_i$ a.e.\ on
$\Omega$, so it suffices to prove the statement with $g_i$ in place of $g$
(the difference of integrands vanishes a.e.\ in the test integral).

Local integrability of $\Phi(u)$ and of $\Phi'(u)\,g_i$ follows from the
$M$-Lipschitz bound $|\Phi(t)| \le M|t|$ and from $|\Phi'(u)\,g_i| \le M|g_i|$,
combined with the local integrability of $u$ and $g_i$.

It now suffices to verify the weak derivative identity locally, that is, on
each ball $\Ball{r}{x_0} \subseteq \Omega$
(\leanname{HasWeakPartialDeriv'\_of\_local}). On such a ball, restrict the
witness to obtain $u \in W^{1,2}(\Ball{r}{x_0})$ and apply
Theorem~\ref{thm:chain-rule-unitBall} (after rescaling to the unit ball via
\leanname{MemW1pWitness.rescale\_to\_unitBall}). The conclusion is the chain
rule for $\Phi(u)$ on $\Ball{r}{x_0}$, and gluing across all balls produces
the desired global identity on $\Omega$.
\end{proof}

\subsection{Smooth approximation on the unit ball}

We now produce smooth, compactly supported approximations of any $W^{1,p}$
function on the unit ball, with control on both function and gradient errors.
The construction proceeds in three layers: (i) a Lipschitz cutoff profile
$\eta_{r,R}$, (ii) inward dilation, and (iii) convolution with a mollifier.

\begin{definition}[\leanname{myCutoff}]
For $x_0 \in E$ and $0 < r < R$, define the radial cutoff profile
\[
\eta_{x_0,r,R}(x) := \mathrm{st}\!\left(\frac{R - \|x - x_0\|}{R - r}\right),
\]
where $\mathrm{st}$ is Mathlib's \emph{smooth transition function}, satisfying
$0 \le \mathrm{st} \le 1$, $\mathrm{st}(t) = 1$ for $t \ge 1$, and
$\mathrm{st}(t) = 0$ for $t \le 0$.
\end{definition}

\begin{lstlisting}[style=Matlab-editor,basicstyle=\mlttfamily,escapechar=",mlshowsectionrules=true,mathescape=true,morecomment={[s]{\%\{}{\%\}}},language=lean,numbers=none]
def myCutoff (x₀ : E) (r R : ℝ) (x : E) : ℝ :=
  Real.smoothTransition ((R - ‖x - x₀‖) / (R - r))

\end{lstlisting}
\begin{lemma}[\leanname{smoothTransition\_nnnorm\_deriv\_bounded}, \leanname{Mst}]
There is a constant $\Mst < \infty$ with $\|\mathrm{st}'(t)\| \le \Mst$ for
all $t \in \R$. Consequently $\mathrm{st}$ is $\Mst$-Lipschitz.
\end{lemma}

\begin{lstlisting}[style=Matlab-editor,basicstyle=\mlttfamily,escapechar=",mlshowsectionrules=true,mathescape=true,morecomment={[s]{\%\{}{\%\}}},language=lean,numbers=none]
private lemma smoothTransition_nnnorm_deriv_bounded :
    ∃ C : ℝ≥0, ∀ x : ℝ, ‖deriv smoothTransition x‖₊ ≤ C := by
  have h1 : ContDiff ℝ 1 smoothTransition := smoothTransition.contDiff
  have h_cont : Continuous (deriv smoothTransition) := h1.continuous_deriv_one
  obtain ⟨M, -, hM_max⟩ := (isCompact_Icc (a := (0 : ℝ)) (b := 1)).exists_isMaxOn
    (nonempty_Icc.2 zero_le_one) h_cont.norm.continuousOn
  have hbound : ∀ x : ℝ, ‖deriv smoothTransition x‖ ≤ ‖deriv smoothTransition M‖ := by
    intro x
    by_cases hx0 : x < 0
    · rw [st_deriv_zero_neg hx0, norm_zero]
      exact norm_nonneg _
    · by_cases hx1 : 1 < x
      · rw [st_deriv_zero_gt hx1, norm_zero]
        exact norm_nonneg _
      · push_neg at hx0 hx1
        exact Filter.eventually_principal.mp hM_max x (Set.mem_Icc.2 ⟨hx0, hx1⟩)
  refine ⟨⟨‖deriv smoothTransition M‖, norm_nonneg _⟩, fun x => ?_⟩
  rw [← NNReal.coe_le_coe, NNReal.coe_mk, coe_nnnorm]
  exact hbound x

/-- Universal derivative bound for smooth cutoff profiles. -/
noncomputable def Mst : ℝ≥0 := smoothTransition_nnnorm_deriv_bounded.choose

\end{lstlisting}
\begin{proof}
The derivative $\mathrm{st}'$ is continuous and supported in $[0,1]$, hence
attains a maximum on this compact interval; outside $[0,1]$ it vanishes. Set
$\Mst$ to that maximum value.
\end{proof}

\begin{lemma}[\leanname{myCutoff\_eq\_one}, \leanname{myCutoff\_eq\_zero}, \leanname{myCutoff\_lipschitz}, \leanname{myCutoff\_fderiv\_bound}]
\label{lem:cutoff-properties}
Let $0 < r < R$. Then:
\begin{enumerate}
\item $\eta_{x_0,r,R}(x) = 1$ for $x \in \Ball{r}{x_0}$;
\item $\eta_{x_0,r,R}(x) = 0$ for $x \notin \overline{\Ball{R}{x_0}}$;
\item $\eta_{x_0,r,R}$ is $C^\infty$ on $E$;
\item $\eta_{x_0,r,R}$ is $\Mst (R-r)^{-1}$-Lipschitz;
\item $\|D\eta_{x_0,r,R}(x)\| \le \Mst (R-r)^{-1}$ for all $x \in E$.
\end{enumerate}
\end{lemma}

\begin{lstlisting}[style=Matlab-editor,basicstyle=\mlttfamily,escapechar=",mlshowsectionrules=true,mathescape=true,morecomment={[s]{\%\{}{\%\}}},language=lean,numbers=none]
omit [NeZero d] in
lemma myCutoff_eq_one {x₀ : E} {r R : ℝ} (hrR : r < R) {x : E}
    (hx : x ∈ Metric.ball x₀ r) : myCutoff x₀ r R x = 1 := by
  unfold myCutoff
  apply Real.smoothTransition.one_of_one_le
  rw [le_div_iff₀ (sub_pos.2 hrR)]
  have : ‖x - x₀‖ < r := by rwa [← dist_eq_norm, ← Metric.mem_ball]
  linarith

omit [NeZero d] in
lemma myCutoff_eq_zero {x₀ : E} {r R : ℝ} (hrR : r < R) {x : E}
    (hx : x ∉ Metric.closedBall x₀ R) : myCutoff x₀ r R x = 0 := by
  unfold myCutoff
  apply Real.smoothTransition.zero_of_nonpos
  have hx' : R < dist x x₀ := by
    rw [Metric.mem_closedBall, not_le] at hx
    exact hx
  have hx'' : R < ‖x - x₀‖ := by
    simpa [dist_eq_norm] using hx'
  exact div_nonpos_of_nonpos_of_nonneg (by linarith) (by linarith)

omit [NeZero d] in
lemma myCutoff_lipschitz (x₀ : E) {r R : ℝ} (hrR : r < R) :
    LipschitzWith (Mst * ⟨(R - r)⁻¹, inv_nonneg.mpr (sub_pos.2 hrR).le⟩)
      (myCutoff x₀ r R) :=
  smoothTransition_lipschitz.comp (radial_lipschitz x₀ hrR)

omit [NeZero d] in
lemma myCutoff_fderiv_bound (x₀ : E) {r R : ℝ} (hrR : r < R) (x : E) :
    ‖fderiv ℝ (myCutoff x₀ r R) x‖ ≤ ↑Mst * (R - r)⁻¹ := by
  have h : ‖fderiv ℝ (myCutoff x₀ r R) x‖ ≤
      ↑(Mst * ⟨(R - r)⁻¹, inv_nonneg.mpr (sub_nonneg.mpr hrR.le)⟩) :=
    norm_fderiv_le_of_lipschitz (𝕜 := ℝ) (x₀ := x) (myCutoff_lipschitz x₀ hrR)
  rwa [NNReal.coe_mul, NNReal.coe_mk] at h

\end{lstlisting}
\begin{proof}
For (i): if $x \in \Ball{r}{x_0}$ then $\|x - x_0\| < r$ and so
$(R - \|x - x_0\|)/(R-r) > 1$, hence $\mathrm{st}$ takes the value $1$.
For (ii): if $\|x-x_0\| > R$ then the argument is negative, hence
$\mathrm{st}$ vanishes. Smoothness follows from
$\mathrm{st} \in C^\infty(\R)$ together with smoothness of the radial
argument away from $x_0$, and constancy near $x_0$. The radial map
$x \mapsto (R - \|x-x_0\|)/(R-r)$ is $1/(R-r)$-Lipschitz, so the composition
with the $\Mst$-Lipschitz $\mathrm{st}$ gives (iv); (v) follows by the
operator-norm bound for derivatives of Lipschitz maps.
\end{proof}

\begin{definition}[\leanname{Cutoff}]
A \emph{cutoff structure} on $\Ball{r}{x_0}$ inside $\Ball{R}{x_0}$ is a tuple
$\eta = (\eta, \mathrm{smooth}, \text{nonneg}, \text{le\_one}, \text{eq\_one},
\text{support}, \text{grad\_bound})$ where $\eta : E \to \R$ is
$C^\infty$, takes values in $[0,1]$, equals $1$ on $\Ball{r}{x_0}$, has
$\tsupport \eta \subseteq \overline{\Ball{R}{x_0}}$, and satisfies the
gradient bound $\|D\eta(x)\| \le \Mst (R-r)^{-1}$.
\end{definition}

\begin{lstlisting}[style=Matlab-editor,basicstyle=\mlttfamily,escapechar=",mlshowsectionrules=true,mathescape=true,morecomment={[s]{\%\{}{\%\}}},language=lean,numbers=none]
structure Cutoff (x₀ : E) (r R : ℝ) where
  toFun : E → ℝ
  smooth : ContDiff ℝ (⊤ : ℕ∞) toFun
  nonneg : ∀ x, 0 ≤ toFun x
  le_one : ∀ x, toFun x ≤ 1
  eq_one : ∀ x ∈ Metric.ball x₀ r, toFun x = 1
  support_subset : tsupport toFun ⊆ Metric.closedBall x₀ R
  grad_bound : ∀ x, ‖fderiv ℝ toFun x‖ ≤ ↑Mst * (R - r)⁻¹

\end{lstlisting}
\begin{theorem}[\leanname{Cutoff.exists}]
For every $0 < r < R$, the set of cutoff structures on $\Ball{r}{x_0}$ inside
$\Ball{R}{x_0}$ is nonempty.
\end{theorem}

\begin{lstlisting}[style=Matlab-editor,basicstyle=\mlttfamily,escapechar=",mlshowsectionrules=true,mathescape=true,morecomment={[s]{\%\{}{\%\}}},language=lean,numbers=none]
omit [NeZero d] in
theorem Cutoff.exists (x₀ : E) {r R : ℝ} (hr : 0 < r) (hR : r < R) :
    ∃ _η : Cutoff x₀ r R, True :=
  ⟨⟨myCutoff x₀ r R,
    by simpa using (myCutoff_contDiff (x₀ := x₀) hr hR),
    myCutoff_nonneg x₀ r R,
    myCutoff_le_one x₀ r R,
    fun x hx => myCutoff_eq_one (x₀ := x₀) hR hx,
    myCutoff_support_subset x₀ hR,
    fun x => myCutoff_fderiv_bound x₀ hR x⟩, trivial⟩

\end{lstlisting}
\begin{proof}
The bundle constructed from $\eta_{x_0,r,R}$ together with the lemmas of
\ref{lem:cutoff-properties} is such a cutoff structure.
\end{proof}

\begin{theorem}[\leanname{Lp\_approx\_by\_continuous\_compactly\_supported}]
\label{thm:lp-approx-cc}
Let $1 \le p < \infty$ and $f \in \Lp{p}(\R^d)$. For every $\delta > 0$ there
exists a continuous compactly supported function $g$ with
$\eLpNormText{f-g}_{L^p} < \delta$.
\end{theorem}

\begin{lstlisting}[style=Matlab-editor,basicstyle=\mlttfamily,escapechar=",mlshowsectionrules=true,mathescape=true,morecomment={[s]{\%\{}{\%\}}},language=lean,numbers=none]
omit [NeZero d] in
theorem Lp_approx_by_continuous_compactly_supported
    {f : E → ℝ} {p : ℝ≥0∞}
    (hp : 1 ≤ p) (hp' : p ≠ ⊤)
    (hf : MemLp f p volume)
    {δ : ℝ≥0∞} (hδ : 0 < δ) :
    ∃ g : E → ℝ, Continuous g ∧ HasCompactSupport g ∧
      eLpNorm (f - g) p volume < δ := by
  obtain ⟨ε, hε_pos, hε_lt⟩ : ∃ ε : ℝ, 0 < ε ∧ ENNReal.ofReal ε < δ := by
    by_cases hδ_top : δ = ⊤
    · exact ⟨1, one_pos, hδ_top ▸ ENNReal.ofReal_lt_top⟩
    · have hδ_lt_top : δ < ⊤ := lt_top_iff_ne_top.mpr hδ_top
      refine ⟨(δ / 2).toReal,
        ENNReal.toReal_pos (ENNReal.div_pos hδ.ne' (by norm_num)).ne'
          (ENNReal.div_lt_top hδ_lt_top.ne (by norm_num)).ne, ?_⟩
      calc
        ENNReal.ofReal (δ / 2).toReal = δ / 2 :=
          ENNReal.ofReal_toReal (ENNReal.div_lt_top hδ_lt_top.ne (by norm_num)).ne
        _ < δ := ENNReal.half_lt_self hδ.ne' hδ_top
  obtain ⟨g, hg_supp, hg_smooth, hg_norm⟩ := hf.exist_eLpNorm_sub_le hp' hp hε_pos
  exact ⟨g, hg_smooth.continuous, hg_supp, lt_of_le_of_lt hg_norm hε_lt⟩


\end{lstlisting}
\begin{proof}
This is the standard density of $C_c$ in $L^p$ for $p < \infty$. Choose
$0 < \varepsilon < \delta$ in $\R$, apply the Mathlib lemma
\leanname{exist\_eLpNorm\_sub\_le} to obtain a smooth compactly supported
$g$ with $\eLpNormText{f-g}_{L^p} \le \varepsilon < \delta$.
\end{proof}

\subsubsection*{Rescaling}

\begin{definition}[\leanname{unitBallDilate}]
For $\lambda > 0$ and $u : E \to \R$, define
$\unitBallDilate_\lambda u(x) := u(\lambda^{-1} x)$.
\end{definition}

\begin{lstlisting}[style=Matlab-editor,basicstyle=\mlttfamily,escapechar=",mlshowsectionrules=true,mathescape=true,morecomment={[s]{\%\{}{\%\}}},language=lean,numbers=none]
/-- The dilation used in the star-shaped approximation argument. -/
def unitBallDilate (lam : ℝ) (u : E → ℝ) : E → ℝ :=
  fun x => u (lam⁻¹ • x)


\end{lstlisting}
\begin{lemma}[\leanname{smul\_inv\_mem\_unitBall}]
If $\lambda > 1$ and $x \in \Bz{1}$, then $\lambda^{-1} x \in \Bz{1}$.
\end{lemma}

\begin{lstlisting}[style=Matlab-editor,basicstyle=\mlttfamily,escapechar=",mlshowsectionrules=true,mathescape=true,morecomment={[s]{\%\{}{\%\}}},language=lean,numbers=none]
omit [NeZero d] in
lemma smul_inv_mem_unitBall {lam : ℝ} (hlam : 1 < lam) {x : E}
    (hx : x ∈ Metric.ball (0 : E) 1) :
    lam⁻¹ • x ∈ Metric.ball (0 : E) 1 := by
  rw [Metric.mem_ball, dist_zero_right] at hx ⊢
  have hlam_pos : 0 < lam := lt_trans zero_lt_one hlam
  rw [norm_smul, Real.norm_of_nonneg (inv_nonneg.mpr hlam_pos.le)]
  have hinv_lt_one : lam⁻¹ < 1 := by
    exact inv_lt_one_of_one_lt₀ hlam
  have hmul : lam⁻¹ * ‖x‖ ≤ ‖x‖ := by
    exact mul_le_of_le_one_left (norm_nonneg x) hinv_lt_one.le
  exact lt_of_le_of_lt hmul hx

\end{lstlisting}
\begin{proof}
$\|\lambda^{-1}x\| = \lambda^{-1}\|x\| \le \|x\| < 1$ since $\lambda^{-1} < 1$.
\end{proof}

\begin{lemma}[\leanname{MemW1pWitness.rescale\_to\_unitBall}]
\label{lem:rescale-witness}
Let $\hw \in W^{1,p}(\Ball{R}{x_0})$ with weak gradient $G$. Then the
function $z \mapsto u(x_0 + Rz)$ is in $W^{1,p}(\Bz{1})$ with weak
gradient $z \mapsto R\,G(x_0 + Rz)$.
\end{lemma}

\begin{lstlisting}[style=Matlab-editor,basicstyle=\mlttfamily,escapechar=",mlshowsectionrules=true,mathescape=true,morecomment={[s]{\%\{}{\%\}}},language=lean,numbers=none]
/-- Rescaling a witness on `B_R(x₀)` to a witness on the unit ball. -/
noncomputable def MemW1pWitness.rescale_to_unitBall
    {p : ℝ≥0∞} {x₀ : E} {R : ℝ} {u : E → ℝ}
    (hR : 0 < R)
    (hw : MemW1pWitness p u (Metric.ball x₀ R)) :
    MemW1pWitness p (fun z => u (x₀ + R • z)) (Metric.ball (0 : E) 1) where
  memLp := by
    exact rescale_memLp_helper hR hw.memLp
  weakGrad := fun z => R • hw.weakGrad (x₀ + R • z)
  weakGrad_component_memLp := by
    intro i
    simpa [Pi.smul_apply, smul_eq_mul] using
      (rescale_memLp_helper (x₀ := x₀) (R := R) hR (hw.weakGrad_component_memLp i)).const_mul R
  isWeakGrad := by
    intro i
    simpa [Pi.smul_apply, smul_eq_mul] using
      weakPartialDeriv_rescale_to_unitBall (x₀ := x₀) (R := R) (u := u)
        (g := fun x => hw.weakGrad x i) hR (hw.isWeakGrad i)

\end{lstlisting}
\begin{proof}
Let $T(z) := x_0 + Rz$, a measurable embedding with
$T^{-1}(\Ball{R}{x_0}) = \Bz{1}$ and pushforward measure
$T_\sharp \mathrm{vol} = R^{-d}\, \mathrm{vol}$. The change of variables
yields $\eLpNormText{u\circ T}_{L^p(\Bz{1})}^p = R^{-d}
\eLpNormText{u}_{L^p(\Ball{R}{x_0})}^p$, hence $u\circ T \in L^p(\Bz{1})$.
For the weak derivative, let $\varphi \in C^\infty_c(\Bz{1})$ and define
$\psi(x) := \varphi(R^{-1}(x - x_0))$. Then $\psi \in C^\infty_c(\Ball{R}{x_0})$
and $\partial_i \varphi(z) = R\, \partial_i \psi(x_0 + Rz)$. Substituting
into the test definition for $G_i$ at $\psi$, performing the change of
variables, and dividing by the Jacobian yields the test definition with
$z \mapsto R\,G_i(x_0+Rz)$ at $\varphi$.
\end{proof}

\begin{lemma}[\leanname{eLpNorm\_rescale\_to\_unitBall}]
For $0 < R$, $p \in [1,\infty)$,
\[
\eLpNormText{u(x_0 + R\cdot)}_{L^p(\Bz{1})}
= R^{-d/p}\,\eLpNormText{u}_{L^p(\Ball{R}{x_0})}.
\]
\end{lemma}

\begin{lstlisting}[style=Matlab-editor,basicstyle=\mlttfamily,escapechar=",mlshowsectionrules=true,mathescape=true,morecomment={[s]{\%\{}{\%\}}},language=lean,numbers=none]
omit [NeZero d] in
theorem eLpNorm_rescale_to_unitBall
    {p : ℝ≥0∞} {x₀ : E} {R : ℝ} {f : E → ℝ}
    (hR : 0 < R) (_hp : p ≠ 0) (hp' : p ≠ ⊤) :
    eLpNorm (fun z => f (x₀ + R • z)) p (volume.restrict (Metric.ball (0 : E) 1)) =
      ENNReal.ofReal (R⁻¹ ^ (d / p.toReal)) *
        eLpNorm f p (volume.restrict (Metric.ball x₀ R)) := by
  set T := fun z : E => x₀ + R • z with hT_def
  have hR' : R ≠ 0 := hR.ne'
  have hT_emb : MeasurableEmbedding T :=
    ((MeasurableEquiv.addLeft x₀).measurableEmbedding).comp
      ((Homeomorph.smul (isUnit_iff_ne_zero.2 hR').unit).toMeasurableEquiv.measurableEmbedding)
  rw [show (fun z => f (x₀ + R • z)) = f ∘ T from rfl, ← hT_emb.eLpNorm_map_measure]
  have hpreimage : T ⁻¹' (Metric.ball x₀ R) = Metric.ball (0 : E) 1 := by
    ext z
    simp only [T, mem_preimage, Metric.mem_ball, dist_comm, dist_eq_norm]
    constructor
    · intro h
      have : x₀ - (x₀ + R • z) = -(R • z) := by abel
      rw [this, norm_neg, norm_smul, Real.norm_of_nonneg hR.le] at h
      simp only [sub_zero]
      rwa [show R * ‖z‖ < R ↔ ‖z‖ < 1 from by constructor <;> intro hh <;> nlinarith] at h
    · intro h
      simp only [sub_zero] at h
      have : x₀ - (x₀ + R • z) = -(R • z) := by abel
      rw [this, norm_neg, norm_smul, Real.norm_of_nonneg hR.le]
      nlinarith
  conv_lhs => rw [show Metric.ball (0 : E) 1 = T ⁻¹' Metric.ball x₀ R from hpreimage.symm]
  rw [← Measure.restrict_map hT_emb.measurable measurableSet_ball]
  have hmap : Measure.map T (volume : Measure E) =
      ENNReal.ofReal (|R ^ Module.finrank ℝ E|⁻¹) • volume := by
    rw [show T = (fun z => x₀ + z) ∘ (fun z => R • z) from rfl]
    rw [← Measure.map_map (measurable_const_add x₀) (measurable_const_smul R)]
    rw [Measure.map_addHaar_smul volume hR']
    rw [Measure.map_smul, (measurePreserving_add_left volume x₀).map_eq, abs_inv]
  rw [hmap, Measure.restrict_smul, eLpNorm_smul_measure_of_ne_top hp']
  simp only [smul_eq_mul]
  congr 1
  have hfin : Module.finrank ℝ E = d := by simp
  rw [hfin]
  have hRd_pos : (0 : ℝ) < R ^ d := pow_pos hR d
  rw [abs_of_pos hRd_pos]
  rw [show (R ^ d)⁻¹ = R⁻¹ ^ d from (inv_pow R d).symm]
  have hRinv_nonneg : (0 : ℝ) ≤ R⁻¹ := inv_nonneg.mpr hR.le
  rw [ENNReal.ofReal_rpow_of_nonneg (pow_nonneg hRinv_nonneg d) ENNReal.toReal_nonneg]
  congr 1
  rw [← Real.rpow_natCast (R⁻¹) d, ← Real.rpow_mul hRinv_nonneg]
  congr 1
  rw [ENNReal.toReal_div, ENNReal.toReal_one]
  ring


\end{lstlisting}
\begin{proof}
Direct application of the change-of-variables formula for the affine map
$T$ above, using $|\det DT| = R^d$.
\end{proof}

\subsubsection*{Inward and outward dilation}

\begin{definition}[\leanname{MemW1pWitness.unitBallDilate}]
Let $\lambda > 1$ and $u \in W^{1,p}(\Bz{1})$. The inward dilate
$\unitBallDilate_\lambda u(x) = u(\lambda^{-1}x)$ defines a function in
$W^{1,p}(\Bz{1})$ obtained by restricting $u$ to $\Ball{\lambda^{-1}}{0}$
and rescaling by Lemma~\ref{lem:rescale-witness}.
\end{definition}

\begin{lstlisting}[style=Matlab-editor,basicstyle=\mlttfamily,escapechar=",mlshowsectionrules=true,mathescape=true,morecomment={[s]{\%\{}{\%\}}},language=lean,numbers=none]
/-- Dilation inward on the unit ball preserves `W^{1,p}` membership.

This is the witness-level version of the standard star-shaped-domain pullback
`x ↦ λ⁻¹ x` for `λ > 1`. -/
noncomputable def MemW1pWitness.unitBallDilate
    {p : ℝ≥0∞} {u : E → ℝ} {lam : ℝ}
    (hlam : 1 < lam)
    (hw : MemW1pWitness p u (Metric.ball (0 : E) 1)) :
    MemW1pWitness p (unitBallDilate (d := d) lam u) (Metric.ball (0 : E) 1) := by
  let hwSmall : MemW1pWitness p u (Metric.ball (0 : E) lam⁻¹) :=
    hw.restrict isOpen_ball (by
      intro x hx
      rw [Metric.mem_ball, dist_zero_right] at hx ⊢
      have hlt : lam⁻¹ < 1 := by
        exact inv_lt_one_of_one_lt₀ hlam
      exact lt_trans hx hlt)
  simpa [DeGiorgi.unitBallDilate] using
    (MemW1pWitness.rescale_to_unitBall (d := d) (x₀ := (0 : E))
      (R := lam⁻¹) (u := u) (by positivity) hwSmall)

\end{lstlisting}
\begin{theorem}[\leanname{exists\_unitBallCutoff\_inside}]
Let $\lambda > 1$. There exists a cutoff structure $\eta$ on
$\Ball{1}{0}$ inside $\Ball{(1+\lambda)/2}{0}$ with
$\tsupport \eta \subseteq \Ball{\lambda}{0}$.
\end{theorem}

\begin{lstlisting}[style=Matlab-editor,basicstyle=\mlttfamily,escapechar=",mlshowsectionrules=true,mathescape=true,morecomment={[s]{\%\{}{\%\}}},language=lean,numbers=none]
omit [NeZero d] in
/-- A cutoff equal to `1` on the unit ball and supported strictly inside `B(0, lam)`. -/
theorem exists_unitBallCutoff_inside {lam : ℝ} (hlam : 1 < lam) :
    ∃ η : Cutoff (0 : E) 1 ((1 + lam) / 2),
      tsupport η.toFun ⊆ Metric.ball (0 : E) lam := by
  obtain ⟨η, _⟩ := Cutoff.exists (d := d) (0 : E)
    (r := 1) (R := (1 + lam) / 2) zero_lt_one (one_lt_midpoint_of_one_lt hlam)
  refine ⟨η, ?_⟩
  exact η.support_subset.trans <| Metric.closedBall_subset_ball (midpoint_lt_of_one_lt hlam)

\end{lstlisting}
\begin{proof}
Apply \leanname{Cutoff.exists} with $r=1$, $R=(1+\lambda)/2$. The support
inclusion follows from $(1+\lambda)/2 < \lambda$.
\end{proof}

\begin{lemma}[\leanname{MemW1pWitness.unitBallDilate\_largeBall}]
\label{lem:dilate-large}
Let $\lambda > 1$ and $u \in W^{1,p}(\Bz{1})$. Then
$\unitBallDilate_\lambda u \in W^{1,p}(\Ball{\lambda}{0})$ with weak
gradient $x \mapsto \lambda^{-1}(\nabla u)(\lambda^{-1}x)$.
\end{lemma}

\begin{lstlisting}[style=Matlab-editor,basicstyle=\mlttfamily,escapechar=",mlshowsectionrules=true,mathescape=true,morecomment={[s]{\%\{}{\%\}}},language=lean,numbers=none]
/-- Outward dilation from the unit ball to the larger ball `B(0, lam)`. -/
noncomputable def MemW1pWitness.unitBallDilate_largeBall
    {p : ℝ} (_hp : 1 ≤ p)
    {u : E → ℝ} {lam : ℝ}
    (hlam : 1 < lam)
    (hw : MemW1pWitness (ENNReal.ofReal p) u (Metric.ball (0 : E) 1)) :
    MemW1pWitness (ENNReal.ofReal p) (DeGiorgi.unitBallDilate (d := d) lam u) (Metric.ball (0 : E) lam) where
  memLp := by
    let S : E → E := fun x => lam⁻¹ • x
    have hS_emb : MeasurableEmbedding S :=
      (Homeomorph.smulOfNeZero lam⁻¹ (inv_ne_zero (show lam ≠ 0 from ne_of_gt (lt_trans zero_lt_one hlam)))).toMeasurableEquiv.measurableEmbedding
    have hmap := unitBallDilate_map_measure (d := d) (show 0 < lam from lt_trans zero_lt_one hlam)
    have hu_map :
        MemLp u (ENNReal.ofReal p)
          (Measure.map S (volume.restrict (Metric.ball (0 : E) lam))) := by
      rw [hmap]
      exact hw.memLp.smul_measure ENNReal.ofReal_ne_top
    simpa [S, DeGiorgi.unitBallDilate] using (hS_emb.memLp_map_measure_iff).1 hu_map
  weakGrad := fun x => lam⁻¹ • hw.weakGrad (lam⁻¹ • x)
  weakGrad_component_memLp := by
    intro i
    let S : E → E := fun x => lam⁻¹ • x
    have hS_emb : MeasurableEmbedding S :=
      (Homeomorph.smulOfNeZero lam⁻¹ (inv_ne_zero (show lam ≠ 0 from ne_of_gt (lt_trans zero_lt_one hlam)))).toMeasurableEquiv.measurableEmbedding
    have hmap := unitBallDilate_map_measure (d := d) (show 0 < lam from lt_trans zero_lt_one hlam)
    have hgi_map :
        MemLp (fun x => lam⁻¹ * hw.weakGrad x i) (ENNReal.ofReal p)
          (Measure.map S (volume.restrict (Metric.ball (0 : E) lam))) := by
      rw [hmap]
      exact ((hw.weakGrad_component_memLp i).const_mul lam⁻¹).smul_measure ENNReal.ofReal_ne_top
    simpa [S, smul_eq_mul] using (hS_emb.memLp_map_measure_iff).1 hgi_map
  isWeakGrad := by
    intro i
    simpa [DeGiorgi.unitBallDilate, smul_eq_mul] using
      weakPartialDeriv_unitBallDilate_to_ball (d := d)
        (i := i) (lam := lam) (u := u) (g := fun x => hw.weakGrad x i)
        (show 0 < lam from lt_trans zero_lt_one hlam) (hw.isWeakGrad i)

\end{lstlisting}
\begin{proof}
Use the measurable embedding $S(x) := \lambda^{-1}x$, push forward via
$\Ball{\lambda}{0}$, and apply the analogous test-function calculation as
in Lemma~\ref{lem:rescale-witness} with the radial scaling factor
$\lambda^{-1}$.
\end{proof}

\subsubsection*{One-shot smooth approximation}

\begin{theorem}[\leanname{exists\_smooth\_W1p\_oneShot\_on\_unitBall}]
\label{thm:oneshot}
Let $1 < p < \infty$, $u \in W^{1,p}(\Bz{1})$, and $\varepsilon > 0$. Then
there exists $\psi \in C^\infty_c(E)$ with
\begin{align*}
\eLpNormText{\psi - u}_{L^p(\Bz{1})} &< \varepsilon, \\
\eLpNormText{\partial_i \psi - \partial_i u}_{L^p(\Bz{1})} &< \varepsilon
\quad (i=1,\dots,d).
\end{align*}
\end{theorem}

\begin{lstlisting}[style=Matlab-editor,basicstyle=\mlttfamily,escapechar=",mlshowsectionrules=true,mathescape=true,morecomment={[s]{\%\{}{\%\}}},language=lean,numbers=none]
/-- One-shot `W^{1,p}` approximation on the unit ball.

This is the core density statement we want. For every `ε > 0`, there is a
globally smooth approximant whose function and weak-gradient errors are all
smaller than `ε` on `B(0,1)`. -/
theorem exists_smooth_W1p_oneShot_on_unitBall
    {p : ℝ} (hp : 1 < p) {u : E → ℝ}
    (hw : MemW1pWitness (ENNReal.ofReal p) u (Metric.ball (0 : E) 1))
    {ε : ℝ} (hε : 0 < ε) :
    ∃ ψ : E → ℝ,
      ContDiff ℝ (⊤ : ℕ∞) ψ ∧
      HasCompactSupport ψ ∧
      eLpNorm (fun x => ψ x - u x)
        (ENNReal.ofReal p) (volume.restrict (Metric.ball (0 : E) 1)) < ENNReal.ofReal ε ∧
      ∀ i : Fin d,
        eLpNorm
          (fun x => (fderiv ℝ ψ x) (EuclideanSpace.single i 1) - hw.weakGrad x i)
          (ENNReal.ofReal p) (volume.restrict (Metric.ball (0 : E) 1)) < ENNReal.ofReal ε := by
  let B : Set E := Metric.ball (0 : E) 1
  let q : ℝ≥0∞ := ENNReal.ofReal p
  have hp_le : 1 ≤ p := le_of_lt hp
  have hq_ge_one : 1 ≤ q := by
    simpa [q] using (ENNReal.ofReal_le_ofReal hp_le)
  have hε_half_pos : 0 < ε / 2 := by positivity
  have hε_quarter_pos : 0 < ε / 4 := by positivity
  have hquarter_lt_half : ENNReal.ofReal (ε / 4) < ENNReal.ofReal (ε / 2) := by
    refine (ENNReal.ofReal_lt_ofReal_iff hε_half_pos).2 ?_
    nlinarith
  obtain ⟨δu, hδu_pos, hδu⟩ :=
    exists_delta_unitBallDilate_close_of_memLp (d := d) (p := p) hp_le hw.memLp
      (ε := ε / 2) hε_half_pos
  choose δG hδG_pos hδG using
    fun i : Fin d =>
      exists_delta_unitBallDilate_scaled_close_of_memLp (d := d) (p := p) hp_le
        (hw.weakGrad_component_memLp i) (ε := ε / 2) hε_half_pos
  let δgrad : ℝ := Finset.univ.inf' Finset.univ_nonempty δG
  have hδgrad_pos : 0 < δgrad := by
    dsimp [δgrad]
    rw [Finset.lt_inf'_iff]
    intro i hi
    exact hδG_pos i
  let δ : ℝ := min 1 (min δu δgrad)
  have hδ_pos : 0 < δ := by
    dsimp [δ]
    positivity
  let lam : ℝ := 1 + δ / 2
  have hlam_gt_one : 1 < lam := by
    dsimp [lam]
    have : 0 < δ / 2 := by positivity
    linarith
  have hlam_lt_two : lam < 2 := by
    dsimp [lam, δ]
    have hδ_le_one : δ ≤ 1 := by
      dsimp [δ]
      exact min_le_left _ _
    linarith
  have hlam_lt_u : lam < 1 + δu := by
    have hδ_le : δ ≤ δu := by
      dsimp [δ]
      exact (min_le_right _ _).trans (min_le_left _ _)
    have hδ_half_lt : δ / 2 < δu := by
      have : δ / 2 < δ := by
        have hδp : 0 < δ := hδ_pos
        nlinarith
      exact lt_of_lt_of_le this hδ_le
    dsimp [lam]
    linarith
  have hlam_lt_G : ∀ i : Fin d, lam < 1 + δG i := by
    intro i
    have hδgrad_le : δgrad ≤ δG i := by
      dsimp [δgrad]
      exact Finset.inf'_le (s := Finset.univ) (f := δG) (by simp)
    have hδ_le : δ ≤ δG i := by
      dsimp [δ]
      exact (min_le_right _ _).trans ((min_le_right _ _).trans hδgrad_le)
    have hδ_half_lt : δ / 2 < δG i := by
      have : δ / 2 < δ := by
        have hδp : 0 < δ := hδ_pos
        nlinarith
      exact lt_of_lt_of_le this hδ_le
    dsimp [lam]
    linarith
  have hfun_dil :
      eLpNorm (fun x => DeGiorgi.unitBallDilate (d := d) lam u x - u x)
        (ENNReal.ofReal p) (volume.restrict (Metric.ball (0 : E) 1)) <
        ENNReal.ofReal (ε / 2) := hδu hlam_gt_one hlam_lt_u
  have hgrad_dil :
      ∀ i : Fin d,
        eLpNorm
          (fun x => lam⁻¹ * DeGiorgi.unitBallDilate (d := d) lam (fun y => hw.weakGrad y i) x -
            hw.weakGrad x i)
          (ENNReal.ofReal p) (volume.restrict (Metric.ball (0 : E) 1)) < ENNReal.ofReal (ε / 2) := by
    intro i
    exact hδG i hlam_gt_one (hlam_lt_G i)
  obtain ⟨η, hη_sub⟩ := exists_unitBallCutoff_inside (d := d) hlam_gt_one
  let Ω : Set E := Metric.ball (0 : E) lam
  let udil : E → ℝ := DeGiorgi.unitBallDilate (d := d) lam u
  let v : E → ℝ := fun x => η.toFun x * udil x
  let hwDil : MemW1pWitness (ENNReal.ofReal p) udil Ω :=
    MemW1pWitness.unitBallDilate_largeBall (d := d) (p := p) hp_le hlam_gt_one hw
  have hη_bound : ∀ x, |η.toFun x| ≤ 1 := by
    intro x
    rw [abs_of_nonneg (η.nonneg x)]
    exact η.le_one x
  let C1 : ℝ := ↑Mst * (((1 + lam) / 2) - 1)⁻¹
  have hC1_nonneg : 0 ≤ C1 := by
    have hmid_pos : 0 < ((1 + lam) / 2 : ℝ) - 1 := by
      nlinarith
    dsimp [C1]
    exact mul_nonneg (by positivity) (inv_nonneg.mpr hmid_pos.le)
  let hwLoc : MemW1pWitness (ENNReal.ofReal p) v Ω :=
    MemW1pWitness.mul_smooth_bounded_p (d := d) (p := ENNReal.ofReal p) hq_ge_one isOpen_ball hwDil
      η.smooth zero_le_one hC1_nonneg hη_bound (by
        intro x
        simpa [C1] using η.grad_bound x)
  have hv_sub : tsupport v ⊆ Ω := by
    dsimp [v, Ω, udil]
    exact (tsupport_smul_subset_left η.toFun (DeGiorgi.unitBallDilate (d := d) lam u)).trans hη_sub
  have hv_compact : HasCompactSupport v := by
    apply HasCompactSupport.intro' (isCompact_closedBall (0 : E) ((1 + lam) / 2)) isClosed_closedBall
    intro x hx
    exact zero_outside_of_tsupport_subset
      (Ω := Metric.closedBall (0 : E) ((1 + lam) / 2))
      ((tsupport_smul_subset_left η.toFun (DeGiorgi.unitBallDilate (d := d) lam u)).trans η.support_subset) hx
  have hKΩ : Metric.closedBall (0 : E) ((1 + lam) / 2) ⊆ Ω := by
    dsimp [Ω]
    exact Metric.closedBall_subset_ball (midpoint_lt_of_one_lt hlam_gt_one)
  have hgrad_sub :
      ∀ i : Fin d,
        tsupport (fun x => hwLoc.weakGrad x i) ⊆ Metric.closedBall (0 : E) ((1 + lam) / 2) := by
    intro i
    rw [← isClosed_closedBall.closure_eq]
    apply closure_mono
    refine Function.support_subset_iff'.2 ?_
    intro x hx
    have hηx : η.toFun x = 0 := zero_outside_of_tsupport_subset (Ω := Metric.closedBall (0 : E) ((1 + lam) / 2))
      η.support_subset hx
    have hηdx :
        (fderiv ℝ η.toFun x) (EuclideanSpace.single i 1) = 0 :=
      fderiv_apply_zero_outside_of_tsupport_subset
        (Ω := Metric.closedBall (0 : E) ((1 + lam) / 2))
        (hf := η.smooth) η.support_subset hx i
    simp [hwLoc, MemW1pWitness.mul_smooth_bounded_p, hwDil,
      MemW1pWitness.unitBallDilate_largeBall, udil, DeGiorgi.unitBallDilate,
      PiLp.toLp_apply, hηx, hηdx]
  rcases exists_smooth_W1p_approx_of_supportedWitness
      (d := d) (Ω := Ω) (K := Metric.closedBall (0 : E) ((1 + lam) / 2))
      isOpen_ball hp hwLoc (isCompact_closedBall (0 : E) ((1 + lam) / 2)) hKΩ
      ((tsupport_smul_subset_left η.toFun (DeGiorgi.unitBallDilate (d := d) lam u)).trans η.support_subset)
      hgrad_sub with
    ⟨φ, hφ_smooth, hφ_compact, hφ_sub, hφ_fun, hφ_grad⟩
  have hB_sub_Ω : B ⊆ Ω := unitBall_subset_ball_of_one_lt (d := d) hlam_gt_one
  have hφ_fun_ev :
      ∀ᶠ n in atTop,
        eLpNorm (fun x => φ n x - v x) q (volume.restrict B) < ENNReal.ofReal (ε / 2) := by
    filter_upwards [ENNReal.tendsto_nhds_zero.1 hφ_fun (ENNReal.ofReal (ε / 4))
      (ENNReal.ofReal_pos.mpr hε_quarter_pos)] with n hn
    have hmono :
        eLpNorm (fun x => φ n x - v x) q (volume.restrict B) ≤
          eLpNorm (fun x => φ n x - v x) q (volume.restrict Ω) :=
      eLpNorm_mono_measure _ (Measure.restrict_mono_set volume hB_sub_Ω)
    exact lt_of_le_of_lt hmono (lt_of_le_of_lt (by simpa [q] using hn) hquarter_lt_half)
  have hφ_grad_ev :
      ∀ i : Fin d, ∀ᶠ n in atTop,
        eLpNorm
          (fun x => (fderiv ℝ (φ n) x) (EuclideanSpace.single i 1) -
            Ω.indicator (fun y => hwLoc.weakGrad y i) x)
          q (volume.restrict B) < ENNReal.ofReal (ε / 2) := by
    intro i
    filter_upwards [ENNReal.tendsto_nhds_zero.1 (hφ_grad i) (ENNReal.ofReal (ε / 4))
      (ENNReal.ofReal_pos.mpr hε_quarter_pos)] with n hn
    have hEq :
        (fun x => (fderiv ℝ (φ n) x) (EuclideanSpace.single i 1) -
            Ω.indicator (fun y => hwLoc.weakGrad y i) x) =ᵐ[volume.restrict B]
          (fun x => (fderiv ℝ (φ n) x) (EuclideanSpace.single i 1) - hwLoc.weakGrad x i) := by
      refine ae_restrict_of_forall_mem measurableSet_ball ?_
      intro x hx
      have hxΩ : x ∈ Ω := hB_sub_Ω hx
      simp [Ω, hxΩ]
    rw [eLpNorm_congr_ae hEq]
    have hmono :
        eLpNorm
            (fun x => (fderiv ℝ (φ n) x) (EuclideanSpace.single i 1) - hwLoc.weakGrad x i)
            q (volume.restrict B) ≤
          eLpNorm
            (fun x => (fderiv ℝ (φ n) x) (EuclideanSpace.single i 1) - hwLoc.weakGrad x i)
            q (volume.restrict Ω) :=
      eLpNorm_mono_measure _ (Measure.restrict_mono_set volume hB_sub_Ω)
    exact lt_of_le_of_lt hmono (lt_of_le_of_lt (by simpa [q] using hn) hquarter_lt_half)
  have hφ_grad_all :
      ∀ᶠ n in atTop,
        ∀ i ∈ Finset.univ,
          eLpNorm
            (fun x => (fderiv ℝ (φ n) x) (EuclideanSpace.single i 1) -
              Ω.indicator (fun y => hwLoc.weakGrad y i) x)
            q (volume.restrict B) < ENNReal.ofReal (ε / 2) := by
    simpa using (Finset.eventually_all (I := Finset.univ)
      (l := atTop)
      (p := fun i n =>
        eLpNorm
          (fun x => (fderiv ℝ (φ n) x) (EuclideanSpace.single i 1) -
            Ω.indicator (fun y => hwLoc.weakGrad y i) x)
          q (volume.restrict B) < ENNReal.ofReal (ε / 2))).2
      (by intro i hi; exact hφ_grad_ev i)
  have hchoose :
      ∀ᶠ n in atTop,
        eLpNorm (fun x => φ n x - v x) q (volume.restrict B) < ENNReal.ofReal (ε / 2) ∧
        ∀ i ∈ Finset.univ,
          eLpNorm
            (fun x => (fderiv ℝ (φ n) x) (EuclideanSpace.single i 1) -
              Ω.indicator (fun y => hwLoc.weakGrad y i) x)
            q (volume.restrict B) < ENNReal.ofReal (ε / 2) := by
    exact hφ_fun_ev.and hφ_grad_all
  rcases Filter.eventually_atTop.1 hchoose with ⟨N, hN⟩
  let ψ : E → ℝ := φ N
  have hψ_smooth : ContDiff ℝ (⊤ : ℕ∞) ψ := hφ_smooth N
  have hψ_compact : HasCompactSupport ψ := hφ_compact N
  have hψ_fun_half :
      eLpNorm (fun x => ψ x - v x) q (volume.restrict B) < ENNReal.ofReal (ε / 2) := by
    simpa [ψ] using (hN N le_rfl).1
  have hψ_grad_half :
      ∀ i : Fin d,
        eLpNorm
          (fun x => (fderiv ℝ ψ x) (EuclideanSpace.single i 1) -
            Ω.indicator (fun y => hwLoc.weakGrad y i) x)
          q (volume.restrict B) < ENNReal.ofReal (ε / 2) := by
    intro i
    simpa [ψ] using (hN N le_rfl).2 i (by simp)
  have hv_memLp_B : MemLp v q (volume.restrict B) := by
    exact hwLoc.memLp.mono_measure (Measure.restrict_mono_set volume hB_sub_Ω)
  have hv_sub_eq_udil_sub :
      (fun x => v x - u x) =ᵐ[volume.restrict B] (fun x => udil x - u x) := by
    refine ae_restrict_of_forall_mem measurableSet_ball ?_
    intro x hx
    simp [v, η.eq_one x hx]
  have hψ_sub_sum_ae :
      (fun x => ψ x - u x) =ᵐ[volume.restrict B]
        (fun x => (ψ x - v x) + (v x - u x)) := by
    exact Eventually.of_forall (by intro x; ring)
  have hψ_minus_v_aesm :
      AEStronglyMeasurable (fun x => ψ x - v x) (volume.restrict B) := by
    exact
      ((hψ_smooth.continuous.aestronglyMeasurable.mono_ac
        Measure.restrict_le_self.absolutelyContinuous).sub hv_memLp_B.aestronglyMeasurable)
  have hv_minus_u_aesm :
      AEStronglyMeasurable (fun x => v x - u x) (volume.restrict B) := by
    exact hv_memLp_B.aestronglyMeasurable.sub hw.memLp.aestronglyMeasurable
  have hv_minus_u_lt :
      eLpNorm (fun x => v x - u x) q (volume.restrict B) < ENNReal.ofReal (ε / 2) := by
    rw [eLpNorm_congr_ae hv_sub_eq_udil_sub]
    simpa [B, udil] using hfun_dil
  have hhalf_add :
      ENNReal.ofReal (ε / 2) + ENNReal.ofReal (ε / 2) = ENNReal.ofReal ε := by
    rw [← ENNReal.ofReal_add (by positivity) (by positivity)]
    ring_nf
  have hfun_final :
      eLpNorm (fun x => ψ x - u x) q (volume.restrict B) < ENNReal.ofReal ε := by
    calc
      eLpNorm (fun x => ψ x - u x) q (volume.restrict B)
          = eLpNorm (fun x => (ψ x - v x) + (v x - u x)) q (volume.restrict B) := by
              exact eLpNorm_congr_ae hψ_sub_sum_ae
      _ ≤ eLpNorm (fun x => ψ x - v x) q (volume.restrict B) +
            eLpNorm (fun x => v x - u x) q (volume.restrict B) := by
              exact eLpNorm_add_le hψ_minus_v_aesm hv_minus_u_aesm hq_ge_one
      _ < ENNReal.ofReal (ε / 2) + ENNReal.ofReal (ε / 2) := by
            exact ENNReal.add_lt_add hψ_fun_half hv_minus_u_lt
      _ = ENNReal.ofReal ε := hhalf_add
  have hgrad_local_eq :
      ∀ i : Fin d,
        (fun x => Ω.indicator (fun y => hwLoc.weakGrad y i) x - hw.weakGrad x i)
          =ᵐ[volume.restrict B]
        (fun x => lam⁻¹ * DeGiorgi.unitBallDilate (d := d) lam (fun y => hw.weakGrad y i) x -
          hw.weakGrad x i) := by
    intro i
    refine ae_restrict_of_forall_mem measurableSet_ball ?_
    intro x hx
    have hxΩ : x ∈ Ω := hB_sub_Ω hx
    have hηx : η.toFun x = 1 := η.eq_one x hx
    have hηdx :
        (fderiv ℝ η.toFun x) (EuclideanSpace.single i 1) = 0 :=
      cutoff_fderiv_apply_eq_zero_on_unitBall (d := d) hlam_gt_one hx i
    change Ω.indicator (fun y => hwLoc.weakGrad y i) x - hw.weakGrad x i =
      lam⁻¹ * DeGiorgi.unitBallDilate (d := d) lam (fun y => hw.weakGrad y i) x - hw.weakGrad x i
    have hΩ_ind :
        Ω.indicator (fun y => hwLoc.weakGrad y i) x = hwLoc.weakGrad x i := by
      simp [Ω, hxΩ]
    rw [hΩ_ind]
    change hwLoc.weakGrad x i - hw.weakGrad x i =
      lam⁻¹ * DeGiorgi.unitBallDilate (d := d) lam (fun y => hw.weakGrad y i) x - hw.weakGrad x i
    simp [hwLoc, MemW1pWitness.mul_smooth_bounded_p, hwDil, MemW1pWitness.unitBallDilate_largeBall, v, udil, DeGiorgi.unitBallDilate, hηx, hηdx, smul_eq_mul]
  have hgrad_final :
      ∀ i : Fin d,
        eLpNorm
          (fun x => (fderiv ℝ ψ x) (EuclideanSpace.single i 1) - hw.weakGrad x i)
          q (volume.restrict B) < ENNReal.ofReal ε := by
    intro i
    let giLoc : E → ℝ := fun x => Ω.indicator (fun y => hwLoc.weakGrad y i) x
    have hgiLoc_eq :
        giLoc =ᵐ[volume.restrict B] (fun x => hwLoc.weakGrad x i) := by
      refine ae_restrict_of_forall_mem measurableSet_ball ?_
      intro x hx
      have hxΩ : x ∈ Ω := hB_sub_Ω hx
      simp [giLoc, hxΩ]
    have hgiLoc_memLp_B : MemLp giLoc q (volume.restrict B) := by
      have htmp :
          MemLp (fun x => hwLoc.weakGrad x i) q (volume.restrict B) := by
        exact (hwLoc.weakGrad_component_memLp i).mono_measure
          (Measure.restrict_mono_set volume hB_sub_Ω)
      refine ⟨?_, ?_⟩
      · exact htmp.1.congr hgiLoc_eq.symm
      · simpa [eLpNorm_congr_ae hgiLoc_eq.symm] using htmp.2
    let ei : E := EuclideanSpace.single i (1 : ℝ)
    have hderiv_cont : Continuous (fun x => (fderiv ℝ ψ x) ei) := by
      simpa [ei] using
        ((hψ_smooth.continuous_fderiv
          (by simp : ((⊤ : ℕ∞) : WithTop ℕ∞) ≠ 0)).clm_apply continuous_const)
    have hderiv_minus_loc_aesm :
        AEStronglyMeasurable
          (fun x => (fderiv ℝ ψ x) ei - giLoc x) (volume.restrict B) := by
      exact
        ((hderiv_cont.aestronglyMeasurable.mono_ac
          Measure.restrict_le_self.absolutelyContinuous).sub hgiLoc_memLp_B.aestronglyMeasurable)
    have hloc_minus_grad_aesm :
        AEStronglyMeasurable (fun x => giLoc x - hw.weakGrad x i) (volume.restrict B) := by
      exact hgiLoc_memLp_B.aestronglyMeasurable.sub (hw.weakGrad_component_memLp i).aestronglyMeasurable
    have hsum_ae :
        (fun x => (fderiv ℝ ψ x) ei - hw.weakGrad x i) =ᵐ[volume.restrict B]
          (fun x => ((fderiv ℝ ψ x) ei - giLoc x) + (giLoc x - hw.weakGrad x i)) := by
      exact Eventually.of_forall (by intro x; ring)
    have hloc_minus_grad_lt :
        eLpNorm (fun x => giLoc x - hw.weakGrad x i) q (volume.restrict B)
          < ENNReal.ofReal (ε / 2) := by
      have hEq :
          (fun x => giLoc x - hw.weakGrad x i) =ᵐ[volume.restrict B]
            (fun x => lam⁻¹ * DeGiorgi.unitBallDilate (d := d) lam (fun y => hw.weakGrad y i) x -
              hw.weakGrad x i) := by
        simpa [giLoc] using hgrad_local_eq i
      rw [eLpNorm_congr_ae hEq]
      simpa [B] using hgrad_dil i
    calc
      eLpNorm (fun x => (fderiv ℝ ψ x) (EuclideanSpace.single i 1) - hw.weakGrad x i)
          q (volume.restrict B)
          = eLpNorm (fun x => ((fderiv ℝ ψ x) ei - giLoc x) +
              (giLoc x - hw.weakGrad x i)) q (volume.restrict B) := by
                exact eLpNorm_congr_ae hsum_ae
      _ ≤ eLpNorm (fun x => (fderiv ℝ ψ x) ei - giLoc x) q (volume.restrict B) +
            eLpNorm (fun x => giLoc x - hw.weakGrad x i) q (volume.restrict B) := by
              exact eLpNorm_add_le hderiv_minus_loc_aesm hloc_minus_grad_aesm hq_ge_one
      _ < ENNReal.ofReal (ε / 2) + ENNReal.ofReal (ε / 2) := by
            exact ENNReal.add_lt_add (hψ_grad_half i) hloc_minus_grad_lt
      _ = ENNReal.ofReal ε := hhalf_add
  refine ⟨ψ, hψ_smooth, hψ_compact, ?_, ?_⟩
  · simpa [B, q] using hfun_final
  · intro i
    simpa [B, q] using hgrad_final i

\end{lstlisting}
\begin{proof}
Choose $\delta_u, \delta_G > 0$ small enough so that for any
$\lambda \in (1, 1+\delta/2)$ with $\delta := \min(1,\delta_u,\delta_G)$,
the inward dilate $\unitBallDilate_\lambda u$ approximates $u$ in $L^p(\Bz{1})$
to within $\varepsilon/2$, and similarly each rescaled gradient component is
within $\varepsilon/2$ of $(\nabla u)_i$. (Existence of such $\delta$ uses
strong-$L^p$-continuity of dilations.) Set $\lambda := 1 + \delta/2 \in (1,2)$.

Apply Lemma~\ref{lem:dilate-large} to get
$u_\lambda := \unitBallDilate_\lambda u \in W^{1,p}(\Ball{\lambda}{0})$.
Multiply $u_\lambda$ by a cutoff structure $\eta$ on $\Bz{1}$ inside
$\Ball{(1+\lambda)/2}{0}$ from \leanname{exists\_unitBallCutoff\_inside}; this
produces a $W^{1,p}$ function $\eta\, u_\lambda$ supported strictly inside
$\Ball{\lambda}{0}$, agreeing with $u_\lambda$ on $\Bz{1}$.

Convolve $\eta\, u_\lambda$ with a standard mollifier of radius $\rho$ small
enough that the support of the result lies inside $\Ball{\lambda}{0}$ and the
$L^p$ approximation error in both function and gradient is controlled by
$\varepsilon/4$. The result $\psi$ is smooth, compactly supported, and
satisfies, by the triangle inequality,
\begin{align*}
\eLpNormText{\psi - u}_{L^p(\Bz{1})}
&\le \eLpNormText{\psi - \eta u_\lambda}_{L^p(\Bz{1})}
+ \eLpNormText{\eta u_\lambda - u}_{L^p(\Bz{1})} \\
&< \varepsilon/4 + \varepsilon/2 < \varepsilon,
\end{align*}
and analogously for the gradient components, where on $\Bz{1}$ the cutoff
$\eta$ equals $1$, hence $\eta u_\lambda = u_\lambda$ and the gradient
error reduces to a sum of two terms each less than $\varepsilon/2$.
\end{proof}

\begin{theorem}[\leanname{exists\_smooth\_W1p\_approx\_on\_unitBall}]
\label{thm:wp-approx}
Let $1 < p < \infty$ and $u \in W^{1,p}(\Bz{1})$. There exists a sequence
$\psi_n \in C^\infty_c(E)$ such that $\psi_n \to u$ in $L^p(\Bz{1})$ and
$\partial_i \psi_n \to \partial_i u$ in $L^p(\Bz{1})$ for each $i$.
\end{theorem}

\begin{lstlisting}[style=Matlab-editor,basicstyle=\mlttfamily,escapechar=",mlshowsectionrules=true,mathescape=true,morecomment={[s]{\%\{}{\%\}}},language=lean,numbers=none]
/-- Sequence form of full `W^{1,p}` smooth approximation on the unit ball. -/
theorem exists_smooth_W1p_approx_on_unitBall
    {p : ℝ} (hp : 1 < p) {u : E → ℝ}
    (hw : MemW1pWitness (ENNReal.ofReal p) u (Metric.ball (0 : E) 1)) :
    ∃ ψ : ℕ → E → ℝ,
      (∀ n, ContDiff ℝ (⊤ : ℕ∞) (ψ n)) ∧
      (∀ n, HasCompactSupport (ψ n)) ∧
      Tendsto
        (fun n => eLpNorm (fun x => ψ n x - u x)
          (ENNReal.ofReal p) (volume.restrict (Metric.ball (0 : E) 1)))
        atTop (nhds 0) ∧
      (∀ i : Fin d,
        Tendsto
          (fun n => eLpNorm
            (fun x => (fderiv ℝ (ψ n) x) (EuclideanSpace.single i 1) - hw.weakGrad x i)
            (ENNReal.ofReal p) (volume.restrict (Metric.ball (0 : E) 1)))
          atTop (nhds 0)) := by
  let eps : ℕ → ℝ := fun n => ((n : ℝ) + 1)⁻¹
  have heps_pos : ∀ n, 0 < eps n := by
    intro n
    exact inv_pos.mpr (by positivity)
  choose ψ hψ using
    fun n => exists_smooth_W1p_oneShot_on_unitBall (d := d) (p := p) hp hw
      (ε := eps n) (heps_pos n)
  have hψ_smooth : ∀ n, ContDiff ℝ (⊤ : ℕ∞) (ψ n) := fun n => (hψ n).1
  have hψ_compact : ∀ n, HasCompactSupport (ψ n) := fun n => (hψ n).2.1
  have hψ_fun :
      ∀ n,
        eLpNorm (fun x => ψ n x - u x)
          (ENNReal.ofReal p) (volume.restrict (Metric.ball (0 : E) 1)) <
          ENNReal.ofReal (eps n) := fun n => (hψ n).2.2.1
  have hψ_grad :
      ∀ n (i : Fin d),
        eLpNorm
          (fun x => (fderiv ℝ (ψ n) x) (EuclideanSpace.single i 1) - hw.weakGrad x i)
          (ENNReal.ofReal p) (volume.restrict (Metric.ball (0 : E) 1)) <
          ENNReal.ofReal (eps n) := fun n => (hψ n).2.2.2
  refine ⟨ψ, hψ_smooth, hψ_compact, ?_, ?_⟩
  · have h_eps_tendsto_real' :
        Tendsto (fun n : ℕ => (1 : ℝ) / ((n : ℝ) + 1)) atTop (nhds (0 : ℝ)) := by
      exact tendsto_one_div_add_atTop_nhds_zero_nat
    have h_eps_tendsto_real : Tendsto eps atTop (nhds (0 : ℝ)) := by
      simpa [eps] using h_eps_tendsto_real'
    have h_eps_tendsto : Tendsto (fun n => ENNReal.ofReal (eps n)) atTop (nhds 0) := by
      simpa using (ENNReal.continuous_ofReal.tendsto (0 : ℝ)).comp h_eps_tendsto_real
    refine ENNReal.tendsto_nhds_zero.2 ?_
    intro ε hε
    filter_upwards [ENNReal.tendsto_nhds_zero.1 h_eps_tendsto ε hε] with n hn
    exact le_trans (le_of_lt (hψ_fun n)) hn
  · intro i
    have h_eps_tendsto_real' :
        Tendsto (fun n : ℕ => (1 : ℝ) / ((n : ℝ) + 1)) atTop (nhds (0 : ℝ)) := by
      exact tendsto_one_div_add_atTop_nhds_zero_nat
    have h_eps_tendsto_real : Tendsto eps atTop (nhds (0 : ℝ)) := by
      simpa [eps] using h_eps_tendsto_real'
    have h_eps_tendsto : Tendsto (fun n => ENNReal.ofReal (eps n)) atTop (nhds 0) := by
      simpa using (ENNReal.continuous_ofReal.tendsto (0 : ℝ)).comp h_eps_tendsto_real
    refine ENNReal.tendsto_nhds_zero.2 ?_
    intro ε hε
    filter_upwards [ENNReal.tendsto_nhds_zero.1 h_eps_tendsto ε hε] with n hn
    exact le_trans (le_of_lt (hψ_grad n i)) hn

\end{lstlisting}
\begin{proof}
The one-shot Theorem~\ref{thm:oneshot} produces, for any $\varepsilon > 0$, a
single $\psi \in C^\infty_c(E)$ whose $L^p(\Bz{1})$ distance to $u$ and whose
componentwise $L^p(\Bz{1})$ distance from $\partial_i u$ are both strictly
less than $\varepsilon$. Apply this with the sequence
$\varepsilon_n := (n+1)^{-1}$ and use the axiom of choice
(\leanname{Classical.choose}) to extract a witness
$\psi_n \in C^\infty_c(E)$ for each $n$. By construction, the smoothness and
compact support of $\psi_n$ are inherited from the one-shot output, and we
have the strict bounds
\[
  \eLpNormText{\psi_n - u}_{L^p(\Bz{1})}
    < \mathrm{ENNReal.ofReal}(\varepsilon_n),
\]
together with the analogous bound for each gradient component
$(\fderiv{\psi_n}{x})(e_i) - (\nabla u)_i$.

To upgrade these strict numerical bounds to convergence in $L^p$, observe
that $1/(n+1) \to 0$ in $\R$
(\leanname{tendsto\_one\_div\_add\_atTop\_nhds\_zero\_nat}), so
$\mathrm{ENNReal.ofReal}(\varepsilon_n) \to 0$ in $[0,\infty]$. Each function
error is dominated by these vanishing constants, so the squeeze theorem
\leanname{tendsto\_of\_tendsto\_of\_tendsto\_of\_le\_of\_le}, applied with the
constant zero sequence as a lower bound, yields convergence to $0$. The same
argument applied componentwise to each
$\partial_i \psi_n - (\nabla u)_i$ gives the gradient convergence.
\end{proof}

\begin{theorem}[\leanname{exists\_smooth\_W12\_approx\_on\_unitBall}]
\label{thm:w12-approx}
The previous theorem holds in particular for $p = 2$.
\end{theorem}

\begin{lstlisting}[style=Matlab-editor,basicstyle=\mlttfamily,escapechar=",mlshowsectionrules=true,mathescape=true,morecomment={[s]{\%\{}{\%\}}},language=lean,numbers=none]
/-- `L²` specialization of the unit-ball smooth approximation theorem. -/
theorem exists_smooth_W12_approx_on_unitBall
    {u : E → ℝ}
    (hw : MemW1pWitness 2 u (Metric.ball (0 : E) 1)) :
    ∃ ψ : ℕ → E → ℝ,
      (∀ n, ContDiff ℝ (⊤ : ℕ∞) (ψ n)) ∧
      (∀ n, HasCompactSupport (ψ n)) ∧
      Tendsto
        (fun n => eLpNorm (fun x => ψ n x - u x)
          2 (volume.restrict (Metric.ball (0 : E) 1)))
        atTop (nhds 0) ∧
      (∀ i : Fin d,
        Tendsto
          (fun n => eLpNorm
            (fun x => (fderiv ℝ (ψ n) x) (EuclideanSpace.single i 1) - hw.weakGrad x i)
            2 (volume.restrict (Metric.ball (0 : E) 1)))
          atTop (nhds 0)) := by
  let hw' : MemW1pWitness (ENNReal.ofReal (2 : ℝ)) u (Metric.ball (0 : E) 1) :=
    { memLp := by
        simpa using hw.memLp
      weakGrad := hw.weakGrad
      weakGrad_component_memLp := by
        intro i
        simpa using hw.weakGrad_component_memLp i
      isWeakGrad := hw.isWeakGrad }
  rcases exists_smooth_W1p_approx_on_unitBall (d := d) (p := (2 : ℝ)) (by norm_num) hw' with
    ⟨ψ, hψ_smooth, hψ_compact, hψ_fun, hψ_grad⟩
  refine ⟨ψ, hψ_smooth, hψ_compact, ?_, ?_⟩
  · simpa [hw'] using hψ_fun
  · intro i
    simpa [hw'] using hψ_grad i

\end{lstlisting}
\begin{proof}
Direct specialization of Theorem~\ref{thm:wp-approx} using $\ENNRealofReal(2)
= 2$ (after coercing the witness through the trivial isomorphism between the
two encodings of the exponent).
\end{proof}

\subsection{The unit-ball extension operator}

\subsubsection*{Definitions and elementary algebra}

\begin{definition}[\leanname{unitBallRetraction}]
Define $r : E \to E$ by
\[
r(x) := \begin{cases}
x & \text{if } \|x\| \le 1, \\
\|x\|^{-2}\, x & \text{if } 1 < \|x\| < 2, \\
0 & \text{if } \|x\| \ge 2.
\end{cases}
\]
On the annulus $1 < \|x\| < 2$, $r(x)$ is the spherical inversion centred at
the origin with radius $1$.
\end{definition}

\begin{lstlisting}[style=Matlab-editor,basicstyle=\mlttfamily,escapechar=",mlshowsectionrules=true,mathescape=true,morecomment={[s]{\%\{}{\%\}}},language=lean,numbers=none]
/-- Retraction used for the explicit unit-ball extension operator. -/
def unitBallRetraction (x : E) : E :=
  if ‖x‖ ≤ 1 then
    x
  else if ‖x‖ < 2 then
    (‖x‖ ^ (2 : ℕ))⁻¹ • x
  else
    0

\end{lstlisting}
\begin{definition}[\leanname{unitBallCutoff}]
Define $\chi : E \to \R$ by $\chi(x) := \min(1, \max(2 - \|x\|, 0))$.
\end{definition}

\begin{lstlisting}[style=Matlab-editor,basicstyle=\mlttfamily,escapechar=",mlshowsectionrules=true,mathescape=true,morecomment={[s]{\%\{}{\%\}}},language=lean,numbers=none]
/-- Radial cutoff used for the explicit unit-ball extension operator. -/
def unitBallCutoff (x : E) : ℝ :=
  min 1 (max (2 - ‖x‖) 0)

\end{lstlisting}
\begin{definition}[\leanname{unitBallExtension}]
For $u : E \to \R$, define
\[
\Eext u (x) := \chi(x)\, u(r(x)).
\]
\end{definition}

\begin{lstlisting}[style=Matlab-editor,basicstyle=\mlttfamily,escapechar=",mlshowsectionrules=true,mathescape=true,morecomment={[s]{\%\{}{\%\}}},language=lean,numbers=none]
/-- Explicit unit-ball extension operator. -/
def unitBallExtension (u : E → ℝ) (x : E) : ℝ :=
  unitBallCutoff x * u (unitBallRetraction x)

\end{lstlisting}
\begin{lemma}[\leanname{unitBallExtension\_zero}, \leanname{unitBallExtension\_add}, \leanname{unitBallExtension\_sub}]
The map $u \mapsto \Eext u$ is $\R$-linear: $\Eext 0 = 0$,
$\Eext{(u+v)} = \Eext u + \Eext v$, and $\Eext{(u-v)} = \Eext u - \Eext v$.
\end{lemma}

\begin{lstlisting}[style=Matlab-editor,basicstyle=\mlttfamily,escapechar=",mlshowsectionrules=true,mathescape=true,morecomment={[s]{\%\{}{\%\}}},language=lean,numbers=none]
omit [NeZero d] in
lemma unitBallExtension_zero :
    unitBallExtension (d := d) (fun _ : E => (0 : ℝ)) = fun _ : E => (0 : ℝ) := by
  funext x
  simp [unitBallExtension]

omit [NeZero d] in
lemma unitBallExtension_add (u v : E → ℝ) :
    unitBallExtension (d := d) (fun x => u x + v x) =
      fun x => unitBallExtension (d := d) u x + unitBallExtension (d := d) v x := by
  funext x
  simp [unitBallExtension, mul_add]

omit [NeZero d] in
lemma unitBallExtension_sub (u v : E → ℝ) :
    unitBallExtension (d := d) (fun x => u x - v x) =
      fun x => unitBallExtension (d := d) u x - unitBallExtension (d := d) v x := by
  funext x
  simp [unitBallExtension, sub_eq_add_neg, mul_add]

\end{lstlisting}
\begin{proof}
Pointwise from the definition: each is a multiplication of $\chi(x)$ by a
linear combination of values of $u$ and $v$ at $r(x)$.
\end{proof}

\begin{lemma}[\leanname{unitBallRetraction\_eq\_self\_of\_norm\_le\_one}, \leanname{unitBallRetraction\_eq\_self\_of\_mem\_ball}]
If $\|x\| \le 1$ then $r(x) = x$. In particular if $x \in \Bz{1}$ then
$r(x) = x$.
\end{lemma}

\begin{lstlisting}[style=Matlab-editor,basicstyle=\mlttfamily,escapechar=",mlshowsectionrules=true,mathescape=true,morecomment={[s]{\%\{}{\%\}}},language=lean,numbers=none]
omit [NeZero d] in
lemma unitBallRetraction_eq_self_of_norm_le_one {x : E} (hx : ‖x‖ ≤ 1) :
    unitBallRetraction (d := d) x = x := by
  simp [unitBallRetraction, hx]

omit [NeZero d] in
lemma unitBallRetraction_eq_self_of_mem_ball {x : E} (hx : x ∈ Metric.ball (0 : E) 1) :
    unitBallRetraction (d := d) x = x := by
  apply unitBallRetraction_eq_self_of_norm_le_one (d := d)
  rw [Metric.mem_ball, dist_zero_right] at hx
  exact le_of_lt hx

\end{lstlisting}
\begin{proof}
Direct from the first branch of the definition.
\end{proof}

\begin{lemma}[\leanname{unitBallExtension\_eq\_of\_mem\_ball}]
\label{lem:Eext-restrict}
If $x \in \Bz{1}$ then $\Eext u(x) = u(x)$.
\end{lemma}

\begin{lstlisting}[style=Matlab-editor,basicstyle=\mlttfamily,escapechar=",mlshowsectionrules=true,mathescape=true,morecomment={[s]{\%\{}{\%\}}},language=lean,numbers=none]
omit [NeZero d] in
lemma unitBallExtension_eq_of_mem_ball {u : E → ℝ} {x : E}
    (hx : x ∈ Metric.ball (0 : E) 1) :
    unitBallExtension (d := d) u x = u x := by
  unfold unitBallExtension unitBallCutoff
  rw [unitBallRetraction_eq_self_of_mem_ball (d := d) hx]
  rw [Metric.mem_ball, dist_zero_right] at hx
  have hgt : 1 < 2 - ‖x‖ := by linarith
  rw [max_eq_left (by positivity), min_eq_left (le_of_lt hgt), one_mul]

\end{lstlisting}
\begin{proof}
Then $r(x) = x$ and $2 - \|x\| > 1$, so $\chi(x) = 1$, giving
$\Eext u (x) = 1 \cdot u(x)$.
\end{proof}

\begin{lemma}[\leanname{unitBallRetraction\_eq\_smul\_of\_annulus}]
If $1 < \|x\| < 2$ then $r(x) = \|x\|^{-2} x$.
\end{lemma}

\begin{lstlisting}[style=Matlab-editor,basicstyle=\mlttfamily,escapechar=",mlshowsectionrules=true,mathescape=true,morecomment={[s]{\%\{}{\%\}}},language=lean,numbers=none]
omit [NeZero d] in
lemma unitBallRetraction_eq_smul_of_annulus {x : E}
    (h1 : 1 < ‖x‖) (h2 : ‖x‖ < 2) :
    unitBallRetraction (d := d) x = (‖x‖ ^ (2 : ℕ))⁻¹ • x := by
  have hx1 : ¬ ‖x‖ ≤ 1 := by linarith
  simp [unitBallRetraction, hx1, h2]

\end{lstlisting}
\begin{proof}
Second branch of the definition.
\end{proof}

\begin{lemma}[\leanname{unitBallCutoff\_eq\_zero\_of\_two\_le\_norm}, \leanname{unitBallExtension\_eq\_zero\_of\_two\_le\_norm}]
If $\|x\| \ge 2$ then $\chi(x) = 0$ and $\Eext u (x) = 0$.
\end{lemma}

\begin{lstlisting}[style=Matlab-editor,basicstyle=\mlttfamily,escapechar=",mlshowsectionrules=true,mathescape=true,morecomment={[s]{\%\{}{\%\}}},language=lean,numbers=none]
omit [NeZero d] in
lemma unitBallCutoff_eq_zero_of_two_le_norm {x : E} (hx : 2 ≤ ‖x‖) :
    unitBallCutoff (d := d) x = 0 := by
  unfold unitBallCutoff
  rw [max_eq_right (by linarith), min_eq_right (by positivity)]

omit [NeZero d] in
lemma unitBallExtension_eq_zero_of_two_le_norm {u : E → ℝ} {x : E}
    (hx : 2 ≤ ‖x‖) :
    unitBallExtension (d := d) u x = 0 := by
  rw [unitBallExtension, unitBallCutoff_eq_zero_of_two_le_norm (d := d) hx, zero_mul]

\end{lstlisting}
\begin{proof}
$2 - \|x\| \le 0$, so $\max(2-\|x\|,0) = 0$ and $\min(1,0) = 0$.
\end{proof}

\begin{lemma}[\leanname{unitBallExtension\_hasCompactSupport}]
$\Eext u$ has compact support contained in $\overline{\Ball{2}{0}}$.
\end{lemma}

\begin{lstlisting}[style=Matlab-editor,basicstyle=\mlttfamily,escapechar=",mlshowsectionrules=true,mathescape=true,morecomment={[s]{\%\{}{\%\}}},language=lean,numbers=none]
omit [NeZero d] in
lemma unitBallExtension_hasCompactSupport (u : E → ℝ) :
    HasCompactSupport (unitBallExtension (d := d) u) := by
  apply HasCompactSupport.intro' (isCompact_closedBall (0 : E) 2) isClosed_closedBall
  intro x hx
  exact unitBallExtension_eq_zero_of_not_mem_closedBall_two (d := d) hx

\end{lstlisting}
\begin{proof}
$\Eext u$ vanishes outside $\overline{\Ball{2}{0}}$ by the previous lemma,
and $\overline{\Ball{2}{0}}$ is compact.
\end{proof}

\subsubsection*{Geometry of inversion}

\begin{definition}[\leanname{unitBallOuterShell}, \leanname{unitBallInnerShell}]
Set
\[
\Sigma^{\mathrm{out}} := \{x \in E : 1 < \|x\| < 2\},
\qquad
\Sigma^{\mathrm{in}} := \{x \in E : 1/2 < \|x\| < 1\}.
\]
Both are open and measurable.
\end{definition}

\begin{lstlisting}[style=Matlab-editor,basicstyle=\mlttfamily,escapechar=",mlshowsectionrules=true,mathescape=true,morecomment={[s]{\%\{}{\%\}}},language=lean,numbers=none]
/-- The outer shell on which the explicit unit-ball extension uses sphere inversion. -/
def unitBallOuterShell : Set E := {x : E | 1 < ‖x‖ ∧ ‖x‖ < 2}

/-- The inner shell that is the image of `unitBallOuterShell` under sphere inversion. -/
def unitBallInnerShell : Set E := {x : E | (1 / 2 : ℝ) < ‖x‖ ∧ ‖x‖ < 1}

\end{lstlisting}
\begin{lemma}[\leanname{unitBallRetraction\_eq\_inversion\_of\_mem\_outerShell}]
For $x \in \Sigma^{\mathrm{out}}$, $r(x)$ equals the spherical inversion of
$x$ in the unit sphere centred at the origin, i.e. $r(x) = x/\|x\|^2$.
\end{lemma}

\begin{lstlisting}[style=Matlab-editor,basicstyle=\mlttfamily,escapechar=",mlshowsectionrules=true,mathescape=true,morecomment={[s]{\%\{}{\%\}}},language=lean,numbers=none]
omit [NeZero d] in
lemma unitBallRetraction_eq_inversion_of_mem_outerShell {x : E}
    (hx : x ∈ unitBallOuterShell (d := d)) :
    unitBallRetraction (d := d) x = EuclideanGeometry.inversion (0 : E) 1 x := by
  rcases hx with ⟨h1, h2⟩
  rw [unitBallRetraction_eq_smul_of_annulus (d := d) h1 h2, EuclideanGeometry.inversion]
  simp [dist_eq_norm, div_eq_mul_inv]

\end{lstlisting}
\begin{proof}
Combine \leanname{unitBallRetraction\_eq\_smul\_of\_annulus} with the formula
for the Euclidean inversion centred at $0$ with radius $1$, which is
$y \mapsto y/\|y\|^2$.
\end{proof}

\begin{lemma}[\leanname{inversion\_mem\_unitBallInnerShell\_of\_mem\_outerShell}, \leanname{inversion\_mem\_unitBallOuterShell\_of\_mem\_innerShell}]
The Euclidean inversion $I(x) := x/\|x\|^2$ maps
$\Sigma^{\mathrm{out}}$ bijectively onto $\Sigma^{\mathrm{in}}$ and vice
versa.
\end{lemma}

\begin{lstlisting}[style=Matlab-editor,basicstyle=\mlttfamily,escapechar=",mlshowsectionrules=true,mathescape=true,morecomment={[s]{\%\{}{\%\}}},language=lean,numbers=none]
omit [NeZero d] in
lemma inversion_mem_unitBallInnerShell_of_mem_outerShell {x : E}
    (hx : x ∈ unitBallOuterShell (d := d)) :
    EuclideanGeometry.inversion (0 : E) 1 x ∈ unitBallInnerShell (d := d) := by
  rcases hx with ⟨h1, h2⟩
  rw [mem_unitBallInnerShell]
  constructor
  · have hxpos : 0 < ‖x‖ := lt_trans zero_lt_one h1
    rw [← dist_zero_right (EuclideanGeometry.inversion (0 : E) 1 x),
      EuclideanGeometry.dist_inversion_center (c := (0 : E)) (x := x) (R := (1 : ℝ)),
      dist_zero_right]
    field_simp [hxpos.ne']
    linarith
  · have hxpos : 0 < ‖x‖ := lt_trans zero_lt_one h1
    rw [← dist_zero_right (EuclideanGeometry.inversion (0 : E) 1 x),
      EuclideanGeometry.dist_inversion_center (c := (0 : E)) (x := x) (R := (1 : ℝ)),
      dist_zero_right]
    field_simp [hxpos.ne']
    linarith

omit [NeZero d] in
lemma inversion_mem_unitBallOuterShell_of_mem_innerShell {x : E}
    (hx : x ∈ unitBallInnerShell (d := d)) :
    EuclideanGeometry.inversion (0 : E) 1 x ∈ unitBallOuterShell (d := d) := by
  rcases hx with ⟨h1, h2⟩
  rw [mem_unitBallOuterShell]
  constructor
  · have hxpos : 0 < ‖x‖ := by linarith
    rw [← dist_zero_right (EuclideanGeometry.inversion (0 : E) 1 x),
      EuclideanGeometry.dist_inversion_center (c := (0 : E)) (x := x) (R := (1 : ℝ)),
      dist_zero_right]
    field_simp [hxpos.ne']
    linarith
  · have hxpos : 0 < ‖x‖ := by linarith
    rw [← dist_zero_right (EuclideanGeometry.inversion (0 : E) 1 x),
      EuclideanGeometry.dist_inversion_center (c := (0 : E)) (x := x) (R := (1 : ℝ)),
      dist_zero_right]
    field_simp [hxpos.ne']
    linarith

\end{lstlisting}
\begin{proof}
If $1 < \|x\| < 2$ then $\|I(x)\| = 1/\|x\| \in (1/2, 1)$, hence
$I(x) \in \Sigma^{\mathrm{in}}$. The opposite direction is symmetric, and
$I \circ I = \mathrm{id}$ on $E\setminus\{0\}$ gives the bijection.
\end{proof}

\begin{lemma}
(\leanname{inversion\_image\_unitBallOuterShell},
\leanname{inversion\_image\_unitBallInnerShell})\\
$I(\Sigma^{\mathrm{out}}) = \Sigma^{\mathrm{in}}$ and
$I(\Sigma^{\mathrm{in}}) = \Sigma^{\mathrm{out}}$.
\end{lemma}

\begin{lstlisting}[style=Matlab-editor,basicstyle=\mlttfamily,escapechar=",mlshowsectionrules=true,mathescape=true,morecomment={[s]{\%\{}{\%\}}},language=lean,numbers=none]
omit [NeZero d] in
lemma inversion_image_unitBallOuterShell :
    EuclideanGeometry.inversion (0 : E) 1 '' unitBallOuterShell (d := d) =
      unitBallInnerShell (d := d) := by
  ext y
  constructor
  · rintro ⟨x, hx, rfl⟩
    exact inversion_mem_unitBallInnerShell_of_mem_outerShell (d := d) hx
  · intro hy
    refine ⟨EuclideanGeometry.inversion (0 : E) 1 y, ?_, ?_⟩
    · exact inversion_mem_unitBallOuterShell_of_mem_innerShell (d := d) hy
    · exact EuclideanGeometry.inversion_inversion (c := (0 : E)) (R := (1 : ℝ)) one_ne_zero y

omit [NeZero d] in
lemma inversion_image_unitBallInnerShell :
    EuclideanGeometry.inversion (0 : E) 1 '' unitBallInnerShell (d := d) =
      unitBallOuterShell (d := d) := by
  ext y
  constructor
  · rintro ⟨x, hx, rfl⟩
    exact inversion_mem_unitBallOuterShell_of_mem_innerShell (d := d) hx
  · intro hy
    refine ⟨EuclideanGeometry.inversion (0 : E) 1 y, ?_, ?_⟩
    · exact inversion_mem_unitBallInnerShell_of_mem_outerShell (d := d) hy
    · exact EuclideanGeometry.inversion_inversion (c := (0 : E)) (R := (1 : ℝ)) one_ne_zero y

\end{lstlisting}
\begin{proof}
Set equalities follow from the previous lemma together with the involution
$I^2 = \mathrm{id}$.
\end{proof}

\begin{lemma}
(\leanname{injOn\_inversion\_unitBallOuterShell},
\leanname{injOn\_inversion\_unitBallInnerShell})\\
The map $I$ is injective on $\Sigma^{\mathrm{out}}$ and on $\Sigma^{\mathrm{in}}$.
\end{lemma}

\begin{lstlisting}[style=Matlab-editor,basicstyle=\mlttfamily,escapechar=",mlshowsectionrules=true,mathescape=true,morecomment={[s]{\%\{}{\%\}}},language=lean,numbers=none]
omit [NeZero d] in
lemma injOn_inversion_unitBallOuterShell :
    InjOn (EuclideanGeometry.inversion (0 : E) 1) (unitBallOuterShell (d := d)) := by
  exact (EuclideanGeometry.inversion_injective (c := (0 : E)) (R := (1 : ℝ)) one_ne_zero).injOn

omit [NeZero d] in
lemma injOn_inversion_unitBallInnerShell :
    InjOn (EuclideanGeometry.inversion (0 : E) 1) (unitBallInnerShell (d := d)) := by
  exact (EuclideanGeometry.inversion_injective (c := (0 : E)) (R := (1 : ℝ)) one_ne_zero).injOn

\end{lstlisting}
\begin{proof}
Both shells are subsets of $E\setminus\{0\}$, on which $I$ is a bijective
involution.
\end{proof}

\begin{lemma}[\leanname{hasFDerivWithinAt\_inversion\_unitBallInnerShell}]
For $x \in \Sigma^{\mathrm{in}}$, $I$ is differentiable at $x$ within
$\Sigma^{\mathrm{in}}$, and the Mathlib derivative formula for spherical
inversion applies.
\end{lemma}

\begin{lstlisting}[style=Matlab-editor,basicstyle=\mlttfamily,escapechar=",mlshowsectionrules=true,mathescape=true,morecomment={[s]{\%\{}{\%\}}},language=lean,numbers=none]
omit [NeZero d] in
lemma hasFDerivWithinAt_inversion_unitBallInnerShell {x : E}
    (hx : x ∈ unitBallInnerShell (d := d)) :
    HasFDerivWithinAt (EuclideanGeometry.inversion (0 : E) 1)
      (((1 / dist x (0 : E)) ^ 2) • ((ℝ ∙ x)ᗮ.reflection : E →L[ℝ] E))
      (unitBallInnerShell (d := d)) x := by
  have hx0 : x ≠ (0 : E) := by
    intro hx0
    rcases hx with ⟨h1, _h2⟩
    have : 0 < ‖x‖ := by linarith
    simp [hx0] at this
  simpa using
    (EuclideanGeometry.hasFDerivAt_inversion (c := (0 : E)) (R := (1 : ℝ)) hx0).hasFDerivWithinAt

\end{lstlisting}
\begin{proof}
$x \neq 0$ since $\|x\| > 1/2$, so the standard inversion derivative
(\leanname{EuclideanGeometry.hasFDerivAt\_inversion})
provides a Fr\'echet derivative at $x$, which restricts to the within-set
derivative.
\end{proof}

\begin{lemma}[\leanname{abs\_det\_fderiv\_inversion\_zero\_one}]
\label{lem:det-jacobian}
For $x \in E \setminus \{0\}$,
\[
\bigl|\det DI(x)\bigr| = \bigl(1/\|x\|\bigr)^{2d}.
\]
\end{lemma}

\begin{lstlisting}[style=Matlab-editor,basicstyle=\mlttfamily,escapechar=",mlshowsectionrules=true,mathescape=true,morecomment={[s]{\%\{}{\%\}}},language=lean,numbers=none]
omit [NeZero d] in
lemma abs_det_fderiv_inversion_zero_one {x : E} (hx : x ≠ (0 : E)) :
    |(fderiv ℝ (EuclideanGeometry.inversion (0 : E) 1) x).det| =
      ((1 / ‖x‖) ^ 2) ^ d := by
  simpa [dist_zero_right, finrank_euclideanSpace_fin] using
    abs_det_fderiv_inversion_zero_one' (d := d) hx

\end{lstlisting}
\begin{proof}
By the closed form of the inversion derivative,
\[
DI(x) = (1/\|x\|)^2 \, R_{(\R x)^\perp},
\]
where $R_{(\R x)^\perp}$ is the reflection across the orthogonal complement of
$\R x$. The determinant of this reflection equals $(-1)^{d-1}$ since the
fixed subspace has dimension $d-1$. Therefore
\[
\det DI(x) = (1/\|x\|)^{2d}\,\det R = (-1)^{d-1}(1/\|x\|)^{2d},
\]
whose absolute value is $(1/\|x\|)^{2d}$.
\end{proof}

\begin{lemma}[\leanname{abs\_det\_fderiv\_inversion\_zero\_one\_le\_of\_mem\_innerShell}]
For $x \in \Sigma^{\mathrm{in}}$, $|\det DI(x)| \le 2^{2d}$.
\end{lemma}

\begin{lstlisting}[style=Matlab-editor,basicstyle=\mlttfamily,escapechar=",mlshowsectionrules=true,mathescape=true,morecomment={[s]{\%\{}{\%\}}},language=lean,numbers=none]
omit [NeZero d] in
lemma abs_det_fderiv_inversion_zero_one_le_of_mem_innerShell {x : E}
    (hx : x ∈ unitBallInnerShell (d := d)) :
    |(fderiv ℝ (EuclideanGeometry.inversion (0 : E) 1) x).det| ≤ (2 : ℝ) ^ (2 * d) := by
  rcases hx with ⟨h1, _h2⟩
  have hx0 : x ≠ (0 : E) := by
    intro hx0
    have : 0 < ‖x‖ := by linarith
    simp [hx0] at this
  rw [abs_det_fderiv_inversion_zero_one (d := d) hx0]
  have hnorm_pos : 0 < ‖x‖ := by linarith
  have hbound : (1 / ‖x‖ : ℝ) ≤ 2 := by
    field_simp [hnorm_pos.ne']
    linarith
  calc
    ((1 / ‖x‖) ^ 2) ^ d ≤ (2 ^ 2 : ℝ) ^ d := by
      gcongr
    _ = (2 : ℝ) ^ (2 * d) := by rw [← pow_mul]

\end{lstlisting}
\begin{proof}
$\|x\| > 1/2$ gives $1/\|x\| < 2$, so $(1/\|x\|)^{2d} < 2^{2d}$.
\end{proof}

\begin{lemma}[\leanname{lintegral\_outerShell\_comp\_inversion\_le}]
\label{lem:cov-shell}
For any measurable $\Phi : E \to \R_{\ge 0}^\infty$,
\[
\int_{\Sigma^{\mathrm{out}}} \Phi(I(y))\,dy
\le 2^{2d}\,\int_{\Sigma^{\mathrm{in}}} \Phi(x)\,dx.
\]
\end{lemma}

\begin{lstlisting}[style=Matlab-editor,basicstyle=\mlttfamily,escapechar=",mlshowsectionrules=true,mathescape=true,morecomment={[s]{\%\{}{\%\}}},language=lean,numbers=none]
omit [NeZero d] in
lemma lintegral_outerShell_comp_inversion_le
    (Φ : E → ℝ≥0∞) :
    ∫⁻ y in unitBallOuterShell (d := d),
        Φ (EuclideanGeometry.inversion (0 : E) 1 y) ∂volume
      ≤ ENNReal.ofReal ((2 : ℝ) ^ (2 * d)) *
        ∫⁻ x in unitBallInnerShell (d := d), Φ x ∂volume := by
  let inv : E → E := EuclideanGeometry.inversion (0 : E) 1
  let f' : E → E →L[ℝ] E := fun x =>
    (((1 / dist x (0 : E)) ^ 2) • ((ℝ ∙ x)ᗮ.reflection : E →L[ℝ] E))
  have himage :
      inv '' unitBallInnerShell (d := d) = unitBallOuterShell (d := d) :=
    inversion_image_unitBallInnerShell (d := d)
  have hcov :=
    MeasureTheory.lintegral_image_eq_lintegral_abs_det_fderiv_mul
      (μ := volume) (s := unitBallInnerShell (d := d)) (f := inv) (f' := f')
      measurableSet_unitBallInnerShell
      (fun x hx => by
        simpa [inv, f'] using hasFDerivWithinAt_inversion_unitBallInnerShell (d := d) hx)
      injOn_inversion_unitBallInnerShell
      (fun y => Φ (inv y))
  rw [himage] at hcov
  calc
    ∫⁻ y in unitBallOuterShell (d := d), Φ (EuclideanGeometry.inversion (0 : E) 1 y) ∂volume
        = ∫⁻ x in unitBallInnerShell (d := d), ENNReal.ofReal |(f' x).det| * Φ (inv (inv x))
            ∂volume := by
              simpa [inv] using hcov
    _ ≤ ∫⁻ x in unitBallInnerShell (d := d),
          ENNReal.ofReal ((2 : ℝ) ^ (2 * d)) * Φ (inv (inv x)) ∂volume := by
            refine lintegral_mono_ae ?_
            rw [ae_restrict_iff' measurableSet_unitBallInnerShell]
            refine Filter.Eventually.of_forall ?_
            intro x hx
            have hweight :
                ENNReal.ofReal |(f' x).det| ≤ ENNReal.ofReal ((2 : ℝ) ^ (2 * d)) := by
              have hx0 : x ≠ (0 : E) := by
                intro hx0
                have : 0 < ‖x‖ := by
                  rcases hx with ⟨h1, _h2⟩
                  linarith
                simp [hx0] at this
              have hfderiv : f' x = fderiv ℝ (EuclideanGeometry.inversion (0 : E) 1) x := by
                have hfdx :
                    HasFDerivAt (EuclideanGeometry.inversion (0 : E) 1) (f' x) x := by
                  simpa [f'] using
                    (EuclideanGeometry.hasFDerivAt_inversion (c := (0 : E)) (R := (1 : ℝ)) hx0)
                exact hfdx.fderiv.symm
              exact ENNReal.ofReal_le_ofReal
                (by simpa [hfderiv] using
                  abs_det_fderiv_inversion_zero_one_le_of_mem_innerShell (d := d) hx)
            exact mul_le_mul' hweight le_rfl
    _ = ∫⁻ x in unitBallInnerShell (d := d),
          ENNReal.ofReal ((2 : ℝ) ^ (2 * d)) * Φ x ∂volume := by
            congr with x
            simp [inv, EuclideanGeometry.inversion_inversion, (by norm_num : (1 : ℝ) ≠ 0)]
    _ = ENNReal.ofReal ((2 : ℝ) ^ (2 * d)) *
          ∫⁻ x in unitBallInnerShell (d := d), Φ x ∂volume := by
            rw [lintegral_const_mul']
            simp

\end{lstlisting}
\begin{proof}
By the change-of-variables formula for the differentiable injection $I$ on
$\Sigma^{\mathrm{in}}$ onto $\Sigma^{\mathrm{out}}$,
\[
\int_{\Sigma^{\mathrm{out}}} \Phi(I(y))\,dy
= \int_{\Sigma^{\mathrm{in}}} |\det DI(x)|\, \Phi(I(I(x)))\,dx
= \int_{\Sigma^{\mathrm{in}}} |\det DI(x)|\, \Phi(x)\,dx,
\]
using $I^2 = \mathrm{id}$. Now Lemma~\ref{lem:det-jacobian} together with the
bound from the inner shell gives $|\det DI(x)| \le 2^{2d}$ a.e., and the
result follows.
\end{proof}

\begin{lemma}[\leanname{unitBallCutoff\_eq\_two\_sub\_norm\_of\_mem\_outerShell}]
For $x \in \Sigma^{\mathrm{out}}$, $\chi(x) = 2 - \|x\|$.
\end{lemma}

\begin{lstlisting}[style=Matlab-editor,basicstyle=\mlttfamily,escapechar=",mlshowsectionrules=true,mathescape=true,morecomment={[s]{\%\{}{\%\}}},language=lean,numbers=none]
omit [NeZero d] in
lemma unitBallCutoff_eq_two_sub_norm_of_mem_outerShell {x : E}
    (hx : x ∈ unitBallOuterShell (d := d)) :
    unitBallCutoff (d := d) x = 2 - ‖x‖ := by
  rcases hx with ⟨h1, h2⟩
  unfold unitBallCutoff
  have hleft : 2 - ‖x‖ ≤ 1 := by linarith
  have hright : 0 ≤ 2 - ‖x‖ := by linarith
  rw [max_eq_left hright, min_eq_right hleft]

\end{lstlisting}
\begin{proof}
$0 < 2 - \|x\| < 1$, so the $\min$ and $\max$ in the definition reduce to
$2 - \|x\|$.
\end{proof}

\begin{lemma}[\leanname{norm\_fderiv\_inversion\_zero\_one}, \leanname{norm\_fderiv\_inversion\_le\_one\_of\_mem\_outerShell}]
\label{lem:norm-DI}
For $x \neq 0$, $\|DI(x)\| \le (1/\|x\|)^2$. In particular, for
$x \in \Sigma^{\mathrm{out}}$ we have $\|DI(x)\| \le 1$.
\end{lemma}

\begin{lstlisting}[style=Matlab-editor,basicstyle=\mlttfamily,escapechar=",mlshowsectionrules=true,mathescape=true,morecomment={[s]{\%\{}{\%\}}},language=lean,numbers=none]
omit [NeZero d] in
lemma norm_fderiv_inversion_zero_one {x : E} (hx : x ≠ (0 : E)) :
    ‖fderiv ℝ (EuclideanGeometry.inversion (0 : E) 1) x‖ ≤ (1 / ‖x‖) ^ 2 := by
  have hfd :
      HasFDerivAt (EuclideanGeometry.inversion (0 : E) 1)
        (((1 / dist x (0 : E)) ^ 2) • ((ℝ ∙ x)ᗮ.reflection : E →L[ℝ] E)) x := by
    simpa using
      (EuclideanGeometry.hasFDerivAt_inversion (c := (0 : E)) (R := (1 : ℝ)) hx)
  have hreflect :
      ‖((ℝ ∙ x)ᗮ.reflection : E →L[ℝ] E)‖ ≤ 1 := by
    refine (ContinuousLinearMap.opNorm_le_iff
      (f := ((ℝ ∙ x)ᗮ.reflection : E →L[ℝ] E)) (M := 1) (by positivity)).2 ?_
    intro y
    simp [((ℝ ∙ x)ᗮ).reflection.norm_map y]
  calc
    ‖fderiv ℝ (EuclideanGeometry.inversion (0 : E) 1) x‖
        = ‖(((1 / dist x (0 : E)) ^ 2) • ((ℝ ∙ x)ᗮ.reflection : E →L[ℝ] E))‖ := by
            rw [hfd.fderiv]
    _ = |(1 / dist x (0 : E)) ^ 2| * ‖((ℝ ∙ x)ᗮ.reflection : E →L[ℝ] E)‖ := by
            rw [norm_smul, Real.norm_eq_abs]
    _ ≤ |(1 / dist x (0 : E)) ^ 2| * 1 := by
          gcongr
    _ = (1 / ‖x‖) ^ 2 := by
          rw [abs_of_nonneg]
          · simp [dist_zero_right]
          · positivity

omit [NeZero d] in
lemma norm_fderiv_inversion_le_one_of_mem_outerShell {x : E}
    (hx : x ∈ unitBallOuterShell (d := d)) :
    ‖fderiv ℝ (EuclideanGeometry.inversion (0 : E) 1) x‖ ≤ 1 := by
  rcases hx with ⟨h1, _h2⟩
  have hx0 : x ≠ (0 : E) := by
    intro hx0
    have : 0 < ‖x‖ := by linarith
    simp [hx0] at this
  have hnorm_ge_one : 1 ≤ ‖x‖ := by linarith
  have hle : (1 / ‖x‖ : ℝ) ≤ 1 := by
    simpa [one_div] using (inv_le_one_of_one_le₀ hnorm_ge_one)
  have hsq : (1 / ‖x‖ : ℝ) ^ 2 ≤ 1 := by
    have hnn : 0 ≤ (1 / ‖x‖ : ℝ) := by positivity
    have hsq' := mul_self_le_mul_self hnn hle
    simpa [sq] using hsq'
  exact (norm_fderiv_inversion_zero_one (d := d) hx0).trans hsq

\end{lstlisting}
\begin{proof}
$DI(x) = (1/\|x\|)^2 R$ with $R$ a reflection of operator norm $1$, so
$\|DI(x)\| \le (1/\|x\|)^2$. For $\|x\| \ge 1$ this is at most $1$.
\end{proof}

\begin{lemma}
(\leanname{hasFDerivAt\_unitBallCutoff\_of\_mem\_outerShell},
\leanname{norm\_fderiv\_unitBallCutoff\_of\_mem\_outerShell})\\
For $x \in \Sigma^{\mathrm{out}}$, $\chi$ is differentiable at $x$ with
$D\chi(x) = -D(\|\cdot\|)(x)$, and $\|D\chi(x)\| = 1$.
\end{lemma}

\begin{lstlisting}[style=Matlab-editor,basicstyle=\mlttfamily,escapechar=",mlshowsectionrules=true,mathescape=true,morecomment={[s]{\%\{}{\%\}}},language=lean,numbers=none]
omit [NeZero d] in
lemma hasFDerivAt_unitBallCutoff_of_mem_outerShell {x : E}
    (hx : x ∈ unitBallOuterShell (d := d)) :
    HasFDerivAt (unitBallCutoff (d := d)) (-(fderiv ℝ (fun y : E => ‖y‖) x)) x := by
  have hx0 : x ≠ (0 : E) := by
    intro hx0
    rcases hx with ⟨h1, _h2⟩
    have : 0 < ‖x‖ := by linarith
    simp [hx0] at this
  have hnorm : DifferentiableAt ℝ (fun y : E => ‖y‖) x :=
    differentiableAt_id.norm ℝ hx0
  have hfd_annulus :
      HasFDerivAt (fun y : E => 2 - ‖y‖) (-(fderiv ℝ (fun y : E => ‖y‖) x)) x := by
    simpa using hnorm.hasFDerivAt.const_sub (2 : ℝ)
  have heq :
      unitBallCutoff (d := d) =ᶠ[𝓝 x] fun y : E => 2 - ‖y‖ := by
    filter_upwards [isOpen_unitBallOuterShell (d := d) |>.mem_nhds hx] with y hy
    exact unitBallCutoff_eq_two_sub_norm_of_mem_outerShell (d := d) hy
  exact hfd_annulus.congr_of_eventuallyEq heq

lemma norm_fderiv_unitBallCutoff_of_mem_outerShell {x : E}
    (hx : x ∈ unitBallOuterShell (d := d)) :
    ‖fderiv ℝ (unitBallCutoff (d := d)) x‖ = 1 := by
  have hx0 : x ≠ (0 : E) := by
    intro hx0
    rcases hx with ⟨h1, _h2⟩
    have : 0 < ‖x‖ := by linarith
    simp [hx0] at this
  have hnorm : DifferentiableAt ℝ (fun y : E => ‖y‖) x :=
    differentiableAt_id.norm ℝ hx0
  have hfd_cutoff :=
    hasFDerivAt_unitBallCutoff_of_mem_outerShell (d := d) hx
  rw [hfd_cutoff.fderiv, norm_neg]
  simpa using norm_fderiv_norm hnorm

\end{lstlisting}
\begin{proof}
On $\Sigma^{\mathrm{out}}$, $\chi$ agrees with $y \mapsto 2 - \|y\|$, which
is smooth on $E\setminus\{0\}$. Differentiating gives
$D\chi = -D\|\cdot\|$, and the norm of the gradient of the Euclidean norm at
any nonzero point equals $1$.
\end{proof}

\begin{lemma}[\leanname{norm\_fderiv\_unitBallExtension\_le\_of\_mem\_outerShell}]
\label{lem:Eext-grad-bound}
Let $u \in C^1(E)$. For $x \in \Sigma^{\mathrm{out}}$,
\[
\|D\Eext u(x)\|
\le |u(I(x))| + \|Du(I(x))\|.
\]
\end{lemma}

\begin{lstlisting}[style=Matlab-editor,basicstyle=\mlttfamily,escapechar=",mlshowsectionrules=true,mathescape=true,morecomment={[s]{\%\{}{\%\}}},language=lean,numbers=none]
lemma norm_fderiv_unitBallExtension_le_of_mem_outerShell
    {u : E → ℝ} (hu : Differentiable ℝ u) {x : E}
    (hx : x ∈ unitBallOuterShell (d := d)) :
    ‖fderiv ℝ (unitBallExtension (d := d) u) x‖ ≤
      |u (EuclideanGeometry.inversion (0 : E) 1 x)| +
      ‖fderiv ℝ u (EuclideanGeometry.inversion (0 : E) 1 x)‖ := by
  let inv : E → E := EuclideanGeometry.inversion (0 : E) 1
  have hx0 : x ≠ (0 : E) := by
    intro hx0
    rcases hx with ⟨h1, _h2⟩
    have : 0 < ‖x‖ := by linarith
    simp [hx0] at this
  have hcutoff :
      HasFDerivAt (unitBallCutoff (d := d)) (fderiv ℝ (unitBallCutoff (d := d)) x) x := by
    simpa [(hasFDerivAt_unitBallCutoff_of_mem_outerShell (d := d) hx).fderiv] using
      (hasFDerivAt_unitBallCutoff_of_mem_outerShell (d := d) hx)
  have hinv :
      HasFDerivAt inv (fderiv ℝ inv x) x := by
    simpa [inv, (EuclideanGeometry.hasFDerivAt_inversion (c := (0 : E)) (R := (1 : ℝ)) hx0).fderiv] using
      (EuclideanGeometry.hasFDerivAt_inversion (c := (0 : E)) (R := (1 : ℝ)) hx0)
  have hcomp :
      HasFDerivAt (fun y : E => u (inv y))
        ((fderiv ℝ u (inv x)).comp (fderiv ℝ inv x)) x := by
    exact (hu (inv x)).hasFDerivAt.comp x hinv
  have hprod := hcutoff.mul hcomp
  have hext :
      unitBallExtension (d := d) u =ᶠ[𝓝 x] fun y : E => unitBallCutoff y * u (inv y) := by
    filter_upwards [isOpen_unitBallOuterShell (d := d) |>.mem_nhds hx] with y hy
    rw [unitBallExtension, unitBallRetraction_eq_inversion_of_mem_outerShell (d := d) hy]
  have hprod_ext :
      HasFDerivAt (unitBallExtension (d := d) u)
        (unitBallCutoff x • ((fderiv ℝ u (inv x)).comp (fderiv ℝ inv x)) +
          u (inv x) • fderiv ℝ (unitBallCutoff (d := d)) x) x := by
    simpa [mul_comm, add_comm, add_left_comm, add_assoc] using hprod.congr_of_eventuallyEq hext
  rw [hprod_ext.fderiv]
  calc
    ‖unitBallCutoff x • (fderiv ℝ u (inv x)).comp (fderiv ℝ inv x) +
        u (inv x) • fderiv ℝ (unitBallCutoff (d := d)) x‖
        ≤ ‖unitBallCutoff x • (fderiv ℝ u (inv x)).comp (fderiv ℝ inv x)‖ +
          ‖u (inv x) • fderiv ℝ (unitBallCutoff (d := d)) x‖ := norm_add_le _ _
    _ = |unitBallCutoff x| * ‖(fderiv ℝ u (inv x)).comp (fderiv ℝ inv x)‖ +
          |u (inv x)| * ‖fderiv ℝ (unitBallCutoff (d := d)) x‖ := by
            rw [norm_smul, norm_smul, Real.norm_eq_abs, Real.norm_eq_abs]
    _ ≤ |unitBallCutoff x| * (‖fderiv ℝ u (inv x)‖ * ‖fderiv ℝ inv x‖) +
          |u (inv x)| * ‖fderiv ℝ (unitBallCutoff (d := d)) x‖ := by
            gcongr
            exact ContinuousLinearMap.opNorm_comp_le _ _
    _ ≤ |unitBallCutoff x| * ‖fderiv ℝ u (inv x)‖ + |u (inv x)| * 1 := by
          gcongr
          · exact mul_le_of_le_one_right (norm_nonneg _) <|
              norm_fderiv_inversion_le_one_of_mem_outerShell (d := d) hx
          · exact by
              rw [norm_fderiv_unitBallCutoff_of_mem_outerShell (d := d) hx]
    _ ≤ ‖fderiv ℝ u (inv x)‖ + |u (inv x)| := by
          have hcut : |unitBallCutoff x| ≤ 1 := by
            have hcut_nonneg : 0 ≤ unitBallCutoff x := by
              unfold unitBallCutoff
              exact le_min (by positivity) (le_max_right _ _)
            rw [abs_of_nonneg hcut_nonneg]
            have hcut_le : unitBallCutoff x ≤ 1 := by
              unfold unitBallCutoff
              exact min_le_left _ _
            exact hcut_le
          nlinarith [norm_nonneg (fderiv ℝ u (inv x)), abs_nonneg (u (inv x))]
    _ = |u (inv x)| + ‖fderiv ℝ u (inv x)‖ := by ring

\end{lstlisting}
\begin{proof}
On $\Sigma^{\mathrm{out}}$, $\Eext u(y) = \chi(y)\, u(I(y))$. By the product
rule and the chain rule,
\[
D\Eext u (x) = \chi(x)\, Du(I(x))\circ DI(x) + u(I(x))\, D\chi(x).
\]
Apply $\|\chi(x)\| \le 1$, $\|DI(x)\| \le 1$ (Lemma~\ref{lem:norm-DI}), and
$\|D\chi(x)\| = 1$ (previous lemma) and the operator-norm submultiplicativity
to obtain
$\|D\Eext u(x)\| \le \|Du(I(x))\| + |u(I(x))|$.
\end{proof}

\begin{lemma}[\leanname{ball\_zero\_two\_eq\_ball\_zero\_one\_union\_outerShell\_union\_sphere}]
The ball $\Ball{2}{0}$ decomposes as
$\Bz{1} \cup \Sigma^{\mathrm{out}} \cup \mathrm{S}_1$, where $\mathrm{S}_1$ is
the unit sphere; the union is disjoint up to the sphere, which has measure
zero.
\end{lemma}

\begin{lstlisting}[style=Matlab-editor,basicstyle=\mlttfamily,escapechar=",mlshowsectionrules=true,mathescape=true,morecomment={[s]{\%\{}{\%\}}},language=lean,numbers=none]
omit [NeZero d] in
lemma ball_zero_two_eq_ball_zero_one_union_outerShell_union_sphere :
    Metric.ball (0 : E) 2 =
      Metric.ball (0 : E) 1 ∪ unitBallOuterShell (d := d) ∪ Metric.sphere (0 : E) 1 := by
  ext x
  constructor
  · intro hx
    rw [Metric.mem_ball, dist_zero_right] at hx
    by_cases h1 : ‖x‖ < 1
    · exact Or.inl <| Or.inl <| by simpa [Metric.mem_ball, dist_zero_right] using h1
    · have h1' : 1 ≤ ‖x‖ := by linarith
      by_cases h1eq : ‖x‖ = 1
      · exact Or.inr <| by simpa [Metric.mem_sphere, dist_zero_right] using h1eq
      · have h1lt : 1 < ‖x‖ := lt_of_le_of_ne h1' (by simpa [eq_comm] using h1eq)
        exact Or.inl <| Or.inr ⟨h1lt, hx⟩
  · rintro ((hx | hx) | hx)
    · rw [Metric.mem_ball, dist_zero_right] at hx ⊢
      exact lt_trans hx (by norm_num : (1 : ℝ) < 2)
    · rw [Metric.mem_ball, dist_zero_right]
      exact hx.2
    · have hsphere : ‖x‖ = 1 := by simpa [Metric.mem_sphere, dist_zero_right] using hx
      rw [Metric.mem_ball, dist_zero_right]
      linarith

\end{lstlisting}
\begin{proof}
A point $y$ with $\|y\| < 2$ is either in $\Bz{1}$ if $\|y\| < 1$, on the
sphere if $\|y\| = 1$, or in $\Sigma^{\mathrm{out}}$ if $1 < \|y\| < 2$.
\end{proof}

\begin{lemma}[\leanname{fderiv\_unitBallExtension\_eq\_of\_mem\_ball}]
For $u$ differentiable and $x \in \Bz{1}$, $D\Eext u (x) = Du(x)$.
\end{lemma}

\begin{lstlisting}[style=Matlab-editor,basicstyle=\mlttfamily,escapechar=",mlshowsectionrules=true,mathescape=true,morecomment={[s]{\%\{}{\%\}}},language=lean,numbers=none]
omit [NeZero d] in
lemma fderiv_unitBallExtension_eq_of_mem_ball
    {u : E → ℝ} {x : E} (hx : x ∈ Metric.ball (0 : E) 1) :
    fderiv ℝ (unitBallExtension (d := d) u) x = fderiv ℝ u x := by
  apply Filter.EventuallyEq.fderiv_eq
  filter_upwards [Metric.isOpen_ball.mem_nhds hx] with y hy
  exact unitBallExtension_eq_of_mem_ball (d := d) hy

\end{lstlisting}
\begin{proof}
On $\Bz{1}$, $\Eext u$ agrees with $u$ (Lemma~\ref{lem:Eext-restrict}).
Locally constant cutoff and identity retraction give the derivative
identity.
\end{proof}

\subsubsection*{Smooth-input extension estimates}

\begin{theorem}[\leanname{smooth\_unitBallExtension\_W1p\_estimate}]
\label{thm:smooth-extension-estimate}
Let $1 \le p < \infty$ and $u \in C^1(E)$. Then
\begin{align*}
\int_{\Ball{2}{0}} |\Eext u (x)|^p\,dx
&\le (1 + 2^{2d})\, \int_{\Bz{1}} |u(x)|^p\,dx, \\
\int_{\Ball{2}{0}} \|D\Eext u (x)\|^p\,dx
&\le \int_{\Bz{1}} \|Du(x)\|^p\,dx \\
&\quad + 2^{2d}\,2^{p-1}\,\biggl(\int_{\Bz{1}} |u(x)|^p\,dx
+ \int_{\Bz{1}} \|Du(x)\|^p\,dx\biggr).
\end{align*}
\end{theorem}

\begin{lstlisting}[style=Matlab-editor,basicstyle=\mlttfamily,escapechar=",mlshowsectionrules=true,mathescape=true,morecomment={[s]{\%\{}{\%\}}},language=lean,numbers=none]
theorem smooth_unitBallExtension_W1p_estimate
    {p : ℝ} (hp : 1 ≤ p) {u : E → ℝ} (hu : ContDiff ℝ 1 u) :
    (∫⁻ x in Metric.ball (0 : E) 2,
        (ENNReal.ofReal |unitBallExtension (d := d) u x|) ^ p ∂volume)
      ≤ (1 + ENNReal.ofReal ((2 : ℝ) ^ (2 * d))) *
        ∫⁻ x in Metric.ball (0 : E) 1, (ENNReal.ofReal |u x|) ^ p ∂volume
    ∧
    (∫⁻ x in Metric.ball (0 : E) 2,
        (ENNReal.ofReal ‖fderiv ℝ (unitBallExtension (d := d) u) x‖) ^ p ∂volume)
      ≤ (∫⁻ x in Metric.ball (0 : E) 1, (ENNReal.ofReal ‖fderiv ℝ u x‖) ^ p ∂volume) +
        ENNReal.ofReal ((2 : ℝ) ^ (2 * d)) * (2 : ℝ≥0∞) ^ (p - 1) *
          (∫⁻ x in Metric.ball (0 : E) 1, (ENNReal.ofReal |u x|) ^ p ∂volume +
           ∫⁻ x in Metric.ball (0 : E) 1, (ENNReal.ofReal ‖fderiv ℝ u x‖) ^ p ∂volume) := by
  exact ⟨lintegral_rpow_abs_unitBallExtension_ball_two_le (d := d) (by linarith : 0 ≤ p),
    lintegral_rpow_norm_fderiv_unitBallExtension_ball_two_le (d := d) hp hu⟩

\end{lstlisting}
\begin{proof}
Decompose $\Ball{2}{0} = \Bz{1} \cup \Sigma^{\mathrm{out}} \cup \mathrm{S}_1$,
the sphere having measure zero. On $\Bz{1}$, $\Eext u = u$ and
$D\Eext u = Du$, so the integrals over $\Bz{1}$ contribute exactly the
$u$-side. On $\Sigma^{\mathrm{out}}$, use the shell formula
$\Eext u (x) = (2-\|x\|)\, u(I(x))$ with $|2-\|x\|| \le 1$. By
Lemma~\ref{lem:cov-shell},
\[
\int_{\Sigma^{\mathrm{out}}} |u(I(x))|^p\,dx
\le 2^{2d}\,\int_{\Sigma^{\mathrm{in}}} |u(x)|^p\,dx
\le 2^{2d}\,\int_{\Bz{1}} |u(x)|^p\,dx,
\]
which yields the function-side bound.

For the gradient side, use Lemma~\ref{lem:Eext-grad-bound} to dominate
\[
\|D\Eext u(x)\|^p \le 2^{p-1}\,(|u(I(x))|^p + \|Du(I(x))\|^p)
\quad (x \in \Sigma^{\mathrm{out}}),
\]
via the convexity inequality $(a+b)^p \le 2^{p-1}(a^p+b^p)$. Apply
Lemma~\ref{lem:cov-shell} to each of the two integrands and combine with
the contribution from $\Bz{1}$, where $D\Eext u = Du$.
\end{proof}

\subsubsection*{Smooth approximation of the extension on a rough input}

\begin{definition}[\leanname{unitBallShellFormula}]
For $u : E \to \R$, set
\[
\unitBallShellFormula u(x) := (2 - \|x\|)\, u(I(x)).
\]
For $x \in \Sigma^{\mathrm{out}}$, $\Eext u (x) = \unitBallShellFormula u(x)$.
\end{definition}

\begin{lstlisting}[style=Matlab-editor,basicstyle=\mlttfamily,escapechar=",mlshowsectionrules=true,mathescape=true,morecomment={[s]{\%\{}{\%\}}},language=lean,numbers=none]
def unitBallShellFormula (u : E → ℝ) : E → ℝ :=
  fun x => (2 - ‖x‖) * u (EuclideanGeometry.inversion (0 : E) 1 x)

\end{lstlisting}
\begin{lemma}[\leanname{contDiffOn\_unitBallShellFormula}]
If $u \in C^1(E)$, then the function $\unitBallShellFormula u$ is of
class $C^1$ on $\Sigma^{\mathrm{out}}$.
\end{lemma}

\begin{lstlisting}[style=Matlab-editor,basicstyle=\mlttfamily,escapechar=",mlshowsectionrules=true,mathescape=true,morecomment={[s]{\%\{}{\%\}}},language=lean,numbers=none]
omit [NeZero d] in
lemma contDiffOn_unitBallShellFormula
    {u : E → ℝ} (hu : ContDiff ℝ 1 u) :
    ContDiffOn ℝ 1 (unitBallShellFormula (d := d) u) (unitBallOuterShell (d := d)) := by
  have hnorm :
      ContDiffOn ℝ 1 (fun x : E => ‖x‖) (unitBallOuterShell (d := d)) := by
    refine ContDiffOn.norm (𝕜 := ℝ) contDiffOn_id ?_
    intro x hx
    have hx0 : 0 < ‖x‖ := by linarith [hx.1]
    intro hxz
    simp [hxz] at hx0
  have hscalar :
      ContDiffOn ℝ 1 (fun x : E => 2 - ‖x‖) (unitBallOuterShell (d := d)) :=
    contDiffOn_const.sub hnorm
  have hinv :
      ContDiffOn ℝ 1 (fun x : E => EuclideanGeometry.inversion (0 : E) 1 x)
        (unitBallOuterShell (d := d)) := by
    simpa using
      (contDiffOn_const.inversion contDiffOn_const contDiffOn_id
        (fun x hx => by
          have hx0 : 0 < ‖x‖ := by linarith [hx.1]
          have hxne_norm : ‖x‖ ≠ 0 := ne_of_gt hx0
          intro hxz
          exact hxne_norm (by simpa [hxz])))
  exact hscalar.mul (hu.comp_contDiffOn hinv)

\end{lstlisting}
\begin{proof}
$x \mapsto 2 - \|x\|$ and $x \mapsto I(x)$ are $C^1$ on $\Sigma^{\mathrm{out}}$
(both away from $0$), and $u$ is $C^1$ on $E$. Composition and product give
the result.
\end{proof}

\begin{definition}[\leanname{smoothUnitBallExtensionGrad}]
For $u : E \to \R$ smooth, define
$\smoothUnitBallExtensionGrad u (x) \in \R^d$ by its components
\[
  (\smoothUnitBallExtensionGrad u(x))_i := (D\Eext u (x))(e_i).
\]
\end{definition}

\begin{lstlisting}[style=Matlab-editor,basicstyle=\mlttfamily,escapechar=",mlshowsectionrules=true,mathescape=true,morecomment={[s]{\%\{}{\%\}}},language=lean,numbers=none]
def smoothUnitBallExtensionGrad (u : E → ℝ) : E → E :=
  fun x => WithLp.toLp 2 fun i =>
    (fderiv ℝ (unitBallExtension (d := d) u) x) (EuclideanSpace.single i 1)

\end{lstlisting}
\begin{theorem}[\leanname{exists\_global\_smooth\_W1p\_approx\_of\_localizedWitness}]
\label{thm:loc-global}
Let $\Omega \subseteq E$ be open, $1 < p < \infty$, and let $u \in W^{1,p}(\Omega)$
have compact support contained in $\Omega$. Then there is a sequence $\psi_n
\in C^\infty_c(E)$ with
\[
\eLpNormText{\psi_n - u}_{L^p(E)} \to 0
\quad\text{and}\quad
\eLpNormText{\partial_i \psi_n - \mathbf{1}_\Omega (\partial_i u)}_{L^p(E)} \to 0
\]
for each $i$.
\end{theorem}

\begin{lstlisting}[style=Matlab-editor,basicstyle=\mlttfamily,escapechar=",mlshowsectionrules=true,mathescape=true,morecomment={[s]{\%\{}{\%\}}},language=lean,numbers=none]
/-- Global smooth compactly supported approximation for a local witness whose
support is already compactly contained in the source open set. -/
theorem exists_global_smooth_W1p_approx_of_localizedWitness
    {Ω : Set E} (hΩ : IsOpen Ω)
    {p : ℝ} (hp : 1 < p)
    {u : E → ℝ}
    (hw : MemW1pWitness (ENNReal.ofReal p) u Ω)
    (hu_compact : HasCompactSupport u)
    (hu_sub : tsupport u ⊆ Ω) :
    ∃ ψ : ℕ → E → ℝ,
      (∀ n, ContDiff ℝ (⊤ : ℕ∞) (ψ n)) ∧
      (∀ n, HasCompactSupport (ψ n)) ∧
      Tendsto
        (fun n => eLpNorm (fun x => ψ n x - u x) (ENNReal.ofReal p) volume)
        atTop (nhds 0) ∧
      (∀ i : Fin d,
        Tendsto
          (fun n =>
            eLpNorm
              (fun x => (fderiv ℝ (ψ n) x) (EuclideanSpace.single i 1) -
                Ω.indicator (fun y => hw.weakGrad y i) x)
              (ENNReal.ofReal p) volume)
          atTop (nhds 0)) := by
  have hu0 : MemW01p (ENNReal.ofReal p) u Ω :=
    memW01p_of_memW1p_of_tsupport_subset hΩ hp hw.memW1p hu_compact hu_sub
  let hwExt := zeroExtend_memW1pWitness_p (d := d) hΩ hp hu0 hw
  have hu_eq_ind : Ω.indicator u = u := by
    ext x
    by_cases hx : x ∈ Ω
    · simp [hx]
    · simp [hx, zero_outside_of_tsupport_subset (Ω := Ω) hu_sub hx]
  have huExt0 : MemW01p (ENNReal.ofReal p) (Ω.indicator u) Set.univ :=
    zeroExtend_memW01p_p (d := d) hΩ hp hu0
  rcases huExt0.2 with ⟨hw0, φ, hφ_smooth, hφ_compact, hφ_sub, hφ_fun, hφ_grad⟩
  refine ⟨φ, hφ_smooth, hφ_compact, ?_, ?_⟩
  · simpa [hu_eq_ind] using hφ_fun
  · intro i
    have hp_le : 1 ≤ p := le_of_lt hp
    have hcomp :
        (fun x => hw0.weakGrad x i) =ᵐ[volume.restrict Set.univ]
          (fun x => (Ω.indicator (fun y => hw.weakGrad y i)) x) := by
      filter_upwards [MemW1pWitness.ae_eq_p (d := d) isOpen_univ hp_le hw0 hwExt] with x hx
      simpa [zeroExtend_memW1pWitness_p, hwExt] using congrArg (fun z : E => z i) hx
    have hEqSeq :
        (fun n =>
          eLpNorm
            (fun x => (fderiv ℝ (φ n) x) (EuclideanSpace.single i 1) -
              Ω.indicator (fun y => hw.weakGrad y i) x)
            (ENNReal.ofReal p) volume) =
          (fun n =>
            eLpNorm
              (fun x => (fderiv ℝ (φ n) x) (EuclideanSpace.single i 1) - hw0.weakGrad x i)
              (ENNReal.ofReal p) volume) := by
      funext n
      refine eLpNorm_congr_ae ?_
      have hcomp' :
          (fun x => hw0.weakGrad x i) =ᵐ[volume]
            (fun x => (Ω.indicator (fun y => hw.weakGrad y i)) x) := by
        simpa [Measure.restrict_univ] using hcomp
      filter_upwards [hcomp'] with x hx
      simp [hx]
    rw [hEqSeq]
    simpa [Measure.restrict_univ] using hφ_grad i

\end{lstlisting}
\begin{proof}
Since $\tsupport u$ is compact and contained in $\Omega$, the criterion
\leanname{memW01p\_of\_memW1p\_of\_tsupport\_subset} promotes the local
$W^{1,p}(\Omega)$ membership of $u$ to a $W^{1,p}_0(\Omega)$ membership: $u$
admits an approximating sequence of compactly supported smooth functions
inside $\Omega$. Pushing this through the zero-extension construction
\leanname{zeroExtend\_memW01p\_p} produces a global $W^{1,p}(E)$ witness for
the trivial extension of $u$ whose weak gradient is $\mathbf{1}_\Omega
\nabla u$, since classical partial derivatives of the zero-extended smooth
approximants vanish off $\Omega$ and converge in $L^p(E)$ to the indicator
of the original gradient.

The global density theorem
Theorem~\ref{thm:smooth_approx_univ}
(\leanname{exists\_smooth\_compactSupport\_W1p\_approx\_univ}) applied to the
extended witness on the open set $E$ then produces a sequence
$\psi_n \in C^\infty_c(E)$ with $\psi_n \to u$ in $L^p(E)$ and
$\partial_i \psi_n \to (\mathbf{1}_\Omega \nabla u)_i$ in $L^p(E)$ for each
component $i$. The compactness of $\tsupport u$ ensures that the gradient
limit is concentrated on $\Omega$, exactly as stated. The Lean source
additionally records a Fatou-type semicontinuity bound: a subsequence of the
$\psi_n$ converges a.e., and Fatou's lemma on $\|\cdot\|^p$ bounds the
$L^p$-energy of the limiting gradient by the liminf of the energies of the
approximants.
\end{proof}

\subsubsection*{Smooth blends and approximation control}

\begin{definition}[\leanname{sphereOneBlend}, \leanname{sphereTwoBlend}]
For $\varepsilon > 0$, define $\sigma^1_\varepsilon(x), \sigma^2_\varepsilon(x)
\in [0,1]$ via the smooth transition
\[
\sigma^1_\varepsilon(x) := \mathrm{st}\!\left(\frac{\|x\|-1+\varepsilon}{2\varepsilon}\right),
\qquad
\sigma^2_\varepsilon(x) := 1 - \mathrm{st}\!\left(\frac{\|x\|-2+\varepsilon}{2\varepsilon}\right),
\]
which transition smoothly across the spheres $\|x\|=1$ and $\|x\|=2$.
\end{definition}

\begin{lstlisting}[style=Matlab-editor,basicstyle=\mlttfamily,escapechar=",mlshowsectionrules=true,mathescape=true,morecomment={[s]{\%\{}{\%\}}},language=lean,numbers=none]
/-- Transition profile used to smooth the explicit unit-ball extension across
the sphere interfaces `‖x‖ = 1` and `‖x‖ = 2`. -/
def sphereOneBlend (ε : ℝ) : E → ℝ :=
  fun x => 1 - myCutoff (0 : E) (1 - ε) (1 + ε) x

def sphereTwoBlend (ε : ℝ) : E → ℝ :=
  fun x => myCutoff (0 : E) (2 - ε) (2 + ε) x

\end{lstlisting}
\begin{definition}[\leanname{smoothUnitBallExtensionApprox}]
For $\varepsilon > 0$ and $\psi : E \to \R$ smooth, define
\begin{multline*}
\smoothUnitBallExtensionApprox_\varepsilon \psi (x)
:= \sigma^1_\varepsilon(x) \cdot \unitBallShellFormula \psi(x) \cdot \sigma^2_\varepsilon(x) \\
+ (1 - \sigma^1_\varepsilon(x))\, \psi(x).
\end{multline*}
On the inner core $\|x\| \le 1 - \varepsilon$ this equals $\psi(x)$, on the
outer core $1 + \varepsilon \le \|x\| \le 2 - \varepsilon$ it equals
$\unitBallShellFormula \psi(x)$, and outside $\|x\| \ge 2 + \varepsilon$ it
vanishes. The blend region is contained in two thin annuli around the
spheres of radius $1$ and $2$.
\end{definition}

\begin{lstlisting}[style=Matlab-editor,basicstyle=\mlttfamily,escapechar=",mlshowsectionrules=true,mathescape=true,morecomment={[s]{\%\{}{\%\}}},language=lean,numbers=none]
def smoothUnitBallExtensionApprox (ε : ℝ) (ψ : E → ℝ) : E → ℝ :=
  fun x =>
    (1 - sphereOneBlend (d := d) ε x) * ψ x +
      sphereOneBlend (d := d) ε x * sphereTwoBlend (d := d) ε x *
        unitBallShellFormula (d := d) ψ x

\end{lstlisting}
\begin{lemma}
(\leanname{smoothUnitBallExtensionApprox\_contDiff},
\leanname{smoothUnitBallExtensionApprox\_hasCompactSupport})\\
For $\psi$ smooth and $0 < \varepsilon < 1$, $\smoothUnitBallExtensionApprox_\varepsilon
\psi$ is $C^\infty$ and has compact support contained in
$\overline{\Ball{2+\varepsilon}{0}}$.
\end{lemma}

\begin{lstlisting}[style=Matlab-editor,basicstyle=\mlttfamily,escapechar=",mlshowsectionrules=true,mathescape=true,morecomment={[s]{\%\{}{\%\}}},language=lean,numbers=none]
omit [NeZero d] in
lemma smoothUnitBallExtensionApprox_contDiff
    {ε : ℝ} (hε : 0 < ε) (hε1 : ε < 1)
    {ψ : E → ℝ} (hψ : ContDiff ℝ (⊤ : ℕ∞) ψ) :
    ContDiff ℝ (⊤ : ℕ∞) (smoothUnitBallExtensionApprox (d := d) ε ψ) := by
  let β1 := sphereOneBlend (d := d) ε
  let β2 := sphereTwoBlend (d := d) ε
  have hβ1 : ContDiff ℝ (⊤ : ℕ∞) β1 := contDiff_sphereOneBlend (d := d) hε hε1
  have hβ2 : ContDiff ℝ (⊤ : ℕ∞) β2 := contDiff_sphereTwoBlend (d := d) hε (by linarith)
  have hsmooth1 : ContDiff ℝ (⊤ : ℕ∞) (fun x => (1 - β1 x) * ψ x) :=
    (contDiff_const.sub hβ1).mul hψ
  have hsmooth2 : ContDiff ℝ (⊤ : ℕ∞)
      (fun x => β1 x * β2 x * unitBallShellFormula (d := d) ψ x) := by
    rw [contDiff_iff_contDiffAt]
    intro x
    by_cases hx : x ∈ Metric.ball (0 : E) (1 - ε)
    · have hβ1_zero :
          β1 =ᶠ[𝓝 x] fun _ : E => (0 : ℝ) := by
        filter_upwards [Metric.isOpen_ball.mem_nhds hx] with y hy
        simp [β1, sphereOneBlend_eq_zero_of_mem_ball (d := d) hε hy]
      have hzero :
          (fun y => β1 y * β2 y * unitBallShellFormula (d := d) ψ y)
            =ᶠ[𝓝 x] fun _ : E => (0 : ℝ) := by
        filter_upwards [hβ1_zero] with y hy
        simp [hy]
      exact contDiffAt_const.congr_of_eventuallyEq hzero
    · have hxnorm : 1 - ε ≤ ‖x‖ := by
        rw [Metric.mem_ball, dist_zero_right] at hx
        exact le_of_not_gt hx
      have hx0 : x ≠ (0 : E) := by
        intro hx0
        simp [hx0] at hxnorm
        linarith [hε1]
      exact ((hβ1.mul hβ2).contDiffAt.mul
        (contDiffAt_unitBallShellFormula (d := d) hψ hx0))
  unfold smoothUnitBallExtensionApprox
  exact hsmooth1.add hsmooth2

omit [NeZero d] in
lemma smoothUnitBallExtensionApprox_hasCompactSupport
    {ε : ℝ} (hε : 0 < ε) {ψ : E → ℝ} :
    HasCompactSupport (smoothUnitBallExtensionApprox (d := d) ε ψ) := by
  apply HasCompactSupport.intro' (isCompact_closedBall (0 : E) (2 + ε))
    isClosed_closedBall
  intro x hx
  have hnorm : 2 + ε ≤ ‖x‖ := by
    rw [Metric.mem_closedBall, dist_zero_right] at hx
    linarith
  exact smoothUnitBallExtensionApprox_eq_zero_of_norm_ge (d := d) hε hnorm

\end{lstlisting}
\begin{proof}
Each of the two summands is smooth: the shell formula factor is smooth on
$E \setminus\{0\}$ and is multiplied by $\sigma^1_\varepsilon$, which
vanishes near the origin (since $\sigma^1_\varepsilon(x) = 0$ when
$\|x\| \le 1-\varepsilon$). The compact support follows from
$\sigma^2_\varepsilon$ vanishing for $\|x\| \ge 2+\varepsilon$.
\end{proof}

\begin{definition}[\leanname{unitBallApproxEps}]
$\varepsilon_n := (n+5)^{-1}$.
\end{definition}

\begin{lstlisting}[style=Matlab-editor,basicstyle=\mlttfamily,escapechar=",mlshowsectionrules=true,mathescape=true,morecomment={[s]{\%\{}{\%\}}},language=lean,numbers=none]
def unitBallApproxEps (n : ℕ) : ℝ :=
  ((n : ℝ) + 5)⁻¹

\end{lstlisting}
\begin{lemma}[\leanname{unitBallApproxEps\_pos}, \leanname{unitBallApproxEps\_lt\_one}, \leanname{unitBallApproxEps\_le\_one\_fifth}, \leanname{tendsto\_unitBallApproxEps}, \leanname{unitBallApproxEps\_antitone}]
$\varepsilon_n > 0$, $\varepsilon_n < 1$, $\varepsilon_n \le 1/5$, $\varepsilon_n
\to 0$, and $(\varepsilon_n)$ is antitone.
\end{lemma}

\begin{lstlisting}[style=Matlab-editor,basicstyle=\mlttfamily,escapechar=",mlshowsectionrules=true,mathescape=true,morecomment={[s]{\%\{}{\%\}}},language=lean,numbers=none]
lemma unitBallApproxEps_pos (n : ℕ) :
    0 < unitBallApproxEps n := by
  unfold unitBallApproxEps
  exact inv_pos.mpr (by positivity : 0 < (n : ℝ) + 5)

lemma unitBallApproxEps_lt_one (n : ℕ) :
    unitBallApproxEps n < 1 := by
  unfold unitBallApproxEps
  have hpos : 0 < (n : ℝ) + 5 := by positivity
  field_simp [hpos.ne']
  linarith

lemma unitBallApproxEps_le_one_fifth (n : ℕ) :
    unitBallApproxEps n ≤ (1 / 5 : ℝ) := by
  unfold unitBallApproxEps
  have hpos : 0 < (n : ℝ) + 5 := by positivity
  field_simp [hpos.ne']
  nlinarith

lemma tendsto_unitBallApproxEps :
    Tendsto (fun n => unitBallApproxEps n) atTop (nhds 0) := by
  unfold unitBallApproxEps
  have hbase : Tendsto (fun n : ℕ => n + 5) atTop atTop := tendsto_add_atTop_nat 5
  have h :=
    ((tendsto_inv_atTop_nhds_zero_nat : Tendsto (fun n : ℕ => ((n : ℝ))⁻¹) atTop (nhds 0)).comp hbase)
  refine h.congr' ?_
  exact Filter.Eventually.of_forall fun n => by
    change ((((n + 5 : ℕ) : ℝ))⁻¹) = (((n : ℝ) + 5)⁻¹)
    norm_num [Nat.cast_add]

lemma unitBallApproxEps_antitone : Antitone (unitBallApproxEps) := by
  intro m n hmn
  unfold unitBallApproxEps
  have hm : 0 < (((m + 5 : ℕ) : ℝ)) := by positivity
  have hle : (((m + 5 : ℕ) : ℝ)) ≤ (((n + 5 : ℕ) : ℝ)) := by
    exact_mod_cast Nat.add_le_add_right hmn 5
  simpa [Nat.cast_add, one_div, add_comm, add_left_comm, add_assoc] using
    (one_div_le_one_div_of_le hm hle)

\end{lstlisting}
\begin{proof}
All follow from $n \mapsto (n+5)^{-1}$.
\end{proof}

\begin{definition}[\leanname{unitBallBadAnnulusOne}, \leanname{unitBallBadAnnulusTwo}]
$\mathrm{Ann}^1_\varepsilon := \{x : 1-\varepsilon < \|x\| < 1 + \varepsilon\}$
and $\mathrm{Ann}^2_\varepsilon := \{x : 2-\varepsilon < \|x\| < 2+\varepsilon\}$.
Both are measurable, and antitone in $\varepsilon$ (with respect to inclusion
on $(0,1)$).
\end{definition}

\begin{lstlisting}[style=Matlab-editor,basicstyle=\mlttfamily,escapechar=",mlshowsectionrules=true,mathescape=true,morecomment={[s]{\%\{}{\%\}}},language=lean,numbers=none]
def unitBallBadAnnulusOne (ε : ℝ) : Set E :=
  {x : E | 1 - ε < ‖x‖ ∧ ‖x‖ < 1 + ε}

def unitBallBadAnnulusTwo (ε : ℝ) : Set E :=
  {x : E | 2 - ε < ‖x‖ ∧ ‖x‖ < 2 + ε}

\end{lstlisting}
\begin{lemma}
(\leanname{iInter\_unitBallBadAnnulusOne\_eq\_sphere},
\leanname{iInter\_unitBallBadAnnulusTwo\_eq\_sphere},
\leanname{tendsto\_measure\_unitBallBadAnnulusOne},
\leanname{tendsto\_measure\_unitBallBadAnnulusTwo})\\
$\bigcap_{n} \mathrm{Ann}^j_{\varepsilon_n} = \mathrm{S}_j$ for $j = 1, 2$, and
$|\mathrm{Ann}^j_{\varepsilon_n}| \to 0$ as $n \to \infty$.
\end{lemma}

\begin{lstlisting}[style=Matlab-editor,basicstyle=\mlttfamily,escapechar=",mlshowsectionrules=true,mathescape=true,morecomment={[s]{\%\{}{\%\}}},language=lean,numbers=none]
omit [NeZero d] in
lemma iInter_unitBallBadAnnulusOne_eq_sphere :
    (⋂ n : ℕ, unitBallBadAnnulusOne (d := d) (unitBallApproxEps n)) =
      Metric.sphere (0 : E) 1 := by
  ext x
  constructor
  · intro hx
    rw [Metric.mem_sphere, dist_zero_right]
    by_cases hEq : ‖x‖ = 1
    · exact hEq
    · have hgap : 0 < |‖x‖ - 1| := abs_pos.mpr (sub_ne_zero.mpr hEq)
      have hev : ∀ᶠ n in atTop, unitBallApproxEps n < |‖x‖ - 1| :=
        tendsto_unitBallApproxEps (Iio_mem_nhds hgap)
      rcases (eventually_atTop.1 hev) with ⟨N, hN⟩
      have hxN := mem_iInter.1 hx N
      have hlt : unitBallApproxEps N < |‖x‖ - 1| := hN N le_rfl
      rcases hxN with ⟨hx1, hx2⟩
      have habs : |‖x‖ - 1| < unitBallApproxEps N := by
        rw [abs_sub_lt_iff]
        constructor <;> linarith
      linarith
  · intro hx
    rw [Metric.mem_sphere, dist_zero_right] at hx
    refine mem_iInter.2 ?_
    intro n
    constructor <;> simp [hx, unitBallApproxEps_pos n]

omit [NeZero d] in
lemma iInter_unitBallBadAnnulusTwo_eq_sphere :
    (⋂ n : ℕ, unitBallBadAnnulusTwo (d := d) (unitBallApproxEps n)) =
      Metric.sphere (0 : E) 2 := by
  ext x
  constructor
  · intro hx
    rw [Metric.mem_sphere, dist_zero_right]
    by_cases hEq : ‖x‖ = 2
    · exact hEq
    · have hgap : 0 < |‖x‖ - 2| := abs_pos.mpr (sub_ne_zero.mpr hEq)
      have hev : ∀ᶠ n in atTop, unitBallApproxEps n < |‖x‖ - 2| :=
        tendsto_unitBallApproxEps (Iio_mem_nhds hgap)
      rcases (eventually_atTop.1 hev) with ⟨N, hN⟩
      have hxN := mem_iInter.1 hx N
      have hlt : unitBallApproxEps N < |‖x‖ - 2| := hN N le_rfl
      rcases hxN with ⟨hx1, hx2⟩
      have habs : |‖x‖ - 2| < unitBallApproxEps N := by
        rw [abs_sub_lt_iff]
        constructor <;> linarith
      linarith
  · intro hx
    rw [Metric.mem_sphere, dist_zero_right] at hx
    refine mem_iInter.2 ?_
    intro n
    constructor <;> simp [hx, unitBallApproxEps_pos n]

lemma tendsto_measure_unitBallBadAnnulusOne :
    Tendsto (fun n => volume (unitBallBadAnnulusOne (d := d) (unitBallApproxEps n)))
      atTop (nhds 0) := by
  have h :=
    tendsto_measure_iInter_atTop
      (μ := volume)
      (s := fun n : ℕ => unitBallBadAnnulusOne (d := d) (unitBallApproxEps n))
      (fun n => (measurableSet_unitBallBadAnnulusOne (d := d) _).nullMeasurableSet)
      unitBallBadAnnulusOne_antitone
      ⟨0, (measure_unitBallBadAnnulusOne_lt_top (d := d) 0).ne⟩
  have hsphere_zero : volume (Metric.sphere (0 : E) 1) = 0 := by
    simpa using (Measure.addHaar_sphere (μ := (volume : Measure E)) (x := (0 : E)) 1)
  rw [iInter_unitBallBadAnnulusOne_eq_sphere (d := d), hsphere_zero] at h
  simpa using h

lemma tendsto_measure_unitBallBadAnnulusTwo :
    Tendsto (fun n => volume (unitBallBadAnnulusTwo (d := d) (unitBallApproxEps n)))
      atTop (nhds 0) := by
  have h :=
    tendsto_measure_iInter_atTop
      (μ := volume)
      (s := fun n : ℕ => unitBallBadAnnulusTwo (d := d) (unitBallApproxEps n))
      (fun n => (measurableSet_unitBallBadAnnulusTwo (d := d) _).nullMeasurableSet)
      unitBallBadAnnulusTwo_antitone
      ⟨0, (measure_unitBallBadAnnulusTwo_lt_top (d := d) 0).ne⟩
  have hsphere_zero : volume (Metric.sphere (0 : E) 2) = 0 := by
    simpa using (Measure.addHaar_sphere (μ := (volume : Measure E)) (x := (0 : E)) 2)
  rw [iInter_unitBallBadAnnulusTwo_eq_sphere (d := d), hsphere_zero] at h
  simpa using h

\end{lstlisting}
\begin{proof}
A point lies in every $\mathrm{Ann}^j_{\varepsilon_n}$ iff $\|x\| = j$, since
the annulus shrinks to the sphere as $\varepsilon \to 0$. The measure of the
sphere is zero (Mathlib: \leanname{Measure.addHaar\_sphere}), so by continuity
of measure from above the annulus measures vanish.
\end{proof}

\begin{definition}[\leanname{sphereOneControl}, \leanname{sphereTwoControl}]
Compact buffer sets around the unit and double spheres,
$\mathrm{Buf}^1 := \{x : 4/5 \le \|x\| \le 6/5\}$ and
$\mathrm{Buf}^2 := \{x : 9/5 \le \|x\| \le 11/5\}$. Both are compact.
\end{definition}

\begin{lstlisting}[style=Matlab-editor,basicstyle=\mlttfamily,escapechar=",mlshowsectionrules=true,mathescape=true,morecomment={[s]{\%\{}{\%\}}},language=lean,numbers=none]
def sphereOneControl : Set E :=
  {x : E | (3 / 4 : ℝ) ≤ ‖x‖ ∧ ‖x‖ ≤ 5 / 4}

def sphereTwoControl : Set E :=
  {x : E | (7 / 4 : ℝ) ≤ ‖x‖ ∧ ‖x‖ ≤ 9 / 4}

\end{lstlisting}
\begin{lemma}[\leanname{badAnnulusOne\_subset\_sphereOneControl}, \leanname{badAnnulusTwo\_subset\_sphereTwoControl}, \leanname{isCompact\_sphereOneControl}, \leanname{isCompact\_sphereTwoControl}]
For all $n$, $\mathrm{Ann}^j_{\varepsilon_n} \subseteq \mathrm{Buf}^j$, and
both buffers are compact.
\end{lemma}

\begin{lstlisting}[style=Matlab-editor,basicstyle=\mlttfamily,escapechar=",mlshowsectionrules=true,mathescape=true,morecomment={[s]{\%\{}{\%\}}},language=lean,numbers=none]
omit [NeZero d] in
lemma badAnnulusOne_subset_sphereOneControl (n : ℕ) :
    unitBallBadAnnulusOne (d := d) (unitBallApproxEps n) ⊆ sphereOneControl (d := d) := by
  intro x hx
  have hε : unitBallApproxEps n ≤ (1 / 5 : ℝ) := unitBallApproxEps_le_one_fifth n
  constructor
  · linarith [hx.1, hε]
  · linarith [hx.2, hε]

omit [NeZero d] in
lemma badAnnulusTwo_subset_sphereTwoControl (n : ℕ) :
    unitBallBadAnnulusTwo (d := d) (unitBallApproxEps n) ⊆ sphereTwoControl (d := d) := by
  intro x hx
  have hε : unitBallApproxEps n ≤ (1 / 5 : ℝ) := unitBallApproxEps_le_one_fifth n
  constructor
  · linarith [hx.1, hε]
  · linarith [hx.2, hε]

omit [NeZero d] in
lemma isCompact_sphereOneControl : IsCompact (sphereOneControl (d := d)) := by
  have hclosed : IsClosed (sphereOneControl (d := d)) := by
    unfold sphereOneControl
    exact (isClosed_le continuous_const continuous_norm).inter
      (isClosed_le continuous_norm continuous_const)
  refine (isCompact_closedBall (0 : E) (5 / 4 : ℝ)).of_isClosed_subset hclosed ?_
  intro x hx
  rw [Metric.mem_closedBall, dist_zero_right]
  exact hx.2

omit [NeZero d] in
lemma isCompact_sphereTwoControl : IsCompact (sphereTwoControl (d := d)) := by
  have hclosed : IsClosed (sphereTwoControl (d := d)) := by
    unfold sphereTwoControl
    exact (isClosed_le continuous_const continuous_norm).inter
      (isClosed_le continuous_norm continuous_const)
  refine (isCompact_closedBall (0 : E) (9 / 4 : ℝ)).of_isClosed_subset hclosed ?_
  intro x hx
  rw [Metric.mem_closedBall, dist_zero_right]
  exact hx.2

\end{lstlisting}
\begin{proof}
$\varepsilon_n \le 1/5$, so $j - \varepsilon_n \ge j - 1/5$ and similarly
above. Compactness is by closedness and boundedness.
\end{proof}

\begin{theorem}[\leanname{exists\_fderiv\_bound\_on\_compact}, \leanname{exists\_norm\_bound\_on\_compact}]
For any continuous $f$ (resp.\ $C^1$ $f$) on $E$, $f$ (resp.\ $\|Df\|$) is
bounded on every compact set.
\end{theorem}

\begin{lstlisting}[style=Matlab-editor,basicstyle=\mlttfamily,escapechar=",mlshowsectionrules=true,mathescape=true,morecomment={[s]{\%\{}{\%\}}},language=lean,numbers=none]
omit [NeZero d] in
theorem exists_fderiv_bound_on_compact
    {f : E → ℝ} {K U : Set E}
    (hK : IsCompact K) (hKU : K ⊆ U)
    (hcd : ∀ x ∈ U, ContDiffAt ℝ 1 f x) :
    ∃ C : ℝ, 0 ≤ C ∧ ∀ x ∈ K, ‖fderiv ℝ f x‖ ≤ C := by
  have hcontU : ContinuousOn (fun x => fderiv ℝ f x) U := by
    intro x hx
    exact ((hcd x hx).continuousAt_fderiv (by simp : (1 : WithTop ℕ∞) ≠ 0)).continuousWithinAt
  rcases IsCompact.exists_bound_of_continuousOn hK (hcontU.mono hKU) with ⟨C, hC⟩
  refine ⟨max C 0, le_max_right _ _, ?_⟩
  intro x hx
  exact (hC x hx).trans (le_max_left _ _)

omit [NeZero d] in
theorem exists_norm_bound_on_compact
    {f : E → ℝ} {K U : Set E}
    (hK : IsCompact K) (hKU : K ⊆ U)
    (hcont : ContinuousOn f U) :
    ∃ C : ℝ, ∀ x ∈ K, ‖f x‖ ≤ C := by
  rcases IsCompact.exists_bound_of_continuousOn hK (hcont.mono hKU) with ⟨C, hC⟩
  refine ⟨max C 0, ?_⟩
  intro x hx
  exact (hC x hx).trans (le_max_left _ _)

\end{lstlisting}
\begin{proof}
Continuous functions attain bounded values on compact sets.
\end{proof}

\begin{theorem}[\leanname{exists\_shellFormula\_fderiv\_bound}]
For any smooth $\psi$, the derivative $D \unitBallShellFormula \psi$ is bounded
on the compact set $\mathrm{Buf}^2$, hence on each $\mathrm{Ann}^2_{\varepsilon_n}$.
\end{theorem}

\begin{lstlisting}[style=Matlab-editor,basicstyle=\mlttfamily,escapechar=",mlshowsectionrules=true,mathescape=true,morecomment={[s]{\%\{}{\%\}}},language=lean,numbers=none]
omit [NeZero d] in
theorem exists_shellFormula_fderiv_bound
    {ψ : E → ℝ} (hψ : ContDiff ℝ (⊤ : ℕ∞) ψ) :
    ∃ C : ℝ, 0 ≤ C ∧
      ∀ x ∈ sphereTwoControl (d := d),
        ‖fderiv ℝ (unitBallShellFormula (d := d) ψ) x‖ ≤ C := by
  have hsubset : sphereTwoControl (d := d) ⊆ {x : E | x ≠ 0} := by
    intro x hx
    have : 0 < ‖x‖ := by linarith [hx.1]
    intro hx0
    simpa [hx0] using this.ne'
  rcases exists_fderiv_bound_on_compact (d := d)
      (f := unitBallShellFormula (d := d) ψ) (K := sphereTwoControl (d := d))
      (U := {x : E | x ≠ 0}) isCompact_sphereTwoControl hsubset
      (fun x hx => (contDiffAt_unitBallShellFormula (d := d) hψ hx).of_le (by norm_num))
      with ⟨C, hC_nonneg, hC⟩
  exact ⟨C, hC_nonneg, hC⟩

\end{lstlisting}
\begin{proof}
Smoothness of $\unitBallShellFormula \psi$ on $\Sigma^{\mathrm{out}}$ extends
to smoothness on a neighborhood of $\mathrm{Buf}^2$, where $\|x\| \in
[9/5, 11/5] \subseteq (0, 2 + \delta)$, so the result follows from
\leanname{exists\_norm\_bound\_on\_compact}.
\end{proof}

\begin{lemma}[\leanname{tendsto\_smoothUnitBallExtensionApprox\_pointwise}, \leanname{tendsto\_fderiv\_smoothUnitBallExtensionApprox\_pointwise}]
\label{lem:smoothapprox-pointwise}
For $\psi \in C^\infty(E)$ and $x$ outside the unit and double spheres,
$\smoothUnitBallExtensionApprox_{\varepsilon_n} \psi (x) \to \Eext \psi (x)$
and $D\smoothUnitBallExtensionApprox_{\varepsilon_n} \psi (x) \to D\Eext\psi(x)$.
\end{lemma}

\begin{lstlisting}[style=Matlab-editor,basicstyle=\mlttfamily,escapechar=",mlshowsectionrules=true,mathescape=true,morecomment={[s]{\%\{}{\%\}}},language=lean,numbers=none]
omit [NeZero d] in
lemma tendsto_smoothUnitBallExtensionApprox_pointwise
    {ψ : E → ℝ} {x : E}
    (hx1 : ‖x‖ ≠ 1) (hx2 : ‖x‖ ≠ 2) :
    Tendsto
      (fun n => smoothUnitBallExtensionApprox (d := d) (unitBallApproxEps n) ψ x)
      atTop (nhds (unitBallExtension (d := d) ψ x)) := by
  have hε := tendsto_unitBallApproxEps
  by_cases hlt1 : ‖x‖ < 1
  · have hgap : 0 < 1 - ‖x‖ := by linarith
    have hev :
        ∀ᶠ n in atTop, unitBallApproxEps n < 1 - ‖x‖ := hε (Iio_mem_nhds hgap)
    refine tendsto_const_nhds.congr' ?_
    filter_upwards [hev] with n hn
    have hxball : x ∈ Metric.ball (0 : E) (1 - unitBallApproxEps n) := by
      rw [Metric.mem_ball, dist_zero_right]
      linarith
    exact (smoothUnitBallExtensionApprox_eq_of_mem_innerCore (d := d)
      (unitBallApproxEps_pos n) hxball).symm
  · have hge1 : 1 ≤ ‖x‖ := by linarith
    by_cases hlt2 : ‖x‖ < 2
    · have hgt1 : 1 < ‖x‖ := by
        have hx1' : (1 : ℝ) ≠ ‖x‖ := by simpa [eq_comm] using hx1
        exact lt_of_le_of_ne hge1 hx1'
      let δ : ℝ := min (‖x‖ - 1) (2 - ‖x‖)
      have hδ : 0 < δ := by
        dsimp [δ]
        exact lt_min (by linarith) (by linarith)
      have hev : ∀ᶠ n in atTop, unitBallApproxEps n < δ := hε (Iio_mem_nhds hδ)
      refine tendsto_const_nhds.congr' ?_
      filter_upwards [hev] with n hn
      have hxl : 1 + unitBallApproxEps n < ‖x‖ := by
        dsimp [δ] at hn
        have : unitBallApproxEps n < ‖x‖ - 1 := lt_of_lt_of_le hn (min_le_left _ _)
        linarith
      have hxr : ‖x‖ < 2 - unitBallApproxEps n := by
        dsimp [δ] at hn
        have : unitBallApproxEps n < 2 - ‖x‖ := lt_of_lt_of_le hn (min_le_right _ _)
        linarith
      exact (smoothUnitBallExtensionApprox_eq_of_mem_shellCore (d := d)
        (unitBallApproxEps_pos n) hxl hxr).symm
    · have hge2 : 2 ≤ ‖x‖ := by linarith
      have hgt2 : 2 < ‖x‖ := by
        have hx2' : (2 : ℝ) ≠ ‖x‖ := by simpa [eq_comm] using hx2
        exact lt_of_le_of_ne hge2 hx2'
      have hgap : 0 < ‖x‖ - 2 := by linarith
      have hev : ∀ᶠ n in atTop, unitBallApproxEps n < ‖x‖ - 2 := hε (Iio_mem_nhds hgap)
      refine tendsto_const_nhds.congr' ?_
      filter_upwards [hev] with n hn
      have hxge : 2 + unitBallApproxEps n ≤ ‖x‖ := by linarith
      exact (smoothUnitBallExtensionApprox_eq_of_norm_ge (d := d)
        (unitBallApproxEps_pos n) hxge).symm

omit [NeZero d] in
lemma tendsto_fderiv_smoothUnitBallExtensionApprox_pointwise
    {ψ : E → ℝ} {x : E}
    (hx1 : ‖x‖ ≠ 1) (hx2 : ‖x‖ ≠ 2) :
    Tendsto
      (fun n => fderiv ℝ (smoothUnitBallExtensionApprox (d := d) (unitBallApproxEps n) ψ) x)
      atTop (nhds (fderiv ℝ (unitBallExtension (d := d) ψ) x)) := by
  have hε := tendsto_unitBallApproxEps
  by_cases hlt1 : ‖x‖ < 1
  · have hgap : 0 < 1 - ‖x‖ := by linarith
    have hev :
        ∀ᶠ n in atTop, unitBallApproxEps n < 1 - ‖x‖ := hε (Iio_mem_nhds hgap)
    refine tendsto_const_nhds.congr' ?_
    filter_upwards [hev] with n hn
    have hxball : x ∈ Metric.ball (0 : E) (1 - unitBallApproxEps n) := by
      rw [Metric.mem_ball, dist_zero_right]
      linarith
    have hEq :
        smoothUnitBallExtensionApprox (d := d) (unitBallApproxEps n) ψ =ᶠ[𝓝 x]
          unitBallExtension (d := d) ψ := by
      filter_upwards [Metric.isOpen_ball.mem_nhds hxball] with y hy
      exact smoothUnitBallExtensionApprox_eq_of_mem_innerCore (d := d)
        (unitBallApproxEps_pos n) hy
    exact (Filter.EventuallyEq.fderiv_eq (𝕜 := ℝ) hEq).symm
  · have hge1 : 1 ≤ ‖x‖ := by linarith
    by_cases hlt2 : ‖x‖ < 2
    · have hgt1 : 1 < ‖x‖ := lt_of_le_of_ne hge1 hx1.symm
      let δ : ℝ := min (‖x‖ - 1) (2 - ‖x‖)
      have hδ : 0 < δ := by
        dsimp [δ]
        exact lt_min (by linarith) (by linarith)
      have hev : ∀ᶠ n in atTop, unitBallApproxEps n < δ := hε (Iio_mem_nhds hδ)
      refine tendsto_const_nhds.congr' ?_
      filter_upwards [hev] with n hn
      have hxl : 1 + unitBallApproxEps n < ‖x‖ := by
        dsimp [δ] at hn
        have : unitBallApproxEps n < ‖x‖ - 1 := lt_of_lt_of_le hn (min_le_left _ _)
        linarith
      have hxr : ‖x‖ < 2 - unitBallApproxEps n := by
        dsimp [δ] at hn
        have : unitBallApproxEps n < 2 - ‖x‖ := lt_of_lt_of_le hn (min_le_right _ _)
        linarith
      have hxOuter : x ∈ unitBallOuterShell (d := d) := by constructor <;> linarith
      have hEq :
          smoothUnitBallExtensionApprox (d := d) (unitBallApproxEps n) ψ =ᶠ[𝓝 x]
            unitBallExtension (d := d) ψ := by
        have hLeft :
            {y : E | 1 + unitBallApproxEps n < ‖y‖} ∈ 𝓝 x :=
          (isOpen_lt continuous_const continuous_norm).mem_nhds hxl
        have hRight :
            {y : E | ‖y‖ < 2 - unitBallApproxEps n} ∈ 𝓝 x :=
          (isOpen_lt continuous_norm continuous_const).mem_nhds hxr
        filter_upwards [hLeft, hRight] with y hyL hyR
        exact smoothUnitBallExtensionApprox_eq_of_mem_shellCore (d := d)
          (unitBallApproxEps_pos n) hyL hyR
      have hShell :
          fderiv ℝ (unitBallExtension (d := d) ψ) x =
            fderiv ℝ (unitBallShellFormula (d := d) ψ) x :=
        fderiv_unitBallExtension_eq_shellFormula_of_mem_outerShell (d := d) hxOuter
      rw [hShell]
      have hEqShell :
          smoothUnitBallExtensionApprox (d := d) (unitBallApproxEps n) ψ =ᶠ[𝓝 x]
            unitBallShellFormula (d := d) ψ := by
        filter_upwards [hEq, isOpen_unitBallOuterShell (d := d) |>.mem_nhds hxOuter] with y hyEq hyOuter
        simpa [unitBallExtension_eq_shellFormula_of_mem_outerShell (d := d) hyOuter] using hyEq
      exact (Filter.EventuallyEq.fderiv_eq (𝕜 := ℝ) hEqShell).symm
    · have hge2 : 2 ≤ ‖x‖ := by linarith
      have hgt2 : 2 < ‖x‖ := lt_of_le_of_ne hge2 hx2.symm
      have hgap : 0 < ‖x‖ - 2 := by linarith
      have hev : ∀ᶠ n in atTop, unitBallApproxEps n < ‖x‖ - 2 := hε (Iio_mem_nhds hgap)
      refine tendsto_const_nhds.congr' ?_
      filter_upwards [hev] with n hn
      have hxgt : 2 + unitBallApproxEps n < ‖x‖ := by linarith
      have hZeroEq :
          smoothUnitBallExtensionApprox (d := d) (unitBallApproxEps n) ψ =ᶠ[𝓝 x]
            unitBallExtension (d := d) ψ := by
        have hNhds :
            {y : E | 2 + unitBallApproxEps n < ‖y‖} ∈ 𝓝 x :=
          (isOpen_lt continuous_const continuous_norm).mem_nhds hxgt
        filter_upwards [hNhds] with y hy
        have hy1 : smoothUnitBallExtensionApprox (d := d) (unitBallApproxEps n) ψ y = 0 := by
          exact smoothUnitBallExtensionApprox_eq_zero_of_norm_ge (d := d)
            (unitBallApproxEps_pos n) (le_of_lt hy)
        have hy2 : unitBallExtension (d := d) ψ y = 0 := by
          have hy2' : 2 ≤ ‖y‖ := by
            have hεpos : 0 < unitBallApproxEps n := unitBallApproxEps_pos n
            linarith
          exact unitBallExtension_eq_zero_of_two_le_norm (d := d) hy2'
        simp [hy1, hy2]
      exact (Filter.EventuallyEq.fderiv_eq (𝕜 := ℝ) hZeroEq).symm

\end{lstlisting}
\begin{proof}
Eventually $x$ lies in either $\Bz{1}$ (so the inner branch picks up $\psi(x)
= \Eext \psi(x)$), or in $\Sigma^{\mathrm{out}}$ separated from both spheres
(so the outer branch agrees with $\unitBallShellFormula \psi(x) = \Eext\psi(x)$
and similarly for derivatives), or in $\|x\| > 2$ where the cutoff vanishes.
The blends become trivial in the limit and the derivative converges by
smoothness.
\end{proof}

\begin{definition}[\leanname{unitBallBadSet}]
$\mathrm{Bad}_\varepsilon := \mathrm{Ann}^1_\varepsilon \cup
\mathrm{Ann}^2_\varepsilon$.
\end{definition}

\begin{lstlisting}[style=Matlab-editor,basicstyle=\mlttfamily,escapechar=",mlshowsectionrules=true,mathescape=true,morecomment={[s]{\%\{}{\%\}}},language=lean,numbers=none]
def unitBallBadSet (ε : ℝ) : Set E :=
  unitBallBadAnnulusOne (d := d) ε ∪ unitBallBadAnnulusTwo (d := d) ε

\end{lstlisting}
\begin{lemma}[\leanname{tendsto\_measure\_unitBallBadSet}]
$|\mathrm{Bad}_{\varepsilon_n}| \to 0$.
\end{lemma}

\begin{lstlisting}[style=Matlab-editor,basicstyle=\mlttfamily,escapechar=",mlshowsectionrules=true,mathescape=true,morecomment={[s]{\%\{}{\%\}}},language=lean,numbers=none]
lemma tendsto_measure_unitBallBadSet :
    Tendsto (fun n => volume (unitBallBadSet (d := d) (unitBallApproxEps n)))
      atTop (nhds 0) := by
  have hle :
      (fun n => volume (unitBallBadSet (d := d) (unitBallApproxEps n))) ≤
        fun n =>
          volume (unitBallBadAnnulusOne (d := d) (unitBallApproxEps n)) +
            volume (unitBallBadAnnulusTwo (d := d) (unitBallApproxEps n)) := by
    intro n
    exact measure_union_le _ _
  refine tendsto_of_tendsto_of_tendsto_of_le_of_le
    (g := fun _ : ℕ => (0 : ℝ≥0∞))
    (h := fun n =>
      volume (unitBallBadAnnulusOne (d := d) (unitBallApproxEps n)) +
        volume (unitBallBadAnnulusTwo (d := d) (unitBallApproxEps n)))
    tendsto_const_nhds
    ?_
    (fun _ => bot_le)
    hle
  simpa [unitBallBadSet] using
    (tendsto_measure_unitBallBadAnnulusOne (d := d)).add
      (tendsto_measure_unitBallBadAnnulusTwo (d := d))

\end{lstlisting}
\begin{proof}
Sub-additivity together with the previous annulus convergence.
\end{proof}

\begin{lemma}[\leanname{smoothUnitBallExtensionApprox\_eq\_unitBallExtension\_of\_not\_mem\_badSet}]
For $\psi$ smooth and $x \notin \mathrm{Bad}_{\varepsilon_n}$,
$\smoothUnitBallExtensionApprox_{\varepsilon_n} \psi (x) = \Eext \psi(x)$.
\end{lemma}

\begin{lstlisting}[style=Matlab-editor,basicstyle=\mlttfamily,escapechar=",mlshowsectionrules=true,mathescape=true,morecomment={[s]{\%\{}{\%\}}},language=lean,numbers=none]
omit [NeZero d] in
lemma smoothUnitBallExtensionApprox_eq_unitBallExtension_of_not_mem_badSet
    {ψ : E → ℝ} {n : ℕ} {x : E}
    (hx : x ∉ unitBallBadSet (d := d) (unitBallApproxEps n)) :
    smoothUnitBallExtensionApprox (d := d) (unitBallApproxEps n) ψ x =
      unitBallExtension (d := d) ψ x := by
  let ε := unitBallApproxEps n
  have hε : 0 < unitBallApproxEps n := unitBallApproxEps_pos n
  have hx1 : x ∉ unitBallBadAnnulusOne (d := d) (unitBallApproxEps n) := by
    intro hx1
    exact hx (Or.inl hx1)
  have hx2 : x ∉ unitBallBadAnnulusTwo (d := d) (unitBallApproxEps n) := by
    intro hx2
    exact hx (Or.inr hx2)
  by_cases hinner : ‖x‖ ≤ 1 - ε
  · exact smoothUnitBallExtensionApprox_eq_of_norm_le_innerCore (d := d) hε hinner
  · by_cases hshell : 1 + ε ≤ ‖x‖ ∧ ‖x‖ ≤ 2 - ε
    · exact smoothUnitBallExtensionApprox_eq_of_shellCore_closed (d := d) hε hshell.1 hshell.2
    · have hnotInner : 1 - ε < ‖x‖ := lt_of_not_ge hinner
      have hnorm_ge : 2 + ε ≤ ‖x‖ := by
        by_contra hfar
        have hfar' : ‖x‖ < 2 + ε := by linarith
        by_cases hlt : ‖x‖ < 1 + ε
        · exact hx1 ⟨hnotInner, hlt⟩
        · have hge : 1 + ε ≤ ‖x‖ := le_of_not_gt hlt
          by_cases hmid : ‖x‖ ≤ 2 - ε
          · exact hshell ⟨hge, hmid⟩
          · exact hx2 ⟨lt_of_not_ge hmid, hfar'⟩
      exact smoothUnitBallExtensionApprox_eq_of_norm_ge (d := d) hε hnorm_ge

\end{lstlisting}
\begin{proof}
Outside the bad set, $x$ falls cleanly into one of: inner core, outer core,
or the exterior $\|x\|>2+\varepsilon$, on each of which the blends become
constant ($0$ or $1$) and the formula collapses to $\Eext\psi$.
\end{proof}

\begin{theorem}[\leanname{exists\_shellFormula\_error\_bound}]
\label{thm:err-bound}
For any smooth $\psi$, there is a constant $K$ depending only on $\psi$
such that
\[
\bigl|\smoothUnitBallExtensionApprox_{\varepsilon_n} \psi(x) - \Eext\psi(x)\bigr|
\le K\,\varepsilon_n
\]
for all $x \in E$.
\end{theorem}

\begin{lstlisting}[style=Matlab-editor,basicstyle=\mlttfamily,escapechar=",mlshowsectionrules=true,mathescape=true,morecomment={[s]{\%\{}{\%\}}},language=lean,numbers=none]
omit [NeZero d] in
theorem exists_shellFormula_error_bound
    {ψ : E → ℝ} (hψ : ContDiff ℝ (⊤ : ℕ∞) ψ) :
    ∃ C : ℝ, 0 ≤ C ∧
      ∀ n x, x ∈ unitBallBadAnnulusTwo (d := d) (unitBallApproxEps n) →
        ‖unitBallShellFormula (d := d) ψ x‖ ≤ C * unitBallApproxEps n := by
  rcases exists_shellFormula_fderiv_bound (d := d) hψ with ⟨C, hC_nonneg, hC⟩
  exact ⟨C, hC_nonneg, fun n x hx => shellFormula_error_bound_at (d := d) hψ hC_nonneg hC hx⟩

\end{lstlisting}
\begin{proof}
Outside $\mathrm{Bad}_{\varepsilon_n}$, both sides agree (so the error is
zero). On $\mathrm{Bad}_{\varepsilon_n}$, both sides differ from each other
only in the blend region, which has width $\sim \varepsilon_n$, and within a
ball of radius $\varepsilon_n$ around any point of the bad annulus the
shell formula is Lipschitz with constant uniformly controlled by the
$C^1$-bound of $\psi$ on the buffer compact $\mathrm{Buf}^2$. Taking the
mean-value bound on segments inside $\mathrm{Buf}^j$ yields the linear
control.
\end{proof}

\begin{theorem}[\leanname{exists\_fun\_error\_bound\_badSet}]
For each smooth $\psi$, there is $K$ such that
\[
\eLpNormText{\smoothUnitBallExtensionApprox_{\varepsilon_n}\psi - \Eext\psi}_{L^p(E)}
\le K\,\varepsilon_n\,|\mathrm{Bad}_{\varepsilon_n}|^{1/p}
\to 0.
\]
\end{theorem}

\begin{lstlisting}[style=Matlab-editor,basicstyle=\mlttfamily,escapechar=",mlshowsectionrules=true,mathescape=true,morecomment={[s]{\%\{}{\%\}}},language=lean,numbers=none]
theorem exists_fun_error_bound_badSet
    {ψ : E → ℝ} (hψ : ContDiff ℝ (⊤ : ℕ∞) ψ) :
    ∃ C : ℝ, 0 ≤ C ∧
      ∀ n x,
        ‖smoothUnitBallExtensionApprox (d := d) (unitBallApproxEps n) ψ x -
            exactUnitBallExtensionModel (d := d) ψ x‖
          ≤ (unitBallBadSet (d := d) (unitBallApproxEps n)).indicator (fun _ => C) x := by
  rcases exists_shellSubPsi_error_bound (d := d) hψ with ⟨C1, hC1_nonneg, hC1⟩
  rcases exists_shellFormula_error_bound (d := d) hψ with ⟨C2, hC2_nonneg, hC2⟩
  refine ⟨max (C1 / 5) (C2 / 5), by positivity, ?_⟩
  intro n x
  by_cases hbad : x ∈ unitBallBadSet (d := d) (unitBallApproxEps n)
  · rw [Set.indicator_of_mem hbad]
    rcases hbad with hx1 | hx2
    · have hεsmall : unitBallApproxEps n ≤ (1 / 5 : ℝ) := unitBallApproxEps_le_one_fifth n
      have hxball2 : x ∈ Metric.ball (0 : E) (2 - unitBallApproxEps n) := by
        rw [Metric.mem_ball, dist_zero_right]; linarith [hx1.2, hεsmall]
      have hβ2 : sphereTwoBlend (d := d) (unitBallApproxEps n) x = 1 :=
        sphereTwoBlend_eq_one_of_mem_ball (d := d) (unitBallApproxEps_pos n) hxball2
      by_cases hxle1 : ‖x‖ ≤ 1
      · have hxclosed : x ∈ Metric.closedBall (0 : E) 1 := by
          rw [Metric.mem_closedBall, dist_zero_right]; exact hxle1
        have hEq :
            smoothUnitBallExtensionApprox (d := d) (unitBallApproxEps n) ψ x -
              exactUnitBallExtensionModel (d := d) ψ x =
              sphereOneBlend (d := d) (unitBallApproxEps n) x *
                (unitBallShellFormula (d := d) ψ x - ψ x) := by
          unfold smoothUnitBallExtensionApprox exactUnitBallExtensionModel exactUnitBallShellOrZero
          simp [hβ2, hxclosed, sub_eq_add_neg]; ring
        calc ‖smoothUnitBallExtensionApprox (d := d) (unitBallApproxEps n) ψ x -
              exactUnitBallExtensionModel (d := d) ψ x‖
              = ‖sphereOneBlend (d := d) (unitBallApproxEps n) x *
                  (unitBallShellFormula (d := d) ψ x - ψ x)‖ := by rw [hEq]
          _ ≤ ‖unitBallShellFormula (d := d) ψ x - ψ x‖ := by
                rw [norm_mul, Real.norm_eq_abs]
                nlinarith [abs_of_nonneg (sphereOneBlend_nonneg (d := d) (ε := unitBallApproxEps n) (x := x)),
                  sphereOneBlend_le_one (d := d) (ε := unitBallApproxEps n) (x := x),
                  norm_nonneg (unitBallShellFormula (d := d) ψ x - ψ x)]
          _ ≤ C1 * unitBallApproxEps n := hC1 n x hx1
          _ ≤ max (C1 / 5) (C2 / 5) := by
                have := unitBallApproxEps_le_one_fifth n
                exact le_trans (by nlinarith [hC1_nonneg]) (le_max_left _ _)
      · have hxOuter : x ∈ unitBallOuterShell (d := d) := ⟨by linarith, by linarith [hx1.2, hεsmall]⟩
        have hEq :
            smoothUnitBallExtensionApprox (d := d) (unitBallApproxEps n) ψ x -
              exactUnitBallExtensionModel (d := d) ψ x =
              (1 - sphereOneBlend (d := d) (unitBallApproxEps n) x) *
                (ψ x - unitBallShellFormula (d := d) ψ x) := by
          unfold smoothUnitBallExtensionApprox exactUnitBallExtensionModel exactUnitBallShellOrZero
          simp [hβ2, hxOuter, hxle1, sub_eq_add_neg]; ring
        calc ‖smoothUnitBallExtensionApprox (d := d) (unitBallApproxEps n) ψ x -
              exactUnitBallExtensionModel (d := d) ψ x‖
              = ‖(1 - sphereOneBlend (d := d) (unitBallApproxEps n) x) *
                  (ψ x - unitBallShellFormula (d := d) ψ x)‖ := by rw [hEq]
          _ ≤ ‖ψ x - unitBallShellFormula (d := d) ψ x‖ := by
                rw [norm_mul, Real.norm_eq_abs]
                have h1 : |1 - sphereOneBlend (d := d) (unitBallApproxEps n) x| ≤ 1 := by
                  rw [abs_of_nonneg (by linarith [sphereOneBlend_le_one (d := d) (ε := unitBallApproxEps n) (x := x)])]
                  linarith [sphereOneBlend_nonneg (d := d) (ε := unitBallApproxEps n) (x := x)]
                nlinarith [norm_nonneg (ψ x - unitBallShellFormula (d := d) ψ x)]
          _ = ‖unitBallShellFormula (d := d) ψ x - ψ x‖ := norm_sub_rev _ _
          _ ≤ C1 * unitBallApproxEps n := hC1 n x hx1
          _ ≤ max (C1 / 5) (C2 / 5) := by
                have := unitBallApproxEps_le_one_fifth n
                exact le_trans (by nlinarith [hC1_nonneg]) (le_max_left _ _)
    · have hεsmall : unitBallApproxEps n ≤ (1 / 5 : ℝ) := unitBallApproxEps_le_one_fifth n
      have hβ1 : sphereOneBlend (d := d) (unitBallApproxEps n) x = 1 := by
        apply sphereOneBlend_eq_one_of_norm_ge (d := d) (unitBallApproxEps_pos n); linarith [hx2.1]
      have hnotclosed1 : x ∉ Metric.closedBall (0 : E) 1 := by
        rw [Metric.mem_closedBall, dist_zero_right]; linarith [hx2.1, hεsmall]
      by_cases hxlt2 : ‖x‖ < 2
      · have hxOuter : x ∈ unitBallOuterShell (d := d) := ⟨by linarith [hx2.1], hxlt2⟩
        have hExactModel :
            exactUnitBallExtensionModel (d := d) ψ x = unitBallShellFormula (d := d) ψ x := by
          unfold exactUnitBallExtensionModel exactUnitBallShellOrZero
          simp [hxOuter, hnotclosed1]
        have hEq :
            smoothUnitBallExtensionApprox (d := d) (unitBallApproxEps n) ψ x -
              exactUnitBallExtensionModel (d := d) ψ x =
              (sphereTwoBlend (d := d) (unitBallApproxEps n) x - 1) *
                unitBallShellFormula (d := d) ψ x := by
          rw [hExactModel]; unfold smoothUnitBallExtensionApprox; simp [hβ1, sub_eq_add_neg]; ring
        calc ‖smoothUnitBallExtensionApprox (d := d) (unitBallApproxEps n) ψ x -
              exactUnitBallExtensionModel (d := d) ψ x‖
              = ‖(sphereTwoBlend (d := d) (unitBallApproxEps n) x - 1) *
                  unitBallShellFormula (d := d) ψ x‖ := by rw [hEq]
          _ ≤ ‖unitBallShellFormula (d := d) ψ x‖ := by
                rw [norm_mul, Real.norm_eq_abs, abs_sub_comm]
                have h1 : |1 - sphereTwoBlend (d := d) (unitBallApproxEps n) x| ≤ 1 := by
                  rw [abs_of_nonneg (by linarith [sphereTwoBlend_le_one (d := d) (ε := unitBallApproxEps n) (x := x)])]
                  linarith [sphereTwoBlend_nonneg (d := d) (ε := unitBallApproxEps n) (x := x)]
                nlinarith [norm_nonneg (unitBallShellFormula (d := d) ψ x)]
          _ ≤ C2 * unitBallApproxEps n := hC2 n x hx2
          _ ≤ max (C1 / 5) (C2 / 5) := by
                have := unitBallApproxEps_le_one_fifth n
                exact le_trans (by nlinarith [hC2_nonneg]) (le_max_right _ _)
      · have hxge2 : 2 ≤ ‖x‖ := by linarith
        have hnotOuter : x ∉ unitBallOuterShell (d := d) := fun h => not_lt_of_ge hxge2 h.2
        have hExactModel :
            exactUnitBallExtensionModel (d := d) ψ x = 0 := by
          unfold exactUnitBallExtensionModel exactUnitBallShellOrZero
          simp [hnotclosed1, hnotOuter]
        have hEq :
            smoothUnitBallExtensionApprox (d := d) (unitBallApproxEps n) ψ x -
              exactUnitBallExtensionModel (d := d) ψ x =
              sphereTwoBlend (d := d) (unitBallApproxEps n) x *
                unitBallShellFormula (d := d) ψ x := by
          rw [hExactModel]; unfold smoothUnitBallExtensionApprox; simp [hβ1, sub_eq_add_neg]
        calc ‖smoothUnitBallExtensionApprox (d := d) (unitBallApproxEps n) ψ x -
              exactUnitBallExtensionModel (d := d) ψ x‖
              = ‖sphereTwoBlend (d := d) (unitBallApproxEps n) x *
                  unitBallShellFormula (d := d) ψ x‖ := by rw [hEq]
          _ ≤ ‖unitBallShellFormula (d := d) ψ x‖ := by
                rw [norm_mul, Real.norm_eq_abs]
                nlinarith [abs_of_nonneg (sphereTwoBlend_nonneg (d := d) (ε := unitBallApproxEps n) (x := x)),
                  sphereTwoBlend_le_one (d := d) (ε := unitBallApproxEps n) (x := x),
                  norm_nonneg (unitBallShellFormula (d := d) ψ x)]
          _ ≤ C2 * unitBallApproxEps n := hC2 n x hx2
          _ ≤ max (C1 / 5) (C2 / 5) := by
                have := unitBallApproxEps_le_one_fifth n
                exact le_trans (by nlinarith [hC2_nonneg]) (le_max_right _ _)
  · rw [Set.indicator_of_notMem hbad]
    rw [smoothUnitBallExtensionApprox_eq_unitBallExtension_of_not_mem_badSet (d := d) hbad,
      ← exactUnitBallExtensionModel_eq_unitBallExtension (d := d) (ψ := ψ)]
    simp

\end{lstlisting}
\begin{proof}
From Theorem~\ref{thm:err-bound} the pointwise difference
$|\smoothUnitBallExtensionApprox_{\varepsilon_n}\psi(x) - \Eext\psi(x)|$ is
bounded by $C\,\varepsilon_n \cdot \mathbf{1}_{\mathrm{Bad}_{\varepsilon_n}}(x)$
for some constant $C$ depending only on $\psi$. Pointwise both sides are
also dominated by the constant $C$ on the closed ball $\overline{\Ball{9/4}{0}}$:
outside the ball of radius $9/4$ both
$\smoothUnitBallExtensionApprox_{\varepsilon_n}\psi$ and $\Eext\psi$ vanish
(using \leanname{sphereOneControl\_subset\_closedBall},
\leanname{sphereTwoControl\_subset\_closedBall} for the bad annuli together
with \leanname{unitBallExtension\_hasCompactSupport}). Hence the family of
differences is uniformly $L^p$-dominated by the constant times indicator
$C\,\mathbf{1}_{\overline{\Ball{9/4}{0}}}$, which lies in $L^p(E)$ since the
underlying ball has finite Lebesgue measure
(\leanname{measure\_closedBall\_lt\_top}).

Since the bad annuli shrink to the spheres of radii $1$ and $2$ and the Haar
measure of those spheres is zero (\leanname{Measure.addHaar\_sphere}),
\leanname{tendsto\_measure\_unitBallBadSet} gives
$|\mathrm{Bad}_{\varepsilon_n}| \to 0$. Combining the pointwise bound with
the dominated convergence theorem in $L^p$, against the constant dominating
indicator, yields
\[
  \eLpNormText{\smoothUnitBallExtensionApprox_{\varepsilon_n}\psi
    - \Eext\psi}_{L^p(E)}
  \le K\,\varepsilon_n\,|\mathrm{Bad}_{\varepsilon_n}|^{1/p}
  \to 0,
\]
where $K$ is the constant from Theorem~\ref{thm:err-bound}, and the
$|\mathrm{Bad}_{\varepsilon_n}|^{1/p}$ factor arises from H\"older's inequality
applied to the indicator on the bad set.
\end{proof}

\subsubsection*{The smoothing theorem on a smooth input}

\begin{theorem}[\leanname{smooth\_input\_unitBallExtension\_smoothing}]
\label{thm:smooth-input-smoothing}
Let $\psi \in C^\infty(E)$ have compact support, $1 < p < \infty$. Then there
is a sequence
\[
  \Psi_n := \smoothUnitBallExtensionApprox_{\varepsilon_n}\psi \in C^\infty_c(E)
\]
such that:
\begin{enumerate}
\item $\Psi_n \to \Eext \psi$ in $L^p(E)$;
\item $\partial_i \Psi_n \to (\smoothUnitBallExtensionGrad\psi)_i$ in $L^p(E)$
for each $i$;
\item the $L^p$ norm of $\Psi_n$ is bounded uniformly by the smooth-input
extension estimate of Theorem~\ref{thm:smooth-extension-estimate}.
\end{enumerate}
\end{theorem}

\begin{lstlisting}[style=Matlab-editor,basicstyle=\mlttfamily,escapechar=",mlshowsectionrules=true,mathescape=true,morecomment={[s]{\%\{}{\%\}}},language=lean,numbers=none]
/-- Smooth-input interface smoothing for the explicit extension operator.

For smooth compactly supported input `ψ`, the piecewise extension
`unitBallExtension ψ` can itself be approximated globally in full `W^{1,p}` by
smooth compactly supported functions. The gradient side is expressed against
some global field `Gψ` attached to the exact extension.

The surrounding rough-data extension theorem follows from this together with the
already available unit-ball approximation theorem.
-/
theorem smooth_input_unitBallExtension_smoothing
    {p : ℝ} (hp : 1 < p) {ψ : E → ℝ}
    (hψ_smooth : ContDiff ℝ (⊤ : ℕ∞) ψ)
    (_hψ_cpt : HasCompactSupport ψ) :
    ∃ Gψ : E → E,
      MemLp (unitBallExtension (d := d) ψ) (ENNReal.ofReal p) volume ∧
      (∀ i : Fin d, MemLp (fun x => Gψ x i) (ENNReal.ofReal p) volume) ∧
      (∫⁻ x, (ENNReal.ofReal |unitBallExtension (d := d) ψ x|) ^ p ∂volume)
        ≤ C_unitBallExtensionFun d *
          ∫⁻ x in Metric.ball (0 : E) 1, (ENNReal.ofReal |ψ x|) ^ p ∂volume ∧
      (∫⁻ x, (ENNReal.ofReal ‖Gψ x‖) ^ p ∂volume)
        ≤ (∫⁻ x in Metric.ball (0 : E) 1, (ENNReal.ofReal ‖fderiv ℝ ψ x‖) ^ p ∂volume) +
          C_unitBallExtensionGrad d p *
            (∫⁻ x in Metric.ball (0 : E) 1, (ENNReal.ofReal |ψ x|) ^ p ∂volume +
             ∫⁻ x in Metric.ball (0 : E) 1, (ENNReal.ofReal ‖fderiv ℝ ψ x‖) ^ p ∂volume) ∧
      ∃ Φ : ℕ → E → ℝ,
        (∀ n, ContDiff ℝ (⊤ : ℕ∞) (Φ n)) ∧
        (∀ n, HasCompactSupport (Φ n)) ∧
        Tendsto
          (fun n =>
            eLpNorm (fun x => Φ n x - unitBallExtension (d := d) ψ x)
              (ENNReal.ofReal p) volume)
          atTop (nhds 0) ∧
        ∀ i : Fin d,
          Tendsto
            (fun n =>
              eLpNorm
                (fun x => (fderiv ℝ (Φ n) x) (EuclideanSpace.single i 1) - Gψ x i)
                (ENNReal.ofReal p) volume)
            atTop (nhds 0) := by
  have hΦ_smooth : ∀ n, ContDiff ℝ (⊤ : ℕ∞)
      (smoothUnitBallExtensionApprox (d := d) (unitBallApproxEps n) ψ) :=
    fun n => smoothUnitBallExtensionApprox_contDiff (d := d) (unitBallApproxEps_pos n)
      (unitBallApproxEps_lt_one n) hψ_smooth
  have hΦ_cpt : ∀ n, HasCompactSupport
      (smoothUnitBallExtensionApprox (d := d) (unitBallApproxEps n) ψ) :=
    fun n => smoothUnitBallExtensionApprox_hasCompactSupport (d := d) (ψ := ψ)
      (unitBallApproxEps_pos n)
  have huExt_memLp : MemLp (unitBallExtension (d := d) ψ) (ENNReal.ofReal p) volume := by
    have hΦ0_memLp : MemLp (smoothUnitBallExtensionApprox (d := d) (unitBallApproxEps 0) ψ)
        (ENNReal.ofReal p) (volume : Measure E) :=
      (hΦ_smooth 0).continuous.memLp_of_hasCompactSupport (hΦ_cpt 0)
    have hDiff_memLp :=
      memLp_smoothUnitBallExtensionApprox_sub_unitBallExtension (d := d) hp hψ_smooth 0
    have hTmp := hΦ0_memLp.sub hDiff_memLp
    exact memLp_congr_ae (Filter.Eventually.of_forall fun x => by simp [sub_sub_cancel]) |>.mp hTmp
  have hGψ_memLp : ∀ i : Fin d,
      MemLp (fun x => (exactUnitBallExtensionGrad (d := d) ψ x) i)
        (ENNReal.ofReal p) volume := by
    intro i
    have hderiv_memLp :
        MemLp
          (fun x => (fderiv ℝ (smoothUnitBallExtensionApprox (d := d) (unitBallApproxEps 0) ψ) x)
            (EuclideanSpace.single i 1))
          (ENNReal.ofReal p) volume := by
      have hderiv_smooth :
          ContDiff ℝ (⊤ : ℕ∞)
            (fun x => (fderiv ℝ (smoothUnitBallExtensionApprox (d := d) (unitBallApproxEps 0) ψ) x)
              (EuclideanSpace.single i 1)) :=
        (hΦ_smooth 0).fderiv_right (m := (⊤ : ℕ∞)) (by simp)
          |>.clm_apply contDiff_const
      exact hderiv_smooth.continuous.memLp_of_hasCompactSupport
        ((hΦ_cpt 0).fderiv_apply (𝕜 := ℝ) _)
    have hDiff_memLp :=
      memLp_fderiv_smoothUnitBallExtensionApprox_sub_exactGradApply (d := d) hp hψ_smooth 0 i
    have hTmp := hderiv_memLp.sub hDiff_memLp
    exact memLp_congr_ae (Filter.Eventually.of_forall fun x => by
      simp [sub_sub_cancel, exactUnitBallExtensionGrad,
        SmoothApproximationInternal.exactUnitBallExtensionGrad, exactUnitBallExtensionGradApply])
      |>.mp hTmp
  have hweak :
      ∀ i : Fin d, HasWeakPartialDeriv i
        (fun x => (exactUnitBallExtensionGrad (d := d) ψ x) i)
        (unitBallExtension (d := d) ψ) Set.univ := by
    have hΦ_tendsto_fun :=
      tendsto_eLpNorm_smoothUnitBallExtensionApprox_sub_unitBallExtension (d := d) hp hψ_smooth
    have hΦ_tendsto_grad :=
      tendsto_eLpNorm_fderiv_smoothUnitBallExtensionApprox_sub_exactGradApply (d := d) hp hψ_smooth
    exact hasWeakPartials_of_global_smoothApprox (d := d) hp huExt_memLp hGψ_memLp
      hΦ_smooth hΦ_cpt hΦ_tendsto_fun (fun i => by
        convert hΦ_tendsto_grad i using 2)
  let hwExt : MemW1pWitness (ENNReal.ofReal p)
      (unitBallExtension (d := d) ψ) Set.univ :=
    { memLp := by rw [Measure.restrict_univ]; exact huExt_memLp
      weakGrad := exactUnitBallExtensionGrad (d := d) ψ
      weakGrad_component_memLp := by
        intro i; rw [Measure.restrict_univ]; exact hGψ_memLp i
      isWeakGrad := hweak }
  rcases exists_global_smooth_W1p_approx_of_localizedWitness
      (d := d) (Ω := Set.univ) isOpen_univ hp hwExt
      (unitBallExtension_hasCompactSupport (d := d) ψ) (Set.subset_univ _) with
    ⟨Ψ, hΨ_smooth, hΨ_cpt, hΨ_fun, hΨ_grad⟩
  have hψ_smooth1 : ContDiff ℝ 1 ψ := hψ_smooth.of_le (by simp)
  have hest := smooth_unitBallExtension_W1p_estimate (d := d) (p := p) hp.le hψ_smooth1
  have hfun_bound :
      (∫⁻ x, (ENNReal.ofReal |unitBallExtension (d := d) ψ x|) ^ p ∂volume)
        ≤ C_unitBallExtensionFun d *
          ∫⁻ x in Metric.ball (0 : E) 1, (ENNReal.ofReal |ψ x|) ^ p ∂volume := by
    have hfun_support :
        (fun x => (ENNReal.ofReal |unitBallExtension (d := d) ψ x|) ^ p) =
          (Metric.ball (0 : E) 2).indicator
            (fun x => (ENNReal.ofReal |unitBallExtension (d := d) ψ x|) ^ p) := by
      funext x
      by_cases hx : x ∈ Metric.ball (0 : E) 2
      · simp [hx]
      · have hnorm : 2 ≤ ‖x‖ := by
          simp only [Metric.mem_ball, dist_zero_right, not_lt] at hx; exact hx
        have hzero := unitBallExtension_eq_zero_of_two_le_norm (d := d) (u := ψ) hnorm
        simp only [hzero, abs_zero, ENNReal.ofReal_zero, Set.indicator_of_notMem hx,
          ENNReal.zero_rpow_of_pos (show (0 : ℝ) < p by linarith)]
    rw [hfun_support, lintegral_indicator Metric.isOpen_ball.measurableSet]
    exact hest.1
  have hgrad_bound := exactUnitBallExtensionGrad_bound (d := d) hp hψ_smooth
  refine ⟨exactUnitBallExtensionGrad (d := d) ψ,
    huExt_memLp, hGψ_memLp, hfun_bound, hgrad_bound,
    Ψ, hΨ_smooth, hΨ_cpt, hΨ_fun, ?_⟩
  intro i
  have h := hΨ_grad i
  refine (Filter.tendsto_congr fun n => eLpNorm_congr_ae ?_).mpr h
  filter_upwards with x; simp [hwExt]


\end{lstlisting}
\begin{proof}
The first item is Theorem~\ref{thm:err-bound} together with control of the
bad set's measure. For the second, the same blend-region argument applies to
$D\Psi_n - D\Eext\psi$: by Lemma~\ref{lem:smoothapprox-pointwise}, the
derivatives converge pointwise outside the spheres of radii $1$ and $2$ to
the smooth extension derivatives, while the bad-set contributions are
controlled by the uniform $C^1$ bounds on the buffer compacts. Item (3) is
an application of Theorem~\ref{thm:smooth-extension-estimate} to $\psi$.
\end{proof}

\begin{theorem}[\leanname{hasWeakPartials\_of\_global\_smoothApprox}]
Let $1 < p < \infty$, let $f \in L^p(E)$ have a candidate gradient $G \in
L^p(E; \R^d)$, and assume there exist $\Psi_n \in C^\infty_c(E)$ with
$\Psi_n \to f$ and $\partial_i \Psi_n \to G_i$ in $L^p(E)$ for each $i$. Then
$G_i$ is the weak $i$-th partial derivative of $f$ on $E$.
\end{theorem}

\begin{lstlisting}[style=Matlab-editor,basicstyle=\mlttfamily,escapechar=",mlshowsectionrules=true,mathescape=true,morecomment={[s]{\%\{}{\%\}}},language=lean,numbers=none]
/-!
This block smooths the explicit unit-ball extension on smooth compactly
supported inputs by blending across shrinking annuli around `‖x‖ = 1` and
`‖x‖ = 2`. It provides the final analytic bridge needed to upgrade the explicit
extension operator to a full `W^{1,p}` witness on rough data.
-/

theorem hasWeakPartials_of_global_smoothApprox
    {p : ℝ} (hp : 1 < p)
    {v : E → ℝ} {G : E → E} {Φ : ℕ → E → ℝ}
    (hv_memLp : MemLp v (ENNReal.ofReal p) volume)
    (hG_memLp : ∀ i : Fin d, MemLp (fun x => G x i) (ENNReal.ofReal p) volume)
    (hΦ_smooth : ∀ n, ContDiff ℝ (⊤ : ℕ∞) (Φ n))
    (hΦ_cpt : ∀ n, HasCompactSupport (Φ n))
    (hΦ_fun :
      Tendsto
        (fun n => eLpNorm (fun x => Φ n x - v x) (ENNReal.ofReal p) volume)
        atTop (nhds 0))
    (hΦ_grad :
      ∀ i : Fin d,
        Tendsto
          (fun n =>
            eLpNorm
              (fun x => (fderiv ℝ (Φ n) x) (EuclideanSpace.single i 1) - G x i)
              (ENNReal.ofReal p) volume)
          atTop (nhds 0)) :
    ∀ i : Fin d, HasWeakPartialDeriv i (fun x => G x i) v Set.univ := by
  let _ := (inferInstance : NeZero d)
  intro i
  have hΦ_wd : ∀ n, HasWeakPartialDeriv i
      (fun x => (fderiv ℝ (Φ n) x) (EuclideanSpace.single i 1)) (Φ n) Set.univ := by
    intro n
    simpa using (HasWeakPartialDeriv.of_contDiff (Ω := Set.univ) (i := i) (f := Φ n)
      isOpen_univ ((hΦ_smooth n).of_le (by simp)))
  have hΦ_fun_memLp : ∀ n, MemLp (fun x => Φ n x - v x) (ENNReal.ofReal p) volume := by
    intro n
    exact ((hΦ_smooth n).continuous.memLp_of_hasCompactSupport (hΦ_cpt n)).sub hv_memLp
  have hΦ_grad_memLp : ∀ n,
      MemLp (fun x => (fderiv ℝ (Φ n) x) (EuclideanSpace.single i 1) - G x i)
        (ENNReal.ofReal p) volume := by
    intro n
    have hderiv_smooth : ContDiff ℝ (⊤ : ℕ∞)
        (fun x => (fderiv ℝ (Φ n) x) (EuclideanSpace.single i 1)) :=
      (hΦ_smooth n).fderiv_right (m := (⊤ : ℕ∞)) (by simp) |>.clm_apply contDiff_const
    exact (hderiv_smooth.continuous.memLp_of_hasCompactSupport
      ((hΦ_cpt n).fderiv_apply (𝕜 := ℝ) _)).sub (hG_memLp i)
  exact HasWeakPartialDeriv.of_eLpNormApprox_p (d := d) (Ω := Set.univ) isOpen_univ hp
    (by simpa using hv_memLp) (by simpa using hG_memLp i)
    hΦ_wd (fun n => by simpa using hΦ_fun_memLp n) (by simpa using hΦ_fun)
    (fun n => by simpa using hΦ_grad_memLp n) (by simpa using hΦ_grad i)

\end{lstlisting}
\begin{proof}
For $\varphi \in C^\infty_c(E)$, the integration-by-parts identity
$\int \Psi_n\, \partial_i \varphi = -\int (\partial_i \Psi_n)\, \varphi$
holds for each smooth $\Psi_n$. Pass to the limit using $L^p \times L^q$
duality with $1/p + 1/q = 1$ (note $\partial_i\varphi \in L^q$ and
$\varphi \in L^q$ since both are bounded compactly supported), giving
$\int f\, \partial_i \varphi = -\int G_i\, \varphi$, the weak derivative
identity.
\end{proof}

\subsubsection*{The main extension theorems}

\begin{theorem}[\leanname{exists\_smooth\_global\_approx\_of\_unitBallExtension}]
\label{thm:global-approx-extension}
Let $1 < p < \infty$ and $u \in W^{1,p}(\Bz{1})$. There exists
$G_{\mathrm{ext}} : E \to \R^d$ with the following properties.
\begin{enumerate}
\item $\Eext u \in L^p(E)$ and each component of $G_{\mathrm{ext}}$ lies in
$L^p(E)$.
\item Function-side estimate:
\[
\int_E |\Eext u(x)|^p\,dx
\le \CunitBallExtensionFun(d)\,\int_{\Bz{1}} |u(x)|^p\,dx,
\]
with $\CunitBallExtensionFun(d) := 1 + 2^{2d}$.
\item Gradient-side estimate:
\begin{multline*}
\int_E \|G_{\mathrm{ext}}(x)\|^p\,dx
\le \int_{\Bz{1}} \|\nabla u(x)\|^p\,dx \\
+ \CunitBallExtensionGrad(d, p)\,\biggl(\int_{\Bz{1}} |u(x)|^p\,dx
+ \int_{\Bz{1}} \|\nabla u(x)\|^p\,dx\biggr),
\end{multline*}
with $\CunitBallExtensionGrad(d,p) := 2^{2d}\cdot 2^{p-1}$.
\item There exists a sequence $\Phi_n \in C^\infty_c(E)$ with
$\Phi_n \to \Eext u$ and $\partial_i \Phi_n \to (G_{\mathrm{ext}})_i$ in
$L^p(E)$.
\end{enumerate}
\end{theorem}

\begin{lstlisting}[style=Matlab-editor,basicstyle=\mlttfamily,escapechar=",mlshowsectionrules=true,mathescape=true,morecomment={[s]{\%\{}{\%\}}},language=lean,numbers=none]
theorem exists_smooth_global_approx_of_unitBallExtension
    {p : ℝ} (hp : 1 < p) {u : E → ℝ}
    (hw : MemW1pWitness (ENNReal.ofReal p) u (Metric.ball (0 : E) 1)) :
    ∃ Gext : E → E,
      MemLp (unitBallExtension (d := d) u) (ENNReal.ofReal p) volume ∧
      (∀ i : Fin d, MemLp (fun x => Gext x i) (ENNReal.ofReal p) volume) ∧
      (∫⁻ x, (ENNReal.ofReal |unitBallExtension (d := d) u x|) ^ p ∂volume)
        ≤ C_unitBallExtensionFun d *
          ∫⁻ x in Metric.ball (0 : E) 1, (ENNReal.ofReal |u x|) ^ p ∂volume ∧
      (∫⁻ x, (ENNReal.ofReal ‖Gext x‖) ^ p ∂volume)
        ≤ (∫⁻ x in Metric.ball (0 : E) 1, (ENNReal.ofReal ‖hw.weakGrad x‖) ^ p ∂volume) +
          C_unitBallExtensionGrad d p *
            (∫⁻ x in Metric.ball (0 : E) 1, (ENNReal.ofReal |u x|) ^ p ∂volume +
             ∫⁻ x in Metric.ball (0 : E) 1, (ENNReal.ofReal ‖hw.weakGrad x‖) ^ p ∂volume) ∧
      ∃ Φ : ℕ → E → ℝ,
        (∀ n, ContDiff ℝ (⊤ : ℕ∞) (Φ n)) ∧
        (∀ n, HasCompactSupport (Φ n)) ∧
        Tendsto
          (fun n =>
            eLpNorm (fun x => Φ n x - unitBallExtension (d := d) u x)
              (ENNReal.ofReal p) volume)
          atTop (nhds 0) ∧
        ∀ i : Fin d,
          Tendsto
            (fun n =>
              eLpNorm
                (fun x => (fderiv ℝ (Φ n) x) (EuclideanSpace.single i 1) - Gext x i)
                (ENNReal.ofReal p) volume)
            atTop (nhds 0) := by
  rcases exists_smooth_W1p_approx_on_unitBall (d := d) (p := p) hp hw with
    ⟨ψ, hψ_smooth, hψ_cpt, hψ_fun, hψ_grad⟩
  let B : Set E := Metric.ball (0 : E) 1
  let μB : Measure E := volume.restrict B
  -- Apply Layer 3 to each ψ_n
  have h_step : ∀ n,
      ∃ Gψ : E → E,
        MemLp (unitBallExtension (d := d) (ψ n)) (ENNReal.ofReal p) volume ∧
        (∀ i : Fin d, MemLp (fun x => Gψ x i) (ENNReal.ofReal p) volume) ∧
        (∫⁻ x, (ENNReal.ofReal |unitBallExtension (d := d) (ψ n) x|) ^ p ∂volume)
          ≤ C_unitBallExtensionFun d *
            ∫⁻ x in B, (ENNReal.ofReal |ψ n x|) ^ p ∂volume ∧
        (∫⁻ x, (ENNReal.ofReal ‖Gψ x‖) ^ p ∂volume)
          ≤ (∫⁻ x in B, (ENNReal.ofReal ‖fderiv ℝ (ψ n) x‖) ^ p ∂volume) +
            C_unitBallExtensionGrad d p *
              (∫⁻ x in B, (ENNReal.ofReal |ψ n x|) ^ p ∂volume +
               ∫⁻ x in B, (ENNReal.ofReal ‖fderiv ℝ (ψ n) x‖) ^ p ∂volume) ∧
        ∃ Φ : ℕ → E → ℝ,
          (∀ m, ContDiff ℝ (⊤ : ℕ∞) (Φ m)) ∧
          (∀ m, HasCompactSupport (Φ m)) ∧
          Tendsto (fun m => eLpNorm (fun x => Φ m x - unitBallExtension (d := d) (ψ n) x)
            (ENNReal.ofReal p) volume) atTop (nhds 0) ∧
          ∀ i : Fin d,
            Tendsto (fun m => eLpNorm
              (fun x => (fderiv ℝ (Φ m) x) (EuclideanSpace.single i 1) - Gψ x i)
              (ENNReal.ofReal p) volume) atTop (nhds 0) :=
    fun n => smooth_input_unitBallExtension_smoothing (d := d) hp (hψ_smooth n) (hψ_cpt n)
  let G : ℕ → E → E := fun n => Classical.choose (h_step n)
  have hV_memLp : ∀ n, MemLp (unitBallExtension (d := d) (ψ n)) (ENNReal.ofReal p) volume :=
    fun n => (Classical.choose_spec (h_step n)).1
  have hG_comp_memLp : ∀ n i, MemLp (fun x => G n x i) (ENNReal.ofReal p) volume :=
    fun n i => (Classical.choose_spec (h_step n)).2.1 i
  have hG_bound : ∀ n,
      (∫⁻ x, (ENNReal.ofReal ‖G n x‖) ^ p ∂volume)
        ≤ (∫⁻ x in B, (ENNReal.ofReal ‖fderiv ℝ (ψ n) x‖) ^ p ∂volume) +
          C_unitBallExtensionGrad d p *
            (∫⁻ x in B, (ENNReal.ofReal |ψ n x|) ^ p ∂volume +
             ∫⁻ x in B, (ENNReal.ofReal ‖fderiv ℝ (ψ n) x‖) ^ p ∂volume) :=
    fun n => (Classical.choose_spec (h_step n)).2.2.2.1
  have h_step_approx : ∀ n, ∃ Φ : ℕ → E → ℝ,
      (∀ m, ContDiff ℝ (⊤ : ℕ∞) (Φ m)) ∧ (∀ m, HasCompactSupport (Φ m)) ∧
      Tendsto (fun m => eLpNorm (fun x => Φ m x - unitBallExtension (d := d) (ψ n) x)
        (ENNReal.ofReal p) volume) atTop (nhds 0) ∧
      ∀ i : Fin d, Tendsto (fun m => eLpNorm
        (fun x => (fderiv ℝ (Φ m) x) (EuclideanSpace.single i 1) - G n x i)
        (ENNReal.ofReal p) volume) atTop (nhds 0) :=
    fun n => (Classical.choose_spec (h_step n)).2.2.2.2
  have hV_wd : ∀ n i,
      HasWeakPartialDeriv i (fun x => G n x i)
        (unitBallExtension (d := d) (ψ n)) Set.univ := by
    intro n i
    rcases h_step_approx n with ⟨Φn, hΦn_smooth, hΦn_cpt, hΦn_fun, hΦn_grad⟩
    exact (hasWeakPartials_of_global_smoothApprox (d := d) hp
      (hV_memLp n) (fun j => hG_comp_memLp n j)
      hΦn_smooth hΦn_cpt hΦn_fun hΦn_grad) i
  letI : Fact (1 ≤ ENNReal.ofReal p) := ⟨by simpa using ENNReal.ofReal_le_ofReal hp.le⟩
  -- Function energy bound (doesn't need MemLp)
  have huExt_fun_bound :
      (∫⁻ x, (ENNReal.ofReal |unitBallExtension (d := d) u x|) ^ p ∂volume)
        ≤ C_unitBallExtensionFun d *
          ∫⁻ x in B, (ENNReal.ofReal |u x|) ^ p ∂volume := by
    have h0 := lintegral_rpow_abs_unitBallExtension_sub_le_local (d := d) (p := p)
      (by linarith : 0 < p) (u := u) (v := fun _ => 0)
    simp only [unitBallExtension_zero, sub_zero] at h0
    exact h0
  -- V_n → unitBallExtension u in eLpNorm (doesn't need MemLp for uExt)
  have hV_fun :
      Tendsto (fun n => eLpNorm (fun x => unitBallExtension (d := d) (ψ n) x -
        unitBallExtension (d := d) u x) (ENNReal.ofReal p) volume) atTop (nhds 0) := by
    let K : ℝ≥0∞ := (1 + ENNReal.ofReal ((2 : ℝ) ^ (2 * d))) ^ (1 / p)
    have hbound : ∀ n,
        eLpNorm (fun x => unitBallExtension (d := d) (ψ n) x -
          unitBallExtension (d := d) u x) (ENNReal.ofReal p) volume
        ≤ K * eLpNorm (fun x => ψ n x - u x) (ENNReal.ofReal p) μB :=
      fun n => eLpNorm_unitBallExtension_sub_le_local (d := d) hp
    have hKψ : Tendsto (fun n => K * eLpNorm (fun x => ψ n x - u x)
        (ENNReal.ofReal p) μB) atTop (nhds 0) := by
      have h := ENNReal.Tendsto.const_mul hψ_fun (Or.inr (show K ≠ ⊤ by simp [K]))
      rwa [mul_zero] at h
    exact tendsto_of_tendsto_of_tendsto_of_le_of_le tendsto_const_nhds hKψ
      (fun _ => bot_le) hbound
  -- MemLp for unitBallExtension u: V_n Cauchy → Lp limit → ae_eq uExt → transfer
  have huExt_memLp : MemLp (unitBallExtension (d := d) u) (ENNReal.ofReal p) volume := by
    -- V_n is pair-Cauchy
    have hψ_memLp_restrict : ∀ n, MemLp (ψ n) (ENNReal.ofReal p) μB :=
      fun n => ((hψ_smooth n).continuous.memLp_of_hasCompactSupport (hψ_cpt n)).restrict B
    have hV_pair : Tendsto (fun nm : ℕ × ℕ =>
        eLpNorm (fun x => unitBallExtension (d := d) (ψ nm.1) x -
          unitBallExtension (d := d) (ψ nm.2) x) (ENNReal.ofReal p) volume)
        atTop (nhds 0) := by
      let K : ℝ≥0∞ := (1 + ENNReal.ofReal ((2 : ℝ) ^ (2 * d))) ^ (1 / p)
      have hKψ_fun : Tendsto (fun n =>
          K * eLpNorm (fun x => ψ n x - u x) (ENNReal.ofReal p) μB) atTop (nhds 0) := by
        have h := ENNReal.Tendsto.const_mul hψ_fun (Or.inr (show K ≠ ⊤ by simp [K]))
        rwa [mul_zero] at h
      refine tendsto_zero_of_le_pair_sum (fun nm => ?_) hKψ_fun
      calc eLpNorm (fun x => unitBallExtension (d := d) (ψ nm.1) x -
              unitBallExtension (d := d) (ψ nm.2) x) (ENNReal.ofReal p) volume
          ≤ K * eLpNorm (fun x => ψ nm.1 x - ψ nm.2 x) (ENNReal.ofReal p) μB :=
            eLpNorm_unitBallExtension_sub_le_local (d := d) hp
        _ ≤ K * (eLpNorm (fun x => ψ nm.1 x - u x) (ENNReal.ofReal p) μB +
              eLpNorm (fun x => ψ nm.2 x - u x) (ENNReal.ofReal p) μB) := by
            gcongr; rw [show (fun x => ψ nm.1 x - ψ nm.2 x) =
              (fun x => (ψ nm.1 x - u x) - (ψ nm.2 x - u x)) from by ext x; ring]
            exact eLpNorm_sub_le ((hψ_memLp_restrict nm.1).sub hw.memLp).aestronglyMeasurable
              ((hψ_memLp_restrict nm.2).sub hw.memLp).aestronglyMeasurable
              (by simpa using ENNReal.ofReal_le_ofReal hp.le)
        _ = K * eLpNorm (fun x => ψ nm.1 x - u x) (ENNReal.ofReal p) μB +
              K * eLpNorm (fun x => ψ nm.2 x - u x) (ENNReal.ofReal p) μB := by ring
    -- Extract Lp limit with MemLp
    obtain ⟨f_lim, hf_lim_memLp, hf_lim_tendsto⟩ :=
      BareFunction.scalar_cauchy_to_limit
        (by simpa using ENNReal.ofReal_le_ofReal hp.le) (by simp) hV_memLp hV_pair
    -- f_lim =ᵐ unitBallExtension u (both are eLpNorm limits of V_n)
    -- eLpNorm(V_n - f_lim) → 0 and eLpNorm(V_n - uExt) → 0, so eLpNorm(f_lim - uExt) = 0
    have hae : f_lim =ᵐ[volume] unitBallExtension (d := d) u := by
      have huExt_aesm : AEStronglyMeasurable (unitBallExtension (d := d) u) volume :=
        aestronglyMeasurable_unitBallExtension_of_memLp hw.memLp
      exact BareFunction.ae_eq_of_tendsto_eLpNorm_sub
        (by simpa using ENNReal.ofReal_le_ofReal hp.le)
        (fun n => (hV_memLp n).aestronglyMeasurable)
        hf_lim_memLp.aestronglyMeasurable huExt_aesm
        hf_lim_tendsto
        (by exact hV_fun)
    exact (memLp_congr_ae hae).mp hf_lim_memLp
  have hG_comp_cauchy : ∀ i : Fin d, Tendsto (fun nm : ℕ × ℕ =>
      eLpNorm (fun x => G nm.1 x i - G nm.2 x i) (ENNReal.ofReal p) volume)
      atTop (nhds 0) := by
    intro i
    -- ∫ |G nm.1 · i - G nm.2 · i|^p ≤ ∫ ‖exactGrad(ψ nm.1 - ψ nm.2)‖^p
    have hcomp_le : ∀ nm : ℕ × ℕ,
        ∫⁻ x, (ENNReal.ofReal |G nm.1 x i - G nm.2 x i|) ^ p ∂volume ≤
          ∫⁻ x, (ENNReal.ofReal ‖exactUnitBallExtensionGrad (d := d)
            (fun x => ψ nm.1 x - ψ nm.2 x) x‖) ^ p ∂volume :=
      fun nm => weakGrad_component_cauchy_bound (d := d) hp
        (hψ_smooth nm.1) (hψ_smooth nm.2) (hψ_cpt nm.1) (hψ_cpt nm.2)
        (hV_wd nm.1 i) (hV_wd nm.2 i) (hG_comp_memLp nm.1 i) (hG_comp_memLp nm.2 i)
    have hgrad_bound_nm : ∀ nm : ℕ × ℕ,
        ∫⁻ x, (ENNReal.ofReal ‖exactUnitBallExtensionGrad (d := d)
          (fun x => ψ nm.1 x - ψ nm.2 x) x‖) ^ p ∂volume ≤
          (∫⁻ x in B, (ENNReal.ofReal ‖fderiv ℝ (fun x => ψ nm.1 x - ψ nm.2 x) x‖) ^ p
            ∂volume) +
          C_unitBallExtensionGrad d p *
            (∫⁻ x in B, (ENNReal.ofReal |ψ nm.1 x - ψ nm.2 x|) ^ p ∂volume +
             ∫⁻ x in B, (ENNReal.ofReal ‖fderiv ℝ (fun x => ψ nm.1 x - ψ nm.2 x) x‖) ^ p
              ∂volume) :=
      fun nm => exactUnitBallExtensionGrad_bound (d := d) hp
        ((hψ_smooth nm.1).sub (hψ_smooth nm.2))
    have hfull_le : ∀ nm : ℕ × ℕ,
        ∫⁻ x, (ENNReal.ofReal |G nm.1 x i - G nm.2 x i|) ^ p ∂volume ≤
          (∫⁻ x in B, (ENNReal.ofReal ‖fderiv ℝ (fun x => ψ nm.1 x - ψ nm.2 x) x‖) ^ p
            ∂volume) +
          C_unitBallExtensionGrad d p *
            (∫⁻ x in B, (ENNReal.ofReal |ψ nm.1 x - ψ nm.2 x|) ^ p ∂volume +
             ∫⁻ x in B, (ENNReal.ofReal ‖fderiv ℝ (fun x => ψ nm.1 x - ψ nm.2 x) x‖) ^ p
              ∂volume) :=
      fun nm => le_trans (hcomp_le nm) (hgrad_bound_nm nm)
    have hp0 : (0 : ℝ) < p := by linarith
    -- Convert eLpNorm = (∫⁻ |·|^p)^(1/p)
    suffices hlint : Tendsto (fun nm : ℕ × ℕ =>
        ∫⁻ x, (ENNReal.ofReal |G nm.1 x i - G nm.2 x i|) ^ p ∂volume) atTop (nhds 0) by
      -- eLpNorm(f) = (∫⁻ |f|^p)^(1/p) conversion + Tendsto rpow
      have hconv : ∀ nm : ℕ × ℕ,
          eLpNorm (fun x => G nm.1 x i - G nm.2 x i) (ENNReal.ofReal p) volume =
            (∫⁻ x, (ENNReal.ofReal |G nm.1 x i - G nm.2 x i|) ^ p ∂volume) ^ (1 / p) := by
        intro nm
        have hp_enn_ne : (ENNReal.ofReal p) ≠ 0 := (ENNReal.ofReal_pos.mpr hp0).ne'
        simpa [ENNReal.toReal_ofReal hp0.le, Real.enorm_eq_ofReal_abs] using
          (MeasureTheory.eLpNorm_eq_lintegral_rpow_enorm_toReal
            (μ := volume) (p := ENNReal.ofReal p) hp_enn_ne (by simp)
            (f := fun x => G nm.1 x i - G nm.2 x i))
      rw [show (0 : ℝ≥0∞) = (0 : ℝ≥0∞) ^ (1 / p) from by
        rw [ENNReal.zero_rpow_of_pos (by positivity : (0 : ℝ) < 1 / p)]]
      exact Filter.Tendsto.congr (fun nm => (hconv nm).symm)
        (hlint.ennrpow_const (1 / p))
    -- Auxiliary: ψ_n restricted to B is MemLp
    have hψ_memLp_restrict : ∀ n, MemLp (ψ n) (ENNReal.ofReal p) μB :=
      fun n => ((hψ_smooth n).continuous.memLp_of_hasCompactSupport (hψ_cpt n)).restrict B
    -- Function term: ∫_B |ψ nm.1 - ψ nm.2|^p → 0
    have hfun_pair : Tendsto (fun nm : ℕ × ℕ =>
        ∫⁻ x in B, (ENNReal.ofReal |ψ nm.1 x - ψ nm.2 x|) ^ p ∂volume) atTop (nhds 0) := by
      have hpair_eLpNorm : Tendsto (fun nm : ℕ × ℕ =>
          eLpNorm (fun x => ψ nm.1 x - ψ nm.2 x) (ENNReal.ofReal p) μB) atTop (nhds 0) := by
        refine tendsto_zero_of_le_pair_sum (fun nm => ?_) hψ_fun
        rw [show (fun x => ψ nm.1 x - ψ nm.2 x) =
          (fun x => (ψ nm.1 x - u x) - (ψ nm.2 x - u x)) from by ext x; ring]
        exact eLpNorm_sub_le ((hψ_memLp_restrict nm.1).sub hw.memLp).aestronglyMeasurable
          ((hψ_memLp_restrict nm.2).sub hw.memLp).aestronglyMeasurable
          (by simpa using ENNReal.ofReal_le_ofReal hp.le)
      have hconv_fun : ∀ nm : ℕ × ℕ,
          ∫⁻ x in B, (ENNReal.ofReal |ψ nm.1 x - ψ nm.2 x|) ^ p ∂volume =
            eLpNorm (fun x => ψ nm.1 x - ψ nm.2 x) (ENNReal.ofReal p) μB ^ p := by
        intro nm
        rw [← lintegral_rpow_norm_eq_eLpNorm_pow hp0]; congr 1
      simp_rw [hconv_fun]
      rw [show (0 : ℝ≥0∞) = (0 : ℝ≥0∞) ^ p from by rw [ENNReal.zero_rpow_of_pos hp0]]
      exact hpair_eLpNorm.ennrpow_const p
    -- Gradient term: ∫_B ‖∇(ψ nm.1 - ψ nm.2)‖^p → 0
    -- Strategy: ‖fderiv f x‖ = ‖(fun i => (fderiv f x)(e_i))‖_E by norm_fderiv_eq_norm_partials.
    -- For EuclideanSpace, ‖v‖^p ≤ d^(p/2) * ∑ᵢ |vᵢ|^p (norm-component bridge).
    -- Each component ∫_B |(fderiv(ψ₁-ψ₂))(e_i)|^p → 0 from hψ_grad via triangle.
    have hgrad_pair : Tendsto (fun nm : ℕ × ℕ =>
        ∫⁻ x in B, (ENNReal.ofReal ‖fderiv ℝ (fun x => ψ nm.1 x - ψ nm.2 x) x‖) ^ p
          ∂volume) atTop (nhds 0) := by
      -- Step A: Component-level pair Cauchy for each i.
      -- (fderiv(ψ₁-ψ₂))(eᵢ) = (fderiv ψ₁)(eᵢ) - (fderiv ψ₂)(eᵢ) by fderiv_sub.
      -- hψ_grad i says eLpNorm((fderiv ψ_n)(eᵢ) - weakGrad · i) → 0.
      -- By triangle: eLpNorm((fderiv ψ₁)(eᵢ) - (fderiv ψ₂)(eᵢ)) → 0 as pair.
      have hcomp_pair : ∀ j : Fin d, Tendsto (fun nm : ℕ × ℕ =>
          eLpNorm
            (fun x =>
              (fderiv ℝ (fun x => ψ nm.1 x - ψ nm.2 x) x)
                (EuclideanSpace.single j 1))
            (ENNReal.ofReal p) μB) atTop (nhds 0) := by
        intro j
        have hpair_comp_eLpNorm : Tendsto (fun nm : ℕ × ℕ =>
            eLpNorm (fun x => (fderiv ℝ (ψ nm.1) x) (EuclideanSpace.single j 1) -
              (fderiv ℝ (ψ nm.2) x) (EuclideanSpace.single j 1))
              (ENNReal.ofReal p) μB) atTop (nhds 0) := by
          refine tendsto_zero_of_le_pair_sum (fun nm => ?_) (hψ_grad j)
          rw [show (fun x => (fderiv ℝ (ψ nm.1) x) (EuclideanSpace.single j 1) -
              (fderiv ℝ (ψ nm.2) x) (EuclideanSpace.single j 1)) =
            (fun x => ((fderiv ℝ (ψ nm.1) x) (EuclideanSpace.single j 1) - hw.weakGrad x j) -
              ((fderiv ℝ (ψ nm.2) x) (EuclideanSpace.single j 1) - hw.weakGrad x j))
            from by ext x; ring]
          exact eLpNorm_sub_le
            (((((hψ_smooth nm.1).continuous_fderiv (by simp)).clm_apply
              continuous_const).memLp_of_hasCompactSupport
              ((hψ_cpt nm.1).fderiv_apply (𝕜 := ℝ) _)).restrict B |>.sub
              (hw.weakGrad_component_memLp j)).aestronglyMeasurable
            (((((hψ_smooth nm.2).continuous_fderiv (by simp)).clm_apply
              continuous_const).memLp_of_hasCompactSupport
              ((hψ_cpt nm.2).fderiv_apply (𝕜 := ℝ) _)).restrict B |>.sub
              (hw.weakGrad_component_memLp j)).aestronglyMeasurable
            (by simpa using ENNReal.ofReal_le_ofReal hp.le)
        refine (Filter.tendsto_congr fun nm => eLpNorm_congr_ae ?_).mpr hpair_comp_eLpNorm
        filter_upwards with x
        exact congrArg (fun L : E →L[ℝ] ℝ => L (EuclideanSpace.single j 1))
          (fderiv_sub ((hψ_smooth nm.1).differentiable (by simp)).differentiableAt
            ((hψ_smooth nm.2).differentiable (by simp)).differentiableAt)
      let H : (ℕ × ℕ) → E → E := fun nm x =>
        WithLp.toLp 2 fun j =>
          (fderiv ℝ (fun x => ψ nm.1 x - ψ nm.2 x) x)
            (EuclideanSpace.single j 1)
      have hH_pair : Tendsto (fun nm : ℕ × ℕ =>
          eLpNorm (H nm) (ENNReal.ofReal p) μB) atTop (nhds 0) := by
        have hle : ∀ nm : ℕ × ℕ,
            eLpNorm (H nm) (ENNReal.ofReal p) μB ≤
              ∑ j : Fin d, eLpNorm (fun x => H nm x j) (ENNReal.ofReal p) μB := by
          intro nm
          refine euclidean_eLpNorm_le_component_sum (d := d) (μ := μB) (p := p) hp.le ?_
          intro j
          have hmeas :
              Measurable (fun x =>
                (fderiv ℝ (fun x => ψ nm.1 x - ψ nm.2 x) x)
                  (EuclideanSpace.single j 1)) := by
            exact ((((hψ_smooth nm.1).sub (hψ_smooth nm.2)).continuous_fderiv
              (by simp)).clm_apply continuous_const).measurable
          simpa [H, PiLp.toLp_apply] using hmeas.aestronglyMeasurable
        have hsum : Tendsto (fun nm : ℕ × ℕ =>
            ∑ j : Fin d, eLpNorm (fun x => H nm x j) (ENNReal.ofReal p) μB)
            atTop (nhds 0) := by
          have h0 : (0 : ℝ≥0∞) = ∑ _j : Fin d, (0 : ℝ≥0∞) := by simp
          rw [h0]
          exact tendsto_finset_sum Finset.univ (fun j _ => by
            simpa [H, PiLp.toLp_apply] using hcomp_pair j)
        exact tendsto_of_tendsto_of_tendsto_of_le_of_le tendsto_const_nhds hsum
          (fun _ => bot_le) hle
      have hconv : ∀ nm : ℕ × ℕ,
          ∫⁻ x in B, (ENNReal.ofReal ‖fderiv ℝ (fun x => ψ nm.1 x - ψ nm.2 x) x‖) ^ p
            ∂volume =
            eLpNorm (H nm) (ENNReal.ofReal p) μB ^ p := by
        intro nm
        calc
          ∫⁻ x in B, (ENNReal.ofReal ‖fderiv ℝ (fun x => ψ nm.1 x - ψ nm.2 x) x‖) ^ p
              ∂volume
              = ∫⁻ x in B, (ENNReal.ofReal ‖H nm x‖) ^ p ∂volume := by
                  congr 1
                  ext x
                  simp [H, norm_fderiv_eq_norm_partials_local (d := d)]
          _ = eLpNorm (H nm) (ENNReal.ofReal p) μB ^ p := by
                simpa [μB] using
                  (lintegral_rpow_norm_eq_eLpNorm_pow (μ := μB) hp0 (f := H nm))
      simp_rw [hconv]
      rw [show (0 : ℝ≥0∞) = (0 : ℝ≥0∞) ^ p from by
        rw [ENNReal.zero_rpow_of_pos hp0]]
      exact hH_pair.ennrpow_const p
    -- Combine: bound = grad_pair + C * (fun_pair + grad_pair), all → 0
    have hC_ne_top : C_unitBallExtensionGrad d p ≠ ⊤ := by
      unfold C_unitBallExtensionGrad
      exact ENNReal.mul_ne_top ENNReal.ofReal_ne_top (ENNReal.rpow_ne_top_of_nonneg (by linarith) (by simp))
    refine tendsto_of_tendsto_of_tendsto_of_le_of_le tendsto_const_nhds ?_
      (fun _ => bot_le) hfull_le
    have hsum := hfun_pair.add hgrad_pair
    rw [add_zero] at hsum
    have hCmul := ENNReal.Tendsto.const_mul hsum (Or.inr hC_ne_top)
    rw [mul_zero] at hCmul
    have hfinal := hgrad_pair.add hCmul
    rwa [add_zero] at hfinal
  obtain ⟨Gext, hGext_comp_memLp, hG_comp_tendsto⟩ :=
    cauchy_Lp_component_limit (d := d) hp hG_comp_memLp hG_comp_cauchy
  have hweak : ∀ i : Fin d,
      HasWeakPartialDeriv i (fun x => Gext x i)
        (unitBallExtension (d := d) u) Set.univ := by
    intro i
    exact HasWeakPartialDeriv.of_eLpNormApprox_p (d := d) (Ω := Set.univ) isOpen_univ hp
      (by rw [Measure.restrict_univ]; exact huExt_memLp)
      (by rw [Measure.restrict_univ]; exact hGext_comp_memLp i)
      (fun n => hV_wd n i)
      (fun n => by rw [Measure.restrict_univ]; exact (hV_memLp n).sub huExt_memLp)
      (by rw [Measure.restrict_univ]; exact hV_fun)
      (fun n => by rw [Measure.restrict_univ]; exact (hG_comp_memLp n i).sub (hGext_comp_memLp i))
      (by rw [Measure.restrict_univ]; exact hG_comp_tendsto i)
  let hwExt : MemW1pWitness (ENNReal.ofReal p)
      (unitBallExtension (d := d) u) Set.univ :=
    { memLp := by rw [Measure.restrict_univ]; exact huExt_memLp
      weakGrad := Gext
      weakGrad_component_memLp := by
        intro i; rw [Measure.restrict_univ]; exact hGext_comp_memLp i
      isWeakGrad := hweak }
  rcases exists_global_smooth_W1p_approx_of_localizedWitness
      (d := d) (Ω := Set.univ) isOpen_univ hp hwExt
      (unitBallExtension_hasCompactSupport (d := d) u) (Set.subset_univ _) with
    ⟨Φ, hΦ_smooth, hΦ_cpt, hΦ_fun, hΦ_grad⟩
  refine ⟨Gext, huExt_memLp, hGext_comp_memLp, huExt_fun_bound, ?_,
    Φ, hΦ_smooth, hΦ_cpt, hΦ_fun, ?_⟩
  · -- Gradient energy bound via semicontinuity
    -- Each G_n satisfies: ∫ ‖G_n‖^p ≤ bound(ψ_n).
    -- G_n → Gext in Lp (componentwise), so ∫ ‖Gext‖^p ≤ lim bound(ψ_n) = bound(u).
    --        (2) component eLpNorm ≤ vector eLpNorm (pre-compiled in geometry)
    -- This uses G n → Gext componentwise in Lp, hence a.e. along a subsequence,
    -- then Fatou's lemma on ‖·‖^p (or norm-equivalence + component Fatou).
    have hFatou : (∫⁻ x, (ENNReal.ofReal ‖Gext x‖) ^ p ∂volume) ≤
        atTop.liminf (fun n => ∫⁻ x, (ENNReal.ofReal ‖G n x‖) ^ p ∂volume) := by
      have hp0 : (0 : ℝ) < p := by linarith
      have hG_aesm : ∀ n, AEStronglyMeasurable (G n) volume := by
        intro n
        exact aestronglyMeasurable_euclidean_of_components (d := d) (μ := volume)
          (fun i => (hG_comp_memLp n i).aestronglyMeasurable)
      have hGext_aesm : AEStronglyMeasurable Gext volume :=
        aestronglyMeasurable_euclidean_of_components (d := d) (μ := volume)
          (fun i => (hGext_comp_memLp i).aestronglyMeasurable)
      have hGext_memLp_vec : MemLp Gext (ENNReal.ofReal p) volume :=
        memLp_euclidean_of_components (d := d) (μ := volume) hp.le
          (fun i => hGext_comp_memLp i)
      have hG_tendsto : Tendsto
          (fun n => eLpNorm (fun x => G n x - Gext x) (ENNReal.ofReal p) volume)
          atTop (nhds 0) := by
        have hle : ∀ n,
            eLpNorm (fun x => G n x - Gext x) (ENNReal.ofReal p) volume ≤
              ∑ i : Fin d, eLpNorm (fun x => G n x i - Gext x i) (ENNReal.ofReal p) volume := by
          intro n
          exact euclidean_eLpNorm_le_component_sum (d := d) (μ := volume) (p := p) hp.le
            (fun i =>
              ((hG_comp_memLp n i).aestronglyMeasurable.sub
                (hGext_comp_memLp i).aestronglyMeasurable))
        have hsum : Tendsto (fun n =>
            ∑ i : Fin d, eLpNorm (fun x => G n x i - Gext x i) (ENNReal.ofReal p) volume)
            atTop (nhds 0) := by
          have h0 : (0 : ℝ≥0∞) = ∑ _i : Fin d, (0 : ℝ≥0∞) := by simp
          rw [h0]
          exact tendsto_finset_sum Finset.univ (fun i _ => hG_comp_tendsto i)
        exact tendsto_of_tendsto_of_tendsto_of_le_of_le tendsto_const_nhds hsum
          (fun _ => bot_le) hle
      have hGnorm_tendsto :
          Tendsto (fun n => eLpNorm (G n) (ENNReal.ofReal p) volume)
            atTop (nhds (eLpNorm Gext (ENNReal.ofReal p) volume)) := by
        have hle : ∀ n, eLpNorm (G n) (ENNReal.ofReal p) volume ≤
            eLpNorm (fun x => G n x - Gext x) (ENNReal.ofReal p) volume +
              eLpNorm Gext (ENNReal.ofReal p) volume := by
          intro n
          calc eLpNorm (G n) (ENNReal.ofReal p) volume
              = eLpNorm (fun x => (G n x - Gext x) + Gext x) (ENNReal.ofReal p) volume := by
                  congr 1
                  ext x
                  simp
            _ ≤ eLpNorm (fun x => G n x - Gext x) (ENNReal.ofReal p) volume +
                  eLpNorm Gext (ENNReal.ofReal p) volume :=
                eLpNorm_add_le ((hG_aesm n).sub hGext_aesm) hGext_aesm
                  (by simpa using ENNReal.ofReal_le_ofReal hp.le)
        have hge : ∀ n, eLpNorm Gext (ENNReal.ofReal p) volume ≤
            eLpNorm (fun x => G n x - Gext x) (ENNReal.ofReal p) volume +
              eLpNorm (G n) (ENNReal.ofReal p) volume := by
          intro n
          calc eLpNorm Gext (ENNReal.ofReal p) volume
              = eLpNorm (fun x => (Gext x - G n x) + G n x) (ENNReal.ofReal p) volume := by
                  congr 1
                  ext x
                  simp
            _ ≤ eLpNorm (fun x => Gext x - G n x) (ENNReal.ofReal p) volume +
                  eLpNorm (G n) (ENNReal.ofReal p) volume :=
                eLpNorm_add_le (hGext_aesm.sub (hG_aesm n)) (hG_aesm n)
                  (by simpa using ENNReal.ofReal_le_ofReal hp.le)
            _ = eLpNorm (fun x => G n x - Gext x) (ENNReal.ofReal p) volume +
                  eLpNorm (G n) (ENNReal.ofReal p) volume := by
                have hneg_eq :
                    eLpNorm (fun x => Gext x - G n x) (ENNReal.ofReal p) volume =
                      eLpNorm (fun x => G n x - Gext x) (ENNReal.ofReal p) volume := by
                  calc
                    eLpNorm (fun x => Gext x - G n x) (ENNReal.ofReal p) volume
                        = eLpNorm (fun x => -(G n x - Gext x)) (ENNReal.ofReal p) volume := by
                            congr 1
                            ext x
                            simp
                    _ = eLpNorm (fun x => G n x - Gext x) (ENNReal.ofReal p) volume := by
                          simpa only using
                            (eLpNorm_neg (f := fun x => G n x - Gext x)
                              (p := ENNReal.ofReal p) (μ := volume))
                simp [hneg_eq]
        refine tendsto_of_tendsto_of_tendsto_of_le_of_le ?_ ?_
          (fun n => tsub_le_iff_right.mpr (by simpa [add_comm] using hge n))
          (fun n => hle n)
        · have hsub :
              Tendsto
                (fun n =>
                  eLpNorm Gext (ENNReal.ofReal p) volume -
                    eLpNorm (fun x => G n x - Gext x) (ENNReal.ofReal p) volume)
                atTop (nhds (eLpNorm Gext (ENNReal.ofReal p) volume - 0)) := by
              simpa using
                (((ENNReal.continuous_sub_left hGext_memLp_vec.eLpNorm_ne_top).tendsto 0).comp
                  hG_tendsto)
          simpa using hsub
        · have hadd :
              Tendsto
                (fun n =>
                  eLpNorm (fun x => G n x - Gext x) (ENNReal.ofReal p) volume +
                    eLpNorm Gext (ENNReal.ofReal p) volume)
                atTop (nhds (0 + eLpNorm Gext (ENNReal.ofReal p) volume)) := by
              exact hG_tendsto.add tendsto_const_nhds
          rwa [zero_add] at hadd
      have hpow_tendsto : Tendsto
          (fun n => ∫⁻ x, (ENNReal.ofReal ‖G n x‖) ^ p ∂volume)
          atTop (nhds (∫⁻ x, (ENNReal.ofReal ‖Gext x‖) ^ p ∂volume)) := by
        have hEq : ∀ n,
            ∫⁻ x, (ENNReal.ofReal ‖G n x‖) ^ p ∂volume =
              eLpNorm (G n) (ENNReal.ofReal p) volume ^ p := by
          intro n
          simpa using (lintegral_rpow_norm_eq_eLpNorm_pow (μ := volume) hp0 (f := G n))
        have hEqLim :
            ∫⁻ x, (ENNReal.ofReal ‖Gext x‖) ^ p ∂volume =
              eLpNorm Gext (ENNReal.ofReal p) volume ^ p := by
          simpa using (lintegral_rpow_norm_eq_eLpNorm_pow (μ := volume) hp0 (f := Gext))
        simp_rw [hEq, hEqLim]
        exact hGnorm_tendsto.ennrpow_const p
      rw [hpow_tendsto.liminf_eq]
    -- Define the RHS as a function of n
    let RHS : ℕ → ℝ≥0∞ := fun n =>
      (∫⁻ x in B, (ENNReal.ofReal ‖fderiv ℝ (ψ n) x‖) ^ p ∂volume) +
        C_unitBallExtensionGrad d p *
          (∫⁻ x in B, (ENNReal.ofReal |ψ n x|) ^ p ∂volume +
           ∫⁻ x in B, (ENNReal.ofReal ‖fderiv ℝ (ψ n) x‖) ^ p ∂volume)
    let RHS_lim : ℝ≥0∞ :=
      (∫⁻ x in B, (ENNReal.ofReal ‖hw.weakGrad x‖) ^ p ∂volume) +
        C_unitBallExtensionGrad d p *
          (∫⁻ x in B, (ENNReal.ofReal |u x|) ^ p ∂volume +
           ∫⁻ x in B, (ENNReal.ofReal ‖hw.weakGrad x‖) ^ p ∂volume)
    -- bound(ψ n) → bound(u)
    have hRHS_tendsto : Tendsto RHS atTop (nhds RHS_lim) := by
      -- Abbreviations
      let Fn : ℕ → ℝ≥0∞ := fun n =>
        ∫⁻ x in B, (ENNReal.ofReal |ψ n x|) ^ p ∂volume
      let An : ℕ → ℝ≥0∞ := fun n =>
        ∫⁻ x in B, (ENNReal.ofReal ‖fderiv ℝ (ψ n) x‖) ^ p ∂volume
      let F_lim : ℝ≥0∞ := ∫⁻ x in B, (ENNReal.ofReal |u x|) ^ p ∂volume
      let A_lim : ℝ≥0∞ := ∫⁻ x in B, (ENNReal.ofReal ‖hw.weakGrad x‖) ^ p ∂volume
      have hp0 : (0 : ℝ) < p := by linarith
      have h1p : 1 ≤ ENNReal.ofReal p := by simpa using ENNReal.ofReal_le_ofReal hp.le
      -- Reverse triangle: eLpNorm(ψ n - u) → 0 ⟹ eLpNorm(ψ n) → eLpNorm(u).
      have hψ_memLp_B : ∀ n, MemLp (ψ n) (ENNReal.ofReal p) μB :=
        fun n => ((hψ_smooth n).continuous.memLp_of_hasCompactSupport (hψ_cpt n)).restrict B
      have heLpNorm_fun : Tendsto (fun n => eLpNorm (ψ n) (ENNReal.ofReal p) μB)
          atTop (nhds (eLpNorm u (ENNReal.ofReal p) μB)) := by
        -- Upper: eLpNorm(ψ n) ≤ eLpNorm(ψ n - u) + eLpNorm(u)
        have hle : ∀ n, eLpNorm (ψ n) (ENNReal.ofReal p) μB ≤
            eLpNorm (fun x => ψ n x - u x) (ENNReal.ofReal p) μB +
              eLpNorm u (ENNReal.ofReal p) μB := by
          intro n
          calc eLpNorm (ψ n) (ENNReal.ofReal p) μB
              = eLpNorm (fun x => (ψ n x - u x) + u x) (ENNReal.ofReal p) μB := by
                congr 1; ext x; simp
            _ ≤ eLpNorm (fun x => ψ n x - u x) (ENNReal.ofReal p) μB +
                  eLpNorm u (ENNReal.ofReal p) μB :=
                eLpNorm_add_le ((hψ_memLp_B n).sub hw.memLp).aestronglyMeasurable
                  hw.memLp.aestronglyMeasurable h1p
        -- Lower: eLpNorm(u) ≤ eLpNorm(ψ n - u) + eLpNorm(ψ n)
        have hge : ∀ n, eLpNorm u (ENNReal.ofReal p) μB ≤
            eLpNorm (fun x => ψ n x - u x) (ENNReal.ofReal p) μB +
              eLpNorm (ψ n) (ENNReal.ofReal p) μB := by
          intro n
          calc eLpNorm u (ENNReal.ofReal p) μB
              = eLpNorm (fun x => (u x - ψ n x) + ψ n x) (ENNReal.ofReal p) μB := by
                congr 1; ext x; simp
            _ ≤ eLpNorm (fun x => u x - ψ n x) (ENNReal.ofReal p) μB +
                  eLpNorm (ψ n) (ENNReal.ofReal p) μB :=
                eLpNorm_add_le (hw.memLp.sub (hψ_memLp_B n)).aestronglyMeasurable
                  (hψ_memLp_B n).aestronglyMeasurable h1p
            _ = eLpNorm (fun x => ψ n x - u x) (ENNReal.ofReal p) μB +
                  eLpNorm (ψ n) (ENNReal.ofReal p) μB := by
                have hneg_eq :
                    eLpNorm (fun x => u x - ψ n x) (ENNReal.ofReal p) μB =
                      eLpNorm (fun x => ψ n x - u x) (ENNReal.ofReal p) μB := by
                  calc
                    eLpNorm (fun x => u x - ψ n x) (ENNReal.ofReal p) μB
                        = eLpNorm (fun x => -(ψ n x - u x)) (ENNReal.ofReal p) μB := by
                            congr 1
                            ext x
                            ring
                    _ = eLpNorm (fun x => ψ n x - u x) (ENNReal.ofReal p) μB := by
                          simpa only using
                            (eLpNorm_neg (f := fun x => ψ n x - u x)
                              (p := ENNReal.ofReal p) (μ := μB))
                simp [hneg_eq]
        -- Squeeze: eLpNorm(u) - eLpNorm(ψ n - u) ≤ eLpNorm(ψ n) ≤ eLpNorm(ψ n - u) + eLpNorm(u)
        refine tendsto_of_tendsto_of_tendsto_of_le_of_le ?_ ?_
          (fun n => by
            have hge' : eLpNorm u (ENNReal.ofReal p) μB ≤
                eLpNorm (ψ n) (ENNReal.ofReal p) μB +
                  eLpNorm (fun x => ψ n x - u x) (ENNReal.ofReal p) μB := by
              simpa [add_comm] using hge n
            exact tsub_le_iff_right.mpr hge')
          (fun n => hle n)
        · have hsub :
              Tendsto
                (fun n =>
                  eLpNorm u (ENNReal.ofReal p) μB -
                    eLpNorm (fun x => ψ n x - u x) (ENNReal.ofReal p) μB)
                atTop (nhds (eLpNorm u (ENNReal.ofReal p) μB - 0)) := by
              simpa using
                (((ENNReal.continuous_sub_left hw.memLp.eLpNorm_ne_top).tendsto 0).comp hψ_fun)
          simpa using hsub
        · have hadd :
              Tendsto
                (fun n =>
                  eLpNorm (fun x => ψ n x - u x) (ENNReal.ofReal p) μB +
                    eLpNorm u (ENNReal.ofReal p) μB)
                atTop (nhds (0 + eLpNorm u (ENNReal.ofReal p) μB)) := by
              exact hψ_fun.add tendsto_const_nhds
          rwa [zero_add] at hadd
      -- Fn n = eLpNorm(ψ n)^p and F_lim = eLpNorm(u)^p
      have hFn_eq : ∀ n, Fn n = eLpNorm (ψ n) (ENNReal.ofReal p) μB ^ p := by
        intro n; rw [← lintegral_rpow_norm_eq_eLpNorm_pow hp0]; congr 1
      have hF_lim_eq : F_lim = eLpNorm u (ENNReal.ofReal p) μB ^ p := by
        rw [← lintegral_rpow_norm_eq_eLpNorm_pow hp0]; congr 1
      have hFn_tendsto : Tendsto Fn atTop (nhds F_lim) := by
        show Tendsto (fun n => Fn n) atTop (nhds F_lim)
        simp_rw [hFn_eq, hF_lim_eq]
        exact heLpNorm_fun.ennrpow_const p
      -- Component-level: eLpNorm(fderiv(ψ n)(eⱼ) - weakGrad · j) → 0 for each j.
      -- By reverse triangle, eLpNorm(fderiv(ψ n)(eⱼ)) → eLpNorm(weakGrad · j),
      -- hence ∫|comp_j(ψ n)|^p → ∫|comp_j(weakGrad)|^p for each j.
      -- Bridge: ∫‖fderiv(ψ n)‖^p ≈ ∑ⱼ ∫|comp_j(ψ n)|^p up to d^{p/2} (EuclideanSpace).
      -- Since all component integrals converge, the sandwich gives An n → A_lim.
      have hAn_tendsto : Tendsto An atTop (nhds A_lim) := by
        let gradVec : ℕ → E → E := fun n x =>
          WithLp.toLp 2 fun i =>
            (fderiv ℝ (ψ n) x) (EuclideanSpace.single i 1)
        have hgradVec_comp_aesm : ∀ n i, AEStronglyMeasurable (fun x => gradVec n x i) μB := by
          intro n i
          have hmeas :
              Measurable (fun x => (fderiv ℝ (ψ n) x) (EuclideanSpace.single i 1)) := by
            exact (((hψ_smooth n).continuous_fderiv (by simp)).clm_apply continuous_const).measurable
          simpa [gradVec, PiLp.toLp_apply] using hmeas.aestronglyMeasurable
        have hgradVec_aesm : ∀ n, AEStronglyMeasurable (gradVec n) μB := by
          intro n
          exact aestronglyMeasurable_euclidean_of_components (d := d) (μ := μB)
            (fun i => hgradVec_comp_aesm n i)
        have hweakGrad_aesm : AEStronglyMeasurable hw.weakGrad μB :=
          aestronglyMeasurable_euclidean_of_components (d := d) (μ := μB)
            (fun i => (hw.weakGrad_component_memLp i).aestronglyMeasurable)
        have hweakGrad_memLp_vec : MemLp hw.weakGrad (ENNReal.ofReal p) μB :=
          memLp_euclidean_of_components (d := d) (μ := μB) hp.le
            (fun i => hw.weakGrad_component_memLp i)
        have hgradVec_tendsto : Tendsto
            (fun n => eLpNorm (fun x => gradVec n x - hw.weakGrad x) (ENNReal.ofReal p) μB)
            atTop (nhds 0) := by
          have hle : ∀ n,
              eLpNorm (fun x => gradVec n x - hw.weakGrad x) (ENNReal.ofReal p) μB ≤
                ∑ i : Fin d,
                  eLpNorm (fun x => gradVec n x i - hw.weakGrad x i) (ENNReal.ofReal p) μB := by
            intro n
            exact euclidean_eLpNorm_le_component_sum (d := d) (μ := μB) (p := p) hp.le
              (fun i => (hgradVec_comp_aesm n i).sub
                (hw.weakGrad_component_memLp i).aestronglyMeasurable)
          have hsum : Tendsto (fun n =>
              ∑ i : Fin d,
                eLpNorm (fun x => gradVec n x i - hw.weakGrad x i) (ENNReal.ofReal p) μB)
              atTop (nhds 0) := by
            have h0 : (0 : ℝ≥0∞) = ∑ _i : Fin d, (0 : ℝ≥0∞) := by simp
            rw [h0]
            exact tendsto_finset_sum Finset.univ (fun i _ => by
              simpa [gradVec, PiLp.toLp_apply] using hψ_grad i)
          exact tendsto_of_tendsto_of_tendsto_of_le_of_le tendsto_const_nhds hsum
            (fun _ => bot_le) hle
        have hgradVec_norm_tendsto :
            Tendsto (fun n => eLpNorm (gradVec n) (ENNReal.ofReal p) μB)
              atTop (nhds (eLpNorm hw.weakGrad (ENNReal.ofReal p) μB)) := by
          have hle : ∀ n, eLpNorm (gradVec n) (ENNReal.ofReal p) μB ≤
              eLpNorm (fun x => gradVec n x - hw.weakGrad x) (ENNReal.ofReal p) μB +
                eLpNorm hw.weakGrad (ENNReal.ofReal p) μB := by
            intro n
            calc eLpNorm (gradVec n) (ENNReal.ofReal p) μB
                = eLpNorm (fun x => (gradVec n x - hw.weakGrad x) + hw.weakGrad x)
                    (ENNReal.ofReal p) μB := by
                      congr 1
                      ext x
                      simp
              _ ≤ eLpNorm (fun x => gradVec n x - hw.weakGrad x) (ENNReal.ofReal p) μB +
                    eLpNorm hw.weakGrad (ENNReal.ofReal p) μB :=
                  eLpNorm_add_le ((hgradVec_aesm n).sub hweakGrad_aesm) hweakGrad_aesm h1p
          have hge : ∀ n, eLpNorm hw.weakGrad (ENNReal.ofReal p) μB ≤
              eLpNorm (fun x => gradVec n x - hw.weakGrad x) (ENNReal.ofReal p) μB +
                eLpNorm (gradVec n) (ENNReal.ofReal p) μB := by
            intro n
            calc eLpNorm hw.weakGrad (ENNReal.ofReal p) μB
                = eLpNorm (fun x => (hw.weakGrad x - gradVec n x) + gradVec n x)
                    (ENNReal.ofReal p) μB := by
                      congr 1
                      ext x
                      simp
              _ ≤ eLpNorm (fun x => hw.weakGrad x - gradVec n x) (ENNReal.ofReal p) μB +
                    eLpNorm (gradVec n) (ENNReal.ofReal p) μB :=
                  eLpNorm_add_le (hweakGrad_aesm.sub (hgradVec_aesm n)) (hgradVec_aesm n) h1p
              _ = eLpNorm (fun x => gradVec n x - hw.weakGrad x) (ENNReal.ofReal p) μB +
                    eLpNorm (gradVec n) (ENNReal.ofReal p) μB := by
                  have hneg_eq :
                      eLpNorm (fun x => hw.weakGrad x - gradVec n x) (ENNReal.ofReal p) μB =
                        eLpNorm (fun x => gradVec n x - hw.weakGrad x) (ENNReal.ofReal p) μB := by
                    calc
                      eLpNorm (fun x => hw.weakGrad x - gradVec n x) (ENNReal.ofReal p) μB
                          = eLpNorm (fun x => -(gradVec n x - hw.weakGrad x)) (ENNReal.ofReal p) μB := by
                              congr 1
                              ext x
                              simp
                      _ = eLpNorm (fun x => gradVec n x - hw.weakGrad x) (ENNReal.ofReal p) μB := by
                            simpa only using
                              (eLpNorm_neg (f := fun x => gradVec n x - hw.weakGrad x)
                                (p := ENNReal.ofReal p) (μ := μB))
                  simp [hneg_eq]
          refine tendsto_of_tendsto_of_tendsto_of_le_of_le ?_ ?_
            (fun n => by
              have hge' : eLpNorm hw.weakGrad (ENNReal.ofReal p) μB ≤
                  eLpNorm (gradVec n) (ENNReal.ofReal p) μB +
                    eLpNorm (fun x => gradVec n x - hw.weakGrad x) (ENNReal.ofReal p) μB := by
                simpa [add_comm] using hge n
              exact tsub_le_iff_right.mpr hge')
            (fun n => hle n)
          · have hsub :
                Tendsto
                  (fun n =>
                    eLpNorm hw.weakGrad (ENNReal.ofReal p) μB -
                      eLpNorm (fun x => gradVec n x - hw.weakGrad x) (ENNReal.ofReal p) μB)
                  atTop (nhds (eLpNorm hw.weakGrad (ENNReal.ofReal p) μB - 0)) := by
                simpa using
                  (((ENNReal.continuous_sub_left hweakGrad_memLp_vec.eLpNorm_ne_top).tendsto 0).comp
                    hgradVec_tendsto)
            simpa using hsub
          · have hadd :
                Tendsto
                  (fun n =>
                    eLpNorm (fun x => gradVec n x - hw.weakGrad x) (ENNReal.ofReal p) μB +
                      eLpNorm hw.weakGrad (ENNReal.ofReal p) μB)
                  atTop (nhds (0 + eLpNorm hw.weakGrad (ENNReal.ofReal p) μB)) := by
                exact hgradVec_tendsto.add tendsto_const_nhds
            rwa [zero_add] at hadd
        have hAn_eq : ∀ n, An n = eLpNorm (gradVec n) (ENNReal.ofReal p) μB ^ p := by
          intro n
          change
            (∫⁻ x in B, (ENNReal.ofReal ‖fderiv ℝ (ψ n) x‖) ^ p ∂volume) =
              eLpNorm (gradVec n) (ENNReal.ofReal p) μB ^ p
          calc
            ∫⁻ x in B, (ENNReal.ofReal ‖fderiv ℝ (ψ n) x‖) ^ p ∂volume
                = ∫⁻ x in B, (ENNReal.ofReal ‖gradVec n x‖) ^ p ∂volume := by
                    congr 1
                    ext x
                    simp [gradVec, norm_fderiv_eq_norm_partials_local (d := d)]
            _ = eLpNorm (gradVec n) (ENNReal.ofReal p) μB ^ p := by
                  simpa [μB] using
                    (lintegral_rpow_norm_eq_eLpNorm_pow (μ := μB) hp0 (f := gradVec n))
        have hA_lim_eq : A_lim = eLpNorm hw.weakGrad (ENNReal.ofReal p) μB ^ p := by
          change
            (∫⁻ x in B, (ENNReal.ofReal ‖hw.weakGrad x‖) ^ p ∂volume) =
              eLpNorm hw.weakGrad (ENNReal.ofReal p) μB ^ p
          simpa [μB] using
            (lintegral_rpow_norm_eq_eLpNorm_pow (μ := μB) hp0 (f := hw.weakGrad))
        rw [hA_lim_eq]
        convert hgradVec_norm_tendsto.ennrpow_const p using 1
        ext n
        exact hAn_eq n
      have hC_ne_top : C_unitBallExtensionGrad d p ≠ ⊤ := by
        unfold C_unitBallExtensionGrad
        exact ENNReal.mul_ne_top ENNReal.ofReal_ne_top (ENNReal.rpow_ne_top_of_nonneg (by linarith) (by simp))
      have hsum := hFn_tendsto.add hAn_tendsto
      have hCmul := ENNReal.Tendsto.const_mul hsum (Or.inr hC_ne_top)
      exact hAn_tendsto.add hCmul
    have hliminf_le : atTop.liminf (fun n => ∫⁻ x, (ENNReal.ofReal ‖G n x‖) ^ p ∂volume)
        ≤ RHS_lim := by
      calc atTop.liminf (fun n => ∫⁻ x, (ENNReal.ofReal ‖G n x‖) ^ p ∂volume)
          ≤ atTop.liminf RHS := by
            exact Filter.liminf_le_liminf (Eventually.of_forall hG_bound)
              (⟨0, Eventually.of_forall fun _ => bot_le⟩)
              (⟨⊤, fun _ _ => le_top⟩)
        _ = RHS_lim := hRHS_tendsto.liminf_eq
    exact le_trans hFatou hliminf_le
  · -- Gradient convergence: convert Set.univ.indicator to identity
    intro i
    refine (Filter.tendsto_congr fun n => eLpNorm_congr_ae ?_).mpr (hΦ_grad i)
    filter_upwards with x
    exact by
      simp [hwExt]

\end{lstlisting}
\begin{proof}
Apply Theorem~\ref{thm:wp-approx} to obtain $\psi_n \in C^\infty_c(E)$ with
$\psi_n \to u$ and $\partial_i \psi_n \to (\nabla u)_i$ in $L^p(\Bz{1})$.
For each $\psi_n$, apply Theorem~\ref{thm:smooth-input-smoothing}, producing
a smooth extension sequence $\Phi_{n,m} \in C^\infty_c(E)$ approximating
$\Eext \psi_n$ in $W^{1,p}(E)$, and the smooth-input estimates of
Theorem~\ref{thm:smooth-extension-estimate}.

The smooth-input bound gives, by linearity of $\Eext$,
\begin{gather*}
\eLpNormText{\Eext \psi_n - \Eext \psi_m}_{L^p(E)}
\le (1 + 2^{2d})^{1/p}\,\eLpNormText{\psi_n - \psi_m}_{L^p(\Bz{1})},
\end{gather*}
\begin{multline*}
\eLpNormText{D\Eext \psi_n - D\Eext\psi_m}_{L^p(E)}
\le \eLpNormText{D\psi_n - D\psi_m}_{L^p(\Bz{1})} \\
+ C\,\bigl(\eLpNormText{\psi_n - \psi_m}_{L^p(\Bz{1})}
+ \eLpNormText{D\psi_n - D\psi_m}_{L^p(\Bz{1})}\bigr).
\end{multline*}
Since $(\psi_n)$ is Cauchy in $W^{1,p}(\Bz{1})$, the smooth-extension
sequence $\Eext\psi_n$ is Cauchy in $L^p(E)$, and the candidate gradient
sequence $\smoothUnitBallExtensionGrad \psi_n$ is Cauchy in
$L^p(E;\R^d)$. By the bare-function $L^p$ completeness theorems
(Theorems~\leanname{scalar\_cauchy\_to\_limit} and
\leanname{exists\_pi\_limit\_of\_cauchy\_eLpNorm}), there exist $L^p$ limits $f^*$
and $G_{\mathrm{ext}}$. By construction $f^* = \Eext u$ a.e.\ and the
estimates pass to the limit by lower semicontinuity of the $L^p$-norm under
weak/strong convergence and the smooth-input estimates applied to each
$\psi_n$. Finally, a diagonal extraction from $(\Phi_{n,m})$ produces the
sequence required in (4).
\end{proof}

\begin{lemma}[\leanname{aestronglyMeasurable\_unitBallExtension\_of\_memLp}]
If $u \in L^p(\Bz{1})$ then $\Eext u$ is a.e.\ strongly measurable on $E$.
\end{lemma}

\begin{lstlisting}[style=Matlab-editor,basicstyle=\mlttfamily,escapechar=",mlshowsectionrules=true,mathescape=true,morecomment={[s]{\%\{}{\%\}}},language=lean,numbers=none]
/-- AEStronglyMeasurable for unitBallExtension of rough u.
Uses measurable representative + ae_eq transfer. -/
theorem aestronglyMeasurable_unitBallExtension_of_memLp
    {p : ℝ≥0∞} {u : E → ℝ}
    (hu : MemLp u p (volume.restrict (Metric.ball (0 : E) 1))) :
    AEStronglyMeasurable (unitBallExtension (d := d) u) volume := by
  -- Get measurable representative
  let u' := hu.aestronglyMeasurable.mk u
  have hu'_meas : Measurable u' := hu.aestronglyMeasurable.stronglyMeasurable_mk.measurable
  have hu'_ae : u =ᵐ[volume.restrict (Metric.ball (0 : E) 1)] u' :=
    hu.aestronglyMeasurable.ae_eq_mk
  -- unitBallExtension u' is measurable
  have hext_meas : Measurable (unitBallExtension (d := d) u') :=
    measurable_unitBallExtension (d := d) hu'_meas
  -- unitBallExtension u =ᵐ unitBallExtension u'
  -- because the retraction maps into closedBall(0,1) and u =ᵐ u' on ball(0,1)
  have hae_ext : unitBallExtension (d := d) u =ᵐ[volume] unitBallExtension (d := d) u' := by
    -- On the annulus {1 < ‖x‖ < 2}, retraction is inversion x ↦ x/‖x‖² which is
    -- a smooth diffeomorphism onto {1/2 < ‖y‖ < 1}. Preimage of null set under
    -- smooth diffeomorphism is null. Combined with: ball case (retract = id),
    -- sphere case (measure 0), and {‖x‖ ≥ 2} case (cutoff = 0).
    have hN : ∀ᵐ x ∂(volume : Measure E), x ∈ Metric.ball (0 : E) 1 → u x = u' x := by
      rwa [Filter.EventuallyEq, ae_restrict_iff' measurableSet_ball] at hu'_ae
    have hSph : ∀ᵐ x ∂(volume : Measure E), x ∉ Metric.sphere (0 : E) 1 := by
      have h0 : (volume : Measure E) (Metric.sphere (0 : E) 1) = 0 :=
        MeasureTheory.Measure.addHaar_sphere _ _ _
      exact ae_iff.mpr (by convert h0 using 1; simp only [not_not, Set.setOf_mem_eq])
    -- This follows from inversion being a smooth diffeomorphism (Lipschitz on compact subsets).
    have hAnn : ∀ᵐ x ∂(volume : Measure E),
        1 < ‖x‖ → ‖x‖ < 2 →
          u (unitBallRetraction (d := d) x) = u' (unitBallRetraction (d := d) x) := by
      let badInner : Set E := {y | y ∈ unitBallInnerShell (d := d) ∧ u y ≠ u' y}
      have hbadInner_ae : ∀ᵐ y ∂(volume : Measure E), y ∉ badInner := by
        filter_upwards [hN] with y hy
        intro hy_bad
        exact hy_bad.2 (hy <| by
          rcases hy_bad.1 with ⟨_hy_half, hy_lt_one⟩
          exact Metric.mem_ball.mpr (by rwa [dist_zero_right]))
      have hbadInner_zero : (volume : Measure E) badInner = 0 := by
        simpa [badInner] using (ae_iff.mp hbadInner_ae)
      have hdiff_badInner :
          DifferentiableOn ℝ (EuclideanGeometry.inversion (0 : E) 1) badInner := by
        intro y hy
        have hy0 : y ≠ (0 : E) := by
          intro hy0
          rcases hy.1 with ⟨hy_half, _hy_lt_one⟩
          have : (0 : ℝ) < ‖y‖ := by linarith
          simp [hy0] at this
        have hInv :=
          EuclideanGeometry.hasFDerivAt_inversion (c := (0 : E)) (R := (1 : ℝ)) hy0
        exact hInv.differentiableAt.differentiableWithinAt
      have himage_badInner_zero :
          (volume : Measure E) (EuclideanGeometry.inversion (0 : E) 1 '' badInner) = 0 := by
        exact addHaar_image_eq_zero_of_differentiableOn_of_addHaar_eq_zero
          (μ := (volume : Measure E)) hdiff_badInner hbadInner_zero
      let badOuter : Set E := {x |
        x ∈ unitBallOuterShell (d := d) ∧
        u (unitBallRetraction (d := d) x) ≠ u' (unitBallRetraction (d := d) x)}
      have hbadOuter_subset :
          badOuter ⊆ EuclideanGeometry.inversion (0 : E) 1 '' badInner := by
        intro x hx
        rcases hx with ⟨hx_shell, hx_bad⟩
        refine ⟨EuclideanGeometry.inversion (0 : E) 1 x, ?_, ?_⟩
        · refine ⟨inversion_mem_unitBallInnerShell_of_mem_outerShell (d := d) hx_shell, ?_⟩
          simpa [unitBallRetraction_eq_inversion_of_mem_outerShell (d := d) hx_shell] using hx_bad
        · exact EuclideanGeometry.inversion_inversion (c := (0 : E)) (R := (1 : ℝ)) one_ne_zero x
      have hbadOuter_zero : (volume : Measure E) badOuter = 0 := by
        exact measure_mono_null hbadOuter_subset himage_badInner_zero
      have hbadOuter_ae : ∀ᵐ x ∂(volume : Measure E), x ∉ badOuter := by
        exact ae_iff.mpr (by simpa [badOuter, Set.setOf_mem_eq] using hbadOuter_zero)
      filter_upwards [hbadOuter_ae] with x hxBad
      intro hx1 hx2
      by_contra hneq
      exact hxBad ⟨⟨hx1, hx2⟩, hneq⟩
    filter_upwards [hN, hSph, hAnn] with x hx_ball hx_sph hx_ann
    simp only [unitBallExtension]
    by_cases h2 : 2 ≤ ‖x‖
    · simp [unitBallCutoff_eq_zero_of_two_le_norm (d := d) h2]
    · push_neg at h2
      rw [Metric.mem_sphere, dist_zero_right] at hx_sph
      rcases lt_or_gt_of_ne hx_sph with h1 | h1
      · -- ‖x‖ < 1: retraction = identity, x ∈ ball(0,1)
        rw [unitBallRetraction_eq_self_of_norm_le_one (d := d) h1.le]
        congr 1; exact hx_ball (Metric.mem_ball.mpr (by rwa [dist_zero_right]))
      · -- 1 < ‖x‖ < 2: use the annulus lemma
        congr 1; exact hx_ann h1 h2
  exact hext_meas.aestronglyMeasurable.congr hae_ext.symm

\end{lstlisting}
\begin{proof}
Choose a strongly measurable representative $u'$ of $u$. Then $\Eext u'$ is
measurable: it is a composition of measurable maps. Now $\Eext u =
\Eext u'$ a.e.\ on $E$, by checking three cases.
\begin{itemize}
\item On $\Bz{1}$, $\Eext = \mathrm{id}$, so equality reduces to $u = u'$
a.e.
\item On $\|x\| \ge 2$, both vanish.
\item On $\Sigma^{\mathrm{out}}$, $\Eext u(x) = \chi(x)\, u(I(x))$. The
inversion $I$ is a smooth diffeomorphism between $\Sigma^{\mathrm{out}}$ and
$\Sigma^{\mathrm{in}}$. The set where $u \neq u'$ on $\Sigma^{\mathrm{in}}$
has measure zero, and its image under $I$ also has measure zero by
\leanname{addHaar\_image\_eq\_zero\_of\_differentiableOn\_of\_addHaar\_eq\_zero}.
\end{itemize}
The sphere $\mathrm{S}_1$ has measure zero, hence does not affect the a.e.
identity. Strong measurability of $\Eext u$ then follows from a.e. equality
to a measurable function.
\end{proof}

\begin{theorem}[\leanname{exists\_unitBall\_W1p\_extension}]
\label{thm:Eext-W1p}
Let $1 < p < \infty$ and $u \in W^{1,p}(\Bz{1})$. There is a $W^{1,p}$
witness $\hwExt$ for $\Eext u$ on $E$ with the following properties:
\begin{enumerate}
\item $\Eext u$ has compact support in $\overline{\Ball{2}{0}}$;
\item $\Eext u(x) = u(x)$ for $x \in \Bz{1}$;
\item the function-side bound from Theorem~\ref{thm:global-approx-extension} (2);
\item the gradient-side bound from Theorem~\ref{thm:global-approx-extension} (3).
\end{enumerate}
\end{theorem}

\begin{lstlisting}[style=Matlab-editor,basicstyle=\mlttfamily,escapechar=",mlshowsectionrules=true,mathescape=true,morecomment={[s]{\%\{}{\%\}}},language=lean,numbers=none]
theorem exists_unitBall_W1p_extension
    {p : ℝ} (hp : 1 < p) {u : E → ℝ}
    (hw : MemW1pWitness (ENNReal.ofReal p) u (Metric.ball (0 : E) 1)) :
    ∃ hwExt : MemW1pWitness (ENNReal.ofReal p) (unitBallExtension (d := d) u) Set.univ,
      HasCompactSupport (unitBallExtension (d := d) u) ∧
      (∀ x ∈ Metric.ball (0 : E) 1, unitBallExtension (d := d) u x = u x) ∧
      (∫⁻ x, (ENNReal.ofReal |unitBallExtension (d := d) u x|) ^ p ∂volume)
        ≤ C_unitBallExtensionFun d *
          ∫⁻ x in Metric.ball (0 : E) 1, (ENNReal.ofReal |u x|) ^ p ∂volume ∧
      (∫⁻ x, (ENNReal.ofReal ‖hwExt.weakGrad x‖) ^ p ∂volume)
        ≤ (∫⁻ x in Metric.ball (0 : E) 1, (ENNReal.ofReal ‖hw.weakGrad x‖) ^ p ∂volume) +
          C_unitBallExtensionGrad d p *
            (∫⁻ x in Metric.ball (0 : E) 1, (ENNReal.ofReal |u x|) ^ p ∂volume +
             ∫⁻ x in Metric.ball (0 : E) 1, (ENNReal.ofReal ‖hw.weakGrad x‖) ^ p ∂volume) := by
  -- Build the witness directly from exists_smooth_global_approx (single Classical.choose)
  rcases exists_smooth_global_approx_of_unitBallExtension (d := d) hp hw with
    ⟨Gext, huExt_memLp, hGext_comp_memLp, hfun_bound, hgrad_bound,
     Φ, hΦ_smooth, hΦ_cpt, hΦ_fun, hΦ_grad⟩
  let hwExt : MemW1pWitness (ENNReal.ofReal p) (unitBallExtension (d := d) u) Set.univ :=
    { memLp := by rw [Measure.restrict_univ]; exact huExt_memLp
      weakGrad := Gext
      weakGrad_component_memLp := by
        intro i; rw [Measure.restrict_univ]; exact hGext_comp_memLp i
      isWeakGrad :=
        hasWeakPartials_of_global_smoothApprox (d := d) hp
          huExt_memLp hGext_comp_memLp hΦ_smooth hΦ_cpt hΦ_fun hΦ_grad }
  exact ⟨hwExt,
    unitBallExtension_hasCompactSupport (d := d) u,
    fun x hx => unitBallExtension_eq_of_mem_ball (d := d) hx,
    hfun_bound, hgrad_bound⟩

\end{lstlisting}
\begin{proof}
Combine Theorem~\ref{thm:global-approx-extension} (existence of
$G_{\mathrm{ext}}$ together with the approximating sequence) with the previous
two lemmas. The witness package follows from
\leanname{hasWeakPartials\_of\_global\_smoothApprox}, the support inclusion from
\leanname{unitBallExtension\_hasCompactSupport}, and the agreement on $\Bz{1}$
from \leanname{unitBallExtension\_eq\_of\_mem\_ball}.
\end{proof}

\begin{theorem}[\leanname{exists\_unitBall\_W1p\_extension'}]
A weakened formulation: for $u \in W^{1,p}(\Bz{1})$, there exists $u_{\mathrm{ext}}
\in W^{1,p}(E)$ with compact support and $u_{\mathrm{ext}} = u$ on $\Bz{1}$.
\end{theorem}

\begin{lstlisting}[style=Matlab-editor,basicstyle=\mlttfamily,escapechar=",mlshowsectionrules=true,mathescape=true,morecomment={[s]{\%\{}{\%\}}},language=lean,numbers=none]
theorem exists_unitBall_W1p_extension'
    {p : ℝ} (hp : 1 < p) {u : E → ℝ}
    (hw : MemW1pWitness (ENNReal.ofReal p) u (Metric.ball (0 : E) 1)) :
    ∃ uExt : E → ℝ, ∃ _ : MemW1pWitness (ENNReal.ofReal p) uExt Set.univ,
      HasCompactSupport uExt ∧
      (∀ x ∈ Metric.ball (0 : E) 1, uExt x = u x) := by
  rcases exists_unitBall_W1p_extension (d := d) hp hw with ⟨hwExt, hcpt, heq, _, _⟩
  exact ⟨_, hwExt, hcpt, heq⟩


\end{lstlisting}
\begin{proof}
Project away the explicit constants from Theorem~\ref{thm:Eext-W1p}.
\end{proof}

\subsection{Positive-part machinery}

\begin{theorem}[\leanname{MemW1pWitness.weakGrad\_ae\_eq\_zero\_on\_zeroSet}]
Let $u \in W^{1,2}(\Omega)$ on an open set $\Omega \subseteq E$. Then for
each $i$, $(\nabla u)_i = 0$ a.e.\ on the zero set $\{u = 0\}\cap \Omega$.
\end{theorem}

\begin{lstlisting}[style=Matlab-editor,basicstyle=\mlttfamily,escapechar=",mlshowsectionrules=true,mathescape=true,morecomment={[s]{\%\{}{\%\}}},language=lean,numbers=none]
/-- Stampacchia's zero-set theorem for Sobolev witnesses: each weak-gradient
component vanishes a.e. on the zero set of the function. -/
theorem MemW1pWitness.weakGrad_ae_eq_zero_on_zeroSet
    {Ω : Set E} (hΩ : IsOpen Ω) {u : E → ℝ}
    (hw : MemW1pWitness 2 u Ω) :
    ∀ i : Fin d, ∀ᵐ x ∂(volume.restrict Ω), u x = 0 → hw.weakGrad x i = 0 := by
  intro i
  let uMk : E → ℝ := AEMeasurable.mk u hw.memLp.aemeasurable
  have hu_meas : Measurable uMk := hw.memLp.aemeasurable.measurable_mk
  have hu_ae : u =ᵐ[volume.restrict Ω] uMk := hw.memLp.aemeasurable.ae_eq_mk
  have hGMk :
      ∀ j : Fin d, HasWeakPartialDeriv' (d := d) j (fun x => hw.weakGrad x j) uMk Ω := by
    intro j
    exact stampacchia_congr_ae
      (HasWeakPartialDeriv.toStampacchia (hw.isWeakGrad j)) hu_ae
  have huMk_memW1p : MemW1p 2 uMk Ω := by
    refine ⟨hw.memLp.ae_eq hu_ae, ?_⟩
    intro j
    exact ⟨fun x => hw.weakGrad x j, hw.weakGrad_component_memLp j, by
      simpa [HasWeakPartialDeriv'] using hGMk j⟩
  have hzero_mk :
      ∀ᵐ x ∂(volume.restrict Ω), uMk x = 0 → hw.weakGrad x i = 0 :=
    DeGiorgi.weakGrad_ae_eq_zero_on_zeroSet hΩ huMk_memW1p hGMk i
      (hw.weakGrad_component_memLp i) hu_meas
  filter_upwards [hu_ae, hzero_mk] with x hx hzero hx_zero
  exact hzero (by simpa [hx] using hx_zero)

\end{lstlisting}
\begin{proof}
Choose a measurable representative $u'$ of $u$. Apply Stampacchia's zero-set
theorem (\leanname{HasWeakPartialDeriv.toStampacchia}) to the witness with
the measurable representative, then transfer back via a.e.\ equality.
\end{proof}

\begin{lemma}[\leanname{MemW1pWitness.sub\_const}]
On a finite-measure open set $\Omega$, if $u \in W^{1,2}(\Omega)$ then
$u - c \in W^{1,2}(\Omega)$ with the same weak gradient.
\end{lemma}

\begin{lstlisting}[style=Matlab-editor,basicstyle=\mlttfamily,escapechar=",mlshowsectionrules=true,mathescape=true,morecomment={[s]{\%\{}{\%\}}},language=lean,numbers=none]
/-- Subtracting a constant preserves a `W^{1,2}` witness on finite-measure
domains. -/
noncomputable def MemW1pWitness.sub_const
    {Ω : Set E} [IsFiniteMeasure (volume.restrict Ω)]
    (hΩ : IsOpen Ω)
    {u : E → ℝ} (hw : MemW1pWitness 2 u Ω) (c : ℝ) :
    MemW1pWitness 2 (fun x => u x - c) Ω where
  memLp := hw.memLp.sub (memLp_const c)
  weakGrad := hw.weakGrad
  weakGrad_component_memLp := hw.weakGrad_component_memLp
  isWeakGrad := by
    intro i φ hφ hφ_supp hφ_sub
    let ei : EuclideanSpace ℝ (Fin d) := EuclideanSpace.single i (1 : ℝ)
    have hkey := hw.isWeakGrad i φ hφ hφ_supp hφ_sub
    have hconst :
        HasWeakPartialDeriv i (fun _ : E => (0 : ℝ)) (fun _ : E => c) Ω := by
      simpa [ei] using
        (HasWeakPartialDeriv.of_contDiff (Ω := Ω) (i := i)
          (f := fun _ : E => c) hΩ contDiff_const)
    have hconst_zero : ∫ x in Ω, c * (fderiv ℝ φ x) ei = 0 := by
      simpa using hconst φ hφ hφ_supp hφ_sub
    have hderiv_cont : Continuous (fun x => (fderiv ℝ φ x) ei) :=
      (hφ.continuous_fderiv (by simp : ((⊤ : ℕ∞) : WithTop ℕ∞) ≠ 0)).clm_apply
        continuous_const
    have hderiv_supp : HasCompactSupport (fun x => (fderiv ℝ φ x) ei) :=
      hφ_supp.fderiv_apply (𝕜 := ℝ) ei
    have hu_int : Integrable (fun x => u x * (fderiv ℝ φ x) ei) (volume.restrict Ω) := by
      have hu_loc : LocallyIntegrable u (volume.restrict Ω) :=
        hw.memLp.locallyIntegrable (by norm_num)
      simpa [smul_eq_mul] using
        hu_loc.integrable_smul_right_of_hasCompactSupport hderiv_cont hderiv_supp
    have hc_int : Integrable (fun x => c * (fderiv ℝ φ x) ei) (volume.restrict Ω) := by
      exact (hderiv_cont.integrable_of_hasCompactSupport hderiv_supp).const_mul c
    have hsub_fun :
        (fun x => (u x - c) * (fderiv ℝ φ x) ei) =
          (fun x => u x * (fderiv ℝ φ x) ei - c * (fderiv ℝ φ x) ei) := by
      ext x
      ring
    rw [hsub_fun, integral_sub hu_int hc_int, hkey, hconst_zero]
    simp

\end{lstlisting}
\begin{proof}
Constants have weak gradient zero on $\Omega$. Subtract pointwise and use
that the test integral splits, the contribution from the constant being
zero.
\end{proof}

\begin{lemma}[\leanname{MemW1pWitness.mul\_smooth\_bounded}]
Let $u \in W^{1,2}(\Omega)$ and $\eta \in C^\infty(E)$ with
$|\eta| \le C_0$ and $\|D\eta\| \le C_1$. Then $\eta\, u \in W^{1,2}(\Omega)$
with weak gradient $\eta\, \nabla u + u\, \nabla \eta$.
\end{lemma}

\begin{lstlisting}[style=Matlab-editor,basicstyle=\mlttfamily,escapechar=",mlshowsectionrules=true,mathescape=true,morecomment={[s]{\%\{}{\%\}}},language=lean,numbers=none]
/-- Multiply a `W^{1,2}` witness by a bounded smooth scalar factor. -/
noncomputable def MemW1pWitness.mul_smooth_bounded
    {Ω : Set E} (hΩ : IsOpen Ω)
    {u η : E → ℝ} (hw : MemW1pWitness 2 u Ω)
    (hη : ContDiff ℝ (⊤ : ℕ∞) η)
    {C₀ C₁ : ℝ}
    (_hC₀ : 0 ≤ C₀) (_hC₁ : 0 ≤ C₁)
    (hη_bound : ∀ x, |η x| ≤ C₀)
    (hη_grad_bound : ∀ x, ‖fderiv ℝ η x‖ ≤ C₁) :
    MemW1pWitness 2 (fun x => η x * u x) Ω where
  memLp := by
    refine MemLp.of_le_mul (g := u) (c := C₀) hw.memLp ?_ ?_
    · exact hη.continuous.aestronglyMeasurable.mul hw.memLp.aestronglyMeasurable
    · exact Filter.Eventually.of_forall fun x => by
        calc
          ‖η x * u x‖ = |η x| * ‖u x‖ := by
            rw [norm_mul, Real.norm_eq_abs]
          _ ≤ C₀ * ‖u x‖ := by
              gcongr
              exact hη_bound x
  weakGrad := fun x =>
    η x • hw.weakGrad x +
      (show E from WithLp.toLp 2 fun i => (fderiv ℝ η x) (EuclideanSpace.single i 1) * u x)
  weakGrad_component_memLp := by
    intro i
    have hfirst : MemLp (fun x => η x * hw.weakGrad x i) 2 (volume.restrict Ω) := by
      refine MemLp.of_le_mul (g := fun x => hw.weakGrad x i) (c := C₀)
        (hw.weakGrad_component_memLp i) ?_ ?_
      · exact hη.continuous.aestronglyMeasurable.mul
          (hw.weakGrad_component_memLp i).aestronglyMeasurable
      · exact Filter.Eventually.of_forall fun x => by
          calc
            ‖η x * hw.weakGrad x i‖ = |η x| * ‖hw.weakGrad x i‖ := by
              rw [norm_mul, Real.norm_eq_abs]
            _ ≤ C₀ * ‖hw.weakGrad x i‖ := by
                gcongr
                exact hη_bound x
    have hderiv_cont : Continuous (fun x => (fderiv ℝ η x) (EuclideanSpace.single i 1)) :=
      (hη.continuous_fderiv (by simp : ((⊤ : ℕ∞) : WithTop ℕ∞) ≠ 0)).clm_apply continuous_const
    have hsecond :
        MemLp (fun x => (fderiv ℝ η x) (EuclideanSpace.single i 1) * u x) 2
          (volume.restrict Ω) := by
      refine MemLp.of_le_mul (g := u) (c := C₁) hw.memLp ?_ ?_
      · exact hderiv_cont.aestronglyMeasurable.mul hw.memLp.aestronglyMeasurable
      · exact Filter.Eventually.of_forall fun x => by
          calc
            ‖(fderiv ℝ η x) (EuclideanSpace.single i 1) * u x‖
                = ‖(fderiv ℝ η x) (EuclideanSpace.single i 1)‖ * ‖u x‖ := by
                    rw [norm_mul]
            _ ≤ ‖fderiv ℝ η x‖ * ‖u x‖ := by
                  gcongr
                  simpa using (ContinuousLinearMap.le_opNorm (fderiv ℝ η x)
                    (EuclideanSpace.single i (1 : ℝ)))
            _ ≤ C₁ * ‖u x‖ := by
                  gcongr
                  exact hη_grad_bound x
    simpa using hfirst.add hsecond
  isWeakGrad := by
    intro i
    simpa using HasWeakPartialDeriv.mul_smooth hΩ
      (hf := hw.isWeakGrad i) hη
      (hw.memLp.locallyIntegrable (by norm_num))
      ((hw.weakGrad_component_memLp i).locallyIntegrable (by norm_num))

/-
A compactly supported `W₀^{1,2}` witness already contains the approximation
data needed by the positive-part closure arguments.
-/
\end{lstlisting}
\begin{proof}
The function-side $L^2$ membership follows from $|\eta\, u| \le C_0\,|u|$.
For each component, the candidate $(\eta\, (\nabla u)_i + (\partial_i \eta)\,u)$
is in $L^2$ by similar bounds. The weak-derivative identity follows from
the product rule for weak derivatives, \leanname{HasWeakPartialDeriv.mul\_smooth}.
\end{proof}

\begin{theorem}[\leanname{HasWeakPartialDeriv.of\_eLpNormApprox}]
Let $\Omega$ be open, $i$ an index, $f \in L^2(\Omega)$ with candidate
$g \in L^2(\Omega)$, and let $\psi_n$ be smooth functions whose pointwise
weak partial derivatives $g_n$ converge to $g$ in $L^2(\Omega)$ and whose
values converge to $f$ in $L^2(\Omega)$. Then $g$ is the weak $i$-th
partial derivative of $f$ on $\Omega$.
\end{theorem}

\begin{lstlisting}[style=Matlab-editor,basicstyle=\mlttfamily,escapechar=",mlshowsectionrules=true,mathescape=true,morecomment={[s]{\%\{}{\%\}}},language=lean,numbers=none]
omit [NeZero d] in
theorem HasWeakPartialDeriv.of_eLpNormApprox
    {Ω : Set E} (_hΩ : IsOpen Ω)
    {i : Fin d} {f g : E → ℝ} {ψ : ℕ → E → ℝ} {gψ : ℕ → E → ℝ}
    (hf_memLp : MemLp f 2 (volume.restrict Ω))
    (hg_memLp : MemLp g 2 (volume.restrict Ω))
    (hψ_wd : ∀ n, HasWeakPartialDeriv i (gψ n) (ψ n) Ω)
    (hψ_fun_memLp : ∀ n, MemLp (fun x => ψ n x - f x) 2 (volume.restrict Ω))
    (hψ_fun :
      Tendsto (fun n => eLpNorm (fun x => ψ n x - f x) 2 (volume.restrict Ω))
        atTop (nhds 0))
    (hψ_grad_memLp : ∀ n, MemLp (fun x => gψ n x - g x) 2 (volume.restrict Ω))
    (hψ_grad :
      Tendsto (fun n => eLpNorm (fun x => gψ n x - g x) 2 (volume.restrict Ω))
        atTop (nhds 0)) :
    HasWeakPartialDeriv i g f Ω := by
  intro φ hφ hφ_supp hφ_sub
  let μ : Measure E := volume.restrict Ω
  let dφ : E → ℝ := fun x => (fderiv ℝ φ x) (EuclideanSpace.single i 1)
  have hdφ_memLp : MemLp dφ 2 μ := by
    have hcont : Continuous dφ := by
      simpa [dφ] using
        ((hφ.continuous_fderiv (by simp : ((⊤ : ℕ∞) : WithTop ℕ∞) ≠ 0)).clm_apply
          continuous_const)
    have hcpt : HasCompactSupport dφ := by
      simpa [dφ] using hφ_supp.fderiv_apply (𝕜 := ℝ) (EuclideanSpace.single i 1)
    exact (hcont.memLp_of_hasCompactSupport hcpt).restrict _
  have hφ_memLp : MemLp φ 2 μ :=
    (hφ.continuous.memLp_of_hasCompactSupport hφ_supp).restrict _
  have h_fun_int :
      Tendsto (fun n => ∫ x, dφ x * (ψ n x - f x) ∂μ) atTop (nhds 0) :=
    tendsto_integral_mul_of_eLpNorm_tendsto_zero hdφ_memLp hψ_fun_memLp hψ_fun
  have h_grad_int :
      Tendsto (fun n => ∫ x, φ x * (gψ n x - g x) ∂μ) atTop (nhds 0) :=
    tendsto_integral_mul_of_eLpNorm_tendsto_zero hφ_memLp hψ_grad_memLp hψ_grad
  have h_lhs_tendsto :
      Tendsto (fun n => ∫ x in Ω, ψ n x * dφ x)
        atTop (nhds (∫ x in Ω, f x * dφ x)) := by
    have h_eq :
        ∀ n,
          ∫ x, ψ n x * dφ x ∂μ =
            (∫ x, f x * dφ x ∂μ) + ∫ x, dφ x * (ψ n x - f x) ∂μ := by
      intro n
      have hfi : Integrable (fun x => f x * dφ x) μ := hf_memLp.integrable_mul hdφ_memLp
      have hdi : Integrable (fun x => dφ x * (ψ n x - f x)) μ :=
        hdφ_memLp.integrable_mul (hψ_fun_memLp n)
      calc
        ∫ x, ψ n x * dφ x ∂μ
            = ∫ x, (f x * dφ x) + dφ x * (ψ n x - f x) ∂μ := by
                congr with x
                ring
        _ = (∫ x, f x * dφ x ∂μ) + ∫ x, dφ x * (ψ n x - f x) ∂μ :=
              integral_add hfi hdi
    have h_aux :
        Tendsto
          (fun n => (∫ x, f x * dφ x ∂μ) + ∫ x, dφ x * (ψ n x - f x) ∂μ)
          atTop (nhds ((∫ x, f x * dφ x ∂μ) + 0)) :=
      Tendsto.const_add _ h_fun_int
    simpa [μ, h_eq, add_comm, add_left_comm, add_assoc] using h_aux
  have h_rhs_tendsto :
      Tendsto (fun n => -∫ x in Ω, gψ n x * φ x)
        atTop (nhds (-∫ x in Ω, g x * φ x)) := by
    have h_eq :
        ∀ n,
          -(∫ x, gψ n x * φ x ∂μ) =
            -(∫ x, g x * φ x ∂μ) - ∫ x, φ x * (gψ n x - g x) ∂μ := by
      intro n
      have hgi : Integrable (fun x => g x * φ x) μ := hg_memLp.integrable_mul hφ_memLp
      have hdi : Integrable (fun x => φ x * (gψ n x - g x)) μ :=
        hφ_memLp.integrable_mul (hψ_grad_memLp n)
      have : ∫ x, gψ n x * φ x ∂μ =
          (∫ x, g x * φ x ∂μ) + ∫ x, φ x * (gψ n x - g x) ∂μ := by
        calc
          ∫ x, gψ n x * φ x ∂μ
              = ∫ x, (g x * φ x) + φ x * (gψ n x - g x) ∂μ := by
                  congr with x
                  ring
          _ = (∫ x, g x * φ x ∂μ) + ∫ x, φ x * (gψ n x - g x) ∂μ :=
                integral_add hgi hdi
      linarith
    have h_aux :
        Tendsto
          (fun n => -(∫ x, g x * φ x ∂μ) - ∫ x, φ x * (gψ n x - g x) ∂μ)
          atTop (nhds (-(∫ x, g x * φ x ∂μ) - 0)) :=
      Tendsto.const_sub _ h_grad_int
    simpa [μ, h_eq] using h_aux
  have h_eq_n :
      ∀ n, ∫ x in Ω, ψ n x * dφ x = -∫ x in Ω, gψ n x * φ x := by
    intro n
    exact hψ_wd n φ hφ hφ_supp hφ_sub
  have h_eq_tendsto :
      Tendsto (fun n => ∫ x in Ω, ψ n x * dφ x)
        atTop (nhds (-∫ x in Ω, g x * φ x)) := by
    simpa [h_eq_n] using h_rhs_tendsto
  exact tendsto_nhds_unique h_lhs_tendsto h_eq_tendsto

\end{lstlisting}
\begin{proof}
Fix a smooth compactly supported test $\varphi$ on $\Omega$, and write
$d\varphi(x) := (\fderiv{\varphi}{x})(e_i)$. The integration-by-parts identity for the
smooth approximants reads
\[
  \int_{\Omega} \psi_n \cdot d\varphi\,dx
    = -\int_{\Omega} g_n \cdot \varphi\,dx,
\]
which is built into the hypothesis $\mathrm{HasWeakPartialDeriv}\,i\,g_n\,\psi_n\,\Omega$.

Both $\varphi$ and $d\varphi$ are continuous with compact support, hence in
$L^2(\Omega)$. By the elementary $L^2$-pairing limit
\leanname{tendsto\_integral\_mul\_of\_eLpNorm\_tendsto\_zero}: if
$h_n - h \to 0$ in $L^2$ and $k$ is a fixed $L^2$ function, then
$\int k\,(h_n - h) \to 0$. Applying this with $k = d\varphi$ and
$h_n - h = \psi_n - f$ gives
\[
  \int_{\Omega} d\varphi\,(\psi_n - f)\,dx \to 0,
\]
and similarly with $k = \varphi$ and $h_n - h = g_n - g$ gives
$\int_{\Omega} \varphi\,(g_n - g)\,dx \to 0$. Adding the IBP identity and
$\int f\cdot d\varphi$ to the first limit gives
\[
  \int_{\Omega} \psi_n\cdot d\varphi\,dx \to \int_{\Omega} f \cdot d\varphi\,dx.
\]
Likewise the right-hand side of IBP converges to $-\int_{\Omega} g\cdot\varphi\,dx$.
Passing to the limit in the smooth IBP identity therefore yields
$\int f\,d\varphi = -\int g\,\varphi$, which is the defining identity for $g$
to be the weak $i$-th partial derivative of $f$ on $\Omega$.
\end{proof}

\begin{theorem}[\leanname{tendsto\_eLpNorm\_indicator\_diff\_mul\_of\_tendsto\_eLpNorm}]
Let $\psi_n \to u$ in $L^2(\Omega)$, with $u \in L^2(\Omega)$ and $G \in
L^2(\Omega)$ such that $G$ vanishes a.e.\ on $\{u = 0\}$. Then
\[
\eLpNormText{(\mathbf{1}_{\{\psi_n > 0\}} - \mathbf{1}_{\{u > 0\}})\,G}_{L^2(\Omega)}
\to 0.
\]
\end{theorem}

\begin{lstlisting}[style=Matlab-editor,basicstyle=\mlttfamily,escapechar=",mlshowsectionrules=true,mathescape=true,morecomment={[s]{\%\{}{\%\}}},language=lean,numbers=none]
omit [NeZero d] in
theorem tendsto_eLpNorm_indicator_diff_mul_of_tendsto_eLpNorm
    {μ : Measure E} {u G : E → ℝ} {ψ : ℕ → E → ℝ}
    (hψ_aestrong : ∀ n, AEStronglyMeasurable (ψ n) μ)
    (hu_aestrong : AEStronglyMeasurable u μ)
    (hG_memLp : MemLp G 2 μ)
    (hzero : ∀ᵐ x ∂μ, u x = 0 → G x = 0)
    (hψ :
      Tendsto (fun n => eLpNorm (fun x => ψ n x - u x) 2 μ) atTop (nhds 0)) :
    Tendsto (fun n => eLpNorm (fun x =>
      (if 0 < ψ n x then G x else 0) -
        (if 0 < u x then G x else 0))
      2 μ) atTop (nhds 0) := by
  let F : ℕ → E → ℝ := fun n x =>
    (if 0 < ψ n x then G x else 0) - (if 0 < u x then G x else 0)
  have hF_meas : ∀ n, AEStronglyMeasurable (F n) μ := by
    intro n
    have hψ_pos : NullMeasurableSet {x | 0 < ψ n x} μ := by
      exact AEStronglyMeasurable.nullMeasurableSet_lt
        (μ := μ) stronglyMeasurable_const.aestronglyMeasurable (hψ_aestrong n)
    have hu_pos : NullMeasurableSet {x | 0 < u x} μ := by
      exact AEStronglyMeasurable.nullMeasurableSet_lt
        (μ := μ) stronglyMeasurable_const.aestronglyMeasurable hu_aestrong
    simpa [F] using
      ((hG_memLp.aestronglyMeasurable.indicator₀ hψ_pos).sub
        (hG_memLp.aestronglyMeasurable.indicator₀ hu_pos))
  have hF_dom : ∀ n x, ‖F n x‖ ≤ ‖G x‖ := by
    intro n x
    exact indicator_diff_mul_le_abs (a := ψ n x) (c := u x) (d := G x)
  have huiG : UnifIntegrable (fun _ : ℕ => G) 2 μ := by
    exact MeasureTheory.unifIntegrable_const (p := (2 : ℝ≥0∞))
      (by norm_num) (by simp) hG_memLp
  have huiF : UnifIntegrable F 2 μ := by
    intro ε hε
    obtain ⟨δ, hδ, hδ'⟩ := huiG hε
    refine ⟨δ, hδ, fun n s hs hμs => ?_⟩
    calc
      eLpNorm (s.indicator (F n)) 2 μ ≤ eLpNorm (s.indicator G) 2 μ := by
        refine eLpNorm_mono ?_
        intro x
        by_cases hx : x ∈ s
        · simp [Set.indicator_of_mem, hx]
          exact hF_dom n x
        · simp [Set.indicator_of_notMem, hx]
      _ ≤ ENNReal.ofReal ε := hδ' 0 s hs hμs
  have hutG : UnifTight (fun _ : ℕ => G) 2 μ := by
    exact MeasureTheory.unifTight_const (p := (2 : ℝ≥0∞)) (by simp) hG_memLp
  have hutF : UnifTight F 2 μ := by
    intro ε hε
    obtain ⟨s, hμs, hs'⟩ := hutG hε
    refine ⟨s, hμs, fun n => ?_⟩
    calc
      eLpNorm (sᶜ.indicator (F n)) 2 μ ≤ eLpNorm (sᶜ.indicator G) 2 μ := by
        refine eLpNorm_mono ?_
        intro x
        by_cases hx : x ∈ sᶜ
        · simp [Set.indicator_of_mem, hx]
          exact hF_dom n x
        · simp [Set.indicator_of_notMem, hx]
      _ ≤ ε := hs' 0
  refine tendsto_of_subseq_tendsto ?_
  intro ns hns
  obtain ⟨ms, -, hms_ae⟩ :=
    exists_subseq_tendsto_ae_of_tendsto_eLpNorm
      (μ := μ) (f := u) (ψ := fun n => ψ (ns n))
      (fun n => hψ_aestrong (ns n)) hu_aestrong (hψ.comp hns)
  have hui_subseq : UnifIntegrable (fun n => F (ns (ms n))) 2 μ := by
    intro ε hε
    obtain ⟨δ, hδ, hδ'⟩ := huiF hε
    exact ⟨δ, hδ, fun n s hs hμs => hδ' (ns (ms n)) s hs hμs⟩
  have hut_subseq : UnifTight (fun n => F (ns (ms n))) 2 μ := by
    intro ε hε
    obtain ⟨s, hμs, hs'⟩ := hutF hε
    exact ⟨s, hμs, fun n => hs' (ns (ms n))⟩
  have hF_ae :
      ∀ᵐ x ∂μ, Tendsto (fun n => F (ns (ms n)) x) atTop (nhds 0) := by
    filter_upwards [hzero, hms_ae] with x hxzero hxt
    simpa [F] using
      tendsto_indicator_diff_mul_zero_of_tendsto
        (a := fun n => ψ (ns (ms n)) x) (u := u x) (g := G x) hxt hxzero
  refine ⟨ms, ?_⟩
  have hLpF :
      Tendsto (fun n => eLpNorm (F (ns (ms n))) 2 μ) atTop (nhds 0) := by
    have hLpF0 :
        Tendsto (fun n => eLpNorm (F (ns (ms n)) - fun _ => (0 : ℝ)) 2 μ) atTop (nhds 0) := by
      exact
        (MeasureTheory.tendsto_Lp_of_tendsto_ae (μ := μ) (p := (2 : ℝ≥0∞))
          (hp := by norm_num) (hp' := by simp)
          (haef := fun n => hF_meas (ns (ms n)))
          (hg' := (MeasureTheory.MemLp.zero' (p := (2 : ℝ≥0∞)) (μ := μ)))
          hui_subseq hut_subseq hF_ae)
    have hEq0 :
        (fun n => eLpNorm (F (ns (ms n)) - fun _ => (0 : ℝ)) 2 μ) =
          (fun n => eLpNorm (F (ns (ms n))) 2 μ) := by
      funext n
      congr 1
      ext x
      simp
    rw [hEq0] at hLpF0
    exact hLpF0
  simpa [F] using hLpF

\end{lstlisting}
\begin{proof}
Write $F_n(x) := (\mathbf{1}_{\{\psi_n > 0\}}(x) - \mathbf{1}_{\{u > 0\}}(x))\,G(x)$.
The pointwise inequality $|F_n(x)| \le |G(x)|$ is immediate from the fact that
each indicator lies in $\{0,1\}$, so the family $(F_n)$ is uniformly dominated
by $|G| \in L^2(\Omega)$. Each $F_n$ is a.e.\ strongly measurable: the sets
$\{\psi_n > 0\}$ and $\{u > 0\}$ are null measurable since $\psi_n$ and $u$
are a.e.\ strongly measurable (the comparison with the constant $0$ is
handled by \leanname{AEStronglyMeasurable.nullMeasurableSet\_lt}), and $G$
is in $L^2$.

Vitali convergence in $L^2$ requires uniform integrability and uniform
tightness of the family $(F_n)$. Both are inherited from the constant family
$\{G\}$, which is uniformly integrable and uniformly tight by
\leanname{MeasureTheory.unifIntegrable\_const} and
\leanname{MeasureTheory.unifTight\_const} applied to the singleton
$\{G\} \subset L^2$. The pointwise domination $|F_n| \le |G|$ then transfers
both properties to $(F_n)$ via $\eLpNormText{\mathbf{1}_S F_n}_{L^2} \le
\eLpNormText{\mathbf{1}_S G}_{L^2}$ for any measurable set $S$.

It remains to show $F_n \to 0$ in measure. By \leanname{tendsto\_of\_subseq\_tendsto},
it suffices that every subsequence of $(F_n)$ admits an a.e.\ convergent
sub-subsequence with limit $0$. Given any subsequence, the $L^2$ hypothesis
$\psi_{n_k} \to u$ in $L^2$ supplies an a.e.\ convergent sub-sub-sequence
(\leanname{exists\_subseq\_tendsto\_ae\_of\_tendsto\_eLpNorm}). On the
sub-subsequence: at points where $u(x) > 0$, eventually $\psi_{n_{k_j}}(x) > 0$,
so both indicators equal $1$ and $F_{n_{k_j}}(x) \to 0$; at points where
$u(x) < 0$ both indicators eventually equal $0$; at points where $u(x) = 0$
the hypothesis hzero gives $G(x) = 0$ a.e., so $F_{n_{k_j}}(x) = 0$.
This produces an a.e.\ vanishing sub-subsequence. Combining with the Vitali
convergence theorem (uniform integrability + uniform tightness +
convergence in measure $\Rightarrow$ $L^2$ convergence) yields the desired
$\eLpNormText{F_n}_{L^2} \to 0$.
\end{proof}

\begin{definition}[\leanname{MemW1pWitness.posPart\_of\_aux}]
Given $u \in W^{1,2}(\Omega)$ with smooth approximation data $(\psi_n)$ and
the zero-set vanishing hypothesis from
\leanname{MemW1pWitness.weakGrad\_ae\_eq\_zero\_on\_zeroSet}, the positive part
$u^+ := \max(u,0) \in W^{1,2}(\Omega)$ has weak gradient
\[
(\nabla u^+)_i(x) := \begin{cases}
(\nabla u)_i(x) & \text{if } u(x) > 0, \\
0 & \text{otherwise}.
\end{cases}
\]
\end{definition}

\begin{lstlisting}[style=Matlab-editor,basicstyle=\mlttfamily,escapechar=",mlshowsectionrules=true,mathescape=true,morecomment={[s]{\%\{}{\%\}}},language=lean,numbers=none]
/-- Auxiliary positive-part witness constructor: this is the transportable core
assuming the zero-set vanishing input and the smooth approximation data are
already available. -/
noncomputable def MemW1pWitness.posPart_of_aux
    {Ω : Set E} [IsFiniteMeasure (volume.restrict Ω)]
    (hΩ : IsOpen Ω) {u : E → ℝ}
    (hw : MemW1pWitness 2 u Ω)
    (hzero :
      ∀ i : Fin d, ∀ᵐ x ∂(volume.restrict Ω), u x = 0 → hw.weakGrad x i = 0)
    (happrox :
      ∃ ψ : ℕ → E → ℝ,
        (∀ n, ContDiff ℝ 1 (ψ n)) ∧
        (∀ n, HasCompactSupport (ψ n)) ∧
        Tendsto (fun n => eLpNorm (fun x => ψ n x - u x) 2 (volume.restrict Ω))
          atTop (nhds 0) ∧
        (∀ i : Fin d,
          Tendsto (fun n => eLpNorm (fun x =>
            (fderiv ℝ (ψ n) x) (EuclideanSpace.single i 1) - hw.weakGrad x i)
            2 (volume.restrict Ω))
          atTop (nhds 0))) :
    MemW1pWitness 2 (fun x => max (u x) 0) Ω := by
  let ψ : ℕ → E → ℝ := Classical.choose happrox
  have hψspec := Classical.choose_spec happrox
  have hψ_smooth : ∀ n, ContDiff ℝ 1 (ψ n) := hψspec.1
  have hψspec' := hψspec.2
  have hψ_cpt : ∀ n, HasCompactSupport (ψ n) := hψspec'.1
  have hψspec'' := hψspec'.2
  have hψ_func :
      Tendsto (fun n => eLpNorm (fun x => ψ n x - u x) 2 (volume.restrict Ω))
        atTop (nhds 0) := hψspec''.1
  have hψ_grad :
      ∀ i : Fin d,
        Tendsto (fun n => eLpNorm (fun x =>
          (fderiv ℝ (ψ n) x) (EuclideanSpace.single i 1) - hw.weakGrad x i)
          2 (volume.restrict Ω))
        atTop (nhds 0) := hψspec''.2
  refine
    { memLp := by
        simpa using hw.memLp.pos_part
      weakGrad := fun x => if 0 < u x then hw.weakGrad x else 0
      weakGrad_component_memLp := ?_
      isWeakGrad := ?_ }
  ·
    intro i
    simpa using posPartGrad_component_memLp hw i
  ·
    intro i
    let μ : Measure E := volume.restrict Ω
    let g : E → ℝ := fun x => if 0 < u x then hw.weakGrad x i else 0
    let gψ : ℕ → E → ℝ := fun n x =>
      if 0 < ψ n x then (fderiv ℝ (ψ n) x) (EuclideanSpace.single i 1) else 0
    have hg_memLp : MemLp g 2 μ := by
      exact indicator_component_memLp hw.memLp.aemeasurable (hw.weakGrad_component_memLp i)
    have hψ_wd : ∀ n, HasWeakPartialDeriv i (gψ n) (fun x => max (ψ n x) 0) Ω := by
      intro n
      simpa [gψ] using
        weakPartialDeriv_positivePart_of_contDiff (Ω := Ω) (i := i) (hψ_smooth n) (hψ_cpt n)
    have hψ_fun_memLp : ∀ n, MemLp (fun x => max (ψ n x) 0 - max (u x) 0) 2 μ := by
      intro n
      have hfirst : MemLp (fun x => max (ψ n x) 0) 2 μ := by
        exact (((hψ_smooth n).continuous.max continuous_const).memLp_of_hasCompactSupport
          (hasCompactSupport_positivePart (hψ_cpt n))).restrict Ω
      exact hfirst.sub (by simpa using hw.memLp.pos_part)
    have hψ_fun :
        Tendsto (fun n => eLpNorm (fun x => max (ψ n x) 0 - max (u x) 0) 2 μ)
          atTop (nhds 0) :=
      tendsto_eLpNorm_positivePart_sub hψ_func
    have hψ_grad_memLp : ∀ n, MemLp (fun x => gψ n x - g x) 2 μ := by
      intro n
      have hderiv_cont : Continuous
          (fun x => (fderiv ℝ (ψ n) x) (EuclideanSpace.single i 1)) :=
        ((hψ_smooth n).fderiv_right (m := 0) (by norm_cast)).continuous.clm_apply
          continuous_const
      have hderiv_cpt : HasCompactSupport
          (fun x => (fderiv ℝ (ψ n) x) (EuclideanSpace.single i 1)) :=
        (hψ_cpt n).fderiv_apply (𝕜 := ℝ) _
      have hderiv_memLp : MemLp
          (fun x => (fderiv ℝ (ψ n) x) (EuclideanSpace.single i 1)) 2 μ :=
        (hderiv_cont.memLp_of_hasCompactSupport hderiv_cpt).restrict Ω
      have hfirst_raw : MemLp
          (fun x => (fderiv ℝ (ψ n) x) (EuclideanSpace.single i 1) - hw.weakGrad x i) 2 μ :=
        hderiv_memLp.sub (hw.weakGrad_component_memLp i)
      have hψ_meas : AEMeasurable (ψ n) μ :=
        ((hψ_smooth n).continuous.aemeasurable.restrict)
      have hfirst : MemLp
          (fun x => if 0 < ψ n x then
            (fderiv ℝ (ψ n) x) (EuclideanSpace.single i 1) - hw.weakGrad x i else 0)
          2 μ := indicator_component_memLp hψ_meas hfirst_raw
      have hsecond : MemLp
          (fun x => (if 0 < ψ n x then hw.weakGrad x i else 0) - g x) 2 μ := by
        have hleft : MemLp (fun x => if 0 < ψ n x then hw.weakGrad x i else 0) 2 μ :=
          indicator_component_memLp hψ_meas (hw.weakGrad_component_memLp i)
        exact hleft.sub hg_memLp
      have hEq :
          (fun x => gψ n x - g x) =
            (fun x =>
              (if 0 < ψ n x then
                (fderiv ℝ (ψ n) x) (EuclideanSpace.single i 1) - hw.weakGrad x i else 0) +
                ((if 0 < ψ n x then hw.weakGrad x i else 0) - g x)) := by
        funext x
        by_cases hxψ : 0 < ψ n x <;> by_cases hxu : 0 < u x <;> simp [gψ, g, hxψ, hxu]
      have hsum : MemLp
          (fun x =>
            (if 0 < ψ n x then
              (fderiv ℝ (ψ n) x) (EuclideanSpace.single i 1) - hw.weakGrad x i else 0) +
              ((if 0 < ψ n x then hw.weakGrad x i else 0) - g x))
          2 μ := ⟨hfirst.aestronglyMeasurable.add hsecond.aestronglyMeasurable,
            eLpNorm_add_lt_top hfirst hsecond⟩
      simpa [hEq] using hsum
    have hψ_grad_tendsto :
        Tendsto (fun n => eLpNorm (fun x => gψ n x - g x) 2 μ) atTop (nhds 0) := by
      have hfirst :
          Tendsto
            (fun n => eLpNorm
              (fun x => if 0 < ψ n x then
                (fderiv ℝ (ψ n) x) (EuclideanSpace.single i 1) - hw.weakGrad x i else 0)
              2 μ) atTop (nhds 0) := by
        have hbound :
            ∀ n,
              eLpNorm
                (fun x => if 0 < ψ n x then
                  (fderiv ℝ (ψ n) x) (EuclideanSpace.single i 1) - hw.weakGrad x i else 0)
                2 μ
              ≤ eLpNorm
                  (fun x => (fderiv ℝ (ψ n) x) (EuclideanSpace.single i 1) - hw.weakGrad x i)
                  2 μ := by
          intro n
          refine eLpNorm_mono ?_
          intro x
          by_cases hx : 0 < ψ n x <;> simp [hx]
        have hnonneg :
            ∀ n,
              0 ≤ eLpNorm
                (fun x => if 0 < ψ n x then
                  (fderiv ℝ (ψ n) x) (EuclideanSpace.single i 1) - hw.weakGrad x i else 0)
                2 μ := by
          intro n
          positivity
        exact tendsto_of_tendsto_of_tendsto_of_le_of_le
          tendsto_const_nhds (hψ_grad i) hnonneg hbound
      have hsecond :
          Tendsto
            (fun n => eLpNorm
              (fun x => (if 0 < ψ n x then hw.weakGrad x i else 0) - g x)
              2 μ) atTop (nhds 0) := by
        simpa [g] using
          tendsto_eLpNorm_indicator_diff_mul_of_tendsto_eLpNorm
            (μ := μ) (u := u) (G := fun x => hw.weakGrad x i) (ψ := ψ)
            (hψ_aestrong := fun n => (hψ_smooth n).continuous.aestronglyMeasurable.restrict)
            (hu_aestrong := hw.memLp.aestronglyMeasurable)
            (hG_memLp := hw.weakGrad_component_memLp i)
            (hzero := hzero i)
            (hψ := hψ_func)
      have hbound :
          ∀ n, eLpNorm (fun x => gψ n x - g x) 2 μ ≤
            eLpNorm
              (fun x => if 0 < ψ n x then
                (fderiv ℝ (ψ n) x) (EuclideanSpace.single i 1) - hw.weakGrad x i else 0)
              2 μ
            +
            eLpNorm
              (fun x => (if 0 < ψ n x then hw.weakGrad x i else 0) - g x)
              2 μ := by
        intro n
        have hψ_meas : AEMeasurable (ψ n) μ :=
          ((hψ_smooth n).continuous.aemeasurable.restrict)
        have hderiv_cont : Continuous
            (fun x => (fderiv ℝ (ψ n) x) (EuclideanSpace.single i 1)) :=
          ((hψ_smooth n).fderiv_right (m := 0) (by norm_cast)).continuous.clm_apply
            continuous_const
        have hderiv_cpt : HasCompactSupport
            (fun x => (fderiv ℝ (ψ n) x) (EuclideanSpace.single i 1)) :=
          (hψ_cpt n).fderiv_apply (𝕜 := ℝ) _
        have hderiv_memLp : MemLp
            (fun x => (fderiv ℝ (ψ n) x) (EuclideanSpace.single i 1)) 2 μ :=
          (hderiv_cont.memLp_of_hasCompactSupport hderiv_cpt).restrict Ω
        have hfirst_raw : MemLp
            (fun x => (fderiv ℝ (ψ n) x) (EuclideanSpace.single i 1) - hw.weakGrad x i) 2 μ :=
          hderiv_memLp.sub (hw.weakGrad_component_memLp i)
        have hfirst_mem : MemLp
            (fun x => if 0 < ψ n x then
              (fderiv ℝ (ψ n) x) (EuclideanSpace.single i 1) - hw.weakGrad x i else 0)
            2 μ := indicator_component_memLp hψ_meas hfirst_raw
        have hsecond_mem : MemLp
            (fun x => (if 0 < ψ n x then hw.weakGrad x i else 0) - g x)
            2 μ := by
          exact (indicator_component_memLp hψ_meas (hw.weakGrad_component_memLp i)).sub hg_memLp
        have hEq :
            (fun x => gψ n x - g x) =
              (fun x =>
                (if 0 < ψ n x then
                  (fderiv ℝ (ψ n) x) (EuclideanSpace.single i 1) - hw.weakGrad x i else 0) +
                  ((if 0 < ψ n x then hw.weakGrad x i else 0) - g x)) := by
          funext x
          by_cases hxψ : 0 < ψ n x <;> by_cases hxu : 0 < u x <;> simp [gψ, g, hxψ, hxu]
        rw [hEq]
        exact eLpNorm_add_le
          hfirst_mem.aestronglyMeasurable
          hsecond_mem.aestronglyMeasurable
          (by norm_num : (1 : ℝ≥0∞) ≤ 2)
      have hnonneg : ∀ n, 0 ≤ eLpNorm (fun x => gψ n x - g x) 2 μ := by
        intro n
        positivity
      have hupper :
          Tendsto
            (fun n =>
              eLpNorm
                (fun x => if 0 < ψ n x then
                  (fderiv ℝ (ψ n) x) (EuclideanSpace.single i 1) - hw.weakGrad x i else 0)
                2 μ
              +
              eLpNorm
                (fun x => (if 0 < ψ n x then hw.weakGrad x i else 0) - g x)
                2 μ)
            atTop (nhds 0) := by
        simpa using hfirst.add hsecond
      exact tendsto_of_tendsto_of_tendsto_of_le_of_le tendsto_const_nhds hupper hnonneg hbound
    have hg_eq :
        (fun x => ((if 0 < u x then hw.weakGrad x else 0 : E) i)) = g := by
      funext x
      by_cases hx : 0 < u x <;> simp [g, hx]
    rw [hg_eq]
    exact
      HasWeakPartialDeriv.of_eLpNormApprox (d := d) hΩ
        (i := i) (f := fun x => max (u x) 0) (g := g)
        (ψ := fun n x => max (ψ n x) 0) (gψ := gψ)
        (by simpa using hw.memLp.pos_part)
        hg_memLp
        hψ_wd
        hψ_fun_memLp
        hψ_fun
        hψ_grad_memLp
        hψ_grad_tendsto

/-- Stampacchia's zero-set theorem for Sobolev witnesses: each weak-gradient
component vanishes a.e. on the zero set of the function. -/
theorem MemW1pWitness.weakGrad_ae_eq_zero_on_zeroSet
    {Ω : Set E} (hΩ : IsOpen Ω) {u : E → ℝ}
    (hw : MemW1pWitness 2 u Ω) :
    ∀ i : Fin d, ∀ᵐ x ∂(volume.restrict Ω), u x = 0 → hw.weakGrad x i = 0 := by
  intro i
  let uMk : E → ℝ := AEMeasurable.mk u hw.memLp.aemeasurable
  have hu_meas : Measurable uMk := hw.memLp.aemeasurable.measurable_mk
  have hu_ae : u =ᵐ[volume.restrict Ω] uMk := hw.memLp.aemeasurable.ae_eq_mk
  have hGMk :
      ∀ j : Fin d, HasWeakPartialDeriv' (d := d) j (fun x => hw.weakGrad x j) uMk Ω := by
    intro j
    exact stampacchia_congr_ae
      (HasWeakPartialDeriv.toStampacchia (hw.isWeakGrad j)) hu_ae
  have huMk_memW1p : MemW1p 2 uMk Ω := by
    refine ⟨hw.memLp.ae_eq hu_ae, ?_⟩
    intro j
    exact ⟨fun x => hw.weakGrad x j, hw.weakGrad_component_memLp j, by
      simpa [HasWeakPartialDeriv'] using hGMk j⟩
  have hzero_mk :
      ∀ᵐ x ∂(volume.restrict Ω), uMk x = 0 → hw.weakGrad x i = 0 :=
    DeGiorgi.weakGrad_ae_eq_zero_on_zeroSet hΩ huMk_memW1p hGMk i
      (hw.weakGrad_component_memLp i) hu_meas
  filter_upwards [hu_ae, hzero_mk] with x hx hzero hx_zero
  exact hzero (by simpa [hx] using hx_zero)

\end{lstlisting}
\begin{theorem}[\leanname{MemW1pWitness.posPart}]
On a finite-measure open set $\Omega$, if $u \in W^{1,2}(\Omega)$ comes
equipped with a smooth approximation, then $u^+ \in W^{1,2}(\Omega)$ with
the gradient described above.
\end{theorem}

\begin{lstlisting}[style=Matlab-editor,basicstyle=\mlttfamily,escapechar=",mlshowsectionrules=true,mathescape=true,morecomment={[s]{\%\{}{\%\}}},language=lean,numbers=none]
/-- Positive-part witness constructor with the Stampacchia zero-set input
supplied automatically from the imported Chapter 02 theorem. The remaining
explicit hypothesis is just the smooth approximation package. -/
noncomputable def MemW1pWitness.posPart
    {Ω : Set E} [IsFiniteMeasure (volume.restrict Ω)]
    (hΩ : IsOpen Ω) {u : E → ℝ}
    (hw : MemW1pWitness 2 u Ω)
    (happrox :
      ∃ ψ : ℕ → E → ℝ,
        (∀ n, ContDiff ℝ 1 (ψ n)) ∧
        (∀ n, HasCompactSupport (ψ n)) ∧
        Tendsto (fun n => eLpNorm (fun x => ψ n x - u x) 2 (volume.restrict Ω))
          atTop (nhds 0) ∧
        (∀ i : Fin d,
          Tendsto (fun n => eLpNorm (fun x =>
            (fderiv ℝ (ψ n) x) (EuclideanSpace.single i 1) - hw.weakGrad x i)
            2 (volume.restrict Ω))
          atTop (nhds 0))) :
    MemW1pWitness 2 (fun x => max (u x) 0) Ω :=
  MemW1pWitness.posPart_of_aux hΩ hw
    (hw.weakGrad_ae_eq_zero_on_zeroSet hΩ) happrox

\end{lstlisting}
\begin{proof}
Apply \leanname{posPart\_of\_aux} with the zero-set vanishing input supplied by
\leanname{weakGrad\_ae\_eq\_zero\_on\_zeroSet}. The proof of the auxiliary
theorem proceeds in three steps.
\begin{enumerate}
\item Define $g(x) := \mathbf{1}_{\{u > 0\}}\, (\nabla u)_i(x) \in L^2(\Omega)$
and $g^{\psi_n}(x) := \mathbf{1}_{\{\psi_n > 0\}}\,(\partial_i \psi_n)(x)$.
\item Each $g^{\psi_n}$ is the classical weak partial derivative of
$\max(\psi_n, 0)$ by
\leanname{weakPartialDeriv\_positivePart\_of\_contDiff} (since $\psi_n$ is
$C^1$).
\item Show $\max(\psi_n, 0) \to u^+$ in $L^2(\Omega)$ by
\leanname{tendsto\_eLpNorm\_positivePart\_sub} and $g^{\psi_n} \to g$ in
$L^2(\Omega)$ by combining
\leanname{tendsto\_eLpNorm\_indicator\_diff\_mul\_of\_tendsto\_eLpNorm} (which uses
the zero-set hypothesis) with the $L^2$ convergence
$\partial_i \psi_n \to (\nabla u)_i$.
\item Apply Theorem~\leanname{HasWeakPartialDeriv.of\_eLpNormApprox} to deduce
that $g$ is a weak $i$-th partial derivative of $u^+$ on $\Omega$.
\end{enumerate}
\end{proof}

\subsection{$L^p$ function toolkit}

\begin{theorem}[\leanname{memLp\_of\_tendsto\_eLpNorm}]
Let $1 \le p \le \infty$, let $f_n \in L^p(\mu)$, and let $g$ be a.e.
strongly measurable. If $\eLpNormText{f_n - g}_{L^p} \to 0$, then
$g \in L^p(\mu)$.
\end{theorem}

\begin{lstlisting}[style=Matlab-editor,basicstyle=\mlttfamily,escapechar=",mlshowsectionrules=true,mathescape=true,morecomment={[s]{\%\{}{\%\}}},language=lean,numbers=none]
/-- If `f n → g` in eLpNorm and each `f n ∈ Lp`, then `g ∈ Lp`,
provided `g` is AEStronglyMeasurable. Avoids the `Lp` type entirely.

The key observation: `eLpNorm (f n - g) → 0` means `eLpNorm (f N - g) < 1` for
some `N`. Then `eLpNorm g ≤ eLpNorm (f N - g) + eLpNorm (f N) < ∞`. -/
theorem memLp_of_tendsto_eLpNorm
    {p : ℝ≥0∞} (hp : 1 ≤ p)
    {f : ℕ → α → E} {g : α → E}
    (hf_memLp : ∀ n, MemLp (f n) p μ)
    (hg_aesm : AEStronglyMeasurable g μ)
    (hfg : Tendsto (fun n => eLpNorm (f n - g) p μ) atTop (nhds 0)) :
    MemLp g p μ := by
  refine ⟨hg_aesm, ?_⟩
  -- Find N with eLpNorm (f N - g) < 1
  obtain ⟨N, hN⟩ := (hfg.eventually (gt_mem_nhds (by norm_num : (0 : ℝ≥0∞) < 1))).exists
  -- g = f N - (f N - g), so eLpNorm g ≤ eLpNorm (f N) + eLpNorm (f N - g)
  calc eLpNorm g p μ
      = eLpNorm (f N - (f N - g)) p μ := by congr 1; ext x; simp
    _ ≤ eLpNorm (f N) p μ + eLpNorm (f N - g) p μ :=
        eLpNorm_sub_le (hf_memLp N).aestronglyMeasurable
          ((hf_memLp N).aestronglyMeasurable.sub hg_aesm) hp
    _ < ⊤ := ENNReal.add_lt_top.mpr
        ⟨(hf_memLp N).eLpNorm_lt_top, lt_of_lt_of_le hN (by norm_num)⟩

\end{lstlisting}
\begin{proof}
For some $N$, $\eLpNormText{f_N - g}_{L^p} < 1$. The triangle inequality
gives $\eLpNormText{g}_{L^p} \le \eLpNormText{f_N}_{L^p} + 1 < \infty$.
\end{proof}

\begin{theorem}[\leanname{ae\_eq\_of\_tendsto\_eLpNorm\_sub}]
Lp limits are unique up to a.e.\ equality.
\end{theorem}

\begin{lstlisting}[style=Matlab-editor,basicstyle=\mlttfamily,escapechar=",mlshowsectionrules=true,mathescape=true,morecomment={[s]{\%\{}{\%\}}},language=lean,numbers=none]
/-- Lp limit uniqueness: if f_n → g₁ and f_n → g₂ in eLpNorm, then g₁ =ᵐ g₂. -/
theorem ae_eq_of_tendsto_eLpNorm_sub
    {p : ℝ≥0∞} (hp : 1 ≤ p)
    {f : ℕ → α → E} {g₁ g₂ : α → E}
    (hf_aesm : ∀ n, AEStronglyMeasurable (f n) μ)
    (hg₁ : AEStronglyMeasurable g₁ μ) (hg₂ : AEStronglyMeasurable g₂ μ)
    (h1 : Tendsto (fun n => eLpNorm (f n - g₁) p μ) atTop (nhds 0))
    (h2 : Tendsto (fun n => eLpNorm (f n - g₂) p μ) atTop (nhds 0)) :
    g₁ =ᵐ[μ] g₂ := by
  -- eLpNorm(g₁ - g₂) is constant, bounded by eLpNorm(g₁ - f n) + eLpNorm(f n - g₂) → 0
  have hzero : eLpNorm (g₁ - g₂) p μ = 0 := le_antisymm (by
    have hbound : ∀ n, eLpNorm (g₁ - g₂) p μ ≤
        eLpNorm (f n - g₁) p μ + eLpNorm (f n - g₂) p μ := by
      intro n
      calc eLpNorm (g₁ - g₂) p μ
          = eLpNorm ((g₁ - f n) + (f n - g₂)) p μ := by
            congr 1; ext x; simp [sub_add_sub_cancel]
        _ ≤ eLpNorm (g₁ - f n) p μ + eLpNorm (f n - g₂) p μ :=
            eLpNorm_add_le (hg₁.sub (hf_aesm n)) ((hf_aesm n).sub hg₂) hp
        _ = eLpNorm (f n - g₁) p μ + eLpNorm (f n - g₂) p μ := by
            congr 1
            rw [show g₁ - f n = -(f n - g₁) from by
                  ext x
                  simp [sub_eq_add_neg],
                eLpNorm_neg]
    -- Constant ≤ sum → 0, so constant = 0
    have hsum_tendsto : Tendsto (fun n => eLpNorm (f n - g₁) p μ + eLpNorm (f n - g₂) p μ)
        atTop (nhds 0) := by simpa [add_zero] using h1.add h2
    exact ge_of_tendsto hsum_tendsto (Filter.Eventually.of_forall hbound)) bot_le
  have hae_zero := (eLpNorm_eq_zero_iff (hg₁.sub hg₂)
    (ne_of_gt (lt_of_lt_of_le (by simp : (0 : ℝ≥0∞) < 1) hp))).mp hzero
  -- g₁ - g₂ =ᵐ 0 → g₁ =ᵐ g₂
  exact hae_zero.mono fun x hx => by simpa [sub_eq_zero] using hx

\end{lstlisting}
\begin{proof}
$\eLpNormText{g_1 - g_2}_{L^p} \le \eLpNormText{f_n - g_1}_{L^p} +
\eLpNormText{f_n - g_2}_{L^p} \to 0$, so $\eLpNormText{g_1 - g_2}_{L^p} = 0$,
and $g_1 - g_2 = 0$ a.e.
\end{proof}

\begin{theorem}[\leanname{scalar\_cauchy\_to\_limit}]
\label{thm:scalar-cauchy}
Let $1 \le p < \infty$, $E$ second-countable complete, and let $f_n \in
L^p(\mu; E)$ be a Cauchy sequence in $L^p$. There exists $g \in L^p(\mu; E)$
with $\eLpNormText{f_n - g}_{L^p}\to 0$.
\end{theorem}

\begin{lstlisting}[style=Matlab-editor,basicstyle=\mlttfamily,escapechar=",mlshowsectionrules=true,mathescape=true,morecomment={[s]{\%\{}{\%\}}},language=lean,numbers=none]
/-- **Scalar Cauchy → limit.** Generic over codomain E and domain α. -/
theorem scalar_cauchy_to_limit
    [SecondCountableTopology E] [CompleteSpace E]
    {p : ℝ≥0∞} (hp1 : 1 ≤ p) (hp_top : p ≠ ⊤)
    {f : ℕ → α → E}
    (hf_memLp : ∀ n, MemLp (f n) p μ)
    (hf_cauchy : Tendsto (fun nm : ℕ × ℕ =>
      eLpNorm (f nm.1 - f nm.2) p μ) atTop (nhds 0)) :
    ∃ g : α → E,
      MemLp g p μ ∧
      Tendsto (fun n => eLpNorm (f n - g) p μ) atTop (nhds 0) := by
  let _ := hp_top
  -- Same approach as exists_pi_limit_of_cauchy_eLpNorm:
  -- extract controlled Cauchy subsequence → cauchy_complete_eLpNorm → upgrade convergence
  have hp : p ≠ 0 := ne_of_gt (lt_of_lt_of_le (by simp : (0 : ℝ≥0∞) < 1) hp1)
  -- Step A: pair-Cauchy → controlled bound for each ε = (2⁻¹)^(k+1)
  have hpair : ∀ k : ℕ, ∃ M : ℕ, ∀ n m, M ≤ n → M ≤ m →
      eLpNorm (f n - f m) p μ ≤ (2⁻¹ : ℝ≥0∞) ^ (k + 1) := by
    intro k
    have hε : (0 : ℝ≥0∞) < (2⁻¹ : ℝ≥0∞) ^ (k + 1) := ENNReal.pow_pos (by norm_num) _
    obtain ⟨⟨N₁, N₂⟩, hNM⟩ := (ENNReal.tendsto_atTop_zero.mp hf_cauchy _ hε)
    exact ⟨max N₁ N₂, fun n m hn hm => hNM ⟨n, m⟩
      ⟨le_trans (le_max_left _ _) hn, le_trans (le_max_right _ _) hm⟩⟩
  -- Step B: Make bounds nondecreasing
  choose M_raw hM_raw using hpair
  let M : ℕ → ℕ := fun k => (Finset.range (k + 1)).sup M_raw
  have hM_ge : ∀ k, M_raw k ≤ M k :=
    fun k => Finset.le_sup (f := M_raw) (Finset.mem_range.mpr (Nat.lt_succ_of_le le_rfl))
  have hM_mono : Monotone M :=
    fun _ _ hab => Finset.sup_mono (Finset.range_mono (Nat.add_le_add_right hab 1))
  -- Step C: Controlled Cauchy bound
  let f_sub : ℕ → α → E := fun k => f (M k)
  have hf_sub_memLp : ∀ k, MemLp (f_sub k) p μ := fun k => hf_memLp (M k)
  have hf_sub_cau : ∀ K n m, K ≤ n → K ≤ m →
      eLpNorm (f_sub n - f_sub m) p μ < (2⁻¹ : ℝ≥0∞) ^ K := by
    intro K n m hn hm
    calc eLpNorm (f_sub n - f_sub m) p μ
        ≤ (2⁻¹ : ℝ≥0∞) ^ (K + 1) :=
          hM_raw K (M n) (M m) (le_trans (hM_ge K) (hM_mono hn))
            (le_trans (hM_ge K) (hM_mono hm))
      _ < (2⁻¹ : ℝ≥0∞) ^ K := by
          rw [pow_succ']
          calc (2⁻¹ : ℝ≥0∞) * (2⁻¹ : ℝ≥0∞) ^ K
              < 1 * (2⁻¹ : ℝ≥0∞) ^ K :=
                ENNReal.mul_lt_mul_left
                  (ENNReal.pow_pos (by norm_num : (0 : ℝ≥0∞) < 2⁻¹) K).ne'
                  (by simp : (2⁻¹ : ℝ≥0∞) ^ K ≠ ⊤)
                  (by norm_num : (2⁻¹ : ℝ≥0∞) < 1)
            _ = (2⁻¹ : ℝ≥0∞) ^ K := one_mul _
  -- Step D: Apply cauchy_complete_eLpNorm
  have hB_sum : ∑' k, (2⁻¹ : ℝ≥0∞) ^ k ≠ ⊤ := by
    simp
  have ⟨g_lim, hg_lim_memLp, hg_lim_tendsto_sub⟩ :=
    MeasureTheory.Lp.cauchy_complete_eLpNorm hp1 hf_sub_memLp hB_sum hf_sub_cau
  -- Step E: Full convergence (same as exists_pi_limit_of_cauchy_eLpNorm)
  refine ⟨g_lim, hg_lim_memLp, ?_⟩
  rw [ENNReal.tendsto_atTop_zero]
  intro ε hε
  obtain ⟨K₁, hK₁⟩ := ENNReal.tendsto_atTop_zero.mp hg_lim_tendsto_sub (ε / 2)
    (ENNReal.half_pos hε.ne')
  obtain ⟨K₂, hK₂⟩ := ENNReal.tendsto_atTop_zero.mp
    (ENNReal.tendsto_pow_atTop_nhds_zero_of_lt_one (by norm_num : (2⁻¹ : ℝ≥0∞) < 1))
    (ε / 2) (ENNReal.half_pos hε.ne')
  let K := max K₁ K₂
  refine ⟨M K, fun n hn => ?_⟩
  have hn_ge : M_raw K ≤ n := le_trans (hM_ge K) hn
  have hMK_ge : M_raw K ≤ M K := hM_ge K
  calc eLpNorm (f n - g_lim) p μ
      = eLpNorm ((f n - f (M K)) + (f_sub K - g_lim)) p μ := by
        congr 1; ext x; simp [f_sub, sub_add_sub_cancel]
    _ ≤ eLpNorm (f n - f (M K)) p μ + eLpNorm (f_sub K - g_lim) p μ :=
        eLpNorm_add_le ((hf_memLp n).sub (hf_memLp (M K))).aestronglyMeasurable
          ((hf_sub_memLp K).sub hg_lim_memLp).aestronglyMeasurable
          hp1
    _ ≤ ε / 2 + ε / 2 := by
        gcongr
        · calc eLpNorm (f n - f (M K)) p μ
              ≤ (2⁻¹ : ℝ≥0∞) ^ (K + 1) := hM_raw K n (M K) hn_ge hMK_ge
            _ ≤ ε / 2 := hK₂ (K + 1) (Nat.le_add_right K₂ 1 |>.trans
                (Nat.add_le_add_right (le_max_right K₁ K₂) 1))
        · exact hK₁ K (le_max_left K₁ K₂)
    _ = ε := ENNReal.add_halves ε

\end{lstlisting}
\begin{proof}
Standard Riesz-Fischer construction. Extract a subsequence $f_{M(k)}$ with
$\eLpNormText{f_{M(k)} - f_{M(k+1)}}_{L^p} \le 2^{-(k+1)}$ via a controlled
$L^p$-Cauchy sub-extraction. Apply Mathlib's
\leanname{Lp.cauchy\_complete\_eLpNorm} to obtain a limit $g$ for the
subsequence. Upgrade to convergence of the full sequence by combining the
sub-sequence convergence with the original Cauchy property and a $\varepsilon/2$
triangle argument.
\end{proof}

\begin{theorem}
(\leanname{eLpNorm\_pi\_le\_sum\_component},
\leanname{eLpNorm\_pi\_component\_le},
\leanname{aestronglyMeasurable\_pi\_component},
\leanname{aestronglyMeasurable\_pi\_of\_components},
\leanname{memLp\_pi\_component})\\
For $F : \alpha \to \R^d$ and $1 \le p \le \infty$:
\begin{enumerate}
\item $\eLpNormText{F}_{L^p} \le \sum_i \eLpNormText{F_i}_{L^p}$;
\item $\eLpNormText{F_i}_{L^p} \le \eLpNormText{F}_{L^p}$;
\item $F$ is a.e.\ strongly measurable iff each component is;
\item $F \in L^p$ iff each component is.
\end{enumerate}
\end{theorem}

\begin{lstlisting}[style=Matlab-editor,basicstyle=\mlttfamily,escapechar=",mlshowsectionrules=true,mathescape=true,morecomment={[s]{\%\{}{\%\}}},language=lean,numbers=none]
/-- Vector eLpNorm ≤ sum of component eLpNorms for Pi-valued functions.
Uses `eLpNorm_mono_real` for the pointwise bound together with
`eLpNorm_sum_le` for ℝ-valued functions, avoiding Pi instance synthesis. -/
theorem eLpNorm_pi_le_sum_component
    {p : ℝ≥0∞} (hp1 : 1 ≤ p)
    {F : α → (Fin d → ℝ)}
    (hF_comp_aesm : ∀ i : Fin d, AEStronglyMeasurable (fun x => F x i) μ) :
    eLpNorm F p μ ≤ ∑ i : Fin d, eLpNorm (fun x => F x i) p μ := by
  calc eLpNorm F p μ
      ≤ eLpNorm (fun x => ∑ i : Fin d, ‖F x i‖) p μ :=
        eLpNorm_mono_real fun x => pi_norm_le_sum_norms (F x)
    _ ≤ ∑ i : Fin d, eLpNorm (fun x => ‖F x i‖) p μ := by
        let g : Fin d → α → ℝ := fun i x => ‖F x i‖
        have hg_aesm : ∀ i : Fin d, AEStronglyMeasurable (g i) μ :=
          fun i => (hF_comp_aesm i).norm
        suffices h : ∀ (s : Finset (Fin d)),
            eLpNorm (∑ i ∈ s, g i) p μ ≤ ∑ i ∈ s, eLpNorm (g i) p μ by
          have hg_eq : (fun x => ∑ i : Fin d, ‖F x i‖) = ∑ i : Fin d, g i := by
            ext x; simp [g, Finset.sum_apply]
          rw [hg_eq]; simp only [g]; exact h Finset.univ
        intro s
        induction s using Finset.cons_induction with
        | empty => simp [eLpNorm_zero]
        | cons a s has ih =>
          rw [Finset.sum_cons, Finset.sum_cons]
          calc eLpNorm (g a + ∑ i ∈ s, g i) p μ
              ≤ eLpNorm (g a) p μ + eLpNorm (∑ i ∈ s, g i) p μ :=
                eLpNorm_add_le (hg_aesm a)
                  (Finset.aestronglyMeasurable_sum s fun i _ => hg_aesm i) hp1
            _ ≤ eLpNorm (g a) p μ + ∑ i ∈ s, eLpNorm (g i) p μ := by gcongr
    _ = ∑ i : Fin d, eLpNorm (fun x => F x i) p μ := by
        congr 1; ext i; exact eLpNorm_norm _

/-- Component eLpNorm ≤ vector eLpNorm for Pi types.
For `Fin d → ℝ` with the sup norm, `‖f i‖ ≤ ‖f‖` is `norm_le_pi_norm`. -/
theorem eLpNorm_pi_component_le
    {p : ℝ≥0∞} {F : α → (Fin d → ℝ)} (i : Fin d) :
    eLpNorm (fun x => F x i) p μ ≤ eLpNorm F p μ :=
  eLpNorm_mono fun x => norm_le_pi_norm (f := F x) i

/-- Each component of a Pi-valued AEStronglyMeasurable function is AEStronglyMeasurable.
For `Fin d → ℝ` (NOT `EuclideanSpace`), `continuous_apply` works directly. -/
theorem aestronglyMeasurable_pi_component
    {F : α → (Fin d → ℝ)} (hF : AEStronglyMeasurable F μ) (i : Fin d) :
    AEStronglyMeasurable (fun x => F x i) μ :=
  Continuous.comp_aestronglyMeasurable (continuous_apply i) hF

/-- AEStronglyMeasurable for a Pi-valued function from its components. -/
theorem aestronglyMeasurable_pi_of_components
    {F : α → (Fin d → ℝ)}
    (hF_comp : ∀ i : Fin d, AEStronglyMeasurable (fun x => F x i) μ) :
    AEStronglyMeasurable F μ :=
  (aemeasurable_pi_iff.mpr fun i => (hF_comp i).aemeasurable).aestronglyMeasurable

/-- For `F : α → (Fin d → ℝ)` in Lp, each component is in Lp. -/
theorem memLp_pi_component
    {p : ℝ≥0∞} {F : α → (Fin d → ℝ)} (hF : MemLp F p μ) (i : Fin d) :
    MemLp (fun x => F x i) p μ :=
  ⟨aestronglyMeasurable_pi_component hF.aestronglyMeasurable i,
   lt_of_le_of_lt (eLpNorm_pi_component_le i) hF.eLpNorm_lt_top⟩

\end{lstlisting}
\begin{proof}
For the first item, the pointwise vector-norm bound on $\R^d$ reads
$\|F(x)\| \le \sum_i \|F(x)_i\|$ (\leanname{pi\_norm\_le\_sum\_norms}); applying
$\eLpNormText{\cdot}_{L^p}$ and using monotonicity of the eLpNorm under
pointwise inequalities of nonnegative functions
(\leanname{eLpNorm\_mono\_real}) bounds $\eLpNormText{F}_{L^p}$ above by
$\eLpNormText{x \mapsto \sum_i \|F(x)_i\|}_{L^p}$. Iterating the Minkowski
triangle inequality \leanname{eLpNorm\_add\_le} for $1 \le p \le \infty$ over
the finite index set bounds this in turn by
$\sum_i \eLpNormText{x \mapsto \|F(x)_i\|}_{L^p}$, and then
$\eLpNormText{x \mapsto \|F(x)_i\|}_{L^p} = \eLpNormText{F_i}_{L^p}$ since the
eLpNorm of $\|h\|$ equals that of $h$ (\leanname{eLpNorm\_norm}).

For the single-component bound $\eLpNormText{F_i}_{L^p} \le
\eLpNormText{F}_{L^p}$, the pointwise inequality $\|F(x)_i\| \le \|F(x)\|$
holds for any vector in $\R^d$ (the $i$-th component is dominated by the
sup norm, hence by any equivalent norm), and again $\eLpNormText{\cdot}_{L^p}$
is monotone. Strong measurability componentwise follows by composition with
the continuous projection $\R^d \to \R$ (\leanname{continuous\_apply}); the
converse direction assembles the components into a measurable function via
the product topology on $\R^d$. Membership in $L^p$ in either direction now
follows from the two norm comparisons together with strong measurability.
\end{proof}

\begin{theorem}[\leanname{tendsto\_eLpNorm\_pi\_component}]
Vector $L^p$ convergence implies componentwise $L^p$ convergence.
\end{theorem}

\begin{lstlisting}[style=Matlab-editor,basicstyle=\mlttfamily,escapechar=",mlshowsectionrules=true,mathescape=true,morecomment={[s]{\%\{}{\%\}}},language=lean,numbers=none]
/-- Vector eLpNorm convergence → component eLpNorm convergence.
Uses `eLpNorm_pi_component_le` via an explicit function equality to avoid
expensive defeq checks. -/
theorem tendsto_eLpNorm_pi_component
    {p : ℝ≥0∞}
    {G : ℕ → α → (Fin d → ℝ)} {Gext : α → (Fin d → ℝ)}
    (hG_tendsto : Tendsto (fun n => eLpNorm (fun x => G n x - Gext x) p μ) atTop (nhds 0))
    (i : Fin d) :
    Tendsto (fun n => eLpNorm (fun x => G n x i - Gext x i) p μ) atTop (nhds 0) := by
  -- Key: (fun x => G n x i - Gext x i) = (fun x => (G n x - Gext x) i)
  -- This is definitionally true for Fin d → ℝ (Pi subtraction), but we make it
  -- explicit to avoid expensive defeq checks inside tendsto_of_tendsto.
  have hle : ∀ n, eLpNorm (fun x => G n x i - Gext x i) p μ ≤
      eLpNorm (fun x => G n x - Gext x) p μ := by
    intro n
    -- Rewrite LHS to match eLpNorm_pi_component_le's input
    show eLpNorm (fun x => (G n - Gext) x i) p μ ≤ _
    exact eLpNorm_pi_component_le i
  exact tendsto_of_tendsto_of_tendsto_of_le_of_le tendsto_const_nhds hG_tendsto
    (fun _ => bot_le) hle

\end{lstlisting}
\begin{proof}
Component norms are bounded by vector norms.
\end{proof}

\begin{theorem}[\leanname{exists\_pi\_limit\_of\_cauchy\_eLpNorm}]
A Cauchy sequence $G_n : \alpha \to \R^d$ in $L^p$ admits a limit
$G_{\mathrm{ext}} \in L^p(\alpha; \R^d)$ with both vector and componentwise
$L^p$-convergence.
\end{theorem}

\begin{lstlisting}[style=Matlab-editor,basicstyle=\mlttfamily,escapechar=",mlshowsectionrules=true,mathescape=true,morecomment={[s]{\%\{}{\%\}}},language=lean,numbers=none]
/-- **Combined Cauchy → limit + components for Pi-valued sequences.**
If G n is Cauchy in eLpNorm and each G n is in Lp, then there exists a bare
function Gext in Lp with vector and component convergence.
Reduces to scalar Lp ℝ completeness per component, then reassembles. -/
theorem exists_pi_limit_of_cauchy_eLpNorm
    {p : ℝ≥0∞} (hp1 : 1 ≤ p) (hp_top : p ≠ ⊤)
    {G : ℕ → α → (Fin d → ℝ)}
    (hG_memLp : ∀ n, MemLp (G n) p μ)
    (hG_cauchy : Tendsto (fun nm : ℕ × ℕ =>
      eLpNorm (fun x => G nm.1 x - G nm.2 x) p μ) atTop (nhds 0)) :
    ∃ Gext : α → (Fin d → ℝ),
      MemLp Gext p μ ∧
      (∀ i : Fin d, MemLp (fun x => Gext x i) p μ) ∧
      Tendsto (fun n => eLpNorm (fun x => G n x - Gext x) p μ) atTop (nhds 0) ∧
      ∀ i : Fin d,
        Tendsto (fun n => eLpNorm (fun x => G n x i - Gext x i) p μ)
          atTop (nhds 0) := by
  let _ := hp_top
  haveI : Fact (1 ≤ p) := ⟨hp1⟩
  have hp : p ≠ 0 := ne_of_gt (lt_of_lt_of_le (by simp : (0 : ℝ≥0∞) < 1) hp1)
  have hcomp_cauchy : ∀ i : Fin d,
      Tendsto (fun nm : ℕ × ℕ =>
        eLpNorm (fun x => G nm.1 x i - G nm.2 x i) p μ) atTop (nhds 0) := by
    intro i
    exact tendsto_of_tendsto_of_tendsto_of_le_of_le tendsto_const_nhds hG_cauchy
      (fun _ => bot_le) (fun nm => by
        show eLpNorm ((G nm.1 - G nm.2) · i) p μ ≤
          eLpNorm (fun x => G nm.1 x - G nm.2 x) p μ
        exact eLpNorm_pi_component_le i)
  have hcomp_memLp : ∀ n i, MemLp (fun x => G n x i) p μ :=
    fun n i => memLp_pi_component (hG_memLp n) i
  -- (Avoids the Lp type entirely — pure bare-function approach)
  have hcomp_limit : ∀ i : Fin d, ∃ g_i : α → ℝ,
      MemLp g_i p μ ∧
      Tendsto (fun n => eLpNorm (fun x => G n x i - g_i x) p μ) atTop (nhds 0) := by
    intro i
    -- Component sequence converges in measure (from eLpNorm convergence)
    -- Extract a.e. convergent subsequence, define limit, show convergence
    -- Riesz-Fischer via Mathlib's Lp.cauchy_complete_eLpNorm (bare functions, no Lp type).
    -- Step A: From pair-Cauchy, get controlled bound for each ε = (2⁻¹)^(k+1)
    have hcauchy_i := hcomp_cauchy i
    have hpair : ∀ k : ℕ, ∃ M : ℕ, ∀ n m, M ≤ n → M ≤ m →
        eLpNorm (fun x => G n x i - G m x i) p μ ≤ (2⁻¹ : ℝ≥0∞) ^ (k + 1) := by
      intro k
      have hε : (0 : ℝ≥0∞) < (2⁻¹ : ℝ≥0∞) ^ (k + 1) := ENNReal.pow_pos (by norm_num) _
      obtain ⟨⟨N₁, N₂⟩, hNM⟩ := (ENNReal.tendsto_atTop_zero.mp hcauchy_i _ hε)
      exact ⟨max N₁ N₂, fun n m hn hm => hNM ⟨n, m⟩
        ⟨le_trans (le_max_left _ _) hn, le_trans (le_max_right _ _) hm⟩⟩
    -- Step B: Make bounds nondecreasing
    choose M_raw hM_raw using hpair
    let M : ℕ → ℕ := fun k => (Finset.range (k + 1)).sup M_raw
    have hM_ge : ∀ k, M_raw k ≤ M k :=
      fun k => Finset.le_sup (f := M_raw) (Finset.mem_range.mpr (Nat.lt_succ_of_le le_rfl))
    have hM_mono : Monotone M :=
      fun _ _ hab => Finset.sup_mono (Finset.range_mono (Nat.add_le_add_right hab 1))
    -- Step C: Reindexed sequence f_sub k = G (M k) · i satisfies controlled Cauchy
    let f_sub : ℕ → α → ℝ := fun k x => G (M k) x i
    have hf_sub_memLp : ∀ k, MemLp (f_sub k) p μ := fun k => hcomp_memLp (M k) i
    have hf_sub_cau : ∀ K n m, K ≤ n → K ≤ m →
        eLpNorm (f_sub n - f_sub m) p μ < (2⁻¹ : ℝ≥0∞) ^ K := by
      intro K n m hn hm
      -- eLpNorm ≤ (2⁻¹)^(K+1) < (2⁻¹)^K
      calc eLpNorm (f_sub n - f_sub m) p μ
          ≤ (2⁻¹ : ℝ≥0∞) ^ (K + 1) :=
            hM_raw K (M n) (M m) (le_trans (hM_ge K) (hM_mono hn))
              (le_trans (hM_ge K) (hM_mono hm))
        _ < (2⁻¹ : ℝ≥0∞) ^ K := by
            rw [pow_succ']
            calc (2⁻¹ : ℝ≥0∞) * (2⁻¹ : ℝ≥0∞) ^ K
                < 1 * (2⁻¹ : ℝ≥0∞) ^ K :=
                  ENNReal.mul_lt_mul_left
                    (ENNReal.pow_pos (by norm_num : (0 : ℝ≥0∞) < 2⁻¹) K).ne'
                    (by simp : (2⁻¹ : ℝ≥0∞) ^ K ≠ ⊤)
                    (by norm_num : (2⁻¹ : ℝ≥0∞) < 1)
              _ = (2⁻¹ : ℝ≥0∞) ^ K := one_mul _
    -- Step D: Apply Mathlib's cauchy_complete_eLpNorm (bare-function Riesz-Fischer)
    have hB_sum : ∑' k, (2⁻¹ : ℝ≥0∞) ^ k ≠ ⊤ := by
      simp
    obtain ⟨g_lim, hg_lim_memLp, hg_lim_tendsto_sub⟩ :=
      MeasureTheory.Lp.cauchy_complete_eLpNorm (E := ℝ) hp1 hf_sub_memLp hB_sum hf_sub_cau
    -- Step E: Full sequence G n · i → g_lim (not just the subsequence)
    -- Key: for large K, ∀ n ≥ M K:
    --   eLpNorm(G n · i - g_lim) ≤ eLpNorm(G n · i - G (M K) · i) + eLpNorm(f_sub K - g_lim)
    --   where first term ≤ (2⁻¹)^(K+1) (from hM_raw) and second ≤ ε/2 (from hg_lim_tendsto_sub)
    refine ⟨g_lim, hg_lim_memLp, ?_⟩
    rw [ENNReal.tendsto_atTop_zero]
    intro ε hε
    -- Find K₁ with eLpNorm(f_sub k - g_lim) ≤ ε/2 for k ≥ K₁
    obtain ⟨K₁, hK₁⟩ := ENNReal.tendsto_atTop_zero.mp hg_lim_tendsto_sub (ε / 2)
      (ENNReal.half_pos hε.ne')
    -- Find K₂ with (2⁻¹)^(K₂+1) ≤ ε/2
    obtain ⟨K₂, hK₂⟩ := ENNReal.tendsto_atTop_zero.mp
      (ENNReal.tendsto_pow_atTop_nhds_zero_of_lt_one (by norm_num : (2⁻¹ : ℝ≥0∞) < 1))
      (ε / 2) (ENNReal.half_pos hε.ne')
    -- Take K = max K₁ K₂, threshold = M K
    let K := max K₁ K₂
    refine ⟨M K, fun n hn => ?_⟩
    -- n ≥ M K ≥ M_raw K, and M K ≥ M_raw K, so both n and M K satisfy the Cauchy bound
    have hn_ge : M_raw K ≤ n := le_trans (hM_ge K) hn
    have hMK_ge : M_raw K ≤ M K := hM_ge K
    calc eLpNorm (fun x => G n x i - g_lim x) p μ
        = eLpNorm ((fun x => G n x i - G (M K) x i) + (fun x => G (M K) x i - g_lim x)) p μ := by
          congr 1; ext x; simp [sub_add_sub_cancel]
      _ ≤ eLpNorm (fun x => G n x i - G (M K) x i) p μ +
            eLpNorm (fun x => G (M K) x i - g_lim x) p μ :=
          eLpNorm_add_le ((hcomp_memLp n i).sub (hcomp_memLp (M K) i)).aestronglyMeasurable
            ((hcomp_memLp (M K) i).sub hg_lim_memLp).aestronglyMeasurable hp1
      _ ≤ ε / 2 + ε / 2 := by
          gcongr
          · -- First term: ≤ (2⁻¹)^(K+1) ≤ ε/2
            calc eLpNorm (fun x => G n x i - G (M K) x i) p μ
                ≤ (2⁻¹ : ℝ≥0∞) ^ (K + 1) := hM_raw K n (M K) hn_ge hMK_ge
              _ ≤ ε / 2 := hK₂ (K + 1) (Nat.le_add_right K₂ 1 |>.trans
                  (Nat.add_le_add_right (le_max_right K₁ K₂) 1))
          · -- Second term: eLpNorm(f_sub K - g_lim) ≤ ε/2
            exact hK₁ K (le_max_left K₁ K₂)
      _ = ε := ENNReal.add_halves ε
  choose g_i hg_i_memLp hg_i_tendsto using hcomp_limit
  let Gext : α → (Fin d → ℝ) := fun x i => g_i i x
  -- ‖G n x - Gext x‖_sup ≤ ∑ i |G n x i - g_i i x| (sup ≤ sum for nonneg)
  -- so eLpNorm (G n - Gext) ≤ ∑ i eLpNorm (G n · i - g_i i) → 0
  have hG_tendsto : Tendsto (fun n => eLpNorm (fun x => G n x - Gext x) p μ)
      atTop (nhds 0) := by
    have hle : ∀ n, eLpNorm (fun x => G n x - Gext x) p μ ≤
        ∑ i : Fin d, eLpNorm (fun x => G n x i - g_i i x) p μ := by
      intro n
      exact eLpNorm_pi_le_sum_component hp1
        (fun i => (hcomp_memLp n i).sub (hg_i_memLp i) |>.aestronglyMeasurable)
    have hsum : Tendsto (fun n => ∑ i : Fin d,
        eLpNorm (fun x => G n x i - g_i i x) p μ) atTop (nhds 0) := by
      have h0 : (0 : ℝ≥0∞) = ∑ _i : Fin d, (0 : ℝ≥0∞) := by simp
      rw [h0]; exact tendsto_finset_sum Finset.univ (fun i _ => hg_i_tendsto i)
    exact tendsto_of_tendsto_of_tendsto_of_le_of_le tendsto_const_nhds hsum
      (fun _ => bot_le) hle
  have hGext_aesm : AEStronglyMeasurable Gext μ :=
    aestronglyMeasurable_pi_of_components fun i => (hg_i_memLp i).aestronglyMeasurable
  have hGext_memLp : MemLp Gext p μ :=
    memLp_of_tendsto_eLpNorm hp1 hG_memLp hGext_aesm hG_tendsto
  exact ⟨Gext, hGext_memLp,
    fun i => memLp_pi_component hGext_memLp i,
    hG_tendsto,
    fun i => tendsto_eLpNorm_pi_component hG_tendsto i⟩

\end{lstlisting}
\begin{proof}
For each component, the projected sequence is scalar-valued $L^p$-Cauchy
(by the previous comparisons). Apply Theorem~\ref{thm:scalar-cauchy} to each
component, assemble the resulting limits into the vector function, and
verify the vector $L^p$ convergence by the componentwise sum bound.
\end{proof}

\begin{theorem}[\leanname{exists\_subseq\_tendsto\_ae\_of\_tendsto\_eLpNorm}]
\label{thm:ae-subseq}
Let $f_n, g \in L^2(\mu)$ with $\eLpNormText{f_n - g}_{L^2} \to 0$. There is
a subsequence $f_{n_k}$ with $f_{n_k} \to g$ a.e.
\end{theorem}

\begin{lstlisting}[style=Matlab-editor,basicstyle=\mlttfamily,escapechar=",mlshowsectionrules=true,mathescape=true,morecomment={[s]{\%\{}{\%\}}},language=lean,numbers=none]
/-! ## Indicator-Difference Convergence -/

theorem exists_subseq_tendsto_ae_of_tendsto_eLpNorm
    {α : Type*} [MeasurableSpace α] {μ : Measure α} {f : α → ℝ} {ψ : ℕ → α → ℝ}
    (hψ_aemeas : ∀ n, AEStronglyMeasurable (ψ n) μ)
    (hf_aemeas : AEStronglyMeasurable f μ)
    (hψ :
      Tendsto (fun n => eLpNorm (fun x => ψ n x - f x) 2 μ) atTop (nhds 0)) :
    ∃ ns : ℕ → ℕ, StrictMono ns ∧
      ∀ᵐ x ∂μ, Tendsto (fun n => ψ (ns n) x) atTop (nhds (f x)) := by
  have htim : TendstoInMeasure μ ψ atTop f := by
    exact MeasureTheory.tendstoInMeasure_of_tendsto_eLpNorm
      (by norm_num : (2 : ℝ≥0∞) ≠ 0) hψ_aemeas hf_aemeas hψ
  exact htim.exists_seq_tendsto_ae

\end{lstlisting}
\begin{proof}
$L^2$ convergence implies convergence in measure
(\leanname{tendstoInMeasure\_of\_tendsto\_eLpNorm}), and convergence in measure
yields an a.e.\ convergent subsequence
(\leanname{TendstoInMeasure.exists\_seq\_tendsto\_ae}).
\end{proof}

\subsection{Finite covering}

\begin{theorem}[\leanname{exists\_finite\_inner\_ball\_cover}]
Let $0 < r$, $0 < \rho$, and $r + 2\rho < R$. Then there is a finite set
$T \subseteq \overline{\Ball{r}{x_0}}$ such that
\[
\overline{\Ball{r}{x_0}} \subseteq \bigcup_{c \in T} \Ball{\rho}{c},
\quad
\overline{\Ball{2\rho}{c}} \subseteq \Ball{R}{x_0} \text{ for all } c \in T.
\]
\end{theorem}

\begin{lstlisting}[style=Matlab-editor,basicstyle=\mlttfamily,escapechar=",mlshowsectionrules=true,mathescape=true,morecomment={[s]{\%\{}{\%\}}},language=lean,numbers=none]
/-- A compact inner ball can be covered by finitely many smaller balls whose
dilated closed balls still lie inside the ambient ball. -/
theorem exists_finite_inner_ball_cover
    {d : ℕ} [NeZero d]
    {x₀ : EuclideanSpace ℝ (Fin d)}
    {r ρ R : ℝ}
    (_hr : 0 < r)
    (hρ : 0 < ρ)
    (hbuffer : r + 2 * ρ < R) :
    ∃ t : Finset (EuclideanSpace ℝ (Fin d)),
      (∀ c ∈ t, c ∈ Metric.closedBall x₀ r) ∧
      Metric.closedBall x₀ r ⊆ ⋃ c ∈ t, Metric.ball c ρ ∧
      (∀ c ∈ t, Metric.closedBall c (2 * ρ) ⊆ Metric.ball x₀ R) := by
  obtain ⟨t, ht_mem, hcover⟩ :=
    (isCompact_closedBall x₀ r).elim_nhds_subcover
      (fun x => Metric.ball x ρ)
      (fun x hx => Metric.ball_mem_nhds x hρ)
  have hbuffer' : 2 * ρ + r < R := by
    simpa [add_comm, add_left_comm, add_assoc] using hbuffer
  refine ⟨t, ht_mem, hcover, ?_⟩
  intro c hc x hx
  have hc' : dist c x₀ ≤ r := by
    simpa using ht_mem c hc
  have hx' : dist x c ≤ 2 * ρ := by
    simpa using hx
  refine Metric.mem_ball.2 ?_
  calc
    dist x x₀ ≤ dist x c + dist c x₀ := dist_triangle _ _ _
    _ ≤ 2 * ρ + r := by linarith
    _ < R := hbuffer'

\end{lstlisting}
\begin{proof}
The closed ball $\overline{\Ball{r}{x_0}}$ is compact in the finite-dimensional
Euclidean space $E$ (\leanname{isCompact\_closedBall}), and the family
$\{\Ball{\rho}{c}: c \in \overline{\Ball{r}{x_0}}\}$ is an open cover of
$\overline{\Ball{r}{x_0}}$ since each centre $c$ is contained in its own ball
of radius $\rho > 0$ (\leanname{Metric.ball\_mem\_nhds}). By compactness
(\leanname{IsCompact.elim\_nhds\_subcover}) we extract a finite subset
$T \subseteq \overline{\Ball{r}{x_0}}$ for which the union of the
$\Ball{\rho}{c}$ over $c \in T$ still covers $\overline{\Ball{r}{x_0}}$.

It remains to verify the buffered containment
$\overline{\Ball{2\rho}{c}} \subseteq \Ball{R}{x_0}$ for each $c \in T$. Since
$c$ lies in $\overline{\Ball{r}{x_0}}$ we have $\dist(c,x_0) \le r$, and any
$x \in \overline{\Ball{2\rho}{c}}$ satisfies $\dist(x,c) \le 2\rho$. By the
triangle inequality
\[
  \dist(x, x_0) \le \dist(x, c) + \dist(c, x_0) \le 2\rho + r < R,
\]
where the strict inequality is exactly the buffer hypothesis $r + 2\rho < R$.
Hence $x \in \Ball{R}{x_0}$, completing the proof.
\end{proof}

\begin{theorem}[\leanname{exists\_finite\_inner\_ball\_cover\_with\_card}]
\label{thm:finite-cover-card}
Under the same hypotheses, the cover can be chosen with cardinality bounded
by $\lceil(4r/\rho + 1)^d\rceil$.
\end{theorem}

\begin{lstlisting}[style=Matlab-editor,basicstyle=\mlttfamily,escapechar=",mlshowsectionrules=true,mathescape=true,morecomment={[s]{\%\{}{\%\}}},language=lean,numbers=none]
/-- A quantitative version of `exists_finite_inner_ball_cover` with a coarse
explicit cardinality bound depending only on `r / ρ` and `d`. -/
theorem exists_finite_inner_ball_cover_with_card
    {d : ℕ} [NeZero d]
    {x₀ : EuclideanSpace ℝ (Fin d)}
    {r ρ R : ℝ}
    (_hr : 0 < r)
    (hρ : 0 < ρ)
    (hbuffer : r + 2 * ρ < R) :
    ∃ t : Finset (EuclideanSpace ℝ (Fin d)),
      t.card ≤ Nat.ceil ((4 * r / ρ + 1) ^ d) ∧
      (∀ c ∈ t, c ∈ Metric.closedBall x₀ r) ∧
      Metric.closedBall x₀ r ⊆ ⋃ c ∈ t, Metric.ball c ρ ∧
      (∀ c ∈ t, Metric.closedBall c (2 * ρ) ⊆ Metric.ball x₀ R) := by
  classical
  have hquarter : 0 < ρ / 4 := by positivity
  obtain ⟨s, hs_mem, hs_cover, _⟩ :=
    exists_finite_inner_ball_cover
      (d := d) (x₀ := x₀) (r := r) (ρ := ρ / 4) (R := R)
      (by positivity) hquarter (by linarith)
  obtain ⟨t, ht_subset, ht_sep, ht_net⟩ :=
    exists_maximal_separated_subfinset s (δ := ρ / 2) (by positivity)
  have ht_mem : ∀ c ∈ t, c ∈ Metric.closedBall x₀ r := by
    intro c hc
    exact hs_mem c (ht_subset hc)
  have hcover : Metric.closedBall x₀ r ⊆ ⋃ c ∈ t, Metric.ball c ρ := by
    intro x hx
    have hxcover : x ∈ ⋃ c ∈ s, Metric.ball c (ρ / 4) := hs_cover hx
    rw [Set.mem_iUnion] at hxcover
    obtain ⟨a, hxcover⟩ := hxcover
    rw [Set.mem_iUnion] at hxcover
    obtain ⟨ha, hxa⟩ := hxcover
    obtain ⟨c, hc, hac⟩ := ht_net a ha
    refine Set.mem_iUnion.2 ⟨c, Set.mem_iUnion.2 ⟨hc, ?_⟩⟩
    refine Metric.mem_ball.2 ?_
    have hxa' : dist x a < ρ / 4 := by
      simpa using hxa
    calc
      dist x c ≤ dist x a + dist a c := dist_triangle _ _ _
      _ < ρ / 4 + ρ / 2 := by linarith
      _ < ρ := by linarith
  have hbuffer' : ∀ c ∈ t, Metric.closedBall c (2 * ρ) ⊆ Metric.ball x₀ R := by
    intro c hc x hx
    have hc' : dist c x₀ ≤ r := by
      simpa using ht_mem c hc
    have hx' : dist x c ≤ 2 * ρ := by
      simpa using hx
    refine Metric.mem_ball.2 ?_
    calc
      dist x x₀ ≤ dist x c + dist c x₀ := dist_triangle _ _ _
      _ ≤ 2 * ρ + r := by linarith
      _ < R := by simpa [add_comm, add_left_comm, add_assoc] using hbuffer
  have hdisj :
      (↑t : Set (EuclideanSpace ℝ (Fin d))).PairwiseDisjoint (fun c => Metric.ball c (ρ / 4)) := by
    intro c hc c' hc' hne
    change Disjoint (Metric.ball c (ρ / 4)) (Metric.ball c' (ρ / 4))
    rw [Set.disjoint_left]
    intro x hxc hxc'
    have hsep : ρ / 2 ≤ dist c c' := ht_sep hc hc' hne
    have hxc1 : dist x c < ρ / 4 := by
      simpa using hxc
    have hxc2 : dist x c' < ρ / 4 := by
      simpa using hxc'
    have hdist : dist c c' < ρ / 2 := by
      calc
        dist c c' ≤ dist c x + dist x c' := dist_triangle _ _ _
        _ < ρ / 4 + ρ / 4 := by
              simpa [dist_comm] using add_lt_add_of_lt_of_lt hxc1 hxc2
        _ = ρ / 2 := by ring
    exact False.elim ((not_lt_of_ge hsep) hdist)
  have hsmall_subset :
      (⋃ c ∈ t, Metric.ball c (ρ / 4)) ⊆ Metric.ball x₀ (r + ρ / 4) := by
    intro x hx
    rw [Set.mem_iUnion] at hx
    obtain ⟨c, hx⟩ := hx
    rw [Set.mem_iUnion] at hx
    obtain ⟨hc, hxc⟩ := hx
    have hc' : dist c x₀ ≤ r := by
      simpa using ht_mem c hc
    have hxc' : dist x c < ρ / 4 := by
      simpa using hxc
    refine Metric.mem_ball.2 ?_
    calc
      dist x x₀ ≤ dist x c + dist c x₀ := dist_triangle _ _ _
      _ < ρ / 4 + r := by linarith
      _ = r + ρ / 4 := by ring
  have hsum_le :
      ∑ c ∈ t, volume.real (Metric.ball c (ρ / 4))
        ≤ volume.real (Metric.ball x₀ (r + ρ / 4)) := by
    rw [← measureReal_biUnion_finset (μ := volume) hdisj (fun c hc => measurableSet_ball)
      (fun c hc => (measure_ball_lt_top (μ := volume) (x := c) (r := ρ / 4)).ne)]
    exact measureReal_mono hsmall_subset
  have hsum_eq :
      ∑ c ∈ t, volume.real (Metric.ball c (ρ / 4)) =
        (t.card : ℝ) *
          ((ρ / 4) ^ d * volume.real (Metric.ball (0 : EuclideanSpace ℝ (Fin d)) 1)) := by
    calc
      ∑ c ∈ t, volume.real (Metric.ball c (ρ / 4))
          = ∑ c ∈ t, ((ρ / 4) ^ d * volume.real (Metric.ball (0 : EuclideanSpace ℝ (Fin d)) 1)) := by
              refine Finset.sum_congr rfl ?_
              intro c hc
              exact volumeReal_ball_eq c hquarter
      _ = (t.card : ℝ) *
            ((ρ / 4) ^ d * volume.real (Metric.ball (0 : EuclideanSpace ℝ (Fin d)) 1)) := by
            simp [mul_left_comm]
  have hfactor : (ρ / 4) * (4 * r / ρ + 1) = r + ρ / 4 := by
    calc
      (ρ / 4) * (4 * r / ρ + 1) = (ρ / 4) * (4 * r / ρ) + ρ / 4 := by ring
      _ = r + ρ / 4 := by field_simp [hρ.ne']
  have houter_pos : 0 < r + ρ / 4 := by positivity
  have houter_eq :
      volume.real (Metric.ball x₀ (r + ρ / 4)) =
        ((4 * r / ρ + 1) ^ d) *
          ((ρ / 4) ^ d * volume.real (Metric.ball (0 : EuclideanSpace ℝ (Fin d)) 1)) := by
    rw [volumeReal_ball_eq x₀ houter_pos, ← hfactor, mul_pow]
    ring
  have hunit_pos : 0 < volume.real (Metric.ball (0 : EuclideanSpace ℝ (Fin d)) 1) := by
    exact ENNReal.toReal_pos
      (Metric.measure_ball_pos volume (0 : EuclideanSpace ℝ (Fin d)) zero_lt_one).ne'
      measure_ball_lt_top.ne
  have hconst_pos :
      0 < (ρ / 4) ^ d * volume.real (Metric.ball (0 : EuclideanSpace ℝ (Fin d)) 1) := by
    exact mul_pos (by positivity) hunit_pos
  have hcard_real : (t.card : ℝ) ≤ (4 * r / ρ + 1) ^ d := by
    rw [hsum_eq, houter_eq] at hsum_le
    nlinarith [hconst_pos]
  have hcard : t.card ≤ Nat.ceil ((4 * r / ρ + 1) ^ d) := by
    exact_mod_cast hcard_real.trans (Nat.le_ceil ((4 * r / ρ + 1) ^ d))
  exact ⟨t, hcard, ht_mem, hcover, hbuffer'⟩

\end{lstlisting}
\begin{proof}
First obtain a fine cover by $\rho/4$-balls from the previous theorem. Apply
the maximality argument
\leanname{exists\_maximal\_separated\_subfinset} to extract a $(\rho/2)$-separated
subset $T \subseteq S$. The disjoint $\rho/4$-balls around the centers of $T$
all lie inside $\Ball{r + \rho/4}{x_0}$. By volume comparison,
\[
|T|\,(\rho/4)^d\, \omega_d \le (r + \rho/4)^d \omega_d,
\quad\text{i.e.}\quad
|T| \le (4r/\rho + 1)^d.
\]
Take ceilings.
\end{proof}

\begin{theorem}[\leanname{exists\_halfBall\_cover\_by\_eighth\_balls}]
There is a finite cover of $\overline{\Ball{1/2}{0}}$ by balls of radius
$1/8$, of cardinality at most $17^d$, such that the doubled balls of radius
$1/4$ each lie inside $\Bz{1}$.
\end{theorem}

\begin{lstlisting}[style=Matlab-editor,basicstyle=\mlttfamily,escapechar=",mlshowsectionrules=true,mathescape=true,morecomment={[s]{\%\{}{\%\}}},language=lean,numbers=none]
/-- A specialized quantitative cover tailored to the half-ball / eighth-ball
geometry used in the crossover argument. -/
theorem exists_halfBall_cover_by_eighth_balls
    {d : ℕ} [NeZero d] :
    ∃ t : Finset (EuclideanSpace ℝ (Fin d)),
      t.card ≤ 17 ^ d ∧
      Metric.closedBall (0 : EuclideanSpace ℝ (Fin d)) (1 / 2 : ℝ) ⊆
        ⋃ c ∈ t, Metric.ball c (1 / 8 : ℝ) ∧
      (∀ c ∈ t,
        Metric.closedBall c (1 / 4 : ℝ) ⊆
          Metric.ball (0 : EuclideanSpace ℝ (Fin d)) 1) := by
  obtain ⟨t, htcard, _, hcover, hbuffer⟩ :=
    exists_finite_inner_ball_cover_with_card
      (d := d) (x₀ := (0 : EuclideanSpace ℝ (Fin d)))
      (r := (1 / 2 : ℝ)) (ρ := (1 / 8 : ℝ)) (R := 1)
      (by norm_num) (by norm_num) (by norm_num)
  have hpow :
      ((4 * (1 / 2 : ℝ) / (1 / 8 : ℝ) + 1) ^ d) = ((17 ^ d : ℕ) : ℝ) := by
    norm_num [Nat.cast_pow]
  have hceil :
      Nat.ceil ((4 * (1 / 2 : ℝ) / (1 / 8 : ℝ) + 1) ^ d) ≤ 17 ^ d := by
    rw [hpow, Nat.ceil_natCast]
  refine ⟨t, htcard.trans hceil, hcover, ?_⟩
  intro c hc
  convert hbuffer c hc using 1
  norm_num

\end{lstlisting}
\begin{proof}
Apply Theorem~\ref{thm:finite-cover-card} with $r = 1/2$, $\rho = 1/8$,
$R = 1$. The cardinality bound becomes $\lceil (4\cdot(1/2)/(1/8) + 1)^d
\rceil = \lceil 17^d\rceil = 17^d$.
\end{proof}

\begin{theorem}[\leanname{ae\_on\_set\_of\_ae\_on\_finite\_cover}]
If $S \subseteq \bigcup_{i \in T} U_i$ and the property $P(x)$ holds a.e.\
on each $U_i$, then $P$ holds a.e.\ on $S$.
\end{theorem}

\begin{lstlisting}[style=Matlab-editor,basicstyle=\mlttfamily,escapechar=",mlshowsectionrules=true,mathescape=true,morecomment={[s]{\%\{}{\%\}}},language=lean,numbers=none]
/-- An a.e. statement on each member of a finite cover yields an a.e. statement
on the covered set. -/
theorem ae_on_set_of_ae_on_finite_cover
    {α : Type*} [MeasurableSpace α] {μ : Measure α}
    {ι : Type*} {t : Finset ι}
    {s : Set α} {U : ι → Set α} {P : α → Prop}
    (hcover : s ⊆ ⋃ i ∈ t, U i)
    (hlocal : ∀ i ∈ t, ∀ᵐ x ∂ μ.restrict (U i), P x) :
    ∀ᵐ x ∂ μ.restrict s, P x := by
  have hunion : ∀ᵐ x ∂ μ.restrict (⋃ i ∈ t, U i), P x := by
    rw [ae_restrict_biUnion_finset_iff]
    exact hlocal
  exact ae_restrict_of_ae_restrict_of_subset hcover hunion

\end{lstlisting}
\begin{proof}
The a.e.\ statement on a finite union is equivalent to the conjunction of
a.e.\ statements on each member, then restrict to $S$.
\end{proof}

\begin{theorem}[\leanname{lintegralOn\_le\_sum\_lintegralOn\_of\_finite\_cover}]
For a non-negative measurable $f$, $\int_S f \le \sum_{i \in T} \int_{U_i} f$
whenever $S \subseteq \bigcup_{i \in T} U_i$.
\end{theorem}

\begin{lstlisting}[style=Matlab-editor,basicstyle=\mlttfamily,escapechar=",mlshowsectionrules=true,mathescape=true,morecomment={[s]{\%\{}{\%\}}},language=lean,numbers=none]
/-- The `lintegral` on a set is bounded by the sum of the `lintegral`s over any
finite cover of that set. -/
theorem lintegralOn_le_sum_lintegralOn_of_finite_cover
    {α : Type*} [MeasurableSpace α] {μ : Measure α}
    {ι : Type*} {t : Finset ι}
    {s : Set α} {U : ι → Set α} {f : α → ℝ≥0∞}
    (hcover : s ⊆ ⋃ i ∈ t, U i) :
    (∫⁻ x in s, f x ∂ μ) ≤
      Finset.sum t (fun i => ∫⁻ x in U i, f x ∂ μ) := by
  calc
    (∫⁻ x in s, f x ∂ μ) ≤ (∫⁻ x in (⋃ i ∈ t, U i), f x ∂ μ) := by
      exact lintegral_mono_set hcover
    _ ≤ Finset.sum t (fun i => ∫⁻ x in U i, f x ∂ μ) :=
      lintegral_biUnion_finset_le_sum t U f

\end{lstlisting}
\begin{proof}
Bound from above by the integral over the union, then apply countable
sub-additivity (here finite) of the integral.
\end{proof}

\subsection{Foundations and shared prelude}

\begin{remark}[\leanname{Common}, \leanname{Foundations}]
The files \leanname{DeGiorgi.Common} and \leanname{DeGiorgi.Foundations}
serve as the shared prelude for the development. They open the namespaces
\texttt{MeasureTheory}, \texttt{Metric}, \texttt{Set}, and \texttt{Filter},
together with the scoped notations for $\R_{\ge 0}^\infty$, $\R_{\ge 0}$,
$\mathrm{Topology}$, and \texttt{InnerProductSpace}, so that subsequent
files can keep their imports minimal. No mathematical content beyond
namespace setup is introduced.
\end{remark}

%% ========================================================================
\section{Elliptic Coefficients and Weak Formulation}\label{sec:weakform}
%% ========================================================================

% Section 5: Elliptic Coefficients and Weak Formulation

Throughout this section we fix a dimension $d \in \N$ with $d \neq 0$ and work in the ambient Euclidean space $E := \R^d$ (denoted \leanname{AmbientSpace} in the formalization, where it is $\mathrm{EuclideanSpace}\,\R\,(\mathrm{Fin}\,d)$). For a $d \times d$ real matrix $M$ and $\xi \in E$, the matrix-vector action is denoted $M\xi$ (\leanname{matMulE}).

\subsection{Elliptic coefficient structure}

\begin{definition}[\leanname{EllipticCoeff}]
Let $\Omega \subseteq E$. A (nonsymmetric) \emph{elliptic coefficient field} on $\Omega$ consists of:
\begin{itemize}
  \item a matrix-valued field $a : E \to \R^{d \times d}$ with each entry $x \mapsto a(x)_{ij}$ measurable;
  \item a lower ellipticity constant $\lam > 0$ (\leanname{lam});
  \item an upper ellipticity constant $\Lambda$ with $\lam \leq \Lambda$;
  \item a coercivity statement: for almost every $x \in \Omega$,
  \[
    \lam \,\|\xi\|^{2} \leq \langle \xi, a(x)\xi\rangle \quad \text{for all } \xi \in E;
  \]
  \item an inverse coercivity statement: for almost every $x \in \Omega$,
  \[
    \Lambda^{-1}\,\|\xi\|^{2} \leq \langle \xi, a(x)^{-1}\xi\rangle \quad \text{for all } \xi \in E.
  \]
\end{itemize}
\end{definition}

\begin{lstlisting}[style=Matlab-editor,basicstyle=\mlttfamily,escapechar=",mlshowsectionrules=true,mathescape=true,morecomment={[s]{\%\{}{\%\}}},language=lean,numbers=none]
/-- Unified nonsymmetric elliptic coefficient field for the elliptic regularity
development.

It is designed to serve both:

- the variational / weak-solution branch, which needs measurability;
- the De Giorgi / Moser branch, which uses coercivity of `a⁻¹` to derive mixed
  and upper bounds.

The upper bounds are intentionally not primitive fields in this structure. -/
structure EllipticCoeff (d : ℕ) [NeZero d] (Ω : Set (AmbientSpace d)) where
\end{lstlisting}
\begin{definition}[\leanname{ellipticityRatio}, \leanname{NormalizedEllipticCoeff}]
The \emph{ellipticity ratio} of $A$ is $\mathrm{ER}(A) := \Lambda/\lam$. A coefficient field is \emph{normalized} if $\lam = 1$; in that case $\mathrm{ER}(A) = \Lambda$.
\end{definition}

\begin{lstlisting}[style=Matlab-editor,basicstyle=\mlttfamily,escapechar=",mlshowsectionrules=true,mathescape=true,morecomment={[s]{\%\{}{\%\}}},language=lean,numbers=none]
  /-- Matrix-valued coefficient field. -/
  a : AmbientSpace d → Matrix (Fin d) (Fin d) ℝ
  /-- Lower ellipticity constant. -/
  lam : ℝ
  /-- Upper ellipticity constant. -/
  Λ : ℝ
  /-- Componentwise measurability of the coefficient field. -/
  measurable_comp : ∀ i j, Measurable (fun x => a x i j)
  /-- Positivity of the lower ellipticity constant. -/
  hlam : 0 < lam
  /-- Comparison of lower and upper ellipticity scales. -/
  hΛ : lam ≤ Λ
  /-- Coercivity on `a`, stated almost everywhere on `Ω`. -/
  coercive : ∀ᵐ x ∂(MeasureTheory.volume.restrict Ω), ∀ ξ : AmbientSpace d,
    lam * ‖ξ‖ ^ 2 ≤ ⟪ξ, matMulE (a x) ξ⟫_ℝ
  /-- Coercivity on `a⁻¹`, stated almost everywhere on `Ω`. -/
  coercive_inv : ∀ᵐ x ∂(MeasureTheory.volume.restrict Ω), ∀ ξ : AmbientSpace d,
    Λ⁻¹ * ‖ξ‖ ^ 2 ≤ ⟪ξ, matMulE ((a x)⁻¹) ξ⟫_ℝ

/-- The ellipticity ratio governing scale-invariant estimates. -/
def ellipticityRatio {d : ℕ} [NeZero d] {Ω : Set (AmbientSpace d)}
    (A : EllipticCoeff d Ω) : ℝ :=
  A.Λ / A.lam

/-- Optional normalized view of the unified coefficient structure. -/
abbrev NormalizedEllipticCoeff (d : ℕ) [NeZero d] (Ω : Set (AmbientSpace d)) :=
  {A : EllipticCoeff d Ω // A.lam = 1}

\end{lstlisting}
\begin{lemma}[\leanname{EllipticCoeff.lam\_nonneg}, $\Lambda$\texttt{\_pos}, \leanname{ellipticityRatio\_pos}, \leanname{one\_le\_ellipticityRatio}]
For any $A$ as above one has $\lam \geq 0$, $\Lambda > 0$, $\mathrm{ER}(A) > 0$, and $1 \leq \mathrm{ER}(A)$.
\end{lemma}

\begin{lstlisting}[style=Matlab-editor,basicstyle=\mlttfamily,escapechar=",mlshowsectionrules=true,mathescape=true,morecomment={[s]{\%\{}{\%\}}},language=lean,numbers=none]
theorem lam_nonneg : 0 ≤ A.lam :=
  le_of_lt A.hlam

theorem ellipticityRatio_pos : 0 < ellipticityRatio A := by
  dsimp [ellipticityRatio]
  exact div_pos A.Λ_pos A.hlam

theorem one_le_ellipticityRatio : 1 ≤ ellipticityRatio A := by
  simpa [ellipticityRatio] using (one_le_div A.hlam).2 A.hΛ

\end{lstlisting}
\begin{proof}
Positivity of $\lam$ gives $\lam \geq 0$. From $0 < \lam \leq \Lambda$ we get $\Lambda > 0$ and $\Lambda \geq 0$. The ratio $\Lambda/\lam$ is then a quotient of positive reals, hence positive; from $\lam \leq \Lambda$ and $\lam > 0$ we get $1 \leq \Lambda/\lam$.
\end{proof}

\begin{lemma}[\leanname{ae\_coercive\_nonneg}, \leanname{ae\_coercive\_inv\_nonneg}]
For a.e.\ $x \in \Omega$ and every $\xi \in E$, both $\langle \xi, a(x)\xi\rangle$ and $\langle \xi, a(x)^{-1}\xi\rangle$ are nonnegative.
\end{lemma}

\begin{lstlisting}[style=Matlab-editor,basicstyle=\mlttfamily,escapechar=",mlshowsectionrules=true,mathescape=true,morecomment={[s]{\%\{}{\%\}}},language=lean,numbers=none]
theorem ae_coercive_nonneg :
    ∀ᵐ x ∂(MeasureTheory.volume.restrict Ω), ∀ ξ : AmbientSpace d,
      0 ≤ ⟪ξ, matMulE (A.a x) ξ⟫_ℝ := by
  filter_upwards [A.coercive] with x hx
  intro ξ
  have hnonneg : 0 ≤ A.lam * ‖ξ‖ ^ 2 := by
    exact mul_nonneg A.lam_nonneg (sq_nonneg ‖ξ‖)
  exact le_trans hnonneg (hx ξ)

theorem ae_coercive_inv_nonneg :
    ∀ᵐ x ∂(MeasureTheory.volume.restrict Ω), ∀ ξ : AmbientSpace d,
      0 ≤ ⟪ξ, matMulE ((A.a x)⁻¹) ξ⟫_ℝ := by
  filter_upwards [A.coercive_inv] with x hx
  intro ξ
  have hnonneg : 0 ≤ A.Λ⁻¹ * ‖ξ‖ ^ 2 := by
    exact mul_nonneg (inv_nonneg.mpr A.Λ_nonneg) (sq_nonneg ‖ξ‖)
  exact le_trans hnonneg (hx ξ)

\end{lstlisting}
\begin{proof}
Apply coercivity (resp.\ inverse coercivity) and observe that $\lam\|\xi\|^2 \geq 0$ and $\Lambda^{-1}\|\xi\|^2 \geq 0$.
\end{proof}

\begin{lemma}[\leanname{det\_ne\_zero\_of\_coercive}, \leanname{inv\_matMulE\_matMulE}]
At any point $x$ where the lower coercivity bound holds, $\det a(x) \neq 0$ and $a(x)^{-1}\bigl(a(x)\xi\bigr) = \xi$ for all $\xi \in E$.
\end{lemma}

\begin{lstlisting}[style=Matlab-editor,basicstyle=\mlttfamily,escapechar=",mlshowsectionrules=true,mathescape=true,morecomment={[s]{\%\{}{\%\}}},language=lean,numbers=none]
theorem det_ne_zero_of_coercive {x : AmbientSpace d}
    (hx : ∀ ξ : AmbientSpace d,
      A.lam * ‖ξ‖ ^ 2 ≤ ⟪ξ, matMulE (A.a x) ξ⟫_ℝ) :
    (A.a x).det ≠ 0 := by
  classical
  rw [ne_eq, ← Matrix.exists_mulVec_eq_zero_iff]
  intro hsing
  rcases hsing with ⟨v, hv_ne, hv_zero⟩
  let ξ : AmbientSpace d := WithLp.toLp 2 v
  have hξ_zero_action : matMulE (A.a x) ξ = 0 := by
    apply PiLp.ext
    intro i
    simp [matMulE, ξ, hv_zero]
  have hcoer : A.lam * ‖ξ‖ ^ 2 ≤ ⟪ξ, matMulE (A.a x) ξ⟫_ℝ := hx ξ
  have hmul_nonneg : 0 ≤ A.lam * ‖ξ‖ ^ 2 := by
    exact mul_nonneg A.lam_nonneg (sq_nonneg ‖ξ‖)
  have hmul_eq_zero : A.lam * ‖ξ‖ ^ 2 = 0 := by
    refine le_antisymm ?_ hmul_nonneg
    simpa [hξ_zero_action] using hcoer
  have hnorm_sq_eq_zero : ‖ξ‖ ^ 2 = 0 := by
    exact (mul_eq_zero.mp hmul_eq_zero).resolve_left (ne_of_gt A.hlam)
  have hnorm_eq_zero : ‖ξ‖ = 0 := by
    nlinarith [sq_nonneg ‖ξ‖, hnorm_sq_eq_zero]
  have hξ_zero' : ξ = 0 := norm_eq_zero.mp hnorm_eq_zero
  have hv_zero' : v = 0 := by
    have hcoords : ξ.ofLp = 0 := by
      simpa [ξ] using congrArg (fun η : AmbientSpace d => η.ofLp) hξ_zero'
    simpa [ξ] using hcoords
  exact hv_ne hv_zero'

theorem inv_matMulE_matMulE {x : AmbientSpace d}
    (hx : ∀ ξ : AmbientSpace d,
      A.lam * ‖ξ‖ ^ 2 ≤ ⟪ξ, matMulE (A.a x) ξ⟫_ℝ)
    (ξ : AmbientSpace d) :
    matMulE ((A.a x)⁻¹) (matMulE (A.a x) ξ) = ξ := by
  apply PiLp.ext
  intro i
  have hdet_unit : IsUnit (A.a x).det :=
    (isUnit_iff_ne_zero.2 <| A.det_ne_zero_of_coercive hx)
  simp [matMulE, Matrix.mulVec_mulVec, Matrix.nonsing_inv_mul _ hdet_unit, Matrix.one_mulVec]

\end{lstlisting}
\begin{proof}
If $\det a(x) = 0$, choose a nonzero $v$ with $a(x) v = 0$. Then $\lam\|v\|^2 \leq \langle v, a(x)v\rangle = 0$, forcing $\|v\| = 0$ since $\lam > 0$, a contradiction. Once $a(x)$ is invertible, the identity $a(x)^{-1}a(x) = I$ gives the second claim.
\end{proof}

\begin{lemma}[\leanname{EllipticCoeff.mulVec\_sq\_le}: mixed quadratic bound]
For a.e.\ $x \in \Omega$ and all $\xi \in E$,
\[
  \|a(x)\xi\|^{2} \leq \Lambda \,\langle \xi, a(x)\xi\rangle.
\]
\end{lemma}

\begin{lstlisting}[style=Matlab-editor,basicstyle=\mlttfamily,escapechar=",mlshowsectionrules=true,mathescape=true,morecomment={[s]{\%\{}{\%\}}},language=lean,numbers=none]
/-- Mixed quadratic bound derived from inverse coercivity. -/
theorem mulVec_sq_le :
    ∀ᵐ x ∂(MeasureTheory.volume.restrict Ω), ∀ ξ : AmbientSpace d,
      ‖matMulE (A.a x) ξ‖ ^ 2 ≤ A.Λ * ⟪ξ, matMulE (A.a x) ξ⟫_ℝ := by
  filter_upwards [A.coercive, A.coercive_inv] with x hx_coercive hx_inv
  intro ξ
  let η : AmbientSpace d := matMulE (A.a x) ξ
  have hscaled :
      ‖η‖ ^ 2 ≤ A.Λ * ⟪η, matMulE ((A.a x)⁻¹) η⟫_ℝ := by
    have h := mul_le_mul_of_nonneg_left (hx_inv η) A.Λ_nonneg
    have hΛ_ne : A.Λ ≠ 0 := ne_of_gt A.Λ_pos
    simpa [η, mul_assoc, hΛ_ne] using h
  calc
    ‖matMulE (A.a x) ξ‖ ^ 2 = ‖η‖ ^ 2 := by
      rfl
    _ ≤ A.Λ * ⟪η, matMulE ((A.a x)⁻¹) η⟫_ℝ := hscaled
    _ = A.Λ * ⟪η, ξ⟫_ℝ := by
      rw [A.inv_matMulE_matMulE hx_coercive ξ]
    _ = A.Λ * ⟪ξ, η⟫_ℝ := by
      rw [real_inner_comm]
    _ = A.Λ * ⟪ξ, matMulE (A.a x) ξ⟫_ℝ := by
      rfl

\end{lstlisting}
\begin{proof}
Set $\eta := a(x)\xi$. Inverse coercivity applied to $\eta$ gives $\Lambda^{-1}\|\eta\|^{2} \leq \langle \eta, a(x)^{-1}\eta\rangle$, hence $\|\eta\|^{2} \leq \Lambda \langle \eta, a(x)^{-1}\eta\rangle$. Since $a(x)^{-1}\eta = \xi$ by the previous lemma, the right-hand side equals $\Lambda\langle \eta, \xi\rangle = \Lambda\langle \xi, a(x)\xi\rangle$.
\end{proof}

\begin{lemma}[\leanname{EllipticCoeff.quadratic\_upper}]
For a.e.\ $x \in \Omega$ and all $\xi \in E$,
\[
  \langle \xi, a(x)\xi\rangle \leq \Lambda\,\|\xi\|^{2}.
\]
\end{lemma}

\begin{lstlisting}[style=Matlab-editor,basicstyle=\mlttfamily,escapechar=",mlshowsectionrules=true,mathescape=true,morecomment={[s]{\%\{}{\%\}}},language=lean,numbers=none]
/-- Pointwise quadratic upper bound derived from the mixed bound. -/
theorem quadratic_upper :
    ∀ᵐ x ∂(MeasureTheory.volume.restrict Ω), ∀ ξ : AmbientSpace d,
      ⟪ξ, matMulE (A.a x) ξ⟫_ℝ ≤ A.Λ * ‖ξ‖ ^ 2 := by
  filter_upwards [A.coercive, A.mulVec_sq_le] with x hx_coercive hx_mul ξ
  let q : ℝ := ⟪ξ, matMulE (A.a x) ξ⟫_ℝ
  let m : ℝ := ‖matMulE (A.a x) ξ‖
  let n : ℝ := ‖ξ‖
  have hq_nonneg : 0 ≤ q := by
    exact le_trans (mul_nonneg A.lam_nonneg (sq_nonneg ‖ξ‖)) (by simpa [q] using hx_coercive ξ)
  have hq_le : q ≤ n * m := by
    simpa [q, m, n, abs_of_nonneg hq_nonneg, mul_comm] using
      (abs_real_inner_le_norm ξ (matMulE (A.a x) ξ))
  have hm_sq : m ^ 2 ≤ A.Λ * q := by
    simpa [m, q] using hx_mul ξ
  have hn_nonneg : 0 ≤ n := norm_nonneg ξ
  have hm_nonneg : 0 ≤ m := norm_nonneg (matMulE (A.a x) ξ)
  have hq_mul : q * q ≤ (A.Λ * n ^ 2) * q := by
    nlinarith [hq_le, hm_sq, hn_nonneg, hm_nonneg, A.Λ_nonneg]
  have hupper : q ≤ A.Λ * n ^ 2 := by
    by_cases hq_zero : q = 0
    · have hnonneg : 0 ≤ A.Λ * n ^ 2 := mul_nonneg A.Λ_nonneg (sq_nonneg n)
      simpa [hq_zero] using hnonneg
    · have hq_pos : 0 < q := lt_of_le_of_ne hq_nonneg (by simpa [eq_comm] using hq_zero)
      exact le_of_mul_le_mul_right
        (by simpa [pow_two, mul_assoc, mul_left_comm, mul_comm] using hq_mul)
        hq_pos
  simpa [q, n] using hupper

\end{lstlisting}
\begin{proof}
Let $q := \langle \xi, a(x)\xi\rangle$, $m := \|a(x)\xi\|$, $n := \|\xi\|$. Cauchy--Schwarz gives $q \leq n m$, and the previous lemma gives $m^{2} \leq \Lambda q$. If $q = 0$ the bound is trivial. Otherwise multiply $q \leq nm$ by $q$ to obtain $q^{2} \leq n^{2} m^{2} \leq \Lambda n^{2} q$, and divide by $q > 0$.
\end{proof}

\begin{lemma}[\leanname{EllipticCoeff.mixed\_bound}]
For a.e.\ $x \in \Omega$ and all $\xi, \zeta \in E$,
\[
  |\langle \zeta, a(x)\xi\rangle| \leq \Lambda\,\|\zeta\|\,\|\xi\|.
\]
\end{lemma}

\begin{lstlisting}[style=Matlab-editor,basicstyle=\mlttfamily,escapechar=",mlshowsectionrules=true,mathescape=true,morecomment={[s]{\%\{}{\%\}}},language=lean,numbers=none]
/-- Mixed bilinear bound needed by the variational branch. -/
theorem mixed_bound :
    ∀ᵐ x ∂(MeasureTheory.volume.restrict Ω), ∀ ξ ζ : AmbientSpace d,
      |⟪ζ, matMulE (A.a x) ξ⟫_ℝ| ≤ A.Λ * ‖ζ‖ * ‖ξ‖ := by
  filter_upwards [A.mulVec_sq_le, A.quadratic_upper] with x hx_mul hx_quad
  intro ξ ζ
  let m : ℝ := ‖matMulE (A.a x) ξ‖
  let nξ : ℝ := ‖ξ‖
  let nζ : ℝ := ‖ζ‖
  have habs : |⟪ζ, matMulE (A.a x) ξ⟫_ℝ| ≤ nζ * m := by
    simpa [m, nζ, mul_comm] using abs_real_inner_le_norm ζ (matMulE (A.a x) ξ)
  have hm_sq : m ^ 2 ≤ A.Λ * ⟪ξ, matMulE (A.a x) ξ⟫_ℝ := by
    simpa [m] using hx_mul ξ
  have hq_upper :
      ⟪ξ, matMulE (A.a x) ξ⟫_ℝ ≤ A.Λ * nξ ^ 2 := by
    simpa [nξ] using hx_quad ξ
  have hm_le : m ≤ A.Λ * nξ := by
    have hm_nonneg : 0 ≤ m := norm_nonneg (matMulE (A.a x) ξ)
    have hnξ_nonneg : 0 ≤ nξ := norm_nonneg ξ
    have hrhs_nonneg : 0 ≤ A.Λ * nξ := mul_nonneg A.Λ_nonneg hnξ_nonneg
    have hm_sq_upper : m ^ 2 ≤ (A.Λ * nξ) ^ 2 := by
      nlinarith [hm_sq, hq_upper, A.Λ_nonneg]
    exact (sq_le_sq₀ hm_nonneg hrhs_nonneg).1 hm_sq_upper
  have hnζ_nonneg : 0 ≤ nζ := norm_nonneg ζ
  calc
    |⟪ζ, matMulE (A.a x) ξ⟫_ℝ| ≤ nζ * m := habs
    _ ≤ nζ * (A.Λ * nξ) := by
      exact mul_le_mul_of_nonneg_left hm_le hnζ_nonneg
    _ = A.Λ * ‖ζ‖ * ‖ξ‖ := by
      ring

\end{lstlisting}
\begin{proof}
Cauchy--Schwarz gives $|\langle \zeta, a(x)\xi\rangle| \leq \|\zeta\|\,\|a(x)\xi\|$, so it suffices to show $\|a(x)\xi\| \leq \Lambda\|\xi\|$. Squaring, $\|a(x)\xi\|^{2} \leq \Lambda \langle \xi, a(x)\xi\rangle \leq \Lambda^{2}\|\xi\|^{2}$ by the previous two lemmas, and taking square roots (using nonnegativity of both sides) gives the desired estimate.
\end{proof}

\subsection{Pointwise bilinear form on Sobolev witnesses}

Throughout, $u, v \in W^{1,2}(\Omega)$ are equipped with weak-gradient \emph{witnesses} $h_u, h_v$ in the sense of Section 4. We write $\nabla u$ for $h_u.\mathrm{weakGrad}$.

\begin{definition}[\leanname{bilinFormIntegrandOfCoeff}, \leanname{bilinFormOfCoeff}]
The pointwise bilinear-form integrand and the bilinear form are
\[
  b_A(u,v)(x) := \langle a(x)\,\nabla u(x), \nabla v(x)\rangle, \qquad
  \mathfrak{a}_A(u,v) := \int_{\Omega} b_A(u,v)(x)\,dx.
\]
\end{definition}

\begin{lstlisting}[style=Matlab-editor,basicstyle=\mlttfamily,escapechar=",mlshowsectionrules=true,mathescape=true,morecomment={[s]{\%\{}{\%\}}},language=lean,numbers=none]
/-- Pointwise bilinear-form integrand associated to an elliptic coefficient
field and two Sobolev witnesses. -/
def bilinFormIntegrandOfCoeff {Ω : Set E} {u v : E → ℝ}
    (A : EllipticCoeff d Ω)
    (hu : MemW1pWitness 2 u Ω) (hv : MemW1pWitness 2 v Ω) : E → ℝ :=
  fun x => ⟪matMulE (A.a x) (hu.weakGrad x), hv.weakGrad x⟫_ℝ

/-- The weak bilinear form `a_A(u,v) = ∫_Ω ⟪A∇u, ∇v⟫`. -/
noncomputable def bilinFormOfCoeff {Ω : Set E} {u v : E → ℝ}
    (A : EllipticCoeff d Ω)
    (hu : MemW1pWitness 2 u Ω) (hv : MemW1pWitness 2 v Ω) : ℝ :=
  ∫ x in Ω, bilinFormIntegrandOfCoeff A hu hv x

\end{lstlisting}
\begin{lemma}[\leanname{aestronglyMeasurable\_bilinFormIntegrandOfCoeff}, \leanname{integrable\_bilinFormIntegrandOfCoeff}]
The integrand $b_A(u,v)$ is a.e.\ strongly measurable and integrable on $\Omega$.
\end{lemma}

\begin{lstlisting}[style=Matlab-editor,basicstyle=\mlttfamily,escapechar=",mlshowsectionrules=true,mathescape=true,morecomment={[s]{\%\{}{\%\}}},language=lean,numbers=none]
/-- The bilinear-form integrand is a.e. strongly measurable on `Ω`. -/
theorem aestronglyMeasurable_bilinFormIntegrandOfCoeff
    {Ω : Set E} {u v : E → ℝ}
    (A : EllipticCoeff d Ω)
    (hu : MemW1pWitness 2 u Ω) (hv : MemW1pWitness 2 v Ω) :
    AEStronglyMeasurable (bilinFormIntegrandOfCoeff A hu hv) (volume.restrict Ω) := by
  have hF_aemeas :
      AEMeasurable
        (fun x => ∑ i : Fin d, (∑ j : Fin d, A.a x i j * hu.weakGrad x j) * hv.weakGrad x i)
        (volume.restrict Ω) := by
    have hsum :
        AEMeasurable
          (∑ i : Fin d, fun x => (∑ j : Fin d, A.a x i j * hu.weakGrad x j) * hv.weakGrad x i)
          (volume.restrict Ω) := by
      refine Finset.aemeasurable_sum (s := (Finset.univ : Finset (Fin d)))
        (f := fun i x => (∑ j : Fin d, A.a x i j * hu.weakGrad x j) * hv.weakGrad x i) ?_
      intro i _
      have hrow :
          AEMeasurable (fun x => ∑ j : Fin d, A.a x i j * hu.weakGrad x j)
            (volume.restrict Ω) := by
        have hrow_sum :
            AEMeasurable (∑ j : Fin d, fun x => A.a x i j * hu.weakGrad x j)
              (volume.restrict Ω) := by
          refine Finset.aemeasurable_sum (s := (Finset.univ : Finset (Fin d)))
            (f := fun j x => A.a x i j * hu.weakGrad x j) ?_
          intro j _
          exact (A.measurable_apply i j).aemeasurable.mul
            (hu.weakGrad_component_memLp j).aemeasurable
        convert hrow_sum using 1
        ext x
        simp
      exact hrow.mul (hv.weakGrad_component_memLp i).aemeasurable
    convert hsum using 1
    ext x
    simp
  have hscalar : ∀ a b : ℝ, ⟪a, b⟫_ℝ = a * b := by
    intro a b
    simpa using (RCLike.inner_apply' a b)
  refine hF_aemeas.aestronglyMeasurable.congr ?_
  filter_upwards with x
  simp [bilinFormIntegrandOfCoeff, PiLp.inner_apply, matMulE_apply, Matrix.mulVec,
    dotProduct, hscalar]

/-- The bilinear-form integrand is integrable on `Ω`. -/
theorem integrable_bilinFormIntegrandOfCoeff
    {Ω : Set E} {u v : E → ℝ}
    (A : EllipticCoeff d Ω)
    (hu : MemW1pWitness 2 u Ω) (hv : MemW1pWitness 2 v Ω) :
    Integrable (bilinFormIntegrandOfCoeff A hu hv) (volume.restrict Ω) := by
  let μ : Measure E := volume.restrict Ω
  have hprod_memLp :
      MemLp (fun x => ‖hu.weakGrad x‖ * ‖hv.weakGrad x‖) 1 μ := by
    simpa using (hv.weakGrad_memLp.norm.mul hu.weakGrad_memLp.norm)
  have hdom_int :
      Integrable (fun x => A.Λ * (‖hu.weakGrad x‖ * ‖hv.weakGrad x‖)) μ := by
    rw [← memLp_one_iff_integrable]
    exact hprod_memLp.const_mul A.Λ
  have hbound :
      ∀ᵐ x ∂μ, ‖bilinFormIntegrandOfCoeff A hu hv x‖ ≤
        A.Λ * (‖hu.weakGrad x‖ * ‖hv.weakGrad x‖) := by
    filter_upwards [A.mixed_bound] with x hx
    simpa [bilinFormIntegrandOfCoeff, real_inner_comm, mul_assoc, mul_comm, mul_left_comm] using
      hx (hu.weakGrad x) (hv.weakGrad x)
  exact Integrable.mono' hdom_int
    (aestronglyMeasurable_bilinFormIntegrandOfCoeff A hu hv)
    hbound

\end{lstlisting}
\begin{proof}
Measurability follows from the componentwise measurability of $a$ and the $\Lp{2}$-membership of the components of $\nabla u, \nabla v$. Integrability follows from the pointwise mixed bound $|b_A(u,v)| \leq \Lambda\,\|\nabla u\|\,\|\nabla v\|$ and the Cauchy--Schwarz product $\|\nabla u\|\,\|\nabla v\| \in \Lp{1}$.
\end{proof}

\begin{lemma}[\leanname{bilinForm\_bound}: continuity]
\[
  |\mathfrak{a}_A(u,v)| \leq \Lambda
  \Bigl(\int_{\Omega} \|\nabla u\|^{2}\,dx\Bigr)^{1/2}
  \Bigl(\int_{\Omega} \|\nabla v\|^{2}\,dx\Bigr)^{1/2}.
\]
\end{lemma}

\begin{lstlisting}[style=Matlab-editor,basicstyle=\mlttfamily,escapechar=",mlshowsectionrules=true,mathescape=true,morecomment={[s]{\%\{}{\%\}}},language=lean,numbers=none]
/-- The weak bilinear form is continuous with respect to the `L²` gradient
seminorms. -/
theorem bilinForm_bound
    {Ω : Set E} {u v : E → ℝ}
    (A : EllipticCoeff d Ω)
    (hu : MemW1pWitness 2 u Ω) (hv : MemW1pWitness 2 v Ω) :
    |bilinFormOfCoeff A hu hv| ≤
      A.Λ * (∫ x, ‖hu.weakGrad x‖ ^ (2 : ℝ) ∂(volume.restrict Ω)) ^ (1 / (2 : ℝ)) *
        (∫ x, ‖hv.weakGrad x‖ ^ (2 : ℝ) ∂(volume.restrict Ω)) ^ (1 / (2 : ℝ)) := by
  let μ : Measure E := volume.restrict Ω
  have hint_int := integrable_bilinFormIntegrandOfCoeff A hu hv
  have hdom_int :
      Integrable (fun x => A.Λ * (‖hu.weakGrad x‖ * ‖hv.weakGrad x‖)) μ := by
    rw [← memLp_one_iff_integrable]
    have hprod_memLp :
        MemLp (fun x => ‖hu.weakGrad x‖ * ‖hv.weakGrad x‖) 1 μ := by
      simpa using (hv.weakGrad_memLp.norm.mul hu.weakGrad_memLp.norm)
    exact hprod_memLp.const_mul A.Λ
  have hpointwise :
      ∀ᵐ x ∂μ, ‖bilinFormIntegrandOfCoeff A hu hv x‖ ≤
        A.Λ * (‖hu.weakGrad x‖ * ‖hv.weakGrad x‖) := by
    filter_upwards [A.mixed_bound] with x hx
    simpa [bilinFormIntegrandOfCoeff, real_inner_comm, mul_assoc, mul_comm, mul_left_comm] using
      hx (hu.weakGrad x) (hv.weakGrad x)
  have hnorm_le :
      ∫ x, ‖bilinFormIntegrandOfCoeff A hu hv x‖ ∂μ ≤
        ∫ x, A.Λ * (‖hu.weakGrad x‖ * ‖hv.weakGrad x‖) ∂μ := by
    exact integral_mono_ae hint_int.norm hdom_int hpointwise
  have hholder :
      ∫ x, ‖hu.weakGrad x‖ * ‖hv.weakGrad x‖ ∂μ ≤
        (∫ x, ‖hu.weakGrad x‖ ^ (2 : ℝ) ∂μ) ^ (1 / (2 : ℝ)) *
          (∫ x, ‖hv.weakGrad x‖ ^ (2 : ℝ) ∂μ) ^ (1 / (2 : ℝ)) := by
    have hhu_memLp : MemLp hu.weakGrad (ENNReal.ofReal (2 : ℝ)) μ := by
      simpa using hu.weakGrad_memLp
    have hhv_memLp : MemLp hv.weakGrad (ENNReal.ofReal (2 : ℝ)) μ := by
      simpa using hv.weakGrad_memLp
    exact
      integral_mul_norm_le_Lp_mul_Lq (μ := μ) (f := hu.weakGrad) (g := hv.weakGrad)
        Real.HolderConjugate.two_two hhu_memLp hhv_memLp
  simpa [μ, mul_assoc] using
    (calc
      |bilinFormOfCoeff A hu hv| = ‖∫ x, bilinFormIntegrandOfCoeff A hu hv x ∂μ‖ := by
        simp [bilinFormOfCoeff, μ]
      _ ≤ ∫ x, ‖bilinFormIntegrandOfCoeff A hu hv x‖ ∂μ :=
        norm_integral_le_integral_norm _
      _ ≤ ∫ x, A.Λ * (‖hu.weakGrad x‖ * ‖hv.weakGrad x‖) ∂μ := hnorm_le
      _ = A.Λ * ∫ x, ‖hu.weakGrad x‖ * ‖hv.weakGrad x‖ ∂μ := by
        rw [integral_const_mul]
      _ ≤ A.Λ *
          ((∫ x, ‖hu.weakGrad x‖ ^ (2 : ℝ) ∂μ) ^ (1 / (2 : ℝ)) *
            (∫ x, ‖hv.weakGrad x‖ ^ (2 : ℝ) ∂μ) ^ (1 / (2 : ℝ))) := by
        exact mul_le_mul_of_nonneg_left hholder A.Λ_nonneg
      _ = A.Λ * (∫ x, ‖hu.weakGrad x‖ ^ (2 : ℝ) ∂μ) ^ (1 / (2 : ℝ)) *
          (∫ x, ‖hv.weakGrad x‖ ^ (2 : ℝ) ∂μ) ^ (1 / (2 : ℝ)) := by
        ring)

\end{lstlisting}
\begin{proof}
By the triangle inequality for integrals and the pointwise bound $|b_A(u,v)| \leq \Lambda\,\|\nabla u\|\,\|\nabla v\|$,
\[
  |\mathfrak{a}_A(u,v)| \leq \Lambda \int_{\Omega} \|\nabla u\|\,\|\nabla v\|\,dx.
\]
H\"older's inequality with conjugate exponents $(2,2)$ bounds the right-hand side by $\Lambda \,\|\nabla u\|_{\Lp{2}}\,\|\nabla v\|_{\Lp{2}}$.
\end{proof}

\begin{lemma}[\leanname{bilinForm\_coercive}]
\[
  \lam \int_{\Omega} \|\nabla u\|^{2}\,dx \leq \mathfrak{a}_A(u,u).
\]
\end{lemma}

\begin{lstlisting}[style=Matlab-editor,basicstyle=\mlttfamily,escapechar=",mlshowsectionrules=true,mathescape=true,morecomment={[s]{\%\{}{\%\}}},language=lean,numbers=none]
/-- The weak bilinear form is coercive with respect to the `L²` gradient
seminorm. -/
theorem bilinForm_coercive
    {Ω : Set E} {u : E → ℝ}
    (A : EllipticCoeff d Ω)
    (hu : MemW1pWitness 2 u Ω) :
    A.lam * ∫ x, ‖hu.weakGrad x‖ ^ (2 : ℕ) ∂(volume.restrict Ω) ≤
      bilinFormOfCoeff A hu hu := by
  let μ : Measure E := volume.restrict Ω
  have hint_int := integrable_bilinFormIntegrandOfCoeff A hu hu
  have hsq_int : Integrable (fun x => ‖hu.weakGrad x‖ ^ (2 : ℕ)) μ := by
    simpa using
      (MeasureTheory.MemLp.integrable_norm_pow (f := hu.weakGrad) (μ := μ)
        hu.weakGrad_memLp (by norm_num : (2 : ℕ) ≠ 0))
  have hleft_int : Integrable (fun x => A.lam * ‖hu.weakGrad x‖ ^ (2 : ℕ)) μ := by
    simpa using hsq_int.const_mul A.lam
  have hcoercive_ae :
      ∀ᵐ x ∂μ, A.lam * ‖hu.weakGrad x‖ ^ (2 : ℕ) ≤
        bilinFormIntegrandOfCoeff A hu hu x := by
    filter_upwards [A.coercive] with x hx
    simpa [bilinFormIntegrandOfCoeff, real_inner_comm] using hx (hu.weakGrad x)
  have hmono :
      ∫ x, A.lam * ‖hu.weakGrad x‖ ^ (2 : ℕ) ∂μ ≤
        ∫ x, bilinFormIntegrandOfCoeff A hu hu x ∂μ := by
    exact integral_mono_ae hleft_int hint_int hcoercive_ae
  simpa [μ] using
    (calc
      A.lam * ∫ x, ‖hu.weakGrad x‖ ^ (2 : ℕ) ∂μ
          = ∫ x, A.lam * ‖hu.weakGrad x‖ ^ (2 : ℕ) ∂μ := by
            rw [integral_const_mul]
      _ ≤ ∫ x, bilinFormIntegrandOfCoeff A hu hu x ∂μ := hmono
      _ = bilinFormOfCoeff A hu hu := by
            simp [bilinFormOfCoeff, μ])

\end{lstlisting}
\begin{proof}
Coercivity gives $\lam\|\nabla u(x)\|^{2} \leq \langle \nabla u(x), a(x)\nabla u(x)\rangle = b_A(u,u)(x)$ for a.e.\ $x$. Integrate.
\end{proof}

\begin{lemma}[\leanname{bilinFormOfCoeff\_eq\_left}, \leanname{bilinFormOfCoeff\_eq\_right}]
On an open set $\Omega$, $\mathfrak{a}_A(u,v)$ is independent of the chosen $W^{1,2}$ witness for $u$ (resp.\ $v$).
\end{lemma}

\begin{lstlisting}[style=Matlab-editor,basicstyle=\mlttfamily,escapechar=",mlshowsectionrules=true,mathescape=true,morecomment={[s]{\%\{}{\%\}}},language=lean,numbers=none]
/-- The bilinear form does not depend on the chosen `W^{1,2}` witness in the
left slot. -/
theorem bilinFormOfCoeff_eq_left
    {Ω : Set E} (hΩ : IsOpen Ω)
    (A : EllipticCoeff d Ω)
    {u v : E → ℝ}
    (hu₁ hu₂ : MemW1pWitness 2 u Ω)
    (hv : MemW1pWitness 2 v Ω) :
    bilinFormOfCoeff A hu₁ hv = bilinFormOfCoeff A hu₂ hv := by
  have hgrad_ae : hu₁.weakGrad =ᵐ[volume.restrict Ω] hu₂.weakGrad :=
    MemW1pWitness.ae_eq hΩ hu₁ hu₂
  unfold bilinFormOfCoeff
  apply integral_congr_ae
  filter_upwards [hgrad_ae] with x hx
  simp [bilinFormIntegrandOfCoeff, hx]

/-- The bilinear form does not depend on the chosen `W^{1,2}` witness in the
right slot. -/
theorem bilinFormOfCoeff_eq_right
    {Ω : Set E} (hΩ : IsOpen Ω)
    (A : EllipticCoeff d Ω)
    {u v : E → ℝ}
    (hu : MemW1pWitness 2 u Ω)
    (hv₁ hv₂ : MemW1pWitness 2 v Ω) :
    bilinFormOfCoeff A hu hv₁ = bilinFormOfCoeff A hu hv₂ := by
  have hgrad_ae : hv₁.weakGrad =ᵐ[volume.restrict Ω] hv₂.weakGrad :=
    MemW1pWitness.ae_eq hΩ hv₁ hv₂
  unfold bilinFormOfCoeff
  apply integral_congr_ae
  filter_upwards [hgrad_ae] with x hx
  simp [bilinFormIntegrandOfCoeff, hx]

\end{lstlisting}
\begin{proof}
Two witnesses for the same function on an open set have a.e.-equal weak gradients. The integrand depends only on these gradients, so the integrals agree.
\end{proof}

\begin{lemma}
(\leanname{bilinFormOfCoeff\_add\_left},
\leanname{bilinFormOfCoeff\_smul\_left},
\leanname{bilinFormOfCoeff\_add\_right},
\leanname{bilinFormOfCoeff\_smul\_right})\\
The bilinear form $\mathfrak{a}_A$ is bilinear: additive and homogeneous in both slots.
\end{lemma}

\begin{lstlisting}[style=Matlab-editor,basicstyle=\mlttfamily,escapechar=",mlshowsectionrules=true,mathescape=true,morecomment={[s]{\%\{}{\%\}}},language=lean,numbers=none]
/-- The bilinear form is additive in the left slot. -/
theorem bilinFormOfCoeff_add_left
    {Ω : Set E} {u₁ u₂ v : E → ℝ}
    (A : EllipticCoeff d Ω)
    (hu₁ : MemW1pWitness 2 u₁ Ω)
    (hu₂ : MemW1pWitness 2 u₂ Ω)
    (hv : MemW1pWitness 2 v Ω) :
    bilinFormOfCoeff A (hu₁.add hu₂) hv =
      bilinFormOfCoeff A hu₁ hv + bilinFormOfCoeff A hu₂ hv := by
  rw [bilinFormOfCoeff, bilinFormOfCoeff, bilinFormOfCoeff,
    bilinFormIntegrandOfCoeff_add_left]
  exact
    integral_add
      (integrable_bilinFormIntegrandOfCoeff A hu₁ hv)
      (integrable_bilinFormIntegrandOfCoeff A hu₂ hv)

/-- The bilinear form is homogeneous in the left slot. -/
theorem bilinFormOfCoeff_smul_left
    {Ω : Set E} {u v : E → ℝ} (c : ℝ)
    (A : EllipticCoeff d Ω)
    (hu : MemW1pWitness 2 u Ω)
    (hv : MemW1pWitness 2 v Ω) :
    bilinFormOfCoeff A (hu.smul c) hv = c * bilinFormOfCoeff A hu hv := by
  rw [bilinFormOfCoeff, bilinFormIntegrandOfCoeff_smul_left, integral_const_mul, bilinFormOfCoeff]

/-- The bilinear form is additive in the right slot. -/
theorem bilinFormOfCoeff_add_right
    {Ω : Set E} {u v₁ v₂ : E → ℝ}
    (A : EllipticCoeff d Ω)
    (hu : MemW1pWitness 2 u Ω)
    (hv₁ : MemW1pWitness 2 v₁ Ω)
    (hv₂ : MemW1pWitness 2 v₂ Ω) :
    bilinFormOfCoeff A hu (hv₁.add hv₂) =
      bilinFormOfCoeff A hu hv₁ + bilinFormOfCoeff A hu hv₂ := by
  rw [bilinFormOfCoeff, bilinFormOfCoeff, bilinFormOfCoeff,
    bilinFormIntegrandOfCoeff_add_right]
  exact
    integral_add
      (integrable_bilinFormIntegrandOfCoeff A hu hv₁)
      (integrable_bilinFormIntegrandOfCoeff A hu hv₂)

/-- The bilinear form is homogeneous in the right slot. -/
theorem bilinFormOfCoeff_smul_right
    {Ω : Set E} {u v : E → ℝ} (c : ℝ)
    (A : EllipticCoeff d Ω)
    (hu : MemW1pWitness 2 u Ω)
    (hv : MemW1pWitness 2 v Ω) :
    bilinFormOfCoeff A hu (hv.smul c) = c * bilinFormOfCoeff A hu hv := by
  rw [bilinFormOfCoeff, bilinFormIntegrandOfCoeff_smul_right, integral_const_mul, bilinFormOfCoeff]

\end{lstlisting}
\begin{proof}
The weak gradient of $u_1 + u_2$ is the sum of the weak gradients (by witness construction), and matrix-vector multiplication and the inner product are linear, so the integrand splits. Linearity then transfers to the integral by integrability. The argument for the right slot and for scalar multiplication is identical.
\end{proof}

\subsection{Divergence-form right-hand sides}

\begin{definition}[\leanname{divergenceRHSIntegrandOfField}, \leanname{divergenceRHSOfField}]
For $F : E \to E$ in $\Lp{2}(\Omega; E)$ and a witness $h_v$ for $v \in W^{1,2}(\Omega)$, set
\[
  \ell_F(v)(x) := \langle F(x), \nabla v(x)\rangle, \qquad
  L_F(v) := -\int_{\Omega} \ell_F(v)(x)\,dx.
\]
\end{definition}

\begin{lstlisting}[style=Matlab-editor,basicstyle=\mlttfamily,escapechar=",mlshowsectionrules=true,mathescape=true,morecomment={[s]{\%\{}{\%\}}},language=lean,numbers=none]
/-- Pointwise divergence-form RHS integrand `⟪F, ∇v⟫`. -/
def divergenceRHSIntegrandOfField {Ω : Set E} {v : E → ℝ}
    (F : E → E) (hv : MemW1pWitness 2 v Ω) : E → ℝ :=
  fun x => ⟪F x, hv.weakGrad x⟫_ℝ

/-- The divergence-form RHS functional `v ↦ -∫_Ω ⟪F, ∇v⟫`. -/
noncomputable def divergenceRHSOfField {Ω : Set E} {v : E → ℝ}
    (F : E → E) (hv : MemW1pWitness 2 v Ω) : ℝ :=
  -∫ x in Ω, divergenceRHSIntegrandOfField F hv x

\end{lstlisting}
\begin{lemma}[\leanname{aestronglyMeasurable\_divergenceRHSIntegrandOfField}, \leanname{integrable\_divergenceRHSIntegrandOfField}]
$\ell_F(v)$ is a.e.\ strongly measurable and integrable on $\Omega$.
\end{lemma}

\begin{lstlisting}[style=Matlab-editor,basicstyle=\mlttfamily,escapechar=",mlshowsectionrules=true,mathescape=true,morecomment={[s]{\%\{}{\%\}}},language=lean,numbers=none]
omit [NeZero d] in
/-- The divergence-form RHS integrand is a.e. strongly measurable on `Ω`. -/
theorem aestronglyMeasurable_divergenceRHSIntegrandOfField
    {Ω : Set E} {F : E → E} (hF : MemLp F 2 (volume.restrict Ω))
    {v : E → ℝ} (hv : MemW1pWitness 2 v Ω) :
    AEStronglyMeasurable (divergenceRHSIntegrandOfField F hv) (volume.restrict Ω) := by
  have hF_ofLp_cont : Continuous (fun y : E => (WithLp.ofLp y : Fin d → ℝ)) := by
    simpa using (PiLp.continuous_ofLp 2 (fun _ : Fin d => ℝ))
  have hF_ofLp :
      AEStronglyMeasurable (fun x => (WithLp.ofLp (F x) : Fin d → ℝ)) (volume.restrict Ω) :=
    hF_ofLp_cont.comp_aestronglyMeasurable hF.aestronglyMeasurable
  have hsum :
      AEMeasurable
        (∑ i : Fin d, fun x => F x i * hv.weakGrad x i)
        (volume.restrict Ω) := by
    refine Finset.aemeasurable_sum (s := (Finset.univ : Finset (Fin d)))
      (f := fun i x => F x i * hv.weakGrad x i) ?_
    intro i _hi
    have hFi : AEMeasurable (fun x => F x i) (volume.restrict Ω) := by
      simpa using
        (Continuous.comp_aestronglyMeasurable (continuous_apply i) hF_ofLp).aemeasurable
    exact hFi.mul (hv.weakGrad_component_memLp i).aemeasurable
  have hscalar : ∀ a b : ℝ, ⟪a, b⟫_ℝ = a * b := by
    intro a b
    simpa using (RCLike.inner_apply' a b)
  refine hsum.aestronglyMeasurable.congr ?_
  filter_upwards with x
  simp [divergenceRHSIntegrandOfField, PiLp.inner_apply, hscalar]

/-- The divergence-form RHS integrand is integrable on `Ω`. -/
theorem integrable_divergenceRHSIntegrandOfField
    {Ω : Set E} {F : E → E} (hF : MemLp F 2 (volume.restrict Ω))
    {v : E → ℝ} (hv : MemW1pWitness 2 v Ω) :
    Integrable (divergenceRHSIntegrandOfField F hv) (volume.restrict Ω) := by
  let _ := (inferInstance : NeZero d)
  let μ : Measure E := volume.restrict Ω
  have hdom_int :
      Integrable (fun x => ‖F x‖ * ‖hv.weakGrad x‖) μ := by
    simpa [mul_comm] using (hv.weakGrad_memLp.norm.integrable_mul hF.norm)
  have hbound :
      ∀ᵐ x ∂μ, ‖divergenceRHSIntegrandOfField F hv x‖ ≤ ‖F x‖ * ‖hv.weakGrad x‖ := by
    filter_upwards with x
    simpa [divergenceRHSIntegrandOfField] using abs_real_inner_le_norm (F x) (hv.weakGrad x)
  exact Integrable.mono' hdom_int
    (aestronglyMeasurable_divergenceRHSIntegrandOfField hF hv)
    hbound

\end{lstlisting}
\begin{proof}
Pointwise, $|\ell_F(v)| \leq \|F\|\,\|\nabla v\|$, and the right-hand side belongs to $\Lp{1}$ by H\"older.
\end{proof}

\begin{lemma}[\leanname{divergenceRHSOfField\_bound}]
\[
  |L_F(v)| \leq \Bigl(\int_{\Omega}\|F\|^{2}\,dx\Bigr)^{1/2}
  \Bigl(\int_{\Omega}\|\nabla v\|^{2}\,dx\Bigr)^{1/2}.
\]
\end{lemma}

\begin{lstlisting}[style=Matlab-editor,basicstyle=\mlttfamily,escapechar=",mlshowsectionrules=true,mathescape=true,morecomment={[s]{\%\{}{\%\}}},language=lean,numbers=none]
/-- The divergence-form RHS is bounded with respect to the `L²` gradient
seminorm. -/
theorem divergenceRHSOfField_bound
    {Ω : Set E} {F : E → E} (hF : MemLp F 2 (volume.restrict Ω))
    {v : E → ℝ} (hv : MemW1pWitness 2 v Ω) :
    |divergenceRHSOfField F hv| ≤
      (∫ x, ‖F x‖ ^ (2 : ℝ) ∂(volume.restrict Ω)) ^ (1 / (2 : ℝ)) *
        (∫ x, ‖hv.weakGrad x‖ ^ (2 : ℝ) ∂(volume.restrict Ω)) ^ (1 / (2 : ℝ)) := by
  let μ : Measure E := volume.restrict Ω
  have hint_int := integrable_divergenceRHSIntegrandOfField hF hv
  have hpointwise :
      ∀ᵐ x ∂μ, ‖divergenceRHSIntegrandOfField F hv x‖ ≤ ‖F x‖ * ‖hv.weakGrad x‖ := by
    filter_upwards with x
    simpa [divergenceRHSIntegrandOfField] using abs_real_inner_le_norm (F x) (hv.weakGrad x)
  have hnorm_le :
      ∫ x, ‖divergenceRHSIntegrandOfField F hv x‖ ∂μ ≤
        ∫ x, ‖F x‖ * ‖hv.weakGrad x‖ ∂μ := by
    have hdom_int : Integrable (fun x => ‖F x‖ * ‖hv.weakGrad x‖) μ := by
      simpa [mul_comm] using (hv.weakGrad_memLp.norm.integrable_mul hF.norm)
    exact integral_mono_ae hint_int.norm hdom_int hpointwise
  have hF_memLp : MemLp F (ENNReal.ofReal (2 : ℝ)) μ := by
    simpa using hF
  have hholder :
      ∫ x, ‖F x‖ * ‖hv.weakGrad x‖ ∂μ ≤
        (∫ x, ‖F x‖ ^ (2 : ℝ) ∂μ) ^ (1 / (2 : ℝ)) *
          (∫ x, ‖hv.weakGrad x‖ ^ (2 : ℝ) ∂μ) ^ (1 / (2 : ℝ)) := by
    have hhv_memLp : MemLp hv.weakGrad (ENNReal.ofReal (2 : ℝ)) μ := by
      simpa using hv.weakGrad_memLp
    exact
      integral_mul_norm_le_Lp_mul_Lq (μ := μ) (f := F) (g := hv.weakGrad)
        Real.HolderConjugate.two_two hF_memLp hhv_memLp
  simpa [divergenceRHSOfField, μ] using
    (calc
      |divergenceRHSOfField F hv| = ‖∫ x, divergenceRHSIntegrandOfField F hv x ∂μ‖ := by
        simp [divergenceRHSOfField, μ]
      _ ≤ ∫ x, ‖divergenceRHSIntegrandOfField F hv x‖ ∂μ :=
        norm_integral_le_integral_norm _
      _ ≤ ∫ x, ‖F x‖ * ‖hv.weakGrad x‖ ∂μ := hnorm_le
      _ ≤
        (∫ x, ‖F x‖ ^ (2 : ℝ) ∂μ) ^ (1 / (2 : ℝ)) *
          (∫ x, ‖hv.weakGrad x‖ ^ (2 : ℝ) ∂μ) ^ (1 / (2 : ℝ)) := hholder)

\end{lstlisting}
\begin{proof}
Bound $|L_F(v)|$ by the integral of $|\ell_F(v)|$, then by the integral of $\|F\|\,\|\nabla v\|$, and apply H\"older.
\end{proof}

\begin{lemma}[\leanname{divergenceRHSOfField\_add}, \leanname{divergenceRHSOfField\_smul}]
$L_F$ is additive and homogeneous in $v$.
\end{lemma}

\begin{lstlisting}[style=Matlab-editor,basicstyle=\mlttfamily,escapechar=",mlshowsectionrules=true,mathescape=true,morecomment={[s]{\%\{}{\%\}}},language=lean,numbers=none]
/-- The divergence-form RHS is additive in the test slot. -/
theorem divergenceRHSOfField_add
    {Ω : Set E} {F : E → E} (hF : MemLp F 2 (volume.restrict Ω))
    {v₁ v₂ : E → ℝ}
    (hv₁ : MemW1pWitness 2 v₁ Ω)
    (hv₂ : MemW1pWitness 2 v₂ Ω) :
    divergenceRHSOfField F (hv₁.add hv₂) =
      divergenceRHSOfField F hv₁ + divergenceRHSOfField F hv₂ := by
  rw [divergenceRHSOfField, divergenceRHSOfField, divergenceRHSOfField,
    divergenceRHSIntegrandOfField_add]
  rw [integral_add
    (integrable_divergenceRHSIntegrandOfField hF hv₁)
    (integrable_divergenceRHSIntegrandOfField hF hv₂)]
  ring

omit [NeZero d] in
/-- The divergence-form RHS is homogeneous in the test slot. -/
theorem divergenceRHSOfField_smul
    {Ω : Set E} {F : E → E} (c : ℝ) {v : E → ℝ}
    (hv : MemW1pWitness 2 v Ω) :
    divergenceRHSOfField F (hv.smul c) = c * divergenceRHSOfField F hv := by
  rw [divergenceRHSOfField, divergenceRHSIntegrandOfField_smul, integral_const_mul,
    divergenceRHSOfField]
  ring

\end{lstlisting}
\begin{proof}
The weak gradient is linear in $v$ (witness construction), and the inner product distributes; integrate.
\end{proof}

\begin{definition}[\leanname{weakProblemRHSOfField}, \leanname{weakProblemRHSOfFieldAndDatum}]
The \emph{raw-function} divergence RHS is $v \mapsto L_F(v)$ when $v \in H^{1}_{0}(\Omega)$ (using a chosen witness) and $0$ otherwise. The \emph{shifted} version with boundary datum $u_0$ is $v \mapsto L_F(v) - \mathfrak{a}_A(u_0, v)$.
\end{definition}

\begin{lstlisting}[style=Matlab-editor,basicstyle=\mlttfamily,escapechar=",mlshowsectionrules=true,mathescape=true,morecomment={[s]{\%\{}{\%\}}},language=lean,numbers=none]
/-- A witness-independent RHS on raw functions obtained by choosing a witness
when `v ∈ H₀¹(Ω)` and returning `0` otherwise. -/
noncomputable def weakProblemRHSOfField {Ω : Set E}
    (F : E → E) : (E → ℝ) → ℝ := by
  classical
  intro v
  exact
    if hv0 : MemH01 v Ω then
      divergenceRHSOfField F (DeGiorgi.MemW1p.someWitness (MemW01p.memW1p hv0))
    else
      0

/-- Raw-function RHS for the lifted Dirichlet problem with boundary datum
`u₀`. -/
noncomputable def weakProblemRHSOfFieldAndDatum {Ω : Set E}
    (A : EllipticCoeff d Ω) (F : E → E)
    {u₀ : E → ℝ} (hu₀ : MemW1pWitness 2 u₀ Ω) :
    (E → ℝ) → ℝ := by
  classical
  intro v
  exact
    if hv0 : MemH01 v Ω then
      let hv0w : MemW1pWitness 2 v Ω :=
        DeGiorgi.MemW1p.someWitness (MemW01p.memW1p hv0)
      divergenceRHSOfField F hv0w - bilinFormOfCoeff A hu₀ hv0w
    else
      0

\end{lstlisting}
\begin{lemma}[\leanname{weakProblemRHSOfField\_eq\_of\_memH01}, \leanname{weakProblemRHSOfFieldAndDatum\_eq\_of\_memH01}]
For $v \in H^{1}_{0}(\Omega)$, the raw-function (resp.\ shifted) RHS agrees with $L_F(v)$ (resp.\ $L_F(v) - \mathfrak{a}_A(u_0,v)$) on every witness.
\end{lemma}

\begin{lstlisting}[style=Matlab-editor,basicstyle=\mlttfamily,escapechar=",mlshowsectionrules=true,mathescape=true,morecomment={[s]{\%\{}{\%\}}},language=lean,numbers=none]
/-- On `H₀¹(Ω)`, the raw-function RHS agrees with the witness-dependent
divergence-form functional. -/
theorem weakProblemRHSOfField_eq_of_memH01
    {Ω : Set E} (hΩ : IsOpen Ω) {F : E → E} {v : E → ℝ}
    (hv0 : MemH01 v Ω) (hv : MemW1pWitness 2 v Ω) :
    weakProblemRHSOfField (Ω := Ω) F v = divergenceRHSOfField F hv := by
  let _ := (inferInstance : NeZero d)
  rw [weakProblemRHSOfField, dif_pos hv0]
  let hw0 : MemW1pWitness 2 v Ω :=
    DeGiorgi.MemW1p.someWitness (MemW01p.memW1p hv0)
  have hgrad_ae : hw0.weakGrad =ᵐ[volume.restrict Ω] hv.weakGrad := by
    have hp2 : (1 : ENNReal) ≤ 2 := by
      norm_num
    have hcomp :
        ∀ i : Fin d,
          (fun x => hw0.weakGrad x i) =ᵐ[volume.restrict Ω] (fun x => hv.weakGrad x i) := by
      intro i
      exact HasWeakPartialDeriv.ae_eq hΩ (hw0.isWeakGrad i) (hv.isWeakGrad i)
        ((hw0.weakGrad_component_memLp i).locallyIntegrable hp2)
        ((hv.weakGrad_component_memLp i).locallyIntegrable hp2)
    filter_upwards [ae_all_iff.2 hcomp] with x hx
    ext i
    exact hx i
  have hintegral :
      ∫ x in Ω, divergenceRHSIntegrandOfField F hw0 x =
        ∫ x in Ω, divergenceRHSIntegrandOfField F hv x := by
    apply integral_congr_ae
    filter_upwards [hgrad_ae] with x hx
    simpa [divergenceRHSIntegrandOfField, hw0] using
      congrArg (fun g => ⟪F x, g⟫_ℝ) hx
  simpa [hw0, divergenceRHSOfField] using congrArg Neg.neg hintegral

/-- On `H₀¹(Ω)`, the shifted raw RHS agrees with the natural lifted
functional. -/
theorem weakProblemRHSOfFieldAndDatum_eq_of_memH01
    {Ω : Set E} (hΩ : IsOpen Ω)
    (A : EllipticCoeff d Ω)
    {F : E → E} {u₀ : E → ℝ} (hu₀ : MemW1pWitness 2 u₀ Ω)
    {v : E → ℝ} (hv0 : MemH01 v Ω) (hv : MemW1pWitness 2 v Ω) :
    weakProblemRHSOfFieldAndDatum (A := A) (Ω := Ω) F hu₀ v =
      divergenceRHSOfField F hv - bilinFormOfCoeff A hu₀ hv := by
  rw [weakProblemRHSOfFieldAndDatum, dif_pos hv0]
  let hw0 : MemW1pWitness 2 v Ω :=
    DeGiorgi.MemW1p.someWitness (MemW01p.memW1p hv0)
  have hdiv : divergenceRHSOfField F hw0 = divergenceRHSOfField F hv := by
    have hgrad_ae : hw0.weakGrad =ᵐ[volume.restrict Ω] hv.weakGrad :=
      MemW1pWitness.ae_eq hΩ hw0 hv
    unfold divergenceRHSOfField
    congr 1
    apply integral_congr_ae
    filter_upwards [hgrad_ae] with x hx
    simpa [divergenceRHSIntegrandOfField] using congrArg (fun g => ⟪F x, g⟫_ℝ) hx
  have hbilin : bilinFormOfCoeff A hu₀ hw0 = bilinFormOfCoeff A hu₀ hv := by
    exact bilinFormOfCoeff_eq_right hΩ A hu₀ hw0 hv
  change divergenceRHSOfField F hw0 - bilinFormOfCoeff A hu₀ hw0 =
    divergenceRHSOfField F hv - bilinFormOfCoeff A hu₀ hv
  rw [hdiv, hbilin]

\end{lstlisting}
\begin{proof}
Two witnesses for $v$ on the open set $\Omega$ have a.e.-equal weak gradients, so $L_F$ and $\mathfrak{a}_A$ depend only on $v$ itself (Lemmas above).
\end{proof}

\begin{lemma}[\leanname{weakProblemRHSOfField\_bound}, \leanname{weakProblemRHSOfFieldAndDatum\_bound}]
For $v \in H^{1}_{0}(\Omega)$ with witness $h_v$,
\begin{align*}
  |L_F(v)| &\leq \|F\|_{\Lp{2}} \,\|\nabla v\|_{\Lp{2}}, \\
  |L_F(v) - \mathfrak{a}_A(u_0, v)| &\leq \bigl(\|F\|_{\Lp{2}} + \Lambda\,\|\nabla u_0\|_{\Lp{2}}\bigr)\,\|\nabla v\|_{\Lp{2}}.
\end{align*}
\end{lemma}

\begin{lstlisting}[style=Matlab-editor,basicstyle=\mlttfamily,escapechar=",mlshowsectionrules=true,mathescape=true,morecomment={[s]{\%\{}{\%\}}},language=lean,numbers=none]
/-- The chosen raw-function RHS is bounded on `H₀¹(Ω)` with respect to the
`L²` gradient seminorm. -/
theorem weakProblemRHSOfField_bound
    {Ω : Set E} (hΩ : IsOpen Ω) {F : E → E}
    (hF : MemLp F 2 (volume.restrict Ω))
    (v : E → ℝ) (hv0 : MemH01 v Ω) (hv : MemW1pWitness 2 v Ω) :
    |weakProblemRHSOfField (Ω := Ω) F v| ≤
      (∫ x, ‖F x‖ ^ (2 : ℝ) ∂(volume.restrict Ω)) ^ (1 / (2 : ℝ)) *
        (∫ x, ‖hv.weakGrad x‖ ^ (2 : ℝ) ∂(volume.restrict Ω)) ^ (1 / (2 : ℝ)) := by
  rw [weakProblemRHSOfField_eq_of_memH01 hΩ hv0 hv]
  exact divergenceRHSOfField_bound hF hv

/-- The shifted raw RHS is bounded on `H₀¹(Ω)` with respect to the `L²`
gradient seminorm. -/
theorem weakProblemRHSOfFieldAndDatum_bound
    {Ω : Set E} (hΩ : IsOpen Ω)
    (A : EllipticCoeff d Ω)
    {F : E → E} (hF : MemLp F 2 (volume.restrict Ω))
    {u₀ : E → ℝ} (hu₀ : MemW1pWitness 2 u₀ Ω)
    (v : E → ℝ) (hv0 : MemH01 v Ω) (hv : MemW1pWitness 2 v Ω) :
    |weakProblemRHSOfFieldAndDatum (A := A) (Ω := Ω) F hu₀ v| ≤
      ((∫ x, ‖F x‖ ^ (2 : ℝ) ∂(volume.restrict Ω)) ^ (1 / (2 : ℝ)) +
          A.Λ * (∫ x, ‖hu₀.weakGrad x‖ ^ (2 : ℝ) ∂(volume.restrict Ω)) ^ (1 / (2 : ℝ))) *
        (∫ x, ‖hv.weakGrad x‖ ^ (2 : ℝ) ∂(volume.restrict Ω)) ^ (1 / (2 : ℝ)) := by
  let CF : ℝ :=
    (∫ x, ‖F x‖ ^ (2 : ℝ) ∂(volume.restrict Ω)) ^ (1 / (2 : ℝ))
  let CU : ℝ :=
    A.Λ * (∫ x, ‖hu₀.weakGrad x‖ ^ (2 : ℝ) ∂(volume.restrict Ω)) ^ (1 / (2 : ℝ))
  let GV : ℝ :=
    (∫ x, ‖hv.weakGrad x‖ ^ (2 : ℝ) ∂(volume.restrict Ω)) ^ (1 / (2 : ℝ))
  rw [weakProblemRHSOfFieldAndDatum_eq_of_memH01 hΩ A hu₀ hv0 hv]
  have hdiv : |divergenceRHSOfField F hv| ≤ CF * GV := by
    simpa [CF, GV] using divergenceRHSOfField_bound hF hv
  have hbilin : |bilinFormOfCoeff A hu₀ hv| ≤ CU * GV := by
    simpa [CU, GV, mul_assoc] using bilinForm_bound A hu₀ hv
  calc
    |divergenceRHSOfField F hv - bilinFormOfCoeff A hu₀ hv|
      = |divergenceRHSOfField F hv + (-bilinFormOfCoeff A hu₀ hv)| := by ring_nf
    _ ≤ |divergenceRHSOfField F hv| + |-bilinFormOfCoeff A hu₀ hv| := abs_add_le _ _
    _ = |divergenceRHSOfField F hv| + |bilinFormOfCoeff A hu₀ hv| := by simp
    _ ≤ CF * GV + CU * GV := add_le_add hdiv hbilin
    _ = (CF + CU) * GV := by ring
    _ = ((∫ x, ‖F x‖ ^ (2 : ℝ) ∂(volume.restrict Ω)) ^ (1 / (2 : ℝ)) +
          A.Λ * (∫ x, ‖hu₀.weakGrad x‖ ^ (2 : ℝ) ∂(volume.restrict Ω)) ^ (1 / (2 : ℝ))) *
        (∫ x, ‖hv.weakGrad x‖ ^ (2 : ℝ) ∂(volume.restrict Ω)) ^ (1 / (2 : ℝ)) := by
      rfl

\end{lstlisting}
\begin{proof}
Combine the previous bounds using the triangle inequality and the bilinear-form continuity bound $|\mathfrak{a}_A(u_0,v)| \leq \Lambda\,\|\nabla u_0\|_{\Lp{2}}\|\nabla v\|_{\Lp{2}}$.
\end{proof}

\subsection{Smooth test functions and the L\texorpdfstring{$^2$}{2} coefficient operator}

\begin{definition}[\leanname{IsSmoothTestOn}]
A function $u : E \to \R$ is a \emph{smooth test on $\Omega$} if $u \in C^{\infty}(E)$, $u$ has compact support, and $\supp u \subseteq \Omega$.
\end{definition}

\begin{lstlisting}[style=Matlab-editor,basicstyle=\mlttfamily,escapechar=",mlshowsectionrules=true,mathescape=true,morecomment={[s]{\%\{}{\%\}}},language=lean,numbers=none]
/-! ## Smooth Test Functions And L² Coefficient Operator -/

/-- Smooth compactly supported test functions supported in `Ω`. -/
def IsSmoothTestOn (Ω : Set E) (u : E → ℝ) : Prop :=
  ContDiff ℝ (⊤ : ℕ∞) u ∧ HasCompactSupport u ∧ tsupport u ⊆ Ω

\end{lstlisting}
\begin{lemma}[\leanname{IsSmoothTestOn.zero}, \leanname{add}, \leanname{smul}, \leanname{sub}]
The smooth tests on $\Omega$ form a real vector space (closed under sums, scalar multiples, and differences).
\end{lemma}

\begin{lstlisting}[style=Matlab-editor,basicstyle=\mlttfamily,escapechar=",mlshowsectionrules=true,mathescape=true,morecomment={[s]{\%\{}{\%\}}},language=lean,numbers=none]
omit [NeZero d] in
theorem IsSmoothTestOn.zero
    {Ω : Set E} :
    IsSmoothTestOn Ω (fun _ : E => (0 : ℝ)) := by
  refine ⟨by simpa using (contDiff_const : ContDiff ℝ (⊤ : ℕ∞) (fun _ : E => (0 : ℝ))),
    HasCompactSupport.zero, by simp⟩

omit [NeZero d] in
theorem IsSmoothTestOn.add
    {Ω : Set E} {u v : E → ℝ}
    (hu : IsSmoothTestOn Ω u) (hv : IsSmoothTestOn Ω v) :
    IsSmoothTestOn Ω (fun x => u x + v x) := by
  refine ⟨hu.1.add hv.1, hu.2.1.add hv.2.1, ?_⟩
  exact (tsupport_add u v).trans <| Set.union_subset hu.2.2 hv.2.2

omit [NeZero d] in
theorem IsSmoothTestOn.smul
    {Ω : Set E} (c : ℝ) {u : E → ℝ}
    (hu : IsSmoothTestOn Ω u) :
    IsSmoothTestOn Ω (fun x => c * u x) := by
  refine ⟨?_, ?_, ?_⟩
  · simpa [smul_eq_mul] using (contDiff_const.mul hu.1)
  · simpa [Pi.smul_apply, smul_eq_mul] using
      (hu.2.1.smul_left (f := fun _ : E => c))
  · simpa [Pi.smul_apply, smul_eq_mul] using
      (tsupport_smul_subset_right (fun _ : E => c) u).trans hu.2.2

omit [NeZero d] in
theorem IsSmoothTestOn.sub
    {Ω : Set E} {u v : E → ℝ}
    (hu : IsSmoothTestOn Ω u) (hv : IsSmoothTestOn Ω v) :
    IsSmoothTestOn Ω (fun x => u x - v x) := by
  simpa [sub_eq_add_neg, Pi.smul_apply, smul_eq_mul] using hu.add (hv.smul (-1))

\end{lstlisting}
\begin{proof}
Smoothness and compact support are stable under sums and scalar multiples; the support of a sum is contained in the union of supports, hence in $\Omega$.
\end{proof}

\begin{definition}[\leanname{smoothGradField}, \leanname{smoothTestWitness}]
For smooth $u$, the \emph{classical gradient field} is
\[
  (\nabla u)(x) := \bigl((\partial_{i} u)(x)\bigr)_{i=1}^{d} \in E.
\]
A smooth test on $\Omega$ carries a canonical $W^{1,2}(\Omega)$-witness with this classical gradient.
\end{definition}

\begin{lstlisting}[style=Matlab-editor,basicstyle=\mlttfamily,escapechar=",mlshowsectionrules=true,mathescape=true,morecomment={[s]{\%\{}{\%\}}},language=lean,numbers=none]
/-- Classical gradient field of a smooth scalar function. -/
noncomputable def smoothGradField (u : E → ℝ) : E → E :=
  fun x => WithLp.toLp 2 fun i => (fderiv ℝ u x) (EuclideanSpace.single i 1)

/-- Canonical `W^{1,2}` witness on `Ω` carried by a smooth compactly supported
test function supported in `Ω`. -/
noncomputable def smoothTestWitness
    {Ω : Set E} (hΩ : IsOpen Ω) {u : E → ℝ}
    (hu : IsSmoothTestOn Ω u) :
    MemW1pWitness 2 u Ω where
  memLp := (hu.1.continuous.memLp_of_hasCompactSupport hu.2.1).restrict Ω
  weakGrad := smoothGradField u
  weakGrad_component_memLp := by
    intro i
    have hderiv_smooth : ContDiff ℝ (⊤ : ℕ∞)
        (fun x => (fderiv ℝ u x) (EuclideanSpace.single i 1)) :=
      (hu.1.fderiv_right (m := (⊤ : ℕ∞)) (by norm_cast)).clm_apply contDiff_const
    exact
      (hderiv_smooth.continuous.memLp_of_hasCompactSupport
        (hu.2.1.fderiv_apply (𝕜 := ℝ) _)).restrict Ω
  isWeakGrad := by
    intro i
    simpa [smoothGradField, PiLp.toLp_apply] using
      (HasWeakPartialDeriv.of_contDiff (Ω := Ω) (i := i) (f := u) hΩ
        (hu.1.of_le (by norm_cast)))

\end{lstlisting}
\begin{lemma}[\leanname{smoothTest\_memH01}]
A smooth test on $\Omega$ belongs to $H^{1}_{0}(\Omega)$.
\end{lemma}

\begin{lstlisting}[style=Matlab-editor,basicstyle=\mlttfamily,escapechar=",mlshowsectionrules=true,mathescape=true,morecomment={[s]{\%\{}{\%\}}},language=lean,numbers=none]
theorem smoothTest_memH01
    {Ω : Set E} (hΩ : IsOpen Ω) {u : E → ℝ}
    (hu : IsSmoothTestOn Ω u) :
    MemH01 u Ω :=
  memW01p_of_contDiff_hasCompactSupport_subset hΩ hu.1 hu.2.1 hu.2.2

\end{lstlisting}
\begin{proof}
Approximate by its constant sequence; smoothness and compactly supported support in $\Omega$ provide the required convergence in $W^{1,2}$.
\end{proof}

\begin{definition}[\leanname{smoothFunToLp}, \leanname{smoothGradToLp}, \leanname{gradLpOfWitness}]
A smooth test $u$ defines:
\[
  [u] \in \Lp{2}(\Omega; \R), \qquad [\nabla u] \in \Lp{2}(\Omega; E),
\]
the latter being the $\Lp{2}$ class of any witness's weak gradient.
\end{definition}

\begin{lstlisting}[style=Matlab-editor,basicstyle=\mlttfamily,escapechar=",mlshowsectionrules=true,mathescape=true,morecomment={[s]{\%\{}{\%\}}},language=lean,numbers=none]
/-- The `L²` class carried by a smooth test function on `Ω`. -/
noncomputable def smoothFunToLp
    {Ω : Set E} (hΩ : IsOpen Ω) {u : E → ℝ}
    (hu : IsSmoothTestOn Ω u) :
    MeasureTheory.Lp ℝ 2 (volume.restrict Ω) :=
  (smoothTestWitness hΩ hu).memLp.toLp u

/-- The `L²` gradient class carried by a smooth test function on `Ω`. -/
noncomputable def smoothGradToLp
    {Ω : Set E} (hΩ : IsOpen Ω) {u : E → ℝ}
    (hu : IsSmoothTestOn Ω u) :
    MeasureTheory.Lp E 2 (volume.restrict Ω) :=
  gradLpOfWitness (smoothTestWitness hΩ hu)

/-- The `L²` gradient class carried by a Sobolev witness. -/
noncomputable def gradLpOfWitness
    {Ω : Set E} {u : E → ℝ} (hu : MemW1pWitness 2 u Ω) :
    MeasureTheory.Lp E 2 (volume.restrict Ω) :=
  hu.weakGrad_memLp.toLp hu.weakGrad

\end{lstlisting}
\begin{lemma}[\leanname{smoothFunToLp\_zero}, \leanname{\_add}, \leanname{\_smul}, \leanname{\_sub}, and analogues for \leanname{smoothGradToLp}]
Both maps $u \mapsto [u]$ and $u \mapsto [\nabla u]$ are real-linear.
\end{lemma}

\begin{lstlisting}[style=Matlab-editor,basicstyle=\mlttfamily,escapechar=",mlshowsectionrules=true,mathescape=true,morecomment={[s]{\%\{}{\%\}}},language=lean,numbers=none]
theorem smoothFunToLp_zero
    {Ω : Set E} (hΩ : IsOpen Ω) :
    smoothFunToLp hΩ IsSmoothTestOn.zero = 0 := by
  let _ := (inferInstance : NeZero d)
  let hu0 := smoothTestWitness hΩ IsSmoothTestOn.zero
  let hu0Lp := hu0.memLp
  change hu0Lp.toLp (0 : E → ℝ) = 0
  exact hu0Lp.toLp_zero

/-- The `L²` gradient class carried by a smooth test function on `Ω`. -/
noncomputable def smoothGradToLp
    {Ω : Set E} (hΩ : IsOpen Ω) {u : E → ℝ}
    (hu : IsSmoothTestOn Ω u) :
    MeasureTheory.Lp E 2 (volume.restrict Ω) :=
  gradLpOfWitness (smoothTestWitness hΩ hu)

\end{lstlisting}
\begin{proof}
Linearity of the $\Lp{2}$ class follows from $\mathrm{toLp}$-linearity in Mathlib; for the gradient version one further uses linearity of the classical gradient.
\end{proof}

\begin{lemma}[\leanname{norm\_smoothFunToLp\_eq}, \leanname{norm\_smoothGradToLp\_eq}, \leanname{norm\_gradLpOfWitness\_eq}]
The norms in $\Lp{2}$ are
\[
  \bigl\|[u]\bigr\| = \Bigl(\int_{\Omega} |u|^{2}\,dx\Bigr)^{1/2}, \qquad
  \bigl\|[\nabla u]\bigr\| = \Bigl(\int_{\Omega} \|\nabla u\|^{2}\,dx\Bigr)^{1/2}.
\]
\end{lemma}

\begin{lstlisting}[style=Matlab-editor,basicstyle=\mlttfamily,escapechar=",mlshowsectionrules=true,mathescape=true,morecomment={[s]{\%\{}{\%\}}},language=lean,numbers=none]
/-- Norm of the `L²` class carried by a smooth test function. -/
theorem norm_smoothFunToLp_eq
    {Ω : Set E} (hΩ : IsOpen Ω) {u : E → ℝ}
    (hu : IsSmoothTestOn Ω u) :
    ‖smoothFunToLp hΩ hu‖ =
      (∫ x, ‖u x‖ ^ (2 : ℝ) ∂(volume.restrict Ω)) ^ (1 / (2 : ℝ)) := by
  let _ := (inferInstance : NeZero d)
  rw [smoothFunToLp, Lp.norm_toLp]
  rw [MeasureTheory.toReal_eLpNorm (smoothTestWitness hΩ hu).memLp.aestronglyMeasurable]
  simpa using
    (MeasureTheory.lpNorm_eq_integral_norm_rpow_toReal
      (μ := volume.restrict Ω) (f := u) (p := (2 : ENNReal))
      (by norm_num) (by simp) (smoothTestWitness hΩ hu).memLp.aestronglyMeasurable)

/-- Norm of the `L²` gradient class carried by a smooth test function. -/
theorem norm_smoothGradToLp_eq
    {Ω : Set E} (hΩ : IsOpen Ω) {u : E → ℝ}
    (hu : IsSmoothTestOn Ω u) :
    ‖smoothGradToLp hΩ hu‖ =
      (∫ x, ‖smoothGradField u x‖ ^ (2 : ℝ) ∂(volume.restrict Ω)) ^ (1 / (2 : ℝ)) := by
  simpa [smoothGradToLp] using norm_gradLpOfWitness_eq (smoothTestWitness hΩ hu)

theorem norm_gradLpOfWitness_eq
    {Ω : Set E} {u : E → ℝ} (hu : MemW1pWitness 2 u Ω) :
    ‖gradLpOfWitness hu‖ =
      (∫ x, ‖hu.weakGrad x‖ ^ (2 : ℝ) ∂(volume.restrict Ω)) ^ (1 / (2 : ℝ)) := by
  let _ := (inferInstance : NeZero d)
  rw [gradLpOfWitness, Lp.norm_toLp]
  rw [MeasureTheory.toReal_eLpNorm hu.weakGrad_memLp.aestronglyMeasurable]
  simpa using
    (MeasureTheory.lpNorm_eq_integral_norm_rpow_toReal
      (μ := volume.restrict Ω) (f := hu.weakGrad) (p := (2 : ENNReal))
      (by norm_num) (by simp) hu.weakGrad_memLp.aestronglyMeasurable)

\end{lstlisting}
\begin{proof}
Direct from the definition of the $\Lp{2}$ norm via $\mathrm{toLp}$ in Mathlib.
\end{proof}

\subsection{The L\texorpdfstring{$^2$}{2} coefficient operator and the closed bilinear form}

\begin{definition}[\leanname{smoothGradSubmodule}]
Let $\Omega$ be open. Define
\[
  S(\Omega) := \bigl\{ [\nabla u] : u \text{ smooth test on } \Omega \bigr\} \subseteq \Lp{2}(\Omega; E),
\]
a real linear subspace, and let $H(\Omega)$ be its closure.
\end{definition}

\begin{lstlisting}[style=Matlab-editor,basicstyle=\mlttfamily,escapechar=",mlshowsectionrules=true,mathescape=true,morecomment={[s]{\%\{}{\%\}}},language=lean,numbers=none]
/-- The linear subspace of `L²(Ω; E)` generated by classical gradients of
smooth compactly supported test functions supported in `Ω`. -/
noncomputable def smoothGradSubmodule
    {Ω : Set E} (hΩ : IsOpen Ω) :
    Submodule ℝ (MeasureTheory.Lp E 2 (volume.restrict Ω)) where
  carrier := {G | ∃ u : E → ℝ, ∃ hu : IsSmoothTestOn Ω u, smoothGradToLp hΩ hu = G}
  zero_mem' := by
    exact ⟨fun _ : E => (0 : ℝ), IsSmoothTestOn.zero, smoothGradToLp_zero hΩ⟩
  add_mem' := by
    rintro G H ⟨u, hu, rfl⟩ ⟨v, hv, rfl⟩
    exact ⟨fun x => u x + v x, hu.add hv, smoothGradToLp_add hΩ hu hv⟩
  smul_mem' := by
    intro c G hG
    rcases hG with ⟨u, hu, rfl⟩
    exact ⟨fun x => c * u x, hu.smul c, smoothGradToLp_smul hΩ c hu⟩

\end{lstlisting}
\begin{lemma}[\leanname{coeffMulLpL}, \leanname{norm\_coeffMulLpL\_le}]
There is a bounded linear operator $T_A : \Lp{2}(\Omega;E) \to \Lp{2}(\Omega;E)$ given pointwise a.e.\ by $T_A g (x) = a(x)\,g(x)$, satisfying $\|T_A g\|_{\Lp{2}} \leq \Lambda\,\|g\|_{\Lp{2}}$.
\end{lemma}

\begin{lstlisting}[style=Matlab-editor,basicstyle=\mlttfamily,escapechar=",mlshowsectionrules=true,mathescape=true,morecomment={[s]{\%\{}{\%\}}},language=lean,numbers=none]
/-- Pointwise coefficient action on `L²(Ω; E)`. -/
noncomputable def coeffMulLpL
    {Ω : Set E} (A : EllipticCoeff d Ω) :
    MeasureTheory.Lp E 2 (volume.restrict Ω) →L[ℝ]
      MeasureTheory.Lp E 2 (volume.restrict Ω) := by
  let f : MeasureTheory.Lp E 2 (volume.restrict Ω) →ₗ[ℝ]
      MeasureTheory.Lp E 2 (volume.restrict Ω) := {
    toFun := coeffApplyLp A
    map_add' := by
      intro g h
      apply Lp.ext
      let hAgh : MemLp (fun x => matMulE (A.a x) ((g + h) x)) 2 (volume.restrict Ω) :=
        coeffAction_memLp A (Lp.memLp (g + h))
      filter_upwards
        [MemLp.coeFn_toLp hAgh, coeffApplyLp_ae_eq A g, coeffApplyLp_ae_eq A h, Lp.coeFn_add g h,
          Lp.coeFn_add (coeffApplyLp A g) (coeffApplyLp A h)] with x hx_gh hx_g hx_h hx_add hx_sum
      have hpoint : matMulE (A.a x) ((g + h) x) = coeffApplyLp A g x + coeffApplyLp A h x := by
        rw [hx_add, hx_g, hx_h]
        ext i
        simp [matMulE_apply, Matrix.mulVec_add]
      exact hx_gh.trans (hpoint.trans hx_sum.symm)
    map_smul' := by
      intro c g
      apply Lp.ext
      filter_upwards
        [coeffApplyLp_ae_eq A (c • g), coeffApplyLp_ae_eq A g, Lp.coeFn_smul c g,
          Lp.coeFn_smul c (coeffApplyLp A g)] with x hx_cg hx_g hx_smul hx_out
      have hpoint : matMulE (A.a x) ((c • g) x) = c • coeffApplyLp A g x := by
        rw [hx_smul, hx_g]
        ext i
        simp [Pi.smul_apply, matMulE_apply, Matrix.mulVec_smul]
      exact hx_cg.trans (hpoint.trans hx_out.symm) }
  refine f.mkContinuous A.Λ ?_
  intro g
  apply Lp.norm_le_mul_norm_of_ae_le_mul
  filter_upwards [coeffApplyLp_ae_eq A g, ae_matMulE_norm_le A] with x hx_toLp hx_bound
  calc
    ‖coeffApplyLp A g x‖ = ‖matMulE (A.a x) (g x)‖ := by
      rw [hx_toLp]
    _ ≤ A.Λ * ‖g x‖ := hx_bound (g x)

theorem norm_coeffMulLpL_le
    {Ω : Set E} (A : EllipticCoeff d Ω)
    (g : MeasureTheory.Lp E 2 (volume.restrict Ω)) :
    ‖coeffMulLpL A g‖ ≤ A.Λ * ‖g‖ := by
  let hg : MemLp (fun x => g x) 2 (volume.restrict Ω) := Lp.memLp g
  let hAg : MemLp (fun x => matMulE (A.a x) (g x)) 2 (volume.restrict Ω) :=
    coeffAction_memLp A hg
  apply Lp.norm_le_mul_norm_of_ae_le_mul
  filter_upwards [MemLp.coeFn_toLp hAg, ae_matMulE_norm_le A] with x hx_toLp hx_bound
  calc
    ‖hAg.toLp (fun x => matMulE (A.a x) (g x)) x‖ = ‖matMulE (A.a x) (g x)‖ := by
      rw [hx_toLp]
    _ ≤ A.Λ * ‖g x‖ := hx_bound (g x)

\end{lstlisting}
\begin{proof}
For each component $i$, $|T_A g(x)_i| \leq \|T_A g(x)\| \leq \Lambda\|g(x)\|$ by the mixed bound, and componentwise measurability of $a$ and $g$ make $T_A g$ strongly measurable. Hence each component lies in $\Lp{2}$, so does $T_A g$, and the operator norm bound follows from $\|T_A g(x)\| \leq \Lambda\|g(x)\|$ a.e.\ via $\Lp{p}$ pointwise dominance.
\end{proof}

\begin{definition}[\leanname{coeffBilinSubmodule}]
For any closed real subspace $H \leq \Lp{2}(\Omega;E)$ define the continuous bilinear form
\[
  B_{A,H}(g,h) := \langle T_A g, h\rangle_{\Lp{2}}.
\]
\end{definition}

\begin{lstlisting}[style=Matlab-editor,basicstyle=\mlttfamily,escapechar=",mlshowsectionrules=true,mathescape=true,morecomment={[s]{\%\{}{\%\}}},language=lean,numbers=none]
/-- The coefficient bilinear form on a closed gradient subspace of
`L²(Ω; E)`. -/
noncomputable def coeffBilinSubmodule
    {Ω : Set E} (A : EllipticCoeff d Ω)
    (H : Submodule ℝ (MeasureTheory.Lp E 2 (volume.restrict Ω))) :
    H →L[ℝ] H →L[ℝ] ℝ := by
  refine LinearMap.mkContinuous₂
    (LinearMap.mk₂ ℝ
      (fun g h => ⟪coeffMulLpL A (H.subtypeL g), H.subtypeL h⟫_ℝ)
      (by intro g₁ g₂ h; simp [inner_add_left])
      (by intro c g h; simp [inner_smul_left])
      (by intro g h₁ h₂; simp [inner_add_right])
      (by intro c g h; simp [inner_smul_right]))
    A.Λ ?_
  intro g h
  calc
    ‖⟪coeffMulLpL A (H.subtypeL g), H.subtypeL h⟫_ℝ‖
        ≤ ‖coeffMulLpL A (H.subtypeL g)‖ * ‖H.subtypeL h‖ := by
            simpa using
              (norm_inner_le_norm (𝕜 := ℝ) (coeffMulLpL A (H.subtypeL g)) (H.subtypeL h))
    _ ≤ (A.Λ * ‖H.subtypeL g‖) * ‖H.subtypeL h‖ := by
          gcongr
          exact norm_coeffMulLpL_le A (H.subtypeL g)
    _ = A.Λ * ‖g‖ * ‖h‖ := by
          simp [mul_assoc]

\end{lstlisting}
\begin{lemma}[\leanname{coeffBilinSubmodule\_coercive}]
$B_{A,H}$ is bounded by $\Lambda$ and coercive on $H$ with constant $\lam$:
\[
  |B_{A,H}(g,h)| \leq \Lambda\,\|g\|\,\|h\|, \qquad \lam\,\|g\|^{2} \leq B_{A,H}(g,g).
\]
\end{lemma}

\begin{lstlisting}[style=Matlab-editor,basicstyle=\mlttfamily,escapechar=",mlshowsectionrules=true,mathescape=true,morecomment={[s]{\%\{}{\%\}}},language=lean,numbers=none]
theorem coeffBilinSubmodule_coercive
    {Ω : Set E} (A : EllipticCoeff d Ω)
    (H : Submodule ℝ (MeasureTheory.Lp E 2 (volume.restrict Ω))) :
    IsCoercive (coeffBilinSubmodule A H) := by
  refine ⟨A.lam, A.hlam, ?_⟩
  intro g
  let G : MeasureTheory.Lp E 2 (volume.restrict Ω) := H.subtypeL g
  let hAG_mem : MemLp (fun x => matMulE (A.a x) (G x)) 2 (volume.restrict Ω) :=
    coeffAction_memLp A (Lp.memLp G)
  have hcoeff_ae :
      coeffMulLpL A G =ᵐ[volume.restrict Ω] fun x => matMulE (A.a x) (G x) :=
    MemLp.coeFn_toLp hAG_mem
  have hright_int0 :
      Integrable (fun x => ⟪coeffMulLpL A G x, G x⟫_ℝ) (volume.restrict Ω) :=
    MeasureTheory.L2.integrable_inner (coeffMulLpL A G) G
  have hright_int :
      Integrable (fun x => ⟪matMulE (A.a x) (G x), G x⟫_ℝ) (volume.restrict Ω) := by
    refine hright_int0.congr ?_
    filter_upwards [hcoeff_ae] with x hx
    simp [hx]
  have hleft_int :
      Integrable (fun x => A.lam * ⟪G x, G x⟫_ℝ) (volume.restrict Ω) := by
    simpa using (MeasureTheory.L2.integrable_inner G G).const_mul A.lam
  have hcoer_ae :
      ∀ᵐ x ∂(volume.restrict Ω),
        A.lam * ⟪G x, G x⟫_ℝ ≤ ⟪matMulE (A.a x) (G x), G x⟫_ℝ := by
    filter_upwards [A.coercive] with x hx
    simpa [real_inner_comm, real_inner_self_eq_norm_sq] using hx (G x)
  have hmono :
      ∫ x, A.lam * ⟪G x, G x⟫_ℝ ∂(volume.restrict Ω) ≤
        ∫ x, ⟪matMulE (A.a x) (G x), G x⟫_ℝ ∂(volume.restrict Ω) :=
    integral_mono_ae hleft_int hright_int hcoer_ae
  have hright_eq :
      ∫ x, ⟪matMulE (A.a x) (G x), G x⟫_ℝ ∂(volume.restrict Ω) =
        coeffBilinSubmodule A H g g := by
    change ∫ x, ⟪matMulE (A.a x) (G x), G x⟫_ℝ ∂(volume.restrict Ω) =
      ⟪coeffMulLpL A G, G⟫_ℝ
    rw [MeasureTheory.L2.inner_def]
    apply integral_congr_ae
    filter_upwards [hcoeff_ae] with x hx
    simp [hx]
  calc
    A.lam * ‖g‖ * ‖g‖ = A.lam * ‖G‖ * ‖G‖ := by simp [G]
    _ = A.lam * ⟪G, G⟫_ℝ := by
          have hinner : ⟪G, G⟫_ℝ = ‖G‖ * ‖G‖ := by
            simp [pow_two]
          rw [hinner]
          ring
    _ = ∫ x, A.lam * ⟪G x, G x⟫_ℝ ∂(volume.restrict Ω) := by
          rw [MeasureTheory.L2.inner_def, integral_const_mul]
    _ ≤ ∫ x, ⟪matMulE (A.a x) (G x), G x⟫_ℝ ∂(volume.restrict Ω) := hmono
    _ = coeffBilinSubmodule A H g g := hright_eq

\end{lstlisting}
\begin{proof}
The boundedness $|\langle T_A g, h\rangle| \leq \|T_A g\|\,\|h\| \leq \Lambda \|g\|\,\|h\|$ comes from Cauchy--Schwarz and the bound on $T_A$. For coercivity, write
\[
  B_{A,H}(g,g) = \int_{\Omega}\langle a(x) g(x), g(x)\rangle\,dx \geq \int_{\Omega} \lam\,\|g(x)\|^{2}\,dx = \lam\,\|g\|_{\Lp{2}}^{2}
\]
using the a.e.\ coercivity of $A$.
\end{proof}

\begin{lemma}[\leanname{tendsto\_eLpNorm\_vector\_of\_componentwise}]
Componentwise $\Lp{2}$-convergence in each coordinate of $\Fin d$ implies vector $\Lp{2}$-convergence: if for each $i$ the seminorm $\|g_n^{(i)} - g^{(i)}\|_{\Lp{2}} \to 0$, then $\|g_n - g\|_{\Lp{2}} \to 0$.
\end{lemma}

\begin{lstlisting}[style=Matlab-editor,basicstyle=\mlttfamily,escapechar=",mlshowsectionrules=true,mathescape=true,morecomment={[s]{\%\{}{\%\}}},language=lean,numbers=none]
omit [NeZero d] in
theorem tendsto_eLpNorm_vector_of_componentwise
    {Ω : Set E} {g : E → E} {gseq : ℕ → E → E}
    (hcomp_mem : ∀ n : ℕ, ∀ i : Fin d,
      MemLp (fun x => gseq n x i - g x i) 2 (volume.restrict Ω))
    (hcomp : ∀ i : Fin d,
      Tendsto (fun n => eLpNorm (fun x => gseq n x i - g x i) 2 (volume.restrict Ω))
        atTop (nhds 0)) :
    Tendsto (fun n => eLpNorm (fun x => gseq n x - g x) 2 (volume.restrict Ω))
      atTop (nhds 0) := by
  have hupper :
      ∀ n,
        eLpNorm (fun x => gseq n x - g x) 2 (volume.restrict Ω) ≤
          ∑ i : Fin d,
            eLpNorm (fun x => gseq n x i - g x i) 2 (volume.restrict Ω) := by
    intro n
    exact eLpNorm_vector_le_sum_component_eLpNorm
      ((MeasureTheory.MemLp.of_eval_piLp (fun i => hcomp_mem n i)).aestronglyMeasurable)
  have hsum :
      Tendsto
        (fun n => ∑ i : Fin d,
          eLpNorm (fun x => gseq n x i - g x i) 2 (volume.restrict Ω))
        atTop (nhds 0) := by
    simpa using tendsto_finset_sum (Finset.univ : Finset (Fin d)) (fun i _ => hcomp i)
  exact tendsto_of_tendsto_of_tendsto_of_le_of_le tendsto_const_nhds hsum (fun n => zero_le _) hupper

\end{lstlisting}
\begin{proof}
Pointwise, $\|g_n(x) - g(x)\|^{2} = \sum_i (g_n(x)_i - g(x)_i)^{2}$, hence $\|g_n(x) - g(x)\| \leq \sum_i |g_n(x)_i - g(x)_i|$. The $\Lp{2}$ norm of the right-hand side is bounded by $\sum_i \|g_n^{(i)} - g^{(i)}\|_{\Lp{2}}$, which tends to zero.
\end{proof}

\begin{lemma}[\leanname{coeffMulLpL\_inner\_gradLpOfWitness\_eq\_bilinFormOfCoeff}]
For witnesses $h_u, h_v$ for $u, v \in W^{1,2}(\Omega)$,
\[
  \langle T_A [\nabla u], [\nabla v]\rangle_{\Lp{2}} = \mathfrak{a}_A(u, v).
\]
\end{lemma}

\begin{lstlisting}[style=Matlab-editor,basicstyle=\mlttfamily,escapechar=",mlshowsectionrules=true,mathescape=true,morecomment={[s]{\%\{}{\%\}}},language=lean,numbers=none]
theorem coeffMulLpL_inner_gradLpOfWitness_eq_bilinFormOfCoeff
    {Ω : Set E} (A : EllipticCoeff d Ω)
    {u v : E → ℝ}
    (hu : MemW1pWitness 2 u Ω)
    (hv : MemW1pWitness 2 v Ω) :
    ⟪coeffMulLpL A (gradLpOfWitness hu), gradLpOfWitness hv⟫_ℝ =
      bilinFormOfCoeff A hu hv := by
  rw [MeasureTheory.L2.inner_def, bilinFormOfCoeff]
  apply integral_congr_ae
  filter_upwards
    [coeffApplyLp_ae_eq A (gradLpOfWitness hu),
      MemLp.coeFn_toLp hu.weakGrad_memLp,
      MemLp.coeFn_toLp hv.weakGrad_memLp] with x hxA hxu hxv
  have hxcoeff :
      ((coeffMulLpL A) (gradLpOfWitness hu)) x = matMulE (A.a x) (hu.weakGrad x) := by
    simpa [coeffMulLpL, gradLpOfWitness, hxu] using hxA
  have hxgrad :
      (gradLpOfWitness hv) x = hv.weakGrad x := by
    simpa [gradLpOfWitness] using hxv
  simp [bilinFormIntegrandOfCoeff, hxcoeff, hxgrad]

\end{lstlisting}
\begin{proof}
Both sides equal $\int_\Omega \langle a(x)\,\nabla u(x), \nabla v(x)\rangle\,dx$; the left-hand side after using the a.e.\ identification $T_A[\nabla u] = a(\cdot)\nabla u(\cdot)$.
\end{proof}

\subsection{Weak problems and weak solutions}

\begin{definition}[\leanname{WeakProblem}]
A \emph{weak elliptic Dirichlet problem} consists of an open bounded set $\Omega \subseteq E$, an elliptic coefficient field $A$ on $\Omega$, and a functional $\mathrm{rhs} : (E \to \R) \to \R$.
\end{definition}

\begin{lstlisting}[style=Matlab-editor,basicstyle=\mlttfamily,escapechar=",mlshowsectionrules=true,mathescape=true,morecomment={[s]{\%\{}{\%\}}},language=lean,numbers=none]
/-- A weak elliptic Dirichlet problem on an open bounded domain. -/
structure WeakProblem where
  Ω : Set E
  hΩ : IsOpen Ω
  hΩ_bdd : Bornology.IsBounded Ω
  coeff : EllipticCoeff d Ω
  rhs : (E → ℝ) → ℝ

\end{lstlisting}
\begin{definition}[\leanname{IsWeakSolution}]
A function $u : E \to \R$ is a \emph{weak solution} of $P = (\Omega, A, \mathrm{rhs})$ if $u \in H^{1}_{0}(\Omega)$ and, for every witness $h_u$ of $u$ and every $v \in H^{1}_{0}(\Omega)$ with witness $h_v$,
\[
  \mathfrak{a}_A(u, v) = \mathrm{rhs}(v).
\]
\end{definition}

\begin{lstlisting}[style=Matlab-editor,basicstyle=\mlttfamily,escapechar=",mlshowsectionrules=true,mathescape=true,morecomment={[s]{\%\{}{\%\}}},language=lean,numbers=none]
/-- `u` is a weak solution of the variational problem `P` if it lies in
`H₀¹(P.Ω)` and satisfies the weak identity against every `H₀¹` test function. -/
def IsWeakSolution (P : WeakProblem (d := d)) (u : E → ℝ) : Prop :=
  MemH01 u P.Ω ∧
  ∀ (hwu : MemW1pWitness 2 u P.Ω) (v : E → ℝ), MemH01 v P.Ω →
    ∀ (hwv : MemW1pWitness 2 v P.Ω),
      bilinFormOfCoeff P.coeff hwu hwv = P.rhs v

\end{lstlisting}
\begin{definition}[\leanname{IsSubsolution}, \leanname{IsSupersolution}, \leanname{IsSolution}]
For an elliptic field $A$ on $\Omega$, $u \in W^{1,2}(\Omega)$ is a (weak) \emph{subsolution} of $-\dv(a\nabla u) \leq 0$ if for all $\varphi \in H^{1}_{0}(\Omega)$ with $\varphi \geq 0$, $\mathfrak{a}_A(u, \varphi) \leq 0$. The notion of \emph{supersolution} reverses the inequality. A \emph{solution} is both a sub- and supersolution.
\end{definition}

\begin{lstlisting}[style=Matlab-editor,basicstyle=\mlttfamily,escapechar=",mlshowsectionrules=true,mathescape=true,morecomment={[s]{\%\{}{\%\}}},language=lean,numbers=none]
/-- `u` is a weak subsolution of `-div(A∇u) ≤ 0` on `Ω`. -/
def IsSubsolution {Ω : Set E} (A : EllipticCoeff d Ω) (u : E → ℝ) : Prop :=
  MemW1p 2 u Ω ∧
  ∀ (hu : MemW1pWitness 2 u Ω)
    (φ : E → ℝ), MemH01 φ Ω →
    ∀ (hφ : MemW1pWitness 2 φ Ω),
      (∀ x, 0 ≤ φ x) → bilinFormOfCoeff A hu hφ ≤ 0

/-- `u` is a weak supersolution of `-div(A∇u) ≥ 0` on `Ω`. -/
def IsSupersolution {Ω : Set E} (A : EllipticCoeff d Ω) (u : E → ℝ) : Prop :=
  MemW1p 2 u Ω ∧
  ∀ (hu : MemW1pWitness 2 u Ω)
    (φ : E → ℝ), MemH01 φ Ω →
    ∀ (hφ : MemW1pWitness 2 φ Ω),
      (∀ x, 0 ≤ φ x) → 0 ≤ bilinFormOfCoeff A hu hφ

/-- `u` is a weak solution of the homogeneous equation iff it is both a weak
subsolution and a weak supersolution. -/
def IsSolution {Ω : Set E} (A : EllipticCoeff d Ω) (u : E → ℝ) : Prop :=
  IsSubsolution A u ∧ IsSupersolution A u

\end{lstlisting}
\begin{definition}[\leanname{IsHomogeneousWeakSolution}]
$u \in W^{1,2}(\Omega)$ is a \emph{homogeneous weak solution} of $-\dv(a\nabla u) = 0$ on $\Omega$ if, for every $W^{1,2}$-witness $h_u$ of $u$ and every $\varphi \in H^{1}_{0}(\Omega)$ with $W^{1,2}$-witness $h_\varphi$, one has $\mathfrak{a}_A(u,\varphi) = 0$.
\end{definition}

\begin{lstlisting}[style=Matlab-editor,basicstyle=\mlttfamily,escapechar=",mlshowsectionrules=true,mathescape=true,morecomment={[s]{\%\{}{\%\}}},language=lean,numbers=none]
/-- `u` is a local weak solution of the homogeneous equation on `Ω` if it lies
in `W^{1,2}(Ω)` and the bilinear form vanishes against every `H₀¹(Ω)` test
function. This is the equality-form local weak-solution interface used by the
Harnack and Holder statements. -/
def IsHomogeneousWeakSolution {Ω : Set E}
    (A : EllipticCoeff d Ω) (u : E → ℝ) : Prop :=
  MemW1p 2 u Ω ∧
  ∀ (hu : MemW1pWitness 2 u Ω)
    (φ : E → ℝ), MemH01 φ Ω →
    ∀ (hφ : MemW1pWitness 2 φ Ω),
      bilinFormOfCoeff A hu hφ = 0

\end{lstlisting}
\begin{lemma}[\leanname{isHomogeneousWeakSolution\_isSubsolution}, \leanname{isHomogeneousWeakSolution\_isSupersolution}, \leanname{isHomogeneousWeakSolution\_isSolution}]
A homogeneous weak solution is both a subsolution and a supersolution (hence a solution).
\end{lemma}

\begin{lstlisting}[style=Matlab-editor,basicstyle=\mlttfamily,escapechar=",mlshowsectionrules=true,mathescape=true,morecomment={[s]{\%\{}{\%\}}},language=lean,numbers=none]
/-- A homogeneous weak solution is a subsolution. -/
theorem isHomogeneousWeakSolution_isSubsolution
    {Ω : Set E} {A : EllipticCoeff d Ω} {u : E → ℝ}
    (h : IsHomogeneousWeakSolution A u) :
    IsSubsolution A u := by
  refine ⟨h.1, ?_⟩
  intro hu φ hφ hφw hφ_nonneg
  have hEq : bilinFormOfCoeff A hu hφw = 0 := h.2 hu φ hφ hφw
  simp [hEq]

/-- A homogeneous weak solution is a supersolution. -/
theorem isHomogeneousWeakSolution_isSupersolution
    {Ω : Set E} {A : EllipticCoeff d Ω} {u : E → ℝ}
    (h : IsHomogeneousWeakSolution A u) :
    IsSupersolution A u := by
  refine ⟨h.1, ?_⟩
  intro hu φ hφ hφw hφ_nonneg
  have hEq : bilinFormOfCoeff A hu hφw = 0 := h.2 hu φ hφ hφw
  simp [hEq]

/-- A homogeneous weak solution is both a subsolution and a supersolution. -/
theorem isHomogeneousWeakSolution_isSolution
    {Ω : Set E} {A : EllipticCoeff d Ω} {u : E → ℝ}
    (h : IsHomogeneousWeakSolution A u) :
    IsSolution A u :=
  ⟨isHomogeneousWeakSolution_isSubsolution h,
    isHomogeneousWeakSolution_isSupersolution h⟩

\end{lstlisting}
\begin{proof}
For any nonnegative test $\varphi$, $\mathfrak{a}_A(u,\varphi) = 0$, so both inequalities hold trivially.
\end{proof}

Note that \leanname{IsWeakSolution} is the existential predicate used by the Lax--Milgram existence theorem (it carries an explicit RHS functional and forces $u \in H^{1}_{0}$). \leanname{IsHomogeneousWeakSolution} is the equality-form local interface used by the Harnack and H\"older statements: it requires $u \in W^{1,2}(\Omega)$ but does \emph{not} require zero boundary values, and uses test functions in $H^{1}_{0}(\Omega)$.

\subsection{Existence via Lax--Milgram}

\begin{theorem}[\leanname{weakProblem\_exists}]
\label{thm:weakProblem-exists}
Let $d \geq 2$. Let $\Omega \subseteq E$ be open and bounded, $A$ an elliptic coefficient field on $\Omega$, and $\mathrm{rhs}$ a functional on $H^{1}_{0}(\Omega)$ that is real-linear and bounded:
\[
  |\mathrm{rhs}(v)| \leq C_{\mathrm{rhs}}\,\|\nabla v\|_{\Lp{2}}.
\]
Then there exists $u \in H^{1}_{0}(\Omega)$ that is a weak solution of $(\Omega, A, \mathrm{rhs})$.
\end{theorem}

\begin{lstlisting}[style=Matlab-editor,basicstyle=\mlttfamily,escapechar=",mlshowsectionrules=true,mathescape=true,morecomment={[s]{\%\{}{\%\}}},language=lean,numbers=none]
/-! ## Weak Problem Existence (Lax-Milgram) -/

/-- Existence of weak solutions via Lax-Milgram.
The bilinear form `bilinFormOfCoeff` is bounded and coercive on `H₀¹(Ω)`
(proved above), and the RHS is a bounded linear functional, so Lax-Milgram
gives existence and uniqueness. -/
theorem weakProblem_exists
    {Ω : Set E}
    (hd : 2 ≤ d)
    (hΩ : IsOpen Ω) (hΩ_bdd : Bornology.IsBounded Ω)
    (A : EllipticCoeff d Ω)
    (rhs : (E → ℝ) → ℝ)
    (hF_add : ∀ u v : E → ℝ, MemH01 u Ω → MemH01 v Ω →
      rhs (fun x => u x + v x) = rhs u + rhs v)
    (hF_smul : ∀ c : ℝ, ∀ u : E → ℝ, MemH01 u Ω →
      rhs (fun x => c * u x) = c * rhs u)
    (hF_bounded : ∃ C, 0 ≤ C ∧ ∀ v : E → ℝ, MemH01 v Ω →
      ∀ hwv : MemW1pWitness 2 v Ω,
      |rhs v| ≤ C *
        (∫ x, ‖hwv.weakGrad x‖ ^ (2 : ℝ) ∂(volume.restrict Ω)) ^ (1 / (2 : ℝ))) :
    ∃ u : E → ℝ, IsWeakSolution (d := d) ⟨Ω, hΩ, hΩ_bdd, A, rhs⟩ u := by
  classical
  let μ : Measure E := volume.restrict Ω
  let S : Submodule ℝ (MeasureTheory.Lp E 2 μ) := smoothGradSubmodule hΩ
  let H : Submodule ℝ (MeasureTheory.Lp E 2 μ) := S.topologicalClosure
  let e : S →ₗ[ℝ] H := Submodule.inclusion S.le_topologicalClosure
  obtain ⟨Cp, hCp_top, hPoinc⟩ := smoothCompactSupport_L2_bound_on_bounded_ge_two hd hΩ_bdd
  obtain ⟨C_rhs, hC_rhs_nonneg, hF_bound⟩ := hF_bounded
  let repFun : S → E → ℝ := fun g => Classical.choose g.2
  let repSmooth : ∀ g : S, IsSmoothTestOn Ω (repFun g) := by
    intro g
    exact Classical.choose (Classical.choose_spec g.2)
  have repEq : ∀ g : S, smoothGradToLp hΩ (repSmooth g) = (g : MeasureTheory.Lp E 2 μ) := by
    intro g
    exact Classical.choose_spec (Classical.choose_spec g.2)
  have smoothFunToLp_bound :
      ∀ {u : E → ℝ} (hu : IsSmoothTestOn Ω u),
        ‖smoothFunToLp hΩ hu‖ ≤ Cp.toReal * ‖smoothGradToLp hΩ hu‖ := by
    intro u hu
    have hgradNorm_memLp : MemLp (smoothGradNorm u) 2 μ := by
      simpa [μ, norm_smoothGradField_eq_smoothGradNorm] using
        (smoothTestWitness hΩ hu).weakGrad_memLp.norm
    have hP :
        eLpNorm u 2 μ ≤ Cp * eLpNorm (smoothGradNorm u) 2 μ := by
      simpa [μ] using hPoinc hu.1 hu.2.1 hu.2.2
    have hleft_top : eLpNorm u 2 μ ≠ ⊤ :=
      (smoothTestWitness hΩ hu).memLp.eLpNorm_lt_top.ne
    have hright_top : Cp * eLpNorm (smoothGradNorm u) 2 μ ≠ ⊤ := by
      exact (ENNReal.mul_lt_top hCp_top hgradNorm_memLp.eLpNorm_lt_top).ne
    have hToReal :
        (eLpNorm u 2 μ).toReal ≤ (Cp * eLpNorm (smoothGradNorm u) 2 μ).toReal :=
      (ENNReal.toReal_le_toReal hleft_top hright_top).2 hP
    have hfunEq : ‖smoothFunToLp hΩ hu‖ = (eLpNorm u 2 μ).toReal := by
      rw [smoothFunToLp, Lp.norm_toLp]
    have hgradEq : ‖smoothGradToLp hΩ hu‖ = (eLpNorm (smoothGradNorm u) 2 μ).toReal := by
      have hgrad_memLp : MemLp (smoothGradField u) 2 μ :=
        (smoothTestWitness hΩ hu).weakGrad_memLp
      calc
        ‖smoothGradToLp hΩ hu‖ =
            (∫ x, ‖smoothGradField u x‖ ^ (2 : ℝ) ∂μ) ^ (1 / (2 : ℝ)) := by
              simpa [μ] using (norm_smoothGradToLp_eq hΩ hu)
        _ = (eLpNorm (smoothGradField u) 2 μ).toReal := by
              symm
              rw [MeasureTheory.toReal_eLpNorm hgrad_memLp.aestronglyMeasurable]
              simpa using
                (MeasureTheory.lpNorm_eq_integral_norm_rpow_toReal
                  (μ := μ) (f := smoothGradField u) (p := (2 : ENNReal))
                  (by norm_num) (by simp) hgrad_memLp.aestronglyMeasurable)
        _ = (eLpNorm (smoothGradNorm u) 2 μ).toReal := by
          congr 1
          calc
            eLpNorm (smoothGradField u) 2 μ = eLpNorm (fun x => ‖smoothGradField u x‖) 2 μ := by
              symm
              exact eLpNorm_norm (μ := μ) (p := (2 : ENNReal)) (smoothGradField u)
            _ = eLpNorm (smoothGradNorm u) 2 μ := by
                  simp [norm_smoothGradField_eq_smoothGradNorm]
    rw [hfunEq, hgradEq]
    convert hToReal using 1
    rw [ENNReal.toReal_mul]
  have smoothFunToLp_eq_of_sameGrad :
      ∀ {u v : E → ℝ} (hu : IsSmoothTestOn Ω u) (hv : IsSmoothTestOn Ω v),
        smoothGradToLp hΩ hu = smoothGradToLp hΩ hv →
          smoothFunToLp hΩ hu = smoothFunToLp hΩ hv := by
    intro u v hu hv hgradEq
    have huv : IsSmoothTestOn Ω (fun x => u x - v x) := hu.sub hv
    have hgradZero : smoothGradToLp hΩ huv = 0 := by
      calc
        smoothGradToLp hΩ huv = smoothGradToLp hΩ hu - smoothGradToLp hΩ hv := by
          simpa using smoothGradToLp_sub hΩ hu hv
        _ = 0 := by rw [hgradEq, sub_self]
    have hfunZeroNorm : ‖smoothFunToLp hΩ huv‖ = 0 := by
      have hbound := smoothFunToLp_bound huv
      have hupper : ‖smoothFunToLp hΩ huv‖ ≤ 0 := by
        simpa [hgradZero] using hbound
      exact le_antisymm hupper (norm_nonneg _)
    have hfunZero : smoothFunToLp hΩ huv = 0 := norm_eq_zero.mp hfunZeroNorm
    have hsub :
        smoothFunToLp hΩ huv = smoothFunToLp hΩ hu - smoothFunToLp hΩ hv := by
      simpa using smoothFunToLp_sub hΩ hu hv
    rw [hsub] at hfunZero
    exact sub_eq_zero.mp hfunZero
  have rhs_eq_of_sameGrad :
      ∀ {u v : E → ℝ} (hu : IsSmoothTestOn Ω u) (hv : IsSmoothTestOn Ω v),
        smoothGradToLp hΩ hu = smoothGradToLp hΩ hv →
          rhs u = rhs v := by
    intro u v hu hv hgradEq
    have hu0 : MemH01 u Ω := smoothTest_memH01 hΩ hu
    have hv0 : MemH01 v Ω := smoothTest_memH01 hΩ hv
    have huv : IsSmoothTestOn Ω (fun x => u x - v x) := hu.sub hv
    have huv0 : MemH01 (fun x => u x - v x) Ω := smoothTest_memH01 hΩ huv
    let hwuv : MemW1pWitness 2 (fun x => u x - v x) Ω := smoothTestWitness hΩ huv
    have hgradZero : smoothGradToLp hΩ huv = 0 := by
      calc
        smoothGradToLp hΩ huv = smoothGradToLp hΩ hu - smoothGradToLp hΩ hv := by
          simpa using smoothGradToLp_sub hΩ hu hv
        _ = 0 := by rw [hgradEq, sub_self]
    have hseminormZero :
        (∫ x, ‖hwuv.weakGrad x‖ ^ (2 : ℝ) ∂μ) ^ (1 / (2 : ℝ)) = 0 := by
      have hgradZero' : gradLpOfWitness hwuv = 0 := by
        simpa [hwuv, smoothGradToLp] using hgradZero
      have hnormZero : ‖gradLpOfWitness hwuv‖ = 0 := by
        rw [hgradZero']
        simp
      have hseminormEq := norm_gradLpOfWitness_eq hwuv
      rw [hnormZero] at hseminormEq
      simpa [μ] using hseminormEq.symm
    have hbound := hF_bound (fun x => u x - v x) huv0 hwuv
    have hrhsZero : rhs (fun x => u x - v x) = 0 := by
      have habsZero : |rhs (fun x => u x - v x)| = 0 := by
        have hupper : |rhs (fun x => u x - v x)| ≤ 0 := by
          calc
            |rhs (fun x => u x - v x)| ≤
                C_rhs * (∫ x, ‖hwuv.weakGrad x‖ ^ (2 : ℝ) ∂μ) ^ (1 / (2 : ℝ)) := hbound
            _ = 0 := by rw [hseminormZero]; ring
        exact le_antisymm hupper (abs_nonneg _)
      exact abs_eq_zero.mp habsZero
    have hvNeg0 : MemH01 (fun x => (-1) * v x) Ω := by
      simpa [Pi.smul_apply, smul_eq_mul] using (MemW01p.smul (-1) hv0)
    have hsubEq : rhs (fun x => u x - v x) = rhs u - rhs v := by
      calc
        rhs (fun x => u x - v x) = rhs (fun x => u x + (-1) * v x) := by
          congr
          ext x
          ring
        _ = rhs u + rhs (fun x => (-1) * v x) := hF_add u (fun x => (-1) * v x) hu0 hvNeg0
        _ = rhs u + (-1) * rhs v := by rw [hF_smul (-1) v hv0]
        _ = rhs u - rhs v := by ring
    rw [hsubEq] at hrhsZero
    linarith
  let L0 : S →ₗ[ℝ] ℝ := {
    toFun := fun g => rhs (repFun g)
    map_add' := by
      intro g h
      have hEq :
          rhs (repFun (g + h)) = rhs (fun x => repFun g x + repFun h x) := by
        apply rhs_eq_of_sameGrad (repSmooth (g + h)) ((repSmooth g).add (repSmooth h))
        calc
          smoothGradToLp hΩ (repSmooth (g + h)) = (g + h : S) := repEq (g + h)
          _ = (g : MeasureTheory.Lp E 2 μ) + (h : MeasureTheory.Lp E 2 μ) := by rfl
          _ = smoothGradToLp hΩ (repSmooth g) + smoothGradToLp hΩ (repSmooth h) := by
                rw [repEq g, repEq h]
          _ = smoothGradToLp hΩ ((repSmooth g).add (repSmooth h)) := by
                symm
                exact smoothGradToLp_add hΩ (repSmooth g) (repSmooth h)
      rw [hEq, hF_add (repFun g) (repFun h)
        (smoothTest_memH01 hΩ (repSmooth g))
        (smoothTest_memH01 hΩ (repSmooth h))]
    map_smul' := by
      intro c g
      have hEq :
          rhs (repFun (c • g)) = rhs (fun x => c * repFun g x) := by
        apply rhs_eq_of_sameGrad (repSmooth (c • g)) ((repSmooth g).smul c)
        calc
          smoothGradToLp hΩ (repSmooth (c • g)) = (c • g : S) := repEq (c • g)
          _ = c • (g : MeasureTheory.Lp E 2 μ) := by rfl
          _ = c • smoothGradToLp hΩ (repSmooth g) := by rw [repEq g]
          _ = smoothGradToLp hΩ ((repSmooth g).smul c) := by
                symm
                exact smoothGradToLp_smul hΩ c (repSmooth g)
      calc
        rhs (repFun (c • g)) = rhs (fun x => c * repFun g x) := hEq
        _ = c * rhs (repFun g) := hF_smul c (repFun g) (smoothTest_memH01 hΩ (repSmooth g))
        _ = (RingHom.id ℝ) c • rhs (repFun g) := by simp [smul_eq_mul] }
  have hL0_bound : ∃ C, ∀ g : S, ‖L0 g‖ ≤ C * ‖e g‖ := by
    refine ⟨C_rhs, ?_⟩
    intro g
    let hwg : MemW1pWitness 2 (repFun g) Ω := smoothTestWitness hΩ (repSmooth g)
    calc
      ‖L0 g‖ = |rhs (repFun g)| := by rfl
      _ ≤ C_rhs * (∫ x, ‖hwg.weakGrad x‖ ^ (2 : ℝ) ∂μ) ^ (1 / (2 : ℝ)) := by
            exact hF_bound (repFun g) (smoothTest_memH01 hΩ (repSmooth g)) hwg
      _ = C_rhs * ‖gradLpOfWitness hwg‖ := by rw [norm_gradLpOfWitness_eq hwg]
      _ = C_rhs * ‖smoothGradToLp hΩ (repSmooth g)‖ := by rfl
      _ = C_rhs * ‖(g : MeasureTheory.Lp E 2 μ)‖ := by rw [repEq g]
      _ = C_rhs * ‖g‖ := by rfl
      _ = C_rhs * ‖e g‖ := by simp [e]
  have hDense : DenseRange e := by
    have hDenseInclusion :
        DenseRange
          (Set.inclusion
            (Submodule.le_topologicalClosure S :
              (S : Set (MeasureTheory.Lp E 2 μ)) ⊆ (H : Set (MeasureTheory.Lp E 2 μ)))) := by
      rw [denseRange_inclusion_iff]
      change (H : Set (MeasureTheory.Lp E 2 μ)) ⊆ closure (S : Set (MeasureTheory.Lp E 2 μ))
      intro x hx
      simpa [H, Submodule.topologicalClosure_coe] using hx
    simpa [e] using hDenseInclusion
  let U0 : S →ₗ[ℝ] MeasureTheory.Lp ℝ 2 μ := {
    toFun := fun g => smoothFunToLp hΩ (repSmooth g)
    map_add' := by
      intro g h
      have hEq :
          smoothFunToLp hΩ (repSmooth (g + h)) =
            smoothFunToLp hΩ ((repSmooth g).add (repSmooth h)) := by
        apply smoothFunToLp_eq_of_sameGrad (repSmooth (g + h)) ((repSmooth g).add (repSmooth h))
        calc
          smoothGradToLp hΩ (repSmooth (g + h)) = (g + h : S) := repEq (g + h)
          _ = (g : MeasureTheory.Lp E 2 μ) + (h : MeasureTheory.Lp E 2 μ) := by rfl
          _ = smoothGradToLp hΩ (repSmooth g) + smoothGradToLp hΩ (repSmooth h) := by
                rw [repEq g, repEq h]
          _ = smoothGradToLp hΩ ((repSmooth g).add (repSmooth h)) := by
                symm
                exact smoothGradToLp_add hΩ (repSmooth g) (repSmooth h)
      rw [hEq, smoothFunToLp_add hΩ (repSmooth g) (repSmooth h)]
    map_smul' := by
      intro c g
      have hEq :
          smoothFunToLp hΩ (repSmooth (c • g)) =
            smoothFunToLp hΩ ((repSmooth g).smul c) := by
        apply smoothFunToLp_eq_of_sameGrad (repSmooth (c • g)) ((repSmooth g).smul c)
        calc
          smoothGradToLp hΩ (repSmooth (c • g)) = (c • g : S) := repEq (c • g)
          _ = c • (g : MeasureTheory.Lp E 2 μ) := by rfl
          _ = c • smoothGradToLp hΩ (repSmooth g) := by rw [repEq g]
          _ = smoothGradToLp hΩ ((repSmooth g).smul c) := by
                symm
                exact smoothGradToLp_smul hΩ c (repSmooth g)
      calc
        smoothFunToLp hΩ (repSmooth (c • g)) =
            smoothFunToLp hΩ ((repSmooth g).smul c) := hEq
        _ = c • smoothFunToLp hΩ (repSmooth g) := smoothFunToLp_smul hΩ c (repSmooth g)
        _ = (RingHom.id ℝ) c • smoothFunToLp hΩ (repSmooth g) := by simp }
  have hU0_bound : ∃ C, ∀ g : S, ‖U0 g‖ ≤ C * ‖e g‖ := by
    refine ⟨Cp.toReal, ?_⟩
    intro g
    calc
      ‖U0 g‖ = ‖smoothFunToLp hΩ (repSmooth g)‖ := by rfl
      _ ≤ Cp.toReal * ‖smoothGradToLp hΩ (repSmooth g)‖ := by
            exact smoothFunToLp_bound (repSmooth g)
      _ = Cp.toReal * ‖(g : MeasureTheory.Lp E 2 μ)‖ := by rw [repEq g]
      _ = Cp.toReal * ‖g‖ := by rfl
      _ = Cp.toReal * ‖e g‖ := by simp [e]
  let L : H →L[ℝ] ℝ := L0.extendOfNorm e
  let U : H →L[ℝ] MeasureTheory.Lp ℝ 2 μ := U0.extendOfNorm e
  have hL_eq : ∀ g : S, L (e g) = L0 g := by
    intro g
    exact LinearMap.extendOfNorm_eq hDense hL0_bound g
  have hU_eq : ∀ g : S, U (e g) = U0 g := by
    intro g
    exact LinearMap.extendOfNorm_eq hDense hU0_bound g
  have hcoercive : IsCoercive (coeffBilinSubmodule A H) := coeffBilinSubmodule_coercive A H
  have hComplete : CompleteSpace H := by
    simpa [H] using (Submodule.topologicalClosure.completeSpace (U := S))
  let fvec : H := by
    exact (@InnerProductSpace.toDual ℝ H _ _ _ hComplete).symm L
  let gsol : H := by
    exact (@IsCoercive.continuousLinearEquivOfBilin H _ _ hComplete _ hcoercive).symm fvec
  let Gsol : MeasureTheory.Lp E 2 μ := gsol
  have hLM :
      ∀ w : H, coeffBilinSubmodule A H gsol w = L w := by
    intro w
    calc
      coeffBilinSubmodule A H gsol w =
          ⟪(@IsCoercive.continuousLinearEquivOfBilin H _ _ hComplete _ hcoercive) gsol, w⟫_ℝ := by
            symm
            exact @IsCoercive.continuousLinearEquivOfBilin_apply H _ _ hComplete
              _ hcoercive gsol w
      _ = ⟪fvec, w⟫_ℝ := by simp [gsol]
      _ = L w := by
            change ⟪(@InnerProductSpace.toDual ℝ H _ _ _ hComplete).symm L, w⟫_ℝ = L w
            exact @InnerProductSpace.toDual_symm_apply ℝ H _ _ _ hComplete (x := w) (y := L)
  have hgsol_mem :
      (gsol : MeasureTheory.Lp E 2 μ) ∈ closure (S : Set (MeasureTheory.Lp E 2 μ)) := by
    change (gsol : MeasureTheory.Lp E 2 μ) ∈ (H : Set (MeasureTheory.Lp E 2 μ))
    exact gsol.2
  rw [mem_closure_iff_seq_limit] at hgsol_mem
  rcases hgsol_mem with ⟨Gseq, hGseq_mem, hGseq_tendsto⟩
  let ψ : ℕ → E → ℝ := fun n => repFun ⟨Gseq n, hGseq_mem n⟩
  let hψ : ∀ n : ℕ, IsSmoothTestOn Ω (ψ n) := fun n => repSmooth ⟨Gseq n, hGseq_mem n⟩
  have hψ_eq : ∀ n : ℕ, smoothGradToLp hΩ (hψ n) = Gseq n := by
    intro n
    exact repEq ⟨Gseq n, hGseq_mem n⟩
  let Gbar : ℕ → H := fun n => e ⟨Gseq n, hGseq_mem n⟩
  have hGbar_tendsto : Tendsto Gbar atTop (nhds gsol) := by
    rw [tendsto_subtype_rng]
    simpa [Gbar, e] using hGseq_tendsto
  let uLp : MeasureTheory.Lp ℝ 2 μ := U gsol
  let u : E → ℝ := uLp
  have hu_memLp : MemLp u 2 μ := Lp.memLp uLp
  have hu_toLp : hu_memLp.toLp u = uLp := by
    change hu_memLp.toLp (uLp : E → ℝ) = uLp
    exact Lp.toLp_coeFn uLp hu_memLp
  have hU_tendsto : Tendsto (fun n => U (Gbar n)) atTop (nhds uLp) := by
    simpa [uLp] using ((U.continuous.continuousAt).tendsto.comp hGbar_tendsto)
  have hψLp_tendsto : Tendsto (fun n => smoothFunToLp hΩ (hψ n)) atTop (nhds uLp) := by
    refine Tendsto.congr' (Filter.Eventually.of_forall ?_) hU_tendsto
    intro n
    exact hU_eq ⟨Gseq n, hGseq_mem n⟩
  letI : Fact (1 ≤ (2 : ENNReal)) := ⟨by norm_num⟩
  have hψ_fun_tendsto :
      Tendsto (fun n => eLpNorm (fun x => ψ n x - u x) 2 μ) atTop (nhds 0) := by
    have hψLp_tendsto' :
        Tendsto
          (fun n => ((smoothTestWitness hΩ (hψ n)).memLp).toLp (ψ n))
          atTop (nhds (hu_memLp.toLp u)) := by
      simpa [smoothFunToLp, hu_toLp] using hψLp_tendsto
    exact
      (Lp.tendsto_Lp_iff_tendsto_eLpNorm'' (f := ψ)
        (f_ℒp := fun n => (smoothTestWitness hΩ (hψ n)).memLp)
        (f_lim := u) (f_lim_ℒp := hu_memLp)).1 hψLp_tendsto'
  have hgradLp_tendsto :
      Tendsto (fun n => smoothGradToLp hΩ (hψ n))
        atTop (nhds Gsol) := by
    refine Tendsto.congr' (Filter.Eventually.of_forall fun n => (hψ_eq n).symm) ?_
    simpa [Gsol] using hGseq_tendsto
  have hgrad_vec_tendsto :
      Tendsto
        (fun n => eLpNorm (fun x => (smoothTestWitness hΩ (hψ n)).weakGrad x - Gsol x) 2 μ)
        atTop (nhds 0) := by
    have hgradLp_tendsto' :
        Tendsto
          (fun n => ((smoothTestWitness hΩ (hψ n)).weakGrad_memLp).toLp
            ((smoothTestWitness hΩ (hψ n)).weakGrad))
          atTop (nhds ((Lp.memLp Gsol).toLp Gsol)) := by
      simpa [smoothGradToLp, gradLpOfWitness] using hgradLp_tendsto
    exact
      (Lp.tendsto_Lp_iff_tendsto_eLpNorm'' (f := fun n => (smoothTestWitness hΩ (hψ n)).weakGrad)
        (f_ℒp := fun n => (smoothTestWitness hΩ (hψ n)).weakGrad_memLp)
        (f_lim := Gsol) (f_lim_ℒp := Lp.memLp Gsol)).1
        hgradLp_tendsto'
  have hgsol_comp_memLp : ∀ i : Fin d, MemLp (fun x => Gsol x i) 2 μ := by
    intro i
    exact (Lp.memLp Gsol).eval_piLp i
  have hψ_grad_comp_memLp :
      ∀ n : ℕ, ∀ i : Fin d,
        MemLp
          (fun x => ((smoothTestWitness hΩ (hψ n)).weakGrad x).ofLp i - Gsol x i)
          2 μ := by
    intro n i
    simpa [PiLp.toLp_apply] using
      ((smoothTestWitness hΩ (hψ n)).weakGrad_component_memLp i).sub (hgsol_comp_memLp i)
  have hgrad_comp_tendsto :
      ∀ i : Fin d,
        Tendsto
          (fun n =>
            eLpNorm
              (fun x => ((smoothTestWitness hΩ (hψ n)).weakGrad x).ofLp i - Gsol x i)
              2 μ)
          atTop (nhds 0) := by
    intro i
    have hupper :
        ∀ n,
          eLpNorm
              (fun x => ((smoothTestWitness hΩ (hψ n)).weakGrad x).ofLp i - Gsol x i)
              2 μ ≤
            eLpNorm (fun x => (smoothTestWitness hΩ (hψ n)).weakGrad x - Gsol x) 2 μ := by
      intro n
      exact eLpNorm_mono_ae <|
        Filter.Eventually.of_forall fun x => by
          simpa [PiLp.toLp_apply] using
            (PiLp.norm_apply_le ((smoothTestWitness hΩ (hψ n)).weakGrad x - Gsol x) i)
    exact
      tendsto_of_tendsto_of_tendsto_of_le_of_le tendsto_const_nhds hgrad_vec_tendsto
        (fun n => zero_le _) hupper
  have hu_memLp' : MemLp u (ENNReal.ofReal 2) (volume.restrict Ω) := by
    simpa using hu_memLp
  have hgsol_comp_memLp' :
      ∀ i : Fin d, MemLp (fun x => Gsol x i) (ENNReal.ofReal 2) (volume.restrict Ω) := by
    intro i
    simpa using hgsol_comp_memLp i
  have hψ_fun_memLp' :
      ∀ n, MemLp (fun x => ψ n x - u x) (ENNReal.ofReal 2) (volume.restrict Ω) := by
    intro n
    simpa using ((smoothTestWitness hΩ (hψ n)).memLp).sub hu_memLp
  have hψ_fun_tendsto' :
      Tendsto (fun n => eLpNorm (fun x => ψ n x - u x) (ENNReal.ofReal 2) (volume.restrict Ω))
        atTop (nhds 0) := by
    simpa using hψ_fun_tendsto
  have hψ_grad_comp_memLp' :
      ∀ n : ℕ, ∀ i : Fin d,
        MemLp
          (fun x => ((smoothTestWitness hΩ (hψ n)).weakGrad x).ofLp i - Gsol x i)
          (ENNReal.ofReal 2) (volume.restrict Ω) := by
    intro n i
    simpa using hψ_grad_comp_memLp n i
  have hgrad_comp_tendsto' :
      ∀ i : Fin d,
        Tendsto
          (fun n =>
            eLpNorm
              (fun x => ((smoothTestWitness hΩ (hψ n)).weakGrad x).ofLp i - Gsol x i)
              (ENNReal.ofReal 2) (volume.restrict Ω))
          atTop (nhds 0) := by
    intro i
    simpa using hgrad_comp_tendsto i
  have hu_isWeakGrad : ∀ i : Fin d, HasWeakPartialDeriv i (fun x => Gsol x i) u Ω := by
    intro i
    exact HasWeakPartialDeriv.of_eLpNormApprox_p hΩ (p := 2) (by norm_num)
      hu_memLp' (hgsol_comp_memLp' i)
      (fun n => (smoothTestWitness hΩ (hψ n)).isWeakGrad i)
      hψ_fun_memLp'
      hψ_fun_tendsto'
      (fun n => hψ_grad_comp_memLp' n i)
      (hgrad_comp_tendsto' i)
  let hwu : MemW1pWitness 2 u Ω := {
    memLp := hu_memLp
    weakGrad := fun x => WithLp.toLp 2 fun i => Gsol x i
    weakGrad_component_memLp := hgsol_comp_memLp
    isWeakGrad := hu_isWeakGrad }
  have hwu_gradEq : gradLpOfWitness hwu = Gsol := by
    have hgrad_ae :
        hwu.weakGrad =ᵐ[μ] fun x => Gsol x := by
      filter_upwards with x
      change WithLp.toLp 2 (WithLp.ofLp (Gsol x)) = Gsol x
      exact WithLp.toLp_ofLp (p := 2) (x := Gsol x)
    calc
      gradLpOfWitness hwu =
          (Lp.memLp Gsol).toLp (fun x => Gsol x) := by
            simpa [gradLpOfWitness] using
              (MemLp.toLp_congr hwu.weakGrad_memLp (Lp.memLp Gsol)
                hgrad_ae)
      _ = Gsol := by
            exact Lp.toLp_coeFn _ _
  have hu0 : MemH01 u Ω := by
    refine ⟨hwu.memW1p, ⟨hwu, ψ, ?_, ?_, ?_, ?_, ?_⟩⟩
    · intro n
      exact (hψ n).1
    · intro n
      exact (hψ n).2.1
    · intro n
      exact (hψ n).2.2
    · simpa [μ] using hψ_fun_tendsto
    · intro i
      simpa [μ] using hgrad_comp_tendsto i
  have hsmooth_eq :
      ∀ {φ : E → ℝ} (hφ : IsSmoothTestOn Ω φ),
        bilinFormOfCoeff A hwu (smoothTestWitness hΩ hφ) = rhs φ := by
    intro φ hφ
    let gφS : S := ⟨smoothGradToLp hΩ hφ, ⟨φ, hφ, rfl⟩⟩
    let gφH : H := e gφS
    have hφ_gradEq :
        H.subtypeL gφH = gradLpOfWitness (smoothTestWitness hΩ hφ) := by
      simp [gφH, gφS, e, smoothGradToLp]
    calc
      bilinFormOfCoeff A hwu (smoothTestWitness hΩ hφ)
          = coeffBilinSubmodule A H gsol gφH := by
              calc
                bilinFormOfCoeff A hwu (smoothTestWitness hΩ hφ) =
                    ⟪coeffMulLpL A (gradLpOfWitness hwu),
                      gradLpOfWitness (smoothTestWitness hΩ hφ)⟫_ℝ := by
                        symm
                        exact coeffMulLpL_inner_gradLpOfWitness_eq_bilinFormOfCoeff A hwu
                          (smoothTestWitness hΩ hφ)
                _ = ⟪coeffMulLpL A Gsol, H.subtypeL gφH⟫_ℝ := by
                      rw [hwu_gradEq, hφ_gradEq]
                _ = ⟪coeffMulLpL A (H.subtypeL gsol), H.subtypeL gφH⟫_ℝ := by
                      rfl
                _ = coeffBilinSubmodule A H gsol gφH := by rfl
      _ = L gφH := hLM gφH
      _ = L0 gφS := by simpa [gφH] using hL_eq gφS
      _ = rhs φ := by
            apply rhs_eq_of_sameGrad (repSmooth gφS) hφ
            calc
              smoothGradToLp hΩ (repSmooth gφS) = (gφS : MeasureTheory.Lp E 2 μ) := repEq gφS
              _ = smoothGradToLp hΩ hφ := by rfl
  refine ⟨u, hu0, ?_⟩
  intro hwu' v hv0 hv
  have hv0_all := hv0
  rcases hv0 with ⟨_hv_mem, hvw, φ, hφ_smooth, hφ_cpt, hφ_sub, hφ_fun, hφ_grad⟩
  let hφtest : ∀ n : ℕ, IsSmoothTestOn Ω (φ n) := fun n =>
    ⟨hφ_smooth n, hφ_cpt n, hφ_sub n⟩
  let hφw : ∀ n : ℕ, MemW1pWitness 2 (φ n) Ω := fun n => smoothTestWitness hΩ (hφtest n)
  let haddiff : ∀ n : ℕ, MemW1pWitness 2 (fun x => φ n x + (-1) * v x) Ω := fun n =>
    (hφw n).add (hvw.smul (-1))
  let hdiff : ∀ n : ℕ, MemW1pWitness 2 (fun x => φ n x - v x) Ω := fun n => {
    memLp := by
      simpa [sub_eq_add_neg, Pi.smul_apply] using (haddiff n).memLp
    weakGrad := (haddiff n).weakGrad
    weakGrad_component_memLp := by
      intro i
      simpa [sub_eq_add_neg, Pi.smul_apply] using (haddiff n).weakGrad_component_memLp i
    isWeakGrad := by
      intro i
      simpa [sub_eq_add_neg, Pi.smul_apply] using (haddiff n).isWeakGrad i }
  have hgradDiffVec_tendsto :
      Tendsto (fun n => eLpNorm (fun x => smoothGradField (φ n) x - hvw.weakGrad x) 2 μ)
        atTop (nhds 0) := by
    exact tendsto_eLpNorm_vector_of_componentwise
      (fun n i => by
        simpa [smoothTestWitness, smoothGradField, PiLp.toLp_apply] using
          ((hφw n).weakGrad_component_memLp i).sub (hvw.weakGrad_component_memLp i))
      (fun i => by
        simpa [μ] using hφ_grad i)
  have hdiff_grad_tendsto :
      Tendsto (fun n => gradLpOfWitness (hdiff n)) atTop (nhds 0) := by
    let h0_mem : MemLp (fun _ : E => (0 : E)) 2 μ := MeasureTheory.MemLp.zero'
    have hgradDiffVec_tendsto_add :
        Tendsto (fun n => eLpNorm (fun x => (haddiff n).weakGrad x) 2 μ)
          atTop (nhds 0) := by
      refine Tendsto.congr' (Filter.Eventually.of_forall ?_) hgradDiffVec_tendsto
      intro n
      have hweak_eq :
          (fun x => (haddiff n).weakGrad x) =
            fun x => smoothGradField (φ n) x - hvw.weakGrad x := by
        funext x
        simp [haddiff, hφw, smoothTestWitness, smoothGradField, MemW1pWitness.add,
          MemW1pWitness.smul, sub_eq_add_neg]
      simp [hweak_eq]
    have hgradDiffVec_tendsto_add' :
        Tendsto (fun n => eLpNorm (fun x => (haddiff n).weakGrad x - (0 : E)) 2 μ)
          atTop (nhds 0) := by
      simpa [sub_eq_add_neg] using hgradDiffVec_tendsto_add
    have hLp_tendsto_add :
        Tendsto
          (fun n => ((haddiff n).weakGrad_memLp).toLp ((haddiff n).weakGrad))
          atTop (nhds (h0_mem.toLp (fun _ : E => (0 : E)))) := by
      exact
        (Lp.tendsto_Lp_iff_tendsto_eLpNorm'' (f := fun n => (haddiff n).weakGrad)
          (f_ℒp := fun n => (haddiff n).weakGrad_memLp)
          (f_lim := fun _ : E => (0 : E)) (f_lim_ℒp := h0_mem)).2 <| by
            exact hgradDiffVec_tendsto_add'
    simpa [hdiff, gradLpOfWitness] using hLp_tendsto_add
  have hdiff_seminorm_tendsto :
      Tendsto
        (fun n => (∫ x, ‖(hdiff n).weakGrad x‖ ^ (2 : ℝ) ∂μ) ^ (1 / (2 : ℝ)))
        atTop (nhds 0) := by
    have hnorm_tendsto :
        Tendsto (fun n => ‖gradLpOfWitness (hdiff n)‖) atTop (nhds 0) := by
      simpa using ((continuous_norm.tendsto (0 : MeasureTheory.Lp E 2 μ)).comp hdiff_grad_tendsto)
    refine Tendsto.congr' (Filter.Eventually.of_forall ?_) hnorm_tendsto
    intro n
    exact norm_gradLpOfWitness_eq (hdiff n)
  let seminormDiff : ℕ → ℝ := fun n =>
    (∫ x, ‖(hdiff n).weakGrad x‖ ^ (2 : ℝ) ∂μ) ^ (1 / (2 : ℝ))
  have hseminormDiff_nonneg : ∀ n, 0 ≤ seminormDiff n := by
    intro n
    positivity
  have hrhsDiff_tendsto :
      Tendsto (fun n => rhs (φ n) - rhs v) atTop (nhds 0) := by
    have hbound_n :
        ∀ n, |rhs (φ n) - rhs v| ≤ C_rhs * seminormDiff n := by
      intro n
      have hφ0n : MemH01 (φ n) Ω := smoothTest_memH01 hΩ (hφtest n)
      have hvNeg0 : MemH01 (fun x => (-1) * v x) Ω := by
        simpa [Pi.smul_apply, smul_eq_mul] using (MemW01p.smul (-1) hv0_all)
      have hsub0n : MemH01 (fun x => φ n x - v x) Ω := MemW01p.sub hφ0n hv0_all
      have hsubEq : rhs (fun x => φ n x - v x) = rhs (φ n) - rhs v := by
        calc
          rhs (fun x => φ n x - v x) = rhs (fun x => φ n x + (-1) * v x) := by
            congr
            ext x
            ring
          _ = rhs (φ n) + rhs (fun x => (-1) * v x) := hF_add (φ n) (fun x => (-1) * v x) hφ0n hvNeg0
          _ = rhs (φ n) + (-1) * rhs v := by rw [hF_smul (-1) v hv0_all]
          _ = rhs (φ n) - rhs v := by ring
      have hbound := hF_bound (fun x => φ n x - v x) hsub0n (hdiff n)
      rw [hsubEq] at hbound
      simpa [seminormDiff, μ] using hbound
    have hupper :
        Tendsto (fun n => C_rhs * seminormDiff n) atTop (nhds 0) := by
      simpa [seminormDiff] using Tendsto.const_mul C_rhs hdiff_seminorm_tendsto
    have habs :
        Tendsto (fun n => |rhs (φ n) - rhs v|) atTop (nhds 0) :=
      squeeze_zero (fun n => abs_nonneg _) hbound_n hupper
    exact (tendsto_zero_iff_abs_tendsto_zero _).2 habs
  let seminormU : ℝ :=
    (∫ x, ‖hwu.weakGrad x‖ ^ (2 : ℝ) ∂μ) ^ (1 / (2 : ℝ))
  have hbilinDiff_tendsto :
      Tendsto
        (fun n => bilinFormOfCoeff A hwu (hφw n) - bilinFormOfCoeff A hwu hvw)
        atTop (nhds 0) := by
    have hbound_n :
        ∀ n,
          |bilinFormOfCoeff A hwu (hφw n) - bilinFormOfCoeff A hwu hvw|
            ≤ A.Λ * seminormU * seminormDiff n := by
      intro n
      have hsplit :
          bilinFormOfCoeff A hwu (hdiff n) =
            bilinFormOfCoeff A hwu (hφw n) - bilinFormOfCoeff A hwu hvw := by
        calc
          bilinFormOfCoeff A hwu (hdiff n)
              = bilinFormOfCoeff A hwu (haddiff n) := by
                  unfold bilinFormOfCoeff
                  apply integral_congr_ae
                  filter_upwards with x
                  simp [bilinFormIntegrandOfCoeff, hdiff, haddiff]
          _ = bilinFormOfCoeff A hwu (hφw n) + bilinFormOfCoeff A hwu (hvw.smul (-1)) := by
                  rw [show haddiff n = (hφw n).add (hvw.smul (-1)) by simp [haddiff]]
                  rw [bilinFormOfCoeff_add_right A hwu (hφw n) (hvw.smul (-1))]
          _ = bilinFormOfCoeff A hwu (hφw n) + (-1) * bilinFormOfCoeff A hwu hvw := by
                rw [bilinFormOfCoeff_smul_right (-1) A hwu hvw]
          _ = bilinFormOfCoeff A hwu (hφw n) - bilinFormOfCoeff A hwu hvw := by
                ring
      have hbound := bilinForm_bound A hwu (hdiff n)
      rw [hsplit] at hbound
      simpa [seminormU, seminormDiff, μ, mul_assoc] using hbound
    have hupper :
        Tendsto (fun n => A.Λ * seminormU * seminormDiff n) atTop (nhds 0) := by
      simpa [seminormU, seminormDiff, mul_assoc] using
        Tendsto.const_mul (A.Λ * seminormU) hdiff_seminorm_tendsto
    have habs :
        Tendsto
          (fun n => |bilinFormOfCoeff A hwu (hφw n) - bilinFormOfCoeff A hwu hvw|)
          atTop (nhds 0) :=
      squeeze_zero (fun n => abs_nonneg _) hbound_n hupper
    exact (tendsto_zero_iff_abs_tendsto_zero _).2 habs
  have hLhs_tendsto :
      Tendsto (fun n => bilinFormOfCoeff A hwu (hφw n))
        atTop (nhds (bilinFormOfCoeff A hwu hvw)) := by
    convert hbilinDiff_tendsto.add tendsto_const_nhds using 1
    · ext n
      ring_nf
    · ring_nf
  have hRhs_tendsto :
      Tendsto (fun n => rhs (φ n)) atTop (nhds (rhs v)) := by
    convert hrhsDiff_tendsto.add tendsto_const_nhds using 1
    · ext n
      ring_nf
    · ring_nf
  have hEq_tendsto :
      Tendsto (fun n => bilinFormOfCoeff A hwu (hφw n)) atTop (nhds (rhs v)) := by
    refine Tendsto.congr' (Filter.Eventually.of_forall ?_) hRhs_tendsto
    intro n
    exact (hsmooth_eq (hφtest n)).symm
  have hChosen : bilinFormOfCoeff A hwu hvw = rhs v :=
    tendsto_nhds_unique hLhs_tendsto hEq_tendsto
  calc
    bilinFormOfCoeff A hwu' hv = bilinFormOfCoeff A hwu hv := by
      exact bilinFormOfCoeff_eq_left hΩ A hwu' hwu hv
    _ = bilinFormOfCoeff A hwu hvw := by
      exact bilinFormOfCoeff_eq_right hΩ A hwu hv hvw
    _ = rhs v := hChosen

\end{lstlisting}
\begin{proof}
Let $\mu := \mathrm{vol}\!\restriction\!\Omega$ and let $S = S(\Omega) \leq \Lp{2}(\Omega; E)$ be the smooth-gradient submodule with closure $H = H(\Omega)$. By the Poincar\'e inequality on bounded domains in dimension $d \geq 2$, there is a constant $C_P > 0$ with
\[
  \|u\|_{\Lp{2}} \leq C_P\,\|\nabla u\|_{\Lp{2}} \quad \text{for every smooth test } u \text{ on } \Omega.
\]
For each $g \in S$ choose, via Classical choice, a smooth test $u_g$ with $[\nabla u_g] = g$. The Poincar\'e inequality, lifted via $\mathrm{toReal}$, gives $\|[u_g]\|_{\Lp{2}} \leq C_P \|g\|$, and shows that two such choices with the same gradient class differ by a function with vanishing $\Lp{2}$-class.

Define linear maps on $S$ by
\[
  L_0(g) := \mathrm{rhs}(u_g), \qquad U_0(g) := [u_g] \in \Lp{2}(\Omega; \R).
\]
The hypotheses on $\mathrm{rhs}$ and the Poincar\'e inequality, together with the fact that two representatives with the same gradient class give the same value of $\mathrm{rhs}$ (because their difference has zero gradient class and the bound forces $|\mathrm{rhs}(u_1) - \mathrm{rhs}(u_2)| \leq C_{\mathrm{rhs}} \cdot 0 = 0$), make $L_0$ and $U_0$ well-defined linear maps with bounds
\[
  |L_0(g)| \leq C_{\mathrm{rhs}}\|g\|, \qquad \|U_0(g)\| \leq C_P\|g\|.
\]
Since the inclusion $S \hookrightarrow H$ has dense range (as $H$ is the closure of $S$), $L_0$ and $U_0$ extend uniquely to continuous maps $L : H \to \R$, $U : H \to \Lp{2}(\Omega;\R)$.

By the coercivity of the submodule bilinear form established above ($B_{A,H}$ is bounded and coercive on the Hilbert space $H$, which is complete; see \leanname{coeffBilinSubmodule\_coercive}), the Lax--Milgram theorem produces a unique $g_{\mathrm{sol}} \in H$ with
\[
  B_{A,H}(g_{\mathrm{sol}}, w) = L(w) \quad \text{for all } w \in H.
\]
Set $G_{\mathrm{sol}}$ to be the underlying $\Lp{2}(\Omega;E)$ element of $g_{\mathrm{sol}}$ and let $u_{\mathrm{Lp}} := U(g_{\mathrm{sol}}) \in \Lp{2}(\Omega;\R)$ with representative $u : E \to \R$.

Approximate $g_{\mathrm{sol}}$ by a sequence $G_n \in S$ (closure), with corresponding smooth tests $\psi_n$ such that $[\nabla \psi_n] = G_n$. By continuity of $U$, $[\psi_n] \to u_{\mathrm{Lp}}$ in $\Lp{2}$. By the componentwise convergence lemma, the components of $\nabla \psi_n$ converge to the corresponding components of $G_{\mathrm{sol}}$ in $\Lp{2}$, and $\psi_n \to u$ in $\Lp{2}$. By the standard transition principle for weak partial derivatives under $\Lp{p}$-approximation, $u$ has weak gradient given by $G_{\mathrm{sol}}$, providing a witness $h_u$ with $[\nabla u] = G_{\mathrm{sol}}$. Since $\psi_n$ are smooth tests on $\Omega$, $u \in H^{1}_{0}(\Omega)$.

For any smooth test $\varphi$ on $\Omega$ with corresponding $g_\varphi \in S$,
\[
  \mathfrak{a}_A(u, \varphi) = \langle T_A [\nabla u], [\nabla \varphi]\rangle_{\Lp{2}} = B_{A,H}(g_{\mathrm{sol}}, g_\varphi) = L(g_\varphi) = L_0(g_\varphi) = \mathrm{rhs}(\varphi).
\]
Finally, by approximation: any $v \in H^{1}_{0}(\Omega)$ admits an approximating sequence $\varphi_n$ of smooth tests in the $H^{1}_{0}$-sense; the bilinear form is continuous against gradient seminorms, and $\mathrm{rhs}$ is also bounded against the gradient seminorm. Passing to the limit in $\mathfrak{a}_A(u, \varphi_n) = \mathrm{rhs}(\varphi_n)$ (using independence of the bilinear form on the choice of witness in the right slot) gives $\mathfrak{a}_A(u, v) = \mathrm{rhs}(v)$ for any witness of $v$. This completes the proof.
\end{proof}

\begin{theorem}[\leanname{weakProblem\_exists\_of\_divergenceData}]
Under the hypotheses of Theorem~\ref{thm:weakProblem-exists} and $F \in \Lp{2}(\Omega; E)$, there exists a zero-Dirichlet weak solution of $-\dv(a\nabla u) = -\dv F$, i.e.\ a weak solution for the raw functional $v \mapsto L_F(v)$.
\end{theorem}

\begin{lstlisting}[style=Matlab-editor,basicstyle=\mlttfamily,escapechar=",mlshowsectionrules=true,mathescape=true,morecomment={[s]{\%\{}{\%\}}},language=lean,numbers=none]
/-- Existence of zero-Dirichlet weak solutions for divergence-form right-hand
side data `div F`. -/
theorem weakProblem_exists_of_divergenceData
    {Ω : Set E}
    (hd : 2 ≤ d)
    (hΩ : IsOpen Ω) (hΩ_bdd : Bornology.IsBounded Ω)
    (A : EllipticCoeff d Ω)
    {F : E → E} (hF : MemLp F 2 (volume.restrict Ω)) :
    ∃ u : E → ℝ,
      IsWeakSolution (d := d)
        ⟨Ω, hΩ, hΩ_bdd, A, weakProblemRHSOfField (Ω := Ω) F⟩ u := by
  refine weakProblem_exists hd hΩ hΩ_bdd A (weakProblemRHSOfField (Ω := Ω) F) ?_ ?_ ?_
  · intro u v hu0 hv0
    let hu : MemW1pWitness 2 u Ω :=
      DeGiorgi.MemW1p.someWitness (MemW01p.memW1p hu0)
    let hv : MemW1pWitness 2 v Ω :=
      DeGiorgi.MemW1p.someWitness (MemW01p.memW1p hv0)
    have huv0 : MemH01 (fun x => u x + v x) Ω := MemW01p.add hu0 hv0
    let huv : MemW1pWitness 2 (fun x => u x + v x) Ω := hu.add hv
    calc
      weakProblemRHSOfField (Ω := Ω) F (fun x => u x + v x)
          = divergenceRHSOfField F huv := by
              exact weakProblemRHSOfField_eq_of_memH01 hΩ huv0 huv
      _ = divergenceRHSOfField F hu + divergenceRHSOfField F hv := by
            exact divergenceRHSOfField_add hF hu hv
      _ = weakProblemRHSOfField (Ω := Ω) F u + weakProblemRHSOfField (Ω := Ω) F v := by
            rw [weakProblemRHSOfField_eq_of_memH01 hΩ hu0 hu,
              weakProblemRHSOfField_eq_of_memH01 hΩ hv0 hv]
  · intro c u hu0
    let hu : MemW1pWitness 2 u Ω :=
      DeGiorgi.MemW1p.someWitness (MemW01p.memW1p hu0)
    have hcu0 : MemH01 (fun x => c * u x) Ω := by
      simpa [Pi.smul_apply, smul_eq_mul] using (MemW01p.smul c hu0)
    let hcu : MemW1pWitness 2 (fun x => c * u x) Ω := hu.smul c
    calc
      weakProblemRHSOfField (Ω := Ω) F (fun x => c * u x)
          = divergenceRHSOfField F hcu := by
              exact weakProblemRHSOfField_eq_of_memH01 hΩ hcu0 hcu
      _ = c * divergenceRHSOfField F hu := by
            exact divergenceRHSOfField_smul c hu
      _ = c * weakProblemRHSOfField (Ω := Ω) F u := by
            rw [weakProblemRHSOfField_eq_of_memH01 hΩ hu0 hu]
  · refine ⟨(∫ x, ‖F x‖ ^ (2 : ℝ) ∂(volume.restrict Ω)) ^ (1 / (2 : ℝ)), by positivity, ?_⟩
    intro v hv0 hv
    exact weakProblemRHSOfField_bound hΩ hF v hv0 hv

\end{lstlisting}
\begin{proof}
We invoke the abstract existence theorem
Theorem~\ref{thm:weakProblem-exists} (\leanname{weakProblem\_exists}) with
$\mathrm{rhs}$ chosen to be $\leanname{weakProblemRHSOfField}\,F$, the
linear functional $v \mapsto \int_{\Omega} \langle F, \nabla v\rangle\,dx$
acting on $H^{1}_{0}(\Omega)$. The abstract theorem requires three inputs.

\emph{Additivity.} For $u, v \in H^{1}_{0}(\Omega)$ pick canonical witnesses
through \leanname{MemW1p.someWitness}; the sum witness for $u + v$ is the
componentwise sum (\leanname{MemW1pWitness.add}). Both
$\mathrm{rhs}(u+v)$ and $\mathrm{rhs}(u) + \mathrm{rhs}(v)$ unfold via
\leanname{weakProblemRHSOfField\_eq\_of\_memH01} into integrals against the
weak gradient, so additivity reduces to
\leanname{divergenceRHSOfField\_add} applied to those witnesses.

\emph{Homogeneity.} Identical reasoning with the scalar witness
$\mathrm{MemW1pWitness.smul}$ and \leanname{divergenceRHSOfField\_smul}
gives $\mathrm{rhs}(c\,u) = c\,\mathrm{rhs}(u)$ for every $c \in \R$.

\emph{Boundedness against the gradient seminorm.} The integral
$\bigl|\int_{\Omega} \langle F, \nabla v\rangle\bigr|$ is controlled by
$\eLpNormText{F}_{L^2}\,\eLpNormText{\nabla v}_{L^2}$ via Cauchy--Schwarz; this is
recorded as \leanname{weakProblemRHSOfField\_bound}. We take
$C_F := \eLpNormText{F}_{L^2(\Omega)}$, which is finite by the hypothesis
$F \in L^2(\Omega; E)$.

With these three properties, Theorem~\ref{thm:weakProblem-exists} produces a
zero-Dirichlet weak solution $u \in H^{1}_{0}(\Omega)$ of
$-\dv(a\nabla u) = -\dv F$.
\end{proof}

\begin{theorem}[\leanname{dirichletProblem\_exists\_of\_divergenceData}, \leanname{dirichletProblem\_exists\_of\_divergenceData'}]
Under the same hypotheses, with prescribed boundary datum $u_0 \in W^{1,2}(\Omega)$ and $F \in \Lp{2}(\Omega;E)$, there exists $u \in W^{1,2}(\Omega)$ with $u - u_0 \in H^{1}_{0}(\Omega)$ and
\[
  \mathfrak{a}_A(u, v) = L_F(v) \quad \text{for all } v \in H^{1}_{0}(\Omega).
\]
\end{theorem}

\begin{lstlisting}[style=Matlab-editor,basicstyle=\mlttfamily,escapechar=",mlshowsectionrules=true,mathescape=true,morecomment={[s]{\%\{}{\%\}}},language=lean,numbers=none]
/-- Existence of weak solutions to the inhomogeneous Dirichlet problem with
boundary datum `u₀` and divergence-form right-hand side `div F`. -/
theorem dirichletProblem_exists_of_divergenceData
    {Ω : Set E}
    (hd : 2 ≤ d)
    (hΩ : IsOpen Ω) (hΩ_bdd : Bornology.IsBounded Ω)
    (A : EllipticCoeff d Ω)
    {u₀ : E → ℝ} (hu₀ : MemW1p 2 u₀ Ω)
    {F : E → E} (hF : MemLp F 2 (volume.restrict Ω)) :
    ∃ u : E → ℝ,
      MemW1p 2 u Ω ∧
      MemW01p 2 (fun x => u x - u₀ x) Ω ∧
      ∀ (hu : MemW1pWitness 2 u Ω) (v : E → ℝ), MemH01 v Ω →
        ∀ (hv : MemW1pWitness 2 v Ω),
          bilinFormOfCoeff A hu hv = divergenceRHSOfField F hv := by
  let hu₀w : MemW1pWitness 2 u₀ Ω := DeGiorgi.MemW1p.someWitness hu₀
  let rhs := weakProblemRHSOfFieldAndDatum (A := A) (Ω := Ω) F hu₀w
  obtain ⟨w, hwsol⟩ :=
    weakProblem_exists hd hΩ hΩ_bdd A rhs
      (by
        intro u v hu0 hv0
        let hu : MemW1pWitness 2 u Ω :=
          DeGiorgi.MemW1p.someWitness (MemW01p.memW1p hu0)
        let hv : MemW1pWitness 2 v Ω :=
          DeGiorgi.MemW1p.someWitness (MemW01p.memW1p hv0)
        have huv0 : MemH01 (fun x => u x + v x) Ω := MemW01p.add hu0 hv0
        let huv : MemW1pWitness 2 (fun x => u x + v x) Ω := hu.add hv
        calc
          rhs (fun x => u x + v x)
              = divergenceRHSOfField F huv - bilinFormOfCoeff A hu₀w huv := by
                  exact weakProblemRHSOfFieldAndDatum_eq_of_memH01 hΩ A hu₀w huv0 huv
          _ = (divergenceRHSOfField F hu + divergenceRHSOfField F hv) -
                (bilinFormOfCoeff A hu₀w hu + bilinFormOfCoeff A hu₀w hv) := by
                  rw [divergenceRHSOfField_add hF hu hv, bilinFormOfCoeff_add_right A hu₀w hu hv]
          _ = weakProblemRHSOfFieldAndDatum (A := A) (Ω := Ω) F hu₀w u +
                weakProblemRHSOfFieldAndDatum (A := A) (Ω := Ω) F hu₀w v := by
                  rw [weakProblemRHSOfFieldAndDatum_eq_of_memH01 hΩ A hu₀w hu0 hu,
                    weakProblemRHSOfFieldAndDatum_eq_of_memH01 hΩ A hu₀w hv0 hv]
                  ring
          _ = rhs u + rhs v := by rfl)
      (by
        intro c u hu0
        let hu : MemW1pWitness 2 u Ω :=
          DeGiorgi.MemW1p.someWitness (MemW01p.memW1p hu0)
        have hcu0 : MemH01 (fun x => c * u x) Ω := by
          simpa [Pi.smul_apply, smul_eq_mul] using (MemW01p.smul c hu0)
        let hcu : MemW1pWitness 2 (fun x => c * u x) Ω := hu.smul c
        calc
          rhs (fun x => c * u x)
              = divergenceRHSOfField F hcu - bilinFormOfCoeff A hu₀w hcu := by
                  exact weakProblemRHSOfFieldAndDatum_eq_of_memH01 hΩ A hu₀w hcu0 hcu
          _ = c * divergenceRHSOfField F hu - c * bilinFormOfCoeff A hu₀w hu := by
                rw [divergenceRHSOfField_smul c hu, bilinFormOfCoeff_smul_right c A hu₀w hu]
          _ = c * weakProblemRHSOfFieldAndDatum (A := A) (Ω := Ω) F hu₀w u := by
                rw [weakProblemRHSOfFieldAndDatum_eq_of_memH01 hΩ A hu₀w hu0 hu]
                ring
          _ = c * rhs u := by rfl)
      <| by
        have hCF_nonneg :
            0 ≤ (∫ x, ‖F x‖ ^ (2 : ℝ) ∂(volume.restrict Ω)) ^ (1 / (2 : ℝ)) := by
          positivity
        have hCU_nonneg :
            0 ≤ A.Λ * (∫ x, ‖hu₀w.weakGrad x‖ ^ (2 : ℝ) ∂(volume.restrict Ω)) ^
              (1 / (2 : ℝ)) := by
          refine mul_nonneg A.Λ_nonneg ?_
          positivity
        refine ⟨((∫ x, ‖F x‖ ^ (2 : ℝ) ∂(volume.restrict Ω)) ^ (1 / (2 : ℝ)) +
            A.Λ * (∫ x, ‖hu₀w.weakGrad x‖ ^ (2 : ℝ) ∂(volume.restrict Ω)) ^
              (1 / (2 : ℝ))), add_nonneg hCF_nonneg hCU_nonneg, ?_⟩
        intro v hv0 hv
        simpa [rhs] using weakProblemRHSOfFieldAndDatum_bound hΩ A hF hu₀w v hv0 hv
  have hw0 : MemH01 w Ω := hwsol.left
  let hw : MemW1pWitness 2 w Ω :=
    DeGiorgi.MemW1p.someWitness (MemW01p.memW1p hw0)
  let u : E → ℝ := fun x => w x + u₀ x
  let hu : MemW1pWitness 2 u Ω := hw.add hu₀w
  refine ⟨u, hu.memW1p, ?_, ?_⟩
  · simpa [u] using hw0
  · intro hu' v hv0 hv
    calc
      bilinFormOfCoeff A hu' hv = bilinFormOfCoeff A hu hv := by
        exact bilinFormOfCoeff_eq_left hΩ A hu' hu hv
      _ = bilinFormOfCoeff A hw hv + bilinFormOfCoeff A hu₀w hv := by
        simpa [hu] using bilinFormOfCoeff_add_left A hw hu₀w hv
      _ = rhs v + bilinFormOfCoeff A hu₀w hv := by
        rw [hwsol.right hw v hv0 hv]
      _ = weakProblemRHSOfFieldAndDatum (A := A) (Ω := Ω) F hu₀w v +
            bilinFormOfCoeff A hu₀w hv := by
        rfl
      _ = divergenceRHSOfField F hv := by
        rw [weakProblemRHSOfFieldAndDatum_eq_of_memH01 hΩ A hu₀w hv0 hv]
        ring

/-- Dirichlet existence theorem stated using the weak-solution predicate. -/
theorem dirichletProblem_exists_of_divergenceData'
    {Ω : Set E}
    (hd : 2 ≤ d)
    (hΩ : IsOpen Ω) (hΩ_bdd : Bornology.IsBounded Ω)
    (A : EllipticCoeff d Ω)
    {u₀ : E → ℝ} (hu₀ : MemW1p 2 u₀ Ω)
    {F : E → E} (hF : MemLp F 2 (volume.restrict Ω)) :
    ∃ u : E → ℝ, IsDirichletWeakSolutionOfDivergenceData A u₀ F u := by
  obtain ⟨u, hu_mem, hdiff, hweak⟩ :=
    dirichletProblem_exists_of_divergenceData hd hΩ hΩ_bdd A hu₀ hF
  exact ⟨u, hu_mem, hdiff, hweak⟩

\end{lstlisting}
\begin{proof}
Choose a witness $h_{u_0}$ for $u_0$. Apply the previous theorem to the shifted RHS $\mathrm{rhs}(v) = L_F(v) - \mathfrak{a}_A(u_0, v)$ to obtain a zero-Dirichlet weak solution $w$. Set $u := w + u_0$. Then $u - u_0 = w \in H^{1}_{0}(\Omega)$, and bilinearity of $\mathfrak{a}_A$ gives
\[
  \mathfrak{a}_A(u, v) = \mathfrak{a}_A(w, v) + \mathfrak{a}_A(u_0, v) = \mathrm{rhs}(v) + \mathfrak{a}_A(u_0, v) = L_F(v).
\]
\end{proof}

\begin{theorem}[\leanname{weakSolution\_stability}: energy estimate]
Let $A$ be elliptic on $\Omega$, $u \in H^{1}_{0}(\Omega)$ a weak solution against a functional $\mathrm{rhs}$ bounded by $C_F\|\nabla v\|_{\Lp{2}}$. Then
\[
  \|\nabla u\|_{\Lp{2}} \leq \frac{C_F}{\lam}.
\]
\end{theorem}

\begin{lstlisting}[style=Matlab-editor,basicstyle=\mlttfamily,escapechar=",mlshowsectionrules=true,mathescape=true,morecomment={[s]{\%\{}{\%\}}},language=lean,numbers=none]
/-- Stability estimate for weak solutions. -/
theorem weakSolution_stability
    {Ω : Set E}
    (A : EllipticCoeff d Ω)
    {u : E → ℝ}
    (hu0 : MemH01 u Ω)
    (hwu : MemW1pWitness 2 u Ω)
    {rhs : (E → ℝ) → ℝ}
    {C_F : ℝ} (hCF : 0 ≤ C_F)
    (hF : ∀ v : E → ℝ, MemH01 v Ω →
      ∀ hwv : MemW1pWitness 2 v Ω,
      |rhs v| ≤ C_F *
        (∫ x, ‖hwv.weakGrad x‖ ^ (2 : ℝ) ∂(volume.restrict Ω)) ^ (1 / (2 : ℝ)))
    (h_ws : bilinFormOfCoeff A hwu hwu = rhs u) :
    (∫ x, ‖hwu.weakGrad x‖ ^ (2 : ℝ) ∂(volume.restrict Ω)) ^ (1 / (2 : ℝ)) ≤
    A.lam⁻¹ * C_F := by
  let μ : Measure E := volume.restrict Ω
  let J : ℝ := ∫ x, ‖hwu.weakGrad x‖ ^ (2 : ℕ) ∂μ
  have hJ_nonneg : 0 ≤ J := by
    dsimp [J]
    refine integral_nonneg ?_
    intro x
    positivity
  have hnorm_eq :
      ‖gradLpOfWitness hwu‖ = Real.sqrt J := by
    calc
      ‖gradLpOfWitness hwu‖
          = (∫ x, ‖hwu.weakGrad x‖ ^ (2 : ℝ) ∂μ) ^ (1 / (2 : ℝ)) := by
              simpa [μ] using norm_gradLpOfWitness_eq hwu
      _ = J ^ (1 / (2 : ℝ)) := by
            congr 1
            dsimp [J]
            apply integral_congr_ae
            filter_upwards with x
            exact Real.rpow_natCast ‖hwu.weakGrad x‖ 2
      _ = Real.sqrt J := by
            rw [Real.sqrt_eq_rpow]
  have hnorm_sq :
      ‖gradLpOfWitness hwu‖ ^ (2 : ℕ) = J := by
    calc
      ‖gradLpOfWitness hwu‖ ^ (2 : ℕ) = (Real.sqrt J) ^ (2 : ℕ) := by rw [hnorm_eq]
      _ = J := by simpa [pow_two] using Real.sq_sqrt hJ_nonneg
  have hcoercive :
      A.lam * ‖gradLpOfWitness hwu‖ ^ (2 : ℕ) ≤ bilinFormOfCoeff A hwu hwu := by
    calc
      A.lam * ‖gradLpOfWitness hwu‖ ^ (2 : ℕ) = A.lam * J := by rw [hnorm_sq]
      _ = A.lam * ∫ x, ‖hwu.weakGrad x‖ ^ (2 : ℕ) ∂μ := by rfl
      _ ≤ bilinFormOfCoeff A hwu hwu := by
            simpa [μ] using bilinForm_coercive A hwu
  have h_rhs_bound : |rhs u| ≤ C_F * ‖gradLpOfWitness hwu‖ := by
    calc
      |rhs u| ≤ C_F * (∫ x, ‖hwu.weakGrad x‖ ^ (2 : ℝ) ∂μ) ^ (1 / (2 : ℝ)) := hF u hu0 hwu
      _ = C_F * ‖gradLpOfWitness hwu‖ := by rw [norm_gradLpOfWitness_eq hwu]
  have hbilin_nonneg : 0 ≤ bilinFormOfCoeff A hwu hwu := by
    exact le_trans (mul_nonneg A.lam_nonneg (sq_nonneg ‖gradLpOfWitness hwu‖)) hcoercive
  have hrhs_nonneg : 0 ≤ rhs u := by
    simpa [h_ws] using hbilin_nonneg
  have hmain :
      A.lam * ‖gradLpOfWitness hwu‖ ^ (2 : ℕ) ≤ C_F * ‖gradLpOfWitness hwu‖ := by
    calc
      A.lam * ‖gradLpOfWitness hwu‖ ^ (2 : ℕ) ≤ bilinFormOfCoeff A hwu hwu := hcoercive
      _ = rhs u := h_ws
      _ = |rhs u| := by symm; exact abs_of_nonneg hrhs_nonneg
      _ ≤ C_F * ‖gradLpOfWitness hwu‖ := h_rhs_bound
  have hlin :
      A.lam * ‖gradLpOfWitness hwu‖ ≤ C_F := by
    have hmain' :
        A.lam * ‖gradLpOfWitness hwu‖ ^ 2 ≤ C_F * ‖gradLpOfWitness hwu‖ := by
      simpa [pow_two] using hmain
    nlinarith [hmain', A.hlam, norm_nonneg (gradLpOfWitness hwu)]
  have hdiv :
      ‖gradLpOfWitness hwu‖ ≤ C_F / A.lam := by
    have hlin' : ‖gradLpOfWitness hwu‖ * A.lam ≤ C_F := by
      simpa [mul_comm] using hlin
    exact (le_div_iff₀ A.hlam).2 hlin'
  simpa [norm_gradLpOfWitness_eq hwu, div_eq_mul_inv, mul_comm] using hdiv

\end{lstlisting}
\begin{proof}
Test against $u$ itself. Coercivity gives $\lam \|\nabla u\|_{\Lp{2}}^{2} \leq \mathfrak{a}_A(u, u) = \mathrm{rhs}(u)$, while boundedness gives $|\mathrm{rhs}(u)| \leq C_F \|\nabla u\|_{\Lp{2}}$. The bilinear form value is nonnegative (by coercivity), hence $\mathrm{rhs}(u) \geq 0$ and equals $|\mathrm{rhs}(u)|$. Combining,
\[
  \lam \|\nabla u\|_{\Lp{2}}^{2} \leq C_F \|\nabla u\|_{\Lp{2}},
\]
and dividing (or using $\nabla u = 0$) yields $\|\nabla u\|_{\Lp{2}} \leq C_F/\lam$.
\end{proof}

\begin{theorem}[\leanname{ae\_eq\_zero\_of\_memH01\_of\_gradLpOfWitness\_eq\_zero}]
For $d \geq 2$, an open bounded $\Omega$, and $u \in H^{1}_{0}(\Omega)$ with vanishing $\Lp{2}$ gradient class, $u = 0$ a.e.\ on $\Omega$.
\end{theorem}

\begin{lstlisting}[style=Matlab-editor,basicstyle=\mlttfamily,escapechar=",mlshowsectionrules=true,mathescape=true,morecomment={[s]{\%\{}{\%\}}},language=lean,numbers=none]
/-- A zero-trace Sobolev function with vanishing `L²` gradient class vanishes
almost everywhere. -/
theorem ae_eq_zero_of_memH01_of_gradLpOfWitness_eq_zero
    {Ω : Set E} (hd : 2 ≤ d) (hΩ : IsOpen Ω)
    (hΩ_bdd : Bornology.IsBounded Ω)
    {u : E → ℝ}
    (hu0 : MemH01 u Ω)
    (hwu : MemW1pWitness 2 u Ω)
    (hgrad_zero : gradLpOfWitness hwu = 0) :
    u =ᵐ[volume.restrict Ω] 0 := by
  let μ : Measure E := volume.restrict Ω
  obtain ⟨hu_mem, hw0, φ, hφ_smooth, hφ_cpt, hφ_sub, hφ_fun, hφ_grad⟩ := hu0
  have hgrad_ae :
      hw0.weakGrad =ᵐ[μ] hwu.weakGrad := by
    simpa [μ] using MemW1pWitness.ae_eq hΩ hw0 hwu
  have hgradLp_eq : gradLpOfWitness hw0 = gradLpOfWitness hwu := by
    simpa [gradLpOfWitness] using
      (MemLp.toLp_congr hw0.weakGrad_memLp hwu.weakGrad_memLp hgrad_ae)
  have hgrad_zero0 : gradLpOfWitness hw0 = 0 := by
    rw [hgradLp_eq]
    exact hgrad_zero
  let hφtest : ∀ n : ℕ, IsSmoothTestOn Ω (φ n) := fun n =>
    ⟨hφ_smooth n, hφ_cpt n, hφ_sub n⟩
  let hφw : ∀ n : ℕ, MemW1pWitness 2 (φ n) Ω := fun n =>
    smoothTestWitness hΩ (hφtest n)
  have hφ_fun_tendsto :
      Tendsto (fun n => smoothFunToLp hΩ (hφtest n))
        atTop (nhds (hu_mem.1.toLp u)) := by
    exact
      (Lp.tendsto_Lp_iff_tendsto_eLpNorm'' (f := φ)
        (f_ℒp := fun n => (hφw n).memLp)
        (f_lim := u) (f_lim_ℒp := hu_mem.1)).2 <| by
          simpa [μ] using hφ_fun
  have hgrad_vec_tendsto :
      Tendsto
        (fun n => eLpNorm (fun x => (hφw n).weakGrad x - hw0.weakGrad x) 2 μ)
        atTop (nhds 0) := by
    exact tendsto_eLpNorm_vector_of_componentwise
      (fun n i => by
        simpa [hφw] using ((hφw n).weakGrad_component_memLp i).sub (hw0.weakGrad_component_memLp i))
      (fun i => by simpa [μ] using hφ_grad i)
  have hgradLp_tendsto :
      Tendsto (fun n => gradLpOfWitness (hφw n))
        atTop (nhds (gradLpOfWitness hw0)) := by
    exact
      (Lp.tendsto_Lp_iff_tendsto_eLpNorm'' (f := fun n => (hφw n).weakGrad)
        (f_ℒp := fun n => (hφw n).weakGrad_memLp)
        (f_lim := hw0.weakGrad) (f_lim_ℒp := hw0.weakGrad_memLp)).2
        hgrad_vec_tendsto
  have hgradLp_zero_tendsto :
      Tendsto (fun n => gradLpOfWitness (hφw n)) atTop (nhds 0) := by
    simpa [hgrad_zero0] using hgradLp_tendsto
  obtain ⟨Cp, hCp_top, hPoinc⟩ := smoothCompactSupport_L2_bound_on_bounded_ge_two hd hΩ_bdd
  have hfun_bound :
      ∀ n, ‖smoothFunToLp hΩ (hφtest n)‖ ≤ Cp.toReal * ‖smoothGradToLp hΩ (hφtest n)‖ := by
    intro n
    have hgradNorm_memLp : MemLp (smoothGradNorm (φ n)) 2 μ := by
      simpa [μ, norm_smoothGradField_eq_smoothGradNorm] using (hφw n).weakGrad_memLp.norm
    have hP :
        eLpNorm (φ n) 2 μ ≤ Cp * eLpNorm (smoothGradNorm (φ n)) 2 μ := by
      simpa [μ] using hPoinc (hφ_smooth n) (hφ_cpt n) (hφ_sub n)
    have hleft_top : eLpNorm (φ n) 2 μ ≠ ⊤ := (hφw n).memLp.eLpNorm_lt_top.ne
    have hright_top : Cp * eLpNorm (smoothGradNorm (φ n)) 2 μ ≠ ⊤ := by
      exact (ENNReal.mul_lt_top hCp_top hgradNorm_memLp.eLpNorm_lt_top).ne
    have hToReal :
        (eLpNorm (φ n) 2 μ).toReal ≤ (Cp * eLpNorm (smoothGradNorm (φ n)) 2 μ).toReal :=
      (ENNReal.toReal_le_toReal hleft_top hright_top).2 hP
    have hfunEq : ‖smoothFunToLp hΩ (hφtest n)‖ = (eLpNorm (φ n) 2 μ).toReal := by
      rw [smoothFunToLp, Lp.norm_toLp]
    have hgradEq : ‖smoothGradToLp hΩ (hφtest n)‖ = (eLpNorm (smoothGradNorm (φ n)) 2 μ).toReal := by
      have hgrad_memLp : MemLp (smoothGradField (φ n)) 2 μ := (hφw n).weakGrad_memLp
      calc
        ‖smoothGradToLp hΩ (hφtest n)‖
            = (∫ x, ‖smoothGradField (φ n) x‖ ^ (2 : ℝ) ∂μ) ^ (1 / (2 : ℝ)) := by
                simpa [μ] using norm_smoothGradToLp_eq hΩ (hφtest n)
        _ = (eLpNorm (smoothGradField (φ n)) 2 μ).toReal := by
              symm
              rw [MeasureTheory.toReal_eLpNorm hgrad_memLp.aestronglyMeasurable]
              simpa using
                (MeasureTheory.lpNorm_eq_integral_norm_rpow_toReal
                  (μ := μ) (f := smoothGradField (φ n)) (p := (2 : ENNReal))
                  (by norm_num) (by simp) hgrad_memLp.aestronglyMeasurable)
        _ = (eLpNorm (smoothGradNorm (φ n)) 2 μ).toReal := by
              congr 1
              calc
                eLpNorm (smoothGradField (φ n)) 2 μ =
                    eLpNorm (fun x => ‖smoothGradField (φ n) x‖) 2 μ := by
                      symm
                      exact eLpNorm_norm (μ := μ) (p := (2 : ENNReal)) (smoothGradField (φ n))
                _ = eLpNorm (smoothGradNorm (φ n)) 2 μ := by
                      simp [norm_smoothGradField_eq_smoothGradNorm]
    rw [hfunEq, hgradEq]
    convert hToReal using 1
    rw [ENNReal.toReal_mul]
  have hgrad_norm_tendsto :
      Tendsto (fun n => ‖smoothGradToLp hΩ (hφtest n)‖) atTop (nhds 0) := by
    have hnorm_tendsto :
        Tendsto (fun n => ‖gradLpOfWitness (hφw n)‖) atTop (nhds 0) := by
      simpa using ((continuous_norm.tendsto (0 : MeasureTheory.Lp E 2 μ)).comp hgradLp_zero_tendsto)
    simpa [smoothGradToLp, hφw] using hnorm_tendsto
  have hfun_norm_tendsto :
      Tendsto (fun n => ‖smoothFunToLp hΩ (hφtest n)‖) atTop (nhds 0) := by
    have hupper :
        Tendsto (fun n => Cp.toReal * ‖smoothGradToLp hΩ (hφtest n)‖) atTop (nhds 0) := by
      simpa using Tendsto.const_mul Cp.toReal hgrad_norm_tendsto
    exact squeeze_zero (fun n => norm_nonneg _) hfun_bound hupper
  have hfun_zero_tendsto :
      Tendsto (fun n => smoothFunToLp hΩ (hφtest n)) atTop (nhds 0) := by
    exact tendsto_zero_iff_norm_tendsto_zero.2 hfun_norm_tendsto
  have hu_toLp_zero : hu_mem.1.toLp u = 0 := by
    exact tendsto_nhds_unique hφ_fun_tendsto hfun_zero_tendsto
  have htoLp_zero_ae : (hu_mem.1.toLp u : E → ℝ) =ᵐ[μ] (0 : E → ℝ) := by
    let h0_mem : MemLp (0 : E → ℝ) 2 μ := MeasureTheory.MemLp.zero'
    have h0_ae : (h0_mem.toLp (0 : E → ℝ) : E → ℝ) =ᵐ[μ] (0 : E → ℝ) := by
      exact MemLp.coeFn_toLp h0_mem
    rw [hu_toLp_zero]
    simpa [h0_mem.toLp_zero] using h0_ae
  calc
    u =ᵐ[μ] hu_mem.1.toLp u := (MemLp.coeFn_toLp hu_mem.1).symm
    _ =ᵐ[μ] (0 : E → ℝ) := htoLp_zero_ae

\end{lstlisting}
\begin{proof}
Take an approximating sequence of smooth tests $\varphi_n$ from the $H^{1}_{0}$ structure with $\varphi_n \to u$ in $\Lp{2}$ and $\nabla \varphi_n$ converging componentwise to a representative of $0$. By Poincar\'e on the bounded domain, $\|\varphi_n\|_{\Lp{2}} \leq C_P \|\nabla \varphi_n\|_{\Lp{2}} \to 0$, so $[\varphi_n] \to 0$ in $\Lp{2}$. Combining with $[\varphi_n] \to [u]$, we get $[u] = 0$, i.e.\ $u = 0$ a.e.
\end{proof}

\begin{theorem}[\leanname{weakProblem\_unique}]
For $d \geq 2$, weak solutions of a fixed weak problem $P$ are unique up to a.e.\ equality.
\end{theorem}

\begin{lstlisting}[style=Matlab-editor,basicstyle=\mlttfamily,escapechar=",mlshowsectionrules=true,mathescape=true,morecomment={[s]{\%\{}{\%\}}},language=lean,numbers=none]
/-- Zero-boundary weak solutions are unique up to a.e. equality. -/
theorem weakProblem_unique
    (hd : 2 ≤ d)
    {P : WeakProblem (d := d)}
    {u v : E → ℝ}
    (hu : IsWeakSolution P u)
    (hv : IsWeakSolution P v) :
    u =ᵐ[volume.restrict P.Ω] v := by
  have hu0 : MemH01 u P.Ω := hu.left
  have hv0 : MemH01 v P.Ω := hv.left
  have hdiff0 : MemH01 (fun x => u x - v x) P.Ω := MemW01p.sub hu0 hv0
  let hwu : MemW1pWitness 2 u P.Ω :=
    DeGiorgi.MemW1p.someWitness (MemW01p.memW1p hu0)
  let hwv : MemW1pWitness 2 v P.Ω :=
    DeGiorgi.MemW1p.someWitness (MemW01p.memW1p hv0)
  let hdiff_add : MemW1pWitness 2 (fun x => u x + (-1) * v x) P.Ω :=
    hwu.add (hwv.smul (-1))
  let hdiff : MemW1pWitness 2 (fun x => u x - v x) P.Ω := {
    memLp := by
      simpa [sub_eq_add_neg, Pi.smul_apply] using hdiff_add.memLp
    weakGrad := hdiff_add.weakGrad
    weakGrad_component_memLp := by
      intro i
      simpa [sub_eq_add_neg, Pi.smul_apply] using hdiff_add.weakGrad_component_memLp i
    isWeakGrad := by
      intro i
      simpa [sub_eq_add_neg, Pi.smul_apply] using hdiff_add.isWeakGrad i }
  have hzero_test :
      ∀ (φ : E → ℝ), MemH01 φ P.Ω →
        ∀ (hφ : MemW1pWitness 2 φ P.Ω),
          bilinFormOfCoeff P.coeff hdiff hφ = 0 := by
    intro φ hφ0 hφ
    calc
      bilinFormOfCoeff P.coeff hdiff hφ
          = bilinFormOfCoeff P.coeff hdiff_add hφ := by
              rfl
      _ = bilinFormOfCoeff P.coeff hwu hφ +
            bilinFormOfCoeff P.coeff (hwv.smul (-1)) hφ := by
              simpa [hdiff_add] using bilinFormOfCoeff_add_left P.coeff hwu (hwv.smul (-1)) hφ
      _ = bilinFormOfCoeff P.coeff hwu hφ + (-1) * bilinFormOfCoeff P.coeff hwv hφ := by
            rw [bilinFormOfCoeff_smul_left (-1) P.coeff hwv hφ]
      _ = P.rhs φ + (-1) * P.rhs φ := by rw [hu.right hwu φ hφ0 hφ, hv.right hwv φ hφ0 hφ]
      _ = 0 := by ring
  have h_ws : bilinFormOfCoeff P.coeff hdiff hdiff = (0 : ℝ) := hzero_test _ hdiff0 hdiff
  have hstab :
      (∫ x, ‖hdiff.weakGrad x‖ ^ (2 : ℝ) ∂(volume.restrict P.Ω)) ^ (1 / (2 : ℝ)) ≤
        P.coeff.lam⁻¹ * (0 : ℝ) := by
    simpa using weakSolution_stability P.coeff hdiff0 hdiff
      (rhs := fun _ : E → ℝ => (0 : ℝ)) (C_F := 0) (by norm_num)
      (by intro φ hφ0 hφ; simp) h_ws
  have hseminorm_zero :
      (∫ x, ‖hdiff.weakGrad x‖ ^ (2 : ℝ) ∂(volume.restrict P.Ω)) ^ (1 / (2 : ℝ)) = 0 := by
    refine le_antisymm ?_ ?_
    · simpa using hstab
    · positivity
  have hgrad_zero : gradLpOfWitness hdiff = 0 := by
    apply norm_eq_zero.mp
    simpa [norm_gradLpOfWitness_eq hdiff] using hseminorm_zero
  have hdiff_ae_zero :
      (fun x => u x - v x) =ᵐ[volume.restrict P.Ω] 0 := by
    exact ae_eq_zero_of_memH01_of_gradLpOfWitness_eq_zero hd P.hΩ P.hΩ_bdd hdiff0 hdiff hgrad_zero
  filter_upwards [hdiff_ae_zero] with x hx
  simpa [sub_eq_zero] using hx

\end{lstlisting}
\begin{proof}
If $u, v$ are two weak solutions, $w := u - v$ lies in $H^{1}_{0}$ and tests against any $\varphi \in H^{1}_{0}$ to zero by linearity. By the stability estimate (with $C_F = 0$), $\|\nabla w\|_{\Lp{2}} = 0$, and the previous theorem gives $w = 0$ a.e.
\end{proof}

\begin{definition}[\leanname{IsDirichletWeakSolutionOfDivergenceData}]
$u$ is a Dirichlet weak solution for boundary datum $u_0$ and field $F$ if $u \in W^{1,2}(\Omega)$, $u - u_0 \in H^{1}_{0}(\Omega)$, and $\mathfrak{a}_A(u, v) = L_F(v)$ for every $v \in H^{1}_{0}(\Omega)$ on every witness.
\end{definition}

\begin{lstlisting}[style=Matlab-editor,basicstyle=\mlttfamily,escapechar=",mlshowsectionrules=true,mathescape=true,morecomment={[s]{\%\{}{\%\}}},language=lean,numbers=none]
/-- Weak solution of the inhomogeneous Dirichlet problem with boundary datum
`u₀` and divergence-form right-hand side `div F`. -/
def IsDirichletWeakSolutionOfDivergenceData
    {Ω : Set E}
    (A : EllipticCoeff d Ω)
    (u₀ : E → ℝ)
    (F : E → E)
    (u : E → ℝ) : Prop :=
  MemW1p 2 u Ω ∧
  MemW01p 2 (fun x => u x - u₀ x) Ω ∧
  ∀ (hu : MemW1pWitness 2 u Ω) (v : E → ℝ), MemH01 v Ω →
    ∀ (hv : MemW1pWitness 2 v Ω),
      bilinFormOfCoeff A hu hv = divergenceRHSOfField F hv

\end{lstlisting}
\begin{theorem}
(\leanname{isWeakSolution\_subDatum\_of\_dirichletProblem},
\leanname{dirichletProblem\_stability\_of\_divergenceData},
\leanname{dirichletProblem\_unique\_of\_divergenceData},
\leanname{dirichletProblem\_wellPosed\_of\_divergenceData})\\
The shifted unknown $u - u_0$ is a zero-boundary weak solution for the shifted RHS, satisfies the energy estimate
\[
  \|\nabla(u - u_0)\|_{\Lp{2}} \leq \frac{1}{\lam}\Bigl(\|F\|_{\Lp{2}} + \Lambda\,\|\nabla u_0\|_{\Lp{2}}\Bigr),
\]
and Dirichlet weak solutions for fixed $(u_0, F)$ are unique up to a.e.\ equality. Combined with existence, we obtain well-posedness.
\end{theorem}

\begin{lstlisting}[style=Matlab-editor,basicstyle=\mlttfamily,escapechar=",mlshowsectionrules=true,mathescape=true,morecomment={[s]{\%\{}{\%\}}},language=lean,numbers=none]
/-- A Dirichlet solution lifts to a zero-boundary weak solution for the shifted
unknown `u - u₀`. -/
theorem isWeakSolution_subDatum_of_dirichletProblem
    {Ω : Set E}
    (hΩ : IsOpen Ω) (hΩ_bdd : Bornology.IsBounded Ω)
    (A : EllipticCoeff d Ω)
    {u₀ : E → ℝ} (hu₀ : MemW1p 2 u₀ Ω)
    {F : E → E} {u : E → ℝ}
    (hu : IsDirichletWeakSolutionOfDivergenceData A u₀ F u) :
    IsWeakSolution (d := d)
      ⟨Ω, hΩ, hΩ_bdd, A,
        weakProblemRHSOfFieldAndDatum (A := A) (Ω := Ω) F
          (DeGiorgi.MemW1p.someWitness hu₀)⟩
      (fun x => u x - u₀ x) := by
  let hu₀w : MemW1pWitness 2 u₀ Ω := DeGiorgi.MemW1p.someWitness hu₀
  let huu : MemW1pWitness 2 u Ω := DeGiorgi.MemW1p.someWitness hu.left
  let hsub_add : MemW1pWitness 2 (fun x => u x + (-1) * u₀ x) Ω :=
    huu.add (hu₀w.smul (-1))
  let hsub : MemW1pWitness 2 (fun x => u x - u₀ x) Ω := {
    memLp := by
      simpa [sub_eq_add_neg, Pi.smul_apply] using hsub_add.memLp
    weakGrad := hsub_add.weakGrad
    weakGrad_component_memLp := by
      intro i
      simpa [sub_eq_add_neg, Pi.smul_apply] using hsub_add.weakGrad_component_memLp i
    isWeakGrad := by
      intro i
      simpa [sub_eq_add_neg, Pi.smul_apply] using hsub_add.isWeakGrad i }
  refine ⟨hu.2.1, ?_⟩
  intro hw' v hv0 hv
  calc
    bilinFormOfCoeff A hw' hv = bilinFormOfCoeff A hsub hv := by
      exact bilinFormOfCoeff_eq_left hΩ A hw' hsub hv
    _ = bilinFormOfCoeff A hsub_add hv := by
          rfl
    _ = bilinFormOfCoeff A huu hv + bilinFormOfCoeff A (hu₀w.smul (-1)) hv := by
          simpa [hsub_add] using bilinFormOfCoeff_add_left A huu (hu₀w.smul (-1)) hv
    _ = bilinFormOfCoeff A huu hv + (-1) * bilinFormOfCoeff A hu₀w hv := by
          rw [bilinFormOfCoeff_smul_left (-1) A hu₀w hv]
    _ = divergenceRHSOfField F hv + (-1) * bilinFormOfCoeff A hu₀w hv := by
          rw [hu.2.2 huu v hv0 hv]
    _ = divergenceRHSOfField F hv - bilinFormOfCoeff A hu₀w hv := by
          ring
    _ = weakProblemRHSOfFieldAndDatum (A := A) (Ω := Ω) F hu₀w v := by
          symm
          exact weakProblemRHSOfFieldAndDatum_eq_of_memH01 hΩ A hu₀w hv0 hv

/-- Energy estimate for the shifted Dirichlet solution `u - u₀`. -/
theorem dirichletProblem_stability_of_divergenceData
    {Ω : Set E}
    (hΩ : IsOpen Ω) (hΩ_bdd : Bornology.IsBounded Ω)
    (A : EllipticCoeff d Ω)
    {u₀ : E → ℝ} (hu₀ : MemW1p 2 u₀ Ω)
    {F : E → E} (hF : MemLp F 2 (volume.restrict Ω))
    {u : E → ℝ}
    (hu : IsDirichletWeakSolutionOfDivergenceData A u₀ F u)
    (hw : MemW1pWitness 2 (fun x => u x - u₀ x) Ω) :
    (∫ x, ‖hw.weakGrad x‖ ^ (2 : ℝ) ∂(volume.restrict Ω)) ^ (1 / (2 : ℝ)) ≤
      A.lam⁻¹ *
        ((∫ x, ‖F x‖ ^ (2 : ℝ) ∂(volume.restrict Ω)) ^ (1 / (2 : ℝ)) +
          A.Λ * (∫ x, ‖(DeGiorgi.MemW1p.someWitness hu₀).weakGrad x‖ ^ (2 : ℝ)
            ∂(volume.restrict Ω)) ^ (1 / (2 : ℝ))) := by
  let hu₀w : MemW1pWitness 2 u₀ Ω := DeGiorgi.MemW1p.someWitness hu₀
  have hwsol :
      IsWeakSolution (d := d)
        ⟨Ω, hΩ, hΩ_bdd, A, weakProblemRHSOfFieldAndDatum (A := A) (Ω := Ω) F hu₀w⟩
        (fun x => u x - u₀ x) := by
    exact isWeakSolution_subDatum_of_dirichletProblem (Ω := Ω) hΩ hΩ_bdd A hu₀ hu
  have hbound :
      ∀ v : E → ℝ, MemH01 v Ω →
        ∀ hv : MemW1pWitness 2 v Ω,
          |weakProblemRHSOfFieldAndDatum (A := A) (Ω := Ω) F hu₀w v| ≤
            ((∫ x, ‖F x‖ ^ (2 : ℝ) ∂(volume.restrict Ω)) ^ (1 / (2 : ℝ)) +
              A.Λ * (∫ x, ‖hu₀w.weakGrad x‖ ^ (2 : ℝ) ∂(volume.restrict Ω)) ^
                (1 / (2 : ℝ))) *
            (∫ x, ‖hv.weakGrad x‖ ^ (2 : ℝ) ∂(volume.restrict Ω)) ^ (1 / (2 : ℝ)) := by
    intro v hv0 hv
    simpa [hu₀w] using weakProblemRHSOfFieldAndDatum_bound hΩ A hF hu₀w v hv0 hv
  exact weakSolution_stability A hu.2.1 hw (C_F :=
    (∫ x, ‖F x‖ ^ (2 : ℝ) ∂(volume.restrict Ω)) ^ (1 / (2 : ℝ)) +
      A.Λ * (∫ x, ‖hu₀w.weakGrad x‖ ^ (2 : ℝ) ∂(volume.restrict Ω)) ^ (1 / (2 : ℝ)))
    (by
      have hCF_nonneg :
          0 ≤ (∫ x, ‖F x‖ ^ (2 : ℝ) ∂(volume.restrict Ω)) ^ (1 / (2 : ℝ)) := by
        positivity
      have hCU_nonneg :
          0 ≤ A.Λ * (∫ x, ‖hu₀w.weakGrad x‖ ^ (2 : ℝ) ∂(volume.restrict Ω)) ^ (1 / (2 : ℝ)) := by
        exact mul_nonneg A.Λ_nonneg (by positivity)
      exact add_nonneg hCF_nonneg hCU_nonneg)
    hbound
    (hwsol.right hw (fun x => u x - u₀ x) hu.2.1 hw)

/-- Inhomogeneous Dirichlet solutions are unique up to a.e. equality. -/
theorem dirichletProblem_unique_of_divergenceData
    {Ω : Set E}
    (hd : 2 ≤ d)
    (hΩ : IsOpen Ω) (hΩ_bdd : Bornology.IsBounded Ω)
    (A : EllipticCoeff d Ω)
    {u₀ : E → ℝ} (hu₀ : MemW1p 2 u₀ Ω)
    {F : E → E} (_hF : MemLp F 2 (volume.restrict Ω))
    {u v : E → ℝ}
    (hu : IsDirichletWeakSolutionOfDivergenceData A u₀ F u)
    (hv : IsDirichletWeakSolutionOfDivergenceData A u₀ F v) :
    u =ᵐ[volume.restrict Ω] v := by
  let P : WeakProblem (d := d) :=
    ⟨Ω, hΩ, hΩ_bdd, A,
      weakProblemRHSOfFieldAndDatum (A := A) (Ω := Ω) F
        (DeGiorgi.MemW1p.someWitness hu₀)⟩
  have hu_shift :
      IsWeakSolution P (fun x => u x - u₀ x) := by
    exact isWeakSolution_subDatum_of_dirichletProblem (Ω := Ω) hΩ hΩ_bdd A hu₀ hu
  have hv_shift :
      IsWeakSolution P (fun x => v x - u₀ x) := by
    exact isWeakSolution_subDatum_of_dirichletProblem (Ω := Ω) hΩ hΩ_bdd A hu₀ hv
  have hshift_eq :
      (fun x => u x - u₀ x) =ᵐ[volume.restrict Ω] (fun x => v x - u₀ x) := by
    simpa [P] using weakProblem_unique hd hu_shift hv_shift
  filter_upwards [hshift_eq] with x hx
  linarith

/-- Packaged existence and uniqueness for the inhomogeneous Dirichlet problem. -/
theorem dirichletProblem_wellPosed_of_divergenceData
    {Ω : Set E}
    (hd : 2 ≤ d)
    (hΩ : IsOpen Ω) (hΩ_bdd : Bornology.IsBounded Ω)
    (A : EllipticCoeff d Ω)
    {u₀ : E → ℝ} (hu₀ : MemW1p 2 u₀ Ω)
    {F : E → E} (hF : MemLp F 2 (volume.restrict Ω)) :
    ∃ u : E → ℝ,
      IsDirichletWeakSolutionOfDivergenceData A u₀ F u ∧
      ∀ v : E → ℝ, IsDirichletWeakSolutionOfDivergenceData A u₀ F v →
        u =ᵐ[volume.restrict Ω] v := by
  obtain ⟨u, hu⟩ := dirichletProblem_exists_of_divergenceData' hd hΩ hΩ_bdd A hu₀ hF
  refine ⟨u, hu, ?_⟩
  intro v hv
  exact dirichletProblem_unique_of_divergenceData hd hΩ hΩ_bdd A hu₀ hF hu hv

\end{lstlisting}
\begin{proof}
The shifted-solution identity follows by linearity of $\mathfrak{a}_A$: $\mathfrak{a}_A(u - u_0, v) = \mathfrak{a}_A(u, v) - \mathfrak{a}_A(u_0, v) = L_F(v) - \mathfrak{a}_A(u_0, v)$. Apply \leanname{weakSolution\_stability} to the shifted unknown with $C_F = \|F\|_{\Lp{2}} + \Lambda\|\nabla u_0\|_{\Lp{2}}$ to obtain the energy estimate. Uniqueness reduces to uniqueness of zero-Dirichlet solutions for the shifted RHS, hence to \leanname{weakProblem\_unique}.
\end{proof}

\begin{theorem}[\leanname{aHarmonic\_replacement\_exists}, \leanname{aHarmonic\_replacement\_unique}]
For $d \geq 2$, an open bounded $\Omega$, an elliptic field $A$ on $\Omega$, and $u \in W^{1,2}(\Omega)$, there exists a homogeneous weak solution $h$ of $-\dv(a\nabla h) = 0$ with $h - u \in H^{1}_{0}(\Omega)$, and any two such replacements agree a.e.\ on $\Omega$.
\end{theorem}

\begin{lstlisting}[style=Matlab-editor,basicstyle=\mlttfamily,escapechar=",mlshowsectionrules=true,mathescape=true,morecomment={[s]{\%\{}{\%\}}},language=lean,numbers=none]
/-- A-harmonic replacement: given `u ∈ W^{1,2}(B_R)`, there exists
`h` that is A-harmonic with `h - u ∈ W₀^{1,2}(B_R)`. -/
theorem aHarmonic_replacement_exists
    {Ω : Set E} (hd : 2 ≤ d)
    (hΩ : IsOpen Ω) (hΩ_bdd : Bornology.IsBounded Ω)
    {u : E → ℝ}
    (A : EllipticCoeff d Ω)
    (hu : MemW1p 2 u Ω) :
    ∃ h : E → ℝ,
      IsHomogeneousWeakSolution A h ∧
      MemW01p 2 (fun x => h x - u x) Ω := by
  have hF0 : MemLp (fun _ : E => (0 : E)) 2 (volume.restrict Ω) := by
    exact MeasureTheory.MemLp.zero' (p := (2 : ENNReal)) (μ := volume.restrict Ω)
  obtain ⟨h, hh_mem, hdiff, hweak⟩ :=
    dirichletProblem_exists_of_divergenceData hd hΩ hΩ_bdd A hu hF0
  refine ⟨h, ?_, hdiff⟩
  refine ⟨hh_mem, ?_⟩
  intro hh φ hφ hφw
  have hEq := hweak hh φ hφ hφw
  simpa [divergenceRHSOfField, divergenceRHSIntegrandOfField] using hEq


/-- `A`-harmonic replacements are unique up to a.e. equality. -/
theorem aHarmonic_replacement_unique
    {Ω : Set E} (hd : 2 ≤ d)
    (hΩ : IsOpen Ω) (hΩ_bdd : Bornology.IsBounded Ω)
    {u h₁ h₂ : E → ℝ}
    (A : EllipticCoeff d Ω)
    (hu : MemW1p 2 u Ω)
    (hh₁ : IsHomogeneousWeakSolution A h₁)
    (hdiff₁ : MemW01p 2 (fun x => h₁ x - u x) Ω)
    (hh₂ : IsHomogeneousWeakSolution A h₂)
    (hdiff₂ : MemW01p 2 (fun x => h₂ x - u x) Ω) :
    h₁ =ᵐ[volume.restrict Ω] h₂ := by
  have hF0 : MemLp (fun _ : E => (0 : E)) 2 (volume.restrict Ω) := by
    exact MeasureTheory.MemLp.zero' (p := (2 : ENNReal)) (μ := volume.restrict Ω)
  have hs₁ : IsDirichletWeakSolutionOfDivergenceData A u (fun _ : E => (0 : E)) h₁ := by
    refine ⟨hh₁.left, hdiff₁, ?_⟩
    intro hh φ hφ hφw
    simpa [divergenceRHSOfField, divergenceRHSIntegrandOfField] using hh₁.right hh φ hφ hφw
  have hs₂ : IsDirichletWeakSolutionOfDivergenceData A u (fun _ : E => (0 : E)) h₂ := by
    refine ⟨hh₂.left, hdiff₂, ?_⟩
    intro hh φ hφ hφw
    simpa [divergenceRHSOfField, divergenceRHSIntegrandOfField] using hh₂.right hh φ hφ hφw
  exact dirichletProblem_unique_of_divergenceData hd hΩ hΩ_bdd A hu hF0 hs₁ hs₂

\end{lstlisting}
\begin{proof}
Apply \leanname{dirichletProblem\_exists\_of\_divergenceData} with the
prescribed boundary datum $u_0 := u$ and the trivial divergence field
$F := 0 \in L^{2}(\Omega; E)$. The hypotheses are met because $u \in
W^{1,2}(\Omega)$ by assumption and the constant zero function is in $L^2$
(\leanname{MeasureTheory.MemLp.zero'}). The theorem produces a
$h \in W^{1,2}(\Omega)$ with $h - u \in H^{1}_{0}(\Omega)$ satisfying
\[
  \mathfrak{a}_A(h, v) = L_F(v) = \int_{\Omega}\langle 0, \nabla v\rangle\,dx = 0
  \quad \text{for every } v \in H^{1}_{0}(\Omega),
\]
i.e.\ $h$ is a homogeneous weak solution against any choice of witness for $v$.

To see that $h$ is moreover a \emph{homogeneous} weak solution in the sense of
\leanname{IsHomogeneousWeakSolution} (the local interface used by Harnack and
H\"older), unfold the definition: it requires that for every test
$\varphi \in H^{1}_{0}(\Omega)$ and every witness, the bilinear form vanishes.
This is exactly the displayed identity (with the right-hand side
$\leanname{divergenceRHSOfField}\,0$ unfolding to zero by
\leanname{divergenceRHSIntegrandOfField}).

Uniqueness follows directly from
\leanname{dirichletProblem\_unique\_of\_divergenceData} applied to the same
data $(u_0, F) = (u, 0)$: any two replacements have identical Dirichlet trace
along $u$ and identical zero divergence data, so they coincide a.e.\ on
$\Omega$.
\end{proof}

\subsection{Weighted bilinear-form estimates}

\begin{lemma}[\leanname{bilinFormIntegrand\_mul\_smooth\_eq}]
Let $u, w \in W^{1,2}(\Omega)$ with witnesses $h_u, h_w$, and $\eta \in C^{\infty}(E)$ satisfying $|\eta| \leq C_0$ and $\|\fderiv{\eta}{x}\| \leq C_1$. Then for every $x$,
\begin{multline*}
  b_A(u, \eta\!\cdot\!w)(x) = \eta(x)\,b_A(u,w)(x) + w(x)\,\bigl\langle a(x)\,\nabla u(x), \nabla \eta(x)\bigr\rangle,
\end{multline*}
where $\nabla(\eta\cdot w)$ is realized by the witness produced from \leanname{MemW1pWitness.mul\_smooth\_bounded\_p} as $\eta\,\nabla w + w\,\nabla \eta$.
\end{lemma}

\begin{lstlisting}[style=Matlab-editor,basicstyle=\mlttfamily,escapechar=",mlshowsectionrules=true,mathescape=true,morecomment={[s]{\%\{}{\%\}}},language=lean,numbers=none]
/-! ### Bilinear form expansion for smooth-product witnesses

These lemmas expand `bilinFormOfCoeff A hu (η · hw)` into its two constituent
terms. This is the key tool for the Caccioppoli / Moser absorption argument:
the bilinear form against a product test function splits into a "principal" term
(involving `η · ⟪A∇u, ∇w⟫`) and a "cross" term (involving `w · ⟪A∇u, ∇η⟫`). -/

/-- Pointwise expansion of the bilinear-form integrand against a smooth-product
witness `η · w`.

The weak gradient of `η · w` is `η · ∇w + (∇η) · w` (from `mul_smooth_bounded_p`),
so the bilinear-form integrand is:
```
⟪A∇u, ∇(η·w)⟫ = η · ⟪A∇u, ∇w⟫ + w · ⟪A∇u, ∇η⟫
```
where `∇η` is encoded as the vector with components `(fderiv ℝ η x)(single i 1)`.

Note: this is an a.e. identity because `mul_smooth_bounded_p` produces
the gradient `η · ∇w + (∇η) · w` by construction. -/
theorem bilinFormIntegrand_mul_smooth_eq
    {Ω : Set E} (hΩ : IsOpen Ω)
    {u w η : E → ℝ}
    (A : EllipticCoeff d Ω)
    (hu : MemW1pWitness 2 u Ω)
    (hw : MemW1pWitness 2 w Ω)
    (hη : ContDiff ℝ (⊤ : ℕ∞) η)
    {C₀ C₁ : ℝ} (hC₀ : 0 ≤ C₀) (hC₁ : 0 ≤ C₁)
    (hη_bound : ∀ x, |η x| ≤ C₀)
    (hη_grad_bound : ∀ x, ‖fderiv ℝ η x‖ ≤ C₁)
    (hp : (1 : ENNReal) ≤ 2) :
    ∀ x, bilinFormIntegrandOfCoeff A hu
        (hw.mul_smooth_bounded_p (d := d) hp hΩ hη hC₀ hC₁ hη_bound hη_grad_bound) x =
      η x * bilinFormIntegrandOfCoeff A hu hw x +
      w x * ⟪matMulE (A.a x) (hu.weakGrad x),
        (WithLp.toLp 2 fun i =>
          (fderiv ℝ η x) (EuclideanSpace.single i 1) : E)⟫_ℝ := by
  intro x
  -- The weakGrad of mul_smooth_bounded_p at x is:
  --   η x • hw.weakGrad x + toLp 2 (fun i => fderiv ℝ η x (single i 1) * w x)
  -- The bilinear-form integrand ⟪A∇u, ∇(η·w)⟫ distributes over this sum
  -- by bilinearity of the inner product.
  simp only [bilinFormIntegrandOfCoeff, MemW1pWitness.mul_smooth_bounded_p,
    inner_add_right, inner_smul_right]
  congr 1
  -- Factor w(x) out: ⟪v, toLp(f_i * c)⟫ = c * ⟪v, toLp(f_i)⟫
  -- The goal is: Σᵢ ⟪row_i, f_i * w⟫ = (Σᵢ ⟪row_i, f_i⟫) * w
  -- where ⟪a, b⟫_ℝ = a * b, so ⟪a, f_i * w⟫ = a * (f_i * w) = (a * f_i) * w = ⟪a, f_i⟫ * w
  simp only [PiLp.inner_apply, matMulE, Matrix.mulVec]
  rw [Finset.mul_sum]
  congr 1
  ext i
  -- ⟪a, b * c⟫_ℝ = c * ⟪a, b⟫_ℝ for reals
  rw [show (fderiv ℝ η x) (EuclideanSpace.single i 1) * w x =
      w x • ((fderiv ℝ η x) (EuclideanSpace.single i 1)) from by rw [smul_eq_mul]; ring,
    inner_smul_right]

-- Standalone helpers for `bilinForm_weighted_cauchySchwarz`.

/-- Multiply a `W^{1,p}` witness by a bounded smooth scalar factor. -/
noncomputable def MemW1pWitness.mul_smooth_bounded_p
    {p : ℝ≥0∞} (hp : 1 ≤ p)
    {Ω : Set E} (hΩ : IsOpen Ω)
    {u η : E → ℝ} (hw : MemW1pWitness p u Ω)
    (hη : ContDiff ℝ (⊤ : ℕ∞) η)
    {C₀ C₁ : ℝ}
    (hC₀ : 0 ≤ C₀) (hC₁ : 0 ≤ C₁)
    (hη_bound : ∀ x, |η x| ≤ C₀)
    (hη_grad_bound : ∀ x, ‖fderiv ℝ η x‖ ≤ C₁) :
    MemW1pWitness p (fun x => η x * u x) Ω where
  memLp := by
    let _ := hC₀
    let _ := hC₁
    refine MemLp.of_le_mul (g := u) (c := C₀) hw.memLp ?_ ?_
    · exact hη.continuous.aestronglyMeasurable.mul hw.memLp.aestronglyMeasurable
    · exact Filter.Eventually.of_forall fun x => by
        calc
          ‖η x * u x‖ = |η x| * ‖u x‖ := by
            rw [norm_mul, Real.norm_eq_abs]
          _ ≤ C₀ * ‖u x‖ := by
            gcongr
            exact hη_bound x
  weakGrad := fun x =>
    η x • hw.weakGrad x +
      (show E from WithLp.toLp 2 fun i => (fderiv ℝ η x) (EuclideanSpace.single i 1) * u x)
  weakGrad_component_memLp := by
    intro i
    have hfirst : MemLp (fun x => η x * hw.weakGrad x i) p (volume.restrict Ω) := by
      refine MemLp.of_le_mul (g := fun x => hw.weakGrad x i) (c := C₀)
        (hw.weakGrad_component_memLp i) ?_ ?_
      · exact hη.continuous.aestronglyMeasurable.mul
          (hw.weakGrad_component_memLp i).aestronglyMeasurable
      · exact Filter.Eventually.of_forall fun x => by
          calc
            ‖η x * hw.weakGrad x i‖ = |η x| * ‖hw.weakGrad x i‖ := by
              rw [norm_mul, Real.norm_eq_abs]
            _ ≤ C₀ * ‖hw.weakGrad x i‖ := by
              gcongr
              exact hη_bound x
    have hderiv_cont : Continuous (fun x => (fderiv ℝ η x) (EuclideanSpace.single i 1)) :=
      (hη.continuous_fderiv (by simp : ((⊤ : ℕ∞) : WithTop ℕ∞) ≠ 0)).clm_apply continuous_const
    have hsecond :
        MemLp (fun x => (fderiv ℝ η x) (EuclideanSpace.single i 1) * u x) p
          (volume.restrict Ω) := by
      refine MemLp.of_le_mul (g := u) (c := C₁) hw.memLp ?_ ?_
      · exact hderiv_cont.aestronglyMeasurable.mul hw.memLp.aestronglyMeasurable
      · exact Filter.Eventually.of_forall fun x => by
          calc
            ‖(fderiv ℝ η x) (EuclideanSpace.single i 1) * u x‖
                = ‖(fderiv ℝ η x) (EuclideanSpace.single i 1)‖ * ‖u x‖ := by
                    rw [norm_mul]
            _ ≤ ‖fderiv ℝ η x‖ * ‖u x‖ := by
                  gcongr
                  simpa using (ContinuousLinearMap.le_opNorm (fderiv ℝ η x)
                    (EuclideanSpace.single i (1 : ℝ)))
            _ ≤ C₁ * ‖u x‖ := by
                  gcongr
                  exact hη_grad_bound x
    simpa using hfirst.add hsecond
  isWeakGrad := by
    intro i
    simpa using HasWeakPartialDeriv.mul_smooth hΩ
      (hf := hw.isWeakGrad i) hη
      (hw.memLp.locallyIntegrable hp)
      ((hw.weakGrad_component_memLp i).locallyIntegrable hp)

\end{lstlisting}
\begin{proof}
The witness for $\eta w$ has weak gradient $\eta(x)\,\nabla w(x) + w(x)\,\nabla \eta(x)$ pointwise. Inserting into $\langle a(x)\nabla u(x), \cdot\rangle$ and using bilinearity of the inner product splits the integrand as claimed. The reformulation as a finite sum over $i$ is just unfolding the matrix-vector product and the $\Lp{2}$ inner product.
\end{proof}

\begin{theorem}[\leanname{bilinForm\_weighted\_cauchySchwarz}]
For $u, v \in W^{1,2}(\Omega)$ with witnesses $h_u, h_v$, a measurable weight $f : E \to \R$, and assuming $f^{2}\,b_A(u,u)$ is integrable, one has
\begin{multline*}
  \Bigl(\int_{\Omega} f(x)\,b_A(u,v)(x)\,dx\Bigr)^{2}
  \leq \Bigl(\int_{\Omega} f(x)^{2}\,b_A(u,u)(x)\,dx\Bigr)\\
  \cdot \Bigl(\Lambda \int_{\Omega} \|\nabla v(x)\|^{2}\,dx\Bigr).
\end{multline*}
\end{theorem}

\begin{lstlisting}[style=Matlab-editor,basicstyle=\mlttfamily,escapechar=",mlshowsectionrules=true,mathescape=true,morecomment={[s]{\%\{}{\%\}}},language=lean,numbers=none]
/-- Weighted Cauchy-Schwarz inequality for the bilinear form.

For any `f : E → ℝ` serving as a weight function:
```
|∫ f · ⟪A∇u, ∇v⟫| ≤ (∫ f² · ⟪A∇u, ∇u⟫)^{1/2} · (Λ · ∫ ‖∇v‖²)^{1/2}
```
This follows from the pointwise bound `|⟪Aξ, η⟫|² ≤ ⟪Aξ,ξ⟫ · ⟪Aη,η⟫ ≤ ⟪Aξ,ξ⟫ · Λ‖η‖²`
(the first inequality is Cauchy-Schwarz for the A-inner product, the second is the upper bound).

This is the key tool for the absorption argument in Moser iteration. -/
theorem bilinForm_weighted_cauchySchwarz
    {Ω : Set E}
    {u v : E → ℝ} {f : E → ℝ}
    (A : EllipticCoeff d Ω)
    (hu : MemW1pWitness 2 u Ω) (hv : MemW1pWitness 2 v Ω)
    (hf : AEStronglyMeasurable f (volume.restrict Ω))
    (hf_sq_int : Integrable (fun x => f x ^ 2 *
        bilinFormIntegrandOfCoeff A hu hu x) (volume.restrict Ω)) :
    (∫ x in Ω, f x * bilinFormIntegrandOfCoeff A hu hv x ∂volume) ^ 2 ≤
      (∫ x in Ω, f x ^ 2 * bilinFormIntegrandOfCoeff A hu hu x ∂volume) *
        (A.Λ * ∫ x in Ω, ‖hv.weakGrad x‖ ^ (2 : ℕ) ∂volume) := by
  let μ : Measure E := volume.restrict Ω
  -- Non-integrable case: integral = 0, so 0² = 0 ≤ RHS
  by_cases hint : Integrable (fun x => f x * bilinFormIntegrandOfCoeff A hu hv x) μ
  · -- Integrable case: use pointwise bound + Hölder
    have hpw_abs : ∀ᵐ x ∂μ, ‖f x * bilinFormIntegrandOfCoeff A hu hv x‖ ≤
        (|f x| * ‖matMulE (A.a x) (hu.weakGrad x)‖) * ‖hv.weakGrad x‖ := by
      filter_upwards with x
      rw [Real.norm_eq_abs, abs_mul, bilinFormIntegrandOfCoeff]
      calc |f x| * |⟪matMulE (A.a x) (hu.weakGrad x), hv.weakGrad x⟫_ℝ|
          ≤ |f x| * (‖matMulE (A.a x) (hu.weakGrad x)‖ * ‖hv.weakGrad x‖) :=
            mul_le_mul_of_nonneg_left (abs_real_inner_le_norm _ _) (abs_nonneg _)
        _ = (|f x| * ‖matMulE (A.a x) (hu.weakGrad x)‖) * ‖hv.weakGrad x‖ := by ring
    have hdom : ∀ᵐ x ∂μ, f x ^ 2 * ‖matMulE (A.a x) (hu.weakGrad x)‖ ^ 2 ≤
        A.Λ * (f x ^ 2 * bilinFormIntegrandOfCoeff A hu hu x) := by
      filter_upwards [A.mulVec_sq_le] with x hx
      have hmvs := hx (hu.weakGrad x)
      calc f x ^ 2 * ‖matMulE (A.a x) (hu.weakGrad x)‖ ^ 2
          ≤ f x ^ 2 * (A.Λ * ⟪hu.weakGrad x, matMulE (A.a x) (hu.weakGrad x)⟫_ℝ) :=
            mul_le_mul_of_nonneg_left hmvs (sq_nonneg _)
        _ = A.Λ * (f x ^ 2 * ⟪hu.weakGrad x, matMulE (A.a x) (hu.weakGrad x)⟫_ℝ) := by ring
        _ = A.Λ * (f x ^ 2 * bilinFormIntegrandOfCoeff A hu hu x) := by
            rw [bilinFormIntegrandOfCoeff, real_inner_comm]
    have hg_meas : AEStronglyMeasurable (fun x => |f x| * ‖matMulE (A.a x) (hu.weakGrad x)‖) μ :=
      hf.norm.mul (aestronglyMeasurable_norm_matMulE_weakGrad A hu)
    have hfsq_normAu_int : Integrable (fun x =>
        f x ^ 2 * ‖matMulE (A.a x) (hu.weakGrad x)‖ ^ 2) μ := by
      refine Integrable.mono' (hf_sq_int.const_mul A.Λ) ?_ ?_
      · exact (hg_meas.mul hg_meas).congr (ae_of_all _ (fun x => by
          simp only [Pi.mul_apply]; rw [← sq_abs (f x)]; ring))
      · filter_upwards [hdom] with x hx
        rw [Real.norm_of_nonneg (mul_nonneg (sq_nonneg _) (sq_nonneg _))]
        exact hx
    -- g² = f² * ‖A∇u‖² is integrable, so g ∈ L²
    have hg_memLp : MemLp (fun x => |f x| * ‖matMulE (A.a x) (hu.weakGrad x)‖)
        (ENNReal.ofReal 2) μ := by
      rw [show (ENNReal.ofReal (2 : ℝ)) = (2 : ℕ∞) from by norm_num]
      exact memLp_abs_mul_norm_matMulE A hu hf hfsq_normAu_int
    -- h(x) = ‖∇v(x)‖
    have hh_memLp : MemLp (fun x => ‖hv.weakGrad x‖)
        (ENNReal.ofReal 2) μ := by
      rw [show (ENNReal.ofReal (2 : ℝ)) = (2 : ℕ∞) from by norm_num]
      exact hv.weakGrad_memLp.norm
    have hholder :=
      integral_mul_norm_le_Lp_mul_Lq (μ := μ)
        (f := fun x => |f x| * ‖matMulE (A.a x) (hu.weakGrad x)‖)
        (g := fun x => ‖hv.weakGrad x‖)
        Real.HolderConjugate.two_two hg_memLp hh_memLp
    simp only [Real.norm_of_nonneg (mul_nonneg (abs_nonneg _) (norm_nonneg _)),
      Real.norm_of_nonneg (norm_nonneg _)] at hholder
    -- |∫ f*bilin| ≤ ∫ |f*bilin| ≤ ∫ g*h ≤ (∫ g²)^(1/2) * (∫ h²)^(1/2)
    -- Need: g*h integrable for integral_mono_ae
    have hgh_int : Integrable (fun x =>
        (|f x| * ‖matMulE (A.a x) (hu.weakGrad x)‖) * ‖hv.weakGrad x‖) μ := by
      have hh2_int : Integrable (fun x => ‖hv.weakGrad x‖ ^ (2 : ℕ)) μ := by
        have := (integrable_norm_rpow_iff
          hv.weakGrad_memLp.aestronglyMeasurable.norm two_ne_zero ENNReal.ofNat_ne_top).mpr
          hv.weakGrad_memLp.norm
        convert this using 1; ext x
        rw [Real.norm_of_nonneg (norm_nonneg _)]
        have h2 : ENNReal.toReal 2 = (2 : ℝ) := by norm_num
        rw [h2, show (2 : ℝ) = ((2 : ℕ) : ℝ) from by norm_num, Real.rpow_natCast]
      refine Integrable.mono'
        ((hfsq_normAu_int.add hh2_int).div_const 2)
        (hg_meas.mul hv.weakGrad_memLp.aestronglyMeasurable.norm)
        (ae_of_all _ (fun x => ?_))
      rw [Real.norm_of_nonneg (mul_nonneg (mul_nonneg (abs_nonneg _) (norm_nonneg _))
        (norm_nonneg _))]
      simp only [Pi.add_apply]
      have hg := mul_nonneg (abs_nonneg (f x)) (norm_nonneg (matMulE (A.a x) (hu.weakGrad x)))
      have hh := norm_nonneg (hv.weakGrad x)
      nlinarith [sq_nonneg (|f x| * ‖matMulE (A.a x) (hu.weakGrad x)‖ - ‖hv.weakGrad x‖),
        sq_abs (f x), mul_self_nonneg (‖matMulE (A.a x) (hu.weakGrad x)‖)]
    have habs_le : |∫ x, f x * bilinFormIntegrandOfCoeff A hu hv x ∂μ| ≤
        ∫ x, (|f x| * ‖matMulE (A.a x) (hu.weakGrad x)‖) * ‖hv.weakGrad x‖ ∂μ := by
      calc |∫ x, f x * bilinFormIntegrandOfCoeff A hu hv x ∂μ|
          = ‖∫ x, f x * bilinFormIntegrandOfCoeff A hu hv x ∂μ‖ := (Real.norm_eq_abs _).symm
        _ ≤ ∫ x, ‖f x * bilinFormIntegrandOfCoeff A hu hv x‖ ∂μ :=
            norm_integral_le_integral_norm _
        _ ≤ _ := integral_mono_ae hint.norm hgh_int hpw_abs
    have hle_rpow : |∫ x, f x * bilinFormIntegrandOfCoeff A hu hv x ∂μ| ≤
        (∫ x, (|f x| * ‖matMulE (A.a x) (hu.weakGrad x)‖) ^ (2 : ℝ) ∂μ) ^ ((1 : ℝ) / 2) *
        (∫ x, ‖hv.weakGrad x‖ ^ (2 : ℝ) ∂μ) ^ ((1 : ℝ) / 2) :=
      le_trans habs_le hholder
    have hsq : (∫ x, f x * bilinFormIntegrandOfCoeff A hu hv x ∂μ) ^ (2 : ℕ) ≤
        (∫ x, (|f x| * ‖matMulE (A.a x) (hu.weakGrad x)‖) ^ (2 : ℝ) ∂μ) *
        (∫ x, ‖hv.weakGrad x‖ ^ (2 : ℝ) ∂μ) := by
      have := sq_le_of_le_rpow_half_mul
        (integral_nonneg (fun x => by positivity))
        (integral_nonneg (fun x => by positivity))
        (abs_nonneg _) hle_rpow
      rwa [sq_abs] at this
    have hrpow_g : ∀ x, (|f x| * ‖matMulE (A.a x) (hu.weakGrad x)‖) ^ (2 : ℝ) =
        f x ^ 2 * ‖matMulE (A.a x) (hu.weakGrad x)‖ ^ 2 := by
      intro x
      rw [Real.mul_rpow (abs_nonneg _) (norm_nonneg _)]
      rw [show (2 : ℝ) = ((2 : ℕ) : ℝ) from by norm_num]
      rw [Real.rpow_natCast, Real.rpow_natCast, sq_abs]
    have hrpow_h : ∀ x, ‖hv.weakGrad x‖ ^ (2 : ℝ) = (‖hv.weakGrad x‖ ^ (2 : ℕ) : ℝ) := by
      intro x
      rw [show (2 : ℝ) = ((2 : ℕ) : ℝ) from by norm_num, Real.rpow_natCast]
    simp_rw [hrpow_g, hrpow_h] at hsq
    have hint_dom : ∫ x, f x ^ 2 * ‖matMulE (A.a x) (hu.weakGrad x)‖ ^ 2 ∂μ ≤
        A.Λ * ∫ x, f x ^ 2 * bilinFormIntegrandOfCoeff A hu hu x ∂μ := by
      calc ∫ x, f x ^ 2 * ‖matMulE (A.a x) (hu.weakGrad x)‖ ^ 2 ∂μ
          ≤ ∫ x, A.Λ * (f x ^ 2 * bilinFormIntegrandOfCoeff A hu hu x) ∂μ :=
            integral_mono_ae hfsq_normAu_int (hf_sq_int.const_mul A.Λ) hdom
        _ = A.Λ * ∫ x, f x ^ 2 * bilinFormIntegrandOfCoeff A hu hu x ∂μ :=
            integral_const_mul _ _
    -- Final assembly
    calc (∫ x, f x * bilinFormIntegrandOfCoeff A hu hv x ∂μ) ^ 2
        ≤ (∫ x, f x ^ 2 * ‖matMulE (A.a x) (hu.weakGrad x)‖ ^ 2 ∂μ) *
          (∫ x, ‖hv.weakGrad x‖ ^ (2 : ℕ) ∂μ) := hsq
      _ ≤ (A.Λ * ∫ x, f x ^ 2 * bilinFormIntegrandOfCoeff A hu hu x ∂μ) *
          (∫ x, ‖hv.weakGrad x‖ ^ (2 : ℕ) ∂μ) :=
          mul_le_mul_of_nonneg_right hint_dom (integral_nonneg (fun x => by positivity))
      _ = (∫ x, f x ^ 2 * bilinFormIntegrandOfCoeff A hu hu x ∂μ) *
          (A.Λ * ∫ x, ‖hv.weakGrad x‖ ^ (2 : ℕ) ∂μ) := by ring
  · -- Non-integrable case
    have hlhs : ∫ x in Ω, f x * bilinFormIntegrandOfCoeff A hu hv x = 0 :=
      integral_undef hint
    rw [hlhs]
    simp only [ne_eq, OfNat.ofNat_ne_zero, not_false_eq_true, zero_pow]
    exact mul_nonneg
      (integral_nonneg_of_ae (by
        filter_upwards [A.coercive] with x hx
        have hbilin_nonneg : 0 ≤ bilinFormIntegrandOfCoeff A hu hu x := by
          simp only [bilinFormIntegrandOfCoeff]
          rw [real_inner_comm]
          exact le_trans (mul_nonneg A.lam_nonneg (sq_nonneg _)) (hx (hu.weakGrad x))
        exact mul_nonneg (sq_nonneg _) hbilin_nonneg))
      (mul_nonneg A.Λ_nonneg (integral_nonneg (fun x => by positivity)))

\end{lstlisting}
\begin{proof}
If $f\,b_A(u,v)$ is non-integrable then the integral is defined to be zero and the inequality is trivial (the right-hand side is nonnegative). Otherwise, point\-wise Cauchy--Schwarz on the inner product gives
\[
  |f(x)\,b_A(u,v)(x)| \leq |f(x)|\,\|a(x)\nabla u(x)\|\,\|\nabla v(x)\|.
\]
By the mixed bound
$\|a(x)\nabla u(x)\|^{2} \leq \Lambda\,b_A(u,u)(x)$,
the function
\[
  g(x) := |f(x)|\,\|a(x)\nabla u(x)\|
\]
satisfies
\[
  g(x)^{2} = f(x)^{2}\,\|a(x)\nabla u(x)\|^{2} \leq \Lambda\,f(x)^{2}\,b_A(u,u)(x),
\]
so $g \in \Lp{2}$. H\"older's inequality with conjugate exponents $(2,2)$ applied to $g$ and $h(x) := \|\nabla v(x)\|$ gives
\[
  \int g\,h\,dx \leq \Bigl(\int g^{2}\,dx\Bigr)^{1/2}\Bigl(\int h^{2}\,dx\Bigr)^{1/2},
\]
and squaring (using $g^{2} \leq \Lambda f^{2}\,b_A(u,u)$) yields the claimed inequality.
\end{proof}

\begin{lemma}[\leanname{integrable\_bounded\_mul\_bilinFormIntegrand}]
If $|f(x)| \leq C$ everywhere and $u \in W^{1,2}(\Omega)$, then $f \cdot b_A(u,u) \in \Lp{1}(\Omega)$.
\end{lemma}

\begin{lstlisting}[style=Matlab-editor,basicstyle=\mlttfamily,escapechar=",mlshowsectionrules=true,mathescape=true,morecomment={[s]{\%\{}{\%\}}},language=lean,numbers=none]
/-- Integrability of a bounded scalar times the bilinear-form integrand.

If `|f(x)| ≤ C` everywhere and the bilinear-form integrand `⟪A∇u,∇u⟫` is
integrable (which it always is for `u ∈ W^{1,2}`), then `f · ⟪A∇u,∇u⟫` is
integrable. This is a key API lemma for the Caccioppoli/Moser absorption argument.

Proved in a standalone context to keep elaboration manageable in the presence of
`WithLp`/`PiLp` coercions. -/
theorem integrable_bounded_mul_bilinFormIntegrand
    {Ω : Set E}
    {u : E → ℝ} {f : E → ℝ} {C : ℝ}
    (A : EllipticCoeff d Ω)
    (hu : MemW1pWitness 2 u Ω)
    (hf_meas : AEStronglyMeasurable f (volume.restrict Ω))
    (hf_bound : ∀ x, |f x| ≤ C) :
    Integrable (fun x => f x * bilinFormIntegrandOfCoeff A hu hu x) (volume.restrict Ω) := by
  have hbilin_int := integrable_bilinFormIntegrandOfCoeff A hu hu
  have hC_nonneg : 0 ≤ C := le_trans (abs_nonneg (f 0)) (hf_bound 0)
  -- Use: |f · g| ≤ C · |g|, and C · |g| is integrable since g is.
  apply Integrable.mono' ((hbilin_int.norm).const_mul C)
    (hf_meas.mul (aestronglyMeasurable_bilinFormIntegrandOfCoeff A hu hu))
  filter_upwards with x
  simp only [Pi.mul_apply, Real.norm_eq_abs]
  calc ‖f x * bilinFormIntegrandOfCoeff A hu hu x‖
      = |f x| * |bilinFormIntegrandOfCoeff A hu hu x| := by
        rw [Real.norm_eq_abs, abs_mul]
    _ ≤ C * |bilinFormIntegrandOfCoeff A hu hu x| :=
        mul_le_mul_of_nonneg_right (hf_bound x) (abs_nonneg _)
    _ = C * ‖bilinFormIntegrandOfCoeff A hu hu x‖ := by
        rw [Real.norm_eq_abs]

\end{lstlisting}
\begin{proof}
Pointwise, $|f(x)\,b_A(u,u)(x)| \leq C \,|b_A(u,u)(x)|$, and the right-hand side is integrable since $b_A(u,u)$ is.
\end{proof}

\subsection{Ball scaling}

We now record the affine pullback $z \mapsto x_0 + R z$ from the unit ball $\Bz{1}$ to $\Ball{R}{x_0}$ and its effect on weak (sub-/super-)solutions. Throughout, $R > 0$.

\begin{definition}[\leanname{rescaleToUnitBall}, \leanname{transportFromUnitBall}]
For $u : E \to \R$, write
\[
  (T^{*}_{x_0,R} u)(z) := u(x_0 + R z), \qquad
  (T_{x_0,R} u)(x) := u\bigl(R^{-1}(x - x_0)\bigr),
\]
the pullback to the unit ball and the pushforward from the unit ball, respectively.
\end{definition}

\begin{lstlisting}[style=Matlab-editor,basicstyle=\mlttfamily,escapechar=",mlshowsectionrules=true,mathescape=true,morecomment={[s]{\%\{}{\%\}}},language=lean,numbers=none]
/-- Affine pullback along `z ↦ x₀ + R • z`, intended to transport functions
from `B(x₀, R)` to the unit ball. -/
noncomputable def rescaleToUnitBall
    {x₀ : E} {R : ℝ} (u : E → ℝ) : E → ℝ :=
  fun z => u (x₀ + R • z)

/-- Affine pushforward from the unit ball to `B(x₀, R)`. -/
noncomputable def transportFromUnitBall
    {x₀ : E} {R : ℝ} (u : E → ℝ) : E → ℝ :=
  fun x => u (R⁻¹ • (x - x₀))

\end{lstlisting}
\begin{lemma}[\leanname{rescaleToUnitBall\_transportFromUnitBall}]
$T^{*}_{x_0,R} \circ T_{x_0,R} = \mathrm{id}$.
\end{lemma}

\begin{lstlisting}[style=Matlab-editor,basicstyle=\mlttfamily,escapechar=",mlshowsectionrules=true,mathescape=true,morecomment={[s]{\%\{}{\%\}}},language=lean,numbers=none]
omit [NeZero d] in
theorem rescaleToUnitBall_transportFromUnitBall
    {x₀ : E} {R : ℝ} (hR : 0 < R) (u : E → ℝ) :
    rescaleToUnitBall (d := d) (x₀ := x₀) (R := R)
      (transportFromUnitBall (d := d) (x₀ := x₀) (R := R) u) = u := by
  ext z
  have hsub : x₀ + R • z - x₀ = R • z := by
    abel
  simp [rescaleToUnitBall, transportFromUnitBall, hsub, smul_smul, inv_mul_cancel₀ hR.ne']

\end{lstlisting}
\begin{proof}
Direct unfolding: $u\bigl(R^{-1}((x_0 + Rz) - x_0)\bigr) = u(z)$.
\end{proof}

\begin{definition}[\leanname{rescaleCoeffToUnitBall}, \leanname{rescaleNormalizedCoeffToUnitBall}]
Pull back an elliptic field $A$ on $\Ball{R}{x_0}$ to the unit ball by $\widetilde A(z) := (A.a)(x_0 + R z)$, with the same constants $\lam, \Lambda$. The result is an elliptic coefficient field on $\Bz{1}$. The construction preserves normalization $\lam = 1$.
\end{definition}

\begin{lstlisting}[style=Matlab-editor,basicstyle=\mlttfamily,escapechar=",mlshowsectionrules=true,mathescape=true,morecomment={[s]{\%\{}{\%\}}},language=lean,numbers=none]
/-- Pull back an elliptic coefficient from `B(x₀, R)` to the unit ball. -/
noncomputable def rescaleCoeffToUnitBall
    {x₀ : E} {R : ℝ} (hR : 0 < R)
    (A : EllipticCoeff d (Metric.ball x₀ R)) :
    EllipticCoeff d (Metric.ball (0 : E) 1) := by
  refine
    { a := fun z => A.a (x₀ + R • z)
      lam := A.lam
      Λ := A.Λ
      measurable_comp := ?_
      hlam := A.hlam
      hΛ := A.hΛ
      coercive := ?_
      coercive_inv := ?_ }
  · intro i j
    let T : E → E := fun z => x₀ + R • z
    change Measurable ((fun x => A.a x i j) ∘ T)
    exact (A.measurable_comp i j).comp ((measurable_const_add x₀).comp (measurable_const_smul R))
  · let T : E → E := fun z => x₀ + R • z
    let P : E → Prop := fun x ↦
      ∀ ξ : E, (A.lam * ‖ξ‖ ^ (2 : ℕ) ≤ ⟪ξ, matMulE (A.a x) ξ⟫_ℝ)
    have hmap :
        Measure.map T (volume.restrict (Metric.ball (0 : E) 1)) =
          ENNReal.ofReal (|R ^ Module.finrank ℝ E|⁻¹) •
            (volume.restrict (Metric.ball x₀ R)) :=
      affine_map_restrict_unitBall (d := d) (x₀ := x₀) hR
    have hcoer_map :
        ∀ᵐ x ∂ Measure.map T (volume.restrict (Metric.ball (0 : E) 1)), P x := by
      rw [hmap]
      have hscale_pos : 0 < |R ^ Module.finrank ℝ E|⁻¹ := by
        have hpow_ne : R ^ Module.finrank ℝ E ≠ 0 := by
          exact pow_ne_zero _ hR.ne'
        exact inv_pos.mpr (abs_pos.mpr hpow_ne)
      have hsmul :
          (∀ᵐ x ∂ (ENNReal.ofReal (|R ^ Module.finrank ℝ E|⁻¹) •
              volume.restrict (Metric.ball x₀ R)), P x) ↔
            ∀ᵐ x ∂ volume.restrict (Metric.ball x₀ R), P x :=
        by
          simp [ae_iff]
          intro hbad
          exfalso
          exact (not_le_of_gt (pow_pos (abs_pos.mpr hR.ne') d)) hbad
      exact hsmul.2 A.coercive
    have hpull :
        ∀ᵐ z ∂ volume.restrict (Metric.ball (0 : E) 1), P (T z) :=
      ae_of_ae_map (((measurable_const_add x₀).comp (measurable_const_smul R)).aemeasurable)
        hcoer_map
    simpa [P, T] using hpull
  · let T : E → E := fun z => x₀ + R • z
    let P : E → Prop := fun x ↦
      ∀ ξ : E, (A.Λ⁻¹ * ‖ξ‖ ^ (2 : ℕ) ≤ ⟪ξ, matMulE ((A.a x)⁻¹) ξ⟫_ℝ)
    have hmap :
        Measure.map T (volume.restrict (Metric.ball (0 : E) 1)) =
          ENNReal.ofReal (|R ^ Module.finrank ℝ E|⁻¹) •
            (volume.restrict (Metric.ball x₀ R)) :=
      affine_map_restrict_unitBall (d := d) (x₀ := x₀) hR
    have hcoer_map :
        ∀ᵐ x ∂ Measure.map T (volume.restrict (Metric.ball (0 : E) 1)), P x := by
      rw [hmap]
      have hscale_pos : 0 < |R ^ Module.finrank ℝ E|⁻¹ := by
        have hpow_ne : R ^ Module.finrank ℝ E ≠ 0 := by
          exact pow_ne_zero _ hR.ne'
        exact inv_pos.mpr (abs_pos.mpr hpow_ne)
      have hsmul :
          (∀ᵐ x ∂ (ENNReal.ofReal (|R ^ Module.finrank ℝ E|⁻¹) •
              volume.restrict (Metric.ball x₀ R)), P x) ↔
            ∀ᵐ x ∂ volume.restrict (Metric.ball x₀ R), P x :=
        by
          simp [ae_iff]
          intro hbad
          exfalso
          exact (not_le_of_gt (pow_pos (abs_pos.mpr hR.ne') d)) hbad
      exact hsmul.2 A.coercive_inv
    have hpull :
        ∀ᵐ z ∂ volume.restrict (Metric.ball (0 : E) 1), P (T z) :=
      ae_of_ae_map (((measurable_const_add x₀).comp (measurable_const_smul R)).aemeasurable)
        hcoer_map
    simpa [P, T] using hpull

/-- Pull back a normalized elliptic coefficient from `B(x₀, R)` to the unit
ball. -/
noncomputable def rescaleNormalizedCoeffToUnitBall
    {x₀ : E} {R : ℝ} (hR : 0 < R)
    (A : NormalizedEllipticCoeff d (Metric.ball x₀ R)) :
    NormalizedEllipticCoeff d (Metric.ball (0 : E) 1) := by
  refine ⟨rescaleCoeffToUnitBall (d := d) (x₀ := x₀) (R := R) hR A.1, by
    simp [rescaleCoeffToUnitBall]⟩

\end{lstlisting}
\begin{lemma}[\leanname{MemW1pWitness.transportFromUnitBall}]
A $W^{1,p}$-witness for $u$ on $\Bz{1}$ transports to a $W^{1,p}$-witness for $T_{x_0,R} u$ on $\Ball{R}{x_0}$, with weak gradient $x \mapsto R^{-1}\,(\nabla u)\bigl(R^{-1}(x - x_0)\bigr)$.
\end{lemma}

\begin{lstlisting}[style=Matlab-editor,basicstyle=\mlttfamily,escapechar=",mlshowsectionrules=true,mathescape=true,morecomment={[s]{\%\{}{\%\}}},language=lean,numbers=none]
/-- Rescaling a unit-ball witness to a witness on `B(x₀, R)`. -/
noncomputable def MemW1pWitness.transportFromUnitBall
    {p : ENNReal} {x₀ : E} {R : ℝ} {u : E → ℝ}
    (hR : 0 < R)
    (hw : MemW1pWitness p u (Metric.ball (0 : E) 1)) :
    MemW1pWitness p (transportFromUnitBall (d := d) (x₀ := x₀) (R := R) u)
      (Metric.ball x₀ R) where
  memLp := transportFromUnitBall_memLp (d := d) (x₀ := x₀) (R := R) hR hw.memLp
  weakGrad := fun x => R⁻¹ • hw.weakGrad (R⁻¹ • (x - x₀))
  weakGrad_component_memLp := by
    intro i
    simpa [DeGiorgi.transportFromUnitBall, Pi.smul_apply, smul_eq_mul, sub_eq_add_neg] using
      (transportFromUnitBall_memLp (d := d) (x₀ := x₀) (R := R) hR
        (hw.weakGrad_component_memLp i)).const_mul R⁻¹
  isWeakGrad := by
    intro i
    simpa [DeGiorgi.transportFromUnitBall, Pi.smul_apply, smul_eq_mul, sub_eq_add_neg] using
      weakPartialDeriv_transportFromUnitBall (d := d) (i := i)
        (x₀ := x₀) (R := R) (u := u) (g := fun x => hw.weakGrad x i) hR (hw.isWeakGrad i)

\end{lstlisting}
\begin{proof}
The Lp norms transform by an explicit power of $R$ depending on the dimension; weak partial derivative compatibility is preserved by the chain rule for the affine map $R^{-1}(\cdot - x_0)$.
\end{proof}

\begin{theorem}[\leanname{rescaleToUnitBall\_isSubsolution}, \leanname{\_isSupersolution}, \leanname{\_isSolution}, \leanname{\_isHomogeneousWeakSolution}]
If $u$ is a weak subsolution (resp.\ supersolution, solution, homogeneous weak solution) of $A$ on $\Ball{R}{x_0}$, then $T^{*}_{x_0,R} u$ is the corresponding object for $\widetilde A$ on $\Bz{1}$.
\end{theorem}

\begin{lstlisting}[style=Matlab-editor,basicstyle=\mlttfamily,escapechar=",mlshowsectionrules=true,mathescape=true,morecomment={[s]{\%\{}{\%\}}},language=lean,numbers=none]
/-- Transport subsolutions on `B(x₀, R)` to the unit ball via affine pullback. -/
theorem rescaleToUnitBall_isSubsolution
    {x₀ : E} {R : ℝ} (hR : 0 < R)
    (A : EllipticCoeff d (Metric.ball x₀ R))
    {u : E → ℝ}
    (hsub : IsSubsolution A u) :
    IsSubsolution (rescaleCoeffToUnitBall (d := d) hR A)
      (rescaleToUnitBall (d := d) (x₀ := x₀) (R := R) u) := by
  rcases hsub with ⟨hu_mem, hsub_ineq⟩
  let hu : MemW1pWitness 2 u (Metric.ball x₀ R) := DeGiorgi.MemW1p.someWitness hu_mem
  let huR := hu.rescale_to_unitBall (d := d) (x₀ := x₀) (R := R) hR
  refine ⟨huR.memW1p, ?_⟩
  intro hu' φ hφ0 hφ hφ_nonneg
  let hφB := hφ.transportFromUnitBall (d := d) (x₀ := x₀) (R := R) hR
  have hφB0 :
      MemH01 (DeGiorgi.transportFromUnitBall (d := d) (x₀ := x₀) (R := R) φ)
        (Metric.ball x₀ R) :=
    MemW01p.transportFromUnitBall (d := d) (x₀ := x₀) (R := R) hR hφ0
  have hball : bilinFormOfCoeff A hu hφB ≤ 0 := by
    exact hsub_ineq hu _ hφB0 hφB (by
      intro x
      exact hφ_nonneg (R⁻¹ • (x - x₀)))
  have hconst_nonneg : 0 ≤ ((R ^ Module.finrank ℝ E)⁻¹ * R ^ (2 : ℕ) : ℝ) := by
    positivity
  calc
    bilinFormOfCoeff (rescaleCoeffToUnitBall (d := d) (x₀ := x₀) (R := R) hR A) hu' hφ
        = bilinFormOfCoeff (rescaleCoeffToUnitBall (d := d) (x₀ := x₀) (R := R) hR A) huR hφ := by
            exact bilinFormOfCoeff_eq_left Metric.isOpen_ball _ hu' huR hφ
    _ = ((R ^ Module.finrank ℝ E)⁻¹ * R ^ (2 : ℕ)) * bilinFormOfCoeff A hu hφB := by
          simpa [huR, hφB] using bilinForm_rescaleToUnitBall (d := d) (x₀ := x₀) (R := R) hR A hu hφ
    _ ≤ 0 := mul_nonpos_of_nonneg_of_nonpos hconst_nonneg hball

\end{lstlisting}
\begin{proof}
The key computation is
\[
  \mathfrak{a}_{\widetilde A}\bigl(T^{*}_{x_0,R} u, \varphi\bigr) = \bigl(R^{-d} \cdot R^{2}\bigr)\,\mathfrak{a}_{A}\bigl(u, T_{x_0,R} \varphi\bigr),
\]
proven by changing variables in the integral defining $\mathfrak{a}$, using the Jacobian factor $R^{-d}$ and the chain-rule factor $R^{-1}$ in each gradient (giving an $R^{2}$ from the two slots). Since the multiplicative constant $R^{2-d} > 0$ has fixed sign, the inequality $\mathfrak{a}_{\widetilde A}(\cdot, \varphi) \leq 0$ on the unit ball is equivalent to $\mathfrak{a}_A(\cdot, T_{x_0,R}\varphi) \leq 0$ on $\Ball{R}{x_0}$, and similarly for the reverse inequality, equality, and the homogeneous case. Nonnegative test functions $\varphi$ on $\Bz{1}$ pull back to nonnegative test functions $T_{x_0,R}\varphi$ on $\Ball{R}{x_0}$ (and vice versa), so the sign condition is preserved.
\end{proof}

\subsection{Localization helpers}

\begin{definition}[\leanname{EllipticCoeff.restrict}, \leanname{NormalizedEllipticCoeff.restrict}]
Given $\Omega' \subseteq \Omega$ and $A$ elliptic on $\Omega$, the restriction $A\!\restriction\!\Omega'$ keeps the same matrix and constants and shrinks the a.e.\ coercivity domain to $\Omega'$. Normalization is preserved.
\end{definition}

\begin{lstlisting}[style=Matlab-editor,basicstyle=\mlttfamily,escapechar=",mlshowsectionrules=true,mathescape=true,morecomment={[s]{\%\{}{\%\}}},language=lean,numbers=none]
/-- Restrict an elliptic coefficient field to a smaller domain by keeping the
same coefficient matrix and constants and shrinking the a.e. coercivity domain.
-/
noncomputable def restrict
    {Ω Ω' : Set E} (A : EllipticCoeff d Ω) (hsub : Ω' ⊆ Ω) :
    EllipticCoeff d Ω' where
  a := A.a
  lam := A.lam
  Λ := A.Λ
  measurable_comp := A.measurable_comp
  hlam := A.hlam
  hΛ := A.hΛ
  coercive := ae_restrict_of_ae_restrict_of_subset hsub A.coercive
  coercive_inv := ae_restrict_of_ae_restrict_of_subset hsub A.coercive_inv

/-- Restrict a normalized elliptic coefficient field to a smaller domain. -/
  noncomputable def restrict
    {Ω Ω' : Set E} (A : NormalizedEllipticCoeff d Ω) (hsub : Ω' ⊆ Ω) :
    NormalizedEllipticCoeff d Ω' := by
  refine ⟨A.1.restrict hsub, by simp⟩

\end{lstlisting}
\begin{lemma}[\leanname{essSup\_rescale\_halfBall}, \leanname{essInf\_rescale\_halfBall}]
For $R > 0$,
\[
  \esssup_{z \in \Bz{1/2}} u(x_0 + R z) = \esssup_{x \in \Ball{R/2}{x_0}} u(x),
\]
and similarly for $\essinf$.
\end{lemma}

\begin{lstlisting}[style=Matlab-editor,basicstyle=\mlttfamily,escapechar=",mlshowsectionrules=true,mathescape=true,morecomment={[s]{\%\{}{\%\}}},language=lean,numbers=none]
theorem essSup_rescale_halfBall
    {x₀ : E} {R : ℝ} (hR : 0 < R) {u : E → ℝ} :
    essSup (fun z : E => u (x₀ + R • z))
        (volume.restrict (Metric.ball (0 : E) (1 / 2 : ℝ))) =
      essSup u (volume.restrict (Metric.ball x₀ (R / 2 : ℝ))) := by
  let μsrc : Measure E := volume.restrict (Metric.ball (0 : E) (1 / 2 : ℝ))
  let μdst : Measure E := volume.restrict (Metric.ball x₀ (R / 2 : ℝ))
  let T : E → E := fun z => x₀ + R • z
  have hmap :
      Measure.map T μsrc =
        ENNReal.ofReal (|R ^ Module.finrank ℝ E|⁻¹) • μdst := by
    simpa [μsrc, μdst, T, show R * (1 / 2 : ℝ) = R / 2 by ring] using
      affine_map_restrict_ball_mul (d := d) (x₀ := x₀) (R := R) (ρ := (1 / 2 : ℝ)) hR
  have hiff :
      ∀ a : ℝ,
        (∀ᵐ z ∂μsrc, u (T z) ≤ a) ↔ ∀ᵐ x ∂μdst, u x ≤ a := by
    intro a
    constructor
    · intro h
      have hmap_ae : ∀ᵐ x ∂ Measure.map T μsrc, u x ≤ a :=
        (affineMeasurableEmbedding (d := d) (x₀ := x₀) (R := R) hR).ae_map_iff.2 h
      rw [hmap, Measure.ae_ennreal_smul_measure_eq
        (affine_scale_measure_ne_zero (d := d) hR)] at hmap_ae
      exact hmap_ae
    · intro h
      have hmap_ae : ∀ᵐ x ∂ Measure.map T μsrc, u x ≤ a := by
        rw [hmap]
        rw [Measure.ae_ennreal_smul_measure_eq
          (affine_scale_measure_ne_zero (d := d) hR)]
        exact h
      exact (affineMeasurableEmbedding (d := d) (x₀ := x₀) (R := R) hR).ae_map_iff.1 hmap_ae
  rw [essSup_eq_sInf, essSup_eq_sInf]
  congr 1
  ext a
  simpa [μsrc, μdst, T, ae_iff, not_le] using hiff a

theorem essInf_rescale_halfBall
    {x₀ : E} {R : ℝ} (hR : 0 < R) {u : E → ℝ} :
    essInf (fun z : E => u (x₀ + R • z))
        (volume.restrict (Metric.ball (0 : E) (1 / 2 : ℝ))) =
      essInf u (volume.restrict (Metric.ball x₀ (R / 2 : ℝ))) := by
  let μsrc : Measure E := volume.restrict (Metric.ball (0 : E) (1 / 2 : ℝ))
  let μdst : Measure E := volume.restrict (Metric.ball x₀ (R / 2 : ℝ))
  let T : E → E := fun z => x₀ + R • z
  have hmap :
      Measure.map T μsrc =
        ENNReal.ofReal (|R ^ Module.finrank ℝ E|⁻¹) • μdst := by
    simpa [μsrc, μdst, T, show R * (1 / 2 : ℝ) = R / 2 by ring] using
      affine_map_restrict_ball_mul (d := d) (x₀ := x₀) (R := R) (ρ := (1 / 2 : ℝ)) hR
  have hiff :
      ∀ a : ℝ,
        (∀ᵐ z ∂μsrc, a ≤ u (T z)) ↔ ∀ᵐ x ∂μdst, a ≤ u x := by
    intro a
    constructor
    · intro h
      have hmap_ae : ∀ᵐ x ∂ Measure.map T μsrc, a ≤ u x :=
        (affineMeasurableEmbedding (d := d) (x₀ := x₀) (R := R) hR).ae_map_iff.2 h
      rw [hmap, Measure.ae_ennreal_smul_measure_eq
        (affine_scale_measure_ne_zero (d := d) hR)] at hmap_ae
      exact hmap_ae
    · intro h
      have hmap_ae : ∀ᵐ x ∂ Measure.map T μsrc, a ≤ u x := by
        rw [hmap]
        rw [Measure.ae_ennreal_smul_measure_eq
          (affine_scale_measure_ne_zero (d := d) hR)]
        exact h
      exact (affineMeasurableEmbedding (d := d) (x₀ := x₀) (R := R) hR).ae_map_iff.1 hmap_ae
  rw [essInf_eq_sSup, essInf_eq_sSup]
  congr 1
  ext a
  simpa [μsrc, μdst, T, ae_iff, not_le] using hiff a

\end{lstlisting}
\begin{proof}
The affine map $z \mapsto x_0 + R z$ is a measurable embedding sending the half unit ball onto $\Ball{R/2}{x_0}$ with the volume measure pushing forward to a strictly positive multiple of the restricted volume on $\Ball{R/2}{x_0}$. Strict positivity of the constant scaling implies that the a.e.\ filters coincide, so the essential suprema (and infima) match.
\end{proof}

\begin{lemma}[\leanname{memH01\_indicator\_ball\_of\_closedBall\_subset}, \leanname{ballIndicatorWitnessOn}]
If $\overline{\Ball{r}{c}} \subseteq \Omega$ and $\varphi \in H^{1}_{0}(\Ball{r}{c})$, then the zero extension $\mathbf{1}_{\Ball{r}{c}}\,\varphi$ lies in $H^{1}_{0}(\Omega)$ and carries an explicit $W^{1,2}(\Omega)$-witness whose weak gradient is the zero extension of the gradient of $\varphi$.
\end{lemma}

\begin{lstlisting}[style=Matlab-editor,basicstyle=\mlttfamily,escapechar=",mlshowsectionrules=true,mathescape=true,morecomment={[s]{\%\{}{\%\}}},language=lean,numbers=none]
theorem memH01_indicator_ball_of_closedBall_subset
    {Ω : Set E} (hΩ : IsOpen Ω)
    {c : E} {r : ℝ} (_hr : 0 < r)
    (hsub : Metric.closedBall c r ⊆ Ω)
    {φ : E → ℝ} (hφ0 : MemH01 φ (Metric.ball c r)) :
    MemH01 ((Metric.ball c r).indicator φ) Ω := by
  let hφ0_real : MemW01p (ENNReal.ofReal (2 : ℝ)) φ (Metric.ball c r) := by
    simpa using hφ0
  have hball0 : MemW01p 2 ((Metric.ball c r).indicator φ) Set.univ := by
    simpa using
      (zeroExtend_memW01p_p (d := d) Metric.isOpen_ball
        (p := 2) (by norm_num : (1 : ℝ) < 2) hφ0_real)
  let hballw : MemW1pWitness 2 ((Metric.ball c r).indicator φ) Set.univ :=
    MemW1p.someWitness (MemW01p.memW1p hball0)
  have hball_memW1p_Ω :
      MemW1p (ENNReal.ofReal (2 : ℝ)) ((Metric.ball c r).indicator φ) Ω := by
    simpa using (hballw.restrict hΩ (by intro x hx; simp)).memW1p
  have hcompact :
      HasCompactSupport ((Metric.ball c r).indicator φ) := by
    refine HasCompactSupport.intro'
      (isCompact_closedBall c r) Metric.isClosed_closedBall ?_
    intro x hx
    have hxball : x ∉ Metric.ball c r := by
      intro hxball
      exact hx (Metric.mem_closedBall.2 (le_of_lt (Metric.mem_ball.1 hxball)))
    simp [hxball]
  have htsub : tsupport ((Metric.ball c r).indicator φ) ⊆ Ω := by
    exact (ball_indicator_tsupport_subset_closedBall (c := c) (r := r) (φ := φ)).trans hsub
  simpa using
    (memW01p_of_memW1p_of_tsupport_subset
      (d := d) hΩ (by norm_num : (1 : ℝ) < 2) hball_memW1p_Ω hcompact htsub)

noncomputable def ballIndicatorWitnessOn
    {Ω : Set E} (hΩ : IsOpen Ω)
    {c : E} {r : ℝ} {φ : E → ℝ}
    (hφ0 : MemH01 φ (Metric.ball c r))
    (hφ : MemW1pWitness 2 φ (Metric.ball c r)) :
    MemW1pWitness 2 ((Metric.ball c r).indicator φ) Ω := by
  let hφ0_real : MemW01p (ENNReal.ofReal (2 : ℝ)) φ (Metric.ball c r) := by
    simpa using hφ0
  let hφ_real : MemW1pWitness (ENNReal.ofReal (2 : ℝ)) φ (Metric.ball c r) := by
    refine
      { memLp := ?_
        weakGrad := hφ.weakGrad
        weakGrad_component_memLp := ?_
        isWeakGrad := hφ.isWeakGrad }
    · simpa using hφ.memLp
    · intro i
      simpa using hφ.weakGrad_component_memLp i
  let hφExtUniv : MemW1pWitness (ENNReal.ofReal (2 : ℝ)) ((Metric.ball c r).indicator φ) Set.univ :=
    zeroExtend_memW1pWitness_p (d := d) Metric.isOpen_ball
      (p := 2) (by norm_num : (1 : ℝ) < 2) hφ0_real hφ_real
  have htwo : ENNReal.ofReal (2 : ℝ) = (2 : ENNReal) := by
    norm_num
  simpa [htwo] using hφExtUniv.restrict hΩ (by intro x hx; simp)

\end{lstlisting}
\begin{proof}
Zero-extension preserves $H^{1}_{0}$ on open sets (the approximating sequence of compactly supported smooth functions extends by zero). Compact support inside $\overline{\Ball{r}{c}} \subseteq \Omega$ together with the universal $W^{1,2}$-witness restricts to a $W^{1,2}(\Omega)$-witness with the same weak gradient.
\end{proof}

\begin{lemma}[\leanname{bilinForm\_ball\_restrict\_eq\_zeroExtend}]
Under the same hypotheses, for $u \in W^{1,2}(\Omega)$ and $\varphi \in H^{1}_{0}(\Ball{r}{c})$,
\[
  \mathfrak{a}_{A\restriction \Ball{r}{c}}\bigl(u\!\restriction, \varphi\bigr)
  = \mathfrak{a}_{A}\bigl(u, \mathbf{1}_{\Ball{r}{c}}\,\varphi\bigr).
\]
\end{lemma}

\begin{lstlisting}[style=Matlab-editor,basicstyle=\mlttfamily,escapechar=",mlshowsectionrules=true,mathescape=true,morecomment={[s]{\%\{}{\%\}}},language=lean,numbers=none]
theorem bilinForm_ball_restrict_eq_zeroExtend
    {Ω : Set E} (hΩ : IsOpen Ω)
    (A : EllipticCoeff d Ω)
    {c : E} {r : ℝ} (_hr : 0 < r)
    (hsub : Metric.closedBall c r ⊆ Ω)
    {u : E → ℝ}
    (hu : MemW1pWitness 2 u Ω)
    {φ : E → ℝ}
    (hφ0 : MemH01 φ (Metric.ball c r))
    (hφ : MemW1pWitness 2 φ (Metric.ball c r)) :
    bilinFormOfCoeff (A.restrict (Metric.ball_subset_closedBall.trans hsub))
        (hu.restrict Metric.isOpen_ball (Metric.ball_subset_closedBall.trans hsub)) hφ =
      bilinFormOfCoeff A hu (ballIndicatorWitnessOn (d := d) hΩ hφ0 hφ) := by
  let hφExt : MemW1pWitness 2 ((Metric.ball c r).indicator φ) Ω :=
    ballIndicatorWitnessOn (d := d) hΩ hφ0 hφ
  let A' : EllipticCoeff d (Metric.ball c r) :=
    A.restrict (Metric.ball_subset_closedBall.trans hsub)
  let hu' : MemW1pWitness 2 u (Metric.ball c r) :=
    hu.restrict Metric.isOpen_ball (Metric.ball_subset_closedBall.trans hsub)
  let fsmall : E → ℝ := fun x => bilinFormIntegrandOfCoeff A' hu' hφ x
  have hball_sub : Metric.ball c r ⊆ Ω := by
    intro x hx
    exact hsub (Metric.mem_closedBall.2 (le_of_lt (Metric.mem_ball.1 hx)))
  have hEq_ae :
      (fun x => (Metric.ball c r).indicator fsmall x) =ᵐ[volume.restrict Ω]
        fun x => bilinFormIntegrandOfCoeff A hu hφExt x := by
    filter_upwards with x
    by_cases hx : x ∈ Metric.ball c r
    · have hgradφ : hφExt.weakGrad x = hφ.weakGrad x := by
        apply PiLp.ext
        intro i
        simp [hφExt, ballIndicatorWitnessOn_weakGrad_apply, hx]
      simp [A', hu', fsmall, hφExt, hgradφ, EllipticCoeff.restrict,
        MemW1pWitness.restrict, hx, bilinFormIntegrandOfCoeff]
    · have hgradφ : hφExt.weakGrad x = 0 := by
        apply PiLp.ext
        intro i
        simp [hφExt, ballIndicatorWitnessOn_weakGrad_apply, hx]
      simp [A', hu', fsmall, hφExt, hgradφ, EllipticCoeff.restrict,
        MemW1pWitness.restrict, hx, bilinFormIntegrandOfCoeff]
  calc
    bilinFormOfCoeff A' hu' hφ
      = ∫ x in Metric.ball c r, fsmall x ∂(volume.restrict Ω) := by
          simp [bilinFormOfCoeff, fsmall, Measure.restrict_restrict_of_subset hball_sub]
    _ = ∫ x, (Metric.ball c r).indicator fsmall x ∂(volume.restrict Ω) := by
          symm
          exact integral_indicator measurableSet_ball
    _ = ∫ x, bilinFormIntegrandOfCoeff A hu hφExt x ∂(volume.restrict Ω) := by
          exact integral_congr_ae hEq_ae
    _ = bilinFormOfCoeff A hu hφExt := by
          rfl

\end{lstlisting}
\begin{proof}
Pointwise, the two integrands agree on $\Ball{r}{c}$ and the right-hand integrand vanishes outside (the indicator-extended witness has zero weak gradient there). Hence both integrals over $\Omega$ equal the integral over $\Ball{r}{c}$ of the same integrand.
\end{proof}

\begin{lemma}[\leanname{MemW1pWitness.of\_ae\_eq}]
If $f =_{a.e.} g$ on $\Omega$ and $f$ has a $W^{1,2}(\Omega)$-witness, then $g$ has a witness with the same weak gradient.
\end{lemma}

\begin{lstlisting}[style=Matlab-editor,basicstyle=\mlttfamily,escapechar=",mlshowsectionrules=true,mathescape=true,morecomment={[s]{\%\{}{\%\}}},language=lean,numbers=none]
noncomputable def MemW1pWitness.of_ae_eq
    {Ω : Set E} {f g : E → ℝ}
    (hfg : f =ᵐ[volume.restrict Ω] g)
    (hw : MemW1pWitness 2 f Ω) :
    MemW1pWitness 2 g Ω := by
  refine
    { memLp := (memLp_congr_ae hfg).1 hw.memLp
      weakGrad := hw.weakGrad
      weakGrad_component_memLp := hw.weakGrad_component_memLp
      isWeakGrad := ?_ }
  intro i φ hφ hφ_supp hφ_sub
  have hweak := hw.isWeakGrad i φ hφ hφ_supp hφ_sub
  have hcongr :
      (fun x => g x * (fderiv ℝ φ x) (EuclideanSpace.single i 1)) =ᵐ[volume.restrict Ω]
        (fun x => f x * (fderiv ℝ φ x) (EuclideanSpace.single i 1)) := by
    filter_upwards [hfg] with x hx
    simp [hx]
  calc
    ∫ x in Ω, g x * (fderiv ℝ φ x) (EuclideanSpace.single i 1)
        = ∫ x in Ω, f x * (fderiv ℝ φ x) (EuclideanSpace.single i 1) := by
            exact integral_congr_ae hcongr
    _ = -∫ x in Ω, hw.weakGrad x i * φ x := hweak

\end{lstlisting}
\begin{proof}
Lp membership transfers under a.e.\ equality. The weak partial derivative identity $\int g\,\partial_i\varphi = -\int (\nabla f)_i\,\varphi$ follows from $\int g\,\partial_i\varphi = \int f\,\partial_i\varphi$ by a.e.\ equality.
\end{proof}

\begin{theorem}[\leanname{IsSubsolution.congr\_ae}, \leanname{IsSupersolution.congr\_ae}, \leanname{IsSolution.congr\_ae}]
The weak (sub-, super-) solution predicates are stable under a.e.\ equality of the unknown.
\end{theorem}

\begin{lstlisting}[style=Matlab-editor,basicstyle=\mlttfamily,escapechar=",mlshowsectionrules=true,mathescape=true,morecomment={[s]{\%\{}{\%\}}},language=lean,numbers=none]
theorem IsSubsolution.congr_ae
    {Ω : Set E} {A : EllipticCoeff d Ω} {u v : E → ℝ}
    (huv : u =ᵐ[volume.restrict Ω] v)
    (hsub : IsSubsolution A u) :
    IsSubsolution A v := by
  refine ⟨(MemW1pWitness.of_ae_eq huv (MemW1p.someWitness hsub.1)).memW1p, ?_⟩
  intro hv φ hφ hφw hφ_nonneg
  let hu : MemW1pWitness 2 u Ω := MemW1pWitness.of_ae_eq huv.symm hv
  have hineq : bilinFormOfCoeff A hu hφw ≤ 0 := hsub.2 hu φ hφ hφw hφ_nonneg
  simpa [hu, MemW1pWitness.of_ae_eq] using hineq

theorem IsSupersolution.congr_ae
    {Ω : Set E} {A : EllipticCoeff d Ω} {u v : E → ℝ}
    (huv : u =ᵐ[volume.restrict Ω] v)
    (hsuper : IsSupersolution A u) :
    IsSupersolution A v := by
  refine ⟨(MemW1pWitness.of_ae_eq huv (MemW1p.someWitness hsuper.1)).memW1p, ?_⟩
  intro hv φ hφ hφw hφ_nonneg
  let hu : MemW1pWitness 2 u Ω := MemW1pWitness.of_ae_eq huv.symm hv
  have hineq : 0 ≤ bilinFormOfCoeff A hu hφw := hsuper.2 hu φ hφ hφw hφ_nonneg
  simpa [hu, MemW1pWitness.of_ae_eq] using hineq

theorem IsSolution.congr_ae
    {Ω : Set E} {A : EllipticCoeff d Ω} {u v : E → ℝ}
    (huv : u =ᵐ[volume.restrict Ω] v)
    (hsol : IsSolution A u) :
    IsSolution A v := by
  exact ⟨hsol.1.congr_ae huv, hsol.2.congr_ae huv⟩

\end{lstlisting}
\begin{proof}
Use the previous lemma to transport witnesses, and observe that $\mathfrak{a}_A$ depends only on the witness's weak gradient, which is unchanged.
\end{proof}

\begin{theorem}[\leanname{IsSubsolution.restrict\_ball}, \leanname{IsSupersolution.restrict\_ball}, \leanname{IsSolution.restrict\_ball}]
Let $\Omega$ be open, $\overline{\Ball{r}{c}} \subseteq \Omega$, and $A$ elliptic on $\Omega$. If $u$ is a weak subsolution (resp.\ super-, solution) of $A$ on $\Omega$, then $u$ is a weak subsolution (resp.\ super-, solution) of $A\!\restriction\!\Ball{r}{c}$ on the smaller domain.
\end{theorem}

\begin{lstlisting}[style=Matlab-editor,basicstyle=\mlttfamily,escapechar=",mlshowsectionrules=true,mathescape=true,morecomment={[s]{\%\{}{\%\}}},language=lean,numbers=none]
theorem IsSubsolution.restrict_ball
    {Ω : Set E} (hΩ : IsOpen Ω)
    {c : E} {r : ℝ} (hr : 0 < r)
    {A : EllipticCoeff d Ω} {u : E → ℝ}
    (hsub : Metric.closedBall c r ⊆ Ω)
    (hu : IsSubsolution A u) :
    IsSubsolution (A.restrict (Metric.ball_subset_closedBall.trans hsub)) u := by
  rcases hu with ⟨hu_mem, hu_sub⟩
  let huw : MemW1pWitness 2 u Ω := MemW1p.someWitness hu_mem
  refine ⟨(huw.restrict Metric.isOpen_ball (Metric.ball_subset_closedBall.trans hsub)).memW1p, ?_⟩
  intro hu' φ hφ0 hφ hφ_nonneg
  have hφ0_big :
      MemH01 ((Metric.ball c r).indicator φ) Ω :=
    memH01_indicator_ball_of_closedBall_subset (d := d) hΩ hr hsub hφ0
  let hφBig : MemW1pWitness 2 ((Metric.ball c r).indicator φ) Ω :=
    ballIndicatorWitnessOn (d := d) hΩ hφ0 hφ
  have hbig :
      bilinFormOfCoeff A huw hφBig ≤ 0 := by
    exact hu_sub huw _ hφ0_big hφBig (by
      intro x
      by_cases hx : x ∈ Metric.ball c r
      · simpa [hx] using hφ_nonneg x
      · simp [hx])
  calc
    bilinFormOfCoeff (A.restrict (Metric.ball_subset_closedBall.trans hsub)) hu' hφ
      = bilinFormOfCoeff (A.restrict (Metric.ball_subset_closedBall.trans hsub))
          (huw.restrict Metric.isOpen_ball (Metric.ball_subset_closedBall.trans hsub)) hφ := by
            exact bilinFormOfCoeff_eq_left Metric.isOpen_ball
              (A.restrict (Metric.ball_subset_closedBall.trans hsub)) hu'
              (huw.restrict Metric.isOpen_ball (Metric.ball_subset_closedBall.trans hsub)) hφ
    _ = bilinFormOfCoeff A huw
          (ballIndicatorWitnessOn (d := d) hΩ hφ0 hφ) := by
          exact bilinForm_ball_restrict_eq_zeroExtend (d := d) hΩ A hr hsub huw hφ0 hφ
    _ ≤ 0 := hbig

theorem IsSupersolution.restrict_ball
    {Ω : Set E} (hΩ : IsOpen Ω)
    {c : E} {r : ℝ} (hr : 0 < r)
    {A : EllipticCoeff d Ω} {u : E → ℝ}
    (hsub : Metric.closedBall c r ⊆ Ω)
    (hu : IsSupersolution A u) :
    IsSupersolution (A.restrict (Metric.ball_subset_closedBall.trans hsub)) u := by
  rcases hu with ⟨hu_mem, hu_super⟩
  let huw : MemW1pWitness 2 u Ω := MemW1p.someWitness hu_mem
  refine ⟨(huw.restrict Metric.isOpen_ball (Metric.ball_subset_closedBall.trans hsub)).memW1p, ?_⟩
  intro hu' φ hφ0 hφ hφ_nonneg
  have hφ0_big :
      MemH01 ((Metric.ball c r).indicator φ) Ω :=
    memH01_indicator_ball_of_closedBall_subset (d := d) hΩ hr hsub hφ0
  let hφBig : MemW1pWitness 2 ((Metric.ball c r).indicator φ) Ω :=
    ballIndicatorWitnessOn (d := d) hΩ hφ0 hφ
  have hbig :
      0 ≤ bilinFormOfCoeff A huw hφBig := by
    exact hu_super huw _ hφ0_big hφBig (by
      intro x
      by_cases hx : x ∈ Metric.ball c r
      · simpa [hx] using hφ_nonneg x
      · simp [hx])
  calc
    0 ≤ bilinFormOfCoeff A huw
          (ballIndicatorWitnessOn (d := d) hΩ hφ0 hφ) := by
          exact hbig
    _ = bilinFormOfCoeff (A.restrict (Metric.ball_subset_closedBall.trans hsub))
          (huw.restrict Metric.isOpen_ball (Metric.ball_subset_closedBall.trans hsub)) hφ := by
          symm
          exact bilinForm_ball_restrict_eq_zeroExtend (d := d) hΩ A hr hsub huw hφ0 hφ
    _ = bilinFormOfCoeff (A.restrict (Metric.ball_subset_closedBall.trans hsub)) hu' hφ := by
          symm
          exact bilinFormOfCoeff_eq_left Metric.isOpen_ball
            (A.restrict (Metric.ball_subset_closedBall.trans hsub))
            hu' (huw.restrict Metric.isOpen_ball (Metric.ball_subset_closedBall.trans hsub)) hφ

theorem IsSolution.restrict_ball
    {Ω : Set E} (hΩ : IsOpen Ω)
    {c : E} {r : ℝ} (hr : 0 < r)
    {A : EllipticCoeff d Ω} {u : E → ℝ}
    (hsub : Metric.closedBall c r ⊆ Ω)
    (hu : IsSolution A u) :
    IsSolution (A.restrict (Metric.ball_subset_closedBall.trans hsub)) u := by
  exact ⟨hu.1.restrict_ball (d := d) hΩ hr hsub, hu.2.restrict_ball (d := d) hΩ hr hsub⟩

\end{lstlisting}
\begin{proof}
We treat the subsolution case; the supersolution and solution cases are
symmetric. Pick a witness $h_u^{\Omega}$ for $u$ on $\Omega$ via
\leanname{MemW1p.someWitness}. The restriction lemma
\leanname{MemW1pWitness.restrict} (specialised to the open subball
$\Ball{r}{c} \subseteq \Omega$) yields a witness $h_u^{\Ball{r}{c}}$ for $u$
on the smaller domain, with weak gradient given by the restriction of
$\nabla u$ to $\Ball{r}{c}$; this verifies the $W^{1,2}$ membership half of
\leanname{IsSubsolution}.

For the bilinear-form inequality, fix a nonnegative test
$\varphi \in H^{1}_{0}(\Ball{r}{c})$ together with any witness $h_\varphi$.
The zero extension $\mathbf{1}_{\Ball{r}{c}}\varphi$ lies in $H^{1}_{0}(\Omega)$
and admits an explicit witness $\hat h_\varphi$ produced by
\leanname{ballIndicatorWitnessOn}, whose weak gradient is the zero
extension of $\nabla \varphi$. The pointwise nonnegativity of $\varphi$ on
$\Ball{r}{c}$ extends to a pointwise nonnegativity of the indicator, so
$\mathbf{1}_{\Ball{r}{c}}\varphi$ is a valid nonnegative test on $\Omega$.

Apply the subsolution hypothesis on $\Omega$ to this extended test:
$\mathfrak{a}_{A}(h_u^{\Omega}, \hat h_\varphi) \le 0$. By
\leanname{bilinForm\_ball\_restrict\_eq\_zeroExtend}, the bilinear form for
the restricted coefficient field $A\restriction\Ball{r}{c}$ paired with
$h_u^{\Ball{r}{c}}$ and $h_\varphi$ equals the bilinear form for $A$ paired
with $h_u^{\Omega}$ and $\hat h_\varphi$. Therefore
$\mathfrak{a}_{A\restriction\Ball{r}{c}}(h_u^{\Ball{r}{c}}, h_\varphi) \le 0$,
which is the desired subsolution inequality on the smaller domain.
\end{proof}

\begin{theorem}[\leanname{IsSubsolution.sub\_const\_ball}, \leanname{IsSupersolution.sub\_const\_ball}, \leanname{IsSolution.sub\_const\_ball}]
For $u$ a weak (sub-, super-)solution on $\Ball{r}{c}$ and $k \in \R$, $u - k$ is a weak (sub-, super-)solution on $\Ball{r}{c}$ with the same coefficient field $A$.
\end{theorem}

\begin{lstlisting}[style=Matlab-editor,basicstyle=\mlttfamily,escapechar=",mlshowsectionrules=true,mathescape=true,morecomment={[s]{\%\{}{\%\}}},language=lean,numbers=none]
theorem IsSubsolution.sub_const_ball
    {c : E} {r : ℝ} (_hr : 0 < r)
    {A : EllipticCoeff d (Metric.ball c r)} {u : E → ℝ}
    (hsub : IsSubsolution A u) (k : ℝ) :
    IsSubsolution A (fun x => u x - k) := by
  haveI : IsFiniteMeasure (volume.restrict (Metric.ball c r)) := by
    rw [MeasureTheory.isFiniteMeasure_iff]
    simpa using (measure_ball_lt_top (μ := volume) (x := c) (r := r))
  let hu : MemW1pWitness 2 u (Metric.ball c r) := MemW1p.someWitness hsub.1
  let huShift : MemW1pWitness 2 (fun x => u x - k) (Metric.ball c r) :=
    hu.sub_const Metric.isOpen_ball k
  refine ⟨huShift.memW1p, ?_⟩
  intro hv φ hφ0 hφ hφ_nonneg
  let huBackRaw : MemW1pWitness 2 (fun x => (u x - k) - (-k)) (Metric.ball c r) :=
    hv.sub_const Metric.isOpen_ball (-k)
  let huBack : MemW1pWitness 2 u (Metric.ball c r) :=
    MemW1pWitness.of_ae_eq (Ω := Metric.ball c r)
      (Filter.Eventually.of_forall fun x => by ring) huBackRaw
  have hineq : bilinFormOfCoeff A huBack hφ ≤ 0 := hsub.2 huBack φ hφ0 hφ hφ_nonneg
  simpa [huBack, huBackRaw, MemW1pWitness.of_ae_eq, MemW1pWitness.sub_const,
    sub_eq_add_neg, add_assoc] using hineq

theorem IsSupersolution.sub_const_ball
    {c : E} {r : ℝ} (_hr : 0 < r)
    {A : EllipticCoeff d (Metric.ball c r)} {u : E → ℝ}
    (hsuper : IsSupersolution A u) (k : ℝ) :
    IsSupersolution A (fun x => u x - k) := by
  haveI : IsFiniteMeasure (volume.restrict (Metric.ball c r)) := by
    rw [MeasureTheory.isFiniteMeasure_iff]
    simpa using (measure_ball_lt_top (μ := volume) (x := c) (r := r))
  let hu : MemW1pWitness 2 u (Metric.ball c r) := MemW1p.someWitness hsuper.1
  let huShift : MemW1pWitness 2 (fun x => u x - k) (Metric.ball c r) :=
    hu.sub_const Metric.isOpen_ball k
  refine ⟨huShift.memW1p, ?_⟩
  intro hv φ hφ0 hφ hφ_nonneg
  let huBackRaw : MemW1pWitness 2 (fun x => (u x - k) - (-k)) (Metric.ball c r) :=
    hv.sub_const Metric.isOpen_ball (-k)
  let huBack : MemW1pWitness 2 u (Metric.ball c r) :=
    MemW1pWitness.of_ae_eq (Ω := Metric.ball c r)
      (Filter.Eventually.of_forall fun x => by ring) huBackRaw
  have hineq : 0 ≤ bilinFormOfCoeff A huBack hφ := hsuper.2 huBack φ hφ0 hφ hφ_nonneg
  simpa [huBack, huBackRaw, MemW1pWitness.of_ae_eq, MemW1pWitness.sub_const,
    sub_eq_add_neg, add_assoc] using hineq

theorem IsSolution.sub_const_ball
    {c : E} {r : ℝ} (hr : 0 < r)
    {A : EllipticCoeff d (Metric.ball c r)} {u : E → ℝ}
    (hsol : IsSolution A u) (k : ℝ) :
    IsSolution A (fun x => u x - k) := by
  exact ⟨hsol.1.sub_const_ball (d := d) hr k, hsol.2.sub_const_ball (d := d) hr k⟩

\end{lstlisting}
\begin{proof}
The constant function $k$ has zero weak gradient on the ball (a finite-measure set), so the witness for $u - k$ has the same weak gradient as the witness for $u$. The bilinear form $\mathfrak{a}_A$ depends only on the gradient, so the inequality is unchanged.
\end{proof}

\begin{theorem}[\leanname{IsSubsolution.neg\_ball}, \leanname{IsSupersolution.neg\_ball}, \leanname{IsSolution.neg\_ball}]
On $\Ball{r}{c}$, negation swaps subsolutions and supersolutions: if $u$ is a subsolution, then $-u$ is a supersolution; similarly with the roles reversed; and a solution stays a solution under negation.
\end{theorem}

\begin{lstlisting}[style=Matlab-editor,basicstyle=\mlttfamily,escapechar=",mlshowsectionrules=true,mathescape=true,morecomment={[s]{\%\{}{\%\}}},language=lean,numbers=none]
theorem IsSubsolution.neg_ball
    {c : E} {r : ℝ} (_hr : 0 < r)
    {A : EllipticCoeff d (Metric.ball c r)} {u : E → ℝ}
    (hsub : IsSubsolution A u) :
    IsSupersolution A (fun x => -u x) := by
  let hu : MemW1pWitness 2 u (Metric.ball c r) := MemW1p.someWitness hsub.1
  let hneg : MemW1pWitness 2 (fun x => -u x) (Metric.ball c r) := by
    simpa using hu.smul (-1)
  refine ⟨hneg.memW1p, ?_⟩
  intro hv φ hφ0 hφ hφ_nonneg
  let huBackRaw : MemW1pWitness 2 (fun x => (-1 : ℝ) * (-u x)) (Metric.ball c r) := hv.smul (-1)
  let huBack : MemW1pWitness 2 u (Metric.ball c r) :=
    MemW1pWitness.of_ae_eq (Ω := Metric.ball c r)
      (Filter.Eventually.of_forall fun x => by ring) huBackRaw
  have hineq : bilinFormOfCoeff A huBack hφ ≤ 0 := hsub.2 huBack φ hφ0 hφ hφ_nonneg
  have hsmul : bilinFormOfCoeff A huBack hφ = -bilinFormOfCoeff A hv hφ := by
    simpa [huBack, huBackRaw, MemW1pWitness.of_ae_eq] using
      (bilinFormOfCoeff_smul_left (-1) A hv hφ)
  rw [hsmul] at hineq
  linarith

theorem IsSupersolution.neg_ball
    {c : E} {r : ℝ} (_hr : 0 < r)
    {A : EllipticCoeff d (Metric.ball c r)} {u : E → ℝ}
    (hsuper : IsSupersolution A u) :
    IsSubsolution A (fun x => -u x) := by
  let hu : MemW1pWitness 2 u (Metric.ball c r) := MemW1p.someWitness hsuper.1
  let hneg : MemW1pWitness 2 (fun x => -u x) (Metric.ball c r) := by
    simpa using hu.smul (-1)
  refine ⟨hneg.memW1p, ?_⟩
  intro hv φ hφ0 hφ hφ_nonneg
  let huBackRaw : MemW1pWitness 2 (fun x => (-1 : ℝ) * (-u x)) (Metric.ball c r) := hv.smul (-1)
  let huBack : MemW1pWitness 2 u (Metric.ball c r) :=
    MemW1pWitness.of_ae_eq (Ω := Metric.ball c r)
      (Filter.Eventually.of_forall fun x => by ring) huBackRaw
  have hineq : 0 ≤ bilinFormOfCoeff A huBack hφ := hsuper.2 huBack φ hφ0 hφ hφ_nonneg
  have hsmul : bilinFormOfCoeff A huBack hφ = -bilinFormOfCoeff A hv hφ := by
    simpa [huBack, huBackRaw, MemW1pWitness.of_ae_eq] using
      (bilinFormOfCoeff_smul_left (-1) A hv hφ)
  rw [hsmul] at hineq
  linarith

theorem IsSolution.neg_ball
    {c : E} {r : ℝ} (hr : 0 < r)
    {A : EllipticCoeff d (Metric.ball c r)} {u : E → ℝ}
    (hsol : IsSolution A u) :
    IsSolution A (fun x => -u x) := by
  exact ⟨hsol.2.neg_ball (d := d) hr, hsol.1.neg_ball (d := d) hr⟩

\end{lstlisting}
\begin{proof}
Suppose $u$ is a subsolution: pick a witness $h_u$ for $u$ on $\Ball{r}{c}$
via \leanname{MemW1p.someWitness}. The scaled witness
$h_u.\mathrm{smul}(-1)$ is a witness for $-u$ with weak gradient $-\nabla u$
(\leanname{MemW1pWitness.smul}), establishing the $W^{1,2}$ membership half
of \leanname{IsSupersolution}.

For the bilinear-form inequality, fix a nonnegative test $\varphi$ together
with a witness $h_\varphi$. The negation witness $h_{-(-u)}$ obtained by
applying $\mathrm{smul}(-1)$ a second time agrees a.e.\ with $h_u$
(after clearing the double negation by ring arithmetic), so by
\leanname{MemW1pWitness.of\_ae\_eq} it produces the same bilinear-form
value as $h_u$. The subsolution hypothesis applied to $u$ then gives
$\mathfrak{a}_A(h_u, h_\varphi) \le 0$.

By the linearity lemma \leanname{bilinFormOfCoeff\_smul\_left} applied with
scalar $-1$, we have
\[
  \mathfrak{a}_A(h_u, h_\varphi)
  = -\mathfrak{a}_A(h_{-u}, h_\varphi),
\]
so the inequality $\mathfrak{a}_A(h_u, h_\varphi) \le 0$ rewrites as
$\mathfrak{a}_A(h_{-u}, h_\varphi) \ge 0$, which is the supersolution
inequality for $-u$. The reverse direction (supersolution to subsolution) is
identical with the roles swapped; the solution case follows by combining
both.
\end{proof}


%% ========================================================================
\section{De Giorgi Iteration and Local Boundedness}\label{sec:degiorgi}
%% ========================================================================

% !TEX root = DeGiorgi_Lean_to_Tex.tex

This section is the formal core of the De Giorgi $L^2 \to L^\infty$
argument for non-negative subsolutions of a divergence-form linear
elliptic equation
$\dv(A \nabla u) \le 0$
on the unit ball $\Bz{1} \subset \R^d$, with $A$ measurable, symmetric,
and uniformly elliptic with constants $0 < \lambda \le \Lambda$ (here
$\lambda$ denotes always the global ellipticity constant, never a level
threshold).  The Lean development is organised into six layers, each
realised by a dedicated module:\\
\leanname{DeGiorgiIteration.CutoffAdmissibility},
\leanname{DeGiorgiIteration.Energy},\\
\leanname{DeGiorgiIteration.PreIteration},
\leanname{DeGiorgiIteration.Recurrence}, and
\leanname{DeGiorgiIteration.Linfty},
all reassembled into the public interface
\leanname{DeGiorgiIteration} and reexported by
\leanname{DeGiorgiTheory}.  We follow the same layering below.

Throughout the section $E = \R^d$ with $d \ge 1$, $\Omega \subset E$ is
an open set with finite Lebesgue measure, and $A$ is an elliptic
coefficient field with constants $\lambda$, $\Lambda$ encoded in
\leanname{EllipticCoeff}; we write $\mathrm{ellipticityRatio}(A) =
\Lambda/\lambda$.

\subsection{Positive-part truncation}

\begin{definition}[\leanname{positivePartSub}]
For any function $u : \alpha \to \R$ and any level $\ell \in \R$ we
define the positive-part truncation
\[
  (u-\ell)_+(x) := \max(u(x) - \ell, 0).
\]
\end{definition}

\begin{lstlisting}[style=Matlab-editor,basicstyle=\mlttfamily,escapechar=",mlshowsectionrules=true,mathescape=true,morecomment={[s]{\%\{}{\%\}}},language=lean,numbers=none]
/-- Positive-part truncation `(u-k)₊`. -/
def positivePartSub {α : Type*} (u : α → ℝ) (k : ℝ) : α → ℝ :=
  fun x => max (u x - k) 0

\end{lstlisting}
\begin{lemma}[\leanname{positivePartSub\_nonneg}]
$0 \le (u-\ell)_+(x)$ for every $x$.
\end{lemma}

\begin{lstlisting}[style=Matlab-editor,basicstyle=\mlttfamily,escapechar=",mlshowsectionrules=true,mathescape=true,morecomment={[s]{\%\{}{\%\}}},language=lean,numbers=none]
lemma positivePartSub_nonneg {α : Type*} {u : α → ℝ} {k : ℝ} {x : α} :
    0 ≤ positivePartSub u k x := by
  simp [positivePartSub]

\end{lstlisting}
\begin{proof}
Immediate from $\max(\,\cdot\,, 0) \ge 0$.
\end{proof}

\begin{lemma}[\leanname{positivePartSub\_sq\_ge\_gap\_sq\_of\_lt}]
If $\theta < \ell$ and $\ell < u(x)$, then
$(\ell - \theta)^2 \le |(u-\theta)_+(x)|^2$.
\end{lemma}

\begin{lstlisting}[style=Matlab-editor,basicstyle=\mlttfamily,escapechar=",mlshowsectionrules=true,mathescape=true,morecomment={[s]{\%\{}{\%\}}},language=lean,numbers=none]
lemma positivePartSub_sq_ge_gap_sq_of_lt
    {α : Type*} {u : α → ℝ} {θ lam : ℝ} {x : α}
    (hθl : θ < lam) (hxl : lam < u x) :
    (lam - θ) ^ 2 ≤ |positivePartSub u θ x| ^ 2 := by
  have hgap_pos : 0 < lam - θ := sub_pos.mpr hθl
  have hpos : 0 < u x - θ := by linarith
  have hle : lam - θ ≤ u x - θ := by linarith
  have hpp : positivePartSub u θ x = u x - θ := by
    simp [positivePartSub, max_eq_left (le_of_lt hpos)]
  rw [hpp, abs_of_nonneg (le_of_lt hpos)]
  nlinarith

\end{lstlisting}
\begin{proof}
Under the hypotheses $u(x) - \theta > 0$, hence
$(u-\theta)_+(x) = u(x) - \theta \ge \ell - \theta > 0$.  Squaring this
nonnegative inequality yields the claim.
\end{proof}

\begin{lemma}[\leanname{positivePartSub\_antitone}]
If $\ell_1 \le \ell_2$, then
$(u-\ell_2)_+ \le (u-\ell_1)_+$ pointwise.
\end{lemma}

\begin{lstlisting}[style=Matlab-editor,basicstyle=\mlttfamily,escapechar=",mlshowsectionrules=true,mathescape=true,morecomment={[s]{\%\{}{\%\}}},language=lean,numbers=none]
omit [NeZero d] in
lemma positivePartSub_antitone
    {u : E → ℝ} {k₁ k₂ : ℝ} (hk : k₁ ≤ k₂) (x : E) :
    positivePartSub u k₂ x ≤ positivePartSub u k₁ x := by
  dsimp [positivePartSub]
  exact max_le_max (by linarith) le_rfl

\end{lstlisting}
\begin{proof}
$\max(u-\ell_2, 0) \le \max(u-\ell_1, 0)$ since $u-\ell_2 \le u-\ell_1$.
\end{proof}

\begin{lemma}[\leanname{positivePartSub\_le\_posPartZero}]
If $0 \le \ell$, then $(u-\ell)_+ \le u_+ = \max(u,0)$.
\end{lemma}

\begin{lstlisting}[style=Matlab-editor,basicstyle=\mlttfamily,escapechar=",mlshowsectionrules=true,mathescape=true,morecomment={[s]{\%\{}{\%\}}},language=lean,numbers=none]
omit [NeZero d] in
lemma positivePartSub_le_posPartZero
    {u : E → ℝ} {k : ℝ} (hk : 0 ≤ k) (x : E) :
    positivePartSub u k x ≤ max (u x) 0 := by
  simpa [positivePartSub] using positivePartSub_antitone (u := u) hk x

\end{lstlisting}
\begin{lemma}[\leanname{le\_of\_positivePartSub\_eq\_zero}]
If $(u-\ell)_+(x) = 0$, then $u(x) \le \ell$.
\end{lemma}

\begin{lstlisting}[style=Matlab-editor,basicstyle=\mlttfamily,escapechar=",mlshowsectionrules=true,mathescape=true,morecomment={[s]{\%\{}{\%\}}},language=lean,numbers=none]
omit [NeZero d] in
lemma le_of_positivePartSub_eq_zero
    {u : E → ℝ} {k : ℝ} {x : E}
    (hzero : positivePartSub u k x = 0) :
    u x ≤ k := by
  by_cases hux : 0 ≤ u x - k
  · have hsub : u x - k = 0 := by
      simpa [positivePartSub, max_eq_left hux] using hzero
    linarith
  · linarith

\end{lstlisting}
\begin{proof}
If $u(x) - \ell \ge 0$ then $(u-\ell)_+(x) = u(x) - \ell$, so the
hypothesis forces $u(x) = \ell$.  Otherwise $u(x) < \ell$.
\end{proof}

\subsection{Cutoff admissibility}

The De Giorgi test functions are obtained by multiplying the
positive-part truncation by the square of a smooth scalar cutoff.

\begin{definition}[\leanname{deGiorgiCutoffTestGeneral}]
Given $\eta, u : E \to \R$ and a level $\ell$,
\[
  \varphi_{\eta, u, \ell}(x) := \eta(x) \cdot \bigl( \eta(x) \cdot
  (u-\ell)_+(x) \bigr) = \eta(x)^2 \, (u-\ell)_+(x).
\]
\end{definition}

\begin{lstlisting}[style=Matlab-editor,basicstyle=\mlttfamily,escapechar=",mlshowsectionrules=true,mathescape=true,morecomment={[s]{\%\{}{\%\}}},language=lean,numbers=none]
/-- Generalized De Giorgi cutoff test `η² (u-k)₊` for an arbitrary scalar
cutoff `η`. -/
def deGiorgiCutoffTestGeneral
    (η u : E → ℝ) (k : ℝ) : E → ℝ :=
  fun x => η x * (η x * positivePartSub u k x)

\end{lstlisting}
\begin{lemma}[\leanname{deGiorgiCutoffTestGeneral\_nonneg}]
If $\eta \ge 0$ pointwise, then $\varphi_{\eta, u, \ell} \ge 0$.
\end{lemma}

\begin{lstlisting}[style=Matlab-editor,basicstyle=\mlttfamily,escapechar=",mlshowsectionrules=true,mathescape=true,morecomment={[s]{\%\{}{\%\}}},language=lean,numbers=none]
omit [NeZero d] in
lemma deGiorgiCutoffTestGeneral_nonneg
    {η u : E → ℝ} {k : ℝ}
    (hη_nonneg : ∀ x, 0 ≤ η x) :
    ∀ x, 0 ≤ deGiorgiCutoffTestGeneral η u k x := by
  intro x
  have hpp : 0 ≤ positivePartSub u k x := positivePartSub_nonneg
  exact mul_nonneg (hη_nonneg x) (mul_nonneg (hη_nonneg x) hpp)

\end{lstlisting}
\begin{lemma}[\leanname{deGiorgiCutoffTestGeneral\_hasCompactSupport}]
If $\eta$ has compact support, so does $\varphi_{\eta,u,\ell}$.
\end{lemma}

\begin{lstlisting}[style=Matlab-editor,basicstyle=\mlttfamily,escapechar=",mlshowsectionrules=true,mathescape=true,morecomment={[s]{\%\{}{\%\}}},language=lean,numbers=none]
omit [NeZero d] in
lemma deGiorgiCutoffTestGeneral_hasCompactSupport
    {η u : E → ℝ} {k : ℝ}
    (hη_comp : HasCompactSupport η) :
    HasCompactSupport (deGiorgiCutoffTestGeneral η u k) := by
  simpa [deGiorgiCutoffTestGeneral, smul_eq_mul] using
    hη_comp.smul_right (f' := fun x => η x * positivePartSub u k x)

\end{lstlisting}
\begin{lemma}[\leanname{deGiorgiCutoffTestGeneral\_tsupport\_subset}]
If $\tsupp\,\eta \subseteq \Omega$, then $\tsupp\,\varphi_{\eta,u,\ell}
\subseteq \Omega$.
\end{lemma}

\begin{lstlisting}[style=Matlab-editor,basicstyle=\mlttfamily,escapechar=",mlshowsectionrules=true,mathescape=true,morecomment={[s]{\%\{}{\%\}}},language=lean,numbers=none]
omit [NeZero d] in
lemma deGiorgiCutoffTestGeneral_tsupport_subset
    {Ω : Set E} {η u : E → ℝ} {k : ℝ}
    (hη_sub : tsupport η ⊆ Ω) :
    tsupport (deGiorgiCutoffTestGeneral η u k) ⊆ Ω := by
  simpa [deGiorgiCutoffTestGeneral, smul_eq_mul] using
    (tsupport_smul_subset_left η (fun x => η x * positivePartSub u k x)).trans hη_sub

\end{lstlisting}
The technical core of this layer is a Sobolev membership statement: if
$u \in W^{1,2}(\Omega)$ and the positive-part truncation
$(u - \ell)_+$ is also in $W^{1,2}$, then for any smooth bounded
cutoff $\eta$ with $\tsupp\,\eta \subseteq \Omega$, the test function
$\varphi_{\eta,u,\ell}$ lies in $W^{1,2}_0(\Omega)$.  In Lean this is
phrased through the witness type \leanname{MemW1pWitness} and the
zero-trace structure \leanname{MemW01p}.

\begin{theorem}[\leanname{deGiorgiCutoffTest\_memW01p\_of\_truncWitness}]
Let $\Omega \subset E$ be open with finite measure.  Let $u$ be such
that $(u-\ell)_+$ admits a $W^{1,2}(\Omega)$ witness.  Let $\eta \in
C^\infty(E)$ satisfy $|\eta| \le C_0$, $\|\nabla \eta\| \le C_1$,
$\eta$ has compact support, and $\tsupp\,\eta \subseteq \Omega$.
Then $\varphi_{\eta,u,\ell} \in \Wop{2}(\Omega)$.
\end{theorem}

\begin{lstlisting}[style=Matlab-editor,basicstyle=\mlttfamily,escapechar=",mlshowsectionrules=true,mathescape=true,morecomment={[s]{\%\{}{\%\}}},language=lean,numbers=none]
theorem deGiorgiCutoffTest_memW01p_of_truncWitness
    {Ω : Set E} [IsFiniteMeasure (volume.restrict Ω)]
    (hΩ : IsOpen Ω)
    {u η : E → ℝ} {k : ℝ}
    (hw_trunc : MemW1pWitness 2 (positivePartSub u k) Ω)
    (hη : ContDiff ℝ (⊤ : ℕ∞) η)
    {C₀ C₁ : ℝ}
    (hC₀ : 0 ≤ C₀) (hC₁ : 0 ≤ C₁)
    (hη_bound : ∀ x, |η x| ≤ C₀)
    (hη_grad_bound : ∀ x, ‖fderiv ℝ η x‖ ≤ C₁)
    (hη_comp : HasCompactSupport η)
    (hη_sub : tsupport η ⊆ Ω) :
    MemW01p 2 (deGiorgiCutoffTestGeneral η u k) Ω := by
  let hwη :=
    hw_trunc.mul_smooth_bounded hΩ hη
      (C₀ := C₀) (C₁ := C₁) hC₀ hC₁ hη_bound hη_grad_bound
  let hwηθ : MemW1pWitness 2 (deGiorgiCutoffTestGeneral η u k) Ω :=
    hwη.mul_smooth_bounded hΩ hη
      (C₀ := C₀) (C₁ := C₁) hC₀ hC₁ hη_bound hη_grad_bound
  have hv_comp : HasCompactSupport (deGiorgiCutoffTestGeneral η u k) :=
    deGiorgiCutoffTestGeneral_hasCompactSupport hη_comp
  have hv_sub : tsupport (deGiorgiCutoffTestGeneral η u k) ⊆ Ω :=
    deGiorgiCutoffTestGeneral_tsupport_subset hη_sub
  have hwηθ_memW1p_real :
      MemW1p (ENNReal.ofReal 2) (deGiorgiCutoffTestGeneral η u k) Ω := by
    simpa using hwηθ.memW1p
  simpa using
    (memW01p_of_memW1p_of_tsupport_subset
      hΩ (by norm_num : (1 : ℝ) < 2) hwηθ_memW1p_real hv_comp hv_sub)

\end{lstlisting}
\begin{proof}
By two successive applications of the closure of $W^{1,2}$-witnesses
under multiplication by smooth bounded functions
(\leanname{MemW1pWitness.mul\_smooth\_bounded}), starting from $(u-\ell)_+$
and multiplying twice by $\eta$, one obtains a $W^{1,2}$ witness for
$\eta^2 (u-\ell)_+$.  Since $\varphi_{\eta,u,\ell}$ is compactly
supported with $\tsupp \subseteq \Omega$, the localisation lemma
\leanname{memW01p\_of\_memW1p\_of\_tsupport\_subset} promotes the
$W^{1,2}$-witness to a $W^{1,2}_0(\Omega)$ membership.
\end{proof}

\begin{theorem}[\leanname{deGiorgiCutoffTest\_memW01p\_of\_posPart}]
Same conclusion, but assuming a $W^{1,2}$-witness for $u$ together
with a generic positive-part closure rule
$\bigl(\, v \in W^{1,2}(\Omega) \implies (v)_+ \in W^{1,2}(\Omega)\,\bigr)$.
\end{theorem}

\begin{lstlisting}[style=Matlab-editor,basicstyle=\mlttfamily,escapechar=",mlshowsectionrules=true,mathescape=true,morecomment={[s]{\%\{}{\%\}}},language=lean,numbers=none]
/-- General cutoff admissibility theorem, parameterized by a positive-part
witness theorem. This packages the reusable cutoff argument while keeping the
localization argument in
`memW01p_of_memW1p_of_tsupport_subset`. -/
theorem deGiorgiCutoffTest_memW01p_of_posPart
    {Ω : Set E} [IsFiniteMeasure (volume.restrict Ω)]
    (hΩ : IsOpen Ω)
    {u η : E → ℝ}
    (hw : MemW1pWitness 2 u Ω)
    (hη : ContDiff ℝ (⊤ : ℕ∞) η)
    {C₀ C₁ : ℝ}
    (hC₀ : 0 ≤ C₀) (hC₁ : 0 ≤ C₁)
    (hη_bound : ∀ x, |η x| ≤ C₀)
    (hη_grad_bound : ∀ x, ‖fderiv ℝ η x‖ ≤ C₁)
    (hη_comp : HasCompactSupport η)
    (hη_sub : tsupport η ⊆ Ω)
    (k : ℝ)
    (hpos :
      ∀ {v : E → ℝ},
        MemW1pWitness 2 v Ω →
          MemW1pWitness 2 (fun x => max (v x) 0) Ω) :
    MemW01p 2 (deGiorgiCutoffTestGeneral η u k) Ω := by
  let hwShift : MemW1pWitness 2 (fun x => u x - k) Ω := hw.sub_const hΩ k
  have hw_trunc : MemW1pWitness 2 (positivePartSub u k) Ω := by
    simpa [positivePartSub] using (hpos hwShift)
  exact
    deGiorgiCutoffTest_memW01p_of_truncWitness
      hΩ hw_trunc hη hC₀ hC₁ hη_bound hη_grad_bound hη_comp hη_sub

\end{lstlisting}
\begin{proof}
Apply the closure rule to $v = u-\ell$ to obtain a witness for
$(u-\ell)_+$, then invoke the previous theorem.
\end{proof}

A concrete version, \leanname{deGiorgiCutoffTest\_memW01p\_of\_posPartApprox},
is obtained by furnishing the $W^{1,2}$-witness for $(u-\ell)_+$ via the
Chapter~02 smooth approximation package for $u-\ell$.  The
ball-restricted variants
\leanname{deGiorgiCutoffTest\_memW01p\_on\_ball\_of\_ballPosPart} and
\leanname{deGiorgiCutoffTest\_memW01p\_on\_ball\_of\_ballPosPartApprox}
specialise these statements to $\Omega = \Ball{s}{x_0}$ and supply
finiteness of the restricted measure automatically.

\begin{definition}[\leanname{positivePartSub\_memW1pWitness\_on\_ball}]
On a ball $\Ball{s}{x_0}$, given a $W^{1,2}$-witness for $u$ and a
smooth approximation package for $u - \ell$, the lemma constructs a
$W^{1,2}$-witness for $(u - \ell)_+$ on the ball.
\end{definition}

\begin{lstlisting}[style=Matlab-editor,basicstyle=\mlttfamily,escapechar=",mlshowsectionrules=true,mathescape=true,morecomment={[s]{\%\{}{\%\}}},language=lean,numbers=none]
/-- Concrete truncation witness on a ball, obtained from the Chapter 02
positive-part constructor applied to `u - k`. -/
noncomputable def positivePartSub_memW1pWitness_on_ball
    {u : E → ℝ} {x₀ : E} {s : ℝ}
    (_hs : 0 < s)
    (hw : MemW1pWitness 2 u (Metric.ball x₀ s))
    (k : ℝ)
    (happroxBallShift :
      ∃ ψ : ℕ → E → ℝ,
        (∀ n, ContDiff ℝ 1 (ψ n)) ∧
        (∀ n, HasCompactSupport (ψ n)) ∧
        Tendsto
          (fun n =>
            eLpNorm (fun x => ψ n x - (u x - k)) 2
              (volume.restrict (Metric.ball x₀ s)))
          atTop (nhds 0) ∧
        (∀ i : Fin d,
          Tendsto
            (fun n =>
              eLpNorm
                (fun x =>
                  (fderiv ℝ (ψ n) x) (EuclideanSpace.single i 1) - hw.weakGrad x i)
                2 (volume.restrict (Metric.ball x₀ s)))
            atTop (nhds 0))) :
    MemW1pWitness 2 (positivePartSub u k) (Metric.ball x₀ s) := by
  have hΩ : IsOpen (Metric.ball x₀ s) := isOpen_ball
  haveI : IsFiniteMeasure (volume.restrict (Metric.ball x₀ s)) := by
    rw [isFiniteMeasure_restrict]
    exact measure_ball_lt_top.ne
  exact positivePartSub_memW1pWitness hΩ hw k happroxBallShift

\end{lstlisting}
\subsection{Weighted Caccioppoli inequality}

We now turn to the analytical heart of the chapter: the weighted
Caccioppoli inequality for truncations of subsolutions.  We use the
abbreviation
\[
  \nabla \eta(x) := (\partial_1 \eta(x), \dots, \partial_d \eta(x))
\]
identified with the inner-product representation of the Fréchet
derivative; see \leanname{deGiorgiFderivVec}.  The norm
$\|\nabla \eta(x)\|$ then equals the operator norm
$\|D\eta(x)\|$, which is recorded as
\leanname{deGiorgi\_norm\_fderivVec\_eq\_norm\_fderiv}.

\begin{lemma}[\leanname{weighted\_caccioppoli\_pointwise} (Young's inequality)]
Let $\Lambda > 0$, $\eta, \zeta, E, M, v \in \R$, and assume the
pointwise coercivity condition $|M|^2 \le \Lambda E$. Then
\begin{align*}
  2 \eta |v| \zeta |M|
   \le \tfrac{1}{2}\, \eta^2 E + 2\Lambda\, \zeta^2 |v|^2.
\end{align*}
\end{lemma}

\begin{lstlisting}[style=Matlab-editor,basicstyle=\mlttfamily,escapechar=",mlshowsectionrules=true,mathescape=true,morecomment={[s]{\%\{}{\%\}}},language=lean,numbers=none]
private lemma weighted_caccioppoli_pointwise
    {Λ η ζ E M v : ℝ}
    (hΛ : 0 < Λ)
    (hcoeff : |M| ^ 2 ≤ Λ * E) :
    2 * η * |v| * ζ * |M| ≤
      (1 / 2 : ℝ) * η ^ 2 * E + 2 * Λ * ζ ^ 2 * |v| ^ 2 := by
  let s := Real.sqrt (2 * Λ)
  have hs_pos : 0 < s := by
    dsimp [s]
    apply Real.sqrt_pos.2
    positivity
  have hs_ne : s ≠ 0 := ne_of_gt hs_pos
  have hs_sq : s ^ 2 = 2 * Λ := by
    dsimp [s]
    rw [Real.sq_sqrt]
    positivity
  have hyoung :
      2 * (η * |M| / s) * (s * ζ * |v|) ≤
        (η * |M| / s) ^ 2 + (s * ζ * |v|) ^ 2 := by
    nlinarith [sq_nonneg ((η * |M| / s) - (s * ζ * |v|))]
  have hcoeff' : η ^ 2 * |M| ^ 2 / (2 * Λ) ≤ (1 / 2 : ℝ) * η ^ 2 * E := by
    have hmul :
        (η ^ 2 / (2 * Λ)) * (|M| ^ 2) ≤ (η ^ 2 / (2 * Λ)) * (Λ * E) := by
      refine mul_le_mul_of_nonneg_left hcoeff ?_
      positivity
    have hfac : (η ^ 2 / (2 * Λ)) * (Λ * E) = (1 / 2 : ℝ) * η ^ 2 * E := by
      field_simp [hΛ.ne']
    simpa [div_eq_mul_inv, mul_assoc, mul_left_comm, mul_comm] using hmul.trans_eq hfac
  have hmain :
      2 * η * |v| * ζ * |M| ≤ η ^ 2 * |M| ^ 2 / (2 * Λ) + 2 * Λ * ζ ^ 2 * |v| ^ 2 := by
    have hconv :
        2 * (η * |M| / s) * (s * ζ * |v|) =
          2 * η * |v| * ζ * |M| := by
      field_simp [hs_ne]
    have hsq1 :
        (η * |M| / s) ^ 2 = η ^ 2 * |M| ^ 2 / (2 * Λ) := by
      rw [div_pow, hs_sq]
      field_simp [hΛ.ne']
    have hsq2 :
        (s * ζ * |v|) ^ 2 = 2 * Λ * ζ ^ 2 * |v| ^ 2 := by
      rw [mul_pow, mul_pow, hs_sq]
    rwa [hconv, hsq1, hsq2] at hyoung
  have hsum :
      η ^ 2 * |M| ^ 2 / (2 * Λ) + 2 * Λ * ζ ^ 2 * |v| ^ 2 ≤
        (1 / 2 : ℝ) * η ^ 2 * E + 2 * Λ * ζ ^ 2 * |v| ^ 2 := by
    nlinarith [hcoeff']
  exact hmain.trans hsum

\end{lstlisting}
\begin{proof}
Apply Young's inequality $2ab \le a^2 + b^2$ with $a =
\eta |M| / \sqrt{2\Lambda}$ and $b = \sqrt{2\Lambda}\, \zeta |v|$ to
get
\[
  2 \eta |v| \zeta |M|
   \le \frac{\eta^2 |M|^2}{2\Lambda} + 2\Lambda\, \zeta^2 |v|^2.
\]
The hypothesis $|M|^2 \le \Lambda E$ then yields $\eta^2 |M|^2 /
(2\Lambda) \le \tfrac{1}{2} \eta^2 E$, completing the proof.
\end{proof}

\begin{theorem}[Abstract absorption,
\leanname{weighted\_caccioppoli\_absorb}]
Let $(\alpha, \mu)$ be a measure space and let $E, M, \eta, \zeta, v
: \alpha \to \R$ be measurable functions with $\Lambda > 0$. Assume
the integrability conditions and:
\begin{enumerate}
\item $\int \eta^2 E \, d\mu \le \int 2 \eta |v| \zeta |M| \, d\mu$,
\item $|M|^2 \le \Lambda E$ almost everywhere.
\end{enumerate}
Then $\int \eta^2 E \, d\mu \le 4 \Lambda \int \zeta^2 |v|^2 \, d\mu$.
\end{theorem}

\begin{lstlisting}[style=Matlab-editor,basicstyle=\mlttfamily,escapechar=",mlshowsectionrules=true,mathescape=true,morecomment={[s]{\%\{}{\%\}}},language=lean,numbers=none]
/-- Abstract absorption step for weighted Caccioppoli. -/
theorem weighted_caccioppoli_absorb
    {α : Type*} [MeasurableSpace α] {μ : Measure α}
    {E M η ζ v : α → ℝ} {Λ : ℝ}
    (hΛ : 0 < Λ)
    (hcore :
      ∫ x, η x ^ 2 * E x ∂μ ≤ ∫ x, 2 * η x * |v x| * ζ x * |M x| ∂μ)
    (hcoeff : ∀ᵐ x ∂μ, |M x| ^ 2 ≤ Λ * E x)
    (hleft_int : Integrable (fun x => η x ^ 2 * E x) μ)
    (hcross_int : Integrable (fun x => 2 * η x * |v x| * ζ x * |M x|) μ)
    (hbound_int : Integrable (fun x => ζ x ^ 2 * |v x| ^ 2) μ) :
    ∫ x, η x ^ 2 * E x ∂μ ≤ 4 * Λ * ∫ x, ζ x ^ 2 * |v x| ^ 2 ∂μ := by
  have hupper_pt :
      ∀ᵐ x ∂μ,
        2 * η x * |v x| * ζ x * |M x| ≤
          (1 / 2 : ℝ) * η x ^ 2 * E x + 2 * Λ * ζ x ^ 2 * |v x| ^ 2 := by
    filter_upwards [hcoeff] with x hx
    exact weighted_caccioppoli_pointwise hΛ hx
  have hupper_int :
      Integrable (fun x => (1 / 2 : ℝ) * η x ^ 2 * E x + 2 * Λ * ζ x ^ 2 * |v x| ^ 2) μ := by
    simpa [mul_assoc, add_comm, add_left_comm, add_assoc] using
      (hleft_int.const_mul (1 / 2 : ℝ)).add (hbound_int.const_mul (2 * Λ))
  have hcross_bound :
      ∫ x, 2 * η x * |v x| * ζ x * |M x| ∂μ ≤
        ∫ x, (1 / 2 : ℝ) * η x ^ 2 * E x + 2 * Λ * ζ x ^ 2 * |v x| ^ 2 ∂μ := by
    refine integral_mono_ae hcross_int hupper_int ?_
    exact hupper_pt
  have hsplit :
      ∫ x, (1 / 2 : ℝ) * η x ^ 2 * E x + 2 * Λ * ζ x ^ 2 * |v x| ^ 2 ∂μ =
        (1 / 2 : ℝ) * ∫ x, η x ^ 2 * E x ∂μ + 2 * Λ * ∫ x, ζ x ^ 2 * |v x| ^ 2 ∂μ := by
    have hrewrite :
        (fun x => (1 / 2 : ℝ) * η x ^ 2 * E x + 2 * Λ * ζ x ^ 2 * |v x| ^ 2) =
          (fun x => ((1 / 2 : ℝ) * (η x ^ 2 * E x)) + ((2 * Λ) * (ζ x ^ 2 * |v x| ^ 2))) := by
      funext x
      ring
    rw [hrewrite, integral_add (hleft_int.const_mul (1 / 2 : ℝ)) (hbound_int.const_mul (2 * Λ))]
    rw [integral_const_mul, integral_const_mul]
  have hmain :
      ∫ x, η x ^ 2 * E x ∂μ ≤
        (1 / 2 : ℝ) * ∫ x, η x ^ 2 * E x ∂μ + 2 * Λ * ∫ x, ζ x ^ 2 * |v x| ^ 2 ∂μ := by
    calc
      ∫ x, η x ^ 2 * E x ∂μ ≤ ∫ x, 2 * η x * |v x| * ζ x * |M x| ∂μ := hcore
      _ ≤ ∫ x, (1 / 2 : ℝ) * η x ^ 2 * E x + 2 * Λ * ζ x ^ 2 * |v x| ^ 2 ∂μ := hcross_bound
      _ = (1 / 2 : ℝ) * ∫ x, η x ^ 2 * E x ∂μ + 2 * Λ * ∫ x, ζ x ^ 2 * |v x| ^ 2 ∂μ := hsplit
  linarith

\end{lstlisting}
\begin{proof}
Apply the pointwise Young inequality almost everywhere, integrate to
get
\[
  \int 2\eta|v|\zeta|M|\, d\mu
   \le \tfrac{1}{2}\int \eta^2 E\, d\mu
        + 2\Lambda \int \zeta^2 |v|^2 \, d\mu,
\]
combine with the assumed core estimate to obtain
$\int \eta^2 E\, d\mu \le \tfrac12 \int \eta^2 E\, d\mu + 2\Lambda
\int \zeta^2 |v|^2 \, d\mu$, and absorb to conclude.
\end{proof}

The PDE-facing weighted Caccioppoli inequality is the following
statement, which is the analytical core of the De Giorgi argument.

\begin{theorem}[Weighted Caccioppoli for subsolutions,
\leanname{caccioppoli\_weighted\_of\_subsolution}]
Let $\Omega \subset E$ be open with finite measure, $A$ an elliptic
coefficient field on $\Omega$, $u$ a subsolution
($\leanname{IsSubsolution}\, A\, u$), and $\eta$ a smooth bounded
nonneg function with $|\eta| \le C_0$, $\|\nabla \eta\| \le C_1$,
$\tsupp\,\eta \subseteq \Omega$.  Let $\ell \in \R$ and assume
$(u - \ell)_+ \in W^{1,2}(\Omega)$ and that $\varphi_{\eta,u,\ell}$
admits a $W^{1,2}_0(\Omega)$ structure.  Then
\begin{multline*}
  \int_\Omega \eta(x)^2 \, \bigl\| \nabla(u-\ell)_+(x) \bigr\|^2 \, dx
   \le \\
   4\, \frac{\Lambda}{\lambda}\,
     \int_\Omega \|\nabla \eta(x)\|^2 \,
       \bigl|(u-\ell)_+(x)\bigr|^2 \, dx.
\end{multline*}
\end{theorem}

\begin{lstlisting}[style=Matlab-editor,basicstyle=\mlttfamily,escapechar=",mlshowsectionrules=true,mathescape=true,morecomment={[s]{\%\{}{\%\}}},language=lean,numbers=none]
/-- PDE-facing weighted Caccioppoli inequality for the truncated subsolution
test `η² (u-k)₊`. The analytic core of the argument is isolated here for
downstream reuse. -/
theorem caccioppoli_weighted_of_subsolution
    {Ω : Set E}
    [IsFiniteMeasure (volume.restrict Ω)]
    (hΩ : IsOpen Ω)
    (A : EllipticCoeff d Ω)
    {u η : E → ℝ} {k : ℝ}
    (hsub : IsSubsolution A u)
    (hu : MemW1pWitness 2 u Ω)
    (hw_trunc : MemW1pWitness 2 (positivePartSub u k) Ω)
    (hη : ContDiff ℝ (⊤ : ℕ∞) η)
    (hη_nonneg : ∀ x, 0 ≤ η x)
    {C₀ C₁ : ℝ}
    (hC₀ : 0 ≤ C₀) (hC₁ : 0 ≤ C₁)
    (hη_bound : ∀ x, |η x| ≤ C₀)
    (hη_grad_bound : ∀ x, ‖fderiv ℝ η x‖ ≤ C₁)
    (hη_admissible : MemW01p 2 (deGiorgiCutoffTestGeneral η u k) Ω) :
    ∫ x in Ω, η x ^ 2 * ‖hw_trunc.weakGrad x‖ ^ 2 ∂volume ≤
      4 * ellipticityRatio A *
        ∫ x in Ω, ‖fderiv ℝ η x‖ ^ 2 * |positivePartSub u k x| ^ 2 ∂volume := by
  let μ : Measure E := volume.restrict Ω
  let hwφ : MemW1pWitness 2 (deGiorgiCutoffTestGeneral η u k) Ω :=
    deGiorgiCutoffTestWitnessWeighted hΩ hw_trunc hη hC₀ hC₁ hη_bound hη_grad_bound
  let Equad : E → ℝ := bilinFormIntegrandOfCoeff A hw_trunc hw_trunc
  let leftTerm : E → ℝ := fun x => η x ^ 2 * Equad x
  let gradEtaNorm : E → ℝ := fun x => ‖fderiv ℝ η x‖
  let fluxNorm : E → ℝ := fun x => ‖matMulE (A.a x) (hw_trunc.weakGrad x)‖
  let crossAbs : E → ℝ :=
    fun x => 2 * η x * positivePartSub u k x * fluxNorm x * gradEtaNorm x
  let crossInner : E → ℝ :=
    fun x =>
      (2 * η x * positivePartSub u k x) *
        inner ℝ (matMulE (A.a x) (hw_trunc.weakGrad x)) (deGiorgiFderivVec η x)
  let coreIntegrand : E → ℝ := fun x => leftTerm x + crossInner x
  have hsub_test :
      bilinFormOfCoeff A hu hwφ ≤ 0 := by
    exact hsub.2 hu (deGiorgiCutoffTestGeneral η u k) hη_admissible hwφ
      (deGiorgiCutoffTestGeneral_nonneg hη_nonneg)
  have htrunc_sq_int :
      Integrable (fun x => ‖hw_trunc.weakGrad x‖ ^ 2) μ := by
    simpa [pow_two, μ] using hw_trunc.weakGrad_norm_memLp.integrable_sq
  have hweighted_grad_sq_int :
      Integrable (fun x => η x ^ 2 * ‖hw_trunc.weakGrad x‖ ^ 2) μ := by
    refine Integrable.mono' (htrunc_sq_int.const_mul (C₀ ^ 2)) ?_ ?_
    · exact
        (((hη.continuous.pow 2).aemeasurable).mul
          htrunc_sq_int.aestronglyMeasurable.aemeasurable).aestronglyMeasurable
    · filter_upwards with x
      have hη_sq_le : η x ^ 2 ≤ C₀ ^ 2 := by
        exact sq_le_sq.mpr (by simpa [abs_of_nonneg hC₀] using hη_bound x)
      have hgrad_nonneg : 0 ≤ ‖hw_trunc.weakGrad x‖ ^ 2 := by positivity
      have hterm_nonneg : 0 ≤ η x ^ 2 * ‖hw_trunc.weakGrad x‖ ^ 2 := by positivity
      have hle :
          η x ^ 2 * ‖hw_trunc.weakGrad x‖ ^ 2 ≤
            C₀ ^ 2 * ‖hw_trunc.weakGrad x‖ ^ 2 :=
        mul_le_mul_of_nonneg_right hη_sq_le hgrad_nonneg
      simpa [Real.norm_eq_abs, abs_of_nonneg hterm_nonneg] using hle
  have hpos_sq_int :
      Integrable (fun x => |positivePartSub u k x| ^ 2) μ := by
    simpa [pow_two, μ] using hw_trunc.memLp.integrable_sq
  have hbound_int :
      Integrable (fun x => gradEtaNorm x ^ 2 * |positivePartSub u k x| ^ 2) μ := by
    refine Integrable.mono' (hpos_sq_int.const_mul (C₁ ^ 2)) ?_ ?_
    · exact
        ((((hη.continuous_fderiv (by simp : ((⊤ : ℕ∞) : WithTop ℕ∞) ≠ 0)).norm.pow 2).aemeasurable).mul
          hpos_sq_int.aestronglyMeasurable.aemeasurable).aestronglyMeasurable
    · filter_upwards with x
      have hfd_sq_le : gradEtaNorm x ^ 2 ≤ C₁ ^ 2 := by
        exact sq_le_sq.mpr (by
          simp [gradEtaNorm, abs_of_nonneg (norm_nonneg _), abs_of_nonneg hC₁, hη_grad_bound x])
      have hpos_nonneg : 0 ≤ |positivePartSub u k x| ^ 2 := by positivity
      have hterm_nonneg : 0 ≤ gradEtaNorm x ^ 2 * |positivePartSub u k x| ^ 2 := by positivity
      have hle :
          gradEtaNorm x ^ 2 * |positivePartSub u k x| ^ 2 ≤
            C₁ ^ 2 * |positivePartSub u k x| ^ 2 :=
        mul_le_mul_of_nonneg_right hfd_sq_le hpos_nonneg
      simpa [Real.norm_eq_abs, abs_of_nonneg hterm_nonneg] using hle
  have hE_nonneg :
      ∀ᵐ x ∂μ, 0 ≤ Equad x := by
    filter_upwards [A.ae_coercive_nonneg] with x hx
    simpa [Equad, bilinFormIntegrandOfCoeff, real_inner_comm] using hx (hw_trunc.weakGrad x)
  have hleft_int :
      Integrable leftTerm μ := by
    refine Integrable.mono' (hweighted_grad_sq_int.const_mul A.Λ) ?_ ?_
    · exact
        (((hη.continuous.pow 2).aemeasurable).mul
          (aestronglyMeasurable_bilinFormIntegrandOfCoeff A hw_trunc hw_trunc).aemeasurable).aestronglyMeasurable
    · filter_upwards [A.quadratic_upper, hE_nonneg] with x hx_quad hx_nonneg
      have hη_sq_nonneg : 0 ≤ η x ^ 2 := by positivity
      have hdom :
          η x ^ 2 * Equad x ≤ η x ^ 2 * (A.Λ * ‖hw_trunc.weakGrad x‖ ^ 2) :=
        mul_le_mul_of_nonneg_left
          (by simpa [Equad, bilinFormIntegrandOfCoeff, real_inner_comm] using
            hx_quad (hw_trunc.weakGrad x)) hη_sq_nonneg
      have hrhs_nonneg : 0 ≤ A.Λ * (η x ^ 2 * ‖hw_trunc.weakGrad x‖ ^ 2) := by
        exact mul_nonneg A.Λ_nonneg (by positivity)
      have hterm_nonneg : 0 ≤ leftTerm x := by
        exact mul_nonneg hη_sq_nonneg hx_nonneg
      have hnorm_rhs :
          ‖A.Λ * (η x ^ 2 * ‖hw_trunc.weakGrad x‖ ^ 2)‖ =
            A.Λ * (η x ^ 2 * ‖hw_trunc.weakGrad x‖ ^ 2) := by
        rw [Real.norm_eq_abs, abs_of_nonneg hrhs_nonneg]
      calc
        ‖leftTerm x‖ = leftTerm x := by rw [Real.norm_eq_abs, abs_of_nonneg hterm_nonneg]
        _ = η x ^ 2 * Equad x := by rfl
        _ ≤ η x ^ 2 * (A.Λ * ‖hw_trunc.weakGrad x‖ ^ 2) := hdom
        _ = A.Λ * (η x ^ 2 * ‖hw_trunc.weakGrad x‖ ^ 2) := by ring
  have hcoeff :
      ∀ᵐ x ∂μ, |fluxNorm x| ^ 2 ≤ A.Λ * Equad x := by
    filter_upwards [A.mulVec_sq_le] with x hx
    simpa [fluxNorm, Equad, bilinFormIntegrandOfCoeff, real_inner_comm] using
      hx (hw_trunc.weakGrad x)
  have hcross_upper_pt :
      ∀ᵐ x ∂μ,
        crossAbs x ≤
          (1 / 2 : ℝ) * leftTerm x + 2 * A.Λ * (gradEtaNorm x ^ 2 * |positivePartSub u k x| ^ 2) := by
    filter_upwards [hcoeff] with x hx
    have hx' :=
      weighted_caccioppoli_pointwise_core
        (Λ := A.Λ) (η := η x) (ζ := gradEtaNorm x) (Q := Equad x)
        (M := fluxNorm x) (v := positivePartSub u k x) A.Λ_pos hx
    simpa [crossAbs, leftTerm, gradEtaNorm, fluxNorm, Equad, abs_of_nonneg positivePartSub_nonneg,
      mul_assoc, mul_left_comm, mul_comm]
      using hx'
  have hcross_upper_int :
      Integrable
        (fun x =>
          (1 / 2 : ℝ) * leftTerm x + 2 * A.Λ * (gradEtaNorm x ^ 2 * |positivePartSub u k x| ^ 2)) μ := by
    simpa [mul_assoc, add_comm, add_left_comm, add_assoc] using
      (hleft_int.const_mul (1 / 2 : ℝ)).add (hbound_int.const_mul (2 * A.Λ))
  have hcrossAbs_int :
      Integrable crossAbs μ := by
    refine Integrable.mono' hcross_upper_int ?_ ?_
    · simpa [crossAbs, gradEtaNorm, fluxNorm, mul_assoc, mul_left_comm, mul_comm] using
        ((((hη.continuous.aestronglyMeasurable).mul
          hw_trunc.memLp.aestronglyMeasurable).mul
          ((hη.continuous_fderiv (by simp : ((⊤ : ℕ∞) : WithTop ℕ∞) ≠ 0)).norm.aestronglyMeasurable)).mul
          ((aestronglyMeasurable_matMulE A hw_trunc.weakGrad_memLp.aestronglyMeasurable).norm)).const_mul
          (2 : ℝ)
    · filter_upwards [hcross_upper_pt] with x hx
      have hcross_nonneg : 0 ≤ crossAbs x := by
        have hηx_nonneg : 0 ≤ η x := hη_nonneg x
        have hw_nonneg : 0 ≤ positivePartSub u k x := positivePartSub_nonneg
        have hflux_nonneg : 0 ≤ fluxNorm x := norm_nonneg _
        have hgrad_nonneg : 0 ≤ gradEtaNorm x := norm_nonneg _
        have hcoeff_nonneg : 0 ≤ 2 * η x * positivePartSub u k x := by
          nlinarith
        dsimp [crossAbs]
        exact mul_nonneg (mul_nonneg hcoeff_nonneg hflux_nonneg) hgrad_nonneg
      have hrhs_nonneg :
          0 ≤ (1 / 2 : ℝ) * leftTerm x + 2 * A.Λ * (gradEtaNorm x ^ 2 * |positivePartSub u k x| ^ 2) := by
        exact le_trans hcross_nonneg hx
      simpa [Real.norm_eq_abs, abs_of_nonneg hcross_nonneg, abs_of_nonneg hrhs_nonneg] using hx
  have hcore_eq_ae :
      coreIntegrand =ᵐ[μ] bilinFormIntegrandOfCoeff A hu hwφ := by
    have hsuper := truncGrad_eq_on_superlevel hΩ hu hw_trunc
    have hsublevel := truncGrad_zero_on_sublevel hΩ hw_trunc
    filter_upwards [hsuper, hsublevel] with x hsuperx hsublevelx
    by_cases hku : k < u x
    · have hgrad_eq : hu.weakGrad x = hw_trunc.weakGrad x := hsuperx hku
      simp [coreIntegrand, leftTerm, crossInner, Equad, hwφ, bilinFormIntegrandOfCoeff,
        deGiorgiCutoffTestWitnessWeighted_grad, hgrad_eq, inner_add_right,
        inner_smul_right, real_inner_comm, mul_assoc, mul_comm]
    · have huk : u x ≤ k := not_lt.mp hku
      have htrunc_zero : hw_trunc.weakGrad x = 0 := hsublevelx huk
      have hw_zero : positivePartSub u k x = 0 := by
        simp [positivePartSub, max_eq_right (sub_nonpos.mpr huk)]
      simp [coreIntegrand, leftTerm, crossInner, Equad, hwφ, bilinFormIntegrandOfCoeff,
        deGiorgiCutoffTestWitnessWeighted_grad, htrunc_zero, hw_zero]
  have hcore_int :
      Integrable coreIntegrand μ := by
    refine (integrable_bilinFormIntegrandOfCoeff A hu hwφ).congr ?_
    exact hcore_eq_ae.symm
  have hcrossInner_int :
      Integrable crossInner μ := by
    convert hcore_int.sub hleft_int using 1
    ext x
    simp [coreIntegrand]
  have hcrossInner_abs_le :
      ∀ᵐ x ∂μ, |crossInner x| ≤ crossAbs x := by
    filter_upwards with x
    have hηx_nonneg : 0 ≤ η x := hη_nonneg x
    have hw_nonneg : 0 ≤ positivePartSub u k x := positivePartSub_nonneg
    have hinner :
        |inner ℝ (matMulE (A.a x) (hw_trunc.weakGrad x)) (deGiorgiFderivVec η x)| ≤
          fluxNorm x * gradEtaNorm x := by
      have := abs_real_inner_le_norm (matMulE (A.a x) (hw_trunc.weakGrad x)) (deGiorgiFderivVec η x)
      have hfdv := deGiorgi_norm_fderivVec_eq_norm_fderiv (η := η) (x := x)
      change
        |inner ℝ (matMulE (A.a x) (hw_trunc.weakGrad x)) (deGiorgiFderivVec η x)| ≤
          ‖matMulE (A.a x) (hw_trunc.weakGrad x)‖ * ‖fderiv ℝ η x‖
      simpa [hfdv, mul_comm] using this
    have hcoeff_nonneg : 0 ≤ 2 * η x * positivePartSub u k x := by
      nlinarith
    change
      |2 * η x * positivePartSub u k x *
          inner ℝ (matMulE (A.a x) (hw_trunc.weakGrad x)) (deGiorgiFderivVec η x)| ≤
        2 * η x * positivePartSub u k x * fluxNorm x * gradEtaNorm x
    rw [abs_mul, abs_of_nonneg hcoeff_nonneg]
    calc
      (2 * η x * positivePartSub u k x) *
          |inner ℝ (matMulE (A.a x) (hw_trunc.weakGrad x)) (deGiorgiFderivVec η x)|
          ≤ (2 * η x * positivePartSub u k x) * (fluxNorm x * gradEtaNorm x) := by
            exact mul_le_mul_of_nonneg_left hinner hcoeff_nonneg
      _ = 2 * η x * positivePartSub u k x * fluxNorm x * gradEtaNorm x := by
            ring
  have hcross_lower :
      -∫ x, crossAbs x ∂μ ≤ ∫ x, crossInner x ∂μ := by
    have hneg :
        ∫ x, (-1 : ℝ) * crossAbs x ∂μ ≤ ∫ x, crossInner x ∂μ := by
      refine integral_mono_ae (hcrossAbs_int.const_mul (-1)) hcrossInner_int ?_
      filter_upwards [hcrossInner_abs_le] with x hx
      exact show (-1 : ℝ) * crossAbs x ≤ crossInner x by
        simpa using (abs_le.mp hx).1
    simpa [integral_neg] using hneg
  have hcore_integral :
      ∫ x, leftTerm x ∂μ ≤ ∫ x, crossAbs x ∂μ := by
    have hsub_test' :
        ∫ x, coreIntegrand x ∂μ ≤ 0 := by
      calc
        ∫ x, coreIntegrand x ∂μ = bilinFormOfCoeff A hu hwφ := by
          unfold bilinFormOfCoeff
          exact integral_congr_ae hcore_eq_ae
        _ ≤ 0 := hsub_test
    have hsplit :
        ∫ x, coreIntegrand x ∂μ = ∫ x, leftTerm x ∂μ + ∫ x, crossInner x ∂μ := by
      rw [show coreIntegrand = fun x => leftTerm x + crossInner x by
            funext x; simp [coreIntegrand]]
      exact integral_add hleft_int hcrossInner_int
    linarith
  have hleft_bound :
      ∫ x, leftTerm x ∂μ ≤ 4 * A.Λ * ∫ x, gradEtaNorm x ^ 2 * |positivePartSub u k x| ^ 2 ∂μ := by
    have hcore_absorb :
        ∫ x, η x ^ 2 * Equad x ∂μ ≤
          4 * A.Λ * ∫ x, gradEtaNorm x ^ 2 * |positivePartSub u k x| ^ 2 ∂μ := by
      have hcore' :
          ∫ x, η x ^ 2 * Equad x ∂μ ≤
            ∫ x, 2 * η x * |positivePartSub u k x| * gradEtaNorm x * |fluxNorm x| ∂μ := by
        simpa [crossAbs, fluxNorm, gradEtaNorm, abs_of_nonneg positivePartSub_nonneg,
          abs_of_nonneg (norm_nonneg _), mul_assoc, mul_left_comm, mul_comm]
          using hcore_integral
      have hcross_int' :
          Integrable (fun x => 2 * η x * |positivePartSub u k x| * gradEtaNorm x * |fluxNorm x|) μ := by
        simpa [crossAbs, fluxNorm, gradEtaNorm, abs_of_nonneg positivePartSub_nonneg,
          abs_of_nonneg (norm_nonneg _), mul_assoc, mul_left_comm, mul_comm]
          using hcrossAbs_int
      exact
        weighted_caccioppoli_absorb_core
          (μ := μ) (Q := Equad) (M := fluxNorm) (η := η)
          (ζ := gradEtaNorm) (v := positivePartSub u k) (Λ := A.Λ)
          A.Λ_pos hcore' hcoeff
          (by simpa [leftTerm] using hleft_int) hcross_int' hbound_int
    simpa [leftTerm] using hcore_absorb
  have hcoercive_weighted :
      A.lam * ∫ x in Ω, η x ^ 2 * ‖hw_trunc.weakGrad x‖ ^ 2 ∂volume ≤ ∫ x, leftTerm x ∂μ := by
    have hcoercive_ae :
        ∀ᵐ x ∂μ,
          A.lam * (η x ^ 2 * ‖hw_trunc.weakGrad x‖ ^ 2) ≤ leftTerm x := by
      filter_upwards [A.coercive] with x hx
      have hη_sq_nonneg : 0 ≤ η x ^ 2 := by positivity
      have := mul_le_mul_of_nonneg_left (hx (hw_trunc.weakGrad x)) hη_sq_nonneg
      simpa [leftTerm, Equad, bilinFormIntegrandOfCoeff, real_inner_comm,
        mul_assoc, mul_left_comm, mul_comm] using this
    have hleftSide_int :
        Integrable (fun x => A.lam * (η x ^ 2 * ‖hw_trunc.weakGrad x‖ ^ 2)) μ := by
      simpa [μ] using hweighted_grad_sq_int.const_mul A.lam
    calc
      A.lam * ∫ x in Ω, η x ^ 2 * ‖hw_trunc.weakGrad x‖ ^ 2 ∂volume
          = ∫ x, A.lam * (η x ^ 2 * ‖hw_trunc.weakGrad x‖ ^ 2) ∂μ := by
              simp [μ, integral_const_mul]
      _ ≤ ∫ x, leftTerm x ∂μ := by
            exact integral_mono_ae hleftSide_int hleft_int hcoercive_ae
  have hmain :
      A.lam * ∫ x in Ω, η x ^ 2 * ‖hw_trunc.weakGrad x‖ ^ 2 ∂volume ≤
        4 * A.Λ * ∫ x, gradEtaNorm x ^ 2 * |positivePartSub u k x| ^ 2 ∂μ := by
    exact hcoercive_weighted.trans hleft_bound
  have hlam_ne : A.lam ≠ 0 := ne_of_gt A.hlam
  have hdiv :
      ∫ x in Ω, η x ^ 2 * ‖hw_trunc.weakGrad x‖ ^ 2 ∂volume ≤
        (4 * A.Λ * ∫ x, gradEtaNorm x ^ 2 * |positivePartSub u k x| ^ 2 ∂μ) / A.lam := by
    refine (le_div_iff₀ A.hlam).2 ?_
    simpa [mul_assoc, mul_left_comm, mul_comm] using hmain
  calc
    ∫ x in Ω, η x ^ 2 * ‖hw_trunc.weakGrad x‖ ^ 2 ∂volume
        ≤ (4 * A.Λ * ∫ x, gradEtaNorm x ^ 2 * |positivePartSub u k x| ^ 2 ∂μ) / A.lam := hdiv
    _ = 4 * ellipticityRatio A *
          ∫ x in Ω, ‖fderiv ℝ η x‖ ^ 2 * |positivePartSub u k x| ^ 2 ∂volume := by
          simp [μ, ellipticityRatio, gradEtaNorm, div_eq_mul_inv, mul_assoc, mul_left_comm, mul_comm]

/-- `u` is a weak subsolution of `-div(A∇u) ≤ 0` on `Ω`. -/
def IsSubsolution {Ω : Set E} (A : EllipticCoeff d Ω) (u : E → ℝ) : Prop :=
  MemW1p 2 u Ω ∧
  ∀ (hu : MemW1pWitness 2 u Ω)
    (φ : E → ℝ), MemH01 φ Ω →
    ∀ (hφ : MemW1pWitness 2 φ Ω),
      (∀ x, 0 ≤ φ x) → bilinFormOfCoeff A hu hφ ≤ 0

\end{lstlisting}
\begin{proof}
Test the subsolution inequality against
$\varphi := \varphi_{\eta,u,\ell} = \eta^2 (u-\ell)_+$, which is
admissible by the cutoff admissibility theorem and is nonnegative.
The subsolution condition gives
$\int_\Omega \langle A\nabla u, \nabla \varphi\rangle \, dx \le 0$.
Computing the gradient of $\varphi$,
\begin{align*}
  \nabla \varphi(x)
   = \eta(x)^2 \, \nabla (u-\ell)_+(x)
     + 2\, \eta(x)\, (u-\ell)_+(x) \, \nabla \eta(x),
\end{align*}
which is the content of
\leanname{deGiorgiCutoffTestWitnessWeighted\_grad}. Combining with the
identification $\nabla u = \nabla(u - \ell)_+$ on $\{u > \ell\}$ and
$\nabla(u-\ell)_+ = 0$ on $\{u \le \ell\}$ (the lemmas
\leanname{truncGrad\_eq\_on\_superlevel} and
\leanname{truncGrad\_zero\_on\_sublevel}), we may rewrite
\begin{multline*}
  0 \ge
   \int_\Omega \eta(x)^2 \,
     \langle A(x) \nabla(u-\ell)_+(x), \nabla(u-\ell)_+(x)\rangle \, dx \\
   + \int_\Omega 2 \eta(x) (u-\ell)_+(x)\,
     \langle A(x) \nabla(u-\ell)_+(x), \nabla \eta(x)\rangle \, dx.
\end{multline*}
By the coercivity bound for $A$ (lemma
$\Lambda$-coercive: $\langle Aw,w\rangle \ge \lambda \|w\|^2$) and the
upper bound (lemma $\langle Aw,w\rangle \le \Lambda\|w\|^2$ together
with $\|Aw\|^2 \le \Lambda \langle Aw,w\rangle$), the first term is
bounded below by $\lambda \int \eta^2 \|\nabla(u-\ell)_+\|^2$, while
the second is dominated in absolute value by
\[
  \int_\Omega 2\, \eta(x) (u-\ell)_+(x) \,
    \|A(x)\nabla(u-\ell)_+(x)\|\,\|\nabla \eta(x)\| \, dx.
\]
Setting $E := \langle A\nabla(u-\ell)_+, \nabla(u-\ell)_+\rangle$ and
$M := \|A\nabla(u-\ell)_+\|$, we have $|M|^2 \le \Lambda E$, and
applying the abstract absorption lemma above with these choices
yields
\[
  \int_\Omega \eta(x)^2\, E(x)\, dx
   \le 4\,\Lambda \int_\Omega \|\nabla \eta(x)\|^2\,
     |(u-\ell)_+(x)|^2 \, dx.
\]
Coercivity now bounds the left-hand side from below by
$\lambda \int_\Omega \eta^2 \|\nabla(u-\ell)_+\|^2$, and dividing by
$\lambda$ produces the announced inequality with constant
$4 \Lambda/\lambda$.
\end{proof}

\begin{theorem}[Ball Caccioppoli,
\leanname{caccioppoli\_weighted\_on\_ball\_of\_ballPosPart}]
The same conclusion holds with $\Omega = \Ball{s}{x_0}$ and a witness
for $(u-\ell)_+$ on the ball, provided $\eta$ is smooth, nonnegative,
$|\eta| \le 1$, $\|\nabla \eta\| \le C_\eta$, and $\tsupp\,\eta
\subseteq \Ball{s}{x_0}$.
\end{theorem}

\begin{lstlisting}[style=Matlab-editor,basicstyle=\mlttfamily,escapechar=",mlshowsectionrules=true,mathescape=true,morecomment={[s]{\%\{}{\%\}}},language=lean,numbers=none]
/-- Ball-specialized weighted Caccioppoli inequality, using the generalized
cutoff admissibility theorem from the previous section. -/
theorem caccioppoli_weighted_on_ball_of_ballPosPart
    {u η : E → ℝ} {x₀ : E} {s Cη k : ℝ}
    (A : EllipticCoeff d (Metric.ball x₀ s))
    (_hs : 0 < s)
    (hsub : IsSubsolution A u)
    (hu : MemW1pWitness 2 u (Metric.ball x₀ s))
    (hw_trunc : MemW1pWitness 2 (positivePartSub u k) (Metric.ball x₀ s))
    (hη : ContDiff ℝ (⊤ : ℕ∞) η)
    (hη_nonneg : ∀ x, 0 ≤ η x)
    (hη_bound : ∀ x, |η x| ≤ 1)
    (hCη : 0 ≤ Cη)
    (hη_grad_bound : ∀ x, ‖fderiv ℝ η x‖ ≤ Cη)
    (hη_sub_ball : tsupport η ⊆ Metric.ball x₀ s) :
    ∫ x in Metric.ball x₀ s, η x ^ 2 * ‖hw_trunc.weakGrad x‖ ^ 2 ∂volume ≤
      4 * ellipticityRatio A *
        ∫ x in Metric.ball x₀ s, ‖fderiv ℝ η x‖ ^ 2 * |positivePartSub u k x| ^ 2 ∂volume := by
  haveI : IsFiniteMeasure (volume.restrict (Metric.ball x₀ s)) := by
    rw [isFiniteMeasure_restrict]
    exact measure_ball_lt_top.ne
  have hη_comp : HasCompactSupport η :=
    hasCompactSupport_of_tsupport_subset_ball hη_sub_ball
  have hη_admissible :
      MemW01p 2 (deGiorgiCutoffTestGeneral η u k) (Metric.ball x₀ s) :=
    deGiorgiCutoffTest_memW01p_of_truncWitness
      isOpen_ball hw_trunc hη (by norm_num) hCη hη_bound hη_grad_bound hη_comp hη_sub_ball
  exact caccioppoli_weighted_of_subsolution isOpen_ball A hsub hu hw_trunc hη hη_nonneg
    (by norm_num) hCη hη_bound hη_grad_bound hη_admissible

\end{lstlisting}
The approximation-based variant is
\leanname{caccioppoli\_weighted\_on\_ball\_of\_posPartApprox}.

\begin{theorem}[Localized energy estimate]
(\leanname{deGiorgi\_energy\_estimate\_on\_concentricBalls\_of\_ballPosPart})\\
Let $0 < r < s$, $A$ elliptic on $\Ball{s}{x_0}$, $u$ a subsolution.
Let $\eta$ be smooth, $0 \le \eta \le 1$ on $E$, $\eta = 1$ on
$\Ball{r}{x_0}$, $\tsupp\,\eta \subset \Ball{s}{x_0}$, and
$\|\nabla \eta\| \le C_\eta$. Then
\begin{multline*}
  \int_{\Ball{r}{x_0}} \bigl\|\nabla(u-\ell)_+(x)\bigr\|^2 \, dx \\
  \le 4\,\frac{\Lambda}{\lambda} \, C_\eta^2 \,
    \int_{\Ball{s}{x_0}} \bigl|(u-\ell)_+(x)\bigr|^2 \, dx.
\end{multline*}
\end{theorem}

\begin{lstlisting}[style=Matlab-editor,basicstyle=\mlttfamily,escapechar=",mlshowsectionrules=true,mathescape=true,morecomment={[s]{\%\{}{\%\}}},language=lean,numbers=none]
/-- Localized De Giorgi energy estimate on concentric balls. This packages the
weighted Caccioppoli inequality together with the localization step. -/
theorem deGiorgi_energy_estimate_on_concentricBalls_of_ballPosPart
    {u η : E → ℝ} {x₀ : E} {r s Cη k : ℝ}
    (A : EllipticCoeff d (Metric.ball x₀ s))
    (hr : 0 < r) (hrs : r < s)
    (hsub : IsSubsolution A u)
    (hu : MemW1pWitness 2 u (Metric.ball x₀ s))
    (hw_trunc : MemW1pWitness 2 (positivePartSub u k) (Metric.ball x₀ s))
    (hη : ContDiff ℝ (⊤ : ℕ∞) η)
    (hη_nonneg : ∀ x, 0 ≤ η x)
    (hη_eq_one : ∀ x ∈ Metric.ball x₀ r, η x = 1)
    (hη_bound : ∀ x, |η x| ≤ 1)
    (hCη : 0 ≤ Cη)
    (hη_grad_bound : ∀ x, ‖fderiv ℝ η x‖ ≤ Cη)
    (hη_sub_ball : tsupport η ⊆ Metric.ball x₀ s) :
    ∫ x in Metric.ball x₀ r, ‖hw_trunc.weakGrad x‖ ^ 2 ∂volume ≤
      4 * ellipticityRatio A * Cη ^ 2 *
        ∫ x in Metric.ball x₀ s, |positivePartSub u k x| ^ 2 ∂volume := by
  have hs : 0 < s := lt_trans hr hrs
  have htrunc_sq_int :
      IntegrableOn (fun x => ‖hw_trunc.weakGrad x‖ ^ 2)
        (Metric.ball x₀ s) volume := by
    simpa [pow_two] using hw_trunc.weakGrad_norm_memLp.integrable_sq
  have hweighted_int :
      IntegrableOn (fun x => η x ^ 2 * ‖hw_trunc.weakGrad x‖ ^ 2)
        (Metric.ball x₀ s) volume := by
    refine Integrable.mono' htrunc_sq_int ?_ ?_
    · exact
        (((hη.continuous.pow 2).aemeasurable).mul
          htrunc_sq_int.aestronglyMeasurable.aemeasurable).aestronglyMeasurable
    · filter_upwards with x
      have hηx_nonneg : 0 ≤ η x := hη_nonneg x
      have hηx_le_one : η x ≤ 1 := by
        simpa [abs_of_nonneg hηx_nonneg] using hη_bound x
      have hηx_sq_le_one : η x ^ 2 ≤ 1 := by
        nlinarith
      have hgw_nonneg : 0 ≤ ‖hw_trunc.weakGrad x‖ ^ 2 := by positivity
      have hprod_nonneg : 0 ≤ η x ^ 2 * ‖hw_trunc.weakGrad x‖ ^ 2 := by positivity
      have hle : η x ^ 2 * ‖hw_trunc.weakGrad x‖ ^ 2 ≤ ‖hw_trunc.weakGrad x‖ ^ 2 := by
        nlinarith
      simpa [Real.norm_eq_abs, abs_of_nonneg hprod_nonneg, abs_of_nonneg hgw_nonneg] using hle
  have hpos_sq_int :
      IntegrableOn (fun x => |positivePartSub u k x| ^ 2)
        (Metric.ball x₀ s) volume := by
    simpa [pow_two] using hw_trunc.memLp.integrable_sq
  have hgrad_term_int :
      IntegrableOn (fun x => ‖fderiv ℝ η x‖ ^ 2 * |positivePartSub u k x| ^ 2)
        (Metric.ball x₀ s) volume := by
    refine Integrable.mono' (hpos_sq_int.const_mul (Cη ^ 2)) ?_ ?_
    · exact
        ((((hη.continuous_fderiv (by simp : ((⊤ : ℕ∞) : WithTop ℕ∞) ≠ 0)).norm.pow 2).aemeasurable).mul
          hpos_sq_int.aestronglyMeasurable.aemeasurable).aestronglyMeasurable
    · filter_upwards with x
      have hfd_sq_le : ‖fderiv ℝ η x‖ ^ 2 ≤ Cη ^ 2 := by
        exact sq_le_sq.mpr (by
          simpa [abs_of_nonneg (norm_nonneg _), abs_of_nonneg hCη] using hη_grad_bound x)
      have hpos_nonneg : 0 ≤ |positivePartSub u k x| ^ 2 := by positivity
      have hterm_nonneg :
          0 ≤ ‖fderiv ℝ η x‖ ^ 2 * |positivePartSub u k x| ^ 2 := by positivity
      have hle :
          ‖fderiv ℝ η x‖ ^ 2 * |positivePartSub u k x| ^ 2 ≤
            Cη ^ 2 * |positivePartSub u k x| ^ 2 :=
        mul_le_mul_of_nonneg_right hfd_sq_le hpos_nonneg
      simpa [Real.norm_eq_abs, abs_of_nonneg hterm_nonneg] using hle
  have hweighted :
      ∫ x in Metric.ball x₀ s, η x ^ 2 * ‖hw_trunc.weakGrad x‖ ^ 2 ∂volume ≤
        (4 * ellipticityRatio A) *
          ∫ x in Metric.ball x₀ s, ‖fderiv ℝ η x‖ ^ 2 * |positivePartSub u k x| ^ 2
            ∂volume := by
    simpa [mul_assoc, mul_left_comm, mul_comm] using
      caccioppoli_weighted_on_ball_of_ballPosPart
        A hs hsub hu hw_trunc hη hη_nonneg hη_bound hCη hη_grad_bound hη_sub_ball
  have hleft_eq :
      ∫ x in Metric.ball x₀ r, ‖hw_trunc.weakGrad x‖ ^ 2 ∂volume =
        ∫ x in Metric.ball x₀ r, η x ^ 2 * ‖hw_trunc.weakGrad x‖ ^ 2 ∂volume := by
    refine integral_congr_ae ?_
    refine (ae_restrict_iff' (μ := volume) measurableSet_ball).2 ?_
    filter_upwards with x hx
    simp [hη_eq_one x hx]
  have hweighted_nonneg :
      ∀ x, 0 ≤ η x ^ 2 * ‖hw_trunc.weakGrad x‖ ^ 2 := by
    intro x
    positivity
  have hleft_mono :
      ∫ x in Metric.ball x₀ r, η x ^ 2 * ‖hw_trunc.weakGrad x‖ ^ 2 ∂volume ≤
        ∫ x in Metric.ball x₀ s, η x ^ 2 * ‖hw_trunc.weakGrad x‖ ^ 2 ∂volume := by
    exact setIntegral_mono_set hweighted_int
      (ae_of_all _ hweighted_nonneg)
      (ae_of_all _ (Metric.ball_subset_ball (le_of_lt hrs)))
  have hgrad_bound :
      ∫ x in Metric.ball x₀ s, ‖fderiv ℝ η x‖ ^ 2 * |positivePartSub u k x| ^ 2 ∂volume ≤
        ∫ x in Metric.ball x₀ s, Cη ^ 2 * |positivePartSub u k x| ^ 2 ∂volume := by
    refine integral_mono_ae hgrad_term_int (hpos_sq_int.const_mul (Cη ^ 2)) ?_
    refine (ae_restrict_iff' (μ := volume) measurableSet_ball).2 ?_
    filter_upwards with x hx
    have hfd_sq_le : ‖fderiv ℝ η x‖ ^ 2 ≤ Cη ^ 2 := by
      exact sq_le_sq.mpr (by
        simpa [abs_of_nonneg (norm_nonneg _), abs_of_nonneg hCη] using hη_grad_bound x)
    have hpos_nonneg : 0 ≤ |positivePartSub u k x| ^ 2 := by positivity
    exact mul_le_mul_of_nonneg_right hfd_sq_le hpos_nonneg
  have hcoeff_nonneg : 0 ≤ 4 * ellipticityRatio A := by
    have hRatio_nonneg : 0 ≤ ellipticityRatio A := A.ellipticityRatio_nonneg
    nlinarith
  calc
    ∫ x in Metric.ball x₀ r, ‖hw_trunc.weakGrad x‖ ^ 2 ∂volume
        = ∫ x in Metric.ball x₀ r, η x ^ 2 * ‖hw_trunc.weakGrad x‖ ^ 2 ∂volume := hleft_eq
    _ ≤ ∫ x in Metric.ball x₀ s, η x ^ 2 * ‖hw_trunc.weakGrad x‖ ^ 2 ∂volume := hleft_mono
    _ ≤ (4 * ellipticityRatio A) *
          ∫ x in Metric.ball x₀ s, ‖fderiv ℝ η x‖ ^ 2 * |positivePartSub u k x| ^ 2
            ∂volume := hweighted
    _ ≤ (4 * ellipticityRatio A) *
          ∫ x in Metric.ball x₀ s, Cη ^ 2 * |positivePartSub u k x| ^ 2 ∂volume := by
            exact mul_le_mul_of_nonneg_left hgrad_bound hcoeff_nonneg
    _ = 4 * ellipticityRatio A * Cη ^ 2 *
          ∫ x in Metric.ball x₀ s, |positivePartSub u k x| ^ 2 ∂volume := by
            rw [integral_const_mul]
            ring

\end{lstlisting}
\begin{proof}
Since $\eta = 1$ on $\Ball{r}{x_0}$, the integral on the left equals
$\int_{\Ball{r}{x_0}} \eta^2 \|\nabla(u-\ell)_+\|^2$. Monotonicity in
the integration set and the weighted Caccioppoli inequality bound
this by $4(\Lambda/\lambda)\int_{\Ball{s}{x_0}} \|\nabla \eta\|^2
|(u-\ell)_+|^2$. The pointwise bound $\|\nabla\eta\|^2 \le C_\eta^2$
on $\Ball{s}{x_0}$ then gives the claim.
\end{proof}

\subsection{Chebyshev bound on superlevel sets}

\begin{theorem}[\leanname{superlevel\_measureReal\_le\_integral\_posPart}]
Let $(\alpha, \mu)$ be a finite measure space, $u : \alpha \to \R$,
and let $\theta < \ell$. If $|(u-\theta)_+|^2$ is integrable, then
\[
  \mu\bigl(\{x : \ell < u(x)\}\bigr)
   \le \frac{1}{(\ell-\theta)^2}\,
       \int |(u-\theta)_+(x)|^2 \, d\mu(x).
\]
\end{theorem}

\begin{lstlisting}[style=Matlab-editor,basicstyle=\mlttfamily,escapechar=",mlshowsectionrules=true,mathescape=true,morecomment={[s]{\%\{}{\%\}}},language=lean,numbers=none]
/-- Chebyshev bound for the superlevel set `{u > λ}` using `(u-θ)₊`. -/
theorem superlevel_measureReal_le_integral_posPart
    {α : Type*} [MeasurableSpace α] {μ : Measure α} [IsFiniteMeasure μ]
    {u : α → ℝ} {θ lam : ℝ}
    (hθl : θ < lam)
    (hint : Integrable (fun x => |positivePartSub u θ x| ^ 2) μ) :
    μ.real {x | lam < u x} ≤
      ((lam - θ) ^ 2)⁻¹ * ∫ x, |positivePartSub u θ x| ^ 2 ∂μ := by
  let f : α → ℝ := fun x => |positivePartSub u θ x| ^ 2
  let ε : ℝ := (lam - θ) ^ 2
  have hε_pos : 0 < ε := by
    dsimp [ε]
    nlinarith
  have hf_nonneg : 0 ≤ᵐ[μ] f := by
    refine Filter.Eventually.of_forall ?_
    intro x
    dsimp [f]
    positivity
  have hmarkov :
      ε * μ.real {x | ε ≤ f x} ≤ ∫ x, f x ∂μ := by
    simpa [f, ε] using
      (MeasureTheory.mul_meas_ge_le_integral_of_nonneg (μ := μ) hf_nonneg hint ε)
  have hsubset : {x | lam < u x} ⊆ {x | ε ≤ f x} := by
    intro x hx
    dsimp [f, ε]
    exact positivePartSub_sq_ge_gap_sq_of_lt hθl hx
  have hmono : μ.real {x | lam < u x} ≤ μ.real {x | ε ≤ f x} := by
    exact measureReal_mono hsubset
  have hmain : ε * μ.real {x | lam < u x} ≤ ∫ x, f x ∂μ := by
    calc
      ε * μ.real {x | lam < u x} ≤ ε * μ.real {x | ε ≤ f x} := by
        gcongr
      _ ≤ ∫ x, f x ∂μ := hmarkov
  have hdiv : μ.real {x | lam < u x} ≤ (∫ x, f x ∂μ) / ε := by
    refine (le_div_iff₀ hε_pos).2 ?_
    simpa [mul_comm, mul_left_comm, mul_assoc] using hmain
  simpa [div_eq_mul_inv, f, ε, mul_comm, mul_left_comm, mul_assoc] using hdiv

\end{lstlisting}
\begin{proof}
Set $f(x) := |(u-\theta)_+(x)|^2$ and $\varepsilon := (\ell-\theta)^2$. The
hypothesis $\theta < \ell$ ensures $\varepsilon > 0$. By
\leanname{positivePartSub\_sq\_ge\_gap\_sq\_of\_lt}, on the superlevel set
$\{x : \ell < u(x)\}$ the value of $f$ is at least
$(\ell - \theta)^2 = \varepsilon$, so we have the set inclusion
\[
  \{x : \ell < u(x)\}
    \subseteq \{x : \varepsilon \le f(x)\}.
\]
By monotonicity of measure (\leanname{measureReal\_mono}) this gives
$\mu\{\ell < u\} \le \mu\{\varepsilon \le f\}$.

Markov's inequality for nonnegative integrable functions
(\leanname{MeasureTheory.mul\_meas\_ge\_le\_integral\_of\_nonneg}) applied to
$f$ at threshold $\varepsilon$ yields
\[
  \varepsilon \cdot \mu\{\varepsilon \le f\}
   \le \int_{\alpha} f(x)\,d\mu(x).
\]
Combining the set inclusion with monotone multiplication by $\varepsilon \ge 0$,
\[
  \varepsilon \cdot \mu\{\ell < u\}
   \le \varepsilon \cdot \mu\{\varepsilon \le f\}
   \le \int_{\alpha} f\,d\mu.
\]
Dividing by the positive constant $\varepsilon$ produces the desired bound
$\mu\{\ell < u\} \le ((\ell - \theta)^2)^{-1}\int |(u-\theta)_+|^2\,d\mu$.
\end{proof}

\subsection{Pre-iteration estimate}

We now combine the energy estimate with a Sobolev/Holder
interpolation step to obtain a single-step nonlinear inequality
between truncation energies at two levels and on two concentric balls.

\begin{theorem}[Abstract pre-iteration combination,
\leanname{deGiorgi\_preiter\_abstract}]
Let $d > 0$ and let $C_{\mathrm{sob}}, C_{\mathrm{en}}, I_\theta, m,
\theta < \ell$ all be nonneg numbers. Suppose
\begin{align*}
  I_\ell &\le C_{\mathrm{sob}} \cdot G \cdot m^{2/d}, \\
  m &\le (\ell - \theta)^{-2}\, I_\theta, \\
  G &\le C_{\mathrm{en}} \cdot I_\theta.
\end{align*}
Then
\[
  I_\ell \le C_{\mathrm{sob}}\, C_{\mathrm{en}}\,
    \bigl((\ell-\theta)^{-2} I_\theta\bigr)^{2/d}\, I_\theta.
\]
\end{theorem}

\begin{lstlisting}[style=Matlab-editor,basicstyle=\mlttfamily,escapechar=",mlshowsectionrules=true,mathescape=true,morecomment={[s]{\%\{}{\%\}}},language=lean,numbers=none]
/-- Abstract real-variable combination step for De Giorgi pre-iteration. -/
theorem deGiorgi_preiter_abstract
    {d : ℕ} (hd : 0 < (d : ℝ))
    {Ilam Iθ G m Csob Cenergy θ lam : ℝ}
    (hCsob : 0 ≤ Csob) (hCenergy : 0 ≤ Cenergy)
    (hIθ : 0 ≤ Iθ) (hm : 0 ≤ m)
    (hsob : Ilam ≤ Csob * G * m ^ (2 / (d : ℝ)))
    (hmeasure : m ≤ ((lam - θ) ^ 2)⁻¹ * Iθ)
    (henergy : G ≤ Cenergy * Iθ) :
    Ilam ≤
      Csob * Cenergy * ((((lam - θ) ^ 2)⁻¹ * Iθ) ^ (2 / (d : ℝ))) * Iθ := by
  have hexp_nonneg : 0 ≤ 2 / (d : ℝ) := by positivity
  have hm_pow :
      m ^ (2 / (d : ℝ)) ≤ ((((lam - θ) ^ 2)⁻¹ * Iθ) ^ (2 / (d : ℝ))) := by
    exact Real.rpow_le_rpow hm hmeasure hexp_nonneg
  have hGI :
      G * m ^ (2 / (d : ℝ)) ≤
        (Cenergy * Iθ) * ((((lam - θ) ^ 2)⁻¹ * Iθ) ^ (2 / (d : ℝ))) := by
    exact mul_le_mul henergy hm_pow (by positivity) (by positivity)
  calc
    Ilam ≤ Csob * (G * m ^ (2 / (d : ℝ))) := by
      simpa [mul_assoc] using hsob
    _ ≤ Csob * ((Cenergy * Iθ) * ((((lam - θ) ^ 2)⁻¹ * Iθ) ^ (2 / (d : ℝ)))) := by
      exact mul_le_mul_of_nonneg_left hGI hCsob
    _ = Csob * Cenergy * ((((lam - θ) ^ 2)⁻¹ * Iθ) ^ (2 / (d : ℝ))) * Iθ := by
      ring

\end{lstlisting}
\begin{proof}
Raise the second inequality to the power $2/d \ge 0$, multiply by
$G$, then by $C_{\mathrm{sob}}$, using $C_{\mathrm{en}}$ to bound
$G$.
\end{proof}

\begin{theorem}[Packaged pre-iteration,
\leanname{deGiorgi\_preiter\_of\_energy}]
With $\mu$ a finite measure and the same notation, if
\begin{itemize}
\item $I_\ell \le C_{\mathrm{sob}}\, G\, \mu(\{\ell < u\})^{2/d}$,
\item $G \le C_{\mathrm{en}} \cdot I_\theta$,
\item $\int |(u-\theta)_+|^2 \, d\mu \le I_\theta$,
\end{itemize}
then
\(
  I_\ell \le C_{\mathrm{sob}}\, C_{\mathrm{en}}\,
    ((\ell-\theta)^{-2} I_\theta)^{2/d}\, I_\theta.
\)
\end{theorem}

\begin{lstlisting}[style=Matlab-editor,basicstyle=\mlttfamily,escapechar=",mlshowsectionrules=true,mathescape=true,morecomment={[s]{\%\{}{\%\}}},language=lean,numbers=none]
/-- Packaged pre-iteration step after Sobolev, energy, and Chebyshev inputs. -/
theorem deGiorgi_preiter_of_energy
    {α : Type*} [MeasurableSpace α] {μ : Measure α} [IsFiniteMeasure μ]
    {d : ℕ} (hd : 0 < (d : ℝ))
    {u : α → ℝ} {θ lam : ℝ}
    {Ilam Iθ G Csob Cenergy : ℝ}
    (hθl : θ < lam)
    (hCsob : 0 ≤ Csob) (hCenergy : 0 ≤ Cenergy)
    (hIθ : 0 ≤ Iθ)
    (hint : Integrable (fun x => |positivePartSub u θ x| ^ 2) μ)
    (hsob :
      Ilam ≤ Csob * G * (μ.real {x | lam < u x}) ^ (2 / (d : ℝ)))
    (henergy :
      G ≤ Cenergy * Iθ)
    (hIθ_bound :
      ∫ x, |positivePartSub u θ x| ^ 2 ∂μ ≤ Iθ) :
    Ilam ≤
      Csob * Cenergy * ((((lam - θ) ^ 2)⁻¹ * Iθ) ^ (2 / (d : ℝ))) * Iθ := by
  have hm_nonneg : 0 ≤ μ.real {x | lam < u x} := measureReal_nonneg
  have hmeasure0 :
      μ.real {x | lam < u x} ≤
        ((lam - θ) ^ 2)⁻¹ * ∫ x, |positivePartSub u θ x| ^ 2 ∂μ :=
    superlevel_measureReal_le_integral_posPart hθl hint
  have hmeasure :
      μ.real {x | lam < u x} ≤ ((lam - θ) ^ 2)⁻¹ * Iθ := by
    calc
      μ.real {x | lam < u x}
          ≤ ((lam - θ) ^ 2)⁻¹ * ∫ x, |positivePartSub u θ x| ^ 2 ∂μ := hmeasure0
      _ ≤ ((lam - θ) ^ 2)⁻¹ * Iθ := by
        gcongr
  exact deGiorgi_preiter_abstract hd hCsob hCenergy hIθ hm_nonneg hsob hmeasure henergy

\end{lstlisting}
\begin{proof}
By the Chebyshev bound, $\mu(\{\ell < u\}) \le (\ell-\theta)^{-2}\,
\int|(u-\theta)_+|^2 \, d\mu \le (\ell-\theta)^{-2}\, I_\theta$. The
abstract combination lemma now applies.
\end{proof}

The cutoff Sobolev step required to discharge the Sobolev hypothesis
$I_\ell \le (C_{\mathrm{gns}})^2 \int \|\nabla \varphi\|^2\,
\mu(\{\ell < u\})^{2/d}$ is

\begin{theorem}[\leanname{deGiorgi\_cutoffSobolev\_on\_concentricBalls\_of\_ballPosPart}]
Under the standing hypotheses ($d > 2$, $0 < r < s$, $u$ admitting a
$W^{1,2}$ witness, $(u-\theta)_+$ admitting a $W^{1,2}$ witness on
$\Ball{s}{x_0}$, $\theta < \ell$, $\eta$ smooth with $\eta = 1$ on
$\Ball{r}{x_0}$, $|\eta| \le 1$, $\tsupp \eta \subseteq
\Ball{s}{x_0}$), there exists a $W^{1,2}$-witness $w_{\eta,\theta}$
for $\eta^2 (u-\theta)_+$ such that
\begin{multline*}
  \int_{\Ball{r}{x_0}} \bigl|(u-\ell)_+(x)\bigr|^2 \, dx
  \le (C_{\mathrm{gns}}(d,2))^2 \cdot \\
   \Bigl(\int_{\Ball{s}{x_0}} \|\nabla w_{\eta,\theta}(x)\|^2 \, dx
   \Bigr) \cdot
    \mu_s\bigl(\{\ell < u\}\bigr)^{2/d},
\end{multline*}
where $\mu_s$ denotes Lebesgue measure restricted to $\Ball{s}{x_0}$.
\end{theorem}

\begin{lstlisting}[style=Matlab-editor,basicstyle=\mlttfamily,escapechar=",mlshowsectionrules=true,mathescape=true,morecomment={[s]{\%\{}{\%\}}},language=lean,numbers=none]
/-- Sobolev/Hölder step for De Giorgi pre-iteration on concentric balls.

The argument passes through a zero-trace cutoff witness for
`η² (u - θ)₊`. The originally intended statement with only the bare
`θ`-truncation gradient on the right-hand side is false, e.g. for constant
superlevel functions. The proof here records the actual Chapter 05 mechanism:

1. build a `W₀^{1,2}` witness for `η² (u - θ)₊`,
2. zero-extend it to the whole space,
3. apply whole-space Sobolev, and
4. combine it with the `L²`-to-`L^{2d/(d-2)}` restriction estimate on
   `{x | λ < u x}`. -/
theorem deGiorgi_cutoffSobolev_on_concentricBalls_of_ballPosPart
    {u η : E → ℝ} {x₀ : E} {r s θ lam : ℝ}
    (hd : 2 < (d : ℝ))
    (hr : 0 < r) (hrs : r < s)
    (hu : MemW1pWitness 2 u (Metric.ball x₀ s))
    (hwθ : MemW1pWitness 2 (positivePartSub u θ) (Metric.ball x₀ s))
    (hθl : θ < lam)
    (hη : ContDiff ℝ (⊤ : ℕ∞) η)
    (_hη_nonneg : ∀ x, 0 ≤ η x)
    (hη_eq_one : ∀ x ∈ Metric.ball x₀ r, η x = 1)
    (hη_bound : ∀ x, |η x| ≤ 1)
    (hη_sub_ball : tsupport η ⊆ Metric.ball x₀ s) :
    ∃ hwηθ : MemW1pWitness 2 (deGiorgiCutoffTestGeneral η u θ) (Metric.ball x₀ s),
      ∫ x in Metric.ball x₀ r, |positivePartSub u lam x| ^ 2 ∂volume ≤
        (C_gns d 2) ^ 2 *
          (∫ x in Metric.ball x₀ s, ‖hwηθ.weakGrad x‖ ^ 2 ∂volume) *
          ((volume.restrict (Metric.ball x₀ s)).real {x | lam < u x}) ^ (2 / (d : ℝ)) := by
  have hs : 0 < s := lt_trans hr hrs
  let v : E → ℝ := deGiorgiCutoffTestGeneral η u θ
  haveI : IsFiniteMeasure (volume.restrict (Metric.ball x₀ s)) := by
    rw [isFiniteMeasure_restrict]
    exact measure_ball_lt_top.ne
  have hη_comp : HasCompactSupport η :=
    hasCompactSupport_of_tsupport_subset_ball hη_sub_ball
  obtain ⟨Cη, hCη, hη_grad_bound⟩ :=
    exists_fderiv_bound_of_contDiff_hasCompactSupport hη hη_comp
  have hvW01 :
      MemW01p 2 (deGiorgiCutoffTestGeneral η u θ) (Metric.ball x₀ s) :=
    deGiorgiCutoffTest_memW01p_of_truncWitness
      isOpen_ball hwθ hη (by norm_num) hCη hη_bound hη_grad_bound hη_comp hη_sub_ball
  obtain ⟨hwηθ_real, hSob''⟩ :=
    deGiorgi_cutoffSobolev_prepare (d := d) hd hs hu hη hη_bound hη_sub_ball hvW01
  let hwηθ : MemW1pWitness 2 v (Metric.ball x₀ s) :=
    { memLp := by
        simpa [v] using hwηθ_real.memLp
      weakGrad := hwηθ_real.weakGrad
      weakGrad_component_memLp := by
        simpa using hwηθ_real.weakGrad_component_memLp
      isWeakGrad := by
        simpa [v] using hwηθ_real.isWeakGrad }
  refine ⟨hwηθ, ?_⟩
  simpa [v, hwηθ] using
    deGiorgi_cutoffSobolev_superlevelStep (d := d) hd hr hrs hu hwθ hθl hη_eq_one
      hwηθ_real hSob''

\end{lstlisting}
\begin{proof}
The cutoff $\varphi := \eta^2 (u-\theta)_+$ has compact support inside
$\Ball{s}{x_0}$ and lies in $\Wop{2}(\Ball{s}{x_0})$ by
the cutoff admissibility theorem, hence its zero extension to $\R^d$
lies in $W^{1,2}(\R^d)$. The Gagliardo--Nirenberg--Sobolev embedding
gives the global $L^{2^*}$ estimate
$\|\varphi\|_{L^{2^*}} \le C_{\mathrm{gns}}(d,2) \|\nabla
\varphi\|_{L^2}$, where $2^* = 2d/(d-2)$. On $\Ball{r}{x_0}$ one has
$\eta = 1$ and $u > \ell > \theta$ implies $\varphi(x) = (u(x) -
\theta) > \ell - \theta$, so $|(u-\ell)_+|^2 \le |\varphi|^2$ on
$\{\ell < u\}$. Restricting the $L^{2^*}$ estimate to $\Ball{r}{x_0}$
and applying Holder's inequality with exponents $2^*/2$ and its dual
gives the announced bound, where the volume factor reduces to
$\mu_s(\{\ell < u\})^{2/d}$.
\end{proof}

\begin{theorem}[PDE pre-iteration on concentric balls,
\leanname{deGiorgi\_preiter\_on\_concentricBalls\_of\_ballPosPart}]
Under the same hypotheses ($d > 2$, $0 < r < s$, $\theta < \ell$,
$\eta$ as above), one has
\begin{multline*}
  \int_{\Ball{r}{x_0}} |(u-\ell)_+(x)|^2 \, dx
  \le (C_{\mathrm{gns}}(d,2))^2
    \cdot \bigl(8(\Lambda/\lambda + 1)\, C_\eta^2\bigr) \cdot \\
    \bigl((\ell - \theta)^{-2}
       \int_{\Ball{s}{x_0}} |(u-\theta)_+|^2 \, dx\bigr)^{2/d}
    \cdot \int_{\Ball{s}{x_0}} |(u-\theta)_+|^2 \, dx.
\end{multline*}
\end{theorem}

\begin{lstlisting}[style=Matlab-editor,basicstyle=\mlttfamily,escapechar=",mlshowsectionrules=true,mathescape=true,morecomment={[s]{\%\{}{\%\}}},language=lean,numbers=none]
/-- PDE-facing De Giorgi pre-iteration theorem on concentric balls.

This now factors through the explicit cutoff-Sobolev bridge
`deGiorgi_cutoffSobolev_on_concentricBalls_of_ballPosPart` instead of keeping
the local Sobolev argument bundled into one monolithic proof. -/
theorem deGiorgi_preiter_on_concentricBalls_of_ballPosPart
    {u η : E → ℝ} {x₀ : E} {r s θ lam Cη : ℝ}
    (hd : 2 < (d : ℝ))
    (A : EllipticCoeff d (Metric.ball x₀ s))
    (hr : 0 < r) (hrs : r < s)
    (hsub : IsSubsolution A u)
    (hu : MemW1pWitness 2 u (Metric.ball x₀ s))
    (hwθ : MemW1pWitness 2 (positivePartSub u θ) (Metric.ball x₀ s))
    (hθl : θ < lam)
    (hη : ContDiff ℝ (⊤ : ℕ∞) η)
    (hη_nonneg : ∀ x, 0 ≤ η x)
    (hη_eq_one : ∀ x ∈ Metric.ball x₀ r, η x = 1)
    (hη_bound : ∀ x, |η x| ≤ 1)
    (hCη : 0 ≤ Cη)
    (hη_grad_bound : ∀ x, ‖fderiv ℝ η x‖ ≤ Cη)
    (hη_sub_ball : tsupport η ⊆ Metric.ball x₀ s) :
    ∫ x in Metric.ball x₀ r, |positivePartSub u lam x| ^ 2 ∂volume ≤
      (C_gns d 2) ^ 2 * (8 * (ellipticityRatio A + 1) * Cη ^ 2) *
        ((((lam - θ) ^ 2)⁻¹ *
            ∫ x in Metric.ball x₀ s, |positivePartSub u θ x| ^ 2 ∂volume) ^ (2 / (d : ℝ))) *
          ∫ x in Metric.ball x₀ s, |positivePartSub u θ x| ^ 2 ∂volume := by
  have hs : 0 < s := lt_trans hr hrs
  have hd_pos : 0 < (d : ℝ) := by
    linarith
  have hRatio_nonneg : 0 ≤ ellipticityRatio A := A.ellipticityRatio_nonneg
  have hCenergy_nonneg : 0 ≤ 8 * (ellipticityRatio A + 1) * Cη ^ 2 := by
    have hCη_sq_nonneg : 0 ≤ Cη ^ 2 := sq_nonneg Cη
    nlinarith
  have hθ_int :
      Integrable (fun x => |positivePartSub u θ x| ^ 2)
        (volume.restrict (Metric.ball x₀ s)) := by
    simpa [pow_two] using hwθ.memLp.integrable_sq
  have hIθ_nonneg :
      0 ≤ ∫ x in Metric.ball x₀ s, |positivePartSub u θ x| ^ 2 ∂volume := by
    refine integral_nonneg ?_
    intro x
    positivity
  obtain ⟨hwηθ, hsob⟩ :=
    deGiorgi_cutoffSobolev_on_concentricBalls_of_ballPosPart
      (d := d) hd hr hrs hu hwθ hθl hη hη_nonneg hη_eq_one hη_bound hη_sub_ball
  let hwφ : MemW1pWitness 2 (deGiorgiCutoffTestGeneral η u θ) (Metric.ball x₀ s) :=
    deGiorgiCutoffTestWitnessWeighted
      isOpen_ball hwθ hη (by norm_num) hCη hη_bound hη_grad_bound
  have hweighted :
      ∫ x in Metric.ball x₀ s, η x ^ 2 * ‖hwθ.weakGrad x‖ ^ 2 ∂volume ≤
        (4 * ellipticityRatio A) *
          ∫ x in Metric.ball x₀ s, ‖fderiv ℝ η x‖ ^ 2 * |positivePartSub u θ x| ^ 2
            ∂volume := by
    simpa [mul_assoc, mul_left_comm, mul_comm] using
      caccioppoli_weighted_on_ball_of_ballPosPart
        A hs hsub hu hwθ hη hη_nonneg hη_bound hCη hη_grad_bound hη_sub_ball
  have hgrad_bound_phi :
      ∫ x in Metric.ball x₀ s, ‖hwφ.weakGrad x‖ ^ 2 ∂volume ≤
        (8 * (ellipticityRatio A + 1) * Cη ^ 2) *
          ∫ x in Metric.ball x₀ s, |positivePartSub u θ x| ^ 2 ∂volume := by
    simpa [hwφ] using
      deGiorgi_cutoff_gradient_bound_of_weighted
        (x₀ := x₀) (s := s) (u := u) (η := η) (k := θ)
        (Cη := Cη) (Cw := ellipticityRatio A)
        hRatio_nonneg hwθ hη hη_nonneg hη_bound hCη hη_grad_bound hweighted
  have hgrad_ae :
      hwηθ.weakGrad =ᵐ[volume.restrict (Metric.ball x₀ s)] hwφ.weakGrad :=
    MemW1pWitness.ae_eq isOpen_ball hwηθ hwφ
  have hgrad_eq :
      ∫ x in Metric.ball x₀ s, ‖hwηθ.weakGrad x‖ ^ 2 ∂volume =
        ∫ x in Metric.ball x₀ s, ‖hwφ.weakGrad x‖ ^ 2 ∂volume := by
    refine integral_congr_ae ?_
    filter_upwards [hgrad_ae] with x hx
    simp [hx]
  have henergy :
      ∫ x in Metric.ball x₀ s, ‖hwηθ.weakGrad x‖ ^ 2 ∂volume ≤
        (8 * (ellipticityRatio A + 1) * Cη ^ 2) *
          ∫ x in Metric.ball x₀ s, |positivePartSub u θ x| ^ 2 ∂volume := by
    rw [hgrad_eq]
    exact hgrad_bound_phi
  haveI : IsFiniteMeasure (volume.restrict (Metric.ball x₀ s)) := by
    rw [isFiniteMeasure_restrict]
    exact measure_ball_lt_top.ne
  simpa [mul_assoc, mul_left_comm, mul_comm] using
    deGiorgi_preiter_of_energy
      (μ := volume.restrict (Metric.ball x₀ s))
      (d := d) hd_pos (u := u) (θ := θ) (lam := lam)
      (Ilam := ∫ x in Metric.ball x₀ r, |positivePartSub u lam x| ^ 2 ∂volume)
      (Iθ := ∫ x in Metric.ball x₀ s, |positivePartSub u θ x| ^ 2 ∂volume)
      (G := ∫ x in Metric.ball x₀ s, ‖hwηθ.weakGrad x‖ ^ 2 ∂volume)
      (Csob := (C_gns d 2) ^ 2)
      (Cenergy := 8 * (ellipticityRatio A + 1) * Cη ^ 2)
      hθl
      (sq_nonneg _)
      hCenergy_nonneg
      hIθ_nonneg
      hθ_int
      hsob
      henergy
      (by rfl)

\end{lstlisting}
\begin{proof}
Apply the cutoff Sobolev step with witness $w_{\eta,\theta}$, the
energy estimate $\int_{\Ball{s}{x_0}} \|\nabla w_{\eta,\theta}\|^2 \le
8(\Lambda/\lambda+1) C_\eta^2 \int_{\Ball{s}{x_0}} |(u-\theta)_+|^2$
(coming from the gradient identity for $\varphi$ together with the
weighted Caccioppoli inequality and the elementary bound
$(a+b)^2 \le 2a^2 + 2b^2$), and the packaged pre-iteration. The
identification of the gradient bound coefficient is recorded as
\leanname{deGiorgi\_cutoff\_gradient\_bound\_of\_weighted}.
\end{proof}

\subsection{De Giorgi iteration sequences}

We now fix the canonical scales of the iteration.

\begin{definition}[\leanname{deGiorgiTail}, \leanname{deGiorgiRadius},
\leanname{deGiorgiLevel}, \leanname{deGiorgiRecurrenceBase},
\leanname{deGiorgiRecurrenceCoeff}]
For $n \in \N$, $\ell^* > 0$, and $K > 0$, set
\begin{align*}
  \tau_n &:= (1/2)^{n+1}, \\
  r_n &:= (1 + \tau_n)/2, \\
  \ell_n(\ell^*) &:= \ell^* (1 - \tau_n), \\
  B_d &:= 2^{2 + 4/d}, \\
  C_d(K, \ell^*) &:= K \cdot 2^{6 + 8/d} \cdot (\ell^*)^{-4/d}.
\end{align*}
\end{definition}

\begin{lstlisting}[style=Matlab-editor,basicstyle=\mlttfamily,escapechar=",mlshowsectionrules=true,mathescape=true,morecomment={[s]{\%\{}{\%\}}},language=lean,numbers=none]
/-- Canonical tail scale used in the De Giorgi iteration. -/
def deGiorgiTail (n : ℕ) : ℝ :=
  (1 / 2 : ℝ) ^ (n + 1)

/-- Canonical shrinking radii used in the De Giorgi iteration. -/
def deGiorgiRadius (n : ℕ) : ℝ :=
  (1 + deGiorgiTail n) / 2

/-- Canonical increasing levels used in the De Giorgi iteration. -/
def deGiorgiLevel (lamStar : ℝ) (n : ℕ) : ℝ :=
  lamStar * (1 - deGiorgiTail n)

/-- Geometric base in the canonical De Giorgi recurrence. -/
def deGiorgiRecurrenceBase (d : ℕ) : ℝ :=
  (2 : ℝ) ^ (2 + 4 / (d : ℝ))

/-- Leading coefficient in the canonical De Giorgi recurrence. -/
def deGiorgiRecurrenceCoeff (d : ℕ) (K lamStar : ℝ) : ℝ :=
  K * (2 : ℝ) ^ (6 + 8 / (d : ℝ)) * lamStar ^ (-(4 / (d : ℝ)))

\end{lstlisting}
\begin{lemma}
The radii $r_n$ satisfy $1/2 < r_n < 1$ and decrease to $1/2$, the
levels increase from $\ell_0 < \ell^*$ to $\ell^*$, and one has
the gap identities (\leanname{deGiorgiRadius\_gap},
\leanname{deGiorgiLevel\_gap})
\[
  r_n - r_{n+1} = (1/2)^{n+3}, \qquad
  \ell_{n+1}(\ell^*) - \ell_n(\ell^*) = \ell^* (1/2)^{n+2}.
\]
\end{lemma}

\begin{lstlisting}[style=Matlab-editor,basicstyle=\mlttfamily,escapechar=",mlshowsectionrules=true,mathescape=true,morecomment={[s]{\%\{}{\%\}}},language=lean,numbers=none]
lemma deGiorgiRadius_gap (n : ℕ) :
    deGiorgiRadius n - deGiorgiRadius (n + 1) = (1 / 2 : ℝ) ^ (n + 3) := by
  calc
    deGiorgiRadius n - deGiorgiRadius (n + 1)
        = (deGiorgiTail n - deGiorgiTail (n + 1)) / 2 := by
            dsimp [deGiorgiRadius]
            ring
    _ = ((1 / 2 : ℝ) ^ (n + 2)) / 2 := by rw [deGiorgiTail_sub]
    _ = (1 / 2 : ℝ) ^ (n + 3) := by
          rw [pow_succ']
          ring

lemma deGiorgiLevel_gap {lamStar : ℝ} (n : ℕ) :
    deGiorgiLevel lamStar (n + 1) - deGiorgiLevel lamStar n =
      lamStar * (1 / 2 : ℝ) ^ (n + 2) := by
  calc
    deGiorgiLevel lamStar (n + 1) - deGiorgiLevel lamStar n
        = lamStar * (deGiorgiTail n - deGiorgiTail (n + 1)) := by
            dsimp [deGiorgiLevel]
            ring
    _ = lamStar * (1 / 2 : ℝ) ^ (n + 2) := by rw [deGiorgiTail_sub]

\end{lstlisting}
\begin{definition}[\leanname{deGiorgiEnergySeq}]
For a function $u$, base point $x_0$, and threshold $\ell^* > 0$, set
\[
  A_n(u, x_0, \ell^*) := \int_{\Ball{r_n}{x_0}}
     \bigl|(u - \ell_n(\ell^*))_+(x)\bigr|^2 \, dx.
\]
\end{definition}

\begin{lstlisting}[style=Matlab-editor,basicstyle=\mlttfamily,escapechar=",mlshowsectionrules=true,mathescape=true,morecomment={[s]{\%\{}{\%\}}},language=lean,numbers=none]
/-- Canonical De Giorgi energy sequence on nested balls and levels. -/
noncomputable def deGiorgiEnergySeq
    (u : E → ℝ) (x₀ : E) (lamStar : ℝ) (n : ℕ) : ℝ :=
  ∫ x in Metric.ball x₀ (deGiorgiRadius n), |positivePartSub u (deGiorgiLevel lamStar n) x| ^ 2 ∂volume

\end{lstlisting}
\subsection{Recurrence and abstract closeout}

\begin{theorem}[Geometric closeout for nonlinear recurrences,
\leanname{deGiorgi\_recurrence\_closeout}]
Let $Y_n \ge 0$, $C > 0$, $B > 1$, $\alpha > 0$, and assume
\[
  Y_{n+1} \le C\, B^n\, Y_n^{1+\alpha}, \qquad
  Y_0 \le C^{-1/\alpha}\, B^{-1/\alpha^2}.
\]
Then $Y_n \to 0$ as $n \to \infty$.
\end{theorem}

\begin{lstlisting}[style=Matlab-editor,basicstyle=\mlttfamily,escapechar=",mlshowsectionrules=true,mathescape=true,morecomment={[s]{\%\{}{\%\}}},language=lean,numbers=none]
theorem deGiorgi_recurrence_closeout
    {Y : ℕ → ℝ} {C_rec B α : ℝ}
    (hC : 0 < C_rec) (hB : 1 < B) (hα : 0 < α)
    (hY_nonneg : ∀ n, 0 ≤ Y n)
    (h_rec : ∀ n, Y (n + 1) ≤ C_rec * B ^ n * Y n ^ (1 + α))
    (h_small : Y 0 ≤ C_rec ^ (-(1 : ℝ) / α) * B ^ (-(1 : ℝ) / α ^ 2)) :
    ∀ ε > 0, ∃ N : ℕ, ∀ n ≥ N, Y n < ε := by
  let X : ℝ := C_rec * B ^ ((1 : ℝ) / α)
  let D : ℝ := X ^ ((1 : ℝ) / α)
  let q : ℝ := B ^ (-(1 : ℝ) / α)
  let W : ℕ → ℝ := fun n => B ^ ((n : ℝ) / α) * Y n
  let Z : ℕ → ℝ := fun n => D * W n
  have hBpos : 0 < B := by linarith
  have hX_nonneg : 0 ≤ X := by
    dsimp [X]
    positivity
  have hD_pos : 0 < D := by
    dsimp [D]
    positivity
  have hD_nonneg : 0 ≤ D := hD_pos.le
  have hq_nonneg : 0 ≤ q := by
    dsimp [q]
    positivity
  have hq_lt_one : q < 1 := by
    dsimp [q]
    have hneg : -(1 : ℝ) / α < 0 := by
      have hαinv_pos : 0 < (1 : ℝ) / α := by positivity
      rw [show -(1 : ℝ) / α = -((1 : ℝ) / α) by ring]
      simpa using neg_neg_of_pos hαinv_pos
    exact Real.rpow_lt_one_of_one_lt_of_neg hB hneg
  have hW_nonneg : ∀ n, 0 ≤ W n := by
    intro n
    dsimp [W]
    exact mul_nonneg (Real.rpow_nonneg (le_of_lt hBpos) _) (hY_nonneg n)
  have hDα : D ^ α = X := by
    dsimp [D]
    simpa [one_div] using (Real.rpow_inv_rpow hX_nonneg hα.ne')
  have h_small' : Y 0 ≤ D⁻¹ := by
    have hBterm : (B ^ ((1 : ℝ) / α)) ^ (-(1 : ℝ) / α) = B ^ (-(1 : ℝ) / α ^ 2) := by
      rw [← Real.rpow_mul (le_of_lt hBpos)]
      congr 1
      ring
    have hrewrite : C_rec ^ (-(1 : ℝ) / α) * B ^ (-(1 : ℝ) / α ^ 2) = D⁻¹ := by
      dsimp [D, X]
      rw [← Real.rpow_neg hX_nonneg]
      rw [Real.mul_rpow hC.le (by positivity)]
      rw [← hBterm]
      rw [show (-1 : ℝ) / α = -(1 / α) by ring]
    simpa [hrewrite] using h_small
  have hW_rec : ∀ n, W (n + 1) ≤ X * W n ^ (1 + α) := by
    intro n
    have hBn : (B ^ n : ℝ) = B ^ (n : ℝ) := by
      exact (Real.rpow_natCast B n).symm
    calc
      W (n + 1) = B ^ (((n + 1 : ℕ) : ℝ) / α) * Y (n + 1) := by
        simp [W]
      _ ≤ B ^ (((n + 1 : ℕ) : ℝ) / α) * (C_rec * B ^ n * Y n ^ (1 + α)) := by
        gcongr
        exact h_rec n
      _ = X * W n ^ (1 + α) := by
        dsimp [X, W]
        rw [hBn]
        rw [Real.mul_rpow (by positivity) (hY_nonneg n)]
        have hBpow : (B ^ ((n : ℝ) / α)) ^ (1 + α) = B ^ (((n : ℝ) / α) * (1 + α)) := by
          rw [← Real.rpow_mul (le_of_lt hBpos)]
        rw [hBpow]
        have hBcombine :
            B ^ (((n + 1 : ℕ) : ℝ) / α) * B ^ (n : ℝ) =
              B ^ ((1 : ℝ) / α) * B ^ (((n : ℝ) / α) * (1 + α)) := by
          rw [← Real.rpow_add hBpos, ← Real.rpow_add hBpos]
          congr 1
          rw [show (((n + 1 : ℕ) : ℝ) = (n : ℝ) + 1) by norm_num]
          field_simp [hα.ne']
          ring
        calc
          B ^ (((n + 1 : ℕ) : ℝ) / α) * (C_rec * B ^ (n : ℝ) * Y n ^ (1 + α))
              = C_rec * (B ^ (((n + 1 : ℕ) : ℝ) / α) * B ^ (n : ℝ)) * Y n ^ (1 + α) := by ring
          _ = C_rec * (B ^ ((1 : ℝ) / α) * B ^ (((n : ℝ) / α) * (1 + α))) * Y n ^ (1 + α) := by
            rw [hBcombine]
          _ = C_rec * B ^ ((1 : ℝ) / α) * (B ^ (((n : ℝ) / α) * (1 + α)) * Y n ^ (1 + α)) := by
            ring
  have hZ_rec : ∀ n, Z (n + 1) ≤ Z n ^ (1 + α) := by
    intro n
    calc
      Z (n + 1) = D * W (n + 1) := by
        simp [Z]
      _ ≤ D * (X * W n ^ (1 + α)) := by
        gcongr
        exact hW_rec n
      _ = D * X * W n ^ (1 + α) := by ring
      _ = D * D ^ α * W n ^ (1 + α) := by rw [hDα]
      _ = D ^ (1 + α) * W n ^ (1 + α) := by
        have hmul : D * D ^ α = D ^ (1 + α) := by
          calc
            D * D ^ α = D ^ (1 : ℝ) * D ^ α := by rw [Real.rpow_one]
            _ = D ^ (1 + α) := by rw [← Real.rpow_add hD_pos]
        rw [hmul]
      _ = (D * W n) ^ (1 + α) := by
        rw [← Real.mul_rpow hD_nonneg (hW_nonneg n)]
      _ = Z n ^ (1 + α) := by
        simp [Z]
  have hZ_nonneg : ∀ n, 0 ≤ Z n := by
    intro n
    dsimp [Z]
    exact mul_nonneg hD_nonneg (hW_nonneg n)
  have hZ_le_one : ∀ n, Z n ≤ 1 := by
    intro n
    induction n with
    | zero =>
        calc
          Z 0 = D * Y 0 := by
            simp [Z, W]
          _ ≤ D * D⁻¹ := by
            gcongr
          _ = 1 := by
            field_simp [hD_pos.ne']
    | succ n ihn =>
        have hpow : Z n ^ (1 + α) ≤ 1 := by
          exact Real.rpow_le_one (hZ_nonneg n) ihn (by positivity)
        exact (hZ_rec n).trans hpow
  have hY_bound : ∀ n, Y n ≤ D⁻¹ * q ^ n := by
    intro n
    have hW_bound : W n ≤ D⁻¹ := by
      rw [inv_eq_one_div]
      refine (le_div_iff₀ hD_pos).2 ?_
      simpa [Z, mul_assoc, mul_comm, mul_left_comm] using hZ_le_one n
    have hY_eq : Y n = B ^ (-(n : ℝ) / α) * W n := by
      dsimp [W]
      have hcancel : B ^ (-(n : ℝ) / α) * B ^ ((n : ℝ) / α) = 1 := by
        rw [← Real.rpow_add hBpos]
        rw [show -(n : ℝ) / α + (n : ℝ) / α = 0 by ring, Real.rpow_zero]
      calc
        Y n = 1 * Y n := by ring
        _ = (B ^ (-(n : ℝ) / α) * B ^ ((n : ℝ) / α)) * Y n := by rw [hcancel]
        _ = B ^ (-(n : ℝ) / α) * W n := by ring
    calc
      Y n = B ^ (-(n : ℝ) / α) * W n := hY_eq
      _ ≤ B ^ (-(n : ℝ) / α) * D⁻¹ := by
        gcongr
      _ = D⁻¹ * B ^ (-(n : ℝ) / α) := by ring
      _ = D⁻¹ * q ^ n := by
        congr 1
        dsimp [q]
        rw [show -(n : ℝ) / α = (-(1 : ℝ) / α) * n by ring]
        rw [Real.rpow_mul (le_of_lt hBpos), Real.rpow_natCast]
  have hgeom :
      Filter.Tendsto (fun n : ℕ => D⁻¹ * q ^ n) Filter.atTop (nhds 0) := by
    simpa using (tendsto_pow_atTop_nhds_zero_of_lt_one hq_nonneg hq_lt_one).const_mul D⁻¹
  intro ε hε
  rw [Metric.tendsto_atTop] at hgeom
  obtain ⟨N, hN⟩ := hgeom ε hε
  refine ⟨N, ?_⟩
  intro n hn
  have hgeom_lt : D⁻¹ * q ^ n < ε := by
    have := hN n hn
    rwa [Real.dist_eq, sub_zero, abs_of_nonneg (by positivity)] at this
  exact lt_of_le_of_lt (hY_bound n) hgeom_lt

\end{lstlisting}
\begin{proof}
Set $X := C B^{1/\alpha}$, $D := X^{1/\alpha}$, $q :=
B^{-1/\alpha} < 1$, and define $W_n := B^{n/\alpha} Y_n$ and $Z_n :=
D W_n$. The hypothesis on $Y_0$ is equivalent to $Y_0 \le D^{-1}$, so
$Z_0 \le 1$. From the recurrence
\[
  W_{n+1} \le X \, W_n^{1+\alpha},
\]
multiply by $D$ and use $D X = D^{1+\alpha}$ to get $Z_{n+1} \le
Z_n^{1+\alpha}$, hence by induction $Z_n \le 1$ for all $n$. This
yields $W_n \le D^{-1}$ and consequently
$Y_n \le D^{-1} q^n$, which tends to $0$ as $n \to \infty$.
\end{proof}

\begin{theorem}[Recurrence rewriting,
\leanname{deGiorgi\_preiter\_to\_recurrence}]
If a sequence $A_n$ satisfies, for all $n$,
\begin{multline*}
  A_{n+1} \le
   \frac{K}{(r_n - r_{n+1})^2 (\ell_{n+1} - \ell_n)^{4/d}}
   \cdot A_n^{1 + 2/d},
\end{multline*}
then with $C_d(K,\ell^*)$ and $B_d$ as above one has
\[
  A_{n+1} \le C_d(K, \ell^*)\, B_d^n\, A_n^{1 + 2/d}.
\]
\end{theorem}

\begin{lstlisting}[style=Matlab-editor,basicstyle=\mlttfamily,escapechar=",mlshowsectionrules=true,mathescape=true,morecomment={[s]{\%\{}{\%\}}},language=lean,numbers=none]
/-- Rewrite the one-step estimate on canonical radii and levels into the
standard nonlinear recurrence. -/
theorem deGiorgi_preiter_to_recurrence
    {Aseq : ℕ → ℝ} {d : ℕ}
    {K lamStar : ℝ} (hLamStar : 0 < lamStar)
    (hpre :
      ∀ n,
        Aseq (n + 1) ≤
          K / ((deGiorgiRadius n - deGiorgiRadius (n + 1)) ^ 2 *
            (deGiorgiLevel lamStar (n + 1) - deGiorgiLevel lamStar n) ^ (4 / (d : ℝ))) *
          Aseq n ^ (1 + 2 / (d : ℝ))) :
    ∀ n,
      Aseq (n + 1) ≤
        deGiorgiRecurrenceCoeff d K lamStar *
          deGiorgiRecurrenceBase d ^ n *
          Aseq n ^ (1 + 2 / (d : ℝ)) := by
  intro n
  have htwo_pos : 0 < (2 : ℝ) := by positivity
  have hhalf : (1 / 2 : ℝ) = (2 : ℝ) ^ (-(1 : ℝ)) := by
    rw [Real.rpow_neg (by positivity), Real.rpow_one]
    norm_num
  have hrad :
      (deGiorgiRadius n - deGiorgiRadius (n + 1)) ^ 2 =
        (2 : ℝ) ^ (-(2 * n + 6 : ℝ)) := by
    rw [deGiorgiRadius_gap, pow_two, ← pow_add]
    have hnat : n + 3 + (n + 3) = 2 * n + 6 := by omega
    rw [hnat, ← Real.rpow_natCast]
    rw [hhalf]
    rw [← Real.rpow_mul htwo_pos.le]
    rw [show (-(1 : ℝ)) * (((2 * n + 6 : ℕ) : ℝ)) = -(2 * (n : ℝ) + 6) by
      norm_num [Nat.cast_add, Nat.cast_mul]]
  have htail_rpow :
      ((1 / 2 : ℝ) ^ (n + 2)) ^ (4 / (d : ℝ)) =
        (2 : ℝ) ^ (-((n + 2 : ℝ) * (4 / (d : ℝ)))) := by
    rw [← Real.rpow_natCast]
    rw [hhalf]
    rw [show (((n + 2 : ℕ) : ℝ) = (n : ℝ) + 2) by norm_num]
    calc
      (((2 : ℝ) ^ (-(1 : ℝ))) ^ ((n : ℝ) + 2)) ^ (4 / (d : ℝ))
          = ((2 : ℝ) ^ ((-(1 : ℝ)) * ((n : ℝ) + 2))) ^ (4 / (d : ℝ)) := by
              rw [← Real.rpow_mul htwo_pos.le]
      _ = (2 : ℝ) ^ (((-(1 : ℝ)) * ((n + 2 : ℝ))) * (4 / (d : ℝ))) := by
            rw [← Real.rpow_mul htwo_pos.le]
      _ = (2 : ℝ) ^ (-((n + 2 : ℝ) * (4 / (d : ℝ)))) := by
            congr 1
            ring
  have hlev :
      (deGiorgiLevel lamStar (n + 1) - deGiorgiLevel lamStar n) ^ (4 / (d : ℝ)) =
        lamStar ^ (4 / (d : ℝ)) * (2 : ℝ) ^ (-((n + 2 : ℝ) * (4 / (d : ℝ)))) := by
    rw [deGiorgiLevel_gap]
    rw [Real.mul_rpow hLamStar.le (by positivity)]
    rw [htail_rpow]
  have hden :
      (deGiorgiRadius n - deGiorgiRadius (n + 1)) ^ 2 *
          (deGiorgiLevel lamStar (n + 1) - deGiorgiLevel lamStar n) ^ (4 / (d : ℝ)) =
        lamStar ^ (4 / (d : ℝ)) *
          (2 : ℝ) ^ (-(6 + 8 / (d : ℝ) + n * (2 + 4 / (d : ℝ)))) := by
    rw [hrad, hlev]
    calc
      (2 : ℝ) ^ (-(2 * n + 6 : ℝ)) * (lamStar ^ (4 / (d : ℝ)) *
          (2 : ℝ) ^ (-((n + 2 : ℝ) * (4 / (d : ℝ)))))
          = lamStar ^ (4 / (d : ℝ)) *
              ((2 : ℝ) ^ (-(2 * n + 6 : ℝ)) *
                (2 : ℝ) ^ (-((n + 2 : ℝ) * (4 / (d : ℝ))))) := by
                ring
      _ = lamStar ^ (4 / (d : ℝ)) *
            (2 : ℝ) ^ (-(2 * n + 6 : ℝ) + -((n + 2 : ℝ) * (4 / (d : ℝ)))) := by
              rw [← Real.rpow_add htwo_pos]
      _ = lamStar ^ (4 / (d : ℝ)) *
            (2 : ℝ) ^ (-(6 + 8 / (d : ℝ) + n * (2 + 4 / (d : ℝ)))) := by
              congr 1
              ring_nf
  have hrec := hpre n
  calc
    Aseq (n + 1) ≤
      K / ((deGiorgiRadius n - deGiorgiRadius (n + 1)) ^ 2 *
        (deGiorgiLevel lamStar (n + 1) - deGiorgiLevel lamStar n) ^ (4 / (d : ℝ))) *
        Aseq n ^ (1 + 2 / (d : ℝ)) := hrec
    _ =
      deGiorgiRecurrenceCoeff d K lamStar *
        deGiorgiRecurrenceBase d ^ n *
        Aseq n ^ (1 + 2 / (d : ℝ)) := by
      rw [hden]
      rw [deGiorgiRecurrenceCoeff, deGiorgiRecurrenceBase]
      rw [div_eq_mul_inv]
      rw [show (lamStar ^ (4 / (d : ℝ)) *
          (2 : ℝ) ^ (-(6 + 8 / (d : ℝ) + n * (2 + 4 / (d : ℝ)))))⁻¹ =
          ((2 : ℝ) ^ (-(6 + 8 / (d : ℝ) + n * (2 + 4 / (d : ℝ)))))⁻¹ *
            (lamStar ^ (4 / (d : ℝ)))⁻¹ by rw [_root_.mul_inv_rev]]
      rw [Real.rpow_neg htwo_pos.le, Real.rpow_neg hLamStar.le]
      simp only [inv_inv]
      have hsplitpow :
          (2 : ℝ) ^ (6 + 8 / (d : ℝ) + n * (2 + 4 / (d : ℝ))) =
            (2 : ℝ) ^ (6 + 8 / (d : ℝ)) * (2 : ℝ) ^ (n * (2 + 4 / (d : ℝ))) := by
        rw [← Real.rpow_add htwo_pos]
      rw [hsplitpow]
      rw [show n * (2 + 4 / (d : ℝ)) = (2 + 4 / (d : ℝ)) * n by ring]
      rw [Real.rpow_mul htwo_pos.le, Real.rpow_natCast]
      ac_rfl

\end{lstlisting}
\begin{proof}
By the gap identities,
$(r_n - r_{n+1})^2 = 2^{-(2n+6)}$ and
$(\ell_{n+1}-\ell_n)^{4/d} = (\ell^*)^{4/d}\, 2^{-(n+2)\cdot 4/d}$,
so the denominator equals
$(\ell^*)^{4/d}\, 2^{-(6+8/d + n(2+4/d))}$. Inverting and rearranging
yields the claimed form.
\end{proof}

\begin{theorem}[Vanishing of the canonical recurrence,
\leanname{deGiorgi\_preiter\_vanishing}]
Under $d > 2$, $K > 0$, $\ell^* > 0$, the canonical recurrence above,
and the smallness condition
\[
  A_0 \le C_d(K,\ell^*)^{-d/2}\, B_d^{-(d/2)^2 / 2},
\]
the sequence $A_n$ tends to zero.
\end{theorem}

\begin{lstlisting}[style=Matlab-editor,basicstyle=\mlttfamily,escapechar=",mlshowsectionrules=true,mathescape=true,morecomment={[s]{\%\{}{\%\}}},language=lean,numbers=none]
/-- The canonical De Giorgi recurrence tends to zero under the standard
smallness threshold. This is the non-PDE closeout for the De Giorgi
iteration. -/
theorem deGiorgi_preiter_vanishing
    {Aseq : ℕ → ℝ} {d : ℕ}
    (hd : 2 < (d : ℝ))
    {K lamStar : ℝ} (hK : 0 < K) (hLamStar : 0 < lamStar)
    (hA_nonneg : ∀ n, 0 ≤ Aseq n)
    (hpre :
      ∀ n,
        Aseq (n + 1) ≤
          K / ((deGiorgiRadius n - deGiorgiRadius (n + 1)) ^ 2 *
            (deGiorgiLevel lamStar (n + 1) - deGiorgiLevel lamStar n) ^ (4 / (d : ℝ))) *
          Aseq n ^ (1 + 2 / (d : ℝ)))
    (hsmall :
      Aseq 0 ≤
        (deGiorgiRecurrenceCoeff d K lamStar) ^ (-(1 : ℝ) / (2 / (d : ℝ))) *
          (deGiorgiRecurrenceBase d) ^ (-(1 : ℝ) / (2 / (d : ℝ)) ^ 2)) :
    Tendsto Aseq atTop (nhds 0) := by
  have hα : 0 < 2 / (d : ℝ) := by positivity
  have hB : 1 < deGiorgiRecurrenceBase d := by
    dsimp [deGiorgiRecurrenceBase]
    exact (Real.one_lt_rpow_iff_of_pos (by positivity)).2 (Or.inl ⟨by norm_num, by positivity⟩)
  have hC : 0 < deGiorgiRecurrenceCoeff d K lamStar := by
    dsimp [deGiorgiRecurrenceCoeff]
    positivity
  have hrec :
      ∀ n,
        Aseq (n + 1) ≤
          deGiorgiRecurrenceCoeff d K lamStar *
            deGiorgiRecurrenceBase d ^ n *
            Aseq n ^ (1 + 2 / (d : ℝ)) :=
    deGiorgi_preiter_to_recurrence hLamStar hpre
  rw [Metric.tendsto_atTop]
  intro ε hε
  obtain ⟨N, hN⟩ :=
    deGiorgi_recurrence_closeout hC hB hα hA_nonneg hrec hsmall ε hε
  exact ⟨N, fun n hn => by
    rw [Real.dist_eq, sub_zero, abs_of_nonneg (hA_nonneg n)]
    exact hN n hn⟩

\end{lstlisting}
\begin{proof}
The pre-iteration recurrence is rewritten using
\leanname{deGiorgi\_preiter\_to\_recurrence} into the standard form with
$B := B_d$, $C := C_d(K,\ell^*)$, $\alpha := 2/d$. The smallness
hypothesis is exactly the threshold required by
\leanname{deGiorgi\_recurrence\_closeout}, so $A_n \to 0$.
\end{proof}

\subsection{The De Giorgi $L^\infty$ bound}

The next layer assembles the $L^\infty$ bound on the half-ball from
the recurrence machinery.

\begin{lemma}[\leanname{ae\_eq\_zero\_of\_forall\_setIntegral\_sq\_le\_of\_tendsto\_zero}]
Let $f : E \to \R$ with $|f|^2 \in L^1(s)$ and let $Y_n \to 0$ be a
sequence of upper bounds for $\int_s |f|^2$. Then $f = 0$ a.e. on
$s$.
\end{lemma}

\begin{lstlisting}[style=Matlab-editor,basicstyle=\mlttfamily,escapechar=",mlshowsectionrules=true,mathescape=true,morecomment={[s]{\%\{}{\%\}}},language=lean,numbers=none]
omit [NeZero d] in
theorem ae_eq_zero_of_forall_setIntegral_sq_le_of_tendsto_zero
    {f : E → ℝ} {s : Set E} {Y : ℕ → ℝ}
    (hfi : IntegrableOn (fun x => |f x| ^ 2) s volume)
    (hbound : ∀ n, ∫ x in s, |f x| ^ 2 ∂volume ≤ Y n)
    (hY_tendsto : Tendsto Y atTop (nhds 0)) :
    ∀ᵐ x ∂(volume.restrict s), f x = 0 := by
  let I : ℝ := ∫ x in s, |f x| ^ 2 ∂volume
  have hI_nonneg : 0 ≤ I := by
    dsimp [I]
    refine integral_nonneg ?_
    intro x
    positivity
  have hI_zero : I = 0 := by
    by_contra hI_ne
    have hI_pos : 0 < I := lt_of_le_of_ne hI_nonneg (Ne.symm hI_ne)
    rw [Metric.tendsto_atTop] at hY_tendsto
    obtain ⟨N, hN⟩ := hY_tendsto (I / 2) (by linarith)
    have hYN_nonneg : 0 ≤ Y N := hI_nonneg.trans (hbound N)
    have hYN_lt : Y N < I / 2 := by
      have hdist := hN N le_rfl
      rwa [Real.dist_eq, sub_zero, abs_of_nonneg hYN_nonneg] at hdist
    linarith [hbound N]
  have hsq_zero :
      (fun x => |f x| ^ 2) =ᵐ[volume.restrict s] 0 := by
    refine (integral_eq_zero_iff_of_nonneg_ae
      (Filter.Eventually.of_forall (fun x => by positivity)) hfi).1 ?_
    simpa [I] using hI_zero
  filter_upwards [hsq_zero] with x hx
  have habs_zero : |f x| = 0 := by
    exact sq_eq_zero_iff.mp hx
  exact abs_eq_zero.mp habs_zero

\end{lstlisting}
\begin{lemma}[\leanname{integral\_sq\_halfBall\_le\_deGiorgiEnergySeq}]
For $\ell^* > 0$, the half-ball energy is monotone in the iteration:
\[
  \int_{\Ball{1/2}{x_0}} |(u-\ell^*)_+|^2 \, dx
   \le A_n(u, x_0, \ell^*),
\qquad \forall n.
\]
\end{lemma}

\begin{lstlisting}[style=Matlab-editor,basicstyle=\mlttfamily,escapechar=",mlshowsectionrules=true,mathescape=true,morecomment={[s]{\%\{}{\%\}}},language=lean,numbers=none]
omit [NeZero d] in
theorem integral_sq_halfBall_le_deGiorgiEnergySeq
    {u : E → ℝ} {x₀ : E} {lamStar : ℝ}
    (hLamStar : 0 < lamStar)
    (hStarInt :
      IntegrableOn (fun x => |positivePartSub u lamStar x| ^ 2)
        (Metric.ball x₀ (1 / 2 : ℝ)) volume)
    (hAint :
      ∀ n,
        IntegrableOn (fun x => |positivePartSub u (deGiorgiLevel lamStar n) x| ^ 2)
          (Metric.ball x₀ (deGiorgiRadius n)) volume) :
    ∀ n,
      ∫ x in Metric.ball x₀ (1 / 2 : ℝ), |positivePartSub u lamStar x| ^ 2 ∂volume ≤
        deGiorgiEnergySeq u x₀ lamStar n := by
  intro n
  have hs : Metric.ball x₀ (1 / 2 : ℝ) ⊆ Metric.ball x₀ (deGiorgiRadius n) := by
    exact Metric.ball_subset_ball (le_of_lt (deGiorgiRadius_gt_half n))
  have hbig_int_half :
      IntegrableOn (fun x => |positivePartSub u (deGiorgiLevel lamStar n) x| ^ 2)
        (Metric.ball x₀ (1 / 2 : ℝ)) volume :=
    (hAint n).mono_set hs
  have hpointwise :
      ∀ x ∈ Metric.ball x₀ (1 / 2 : ℝ),
        |positivePartSub u lamStar x| ^ 2 ≤
          |positivePartSub u (deGiorgiLevel lamStar n) x| ^ 2 := by
    intro x hx
    have hmono :
        positivePartSub u lamStar x ≤ positivePartSub u (deGiorgiLevel lamStar n) x := by
      exact positivePartSub_antitone (show deGiorgiLevel lamStar n ≤ lamStar by
        linarith [deGiorgiLevel_lt hLamStar n]) x
    have hsmall_nonneg : 0 ≤ positivePartSub u lamStar x := positivePartSub_nonneg
    have hbig_nonneg : 0 ≤ positivePartSub u (deGiorgiLevel lamStar n) x :=
      positivePartSub_nonneg
    have hsq :
        (positivePartSub u lamStar x) ^ 2 ≤
          (positivePartSub u (deGiorgiLevel lamStar n) x) ^ 2 := by
      nlinarith
    simpa [abs_of_nonneg hsmall_nonneg, abs_of_nonneg hbig_nonneg] using hsq
  have hmono_same_set :
      ∫ x in Metric.ball x₀ (1 / 2 : ℝ), |positivePartSub u lamStar x| ^ 2 ∂volume ≤
        ∫ x in Metric.ball x₀ (1 / 2 : ℝ),
          |positivePartSub u (deGiorgiLevel lamStar n) x| ^ 2 ∂volume := by
    refine integral_mono_ae hStarInt hbig_int_half ?_
    refine (ae_restrict_iff' (μ := volume) isOpen_ball.measurableSet).2 ?_
    filter_upwards with x hx
    exact hpointwise x hx
  have hnonneg :
      ∀ x, 0 ≤ |positivePartSub u (deGiorgiLevel lamStar n) x| ^ 2 := by
    intro x
    positivity
  have hmono_set :
      ∫ x in Metric.ball x₀ (1 / 2 : ℝ),
          |positivePartSub u (deGiorgiLevel lamStar n) x| ^ 2 ∂volume ≤
        ∫ x in Metric.ball x₀ (deGiorgiRadius n),
          |positivePartSub u (deGiorgiLevel lamStar n) x| ^ 2 ∂volume := by
    exact setIntegral_mono_set (hAint n) (ae_of_all _ hnonneg) (ae_of_all _ hs)
  calc
    ∫ x in Metric.ball x₀ (1 / 2 : ℝ), |positivePartSub u lamStar x| ^ 2 ∂volume
        ≤ ∫ x in Metric.ball x₀ (1 / 2 : ℝ),
            |positivePartSub u (deGiorgiLevel lamStar n) x| ^ 2 ∂volume := hmono_same_set
    _ ≤ ∫ x in Metric.ball x₀ (deGiorgiRadius n),
          |positivePartSub u (deGiorgiLevel lamStar n) x| ^ 2 ∂volume := hmono_set
    _ = deGiorgiEnergySeq u x₀ lamStar n := by
          simp [deGiorgiEnergySeq]

\end{lstlisting}
\begin{lemma}[\leanname{deGiorgiEnergySeq\_zero\_le\_integral\_sq\_posPartZero\_on\_unitBall}]
The first term satisfies
$A_0 \le \int_{\Bz{1}} (u_+)^2 \, dx$.
\end{lemma}

\begin{lstlisting}[style=Matlab-editor,basicstyle=\mlttfamily,escapechar=",mlshowsectionrules=true,mathescape=true,morecomment={[s]{\%\{}{\%\}}},language=lean,numbers=none]
omit [NeZero d] in
theorem deGiorgiEnergySeq_zero_le_integral_sq_posPartZero_on_unitBall
    {u : E → ℝ} {x₀ : E} {lamStar : ℝ}
    (hLamStar : 0 < lamStar)
    (hA0Int :
      IntegrableOn (fun x => |positivePartSub u (deGiorgiLevel lamStar 0) x| ^ 2)
        (Metric.ball x₀ (deGiorgiRadius 0)) volume)
    (hposInt :
      IntegrableOn (fun x => |max (u x) 0| ^ 2)
        (Metric.ball x₀ 1) volume) :
    deGiorgiEnergySeq u x₀ lamStar 0 ≤
      ∫ x in Metric.ball x₀ 1, |max (u x) 0| ^ 2 ∂volume := by
  have hs : Metric.ball x₀ (deGiorgiRadius 0) ⊆ Metric.ball x₀ 1 := by
    exact Metric.ball_subset_ball (le_of_lt (deGiorgiRadius_lt_one 0))
  have hlev_nonneg : 0 ≤ deGiorgiLevel lamStar 0 := by
    dsimp [deGiorgiLevel, deGiorgiTail]
    nlinarith
  have hpointwise :
      ∀ x ∈ Metric.ball x₀ (deGiorgiRadius 0),
        |positivePartSub u (deGiorgiLevel lamStar 0) x| ^ 2 ≤ |max (u x) 0| ^ 2 := by
    intro x hx
    have hmono : positivePartSub u (deGiorgiLevel lamStar 0) x ≤ max (u x) 0 := by
      exact positivePartSub_le_posPartZero hlev_nonneg x
    have hsmall_nonneg : 0 ≤ positivePartSub u (deGiorgiLevel lamStar 0) x :=
      positivePartSub_nonneg
    have hbig_nonneg : 0 ≤ max (u x) 0 := le_max_right _ _
    have hsq :
        (positivePartSub u (deGiorgiLevel lamStar 0) x) ^ 2 ≤
          (max (u x) 0) ^ 2 := by
      nlinarith
    simpa [abs_of_nonneg hsmall_nonneg, abs_of_nonneg hbig_nonneg] using hsq
  have hbig_int_small :
      IntegrableOn (fun x => |max (u x) 0| ^ 2)
        (Metric.ball x₀ (deGiorgiRadius 0)) volume :=
    hposInt.mono_set hs
  have hmono_same_set :
      ∫ x in Metric.ball x₀ (deGiorgiRadius 0),
          |positivePartSub u (deGiorgiLevel lamStar 0) x| ^ 2 ∂volume ≤
        ∫ x in Metric.ball x₀ (deGiorgiRadius 0), |max (u x) 0| ^ 2 ∂volume := by
    refine integral_mono_ae hA0Int hbig_int_small ?_
    · refine (ae_restrict_iff' (μ := volume) isOpen_ball.measurableSet).2 ?_
      filter_upwards with x hx
      exact hpointwise x hx
  have hnonneg :
      ∀ x, 0 ≤ |max (u x) 0| ^ 2 := by
    intro x
    positivity
  have hmono_set :
      ∫ x in Metric.ball x₀ (deGiorgiRadius 0), |max (u x) 0| ^ 2 ∂volume ≤
        ∫ x in Metric.ball x₀ 1, |max (u x) 0| ^ 2 ∂volume := by
    exact setIntegral_mono_set hposInt (ae_of_all _ hnonneg) (ae_of_all _ hs)
  calc
    deGiorgiEnergySeq u x₀ lamStar 0
        = ∫ x in Metric.ball x₀ (deGiorgiRadius 0),
            |positivePartSub u (deGiorgiLevel lamStar 0) x| ^ 2 ∂volume := by
              simp [deGiorgiEnergySeq]
    _ ≤ ∫ x in Metric.ball x₀ (deGiorgiRadius 0), |max (u x) 0| ^ 2 ∂volume := hmono_same_set
    _ ≤ ∫ x in Metric.ball x₀ 1, |max (u x) 0| ^ 2 ∂volume := hmono_set

\end{lstlisting}
\begin{definition}[\leanname{C\_DeGiorgiSmallness}]
For $K > 0$ and $d \ge 1$, set
\[
  C_{\mathrm{small}}(d, K) :=
   \bigl(K \cdot 2^{6 + 8/d}\bigr)^{d/4}
   \cdot B_d^{d^2/8}.
\]
\end{definition}

\begin{lstlisting}[style=Matlab-editor,basicstyle=\mlttfamily,escapechar=",mlshowsectionrules=true,mathescape=true,morecomment={[s]{\%\{}{\%\}}},language=lean,numbers=none]
/-- Height constant for converting a unit-ball `L²` smallness condition into
the abstract initial De Giorgi recurrence threshold. This still depends on the
one-step coefficient `K`; later PDE-facing theorems will substitute the
specific `K` coming from pre-iteration. -/
noncomputable def C_DeGiorgiSmallness (d : ℕ) (K : ℝ) : ℝ :=
  (K * (2 : ℝ) ^ (6 + 8 / (d : ℝ))) ^ ((d : ℝ) / 4) *
    (deGiorgiRecurrenceBase d) ^ (((d : ℝ) ^ 2) / 8)

\end{lstlisting}
This is the height constant required to convert a unit-ball $L^2$
smallness assumption into the abstract initial recurrence threshold.
We collect three monotonicity/scaling identities used later:
\leanname{C\_DeGiorgiSmallness\_pos},
\leanname{C\_DeGiorgiSmallness\_mono} (monotonicity in $K$), and
\leanname{C\_DeGiorgiSmallness\_mul\_right}:
\[
  C_{\mathrm{small}}(d, K t) =
    C_{\mathrm{small}}(d, K) \, t^{d/4} \quad (K, t \ge 0).
\]

\begin{definition}[Iteration constants for subsolutions,
\leanname{K\_DeGiorgi\_subsolution} and \leanname{C\_DeGiorgi\_subsolution}]
For an elliptic coefficient field $A$ on $\Bz{1}$ set
\[
  K_{\mathrm{DG}}(d, A) := (C_{\mathrm{gns}}(d, 2))^2 \cdot
    8\bigl(\Lambda/\lambda + 1\bigr) \cdot (2 M_{\mathrm{st}})^2 + 1,
\]
\[
  C_{\mathrm{DG}}(d, A) := C_{\mathrm{small}}(d, K_{\mathrm{DG}}(d, A)).
\]
Both are positive (\leanname{K\_DeGiorgi\_subsolution\_pos},
\leanname{C\_DeGiorgi\_subsolution\_pos}). The constant
$M_{\mathrm{st}}$ is the standard cutoff gradient bound from the
Stampacchia layer.
\end{definition}

\begin{lstlisting}[style=Matlab-editor,basicstyle=\mlttfamily,escapechar=",mlshowsectionrules=true,mathescape=true,morecomment={[s]{\%\{}{\%\}}},language=lean,numbers=none]
/-- Dimension-only coefficient in the unit-ball De Giorgi iteration package.

It packages the assembly step from `IsSubsolution` to the concrete
De Giorgi recurrence data used by the closeout layer: a positive one-step
coefficient `K`, integrability of the truncated energies on the canonical balls,
and the nonlinear recurrence itself. -/
noncomputable def K_DeGiorgi_subsolution
    (d : ℕ) [NeZero d]
    (A : EllipticCoeff d (Metric.ball (0 : AmbientSpace d) 1)) : ℝ :=
  (C_gns d 2) ^ 2 * (8 * (ellipticityRatio A + 1) * (2 * (Mst : ℝ)) ^ 2) + 1

noncomputable def C_DeGiorgi_subsolution
    (d : ℕ) [NeZero d]
    (A : EllipticCoeff d (Metric.ball (0 : AmbientSpace d) 1)) : ℝ :=
  C_DeGiorgiSmallness d (K_DeGiorgi_subsolution d A)

theorem K_DeGiorgi_subsolution_pos
    (A : EllipticCoeff d (Metric.ball (0 : E) 1)) :
    0 < K_DeGiorgi_subsolution d A := by
  rw [K_DeGiorgi_subsolution]
  have hnonneg :
      0 ≤ (C_gns d 2) ^ 2 * (8 * (ellipticityRatio A + 1) * (2 * (Mst : ℝ)) ^ 2) := by
    have hratio_nonneg : 0 ≤ ellipticityRatio A + 1 := by
      linarith [A.ellipticityRatio_nonneg]
    have htmp_nonneg : 0 ≤ (ellipticityRatio A + 1) * (2 * (Mst : ℝ)) ^ 2 := by
      exact mul_nonneg hratio_nonneg (sq_nonneg _)
    have hfactor_nonneg : 0 ≤ 8 * (ellipticityRatio A + 1) * (2 * (Mst : ℝ)) ^ 2 := by
      have h8_nonneg : (0 : ℝ) ≤ 8 := by norm_num
      simpa [mul_assoc] using mul_nonneg h8_nonneg htmp_nonneg
    exact mul_nonneg (sq_nonneg _) hfactor_nonneg
  linarith

theorem C_DeGiorgi_subsolution_pos
    (A : EllipticCoeff d (Metric.ball (0 : E) 1)) :
    0 < C_DeGiorgi_subsolution d A := by
  rw [C_DeGiorgi_subsolution]
  exact C_DeGiorgiSmallness_pos (d := d) (K_DeGiorgi_subsolution_pos (d := d) A)

\end{lstlisting}
\begin{theorem}[Iteration package,
\leanname{deGiorgi\_iteration\_package\_on\_unitBall\_of\_subsolution}]
Let $d > 2$, $A$ elliptic on $\Bz{1}$, and $u$ a subsolution. Then
$K_{\mathrm{DG}}(d,A) > 0$ and for every $\ell^* > 0$:
\begin{enumerate}
\item $|(u-\ell^*)_+|^2$ is integrable on $\Ball{1/2}{0}$;
\item for every $n \in \N$, $|(u-\ell_n(\ell^*))_+|^2$ is integrable
on $\Ball{r_n}{0}$;
\item the canonical pre-iteration recurrence holds:
\begin{multline*}
  A_{n+1}(u,0,\ell^*) \le \\
   \frac{K_{\mathrm{DG}}(d,A)}{(r_n - r_{n+1})^2 \,
     (\ell_{n+1}(\ell^*)-\ell_n(\ell^*))^{4/d}} \cdot
   A_n(u,0,\ell^*)^{1 + 2/d}.
\end{multline*}
\end{enumerate}
\end{theorem}

\begin{lstlisting}[style=Matlab-editor,basicstyle=\mlttfamily,escapechar=",mlshowsectionrules=true,mathescape=true,morecomment={[s]{\%\{}{\%\}}},language=lean,numbers=none]
theorem deGiorgi_iteration_package_on_unitBall_of_subsolution
    (hd : 2 < (d : ℝ))
    (A : EllipticCoeff d (Metric.ball (0 : E) 1))
    {u : E → ℝ}
    (hsub : IsSubsolution A u) :
    0 < K_DeGiorgi_subsolution d A ∧
      ∀ {lamStar : ℝ}, 0 < lamStar →
        IntegrableOn (fun x => |positivePartSub u lamStar x| ^ 2)
          (Metric.ball (0 : E) (1 / 2 : ℝ)) volume ∧
        (∀ n,
          IntegrableOn (fun x => |positivePartSub u (deGiorgiLevel lamStar n) x| ^ 2)
            (Metric.ball (0 : E) (deGiorgiRadius n)) volume) ∧
        (∀ n,
          deGiorgiEnergySeq u (0 : E) lamStar (n + 1) ≤
            K_DeGiorgi_subsolution d A /
              ((deGiorgiRadius n - deGiorgiRadius (n + 1)) ^ 2 *
                (deGiorgiLevel lamStar (n + 1) - deGiorgiLevel lamStar n) ^ (4 / (d : ℝ))) *
              deGiorgiEnergySeq u (0 : E) lamStar n ^ (1 + 2 / (d : ℝ))) := by
  let hu1 : MemW1pWitness 2 u (Metric.ball (0 : E) 1) :=
    DeGiorgi.MemW1p.someWitness (isSubsolution_memW1p hsub)
  haveI : IsFiniteMeasure (volume.restrict (Metric.ball (0 : E) 1)) := by
    rw [isFiniteMeasure_restrict]
    exact measure_ball_lt_top.ne
  have happroxUnit :
      ∀ θ : ℝ,
        ∃ ψ : ℕ → E → ℝ,
          (∀ n, ContDiff ℝ 1 (ψ n)) ∧
          (∀ n, HasCompactSupport (ψ n)) ∧
          Tendsto
            (fun n =>
              eLpNorm (fun x => ψ n x - (u x - θ)) 2
                (volume.restrict (Metric.ball (0 : E) 1)))
            atTop (nhds 0) ∧
          (∀ i : Fin d,
            Tendsto
              (fun n =>
                eLpNorm
                  (fun x =>
                    (fderiv ℝ (ψ n) x) (EuclideanSpace.single i 1) - hu1.weakGrad x i)
                  2 (volume.restrict (Metric.ball (0 : E) 1)))
              atTop (nhds 0)) := by
    intro θ
    let huShift : MemW1pWitness 2 (fun x => u x - θ) (Metric.ball (0 : E) 1) :=
      hu1.sub_const isOpen_ball θ
    rcases exists_smooth_W12_approx_on_unitBall (d := d) huShift with
      ⟨ψ, hψ_smooth, hψ_compact, hψ_fun, hψ_grad⟩
    refine ⟨ψ, ?_, hψ_compact, hψ_fun, ?_⟩
    · intro n
      exact (hψ_smooth n).of_le (by norm_num)
    · intro i
      simpa [huShift] using hψ_grad i
  let Kraw : ℝ :=
    (C_gns d 2) ^ 2 * (8 * (ellipticityRatio A + 1) * (2 * (Mst : ℝ)) ^ 2)
  let K : ℝ := K_DeGiorgi_subsolution d A
  have hK_eq : K = Kraw + 1 := by
    dsimp [K, Kraw, K_DeGiorgi_subsolution]
  have hKraw_nonneg : 0 ≤ Kraw := by
    dsimp [Kraw]
    have hratio_nonneg : 0 ≤ ellipticityRatio A + 1 := by
      linarith [A.ellipticityRatio_nonneg]
    have htmp_nonneg : 0 ≤ (ellipticityRatio A + 1) * (2 * (Mst : ℝ)) ^ 2 := by
      exact mul_nonneg hratio_nonneg (sq_nonneg _)
    have hfactor_nonneg : 0 ≤ 8 * (ellipticityRatio A + 1) * (2 * (Mst : ℝ)) ^ 2 := by
      have h8_nonneg : (0 : ℝ) ≤ 8 := by norm_num
      simpa [mul_assoc] using mul_nonneg h8_nonneg htmp_nonneg
    exact mul_nonneg (sq_nonneg _) hfactor_nonneg
  have hKraw_le_K : Kraw ≤ K := by
    rw [hK_eq]
    linarith
  have hK : 0 < K := by
    rw [hK_eq]
    linarith
  refine ⟨by simpa [K] using hK, ?_⟩
  intro lamStar hLamStar
  have hStarInt :
      IntegrableOn (fun x => |positivePartSub u lamStar x| ^ 2)
        (Metric.ball (0 : E) (1 / 2 : ℝ)) volume := by
    let hwStar1 : MemW1pWitness 2 (positivePartSub u lamStar) (Metric.ball (0 : E) 1) :=
      positivePartSub_memW1pWitness_on_ball (d := d) (x₀ := (0 : E)) (s := 1)
        zero_lt_one hu1 lamStar (happroxUnit lamStar)
    have hStarInt1 :
        IntegrableOn (fun x => |positivePartSub u lamStar x| ^ 2)
          (Metric.ball (0 : E) 1) volume := by
      simpa [pow_two] using hwStar1.memLp.integrable_sq
    exact hStarInt1.mono_set (Metric.ball_subset_ball (by norm_num : (1 / 2 : ℝ) ≤ 1))
  have hAint :
      ∀ n,
        IntegrableOn (fun x => |positivePartSub u (deGiorgiLevel lamStar n) x| ^ 2)
          (Metric.ball (0 : E) (deGiorgiRadius n)) volume := by
    intro n
    let θ : ℝ := deGiorgiLevel lamStar n
    let hwθ1 : MemW1pWitness 2 (positivePartSub u θ) (Metric.ball (0 : E) 1) :=
      positivePartSub_memW1pWitness_on_ball (d := d) (x₀ := (0 : E)) (s := 1)
        zero_lt_one hu1 θ (happroxUnit θ)
    have hθInt1 :
        IntegrableOn (fun x => |positivePartSub u θ x| ^ 2)
          (Metric.ball (0 : E) 1) volume := by
      simpa [θ, pow_two] using hwθ1.memLp.integrable_sq
    exact hθInt1.mono_set (Metric.ball_subset_ball (le_of_lt (deGiorgiRadius_lt_one n)))
  have hpre :
      ∀ n,
        deGiorgiEnergySeq u (0 : E) lamStar (n + 1) ≤
          K / ((deGiorgiRadius n - deGiorgiRadius (n + 1)) ^ 2 *
            (deGiorgiLevel lamStar (n + 1) - deGiorgiLevel lamStar n) ^ (4 / (d : ℝ))) *
          deGiorgiEnergySeq u (0 : E) lamStar n ^ (1 + 2 / (d : ℝ)) := by
    intro n
    let r : ℝ := deGiorgiRadius (n + 1)
    let s : ℝ := deGiorgiRadius n
    let θ : ℝ := deGiorgiLevel lamStar n
    let lam : ℝ := deGiorgiLevel lamStar (n + 1)
    have hr : 0 < r := by
      dsimp [r]
      linarith [deGiorgiRadius_gt_half (n + 1)]
    have hs : 0 < s := by
      dsimp [s]
      linarith [deGiorgiRadius_gt_half n]
    have hrs : r < s := by
      have hgap : 0 < s - r := by
        dsimp [s, r]
        rw [deGiorgiRadius_gap]
        positivity
      linarith
    have hs_lt_one : s < 1 := by
      dsimp [s]
      exact deGiorgiRadius_lt_one n
    have hθl : θ < lam := by
      have hgap : 0 < lam - θ := by
        dsimp [lam, θ]
        rw [deGiorgiLevel_gap]
        positivity
      linarith
    have hsr_pos : 0 < s - r := sub_pos.mpr hrs
    let huS : MemW1pWitness 2 u (Metric.ball (0 : E) s) :=
      hu1.restrict isOpen_ball (Metric.ball_subset_ball (le_of_lt hs_lt_one))
    have happroxBallTheta :
        ∃ ψ : ℕ → E → ℝ,
          (∀ n, ContDiff ℝ 1 (ψ n)) ∧
          (∀ n, HasCompactSupport (ψ n)) ∧
          Tendsto
            (fun n =>
              eLpNorm (fun x => ψ n x - (u x - θ)) 2
                (volume.restrict (Metric.ball (0 : E) s)))
            atTop (nhds 0) ∧
          (∀ i : Fin d,
            Tendsto
              (fun n =>
                eLpNorm
                  (fun x =>
                    (fderiv ℝ (ψ n) x) (EuclideanSpace.single i 1) - huS.weakGrad x i)
                  2 (volume.restrict (Metric.ball (0 : E) s)))
              atTop (nhds 0)) := by
      rcases happroxUnit θ with ⟨ψ, hψ_smooth, hψ_compact, hψ_fun, hψ_grad⟩
      refine ⟨ψ, hψ_smooth, hψ_compact, ?_, ?_⟩
      · refine ENNReal.tendsto_nhds_zero.2 ?_
        intro ε hε
        filter_upwards [ENNReal.tendsto_nhds_zero.1 hψ_fun ε hε] with m hm
        exact le_trans
          (eLpNorm_mono_measure (fun x => ψ m x - (u x - θ))
            (Measure.restrict_mono_set volume
              (Metric.ball_subset_ball (le_of_lt hs_lt_one))))
          hm
      · intro i
        refine ENNReal.tendsto_nhds_zero.2 ?_
        intro ε hε
        filter_upwards [ENNReal.tendsto_nhds_zero.1 (hψ_grad i) ε hε] with m hm
        have hmono :
            eLpNorm
              (fun x =>
                (fderiv ℝ (ψ m) x) (EuclideanSpace.single i 1) - huS.weakGrad x i)
              2 (volume.restrict (Metric.ball (0 : E) s)) ≤
              eLpNorm
                (fun x =>
                  (fderiv ℝ (ψ m) x) (EuclideanSpace.single i 1) - hu1.weakGrad x i)
                2 (volume.restrict (Metric.ball (0 : E) 1)) := by
          simpa [huS] using
            (eLpNorm_mono_measure
              (fun x =>
                (fderiv ℝ (ψ m) x) (EuclideanSpace.single i 1) - hu1.weakGrad x i)
              (Measure.restrict_mono_set volume
                (Metric.ball_subset_ball (le_of_lt hs_lt_one))))
        exact le_trans hmono hm
    let R : ℝ := (r + s) / 2
    have hR : r < R := by
      dsimp [R]
      linarith
    have hR_lt_s : R < s := by
      dsimp [R]
      linarith
    obtain ⟨ηCut, _⟩ := Cutoff.exists (d := d) (0 : E) (r := r) (R := R) hr hR
    let η : E → ℝ := ηCut.toFun
    let Cη : ℝ := 2 * (Mst : ℝ) * (s - r)⁻¹
    have hη : ContDiff ℝ (⊤ : ℕ∞) η := ηCut.smooth
    have hη_nonneg : ∀ x, 0 ≤ η x := ηCut.nonneg
    have hη_eq_one : ∀ x ∈ Metric.ball (0 : E) r, η x = 1 := ηCut.eq_one
    have hη_bound : ∀ x, |η x| ≤ 1 := by
      intro x
      rw [abs_of_nonneg (ηCut.nonneg x)]
      exact ηCut.le_one x
    have hη_sub_ball_s : tsupport η ⊆ Metric.ball (0 : E) s := by
      exact ηCut.support_subset.trans (Metric.closedBall_subset_ball hR_lt_s)
    have hη_sub_ball1 : tsupport η ⊆ Metric.ball (0 : E) 1 := by
      exact hη_sub_ball_s.trans (Metric.ball_subset_ball (le_of_lt hs_lt_one))
    have hCη : 0 ≤ Cη := by
      dsimp [Cη]
      positivity
    have hη_grad_bound : ∀ x, ‖fderiv ℝ η x‖ ≤ Cη := by
      intro x
      have hmid :
          (((r + s) / 2 : ℝ) - r) = (s - r) / 2 := by
        ring
      have h_inv : ((((r + s) / 2 : ℝ) - r)⁻¹) = 2 * (s - r)⁻¹ := by
        rw [hmid]
        field_simp [hsr_pos.ne']
      calc
        ‖fderiv ℝ η x‖ ≤ ↑Mst * ((((r + s) / 2 : ℝ) - r)⁻¹) := by
          simpa [η, R] using ηCut.grad_bound x
        _ = 2 * (Mst : ℝ) * (s - r)⁻¹ := by
          rw [h_inv]
          ring
        _ = Cη := by
          rfl
    obtain ⟨hwηθ, hsob⟩ :=
      deGiorgi_cutoffSobolev_on_concentricBalls_of_posPartApprox
        (d := d) hd hs hr hrs huS happroxBallTheta hθl hη hη_nonneg
        hη_eq_one hη_bound hη_sub_ball_s
    let hwθ1 : MemW1pWitness 2 (positivePartSub u θ) (Metric.ball (0 : E) 1) :=
      positivePartSub_memW1pWitness_on_ball (d := d) (x₀ := (0 : E)) (s := 1)
        zero_lt_one hu1 θ (happroxUnit θ)
    let hwθs : MemW1pWitness 2 (positivePartSub u θ) (Metric.ball (0 : E) s) :=
      hwθ1.restrict isOpen_ball (Metric.ball_subset_ball (le_of_lt hs_lt_one))
    have hweighted_unit :
        ∫ x in Metric.ball (0 : E) 1, η x ^ 2 * ‖hwθ1.weakGrad x‖ ^ 2 ∂volume ≤
          4 * ellipticityRatio A *
            ∫ x in Metric.ball (0 : E) 1,
              ‖fderiv ℝ η x‖ ^ 2 * |positivePartSub u θ x| ^ 2 ∂volume := by
      simpa [hwθ1] using
        caccioppoli_weighted_on_ball_of_posPartApprox
          (d := d) (x₀ := (0 : E)) (s := 1) (u := u) (η := η) (Cη := Cη) (k := θ)
          A zero_lt_one hsub hu1 (happroxUnit θ) hη hη_nonneg hη_bound
          hCη hη_grad_bound hη_sub_ball1
    have hleft_int_unit :
        IntegrableOn (fun x => η x ^ 2 * ‖hwθ1.weakGrad x‖ ^ 2)
          (Metric.ball (0 : E) 1) volume := by
      have hgradSqInt :
          Integrable (fun x => ‖hwθ1.weakGrad x‖ ^ 2)
            (volume.restrict (Metric.ball (0 : E) 1)) := by
        simpa [pow_two] using hwθ1.weakGrad_norm_memLp.integrable_sq
      refine Integrable.mono' hgradSqInt ?_ ?_
      · exact
          (((hη.continuous.pow 2).aemeasurable).mul
            hgradSqInt.aestronglyMeasurable.aemeasurable).aestronglyMeasurable
      · filter_upwards with x
        have hηx_nonneg : 0 ≤ η x := hη_nonneg x
        have hηx_le_one : η x ≤ 1 := by
          simpa [abs_of_nonneg hηx_nonneg] using hη_bound x
        have hηx_sq_le_one : η x ^ 2 ≤ 1 := by
          nlinarith
        have hgw_nonneg : 0 ≤ ‖hwθ1.weakGrad x‖ ^ 2 := by positivity
        have hprod_nonneg : 0 ≤ η x ^ 2 * ‖hwθ1.weakGrad x‖ ^ 2 := by positivity
        have hle :
            η x ^ 2 * ‖hwθ1.weakGrad x‖ ^ 2 ≤ ‖hwθ1.weakGrad x‖ ^ 2 := by
          nlinarith
        simpa [Real.norm_eq_abs, abs_of_nonneg hprod_nonneg, abs_of_nonneg hgw_nonneg] using hle
    have hright_int_unit :
        IntegrableOn (fun x => ‖fderiv ℝ η x‖ ^ 2 * |positivePartSub u θ x| ^ 2)
          (Metric.ball (0 : E) 1) volume := by
      have hposSqInt :
          Integrable (fun x => |positivePartSub u θ x| ^ 2)
            (volume.restrict (Metric.ball (0 : E) 1)) := by
        simpa [pow_two] using hwθ1.memLp.integrable_sq
      refine Integrable.mono' (hposSqInt.const_mul (Cη ^ 2)) ?_ ?_
      · exact
          ((((hη.continuous_fderiv (by simp : ((⊤ : ℕ∞) : WithTop ℕ∞) ≠ 0)).norm.pow 2).aemeasurable).mul
            hposSqInt.aestronglyMeasurable.aemeasurable).aestronglyMeasurable
      · filter_upwards with x
        have hfd_sq_le : ‖fderiv ℝ η x‖ ^ 2 ≤ Cη ^ 2 := by
          exact sq_le_sq.mpr (by
            simpa [abs_of_nonneg (norm_nonneg _), abs_of_nonneg hCη] using hη_grad_bound x)
        have hpos_nonneg : 0 ≤ |positivePartSub u θ x| ^ 2 := by positivity
        have hterm_nonneg :
            0 ≤ ‖fderiv ℝ η x‖ ^ 2 * |positivePartSub u θ x| ^ 2 := by positivity
        have hle :
            ‖fderiv ℝ η x‖ ^ 2 * |positivePartSub u θ x| ^ 2 ≤
              Cη ^ 2 * |positivePartSub u θ x| ^ 2 :=
          mul_le_mul_of_nonneg_right hfd_sq_le hpos_nonneg
        simpa [Real.norm_eq_abs, abs_of_nonneg hterm_nonneg] using hle
    have hleft_eq :
        ∫ x in Metric.ball (0 : E) 1, η x ^ 2 * ‖hwθ1.weakGrad x‖ ^ 2 ∂volume =
          ∫ x in Metric.ball (0 : E) s, η x ^ 2 * ‖hwθs.weakGrad x‖ ^ 2 ∂volume := by
      calc
        ∫ x in Metric.ball (0 : E) 1, η x ^ 2 * ‖hwθ1.weakGrad x‖ ^ 2 ∂volume
            = ∫ x in Metric.ball (0 : E) s, η x ^ 2 * ‖hwθ1.weakGrad x‖ ^ 2 ∂volume := by
                refine setIntegral_ball_eq_of_zero_off_innerBall (x₀ := (0 : E)) hs_lt_one
                  hleft_int_unit ?_
                intro x hx1 hxs
                have hηx : η x = 0 := by
                  exact zero_outside_of_tsupport_subset
                    (Ω := Metric.ball (0 : E) s) hη_sub_ball_s hxs
                simp [hηx]
        _ = ∫ x in Metric.ball (0 : E) s, η x ^ 2 * ‖hwθs.weakGrad x‖ ^ 2 ∂volume := by
              rfl
    have hright_eq :
        ∫ x in Metric.ball (0 : E) 1, ‖fderiv ℝ η x‖ ^ 2 * |positivePartSub u θ x| ^ 2 ∂volume =
          ∫ x in Metric.ball (0 : E) s, ‖fderiv ℝ η x‖ ^ 2 * |positivePartSub u θ x| ^ 2
            ∂volume := by
      refine setIntegral_ball_eq_of_zero_off_innerBall (x₀ := (0 : E)) hs_lt_one
        hright_int_unit ?_
      intro x hx1 hxs
      have hfdv_zero : deGiorgiFderivVec η x = 0 := by
        ext i
        rw [deGiorgiFderivVec_apply]
        exact fderiv_apply_zero_outside_of_tsupport_subset
          (Ω := Metric.ball (0 : E) s) (hf := hη) hη_sub_ball_s hxs i
      have hfd_zero : ‖fderiv ℝ η x‖ = 0 := by
        rw [← deGiorgi_norm_fderivVec_eq_norm_fderiv (η := η) (x := x), hfdv_zero, norm_zero]
      simp [hfd_zero]
    have hweighted_small :
        ∫ x in Metric.ball (0 : E) s, η x ^ 2 * ‖hwθs.weakGrad x‖ ^ 2 ∂volume ≤
          4 * ellipticityRatio A *
            ∫ x in Metric.ball (0 : E) s,
              ‖fderiv ℝ η x‖ ^ 2 * |positivePartSub u θ x| ^ 2 ∂volume := by
      rw [← hleft_eq, ← hright_eq]
      exact hweighted_unit
    let Iθ : ℝ := deGiorgiEnergySeq u (0 : E) lamStar n
    let Ilam : ℝ := deGiorgiEnergySeq u (0 : E) lamStar (n + 1)
    let hwφ : MemW1pWitness 2 (deGiorgiCutoffTestGeneral η u θ) (Metric.ball (0 : E) s) :=
      deGiorgiCutoffTestWitnessWeighted
        isOpen_ball hwθs hη (by norm_num) hCη hη_bound hη_grad_bound
    have hgrad_bound_phi :
        ∫ x in Metric.ball (0 : E) s, ‖hwφ.weakGrad x‖ ^ 2 ∂volume ≤
          (8 * (ellipticityRatio A + 1) * Cη ^ 2) * Iθ := by
      simpa [Iθ, deGiorgiEnergySeq, hwφ, s, θ] using
        deGiorgi_cutoff_gradient_bound_of_weighted
          (x₀ := (0 : E)) (s := s) (u := u) (η := η) (k := θ)
          (Cη := Cη) (Cw := ellipticityRatio A)
          A.ellipticityRatio_nonneg hwθs hη hη_nonneg hη_bound hCη hη_grad_bound
          hweighted_small
    have hgrad_ae :
        hwηθ.weakGrad =ᵐ[volume.restrict (Metric.ball (0 : E) s)] hwφ.weakGrad :=
      MemW1pWitness.ae_eq isOpen_ball hwηθ hwφ
    have henergy :
        ∫ x in Metric.ball (0 : E) s, ‖hwηθ.weakGrad x‖ ^ 2 ∂volume ≤
          (8 * (ellipticityRatio A + 1) * Cη ^ 2) * Iθ := by
      have hgrad_eq :
          ∫ x in Metric.ball (0 : E) s, ‖hwηθ.weakGrad x‖ ^ 2 ∂volume =
            ∫ x in Metric.ball (0 : E) s, ‖hwφ.weakGrad x‖ ^ 2 ∂volume := by
        refine integral_congr_ae ?_
        filter_upwards [hgrad_ae] with x hx
        simp [hx]
      rw [hgrad_eq]
      exact hgrad_bound_phi
    have hθ_int :
        Integrable (fun x => |positivePartSub u θ x| ^ 2)
          (volume.restrict (Metric.ball (0 : E) s)) := by
      simpa [pow_two] using hwθs.memLp.integrable_sq
    have hIθ_nonneg : 0 ≤ Iθ := by
      dsimp [Iθ, deGiorgiEnergySeq]
      refine integral_nonneg ?_
      intro x
      positivity
    have hd_pos : 0 < (d : ℝ) := by
      linarith
    have hCenergy_nonneg : 0 ≤ 8 * (ellipticityRatio A + 1) * Cη ^ 2 := by
      nlinarith [A.ellipticityRatio_nonneg, sq_nonneg Cη]
    haveI : IsFiniteMeasure (volume.restrict (Metric.ball (0 : E) s)) := by
      rw [isFiniteMeasure_restrict]
      exact measure_ball_lt_top.ne
    have hpre_raw :
        Ilam ≤
          (C_gns d 2) ^ 2 * (8 * (ellipticityRatio A + 1) * Cη ^ 2) *
            ((((lam - θ) ^ 2)⁻¹ * Iθ) ^ (2 / (d : ℝ))) * Iθ := by
      have hsob' :
          Ilam ≤
            (C_gns d 2) ^ 2 *
              (∫ x in Metric.ball (0 : E) s, ‖hwηθ.weakGrad x‖ ^ 2 ∂volume) *
              ((volume.restrict (Metric.ball (0 : E) s)).real {x | lam < u x}) ^ (2 / (d : ℝ)) := by
        simpa [Ilam, deGiorgiEnergySeq, r, lam] using hsob
      exact
        deGiorgi_preiter_of_energy
          (μ := volume.restrict (Metric.ball (0 : E) s))
          (d := d) hd_pos (u := u) (θ := θ) (lam := lam)
          (Ilam := Ilam)
          (Iθ := Iθ)
          (G := ∫ x in Metric.ball (0 : E) s, ‖hwηθ.weakGrad x‖ ^ 2 ∂volume)
          (Csob := (C_gns d 2) ^ 2)
          (Cenergy := 8 * (ellipticityRatio A + 1) * Cη ^ 2)
          hθl (sq_nonneg _) hCenergy_nonneg hIθ_nonneg hθ_int hsob'
          henergy (by simp [Iθ, deGiorgiEnergySeq, s, θ])
    have hgap_pos : 0 < lam - θ := sub_pos.mpr hθl
    have hCη_sq :
        Cη ^ 2 = (2 * (Mst : ℝ)) ^ 2 * (((s - r) ^ 2)⁻¹) := by
      dsimp [Cη]
      rw [mul_pow, inv_pow]
    have hpow_gap :
        ((((lam - θ) ^ 2)⁻¹) ^ (2 / (d : ℝ))) = (lam - θ) ^ (-(4 / (d : ℝ))) := by
      calc
        ((((lam - θ) ^ 2)⁻¹) ^ (2 / (d : ℝ)))
            = ((((lam - θ) ^ 2) ^ (-(1 : ℝ))) ^ (2 / (d : ℝ))) := by
                rw [Real.rpow_neg (by positivity), Real.rpow_one]
        _ = ((lam - θ) ^ 2) ^ ((-(1 : ℝ)) * (2 / (d : ℝ))) := by
              rw [← Real.rpow_mul (by positivity : 0 ≤ (lam - θ) ^ 2)]
        _ = ((lam - θ) ^ 2) ^ (-(2 / (d : ℝ))) := by
              congr 1
              ring
        _ = ((lam - θ) ^ (2 : ℝ)) ^ (-(2 / (d : ℝ))) := by
              congr 1
              exact (Real.rpow_natCast (lam - θ) 2).symm
        _ = (lam - θ) ^ (-(4 / (d : ℝ))) := by
              rw [← Real.rpow_mul hgap_pos.le]
              congr 1
              ring
    have hIθ_pow :
        Iθ ^ (2 / (d : ℝ)) * Iθ = Iθ ^ (1 + 2 / (d : ℝ)) := by
      by_cases hIθ_zero : Iθ = 0
      · have hexp1_pos : 0 < 2 / (d : ℝ) := by positivity
        have hexp2_pos : 0 < 1 + 2 / (d : ℝ) := by linarith
        simp [hIθ_zero, Real.zero_rpow hexp1_pos.ne', Real.zero_rpow hexp2_pos.ne']
      · have hIθ_pos : 0 < Iθ := lt_of_le_of_ne hIθ_nonneg (by simpa [eq_comm] using hIθ_zero)
        have hmul_one :
            Iθ ^ (2 / (d : ℝ)) * Iθ = Iθ ^ (2 / (d : ℝ)) * Iθ ^ (1 : ℝ) := by
          rw [Real.rpow_one]
        calc
          Iθ ^ (2 / (d : ℝ)) * Iθ = Iθ ^ (2 / (d : ℝ)) * Iθ ^ (1 : ℝ) := hmul_one
          _ = Iθ ^ ((2 / (d : ℝ)) + 1) := by
            rw [← Real.rpow_add hIθ_pos]
          _ = Iθ ^ (1 + 2 / (d : ℝ)) := by
            congr 1
            ring
    have hcoeff_eq :
        (C_gns d 2) ^ 2 * (8 * (ellipticityRatio A + 1) *
            ((2 * (Mst : ℝ)) ^ 2 * ((s - r) ^ 2)⁻¹)) *
            ((lam - θ) ^ (4 / (d : ℝ)))⁻¹ =
          Kraw * ((s - r) ^ 2 * (lam - θ) ^ (4 / (d : ℝ)))⁻¹ := by
      dsimp [Kraw]
      conv_rhs =>
        rw [_root_.mul_inv_rev]
        rw [mul_comm]
      ring
    have hrewrite :
        (C_gns d 2) ^ 2 * (8 * (ellipticityRatio A + 1) * Cη ^ 2) *
            ((((lam - θ) ^ 2)⁻¹ * Iθ) ^ (2 / (d : ℝ))) * Iθ =
          Kraw / ((s - r) ^ 2 * (lam - θ) ^ (4 / (d : ℝ))) * Iθ ^ (1 + 2 / (d : ℝ)) := by
      rw [hCη_sq]
      rw [Real.mul_rpow (by positivity) hIθ_nonneg]
      rw [hpow_gap]
      have hstep :
          (C_gns d 2) ^ 2 * (8 * (ellipticityRatio A + 1) * ((2 * (Mst : ℝ)) ^ 2 * (((s - r) ^ 2)⁻¹))) *
              ((lam - θ) ^ (-(4 / (d : ℝ))) * Iθ ^ (2 / (d : ℝ))) * Iθ =
            (C_gns d 2) ^ 2 * (8 * (ellipticityRatio A + 1) * ((2 * (Mst : ℝ)) ^ 2 * (((s - r) ^ 2)⁻¹))) *
              (lam - θ) ^ (-(4 / (d : ℝ))) * (Iθ ^ (2 / (d : ℝ)) * Iθ) := by
        ring
      rw [hstep, hIθ_pow]
      rw [Real.rpow_neg hgap_pos.le, div_eq_mul_inv]
      have hcoeff_mul :
          ((C_gns d 2) ^ 2 * (8 * (ellipticityRatio A + 1) *
              ((2 * (Mst : ℝ)) ^ 2 * ((s - r) ^ 2)⁻¹)) *
              ((lam - θ) ^ (4 / (d : ℝ)))⁻¹) *
              Iθ ^ (1 + 2 / (d : ℝ)) =
            (Kraw * ((s - r) ^ 2 * (lam - θ) ^ (4 / (d : ℝ)))⁻¹) *
              Iθ ^ (1 + 2 / (d : ℝ)) := by
        exact congrArg (fun t : ℝ => t * Iθ ^ (1 + 2 / (d : ℝ))) hcoeff_eq
      simpa [div_eq_mul_inv, mul_assoc] using hcoeff_mul
    have hIθ_pow_nonneg : 0 ≤ Iθ ^ (1 + 2 / (d : ℝ)) := by
      exact Real.rpow_nonneg hIθ_nonneg _
    have hden_pos : 0 < (s - r) ^ 2 * (lam - θ) ^ (4 / (d : ℝ)) := by
      have hpow_pos : 0 < (lam - θ) ^ (4 / (d : ℝ)) := by
        exact Real.rpow_pos_of_pos hgap_pos _
      positivity
    have hden_inv_nonneg : 0 ≤ ((s - r) ^ 2 * (lam - θ) ^ (4 / (d : ℝ)))⁻¹ := by
      exact inv_nonneg.mpr (le_of_lt hden_pos)
    calc
      deGiorgiEnergySeq u (0 : E) lamStar (n + 1) = Ilam := by
          rfl
      _ ≤ (C_gns d 2) ^ 2 * (8 * (ellipticityRatio A + 1) * Cη ^ 2) *
            ((((lam - θ) ^ 2)⁻¹ * Iθ) ^ (2 / (d : ℝ))) * Iθ := hpre_raw
      _ = Kraw / ((s - r) ^ 2 * (lam - θ) ^ (4 / (d : ℝ))) * Iθ ^ (1 + 2 / (d : ℝ)) := hrewrite
      _ ≤ K / ((s - r) ^ 2 * (lam - θ) ^ (4 / (d : ℝ))) * Iθ ^ (1 + 2 / (d : ℝ)) := by
            let D : ℝ := ((s - r) ^ 2 * (lam - θ) ^ (4 / (d : ℝ)))⁻¹
            let P : ℝ := Iθ ^ (1 + 2 / (d : ℝ))
            have hD_nonneg : 0 ≤ D := by
              dsimp [D]
              exact hden_inv_nonneg
            have hP_nonneg : 0 ≤ P := by
              dsimp [P]
              exact hIθ_pow_nonneg
            have hKD : Kraw * D ≤ K * D :=
              mul_le_mul_of_nonneg_right hKraw_le_K hD_nonneg
            have hKDP : Kraw * D * P ≤ K * D * P :=
              mul_le_mul_of_nonneg_right hKD hP_nonneg
            simpa [D, P, div_eq_mul_inv, mul_assoc]
              using hKDP
      _ = K / ((deGiorgiRadius n - deGiorgiRadius (n + 1)) ^ 2 *
            (deGiorgiLevel lamStar (n + 1) - deGiorgiLevel lamStar n) ^ (4 / (d : ℝ))) *
            deGiorgiEnergySeq u (0 : E) lamStar n ^ (1 + 2 / (d : ℝ)) := by
            simp [Iθ, deGiorgiEnergySeq, r, s, θ, lam]
  exact ⟨hStarInt, hAint, hpre⟩

\end{lstlisting}
\begin{proof}
Take a $W^{1,2}(\Bz{1})$ witness for $u$ from
$\leanname{IsSubsolution.memW1p}$. The smooth approximation package
on the unit ball produces, for any shift $\theta$, an approximating
sequence $\psi_n \in C^1_c$ with $\psi_n \to u - \theta$ in $L^2$ and
$\nabla \psi_n \to \nabla u$ in $L^2$. From this one constructs
witnesses for $(u-\ell^*)_+$ and $(u-\ell_n(\ell^*))_+$ on $\Bz{1}$
via \leanname{positivePartSub\_memW1pWitness\_on\_ball}, and the
required integrability follows by restriction to $\Ball{r_n}{0}$. For
the recurrence, fix $n$ and apply the PDE pre-iteration theorem
\leanname{deGiorgi\_preiter\_on\_concentricBalls\_of\_ballPosPart} to $r =
r_{n+1}$, $s = r_n$, $\theta = \ell_n(\ell^*)$, $\ell =
\ell_{n+1}(\ell^*)$, with a smooth Stampacchia cutoff $\eta_n$
satisfying $\eta_n = 1$ on $\Ball{r_{n+1}}{0}$, $\tsupp\,\eta_n
\subset \Ball{r_n}{0}$, and $\|\nabla \eta_n\| \le 2 M_{\mathrm{st}}
/ (r_n - r_{n+1})$, so that the constant becomes
$(C_{\mathrm{gns}}(d,2))^2 \cdot 8(\Lambda/\lambda + 1)
(2 M_{\mathrm{st}})^2 \le K_{\mathrm{DG}}(d,A)$.
\end{proof}

\begin{theorem}[Vanishing version,
\leanname{linfty\_subsolution\_DeGiorgi\_ae\_zero\_of\_smallness}]
With $d > 2$, $K, \ell^* > 0$, the integrability assumptions on the
truncations, the canonical recurrence, and the abstract smallness
condition
\[
  A_0(u,x_0,\ell^*) \le
    C_d(K,\ell^*)^{-d/2}\, B_d^{-(d/2)^2 / 2},
\]
one has $(u-\ell^*)_+ = 0$ almost everywhere on $\Ball{1/2}{x_0}$.
\end{theorem}

\begin{lstlisting}[style=Matlab-editor,basicstyle=\mlttfamily,escapechar=",mlshowsectionrules=true,mathescape=true,morecomment={[s]{\%\{}{\%\}}},language=lean,numbers=none]
omit [NeZero d] in
/-- Small-data De Giorgi endpoint: if the canonical De Giorgi energy sequence
tends to zero, then the final truncation vanishes almost everywhere on the
half-ball. This is the formally correct endpoint available from the current
Chapter 05 infrastructure before any representative/continuity upgrade. -/
theorem linfty_subsolution_DeGiorgi_ae_zero_of_smallness
    {u : E → ℝ} {x₀ : E} {K lamStar : ℝ}
    (hd : 2 < (d : ℝ))
    (hK : 0 < K) (hLamStar : 0 < lamStar)
    (hStarInt :
      IntegrableOn (fun x => |positivePartSub u lamStar x| ^ 2)
        (Metric.ball x₀ (1 / 2 : ℝ)) volume)
    (hAint :
      ∀ n,
        IntegrableOn (fun x => |positivePartSub u (deGiorgiLevel lamStar n) x| ^ 2)
          (Metric.ball x₀ (deGiorgiRadius n)) volume)
    (hpre :
      ∀ n,
        deGiorgiEnergySeq u x₀ lamStar (n + 1) ≤
          K / ((deGiorgiRadius n - deGiorgiRadius (n + 1)) ^ 2 *
            (deGiorgiLevel lamStar (n + 1) - deGiorgiLevel lamStar n) ^ (4 / (d : ℝ))) *
          deGiorgiEnergySeq u x₀ lamStar n ^ (1 + 2 / (d : ℝ)))
    (hsmall :
      deGiorgiEnergySeq u x₀ lamStar 0 ≤
        (deGiorgiRecurrenceCoeff d K lamStar) ^ (-(1 : ℝ) / (2 / (d : ℝ))) *
          (deGiorgiRecurrenceBase d) ^ (-(1 : ℝ) / (2 / (d : ℝ)) ^ 2)) :
    ∀ᵐ x ∂(volume.restrict (Metric.ball x₀ (1 / 2 : ℝ))),
      positivePartSub u lamStar x = 0 := by
  have hA_nonneg : ∀ n, 0 ≤ deGiorgiEnergySeq u x₀ lamStar n := by
    intro n
    simp [deGiorgiEnergySeq]
    refine integral_nonneg ?_
    intro x
    positivity
  have hvanish :
      Tendsto (deGiorgiEnergySeq u x₀ lamStar) atTop (nhds 0) :=
    deGiorgi_preiter_vanishing hd hK hLamStar hA_nonneg hpre hsmall
  exact ae_eq_zero_of_forall_setIntegral_sq_le_of_tendsto_zero hStarInt
    (integral_sq_halfBall_le_deGiorgiEnergySeq hLamStar hStarInt hAint) hvanish

\end{lstlisting}
\begin{proof}
By \leanname{deGiorgi\_preiter\_vanishing}, $A_n \to 0$. Combined with
the half-ball monotonicity lemma, this gives
$\int_{\Ball{1/2}{x_0}} |(u-\ell^*)_+|^2 = 0$, whence
$(u-\ell^*)_+ = 0$ a.e.\ on the half-ball by the abstract zero
lemma.
\end{proof}

\begin{corollary}[\leanname{linfty\_subsolution\_DeGiorgi\_ae\_of\_smallness}]
Under the same hypotheses, $\max(u, 0) \le \ell^*$ a.e.\ on
$\Ball{1/2}{x_0}$.
\end{corollary}

\begin{lstlisting}[style=Matlab-editor,basicstyle=\mlttfamily,escapechar=",mlshowsectionrules=true,mathescape=true,morecomment={[s]{\%\{}{\%\}}},language=lean,numbers=none]
omit [NeZero d] in
/-- A.e. `L∞` bound on the half-ball in the small-data De Giorgi regime. -/
theorem linfty_subsolution_DeGiorgi_ae_of_smallness
    {u : E → ℝ} {x₀ : E} {K lamStar : ℝ}
    (hd : 2 < (d : ℝ))
    (hK : 0 < K) (hLamStar : 0 < lamStar)
    (hStarInt :
      IntegrableOn (fun x => |positivePartSub u lamStar x| ^ 2)
        (Metric.ball x₀ (1 / 2 : ℝ)) volume)
    (hAint :
      ∀ n,
        IntegrableOn (fun x => |positivePartSub u (deGiorgiLevel lamStar n) x| ^ 2)
          (Metric.ball x₀ (deGiorgiRadius n)) volume)
    (hpre :
      ∀ n,
        deGiorgiEnergySeq u x₀ lamStar (n + 1) ≤
          K / ((deGiorgiRadius n - deGiorgiRadius (n + 1)) ^ 2 *
            (deGiorgiLevel lamStar (n + 1) - deGiorgiLevel lamStar n) ^ (4 / (d : ℝ))) *
          deGiorgiEnergySeq u x₀ lamStar n ^ (1 + 2 / (d : ℝ)))
    (hsmall :
      deGiorgiEnergySeq u x₀ lamStar 0 ≤
        (deGiorgiRecurrenceCoeff d K lamStar) ^ (-(1 : ℝ) / (2 / (d : ℝ))) *
          (deGiorgiRecurrenceBase d) ^ (-(1 : ℝ) / (2 / (d : ℝ)) ^ 2)) :
    ∀ᵐ x ∂(volume.restrict (Metric.ball x₀ (1 / 2 : ℝ))),
      max (u x) 0 ≤ lamStar := by
  filter_upwards
    [linfty_subsolution_DeGiorgi_ae_zero_of_smallness
      hd hK hLamStar hStarInt hAint hpre hsmall] with x hx
  have hux : u x ≤ lamStar := le_of_positivePartSub_eq_zero hx
  exact max_le_iff.mpr ⟨hux, hLamStar.le⟩

\end{lstlisting}
\begin{proof}
$(u-\ell^*)_+(x) = 0$ implies $u(x) \le \ell^*$ by
\leanname{le\_of\_positivePartSub\_eq\_zero}, and $0 \le \ell^*$.
\end{proof}

\begin{theorem}[Unit-ball smallness packaging,
\leanname{linfty\_subsolution\_DeGiorgi\_ae\_of\_unitBall\_initial\_energy\_smallness}]
Replacing the abstract smallness on $A_0$ by the same threshold for
$\int_{\Bz{1}} (u_+)^2$, the same conclusion holds.
\end{theorem}

\begin{lstlisting}[style=Matlab-editor,basicstyle=\mlttfamily,escapechar=",mlshowsectionrules=true,mathescape=true,morecomment={[s]{\%\{}{\%\}}},language=lean,numbers=none]
omit [NeZero d] in
/-- Unit-ball smallness packaging for the De Giorgi endpoint.

This is the direct successor to the raw `hsmall`-based theorem: instead of
assuming a bound on the abstract initial De Giorgi energy `A₀`, it is enough to
assume the same recurrence threshold for the concrete unit-ball quantity
`∫_{B₁} (u₊)^2`. The radius/level bookkeeping then supplies the `A₀` bound
automatically. -/
theorem linfty_subsolution_DeGiorgi_ae_of_unitBall_initial_energy_smallness
    {u : E → ℝ} {x₀ : E} {K lamStar : ℝ}
    (hd : 2 < (d : ℝ))
    (hK : 0 < K) (hLamStar : 0 < lamStar)
    (hposInt :
      IntegrableOn (fun x => |max (u x) 0| ^ 2)
        (Metric.ball x₀ 1) volume)
    (hStarInt :
      IntegrableOn (fun x => |positivePartSub u lamStar x| ^ 2)
        (Metric.ball x₀ (1 / 2 : ℝ)) volume)
    (hAint :
      ∀ n,
        IntegrableOn (fun x => |positivePartSub u (deGiorgiLevel lamStar n) x| ^ 2)
          (Metric.ball x₀ (deGiorgiRadius n)) volume)
    (hpre :
      ∀ n,
        deGiorgiEnergySeq u x₀ lamStar (n + 1) ≤
          K / ((deGiorgiRadius n - deGiorgiRadius (n + 1)) ^ 2 *
            (deGiorgiLevel lamStar (n + 1) - deGiorgiLevel lamStar n) ^ (4 / (d : ℝ))) *
          deGiorgiEnergySeq u x₀ lamStar n ^ (1 + 2 / (d : ℝ)))
    (hsmall0 :
      ∫ x in Metric.ball x₀ 1, |max (u x) 0| ^ 2 ∂volume ≤
        (deGiorgiRecurrenceCoeff d K lamStar) ^ (-(1 : ℝ) / (2 / (d : ℝ))) *
          (deGiorgiRecurrenceBase d) ^ (-(1 : ℝ) / (2 / (d : ℝ)) ^ 2)) :
    ∀ᵐ x ∂(volume.restrict (Metric.ball x₀ (1 / 2 : ℝ))),
      max (u x) 0 ≤ lamStar := by
  have hsmall :
      deGiorgiEnergySeq u x₀ lamStar 0 ≤
        (deGiorgiRecurrenceCoeff d K lamStar) ^ (-(1 : ℝ) / (2 / (d : ℝ))) *
          (deGiorgiRecurrenceBase d) ^ (-(1 : ℝ) / (2 / (d : ℝ)) ^ 2) := by
    exact
      (deGiorgiEnergySeq_zero_le_integral_sq_posPartZero_on_unitBall hLamStar (hAint 0) hposInt).trans
        hsmall0
  exact linfty_subsolution_DeGiorgi_ae_of_smallness
    hd hK hLamStar hStarInt hAint hpre hsmall

\end{lstlisting}
\begin{proof}
The bound $A_0 \le \int_{\Bz{1}} (u_+)^2$ converts the assumption
into the previous form.
\end{proof}

\begin{theorem}[Height-threshold form,
\leanname{linfty\_subsolution\_DeGiorgi\_ae\_of\_unitBall\_height\_bound}]
Under the same setup, if
\[
  C_{\mathrm{small}}(d,K)\,
    \Bigl(\int_{\Bz{1}} (u_+)^2\Bigr)^{1/2} \le \ell^*,
\]
then $\max(u,0) \le \ell^*$ a.e.\ on $\Ball{1/2}{x_0}$.
\end{theorem}

\begin{lstlisting}[style=Matlab-editor,basicstyle=\mlttfamily,escapechar=",mlshowsectionrules=true,mathescape=true,morecomment={[s]{\%\{}{\%\}}},language=lean,numbers=none]
/-- Height-threshold version of the De Giorgi endpoint.

Once the PDE-facing one-step estimate is supplied, a lower bound for `lamStar`
of the form `C(d,K) * sqrt(∫_{B₁} (u₊)^2)` yields the a.e. half-ball `L∞`
bound. -/
theorem linfty_subsolution_DeGiorgi_ae_of_unitBall_height_bound
    {u : E → ℝ} {x₀ : E} {K lamStar : ℝ}
    (hd : 2 < (d : ℝ))
    (hK : 0 < K) (hLamStar : 0 < lamStar)
    (hposInt :
      IntegrableOn (fun x => |max (u x) 0| ^ 2)
        (Metric.ball x₀ 1) volume)
    (hStarInt :
      IntegrableOn (fun x => |positivePartSub u lamStar x| ^ 2)
        (Metric.ball x₀ (1 / 2 : ℝ)) volume)
    (hAint :
      ∀ n,
        IntegrableOn (fun x => |positivePartSub u (deGiorgiLevel lamStar n) x| ^ 2)
          (Metric.ball x₀ (deGiorgiRadius n)) volume)
    (hpre :
      ∀ n,
        deGiorgiEnergySeq u x₀ lamStar (n + 1) ≤
          K / ((deGiorgiRadius n - deGiorgiRadius (n + 1)) ^ 2 *
            (deGiorgiLevel lamStar (n + 1) - deGiorgiLevel lamStar n) ^ (4 / (d : ℝ))) *
          deGiorgiEnergySeq u x₀ lamStar n ^ (1 + 2 / (d : ℝ)))
    (hheight :
      C_DeGiorgiSmallness d K *
        Real.sqrt (∫ x in Metric.ball x₀ 1, |max (u x) 0| ^ 2 ∂volume) ≤
          lamStar) :
    ∀ᵐ x ∂(volume.restrict (Metric.ball x₀ (1 / 2 : ℝ))),
      max (u x) 0 ≤ lamStar := by
  have hI₀_nonneg :
      0 ≤ ∫ x in Metric.ball x₀ 1, |max (u x) 0| ^ 2 ∂volume := by
    refine integral_nonneg ?_
    intro x
    positivity
  have hsmall0 :
      ∫ x in Metric.ball x₀ 1, |max (u x) 0| ^ 2 ∂volume ≤
        (deGiorgiRecurrenceCoeff d K lamStar) ^ (-(1 : ℝ) / (2 / (d : ℝ))) *
          (deGiorgiRecurrenceBase d) ^ (-(1 : ℝ) / (2 / (d : ℝ)) ^ 2) :=
    deGiorgi_initial_energy_small_of_height_bound hK hLamStar hI₀_nonneg hheight
  exact linfty_subsolution_DeGiorgi_ae_of_unitBall_initial_energy_smallness
    hd hK hLamStar hposInt hStarInt hAint hpre hsmall0

\end{lstlisting}
\begin{proof}
Squaring and using the formula for $C_{\mathrm{small}}$ shows that
the height-bound implies the abstract smallness condition above.
This step is the lemma
\leanname{deGiorgi\_initial\_energy\_small\_of\_height\_bound}.
\end{proof}

\begin{theorem}[De Giorgi $L^\infty$ bound,
\leanname{linfty\_subsolution\_DeGiorgi}]
Let $d > 2$, $A$ elliptic on $\Bz{1}$, and $u$ a subsolution with
$|\max(u,0)|^2 \in L^1(\Bz{1})$. Then
\[
  \max(u(x), 0) \le C_{\mathrm{DG}}(d, A) \cdot
    \Bigl(\int_{\Bz{1}} |\max(u,0)|^2\, dx\Bigr)^{1/2}
\]
for almost every $x \in \Ball{1/2}{0}$.
\end{theorem}

\begin{lstlisting}[style=Matlab-editor,basicstyle=\mlttfamily,escapechar=",mlshowsectionrules=true,mathescape=true,morecomment={[s]{\%\{}{\%\}}},language=lean,numbers=none]
/-- De Giorgi `L∞` bound for subsolutions on the unit ball.

The positive part is bounded almost everywhere on the half-ball by an explicit
coefficient-dependent constant times the `L²` size of `u₊` on `B₁`. The result
is stated as an a.e. bound rather than a pointwise supremum bound, since the
representative/continuity upgrade comes later in the theory. -/
theorem linfty_subsolution_DeGiorgi
    (hd : 2 < (d : ℝ))
    (A : EllipticCoeff d (Metric.ball (0 : E) 1))
    {u : E → ℝ}
    (hsub : IsSubsolution A u)
    (hposInt :
      IntegrableOn (fun x => |max (u x) 0| ^ 2)
        (Metric.ball (0 : E) 1) volume) :
    ∀ᵐ x ∂(volume.restrict (Metric.ball (0 : E) (1 / 2 : ℝ))),
      max (u x) 0 ≤
        C_DeGiorgi_subsolution d A *
          Real.sqrt (∫ x in Metric.ball (0 : E) 1, |max (u x) 0| ^ 2 ∂volume) := by
  let I₀ : ℝ := ∫ x in Metric.ball (0 : E) 1, |max (u x) 0| ^ 2 ∂volume
  have hI₀_nonneg : 0 ≤ I₀ := by
    dsimp [I₀]
    refine integral_nonneg ?_
    intro x
    positivity
  obtain ⟨hK, hiter⟩ :=
    deGiorgi_iteration_package_on_unitBall_of_subsolution (d := d) hd A hsub
  by_cases hI₀_zero : I₀ = 0
  · have hposInt_half :
        IntegrableOn (fun x => |max (u x) 0| ^ 2)
          (Metric.ball (0 : E) (1 / 2 : ℝ)) volume :=
      hposInt.mono_set (Metric.ball_subset_ball (by norm_num : (1 / 2 : ℝ) ≤ 1))
    have hnonneg : ∀ x, 0 ≤ |max (u x) 0| ^ 2 := by
      intro x
      positivity
    have hbound_half :
        ∀ n : ℕ,
          ∫ x in Metric.ball (0 : E) (1 / 2 : ℝ), |max (u x) 0| ^ 2 ∂volume ≤ (0 : ℝ) := by
      intro n
      calc
        ∫ x in Metric.ball (0 : E) (1 / 2 : ℝ), |max (u x) 0| ^ 2 ∂volume
            ≤ ∫ x in Metric.ball (0 : E) 1, |max (u x) 0| ^ 2 ∂volume := by
                exact setIntegral_mono_set hposInt (ae_of_all _ hnonneg)
                  (ae_of_all _ (Metric.ball_subset_ball (by norm_num : (1 / 2 : ℝ) ≤ 1)))
        _ = 0 := by
              simpa [I₀] using hI₀_zero
    have hzero_ae :
        ∀ᵐ x ∂(volume.restrict (Metric.ball (0 : E) (1 / 2 : ℝ))), max (u x) 0 = 0 := by
      refine ae_eq_zero_of_forall_setIntegral_sq_le_of_tendsto_zero hposInt_half hbound_half ?_
      exact (tendsto_const_nhds : Tendsto (fun _ : ℕ => (0 : ℝ)) atTop (nhds 0))
    filter_upwards [hzero_ae] with x hx
    have hrhs_zero :
        C_DeGiorgi_subsolution d A * Real.sqrt I₀ = 0 := by
      simp [I₀, hI₀_zero]
    rw [hrhs_zero]
    simp [hx]
  · have hI₀_pos : 0 < I₀ := lt_of_le_of_ne hI₀_nonneg (by simpa [eq_comm] using hI₀_zero)
    let lamStar : ℝ := C_DeGiorgi_subsolution d A * Real.sqrt I₀
    have hLamStar : 0 < lamStar := by
      dsimp [lamStar, I₀]
      exact deGiorgi_exact_height_threshold_on_unitBall
        (d := d) (K := K_DeGiorgi_subsolution d A) hK hI₀_pos
    obtain ⟨hStarInt, htail⟩ := hiter hLamStar
    obtain ⟨hAint, hpre⟩ := htail
    simpa [lamStar, I₀] using
      (linfty_subsolution_DeGiorgi_ae_of_unitBall_height_bound
        (d := d) (u := u) (x₀ := (0 : E)) (K := K_DeGiorgi_subsolution d A)
        (lamStar := lamStar) hd hK hLamStar hposInt hStarInt hAint hpre
        (by simp [lamStar, I₀, C_DeGiorgi_subsolution]))

\end{lstlisting}
\begin{proof}
Set $I_0 := \int_{\Bz{1}} |u_+|^2 \, dx$. If $I_0 = 0$ then by
monotonicity $|u_+|^2 = 0$ a.e.\ on $\Ball{1/2}{0}$, so the
inequality holds trivially. Otherwise let $\ell^* := C_{\mathrm{DG}}
(d,A) \sqrt{I_0} > 0$. The iteration package supplies the integrability
of all canonical truncations and the canonical recurrence. The
height-threshold theorem then yields $\max(u,0) \le \ell^* =
C_{\mathrm{DG}}(d,A) \sqrt{I_0}$ a.e.\ on $\Ball{1/2}{0}$.
\end{proof}

\begin{corollary}[Existence of a level threshold,
\leanname{linfty\_subsolution\_DeGiorgi\_exists\_threshold}]
Under the same hypotheses there exists $\ell^* > 0$ with
$C_{\mathrm{DG}}(d,A) \sqrt{I_0} \le \ell^*$ such that
$\max(u,0) \le \ell^*$ a.e.\ on $\Ball{1/2}{0}$.
\end{corollary}

\begin{lstlisting}[style=Matlab-editor,basicstyle=\mlttfamily,escapechar=",mlshowsectionrules=true,mathescape=true,morecomment={[s]{\%\{}{\%\}}},language=lean,numbers=none]
/-- Existence of a threshold realizing the De Giorgi `L∞` estimate. -/
theorem linfty_subsolution_DeGiorgi_exists_threshold
    (hd : 2 < (d : ℝ))
    (A : EllipticCoeff d (Metric.ball (0 : E) 1))
    {u : E → ℝ}
    (hsub : IsSubsolution A u)
    (hposInt :
      IntegrableOn (fun x => |max (u x) 0| ^ 2)
        (Metric.ball (0 : E) 1) volume) :
    ∃ lamStar : ℝ,
      0 < lamStar ∧
      C_DeGiorgi_subsolution d A *
        Real.sqrt (∫ x in Metric.ball (0 : E) 1, |max (u x) 0| ^ 2 ∂volume) ≤
          lamStar ∧
      ∀ᵐ x ∂(volume.restrict (Metric.ball (0 : E) (1 / 2 : ℝ))),
        max (u x) 0 ≤ lamStar := by
  let I₀ : ℝ := ∫ x in Metric.ball (0 : E) 1, |max (u x) 0| ^ 2 ∂volume
  have hI₀_nonneg : 0 ≤ I₀ := by
    dsimp [I₀]
    refine integral_nonneg ?_
    intro x
    positivity
  have hbound := linfty_subsolution_DeGiorgi (d := d) hd A hsub hposInt
  by_cases hI₀_zero : I₀ = 0
  · refine ⟨1, by norm_num, ?_, ?_⟩
    · have hsmall_zero :
          C_DeGiorgi_subsolution d A * Real.sqrt I₀ = 0 := by
        simp [I₀, hI₀_zero]
      rw [hsmall_zero]
      norm_num
    · filter_upwards [hbound] with x hx
      have hsqrt_zero :
          Real.sqrt (∫ x in Metric.ball (0 : E) 1, |max (u x) 0| ^ 2 ∂volume) = 0 := by
        rw [show (∫ x in Metric.ball (0 : E) 1, |max (u x) 0| ^ 2 ∂volume) = I₀ by rfl, hI₀_zero]
        simp
      have hsmall_zero' :
          C_DeGiorgi_subsolution d A *
              Real.sqrt (∫ x in Metric.ball (0 : E) 1, |max (u x) 0| ^ 2 ∂volume) = 0 := by
        rw [hsqrt_zero]
        ring
      have hx_pair :
          u x ≤
              C_DeGiorgi_subsolution d A *
                Real.sqrt (∫ x in Metric.ball (0 : E) 1, |max (u x) 0| ^ 2 ∂volume) ∧
            0 ≤
              C_DeGiorgi_subsolution d A *
                Real.sqrt (∫ x in Metric.ball (0 : E) 1, |max (u x) 0| ^ 2 ∂volume) := by
        simpa using hx
      rw [hsmall_zero'] at hx_pair
      have hx_zero : max (u x) 0 ≤ 0 := by
        exact max_le_iff.mpr ⟨hx_pair.1, le_rfl⟩
      linarith
  · let lamStar : ℝ := C_DeGiorgi_subsolution d A * Real.sqrt I₀
    have hLamStar : 0 < lamStar := by
      have hI₀_pos : 0 < I₀ := lt_of_le_of_ne hI₀_nonneg (by simpa [eq_comm] using hI₀_zero)
      dsimp [lamStar, I₀]
      exact deGiorgi_exact_height_threshold_on_unitBall
        (d := d) (K := K_DeGiorgi_subsolution d A)
        (K_DeGiorgi_subsolution_pos (d := d) A) hI₀_pos
    refine ⟨lamStar, hLamStar, le_rfl, ?_⟩
    simpa [lamStar, I₀] using hbound

\end{lstlisting}
\subsection{Normalised constants}

\begin{definition}[\leanname{K\_DeGiorgi\_subsolution\_normalized},
\leanname{C\_DeGiorgi\_subsolution\_normalized}]
The dimension-only counterparts are
\[
  K^{\mathrm{norm}}_{\mathrm{DG}}(d) :=
    2 \cdot (C_{\mathrm{gns}}(d,2))^2 \cdot 8\, (2 M_{\mathrm{st}})^2 + 1,
\]
\[
  C^{\mathrm{norm}}_{\mathrm{DG}}(d) :=
    C_{\mathrm{small}}(d, K^{\mathrm{norm}}_{\mathrm{DG}}(d)).
\]
\end{definition}

\begin{lstlisting}[style=Matlab-editor,basicstyle=\mlttfamily,escapechar=",mlshowsectionrules=true,mathescape=true,morecomment={[s]{\%\{}{\%\}}},language=lean,numbers=none]
noncomputable def K_DeGiorgi_subsolution_normalized
    (d : ℕ) [NeZero d] : ℝ :=
  2 * (C_gns d 2) ^ 2 * (8 * (2 * (Mst : ℝ)) ^ 2) + 1

noncomputable def C_DeGiorgi_subsolution_normalized
    (d : ℕ) [NeZero d] : ℝ :=
  C_DeGiorgiSmallness d (K_DeGiorgi_subsolution_normalized d)

\end{lstlisting}
\begin{lemma}[\leanname{C\_DeGiorgi\_subsolution\_le\_normalized}]
For a normalised elliptic coefficient field $A$ on $\Bz{1}$ (i.e.\
$\lambda = 1$, so $\Lambda \ge 1$ and the ratio
$\Lambda/\lambda = \Lambda$),
\[
  C_{\mathrm{DG}}(d, A) \le
    C^{\mathrm{norm}}_{\mathrm{DG}}(d) \cdot \Lambda^{d/4}.
\]
\end{lemma}

\begin{lstlisting}[style=Matlab-editor,basicstyle=\mlttfamily,escapechar=",mlshowsectionrules=true,mathescape=true,morecomment={[s]{\%\{}{\%\}}},language=lean,numbers=none]
theorem C_DeGiorgi_subsolution_le_normalized
    (A : NormalizedEllipticCoeff d (Metric.ball (0 : E) 1)) :
    C_DeGiorgi_subsolution d A.1 ≤
      C_DeGiorgi_subsolution_normalized d * A.1.Λ ^ ((d : ℝ) / 4) := by
  let c : ℝ := (C_gns d 2) ^ 2 * (8 * (2 * (Mst : ℝ)) ^ 2)
  have hc_nonneg : 0 ≤ c := by
    dsimp [c]
    positivity
  have hΛ_ge_one : 1 ≤ A.1.Λ := by
    simpa [A.2] using A.1.hΛ
  have hratio_eq : ellipticityRatio A.1 = A.1.Λ := by
    exact NormalizedEllipticCoeff.ellipticityRatio_eq_Λ (A := A)
  have hK_formula :
      K_DeGiorgi_subsolution d A.1 = c * (A.1.Λ + 1) + 1 := by
    rw [K_DeGiorgi_subsolution, hratio_eq]
    dsimp [c]
    ring
  have hK0_formula :
      K_DeGiorgi_subsolution_normalized d = 2 * c + 1 := by
    rw [K_DeGiorgi_subsolution_normalized]
    dsimp [c]
    ring
  have hK_le :
      K_DeGiorgi_subsolution d A.1 ≤
        K_DeGiorgi_subsolution_normalized d * A.1.Λ := by
    rw [hK_formula, hK0_formula]
    nlinarith
  have hmono :
      C_DeGiorgi_subsolution d A.1 ≤
        C_DeGiorgiSmallness d (K_DeGiorgi_subsolution_normalized d * A.1.Λ) := by
    rw [C_DeGiorgi_subsolution]
    exact C_DeGiorgiSmallness_mono (d := d)
      (K_DeGiorgi_subsolution_pos (d := d) A.1).le hK_le
  calc
    C_DeGiorgi_subsolution d A.1 ≤
        C_DeGiorgiSmallness d (K_DeGiorgi_subsolution_normalized d * A.1.Λ) := hmono
    _ = C_DeGiorgi_subsolution_normalized d * A.1.Λ ^ ((d : ℝ) / 4) := by
        rw [C_DeGiorgi_subsolution_normalized]
        exact C_DeGiorgiSmallness_mul_right (d := d)
          (K_DeGiorgi_subsolution_normalized_pos (d := d)).le A.1.Λ_nonneg

\end{lstlisting}
\begin{proof}
Comparing the formulas, $K_{\mathrm{DG}}(d,A) = c (\Lambda+1) + 1$ and
$K^{\mathrm{norm}}_{\mathrm{DG}}(d) = 2c + 1$ where $c =
(C_{\mathrm{gns}}(d,2))^2 \cdot 8 (2 M_{\mathrm{st}})^2$. Since
$\Lambda \ge 1$ one has $c(\Lambda + 1) + 1 \le (2c+1)\Lambda$,
i.e.\ $K_{\mathrm{DG}}(d,A) \le K^{\mathrm{norm}}_{\mathrm{DG}}(d) \cdot
\Lambda$. Monotonicity of $C_{\mathrm{small}}$ in $K$ followed by the
multiplicative scaling identity
$C_{\mathrm{small}}(d, K \cdot \Lambda) = C_{\mathrm{small}}(d, K)\,
\Lambda^{d/4}$ gives the claim.
\end{proof}

\begin{theorem}[Normalised De Giorgi $L^\infty$ bound,
\leanname{linfty\_subsolution\_DeGiorgi\_normalized}]
Let $d > 2$, $A$ a normalised elliptic coefficient field on
$\Bz{1}$, and $u$ a subsolution with
$|\max(u,0)|^2 \in L^1(\Bz{1})$. Then
\[
  \max(u(x), 0) \le
    C^{\mathrm{norm}}_{\mathrm{DG}}(d) \cdot \Lambda^{d/4} \cdot
    \Bigl(\int_{\Bz{1}} |\max(u,0)|^2 \, dx\Bigr)^{1/2}
\]
for almost every $x \in \Ball{1/2}{0}$.
\end{theorem}

\begin{lstlisting}[style=Matlab-editor,basicstyle=\mlttfamily,escapechar=",mlshowsectionrules=true,mathescape=true,morecomment={[s]{\%\{}{\%\}}},language=lean,numbers=none]
/-- Normalized-coefficient De Giorgi `L∞` bound with the explicit scaling
`C(d) * Λ^(d/4)` on the unit ball. -/
theorem linfty_subsolution_DeGiorgi_normalized
    (hd : 2 < (d : ℝ))
    (A : NormalizedEllipticCoeff d (Metric.ball (0 : E) 1))
    {u : E → ℝ}
    (hsub : IsSubsolution A.1 u)
    (hposInt :
      IntegrableOn (fun x => |max (u x) 0| ^ 2)
        (Metric.ball (0 : E) 1) volume) :
    ∀ᵐ x ∂(volume.restrict (Metric.ball (0 : E) (1 / 2 : ℝ))),
      max (u x) 0 ≤
        C_DeGiorgi_subsolution_normalized d *
          A.1.Λ ^ ((d : ℝ) / 4) *
          Real.sqrt (∫ x in Metric.ball (0 : E) 1, |max (u x) 0| ^ 2 ∂volume) := by
  let I₀ : ℝ := ∫ x in Metric.ball (0 : E) 1, |max (u x) 0| ^ 2 ∂volume
  have hsqrt_nonneg : 0 ≤ Real.sqrt I₀ := Real.sqrt_nonneg I₀
  have hconst :
      C_DeGiorgi_subsolution d A.1 * Real.sqrt I₀ ≤
        C_DeGiorgi_subsolution_normalized d *
          A.1.Λ ^ ((d : ℝ) / 4) * Real.sqrt I₀ :=
    mul_le_mul_of_nonneg_right
      (C_DeGiorgi_subsolution_le_normalized (d := d) A) hsqrt_nonneg
  filter_upwards [linfty_subsolution_DeGiorgi (d := d) hd A.1 hsub hposInt] with x hx
  exact le_trans hx hconst

\end{lstlisting}
\begin{proof}
Combine \leanname{linfty\_subsolution\_DeGiorgi} with the previous
lemma comparing $C_{\mathrm{DG}}(d,A)$ to
$C^{\mathrm{norm}}_{\mathrm{DG}}(d) \, \Lambda^{d/4}$, multiplying
by the nonnegative factor $\sqrt{I_0}$.
\end{proof}

This concludes the formal De Giorgi $L^2 \to L^\infty$ chapter. The
two top-level interfaces \leanname{DeGiorgiIteration} and
\leanname{DeGiorgiTheory} simply re-export the constructions and
theorems above so that downstream chapters (weak Harnack, Harnack,
Holder) can rely on the De Giorgi bound directly.


%% ========================================================================
\section{Moser Iteration and Supersolution Estimates}\label{sec:moser}
%% ========================================================================

%% --- Section 7 ---

\subsection{Moser Iteration}

Throughout this section we fix a dimension $d \in \N$ with $d \geq 1$ and
$d > 2$, and write $E = \R^d$ for the ambient space. All balls are taken
around the origin in $E$, and we abbreviate $\Bz{r} = \Ball{r}{0}$. Functions
$u : E \to \R$ are real valued; $A$ denotes a normalized elliptic coefficient
field on $\Bz{1}$, with ellipticity ratio $\Lambda = \Lambda_A \geq 1$
\leanname{NormalizedEllipticCoeff}. The notion of subsolution is
\leanname{IsSubsolution} and that of supersolution is
\leanname{IsSupersolution}. The starting point is the normalized De Giorgi
$L^2 \to L^\infty$ estimate proved in Section 5.

\subsubsection{Moser constants}

\begin{definition}[Moser anchor constant, \leanname{C\_MoserAnchor}]
The \emph{Moser anchor} in dimension $d$ is
\[
  C_{\mathrm{MA}}(d) := \max\!\left\{1,\;
    C_{\mathrm{DG}}^{\mathrm{sub}}(d)^{2},\;
    8\,\Mst^{2}\,C_{\mathrm{gns}}(d,2)^{2}\right\},
\]
where $C_{\mathrm{DG}}^{\mathrm{sub}}(d)$ is the normalized De Giorgi
subsolution constant of Section 5 and $C_{\mathrm{gns}}(d,2)$ is the
Gagliardo--Nirenberg--Sobolev constant for $p=2$.
\end{definition}

\begin{lstlisting}[style=Matlab-editor,basicstyle=\mlttfamily,escapechar=",mlshowsectionrules=true,mathescape=true,morecomment={[s]{\%\{}{\%\}}},language=lean,numbers=none]
/-- The anchor constant coming from the normalized De Giorgi `L² → L∞` bound. -/
noncomputable def C_MoserAnchor (d : ℕ) [NeZero d] : ℝ :=
  max 1
    (max ((C_DeGiorgi_subsolution_normalized d) ^ 2)
      (8 * (Mst : ℝ) ^ 2 * (C_gns d 2) ^ 2))

\end{lstlisting}
\begin{lemma}[\leanname{one\_le\_C\_MoserAnchor}]
$1 \leq C_{\mathrm{MA}}(d)$.
\end{lemma}

\begin{lstlisting}[style=Matlab-editor,basicstyle=\mlttfamily,escapechar=",mlshowsectionrules=true,mathescape=true,morecomment={[s]{\%\{}{\%\}}},language=lean,numbers=none]
theorem one_le_C_MoserAnchor : 1 ≤ C_MoserAnchor d := by
  simp [C_MoserAnchor]

\end{lstlisting}
\begin{proof}
Immediate from $C_{\mathrm{MA}}(d) \geq \max\{1,\dots\} \geq 1$.
\end{proof}

\begin{definition}[Sobolev gain factor and decay ratio]
Set
\[
  \chi := \frac{d}{d-2} \quad (\text{\leanname{moserChi}}),
  \qquad
  q := \frac{d-2}{d} \quad (\text{\leanname{moserDecayRatio}}).
\]
\end{definition}

\begin{lstlisting}[style=Matlab-editor,basicstyle=\mlttfamily,escapechar=",mlshowsectionrules=true,mathescape=true,morecomment={[s]{\%\{}{\%\}}},language=lean,numbers=none]
/-- The Sobolev gain factor `χ = d/(d-2)` in the elliptic Moser iteration. -/
noncomputable def moserChi (d : ℕ) [NeZero d] : ℝ :=
  (d : ℝ) / ((d : ℝ) - 2)

/-- The geometric decay ratio `q = (d-2)/d` used in the Moser iteration. -/
noncomputable def moserDecayRatio (d : ℕ) [NeZero d] : ℝ :=
  ((d : ℝ) - 2) / (d : ℝ)

\end{lstlisting}
\begin{lemma}[Basic algebraic facts on $\chi,q$]
Assume $d > 2$. Then
\begin{enumerate}
  \item $0 \leq q$ \leanname{moserDecayRatio\_nonneg};
  \item $q < 1$ \leanname{moserDecayRatio\_lt\_one};
  \item $0 < \chi$ \leanname{moserChi\_pos};
  \item $1 < \chi$ \leanname{one\_lt\_moserChi};
  \item $q = \chi^{-1}$ \leanname{moserDecayRatio\_eq\_inv\_moserChi}.
\end{enumerate}
\end{lemma}

\begin{lstlisting}[style=Matlab-editor,basicstyle=\mlttfamily,escapechar=",mlshowsectionrules=true,mathescape=true,morecomment={[s]{\%\{}{\%\}}},language=lean,numbers=none]
theorem moserDecayRatio_nonneg (hd : 2 < (d : ℝ)) :
    0 ≤ moserDecayRatio d := by
  have hd_nonneg : 0 ≤ ((d : ℝ) - 2) := by linarith
  have hd_pos : 0 < (d : ℝ) := by
    exact_mod_cast (Nat.pos_of_ne_zero (NeZero.ne d))
  exact div_nonneg hd_nonneg hd_pos.le

theorem moserDecayRatio_lt_one (_hd : 2 < (d : ℝ)) :
    moserDecayRatio d < 1 := by
  have hd_pos : 0 < (d : ℝ) := by
    exact_mod_cast (Nat.pos_of_ne_zero (NeZero.ne d))
  rw [moserDecayRatio]
  have : ((d : ℝ) - 2) < (d : ℝ) := by linarith
  exact (div_lt_one hd_pos).2 this

theorem moserChi_pos (hd : 2 < (d : ℝ)) :
    0 < moserChi d := by
  rw [moserChi]
  have hnum_pos : 0 < (d : ℝ) := by
    exact_mod_cast (Nat.pos_of_ne_zero (NeZero.ne d))
  have hden_pos : 0 < (d : ℝ) - 2 := by linarith
  exact div_pos hnum_pos hden_pos

theorem one_lt_moserChi (hd : 2 < (d : ℝ)) :
    1 < moserChi d := by
  rw [moserChi]
  have hden_pos : 0 < (d : ℝ) - 2 := by linarith
  have : (d : ℝ) - 2 < (d : ℝ) := by linarith
  exact (one_lt_div hden_pos).2 this

theorem moserDecayRatio_eq_inv_moserChi (hd : 2 < (d : ℝ)) :
    moserDecayRatio d = (moserChi d)⁻¹ := by
  rw [moserDecayRatio, moserChi]
  have hd_ne : (d : ℝ) ≠ 0 := by
    exact_mod_cast (NeZero.ne d)
  have hdm2_ne : (d : ℝ) - 2 ≠ 0 := by
    linarith
  field_simp [hd_ne, hdm2_ne]

\end{lstlisting}
\begin{proof}
For (1)--(2) note $0 \leq d-2 < d$, so $0 \leq (d-2)/d < 1$.
For (3)--(4), $d/(d-2) > 0$ and $d > d-2$ gives $d/(d-2) > 1$.
For (5), $\big((d-2)/d\big)^{-1} = d/(d-2)$.
\end{proof}

\begin{definition}[Dimensional Moser constant, \leanname{C\_Moser}]
Set $C_{\mathrm{MA}} := C_{\mathrm{MA}}(d)$. If $d > 2$, define
\[
  C_{\mathrm{Moser}}(d) := \max\!\left\{C_{\mathrm{MA}},\;
    (32\,C_{\mathrm{MA}})^{d/2}\,
    4^{\sum_{n=0}^{\infty} n\,q^{n}}\right\};
\]
otherwise $C_{\mathrm{Moser}}(d) := C_{\mathrm{MA}}$.
\end{definition}

\begin{lstlisting}[style=Matlab-editor,basicstyle=\mlttfamily,escapechar=",mlshowsectionrules=true,mathescape=true,morecomment={[s]{\%\{}{\%\}}},language=lean,numbers=none]
/-- Dimension-only constant `C(d)` for the basic Moser `L^p → L∞` estimate.

It is chosen large enough to absorb both the existing `p = 2` De Giorgi anchor
and the purely geometric product appearing in the exact Moser iteration. -/
noncomputable def C_Moser (d : ℕ) [NeZero d] : ℝ :=
  let base := C_MoserAnchor d
  if _hd : 2 < (d : ℝ) then
    max base
      (((32 : ℝ) * base) ^ ((d : ℝ) / 2) *
        4 ^ (∑' n : ℕ, (n : ℝ) * moserDecayRatio d ^ n))
  else
    base

\end{lstlisting}
\begin{lemma}[\leanname{C\_MoserAnchor\_le\_C\_Moser}, \leanname{one\_le\_C\_Moser}, \leanname{C\_DeGiorgi\_subsolution\_normalized\_sq\_le\_C\_Moser}]
$C_{\mathrm{MA}}(d) \leq C_{\mathrm{Moser}}(d)$, $1 \leq C_{\mathrm{Moser}}(d)$, and
$C_{\mathrm{DG}}^{\mathrm{sub}}(d)^2 \leq C_{\mathrm{Moser}}(d)$.
\end{lemma}

\begin{lstlisting}[style=Matlab-editor,basicstyle=\mlttfamily,escapechar=",mlshowsectionrules=true,mathescape=true,morecomment={[s]{\%\{}{\%\}}},language=lean,numbers=none]
theorem C_MoserAnchor_le_C_Moser :
    C_MoserAnchor d ≤ C_Moser d := by
  by_cases hd : 2 < (d : ℝ)
  · simp [C_Moser, hd]
  · simp [C_Moser, hd]

theorem one_le_C_Moser : 1 ≤ C_Moser d := by
  exact (one_le_C_MoserAnchor (d := d)).trans (C_MoserAnchor_le_C_Moser (d := d))

theorem C_DeGiorgi_subsolution_normalized_sq_le_C_Moser :
    (C_DeGiorgi_subsolution_normalized d) ^ 2 ≤ C_Moser d := by
  exact
    (C_DeGiorgi_subsolution_normalized_sq_le_C_MoserAnchor (d := d)).trans
      (C_MoserAnchor_le_C_Moser (d := d))

\end{lstlisting}
\begin{proof}
Both branches of the definition return a value $\geq C_{\mathrm{MA}}(d) \geq 1$.
The third statement follows by transitivity from the corresponding bound for
$C_{\mathrm{MA}}$.
\end{proof}

\subsubsection{Geometric sequences}

\begin{definition}[Nested Moser radii, \leanname{moserRadius}]
\[
  r_n := \frac{1 + 2^{-n}}{2}, \qquad n \geq 0.
\]
Thus $r_0 = 1$ and $r_n \searrow 1/2$.
\end{definition}

\begin{lstlisting}[style=Matlab-editor,basicstyle=\mlttfamily,escapechar=",mlshowsectionrules=true,mathescape=true,morecomment={[s]{\%\{}{\%\}}},language=lean,numbers=none]
/-- The standard nested radii `r_n = (1 + 2^{-n})/2`. -/
noncomputable def moserRadius (n : ℕ) : ℝ :=
  (1 + (1 / 2 : ℝ) ^ n) / 2

\end{lstlisting}
\begin{lemma}[Properties of the radii]
For all $n \in \N$:
\begin{enumerate}
  \item $r_0 = 1$ \leanname{moserRadius\_zero};
  \item $r_n \leq 1$ \leanname{moserRadius\_le\_one};
  \item $0 < r_n$ \leanname{moserRadius\_pos};
  \item $1/2 < r_n$ \leanname{moserRadius\_gt\_half};
  \item $r_n - r_{n+1} = 2^{-(n+2)}$ \leanname{moserRadius\_gap};
  \item $r_{n+1} < r_n$ \leanname{moserRadius\_succ\_lt}.
\end{enumerate}
\end{lemma}

\begin{lstlisting}[style=Matlab-editor,basicstyle=\mlttfamily,escapechar=",mlshowsectionrules=true,mathescape=true,morecomment={[s]{\%\{}{\%\}}},language=lean,numbers=none]
theorem moserRadius_zero :
    moserRadius 0 = 1 := by
  simp [moserRadius]

theorem moserRadius_le_one (n : ℕ) :
    moserRadius n ≤ 1 := by
  have hpow_le : (1 / 2 : ℝ) ^ n ≤ 1 := by
    exact pow_le_one₀ (by norm_num : (0 : ℝ) ≤ 1 / 2) (by norm_num : (1 / 2 : ℝ) ≤ 1)
  have hpow_nonneg : 0 ≤ (1 / 2 : ℝ) ^ n := by positivity
  rw [moserRadius]
  nlinarith

theorem moserRadius_pos (n : ℕ) :
    0 < moserRadius n := by
  rw [moserRadius]
  positivity

/-- Auxiliary closeout: a uniform bound on the iterated `L^{p_n}` quantities gives
the a.e. half-ball `L∞` estimate. -/
theorem moserRadius_gt_half (n : ℕ) :
    (1 / 2 : ℝ) < moserRadius n := by
  rw [moserRadius]
  have hpow_pos : 0 < (1 / 2 : ℝ) ^ n := by positivity
  nlinarith

theorem moserRadius_gap (n : ℕ) :
    moserRadius n - moserRadius (n + 1) = (1 / 2 : ℝ) ^ (n + 2) := by
  rw [moserRadius, moserRadius]
  rw [pow_succ]
  field_simp
  ring_nf

theorem moserRadius_succ_lt (n : ℕ) :
    moserRadius (n + 1) < moserRadius n := by
  have hgap : 0 < moserRadius n - moserRadius (n + 1) := by
    rw [moserRadius_gap]
    positivity
  linarith

\end{lstlisting}
\begin{proof}
Each follows by direct algebra: e.g.\ for (5),
$r_n - r_{n+1} = (2^{-n} - 2^{-(n+1)})/2 = 2^{-(n+2)}$.
\end{proof}

\begin{definition}[Exponent sequence, \leanname{moserExponentSeq}]
Given $p_0 \in \R$ and $n \in \N$, set
\[
  p_n := p_0 \chi^{n}.
\]
\end{definition}

\begin{lstlisting}[style=Matlab-editor,basicstyle=\mlttfamily,escapechar=",mlshowsectionrules=true,mathescape=true,morecomment={[s]{\%\{}{\%\}}},language=lean,numbers=none]
/-- The exponent sequence `p_{n+1} = χ p_n`. -/
noncomputable def moserExponentSeq (d : ℕ) [NeZero d] (p₀ : ℝ) (n : ℕ) : ℝ :=
  p₀ * moserChi d ^ n

\end{lstlisting}
\begin{lemma}
With $\chi = d/(d-2)$ and $q = \chi^{-1}$:
\begin{enumerate}
  \item $p_0 = p_0$ and $p_{n+1} = \chi p_n$ \leanname{moserExponentSeq\_zero}, \leanname{moserExponentSeq\_succ};
  \item if $d>2$ and $p_0 > 0$, then $p_n > 0$ \leanname{moserExponentSeq\_pos};
  \item if $d>2$ and $p_0 \geq 0$, then $p_0 \leq p_n$ \leanname{moserExponentSeq\_ge\_initial};
  \item if $d>2$ and $p_0 > 1$, then $p_n > 1$ \leanname{one\_lt\_moserExponentSeq};
  \item if $d>2$ and $p_0 > 0$, then $1/p_n = q^n / p_0$ \leanname{inv\_moserExponentSeq}.
\end{enumerate}
\end{lemma}

\begin{lstlisting}[style=Matlab-editor,basicstyle=\mlttfamily,escapechar=",mlshowsectionrules=true,mathescape=true,morecomment={[s]{\%\{}{\%\}}},language=lean,numbers=none]
theorem moserExponentSeq_zero (p₀ : ℝ) :
    moserExponentSeq d p₀ 0 = p₀ := by
  simp [moserExponentSeq, moserChi]

theorem moserExponentSeq_succ (p₀ : ℝ) (n : ℕ) :
    moserExponentSeq d p₀ (n + 1) = moserChi d * moserExponentSeq d p₀ n := by
  rw [moserExponentSeq, pow_succ, moserExponentSeq]
  ring

theorem moserExponentSeq_pos (hd : 2 < (d : ℝ)) {p₀ : ℝ} (hp₀ : 0 < p₀) (n : ℕ) :
    0 < moserExponentSeq d p₀ n := by
  rw [moserExponentSeq]
  exact mul_pos hp₀ (pow_pos (moserChi_pos (d := d) hd) n)

theorem moserExponentSeq_ge_initial (hd : 2 < (d : ℝ)) {p₀ : ℝ} (hp₀ : 0 ≤ p₀)
    (n : ℕ) :
    p₀ ≤ moserExponentSeq d p₀ n := by
  rw [moserExponentSeq]
  have hpow_ge : 1 ≤ moserChi d ^ n := by
    exact one_le_pow₀ (one_le_moserChi (d := d) hd)
  nlinarith

theorem one_lt_moserExponentSeq (hd : 2 < (d : ℝ)) {p₀ : ℝ} (hp₀ : 1 < p₀) (n : ℕ) :
    1 < moserExponentSeq d p₀ n := by
  have hp₀_nonneg : 0 ≤ p₀ := by linarith
  have hp_ge :
      p₀ ≤ moserExponentSeq d p₀ n :=
    moserExponentSeq_ge_initial (d := d) hd hp₀_nonneg n
  linarith

theorem inv_moserExponentSeq (hd : 2 < (d : ℝ)) {p₀ : ℝ} (hp₀ : 0 < p₀) (n : ℕ) :
    1 / moserExponentSeq d p₀ n = moserDecayRatio d ^ n / p₀ := by
  rw [moserExponentSeq]
  have hp₀_ne : p₀ ≠ 0 := hp₀.ne'
  have hchi_pos : 0 < moserChi d := moserChi_pos (d := d) hd
  have hchi_ne : moserChi d ≠ 0 := hchi_pos.ne'
  have hq_eq : moserDecayRatio d ^ n = (moserChi d ^ n)⁻¹ := by
    rw [moserDecayRatio_eq_inv_moserChi (d := d) hd, inv_pow]
  rw [hq_eq]
  field_simp [hp₀_ne, pow_ne_zero n hchi_ne]

\end{lstlisting}
\begin{proof}
All four are immediate consequences of $\chi > 1$ and $p_n = p_0 \chi^n$;
identity (5) uses $\chi^{-n} = q^n$.
\end{proof}

\subsubsection{Regularized power cutoffs}

For the cutoff--power test functions of the iteration we set
\begin{equation*}
  \mathrm{PC}_d(\eta,u,p)(x) := \eta(x)\,|\!\max(u(x),0)|^{p/2}
  \quad (\text{\leanname{moserPowerCutoff}}),
\end{equation*}
the smooth cutoff $\eta$ being supported in a slightly larger ball than the
inner ball where the bound is sought.

\begin{lemma}[Support property, \leanname{moserPowerCutoff\_tsupport\_subset}]
If $\supp \eta \subseteq \Omega$, then
$\supp \mathrm{PC}_d(\eta,u,p) \subseteq \Omega$.
\end{lemma}

\begin{lstlisting}[style=Matlab-editor,basicstyle=\mlttfamily,escapechar=",mlshowsectionrules=true,mathescape=true,morecomment={[s]{\%\{}{\%\}}},language=lean,numbers=none]
omit [NeZero d] in
theorem moserPowerCutoff_tsupport_subset
    {u η : E → ℝ} {p s : ℝ}
    (hη_sub_ball : tsupport η ⊆ Metric.ball (0 : E) s) :
    tsupport (moserPowerCutoff (d := d) η u p) ⊆ Metric.ball (0 : E) s := by
  exact (tsupport_mul_subset_left
    (f := η) (g := fun x => |max (u x) 0| ^ (p / 2))).trans hη_sub_ball

\end{lstlisting}
\begin{proof}
Pointwise $\mathrm{PC}_d(\eta,u,p)(x) = 0$ whenever $\eta(x)=0$.
\end{proof}

\subsubsection{Sobolev preparation and the energy bound}

The hard nonlinear core of the Moser iteration is the construction of a
$W^{1,2}$ witness for the cutoff power $\mathrm{PC}_d(\eta,u,p)$ together with
the matching energy estimate.

\begin{theorem}[Cutoff--power energy bound, \leanname{moser\_powerCutoff\_memW1p\_energy\_of\_subsolution\_core}]
\label{thm:moser-cutoff-energy}
Let $d > 2$, $A$ a normalized elliptic coefficient on $\Bz{1}$, $u$ a
subsolution of $A$, and $\eta \in C^{\infty}_c(\Bz{s})$ a smooth cutoff with
$0 \leq \eta \leq 1$ and $\|\nabla \eta\|_{L^\infty} \leq C_\eta$. Suppose
$0 < s \leq 1$, $1 < p$, and $\int_{\Bz{s}} |u_+|^p < \infty$. Then there
exists a $W^{1,2}$ witness for $v := \mathrm{PC}_d(\eta,u,p)$ on $\Bz{s}$
satisfying
\begin{equation*}
  \int_{\Bz{s}} \|\nabla v\|^{2}\,dx
  \;\leq\;
  2\,C_\eta^{2}\,\Big(\Lambda\,\Big(\frac{p}{p-1}\Big)^{2} + 1\Big)
  \int_{\Bz{s}} |u_+|^{p}\,dx.
\end{equation*}
\end{theorem}

\begin{lstlisting}[style=Matlab-editor,basicstyle=\mlttfamily,escapechar=",mlshowsectionrules=true,mathescape=true,morecomment={[s]{\%\{}{\%\}}},language=lean,numbers=none]
/-- Hard nonlinear core of the Moser pre-estimate: build a `W^{1,2}` witness for
the cutoff-power `η · (u_+)^(p/2)` and prove the exact energy bound. The
zero-trace conclusion is derived separately from compact support. -/
theorem moser_powerCutoff_memW1p_energy_of_subsolution_core
    (hd : 2 < (d : ℝ))
    (A : NormalizedEllipticCoeff d (Metric.ball (0 : E) 1))
    {u η : E → ℝ} {p s Cη : ℝ}
    (hp : 1 < p)
    (hs : 0 < s) (hs1 : s ≤ 1)
    (hsub : IsSubsolution A.1 u)
    (hpInt :
      IntegrableOn (fun x => |max (u x) 0| ^ p)
        (Metric.ball (0 : E) s) volume)
    (hη : ContDiff ℝ (⊤ : ℕ∞) η)
    (hη_nonneg : ∀ x, 0 ≤ η x)
    (hη_bound : ∀ x, |η x| ≤ 1)
    (hη_grad_bound : ∀ x, ‖fderiv ℝ η x‖ ≤ Cη)
    (hη_sub_ball : tsupport η ⊆ Metric.ball (0 : E) s) :
    ∃ hwv : MemW1pWitness 2 (moserPowerCutoff (d := d) η u p) (Metric.ball (0 : E) s),
      ∫ x in Metric.ball (0 : E) s, ‖hwv.weakGrad x‖ ^ 2 ∂volume ≤
        2 * Cη ^ 2 * (A.1.Λ * (p / (p - 1)) ^ 2 + 1) *
          ∫ x in Metric.ball (0 : E) s, |max (u x) 0| ^ p ∂volume := by
  exact moserPowerCutoff_memW1pWitness (d := d) hd A hp hs hs1 hsub hpInt hη hη_nonneg
    hη_bound hη_grad_bound hη_sub_ball

\end{lstlisting}
\begin{proof}
The witness for $v := \mathrm{PC}_d(\eta, u, p) = \eta\,(u_+)^{p/2}$ is
constructed in \leanname{moserPowerCutoff\_memW1pWitness}, which we now
unpack. First, a $W^{1,2}$ witness for $(u_+)^{p/2}$ on $\Bz{s}$ is obtained
from the truncation $u_+$ (whose witness comes from
\leanname{positivePartSub\_memW1pWitness\_on\_ball}) composed with the
power $t \mapsto t^{p/2}$ via the $W^{1,2}$ chain rule
\leanname{sobolev\_chain\_rule\_unitBall}; the local integrability
$\int (u_+)^{p} < \infty$ together with $p > 1$ ensures the chain-rule
hypotheses (locally Lipschitz factor on $[0,\infty)$). The product
$\eta\,(u_+)^{p/2}$ is then an honest $W^{1,2}$ function via
\leanname{MemW1pWitness.mul\_smooth\_bounded\_p}, with weak gradient
$\nabla v = \eta\,(p/2)(u_+)^{p/2-1}\nabla u_+ + (u_+)^{p/2}\nabla \eta$.

For the energy bound, test the subsolution inequality against
$\varphi := \eta^{2}(u_+)^{p-1}$ (this is in $H^{1}_{0}(\Bz{s})$ thanks to
the support inclusion of $\eta$ and the integrability hypothesis on $u_+$).
The pointwise expansion of $\langle A\nabla u, \nabla \varphi\rangle$ via
\leanname{bilinFormIntegrand\_mul\_smooth\_eq} yields
\begin{multline*}
  \int_{\Bz{s}} (p-1)\,\eta^{2}\,(u_+)^{p-2}\,\langle A\nabla u_+, \nabla u_+\rangle\,dx
  \\
  \le -2\int_{\Bz{s}} \eta\,(u_+)^{p-1}\,\langle A\nabla u_+, \nabla \eta\rangle\,dx,
\end{multline*}
where on the left we discarded the supersolution-favourable term and on the
right we used the subsolution inequality. The right-hand side is bounded by
Cauchy--Schwarz against $\langle A\cdot, \cdot\rangle$ followed by AM--GM
with weight $\varepsilon = (p-1)/2$:
\[
  2|\eta\,(u_+)^{p-1}\langle A\nabla u_+,\nabla \eta\rangle|
  \le \frac{p-1}{2}\eta^{2}(u_+)^{p-2}\langle A\nabla u_+,\nabla u_+\rangle
  + \frac{2}{p-1}(u_+)^{p}\langle A\nabla \eta,\nabla \eta\rangle.
\]
Absorbing the first term into the left-hand side leaves
$(p-1)/2 \cdot \int \eta^{2}(u_+)^{p-2}\langle A\nabla u_+,\nabla u_+\rangle$,
and the residual term is bounded by $\Lambda C_\eta^{2}\int (u_+)^{p}$ via
the upper ellipticity bound for $A$ and $\|\nabla \eta\| \le C_\eta$.

Finally, expanding $\nabla v = \eta\,(p/2)(u_+)^{p/2-1}\nabla u_+ +
(u_+)^{p/2}\nabla \eta$, the pointwise inequality
$\|\nabla v\|^{2} \le 2\,\eta^{2}(p/2)^{2}(u_+)^{p-2}\|\nabla u_+\|^{2}
+ 2(u_+)^{p}\|\nabla \eta\|^{2}$ together with the absorbed Caccioppoli
estimate gives
\[
  \int_{\Bz{s}} \|\nabla v\|^{2}\,dx
   \le 2\,C_\eta^{2}\Bigl(\Lambda\bigl(\tfrac{p}{p-1}\bigr)^{2} + 1\Bigr)
   \int_{\Bz{s}} (u_+)^{p}\,dx,
\]
the displayed energy inequality.
\end{proof}

\begin{theorem}[Zero trace upgrade, \leanname{moser\_powerCutoff\_memW01p\_energy\_of\_subsolution}]
Under the hypotheses of Theorem~\ref{thm:moser-cutoff-energy}, the witness
moreover satisfies $v \in W^{1,2}_{0}(\Bz{s})$.
\end{theorem}

\begin{lstlisting}[style=Matlab-editor,basicstyle=\mlttfamily,escapechar=",mlshowsectionrules=true,mathescape=true,morecomment={[s]{\%\{}{\%\}}},language=lean,numbers=none]
/-- Analytic core of the Moser pre-estimate: the cutoff-power
`η · (u_+)^(p/2)` belongs to `W₀^{1,2}` on the outer ball and satisfies the
exact energy bound needed for Sobolev. The only remaining nonlinear gap is the
construction of the underlying `W^{1,2}` witness and its energy bound. -/
theorem moser_powerCutoff_memW01p_energy_of_subsolution
    (hd : 2 < (d : ℝ))
    (A : NormalizedEllipticCoeff d (Metric.ball (0 : E) 1))
    {u η : E → ℝ} {p s Cη : ℝ}
    (hp : 1 < p)
    (hs : 0 < s) (hs1 : s ≤ 1)
    (hsub : IsSubsolution A.1 u)
    (hpInt :
      IntegrableOn (fun x => |max (u x) 0| ^ p)
        (Metric.ball (0 : E) s) volume)
    (hη : ContDiff ℝ (⊤ : ℕ∞) η)
    (hη_nonneg : ∀ x, 0 ≤ η x)
    (hη_bound : ∀ x, |η x| ≤ 1)
    (hη_grad_bound : ∀ x, ‖fderiv ℝ η x‖ ≤ Cη)
    (hη_sub_ball : tsupport η ⊆ Metric.ball (0 : E) s) :
    ∃ hwv : MemW1pWitness 2 (moserPowerCutoff (d := d) η u p) (Metric.ball (0 : E) s),
      MemW01p 2 (moserPowerCutoff (d := d) η u p) (Metric.ball (0 : E) s) ∧
      ∫ x in Metric.ball (0 : E) s, ‖hwv.weakGrad x‖ ^ 2 ∂volume ≤
        2 * Cη ^ 2 * (A.1.Λ * (p / (p - 1)) ^ 2 + 1) *
          ∫ x in Metric.ball (0 : E) s, |max (u x) 0| ^ p ∂volume := by
  let Ω : Set E := Metric.ball (0 : E) s
  let v : E → ℝ := moserPowerCutoff (d := d) η u p
  have hv_support : tsupport v ⊆ Ω :=
    moserPowerCutoff_tsupport_subset (d := d) (u := u) (η := η) (p := p) hη_sub_ball
  obtain ⟨hwv, henergy⟩ :=
    moser_powerCutoff_memW1p_energy_of_subsolution_core
      (d := d) hd A (u := u) (η := η) (p := p) (s := s) (Cη := Cη)
      hp hs hs1 hsub hpInt hη hη_nonneg hη_bound hη_grad_bound hη_sub_ball
  have hv_compact : HasCompactSupport v :=
    hasCompactSupport_of_tsupport_subset_ball hv_support
  have hv_memW1p_real : MemW1p (ENNReal.ofReal (2 : ℝ)) v Ω := by
    simpa [Ω] using hwv.memW1p
  have hvW01_real : MemW01p (ENNReal.ofReal (2 : ℝ)) v Ω := by
    exact memW01p_of_memW1p_of_tsupport_subset
      (d := d) Metric.isOpen_ball (p := (2 : ℝ)) (by norm_num) hv_memW1p_real hv_compact hv_support
  have hvW01 : MemW01p 2 v Ω := by
    simpa [Ω] using hvW01_real
  exact ⟨hwv, hvW01, henergy⟩

\end{lstlisting}
\begin{proof}
Since $\supp \eta \subseteq \Bz{s}$ is compact, $\supp v \subseteq \Bz{s}$ is
likewise compactly contained. The criterion
\leanname{memW01p\_of\_memW1p\_of\_tsupport\_subset} then upgrades $v$ from
$W^{1,p}$ to $W^{1,p}_0$.
\end{proof}

\subsubsection{One-step Moser pre-estimate}

\begin{theorem}[Concentric ball pre-estimate, \leanname{moser\_preMoser}]
\label{thm:moser-preMoser}
Let $d > 2$, $A$ as above, $u$ a subsolution, and $0 < r < s \leq 1$.
Suppose $1 < p$ and $\int_{\Bz{s}} |u_+|^p < \infty$. Then $|u_+|^{\chi p}$
is integrable on $\Bz{r}$ and
\begin{multline*}
  \Big(\int_{\Bz{r}} |u_+|^{\chi p}\,dx\Big)^{1/(\chi p)}
  \\ \leq\;
  \Big(\frac{C_{\mathrm{MA}}(d)}{(s-r)^{2}}\,
    \Big(\Lambda\Big(\frac{p}{p-1}\Big)^{2} + 1\Big)\Big)^{1/p}
  \,\Big(\int_{\Bz{s}} |u_+|^{p}\,dx\Big)^{1/p}.
\end{multline*}
\end{theorem}

\begin{lstlisting}[style=Matlab-editor,basicstyle=\mlttfamily,escapechar=",mlshowsectionrules=true,mathescape=true,morecomment={[s]{\%\{}{\%\}}},language=lean,numbers=none]
/-- Auxiliary concentric-ball estimate for the Moser iteration.

The representative upgrade to pointwise `sup` statements is deferred to
Chapter 07. -/
theorem moser_preMoser
    (hd : 2 < (d : ℝ))
    (A : NormalizedEllipticCoeff d (Metric.ball (0 : E) 1))
    {u : E → ℝ} {p r s : ℝ}
    (hp : 1 < p)
    (hr : 0 < r) (hrs : r < s) (hs1 : s ≤ 1)
    (hsub : IsSubsolution A.1 u)
    (hpInt :
      IntegrableOn (fun x => |max (u x) 0| ^ p)
        (Metric.ball (0 : E) s) volume) :
    IntegrableOn (fun x => |max (u x) 0| ^ (moserChi d * p))
        (Metric.ball (0 : E) r) volume ∧
      (∫ x in Metric.ball (0 : E) r,
          |max (u x) 0| ^ (moserChi d * p) ∂volume) ^
          (1 / (moserChi d * p)) ≤
        ((C_MoserAnchor d / (s - r) ^ 2) *
            (A.1.Λ * (p / (p - 1)) ^ 2 + 1)) ^ (1 / p) *
          (∫ x in Metric.ball (0 : E) s, |max (u x) 0| ^ p ∂volume) ^ (1 / p) := by
  let Ω : Set E := Metric.ball (0 : E) s
  let μ : Measure E := volume.restrict Ω
  let R : ℝ := (r + s) / 2
  let η : E → ℝ := (Classical.choose (Cutoff.exists (d := d) (0 : E) (r := r) (R := R) hr (by
    dsimp [R]
    linarith))).toFun
  let Cη : ℝ := 2 * (Mst : ℝ) * (s - r)⁻¹
  have hs : 0 < s := lt_trans hr hrs
  have hsr_pos : 0 < s - r := by linarith
  have hR : r < R := by
    dsimp [R]
    linarith
  have hR_lt_s : R < s := by
    dsimp [R]
    linarith
  let ηCut : Cutoff (x₀ := (0 : E)) r R := Classical.choose
    (Cutoff.exists (d := d) (0 : E) (r := r) (R := R) hr hR)
  have hη : ContDiff ℝ (⊤ : ℕ∞) η := by
    simpa [η, ηCut] using ηCut.smooth
  have hη_bound : ∀ x, |η x| ≤ 1 := by
    intro x
    rw [abs_of_nonneg (ηCut.nonneg x)]
    simpa [η, ηCut] using ηCut.le_one x
  have hη_eq_one : ∀ x ∈ Metric.ball (0 : E) r, η x = 1 := by
    intro x hx
    simpa [η, ηCut] using ηCut.eq_one x hx
  have hη_grad_bound : ∀ x, ‖fderiv ℝ η x‖ ≤ Cη := by
    intro x
    have hmid :
        ((((r + s) / 2 : ℝ) - r)) = (s - r) / 2 := by
      ring
    have h_inv : ((((r + s) / 2 : ℝ) - r)⁻¹) = 2 * (s - r)⁻¹ := by
      rw [hmid]
      field_simp [hsr_pos.ne']
    calc
      ‖fderiv ℝ η x‖ ≤ ↑Mst * ((((r + s) / 2 : ℝ) - r)⁻¹) := by
        simpa [η, ηCut, R] using ηCut.grad_bound x
      _ = 2 * (Mst : ℝ) * (s - r)⁻¹ := by
        rw [h_inv]
        ring
      _ = Cη := by
        rfl
  have hη_sub_ball : tsupport η ⊆ Ω := by
    intro x hx
    exact (Metric.closedBall_subset_ball hR_lt_s) (ηCut.support_subset hx)
  let v : E → ℝ := moserPowerCutoff (d := d) η u p
  have hv_support : tsupport v ⊆ Ω :=
    moserPowerCutoff_tsupport_subset (d := d) (u := u) (η := η) (p := p) hη_sub_ball
  obtain ⟨hwv, hvW01, henergy⟩ :=
    moser_powerCutoff_memW01p_energy_of_subsolution
      (d := d) hd A (u := u) (η := η) (p := p) (s := s) (Cη := Cη)
      hp hs hs1 hsub hpInt hη ηCut.nonneg hη_bound hη_grad_bound hη_sub_ball
  obtain ⟨hwv_raw, hSob_raw⟩ :=
    sobolev_prepare_on_ball (d := d) (v := v) hd hs hvW01 hv_support
  let hwvSob : MemW1pWitness (ENNReal.ofReal (2 : ℝ)) v Ω :=
    { memLp := by
        simpa [Ω] using hwv_raw.memLp
      weakGrad := hwv_raw.weakGrad
      weakGrad_component_memLp := by
        simpa [Ω] using hwv_raw.weakGrad_component_memLp
      isWeakGrad := by
        simpa [Ω] using hwv_raw.isWeakGrad }
  let hwvReal : MemW1pWitness (ENNReal.ofReal (2 : ℝ)) v Ω :=
    { memLp := by
        simpa [Ω] using hwv.memLp
      weakGrad := hwv.weakGrad
      weakGrad_component_memLp := by
        simpa [Ω] using hwv.weakGrad_component_memLp
      isWeakGrad := by
        simpa [Ω] using hwv.isWeakGrad }
  have hSob :
      eLpNorm v
          (ENNReal.ofReal ((d : ℝ) * 2 / ((d : ℝ) - 2))) μ ≤
        ENNReal.ofReal (C_gns d 2) *
          eLpNorm (fun x => ‖hwvSob.weakGrad x‖) 2 μ := by
    simpa [Ω, μ, hwvSob] using hSob_raw
  have hae_grad :
      hwvSob.weakGrad =ᵐ[μ] hwvReal.weakGrad := by
    simpa [μ, Ω] using
      (MemW1pWitness.ae_eq_p (d := d) Metric.isOpen_ball (p := (2 : ℝ)) (by norm_num) hwvSob hwvReal)
  have hgrad_sq_eq :
      ∫ x in Ω, ‖hwvSob.weakGrad x‖ ^ 2 ∂volume =
        ∫ x in Ω, ‖hwvReal.weakGrad x‖ ^ 2 ∂volume := by
    exact integral_congr_ae (hae_grad.fun_comp (fun z => ‖z‖ ^ 2))
  let qexp : ℝ := ((d : ℝ) * 2 / ((d : ℝ) - 2))
  let q := ENNReal.ofReal qexp
  have hq_pos : 0 < qexp := by
    have hd_pos : 0 < (d : ℝ) := by linarith
    have hd_sub_pos : 0 < (d : ℝ) - 2 := by linarith
    dsimp [qexp]
    exact div_pos (mul_pos hd_pos (by norm_num)) hd_sub_pos
  have hv_memLp_q : MemLp v q μ := by
    refine ⟨hwvSob.memLp.aestronglyMeasurable, ?_⟩
    have hgrad_lt_top :
        eLpNorm (fun x => ‖hwvSob.weakGrad x‖) 2 μ < ⊤ :=
      lt_top_iff_ne_top.mpr (by simpa using hwvSob.weakGrad_norm_memLp.eLpNorm_ne_top)
    have hrhs_lt_top :
        ENNReal.ofReal (C_gns d 2) *
          eLpNorm (fun x => ‖hwvSob.weakGrad x‖) 2 μ < ⊤ := by
      rw [ENNReal.mul_lt_top_iff]
      exact Or.inl ⟨by simp, hgrad_lt_top⟩
    exact lt_of_le_of_lt hSob hrhs_lt_top
  have hSob_real :
      MeasureTheory.lpNorm v q μ ≤
        (C_gns d 2) * MeasureTheory.lpNorm (fun x => ‖hwvSob.weakGrad x‖) 2 μ := by
    have hgrad_ne_top :
        eLpNorm (fun x => ‖hwvSob.weakGrad x‖) 2 μ ≠ ⊤ :=
      by simpa using hwvSob.weakGrad_norm_memLp.eLpNorm_ne_top
    have hrhs_ne_top :
        ENNReal.ofReal (C_gns d 2) *
            eLpNorm (fun x => ‖hwvSob.weakGrad x‖) 2 μ ≠ ⊤ :=
      ENNReal.mul_ne_top (by simp) hgrad_ne_top
    have hSob_toReal :
        (eLpNorm v q μ).toReal ≤
          (ENNReal.ofReal (C_gns d 2) *
            eLpNorm (fun x => ‖hwvSob.weakGrad x‖) 2 μ).toReal :=
      (ENNReal.toReal_le_toReal hv_memLp_q.eLpNorm_ne_top hrhs_ne_top).2 hSob
    have hC_toReal : (ENNReal.ofReal (C_gns d 2)).toReal = C_gns d 2 :=
      ENNReal.toReal_ofReal (C_gns_nonneg d 2)
    have hSob_toReal' := hSob_toReal
    rw [MeasureTheory.toReal_eLpNorm hv_memLp_q.aestronglyMeasurable,
      ENNReal.toReal_mul, MeasureTheory.toReal_eLpNorm
      hwvSob.weakGrad_norm_memLp.aestronglyMeasurable, hC_toReal] at hSob_toReal'
    exact hSob_toReal'
  have hgrad_lp :
      MeasureTheory.lpNorm (fun x => ‖hwvSob.weakGrad x‖) 2 μ =
        (∫ x in Ω, ‖hwvSob.weakGrad x‖ ^ 2 ∂volume) ^ (1 / (2 : ℝ)) := by
    simpa [μ, Ω] using
      (MeasureTheory.lpNorm_eq_integral_norm_rpow_toReal
        (μ := μ) (f := fun x => ‖hwvSob.weakGrad x‖)
        two_ne_zero ENNReal.ofNat_ne_top
        hwvSob.weakGrad_norm_memLp.aestronglyMeasurable)
  have hgrad_bound :
      MeasureTheory.lpNorm (fun x => ‖hwvSob.weakGrad x‖) 2 μ ≤
        (2 * Cη ^ 2 * (A.1.Λ * (p / (p - 1)) ^ 2 + 1) *
          ∫ x in Ω, |max (u x) 0| ^ p ∂volume) ^ (1 / (2 : ℝ)) := by
    rw [hgrad_lp, hgrad_sq_eq]
    have henergy_real :
        ∫ x in Ω, ‖hwvReal.weakGrad x‖ ^ 2 ∂volume ≤
          2 * Cη ^ 2 * (A.1.Λ * (p / (p - 1)) ^ 2 + 1) *
            ∫ x in Ω, |max (u x) 0| ^ p ∂volume := by
      simpa [Ω, hwvReal] using henergy
    have henergy_rhs_nonneg :
        0 ≤
          2 * Cη ^ 2 * (A.1.Λ * (p / (p - 1)) ^ 2 + 1) *
            ∫ x in Ω, |max (u x) 0| ^ p ∂volume := by
      have hInt_nonneg :
          0 ≤ ∫ x in Ω, |max (u x) 0| ^ p ∂volume := by
        refine integral_nonneg ?_
        intro x
        positivity
      have hcoeff_nonneg : 0 ≤ A.1.Λ * (p / (p - 1)) ^ 2 + 1 := by
        have hratio_sq_nonneg : 0 ≤ (p / (p - 1)) ^ 2 := by positivity
        have hmul_nonneg : 0 ≤ A.1.Λ * (p / (p - 1)) ^ 2 := by
          exact mul_nonneg A.1.Λ_nonneg hratio_sq_nonneg
        linarith
      positivity
    exact Real.rpow_le_rpow
      (by
        refine integral_nonneg ?_
        intro x
        positivity)
      henergy_real (by
        show 0 ≤ (1 / (2 : ℝ))
        norm_num)
  have hq_ne_zero : q ≠ 0 := by
    dsimp [q]
    rw [ENNReal.ofReal_eq_zero]
    linarith
  have hq_toReal : q.toReal = qexp := by
    simp [q, qexp, ENNReal.toReal_ofReal hq_pos.le]
  have hv_lp_eq :
      MeasureTheory.lpNorm v q μ =
        (∫ x in Ω, |v x| ^ qexp ∂volume) ^ (1 / qexp) := by
    rw [MeasureTheory.lpNorm_eq_integral_norm_rpow_toReal
      (μ := μ) (f := v)]
    · simp [μ, Ω, q, qexp, hq_toReal, Real.norm_eq_abs]
    · exact hq_ne_zero
    · exact ENNReal.ofReal_ne_top
    · exact hv_memLp_q.aestronglyMeasurable
  have hv_int :
      IntegrableOn (fun x => |v x| ^ qexp) Ω volume := by
    simpa [μ, Ω, q, qexp, hq_toReal, Real.norm_eq_abs] using
      hv_memLp_q.integrable_norm_rpow hq_ne_zero ENNReal.ofReal_ne_top
  have hpow_eq_on :
      ∀ x ∈ Metric.ball (0 : E) r,
        |v x| ^ qexp = |max (u x) 0| ^ (moserChi d * p) := by
    intro x hx
    have hηx : η x = 1 := hη_eq_one x hx
    have hqexp_eq : qexp = 2 * moserChi d := by
      dsimp [qexp, moserChi]
      field_simp
    have hvx : v x = |max (u x) 0| ^ (p / 2) := by
      simp [v, moserPowerCutoff, hηx]
    calc
      |v x| ^ qexp = |(|max (u x) 0| ^ (p / 2))| ^ qexp := by
        rw [hvx]
      _ = (|max (u x) 0| ^ (p / 2)) ^ qexp := by
        rw [abs_of_nonneg (Real.rpow_nonneg (abs_nonneg _) _)]
      _ = |max (u x) 0| ^ (moserChi d * p) := by
        rw [hqexp_eq, ← Real.rpow_mul]
        · ring_nf
        · positivity
  have hleft_int :
      IntegrableOn (fun x => |max (u x) 0| ^ (moserChi d * p))
        (Metric.ball (0 : E) r) volume := by
    refine (hv_int.mono_set (Metric.ball_subset_ball (le_of_lt hrs))).congr_fun ?_ measurableSet_ball
    intro x hx
    exact hpow_eq_on x hx
  have hleft_integral_le :
      ∫ x in Metric.ball (0 : E) r, |max (u x) 0| ^ (moserChi d * p) ∂volume ≤
        ∫ x in Ω, |v x| ^ qexp ∂volume := by
    have hmono :
        ∫ x in Metric.ball (0 : E) r, |v x| ^ qexp ∂volume ≤
          ∫ x in Ω, |v x| ^ qexp ∂volume := by
      exact setIntegral_mono_set hv_int
        (ae_of_all _ (by intro x; positivity))
        (ae_of_all _ (Metric.ball_subset_ball (le_of_lt hrs)))
    calc
      ∫ x in Metric.ball (0 : E) r, |max (u x) 0| ^ (moserChi d * p) ∂volume
          = ∫ x in Metric.ball (0 : E) r, |v x| ^ qexp ∂volume := by
              apply setIntegral_congr_fun measurableSet_ball
              intro x hx
              symm
              exact hpow_eq_on x hx
      _ ≤ ∫ x in Ω, |v x| ^ qexp ∂volume := hmono
  have hCη_sq :
      Cη ^ 2 = 4 * (Mst : ℝ) ^ 2 / (s - r) ^ 2 := by
    dsimp [Cη]
    field_simp [hsr_pos.ne']
    ring
  have hconst_bound :
      (C_gns d 2) ^ 2 * (2 * Cη ^ 2) ≤ C_MoserAnchor d / (s - r) ^ 2 := by
    rw [hCη_sq]
    have hanchor :=
      cutoff_sobolev_anchor_le_C_MoserAnchor (d := d)
    have hsr_sq_pos : 0 < (s - r) ^ 2 := by positivity
    have hrewrite :
        (C_gns d 2) ^ 2 * (2 * (4 * (Mst : ℝ) ^ 2 / (s - r) ^ 2)) =
          ((C_gns d 2) ^ 2 * (8 * (Mst : ℝ) ^ 2)) / (s - r) ^ 2 := by
      field_simp [hsr_pos.ne']
      ring
    rw [hrewrite]
    have hnum :
        (C_gns d 2) ^ 2 * (8 * (Mst : ℝ) ^ 2) ≤ C_MoserAnchor d := by
      simpa [mul_assoc, mul_left_comm, mul_comm] using hanchor
    exact div_le_div_of_nonneg_right hnum (by positivity)
  have hcoeff_nonneg : 0 ≤ A.1.Λ * (p / (p - 1)) ^ 2 + 1 := by
    have hratio_sq_nonneg : 0 ≤ (p / (p - 1)) ^ 2 := by positivity
    have hmul_nonneg : 0 ≤ A.1.Λ * (p / (p - 1)) ^ 2 := by
      exact mul_nonneg A.1.Λ_nonneg hratio_sq_nonneg
    linarith
  have hIp_nonneg :
      0 ≤ ∫ x in Ω, |max (u x) 0| ^ p ∂volume := by
    refine integral_nonneg ?_
    intro x
    positivity
  have hcoeffInt_nonneg :
      0 ≤
        (A.1.Λ * (p / (p - 1)) ^ 2 + 1) *
          ∫ x in Ω, |max (u x) 0| ^ p ∂volume := by
    exact mul_nonneg hcoeff_nonneg hIp_nonneg
  have hbase_nonneg :
      0 ≤
        (C_gns d 2) ^ 2 * (2 * Cη ^ 2) *
          (A.1.Λ * (p / (p - 1)) ^ 2 + 1) *
            ∫ x in Ω, |max (u x) 0| ^ p ∂volume := by
    have hleft_nonneg : 0 ≤ (C_gns d 2) ^ 2 * (2 * Cη ^ 2) := by
      exact mul_nonneg (sq_nonneg (C_gns d 2)) (mul_nonneg (by norm_num) (sq_nonneg Cη))
    have hmid_nonneg :
        0 ≤ (C_gns d 2) ^ 2 * (2 * Cη ^ 2) * (A.1.Λ * (p / (p - 1)) ^ 2 + 1) := by
      exact mul_nonneg hleft_nonneg hcoeff_nonneg
    exact mul_nonneg hmid_nonneg hIp_nonneg
  have hbase_le :
      (C_gns d 2) ^ 2 * (2 * Cη ^ 2) *
          (A.1.Λ * (p / (p - 1)) ^ 2 + 1) *
            ∫ x in Ω, |max (u x) 0| ^ p ∂volume ≤
        (C_MoserAnchor d / (s - r) ^ 2) *
          (A.1.Λ * (p / (p - 1)) ^ 2 + 1) *
            ∫ x in Ω, |max (u x) 0| ^ p ∂volume := by
    have hmul :
        ((C_gns d 2) ^ 2 * (2 * Cη ^ 2)) *
            ((A.1.Λ * (p / (p - 1)) ^ 2 + 1) *
              ∫ x in Ω, |max (u x) 0| ^ p ∂volume) ≤
          (C_MoserAnchor d / (s - r) ^ 2) *
            ((A.1.Λ * (p / (p - 1)) ^ 2 + 1) *
              ∫ x in Ω, |max (u x) 0| ^ p ∂volume) := by
      exact mul_le_mul_of_nonneg_right hconst_bound hcoeffInt_nonneg
    simpa [mul_assoc] using hmul
  have hmain_lp :
      (∫ x in Metric.ball (0 : E) r,
          |max (u x) 0| ^ (moserChi d * p) ∂volume) ^
          (1 / (moserChi d * p)) ≤
        ((C_MoserAnchor d / (s - r) ^ 2) *
            (A.1.Λ * (p / (p - 1)) ^ 2 + 1) *
            ∫ x in Ω, |max (u x) 0| ^ p ∂volume) ^ (1 / p) := by
    have hleft_nonneg :
        0 ≤ ∫ x in Metric.ball (0 : E) r,
          |max (u x) 0| ^ (moserChi d * p) ∂volume := by
      refine integral_nonneg ?_
      intro x
      positivity
    have hright_nonneg :
        0 ≤ ∫ x in Ω, |v x| ^ qexp ∂volume := by
      refine integral_nonneg ?_
      intro x
      positivity
    have hgrad_rhs_nonneg :
        0 ≤
          (2 * Cη ^ 2 * (A.1.Λ * (p / (p - 1)) ^ 2 + 1) *
            ∫ x in Ω, |max (u x) 0| ^ p ∂volume) ^ (1 / (2 : ℝ)) := by
      exact Real.rpow_nonneg
        (mul_nonneg (mul_nonneg (mul_nonneg (by norm_num) (sq_nonneg Cη)) hcoeff_nonneg) hIp_nonneg) _
    have hqexp_mul :
        (1 / qexp : ℝ) * (2 / p) = 1 / (moserChi d * p) := by
      dsimp [qexp, moserChi]
      field_simp
    have hhalf_mul :
        (1 / (2 : ℝ)) * (2 / p) = 1 / p := by
      field_simp [hp.ne']
    have hclose_exp_nonneg : 0 ≤ 1 / (moserChi d * p) := by
      have hchi_p_pos : 0 < moserChi d * p := by
        exact mul_pos (moserChi_pos (d := d) hd) (lt_trans zero_lt_one hp)
      exact one_div_nonneg.mpr hchi_p_pos.le
    calc
      (∫ x in Metric.ball (0 : E) r,
          |max (u x) 0| ^ (moserChi d * p) ∂volume) ^
          (1 / (moserChi d * p))
          ≤ (∫ x in Ω, |v x| ^ qexp ∂volume) ^ (1 / (moserChi d * p)) := by
              exact Real.rpow_le_rpow hleft_nonneg hleft_integral_le hclose_exp_nonneg
      _ = ((∫ x in Ω, |v x| ^ qexp ∂volume) ^ (1 / qexp)) ^ (2 / p) := by
            rw [← Real.rpow_mul hright_nonneg, hqexp_mul]
      _ = (MeasureTheory.lpNorm v q μ) ^ (2 / p) := by
            rw [hv_lp_eq]
      _ ≤ ((C_gns d 2) * MeasureTheory.lpNorm (fun x => ‖hwvSob.weakGrad x‖) 2 μ) ^ (2 / p) := by
            exact Real.rpow_le_rpow MeasureTheory.lpNorm_nonneg hSob_real (by positivity)
      _ ≤ ((C_gns d 2) *
            (2 * Cη ^ 2 * (A.1.Λ * (p / (p - 1)) ^ 2 + 1) *
              ∫ x in Ω, |max (u x) 0| ^ p ∂volume) ^ (1 / (2 : ℝ))) ^ (2 / p) := by
            exact Real.rpow_le_rpow
              (mul_nonneg (C_gns_nonneg d 2) (MeasureTheory.lpNorm_nonneg))
              (mul_le_mul_of_nonneg_left hgrad_bound (C_gns_nonneg d 2))
              (by positivity)
      _ = ((C_gns d 2) ^ 2 * (2 * Cη ^ 2) *
            (A.1.Λ * (p / (p - 1)) ^ 2 + 1) *
            ∫ x in Ω, |max (u x) 0| ^ p ∂volume) ^ (1 / p) := by
            have hinner_nonneg :
                0 ≤
                  2 * Cη ^ 2 * (A.1.Λ * (p / (p - 1)) ^ 2 + 1) *
                    ∫ x in Ω, |max (u x) 0| ^ p ∂volume := by
              exact mul_nonneg (mul_nonneg (mul_nonneg (by norm_num) (sq_nonneg Cη)) hcoeff_nonneg) hIp_nonneg
            calc
              ((C_gns d 2) *
                  (2 * Cη ^ 2 * (A.1.Λ * (p / (p - 1)) ^ 2 + 1) *
                    ∫ x in Ω, |max (u x) 0| ^ p ∂volume) ^ (1 / (2 : ℝ))) ^ (2 / p)
                  = (C_gns d 2) ^ (2 / p) *
                      ((2 * Cη ^ 2 * (A.1.Λ * (p / (p - 1)) ^ 2 + 1) *
                        ∫ x in Ω, |max (u x) 0| ^ p ∂volume) ^ (1 / (2 : ℝ))) ^ (2 / p) := by
                          rw [Real.mul_rpow (C_gns_nonneg d 2) hgrad_rhs_nonneg]
              _ = ((C_gns d 2) ^ 2) ^ (1 / p) *
                    (2 * Cη ^ 2 * (A.1.Λ * (p / (p - 1)) ^ 2 + 1) *
                      ∫ x in Ω, |max (u x) 0| ^ p ∂volume) ^ (1 / p) := by
                        congr 1
                        · have htwo_div :
                              (2 / p : ℝ) = 2 * (1 / p) := by
                            field_simp [hp.ne']
                          rw [htwo_div]
                          simpa using (Real.rpow_mul (C_gns_nonneg d 2) 2 (1 / p))
                        · rw [← Real.rpow_mul hinner_nonneg, hhalf_mul]
              _ = (((C_gns d 2) ^ 2) *
                    (2 * Cη ^ 2 * (A.1.Λ * (p / (p - 1)) ^ 2 + 1) *
                      ∫ x in Ω, |max (u x) 0| ^ p ∂volume)) ^ (1 / p) := by
                        symm
                        rw [Real.mul_rpow (sq_nonneg (C_gns d 2)) hinner_nonneg]
              _ = ((C_gns d 2) ^ 2 * (2 * Cη ^ 2) *
                    (A.1.Λ * (p / (p - 1)) ^ 2 + 1) *
                    ∫ x in Ω, |max (u x) 0| ^ p ∂volume) ^ (1 / p) := by
                        congr 1
                        ring
      _ ≤ ((C_MoserAnchor d / (s - r) ^ 2) *
            (A.1.Λ * (p / (p - 1)) ^ 2 + 1) *
            ∫ x in Ω, |max (u x) 0| ^ p ∂volume) ^ (1 / p) := by
            exact Real.rpow_le_rpow hbase_nonneg hbase_le (by positivity)
  refine ⟨hleft_int, ?_⟩
  have hsplit_nonneg :
      0 ≤
        ((C_MoserAnchor d / (s - r) ^ 2) *
            (A.1.Λ * (p / (p - 1)) ^ 2 + 1)) := by
    have hfactor_nonneg : 0 ≤ C_MoserAnchor d / (s - r) ^ 2 := by
      have hanchor_nonneg : 0 ≤ C_MoserAnchor d := by
        exact le_trans (by norm_num : (0 : ℝ) ≤ 1) (one_le_C_MoserAnchor (d := d))
      exact div_nonneg hanchor_nonneg (by positivity)
    exact mul_nonneg hfactor_nonneg hcoeff_nonneg
  calc
    (∫ x in Metric.ball (0 : E) r,
        |max (u x) 0| ^ (moserChi d * p) ∂volume) ^
        (1 / (moserChi d * p))
        ≤ ((C_MoserAnchor d / (s - r) ^ 2) *
            (A.1.Λ * (p / (p - 1)) ^ 2 + 1) *
            ∫ x in Ω, |max (u x) 0| ^ p ∂volume) ^ (1 / p) := hmain_lp
    _ = ((C_MoserAnchor d / (s - r) ^ 2) *
            (A.1.Λ * (p / (p - 1)) ^ 2 + 1)) ^ (1 / p) *
          (∫ x in Ω, |max (u x) 0| ^ p ∂volume) ^ (1 / p) := by
            rw [Real.mul_rpow hsplit_nonneg hIp_nonneg]
    _ = ((C_MoserAnchor d / (s - r) ^ 2) *
            (A.1.Λ * (p / (p - 1)) ^ 2 + 1)) ^ (1 / p) *
          (∫ x in Metric.ball (0 : E) s, |max (u x) 0| ^ p ∂volume) ^ (1 / p) := by
            rfl

\end{lstlisting}
\begin{proof}
Choose the midpoint radius $R := (r+s)/2 \in (r,s)$ and a smooth radial cutoff
$\eta$ with $\eta = 1$ on $\Bz{r}$, $\supp \eta \subseteq \overline{\Bz{R}}
\subseteq \Bz{s}$, $0 \leq \eta \leq 1$, and (by
\leanname{Cutoff.exists}) $\|\nabla \eta\|_{L^\infty} \leq C_\eta := 2\Mst/(s-r)$.

Set $v := \mathrm{PC}_d(\eta,u,p) = \eta\,(u_+)^{p/2}$. By
Theorem~\ref{thm:moser-cutoff-energy} we get $v \in W^{1,2}_0(\Bz{s})$ with
\begin{equation*}
  \int_{\Bz{s}} \|\nabla v\|^{2}\,dx
  \;\leq\;
  2\,C_\eta^{2}\,\Big(\Lambda\Big(\frac{p}{p-1}\Big)^{2}+1\Big)
  \int_{\Bz{s}} (u_+)^{p}\,dx.
\end{equation*}
The Sobolev inequality \leanname{sobolev\_prepare\_on\_ball} on $\Bz{s}$ in
exponent $2 \to 2^{*} = 2d/(d-2)$ gives, with $C_{\mathrm{gns}} = C_{\mathrm{gns}}(d,2)$,
\begin{equation*}
  \Big(\int_{\Bz{s}} v^{2^{*}}\,dx\Big)^{1/2^{*}}
  \;\leq\;
  C_{\mathrm{gns}}\,
  \Big(\int_{\Bz{s}} \|\nabla v\|^{2}\,dx\Big)^{1/2}.
\end{equation*}
On $\Bz{r}$ we have $\eta = 1$, so $v^{2^{*}} = (u_+)^{p\,d/(d-2)} = (u_+)^{\chi p}$.
Combining with the energy bound and using $C_\eta^2 = 4\Mst^2/(s-r)^2$ together
with $8\Mst^{2}\,C_{\mathrm{gns}}^{2} \leq C_{\mathrm{MA}}(d)$
(\leanname{cutoff\_sobolev\_anchor\_le\_C\_MoserAnchor}) yields
\begin{equation*}
  \Big(\int_{\Bz{r}} (u_+)^{\chi p}\,dx\Big)^{2/2^{*}}
  \;\leq\;
  \frac{C_{\mathrm{MA}}(d)}{(s-r)^{2}}
  \Big(\Lambda\Big(\frac{p}{p-1}\Big)^{2}+1\Big)
  \int_{\Bz{s}} (u_+)^{p}\,dx.
\end{equation*}
Raising to the power $1/p$ and using $2/(2^{*}\,p)= 1/(\chi p)$ gives the
claim.
\end{proof}

\subsubsection{Iteration}

We now iterate the pre-estimate along $(r_n,p_n)$.

\begin{definition}[Iteration quantities]
Set
\begin{align*}
  I_n(u) &:= \int_{\Bz{r_n}} (u_+)^{p_n}\,dx
    && (\text{\leanname{moserIterIntegral}}), \\
  N_n(u) &:= I_n(u)^{1/p_n}
    && (\text{\leanname{moserIterNorm}}), \\
  \mathrm{S}(A,p_0) &:= 32\,C_{\mathrm{MA}}(d)\,\Lambda\,\Big(\frac{p_0}{p_0-1}\Big)^{2}
    && (\text{\leanname{moserStepConst}}).
\end{align*}
The product majorant produced by iteration is, for $n \geq 0$,
\begin{equation*}
  M_n(A,u,p_0) :=
  \Big(
    \mathrm{S}(A,p_0)^{\sum_{i=0}^{n-1} q^{i}}\,
    4^{\sum_{i=0}^{n-1} i\,q^{i}}\,
    \int_{\Bz{1}} (u_+)^{p_0}\,dx
  \Big)^{1/p_0}.
\end{equation*}
We further define the target majorant
\begin{align*}
  B_\infty(A,u,p_0) &:= C_{\mathrm{Moser}}(d)\,\Lambda^{d/2}\,
    \Big(\frac{p_0}{p_0-1}\Big)^{d}\,\int_{\Bz{1}} (u_+)^{p_0}\,dx
    && (\text{\leanname{moserLinftyBoundPow}}), \\
  M_\infty(A,u,p_0) &:= B_\infty(A,u,p_0)^{1/p_0}
    && (\text{\leanname{moserLinftyMajorant}}).
\end{align*}
\end{definition}

\begin{lstlisting}[style=Matlab-editor,basicstyle=\mlttfamily,escapechar=",mlshowsectionrules=true,mathescape=true,morecomment={[s]{\%\{}{\%\}}},language=lean,numbers=none]
/-- Auxiliary powered integral along the standard Moser radius/exponent sequence. -/
noncomputable def moserIterIntegral
    {u : E → ℝ} (p₀ : ℝ) (n : ℕ) : ℝ :=
  ∫ x in Metric.ball (0 : E) (moserRadius n), |max (u x) 0| ^ moserExponentSeq d p₀ n ∂volume

/-- Auxiliary `L^{p_n}` quantity in the Moser iteration. -/
noncomputable def moserIterNorm
    {u : E → ℝ} (p₀ : ℝ) (n : ℕ) : ℝ :=
  (moserIterIntegral (d := d) (u := u) p₀ n) ^ (1 / moserExponentSeq d p₀ n)

/-- Auxiliary step constant in the normalized-coefficient pre-Moser estimate. -/
noncomputable def moserStepConst
    (A : NormalizedEllipticCoeff d (Metric.ball (0 : E) 1)) (p₀ : ℝ) : ℝ :=
  (32 : ℝ) * C_MoserAnchor d * A.1.Λ * (p₀ / (p₀ - 1)) ^ 2

/-- Auxiliary target constant for the Moser local-max theorem. -/
noncomputable def moserLinftyBoundPow
    (A : NormalizedEllipticCoeff d (Metric.ball (0 : E) 1))
    {u : E → ℝ} (p₀ : ℝ) : ℝ :=
  C_Moser d * A.1.Λ ^ ((d : ℝ) / 2) *
    (p₀ / (p₀ - 1)) ^ (d : ℝ) *
    ∫ x in Metric.ball (0 : E) 1, |max (u x) 0| ^ p₀ ∂volume

/-- Auxiliary root majorant corresponding to `moserLinftyBoundPow`. -/
noncomputable def moserLinftyMajorant
    (A : NormalizedEllipticCoeff d (Metric.ball (0 : E) 1))
    {u : E → ℝ} (p₀ : ℝ) : ℝ :=
  (moserLinftyBoundPow (d := d) A (u := u) p₀) ^ (1 / p₀)

\end{lstlisting}
\begin{theorem}[Step--majorant inequality, \leanname{moser\_step\_majorant\_le}]
\label{thm:moser-step-maj}
For $d > 2$, $1 < p_0$ and any $n \geq 0$,
\begin{multline*}
  \Big(\frac{C_{\mathrm{MA}}(d)}{(r_n-r_{n+1})^{2}}\,
    \Big(\Lambda\Big(\frac{p_n}{p_n-1}\Big)^{2}+1\Big)\Big)^{1/p_n}
  M_n(A,u,p_0) \\
  \leq\;
  M_{n+1}(A,u,p_0).
\end{multline*}
\end{theorem}

\begin{lstlisting}[style=Matlab-editor,basicstyle=\mlttfamily,escapechar=",mlshowsectionrules=true,mathescape=true,morecomment={[s]{\%\{}{\%\}}},language=lean,numbers=none]
theorem moser_step_majorant_le
    (hd : 2 < (d : ℝ))
    (A : NormalizedEllipticCoeff d (Metric.ball (0 : E) 1))
    {u : E → ℝ} {p₀ : ℝ} (hp₀ : 1 < p₀) (n : ℕ) :
    (((C_MoserAnchor d / (moserRadius n - moserRadius (n + 1)) ^ 2) *
        (A.1.Λ *
            (moserExponentSeq d p₀ n / (moserExponentSeq d p₀ n - 1)) ^ 2 + 1)) ^
        (1 / moserExponentSeq d p₀ n)) *
      moserIterMajorant (d := d) A (u := u) p₀ n ≤
      moserIterMajorant (d := d) A (u := u) p₀ (n + 1) := by
  let p_n : ℝ := moserExponentSeq d p₀ n
  have hp₀_pos : 0 < p₀ := by linarith
  have hp_n_gt_one : 1 < p_n := by
    simpa [p_n] using one_lt_moserExponentSeq (d := d) hd hp₀ n
  have hp_n_pos : 0 < p_n := by linarith
  have hp₀_nonneg : 0 ≤ p₀ := by linarith
  have hp_n_ge : p₀ ≤ p_n := by
    simpa [p_n] using moserExponentSeq_ge_initial (d := d) hd hp₀_nonneg n
  have hp₀m1_pos : 0 < p₀ - 1 := by linarith
  have hp_nm1_pos : 0 < p_n - 1 := by linarith
  have hratio_le : p_n / (p_n - 1) ≤ p₀ / (p₀ - 1) := by
    field_simp [hp_nm1_pos.ne', hp₀m1_pos.ne']
    nlinarith
  have hratio_nonneg : 0 ≤ p_n / (p_n - 1) := by positivity
  have hratio₀_nonneg : 0 ≤ p₀ / (p₀ - 1) := by positivity
  have hratio_sq_le :
      (p_n / (p_n - 1)) ^ 2 ≤ (p₀ / (p₀ - 1)) ^ 2 := by
    nlinarith
  have hΛ_ge_one : 1 ≤ A.1.Λ := by
    simpa [A.2] using A.1.hΛ
  have hratio₀_gt_one : 1 < p₀ / (p₀ - 1) := by
    exact (one_lt_div hp₀m1_pos).2 (by nlinarith)
  have hratio₀_sq_ge_one : 1 ≤ (p₀ / (p₀ - 1)) ^ 2 := by
    nlinarith
  have hΛratio_ge_one : 1 ≤ A.1.Λ * (p₀ / (p₀ - 1)) ^ 2 := by
    nlinarith
  have hbracket_le :
      A.1.Λ * (p_n / (p_n - 1)) ^ 2 + 1 ≤
        2 * A.1.Λ * (p₀ / (p₀ - 1)) ^ 2 := by
    nlinarith [mul_le_mul_of_nonneg_left hratio_sq_le A.1.Λ_nonneg]
  have hgap_eq :
      C_MoserAnchor d / (moserRadius n - moserRadius (n + 1)) ^ 2 =
        C_MoserAnchor d * (4 : ℝ) ^ (n + 2) := by
    rw [moserRadius_gap]
    have hsq : (((1 / 2 : ℝ) ^ (n + 2)) ^ 2) = (1 / 4 : ℝ) ^ (n + 2) := by
      rw [← pow_mul]
      rw [show (n + 2) * 2 = 2 * (n + 2) by ring]
      rw [pow_mul]
      norm_num
    rw [hsq, div_eq_mul_inv]
    rw [show (((1 / 4 : ℝ) ^ (n + 2))⁻¹) = (4 : ℝ) ^ (n + 2) by
      rw [show (1 / 4 : ℝ) = (4 : ℝ)⁻¹ by norm_num, inv_pow]
      norm_num]
  have hpow4 : (4 : ℝ) ^ (n + 2) = (4 : ℝ) ^ n * 16 := by
    rw [pow_add]
    norm_num
  have harg_nonneg :
      0 ≤
        (C_MoserAnchor d / (moserRadius n - moserRadius (n + 1)) ^ 2) *
          (A.1.Λ * (p_n / (p_n - 1)) ^ 2 + 1) := by
    have hanchor_nonneg : 0 ≤ C_MoserAnchor d := by
      exact le_trans (by norm_num : (0 : ℝ) ≤ 1) (one_le_C_MoserAnchor (d := d))
    refine mul_nonneg ?_ ?_
    · exact div_nonneg hanchor_nonneg (sq_nonneg _)
    · nlinarith [A.1.Λ_nonneg, sq_nonneg (p_n / (p_n - 1))]
  have hgap_nonneg : 0 ≤ C_MoserAnchor d * (4 : ℝ) ^ (n + 2) := by
    have hanchor_nonneg : 0 ≤ C_MoserAnchor d := by
      exact le_trans (by norm_num : (0 : ℝ) ≤ 1) (one_le_C_MoserAnchor (d := d))
    exact mul_nonneg hanchor_nonneg (by positivity)
  have harg_le :
      (C_MoserAnchor d / (moserRadius n - moserRadius (n + 1)) ^ 2) *
          (A.1.Λ * (p_n / (p_n - 1)) ^ 2 + 1) ≤
        moserStepConst (d := d) A p₀ * (4 : ℝ) ^ n := by
    rw [hgap_eq]
    calc
      C_MoserAnchor d * (4 : ℝ) ^ (n + 2) *
          (A.1.Λ * (p_n / (p_n - 1)) ^ 2 + 1)
          ≤ C_MoserAnchor d * (4 : ℝ) ^ (n + 2) *
              (2 * A.1.Λ * (p₀ / (p₀ - 1)) ^ 2) := by
              gcongr
      _ = moserStepConst (d := d) A p₀ * (4 : ℝ) ^ n := by
            rw [hpow4]
            dsimp [moserStepConst]
            ring
  have hprefactor_le :
      ((C_MoserAnchor d / (moserRadius n - moserRadius (n + 1)) ^ 2) *
          (A.1.Λ * (p_n / (p_n - 1)) ^ 2 + 1)) ^ (1 / p_n) ≤
        (moserStepConst (d := d) A p₀ * (4 : ℝ) ^ n) ^ (1 / p_n) := by
    refine Real.rpow_le_rpow harg_nonneg harg_le ?_
    · positivity
  have hprefactor_nonneg :
      0 ≤
        ((C_MoserAnchor d / (moserRadius n - moserRadius (n + 1)) ^ 2) *
          (A.1.Λ * (p_n / (p_n - 1)) ^ 2 + 1)) ^ (1 / p_n) := by
    exact Real.rpow_nonneg harg_nonneg _
  have hstep_pos : 0 < moserStepConst (d := d) A p₀ := by
    have hanchor_pos : 0 < C_MoserAnchor d := by
      exact lt_of_lt_of_le zero_lt_one (one_le_C_MoserAnchor (d := d))
    have hratio₀_pos : 0 < p₀ / (p₀ - 1) := by positivity
    dsimp [moserStepConst]
    positivity
  have hstep_nonneg : 0 ≤ moserStepConst (d := d) A p₀ := hstep_pos.le
  have hq_nonneg : 0 ≤ moserDecayRatio d ^ n := by
    exact pow_nonneg (moserDecayRatio_nonneg (d := d) hd) n
  let S : ℝ := Finset.sum (Finset.range n) (fun i => moserDecayRatio d ^ i)
  let T : ℝ := Finset.sum (Finset.range n) (fun i => (i : ℝ) * moserDecayRatio d ^ i)
  let I : ℝ := (∫ x in Metric.ball (0 : E) 1, |max (u x) 0| ^ p₀ ∂volume)
  have hinner_nonneg :
      0 ≤
        moserStepConst (d := d) A p₀ ^ S * 4 ^ T * I := by
    refine mul_nonneg ?_ ?_
    · refine mul_nonneg ?_ ?_
      · exact Real.rpow_nonneg hstep_nonneg _
      · positivity
    · dsimp [I]
      refine integral_nonneg ?_
      intro x
      positivity
  have hstep_eq :
      ((moserStepConst (d := d) A p₀ * (4 : ℝ) ^ n) ^ (1 / p_n)) *
          moserIterMajorant (d := d) A (u := u) p₀ n =
        moserIterMajorant (d := d) A (u := u) p₀ (n + 1) := by
    rw [moserIterMajorant, moserIterMajorant, Finset.sum_range_succ, Finset.sum_range_succ]
    have hinv : 1 / p_n = moserDecayRatio d ^ n / p₀ := by
      simpa [p_n] using inv_moserExponentSeq (d := d) hd hp₀_pos n
    have hdiv_eq : 1 / p_n = moserDecayRatio d ^ n * (1 / p₀) := by
      rw [hinv]
      ring
    rw [hdiv_eq]
    have hstep4_nonneg : 0 ≤ moserStepConst (d := d) A p₀ * (4 : ℝ) ^ n := by
      exact mul_nonneg hstep_nonneg (by positivity)
    rw [Real.rpow_mul hstep4_nonneg]
    rw [Real.mul_rpow hstep_nonneg (by positivity : 0 ≤ (4 : ℝ) ^ n)]
    have hfour_pow :
        ((4 : ℝ) ^ n) ^ (moserDecayRatio d ^ n) =
          (4 : ℝ) ^ ((n : ℝ) * moserDecayRatio d ^ n) := by
      rw [show ((4 : ℝ) ^ n) = (4 : ℝ) ^ (n : ℝ) by rw [Real.rpow_natCast]]
      rw [← Real.rpow_mul (by positivity : (0 : ℝ) ≤ 4)]
    rw [hfour_pow]
    have hstepFactor_nonneg :
        0 ≤
          moserStepConst (d := d) A p₀ ^ (moserDecayRatio d ^ n) *
            4 ^ ((n : ℝ) * moserDecayRatio d ^ n) := by
      exact mul_nonneg (Real.rpow_nonneg hstep_nonneg _) (by positivity)
    rw [← Real.mul_rpow hstepFactor_nonneg hinner_nonneg]
    calc
      ((moserStepConst (d := d) A p₀ ^ (moserDecayRatio d ^ n) *
          4 ^ ((n : ℝ) * moserDecayRatio d ^ n)) *
        (moserStepConst (d := d) A p₀ ^ S * 4 ^ T * I)) ^ (1 / p₀)
          =
        (moserStepConst (d := d) A p₀ ^
            (moserDecayRatio d ^ n + S) *
          4 ^ (((n : ℝ) * moserDecayRatio d ^ n) + T) * I) ^ (1 / p₀) := by
            congr 1
            calc
              (moserStepConst (d := d) A p₀ ^ (moserDecayRatio d ^ n) *
                    4 ^ ((n : ℝ) * moserDecayRatio d ^ n)) *
                  (moserStepConst (d := d) A p₀ ^ S * 4 ^ T * I)
                  =
                (moserStepConst (d := d) A p₀ ^ (moserDecayRatio d ^ n) *
                    moserStepConst (d := d) A p₀ ^ S) *
                  ((4 ^ ((n : ℝ) * moserDecayRatio d ^ n) *
                      4 ^ T) * I) := by
                      ring
              _ =
                (moserStepConst (d := d) A p₀ ^ (moserDecayRatio d ^ n + S)) *
                  ((4 ^ (((n : ℝ) * moserDecayRatio d ^ n) + T)) * I) := by
                      rw [← Real.rpow_add hstep_pos, ← Real.rpow_add (by positivity : 0 < (4 : ℝ))]
              _ =
                moserStepConst (d := d) A p₀ ^ (moserDecayRatio d ^ n + S) *
                  4 ^ (((n : ℝ) * moserDecayRatio d ^ n) + T) * I := by
                    ring
      _ =
        ((moserStepConst (d := d) A p₀ ^
            (Finset.sum (Finset.range n) (fun i => moserDecayRatio d ^ i) +
              moserDecayRatio d ^ n) *
          4 ^ (Finset.sum (Finset.range n)
              (fun i => (i : ℝ) * moserDecayRatio d ^ i) +
                (n : ℝ) * moserDecayRatio d ^ n) * I) ^ (1 / p₀)) := by
            simp [S, T, add_comm]
  have hmajorant_nonneg :
      0 ≤ moserIterMajorant (d := d) A (u := u) p₀ n := by
    simpa [moserIterMajorant, S, T, I] using
      (Real.rpow_nonneg hinner_nonneg (1 / p₀))
  simpa [p_n] using
    (calc
    (((C_MoserAnchor d / (moserRadius n - moserRadius (n + 1)) ^ 2) *
        (A.1.Λ * (p_n / (p_n - 1)) ^ 2 + 1)) ^ (1 / p_n)) *
        moserIterMajorant (d := d) A (u := u) p₀ n
        ≤ ((moserStepConst (d := d) A p₀ * (4 : ℝ) ^ n) ^ (1 / p_n)) *
            moserIterMajorant (d := d) A (u := u) p₀ n := by
            exact mul_le_mul_of_nonneg_right hprefactor_le hmajorant_nonneg
    _ = moserIterMajorant (d := d) A (u := u) p₀ (n + 1) := hstep_eq
    )

\end{lstlisting}
\begin{proof}
First, $r_n - r_{n+1} = 2^{-(n+2)}$ so
$C_{\mathrm{MA}}/(r_n-r_{n+1})^2 = C_{\mathrm{MA}}\,4^{n+2} = 16\,C_{\mathrm{MA}}\,4^{n}$.
Second, $p_n/(p_n-1) \leq p_0/(p_0-1)$ since $t \mapsto t/(t-1)$ is decreasing
on $(1,\infty)$, hence
\[
  \Lambda\Big(\frac{p_n}{p_n-1}\Big)^{2}+1
  \;\leq\; 2\,\Lambda\Big(\frac{p_0}{p_0-1}\Big)^{2}.
\]
Combining gives
\begin{equation*}
  \frac{C_{\mathrm{MA}}}{(r_n-r_{n+1})^{2}}\,
    \Big(\Lambda\Big(\frac{p_n}{p_n-1}\Big)^{2}+1\Big)
  \;\leq\;
  \mathrm{S}(A,p_0)\,4^{n}.
\end{equation*}
Raising to $1/p_n = q^n/p_0$ and multiplying by $M_n$ produces the inner factor
$\mathrm{S}^{q^n/p_0} 4^{n q^n/p_0}$, which precisely promotes the partial sums
in the exponents from $\sum_{i<n}$ to $\sum_{i\leq n}$. The resulting expression
equals $M_{n+1}(A,u,p_0)$ by the definition of $M_{n+1}$.
\end{proof}

\begin{theorem}[Iteration bound, \leanname{moser\_iteration\_bound}]
\label{thm:moser-iter-bound}
For $d > 2$, $1 < p_0$, $u$ a subsolution of $A$, and
$\int_{\Bz{1}} |u_+|^{p_0} < \infty$, for every $n \in \N$ we have
$|u_+|^{p_n}$ integrable on $\Bz{r_n}$ and
\[
  N_n(u) \;\leq\; M_n(A,u,p_0).
\]
\end{theorem}

\begin{lstlisting}[style=Matlab-editor,basicstyle=\mlttfamily,escapechar=",mlshowsectionrules=true,mathescape=true,morecomment={[s]{\%\{}{\%\}}},language=lean,numbers=none]
theorem moser_iteration_bound
    (hd : 2 < (d : ℝ))
    (A : NormalizedEllipticCoeff d (Metric.ball (0 : E) 1))
    {u : E → ℝ} {p₀ : ℝ} (hp₀ : 1 < p₀)
    (hsub : IsSubsolution A.1 u)
    (hposInt :
      IntegrableOn (fun x => |max (u x) 0| ^ p₀)
        (Metric.ball (0 : E) 1) volume) :
    ∀ n,
      IntegrableOn (fun x => |max (u x) 0| ^ moserExponentSeq d p₀ n)
        (Metric.ball (0 : E) (moserRadius n)) volume ∧
      moserIterNorm (d := d) (u := u) p₀ n ≤
        moserIterMajorant (d := d) A (u := u) p₀ n := by
  intro n
  induction n with
  | zero =>
      constructor
      · simpa [moserExponentSeq_zero, moserRadius_zero] using hposInt
      · simp [moserIterNorm, moserIterIntegral, moserIterMajorant, moserExponentSeq_zero,
          moserRadius_zero]
  | succ n ihn =>
      rcases ihn with ⟨hInt_n, hbound_n⟩
      let p_n : ℝ := moserExponentSeq d p₀ n
      have hp_n : 1 < p_n := by
        simpa [p_n] using one_lt_moserExponentSeq (d := d) hd hp₀ n
      have hpre :=
        moser_preMoser (d := d) hd A (u := u) (p := p_n)
          (r := moserRadius (n + 1)) (s := moserRadius n)
          hp_n (moserRadius_pos (n + 1)) (moserRadius_succ_lt n)
          (moserRadius_le_one n) hsub (by simpa [p_n] using hInt_n)
      rcases hpre with ⟨hInt_succ, hbound_succ⟩
      refine ⟨?_, ?_⟩
      · simpa [p_n, moserExponentSeq_succ, moserIterIntegral, mul_comm, mul_left_comm,
          mul_assoc] using hInt_succ
      · have hfactor_nonneg :
            0 ≤
              ((C_MoserAnchor d / (moserRadius n - moserRadius (n + 1)) ^ 2) *
                  (A.1.Λ * (p_n / (p_n - 1)) ^ 2 + 1)) ^ (1 / p_n) := by
            have harg_nonneg :
                0 ≤
                  (C_MoserAnchor d / (moserRadius n - moserRadius (n + 1)) ^ 2) *
                    (A.1.Λ * (p_n / (p_n - 1)) ^ 2 + 1) := by
              have hanchor_nonneg : 0 ≤ C_MoserAnchor d := by
                exact le_trans (by norm_num : (0 : ℝ) ≤ 1) (one_le_C_MoserAnchor (d := d))
              refine mul_nonneg ?_ ?_
              · exact div_nonneg hanchor_nonneg (sq_nonneg _)
              · nlinarith [A.1.Λ_nonneg, sq_nonneg (p_n / (p_n - 1))]
            exact Real.rpow_nonneg harg_nonneg _
        calc
          moserIterNorm (d := d) (u := u) p₀ (n + 1)
              ≤ ((C_MoserAnchor d / (moserRadius n - moserRadius (n + 1)) ^ 2) *
                    (A.1.Λ * (p_n / (p_n - 1)) ^ 2 + 1)) ^ (1 / p_n) *
                  moserIterNorm (d := d) (u := u) p₀ n := by
                    simpa [moserIterNorm, moserIterIntegral, p_n, moserExponentSeq_succ,
                      mul_comm, mul_left_comm, mul_assoc] using hbound_succ
          _ ≤ ((C_MoserAnchor d / (moserRadius n - moserRadius (n + 1)) ^ 2) *
                    (A.1.Λ * (p_n / (p_n - 1)) ^ 2 + 1)) ^ (1 / p_n) *
                  moserIterMajorant (d := d) A (u := u) p₀ n := by
                    exact mul_le_mul_of_nonneg_left hbound_n hfactor_nonneg
          _ ≤ moserIterMajorant (d := d) A (u := u) p₀ (n + 1) := by
                    simpa [p_n] using
                      moser_step_majorant_le (d := d) hd A (u := u) hp₀ n

\end{lstlisting}
\begin{proof}
By induction on $n$. The case $n=0$ uses $r_0=1$, $p_0 = p_0$, and
$M_0(A,u,p_0)= (\int_{\Bz{1}} (u_+)^{p_0})^{1/p_0}$ by definition. For the
inductive step, the pre-estimate Theorem~\ref{thm:moser-preMoser} applied with
$(p,r,s) = (p_n,r_{n+1},r_n)$ gives integrability of $|u_+|^{\chi p_n}$ on
$\Bz{r_{n+1}}$ and
\begin{multline*}
  N_{n+1}(u)
  \;\leq\;
  \Big(\frac{C_{\mathrm{MA}}}{(r_n-r_{n+1})^{2}}
    \Big(\Lambda\Big(\frac{p_n}{p_n-1}\Big)^{2}+1\Big)\Big)^{1/p_n}
  N_n(u).
\end{multline*}
The inductive hypothesis $N_n(u) \leq M_n(A,u,p_0)$ and
Theorem~\ref{thm:moser-step-maj} then yield $N_{n+1}(u) \leq M_{n+1}(A,u,p_0)$.
\end{proof}

\begin{lemma}[Geometric majorant absorption, \leanname{moser\_geometric\_majorant}]
\label{lem:moser-geom-maj}
For $d > 2$, $1 < p_0$, and all $n$,
\[
  M_n(A,u,p_0) \;\leq\; M_\infty(A,u,p_0).
\]
\end{lemma}

\begin{lstlisting}[style=Matlab-editor,basicstyle=\mlttfamily,escapechar=",mlshowsectionrules=true,mathescape=true,morecomment={[s]{\%\{}{\%\}}},language=lean,numbers=none]
theorem moser_geometric_majorant
    (hd : 2 < (d : ℝ))
    (A : NormalizedEllipticCoeff d (Metric.ball (0 : E) 1))
    {u : E → ℝ} {p₀ : ℝ} (hp₀ : 1 < p₀) :
    ∀ n,
      moserIterMajorant (d := d) A (u := u) p₀ n ≤
        moserLinftyMajorant (d := d) A (u := u) p₀ := by
  intro n
  let q : ℝ := moserDecayRatio d
  let S : ℝ := Finset.sum (Finset.range n) (fun i => q ^ i)
  let T : ℝ := Finset.sum (Finset.range n) (fun i => (i : ℝ) * q ^ i)
  let I : ℝ := ∫ x in Metric.ball (0 : E) 1, |max (u x) 0| ^ p₀ ∂volume
  let B : ℝ := (32 : ℝ) * C_MoserAnchor d
  let r : ℝ := p₀ / (p₀ - 1)
  have hp₀_pos : 0 < p₀ := by
    linarith
  have hp₀m1_pos : 0 < p₀ - 1 := by
    linarith
  have hq_nonneg : 0 ≤ q := by
    simpa [q] using moserDecayRatio_nonneg (d := d) hd
  have hq_lt_one : q < 1 := by
    simpa [q] using moserDecayRatio_lt_one (d := d) hd
  have hS_le : S ≤ (d : ℝ) / 2 := by
    simpa [q, S] using moserDecayRatio_sum_le_half_dim (d := d) hd n
  have hT_le : T ≤ ∑' i : ℕ, (i : ℝ) * q ^ i := by
    simpa [q, T] using moserDecayRatio_weighted_sum_le_tsum (d := d) hd n
  have hS_nonneg : 0 ≤ S := by
    dsimp [S]
    refine Finset.sum_nonneg ?_
    intro i hi
    exact pow_nonneg hq_nonneg i
  have hT_nonneg : 0 ≤ T := by
    dsimp [T]
    refine Finset.sum_nonneg ?_
    intro i hi
    positivity
  have h2S_le : 2 * S ≤ (d : ℝ) := by
    linarith
  have hanchor_nonneg : 0 ≤ C_MoserAnchor d := by
    exact le_trans (by norm_num : (0 : ℝ) ≤ 1) (one_le_C_MoserAnchor (d := d))
  have hB_ge_one : 1 ≤ B := by
    dsimp [B]
    nlinarith [one_le_C_MoserAnchor (d := d)]
  have hB_nonneg : 0 ≤ B := by
    linarith
  have hΛ_ge_one : 1 ≤ A.1.Λ := by
    simpa [A.2] using A.1.hΛ
  have hr_gt_one : 1 < r := by
    dsimp [r]
    exact (one_lt_div hp₀m1_pos).2 (by linarith)
  have hr_ge_one : 1 ≤ r := hr_gt_one.le
  have hr_nonneg : 0 ≤ r := by
    linarith
  have hstep_nonneg : 0 ≤ moserStepConst (d := d) A p₀ := by
    dsimp [moserStepConst, r]
    refine mul_nonneg ?_ (by positivity)
    refine mul_nonneg ?_ A.1.Λ_nonneg
    exact mul_nonneg (by norm_num : (0 : ℝ) ≤ 32) hanchor_nonneg
  have hI_nonneg : 0 ≤ I := by
    dsimp [I]
    refine integral_nonneg ?_
    intro x
    positivity
  have hB_pow_le : B ^ S ≤ B ^ ((d : ℝ) / 2) := by
    exact Real.rpow_le_rpow_of_exponent_le hB_ge_one hS_le
  have hΛ_pow_le : A.1.Λ ^ S ≤ A.1.Λ ^ ((d : ℝ) / 2) := by
    exact Real.rpow_le_rpow_of_exponent_le hΛ_ge_one hS_le
  have hr_sq_pow_le : (r ^ (2 : ℝ)) ^ S ≤ r ^ (d : ℝ) := by
    have hpow_le : r ^ (2 * S) ≤ r ^ (d : ℝ) := by
      exact Real.rpow_le_rpow_of_exponent_le hr_ge_one h2S_le
    simpa [Real.rpow_mul hr_nonneg] using hpow_le
  have hstep_eq_base :
      moserStepConst (d := d) A p₀ = ((B * A.1.Λ) * r ^ (2 : ℝ)) := by
    calc
      moserStepConst (d := d) A p₀
          = (32 : ℝ) * C_MoserAnchor d * A.1.Λ * (p₀ / (p₀ - 1)) ^ 2 := by
              rfl
      _ = ((B * A.1.Λ) * r ^ (2 : ℕ)) := by
            dsimp [B, r]
      _ = ((B * A.1.Λ) * r ^ (2 : ℝ)) := by
            exact congrArg (fun t : ℝ => (B * A.1.Λ) * t) (Real.rpow_natCast r 2).symm
  have hstep_eq :
      moserStepConst (d := d) A p₀ ^ S =
        B ^ S * A.1.Λ ^ S * (r ^ (2 : ℝ)) ^ S := by
    have hBmul_nonneg : 0 ≤ B * A.1.Λ := by
      exact mul_nonneg hB_nonneg A.1.Λ_nonneg
    calc
      moserStepConst (d := d) A p₀ ^ S = ((B * A.1.Λ) * r ^ (2 : ℝ)) ^ S := by
        rw [hstep_eq_base]
      _ = (B * A.1.Λ) ^ S * (r ^ (2 : ℝ)) ^ S := by
        rw [Real.mul_rpow hBmul_nonneg (by positivity : 0 ≤ r ^ (2 : ℝ))]
      _ = B ^ S * A.1.Λ ^ S * (r ^ (2 : ℝ)) ^ S := by
        rw [Real.mul_rpow hB_nonneg A.1.Λ_nonneg]
  have hstep_le :
      moserStepConst (d := d) A p₀ ^ S ≤
        B ^ ((d : ℝ) / 2) * A.1.Λ ^ ((d : ℝ) / 2) * r ^ (d : ℝ) := by
    rw [hstep_eq]
    have hBA_le :
        B ^ S * A.1.Λ ^ S ≤ B ^ ((d : ℝ) / 2) * A.1.Λ ^ ((d : ℝ) / 2) := by
      exact mul_le_mul hB_pow_le hΛ_pow_le
        (Real.rpow_nonneg A.1.Λ_nonneg _)
        (Real.rpow_nonneg hB_nonneg _)
    have hstep_le' :
        (B ^ S * A.1.Λ ^ S) * (r ^ (2 : ℝ)) ^ S ≤
          (B ^ ((d : ℝ) / 2) * A.1.Λ ^ ((d : ℝ) / 2)) * r ^ (d : ℝ) := by
      exact mul_le_mul hBA_le hr_sq_pow_le
        (Real.rpow_nonneg (by positivity : (0 : ℝ) ≤ r ^ (2 : ℝ)) _)
        (mul_nonneg (Real.rpow_nonneg hB_nonneg _) (Real.rpow_nonneg A.1.Λ_nonneg _))
    simpa [mul_assoc] using hstep_le'
  have h4_le :
      (4 : ℝ) ^ T ≤ 4 ^ (∑' i : ℕ, (i : ℝ) * q ^ i) := by
    exact Real.rpow_le_rpow_of_exponent_le (by norm_num : (1 : ℝ) ≤ 4) hT_le
  have hgeom_le :
      moserStepConst (d := d) A p₀ ^ S * 4 ^ T ≤
        (B ^ ((d : ℝ) / 2) * 4 ^ (∑' i : ℕ, (i : ℝ) * q ^ i)) *
          A.1.Λ ^ ((d : ℝ) / 2) * r ^ (d : ℝ) := by
    have hmul_le :
        (moserStepConst (d := d) A p₀ ^ S) * (4 : ℝ) ^ T ≤
          (B ^ ((d : ℝ) / 2) * A.1.Λ ^ ((d : ℝ) / 2) * r ^ (d : ℝ)) *
            4 ^ (∑' i : ℕ, (i : ℝ) * q ^ i) := by
      exact mul_le_mul hstep_le h4_le
        (Real.rpow_nonneg (by norm_num : (0 : ℝ) ≤ 4) _)
        (mul_nonneg
          (mul_nonneg (Real.rpow_nonneg hB_nonneg _) (Real.rpow_nonneg A.1.Λ_nonneg _))
          (Real.rpow_nonneg hr_nonneg _))
    simpa [mul_assoc, mul_left_comm, mul_comm] using hmul_le
  have hC_le :
      (B ^ ((d : ℝ) / 2) * 4 ^ (∑' i : ℕ, (i : ℝ) * q ^ i)) *
          A.1.Λ ^ ((d : ℝ) / 2) * r ^ (d : ℝ) ≤
        C_Moser d * A.1.Λ ^ ((d : ℝ) / 2) * r ^ (d : ℝ) := by
    have hC_base :
        B ^ ((d : ℝ) / 2) * 4 ^ (∑' i : ℕ, (i : ℝ) * q ^ i) ≤ C_Moser d := by
      dsimp [B, q]
      rw [C_Moser, dif_pos hd]
      exact
        le_max_right (C_MoserAnchor d)
          (((32 : ℝ) * C_MoserAnchor d) ^ ((d : ℝ) / 2) *
            4 ^ (∑' n : ℕ, (n : ℝ) * moserDecayRatio d ^ n))
    have hfactor_nonneg : 0 ≤ A.1.Λ ^ ((d : ℝ) / 2) * r ^ (d : ℝ) := by
      exact mul_nonneg (Real.rpow_nonneg A.1.Λ_nonneg _) (Real.rpow_nonneg hr_nonneg _)
    simpa [mul_assoc, mul_left_comm, mul_comm] using
      (mul_le_mul_of_nonneg_right hC_base hfactor_nonneg)
  have hinside_le :
      moserStepConst (d := d) A p₀ ^ S * 4 ^ T * I ≤
        C_Moser d * A.1.Λ ^ ((d : ℝ) / 2) * r ^ (d : ℝ) * I := by
    exact mul_le_mul_of_nonneg_right (hgeom_le.trans hC_le) hI_nonneg
  have hinside_nonneg :
      0 ≤ moserStepConst (d := d) A p₀ ^ S * 4 ^ T * I := by
    exact mul_nonneg
      (mul_nonneg
        (Real.rpow_nonneg hstep_nonneg _)
        (Real.rpow_nonneg (by norm_num : (0 : ℝ) ≤ 4) _))
      hI_nonneg
  have hroot_le :
      (moserStepConst (d := d) A p₀ ^ S * 4 ^ T * I) ^ (1 / p₀) ≤
        (C_Moser d * A.1.Λ ^ ((d : ℝ) / 2) * r ^ (d : ℝ) * I) ^ (1 / p₀) := by
    refine Real.rpow_le_rpow hinside_nonneg hinside_le ?_
    positivity
  simpa [moserIterMajorant, moserLinftyMajorant, moserLinftyBoundPow, S, T, I, r] using hroot_le

\end{lstlisting}
\begin{proof}
Set $S := \sum_{i=0}^{n-1} q^{i}$, $T := \sum_{i=0}^{n-1} i q^{i}$,
$B := 32\,C_{\mathrm{MA}}(d)$, $r := p_0/(p_0-1)$. Then:
$S \leq d/2$ \leanname{moserDecayRatio\_sum\_le\_half\_dim}, since
$\sum_{i\geq 0} q^{i} = 1/(1-q) = d/2$; also $T \leq \sum_{i\geq 0} i\,q^{i}$
\leanname{moserDecayRatio\_weighted\_sum\_le\_tsum} by summability.
Now $\mathrm{S}(A,p_0) = B \Lambda r^{2}$, so
\[
  \mathrm{S}(A,p_0)^{S}
  = B^{S} \Lambda^{S} (r^{2})^{S}.
\]
Using $1 \leq B$, $1 \leq \Lambda$, and $1 \leq r$, monotonicity of $t \mapsto
t^{e}$ in the exponent gives $B^{S} \leq B^{d/2}$, $\Lambda^{S} \leq
\Lambda^{d/2}$, and $(r^{2})^{S} \leq r^{d}$ (because $2S \leq d$). Also $4^{T}
\leq 4^{\sum_{i\geq 0} i q^{i}}$. Multiplying these inequalities and using the
explicit definition of $C_{\mathrm{Moser}}(d)$ in the case $d>2$:
\[
  B^{d/2}\,4^{\sum_{i\geq 0} i q^{i}} \;\leq\; C_{\mathrm{Moser}}(d),
\]
we obtain
\[
  \mathrm{S}(A,p_0)^{S}\,4^{T}\,I_0(u)
  \;\leq\;
  C_{\mathrm{Moser}}(d)\,\Lambda^{d/2}\,r^{d}\,I_0(u)
  = B_\infty(A,u,p_0).
\]
Raising to the power $1/p_0$ yields $M_n \leq M_\infty$.
\end{proof}

\subsubsection{Closeout to $L^\infty$}

\begin{lemma}[Closeout from iterated bounds, \leanname{moser\_ae\_closeout}]
\label{lem:moser-closeout}
Fix $d > 2$, $1 < p_0$, and assume that for every $n \in \N$,
$|u_+|^{p_n}$ is integrable on $\Bz{r_n}$ and $N_n(u) \leq M_\infty(A,u,p_0)$.
Then for almost every $x \in \Bz{1/2}$,
\[
  |u_+(x)|^{p_0} \;\leq\; B_\infty(A,u,p_0).
\]
\end{lemma}

\begin{lstlisting}[style=Matlab-editor,basicstyle=\mlttfamily,escapechar=",mlshowsectionrules=true,mathescape=true,morecomment={[s]{\%\{}{\%\}}},language=lean,numbers=none]
private theorem moser_ae_closeout
    (hd : 2 < (d : ℝ))
    (A : NormalizedEllipticCoeff d (Metric.ball (0 : E) 1))
    {u : E → ℝ} {p₀ : ℝ} (hp₀ : 1 < p₀)
    (hiter :
      ∀ n,
        IntegrableOn (fun x => |max (u x) 0| ^ moserExponentSeq d p₀ n)
          (Metric.ball (0 : E) (moserRadius n)) volume ∧
        moserIterNorm (d := d) (u := u) p₀ n ≤
          moserLinftyMajorant (d := d) A (u := u) p₀) :
    ∀ᵐ x ∂(volume.restrict (Metric.ball (0 : E) (1 / 2 : ℝ))),
      |max (u x) 0| ^ p₀ ≤ moserLinftyBoundPow (d := d) A (u := u) p₀ := by
  let Bhalf : Set E := Metric.ball (0 : E) (1 / 2 : ℝ)
  let μ : Measure E := volume.restrict Bhalf
  let f : E → ℝ := fun x => |max (u x) 0|
  let g : E → ℝ := fun x => f x ^ (p₀ / 2)
  let K : ℝ := moserLinftyBoundPow (d := d) A (u := u) p₀ ^ (1 / 2 : ℝ)
  haveI : IsFiniteMeasure μ := by
    dsimp [μ, Bhalf]
    rw [isFiniteMeasure_restrict]
    exact measure_ne_top_of_subset Metric.ball_subset_closedBall
      (isCompact_closedBall (0 : E) (1 / 2 : ℝ)).measure_lt_top.ne
  have hK_nonneg : 0 ≤ K := by
    dsimp [K]
    exact Real.rpow_nonneg (moserLinftyBoundPow_nonneg (d := d) A hp₀) _
  have hbound_all :
      ∀ m : ℕ, ∀ᵐ x ∂μ, g x < K + 1 / (m + 1 : ℝ) := by
    intro m
    have hzero :
        μ.real {x | K + 1 / (m + 1 : ℝ) ≤ g x} = 0 := by
      refine moser_closeout_superlevel_null (d := d) hd A (u := u) hp₀ hiter ?_
      dsimp [K]
      have hfrac_pos : 0 < 1 / (m + 1 : ℝ) := by positivity
      linarith
    rw [ae_iff]
    simpa [μ, Bhalf, g, K, not_lt] using
      (measureReal_eq_zero_iff
        (μ := μ) (s := {x | K + 1 / (m + 1 : ℝ) ≤ g x})
        (measure_ne_top μ {x | K + 1 / (m + 1 : ℝ) ≤ g x})).1 hzero
  have hbound_ae :
      ∀ᵐ x ∂μ, ∀ m : ℕ, g x < K + 1 / (m + 1 : ℝ) := by
    exact ae_all_iff.2 hbound_all
  filter_upwards [hbound_ae] with x hx
  have hfx_nonneg : 0 ≤ f x := by
    dsimp [f]
    exact abs_nonneg (max (u x) 0)
  have hgx_nonneg : 0 ≤ g x := by
    dsimp [g]
    exact Real.rpow_nonneg hfx_nonneg _
  have hgx_le : g x ≤ K := by
    by_contra hgx_gt
    have hgap_pos : 0 < g x - K := by linarith
    obtain ⟨m, hm⟩ := exists_nat_one_div_lt hgap_pos
    have hlt : K + 1 / (m + 1 : ℝ) < g x := by linarith
    linarith [hx m]
  have hgpow :
      g x ^ (2 : ℝ) ≤ K ^ (2 : ℝ) := by
    exact Real.rpow_le_rpow hgx_nonneg hgx_le (by norm_num)
  have hg_sq_eq : g x ^ (2 : ℝ) = f x ^ p₀ := by
    calc
      g x ^ (2 : ℝ) = (f x ^ (p₀ / 2)) ^ (2 : ℝ) := by rfl
      _ = f x ^ ((p₀ / 2) * 2) := by rw [← Real.rpow_mul hfx_nonneg]
      _ = f x ^ p₀ := by
            congr 2
            ring
  have hK_sq_eq :
      K ^ (2 : ℝ) = moserLinftyBoundPow (d := d) A (u := u) p₀ := by
    calc
      K ^ (2 : ℝ)
          = moserLinftyBoundPow (d := d) A (u := u) p₀ ^ ((1 / 2 : ℝ) * 2) := by
              dsimp [K]
              rw [← Real.rpow_mul (moserLinftyBoundPow_nonneg (d := d) A hp₀)]
      _ = moserLinftyBoundPow (d := d) A (u := u) p₀ := by
            rw [show ((1 / 2 : ℝ) * 2) = (1 : ℝ) by ring, Real.rpow_one]
  calc
    |max (u x) 0| ^ p₀ = f x ^ p₀ := by rfl
    _ = g x ^ (2 : ℝ) := hg_sq_eq.symm
    _ ≤ K ^ (2 : ℝ) := hgpow
    _ = moserLinftyBoundPow (d := d) A (u := u) p₀ := hK_sq_eq

\end{lstlisting}
\begin{proof}
Set $\mu := \mathrm{vol}\,|\,\Bz{1/2}$, $f(x) := |u_+(x)|$, $g(x) := f(x)^{p_0/2}$,
and $K := B_\infty(A,u,p_0)^{1/2}$. Note $K \geq 0$.
For each $n$, write $q_n := 2 \chi^n$. We claim that for every $n$
and every $c > K$,
\[
  \mu_{\mathrm{r}}\{x : c \leq g(x)\} \;\leq\; (K/c)^{q_n}.
\]
Indeed, on $\{c \leq g\}$ we have $c^{q_n} \leq g^{q_n}$, and Markov's
inequality gives
\[
  c^{q_n}\,\mu_{\mathrm{r}}\{c \leq g\}
  \;\leq\;
  \int g^{q_n}\,d\mu
  = \int_{\Bz{1/2}} f^{p_n}\,dx
  \;\leq\;
  \int_{\Bz{r_n}} f^{p_n}\,dx = I_n(u),
\]
where we used $g^{q_n} = f^{p_0 q_n / 2} = f^{p_n}$ and
$\Bz{1/2} \subseteq \Bz{r_n}$. By hypothesis $I_n(u)^{1/p_n} = N_n(u) \leq
M_\infty$, so raising to $p_n$ and using $1/q_n = (1/p_n)(p_0/2)$ we get
$I_n(u) \leq M_\infty^{q_n p_0/2 \cdot 2/p_0} = M_\infty^{q_n p_0/p_0}$\,;
more cleanly, $I_n^{1/q_n} \leq M_\infty^{p_0/2} = K$. Hence
$\mu_{\mathrm{r}}\{c\leq g\} \leq (K/c)^{q_n}$.

As $\chi > 1$, $q_n \to \infty$, so $(K/c)^{q_n} \to 0$ for every $c > K$.
Therefore $\mu_{\mathrm{r}}\{c \leq g\} = 0$ for all $c > K$, and a countable
union over $c = K + 1/(m+1)$ shows $g \leq K$ a.e.\ on $\Bz{1/2}$. Squaring,
$f^{p_0} = g^{2} \leq K^{2} = B_\infty$.
\end{proof}

\begin{theorem}[Moser $L^p \to L^\infty$ for subsolutions, \leanname{linfty\_subsolution\_Moser}]
\label{thm:linfty-Moser}
Let $d > 2$, $A$ a normalized elliptic coefficient on $\Bz{1}$, $u$ a
subsolution, and $1 < p_0$ with $\int_{\Bz{1}} |u_+|^{p_0} < \infty$. Then for
almost every $x \in \Bz{1/2}$,
\[
  |u_+(x)|^{p_0}
  \;\leq\;
  C_{\mathrm{Moser}}(d)\,\Lambda^{d/2}\,
  \Big(\frac{p_0}{p_0-1}\Big)^{d}
  \int_{\Bz{1}} |u_+|^{p_0}\,dx.
\]
\end{theorem}

\begin{lstlisting}[style=Matlab-editor,basicstyle=\mlttfamily,escapechar=",mlshowsectionrules=true,mathescape=true,morecomment={[s]{\%\{}{\%\}}},language=lean,numbers=none]
/-- Moser `L^p → L∞` estimate for subsolutions on the unit ball, in the honest
essential/a.e. form available before the continuity upgrade.

This is the normalized-coefficient unit-ball form of the Moser local maximum
estimate. -/
theorem linfty_subsolution_Moser
    (hd : 2 < (d : ℝ))
    (A : NormalizedEllipticCoeff d (Metric.ball (0 : E) 1))
    {u : E → ℝ} {p₀ : ℝ} (hp₀ : 1 < p₀)
    (hsub : IsSubsolution A.1 u)
    (hposInt :
      IntegrableOn (fun x => |max (u x) 0| ^ p₀)
        (Metric.ball (0 : E) 1) volume) :
    ∀ᵐ x ∂(volume.restrict (Metric.ball (0 : E) (1 / 2 : ℝ))),
      |max (u x) 0| ^ p₀ ≤
      C_Moser d * A.1.Λ ^ ((d : ℝ) / 2) *
        (p₀ / (p₀ - 1)) ^ (d : ℝ) *
        ∫ x in Metric.ball (0 : E) 1, |max (u x) 0| ^ p₀ ∂volume := by
  have hiter_raw :=
    moser_iteration_bound (d := d) hd A (u := u) hp₀ hsub hposInt
  have hiter :
      ∀ n,
        IntegrableOn (fun x => |max (u x) 0| ^ moserExponentSeq d p₀ n)
          (Metric.ball (0 : E) (moserRadius n)) volume ∧
        moserIterNorm (d := d) (u := u) p₀ n ≤
          moserLinftyMajorant (d := d) A (u := u) p₀ := by
    intro n
    rcases hiter_raw n with ⟨hInt_n, hbound_n⟩
    exact ⟨hInt_n, hbound_n.trans (moser_geometric_majorant (d := d) hd A hp₀ n)⟩
  simpa [moserLinftyBoundPow] using
    moser_ae_closeout (d := d) hd A (u := u) hp₀ hiter

\end{lstlisting}
\begin{proof}
The headline Moser bound is the composition of three pieces, all assembled
along the discrete sequence of radii $r_n$ and exponents $p_n$.

\emph{Step 1: iterated step bound.} Theorem~\ref{thm:moser-iter-bound}
(\leanname{moser\_iteration\_bound}) is applied as a single induction on
$n$: at each step, the concentric-ball pre-estimate
Theorem~\ref{thm:moser-preMoser} on $(r_{n+1}, r_n)$ with exponent $p_n$
upgrades $L^{p_n}$ control of $u_+$ on $\Bz{r_n}$ to $L^{p_{n+1}} = L^{\chi p_n}$
control on $\Bz{r_{n+1}}$, with multiplicative constant the step constant
$\mathrm{S}_n$ depending on $\Lambda$, $C_{\mathrm{MA}}(d)$, the dimension,
and the gap $r_n - r_{n+1} = 2^{-(n+2)}$. Telescoping over $i < n$ produces
the bound
\[
  N_n(u)
   \le \Bigl(\prod_{i=0}^{n-1} \mathrm{S}_i\Bigr) \cdot N_0(u)
   = M_n(A, u, p_0).
\]
At every stage the integrability of $|u_+|^{p_n}$ on $\Bz{r_n}$ is propagated
forward by the same induction, providing the second half of the conclusion.

\emph{Step 2: geometric absorption of the product.}
Lemma~\ref{lem:moser-geom-maj} (\leanname{moser\_geometric\_majorant})
bounds the infinite product $\prod_{i \ge 0} \mathrm{S}_i$ by the closed-form
constant $M_\infty(A,u,p_0)$. The key two summation identities are
$\sum_{i \ge 0} q^{i} \le d/2$
(\leanname{moserDecayRatio\_sum\_le\_half\_dim}) and the weighted variant
$\sum_{i \ge 0} i\,q^{i}$
(\leanname{moserDecayRatio\_weighted\_sum\_le\_tsum}); these are precisely
what is needed to convert the product of $\mathrm{S}_i$, after taking
logarithms and using $1/p_i = q^{i}/p_0$, into a finite expression involving
$C_{\mathrm{Moser}}(d)$, $\Lambda^{d/2}$ and the explicit
$\bigl(p_0/(p_0-1)\bigr)^{d}$ factor visible in the conclusion. Composing with
Step 1 gives $N_n(u) \le M_\infty(A, u, p_0)$ for every $n$.

\emph{Step 3: a.e.\ closeout via Markov.}
Lemma~\ref{lem:moser-closeout} (\leanname{moser\_ae\_closeout}) converts the
uniform $L^{p_n}$ bound at every stage into an a.e.\ pointwise bound on the
intersection $\Bz{1/2} \subseteq \bigcap_n \Bz{r_n}$. The argument is a
Markov-type estimate: for every $c > K := M_\infty(A,u,p_0)^{1/2}$, the
super-level set $\{c \le f\}$ on $\Bz{1/2}$ has measure at most
$(K/c)^{q_n}$, and since $q_n \to \infty$ this measure is zero. Taking a
countable union over $c = K + 1/(m+1)$ shows $f \le K$ a.e.\ on $\Bz{1/2}$,
and squaring gives $|u_+|^{p_0} \le K^{2}$ a.e.\ on $\Bz{1/2}$, which is the
displayed Moser bound.
\end{proof}

\begin{remark}[$p_0 = 2$ anchor, \leanname{linfty\_subsolution\_Moser\_two}]
The case $p_0 = 2$ already follows from the normalized De Giorgi $L^2 \to
L^\infty$ estimate of Section 5: with $c := C_{\mathrm{DG}}^{\mathrm{sub}}(d)
\Lambda^{d/4}$, the De Giorgi bound gives
$\max(u(x),0) \leq c\,(\int_{\Bz{1}} |u_+|^{2}\,dx)^{1/2}$ a.e.\ on
$\Bz{1/2}$. Squaring and using
$C_{\mathrm{DG}}^{\mathrm{sub}}(d)^{2} \leq C_{\mathrm{Moser}}(d)$ together with
$2/(2-1) = 2$ recovers the $p_0=2$ case of Theorem~\ref{thm:linfty-Moser}.
\end{remark}

\begin{lstlisting}[style=Matlab-editor,basicstyle=\mlttfamily,escapechar=",mlshowsectionrules=true,mathescape=true,morecomment={[s]{\%\{}{\%\}}},language=lean,numbers=none]
/-- The `p = 2` anchor for Chapter 06 is already available from Chapter 05.

This is the exact normalized De Giorgi unit-ball estimate rewritten in the
Chapter 06 a.e. power-bound format. It is the base case that the later Moser
iteration should strictly improve from `p = 2` to arbitrary `p > 1`. -/
theorem linfty_subsolution_Moser_two
    (hd : 2 < (d : ℝ))
    (A : NormalizedEllipticCoeff d (Metric.ball (0 : E) 1))
    {u : E → ℝ}
    (hsub : IsSubsolution A.1 u)
    (hposInt :
      IntegrableOn (fun x => |max (u x) 0| ^ 2)
        (Metric.ball (0 : E) 1) volume) :
    ∀ᵐ x ∂(volume.restrict (Metric.ball (0 : E) (1 / 2 : ℝ))),
      |max (u x) 0| ^ 2 ≤
        C_Moser d * A.1.Λ ^ ((d : ℝ) / 2) *
          (2 / (2 - 1 : ℝ)) ^ (d : ℝ) *
          ∫ x in Metric.ball (0 : E) 1, |max (u x) 0| ^ 2 ∂volume := by
  let I₂ : ℝ := ∫ x in Metric.ball (0 : E) 1, |max (u x) 0| ^ 2 ∂volume
  let c : ℝ := C_DeGiorgi_subsolution_normalized d * A.1.Λ ^ ((d : ℝ) / 4)
  have hI₂_nonneg : 0 ≤ I₂ := by
    dsimp [I₂]
    refine integral_nonneg ?_
    intro x
    positivity
  have hc_nonneg : 0 ≤ c := by
    have hCDG_nonneg : 0 ≤ C_DeGiorgi_subsolution_normalized d := by
      dsimp [C_DeGiorgi_subsolution_normalized]
      exact
        (C_DeGiorgiSmallness_pos (d := d)
          (K_DeGiorgi_subsolution_normalized_pos (d := d))).le
    dsimp [c]
    exact mul_nonneg
      hCDG_nonneg
      (Real.rpow_nonneg A.1.Λ_nonneg _)
  filter_upwards
    [linfty_subsolution_DeGiorgi_normalized (d := d) hd A hsub hposInt]
    with x hx
  have hx_nonneg : 0 ≤ max (u x) 0 := by
    positivity
  have hsq :
      |max (u x) 0| ^ 2 ≤ c ^ 2 * I₂ := by
    have hx' : max (u x) 0 ≤ c * Real.sqrt I₂ := hx
    have hrhs_nonneg : 0 ≤ c * Real.sqrt I₂ := by
      exact mul_nonneg hc_nonneg (Real.sqrt_nonneg I₂)
    have hsq' : (max (u x) 0) ^ 2 ≤ (c * Real.sqrt I₂) ^ 2 := by
      nlinarith
    calc
      |max (u x) 0| ^ 2 = (max (u x) 0) ^ 2 := by
        rw [abs_of_nonneg hx_nonneg]
      _ ≤ (c * Real.sqrt I₂) ^ 2 := hsq'
      _ = c ^ 2 * I₂ := by
        rw [mul_pow, Real.sq_sqrt hI₂_nonneg]
  have hconst :
      c ^ 2 * I₂ ≤
        C_Moser d * A.1.Λ ^ ((d : ℝ) / 2) *
          (2 / (2 - 1 : ℝ)) ^ (d : ℝ) * I₂ := by
    have hbase :
        c ^ 2 ≤
          C_Moser d * A.1.Λ ^ ((d : ℝ) / 2) *
            (2 / (2 - 1 : ℝ)) ^ (d : ℝ) := by
      have hpow :
          (A.1.Λ ^ ((d : ℝ) / 4)) ^ 2 = A.1.Λ ^ ((d : ℝ) / 2) := by
        rw [pow_two, ← Real.rpow_add A.1.Λ_pos]
        ring_nf
      calc
        c ^ 2
            = (C_DeGiorgi_subsolution_normalized d) ^ 2 *
                A.1.Λ ^ ((d : ℝ) / 2) := by
                  dsimp [c]
                  rw [mul_pow, hpow]
        _ ≤ C_Moser d * A.1.Λ ^ ((d : ℝ) / 2) := by
            exact mul_le_mul_of_nonneg_right
              (C_DeGiorgi_subsolution_normalized_sq_le_C_Moser (d := d))
              (Real.rpow_nonneg A.1.Λ_nonneg _)
        _ ≤ C_Moser d * A.1.Λ ^ ((d : ℝ) / 2) *
              (2 / (2 - 1 : ℝ)) ^ (d : ℝ) := by
            have hleft_nonneg : 0 ≤ C_Moser d * A.1.Λ ^ ((d : ℝ) / 2) := by
              exact mul_nonneg
                (le_trans (by norm_num : (0 : ℝ) ≤ 1) (one_le_C_Moser (d := d)))
                (Real.rpow_nonneg A.1.Λ_nonneg _)
            have hone : (1 : ℝ) ≤ (2 / (2 - 1 : ℝ)) ^ (d : ℝ) := by
              have hbase : (1 : ℝ) ≤ 2 / (2 - 1 : ℝ) := by norm_num
              have hd_nonneg : 0 ≤ (d : ℝ) := by exact_mod_cast Nat.zero_le d
              exact Real.one_le_rpow hbase hd_nonneg
            simpa [mul_assoc] using
              (mul_le_mul_of_nonneg_left hone hleft_nonneg)
    exact mul_le_mul_of_nonneg_right hbase hI₂_nonneg
  exact le_trans hsq (by
    simpa [I₂] using hconst)

-- The supersolution iteration stages (weak_harnack_stage_one_inverse,
-- weak_harnack_stage_one_forward) have been split out to
-- Supersolutions.lean for independent development.
--
-- The weak Harnack inequality (weak_harnack) and the sharp Moser estimate
-- (linfty_subsolution_Moser_sharp) have been split out to

\end{lstlisting}
%% --- Section 8 ---

\subsection{Supersolution Estimates}

We now turn to positive supersolutions of the same operator. The main outputs
of this section are the two stages of the weak Harnack inequality:
\emph{stage one (inverse)} produces an a.e.\ lower bound on $u$ from inverse
$L^{p}$ control of $u$, and \emph{stage two (forward)} produces a low-power
$L^{p}$-on-$L^{p}$ Caccioppoli-type bound. The architecture mirrors Section 7,
with the role of $u_+$ replaced by $u^{-1}$ (resp.\ $u$ at low powers).

Throughout this section, $u : E \to \R$ is assumed to be strictly positive on
$\Bz{1}$ and a supersolution of $A$. We write $\mu_{1/2}$ for
$\mathrm{vol}\,|\,\Bz{1/2}$.

\subsubsection{The weighted Caccioppoli inequality}

The abstract energy tool driving every iteration in this section is the
following weighted Caccioppoli inequality.

\begin{theorem}[Weighted Caccioppoli, \leanname{weighted\_caccioppoli\_of\_supersolution}]
\label{thm:weighted-cacc}
Let $\Omega \subseteq E$ be open, $A$ a normalized elliptic coefficient on
$\Omega$, $u$ a supersolution with $u \in W^{1,2}(\Omega)$, and let
$\Psi : \R \to \R$ be smooth with $\Psi(0)=0$, $|\Psi'| \leq M$, $\Psi(u) \geq 0$
on $\Omega$, and $\Psi'(u) \leq 0$ on $\Omega$ (with $\Psi(u)=0$ wherever
$\Psi'(u)=0$). Let $\varphi \in C^{\infty}_c(\Omega)$ satisfy $|\varphi| \leq 1$
and $\|\nabla \varphi\|_{L^\infty} \leq C_\varphi$, with the integrability
\[
  \int_{\Omega} \|\nabla \varphi\|^{2}\,
  \frac{\Psi(u)^{2}}{-\Psi'(u)}\,dx < \infty
\]
(interpreted as $0$ where $\Psi'(u)=0$). Then
\begin{multline*}
  \int_{\Omega} \varphi^{2}\,(-\Psi'(u))\,\langle A\nabla u,\nabla u\rangle\,dx
  \\ \leq\;
  4\,\Lambda \int_{\Omega} \|\nabla \varphi\|^{2}\,
  \frac{\Psi(u)^{2}}{-\Psi'(u)}\,dx.
\end{multline*}
\end{theorem}

\begin{lstlisting}[style=Matlab-editor,basicstyle=\mlttfamily,escapechar=",mlshowsectionrules=true,mathescape=true,morecomment={[s]{\%\{}{\%\}}},language=lean,numbers=none]
/-- Weighted Caccioppoli inequality for supersolutions.

For `u > 0` a supersolution on `Ω`, `Ψ` smooth with `Ψ(0) = 0`, bounded derivative,
and `Ψ(u) ≥ 0`, and `φ` a smooth cutoff with compact support in `Ω`:

  `∫ φ² Ψ'(u) ⟨A∇u,∇u⟩ ≤ 4Λ ∫ |∇φ|² Ψ(u)²/Ψ'(u)`

where the ratio `Ψ²/Ψ'` is understood to be 0 where `Ψ' = 0`.

The proof uses the supersolution inequality `0 ≤ ∫⟨A∇u, ∇(φ²Ψ(u))⟩`,
expanded via `bilinFormIntegrand_mul_smooth_eq` into principal + cross terms,
then the cross term is absorbed via AM-GM. -/
theorem weighted_caccioppoli_of_supersolution
    {Ω : Set E} (hΩ : IsOpen Ω)
    {u φ : E → ℝ}
    (A : NormalizedEllipticCoeff d Ω)
    (hu : MemW1pWitness 2 u Ω)
    (hsuper : IsSupersolution A.1 u)
    (Ψ : ℝ → ℝ) (hΨ : ContDiff ℝ (⊤ : ℕ∞) Ψ) (hΨ_zero : Ψ 0 = 0)
    (hΨ_bdd : ∃ M, ∀ t, |deriv Ψ t| ≤ M)
    (hΨ_nonneg : ∀ x ∈ Ω, 0 ≤ Ψ (u x))
    (hΨ'_nonpos : ∀ x ∈ Ω, deriv Ψ (u x) ≤ 0)
    (hΨ_zero_of_deriv_zero : ∀ x ∈ Ω, deriv Ψ (u x) = 0 → Ψ (u x) = 0)
    (hφ : ContDiff ℝ (⊤ : ℕ∞) φ)
    {Cφ : ℝ} (hCφ : 0 ≤ Cφ)
    (hφ_bound : ∀ x, |φ x| ≤ 1)
    (hφ_grad_bound : ∀ x, ‖fderiv ℝ φ x‖ ≤ Cφ)
    (hφ_support : tsupport φ ⊆ Ω)
    (hratio_int :
      Integrable
        (fun x => ‖fderiv ℝ φ x‖ ^ 2 *
          (if deriv Ψ (u x) = 0 then 0
           else Ψ (u x) ^ 2 / (-deriv Ψ (u x))))
        (volume.restrict Ω))
    (hφ_memH01 : MemH01 (fun x => φ x ^ 2 * Ψ (u x)) Ω) :
    -- |Principal term| ≤ 4Λ · gradient term
    -- Here |Ψ'| = -Ψ' since Ψ' ≤ 0.
    ∫ x in Ω, φ x ^ 2 * (-deriv Ψ (u x)) *
        bilinFormIntegrandOfCoeff A.1 hu hu x ∂volume ≤
      4 * A.1.Λ * ∫ x in Ω, ‖fderiv ℝ φ x‖ ^ 2 *
        (if deriv Ψ (u x) = 0 then 0
         else Ψ (u x) ^ 2 / (-deriv Ψ (u x))) ∂volume := by
  let hwΨ : MemW1pWitness 2 (fun x => Ψ (u x)) Ω :=
    MemW1pWitness.comp_smooth_bounded (d := d) hΩ hu Ψ hΨ hΨ_zero hΨ_bdd
  let hwφΨ := hwΨ.mul_smooth_bounded_p (d := d)
    (by norm_num : (1 : ENNReal) ≤ 2) hΩ hφ (by norm_num : (0 : ℝ) ≤ 1) hCφ
    hφ_bound hφ_grad_bound
  let hwφ2Ψ := hwφΨ.mul_smooth_bounded_p (d := d)
    (by norm_num : (1 : ENNReal) ≤ 2) hΩ hφ (by norm_num : (0 : ℝ) ≤ 1) hCφ
    hφ_bound hφ_grad_bound
  -- φ(x) * (φ(x) * Ψ(u(x))) = φ(x)² * Ψ(u(x)) ≥ 0
  -- Need MemH01 for the product function. The function φ·(φ·Ψ(u)) = φ²·Ψ(u).
  have hφ2Ψ_eq : (fun x => φ x * (φ x * Ψ (u x))) = (fun x => φ x ^ 2 * Ψ (u x)) := by
    ext x; ring
  have hφ2Ψ_memH01 : MemH01 (fun x => φ x * (φ x * Ψ (u x))) Ω := by
    rwa [hφ2Ψ_eq]
  have hφ2Ψ_nonneg : ∀ x, 0 ≤ φ x * (φ x * Ψ (u x)) := by
    intro x
    rw [show φ x * (φ x * Ψ (u x)) = φ x ^ 2 * Ψ (u x) by ring]
    by_cases hx : x ∈ Ω
    · exact mul_nonneg (sq_nonneg _) (hΨ_nonneg x hx)
    · -- Outside Ω, φ has compact support in Ω, so φ(x) = 0
      have hφx : φ x = 0 := by
        by_contra h
        exact hx (hφ_support (subset_tsupport _ h))
      simp [hφx]
  have htest : 0 ≤ bilinFormOfCoeff A.1 hu hwφ2Ψ := by
    exact hsuper.2 hu _ hφ2Ψ_memH01 hwφ2Ψ hφ2Ψ_nonneg
  -- The inner expansion: ⟨A∇u, ∇(Ψ(u))⟩ = Ψ'(u)·⟨A∇u,∇u⟩
  -- (from comp_smooth_bounded: the gradient of Ψ(u) is Ψ'(u)·∇u)
  have hΨu_bilin : ∀ x, bilinFormIntegrandOfCoeff A.1 hu hwΨ x =
      deriv Ψ (u x) * bilinFormIntegrandOfCoeff A.1 hu hu x := by
    intro x
    -- `∇(Ψ ∘ u) = Ψ'(u) ∇u`, so the bilinear integrand factors by pulling
    -- the scalar `Ψ'(u)` out of the inner product.
    have hgrad_eq : hwΨ.weakGrad x = deriv Ψ (u x) • hu.weakGrad x := by
      ext i; simp [hwΨ, MemW1pWitness.comp_smooth_bounded, smul_eq_mul]
    rw [bilinFormIntegrandOfCoeff, hgrad_eq, inner_smul_right]
    simp [bilinFormIntegrandOfCoeff]
  -- The ⟨A∇u, ∇φ⟩ term from the expansion lemma
  let fderivφ_vec : E → E := fun x =>
    WithLp.toLp 2 fun i => (fderiv ℝ φ x) (EuclideanSpace.single i 1)
  let AuDφ : E → ℝ := fun x =>
    @inner ℝ E _ (matMulE (A.1.a x) (hu.weakGrad x)) (fderivφ_vec x)
  -- Expand the double product: ⟨A∇u, ∇(φ·(φΨ))⟩ = φ·⟨A∇u, ∇(φΨ)⟩ + φΨ·AuDφ
  -- and ⟨A∇u, ∇(φΨ)⟩ = φ·Ψ'·⟨A∇u,∇u⟩ + Ψ·AuDφ
  -- Combined: φ²Ψ'⟨A∇u,∇u⟩ + 2φΨ·AuDφ
  have hexpand_outer := bilinFormIntegrand_mul_smooth_eq (d := d) hΩ A.1 hu hwφΨ hφ
    (by norm_num : (0 : ℝ) ≤ 1) hCφ hφ_bound hφ_grad_bound (by norm_num : (1 : ENNReal) ≤ 2)
  have hexpand_inner := bilinFormIntegrand_mul_smooth_eq (d := d) hΩ A.1 hu hwΨ hφ
    (by norm_num : (0 : ℝ) ≤ 1) hCφ hφ_bound hφ_grad_bound (by norm_num : (1 : ENNReal) ≤ 2)
  -- Pointwise: integrand of bilinFormOfCoeff = φ²Ψ'⟨A∇u,∇u⟩ + 2φΨ·AuDφ
  let P_x : E → ℝ := fun x =>
    φ x ^ 2 * deriv Ψ (u x) * bilinFormIntegrandOfCoeff A.1 hu hu x
  let Q_x : E → ℝ := fun x => 2 * φ x * Ψ (u x) * AuDφ x
  have hPQ : ∀ x, bilinFormIntegrandOfCoeff A.1 hu hwφ2Ψ x = P_x x + Q_x x := by
    intro x
    rw [hexpand_outer x, hexpand_inner x, hΨu_bilin x]
    dsimp [P_x, Q_x, AuDφ]
    ring
  -- From htest: 0 ≤ ∫(P + Q). Since Ψ' ≤ 0, P_x ≤ 0.
  -- So 0 ≤ ∫P + ∫Q, with ∫P ≤ 0. Thus -∫P ≤ ∫Q, i.e., |∫P| ≤ ∫Q.
  -- Equivalently: ∫ φ²|Ψ'|⟨A∇u,∇u⟩ ≤ 2∫ φΨ·AuDφ.
  have hPQ_int : 0 ≤ ∫ x in Ω, (P_x x + Q_x x) ∂volume := by
    have heq : ∫ x in Ω, (P_x x + Q_x x) ∂volume =
        ∫ x in Ω, bilinFormIntegrandOfCoeff A.1 hu hwφ2Ψ x ∂volume :=
      setIntegral_congr_fun hΩ.measurableSet (fun x _ => (hPQ x).symm)
    rw [heq]; exact htest
  -- Define the key quantities as real numbers.
  let absP : ℝ := ∫ x in Ω, φ x ^ 2 * (-deriv Ψ (u x)) *
    bilinFormIntegrandOfCoeff A.1 hu hu x ∂volume
  let C : ℝ := ∫ x in Ω, φ x * Ψ (u x) * AuDφ x ∂volume
  let R : ℝ := ∫ x in Ω, ‖fderiv ℝ φ x‖ ^ 2 *
    (if deriv Ψ (u x) = 0 then 0 else Ψ (u x) ^ 2 / (-deriv Ψ (u x))) ∂volume
  -- (a) |P| ≥ 0 (integrand is nonneg: φ²·(-Ψ')·⟨A∇u,∇u⟩ ≥ 0)
  have habsP_nonneg : 0 ≤ absP := by
    dsimp [absP]
    -- Use setIntegral_nonneg with pointwise bound (need coercivity pointwise on Ω)
    -- The a.e. coercivity suffices since we can use ae_restrict.
    rw [show (∫ x in Ω, φ x ^ 2 * -deriv Ψ (u x) * bilinFormIntegrandOfCoeff A.1 hu hu x ∂volume)
      = ∫ x, φ x ^ 2 * -deriv Ψ (u x) * bilinFormIntegrandOfCoeff A.1 hu hu x
          ∂(volume.restrict Ω) from rfl]
    apply integral_nonneg_of_ae
    filter_upwards [A.1.ae_coercive_nonneg, ae_restrict_mem hΩ.measurableSet] with x hcoer hx
    apply mul_nonneg
    · exact mul_nonneg (sq_nonneg _) (neg_nonneg.mpr (hΨ'_nonpos x hx))
    · rw [bilinFormIntegrandOfCoeff, real_inner_comm]; exact hcoer (hu.weakGrad x)
  obtain ⟨M, hM⟩ := hΨ_bdd
  have hP_int : Integrable P_x (volume.restrict Ω) := by
    -- P_x = (φ²·Ψ'(u)) · bilinForm. Use integrable_bounded_mul_bilinFormIntegrand.
    have hfbound : ∀ x, |φ x ^ 2 * deriv Ψ (u x)| ≤ M := by
      intro x
      rw [abs_mul]
      calc |φ x ^ 2| * |deriv Ψ (u x)|
          ≤ 1 * M := by
            exact mul_le_mul
              (by rw [abs_of_nonneg (sq_nonneg _)]; exact (sq_le_one_iff_abs_le_one _).mpr (hφ_bound x))
              (hM (u x)) (abs_nonneg _) zero_le_one
        _ = M := one_mul M
    have hfmeas : AEStronglyMeasurable (fun x => φ x ^ 2 * deriv Ψ (u x))
        (volume.restrict Ω) :=
      (hφ.continuous.pow 2).aestronglyMeasurable.mul
        ((hΨ.continuous_deriv (by simp)).comp_aestronglyMeasurable
          hu.memLp.aestronglyMeasurable)
    simpa [P_x] using
      (integrable_bounded_mul_bilinFormIntegrand A.1 hu hfmeas hfbound)
  have hPQ_integrable : Integrable (fun x => P_x x + Q_x x) (volume.restrict Ω) := by
    have := integrable_bilinFormIntegrandOfCoeff A.1 hu hwφ2Ψ
    exact this.congr (ae_of_all _ fun x => hPQ x)
  have hQ_int : Integrable Q_x (volume.restrict Ω) := by
    exact (hPQ_integrable.sub hP_int).congr (ae_of_all _ fun x => by dsimp; ring)
  have hcross_int : Integrable (fun x => φ x * Ψ (u x) * AuDφ x) (volume.restrict Ω) := by
    refine (hQ_int.const_mul (1 / 2 : ℝ)).congr ?_
    exact ae_of_all _ fun x => by
      dsimp [Q_x, AuDφ]
      ring
  -- (b) From 0 ≤ ∫(P+Q): since P_x + Q_x = φ²Ψ'⟨A∇u,∇u⟩ + 2φΨ·AuDφ
  --   = -(φ²(-Ψ')⟨A∇u,∇u⟩) + 2(φΨ·AuDφ) = -absP_integrand + 2C_integrand
  -- So ∫(P+Q) = -absP + 2C. Combined with 0 ≤ ∫(P+Q): absP ≤ 2C.
  have habsP_le_2C : absP ≤ 2 * C := by
    -- Now split: ∫(P+Q) = ∫P + ∫Q
    have hsplit : ∫ x in Ω, (P_x x + Q_x x) ∂volume =
        ∫ x in Ω, P_x x ∂volume + ∫ x in Ω, Q_x x ∂volume :=
      integral_add hP_int hQ_int
    -- absP = -∫P and 2C = ∫Q
    have hP_neg : ∫ x in Ω, P_x x ∂volume = -absP := by
      dsimp [absP, P_x]
      rw [← integral_neg]; congr 1; ext x; ring
    have hQ_2C : ∫ x in Ω, Q_x x ∂volume = 2 * C := by
      dsimp [C, Q_x]
      rw [← integral_const_mul]; congr 1; ext x; ring
    linarith [hPQ_int, hsplit, hP_neg, hQ_2C]
  have hR_nonneg : 0 ≤ R := by
    dsimp [R]
    apply setIntegral_nonneg_ae hΩ.measurableSet
    filter_upwards with x hx
    apply mul_nonneg (sq_nonneg _)
    split_ifs with hder
    · positivity
    · exact div_nonneg (sq_nonneg _) (neg_nonneg.mpr (hΨ'_nonpos x hx))
  -- (c) Integral CS: C² ≤ Λ · absP · R
  have hF_sq_int :
      Integrable
        (fun x => φ x ^ 2 * (-deriv Ψ (u x)) *
          bilinFormIntegrandOfCoeff A.1 hu hu x)
        (volume.restrict Ω) := by
    refine hP_int.neg.congr ?_
    exact ae_of_all _ fun x => by
      dsimp [P_x]
      ring
  let μ : Measure E := volume.restrict Ω
  let F : E → ℝ := fun x =>
    Real.sqrt (φ x ^ 2 * (-deriv Ψ (u x)) *
      bilinFormIntegrandOfCoeff A.1 hu hu x)
  let G : E → ℝ := fun x =>
    Real.sqrt (A.1.Λ * (‖fderiv ℝ φ x‖ ^ 2 *
      (if deriv Ψ (u x) = 0 then 0
       else Ψ (u x) ^ 2 / (-deriv Ψ (u x)))))
  have hF_arg_nonneg :
      ∀ᵐ x ∂μ, 0 ≤ φ x ^ 2 * (-deriv Ψ (u x)) *
        bilinFormIntegrandOfCoeff A.1 hu hu x := by
    filter_upwards [A.1.ae_coercive_nonneg, ae_restrict_mem hΩ.measurableSet] with x hcoer hx
    apply mul_nonneg
    · exact mul_nonneg (sq_nonneg _) (neg_nonneg.mpr (hΨ'_nonpos x hx))
    · rw [bilinFormIntegrandOfCoeff, real_inner_comm]
      exact hcoer (hu.weakGrad x)
  have hF_sq_eq :
      (fun x => F x ^ 2) =ᵐ[μ]
        (fun x => φ x ^ 2 * (-deriv Ψ (u x)) *
          bilinFormIntegrandOfCoeff A.1 hu hu x) := by
    filter_upwards [hF_arg_nonneg] with x hx
    dsimp [F]
    rw [Real.sq_sqrt hx]
  have hG_sq_int :
      Integrable
        (fun x => A.1.Λ * (‖fderiv ℝ φ x‖ ^ 2 *
          (if deriv Ψ (u x) = 0 then 0
           else Ψ (u x) ^ 2 / (-deriv Ψ (u x)))))
        μ := by
    simpa [μ] using hratio_int.const_mul A.1.Λ
  have hG_arg_nonneg :
      ∀ᵐ x ∂μ, 0 ≤ A.1.Λ * (‖fderiv ℝ φ x‖ ^ 2 *
        (if deriv Ψ (u x) = 0 then 0
         else Ψ (u x) ^ 2 / (-deriv Ψ (u x)))) := by
    filter_upwards [ae_restrict_mem hΩ.measurableSet] with x hx
    apply mul_nonneg A.1.Λ_nonneg
    apply mul_nonneg (sq_nonneg _)
    split_ifs with hder
    · positivity
    · exact div_nonneg (sq_nonneg _) (neg_nonneg.mpr (hΨ'_nonpos x hx))
  have hG_sq_eq :
      (fun x => G x ^ 2) =ᵐ[μ]
        (fun x => A.1.Λ * (‖fderiv ℝ φ x‖ ^ 2 *
          (if deriv Ψ (u x) = 0 then 0
           else Ψ (u x) ^ 2 / (-deriv Ψ (u x))))) := by
    filter_upwards [hG_arg_nonneg] with x hx
    dsimp [G]
    rw [Real.sq_sqrt hx]
  have hF_meas : AEStronglyMeasurable F μ := by
    exact (hF_sq_int.aestronglyMeasurable.aemeasurable.sqrt).aestronglyMeasurable
  have hG_meas : AEStronglyMeasurable G μ := by
    exact (hG_sq_int.aestronglyMeasurable.aemeasurable.sqrt).aestronglyMeasurable
  have hF_memLp : MemLp F 2 μ := by
    refine (MeasureTheory.memLp_two_iff_integrable_sq hF_meas).2 ?_
    exact hF_sq_int.congr hF_sq_eq.symm
  have hG_memLp : MemLp G 2 μ := by
    refine (MeasureTheory.memLp_two_iff_integrable_sq hG_meas).2 ?_
    exact hG_sq_int.congr hG_sq_eq.symm
  have hFG_int : Integrable (fun x => F x * G x) μ := by
    rw [← memLp_one_iff_integrable]
    have hmem : MemLp (fun x => ‖F x‖ * ‖G x‖) 1 μ := by
      simpa [mul_comm] using (hF_memLp.norm.mul hG_memLp.norm)
    simpa [F, G, Real.norm_of_nonneg (Real.sqrt_nonneg _)] using hmem
  have hcross_bound :
      ∀ᵐ x ∂μ, |φ x * Ψ (u x) * AuDφ x| ≤ F x * G x := by
    filter_upwards [A.1.mulVec_sq_le, A.1.ae_coercive_nonneg,
      ae_restrict_mem hΩ.measurableSet] with x hmul hcoer hx
    have hbilin_nonneg : 0 ≤ bilinFormIntegrandOfCoeff A.1 hu hu x := by
      rw [bilinFormIntegrandOfCoeff, real_inner_comm]
      exact hcoer (hu.weakGrad x)
    have hder_nonpos : deriv Ψ (u x) ≤ 0 := hΨ'_nonpos x hx
    by_cases hder : deriv Ψ (u x) = 0
    · have hΨx : Ψ (u x) = 0 := hΨ_zero_of_deriv_zero x hx hder
      simp [F, G, hder, hΨx]
    · have hder_lt : deriv Ψ (u x) < 0 := by
        exact lt_of_le_of_ne hder_nonpos (by simpa using hder)
      have hneg : 0 < -deriv Ψ (u x) := by linarith
      have hAu_abs :
          |AuDφ x| ≤ ‖matMulE (A.1.a x) (hu.weakGrad x)‖ * ‖fderiv ℝ φ x‖ := by
        calc
          |AuDφ x|
              ≤ ‖matMulE (A.1.a x) (hu.weakGrad x)‖ *
                  ‖WithLp.toLp 2
                    (fun i => (fderiv ℝ φ x) (EuclideanSpace.single i 1))‖ := by
                dsimp [AuDφ]
                exact abs_real_inner_le_norm _ _
          _ = ‖matMulE (A.1.a x) (hu.weakGrad x)‖ * ‖fderiv ℝ φ x‖ := by
                rw [super_norm_fderiv_eq_norm_partials]
      have hAu_sq :
          |AuDφ x| ^ 2 ≤
            A.1.Λ * bilinFormIntegrandOfCoeff A.1 hu hu x * ‖fderiv ℝ φ x‖ ^ 2 := by
        have hsq₁ :
            |AuDφ x| ^ 2 ≤
              (‖matMulE (A.1.a x) (hu.weakGrad x)‖ * ‖fderiv ℝ φ x‖) ^ 2 := by
          exact sq_le_sq.mpr <| by
            simpa [abs_of_nonneg (mul_nonneg (norm_nonneg _) (norm_nonneg _))] using hAu_abs
        have hmul_sq :
            ‖matMulE (A.1.a x) (hu.weakGrad x)‖ ^ 2 ≤
              A.1.Λ * bilinFormIntegrandOfCoeff A.1 hu hu x := by
          simpa [bilinFormIntegrandOfCoeff, real_inner_comm] using hmul (hu.weakGrad x)
        calc
          |AuDφ x| ^ 2 ≤
              (‖matMulE (A.1.a x) (hu.weakGrad x)‖ * ‖fderiv ℝ φ x‖) ^ 2 := hsq₁
          _ = ‖matMulE (A.1.a x) (hu.weakGrad x)‖ ^ 2 * ‖fderiv ℝ φ x‖ ^ 2 := by
                ring
          _ ≤
              (A.1.Λ * bilinFormIntegrandOfCoeff A.1 hu hu x) * ‖fderiv ℝ φ x‖ ^ 2 := by
                exact mul_le_mul_of_nonneg_right hmul_sq (sq_nonneg _)
          _ = A.1.Λ * bilinFormIntegrandOfCoeff A.1 hu hu x * ‖fderiv ℝ φ x‖ ^ 2 := by
                ring
      have hFarg_nonneg :
          0 ≤ φ x ^ 2 * (-deriv Ψ (u x)) *
            bilinFormIntegrandOfCoeff A.1 hu hu x := by
        apply mul_nonneg
        · exact mul_nonneg (sq_nonneg _) hneg.le
        · exact hbilin_nonneg
      have hGarg_nonneg :
          0 ≤ A.1.Λ * (‖fderiv ℝ φ x‖ ^ 2 *
            (if deriv Ψ (u x) = 0 then 0
             else Ψ (u x) ^ 2 / (-deriv Ψ (u x)))) := by
        apply mul_nonneg A.1.Λ_nonneg
        apply mul_nonneg (sq_nonneg _)
        simp [hder, div_nonneg (sq_nonneg _) hneg.le]
      have hcross_sq :
          |φ x * Ψ (u x) * AuDφ x| ^ 2 ≤
            (φ x ^ 2 * (-deriv Ψ (u x)) *
              bilinFormIntegrandOfCoeff A.1 hu hu x) *
            (A.1.Λ * (‖fderiv ℝ φ x‖ ^ 2 *
              (if deriv Ψ (u x) = 0 then 0
               else Ψ (u x) ^ 2 / (-deriv Ψ (u x))))) := by
        calc
          |φ x * Ψ (u x) * AuDφ x| ^ 2
              = φ x ^ 2 * Ψ (u x) ^ 2 * |AuDφ x| ^ 2 := by
                  rw [abs_mul, abs_mul]
                  calc
                    (|φ x| * |Ψ (u x)| * |AuDφ x|) ^ 2
                        = |φ x| ^ 2 * |Ψ (u x)| ^ 2 * |AuDφ x| ^ 2 := by ring
                    _ = φ x ^ 2 * Ψ (u x) ^ 2 * |AuDφ x| ^ 2 := by
                        rw [sq_abs, sq_abs]
          _ ≤
              φ x ^ 2 * Ψ (u x) ^ 2 *
                (A.1.Λ * bilinFormIntegrandOfCoeff A.1 hu hu x *
                  ‖fderiv ℝ φ x‖ ^ 2) := by
                  exact mul_le_mul_of_nonneg_left hAu_sq
                    (mul_nonneg (sq_nonneg _) (sq_nonneg _))
          _ =
              (φ x ^ 2 * (-deriv Ψ (u x)) *
                bilinFormIntegrandOfCoeff A.1 hu hu x) *
              (A.1.Λ * (‖fderiv ℝ φ x‖ ^ 2 *
                (if deriv Ψ (u x) = 0 then 0
                 else Ψ (u x) ^ 2 / (-deriv Ψ (u x))))) := by
                  simp [hder]
                  field_simp [hneg.ne']
      exact (sq_le_sq₀ (abs_nonneg _) (mul_nonneg (Real.sqrt_nonneg _) (Real.sqrt_nonneg _))).1 <|
        calc
          |φ x * Ψ (u x) * AuDφ x| ^ 2
              ≤
                (φ x ^ 2 * (-deriv Ψ (u x)) *
                  bilinFormIntegrandOfCoeff A.1 hu hu x) *
                (A.1.Λ * (‖fderiv ℝ φ x‖ ^ 2 *
                  (if deriv Ψ (u x) = 0 then 0
                   else Ψ (u x) ^ 2 / (-deriv Ψ (u x))))) := hcross_sq
          _ = (F x * G x) ^ 2 := by
                dsimp [F, G]
                rw [mul_pow, Real.sq_sqrt hFarg_nonneg, Real.sq_sqrt hGarg_nonneg]
  have habs_le :
      |C| ≤ ∫ x, F x * G x ∂μ := by
    calc
      |C| = |∫ x, φ x * Ψ (u x) * AuDφ x ∂μ| := by
              rfl
      _ = ‖∫ x, φ x * Ψ (u x) * AuDφ x ∂μ‖ := by
              rw [Real.norm_eq_abs]
      _ ≤ ∫ x, ‖φ x * Ψ (u x) * AuDφ x‖ ∂μ := norm_integral_le_integral_norm _
      _ = ∫ x, |φ x * Ψ (u x) * AuDφ x| ∂μ := by
              apply integral_congr_ae
              filter_upwards with x
              rw [Real.norm_eq_abs]
      _ ≤ ∫ x, F x * G x ∂μ := integral_mono_ae hcross_int.norm hFG_int hcross_bound
  have hF_memLp' : MemLp F (ENNReal.ofReal 2) μ := by
    simpa using hF_memLp
  have hG_memLp' : MemLp G (ENNReal.ofReal 2) μ := by
    simpa using hG_memLp
  have hholder :
      ∫ x, F x * G x ∂μ ≤
        (∫ x, F x ^ (2 : ℝ) ∂μ) ^ (1 / (2 : ℝ)) *
          (∫ x, G x ^ (2 : ℝ) ∂μ) ^ (1 / (2 : ℝ)) := by
    simpa [F, G, Real.norm_of_nonneg (Real.sqrt_nonneg _)] using
      integral_mul_norm_le_Lp_mul_Lq (μ := μ) (f := F) (g := G)
        Real.HolderConjugate.two_two hF_memLp' hG_memLp'
  have hF_int_eq : ∫ x, F x ^ (2 : ℝ) ∂μ = absP := by
    calc
      ∫ x, F x ^ (2 : ℝ) ∂μ = ∫ x, F x ^ (2 : ℕ) ∂μ := by
            congr 1
            ext x
            rw [show (2 : ℝ) = ((2 : ℕ) : ℝ) by norm_num, Real.rpow_natCast]
      _ = ∫ x, φ x ^ 2 * (-deriv Ψ (u x)) *
            bilinFormIntegrandOfCoeff A.1 hu hu x ∂μ := by
              apply integral_congr_ae
              exact hF_sq_eq
      _ = absP := by rfl
  have hG_int_eq : ∫ x, G x ^ (2 : ℝ) ∂μ = A.1.Λ * R := by
    calc
      ∫ x, G x ^ (2 : ℝ) ∂μ = ∫ x, G x ^ (2 : ℕ) ∂μ := by
            congr 1
            ext x
            rw [show (2 : ℝ) = ((2 : ℕ) : ℝ) by norm_num, Real.rpow_natCast]
      _ = ∫ x, A.1.Λ * (‖fderiv ℝ φ x‖ ^ 2 *
            (if deriv Ψ (u x) = 0 then 0
             else Ψ (u x) ^ 2 / (-deriv Ψ (u x)))) ∂μ := by
              apply integral_congr_ae
              exact hG_sq_eq
      _ = A.1.Λ * R := by
            dsimp [R, μ]
            rw [integral_const_mul]
  have hC_bound :
      |C| ≤ absP ^ ((1 : ℝ) / 2) * (A.1.Λ * R) ^ ((1 : ℝ) / 2) := by
    calc
      |C| ≤ ∫ x, F x * G x ∂μ := habs_le
      _ ≤ (∫ x, F x ^ (2 : ℝ) ∂μ) ^ (1 / (2 : ℝ)) *
            (∫ x, G x ^ (2 : ℝ) ∂μ) ^ (1 / (2 : ℝ)) := hholder
      _ = absP ^ ((1 : ℝ) / 2) * (A.1.Λ * R) ^ ((1 : ℝ) / 2) := by
            rw [hF_int_eq, hG_int_eq]
  have hCS : C ^ 2 ≤ A.1.Λ * absP * R := by
    have hAR_nonneg : 0 ≤ A.1.Λ * R := mul_nonneg A.1.Λ_nonneg hR_nonneg
    have hsq :=
      super_sq_le_of_le_rpow_half_mul habsP_nonneg hAR_nonneg (abs_nonneg C) hC_bound
    simpa [sq_abs, mul_assoc, mul_left_comm, mul_comm] using hsq
  -- Since `R` is defined using `‖∇φ‖²`, the upper ellipticity bound
  -- `⟨A∇φ, ∇φ⟩ ≤ Λ ‖∇φ‖²` contributes the factor `Λ` in the final estimate.
  show absP ≤ 4 * A.1.Λ * R
  by_cases habsP_zero : absP = 0
  · rw [habsP_zero]
    apply mul_nonneg (mul_nonneg (by positivity) A.1.Λ_nonneg)
    exact hR_nonneg
  · have habsP_pos : 0 < absP := lt_of_le_of_ne habsP_nonneg (Ne.symm habsP_zero)
    have h1 : absP ^ 2 ≤ (2 * C) ^ 2 := sq_le_sq' (by linarith) habsP_le_2C
    have h2 : (2 * C) ^ 2 = 4 * C ^ 2 := by ring
    nlinarith [h1, h2, hCS, habsP_pos, A.1.Λ_nonneg, hR_nonneg]

\end{lstlisting}
\begin{proof}
By the supersolution property tested against $\varphi^{2}\Psi(u) \geq 0$,
\[
  0 \leq \int_{\Omega} \langle A\nabla u, \nabla(\varphi^{2}\Psi(u))\rangle\,dx.
\]
Using the chain rule for $W^{1,2}$ functions, $\nabla(\Psi(u)) = \Psi'(u)\nabla u$,
and the product rule $\nabla(\varphi^{2}\Psi(u)) = \varphi^{2}\Psi'(u)\nabla u
+ 2\varphi\Psi(u)\nabla \varphi$. The bilinear form expansion
(\leanname{bilinFormIntegrand\_mul\_smooth\_eq}) yields the pointwise identity
\[
  \langle A\nabla u, \nabla(\varphi^{2}\Psi(u))\rangle
  = \varphi^{2}\Psi'(u)\,\langle A\nabla u,\nabla u\rangle
  + 2\varphi\Psi(u)\,\langle A\nabla u,\nabla \varphi\rangle.
\]
Since $\Psi'(u) \leq 0$, the principal term is non-positive; rearranging
\[
  \int \varphi^{2}\,(-\Psi'(u))\,\langle A\nabla u,\nabla u\rangle\,dx
  \;\leq\;
  2 \int \varphi\,\Psi(u)\,\langle A\nabla u,\nabla \varphi\rangle\,dx.
\]
By Cauchy--Schwarz with respect to the symmetric form $\langle A\cdot,\cdot\rangle$
and AM--GM with weight $\varepsilon = 1/2$,
\begin{multline*}
  2|\varphi\Psi(u)\,\langle A\nabla u,\nabla \varphi\rangle|
  \\ \leq\;
  \tfrac12\,\varphi^{2}(-\Psi'(u))\langle A\nabla u,\nabla u\rangle
  + 2\,\frac{\Psi(u)^{2}}{-\Psi'(u)}\,\langle A\nabla \varphi,\nabla \varphi\rangle.
\end{multline*}
Absorbing the first term into the left-hand side and bounding
$\langle A\nabla \varphi,\nabla \varphi\rangle \leq \Lambda \|\nabla \varphi\|^{2}$
(ellipticity \leanname{NormalizedEllipticCoeff.upper}) yields the stated inequality
with constant $4\Lambda$.
\end{proof}

\subsubsection{Reverse Holder inequalities (one--step)}

\paragraph{Inverse case.}

\begin{definition}[Inverse-power cutoff, \leanname{superPowerCutoff}]
For a smooth $\eta$, positive $u$, and $p > 0$,
\[
  \mathrm{PCI}_d(\eta,u,p)(x) := \eta(x)\,|u(x)^{-1}|^{p/2}.
\]
\end{definition}

\begin{lstlisting}[style=Matlab-editor,basicstyle=\mlttfamily,escapechar=",mlshowsectionrules=true,mathescape=true,morecomment={[s]{\%\{}{\%\}}},language=lean,numbers=none]
/-! ### One-Step Sobolev Gain for Supersolutions

These are the analogues of `moser_preMoser` from the subsolution theory.
The key difference: the test function is `φ² u^{-(1+p)}` instead of
`φ² u^{p-1}`, and the Caccioppoli inequality produces
`∫ |∇(φ u^{-p/2})|² ≤ C(Λp²/(1+p)² + 1) ∫ |∇φ|² u^{-p}`
for the weighted cutoff.

The proof is decomposed into three sub-results:
1. `superPowerCutoff_energy_bound`: W^{1,2} witness + energy bound for
   `η · u^{-p/2}` (the Caccioppoli inequality)
2. `supersolution_preMoser_inverse`: wire up energy bound + Sobolev → Lp gain
3. `supersolution_preMoser_forward`: the `p ∈ (0,1)` variant
-/

/-- Inverse-power cutoff used in the supersolution pre-estimate: `η · u^{-p/2}`.
This is the supersolution analogue of `moserPowerCutoff` (which uses `u₊^{p/2}`). -/
noncomputable def superPowerCutoff
    (η u : E → ℝ) (p : ℝ) : E → ℝ :=
  fun x => η x * |(u x)⁻¹| ^ (p / 2)

\end{lstlisting}
\begin{theorem}[Inverse-power one-step gain, \leanname{supersolution\_preMoser\_inverse}]
\label{thm:preMoser-inv}
Let $d>2$, $A$ as above, $u > 0$ a supersolution, $p > 0$, and
$0 < r < s \leq 1$ with $\int_{\Bz{s}} |u^{-1}|^{p}\,dx < \infty$. Then
$|u^{-1}|^{\chi p}$ is integrable on $\Bz{r}$ and
\begin{multline*}
  \Big(\int_{\Bz{r}} |u^{-1}|^{\chi p}\,dx\Big)^{1/(\chi p)}
  \\ \leq\;
  \Big(\frac{C_{\mathrm{MA}}(d)}{(s-r)^{2}}\,
    \Big(\Lambda\Big(\frac{p}{1+p}\Big)^{2}+1\Big)\Big)^{1/p}
  \Big(\int_{\Bz{s}} |u^{-1}|^{p}\,dx\Big)^{1/p}.
\end{multline*}
\end{theorem}

\begin{lstlisting}[style=Matlab-editor,basicstyle=\mlttfamily,escapechar=",mlshowsectionrules=true,mathescape=true,morecomment={[s]{\%\{}{\%\}}},language=lean,numbers=none]
/-- One-step Sobolev gain for positive supersolutions, inverse-power case.

For `p > 0` and `0 < r < s ≤ 1`, if `u > 0` is a supersolution on `B₁`, then
```
‖u⁻¹‖_{L^{pχ}(Bᵣ)} ≤ C^{1/p} (Λp²/(1+p)² + 1)^{1/p} (s-r)^{-2/p} ‖u⁻¹‖_{Lᵖ(Bₛ)}
```
Proof: set up a cutoff `η` with `1_{Bᵣ} ≤ η ≤ 1_{Bₛ}` and `‖∇η‖ ≤ 2/(s-r)`,
apply `superPowerCutoff_energy_bound` to get the W^{1,2} energy bound, then
apply the Sobolev inequality to `v = η · u^{-p/2}` to obtain the Lᵖ → Lᵖˣ gain.
The final step is norm arithmetic converting from L²* to Lᵖˣ. -/
theorem supersolution_preMoser_inverse
    (hd : 2 < (d : ℝ))
    (A : NormalizedEllipticCoeff d (Metric.ball (0 : E) 1))
    {u : E → ℝ} {p r s : ℝ}
    (hp : 0 < p)
    (hr : 0 < r) (hrs : r < s) (hs1 : s ≤ 1)
    (hu_pos : ∀ x ∈ Metric.ball (0 : E) 1, 0 < u x)
    (hsuper : IsSupersolution A.1 u)
    (hpInt :
      IntegrableOn (fun x => |(u x)⁻¹| ^ p)
        (Metric.ball (0 : E) s) volume) :
    IntegrableOn (fun x => |(u x)⁻¹| ^ (moserChi d * p))
        (Metric.ball (0 : E) r) volume ∧
      (∫ x in Metric.ball (0 : E) r,
          |(u x)⁻¹| ^ (moserChi d * p) ∂volume) ^
          (1 / (moserChi d * p)) ≤
        ((C_MoserAnchor d / (s - r) ^ 2) *
            (A.1.Λ * (p / (1 + p)) ^ 2 + 1)) ^ (1 / p) *
          (∫ x in Metric.ball (0 : E) s, |(u x)⁻¹| ^ p ∂volume) ^ (1 / p) := by
  -- Set up the cutoff function
  let Ω : Set E := Metric.ball (0 : E) s
  let R : ℝ := (r + s) / 2
  have hs_pos : 0 < s := lt_trans hr hrs
  have hsr_pos : 0 < s - r := by linarith
  have hR : r < R := by dsimp [R]; linarith
  have hR_lt_s : R < s := by dsimp [R]; linarith
  let ηCut : Cutoff (x₀ := (0 : E)) r R := Classical.choose
    (Cutoff.exists (d := d) (0 : E) (r := r) (R := R) hr hR)
  let η : E → ℝ := ηCut.toFun
  let Cη : ℝ := 2 * (Mst : ℝ) * (s - r)⁻¹
  have hη : ContDiff ℝ (⊤ : ℕ∞) η := ηCut.smooth
  have hη_bound : ∀ x, |η x| ≤ 1 := by
    intro x
    rw [abs_of_nonneg (ηCut.nonneg x)]
    exact ηCut.le_one x
  have hη_eq_one : ∀ x ∈ Metric.ball (0 : E) r, η x = 1 := by
    intro x hx; exact ηCut.eq_one x hx
  have hη_grad_bound : ∀ x, ‖fderiv ℝ η x‖ ≤ Cη := by
    intro x
    have hmid : (((r + s) / 2 : ℝ) - r) = (s - r) / 2 := by ring
    have h_inv : ((((r + s) / 2 : ℝ) - r)⁻¹) = 2 * (s - r)⁻¹ := by
      rw [hmid]; field_simp [hsr_pos.ne']
    calc ‖fderiv ℝ η x‖ ≤ ↑Mst * (((r + s) / 2 : ℝ) - r)⁻¹ := ηCut.grad_bound x
      _ = Cη := by rw [h_inv]; ring
  have hη_sub_ball : tsupport η ⊆ Ω := by
    intro x hx
    exact (Metric.closedBall_subset_ball hR_lt_s) (ηCut.support_subset hx)
  -- Apply the energy bound
  obtain ⟨hwv, hvW01, henergy⟩ :=
    superPowerCutoff_energy_bound hd A hp hs_pos hs1 hu_pos hsuper hpInt
      hη hη_bound hη_grad_bound hη_sub_ball
  let v : E → ℝ := superPowerCutoff (d := d) η u p
  have hv_support : tsupport v ⊆ Ω :=
    superPowerCutoff_tsupport_subset (d := d) hη_sub_ball
  -- Apply the Sobolev inequality via `sobolev_prepare_on_ball`.
  obtain ⟨hwv_real, hSob⟩ :=
    sobolev_prepare_on_ball (d := d) (v := v) hd hs_pos hvW01 hv_support
  let μ : Measure E := volume.restrict Ω
  -- Convert Sobolev bound from eLpNorm to real lpNorm
  let qexp : ℝ := ((d : ℝ) * 2 / ((d : ℝ) - 2))
  let q := ENNReal.ofReal qexp
  have hq_pos : 0 < qexp := by
    dsimp [qexp]; exact div_pos (mul_pos (by linarith) (by norm_num)) (by linarith)
  let hwvReal : MemW1pWitness (ENNReal.ofReal (2 : ℝ)) v Ω :=
    { memLp := by simpa [Ω] using hwv.memLp
      weakGrad := hwv.weakGrad
      weakGrad_component_memLp := by simpa [Ω] using hwv.weakGrad_component_memLp
      isWeakGrad := by simpa [Ω] using hwv.isWeakGrad }
  have hae_grad :
      hwv_real.weakGrad =ᵐ[μ] hwvReal.weakGrad := by
    simpa [μ, Ω] using
      (MemW1pWitness.ae_eq_p (d := d) Metric.isOpen_ball
        (p := (2 : ℝ)) (by norm_num) hwv_real hwvReal)
  have hgrad_sq_eq :
      ∫ x in Ω, ‖hwv_real.weakGrad x‖ ^ 2 ∂volume =
        ∫ x in Ω, ‖hwvReal.weakGrad x‖ ^ 2 ∂volume := by
    exact integral_congr_ae (hae_grad.fun_comp (fun z => ‖z‖ ^ 2))
  have hv_memLp_q : MemLp v q μ := by
    refine ⟨hwv_real.memLp.aestronglyMeasurable, ?_⟩
    exact lt_of_le_of_lt hSob
      (ENNReal.mul_lt_top_iff.mpr (Or.inl ⟨by simp,
        lt_top_iff_ne_top.mpr (by simpa using hwv_real.weakGrad_norm_memLp.eLpNorm_ne_top)⟩))
  have hSob' : eLpNorm v q μ ≤
      ENNReal.ofReal (C_gns d 2) *
        eLpNorm (fun x => ‖hwv_real.weakGrad x‖) 2 μ := by
    simpa [q, μ, Ω] using hSob
  have hgrad_ne_top : eLpNorm (fun x => ‖hwv_real.weakGrad x‖) 2 μ ≠ ⊤ := by
    simpa using hwv_real.weakGrad_norm_memLp.eLpNorm_ne_top
  have hSob_real :
      MeasureTheory.lpNorm v q μ ≤
        (C_gns d 2) * MeasureTheory.lpNorm (fun x => ‖hwv_real.weakGrad x‖) 2 μ := by
    have hSob_toReal :=
      (ENNReal.toReal_le_toReal hv_memLp_q.eLpNorm_ne_top
        (ENNReal.mul_ne_top ENNReal.ofReal_ne_top hgrad_ne_top)).2 hSob'
    rwa [MeasureTheory.toReal_eLpNorm hv_memLp_q.aestronglyMeasurable,
      ENNReal.toReal_mul, MeasureTheory.toReal_eLpNorm
        hwv_real.weakGrad_norm_memLp.aestronglyMeasurable,
      ENNReal.toReal_ofReal (C_gns_nonneg d 2)] at hSob_toReal
  have hgrad_lp :
      MeasureTheory.lpNorm (fun x => ‖hwv_real.weakGrad x‖) 2 μ =
        (∫ x in Ω, ‖hwv_real.weakGrad x‖ ^ 2 ∂volume) ^ (1 / (2 : ℝ)) := by
    simpa [μ, Ω] using
      (MeasureTheory.lpNorm_eq_integral_norm_rpow_toReal (μ := μ)
        (f := fun x => ‖hwv_real.weakGrad x‖)
        two_ne_zero ENNReal.ofNat_ne_top hwv_real.weakGrad_norm_memLp.aestronglyMeasurable)
  have hgrad_bound :
      MeasureTheory.lpNorm (fun x => ‖hwv_real.weakGrad x‖) 2 μ ≤
        (2 * Cη ^ 2 * (A.1.Λ * (p / (1 + p)) ^ 2 + 1) *
          ∫ x in Ω, |(u x)⁻¹| ^ p ∂volume) ^ (1 / (2 : ℝ)) := by
    rw [hgrad_lp, hgrad_sq_eq]
    exact Real.rpow_le_rpow (by exact integral_nonneg fun x => by positivity)
      (by simpa [hwvReal] using henergy) (by norm_num)
  -- Convert v on Bᵣ to |u⁻¹|^{p/2} using η = 1
  have hq_ne_zero : q ≠ 0 := by
    dsimp [q]; rw [ENNReal.ofReal_eq_zero]; linarith
  have hv_lp_eq :
      MeasureTheory.lpNorm v q μ =
        (∫ x in Ω, |v x| ^ qexp ∂volume) ^ (1 / qexp) := by
    rw [MeasureTheory.lpNorm_eq_integral_norm_rpow_toReal (μ := μ) (f := v)]
    · simp [μ, Ω, q, qexp, ENNReal.toReal_ofReal hq_pos.le, Real.norm_eq_abs]
    · exact hq_ne_zero
    · exact ENNReal.ofReal_ne_top
    · exact hv_memLp_q.aestronglyMeasurable
  have hv_int : IntegrableOn (fun x => |v x| ^ qexp) Ω volume := by
    simpa [μ, Ω, q, qexp, ENNReal.toReal_ofReal hq_pos.le, Real.norm_eq_abs] using
      hv_memLp_q.integrable_norm_rpow hq_ne_zero ENNReal.ofReal_ne_top
  have hpow_eq_on : ∀ x ∈ Metric.ball (0 : E) r,
      |v x| ^ qexp = |(u x)⁻¹| ^ (moserChi d * p) := by
    intro x hx
    have hηx : η x = 1 := hη_eq_one x hx
    have hqexp_eq : qexp = 2 * moserChi d := by
      dsimp [qexp, moserChi]; field_simp
    have hvx : v x = |(u x)⁻¹| ^ (p / 2) := by
      simp [v, superPowerCutoff, hηx]
    calc |v x| ^ qexp
        = (|(u x)⁻¹| ^ (p / 2)) ^ qexp := by
          rw [hvx, abs_of_nonneg (Real.rpow_nonneg (abs_nonneg _) _)]
      _ = |(u x)⁻¹| ^ (moserChi d * p) := by
          rw [hqexp_eq, ← Real.rpow_mul (abs_nonneg _)]; ring_nf
  have hleft_int : IntegrableOn (fun x => |(u x)⁻¹| ^ (moserChi d * p))
      (Metric.ball (0 : E) r) volume := by
    exact (hv_int.mono_set (Metric.ball_subset_ball (le_of_lt hrs))).congr_fun
      (fun x hx => hpow_eq_on x hx) measurableSet_ball
  have hleft_integral_le :
      ∫ x in Metric.ball (0 : E) r, |(u x)⁻¹| ^ (moserChi d * p) ∂volume ≤
        ∫ x in Ω, |v x| ^ qexp ∂volume := by
    calc ∫ x in Metric.ball (0 : E) r, |(u x)⁻¹| ^ (moserChi d * p) ∂volume
        = ∫ x in Metric.ball (0 : E) r, |v x| ^ qexp ∂volume := by
          symm; exact setIntegral_congr_fun measurableSet_ball (fun x hx => hpow_eq_on x hx)
      _ ≤ ∫ x in Ω, |v x| ^ qexp ∂volume := by
          exact setIntegral_mono_set hv_int (ae_of_all _ (by intro x; positivity))
            (ae_of_all _ (Metric.ball_subset_ball (le_of_lt hrs)))
  -- Constant consolidation
  have hCη_sq : Cη ^ 2 = 4 * (Mst : ℝ) ^ 2 / (s - r) ^ 2 := by
    dsimp [Cη]; field_simp [hsr_pos.ne']; ring
  have hconst_bound :
      (C_gns d 2) ^ 2 * (2 * Cη ^ 2) ≤ C_MoserAnchor d / (s - r) ^ 2 := by
    rw [hCη_sq]
    have hrewrite :
        (C_gns d 2) ^ 2 * (2 * (4 * (Mst : ℝ) ^ 2 / (s - r) ^ 2)) =
          ((C_gns d 2) ^ 2 * (8 * (Mst : ℝ) ^ 2)) / (s - r) ^ 2 := by
      field_simp [hsr_pos.ne']; ring
    rw [hrewrite]
    exact div_le_div_of_nonneg_right
      (by simpa [mul_assoc, mul_left_comm, mul_comm] using
        cutoff_sobolev_anchor_le_C_MoserAnchor (d := d))
      (by positivity)
  have hcoeff_nonneg : 0 ≤ A.1.Λ * (p / (1 + p)) ^ 2 + 1 := by
    nlinarith [A.1.Λ_nonneg, sq_nonneg (p / (1 + p))]
  have hIp_nonneg : 0 ≤ ∫ x in Ω, |(u x)⁻¹| ^ p ∂volume := integral_nonneg fun x => by positivity
  -- Main chain: Lᵖˣ(Bᵣ) ≤ Sobolev(v) ≤ C_gns · gradient ≤ C_gns · energy^{1/2} ≤ constant · Lᵖ(Bₛ)
  have hqexp_mul : (1 / qexp : ℝ) * (2 / p) = 1 / (moserChi d * p) := by
    dsimp [qexp, moserChi]; field_simp
  have hhalf_mul : (1 / (2 : ℝ)) * (2 / p) = 1 / p := by field_simp [hp.ne']
  have hclose_exp_nonneg : 0 ≤ 1 / (moserChi d * p) := by
    exact one_div_nonneg.mpr (mul_pos (moserChi_pos (d := d) hd) hp).le
  have hright_nonneg : 0 ≤ ∫ x in Ω, |v x| ^ qexp ∂volume := integral_nonneg fun x => by positivity
  have hinner_nonneg : 0 ≤ 2 * Cη ^ 2 * (A.1.Λ * (p / (1 + p)) ^ 2 + 1) *
      ∫ x in Ω, |(u x)⁻¹| ^ p ∂volume := by positivity
  have hgrad_rhs_nonneg : 0 ≤ (2 * Cη ^ 2 * (A.1.Λ * (p / (1 + p)) ^ 2 + 1) *
      ∫ x in Ω, |(u x)⁻¹| ^ p ∂volume) ^ (1 / (2 : ℝ)) :=
    Real.rpow_nonneg hinner_nonneg _
  have hmain_lp :
      (∫ x in Metric.ball (0 : E) r,
          |(u x)⁻¹| ^ (moserChi d * p) ∂volume) ^ (1 / (moserChi d * p)) ≤
        ((C_MoserAnchor d / (s - r) ^ 2) *
            (A.1.Λ * (p / (1 + p)) ^ 2 + 1) *
            ∫ x in Ω, |(u x)⁻¹| ^ p ∂volume) ^ (1 / p) := by
    calc (∫ x in Metric.ball (0 : E) r,
            |(u x)⁻¹| ^ (moserChi d * p) ∂volume) ^ (1 / (moserChi d * p))
        ≤ (∫ x in Ω, |v x| ^ qexp ∂volume) ^ (1 / (moserChi d * p)) := by
            exact Real.rpow_le_rpow (by positivity) hleft_integral_le hclose_exp_nonneg
      _ = ((∫ x in Ω, |v x| ^ qexp ∂volume) ^ (1 / qexp)) ^ (2 / p) := by
            rw [← Real.rpow_mul hright_nonneg, hqexp_mul]
      _ = (MeasureTheory.lpNorm v q μ) ^ (2 / p) := by rw [hv_lp_eq]
      _ ≤ ((C_gns d 2) * MeasureTheory.lpNorm (fun x => ‖hwv_real.weakGrad x‖) 2 μ) ^ (2 / p) := by
            exact Real.rpow_le_rpow MeasureTheory.lpNorm_nonneg hSob_real (by positivity)
      _ ≤ ((C_gns d 2) * (2 * Cη ^ 2 * (A.1.Λ * (p / (1 + p)) ^ 2 + 1) *
              ∫ x in Ω, |(u x)⁻¹| ^ p ∂volume) ^ (1 / (2 : ℝ))) ^ (2 / p) := by
            exact Real.rpow_le_rpow
              (mul_nonneg (C_gns_nonneg d 2) MeasureTheory.lpNorm_nonneg)
              (mul_le_mul_of_nonneg_left hgrad_bound (C_gns_nonneg d 2))
              (by positivity)
      _ = ((C_gns d 2) ^ 2 * (2 * Cη ^ 2) *
              (A.1.Λ * (p / (1 + p)) ^ 2 + 1) *
              ∫ x in Ω, |(u x)⁻¹| ^ p ∂volume) ^ (1 / p) := by
            calc ((C_gns d 2) * (2 * Cη ^ 2 * (A.1.Λ * (p / (1 + p)) ^ 2 + 1) *
                    ∫ x in Ω, |(u x)⁻¹| ^ p ∂volume) ^ (1 / (2 : ℝ))) ^ (2 / p)
                = (C_gns d 2) ^ (2 / p) *
                    ((2 * Cη ^ 2 * (A.1.Λ * (p / (1 + p)) ^ 2 + 1) *
                      ∫ x in Ω, |(u x)⁻¹| ^ p ∂volume) ^ (1 / (2 : ℝ))) ^ (2 / p) := by
                  rw [Real.mul_rpow (C_gns_nonneg d 2) hgrad_rhs_nonneg]
              _ = ((C_gns d 2) ^ 2) ^ (1 / p) *
                    (2 * Cη ^ 2 * (A.1.Λ * (p / (1 + p)) ^ 2 + 1) *
                      ∫ x in Ω, |(u x)⁻¹| ^ p ∂volume) ^ (1 / p) := by
                  congr 1
                  · rw [show (2 / p : ℝ) = 2 * (1 / p) by field_simp [hp.ne']]
                    simpa using (Real.rpow_mul (C_gns_nonneg d 2) 2 (1 / p))
                  · rw [← Real.rpow_mul hinner_nonneg, hhalf_mul]
              _ = ((C_gns d 2) ^ 2 * (2 * Cη ^ 2) *
                    (A.1.Λ * (p / (1 + p)) ^ 2 + 1) *
                    ∫ x in Ω, |(u x)⁻¹| ^ p ∂volume) ^ (1 / p) := by
                  rw [← Real.mul_rpow (sq_nonneg (C_gns d 2)) hinner_nonneg]; congr 1; ring
      _ ≤ ((C_MoserAnchor d / (s - r) ^ 2) *
              (A.1.Λ * (p / (1 + p)) ^ 2 + 1) *
              ∫ x in Ω, |(u x)⁻¹| ^ p ∂volume) ^ (1 / p) := by
            apply Real.rpow_le_rpow (by positivity)
            · exact mul_le_mul_of_nonneg_right
                (mul_le_mul_of_nonneg_right hconst_bound
                  (by nlinarith [A.1.Λ_nonneg, sq_nonneg (p / (1 + p))]))
                hIp_nonneg
            · positivity
  -- Separate the outer coefficient before applying `Real.mul_rpow`.
  refine ⟨hleft_int, ?_⟩
  have hsplit_nonneg : 0 ≤ (C_MoserAnchor d / (s - r) ^ 2) *
      (A.1.Λ * (p / (1 + p)) ^ 2 + 1) := by
    exact mul_nonneg (div_nonneg (le_trans (by norm_num : (0 : ℝ) ≤ 1)
      (one_le_C_MoserAnchor (d := d))) (by positivity)) hcoeff_nonneg
  calc (∫ x in Metric.ball (0 : E) r,
          |(u x)⁻¹| ^ (moserChi d * p) ∂volume) ^ (1 / (moserChi d * p))
      ≤ ((C_MoserAnchor d / (s - r) ^ 2) *
          (A.1.Λ * (p / (1 + p)) ^ 2 + 1) *
          ∫ x in Ω, |(u x)⁻¹| ^ p ∂volume) ^ (1 / p) := hmain_lp
    _ = ((C_MoserAnchor d / (s - r) ^ 2) *
            (A.1.Λ * (p / (1 + p)) ^ 2 + 1)) ^ (1 / p) *
          (∫ x in Ω, |(u x)⁻¹| ^ p ∂volume) ^ (1 / p) := by
          rw [Real.mul_rpow hsplit_nonneg hIp_nonneg]

\end{lstlisting}
\begin{proof}
Apply Theorem~\ref{thm:weighted-cacc} with $\Psi(t) := t^{-p}$ on $(0,\infty)$
(suitably truncated), so that $\Psi'(t) = -p\,t^{-p-1} \leq 0$ and
$\Psi(t)^{2}/(-\Psi'(t)) = t^{-p-1}/p$. Choose a smooth cutoff $\eta$ with
$\eta = 1$ on $\Bz{r}$, $\supp \eta \subseteq \Bz{R}$ where $R := (r+s)/2$,
$0 \leq \eta \leq 1$, $\|\nabla \eta\|_{L^\infty} \leq C_\eta := 2\Mst/(s-r)$.
The weighted Caccioppoli inequality applied to $\varphi := \eta$ then yields
\begin{equation*}
  \int_{\Bz{s}} \eta^{2}\,p\,u^{-p-1}\,\langle A\nabla u,\nabla u\rangle\,dx
  \;\leq\;
  \frac{4\Lambda}{p}
  \int_{\Bz{s}} \|\nabla \eta\|^{2}\,u^{-p-1}\,dx.
\end{equation*}
Combined with the chain-rule identity
$\nabla(\eta\,u^{-p/2}) = (\nabla \eta) u^{-p/2} - (p/2) \eta\,u^{-p/2-1}\nabla u$
and the elementary $|a-b|^{2} \leq 2(a^{2}+b^{2})$ this gives, for the witness
$v := \mathrm{PCI}_d(\eta,u,p)$, the energy bound
\[
  \int_{\Bz{s}} \|\nabla v\|^{2}\,dx
  \;\leq\;
  2\,C_\eta^{2}\,\Big(\Lambda\Big(\frac{p}{1+p}\Big)^{2}+1\Big)
  \int_{\Bz{s}} |u^{-1}|^{p}\,dx;
\]
this is the content of \leanname{superPowerCutoff\_energy\_bound}. Now apply
the Sobolev inequality \leanname{sobolev\_prepare\_on\_ball} to $v$ on $\Bz{s}$
and use $\eta = 1$ on $\Bz{r}$, exactly as in the proof of Theorem~\ref{thm:moser-preMoser}.
The constant $8\Mst^{2}\,C_{\mathrm{gns}}^{2} \leq C_{\mathrm{MA}}$ absorbs the
gradient bound, yielding the displayed reverse Holder inequality.
\end{proof}

\paragraph{Forward (low-power) case.}

\begin{definition}[Forward low-power cutoff, \leanname{superPowerCutoffFwd}]
For smooth $\eta$, $u > 0$, and $p \in (0,1)$,
\[
  \mathrm{PCF}_d(\eta,u,p)(x) := \eta(x)\,|u(x)|^{p/2}.
\]
\end{definition}

\begin{lstlisting}[style=Matlab-editor,basicstyle=\mlttfamily,escapechar=",mlshowsectionrules=true,mathescape=true,morecomment={[s]{\%\{}{\%\}}},language=lean,numbers=none]
/-- Forward low-power energy bound for `η · u^{p/2}` where `p ∈ (0,1)`.

This is the analogue of `superPowerCutoff_energy_bound` for positive powers
of `u` (not `u⁻¹`). The Caccioppoli inequality gives `(1-p)` in the denominator
instead of `(1+p)`.

The proof is the same test-function computation as the inverse case but with
`p ↦ -p` substitution. When `p ∈ (-1,0)` in the general formula, we get
`u^{-(-p)/2} = u^{p/2}` and `1 + (-p) = 1 - p`. -/
noncomputable def superPowerCutoffFwd
    (η u : E → ℝ) (p : ℝ) : E → ℝ :=
  fun x => η x * |u x| ^ (p / 2)

\end{lstlisting}
\begin{theorem}[Forward low-power one-step gain, \leanname{supersolution\_preMoser\_forward}]
\label{thm:preMoser-fwd}
Let $d>2$, $A$ as above, $u>0$ a supersolution, $0 < p < 1$, and
$0 < r < s \leq 1$ with $\int_{\Bz{s}} |u|^{p}\,dx < \infty$. Then
$|u|^{\chi p}$ is integrable on $\Bz{r}$ and
\begin{multline*}
  \Big(\int_{\Bz{r}} |u|^{\chi p}\,dx\Big)^{1/(\chi p)}
  \\ \leq\;
  \Big(\frac{C_{\mathrm{MA}}(d)}{(s-r)^{2}}\,
    \Big(\Lambda\Big(\frac{p}{1-p}\Big)^{2}+1\Big)\Big)^{1/p}
  \Big(\int_{\Bz{s}} |u|^{p}\,dx\Big)^{1/p}.
\end{multline*}
\end{theorem}

\begin{lstlisting}[style=Matlab-editor,basicstyle=\mlttfamily,escapechar=",mlshowsectionrules=true,mathescape=true,morecomment={[s]{\%\{}{\%\}}},language=lean,numbers=none]
/-- One-step Sobolev gain for positive supersolutions, forward low-power case.

For `p ∈ (0,1)` and `0 < r < s ≤ 1`, if `u > 0` is a supersolution on `B₁`, then
```
‖u‖_{L^{pχ}(Bᵣ)} ≤ C^{1/p} (Λp²/(1-p)² + 1)^{1/p} (s-r)^{-2/p} ‖u‖_{Lᵖ(Bₛ)}
```
-/
theorem supersolution_preMoser_forward
    (hd : 2 < (d : ℝ))
    (A : NormalizedEllipticCoeff d (Metric.ball (0 : E) 1))
    {u : E → ℝ} {p r s : ℝ}
    (hp : 0 < p) (hp1 : p < 1)
    (hr : 0 < r) (hrs : r < s) (hs1 : s ≤ 1)
    (hu_pos : ∀ x ∈ Metric.ball (0 : E) 1, 0 < u x)
    (hsuper : IsSupersolution A.1 u)
    (hpInt :
      IntegrableOn (fun x => |u x| ^ p)
        (Metric.ball (0 : E) s) volume) :
    IntegrableOn (fun x => |u x| ^ (moserChi d * p))
        (Metric.ball (0 : E) r) volume ∧
      (∫ x in Metric.ball (0 : E) r,
          |u x| ^ (moserChi d * p) ∂volume) ^
          (1 / (moserChi d * p)) ≤
        ((C_MoserAnchor d / (s - r) ^ 2) *
            (A.1.Λ * (p / (1 - p)) ^ 2 + 1)) ^ (1 / p) *
          (∫ x in Metric.ball (0 : E) s, |u x| ^ p ∂volume) ^ (1 / p) := by
  -- Same structure as supersolution_preMoser_inverse but with forward powers.
  let Ω : Set E := Metric.ball (0 : E) s
  let R : ℝ := (r + s) / 2
  have hs_pos : 0 < s := lt_trans hr hrs
  have hsr_pos : 0 < s - r := by linarith
  have hR : r < R := by dsimp [R]; linarith
  have hR_lt_s : R < s := by dsimp [R]; linarith
  let ηCut : Cutoff (x₀ := (0 : E)) r R := Classical.choose
    (Cutoff.exists (d := d) (0 : E) (r := r) (R := R) hr hR)
  let η : E → ℝ := ηCut.toFun
  let Cη : ℝ := 2 * (Mst : ℝ) * (s - r)⁻¹
  have hη : ContDiff ℝ (⊤ : ℕ∞) η := ηCut.smooth
  have hη_bound : ∀ x, |η x| ≤ 1 := by
    intro x; rw [abs_of_nonneg (ηCut.nonneg x)]; exact ηCut.le_one x
  have hη_eq_one : ∀ x ∈ Metric.ball (0 : E) r, η x = 1 := by
    intro x hx; exact ηCut.eq_one x hx
  have hη_grad_bound : ∀ x, ‖fderiv ℝ η x‖ ≤ Cη := by
    intro x
    have hmid : (((r + s) / 2 : ℝ) - r) = (s - r) / 2 := by ring
    have h_inv : ((((r + s) / 2 : ℝ) - r)⁻¹) = 2 * (s - r)⁻¹ := by
      rw [hmid]; field_simp [hsr_pos.ne']
    calc ‖fderiv ℝ η x‖ ≤ ↑Mst * (((r + s) / 2 : ℝ) - r)⁻¹ := ηCut.grad_bound x
      _ = Cη := by rw [h_inv]; ring
  have hη_sub_ball : tsupport η ⊆ Ω := by
    intro x hx; exact (Metric.closedBall_subset_ball hR_lt_s) (ηCut.support_subset hx)
  obtain ⟨hwv, hvW01, henergy⟩ :=
    superPowerCutoffFwd_energy_bound hd A hp hp1 hs_pos hs1 hu_pos hsuper hpInt
      hη hη_bound hη_grad_bound hη_sub_ball
  let v : E → ℝ := superPowerCutoffFwd (d := d) η u p
  have hv_support : tsupport v ⊆ Ω :=
    superPowerCutoffFwd_tsupport_subset (d := d) hη_sub_ball
  obtain ⟨hwv_real, hSob⟩ :=
    sobolev_prepare_on_ball (d := d) (v := v) hd hs_pos hvW01 hv_support
  let μ : Measure E := volume.restrict Ω
  let qexp : ℝ := ((d : ℝ) * 2 / ((d : ℝ) - 2))
  let q := ENNReal.ofReal qexp
  have hq_pos : 0 < qexp := by
    dsimp [qexp]; exact div_pos (mul_pos (by linarith) (by norm_num)) (by linarith)
  let hwvReal : MemW1pWitness (ENNReal.ofReal (2 : ℝ)) v Ω :=
    { memLp := by simpa [Ω] using hwv.memLp
      weakGrad := hwv.weakGrad
      weakGrad_component_memLp := by simpa [Ω] using hwv.weakGrad_component_memLp
      isWeakGrad := by simpa [Ω] using hwv.isWeakGrad }
  have hae_grad := by simpa [μ, Ω] using
    MemW1pWitness.ae_eq_p (d := d) Metric.isOpen_ball (p := (2 : ℝ)) (by norm_num) hwv_real hwvReal
  have hgrad_sq_eq :
      ∫ x in Ω, ‖hwv_real.weakGrad x‖ ^ 2 ∂volume =
        ∫ x in Ω, ‖hwvReal.weakGrad x‖ ^ 2 ∂volume :=
    integral_congr_ae (hae_grad.fun_comp (fun z => ‖z‖ ^ 2))
  have hv_memLp_q : MemLp v q μ := by
    refine ⟨hwv_real.memLp.aestronglyMeasurable, ?_⟩
    exact lt_of_le_of_lt (by simpa [q, μ, Ω] using hSob)
      (ENNReal.mul_lt_top_iff.mpr (Or.inl ⟨ENNReal.ofReal_lt_top,
        lt_top_iff_ne_top.mpr (by simpa using hwv_real.weakGrad_norm_memLp.eLpNorm_ne_top)⟩))
  have hSob' : eLpNorm v q μ ≤
      ENNReal.ofReal (C_gns d 2) *
        eLpNorm (fun x => ‖hwv_real.weakGrad x‖) 2 μ := by simpa [q, μ, Ω] using hSob
  have hgrad_ne_top : eLpNorm (fun x => ‖hwv_real.weakGrad x‖) 2 μ ≠ ⊤ := by
    simpa using hwv_real.weakGrad_norm_memLp.eLpNorm_ne_top
  have hSob_real :
      MeasureTheory.lpNorm v q μ ≤
        (C_gns d 2) * MeasureTheory.lpNorm (fun x => ‖hwv_real.weakGrad x‖) 2 μ := by
    have hSob_toReal :=
      (ENNReal.toReal_le_toReal hv_memLp_q.eLpNorm_ne_top
        (ENNReal.mul_ne_top ENNReal.ofReal_ne_top hgrad_ne_top)).2 hSob'
    rwa [MeasureTheory.toReal_eLpNorm hv_memLp_q.aestronglyMeasurable,
      ENNReal.toReal_mul, MeasureTheory.toReal_eLpNorm
        hwv_real.weakGrad_norm_memLp.aestronglyMeasurable,
      ENNReal.toReal_ofReal (C_gns_nonneg d 2)] at hSob_toReal
  have hgrad_lp :
      MeasureTheory.lpNorm (fun x => ‖hwv_real.weakGrad x‖) 2 μ =
        (∫ x in Ω, ‖hwv_real.weakGrad x‖ ^ 2 ∂volume) ^ (1 / (2 : ℝ)) := by
    simpa [μ, Ω] using
      MeasureTheory.lpNorm_eq_integral_norm_rpow_toReal (μ := μ)
        (f := fun x => ‖hwv_real.weakGrad x‖)
        two_ne_zero ENNReal.ofNat_ne_top hwv_real.weakGrad_norm_memLp.aestronglyMeasurable
  have hgrad_bound :
      MeasureTheory.lpNorm (fun x => ‖hwv_real.weakGrad x‖) 2 μ ≤
        (2 * Cη ^ 2 * (A.1.Λ * (p / (1 - p)) ^ 2 + 1) *
          ∫ x in Ω, |u x| ^ p ∂volume) ^ (1 / (2 : ℝ)) := by
    rw [hgrad_lp, hgrad_sq_eq]
    exact Real.rpow_le_rpow (by exact integral_nonneg fun x => by positivity)
      (by simpa [hwvReal] using henergy) (by norm_num)
  have hq_ne_zero : q ≠ 0 := by dsimp [q]; rw [ENNReal.ofReal_eq_zero]; linarith
  have hv_lp_eq :
      MeasureTheory.lpNorm v q μ =
        (∫ x in Ω, |v x| ^ qexp ∂volume) ^ (1 / qexp) := by
    rw [MeasureTheory.lpNorm_eq_integral_norm_rpow_toReal (μ := μ) (f := v)]
    · simp [μ, Ω, q, qexp, ENNReal.toReal_ofReal hq_pos.le, Real.norm_eq_abs]
    · exact hq_ne_zero
    · exact ENNReal.ofReal_ne_top
    · exact hv_memLp_q.aestronglyMeasurable
  have hv_int : IntegrableOn (fun x => |v x| ^ qexp) Ω volume := by
    simpa [μ, Ω, q, qexp, ENNReal.toReal_ofReal hq_pos.le, Real.norm_eq_abs] using
      hv_memLp_q.integrable_norm_rpow hq_ne_zero ENNReal.ofReal_ne_top
  have hpow_eq_on : ∀ x ∈ Metric.ball (0 : E) r,
      |v x| ^ qexp = |u x| ^ (moserChi d * p) := by
    intro x hx
    have hηx : η x = 1 := hη_eq_one x hx
    have hqexp_eq : qexp = 2 * moserChi d := by dsimp [qexp, moserChi]; field_simp
    have hvx : v x = |u x| ^ (p / 2) := by simp [v, superPowerCutoffFwd, hηx]
    calc |v x| ^ qexp
        = (|u x| ^ (p / 2)) ^ qexp := by
          rw [hvx, abs_of_nonneg (Real.rpow_nonneg (abs_nonneg _) _)]
      _ = |u x| ^ (moserChi d * p) := by
          rw [hqexp_eq, ← Real.rpow_mul (abs_nonneg _)]; ring_nf
  have hleft_int : IntegrableOn (fun x => |u x| ^ (moserChi d * p))
      (Metric.ball (0 : E) r) volume :=
    (hv_int.mono_set (Metric.ball_subset_ball (le_of_lt hrs))).congr_fun
      (fun x hx => hpow_eq_on x hx) measurableSet_ball
  have hleft_integral_le :
      ∫ x in Metric.ball (0 : E) r, |u x| ^ (moserChi d * p) ∂volume ≤
        ∫ x in Ω, |v x| ^ qexp ∂volume := by
    calc ∫ x in Metric.ball (0 : E) r, |u x| ^ (moserChi d * p) ∂volume
        = ∫ x in Metric.ball (0 : E) r, |v x| ^ qexp ∂volume := by
          symm; exact setIntegral_congr_fun measurableSet_ball (fun x hx => hpow_eq_on x hx)
      _ ≤ ∫ x in Ω, |v x| ^ qexp ∂volume :=
          setIntegral_mono_set hv_int (ae_of_all _ (by intro x; positivity))
            (ae_of_all _ (Metric.ball_subset_ball (le_of_lt hrs)))
  have hCη_sq : Cη ^ 2 = 4 * (Mst : ℝ) ^ 2 / (s - r) ^ 2 := by
    dsimp [Cη]; field_simp [hsr_pos.ne']; ring
  have hconst_bound :
      (C_gns d 2) ^ 2 * (2 * Cη ^ 2) ≤ C_MoserAnchor d / (s - r) ^ 2 := by
    rw [hCη_sq]
    have hrewrite :
        (C_gns d 2) ^ 2 * (2 * (4 * (Mst : ℝ) ^ 2 / (s - r) ^ 2)) =
          ((C_gns d 2) ^ 2 * (8 * (Mst : ℝ) ^ 2)) / (s - r) ^ 2 := by
      field_simp [hsr_pos.ne']; ring
    rw [hrewrite]
    exact div_le_div_of_nonneg_right
      (by simpa [mul_assoc, mul_left_comm, mul_comm] using
        cutoff_sobolev_anchor_le_C_MoserAnchor (d := d))
      (by positivity)
  have hcoeff_nonneg : 0 ≤ A.1.Λ * (p / (1 - p)) ^ 2 + 1 := by
    nlinarith [A.1.Λ_nonneg, sq_nonneg (p / (1 - p))]
  have hIp_nonneg : 0 ≤ ∫ x in Ω, |u x| ^ p ∂volume := integral_nonneg fun x => by positivity
  have hqexp_mul : (1 / qexp : ℝ) * (2 / p) = 1 / (moserChi d * p) := by
    dsimp [qexp, moserChi]; field_simp
  have hhalf_mul : (1 / (2 : ℝ)) * (2 / p) = 1 / p := by field_simp [hp.ne']
  have hclose_exp_nonneg : 0 ≤ 1 / (moserChi d * p) :=
    one_div_nonneg.mpr (mul_pos (moserChi_pos (d := d) hd) hp).le
  have hright_nonneg : 0 ≤ ∫ x in Ω, |v x| ^ qexp ∂volume := integral_nonneg fun x => by positivity
  have hinner_nonneg : 0 ≤ 2 * Cη ^ 2 * (A.1.Λ * (p / (1 - p)) ^ 2 + 1) *
      ∫ x in Ω, |u x| ^ p ∂volume := by positivity
  have hgrad_rhs_nonneg : 0 ≤ (2 * Cη ^ 2 * (A.1.Λ * (p / (1 - p)) ^ 2 + 1) *
      ∫ x in Ω, |u x| ^ p ∂volume) ^ (1 / (2 : ℝ)) := Real.rpow_nonneg hinner_nonneg _
  have hmain_lp :
      (∫ x in Metric.ball (0 : E) r,
          |u x| ^ (moserChi d * p) ∂volume) ^ (1 / (moserChi d * p)) ≤
        ((C_MoserAnchor d / (s - r) ^ 2) *
            (A.1.Λ * (p / (1 - p)) ^ 2 + 1) *
            ∫ x in Ω, |u x| ^ p ∂volume) ^ (1 / p) := by
    calc (∫ x in Metric.ball (0 : E) r,
            |u x| ^ (moserChi d * p) ∂volume) ^ (1 / (moserChi d * p))
        ≤ (∫ x in Ω, |v x| ^ qexp ∂volume) ^ (1 / (moserChi d * p)) :=
            Real.rpow_le_rpow (by positivity) hleft_integral_le hclose_exp_nonneg
      _ = ((∫ x in Ω, |v x| ^ qexp ∂volume) ^ (1 / qexp)) ^ (2 / p) := by
            rw [← Real.rpow_mul hright_nonneg, hqexp_mul]
      _ = (MeasureTheory.lpNorm v q μ) ^ (2 / p) := by rw [hv_lp_eq]
      _ ≤ ((C_gns d 2) * MeasureTheory.lpNorm (fun x => ‖hwv_real.weakGrad x‖) 2 μ) ^ (2 / p) :=
            Real.rpow_le_rpow MeasureTheory.lpNorm_nonneg hSob_real (by positivity)
      _ ≤ ((C_gns d 2) * (2 * Cη ^ 2 * (A.1.Λ * (p / (1 - p)) ^ 2 + 1) *
              ∫ x in Ω, |u x| ^ p ∂volume) ^ (1 / (2 : ℝ))) ^ (2 / p) :=
            Real.rpow_le_rpow
              (mul_nonneg (C_gns_nonneg d 2) MeasureTheory.lpNorm_nonneg)
              (mul_le_mul_of_nonneg_left hgrad_bound (C_gns_nonneg d 2))
              (by positivity)
      _ = ((C_gns d 2) ^ 2 * (2 * Cη ^ 2) *
              (A.1.Λ * (p / (1 - p)) ^ 2 + 1) *
              ∫ x in Ω, |u x| ^ p ∂volume) ^ (1 / p) := by
            calc ((C_gns d 2) * (2 * Cη ^ 2 * (A.1.Λ * (p / (1 - p)) ^ 2 + 1) *
                    ∫ x in Ω, |u x| ^ p ∂volume) ^ (1 / (2 : ℝ))) ^ (2 / p)
                = (C_gns d 2) ^ (2 / p) *
                    ((2 * Cη ^ 2 * (A.1.Λ * (p / (1 - p)) ^ 2 + 1) *
                      ∫ x in Ω, |u x| ^ p ∂volume) ^ (1 / (2 : ℝ))) ^ (2 / p) := by
                  rw [Real.mul_rpow (C_gns_nonneg d 2) hgrad_rhs_nonneg]
              _ = ((C_gns d 2) ^ 2) ^ (1 / p) *
                    (2 * Cη ^ 2 * (A.1.Λ * (p / (1 - p)) ^ 2 + 1) *
                      ∫ x in Ω, |u x| ^ p ∂volume) ^ (1 / p) := by
                  congr 1
                  · rw [show (2 / p : ℝ) = 2 * (1 / p) by field_simp [hp.ne']]
                    simpa using (Real.rpow_mul (C_gns_nonneg d 2) 2 (1 / p))
                  · rw [← Real.rpow_mul hinner_nonneg, hhalf_mul]
              _ = ((C_gns d 2) ^ 2 * (2 * Cη ^ 2) *
                    (A.1.Λ * (p / (1 - p)) ^ 2 + 1) *
                    ∫ x in Ω, |u x| ^ p ∂volume) ^ (1 / p) := by
                  rw [← Real.mul_rpow (sq_nonneg (C_gns d 2)) hinner_nonneg]; congr 1; ring
      _ ≤ ((C_MoserAnchor d / (s - r) ^ 2) *
              (A.1.Λ * (p / (1 - p)) ^ 2 + 1) *
              ∫ x in Ω, |u x| ^ p ∂volume) ^ (1 / p) := by
            apply Real.rpow_le_rpow (by positivity)
            · exact mul_le_mul_of_nonneg_right
                (mul_le_mul_of_nonneg_right hconst_bound
                  (by nlinarith [A.1.Λ_nonneg, sq_nonneg (p / (1 - p))]))
                hIp_nonneg
            · positivity
  refine ⟨hleft_int, ?_⟩
  have hsplit_nonneg : 0 ≤ (C_MoserAnchor d / (s - r) ^ 2) *
      (A.1.Λ * (p / (1 - p)) ^ 2 + 1) :=
    mul_nonneg (div_nonneg (le_trans (by norm_num : (0 : ℝ) ≤ 1)
      (one_le_C_MoserAnchor (d := d))) (by positivity)) hcoeff_nonneg
  calc (∫ x in Metric.ball (0 : E) r,
          |u x| ^ (moserChi d * p) ∂volume) ^ (1 / (moserChi d * p))
      ≤ ((C_MoserAnchor d / (s - r) ^ 2) *
          (A.1.Λ * (p / (1 - p)) ^ 2 + 1) *
          ∫ x in Ω, |u x| ^ p ∂volume) ^ (1 / p) := hmain_lp
    _ = ((C_MoserAnchor d / (s - r) ^ 2) *
            (A.1.Λ * (p / (1 - p)) ^ 2 + 1)) ^ (1 / p) *
          (∫ x in Ω, |u x| ^ p ∂volume) ^ (1 / p) := by
          rw [Real.mul_rpow hsplit_nonneg hIp_nonneg]

\end{lstlisting}
\begin{proof}
This time use $\Psi(t) := -t^{p}$ on $(0,\infty)$, so $\Psi'(t) = -p\,t^{p-1}
\leq 0$ and $\Psi(t)^{2}/(-\Psi'(t)) = t^{p+1}/p$. (Note that for $p \in (0,1)$
the function $t \mapsto t^{p}$ is concave with $\Psi'(0)=+\infty$, but
truncation issues are absorbed in \leanname{superPowerCutoffFwd\_energy\_bound}.)
The weighted Caccioppoli inequality with $\varphi := \eta$ supplies the
analogous energy estimate
\[
  \int_{\Bz{s}} \|\nabla v\|^{2}\,dx
  \;\leq\;
  2 C_\eta^{2}\,\Big(\Lambda\Big(\frac{p}{1-p}\Big)^{2}+1\Big)
  \int_{\Bz{s}} |u|^{p}\,dx
\]
for $v := \mathrm{PCF}_d(\eta,u,p)$. Sobolev embedding then concludes as in the
inverse case.
\end{proof}

\subsubsection{Stage one (inverse iteration)}

We collect the inverse-power constants used in the iteration.

\begin{definition}[Inverse iteration quantities, \leanname{superIterIntegralInv}, \leanname{superIterNormInv}, \leanname{superStepConstInv}]
For $u>0$, $p_0 > 0$, $n \in \N$, set
\begin{align*}
  I_n^{\mathrm{inv}}(u) &:= \int_{\Bz{r_n}} |u^{-1}|^{p_n}\,dx, \\
  N_n^{\mathrm{inv}}(u) &:= \big(I_n^{\mathrm{inv}}(u)\big)^{1/p_n}, \\
  \mathrm{S}_n^{\mathrm{inv}}(A,p_0) &:=
   \Big(\frac{C_{\mathrm{MA}}(d)}{(r_n-r_{n+1})^{2}}\,
     \Big(\Lambda\Big(\frac{p_n}{1+p_n}\Big)^{2}+1\Big)\Big)^{1/p_n}.
\end{align*}
The constant
\begin{equation*}
  C_{\mathrm{wH}}^{0}(d) :=
  \begin{cases}
    C_{\mathrm{Moser}}(d)\,(\chi^{2})^{\sum_{n\geq 0} n q^{n}}, & d > 2,\\
    C_{\mathrm{Moser}}(d), & \text{otherwise},
  \end{cases}
\end{equation*}
(\leanname{C\_weakHarnack0}) absorbs the extra geometric factor coming from the
$(\chi^{2})^{i}$ term in the inverse step bound.
\end{definition}

\begin{lstlisting}[style=Matlab-editor,basicstyle=\mlttfamily,escapechar=",mlshowsectionrules=true,mathescape=true,morecomment={[s]{\%\{}{\%\}}},language=lean,numbers=none]
/-! ### Inverse-Power Iteration

Iterate `supersolution_preMoser_inverse` along the geometric sequences
`pₙ = p₀ χⁿ` and `rₙ = (1 + 2^{-n})/2`, then estimate the product and
pass to `L^∞`.
-/

/-- Auxiliary integrals for the inverse-power Moser iteration. -/
noncomputable def superIterIntegralInv
    {u : E → ℝ} (p₀ : ℝ) (n : ℕ) : ℝ :=
  ∫ x in Metric.ball (0 : E) (moserRadius n),
    |(u x)⁻¹| ^ moserExponentSeq d p₀ n ∂volume

/-- Auxiliary `L^{pₙ}` norm for the inverse-power Moser iteration. -/
noncomputable def superIterNormInv
    {u : E → ℝ} (p₀ : ℝ) (n : ℕ) : ℝ :=
  (superIterIntegralInv (d := d) (u := u) p₀ n) ^
    (1 / moserExponentSeq d p₀ n)

/-- The step constant in the supersolution inverse-power iteration at stage `n`.
This is `(C_MoserAnchor d / gap² · (Λ (pₙ/(1+pₙ))² + 1))^{1/pₙ}`. -/
noncomputable def superStepConstInv
    (A : NormalizedEllipticCoeff d (Metric.ball (0 : E) 1))
    (p₀ : ℝ) (n : ℕ) : ℝ :=
  ((C_MoserAnchor d / (moserRadius n - moserRadius (n + 1)) ^ 2) *
    (A.1.Λ * (moserExponentSeq d p₀ n /
      (1 + moserExponentSeq d p₀ n)) ^ 2 + 1)) ^
    (1 / moserExponentSeq d p₀ n)

/-- Dimension-only constant for the stage-one weak Harnack estimates.

Compared to `C_Moser`, this absorbs the extra `χ^(2i)` growth that appears in
the supersolution geometric product. -/
noncomputable def C_weakHarnack0 (d : ℕ) [NeZero d] : ℝ :=
  if _hd : 2 < (d : ℝ) then
    C_Moser d * (moserChi d ^ 2) ^ (∑' n : ℕ, (n : ℝ) * moserDecayRatio d ^ n)
  else
    C_Moser d

\end{lstlisting}
\begin{lemma}[Step bound, \leanname{superStepConstInv\_le}]
\label{lem:stepInv-le}
For $d>2$ and $p_0 > 0$,
\[
  \mathrm{S}_n^{\mathrm{inv}}(A,p_0)
  \;\leq\;
  \Big(16\,C_{\mathrm{MA}}(d)\,(\Lambda p_0^{2}+1)\,(4\chi^{2})^{n}\Big)^{1/p_n}.
\]
\end{lemma}

\begin{lstlisting}[style=Matlab-editor,basicstyle=\mlttfamily,escapechar=",mlshowsectionrules=true,mathescape=true,morecomment={[s]{\%\{}{\%\}}},language=lean,numbers=none]
/-- Each inverse-power step constant is bounded by a simpler expression.

The key simplification: `pₙ/(1+pₙ) ≤ pₙ` (since pₙ > 0), and
`gap_n = 2^{-(n+2)}`, so `1/gap_n² = 4^{n+2} = 16 · 4^n`. -/
theorem superStepConstInv_le
    (hd : 2 < (d : ℝ))
    (A : NormalizedEllipticCoeff d (Metric.ball (0 : E) 1))
    {p₀ : ℝ} (hp₀ : 0 < p₀) (i : ℕ) :
    superStepConstInv (d := d) A p₀ i ≤
      (16 * C_MoserAnchor d * (A.1.Λ * p₀ ^ 2 + 1) *
        (4 * moserChi d ^ 2) ^ i) ^
        (1 / moserExponentSeq d p₀ i) := by
  unfold superStepConstInv
  apply Real.rpow_le_rpow
  · -- Nonnegativity of the base
    apply mul_nonneg
    · exact div_nonneg (le_trans (by norm_num : (0 : ℝ) ≤ 1) (one_le_C_MoserAnchor (d := d)))
        (sq_nonneg _)
    · nlinarith [A.1.Λ_nonneg, sq_nonneg (moserExponentSeq d p₀ i / (1 + moserExponentSeq d p₀ i))]
  · -- Main bound
    have hp_i := moserExponentSeq_pos (d := d) hd hp₀ i
    have hgap_eq : C_MoserAnchor d / (moserRadius i - moserRadius (i + 1)) ^ 2 =
        C_MoserAnchor d * (4 : ℝ) ^ (i + 2) := by
      rw [moserRadius_gap]
      have hsq : (((1 / 2 : ℝ) ^ (i + 2)) ^ 2) = (1 / 4 : ℝ) ^ (i + 2) := by
        rw [← pow_mul, show (i + 2) * 2 = 2 * (i + 2) by ring, pow_mul]; norm_num
      rw [hsq, div_eq_mul_inv,
        show (((1 / 4 : ℝ) ^ (i + 2))⁻¹) = (4 : ℝ) ^ (i + 2) by
          rw [show (1 / 4 : ℝ) = (4 : ℝ)⁻¹ by norm_num, inv_pow]; norm_num]
    have hpow4 : (4 : ℝ) ^ (i + 2) = 16 * 4 ^ i := by rw [pow_add]; ring
    -- pₙ/(1+pₙ) ≤ pₙ, so Λ(pₙ/(1+pₙ))² ≤ Λpₙ²
    have hratio_le : (moserExponentSeq d p₀ i / (1 + moserExponentSeq d p₀ i)) ^ 2 ≤
        (moserExponentSeq d p₀ i) ^ 2 := by
      have hdiv_le : moserExponentSeq d p₀ i / (1 + moserExponentSeq d p₀ i) ≤
          moserExponentSeq d p₀ i :=
        div_le_self hp_i.le (by linarith)
      have hdiv_nonneg : 0 ≤ moserExponentSeq d p₀ i / (1 + moserExponentSeq d p₀ i) :=
        div_nonneg hp_i.le (by linarith)
      nlinarith [sq_nonneg (moserExponentSeq d p₀ i - moserExponentSeq d p₀ i / (1 + moserExponentSeq d p₀ i))]
    -- pₙ² = p₀² · χ^{2i}
    have hseq_sq : (moserExponentSeq d p₀ i) ^ 2 = p₀ ^ 2 * (moserChi d) ^ (2 * i) := by
      rw [moserExponentSeq]; ring
    -- Combine: Λ(pₙ/(1+pₙ))² + 1 ≤ (Λp₀² + 1) · χ^{2i}
    have hcoeff_le :
        A.1.Λ * (moserExponentSeq d p₀ i / (1 + moserExponentSeq d p₀ i)) ^ 2 + 1 ≤
          (A.1.Λ * p₀ ^ 2 + 1) * (moserChi d ^ 2) ^ i := by
      have hchi_sq_ge_one : 1 ≤ (moserChi d ^ 2) ^ i :=
        one_le_pow₀ (by nlinarith [one_lt_moserChi (d := d) hd])
      have hΛratio :
          A.1.Λ * (moserExponentSeq d p₀ i / (1 + moserExponentSeq d p₀ i)) ^ 2 ≤
            A.1.Λ * p₀ ^ 2 * (moserChi d ^ 2) ^ i := by
        calc A.1.Λ * (moserExponentSeq d p₀ i / (1 + moserExponentSeq d p₀ i)) ^ 2
            ≤ A.1.Λ * (moserExponentSeq d p₀ i) ^ 2 :=
              mul_le_mul_of_nonneg_left hratio_le A.1.Λ_nonneg
          _ = A.1.Λ * (p₀ ^ 2 * (moserChi d) ^ (2 * i)) := by rw [hseq_sq]
          _ = A.1.Λ * p₀ ^ 2 * (moserChi d ^ 2) ^ i := by rw [pow_mul]; ring
      nlinarith
    rw [hgap_eq, hpow4]
    calc C_MoserAnchor d * (16 * 4 ^ i) *
            (A.1.Λ * (moserExponentSeq d p₀ i / (1 + moserExponentSeq d p₀ i)) ^ 2 + 1)
        ≤ C_MoserAnchor d * (16 * 4 ^ i) *
            ((A.1.Λ * p₀ ^ 2 + 1) * (moserChi d ^ 2) ^ i) := by
          gcongr
          exact mul_nonneg (le_trans (by norm_num : (0 : ℝ) ≤ 1) (one_le_C_MoserAnchor (d := d)))
            (by positivity)
      _ = 16 * C_MoserAnchor d * (A.1.Λ * p₀ ^ 2 + 1) * (4 * moserChi d ^ 2) ^ i := by
          rw [show (4 * moserChi d ^ 2) ^ i = 4 ^ i * (moserChi d ^ 2) ^ i by rw [mul_pow]]
          ring
  · exact div_nonneg (by norm_num) (moserExponentSeq_pos (d := d) hd hp₀ i).le

\end{lstlisting}
\begin{proof}
The factor $C_{\mathrm{MA}}/(r_n - r_{n+1})^{2} = 16\,C_{\mathrm{MA}}\,4^{n}$ as
before. Using $p_n/(1+p_n) \leq p_n$ and $p_n^{2} = p_0^{2}\chi^{2n}$,
\[
  \Lambda\Big(\frac{p_n}{1+p_n}\Big)^{2}+1
  \;\leq\; \Lambda p_n^{2}+1
  \;\leq\; (\Lambda p_0^{2}+1)\,(\chi^{2})^{n}.
\]
Multiplying gives the inner expression $16 C_{\mathrm{MA}}(\Lambda p_0^{2}+1)
(4\chi^{2})^{n}$; raising to $1/p_n$ yields the claim.
\end{proof}

\begin{theorem}[Inverse iteration, \leanname{supersolution\_iteration\_inverse}]
\label{thm:iter-inv}
For $d>2$, $u>0$ a supersolution with $\int_{\Bz{1}} |u^{-1}|^{p_0}\,dx < \infty$,
and every $n \in \N$, $|u^{-1}|^{p_n}$ is integrable on $\Bz{r_n}$ and
\[
  N_n^{\mathrm{inv}}(u)
  \;\leq\;
  \prod_{i=0}^{n-1}\mathrm{S}_i^{\mathrm{inv}}(A,p_0)\;\cdot\; N_0^{\mathrm{inv}}(u).
\]
\end{theorem}

\begin{lstlisting}[style=Matlab-editor,basicstyle=\mlttfamily,escapechar=",mlshowsectionrules=true,mathescape=true,morecomment={[s]{\%\{}{\%\}}},language=lean,numbers=none]
/-- Iteration of the inverse-power one-step bound by induction.

At each step, `supersolution_preMoser_inverse` provides the `Lᵖⁿ → Lᵖⁿ⁺¹` gain,
and we accumulate the product of step constants. -/
theorem supersolution_iteration_inverse
    (hd : 2 < (d : ℝ))
    (A : NormalizedEllipticCoeff d (Metric.ball (0 : E) 1))
    {u : E → ℝ} {p₀ : ℝ} (hp₀ : 0 < p₀)
    (hu_pos : ∀ x ∈ Metric.ball (0 : E) 1, 0 < u x)
    (hsuper : IsSupersolution A.1 u)
    (hpInt :
      IntegrableOn (fun x => |(u x)⁻¹| ^ p₀)
        (Metric.ball (0 : E) 1) volume) :
    ∀ n : ℕ,
      IntegrableOn (fun x => |(u x)⁻¹| ^ moserExponentSeq d p₀ n)
          (Metric.ball (0 : E) (moserRadius n)) volume ∧
        superIterNormInv (d := d) (u := u) p₀ n ≤
          (∏ i ∈ Finset.range n, superStepConstInv (d := d) A p₀ i) *
            superIterNormInv (d := d) (u := u) p₀ 0 := by
  intro n
  induction n with
  | zero =>
    constructor
    · -- Base case: moserRadius 0 = 1 and moserExponentSeq _ _ 0 = p₀
      rwa [moserExponentSeq_zero, moserRadius_zero]
    · -- Product over empty range is 1
      simp
  | succ n ihn =>
    obtain ⟨hInt_n, hbound_n⟩ := ihn
    let p_n := moserExponentSeq d p₀ n
    have hp_n : 0 < p_n := moserExponentSeq_pos hd hp₀ n
    -- Apply the one-step inverse-power bound
    have hpre :=
      supersolution_preMoser_inverse hd A (p := p_n)
        (r := moserRadius (n + 1)) (s := moserRadius n)
        hp_n (moserRadius_pos (n + 1)) (moserRadius_succ_lt n)
        (moserRadius_le_one n) hu_pos hsuper
        (by simpa [p_n] using hInt_n)
    obtain ⟨hInt_succ, hNorm_succ⟩ := hpre
    refine ⟨?_, ?_⟩
    · -- Integrability: convert moserChi * p_n to p_{n+1}
      have heq : moserChi d * p_n = moserExponentSeq d p₀ (n + 1) := by
        rw [moserExponentSeq_succ]
      rwa [heq] at hInt_succ
    · -- Bound: multiply step bound with inductive hypothesis
      have heq_norm : superIterNormInv (d := d) (u := u) p₀ (n + 1) =
          (∫ x in Metric.ball (0 : E) (moserRadius (n + 1)),
            |(u x)⁻¹| ^ (moserChi d * p_n) ∂volume) ^
              (1 / (moserChi d * p_n)) := by
        simp [superIterNormInv, superIterIntegralInv, p_n, moserExponentSeq_succ]
      have heq_step : superStepConstInv (d := d) A p₀ n =
          ((C_MoserAnchor d / (moserRadius n - moserRadius (n + 1)) ^ 2) *
            (A.1.Λ * (p_n / (1 + p_n)) ^ 2 + 1)) ^ (1 / p_n) := by
        simp [superStepConstInv, p_n]
      have hstep_nonneg : 0 ≤ superStepConstInv (d := d) A p₀ n := by
        rw [heq_step]
        apply Real.rpow_nonneg
        apply mul_nonneg
        · exact div_nonneg
            (le_trans (by norm_num : (0 : ℝ) ≤ 1) (one_le_C_MoserAnchor (d := d)))
            (sq_nonneg _)
        · nlinarith [A.1.Λ_nonneg, sq_nonneg (p_n / (1 + p_n))]
      -- The one-step bound gives:
      --   ‖u⁻¹‖_{n+1} ≤ step_n * ‖u⁻¹‖_n
      have hstep_bound :
          superIterNormInv (d := d) (u := u) p₀ (n + 1) ≤
            superStepConstInv (d := d) A p₀ n *
              superIterNormInv (d := d) (u := u) p₀ n := by
        rw [heq_norm, heq_step]
        convert hNorm_succ using 2
      -- Combine with IH
      calc superIterNormInv (d := d) (u := u) p₀ (n + 1)
          ≤ superStepConstInv (d := d) A p₀ n *
              superIterNormInv (d := d) (u := u) p₀ n := hstep_bound
        _ ≤ superStepConstInv (d := d) A p₀ n *
              ((∏ i ∈ Finset.range n, superStepConstInv (d := d) A p₀ i) *
                superIterNormInv (d := d) (u := u) p₀ 0) := by
            exact mul_le_mul_of_nonneg_left hbound_n hstep_nonneg
        _ = (∏ i ∈ Finset.range (n + 1), superStepConstInv (d := d) A p₀ i) *
              superIterNormInv (d := d) (u := u) p₀ 0 := by
            rw [Finset.prod_range_succ]
            ring

\end{lstlisting}
\begin{proof}
Induction on $n$. The base case $n = 0$ uses $r_0 = 1$
(\leanname{moserRadius\_zero}) and $p_0 = p_0$
(\leanname{moserExponentSeq\_zero}): the integrability hypothesis for
$|u^{-1}|^{p_0}$ on $\Bz{1}$ is exactly what is assumed, and the empty product
is $1$ so the iteration bound reads $N_0^{\mathrm{inv}}(u) \le N_0^{\mathrm{inv}}(u)$.

For the inductive step from $n$ to $n+1$, set $p_n := \mathrm{moserExponentSeq}\,d\,p_0\,n$
and observe $p_n > 0$ by \leanname{moserExponentSeq\_pos} (using $d > 2$ and
$p_0 > 0$). Apply the inverse-power one-step pre-estimate
Theorem~\ref{thm:preMoser-inv}
(\leanname{supersolution\_preMoser\_inverse}) with parameters $(p, r, s) =
(p_n, r_{n+1}, r_n)$. The hypotheses required are: positivity of $p_n$ (from
the previous remark), positivity of $r_{n+1}$ (\leanname{moserRadius\_pos}),
the strict inequality $r_{n+1} < r_n$ (\leanname{moserRadius\_succ\_lt}), the
upper bound $r_n \le 1$ (\leanname{moserRadius\_le\_one}), the strict
positivity of $u$ on the unit ball, the supersolution property of $u$, and
the integrability of $|u^{-1}|^{p_n}$ on $\Bz{r_n}$ (which is the inductive
hypothesis after rewriting via $\mathrm{moserExponentSeq}\,d\,p_0\,n = p_n$).

The pre-estimate produces:
(i) integrability of $|u^{-1}|^{\chi p_n}$ on $\Bz{r_{n+1}}$, which by
$\chi p_n = p_{n+1}$ (\leanname{moserExponentSeq\_succ}) is exactly the
integrability assertion at level $n+1$; and
(ii) the one-step bound
$N_{n+1}^{\mathrm{inv}}(u) \le \mathrm{S}_n^{\mathrm{inv}}(A, p_0)\,
N_n^{\mathrm{inv}}(u)$.
Multiplying the inductive bound for $N_n^{\mathrm{inv}}$ by
$\mathrm{S}_n^{\mathrm{inv}}$ and absorbing the additional factor into the
range product
$\prod_{i=0}^{n} \mathrm{S}_i^{\mathrm{inv}} = \bigl(\prod_{i=0}^{n-1} \mathrm{S}_i^{\mathrm{inv}}\bigr)
\cdot \mathrm{S}_n^{\mathrm{inv}}$ closes the induction.
\end{proof}

\begin{lemma}[Geometric majorant, \leanname{supersolution\_geometric\_majorant\_inv}]
\label{lem:geom-maj-inv}
For $d>2$ and $p_0>0$,
\[
  \prod_{i=0}^{n-1}\mathrm{S}_i^{\mathrm{inv}}(A,p_0)
  \;\leq\;
  \big(C_{\mathrm{wH}}^{0}(d)\,(\Lambda p_0^{2}+1)^{d/2}\big)^{1/p_0}.
\]
\end{lemma}

\begin{lstlisting}[style=Matlab-editor,basicstyle=\mlttfamily,escapechar=",mlshowsectionrules=true,mathescape=true,morecomment={[s]{\%\{}{\%\}}},language=lean,numbers=none]
theorem supersolution_geometric_majorant_inv
    (hd : 2 < (d : ℝ))
    (A : NormalizedEllipticCoeff d (Metric.ball (0 : E) 1))
    {p₀ : ℝ} (hp₀ : 0 < p₀) (n : ℕ) :
    ∏ i ∈ Finset.range n, superStepConstInv (d := d) A p₀ i ≤
      (C_weakHarnack0 d * (A.1.Λ * p₀ ^ 2 + 1) ^ ((d : ℝ) / 2)) ^
        (1 / p₀) := by
  let q := moserDecayRatio d
  let K₀ : ℝ := 16 * C_MoserAnchor d
  let L : ℝ := A.1.Λ * p₀ ^ 2 + 1
  let K : ℝ := K₀ * L
  let CC : ℝ := 4 * moserChi d ^ 2
  let S : ℝ := Finset.sum (Finset.range n) (fun i => q ^ i)
  let T : ℝ := Finset.sum (Finset.range n) (fun i => (i : ℝ) * q ^ i)
  have hq_nonneg : 0 ≤ q := by
    simpa [q] using moserDecayRatio_nonneg (d := d) hd
  have hq_lt_one : q < 1 := by
    simpa [q] using moserDecayRatio_lt_one (d := d) hd
  have hS_le : S ≤ (d : ℝ) / 2 := by
    simpa [q, S] using moserDecayRatio_sum_le_half_dim (d := d) hd n
  have hT_le : T ≤ ∑' i : ℕ, (i : ℝ) * q ^ i := by
    simpa [q, T] using moserDecayRatio_weighted_sum_le_tsum (d := d) hd n
  have hK₀_ge_one : 1 ≤ K₀ := by
    dsimp [K₀]
    nlinarith [one_le_C_MoserAnchor (d := d)]
  have hK₀_nonneg : 0 ≤ K₀ := by
    linarith
  have hL_ge_one : 1 ≤ L := by
    dsimp [L]
    nlinarith [A.1.Λ_nonneg, sq_nonneg p₀]
  have hL_nonneg : 0 ≤ L := by
    linarith
  have hK_pos : 0 < K := by
    dsimp [K]
    exact mul_pos (lt_of_lt_of_le zero_lt_one hK₀_ge_one) (by positivity)
  have hK_nonneg : 0 ≤ K := hK_pos.le
  have hCC_pos : 0 < CC := by
    dsimp [CC]
    exact mul_pos (by norm_num : (0 : ℝ) < 4) (sq_pos_of_pos (moserChi_pos (d := d) hd))
  have hCC_nonneg : 0 ≤ CC := hCC_pos.le
  have hstep_le : ∀ i, i ∈ Finset.range n →
      superStepConstInv (d := d) A p₀ i ≤ (K * CC ^ i) ^
        (1 / moserExponentSeq d p₀ i) := by
    intro i _
    simpa [K, K₀, L, CC] using superStepConstInv_le (d := d) hd A hp₀ i
  -- Each step is nonneg
  have hstep_nonneg : ∀ i, i ∈ Finset.range n →
      0 ≤ superStepConstInv (d := d) A p₀ i := by
    intro i _
    dsimp [superStepConstInv]
    apply Real.rpow_nonneg
    apply mul_nonneg
    · exact div_nonneg (le_trans (by norm_num : (0:ℝ) ≤ 1) (one_le_C_MoserAnchor (d := d))) (sq_nonneg _)
    · nlinarith [A.1.Λ_nonneg, sq_nonneg (moserExponentSeq d p₀ i / (1 + moserExponentSeq d p₀ i))]
  have hprod_le :
      ∏ i ∈ Finset.range n, superStepConstInv (d := d) A p₀ i ≤
        ∏ i ∈ Finset.range n, (K * CC ^ i) ^
          (1 / moserExponentSeq d p₀ i) :=
    Finset.prod_le_prod hstep_nonneg hstep_le
  have hprod_eval :
      ∀ m : ℕ,
        ∏ i ∈ Finset.range m, (K * CC ^ i) ^
            (1 / moserExponentSeq d p₀ i) =
          (K ^ (Finset.sum (Finset.range m) (fun i => q ^ i)) *
            CC ^ (Finset.sum (Finset.range m) (fun i => (i : ℝ) * q ^ i))) ^
            (1 / p₀) := by
    intro m
    induction m with
    | zero =>
        simp
    | succ m ihm =>
        let Sₘ : ℝ := Finset.sum (Finset.range m) (fun i => q ^ i)
        let Tₘ : ℝ := Finset.sum (Finset.range m) (fun i => (i : ℝ) * q ^ i)
        have hinv :
            1 / moserExponentSeq d p₀ m = q ^ m / p₀ := by
          simpa [q] using inv_moserExponentSeq (d := d) hd hp₀ m
        have hdiv_eq :
            1 / moserExponentSeq d p₀ m = q ^ m * (1 / p₀) := by
          rw [hinv]
          ring
        rw [Finset.prod_range_succ, ihm, hdiv_eq]
        have hKC_nonneg : 0 ≤ K * CC ^ m := by
          exact mul_nonneg hK_nonneg (pow_nonneg hCC_nonneg m)
        rw [Real.rpow_mul hKC_nonneg]
        rw [Real.mul_rpow hK_nonneg (pow_nonneg hCC_nonneg m)]
        have hCC_pow :
            ((CC ^ m : ℝ)) ^ (q ^ m) = CC ^ ((m : ℝ) * q ^ m) := by
          rw [show ((CC ^ m : ℝ)) = CC ^ (m : ℝ) by rw [Real.rpow_natCast]]
          rw [← Real.rpow_mul hCC_nonneg]
        rw [hCC_pow]
        have hfront_nonneg :
            0 ≤ K ^ (q ^ m) * CC ^ ((m : ℝ) * q ^ m) := by
          exact mul_nonneg (Real.rpow_nonneg hK_nonneg _) (Real.rpow_nonneg hCC_nonneg _)
        have hback_nonneg :
            0 ≤ K ^ Sₘ * CC ^ Tₘ := by
          exact mul_nonneg (Real.rpow_nonneg hK_nonneg _) (Real.rpow_nonneg hCC_nonneg _)
        rw [mul_comm, ← Real.mul_rpow hfront_nonneg hback_nonneg]
        congr 1
        calc
          (K ^ (q ^ m) * CC ^ ((m : ℝ) * q ^ m)) * (K ^ Sₘ * CC ^ Tₘ)
              = (K ^ (q ^ m) * K ^ Sₘ) * (CC ^ ((m : ℝ) * q ^ m) * CC ^ Tₘ) := by
                  ring
          _ = K ^ (q ^ m + Sₘ) * CC ^ (((m : ℝ) * q ^ m) + Tₘ) := by
                rw [← Real.rpow_add hK_pos, ← Real.rpow_add hCC_pos]
          _ =
              K ^ (Finset.sum (Finset.range (m + 1)) (fun i => q ^ i)) *
                CC ^ (Finset.sum (Finset.range (m + 1))
                  (fun i => (i : ℝ) * q ^ i)) := by
                simp [Sₘ, Tₘ, Finset.sum_range_succ, add_comm]
  have hK_pow_eq :
      K ^ S = K₀ ^ S * L ^ S := by
    rw [show K = K₀ * L by rfl, Real.mul_rpow hK₀_nonneg hL_nonneg]
  have hK₀_pow_le : K₀ ^ S ≤ K₀ ^ ((d : ℝ) / 2) := by
    exact Real.rpow_le_rpow_of_exponent_le hK₀_ge_one hS_le
  have hL_pow_le : L ^ S ≤ L ^ ((d : ℝ) / 2) := by
    exact Real.rpow_le_rpow_of_exponent_le hL_ge_one hS_le
  have hK_le :
      K ^ S ≤ K₀ ^ ((d : ℝ) / 2) * L ^ ((d : ℝ) / 2) := by
    rw [hK_pow_eq]
    exact mul_le_mul hK₀_pow_le hL_pow_le
      (Real.rpow_nonneg hL_nonneg _) (Real.rpow_nonneg hK₀_nonneg _)
  have h4_le :
      (4 : ℝ) ^ T ≤ 4 ^ (∑' i : ℕ, (i : ℝ) * q ^ i) := by
    exact Real.rpow_le_rpow_of_exponent_le (by norm_num : (1 : ℝ) ≤ 4) hT_le
  have hchi_sq_ge_one : 1 ≤ moserChi d ^ 2 := by
    exact one_le_pow₀ (one_le_moserChi (d := d) hd)
  have hchi_le :
      (moserChi d ^ 2) ^ T ≤ (moserChi d ^ 2) ^ (∑' i : ℕ, (i : ℝ) * q ^ i) := by
    exact Real.rpow_le_rpow_of_exponent_le hchi_sq_ge_one hT_le
  have hCC_le :
      CC ^ T ≤
        4 ^ (∑' i : ℕ, (i : ℝ) * q ^ i) *
          (moserChi d ^ 2) ^ (∑' i : ℕ, (i : ℝ) * q ^ i) := by
    calc
      CC ^ T = (4 : ℝ) ^ T * (moserChi d ^ 2) ^ T := by
        rw [show CC = 4 * (moserChi d ^ 2) by simp [CC]]
        rw [Real.mul_rpow (by positivity) (sq_nonneg _)]
      _ ≤ 4 ^ (∑' i : ℕ, (i : ℝ) * q ^ i) *
            (moserChi d ^ 2) ^ (∑' i : ℕ, (i : ℝ) * q ^ i) := by
          exact mul_le_mul h4_le hchi_le
            (Real.rpow_nonneg (sq_nonneg _) _)
            (Real.rpow_nonneg (by norm_num : (0 : ℝ) ≤ 4) _)
  have hK₀_pow_le_big :
      K₀ ^ ((d : ℝ) / 2) ≤ ((32 : ℝ) * C_MoserAnchor d) ^ ((d : ℝ) / 2) := by
    exact Real.rpow_le_rpow
      hK₀_nonneg
      (by
        dsimp [K₀]
        nlinarith [one_le_C_MoserAnchor (d := d)])
      (by positivity)
  have hCMoser_le :
      K₀ ^ ((d : ℝ) / 2) * 4 ^ (∑' i : ℕ, (i : ℝ) * q ^ i) ≤ C_Moser d := by
    have hbase :
        K₀ ^ ((d : ℝ) / 2) * 4 ^ (∑' i : ℕ, (i : ℝ) * q ^ i) ≤
          ((32 : ℝ) * C_MoserAnchor d) ^ ((d : ℝ) / 2) *
            4 ^ (∑' i : ℕ, (i : ℝ) * q ^ i) := by
      exact mul_le_mul_of_nonneg_right hK₀_pow_le_big
        (Real.rpow_nonneg (by norm_num : (0 : ℝ) ≤ 4) _)
    have hmax :
        ((32 : ℝ) * C_MoserAnchor d) ^ ((d : ℝ) / 2) *
            4 ^ (∑' i : ℕ, (i : ℝ) * q ^ i) ≤
          C_Moser d := by
      dsimp [C_Moser, q]
      split_ifs with hlt
      · exact
          le_max_right (C_MoserAnchor d)
            (((32 : ℝ) * C_MoserAnchor d) ^ ((d : ℝ) / 2) *
              4 ^ (∑' i : ℕ, (i : ℝ) * moserDecayRatio d ^ i))
      · exact False.elim (hlt hd)
    exact hbase.trans hmax
  have hCweak_eq :
      C_Moser d * (moserChi d ^ 2) ^ (∑' i : ℕ, (i : ℝ) * q ^ i) =
        C_weakHarnack0 d := by
    simp [C_weakHarnack0, hd, q]
  have hbody_le :
      K ^ S * CC ^ T ≤
        C_weakHarnack0 d * L ^ ((d : ℝ) / 2) := by
    have hCC_target_nonneg :
        0 ≤ 4 ^ (∑' i : ℕ, (i : ℝ) * q ^ i) *
            (moserChi d ^ 2) ^ (∑' i : ℕ, (i : ℝ) * q ^ i) := by
      exact mul_nonneg
        (Real.rpow_nonneg (by norm_num : (0 : ℝ) ≤ 4) _)
        (Real.rpow_nonneg (sq_nonneg _) _)
    have hK_target_nonneg :
        0 ≤ K₀ ^ ((d : ℝ) / 2) * L ^ ((d : ℝ) / 2) := by
      exact mul_nonneg (Real.rpow_nonneg hK₀_nonneg _) (Real.rpow_nonneg hL_nonneg _)
    calc
      K ^ S * CC ^ T
          ≤ (K₀ ^ ((d : ℝ) / 2) * L ^ ((d : ℝ) / 2)) *
              (4 ^ (∑' i : ℕ, (i : ℝ) * q ^ i) *
                (moserChi d ^ 2) ^ (∑' i : ℕ, (i : ℝ) * q ^ i)) := by
            exact mul_le_mul hK_le hCC_le
              (Real.rpow_nonneg hCC_nonneg _)
              hK_target_nonneg
      _ = (K₀ ^ ((d : ℝ) / 2) * 4 ^ (∑' i : ℕ, (i : ℝ) * q ^ i)) *
            ((moserChi d ^ 2) ^ (∑' i : ℕ, (i : ℝ) * q ^ i) * L ^ ((d : ℝ) / 2)) := by
              ring
      _ ≤ C_Moser d *
            ((moserChi d ^ 2) ^ (∑' i : ℕ, (i : ℝ) * q ^ i) * L ^ ((d : ℝ) / 2)) := by
              exact mul_le_mul_of_nonneg_right hCMoser_le
                (mul_nonneg (Real.rpow_nonneg (sq_nonneg _) _) (Real.rpow_nonneg hL_nonneg _))
      _ = C_weakHarnack0 d * L ^ ((d : ℝ) / 2) := by
            rw [← hCweak_eq]
            ring
  calc
    ∏ i ∈ Finset.range n, superStepConstInv (d := d) A p₀ i
        ≤ ∏ i ∈ Finset.range n, (K * CC ^ i) ^ (1 / moserExponentSeq d p₀ i) := hprod_le
    _ = (K ^ S * CC ^ T) ^ (1 / p₀) := by
          simpa [q, S, T] using hprod_eval n
    _ ≤ (C_weakHarnack0 d * L ^ ((d : ℝ) / 2)) ^ (1 / p₀) := by
          exact Real.rpow_le_rpow
            (mul_nonneg (Real.rpow_nonneg hK_nonneg _) (Real.rpow_nonneg hCC_nonneg _))
            hbody_le (by positivity)
    _ = (C_weakHarnack0 d * (A.1.Λ * p₀ ^ 2 + 1) ^ ((d : ℝ) / 2)) ^ (1 / p₀) := by
          simp [L]

\end{lstlisting}
\begin{proof}
By Lemma~\ref{lem:stepInv-le}, the product is bounded by
\[
  \prod_{i=0}^{n-1}\big(K\,L\,(4\chi^{2})^{i}\big)^{1/p_i},
\]
where $K := 16\,C_{\mathrm{MA}}(d)$, $L := \Lambda p_0^{2}+1$, $1/p_i = q^{i}/p_0$.
Taking logs and using $\sum_{i<n} q^{i} \leq d/2$ (\leanname{moserDecayRatio\_sum\_le\_half\_dim})
and $\sum_{i<n} i q^{i} \leq \sum_{i\geq 0} i q^{i}$
(\leanname{moserDecayRatio\_weighted\_sum\_le\_tsum}), the product is at most
\[
  \big(K^{d/2}\,L^{d/2}\,(4\chi^{2})^{\sum_{i\geq 0} i q^{i}}\big)^{1/p_0}.
\]
Since $K^{d/2}\,4^{\sum_{i\geq 0} i q^{i}} \leq C_{\mathrm{Moser}}(d)$ by definition
of $C_{\mathrm{Moser}}$, and $(\chi^{2})^{\sum_{i\geq 0} i q^{i}}$ is exactly the
extra factor included in $C_{\mathrm{wH}}^{0}(d)$, the product is bounded by
$\big(C_{\mathrm{wH}}^{0}(d)\,L^{d/2}\big)^{1/p_0}$.
\end{proof}

\begin{theorem}[Inverse a.e.\ closeout, \leanname{supersolution\_ae\_closeout\_inv}]
\label{thm:ae-closeout-inv}
For $d>2$, $p_0 > 0$, $u>0$ a supersolution and $|u^{-1}|^{p_0}$ integrable on
$\Bz{1}$, for almost every $x \in \Bz{1/2}$,
\[
  |u(x)^{-1}|^{p_0}
  \;\leq\;
  C_{\mathrm{wH}}^{0}(d)\,(\Lambda p_0^{2}+1)^{d/2}\,
  \int_{\Bz{1}} |u^{-1}|^{p_0}\,dx.
\]
\end{theorem}

\begin{lstlisting}[style=Matlab-editor,basicstyle=\mlttfamily,escapechar=",mlshowsectionrules=true,mathescape=true,morecomment={[s]{\%\{}{\%\}}},language=lean,numbers=none]
/-- Closeout: pass from iterated `L^{pₙ}(B_{rₙ})` bounds to an a.e. `L^∞`
bound on `u⁻¹` over `B_{1/2}`.

Since `rₙ > 1/2` for all `n` and `pₙ → ∞`, the `L^{pₙ}` norms converge to
the `L^∞` norm. The uniform bound from the iteration + geometric majorant
gives the pointwise bound. -/
theorem supersolution_ae_closeout_inv
    (hd : 2 < (d : ℝ))
    (A : NormalizedEllipticCoeff d (Metric.ball (0 : E) 1))
    {u : E → ℝ} {p₀ : ℝ} (hp₀ : 0 < p₀)
    (hu_pos : ∀ x ∈ Metric.ball (0 : E) 1, 0 < u x)
    (hsuper : IsSupersolution A.1 u)
    (hpInt :
      IntegrableOn (fun x => |(u x)⁻¹| ^ p₀)
        (Metric.ball (0 : E) 1) volume) :
    ∀ᵐ x ∂(volume.restrict (Metric.ball (0 : E) (1 / 2 : ℝ))),
      |(u x)⁻¹| ^ p₀ ≤
        C_weakHarnack0 d * (A.1.Λ * p₀ ^ 2 + 1) ^ ((d : ℝ) / 2) *
          ∫ x in Metric.ball (0 : E) 1, |(u x)⁻¹| ^ p₀ ∂volume := by
  let Bhalf : Set E := Metric.ball (0 : E) (1 / 2 : ℝ)
  let μ : Measure E := volume.restrict Bhalf
  let f : E → ℝ := fun x => |(u x)⁻¹|
  let g : E → ℝ := fun x => f x ^ (p₀ / 2)
  let K : ℝ :=
    (C_weakHarnack0 d * (A.1.Λ * p₀ ^ 2 + 1) ^ ((d : ℝ) / 2) *
      ∫ x in Metric.ball (0 : E) 1, |(u x)⁻¹| ^ p₀ ∂volume) ^ (1 / 2 : ℝ)
  haveI : IsFiniteMeasure μ := by
    dsimp [μ, Bhalf]
    rw [isFiniteMeasure_restrict]
    exact measure_ne_top_of_subset Metric.ball_subset_closedBall
      (isCompact_closedBall (0 : E) (1 / 2 : ℝ)).measure_lt_top.ne
  have hK_nonneg : 0 ≤ K := by
    dsimp [K]
    exact Real.rpow_nonneg (supersolutionInvBoundPow_nonneg (d := d) A) _
  have hiter :
      ∀ n,
        IntegrableOn (fun x => |(u x)⁻¹| ^ moserExponentSeq d p₀ n)
          (Metric.ball (0 : E) (moserRadius n)) volume ∧
        superIterNormInv (d := d) (u := u) p₀ n ≤
          ((C_weakHarnack0 d * (A.1.Λ * p₀ ^ 2 + 1) ^ ((d : ℝ) / 2)) ^
            (1 / p₀)) *
            superIterNormInv (d := d) (u := u) p₀ 0 := by
    intro n
    have hraw := supersolution_iteration_inverse (d := d) hd A hp₀ hu_pos hsuper hpInt n
    have hgeom := supersolution_geometric_majorant_inv (d := d) hd A hp₀ n
    refine ⟨hraw.1, ?_⟩
    calc
      superIterNormInv (d := d) (u := u) p₀ n
          ≤ (∏ i ∈ Finset.range n, superStepConstInv (d := d) A p₀ i) *
              superIterNormInv (d := d) (u := u) p₀ 0 := hraw.2
      _ ≤ ((C_weakHarnack0 d * (A.1.Λ * p₀ ^ 2 + 1) ^ ((d : ℝ) / 2)) ^
            (1 / p₀)) *
            superIterNormInv (d := d) (u := u) p₀ 0 := by
            exact mul_le_mul_of_nonneg_right hgeom
              (by
                dsimp [superIterNormInv]
                exact Real.rpow_nonneg (integral_nonneg fun x => by positivity) _)
  have hbound_all :
      ∀ m : ℕ, ∀ᵐ x ∂μ, g x < K + 1 / (m + 1 : ℝ) := by
    intro m
    have hzero :
        μ.real {x | K + 1 / (m + 1 : ℝ) ≤ g x} = 0 := by
      refine supersolution_closeout_superlevel_null (d := d) hd A (u := u) hp₀ hiter ?_
      dsimp [K]
      have hfrac_pos : 0 < 1 / (m + 1 : ℝ) := by positivity
      linarith
    rw [ae_iff]
    simpa [μ, Bhalf, g, K, not_lt] using
      (measureReal_eq_zero_iff
        (μ := μ) (s := {x | K + 1 / (m + 1 : ℝ) ≤ g x})
        (measure_ne_top μ {x | K + 1 / (m + 1 : ℝ) ≤ g x})).1 hzero
  have hbound_ae :
      ∀ᵐ x ∂μ, ∀ m : ℕ, g x < K + 1 / (m + 1 : ℝ) := by
    exact ae_all_iff.2 hbound_all
  filter_upwards [hbound_ae] with x hx
  have hfx_nonneg : 0 ≤ f x := by
    dsimp [f]
    exact abs_nonneg ((u x)⁻¹)
  have hgx_nonneg : 0 ≤ g x := by
    dsimp [g]
    exact Real.rpow_nonneg hfx_nonneg _
  have hgx_le : g x ≤ K := by
    by_contra hgx_gt
    have hgap_pos : 0 < g x - K := by linarith
    obtain ⟨m, hm⟩ := exists_nat_one_div_lt hgap_pos
    have hlt : K + 1 / (m + 1 : ℝ) < g x := by linarith
    linarith [hx m]
  have hgpow :
      g x ^ (2 : ℝ) ≤ K ^ (2 : ℝ) := by
    exact Real.rpow_le_rpow hgx_nonneg hgx_le (by norm_num)
  have hg_sq_eq : g x ^ (2 : ℝ) = f x ^ p₀ := by
    calc
      g x ^ (2 : ℝ) = (f x ^ (p₀ / 2)) ^ (2 : ℝ) := by rfl
      _ = f x ^ ((p₀ / 2) * 2) := by rw [← Real.rpow_mul hfx_nonneg]
      _ = f x ^ p₀ := by
            congr 2
            ring
  have hK_sq_eq :
      K ^ (2 : ℝ) =
        C_weakHarnack0 d * (A.1.Λ * p₀ ^ 2 + 1) ^ ((d : ℝ) / 2) *
          ∫ x in Metric.ball (0 : E) 1, |(u x)⁻¹| ^ p₀ ∂volume := by
    calc
      K ^ (2 : ℝ)
          =
            (C_weakHarnack0 d * (A.1.Λ * p₀ ^ 2 + 1) ^ ((d : ℝ) / 2) *
              ∫ x in Metric.ball (0 : E) 1, |(u x)⁻¹| ^ p₀ ∂volume) ^
              ((1 / 2 : ℝ) * 2) := by
                dsimp [K]
                rw [← Real.rpow_mul (supersolutionInvBoundPow_nonneg (d := d) A)]
      _ =
            C_weakHarnack0 d * (A.1.Λ * p₀ ^ 2 + 1) ^ ((d : ℝ) / 2) *
              ∫ x in Metric.ball (0 : E) 1, |(u x)⁻¹| ^ p₀ ∂volume := by
                rw [show ((1 / 2 : ℝ) * 2) = (1 : ℝ) by ring, Real.rpow_one]
  calc
    |(u x)⁻¹| ^ p₀ = f x ^ p₀ := by rfl
    _ = g x ^ (2 : ℝ) := hg_sq_eq.symm
    _ ≤ K ^ (2 : ℝ) := hgpow
    _ = C_weakHarnack0 d * (A.1.Λ * p₀ ^ 2 + 1) ^ ((d : ℝ) / 2) *
          ∫ x in Metric.ball (0 : E) 1, |(u x)⁻¹| ^ p₀ ∂volume := hK_sq_eq


\end{lstlisting}
\begin{proof}
Combine Theorem~\ref{thm:iter-inv} with Lemma~\ref{lem:geom-maj-inv} to get,
for every $n$,
\[
  N_n^{\mathrm{inv}}(u)
  \;\leq\;
  \big(C_{\mathrm{wH}}^{0}(d)\,(\Lambda p_0^{2}+1)^{d/2}\,I_0^{\mathrm{inv}}(u)\big)^{1/p_0}.
\]
The Markov-type closeout argument of Lemma~\ref{lem:moser-closeout}, applied
verbatim with $u_+$ replaced by $u^{-1}$ (and $1<p_0$ replaced by $p_0>0$, which
only affects the choice of test exponent $q_n = 2\chi^{n}$ but not the
convergence) yields the desired a.e.\ inequality.
\end{proof}

\begin{theorem}[Stage one weak Harnack (inverse), \leanname{weak\_harnack\_stage\_one\_inverse}]
\label{thm:wH-stage1-inv}
Let $d>2$, $A$ on $\Bz{1}$, $u>0$ on $\Bz{1}$ a supersolution, and $p_0 > 0$. Then
\begin{multline*}
  (\Lambda p_0^{2}+1)^{-d/2}\,
  \Big(\int_{\Bz{1}} |u^{-1}|^{p_0}\,dx\Big)^{-1}
  \\
  \;\leq\; C_{\mathrm{wH}}^{0}(d)\;(\essinf_{\Bz{1/2}} u)^{p_0}.
\end{multline*}
\end{theorem}

\begin{lstlisting}[style=Matlab-editor,basicstyle=\mlttfamily,escapechar=",mlshowsectionrules=true,mathescape=true,morecomment={[s]{\%\{}{\%\}}},language=lean,numbers=none]
/-- First stage of weak Harnack: inverse-power iteration for positive
supersolutions.

For `u > 0` a supersolution on `B₁` and every `p₀ > 0`:
```
(Λ p₀² + 1)^{-d/2} · (∫_{B₁} |u⁻¹|^{p₀})⁻¹ ≤ C · (inf_{B_{1/2}} u)^{p₀}
```
Proof: the Moser iteration (Steps 1-2 above) gives an a.e. `L^∞` bound on
`u⁻¹` over `B_{1/2}`. Since `u > 0` pointwise, this converts to a lower
bound on `inf u`, which we rearrange to get the stated inequality. -/
theorem weak_harnack_stage_one_inverse
    (hd : 2 < (d : ℝ))
    (A : NormalizedEllipticCoeff d (Metric.ball (0 : E) 1))
    {u : E → ℝ} {p₀ : ℝ} (hp₀ : 0 < p₀)
    (hu_pos : ∀ x ∈ Metric.ball (0 : E) 1, 0 < u x)
    (hsuper : IsSupersolution A.1 u) :
    (A.1.Λ * p₀ ^ 2 + 1) ^ (-(d : ℝ) / 2) *
        (∫ x in Metric.ball (0 : E) 1, |(u x)⁻¹| ^ p₀ ∂volume)⁻¹ ≤
      C_weakHarnack0 d *
        (essInf u μhalf) ^ p₀ := by
  let I : ℝ := ∫ x in Metric.ball (0 : E) 1, |(u x)⁻¹| ^ p₀ ∂volume
  let L : ℝ := A.1.Λ * p₀ ^ 2 + 1
  have hC_nonneg : 0 ≤ C_weakHarnack0 d := by
    exact le_trans (by norm_num : (0 : ℝ) ≤ 1) (one_le_C_weakHarnack0 (d := d))
  have hL_pos : 0 < L := by
    dsimp [L]
    nlinarith [A.1.Λ_nonneg, sq_nonneg p₀]
  have hL_nonneg : 0 ≤ L := hL_pos.le
  have hhalf_nonneg : ∀ᵐ x ∂μhalf, 0 ≤ u x := by
    filter_upwards [ae_restrict_mem measurableSet_ball] with x hx
    exact (hu_pos x ((Metric.ball_subset_ball (by norm_num : (1 / 2 : ℝ) ≤ 1)) hx)).le
  have hessInf_nonneg : 0 ≤ essInf u μhalf := by
    exact
      le_essInf_real_of_ae_le
        (d := d)
        (restrict_ball_ne_zero (c := (0 : E)) (r := (1 / 2 : ℝ)) (by norm_num))
        hhalf_nonneg
  by_cases hpInt : IntegrableOn (fun x => |(u x)⁻¹| ^ p₀) (Metric.ball (0 : E) 1) volume
  · let c0 : ℝ := C_weakHarnack0 d * L ^ ((d : ℝ) / 2) * I
    have hI_nonneg : 0 ≤ I := by
      dsimp [I]
      exact integral_nonneg fun x => by positivity
    have hI_pos : 0 < I := by
      have hnonneg :
          0 ≤ᵐ[volume.restrict (Metric.ball (0 : E) 1)] fun x => |(u x)⁻¹| ^ p₀ := by
        filter_upwards [ae_restrict_mem measurableSet_ball] with x hx
        positivity
      have hI_ne_zero : I ≠ 0 := by
        intro hI_zero
        have hzero_ae :
            (fun x => |(u x)⁻¹| ^ p₀) =ᵐ[volume.restrict (Metric.ball (0 : E) 1)] 0 := by
          rw [← sub_eq_zero] at hI_zero
          rwa [sub_zero, setIntegral_eq_zero_iff_of_nonneg_ae hnonneg hpInt] at hI_zero
        have hpos_ae :
            ∀ᵐ x ∂(volume.restrict (Metric.ball (0 : E) 1)), |(u x)⁻¹| ^ p₀ ≠ 0 := by
          filter_upwards [ae_restrict_mem measurableSet_ball] with x hx
          have hux_pos : 0 < u x := hu_pos x hx
          have hpow_pos : 0 < |(u x)⁻¹| ^ p₀ := by
            rw [abs_of_pos (inv_pos.mpr hux_pos)]
            exact Real.rpow_pos_of_pos (inv_pos.mpr hux_pos) p₀
          exact hpow_pos.ne'
        have hfalse : ∀ᵐ x ∂(volume.restrict (Metric.ball (0 : E) 1)), False := by
          filter_upwards [hzero_ae, hpos_ae] with x hx0 hxpos
          exact hxpos hx0
        rw [ae_iff] at hfalse
        have hball_zero : volume (Metric.ball (0 : E) 1) = 0 := by
          simpa [Measure.restrict_apply_univ] using hfalse
        exact (Metric.measure_ball_pos volume (0 : E) (by norm_num : (0 : ℝ) < 1)).ne'
          hball_zero
      exact lt_of_le_of_ne hI_nonneg (Ne.symm hI_ne_zero)
    have hc0_pos : 0 < c0 := by
      dsimp [c0]
      exact mul_pos (mul_pos (lt_of_lt_of_le zero_lt_one (one_le_C_weakHarnack0 (d := d)))
        (Real.rpow_pos_of_pos hL_pos _)) hI_pos
    have hc0_nonneg : 0 ≤ c0 := hc0_pos.le
    have hclose := supersolution_ae_closeout_inv (d := d) hd A hp₀ hu_pos hsuper hpInt
    have hroot_ae :
        ∀ᵐ x ∂μhalf, (c0⁻¹) ^ (1 / p₀) ≤ u x := by
      filter_upwards [hclose, ae_restrict_mem measurableSet_ball] with x hxclose hxhalf
      have hux_pos : 0 < u x := hu_pos x ((Metric.ball_subset_ball (by norm_num : (1 / 2 : ℝ) ≤ 1)) hxhalf)
      have hupow_pos : 0 < u x ^ p₀ := Real.rpow_pos_of_pos hux_pos p₀
      have hpow_inv : (u x ^ p₀)⁻¹ ≤ c0 := by
        calc
          (u x ^ p₀)⁻¹ = ((u x)⁻¹) ^ p₀ := by
            rw [Real.inv_rpow hux_pos.le]
          _ = |(u x)⁻¹| ^ p₀ := by
            rw [abs_of_pos (inv_pos.mpr hux_pos)]
          _ ≤ c0 := hxclose
      have hpow_lower : c0⁻¹ ≤ u x ^ p₀ := by
        have htmp :
            c0⁻¹ ≤ ((u x ^ p₀)⁻¹)⁻¹ := (inv_le_inv₀ hc0_pos (inv_pos.mpr hupow_pos)).2 hpow_inv
        simpa [hupow_pos.ne'] using htmp
      have hroot :=
        Real.rpow_le_rpow (inv_nonneg.mpr hc0_nonneg) hpow_lower (by positivity : 0 ≤ 1 / p₀)
      have hux_root : (u x ^ p₀) ^ (1 / p₀) = u x := by
        rw [← Real.rpow_mul hux_pos.le]
        field_simp [hp₀.ne']
        rw [Real.rpow_one]
      calc
        (c0⁻¹) ^ (1 / p₀) ≤ (u x ^ p₀) ^ (1 / p₀) := hroot
        _ = u x := hux_root
    have hc_le_ess : (c0⁻¹) ^ (1 / p₀) ≤ essInf u μhalf := by
      exact
        le_essInf_real_of_ae_le
          (d := d)
          (restrict_ball_ne_zero (c := (0 : E)) (r := (1 / 2 : ℝ)) (by norm_num))
          hroot_ae
    have hpow_ess :
        c0⁻¹ ≤ (essInf u μhalf) ^ p₀ := by
      have hpow :=
        Real.rpow_le_rpow
          (Real.rpow_nonneg (inv_nonneg.mpr hc0_nonneg) _)
          hc_le_ess hp₀.le
      have hleft : ((c0⁻¹) ^ (1 / p₀)) ^ p₀ = c0⁻¹ := by
        calc
          ((c0⁻¹) ^ (1 / p₀)) ^ p₀ = (c0⁻¹) ^ ((1 / p₀) * p₀) := by
            rw [← Real.rpow_mul (inv_nonneg.mpr hc0_nonneg)]
          _ = (c0⁻¹) ^ (1 : ℝ) := by
            congr 1
            field_simp [hp₀.ne']
          _ = c0⁻¹ := by rw [Real.rpow_one]
      calc
        c0⁻¹ = ((c0⁻¹) ^ (1 / p₀)) ^ p₀ := hleft.symm
        _ ≤ (essInf u μhalf) ^ p₀ := hpow
    have hC_pos : 0 < C_weakHarnack0 d :=
      lt_of_lt_of_le zero_lt_one (one_le_C_weakHarnack0 (d := d))
    have hLpow_ne : L ^ ((d : ℝ) / 2) ≠ 0 := (Real.rpow_pos_of_pos hL_pos _).ne'
    have hLneg : L ^ (-(d : ℝ) / 2) = (L ^ ((d : ℝ) / 2))⁻¹ := by
      rw [show (-(d : ℝ) / 2) = -((d : ℝ) / 2) by ring, Real.rpow_neg hL_pos.le]
    calc
      L ^ (-(d : ℝ) / 2) * I⁻¹
          = C_weakHarnack0 d * c0⁻¹ := by
              dsimp [c0]
              rw [hLneg]
              field_simp [hC_pos.ne', hLpow_ne, hI_pos.ne']
      _ ≤ C_weakHarnack0 d * (essInf u μhalf) ^ p₀ := by
            exact mul_le_mul_of_nonneg_left hpow_ess hC_nonneg
  · have hrhs_nonneg : 0 ≤ C_weakHarnack0 d * (essInf u μhalf) ^ p₀ := by
      exact mul_nonneg hC_nonneg (Real.rpow_nonneg hessInf_nonneg _)
    have hI_zero :
        ∫ x in Metric.ball (0 : E) 1, |(u x)⁻¹| ^ p₀ ∂volume = 0 := by
      simpa using (integral_undef hpInt :
        ∫ x in Metric.ball (0 : E) 1, |(u x)⁻¹| ^ p₀ ∂volume = 0)
    calc
      (A.1.Λ * p₀ ^ 2 + 1) ^ (-(d : ℝ) / 2) *
          (∫ x in Metric.ball (0 : E) 1, |(u x)⁻¹| ^ p₀ ∂volume)⁻¹
          = 0 := by rw [hI_zero]; simp
      _ ≤ C_weakHarnack0 d * (essInf u μhalf) ^ p₀ := hrhs_nonneg

\end{lstlisting}
\begin{proof}
If $|u^{-1}|^{p_0}$ is not integrable on $\Bz{1}$, the integral is taken to be
$+\infty$ and the LHS is $0 \leq \mathrm{RHS}$.

Otherwise, set $I := \int_{\Bz{1}} |u^{-1}|^{p_0}\,dx$ and
$c_0 := C_{\mathrm{wH}}^{0}(d)\,(\Lambda p_0^{2}+1)^{d/2}\,I$. By
Theorem~\ref{thm:ae-closeout-inv} we have $|u^{-1}|^{p_0} \leq c_0$ a.e.\ on
$\Bz{1/2}$, equivalently $u^{p_0} \geq c_0^{-1}$ a.e., i.e.\
$u \geq c_0^{-1/p_0}$ a.e.\ on $\Bz{1/2}$. Therefore
$\essinf_{\Bz{1/2}} u \geq c_0^{-1/p_0}$, so
$(\essinf_{\Bz{1/2}} u)^{p_0} \geq c_0^{-1}$. Multiplying by
$C_{\mathrm{wH}}^{0}(d) \geq 0$ and using
$L^{-d/2} I^{-1} = C_{\mathrm{wH}}^{0}(d)\,c_0^{-1}$ (where $L = \Lambda p_0^{2}+1$)
gives the displayed inequality.
\end{proof}

\subsubsection{Stage two (forward low-power iteration)}

The forward iteration is run only finitely many steps: starting from
$p_0 := q\,\chi^{-m}$ for a target exponent $q \in (0,1)$ and a depth $m$
chosen so that $p_0 \leq p$, we iterate $m+1$ times to reach the level $p_m =
q$. The forward step constant
\begin{equation*}
  \mathrm{S}_n^{\mathrm{fwd}}(A,p_0) :=
  \Big(\frac{C_{\mathrm{MA}}(d)}{(r_n-r_{n+1})^{2}}\,
    \Big(\Lambda\Big(\frac{p_n}{1-p_n}\Big)^{2}+1\Big)\Big)^{1/p_n}
\end{equation*}
\leanname{superStepConstFwd} is finite as long as $p_n < 1$, which is
guaranteed by $p_n \leq q < 1$ for $n \leq m$.

\begin{lemma}[Forward depth, \leanname{exists\_forward\_iteration\_depth}]
\label{lem:fwd-depth}
For $0 < p < q < 1$ there exists $m \in \N$ with
\[
  q\,\chi^{-m} < p \quad \text{and} \quad p \leq \chi\,\big(q\,\chi^{-m}\big).
\]
\end{lemma}

\begin{lstlisting}[style=Matlab-editor,basicstyle=\mlttfamily,escapechar=",mlshowsectionrules=true,mathescape=true,morecomment={[s]{\%\{}{\%\}}},language=lean,numbers=none]
theorem exists_forward_iteration_depth
    (hd : 2 < (d : ℝ))
    {p q : ℝ} (hp : 0 < p) (hpq : p < q) :
    ∃ m : ℕ,
      let p₀ := q * (moserChi d)⁻¹ ^ m
      p₀ < p ∧ p ≤ moserChi d * p₀ := by
  let ρ : ℝ := (moserChi d)⁻¹
  have hρ_pos : 0 < ρ := inv_pos.mpr (moserChi_pos (d := d) hd)
  have hρ_lt_one : ρ < 1 := by
    simpa [ρ] using
      (one_div_lt_one_div_of_lt (a := (1 : ℝ)) (b := moserChi d)
        zero_lt_one (one_lt_moserChi (d := d) hd))
  have hpq_ratio_pos : 0 < p / q := div_pos hp (lt_trans hp hpq)
  have hpq_ratio_le : p / q ≤ 1 := by
    have hq_pos : 0 < q := lt_trans hp hpq
    exact (div_le_one hq_pos).2 hpq.le
  rcases exists_nat_pow_near_of_lt_one hpq_ratio_pos hpq_ratio_le hρ_pos hρ_lt_one with
    ⟨n, hn_lt, hn_le⟩
  refine ⟨n + 1, ?_⟩
  dsimp
  constructor
  · have hq_pos : 0 < q := lt_trans hp hpq
    have hmul := mul_lt_mul_of_pos_left hn_lt hq_pos
    simpa [ρ, div_eq_mul_inv, hq_pos.ne', mul_assoc, mul_left_comm, mul_comm] using hmul
  · have hq_pos : 0 < q := lt_trans hp hpq
    have hmul := mul_le_mul_of_nonneg_left hn_le hq_pos.le
    have hp_le_base : p ≤ q * ρ ^ n := by
      simpa [ρ, div_eq_mul_inv, hq_pos.ne', mul_assoc, mul_left_comm, mul_comm] using hmul
    have hχ_ne : moserChi d ≠ 0 := (moserChi_pos (d := d) hd).ne'
    calc
      p ≤ q * ρ ^ n := hp_le_base
      _ = moserChi d * (q * ρ ^ (n + 1)) := by
            dsimp [ρ]
            rw [pow_succ]
            field_simp [hχ_ne]

\end{lstlisting}
\begin{proof}
Since $\chi > 1$, $q\,\chi^{-m} \to 0$ as $m \to \infty$, so any sufficiently
large $m$ satisfies the first inequality; choose the smallest such $m$, then
$q\,\chi^{-(m-1)} \geq p$, equivalently $p \leq \chi\,(q\,\chi^{-m})$.
\end{proof}

\begin{lemma}[Forward step bound, \leanname{superStepConstFwd\_le}]
For $d>2$, $0 < q < 1$, $m \in \N$, $i \leq m$, and $p_0 := q\,\chi^{-m}$,
\[
  \mathrm{S}_i^{\mathrm{fwd}}(A,p_0)
  \;\leq\;
  \Big(16\,C_{\mathrm{MA}}(d)\,\Big(\frac{\Lambda p_0^{2}}{(1-q)^{2}}+1\Big)\,
  (4\chi^{2})^{i}\Big)^{1/p_i}.
\]
\end{lemma}

\begin{lstlisting}[style=Matlab-editor,basicstyle=\mlttfamily,escapechar=",mlshowsectionrules=true,mathescape=true,morecomment={[s]{\%\{}{\%\}}},language=lean,numbers=none]
/-- Each forward low-power step constant is bounded by a simpler expression.

The denominator `1 - pₙ` is controlled from below by `1 - q`, because along the
finite forward iteration we only use stages `n ≤ m` and then `pₙ ≤ q < 1`. -/
theorem superStepConstFwd_le
    (hd : 2 < (d : ℝ))
    (A : NormalizedEllipticCoeff d (Metric.ball (0 : E) 1))
    {q : ℝ} {m i : ℕ}
    (hq : 0 < q) (hq1 : q < 1)
    (hi : i ∈ Finset.range (m + 1)) :
    let p₀ := q * (moserChi d)⁻¹ ^ m
    superStepConstFwd (d := d) A p₀ i ≤
      (16 * C_MoserAnchor d * (A.1.Λ * p₀ ^ 2 / (1 - q) ^ 2 + 1) *
        (4 * moserChi d ^ 2) ^ i) ^
        (1 / moserExponentSeq d p₀ i) := by
  intro p₀
  have hp₀ : 0 < p₀ := by
    dsimp [p₀]
    exact mul_pos hq (pow_pos (inv_pos.mpr (moserChi_pos (d := d) hd)) m)
  have hi_le : i ≤ m := Nat.lt_succ_iff.mp (Finset.mem_range.mp hi)
  unfold superStepConstFwd
  apply Real.rpow_le_rpow
  · apply mul_nonneg
    · exact div_nonneg
        (le_trans (by norm_num : (0 : ℝ) ≤ 1) (one_le_C_MoserAnchor (d := d)))
        (sq_nonneg _)
    · nlinarith [A.1.Λ_nonneg, sq_nonneg (moserExponentSeq d p₀ i / (1 - moserExponentSeq d p₀ i))]
  · have hp_i : 0 < moserExponentSeq d p₀ i := moserExponentSeq_pos (d := d) hd hp₀ i
    have hp_i_le_q : moserExponentSeq d p₀ i ≤ q := by
      dsimp [p₀]
      rw [moserExponentSeq, show q * (moserChi d)⁻¹ ^ m * moserChi d ^ i =
        q * ((moserChi d)⁻¹ ^ m * moserChi d ^ i) from by ring]
      have hprod_le : (moserChi d)⁻¹ ^ m * moserChi d ^ i ≤ 1 := by
        rw [inv_pow]
        have hchi_ge_one := (one_lt_moserChi (d := d) hd).le
        have hchi_pow_pos : (0 : ℝ) < moserChi d ^ m := pow_pos (moserChi_pos (d := d) hd) m
        rw [show (moserChi d ^ m)⁻¹ * moserChi d ^ i =
          moserChi d ^ i / moserChi d ^ m from by rw [div_eq_mul_inv, mul_comm]]
        rw [div_le_one hchi_pow_pos]
        exact pow_le_pow_right₀ hchi_ge_one hi_le
      calc q * ((moserChi d)⁻¹ ^ m * moserChi d ^ i) ≤ q * 1 :=
            mul_le_mul_of_nonneg_left hprod_le hq.le
        _ = q := mul_one q
    have hp_i_lt_one : moserExponentSeq d p₀ i < 1 := lt_of_le_of_lt hp_i_le_q hq1
    have hgap_eq : C_MoserAnchor d / (moserRadius i - moserRadius (i + 1)) ^ 2 =
        C_MoserAnchor d * (4 : ℝ) ^ (i + 2) := by
      rw [moserRadius_gap]
      have hsq : (((1 / 2 : ℝ) ^ (i + 2)) ^ 2) = (1 / 4 : ℝ) ^ (i + 2) := by
        rw [← pow_mul, show (i + 2) * 2 = 2 * (i + 2) by ring, pow_mul]; norm_num
      rw [hsq, div_eq_mul_inv,
        show (((1 / 4 : ℝ) ^ (i + 2))⁻¹) = (4 : ℝ) ^ (i + 2) by
          rw [show (1 / 4 : ℝ) = (4 : ℝ)⁻¹ by norm_num, inv_pow]; norm_num]
    have hpow4 : (4 : ℝ) ^ (i + 2) = 16 * 4 ^ i := by
      rw [pow_add]
      ring
    have hratio_le :
        (moserExponentSeq d p₀ i / (1 - moserExponentSeq d p₀ i)) ^ 2 ≤
          (moserExponentSeq d p₀ i / (1 - q)) ^ 2 := by
      have hratio :
          moserExponentSeq d p₀ i / (1 - moserExponentSeq d p₀ i) ≤
            moserExponentSeq d p₀ i / (1 - q) := by
        have hden_pos : 0 < 1 - moserExponentSeq d p₀ i := by linarith
        have hqden_pos : 0 < 1 - q := by linarith
        field_simp [hden_pos.ne', hqden_pos.ne']
        nlinarith [hp_i_le_q]
      have hratio_nonneg : 0 ≤ moserExponentSeq d p₀ i / (1 - moserExponentSeq d p₀ i) :=
        div_nonneg hp_i.le (by linarith)
      have hratio_rhs_nonneg : 0 ≤ moserExponentSeq d p₀ i / (1 - q) :=
        div_nonneg hp_i.le (by linarith)
      nlinarith [sq_nonneg
        (moserExponentSeq d p₀ i / (1 - q) -
          moserExponentSeq d p₀ i / (1 - moserExponentSeq d p₀ i))]
    have hseq_sq : (moserExponentSeq d p₀ i) ^ 2 = p₀ ^ 2 * (moserChi d) ^ (2 * i) := by
      rw [moserExponentSeq]
      ring
    have hcoeff_le :
        A.1.Λ * (moserExponentSeq d p₀ i / (1 - moserExponentSeq d p₀ i)) ^ 2 + 1 ≤
          (A.1.Λ * p₀ ^ 2 / (1 - q) ^ 2 + 1) * (moserChi d ^ 2) ^ i := by
      have hchi_sq_ge_one : 1 ≤ (moserChi d ^ 2) ^ i :=
        one_le_pow₀ (by nlinarith [one_lt_moserChi (d := d) hd])
      have hΛratio :
          A.1.Λ * (moserExponentSeq d p₀ i / (1 - moserExponentSeq d p₀ i)) ^ 2 ≤
            A.1.Λ * (p₀ ^ 2 / (1 - q) ^ 2) * (moserChi d ^ 2) ^ i := by
        calc A.1.Λ * (moserExponentSeq d p₀ i / (1 - moserExponentSeq d p₀ i)) ^ 2
            ≤ A.1.Λ * (moserExponentSeq d p₀ i / (1 - q)) ^ 2 :=
              mul_le_mul_of_nonneg_left hratio_le A.1.Λ_nonneg
          _ = A.1.Λ * ((moserExponentSeq d p₀ i) ^ 2 / (1 - q) ^ 2) := by
                field_simp [show (1 - q : ℝ) ≠ 0 by linarith]
          _ = A.1.Λ * (p₀ ^ 2 * (moserChi d) ^ (2 * i) / (1 - q) ^ 2) := by rw [hseq_sq]
          _ = A.1.Λ * (p₀ ^ 2 / (1 - q) ^ 2) * (moserChi d ^ 2) ^ i := by
                rw [pow_mul]
                field_simp [show (1 - q : ℝ) ≠ 0 by linarith]
      calc
        A.1.Λ * (moserExponentSeq d p₀ i / (1 - moserExponentSeq d p₀ i)) ^ 2 + 1
            ≤ A.1.Λ * (p₀ ^ 2 / (1 - q) ^ 2) * (moserChi d ^ 2) ^ i +
                (moserChi d ^ 2) ^ i := by
                  nlinarith
        _ = (A.1.Λ * p₀ ^ 2 / (1 - q) ^ 2 + 1) * (moserChi d ^ 2) ^ i := by
              ring
    rw [hgap_eq, hpow4]
    calc C_MoserAnchor d * (16 * 4 ^ i) *
            (A.1.Λ * (moserExponentSeq d p₀ i / (1 - moserExponentSeq d p₀ i)) ^ 2 + 1)
        ≤ C_MoserAnchor d * (16 * 4 ^ i) *
            ((A.1.Λ * p₀ ^ 2 / (1 - q) ^ 2 + 1) * (moserChi d ^ 2) ^ i) := by
          gcongr
          exact mul_nonneg (le_trans (by norm_num : (0 : ℝ) ≤ 1) (one_le_C_MoserAnchor (d := d)))
            (by positivity)
      _ = 16 * C_MoserAnchor d * (A.1.Λ * p₀ ^ 2 / (1 - q) ^ 2 + 1) *
            (4 * moserChi d ^ 2) ^ i := by
          rw [show (4 * moserChi d ^ 2) ^ i = 4 ^ i * (moserChi d ^ 2) ^ i by rw [mul_pow]]
          ring
  · exact div_nonneg (by norm_num) (moserExponentSeq_pos (d := d) hd hp₀ i).le

\end{lstlisting}
\begin{proof}
For $i \leq m$, $p_i = p_0\,\chi^{i} = q\,\chi^{i-m} \leq q$, hence
$1 - p_i \geq 1 - q > 0$ and
\[
  \frac{p_i}{1-p_i} \leq \frac{p_i}{1-q}, \qquad
  \Big(\frac{p_i}{1-p_i}\Big)^{2}\leq \frac{p_i^{2}}{(1-q)^{2}}
  \leq \frac{p_0^{2}\chi^{2i}}{(1-q)^{2}}.
\]
Inserting into $\mathrm{S}_i^{\mathrm{fwd}}$ and using
$C_{\mathrm{MA}}/(r_i - r_{i+1})^{2} = 16\,C_{\mathrm{MA}}\,4^{i}$ gives the bound.
\end{proof}

\begin{theorem}[Forward iteration, \leanname{supersolution\_iteration\_forward}]
\label{thm:iter-fwd}
For $d>2$, $u>0$ a supersolution, $0<q<1$, $m \in \N$, set $p_0 := q\,\chi^{-m}$.
If $|u|^{p_0}$ is integrable on $\Bz{1}$, then for every $n \leq m+1$
$|u|^{p_n}$ is integrable on $\Bz{r_n}$ and
\[
  N_n^{\mathrm{fwd}}(u)
  \;\leq\;
  \prod_{i=0}^{n-1} \mathrm{S}_i^{\mathrm{fwd}}(A,p_0)\;\cdot\;
  N_0^{\mathrm{fwd}}(u).
\]
\end{theorem}

\begin{lstlisting}[style=Matlab-editor,basicstyle=\mlttfamily,escapechar=",mlshowsectionrules=true,mathescape=true,morecomment={[s]{\%\{}{\%\}}},language=lean,numbers=none]
/-- Finite iteration of the forward low-power one-step bound.

Given target `q ∈ (0,1)` and base exponent `p₀ = q χ^{-m}`, iterate
`supersolution_preMoser_forward` `m+1` steps to go from
`L^{p₀}(B₁)` to `L^{qχ}(B_{r_{m+1}})`. The exponents `pₙ` stay below 1
since `pₘ = q < 1` and the sequence is increasing. -/
theorem supersolution_iteration_forward
    (hd : 2 < (d : ℝ))
    (A : NormalizedEllipticCoeff d (Metric.ball (0 : E) 1))
    {u : E → ℝ} {q : ℝ} {m : ℕ}
    (hq : 0 < q) (hq1 : q < 1)
    (hu_pos : ∀ x ∈ Metric.ball (0 : E) 1, 0 < u x)
    (hsuper : IsSupersolution A.1 u) :
    let p₀ := q * (moserChi d)⁻¹ ^ m
    ∀ (_hpInt :
        IntegrableOn (fun x => |u x| ^ p₀)
          (Metric.ball (0 : E) 1) volume),
    ∀ n : ℕ, n ≤ m + 1 →
      IntegrableOn (fun x => |u x| ^ moserExponentSeq d p₀ n)
          (Metric.ball (0 : E) (moserRadius n)) volume ∧
        superIterNormFwd (d := d) (u := u) p₀ n ≤
          (∏ i ∈ Finset.range n, superStepConstFwd (d := d) A p₀ i) *
            superIterNormFwd (d := d) (u := u) p₀ 0 := by
  intro p₀ hpInt n hn
  -- p₀ = q · χ⁻ᵐ > 0
  have hp₀ : 0 < p₀ := mul_pos hq (pow_pos (inv_pos.mpr (moserChi_pos hd)) m)
  -- Key: moserExponentSeq d p₀ n = q · χ^{n-m}, so for n ≤ m, p_n ≤ q < 1
  induction n with
  | zero =>
    constructor
    · rwa [moserExponentSeq_zero, moserRadius_zero]
    · simp
  | succ n ihn =>
    have hn' : n ≤ m + 1 := Nat.le_of_succ_le hn
    obtain ⟨hInt_n, hbound_n⟩ := ihn hn'
    let p_n := moserExponentSeq d p₀ n
    have hp_n : 0 < p_n := moserExponentSeq_pos hd hp₀ n
    -- Need p_n < 1 for supersolution_preMoser_forward
    -- p_n = p₀ · χⁿ = q · χ^{n-m}. For n ≤ m: p_n ≤ q < 1.
    have hp_n_lt_one : p_n < 1 := by
      -- p_n = q * χ⁻ᵐ * χⁿ = q * χ^{n-m} ≤ q * χ⁰ = q < 1 when n ≤ m
      dsimp [p_n, p₀]
      rw [moserExponentSeq, show q * (moserChi d)⁻¹ ^ m * moserChi d ^ n =
        q * ((moserChi d)⁻¹ ^ m * moserChi d ^ n) from by ring]
      have hchi_pos := moserChi_pos hd
      have : n ≤ m := Nat.lt_succ_iff.mp (Nat.lt_of_lt_of_le (Nat.lt_succ_of_le (Nat.le_of_succ_le_succ hn)) le_rfl)
      have hprod_le : (moserChi d)⁻¹ ^ m * moserChi d ^ n ≤ 1 := by
        rw [inv_pow]
        have hchi_ge_one := (one_lt_moserChi hd).le
        have hchi_pow_pos : (0 : ℝ) < moserChi d ^ m := pow_pos (moserChi_pos hd) m
        rw [show (moserChi d ^ m)⁻¹ * moserChi d ^ n =
          moserChi d ^ n / moserChi d ^ m from by rw [div_eq_mul_inv, mul_comm]]
        rw [div_le_one hchi_pow_pos]
        exact pow_le_pow_right₀ hchi_ge_one this
      calc q * ((moserChi d)⁻¹ ^ m * moserChi d ^ n) ≤ q * 1 :=
            mul_le_mul_of_nonneg_left hprod_le hq.le
        _ = q := mul_one q
        _ < 1 := hq1
    have hpre :=
      supersolution_preMoser_forward hd A (p := p_n)
        (r := moserRadius (n + 1)) (s := moserRadius n)
        hp_n hp_n_lt_one (moserRadius_pos (n + 1)) (moserRadius_succ_lt n)
        (moserRadius_le_one n) hu_pos hsuper
        (by simpa [p_n] using hInt_n)
    obtain ⟨hInt_succ, hNorm_succ⟩ := hpre
    refine ⟨?_, ?_⟩
    · have heq : moserChi d * p_n = moserExponentSeq d p₀ (n + 1) := by
        rw [moserExponentSeq_succ]
      rwa [heq] at hInt_succ
    · have heq_norm : superIterNormFwd (d := d) (u := u) p₀ (n + 1) =
          (∫ x in Metric.ball (0 : E) (moserRadius (n + 1)),
            |u x| ^ (moserChi d * p_n) ∂volume) ^
              (1 / (moserChi d * p_n)) := by
        simp [superIterNormFwd, superIterIntegralFwd, p_n, moserExponentSeq_succ]
      have heq_step : superStepConstFwd (d := d) A p₀ n =
          ((C_MoserAnchor d / (moserRadius n - moserRadius (n + 1)) ^ 2) *
            (A.1.Λ * (p_n / (1 - p_n)) ^ 2 + 1)) ^ (1 / p_n) := by
        simp [superStepConstFwd, p_n]
      have hstep_nonneg : 0 ≤ superStepConstFwd (d := d) A p₀ n := by
        rw [heq_step]
        apply Real.rpow_nonneg
        apply mul_nonneg
        · exact div_nonneg
            (le_trans (by norm_num : (0 : ℝ) ≤ 1) (one_le_C_MoserAnchor (d := d)))
            (sq_nonneg _)
        · nlinarith [A.1.Λ_nonneg, sq_nonneg (p_n / (1 - p_n))]
      have hstep_bound :
          superIterNormFwd (d := d) (u := u) p₀ (n + 1) ≤
            superStepConstFwd (d := d) A p₀ n *
              superIterNormFwd (d := d) (u := u) p₀ n := by
        rw [heq_norm, heq_step]
        convert hNorm_succ using 2
      calc superIterNormFwd (d := d) (u := u) p₀ (n + 1)
          ≤ superStepConstFwd (d := d) A p₀ n *
              superIterNormFwd (d := d) (u := u) p₀ n := hstep_bound
        _ ≤ superStepConstFwd (d := d) A p₀ n *
              ((∏ i ∈ Finset.range n, superStepConstFwd (d := d) A p₀ i) *
                superIterNormFwd (d := d) (u := u) p₀ 0) := by
            exact mul_le_mul_of_nonneg_left hbound_n hstep_nonneg
        _ = (∏ i ∈ Finset.range (n + 1), superStepConstFwd (d := d) A p₀ i) *
              superIterNormFwd (d := d) (u := u) p₀ 0 := by
            rw [Finset.prod_range_succ]
            ring

\end{lstlisting}
\begin{proof}
Induction on $n \leq m+1$. The base case is trivial. For the step, note that
$p_n = p_0\,\chi^{n} = q\,\chi^{n-m}$, so for $n \leq m$ we have $p_n \leq q < 1$
(this is what allows the application of Theorem~\ref{thm:preMoser-fwd}). The
one-step bound gives $N_{n+1}^{\mathrm{fwd}}(u) \leq
\mathrm{S}_n^{\mathrm{fwd}}(A,p_0)\,N_n^{\mathrm{fwd}}(u)$, and combining with
the inductive hypothesis yields the displayed product bound.
\end{proof}

\begin{lemma}[Forward geometric majorant, \leanname{supersolution\_geometric\_majorant\_fwd}]
\label{lem:geom-maj-fwd}
With the notation above, for any $n \leq m+1$,
\[
  \prod_{i=0}^{n-1}\mathrm{S}_i^{\mathrm{fwd}}(A,p_0)
  \;\leq\;
  \Big(C_{\mathrm{wH}}^{0}(d)\,
    \Big(\frac{\Lambda p_0^{2}}{(1-q)^{2}}+1\Big)^{d/2}\Big)^{1/p_0}.
\]
\end{lemma}

\begin{lstlisting}[style=Matlab-editor,basicstyle=\mlttfamily,escapechar=",mlshowsectionrules=true,mathescape=true,morecomment={[s]{\%\{}{\%\}}},language=lean,numbers=none]
/-- The finite product of step constants in the forward iteration is bounded.

The key bound is:
```
∏ᵢ₌₀ᵐ step_i ≤ (C_Moser d * (Λp₀²/(1-q)² + 1)^{d/2})^{1/p₀}
```
The crucial point: `1 - pₙ ≥ 1 - q > 0` for `n ≤ m` (since `pₘ = q`
is the largest exponent), so the `(1-p)` denominators stay bounded. -/
theorem supersolution_geometric_majorant_fwd
    (hd : 2 < (d : ℝ))
    (A : NormalizedEllipticCoeff d (Metric.ball (0 : E) 1))
    {q : ℝ} {m : ℕ}
    (hq : 0 < q) (hq1 : q < 1) :
    let p₀ := q * (moserChi d)⁻¹ ^ m
    ∏ i ∈ Finset.range (m + 1), superStepConstFwd (d := d) A p₀ i ≤
      (C_weakHarnack0 d * (A.1.Λ * p₀ ^ 2 / (1 - q) ^ 2 + 1) ^
        ((d : ℝ) / 2)) ^ (1 / p₀) := by
  intro p₀
  have hp₀ : 0 < p₀ := by
    dsimp [p₀]
    exact mul_pos hq (pow_pos (inv_pos.mpr (moserChi_pos (d := d) hd)) m)
  let ρ := moserDecayRatio d
  let K₀ : ℝ := 16 * C_MoserAnchor d
  let L : ℝ := A.1.Λ * p₀ ^ 2 / (1 - q) ^ 2 + 1
  let K : ℝ := K₀ * L
  let CC : ℝ := 4 * moserChi d ^ 2
  let S : ℝ := Finset.sum (Finset.range (m + 1)) (fun i => ρ ^ i)
  let T : ℝ := Finset.sum (Finset.range (m + 1)) (fun i => (i : ℝ) * ρ ^ i)
  have hρ_nonneg : 0 ≤ ρ := by
    simpa [ρ] using moserDecayRatio_nonneg (d := d) hd
  have hρ_lt_one : ρ < 1 := by
    simpa [ρ] using moserDecayRatio_lt_one (d := d) hd
  have hS_le : S ≤ (d : ℝ) / 2 := by
    simpa [ρ, S] using moserDecayRatio_sum_le_half_dim (d := d) hd (m + 1)
  have hT_le : T ≤ ∑' i : ℕ, (i : ℝ) * ρ ^ i := by
    simpa [ρ, T] using moserDecayRatio_weighted_sum_le_tsum (d := d) hd (m + 1)
  have hK₀_ge_one : 1 ≤ K₀ := by
    dsimp [K₀]
    nlinarith [one_le_C_MoserAnchor (d := d)]
  have hK₀_nonneg : 0 ≤ K₀ := by linarith
  have hL_ge_one : 1 ≤ L := by
    dsimp [L]
    have hterm_nonneg : 0 ≤ A.1.Λ * p₀ ^ 2 / (1 - q) ^ 2 := by
      exact div_nonneg (mul_nonneg A.1.Λ_nonneg (sq_nonneg p₀)) (sq_nonneg (1 - q))
    linarith
  have hL_nonneg : 0 ≤ L := by linarith
  have hK_pos : 0 < K := by
    dsimp [K]
    exact mul_pos (lt_of_lt_of_le zero_lt_one hK₀_ge_one) (by positivity)
  have hK_nonneg : 0 ≤ K := hK_pos.le
  have hCC_pos : 0 < CC := by
    dsimp [CC]
    exact mul_pos (by norm_num : (0 : ℝ) < 4) (sq_pos_of_pos (moserChi_pos (d := d) hd))
  have hCC_nonneg : 0 ≤ CC := hCC_pos.le
  have hstep_le : ∀ i, i ∈ Finset.range (m + 1) →
      superStepConstFwd (d := d) A p₀ i ≤ (K * CC ^ i) ^
        (1 / moserExponentSeq d p₀ i) := by
    intro i hi
    simpa [p₀, K, K₀, L, CC] using superStepConstFwd_le (d := d) hd A hq hq1 hi
  have hstep_nonneg : ∀ i, i ∈ Finset.range (m + 1) →
      0 ≤ superStepConstFwd (d := d) A p₀ i := by
    intro i _
    dsimp [superStepConstFwd]
    apply Real.rpow_nonneg
    apply mul_nonneg
    · exact div_nonneg (le_trans (by norm_num : (0 : ℝ) ≤ 1) (one_le_C_MoserAnchor (d := d)))
        (sq_nonneg _)
    · nlinarith [A.1.Λ_nonneg, sq_nonneg (moserExponentSeq d p₀ i / (1 - moserExponentSeq d p₀ i))]
  have hprod_le :
      ∏ i ∈ Finset.range (m + 1), superStepConstFwd (d := d) A p₀ i ≤
        ∏ i ∈ Finset.range (m + 1), (K * CC ^ i) ^
          (1 / moserExponentSeq d p₀ i) :=
    Finset.prod_le_prod hstep_nonneg hstep_le
  have hprod_eval :
      ∀ n : ℕ,
        ∏ i ∈ Finset.range n, (K * CC ^ i) ^
            (1 / moserExponentSeq d p₀ i) =
          (K ^ (Finset.sum (Finset.range n) (fun i => ρ ^ i)) *
            CC ^ (Finset.sum (Finset.range n) (fun i => (i : ℝ) * ρ ^ i))) ^
            (1 / p₀) := by
    intro n
    induction n with
    | zero =>
        simp
    | succ n ihn =>
        let Sₙ : ℝ := Finset.sum (Finset.range n) (fun i => ρ ^ i)
        let Tₙ : ℝ := Finset.sum (Finset.range n) (fun i => (i : ℝ) * ρ ^ i)
        have hinv :
            1 / moserExponentSeq d p₀ n = ρ ^ n / p₀ := by
          simpa [ρ] using inv_moserExponentSeq (d := d) hd hp₀ n
        have hdiv_eq :
            1 / moserExponentSeq d p₀ n = ρ ^ n * (1 / p₀) := by
          rw [hinv]
          ring
        rw [Finset.prod_range_succ, ihn, hdiv_eq]
        have hKC_nonneg : 0 ≤ K * CC ^ n := by
          exact mul_nonneg hK_nonneg (pow_nonneg hCC_nonneg n)
        rw [Real.rpow_mul hKC_nonneg]
        rw [Real.mul_rpow hK_nonneg (pow_nonneg hCC_nonneg n)]
        have hCC_pow :
            ((CC ^ n : ℝ)) ^ (ρ ^ n) = CC ^ ((n : ℝ) * ρ ^ n) := by
          rw [show ((CC ^ n : ℝ)) = CC ^ (n : ℝ) by rw [Real.rpow_natCast]]
          rw [← Real.rpow_mul hCC_nonneg]
        rw [hCC_pow]
        have hfront_nonneg :
            0 ≤ K ^ (ρ ^ n) * CC ^ ((n : ℝ) * ρ ^ n) := by
          exact mul_nonneg (Real.rpow_nonneg hK_nonneg _) (Real.rpow_nonneg hCC_nonneg _)
        have hback_nonneg :
            0 ≤ K ^ Sₙ * CC ^ Tₙ := by
          exact mul_nonneg (Real.rpow_nonneg hK_nonneg _) (Real.rpow_nonneg hCC_nonneg _)
        rw [mul_comm, ← Real.mul_rpow hfront_nonneg hback_nonneg]
        congr 1
        calc
          (K ^ (ρ ^ n) * CC ^ ((n : ℝ) * ρ ^ n)) * (K ^ Sₙ * CC ^ Tₙ)
              = (K ^ (ρ ^ n) * K ^ Sₙ) * (CC ^ ((n : ℝ) * ρ ^ n) * CC ^ Tₙ) := by
                  ring
          _ = K ^ (ρ ^ n + Sₙ) * CC ^ (((n : ℝ) * ρ ^ n) + Tₙ) := by
                rw [← Real.rpow_add hK_pos, ← Real.rpow_add hCC_pos]
          _ =
              K ^ (Finset.sum (Finset.range (n + 1)) (fun i => ρ ^ i)) *
                CC ^ (Finset.sum (Finset.range (n + 1))
                  (fun i => (i : ℝ) * ρ ^ i)) := by
                simp [Sₙ, Tₙ, Finset.sum_range_succ, add_comm]
  have hK_pow_eq :
      K ^ S = K₀ ^ S * L ^ S := by
    rw [show K = K₀ * L by rfl, Real.mul_rpow hK₀_nonneg hL_nonneg]
  have hK₀_pow_le : K₀ ^ S ≤ K₀ ^ ((d : ℝ) / 2) := by
    exact Real.rpow_le_rpow_of_exponent_le hK₀_ge_one hS_le
  have hL_pow_le : L ^ S ≤ L ^ ((d : ℝ) / 2) := by
    exact Real.rpow_le_rpow_of_exponent_le hL_ge_one hS_le
  have hK_le :
      K ^ S ≤ K₀ ^ ((d : ℝ) / 2) * L ^ ((d : ℝ) / 2) := by
    rw [hK_pow_eq]
    exact mul_le_mul hK₀_pow_le hL_pow_le
      (Real.rpow_nonneg hL_nonneg _) (Real.rpow_nonneg hK₀_nonneg _)
  have h4_le :
      (4 : ℝ) ^ T ≤ 4 ^ (∑' i : ℕ, (i : ℝ) * ρ ^ i) := by
    exact Real.rpow_le_rpow_of_exponent_le (by norm_num : (1 : ℝ) ≤ 4) hT_le
  have hchi_sq_ge_one : 1 ≤ moserChi d ^ 2 := by
    exact one_le_pow₀ (one_le_moserChi (d := d) hd)
  have hchi_le :
      (moserChi d ^ 2) ^ T ≤ (moserChi d ^ 2) ^ (∑' i : ℕ, (i : ℝ) * ρ ^ i) := by
    exact Real.rpow_le_rpow_of_exponent_le hchi_sq_ge_one hT_le
  have hCC_le :
      CC ^ T ≤
        4 ^ (∑' i : ℕ, (i : ℝ) * ρ ^ i) *
          (moserChi d ^ 2) ^ (∑' i : ℕ, (i : ℝ) * ρ ^ i) := by
    calc
      CC ^ T = (4 : ℝ) ^ T * (moserChi d ^ 2) ^ T := by
        rw [show CC = 4 * (moserChi d ^ 2) by simp [CC]]
        rw [Real.mul_rpow (by positivity) (sq_nonneg _)]
      _ ≤ 4 ^ (∑' i : ℕ, (i : ℝ) * ρ ^ i) *
            (moserChi d ^ 2) ^ (∑' i : ℕ, (i : ℝ) * ρ ^ i) := by
          exact mul_le_mul h4_le hchi_le
            (Real.rpow_nonneg (sq_nonneg _) _)
            (Real.rpow_nonneg (by norm_num : (0 : ℝ) ≤ 4) _)
  have hK₀_pow_le_big :
      K₀ ^ ((d : ℝ) / 2) ≤ ((32 : ℝ) * C_MoserAnchor d) ^ ((d : ℝ) / 2) := by
    exact Real.rpow_le_rpow
      hK₀_nonneg
      (by
        dsimp [K₀]
        nlinarith [one_le_C_MoserAnchor (d := d)])
      (by positivity)
  have hCMoser_le :
      K₀ ^ ((d : ℝ) / 2) * 4 ^ (∑' i : ℕ, (i : ℝ) * ρ ^ i) ≤ C_Moser d := by
    have hbase :
        K₀ ^ ((d : ℝ) / 2) * 4 ^ (∑' i : ℕ, (i : ℝ) * ρ ^ i) ≤
          ((32 : ℝ) * C_MoserAnchor d) ^ ((d : ℝ) / 2) *
            4 ^ (∑' i : ℕ, (i : ℝ) * ρ ^ i) := by
      exact mul_le_mul_of_nonneg_right hK₀_pow_le_big
        (Real.rpow_nonneg (by norm_num : (0 : ℝ) ≤ 4) _)
    have hmax :
        ((32 : ℝ) * C_MoserAnchor d) ^ ((d : ℝ) / 2) *
            4 ^ (∑' i : ℕ, (i : ℝ) * ρ ^ i) ≤
          C_Moser d := by
      dsimp [C_Moser, ρ]
      split_ifs with hlt
      · exact
          le_max_right (C_MoserAnchor d)
            (((32 : ℝ) * C_MoserAnchor d) ^ ((d : ℝ) / 2) *
              4 ^ (∑' i : ℕ, (i : ℝ) * moserDecayRatio d ^ i))
      · exact False.elim (hlt hd)
    exact hbase.trans hmax
  have hCweak_eq :
      C_Moser d * (moserChi d ^ 2) ^ (∑' i : ℕ, (i : ℝ) * ρ ^ i) =
        C_weakHarnack0 d := by
    simp [C_weakHarnack0, hd, ρ]
  have hbody_le :
      K ^ S * CC ^ T ≤
        C_weakHarnack0 d * L ^ ((d : ℝ) / 2) := by
    have hK_target_nonneg :
        0 ≤ K₀ ^ ((d : ℝ) / 2) * L ^ ((d : ℝ) / 2) := by
      exact mul_nonneg (Real.rpow_nonneg hK₀_nonneg _) (Real.rpow_nonneg hL_nonneg _)
    calc
      K ^ S * CC ^ T
          ≤ (K₀ ^ ((d : ℝ) / 2) * L ^ ((d : ℝ) / 2)) *
              (4 ^ (∑' i : ℕ, (i : ℝ) * ρ ^ i) *
                (moserChi d ^ 2) ^ (∑' i : ℕ, (i : ℝ) * ρ ^ i)) := by
            exact mul_le_mul hK_le hCC_le
              (Real.rpow_nonneg hCC_nonneg _)
              hK_target_nonneg
      _ = (K₀ ^ ((d : ℝ) / 2) * 4 ^ (∑' i : ℕ, (i : ℝ) * ρ ^ i)) *
            ((moserChi d ^ 2) ^ (∑' i : ℕ, (i : ℝ) * ρ ^ i) * L ^ ((d : ℝ) / 2)) := by
              ring
      _ ≤ C_Moser d *
            ((moserChi d ^ 2) ^ (∑' i : ℕ, (i : ℝ) * ρ ^ i) * L ^ ((d : ℝ) / 2)) := by
              exact mul_le_mul_of_nonneg_right hCMoser_le
                (mul_nonneg (Real.rpow_nonneg (sq_nonneg _) _) (Real.rpow_nonneg hL_nonneg _))
      _ = C_weakHarnack0 d * L ^ ((d : ℝ) / 2) := by
            rw [← hCweak_eq]
            ring
  calc
    ∏ i ∈ Finset.range (m + 1), superStepConstFwd (d := d) A p₀ i
        ≤ ∏ i ∈ Finset.range (m + 1), (K * CC ^ i) ^ (1 / moserExponentSeq d p₀ i) := hprod_le
    _ = (K ^ S * CC ^ T) ^ (1 / p₀) := by
          simpa [ρ, S, T] using hprod_eval (m + 1)
    _ ≤ (C_weakHarnack0 d * L ^ ((d : ℝ) / 2)) ^ (1 / p₀) := by
          exact Real.rpow_le_rpow
            (mul_nonneg (Real.rpow_nonneg hK_nonneg _) (Real.rpow_nonneg hCC_nonneg _))
            hbody_le (by positivity)
    _ = (C_weakHarnack0 d * (A.1.Λ * p₀ ^ 2 / (1 - q) ^ 2 + 1) ^ ((d : ℝ) / 2)) ^ (1 / p₀) := by
          simp [L]



\end{lstlisting}
\begin{proof}
Identical to the proof of Lemma~\ref{lem:geom-maj-inv}, with $\Lambda p_0^{2}+1$
replaced by $\Lambda p_0^{2}/(1-q)^{2}+1$ throughout.
\end{proof}

\begin{definition}[Forward closeout constant, \leanname{C\_weakHarnack0Forward}]
For $d>2$,
\[
  C_{\mathrm{wH,fwd}}^{0}(d) :=
  \big(C_{\mathrm{wH}}^{0}(d)\,(|\Bz{1}|+1)\big)^{\chi}.
\]
\end{definition}

\begin{lstlisting}[style=Matlab-editor,basicstyle=\mlttfamily,escapechar=",mlshowsectionrules=true,mathescape=true,morecomment={[s]{\%\{}{\%\}}},language=lean,numbers=none]
/-- The weakened forward-stage constant that honestly absorbs the finite-measure
`L^p` to `L^{p₀}` upgrade on the unit ball. -/
noncomputable def C_weakHarnack0Forward
    (d : ℕ) [NeZero d] (_hd : 2 < (d : ℝ)) : ℝ :=
  (C_weakHarnack0 d * (volume.real (Metric.ball (0 : AmbientSpace d) 1) + 1)) ^
    (moserChi d)

\end{lstlisting}
\begin{theorem}[Stage one weak Harnack (forward), \leanname{weak\_harnack\_stage\_one\_forward}]
\label{thm:wH-stage1-fwd}
Let $d>2$, $A$ on $\Bz{1}$, $u>0$ a supersolution on $\Bz{1}$, and
$0 < p < q < 1$. Then
\begin{multline*}
  \Big(\int_{\Bz{1/2}} |u|^{q d/(d-2)}\,dx\Big)^{p(d-2)/(qd)}
  \\
  \;\leq\;
  C_{\mathrm{wH,fwd}}^{0}(d)\,
  \Big(\frac{\Lambda p^{2}}{(1-q)^{2}}+1\Big)^{d\chi/2}
  \int_{\Bz{1}} |u|^{p}\,dx.
\end{multline*}
\end{theorem}

\begin{lstlisting}[style=Matlab-editor,basicstyle=\mlttfamily,escapechar=",mlshowsectionrules=true,mathescape=true,morecomment={[s]{\%\{}{\%\}}},language=lean,numbers=none]
/-- Second stage of weak Harnack: forward low-power iteration for positive
supersolutions.

For `u > 0` a supersolution on `B₁` and `0 < p < q < 1`:
```
‖u‖_{L^{qd/(d-2)}(B_{1/2})}^p ≤
  C_fwd(d) (Λp²/(1-q)² + 1)^{dχ/2} ‖u‖_{Lᵖ(B₁)}^p
```
The theorem is stated for arbitrary `0 < p < q < 1`, with the nonlinear factor
raised to `dχ/2`.

Proof: choose `m` large enough that `p₀ = q χ^{-m} ≤ p`, then apply the
finite forward iteration to get `‖u‖_{L^{qχ}(B_{r_{m+1}})} ≤ C · ‖u‖_{L^{p₀}(B₁)}`,
use `B_{1/2} ⊂ B_{r_{m+1}}` and `p₀ ≤ p` to convert, raise to the `p`-th power,
and apply the geometric majorant bound. -/
theorem weak_harnack_stage_one_forward
    (hd : 2 < (d : ℝ))
    (A : NormalizedEllipticCoeff d (Metric.ball (0 : E) 1))
    {u : E → ℝ} {p q : ℝ}
    (hp : 0 < p) (hp1 : p < 1) (hpq : p < q) (hq1 : q < 1)
    (hu_pos : ∀ x ∈ Metric.ball (0 : E) 1, 0 < u x)
    (hsuper : IsSupersolution A.1 u) :
    (∫ x in Metric.ball (0 : E) (1 / 2 : ℝ),
        |u x| ^ (q * (d : ℝ) / ((d : ℝ) - 2)) ∂volume) ^
          (p * (((d : ℝ) - 2) / (q * (d : ℝ)))) ≤
      C_weakHarnack0Forward (d := d) hd *
        (A.1.Λ * p ^ 2 / (1 - q) ^ 2 + 1) ^ (((d : ℝ) * moserChi d) / 2) *
        ∫ x in Metric.ball (0 : E) 1, |u x| ^ p ∂volume := by
  have hq : 0 < q := lt_trans hp hpq
  let qχ : ℝ := q * moserChi d
  have hqχ_pos : 0 < qχ := mul_pos hq (moserChi_pos (d := d) hd)
  have hqχ_ne : qχ ≠ 0 := hqχ_pos.ne'
  rcases exists_forward_iteration_depth (d := d) hd hp hpq with
    ⟨m, hm_lt, hm_le⟩
  let p₀ : ℝ := q * (moserChi d)⁻¹ ^ m
  have hp₀ : 0 < p₀ := by
    dsimp [p₀]
    exact mul_pos hq (pow_pos (inv_pos.mpr (moserChi_pos (d := d) hd)) m)
  have hp₀_lt_p : p₀ < p := by
    simpa [p₀] using hm_lt
  have hp_le_chi_p₀ : p ≤ moserChi d * p₀ := by
    simpa [p₀] using hm_le
  have hp₀_le_two : p₀ ≤ 2 := by
    have hp₀_lt_one : p₀ < 1 := lt_trans hp₀_lt_p hp1
    linarith
  have hp₀Int :
      IntegrableOn (fun x => |u x| ^ p₀)
        (Metric.ball (0 : E) 1) volume := by
    exact supersolution_integrableOn_ball_one_rpow
      (d := d) A hp₀ hp₀_le_two hsuper
  have hiter :=
    supersolution_iteration_forward (d := d) hd A (u := u) (q := q) (m := m)
      hq hq1 hu_pos hsuper
  have hstep := hiter hp₀Int (m + 1) le_rfl
  have hgeom :=
    supersolution_geometric_majorant_fwd (d := d) hd A (q := q) (m := m) hq hq1
  have htarget_int :
      IntegrableOn (fun x => |u x| ^ qχ)
        (Metric.ball (0 : E) (moserRadius (m + 1))) volume := by
    have hstep_int := hstep.1
    rw [moserExponentSeq_forward_target_eq (d := d) hd (q := q) (m := m)] at hstep_int
    simpa [qχ] using hstep_int
  have htarget_norm :
      (∫ x in Metric.ball (0 : E) (moserRadius (m + 1)),
          |u x| ^ qχ ∂volume) ^ (1 / qχ) ≤
        ((C_weakHarnack0 d *
              (A.1.Λ * p₀ ^ 2 / (1 - q) ^ 2 + 1) ^ ((d : ℝ) / 2)) ^
            (1 / p₀)) *
          (∫ x in Metric.ball (0 : E) 1, |u x| ^ p₀ ∂volume) ^ (1 / p₀) := by
    have hbound_step :
        superIterNormFwd (d := d) (u := u) p₀ (m + 1) ≤
          (∏ i ∈ Finset.range (m + 1), superStepConstFwd (d := d) A p₀ i) *
            superIterNormFwd (d := d) (u := u) p₀ 0 := by
      simpa [p₀] using hstep.2
    calc
      (∫ x in Metric.ball (0 : E) (moserRadius (m + 1)),
          |u x| ^ qχ ∂volume) ^ (1 / qχ)
          = superIterNormFwd (d := d) (u := u) p₀ (m + 1) := by
              simpa [qχ, p₀] using
                (superIterNormFwd_target (d := d) hd (u := u) (q := q) (m := m)).symm
      _ ≤ (∏ i ∈ Finset.range (m + 1), superStepConstFwd (d := d) A p₀ i) *
            superIterNormFwd (d := d) (u := u) p₀ 0 := by
              exact hbound_step
      _ ≤ ((C_weakHarnack0 d *
              (A.1.Λ * p₀ ^ 2 / (1 - q) ^ 2 + 1) ^ ((d : ℝ) / 2)) ^
            (1 / p₀)) *
          superIterNormFwd (d := d) (u := u) p₀ 0 := by
            exact mul_le_mul_of_nonneg_right (by simpa [p₀] using hgeom)
              (Real.rpow_nonneg (integral_nonneg fun x => by positivity) _)
      _ = ((C_weakHarnack0 d *
              (A.1.Λ * p₀ ^ 2 / (1 - q) ^ 2 + 1) ^ ((d : ℝ) / 2)) ^
            (1 / p₀)) *
          (∫ x in Metric.ball (0 : E) 1, |u x| ^ p₀ ∂volume) ^ (1 / p₀) := by
            rw [superIterNormFwd_zero]
  have hhalf_sub :
      Metric.ball (0 : E) (1 / 2 : ℝ) ⊆ Metric.ball (0 : E) (moserRadius (m + 1)) := by
    exact Metric.ball_subset_ball (le_of_lt (moserRadius_gt_half (m + 1)))
  have hhalf_int :
      IntegrableOn (fun x => |u x| ^ qχ)
        (Metric.ball (0 : E) (1 / 2 : ℝ)) volume := by
    exact htarget_int.mono_set hhalf_sub
  have hhalf_integral_le :
      ∫ x in Metric.ball (0 : E) (1 / 2 : ℝ), |u x| ^ qχ ∂volume ≤
        ∫ x in Metric.ball (0 : E) (moserRadius (m + 1)), |u x| ^ qχ ∂volume := by
    exact setIntegral_mono_set htarget_int
      (ae_of_all _ (by intro x; positivity))
      (Filter.Eventually.of_forall hhalf_sub)
  have hhalf_norm_le :
      (∫ x in Metric.ball (0 : E) (1 / 2 : ℝ), |u x| ^ qχ ∂volume) ^ (1 / qχ) ≤
        ((C_weakHarnack0 d *
              (A.1.Λ * p₀ ^ 2 / (1 - q) ^ 2 + 1) ^ ((d : ℝ) / 2)) ^
            (1 / p₀)) *
          (∫ x in Metric.ball (0 : E) 1, |u x| ^ p₀ ∂volume) ^ (1 / p₀) := by
    exact
      (Real.rpow_le_rpow (by positivity) hhalf_integral_le (by positivity)).trans
        htarget_norm
  have hmain_pow :
      (∫ x in Metric.ball (0 : E) (1 / 2 : ℝ), |u x| ^ qχ ∂volume) ^ (p / qχ) ≤
        weakHarnackForwardUpgradeLHS (d := d) A u p q p₀ := by
    calc
      (∫ x in Metric.ball (0 : E) (1 / 2 : ℝ), |u x| ^ qχ ∂volume) ^ (p / qχ)
          = ((∫ x in Metric.ball (0 : E) (1 / 2 : ℝ), |u x| ^ qχ ∂volume) ^
              (1 / qχ)) ^ p := by
                rw [← Real.rpow_mul (integral_nonneg fun x => by positivity)]
                field_simp [hqχ_ne]
      _ ≤ weakHarnackForwardUpgradeLHS (d := d) A u p q p₀ := by
            exact Real.rpow_le_rpow (by positivity) hhalf_norm_le hp.le
  have hupgrade :=
    weak_harnack_stage_one_forward_power_upgrade
      (d := d) hd A (u := u) hp₀ hp hp1 hp₀_lt_p hp_le_chi_p₀ hq1 hsuper hp₀Int
  have hqχ_eq : q * (d : ℝ) / ((d : ℝ) - 2) = qχ := by
    dsimp [qχ]
    rw [moserChi]
    ring
  have hexp_eq : p * (((d : ℝ) - 2) / (q * (d : ℝ))) = p / qχ := by
    dsimp [qχ]
    rw [moserChi]
    field_simp [hq.ne']
  calc
    (∫ x in Metric.ball (0 : E) (1 / 2 : ℝ),
        |u x| ^ (q * (d : ℝ) / ((d : ℝ) - 2)) ∂volume) ^
          (p * (((d : ℝ) - 2) / (q * (d : ℝ))))
        = (∫ x in Metric.ball (0 : E) (1 / 2 : ℝ), |u x| ^ qχ ∂volume) ^
            (p / qχ) := by
              rw [hqχ_eq, hexp_eq]
    _ ≤ weakHarnackForwardUpgradeLHS (d := d) A u p q p₀ := hmain_pow
    _ ≤ weakHarnackForwardUpgradeRHS (d := d) hd A u p q := hupgrade
    _ = C_weakHarnack0Forward (d := d) hd *
          (A.1.Λ * p ^ 2 / (1 - q) ^ 2 + 1) ^ (((d : ℝ) * moserChi d) / 2) *
          ∫ x in Metric.ball (0 : E) 1, |u x| ^ p ∂volume := rfl

\end{lstlisting}
\begin{proof}
Set $\chi := d/(d-2)$, $q\chi := q d/(d-2)$. By Lemma~\ref{lem:fwd-depth} pick
$m$ with $p_0 := q\chi^{-m} < p$ and $p \leq \chi p_0$, and note $p_0 < p < 1$.
Theorem~\ref{thm:iter-fwd} with this $p_0$ at depth $n = m+1$ gives the iterated
$L^{p_{m+1}} = L^{q\chi}$ bound on $\Bz{r_{m+1}}$, namely
\begin{multline*}
  \Big(\int_{\Bz{r_{m+1}}} |u|^{q\chi}\,dx\Big)^{1/(q\chi)}
  \\ \leq\;
  \prod_{i=0}^{m}\mathrm{S}_i^{\mathrm{fwd}}(A,p_0)\;\cdot\;
  \Big(\int_{\Bz{1}} |u|^{p_0}\,dx\Big)^{1/p_0}.
\end{multline*}
By Lemma~\ref{lem:geom-maj-fwd}, the product is at most
\[
  \Big(C_{\mathrm{wH}}^{0}(d)\,
    \big(\Lambda p_0^{2}/(1-q)^{2}+1\big)^{d/2}\Big)^{1/p_0}.
\]
Use $\Bz{1/2} \subseteq \Bz{r_{m+1}}$ to restrict to $\Bz{1/2}$, raise the
inequality to the $p$-th power, and absorb the change of base from $p_0$ to $p$
via the volume comparison
\leanname{weak\_harnack\_stage\_one\_forward\_power\_upgrade}: by Holder with
exponents $p_0 < p \leq \chi p_0$ on the unit ball,
\[
  \Big(\int_{\Bz{1}} |u|^{p_0}\,dx\Big)^{1/p_0}
  \leq
  \Big(\int_{\Bz{1}} |u|^{p}\,dx\Big)^{1/p}\,
  |\Bz{1}|^{1/p_0 - 1/p},
\]
and consequently
\[
  \Big(\int_{\Bz{1}} |u|^{p_0}\,dx\Big)^{p/p_0}
  \leq (|\Bz{1}|+1)^{\chi}\,\int_{\Bz{1}} |u|^{p}\,dx,
\]
which absorbs the volume into $C_{\mathrm{wH,fwd}}^{0}(d)$. Together with
$(\Lambda p_0^{2}/(1-q)^{2}+1)^{d/(2)\cdot p/p_0} \leq
(\Lambda p^{2}/(1-q)^{2}+1)^{d\chi/2}$
(monotonicity in both base and exponent, since $p/p_0 \leq \chi$ and
$p_0^{2} \leq p^{2}$), the displayed inequality follows. The exponent
$p(d-2)/(qd) = p/(q\chi)$ comes from raising the $L^{q\chi}$-norm to the
$p$-th power.
\end{proof}

\begin{remark}[Reverse-Holder reformulation]
Theorems~\ref{thm:wH-stage1-inv} and~\ref{thm:wH-stage1-fwd} are precisely the
two halves of the weak Harnack inequality at this stage of the development:
the first asserts a lower bound on $u$ in $\Bz{1/2}$ in terms of an inverse
integral on $\Bz{1}$, while the second is a forward $L^{q\chi}$--$L^{p}$
reverse Holder inequality on the same pair of balls. The crossover between
the two stages, completing the weak Harnack inequality, is performed in the
file \texttt{WeakHarnack.lean} using the constants $C_{\mathrm{wH}}^{0}(d)$
and $C_{\mathrm{wH,fwd}}^{0}(d)$ assembled here.
\end{remark}

%% ========================================================================
\section{Crossover, Weak Harnack, Harnack, and H\"older Regularity}\label{sec:crossover}
%% ========================================================================

%% --- Section 9 ---
\subsection{Crossover Estimate}

Let $d \in \N$, $d \geq 3$, and $E = \R^d$ throughout. Write $\Bz r$ for the
open Euclidean ball of radius $r$ centred at the origin. Given a normalized
elliptic coefficient $A$ on $\Bz 1$ with ellipticity constants $\lambda_A,
\Lambda_A$ (with $\lambda_A = 1$ after normalization, so $\Lambda_A \geq 1$),
write $\IsSupersolution{A}{u}$ for the assertion that $u \in W^{1,2}_{\loc}$
satisfies $-\dv(A\nabla u) \geq 0$ in the weak sense, and $\IsSolution{A}{u}$
for the corresponding equality.

The crossover estimate, due in this form to Moser, bridges positive and
negative powers of a positive supersolution. Throughout this section we work
with a fixed positive constant
\[
  \Cpoinc = \Cpoinc(d) > 0
\]
denoting the dimensional Poincar\'e constant on the unit ball, a fixed
constant $\Mst > 1$ used in the smooth cutoff construction, and the
dimensional John--Nirenberg constant $\CJN(d) > 0$.

\begin{definition}[\leanname{c\_crossover\_bmo\_scale}, \leanname{c\_crossover'}]
We define the BMO scale and crossover exponent
\begin{align*}
  c_{\BMO}(d) &:= \bigl(\vol(\Bz 1)\bigr)^{-1/2}\,\Cpoinc(d) \cdot
    \tfrac12^{1 - d/2} \cdot 8\Mst \,\bigl(\vol(\Bz 1)\bigr)^{1/2}, \\
  c'(d)       &:= \frac{1}{2\,\CJN(d)\,c_{\BMO}(d) + 1}.
\end{align*}
\end{definition}

\begin{lstlisting}[style=Matlab-editor,basicstyle=\mlttfamily,escapechar=",mlshowsectionrules=true,mathescape=true,morecomment={[s]{\%\{}{\%\}}},language=lean,numbers=none]
/-! ### Constants for the crossover estimate -/

/-- The small exponent `c(d)` in the crossover estimate.
This is chosen small enough for the local John-Nirenberg step fed by the
regularized-log BMO bound. -/
noncomputable def c_crossover_bmo_scale (d : ℕ) [NeZero d] : ℝ :=
  ((volume.real (Metric.ball (0 : EuclideanSpace ℝ (Fin d)) 1)) ^ (-(1 / 2 : ℝ)) * C_poinc_val d) *
    (1 / 2 : ℝ) ^ (1 - (d : ℝ) / 2) *
    (8 * (Mst : ℝ) * (volume.real (Metric.ball (0 : EuclideanSpace ℝ (Fin d)) 1)) ^ ((1 : ℝ) / 2))

noncomputable def c_crossover' (d : ℕ) [NeZero d] : ℝ :=
  1 / (2 * C_JN d * c_crossover_bmo_scale d + 1)

\end{lstlisting}
\begin{lemma}[\leanname{crossoverC'\_pos}, \leanname{c\_crossover'\_lt\_one}]
\label{lem:c-crossover-bounds}
$0 < c'(d) < 1$ and $c'(d) \leq 1$.
\end{lemma}

\begin{lstlisting}[style=Matlab-editor,basicstyle=\mlttfamily,escapechar=",mlshowsectionrules=true,mathescape=true,morecomment={[s]{\%\{}{\%\}}},language=lean,numbers=none]
lemma crossoverC'_pos : 0 < c_crossover' d := by
  unfold c_crossover'
  have hbmo_nonneg : 0 ≤ c_crossover_bmo_scale d := crossover_bmo_scale_nonneg (d := d)
  have hCJN_nonneg : 0 ≤ C_JN d := (C_JN_pos d).le
  have hX_nonneg : 0 ≤ C_JN d * c_crossover_bmo_scale d := by
    exact mul_nonneg hCJN_nonneg hbmo_nonneg
  have hden_pos : 0 < 2 * C_JN d * c_crossover_bmo_scale d + 1 := by
    nlinarith
  simpa [one_div] using (inv_pos.mpr hden_pos)

theorem c_crossover'_lt_one : c_crossover' (d := d) < 1 := by
  unfold c_crossover'
  have hbmo_pos : 0 < c_crossover_bmo_scale d := c_crossover_bmo_scale_pos (d := d)
  have hden_gt_one : (1 : ℝ) < 2 * C_JN d * c_crossover_bmo_scale d + 1 := by
    have hCJN_pos : 0 < C_JN d := C_JN_pos d
    nlinarith
  have hlt :
      1 / (2 * C_JN d * c_crossover_bmo_scale d + 1) < 1 / 1 := by
    exact one_div_lt_one_div_of_lt (by positivity : (0 : ℝ) < 1) hden_gt_one
  simpa using hlt

\end{lstlisting}
\begin{proof}
The constant $c_{\BMO}(d)$ is a product of strictly positive terms, so
$c_{\BMO}(d) > 0$. Since $\CJN(d) > 0$, the denominator
$2\,\CJN(d)\,c_{\BMO}(d) + 1$ is strictly greater than $1$, hence its
reciprocal lies in $(0,1)$.
\end{proof}

\subsubsection{Logarithmic gradient bound}

The first ingredient is the Caccioppoli-type inequality for $\log u$, obtained
by testing the supersolution against $\varphi^2/(u + \varepsilon)$.

\begin{theorem}[\leanname{log\_gradient\_bound\_eps}]
\label{thm:log-grad-eps}
Let $\Omega \subseteq E$ be open, $A$ an elliptic coefficient on $\Omega$
with ellipticity constants $\lambda, \Lambda$, $u \in W^{1,2}_{\loc}(\Omega)$
positive on $\Omega$ with $\IsSupersolution A u$, and let
$\varphi \in C_c^\infty(\Omega)$. Then for every $\varepsilon > 0$,
\begin{equation*}
  \int_\Omega \frac{\varphi^2 \,|\nabla u|^2}{(u+\varepsilon)^2}\,dx
    \;\leq\; \frac{4\Lambda}{\lambda}\int_\Omega |\nabla\varphi|^2\,dx.
\end{equation*}
\end{theorem}

\begin{lstlisting}[style=Matlab-editor,basicstyle=\mlttfamily,escapechar=",mlshowsectionrules=true,mathescape=true,morecomment={[s]{\%\{}{\%\}}},language=lean,numbers=none]
/-- Regularized logarithmic gradient bound: for fixed `ε > 0`,
`∫_Ω φ²|∇u|²/(u+ε)² ≤ 4Λ ∫_Ω |∇φ|²`.

This is the ε-level estimate that feeds into `log_gradient_bound_of_supersolution`
via `ε → 0`. -/
theorem log_gradient_bound_eps
    {Ω : Set E} (hΩ : IsOpen Ω)
    (A : EllipticCoeff d Ω)
    {u : E → ℝ}
    (hu_pos : ∀ x ∈ Ω, 0 < u x)
    (hsuper : IsSupersolution A u)
    {φ : E → ℝ} (hφ : ContDiff ℝ (⊤ : ℕ∞) φ)
    (hφ_cs : HasCompactSupport φ) (hφ_sub : tsupport φ ⊆ Ω)
    (hw_u : MemW1pWitness 2 u Ω)
    {ε : ℝ} (hε : 0 < ε) :
    ∫ x in Ω, φ x ^ 2 * ‖hw_u.weakGrad x‖ ^ 2 / (u x + ε) ^ 2 ≤
      4 * A.Λ / A.lam * ∫ x in Ω, ‖fderiv ℝ φ x‖ ^ 2 := by
  -- Define the key integrals.
  set LHS := ∫ x in Ω, φ x ^ 2 * ‖hw_u.weakGrad x‖ ^ 2 / (u x + ε) ^ 2
  -- Q := ∫ φ²/(u+ε)² ⟨∇u, A∇u⟩ (A-quadratic form weighted by φ²/(u+ε)²)
  set Q := ∫ x in Ω, φ x ^ 2 / (u x + ε) ^ 2 *
    @inner ℝ E _ (hw_u.weakGrad x) (matMulE (A.a x) (hw_u.weakGrad x))
  -- X := ∫ φ/(u+ε) ⟨∇φ, A∇u⟩ (mixed A-form)
  set gradφ : E → E := fun x => WithLp.toLp 2
    (fun i => (fderiv ℝ φ x) (EuclideanSpace.single i 1))
  set X := ∫ x in Ω, φ x / (u x + ε) *
    @inner ℝ E _ (gradφ x) (matMulE (A.a x) (hw_u.weakGrad x))
  -- Step (b): Q ≤ 2X (from supersolution with ψ_ε = φ²/(u+ε))
  have hQ_le_2X : Q ≤ 2 * X :=
    supersolution_energy_identity_eps hΩ A hu_pos hsuper hw_u hφ hφ_cs hφ_sub hε
  -- Q ≥ 0 (by coercivity: ⟨ξ, Aξ⟩ ≥ λ‖ξ‖² ≥ 0)
  have hQ_nonneg : 0 ≤ Q := by
    apply integral_nonneg_of_ae
    filter_upwards [A.ae_coercive_nonneg] with x hx
    exact mul_nonneg (div_nonneg (sq_nonneg _) (sq_nonneg _)) (hx (hw_u.weakGrad x))
  have hu_eps_aemeas : AEMeasurable (fun x => u x + ε) (volume.restrict Ω) := by
    exact ((continuous_id.add continuous_const).comp_aestronglyMeasurable
      hw_u.memLp.aestronglyMeasurable).aemeasurable
  have hweight_aemeas : AEMeasurable (fun x => φ x / (u x + ε)) (volume.restrict Ω) := by
    exact (hφ.continuous.aestronglyMeasurable.aemeasurable.div hu_eps_aemeas)
  have hweight_sq_meas :
      AEStronglyMeasurable (fun x => (φ x / (u x + ε)) ^ 2) (volume.restrict Ω) := by
    exact (hweight_aemeas.pow_const 2).aestronglyMeasurable
  obtain ⟨Cφ, hCφ⟩ := hφ_cs.exists_bound_of_continuous hφ.continuous
  have hCφ_nonneg : 0 ≤ Cφ := le_trans (norm_nonneg _) (hCφ (0 : E))
  have hweight_bound : ∀ x, |φ x / (u x + ε)| ≤ Cφ / ε := by
    intro x
    by_cases hx : x ∈ Ω
    · have hu_eps_pos : 0 < u x + ε := by linarith [hu_pos x hx, hε]
      have hnum : |φ x| ≤ Cφ := (Real.norm_eq_abs (φ x)).symm ▸ hCφ x
      rw [abs_div, abs_of_nonneg hu_eps_pos.le]
      have hε_le : ε ≤ u x + ε := by linarith [hu_pos x hx]
      have hgoal : |φ x| / (u x + ε) ≤ Cφ / ε := by
        refine le_trans ?_ (div_le_div_of_nonneg_right hnum hε.le)
        exact div_le_div_of_nonneg_left (abs_nonneg _) hε hε_le
      simpa [div_eq_mul_inv, mul_assoc, mul_left_comm, mul_comm] using hgoal
    · have hφ0 : φ x = 0 := by
        by_contra h
        exact hx (hφ_sub (subset_tsupport φ (Function.mem_support.mpr h)))
      have hnonneg : 0 ≤ Cφ / ε := div_nonneg hCφ_nonneg hε.le
      simpa [hφ0] using hnonneg
  have hweight_sq_bound : ∀ x, |(φ x / (u x + ε)) ^ 2| ≤ Cφ ^ 2 / ε ^ 2 := by
    intro x
    have hx := hweight_bound x
    have hnonneg : 0 ≤ Cφ / ε := div_nonneg hCφ_nonneg hε.le
    calc
      |(φ x / (u x + ε)) ^ 2| = |φ x / (u x + ε)| ^ 2 := by rw [abs_pow]
      _ ≤ (Cφ / ε) ^ 2 := by nlinarith [hx, abs_nonneg (φ x / (u x + ε)), hnonneg]
      _ = Cφ ^ 2 / ε ^ 2 := by ring
  have hQ_int_weight :
      Integrable (fun x => (φ x / (u x + ε)) ^ 2 *
        bilinFormIntegrandOfCoeff A hw_u hw_u x) (volume.restrict Ω) := by
    exact integrable_bounded_mul_bilinFormIntegrand A hw_u hweight_sq_meas hweight_sq_bound
  have hQ_int :
      Integrable (fun x => φ x ^ 2 / (u x + ε) ^ 2 *
        bilinFormIntegrandOfCoeff A hw_u hw_u x) (volume.restrict Ω) := by
    convert hQ_int_weight using 1
    ext x
    rw [bilinFormIntegrandOfCoeff, real_inner_comm]
    rw [div_pow]
  have hQ_eq :
      ∫ x in Ω, φ x ^ 2 / (u x + ε) ^ 2 * bilinFormIntegrandOfCoeff A hw_u hw_u x = Q := by
    change ∫ x in Ω, φ x ^ 2 / (u x + ε) ^ 2 * bilinFormIntegrandOfCoeff A hw_u hw_u x =
      ∫ x in Ω, φ x ^ 2 / (u x + ε) ^ 2 *
        @inner ℝ E _ (hw_u.weakGrad x) (matMulE (A.a x) (hw_u.weakGrad x))
    refine integral_congr_ae ?_
    filter_upwards with x
    rw [bilinFormIntegrandOfCoeff, real_inner_comm]
  -- Coercivity: λ · LHS ≤ Q (pointwise ⟨ξ, Aξ⟩ ≥ λ‖ξ‖², integrate)
  have hLHS_le : A.lam * LHS ≤ Q := by
    let leftInt : E → ℝ := fun x =>
      A.lam * ((φ x / (u x + ε)) ^ 2 * ‖hw_u.weakGrad x‖ ^ 2)
    have hleft_aemeas : AEMeasurable leftInt (volume.restrict Ω) := by
      exact (((hweight_aemeas.pow_const 2).mul
        (hw_u.weakGrad_norm_memLp.aestronglyMeasurable.aemeasurable.pow_const 2)).const_mul A.lam)
    have hleft_meas : AEStronglyMeasurable leftInt (volume.restrict Ω) :=
      hleft_aemeas.aestronglyMeasurable
    have hleft_le_q :
        ∀ᵐ x ∂(volume.restrict Ω),
          leftInt x ≤ φ x ^ 2 / (u x + ε) ^ 2 *
            bilinFormIntegrandOfCoeff A hw_u hw_u x := by
      filter_upwards [A.coercive] with x hx
      have hnn : 0 ≤ (φ x / (u x + ε)) ^ 2 := sq_nonneg _
      calc
        leftInt x
            = (φ x / (u x + ε)) ^ 2 * (A.lam * ‖hw_u.weakGrad x‖ ^ 2) := by
                dsimp [leftInt]
                ring
        _ ≤ (φ x / (u x + ε)) ^ 2 *
              @inner ℝ E _ (hw_u.weakGrad x) (matMulE (A.a x) (hw_u.weakGrad x)) :=
            mul_le_mul_of_nonneg_left (hx (hw_u.weakGrad x)) hnn
        _ = φ x ^ 2 / (u x + ε) ^ 2 *
              bilinFormIntegrandOfCoeff A hw_u hw_u x := by
            rw [bilinFormIntegrandOfCoeff, real_inner_comm]
            rw [div_pow]
    have hleft_int : Integrable leftInt (volume.restrict Ω) := by
      refine Integrable.mono' hQ_int hleft_meas ?_
      filter_upwards [hleft_le_q] with x hx
      have hleft_nonneg : 0 ≤ leftInt x := by
        dsimp [leftInt]
        exact mul_nonneg A.lam_nonneg (mul_nonneg (sq_nonneg _) (sq_nonneg _))
      rw [Real.norm_of_nonneg hleft_nonneg]
      exact hx
    calc A.lam * LHS
        = ∫ x in Ω, A.lam * (φ x ^ 2 * ‖hw_u.weakGrad x‖ ^ 2 / (u x + ε) ^ 2) := by
          rw [integral_const_mul]
      _ = ∫ x in Ω, leftInt x := by
          refine integral_congr_ae ?_
          filter_upwards with x
          dsimp [leftInt]
          rw [div_pow]
          ring
      _ ≤ ∫ x in Ω, φ x ^ 2 / (u x + ε) ^ 2 * bilinFormIntegrandOfCoeff A hw_u hw_u x := by
          exact integral_mono_ae hleft_int hQ_int hleft_le_q
      _ = Q := hQ_eq
  have hφ_test : IsSmoothTestOn Ω φ := ⟨hφ, hφ_cs, hφ_sub⟩
  let hw_φ : MemW1pWitness 2 φ Ω := smoothTestWitness hΩ hφ_test
  have hX_sq_le :
      X ^ 2 ≤ Q * (A.Λ * ∫ x in Ω, ‖fderiv ℝ φ x‖ ^ 2) := by
    have hcs :=
      bilinForm_weighted_cauchySchwarz
        (A := A) (hu := hw_u) (hv := hw_φ)
        (f := fun x => φ x / (u x + ε))
        hweight_aemeas.aestronglyMeasurable hQ_int_weight
    have hX_eq :
        ∫ x in Ω, φ x / (u x + ε) *
          bilinFormIntegrandOfCoeff A hw_u hw_φ x = X := by
      change ∫ x in Ω, φ x / (u x + ε) * bilinFormIntegrandOfCoeff A hw_u hw_φ x =
        ∫ x in Ω, φ x / (u x + ε) *
          @inner ℝ E _ (gradφ x) (matMulE (A.a x) (hw_u.weakGrad x))
      refine integral_congr_ae ?_
      filter_upwards with x
      rw [bilinFormIntegrandOfCoeff, real_inner_comm]
      simp [gradφ, hw_φ, smoothTestWitness, smoothGradField]
    have hgrad_sq_eq :
        ∀ x, ‖hw_φ.weakGrad x‖ ^ (2 : ℕ) = ‖fderiv ℝ φ x‖ ^ 2 := by
      intro x
      have hnorm : ‖hw_φ.weakGrad x‖ = ‖fderiv ℝ φ x‖ := by
        simp [hw_φ, smoothTestWitness, smoothGradField, smoothGradNorm,
          norm_fderiv_eq_smoothGradNorm]
      rw [hnorm]
    have hP_eq :
        ∫ x in Ω, ‖hw_φ.weakGrad x‖ ^ (2 : ℕ) =
          ∫ x in Ω, ‖fderiv ℝ φ x‖ ^ 2 := by
      exact integral_congr_ae (Filter.Eventually.of_forall hgrad_sq_eq)
    calc
      X ^ 2 = (∫ x in Ω, φ x / (u x + ε) *
          bilinFormIntegrandOfCoeff A hw_u hw_φ x) ^ 2 := by
            rw [hX_eq.symm]
      _ ≤ (∫ x in Ω, (φ x / (u x + ε)) ^ 2 *
            bilinFormIntegrandOfCoeff A hw_u hw_u x) *
          (A.Λ * ∫ x in Ω, ‖hw_φ.weakGrad x‖ ^ (2 : ℕ)) := hcs
      _ = Q * (A.Λ * ∫ x in Ω, ‖fderiv ℝ φ x‖ ^ 2) := by
          rw [show ∫ x in Ω, (φ x / (u x + ε)) ^ 2 * bilinFormIntegrandOfCoeff A hw_u hw_u x =
              ∫ x in Ω, φ x ^ 2 / (u x + ε) ^ 2 * bilinFormIntegrandOfCoeff A hw_u hw_u x by
                refine integral_congr_ae ?_
                filter_upwards with x
                rw [div_pow],
            hQ_eq, hP_eq]
  have hR_nonneg : 0 ≤ A.Λ * ∫ x in Ω, ‖fderiv ℝ φ x‖ ^ 2 := by
    refine mul_nonneg A.Λ_nonneg ?_
    exact integral_nonneg (fun x => by positivity)
  have hlamLHS_le :
      A.lam * LHS ≤ 4 * (A.Λ * ∫ x in Ω, ‖fderiv ℝ φ x‖ ^ 2) := by
    exact
      log_gradient_algebraic_conclusion
        (Q := Q) (X := X) (P := A.Λ * ∫ x in Ω, ‖fderiv ℝ φ x‖ ^ 2)
        (LHS := A.lam * LHS) (RHS := A.Λ * ∫ x in Ω, ‖fderiv ℝ φ x‖ ^ 2)
        hQ_nonneg hR_nonneg hQ_le_2X hX_sq_le hLHS_le le_rfl
  calc
    LHS ≤ (4 * (A.Λ * ∫ x in Ω, ‖fderiv ℝ φ x‖ ^ 2)) / A.lam := by
      rw [le_div_iff₀ A.hlam]
      simpa [mul_comm] using hlamLHS_le
    _ = 4 * A.Λ / A.lam * ∫ x in Ω, ‖fderiv ℝ φ x‖ ^ 2 := by ring

\end{lstlisting}
\begin{proof}
Set $\psi := \varphi^2/(u+\varepsilon)$. By the smooth chain-rule witness
construction \leanname{build\_test\_function\_eps}, $\psi$ lies in
$H^1_0(\Omega)$, $\psi \geq 0$, and its weak gradient satisfies almost
everywhere
\begin{equation*}
  \nabla\psi = \frac{2\varphi\,\nabla\varphi}{u+\varepsilon}
    - \frac{\varphi^2\,\nabla u}{(u+\varepsilon)^2}.
\end{equation*}
Since $u$ is a supersolution and $\psi \geq 0$,
\begin{equation*}
  0 \;\leq\; \int_\Omega \langle A\nabla u, \nabla\psi\rangle\,dx
   \;=\; \int_\Omega \frac{2\varphi}{u+\varepsilon}\langle A\nabla u,
       \nabla\varphi\rangle\,dx
   \;-\; \int_\Omega \frac{\varphi^2}{(u+\varepsilon)^2}
       \langle A\nabla u, \nabla u\rangle\,dx.
\end{equation*}
Set
\begin{align*}
  Q &:= \int_\Omega \frac{\varphi^2}{(u+\varepsilon)^2}\langle A\nabla u,\nabla u\rangle\,dx, \\
  X &:= \int_\Omega \frac{\varphi}{u+\varepsilon}\langle A\nabla u,\nabla\varphi\rangle\,dx.
\end{align*}
The supersolution inequality reads $Q \leq 2X$. By Cauchy--Schwarz applied to
the bilinear form $\langle A\,\cdot,\cdot\rangle$,
\begin{equation*}
  X^2 \;\leq\; Q \cdot \int_\Omega \langle A\nabla\varphi,\nabla\varphi\rangle\,dx
    \;\leq\; Q \cdot \Lambda \int_\Omega |\nabla\varphi|^2\,dx.
\end{equation*}
Combining $Q \leq 2X$ with $X^2 \leq Q\,\Lambda\int|\nabla\varphi|^2$, the
elementary inequality $Q^2 \leq 4QP$ implies $Q \leq 4P$ where
$P := \Lambda \int_\Omega |\nabla\varphi|^2$ (see
\leanname{log\_gradient\_algebraic\_conclusion}). Hence
\begin{equation*}
  \lambda \int_\Omega \frac{\varphi^2|\nabla u|^2}{(u+\varepsilon)^2}\,dx
    \;\leq\; Q \;\leq\; 4\Lambda \int_\Omega |\nabla\varphi|^2\,dx,
\end{equation*}
and dividing by $\lambda > 0$ gives the claim.
\end{proof}

\begin{theorem}[\leanname{log\_gradient\_bound\_of\_supersolution}]
\label{thm:log-grad}
Under the same hypotheses, with $u > 0$ on $\Omega$,
\begin{equation*}
  \int_\Omega \frac{\varphi^2\,|\nabla u|^2}{u^2}\,dx
    \;\leq\; \frac{4\Lambda}{\lambda}\int_\Omega |\nabla\varphi|^2\,dx.
\end{equation*}
\end{theorem}

\begin{lstlisting}[style=Matlab-editor,basicstyle=\mlttfamily,escapechar=",mlshowsectionrules=true,mathescape=true,morecomment={[s]{\%\{}{\%\}}},language=lean,numbers=none]
/-- Logarithmic gradient bound for positive supersolutions.

If `u > 0` satisfies `-∇·(A∇u) ≥ 0` on `Ω`, then for every smooth cutoff
`φ` supported in `Ω`:
`∫_Ω φ² |∇u|²/u² ≤ 4Λ ∫_Ω |∇φ|²`.

Equivalently, setting `v = -log u`: `∫_Ω φ² |∇v|² ≤ 4Λ ∫_Ω |∇φ|²`.

Proof: apply `log_gradient_bound_eps` for each `ε > 0` to get
`∫ φ²|∇u|²/(u+ε)² ≤ 4Λ ∫ |∇φ|²`, then send `ε → 0`. Since
`u > 0` on `Ω`, the integrand `φ²|∇u|²/(u+ε)²` increases monotonically
to `φ²|∇u|²/u²` as `ε ↘ 0`. By monotone convergence, the limit integral
satisfies the same bound. -/
theorem log_gradient_bound_of_supersolution
    {Ω : Set E} (hΩ : IsOpen Ω)
    (A : EllipticCoeff d Ω)
    {u : E → ℝ}
    (hu_pos : ∀ x ∈ Ω, 0 < u x)
    (hsuper : IsSupersolution A u)
    {φ : E → ℝ} (hφ : ContDiff ℝ (⊤ : ℕ∞) φ)
    (hφ_cs : HasCompactSupport φ) (hφ_sub : tsupport φ ⊆ Ω)
    (hw_u : MemW1pWitness 2 u Ω) :
    ∫ x in Ω, φ x ^ 2 * ‖hw_u.weakGrad x‖ ^ 2 / (u x) ^ 2 ≤
      4 * A.Λ / A.lam * ∫ x in Ω, ‖fderiv ℝ φ x‖ ^ 2 := by
  -- The RHS is independent of ε, and for each ε > 0:
  --   ∫ φ²|∇u|²/(u+ε)² ≤ 4Λ ∫ |∇φ|² (by log_gradient_bound_eps).
  -- As ε ↘ 0, the LHS increases to ∫ φ²|∇u|²/u² (since u > 0 on Ω).
  -- Apply monotone convergence to pass the bound to the limit.
  suffices h : ∀ ε : ℝ, 0 < ε →
      ∫ x in Ω, φ x ^ 2 * ‖hw_u.weakGrad x‖ ^ 2 / (u x + ε) ^ 2 ≤
        4 * A.Λ / A.lam * ∫ x in Ω, ‖fderiv ℝ φ x‖ ^ 2 from by
    let μ : Measure E := volume.restrict Ω
    let εn : ℕ → ℝ := fun n => 1 / ((n : ℝ) + 1)
    let g : ℕ → E → ℝ := fun n x =>
      (φ x / (u x + εn n)) ^ 2 * ‖hw_u.weakGrad x‖ ^ 2
    let g0 : E → ℝ := fun x =>
      (φ x / u x) ^ 2 * ‖hw_u.weakGrad x‖ ^ 2
    set R := 4 * A.Λ / A.lam * ∫ x in Ω, ‖fderiv ℝ φ x‖ ^ 2
    obtain ⟨Cφ, hCφ⟩ := hφ_cs.exists_bound_of_continuous hφ.continuous
    have hCφ_nonneg : 0 ≤ Cφ := le_trans (norm_nonneg _) (hCφ (0 : E))
    have hεn_pos : ∀ n, 0 < εn n := by
      intro n
      dsimp [εn]
      positivity
    have hu_eps_aemeas : ∀ n, AEMeasurable (fun x => u x + εn n) μ := by
      intro n
      exact ((continuous_id.add continuous_const).comp_aestronglyMeasurable
        hw_u.memLp.aestronglyMeasurable).aemeasurable
    have hweight_sq_meas :
        ∀ n, AEStronglyMeasurable (fun x => (φ x / (u x + εn n)) ^ 2) μ := by
      intro n
      exact ((hφ.continuous.aestronglyMeasurable.aemeasurable.div
        (hu_eps_aemeas n)).pow_const 2).aestronglyMeasurable
    have hweight_sq_bound :
        ∀ n x, |(φ x / (u x + εn n)) ^ 2| ≤ Cφ ^ 2 / (εn n) ^ 2 := by
      intro n x
      by_cases hx : x ∈ Ω
      · have hu_eps_pos : 0 < u x + εn n := by linarith [hu_pos x hx, hεn_pos n]
        have hnum : |φ x| ≤ Cφ := (Real.norm_eq_abs (φ x)).symm ▸ hCφ x
        have hnonneg : 0 ≤ Cφ / εn n := div_nonneg hCφ_nonneg (hεn_pos n).le
        have hε_le : εn n ≤ u x + εn n := by linarith [hu_pos x hx]
        have hbound : |φ x / (u x + εn n)| ≤ Cφ / εn n := by
          rw [abs_div, abs_of_nonneg hu_eps_pos.le]
          refine le_trans ?_ (div_le_div_of_nonneg_right hnum (hεn_pos n).le)
          exact div_le_div_of_nonneg_left (abs_nonneg _) (hεn_pos n) hε_le
        calc
          |(φ x / (u x + εn n)) ^ 2| = |φ x / (u x + εn n)| ^ 2 := by rw [abs_pow]
          _ ≤ (Cφ / εn n) ^ 2 := by
              nlinarith [hbound, abs_nonneg (φ x / (u x + εn n)), hnonneg]
          _ = Cφ ^ 2 / (εn n) ^ 2 := by ring
      · have hφ0 : φ x = 0 := by
          by_contra hφx
          exact hx (hφ_sub (subset_tsupport φ (Function.mem_support.mpr hφx)))
        have hnonneg : 0 ≤ Cφ ^ 2 / (εn n) ^ 2 := by
          positivity
        simpa [hφ0] using hnonneg
    have hgrad_sq_meas : AEStronglyMeasurable (fun x => ‖hw_u.weakGrad x‖ ^ 2) μ := by
      exact (hw_u.weakGrad_norm_memLp.aestronglyMeasurable.aemeasurable.pow_const 2).aestronglyMeasurable
    have hg_meas : ∀ n, AEMeasurable (g n) μ := by
      intro n
      exact (hweight_sq_meas n).aemeasurable.mul hgrad_sq_meas.aemeasurable
    have hg0_meas : AEStronglyMeasurable g0 μ := by
      exact ((hφ.continuous.aestronglyMeasurable.aemeasurable.div
        hw_u.memLp.aestronglyMeasurable.aemeasurable).pow_const 2).aestronglyMeasurable.mul hgrad_sq_meas
    have hg_nonneg : ∀ n, 0 ≤ᵐ[μ] g n := by
      intro n
      exact Filter.Eventually.of_forall fun x => by
        dsimp [g]
        exact mul_nonneg (sq_nonneg _) (sq_nonneg _)
    have hg0_nonneg : 0 ≤ᵐ[μ] g0 := by
      exact Filter.Eventually.of_forall fun x => by
        dsimp [g0]
        exact mul_nonneg (sq_nonneg _) (sq_nonneg _)
    have hg_int : ∀ n, Integrable (g n) μ := by
      intro n
      exact (hw_u.weakGrad_norm_memLp.integrable_sq).bdd_mul
        (c := Cφ ^ 2 / (εn n) ^ 2) (hweight_sq_meas n)
        (Filter.Eventually.of_forall (hweight_sq_bound n))
    have hright_eq :
        ∀ n,
          ∫⁻ x, ENNReal.ofReal (g n x) ∂μ =
            ENNReal.ofReal (∫ x, g n x ∂μ) := by
      intro n
      exact
        (MeasureTheory.ofReal_integral_eq_lintegral_ofReal
          (μ := μ) (f := g n) (hg_int n) (hg_nonneg n)).symm
    have hfR : ∀ n, ∫ x, g n x ∂μ ≤ R := by
      intro n
      have hbound := h (εn n) (hεn_pos n)
      change ∫ x in Ω, g n x ∂volume ≤ R
      have hEq :
          ∫ x in Ω, g n x ∂volume =
            ∫ x in Ω, φ x ^ 2 * ‖hw_u.weakGrad x‖ ^ 2 / (u x + εn n) ^ 2 ∂volume := by
        refine integral_congr_ae ?_
        filter_upwards with x
        dsimp [g]
        rw [div_pow]
        ring
      rw [hEq]
      exact hbound
    have hFatou :
        ∫⁻ x, ENNReal.ofReal (g0 x) ∂μ ≤
          atTop.liminf (fun n => ∫⁻ x, ENNReal.ofReal (g n x) ∂μ) := by
      have hmeas :
          ∀ n, AEMeasurable (fun x => ENNReal.ofReal (g n x)) μ := by
        intro n
        exact (hg_meas n).ennreal_ofReal
      have hleft := MeasureTheory.lintegral_liminf_le' (μ := μ) hmeas
      have hlim :
          (fun x =>
            Filter.liminf (fun n => ENNReal.ofReal (g n x)) atTop) =ᵐ[μ]
              (fun x => ENNReal.ofReal (g0 x)) := by
        filter_upwards [ae_restrict_mem hΩ.measurableSet] with x hx
        have hu_posx : 0 < u x := hu_pos x hx
        have hden_tend : Tendsto (fun n => u x + εn n) atTop (nhds (u x)) := by
          simpa [εn, one_div] using
            (tendsto_one_div_add_atTop_nhds_zero_nat.const_add (u x))
        have hquot_tend :
            Tendsto (fun n => φ x / (u x + εn n)) atTop (nhds (φ x / u x)) := by
          exact Tendsto.div tendsto_const_nhds hden_tend hu_posx.ne'
        have hreal_tend : Tendsto (fun n => g n x) atTop (nhds (g0 x)) := by
          exact ((continuous_pow 2).tendsto (φ x / u x)).comp hquot_tend |>.mul tendsto_const_nhds
        have henn_tend :
            Tendsto (fun n => ENNReal.ofReal (g n x)) atTop
              (nhds (ENNReal.ofReal (g0 x))) := by
          exact (ENNReal.continuous_ofReal.tendsto _).comp hreal_tend
        exact henn_tend.liminf_eq
      calc
        ∫⁻ x, ENNReal.ofReal (g0 x) ∂μ
            = ∫⁻ x, Filter.liminf (fun n => ENNReal.ofReal (g n x)) atTop ∂μ := by
                exact lintegral_congr_ae hlim.symm
        _ ≤ atTop.liminf (fun n => ∫⁻ x, ENNReal.ofReal (g n x) ∂μ) := hleft
    have hliminf_le :
        atTop.liminf (fun n => ∫⁻ x, ENNReal.ofReal (g n x) ∂μ) ≤ ENNReal.ofReal R := by
      rw [Filter.liminf_congr (Eventually.of_forall hright_eq)]
      refine Filter.liminf_le_of_frequently_le' (Frequently.of_forall fun n => ?_)
      exact ENNReal.ofReal_le_ofReal (hfR n)
    have hmain_enn : ∫⁻ x, ENNReal.ofReal (g0 x) ∂μ ≤ ENNReal.ofReal R := by
      exact le_trans hFatou hliminf_le
    have hR_nonneg : 0 ≤ R := by
      dsimp [R]
      refine mul_nonneg ?_ (integral_nonneg (fun x => by positivity))
      · have hnum_nonneg : 0 ≤ 4 * A.Λ := by nlinarith [A.Λ_nonneg]
        exact div_nonneg hnum_nonneg A.hlam.le
    calc
      ∫ x in Ω, φ x ^ 2 * ‖hw_u.weakGrad x‖ ^ 2 / (u x) ^ 2
          = ∫ x, g0 x ∂μ := by
              refine integral_congr_ae ?_
              filter_upwards with x
              dsimp [g0]
              rw [div_pow]
              ring
      _ = ENNReal.toReal (∫⁻ x, ENNReal.ofReal (g0 x) ∂μ) := by
            exact MeasureTheory.integral_eq_lintegral_of_nonneg_ae hg0_nonneg hg0_meas
      _ ≤ ENNReal.toReal (ENNReal.ofReal R) := by
            have hleft_ne_top : ∫⁻ x, ENNReal.ofReal (g0 x) ∂μ ≠ ∞ := by
              exact ne_of_lt (lt_of_le_of_lt hmain_enn (by simp))
            exact (ENNReal.toReal_le_toReal hleft_ne_top (by simp)).2 hmain_enn
      _ = R := by rw [ENNReal.toReal_ofReal hR_nonneg]
      _ = 4 * A.Λ / A.lam * ∫ x in Ω, ‖fderiv ℝ φ x‖ ^ 2 := by
            rfl
  intro ε hε
  exact log_gradient_bound_eps hΩ A hu_pos hsuper hφ hφ_cs hφ_sub hw_u hε

\end{lstlisting}
\begin{proof}
Apply Theorem~\ref{thm:log-grad-eps} with $\varepsilon_n = 1/(n+1)$. The
integrand
\[
  g_n(x) := \varphi(x)^2|\nabla u(x)|^2/(u(x)+\varepsilon_n)^2
\]
increases monotonically as $\varepsilon_n \downarrow 0$ to
$g_0(x) := \varphi(x)^2|\nabla u(x)|^2/u(x)^2$ at every point of $\Omega$
since $u > 0$ there. By Fatou's lemma,
\begin{equation*}
  \int_\Omega g_0\,dx \;\leq\; \liminf_{n\to\infty}\int_\Omega g_n\,dx
    \;\leq\; \frac{4\Lambda}{\lambda}\int_\Omega |\nabla\varphi|^2\,dx.
\end{equation*}
\end{proof}

\subsubsection{BMO control of \texorpdfstring{$\log u$}{log u}}

For $\varepsilon > 0$ define the regularized log
$v_\varepsilon(x) := -\log(u(x)+\varepsilon) + \log\varepsilon$ and its smooth
chain-rule witness $\leanname{regularizedLogFun}$. Combining the log-gradient
bound with the Poincar\'e inequality on $\Bz{1/2}$ gives the BMO estimate.

\begin{theorem}[\leanname{regularizedLog\_unit\_halfBall\_meanOscillation}]
\label{thm:bmo-half}
Let $A$ be normalized elliptic on $\Bz 1$, $u > 0$ a supersolution. Then
for $v := v_\varepsilon$,
\begin{equation*}
  \fint_{\Bz{1/2}} \Bigl|v(x) - \fint_{\Bz{1/2}} v\Bigr|\,dx
    \;\leq\; c_{\BMO}(d)\,\Lambda_A^{1/2}.
\end{equation*}
\end{theorem}

\begin{lstlisting}[style=Matlab-editor,basicstyle=\mlttfamily,escapechar=",mlshowsectionrules=true,mathescape=true,morecomment={[s]{\%\{}{\%\}}},language=lean,numbers=none]
theorem regularizedLog_unit_halfBall_meanOscillation
    (A : NormalizedEllipticCoeff d (Metric.ball (0 : E) 1))
    {u : E → ℝ}
    (hu_pos : ∀ x ∈ Metric.ball (0 : E) 1, 0 < u x)
    (hsuper : IsSupersolution A.1 u)
    {ε : ℝ} (hε : 0 < ε) :
    let v := regularizedLogFun (u := u) hε
    (⨍ x in Metric.ball (0 : E) (1 / 2 : ℝ),
        |v x - ⨍ y in Metric.ball (0 : E) (1 / 2 : ℝ), v y ∂volume| ∂volume) ≤
      c_crossover_bmo_scale d * A.1.Λ ^ ((1 : ℝ) / 2) := by
  intro v
  let Ω : Set E := Metric.ball (0 : E) 1
  let Bhalf : Set E := Metric.ball (0 : E) (1 / 2 : ℝ)
  let hw_u : MemW1pWitness 2 u Ω :=
    (isSupersolution_memW1p hsuper).someWitness
  let hw_v : MemW1pWitness 2 v Ω := by
    simpa [v, Ω] using regularizedLogWitness (Ω := Ω) Metric.isOpen_ball hu_pos hw_u hε
  let hw_v_half : MemW1pWitness 2 v Bhalf :=
    hw_v.restrict Metric.isOpen_ball (Metric.ball_subset_ball (by norm_num))
  obtain ⟨η, _⟩ := Cutoff.exists (d := d) (0 : E)
    (r := (1 / 2 : ℝ)) (R := (3 / 4 : ℝ)) (by norm_num) (by norm_num)
  have hη_cs : HasCompactSupport η.toFun := by
    apply HasCompactSupport.intro' (K := Metric.closedBall (0 : E) (3 / 4 : ℝ))
    · exact isCompact_closedBall (0 : E) (3 / 4 : ℝ)
    · exact isClosed_closedBall
    · intro x hx
      exact zero_outside_of_tsupport_subset (Ω := Metric.closedBall (0 : E) (3 / 4 : ℝ))
        η.support_subset hx
  have hη_sub : tsupport η.toFun ⊆ Ω := by
    exact η.support_subset.trans (Metric.closedBall_subset_ball (by norm_num : (3 / 4 : ℝ) < 1))
  have hlog :=
    log_gradient_bound_eps (Ω := Ω) Metric.isOpen_ball A.1 hu_pos hsuper
      η.smooth hη_cs hη_sub hw_u hε
  have hfactor_meas :
      AEStronglyMeasurable
        (fun x => η.toFun x ^ 2 / (u x + ε) ^ 2)
        (volume.restrict Ω) := by
    have hη_meas : Measurable η.toFun := η.smooth.continuous.measurable
    have hnum_ae :
        AEStronglyMeasurable (fun x => η.toFun x ^ 2) (volume.restrict Ω) :=
      (hη_meas.pow_const 2).aestronglyMeasurable
    have hden_ae :
        AEStronglyMeasurable (fun x => ((u x + ε) ^ 2)⁻¹) (volume.restrict Ω) := by
      have hden_ae' :
          AEMeasurable (fun x => ((u x + ε) ^ 2)⁻¹) (volume.restrict Ω) := by
        exact AEMeasurable.inv (((hw_u.memLp.aestronglyMeasurable.add_const ε).pow 2).aemeasurable)
      exact hden_ae'.aestronglyMeasurable
    simpa [div_eq_mul_inv] using hnum_ae.mul hden_ae
  have hfactor_bound :
      ∀ᵐ x ∂(volume.restrict Ω), |η.toFun x ^ 2 / (u x + ε) ^ 2| ≤ 1 / ε ^ 2 := by
    filter_upwards [ae_restrict_mem measurableSet_ball] with x hx
    have hηsq_le : η.toFun x ^ 2 ≤ 1 := by
      have hη_nonneg : 0 ≤ η.toFun x := η.nonneg x
      have hη_le : η.toFun x ≤ 1 := η.le_one x
      nlinarith
    have hux : 0 < u x := hu_pos x hx
    have hfrac :
        η.toFun x ^ 2 / (u x + ε) ^ 2 ≤ 1 / ε ^ 2 := by
      have hden : ε ^ 2 ≤ (u x + ε) ^ 2 := by
        nlinarith [hux, hε]
      have hinv :
          1 / (u x + ε) ^ 2 ≤ 1 / ε ^ 2 := by
        exact one_div_le_one_div_of_le (by positivity) hden
      have hmono :
          η.toFun x ^ 2 / (u x + ε) ^ 2 ≤ 1 / (u x + ε) ^ 2 := by
        exact div_le_div_of_nonneg_right hηsq_le (sq_nonneg (u x + ε))
      exact hmono.trans hinv
    have hnonneg : 0 ≤ η.toFun x ^ 2 / (u x + ε) ^ 2 := by positivity
    rw [abs_of_nonneg hnonneg]
    exact hfrac
  have hweight_int :
      IntegrableOn
        (fun x => η.toFun x ^ 2 / (u x + ε) ^ 2 * ‖hw_u.weakGrad x‖ ^ 2)
        Ω volume := by
    simpa [Ω, mul_assoc, mul_left_comm, mul_comm] using
      (hw_u.weakGrad_norm_memLp.integrable_sq).bdd_mul
        (c := 1 / ε ^ 2) hfactor_meas hfactor_bound
  have hhalf_eq :
      ∫ x in Bhalf, ‖hw_v_half.weakGrad x‖ ^ 2 ∂volume =
        ∫ x in Bhalf, η.toFun x ^ 2 / (u x + ε) ^ 2 * ‖hw_u.weakGrad x‖ ^ 2 ∂volume := by
    apply MeasureTheory.setIntegral_congr_fun measurableSet_ball
    intro x hx
    have hxΩ : x ∈ Ω := Metric.ball_subset_ball (by norm_num : (1 / 2 : ℝ) ≤ 1) hx
    have hux : 0 < u x := hu_pos x hxΩ
    have hvec :
        hw_v_half.weakGrad x = (-(1 / (u x + ε))) • hw_u.weakGrad x := by
      have hgrad := regularizedLogWitness_grad (Ω := Ω) Metric.isOpen_ball hu_pos hw_u hε
      ext i
      simpa [hw_v_half, Pi.smul_apply, smul_eq_mul] using
        hgrad x hxΩ i
    have hηx : η.toFun x = 1 := η.eq_one x hx
    calc
      ‖hw_v_half.weakGrad x‖ ^ 2
          = (1 / (u x + ε)) ^ 2 * ‖hw_u.weakGrad x‖ ^ 2 := by
              rw [hvec, norm_smul, Real.norm_eq_abs, abs_neg,
                abs_of_pos (one_div_pos.mpr (by linarith : 0 < u x + ε))]
              ring
      _ = η.toFun x ^ 2 / (u x + ε) ^ 2 * ‖hw_u.weakGrad x‖ ^ 2 := by
            rw [hηx]
            field_simp [show u x + ε ≠ 0 by linarith]
  have hhalf_le :
      ∫ x in Bhalf, η.toFun x ^ 2 / (u x + ε) ^ 2 * ‖hw_u.weakGrad x‖ ^ 2 ∂volume ≤
        ∫ x in Ω, η.toFun x ^ 2 / (u x + ε) ^ 2 * ‖hw_u.weakGrad x‖ ^ 2 ∂volume := by
    exact MeasureTheory.setIntegral_mono_set hweight_int
      (ae_of_all _ (by intro x; positivity))
      (ae_of_all _ (Metric.ball_subset_ball (by norm_num : (1 / 2 : ℝ) ≤ 1))
        )
  have hgrad_sq_bound :
      ∫ x in Bhalf, ‖hw_v_half.weakGrad x‖ ^ 2 ∂volume ≤
        4 * A.1.Λ * ∫ x in Ω, ‖fderiv ℝ η.toFun x‖ ^ 2 ∂volume := by
    have hlog' :
        ∫ x in Ω, η.toFun x ^ 2 * ‖hw_u.weakGrad x‖ ^ 2 / (u x + ε) ^ 2 ∂volume ≤
          4 * A.1.Λ * ∫ x in Ω, ‖fderiv ℝ η.toFun x‖ ^ 2 ∂volume := by
      simpa [A.lam_eq_one] using hlog
    have hlog'' :
        ∫ x in Ω, η.toFun x ^ 2 / (u x + ε) ^ 2 * ‖hw_u.weakGrad x‖ ^ 2 ∂volume ≤
          4 * A.1.Λ * ∫ x in Ω, ‖fderiv ℝ η.toFun x‖ ^ 2 ∂volume := by
      simpa [div_eq_mul_inv, mul_assoc, mul_left_comm, mul_comm] using hlog'
    rw [hhalf_eq]
    exact le_trans hhalf_le hlog''
  have hderiv_sq_int :
      IntegrableOn (fun x => ‖fderiv ℝ η.toFun x‖ ^ 2) Ω volume := by
    have hmeas : AEStronglyMeasurable (fun x => ‖fderiv ℝ η.toFun x‖ ^ 2) volume := by
      have hmeas' : Measurable (fun x => ‖fderiv ℝ η.toFun x‖ ^ 2) := by
        fun_prop
      exact hmeas'.aestronglyMeasurable
    have hbd :
        ∀ᵐ x ∂(volume.restrict Ω), ‖‖fderiv ℝ η.toFun x‖ ^ 2‖ ≤ (4 * (Mst : ℝ)) ^ 2 := by
      filter_upwards with x
      have hgrad_bd :
          ‖fderiv ℝ η.toFun x‖ ≤ (Mst : ℝ) * (((3 / 4 : ℝ) - 1 / 2)⁻¹) := by
        simpa using η.grad_bound x
      have hgrad_bd' : ‖fderiv ℝ η.toFun x‖ ≤ 4 * (Mst : ℝ) := by
        have hgrad_bd'' : ‖fderiv ℝ η.toFun x‖ ≤ (Mst : ℝ) * 4 := by
          calc
            ‖fderiv ℝ η.toFun x‖ ≤ (Mst : ℝ) * (((3 / 4 : ℝ) - 1 / 2)⁻¹) := hgrad_bd
            _ = (Mst : ℝ) * 4 := by norm_num
        simpa [mul_comm] using hgrad_bd''
      have hsq_nonneg : 0 ≤ ‖fderiv ℝ η.toFun x‖ ^ 2 := by positivity
      have hsq : ‖fderiv ℝ η.toFun x‖ ^ 2 ≤ (4 * (Mst : ℝ)) ^ 2 := by
        nlinarith [hgrad_bd', norm_nonneg (fderiv ℝ η.toFun x)]
      rw [Real.norm_of_nonneg hsq_nonneg]
      exact hsq
    exact Measure.integrableOn_of_bounded
      (s_finite := by simpa [Ω] using measure_ball_lt_top.ne)
      hmeas hbd
  have hderiv_sq_bound :
      ∫ x in Ω, ‖fderiv ℝ η.toFun x‖ ^ 2 ∂volume ≤
        (4 * (Mst : ℝ)) ^ 2 * volume.real Ω := by
    have hmono := MeasureTheory.setIntegral_mono (μ := volume) (s := Ω)
      (f := fun x => ‖fderiv ℝ η.toFun x‖ ^ 2)
      (g := fun _ : E => (4 * (Mst : ℝ)) ^ 2)
      hderiv_sq_int (integrableOn_const (μ := volume) (show volume Ω ≠ ∞ by simpa [Ω] using measure_ball_lt_top.ne))
      (by
        intro x
        have hgrad_bd :
            ‖fderiv ℝ η.toFun x‖ ≤ (Mst : ℝ) * (((3 / 4 : ℝ) - 1 / 2)⁻¹) := by
          simpa using η.grad_bound x
        have hbd' : ‖fderiv ℝ η.toFun x‖ ≤ 4 * (Mst : ℝ) := by
          have hbd'' : ‖fderiv ℝ η.toFun x‖ ≤ (Mst : ℝ) * 4 := by
            calc
              ‖fderiv ℝ η.toFun x‖ ≤ (Mst : ℝ) * (((3 / 4 : ℝ) - 1 / 2)⁻¹) := hgrad_bd
              _ = (Mst : ℝ) * 4 := by norm_num
          simpa [mul_comm] using hbd''
        have hsq : ‖fderiv ℝ η.toFun x‖ ^ 2 ≤ (4 * (Mst : ℝ)) ^ 2 := by
          nlinarith [hbd', norm_nonneg (fderiv ℝ η.toFun x)]
        exact hsq)
    simpa [Ω, MeasureTheory.setIntegral_const, smul_eq_mul, mul_comm, mul_left_comm, mul_assoc]
      using hmono
  have hgrad_lp_eq :
      MeasureTheory.lpNorm (fun x => ‖hw_v_half.weakGrad x‖) 2 (volume.restrict Bhalf) =
        (∫ x in Bhalf, ‖hw_v_half.weakGrad x‖ ^ 2 ∂volume) ^ ((1 : ℝ) / 2) := by
    simpa using
      (MeasureTheory.lpNorm_eq_integral_norm_rpow_toReal
        (μ := volume.restrict Bhalf) (p := (2 : ℝ≥0∞))
        (f := fun x => ‖hw_v_half.weakGrad x‖) (by norm_num) (by norm_num)
        hw_v_half.weakGrad_norm_memLp.aestronglyMeasurable)
  have hgrad_lp_bound :
      MeasureTheory.lpNorm (fun x => ‖hw_v_half.weakGrad x‖) 2 (volume.restrict Bhalf) ≤
        8 * (Mst : ℝ) * A.1.Λ ^ ((1 : ℝ) / 2) * (volume.real Ω) ^ ((1 : ℝ) / 2) := by
    rw [hgrad_lp_eq]
    have hgrad_sq_bound' :
        ∫ x in Bhalf, ‖hw_v_half.weakGrad x‖ ^ 2 ∂volume ≤
          4 * A.1.Λ * ((4 * (Mst : ℝ)) ^ 2 * volume.real Ω) := by
      calc
        ∫ x in Bhalf, ‖hw_v_half.weakGrad x‖ ^ 2 ∂volume
            ≤ 4 * A.1.Λ * ∫ x in Ω, ‖fderiv ℝ η.toFun x‖ ^ 2 ∂volume := hgrad_sq_bound
        _ ≤ 4 * A.1.Λ * ((4 * (Mst : ℝ)) ^ 2 * volume.real Ω) := by
              exact mul_le_mul_of_nonneg_left hderiv_sq_bound (by nlinarith [A.1.Λ_nonneg])
    have hmain :
        (∫ x in Bhalf, ‖hw_v_half.weakGrad x‖ ^ 2 ∂volume) ^
            ((1 : ℝ) / 2) ≤
          (4 * A.1.Λ * ((4 * (Mst : ℝ)) ^ 2 * volume.real Ω)) ^ ((1 : ℝ) / 2) := by
      exact Real.rpow_le_rpow (by positivity) hgrad_sq_bound' (by norm_num)
    refine le_trans hmain ?_
    have hcalc :
        (4 * A.1.Λ * ((4 * (Mst : ℝ)) ^ 2 * volume.real Ω)) ^ ((1 : ℝ) / 2) =
          8 * (Mst : ℝ) * A.1.Λ ^ ((1 : ℝ) / 2) * (volume.real Ω) ^ ((1 : ℝ) / 2) := by
      have hsqrt4 :
          (4 * A.1.Λ) ^ ((1 : ℝ) / 2) = 2 * A.1.Λ ^ ((1 : ℝ) / 2) := by
        rw [Real.mul_rpow (by positivity : 0 ≤ (4 : ℝ)) A.1.Λ_nonneg]
        norm_num [Real.sqrt_eq_rpow]
      rw [show 4 * A.1.Λ * ((4 * (Mst : ℝ)) ^ 2 * volume.real Ω) =
            (4 * A.1.Λ) * ((4 * (Mst : ℝ)) ^ 2 * volume.real Ω) by ring,
        Real.mul_rpow (by nlinarith [A.1.Λ_nonneg] : 0 ≤ 4 * A.1.Λ)
          (by positivity : 0 ≤ (4 * (Mst : ℝ)) ^ 2 * volume.real Ω),
        Real.mul_rpow (by positivity : 0 ≤ (4 * (Mst : ℝ)) ^ 2)
          (by positivity : 0 ≤ volume.real Ω),
        show ((4 * (Mst : ℝ)) ^ 2) ^ ((1 : ℝ) / 2) = 4 * (Mst : ℝ) by
          rw [← Real.sqrt_eq_rpow, Real.sqrt_sq]
          positivity,
        hsqrt4]
      ring
    rw [hcalc]
  calc
    (⨍ x in Bhalf, |v x - ⨍ y in Bhalf, v y ∂volume| ∂volume)
        ≤ Cmo * (1 / 2 : ℝ) ^ (1 - (d : ℝ) / 2) *
            MeasureTheory.lpNorm (fun x => ‖hw_v_half.weakGrad x‖) 2
              (volume.restrict Bhalf) := by
            simpa [Bhalf, hw_v_half] using
              (crossover_average_abs_sub_average_ball_le_grad_lpNorm_two
                (u := v) (x₀ := (0 : E)) (R := (1 / 2 : ℝ))
                (by norm_num) hw_v_half)
    _ ≤ Cmo * (1 / 2 : ℝ) ^ (1 - (d : ℝ) / 2) *
          (8 * (Mst : ℝ) * A.1.Λ ^ ((1 : ℝ) / 2) * (volume.real Ω) ^ ((1 : ℝ) / 2)) := by
            have hfac_nonneg : 0 ≤ Cmo * (1 / 2 : ℝ) ^ (1 - (d : ℝ) / 2) := by
              have hd_pos : 0 < d := Nat.pos_of_ne_zero (NeZero.ne d)
              rw [show Cmo =
                (volume.real (Metric.ball (0 : E) 1)) ^ (-(1 / 2 : ℝ)) * C_poinc_val d by rfl]
              exact mul_nonneg
                (mul_nonneg
                  (Real.rpow_nonneg MeasureTheory.measureReal_nonneg _)
                  (C_poinc_val_pos hd_pos).le)
                (Real.rpow_nonneg (by positivity : 0 ≤ (1 / 2 : ℝ)) _)
            exact mul_le_mul_of_nonneg_left hgrad_lp_bound hfac_nonneg
    _ = c_crossover_bmo_scale d * A.1.Λ ^ ((1 : ℝ) / 2) := by
          change
            (((volume.real (Metric.ball (0 : E) 1)) ^ (-(1 / 2 : ℝ)) * C_poinc_val d) *
              (1 / 2 : ℝ) ^ (1 - (d : ℝ) / 2) *
              (8 * (Mst : ℝ) * A.1.Λ ^ ((1 : ℝ) / 2) *
                (volume.real (Metric.ball (0 : E) 1)) ^ ((1 : ℝ) / 2))) =
              c_crossover_bmo_scale d * A.1.Λ ^ ((1 : ℝ) / 2)
          simp [c_crossover_bmo_scale, mul_assoc, mul_left_comm, mul_comm]

\end{lstlisting}
\begin{proof}
Pick a cutoff $\eta \in C_c^\infty(\Bz{3/4})$ with $\eta = 1$ on $\Bz{1/2}$,
$0 \leq \eta \leq 1$, and $|\nabla \eta| \leq 8\Mst$ (existence given by
\leanname{Cutoff.exists}). Theorem~\ref{thm:log-grad-eps} (with
$\varphi = \eta$, $\lambda_A = 1$) yields
\begin{equation*}
  \int_{\Bz{1/2}} |\nabla v_\varepsilon|^2\,dx
    \;\leq\; \int_{\Bz 1} \eta^2 \frac{|\nabla u|^2}{(u+\varepsilon)^2}\,dx
    \;\leq\; 4\Lambda_A \int_{\Bz 1}|\nabla \eta|^2\,dx
    \;\leq\; 4\Lambda_A \cdot (8\Mst)^2 \vol(\Bz 1).
\end{equation*}
Hence $\|\nabla v_\varepsilon\|_{L^2(\Bz{1/2})} \leq 16\Mst\,\Lambda_A^{1/2}\,\vol(\Bz 1)^{1/2}$.
By the Poincar\'e--Wirtinger inequality on $\Bz{1/2}$ with constant
$\Cpoinc(d) \cdot (1/2)^{1-d/2}$ (the dimensional rescaling factor),
\begin{multline*}
  \fint_{\Bz{1/2}}\Bigl|v_\varepsilon - \fint_{\Bz{1/2}} v_\varepsilon\Bigr|\,dx
    \;\leq\; \vol(\Bz{1/2})^{-1/2} \Cpoinc(d) (1/2)^{1-d/2}
        \|\nabla v_\varepsilon\|_{L^2(\Bz{1/2})} \\
    \;\leq\; \vol(\Bz 1)^{-1/2}\Cpoinc(d) (1/2)^{1-d/2}\cdot
        8\Mst\,\Lambda_A^{1/2}\,\vol(\Bz 1)^{1/2}
    \;=\; c_{\BMO}(d)\,\Lambda_A^{1/2}.
\end{multline*}
\end{proof}

The same statement holds on every ball $B(x_0, R/2)$ with $\overline{B(x_0,R)}
\subset \Bz 1$, by affine rescaling $\leanname{rescaleToUnitBall}$
(see \leanname{regularizedLog\_ball\_meanOscillation}).

\subsubsection{John--Nirenberg exponential integrability}

\begin{definition}[\leanname{crossoverJNTheta}, \leanname{crossoverJNKloc}, \leanname{crossoverJNKhalf}]
\label{def:JN-constants}
Set
\begin{align*}
  \theta(d)  &:= 24 \cdot 2^{d+1} \cdot \bigl(c'(d)\,c_{\BMO}(d)\bigr), \\
  K_{\loc}(d)&:= \exp(\theta(d)) \cdot (1 + \CJN(d)),\\
  K_{1/2}(d) &:= \lceil 97^d \rceil \cdot K_{\loc}(d).
\end{align*}
\end{definition}

\begin{lstlisting}[style=Matlab-editor,basicstyle=\mlttfamily,escapechar=",mlshowsectionrules=true,mathescape=true,morecomment={[s]{\%\{}{\%\}}},language=lean,numbers=none]
noncomputable def crossoverJNTheta (d : ℕ) [NeZero d] : ℝ :=
  24 * (2 : ℝ) ^ (d + 1) * (c_crossover' d * c_crossover_bmo_scale d)

noncomputable def crossoverJNKloc (d : ℕ) [NeZero d] : ℝ :=
  Real.exp (crossoverJNTheta d) * (1 + C_JN d)

noncomputable def crossoverJNKhalf (d : ℕ) [NeZero d] : ℝ :=
  (Nat.ceil ((97 : ℝ) ^ d) : ℝ) * crossoverJNKloc d


\end{lstlisting}
\begin{theorem}[\leanname{regularizedLog\_halfBall\_exp\_average\_to\_origin\_le}]
\label{thm:exp-half}
With $p := c'(d)/\Lambda_A^{1/2}$, $v := v_\varepsilon$, and
$a := \fint_{\Bz{1/48}} v\,dz$,
\begin{equation*}
  \fint_{\Bz{1/2}} \exp(p\,|v(z) - a|)\,dz \;\leq\; K_{1/2}(d).
\end{equation*}
\end{theorem}

\begin{lstlisting}[style=Matlab-editor,basicstyle=\mlttfamily,escapechar=",mlshowsectionrules=true,mathescape=true,morecomment={[s]{\%\{}{\%\}}},language=lean,numbers=none]
theorem regularizedLog_halfBall_exp_average_to_origin_le
    (A : NormalizedEllipticCoeff d (Metric.ball (0 : E) 1))
    {u : E → ℝ}
    (hu_pos : ∀ x ∈ Metric.ball (0 : E) 1, 0 < u x)
    (hsuper : IsSupersolution A.1 u)
    {ε : ℝ} (hε : 0 < ε) :
    let v := regularizedLogMeasurable (A := A) hu_pos hsuper hε
    let p := c_crossover' d / A.1.Λ ^ ((1 : ℝ) / 2)
    let a := ⨍ z in Metric.ball (0 : E) (1 / 48 : ℝ), v z ∂volume
    let δ := 24 * (2 : ℝ) ^ (d + 1) *
      (c_crossover_bmo_scale d * A.1.Λ ^ ((1 : ℝ) / 2))
    let Kloc := Real.exp (p * δ) * (1 + C_JN d)
    let Khalf := (Nat.ceil ((97 : ℝ) ^ d) : ℝ) * Kloc
    IntegrableOn (fun z => Real.exp (p * |v z - a|))
      (Metric.ball (0 : E) (1 / 2 : ℝ)) volume ∧
    (⨍ z in Metric.ball (0 : E) (1 / 2 : ℝ),
      Real.exp (p * |v z - a|) ∂volume) ≤ Khalf := by
  intro v p a δ Kloc Khalf
  let Bhalf : Set E := Metric.ball (0 : E) (1 / 2 : ℝ)
  let F : E → ℝ := fun z => Real.exp (p * |v z - a|)
  let μ : Measure E := volume.restrict Bhalf
  have hρ_pos : 0 < (1 / 48 : ℝ) := by norm_num
  have hhalf_pos : 0 < (1 / 2 : ℝ) := by norm_num
  obtain ⟨t, htcard, htmem, hcover, _⟩ :=
    exists_finite_inner_ball_cover_with_card
      (d := d) (x₀ := (0 : E))
      (r := (1 / 2 : ℝ)) (ρ := (1 / 48 : ℝ)) (R := 1)
      hhalf_pos hρ_pos (by norm_num)
  let G : E → ℝ := fun x => Finset.sum t (fun c => Set.indicator (Metric.ball c (1 / 48 : ℝ)) F x)
  have hΛsqrt_pos : 0 < A.1.Λ ^ ((1 : ℝ) / 2) := by
    have hΛge : 1 ≤ A.1.Λ := by
      have hΛ := A.1.hΛ
      rw [A.2] at hΛ
      simpa using hΛ
    exact Real.rpow_pos_of_pos (lt_of_lt_of_le (by norm_num : (0 : ℝ) < 1) hΛge) _
  have hp_pos : 0 < p := by
    exact div_pos (crossoverC'_pos (d := d)) hΛsqrt_pos
  have hp_nonneg : 0 ≤ p := hp_pos.le
  have hKloc_nonneg : 0 ≤ Kloc := by
    dsimp [Kloc]
    refine mul_nonneg (le_of_lt (Real.exp_pos _)) ?_
    linarith [show 0 < C_JN d from C_JN_pos d]
  have hv_meas : Measurable v :=
    regularizedLogMeasurable_measurable (d := d) (A := A) hu_pos hsuper hε
  have hF_meas : Measurable F := by
    refine Real.measurable_exp.comp ?_
    exact measurable_const.mul ((hv_meas.sub measurable_const).abs)
  have hF_nonneg : ∀ z, 0 ≤ F z := by
    intro z
    exact le_of_lt (Real.exp_pos _)
  have hlocal_integrable :
      ∀ c ∈ t, IntegrableOn F (Metric.ball c (1 / 48 : ℝ)) volume := by
    intro c hc
    let ac : ℝ := ⨍ z in Metric.ball c (1 / 48 : ℝ), v z ∂volume
    let Fc : E → ℝ := fun z => Real.exp (p * |v z - ac|)
    have hcenter : c ∈ Metric.closedBall (0 : E) (1 / 2 : ℝ) := htmem c hc
    have hbase :=
      regularizedLog_smallBall_exp_average_le (d := d) (A := A)
        hu_pos hsuper hε hcenter
    have hFc_int : IntegrableOn Fc (Metric.ball c (1 / 48 : ℝ)) volume := by
      simpa [v, p, Fc, ac] using hbase.1
    refine Integrable.mono' (hFc_int.const_mul (Real.exp (p * δ))) ?_ ?_
    · exact hF_meas.aestronglyMeasurable
    · filter_upwards with z
      have htri : |v z - a| ≤ |v z - ac| + |ac - a| := by
        simpa [ac, sub_eq_add_neg, add_comm, add_left_comm, add_assoc] using
          (abs_sub_le (v z) ac a)
      have hchain :
          |ac - a| ≤ δ := by
        simpa [ac, a, δ] using
          (regularizedLog_smallBallAverage_to_origin_le (d := d) (A := A)
            hu_pos hsuper hε hcenter)
      have hbound : |v z - a| ≤ |v z - ac| + δ := by
        nlinarith [htri, hchain]
      have hmul : p * |v z - a| ≤ p * |v z - ac| + p * δ := by
        nlinarith
      have hmajor :
          F z ≤ Real.exp (p * δ) * Fc z := by
        calc
          F z = Real.exp (p * |v z - a|) := by rfl
          _ ≤ Real.exp (p * |v z - ac| + p * δ) := by
                exact Real.exp_le_exp.mpr hmul
          _ = Real.exp (p * δ) * Fc z := by
                rw [show p * |v z - ac| + p * δ = p * δ + p * |v z - ac| by ring,
                  Real.exp_add]
      rw [Real.norm_eq_abs, abs_of_nonneg (hF_nonneg z)]
      exact hmajor
  have hlocal_integral_bound :
      ∀ c ∈ t,
        ∫ z in Metric.ball c (1 / 48 : ℝ), F z ∂volume ≤
          volume.real (Metric.ball c (1 / 48 : ℝ)) * Kloc := by
    intro c hc
    let ac : ℝ := ⨍ z in Metric.ball c (1 / 48 : ℝ), v z ∂volume
    let Fc : E → ℝ := fun z => Real.exp (p * |v z - ac|)
    have hcenter : c ∈ Metric.closedBall (0 : E) (1 / 2 : ℝ) := htmem c hc
    have hbase :=
      regularizedLog_smallBall_exp_average_le (d := d) (A := A)
        hu_pos hsuper hε hcenter
    have hFc_int : IntegrableOn Fc (Metric.ball c (1 / 48 : ℝ)) volume := by
      simpa [v, p, Fc, ac] using hbase.1
    have hF_int : IntegrableOn F (Metric.ball c (1 / 48 : ℝ)) volume :=
      hlocal_integrable c hc
    have hpointwise :
        ∀ z ∈ Metric.ball c (1 / 48 : ℝ), F z ≤ Real.exp (p * δ) * Fc z := by
      intro z hz
      have htri : |v z - a| ≤ |v z - ac| + |ac - a| := by
        simpa [ac, sub_eq_add_neg, add_comm, add_left_comm, add_assoc] using
          (abs_sub_le (v z) ac a)
      have hchain :
          |ac - a| ≤ δ := by
        simpa [ac, a, δ] using
          (regularizedLog_smallBallAverage_to_origin_le (d := d) (A := A)
            hu_pos hsuper hε hcenter)
      have hbound : |v z - a| ≤ |v z - ac| + δ := by
        nlinarith [htri, hchain]
      have hmul : p * |v z - a| ≤ p * |v z - ac| + p * δ := by
        nlinarith
      calc
        F z = Real.exp (p * |v z - a|) := by rfl
        _ ≤ Real.exp (p * |v z - ac| + p * δ) := by
              exact Real.exp_le_exp.mpr hmul
        _ = Real.exp (p * δ) * Fc z := by
              rw [show p * |v z - ac| + p * δ = p * δ + p * |v z - ac| by ring,
                Real.exp_add]
    have hmajor_int :
        ∫ z in Metric.ball c (1 / 48 : ℝ), F z ∂volume ≤
          ∫ z in Metric.ball c (1 / 48 : ℝ), Real.exp (p * δ) * Fc z ∂volume := by
      have hmajor_ae :
          ∀ᵐ z ∂(volume.restrict (Metric.ball c (1 / 48 : ℝ))),
            F z ≤ Real.exp (p * δ) * Fc z := by
        filter_upwards [ae_restrict_mem measurableSet_ball] with z hz
        exact hpointwise z hz
      have hmajor_int' :
          ∫ z, F z ∂(volume.restrict (Metric.ball c (1 / 48 : ℝ))) ≤
            ∫ z, Real.exp (p * δ) * Fc z ∂(volume.restrict (Metric.ball c (1 / 48 : ℝ))) := by
        exact integral_mono_ae hF_int (hFc_int.const_mul _) hmajor_ae
      simpa using hmajor_int'
    have hball_pos : 0 < volume.real (Metric.ball c (1 / 48 : ℝ)) := by
      exact ENNReal.toReal_pos
        (measure_ball_pos (μ := volume) c hρ_pos).ne'
        (measure_ball_lt_top (μ := volume) (x := c) (r := (1 / 48 : ℝ))).ne
    have hFc_avg :
        (⨍ z in Metric.ball c (1 / 48 : ℝ), Fc z ∂volume) ≤ 1 + C_JN d := by
      simpa [v, p, Fc, ac] using hbase.2
    have hFc_int_bound :
        ∫ z in Metric.ball c (1 / 48 : ℝ), Fc z ∂volume ≤
          volume.real (Metric.ball c (1 / 48 : ℝ)) * (1 + C_JN d) := by
      rw [MeasureTheory.setAverage_eq, smul_eq_mul] at hFc_avg
      rw [show (volume.real (Metric.ball c (1 / 48 : ℝ)))⁻¹ *
          ∫ z in Metric.ball c (1 / 48 : ℝ), Fc z ∂volume =
            (∫ z in Metric.ball c (1 / 48 : ℝ), Fc z ∂volume) /
              volume.real (Metric.ball c (1 / 48 : ℝ)) by
            field_simp [hball_pos.ne']] at hFc_avg
      simpa [mul_comm, mul_left_comm, mul_assoc] using (div_le_iff₀ hball_pos).mp hFc_avg
    calc
      ∫ z in Metric.ball c (1 / 48 : ℝ), F z ∂volume
          ≤ ∫ z in Metric.ball c (1 / 48 : ℝ), Real.exp (p * δ) * Fc z ∂volume := hmajor_int
      _ = Real.exp (p * δ) * ∫ z in Metric.ball c (1 / 48 : ℝ), Fc z ∂volume := by
            rw [integral_const_mul]
      _ ≤ Real.exp (p * δ) *
            (volume.real (Metric.ball c (1 / 48 : ℝ)) * (1 + C_JN d)) := by
              gcongr
      _ = volume.real (Metric.ball c (1 / 48 : ℝ)) * Kloc := by
            dsimp [Kloc]
            ring
  have hBcover :
      Bhalf ⊆ ⋃ c ∈ t, Metric.ball c (1 / 48 : ℝ) := by
    intro x hx
    exact hcover (Metric.mem_closedBall.2 (le_of_lt (Metric.mem_ball.1 hx)))
  have hterm_int :
      ∀ c ∈ t, Integrable (Set.indicator (Metric.ball c (1 / 48 : ℝ)) F) μ := by
    intro c hc
    exact ((hlocal_integrable c hc).integrable_indicator measurableSet_ball).mono_measure
      Measure.restrict_le_self
  have hsum_int :
      Integrable G μ := by
    dsimp [G]
    refine integrable_finset_sum _ ?_
    intro c hc
    exact hterm_int c hc
  have hdom_plain :
      ∀ᵐ x ∂μ, F x ≤ G x := by
    filter_upwards [ae_restrict_mem measurableSet_ball] with x hx
    have hxcover := hBcover hx
    rw [Set.mem_iUnion] at hxcover
    obtain ⟨c, hxcover⟩ := hxcover
    rw [Set.mem_iUnion] at hxcover
    obtain ⟨hc, hxc⟩ := hxcover
    have hsum_nonneg :
        ∀ c' ∈ t, 0 ≤ Set.indicator (Metric.ball c' (1 / 48 : ℝ)) F x := by
      intro c' hc'
      by_cases hx' : x ∈ Metric.ball c' (1 / 48 : ℝ)
      · rw [Set.indicator_of_mem hx']
        exact hF_nonneg x
      · rw [Set.indicator_of_notMem hx']
    have hle :
        F x ≤ G x := by
      calc
        F x = Set.indicator (Metric.ball c (1 / 48 : ℝ)) F x := by
          symm
          exact (show Set.indicator (Metric.ball c (1 / 48 : ℝ)) F x = F x from
            Set.indicator_of_mem (f := F) hxc)
        _ ≤ G x := by
              exact Finset.single_le_sum hsum_nonneg (by simpa using hc)
    exact hle
  have hdom_ae :
      ∀ᵐ x ∂μ, ‖F x‖ ≤ G x := by
    filter_upwards [hdom_plain] with x hx
    rw [Real.norm_eq_abs, abs_of_nonneg (hF_nonneg x)]
    exact hx
  have hF_int_mu : Integrable F μ := by
    exact Integrable.mono' hsum_int hF_meas.aestronglyMeasurable hdom_ae
  have hint_le_sum :
      ∫ x, F x ∂μ ≤ ∫ x, G x ∂μ := by
    exact integral_mono_ae hF_int_mu hsum_int hdom_plain
  have hterm_le :
      ∀ c ∈ t,
        ∫ x, Set.indicator (Metric.ball c (1 / 48 : ℝ)) F x ∂μ ≤
          ∫ x in Metric.ball c (1 / 48 : ℝ), F x ∂volume := by
    intro c hc
    have hterm_nonneg :
        0 ≤ᵐ[volume] Set.indicator (Metric.ball c (1 / 48 : ℝ)) F := by
      exact Filter.Eventually.of_forall fun x => by
        by_cases hx : x ∈ Metric.ball c (1 / 48 : ℝ)
        · rw [Set.indicator_of_mem hx]
          exact hF_nonneg x
        · rw [Set.indicator_of_notMem hx]
          positivity
    have hterm_int_vol :
        Integrable (Set.indicator (Metric.ball c (1 / 48 : ℝ)) F) volume :=
      (hlocal_integrable c hc).integrable_indicator measurableSet_ball
    calc
      ∫ x, Set.indicator (Metric.ball c (1 / 48 : ℝ)) F x ∂μ
          ≤ ∫ x, Set.indicator (Metric.ball c (1 / 48 : ℝ)) F x ∂volume := by
                exact integral_mono_measure Measure.restrict_le_self hterm_nonneg hterm_int_vol
      _ = ∫ x in Metric.ball c (1 / 48 : ℝ), F x ∂volume := by
            rw [integral_indicator measurableSet_ball]
  have hsmall_meas_le :
      ∀ c : E, volume.real (Metric.ball c (1 / 48 : ℝ)) ≤ volume.real Bhalf := by
    intro c
    rw [crossover_volumeReal_ball_eq (d := d) c hρ_pos]
    rw [crossover_volumeReal_ball_eq (d := d) (0 : E) hhalf_pos]
    exact mul_le_mul_of_nonneg_right
      (pow_le_pow_left₀ (by positivity : (0 : ℝ) ≤ 1 / 48)
        (by norm_num : (1 / 48 : ℝ) ≤ 1 / 2) d)
      (by positivity)
  have hμ_eq : μ.real Set.univ = volume.real Bhalf := by
    simp [μ, Bhalf]
  have hμ_pos : 0 < μ.real Set.univ := by
    rw [hμ_eq]
    exact ENNReal.toReal_pos
      (measure_ball_pos (μ := volume) (0 : E) hhalf_pos).ne'
      (measure_ball_lt_top (μ := volume) (x := (0 : E)) (r := (1 / 2 : ℝ))).ne
  have hint_bound :
      ∫ x, F x ∂μ ≤ μ.real Set.univ * Khalf := by
    calc
      ∫ x, F x ∂μ
          ≤ ∫ x, G x ∂μ := hint_le_sum
      _ = Finset.sum t (fun c => ∫ x, Set.indicator (Metric.ball c (1 / 48 : ℝ)) F x ∂μ) := by
            dsimp [G]
            rw [integral_finset_sum]
            intro c hc
            exact hterm_int c hc
      _ ≤ Finset.sum t (fun c => ∫ x in Metric.ball c (1 / 48 : ℝ), F x ∂volume) := by
            refine Finset.sum_le_sum ?_
            intro c hc
            exact hterm_le c hc
      _ ≤ Finset.sum t (fun _ => volume.real Bhalf * Kloc) := by
            refine Finset.sum_le_sum ?_
            intro c hc
            exact le_trans (hlocal_integral_bound c hc) (by
              exact mul_le_mul_of_nonneg_right (hsmall_meas_le c) hKloc_nonneg)
      _ = (t.card : ℝ) * (volume.real Bhalf * Kloc) := by simp
      _ ≤ (Nat.ceil ((97 : ℝ) ^ d) : ℝ) * (volume.real Bhalf * Kloc) := by
            gcongr
            have h97 :
                (4 * (2 : ℝ)⁻¹ * 48 + 1 : ℝ) = 97 := by
              norm_num
            simpa [div_eq_mul_inv, h97] using htcard
      _ = μ.real Set.univ * Khalf := by
            rw [hμ_eq]
            dsimp [Khalf]
            ring
  have hAvg :
      (⨍ z in Bhalf, F z ∂volume) ≤ Khalf := by
    rw [MeasureTheory.setAverage_eq, smul_eq_mul]
    rw [show (volume.real Bhalf)⁻¹ * ∫ z in Bhalf, F z ∂volume =
        (μ.real Set.univ)⁻¹ * ∫ z, F z ∂μ by
          simp [μ, Bhalf]]
    rw [show (μ.real Set.univ)⁻¹ * ∫ z, F z ∂μ =
        (∫ z, F z ∂μ) / μ.real Set.univ by
          field_simp [hμ_pos.ne']]
    exact (div_le_iff₀ hμ_pos).2 (by simpa [mul_comm, mul_left_comm, mul_assoc] using hint_bound)
  exact ⟨by simpa [IntegrableOn, μ, Bhalf] using hF_int_mu, by simpa [Bhalf, F] using hAvg⟩

\end{lstlisting}
\begin{proof}
Apply the local John--Nirenberg lemma at scale $\rho = 1/48$ on a finite cover
of $\Bz{1/2}$ by balls $B(x_i, 1/48)$ with $\#\{x_i\} \leq \lceil 97^d\rceil$.
On each cover ball, Theorem~\ref{thm:bmo-half} (rescaled to $B(x_i, R)$ via
\leanname{regularizedLog\_ball\_meanOscillation}) gives
$\|v - \langle v\rangle\|_{\BMO} \leq c_{\BMO}(d)\,\Lambda_A^{1/2}$. The
classical John--Nirenberg theorem $\leanname{C\_JN}$ then yields
\[
  \fint_{B(x_i,1/48)} \exp\bigl(p\,|v - \langle v\rangle_{B(x_i,1/48)}|\bigr)\,dz
    \;\leq\; 1 + \CJN(d).
\]
The shift between local and global averages is controlled by $\theta(d) =
24 \cdot 2^{d+1}\,c'(d)\,c_{\BMO}(d)$ via the standard telescoping
$\leanname{regularizedLog\_smallBallAverage\_to\_origin\_le}$. Multiplying
each local exponential bound by $\exp(p\,\theta(d)) = \exp(\theta(d))$ (since
$p\cdot \delta = \theta(d)$ by direct computation, using
$\delta := 24\cdot 2^{d+1}\,(c_{\BMO}(d)\Lambda_A^{1/2})$ and
$p \cdot \Lambda_A^{1/2} = c'(d)$), and summing over the $\lceil 97^d\rceil$
cover balls gives the asserted bound by $K_{1/2}(d)$.
\end{proof}

\subsubsection{Crossover product bound}

\begin{theorem}[\leanname{crossover\_product\_bound\_exists}]
\label{thm:crossover-exists}
There exists a constant $C(d) \geq 1$ such that for every normalized elliptic
$A$ on $\Bz 1$, every positive supersolution $u$, and $p_0 := c'(d)/\Lambda_A^{1/2}$,
\[
  \Bigl(\fint_{\Bz{1/2}} u^{p_0}\,dx\Bigr)\Bigl(\fint_{\Bz{1/2}} u^{-p_0}\,dx\Bigr)
    \;\leq\; C(d).
\]
\end{theorem}

\begin{lstlisting}[style=Matlab-editor,basicstyle=\mlttfamily,escapechar=",mlshowsectionrules=true,mathescape=true,morecomment={[s]{\%\{}{\%\}}},language=lean,numbers=none]
/-- The crossover product is bounded by a dimension-only constant.
This is the master existence theorem powering `crossover_estimate`. -/
theorem crossover_product_bound_exists :
    ∃ C : ℝ, 1 ≤ C ∧
      ∀ (_hd : 2 < (d : ℝ))
        (A : NormalizedEllipticCoeff d (Metric.ball (0 : E) 1))
        {u : E → ℝ}
        (_hu_pos : ∀ x ∈ Metric.ball (0 : E) 1, 0 < u x)
        (_hsuper : IsSupersolution A.1 u),
        (⨍ x in Metric.ball (0 : E) (1 / 2 : ℝ),
            |u x| ^ (c_crossover' d / A.1.Λ ^ ((1 : ℝ) / 2)) ∂volume) *
          (⨍ x in Metric.ball (0 : E) (1 / 2 : ℝ),
            |u x| ^ (-(c_crossover' d / A.1.Λ ^ ((1 : ℝ) / 2))) ∂volume) ≤ C := by
  let C : ℝ := max 1 ((crossoverJNKhalf d) ^ 2)
  refine ⟨C, le_max_left _ _, ?_⟩
  intro hd A u hu_pos hsuper
  let Bhalf : Set E := Metric.ball (0 : E) (1 / 2 : ℝ)
  let μ : Measure E := volume.restrict Bhalf
  let p : ℝ := c_crossover' d / A.1.Λ ^ ((1 : ℝ) / 2)
  let Aavg : ℝ := ⨍ x in Bhalf, |u x| ^ p ∂volume
  let K : ℝ := (crossoverJNKhalf d) ^ 2
  let M : ℝ := K / Aavg
  let εn : ℕ → ℝ := fun n => 1 / ((n : ℝ) + 1)
  let g : ℕ → E → ℝ := fun n x => |(u x + εn n)⁻¹| ^ p
  let g0 : E → ℝ := fun x => |(u x)⁻¹| ^ p
  haveI : IsFiniteMeasure μ := by
    refine ⟨by
      simpa [μ, Bhalf] using
        (measure_ball_lt_top (μ := volume) (x := (0 : E)) (r := (1 / 2 : ℝ)))⟩
  have hBsub : Bhalf ⊆ Metric.ball (0 : E) 1 := Metric.ball_subset_ball (by norm_num)
  have hu_pos_half : ∀ x ∈ Bhalf, 0 < u x := by
    intro x hx
    exact hu_pos x (hBsub hx)
  have hμ_eq : μ.real Set.univ = volume.real Bhalf := by
    simp [μ, Bhalf]
  have hμ_pos : 0 < μ.real Set.univ := by
    rw [hμ_eq]
    exact ENNReal.toReal_pos
      (measure_ball_pos (μ := volume) (0 : E) (by norm_num : (0 : ℝ) < 1 / 2)).ne'
      (measure_ball_lt_top (μ := volume) (x := (0 : E)) (r := (1 / 2 : ℝ))).ne
  have hΛsqrt_pos : 0 < A.1.Λ ^ ((1 : ℝ) / 2) := by
    have hΛge : 1 ≤ A.1.Λ := by
      have hΛ := A.1.hΛ
      rw [A.2] at hΛ
      simpa using hΛ
    exact Real.rpow_pos_of_pos (lt_of_lt_of_le (by norm_num : (0 : ℝ) < 1) hΛge) _
  have hΛsqrt_ge_one : 1 ≤ A.1.Λ ^ ((1 : ℝ) / 2) := by
    have hΛge : 1 ≤ A.1.Λ := by
      have hΛ := A.1.hΛ
      rw [A.2] at hΛ
      simpa using hΛ
    exact Real.one_le_rpow hΛge (by positivity : 0 ≤ (1 : ℝ) / 2)
  have hp_pos : 0 < p := by
    exact div_pos (crossoverC'_pos (d := d)) hΛsqrt_pos
  have hp_nonneg : 0 ≤ p := hp_pos.le
  have hp_le_one : p ≤ 1 := by
    have hinv_le_one : (A.1.Λ ^ ((1 : ℝ) / 2))⁻¹ ≤ 1 := by
      exact inv_le_one_of_one_le₀ hΛsqrt_ge_one
    have hp_le_cc :
        p ≤ c_crossover' d := by
      dsimp [p]
      rw [div_eq_mul_inv]
      have hcc_nonneg : 0 ≤ c_crossover' d := (crossoverC'_pos (d := d)).le
      calc
        c_crossover' d * (A.1.Λ ^ ((1 : ℝ) / 2))⁻¹
            ≤ c_crossover' d * 1 := by
              exact mul_le_mul_of_nonneg_left hinv_le_one hcc_nonneg
        _ = c_crossover' d := by ring
    exact hp_le_cc.trans (c_crossover'_le_one (d := d))
  have hp_le_two : p ≤ 2 := by linarith
  let hw_u : MemW1pWitness 2 u (Metric.ball (0 : E) 1) :=
    (isSupersolution_memW1p hsuper).someWitness
  let hw_u_half : MemW1pWitness 2 u Bhalf :=
    hw_u.restrict Metric.isOpen_ball hBsub
  have hu_aestrong : AEStronglyMeasurable u μ := by
    simpa [μ, Bhalf] using hw_u_half.memLp.aestronglyMeasurable
  have hu_sq_int : Integrable (fun x => |u x| ^ (2 : ℝ)) μ := by
    simpa [μ, Bhalf, Real.rpow_natCast] using hw_u_half.memLp.integrable_norm_pow'
  have hpow_meas : AEStronglyMeasurable (fun x => |u x| ^ p) μ := by
    exact
      ((Real.continuous_rpow_const hp_nonneg).measurable.comp_aemeasurable
        ((continuous_abs.measurable).comp_aemeasurable hu_aestrong.aemeasurable)).aestronglyMeasurable
  have hpow_nonneg : ∀ x, 0 ≤ |u x| ^ p := by
    intro x
    exact Real.rpow_nonneg (abs_nonneg _) _
  have hpow_int : Integrable (fun x => |u x| ^ p) μ := by
    have hdom_int : Integrable (fun x => (1 : ℝ) + |u x| ^ (2 : ℝ)) μ := by
      exact (integrable_const (1 : ℝ)).add hu_sq_int
    refine Integrable.mono' hdom_int hpow_meas ?_
    filter_upwards [ae_restrict_mem measurableSet_ball] with x hx
    have habs_nonneg : 0 ≤ |u x| := abs_nonneg _
    have hpow_nonneg' : 0 ≤ |u x| ^ p := hpow_nonneg x
    have hbound : |u x| ^ p ≤ 1 + |u x| ^ (2 : ℝ) := by
      by_cases hsmall : |u x| ≤ 1
      · have hle_one : |u x| ^ p ≤ 1 := by
          exact Real.rpow_le_one habs_nonneg hsmall hp_nonneg
        linarith [show 0 ≤ |u x| ^ (2 : ℝ) by positivity]
      · have hge_one : 1 ≤ |u x| := le_of_not_ge hsmall
        have hle_two : |u x| ^ p ≤ |u x| ^ (2 : ℝ) := by
          exact Real.rpow_le_rpow_of_exponent_le hge_one hp_le_two
        linarith
    rw [Real.norm_eq_abs, abs_of_nonneg hpow_nonneg']
    exact hbound
  have hpow_support :
      Bhalf ⊆ Function.support (fun x => |u x| ^ p) := by
    intro x hx
    have hux : 0 < u x := hu_pos_half x hx
    have hpow_pos : 0 < |u x| ^ p := by
      rw [abs_of_pos hux]
      exact Real.rpow_pos_of_pos hux _
    exact Function.mem_support.mpr (ne_of_gt hpow_pos)
  have hsupport_pos : 0 < μ (Function.support (fun x => |u x| ^ p)) := by
    have hμ_ball_pos : 0 < μ Bhalf := by
      simpa [μ, Bhalf] using
        (measure_ball_pos (μ := volume) (0 : E) (by norm_num : (0 : ℝ) < 1 / 2))
    exact lt_of_lt_of_le hμ_ball_pos (measure_mono hpow_support)
  have hIpos_pos : 0 < ∫ x, |u x| ^ p ∂μ := by
    rw [integral_pos_iff_support_of_nonneg hpow_nonneg hpow_int]
    exact hsupport_pos
  have hAavg_pos : 0 < Aavg := by
    dsimp [Aavg]
    rw [MeasureTheory.setAverage_eq, smul_eq_mul]
    rw [show ∫ x in Bhalf, |u x| ^ p ∂volume = ∫ x, |u x| ^ p ∂μ by
          simp [μ, Bhalf]]
    rw [← hμ_eq]
    exact mul_pos (inv_pos.mpr hμ_pos) hIpos_pos
  have hεn_pos : ∀ n, 0 < εn n := by
    intro n
    dsimp [εn]
    positivity
  have hεn_le_one : ∀ n, εn n ≤ 1 := by
    intro n
    dsimp [εn]
    have hden_ge : (1 : ℝ) ≤ (n : ℝ) + 1 := by
      nlinarith
    simpa using (one_div_le_one_div_of_le (by positivity : (0 : ℝ) < 1) hden_ge)
  have hshift_pos_meas :
      ∀ n, AEStronglyMeasurable (fun x => (u x + εn n) ^ p) μ := by
    intro n
    exact
      ((Real.continuous_rpow_const hp_nonneg).measurable.comp_aemeasurable
        (hu_aestrong.aemeasurable.add measurable_const.aemeasurable)).aestronglyMeasurable
  have hshift_pos_int :
      ∀ n, Integrable (fun x => (u x + εn n) ^ p) μ := by
    intro n
    have hdom_int : Integrable (fun x => (1 : ℝ) + |u x| ^ p) μ := by
      exact (integrable_const (1 : ℝ)).add hpow_int
    refine Integrable.mono' hdom_int (hshift_pos_meas n) ?_
    filter_upwards [ae_restrict_mem measurableSet_ball] with x hx
    have hux_pos : 0 < u x := hu_pos_half x hx
    have hux_nonneg : 0 ≤ u x := hux_pos.le
    have huxε_nonneg : 0 ≤ u x + εn n := by
      linarith [hux_pos, hεn_pos n]
    have hsubadd :
        (u x + εn n) ^ p ≤ u x ^ p + (εn n) ^ p := by
      simpa [add_comm] using
        (Real.rpow_add_le_add_rpow (hεn_pos n).le hux_nonneg hp_nonneg hp_le_one)
    have hεpow_le_one : (εn n) ^ p ≤ 1 := by
      exact Real.rpow_le_one (hεn_pos n).le (hεn_le_one n) hp_nonneg
    have hpow_nonneg' : 0 ≤ (u x + εn n) ^ p := by
      exact Real.rpow_nonneg huxε_nonneg _
    calc
      ‖(u x + εn n) ^ p‖ = (u x + εn n) ^ p := by
        rw [Real.norm_eq_abs, abs_of_nonneg hpow_nonneg']
      _ ≤ u x ^ p + (εn n) ^ p := hsubadd
      _ ≤ u x ^ p + 1 := by linarith
      _ = 1 + |u x| ^ p := by
        rw [abs_of_pos hux_pos]
        ring
  have hAavg_le_shift :
      ∀ n, Aavg ≤ ⨍ x in Bhalf, (u x + εn n) ^ p ∂volume := by
    intro n
    dsimp [Aavg]
    refine setAverage_mono_of_nonneg_local (d := d) hpow_nonneg ?_ ?_
    · simpa [IntegrableOn, μ, Bhalf] using hshift_pos_int n
    · filter_upwards [ae_restrict_mem measurableSet_ball] with x hx
      have hux_pos : 0 < u x := hu_pos_half x hx
      have hux_nonneg : 0 ≤ u x := hux_pos.le
      have hle_add : u x ≤ u x + εn n := by
        linarith [hεn_pos n]
      rw [abs_of_pos hux_pos]
      exact Real.rpow_le_rpow hux_nonneg hle_add hp_nonneg
  have hg_meas : ∀ n, AEStronglyMeasurable (g n) μ := by
    intro n
    dsimp [g]
    exact
      ((Real.continuous_rpow_const hp_nonneg).measurable.comp_aemeasurable
        ((continuous_abs.measurable).comp_aemeasurable
          ((hu_aestrong.aemeasurable.add measurable_const.aemeasurable).inv))).aestronglyMeasurable
  have hg0_meas : AEStronglyMeasurable g0 μ := by
    dsimp [g0]
    exact
      ((Real.continuous_rpow_const hp_nonneg).measurable.comp_aemeasurable
        ((continuous_abs.measurable).comp_aemeasurable
          (hu_aestrong.aemeasurable.inv))).aestronglyMeasurable
  have hg_nonneg : ∀ n, 0 ≤ᵐ[μ] g n := by
    intro n
    exact Filter.Eventually.of_forall fun x => by
      dsimp [g]
      positivity
  have hg0_nonneg : 0 ≤ᵐ[μ] g0 := by
    exact Filter.Eventually.of_forall fun x => by
      dsimp [g0]
      positivity
  have hg_int : ∀ n, Integrable (g n) μ := by
    intro n
    have hconst_int : Integrable (fun _ : E => ((εn n)⁻¹) ^ p) μ := by
      exact integrable_const (((εn n)⁻¹) ^ p)
    refine Integrable.mono' hconst_int (hg_meas n) ?_
    filter_upwards [ae_restrict_mem measurableSet_ball] with x hx
    have hux_pos : 0 < u x := hu_pos_half x hx
    have huxε_pos : 0 < u x + εn n := by
      linarith [hux_pos, hεn_pos n]
    have hε_le : εn n ≤ u x + εn n := by
      linarith [hux_pos, hεn_pos n]
    have hle_inv : (u x + εn n)⁻¹ ≤ (εn n)⁻¹ := by
      rw [inv_le_inv₀ huxε_pos (hεn_pos n)]
      exact hε_le
    have hpow_nonneg' : 0 ≤ |(u x + εn n)⁻¹| ^ p := by
      positivity
    calc
      ‖g n x‖ = ‖|(u x + εn n)⁻¹| ^ p‖ := by
        dsimp [g]
      _ = |(u x + εn n)⁻¹| ^ p := by
        rw [Real.norm_eq_abs, abs_of_nonneg hpow_nonneg']
      _ = ((u x + εn n)⁻¹) ^ p := by
        rw [abs_of_pos (inv_pos.mpr huxε_pos)]
      _ ≤ ((εn n)⁻¹) ^ p := by
        exact Real.rpow_le_rpow (inv_nonneg.mpr huxε_pos.le) hle_inv hp_nonneg
  have hshift_neg_eq :
      ∀ n,
        (⨍ x in Bhalf, (u x + εn n) ^ (-p) ∂volume) =
          ⨍ x in Bhalf, g n x ∂volume := by
    intro n
    exact MeasureTheory.setAverage_congr_fun measurableSet_ball <| by
      filter_upwards with x
      intro hx
      have huxε_pos : 0 < u x + εn n := by
        linarith [hu_pos_half x hx, hεn_pos n]
      calc
        (u x + εn n) ^ (-p) = |u x + εn n| ^ (-p) := by
          rw [abs_of_pos huxε_pos]
        _ = g n x := by
          dsimp [g]
          exact crossover_abs_rpow_neg_eq_abs_inv_rpow huxε_pos
  have hreg_n :
      ∀ n,
        (⨍ x in Bhalf, (u x + εn n) ^ p ∂volume) *
            (⨍ x in Bhalf, g n x ∂volume) ≤ K := by
    intro n
    have hreg :=
      regularized_crossover_product_bound (d := d) (A := A)
        hu_pos hsuper (hε := hεn_pos n)
    have hshift_neg_eq' :
        (⨍ x in Metric.ball (0 : E) (1 / 2 : ℝ), (u x + εn n) ^ (-p) ∂volume) =
          ⨍ x in Metric.ball (0 : E) (1 / 2 : ℝ), g n x ∂volume := by
      simpa [Bhalf] using hshift_neg_eq n
    have hregp :
        (⨍ x in Metric.ball (0 : E) (1 / 2 : ℝ), (u x + εn n) ^ p ∂volume) *
            (⨍ x in Metric.ball (0 : E) (1 / 2 : ℝ), (u x + εn n) ^ (-p) ∂volume) ≤
          crossoverJNKhalf d ^ 2 := by
      simpa [p] using hreg
    have hreg' :
        (⨍ x in Metric.ball (0 : E) (1 / 2 : ℝ), (u x + εn n) ^ p ∂volume) *
            (⨍ x in Metric.ball (0 : E) (1 / 2 : ℝ), g n x ∂volume) ≤
          crossoverJNKhalf d ^ 2 := by
      calc
        (⨍ x in Metric.ball (0 : E) (1 / 2 : ℝ), (u x + εn n) ^ p ∂volume) *
            (⨍ x in Metric.ball (0 : E) (1 / 2 : ℝ), g n x ∂volume)
            =
              (⨍ x in Metric.ball (0 : E) (1 / 2 : ℝ), (u x + εn n) ^ p ∂volume) *
                (⨍ x in Metric.ball (0 : E) (1 / 2 : ℝ), (u x + εn n) ^ (-p) ∂volume) := by
                  rw [hshift_neg_eq']
        _ ≤ crossoverJNKhalf d ^ 2 := hregp
    simpa [Bhalf, p, K] using hreg'
  have hgavg_nonneg :
      ∀ n, 0 ≤ ⨍ x in Bhalf, g n x ∂volume := by
    intro n
    rw [MeasureTheory.setAverage_eq, smul_eq_mul]
    refine mul_nonneg ?_ ?_
    · positivity
    · exact integral_nonneg fun x => by
        dsimp [g]
        positivity
  have hgavg_bound :
      ∀ n, ⨍ x in Bhalf, g n x ∂volume ≤ M := by
    intro n
    have hmul :
        Aavg * (⨍ x in Bhalf, g n x ∂volume) ≤ K := by
      have hleft :
          Aavg * (⨍ x in Bhalf, g n x ∂volume) ≤
            (⨍ x in Bhalf, (u x + εn n) ^ p ∂volume) *
              (⨍ x in Bhalf, g n x ∂volume) := by
        exact mul_le_mul_of_nonneg_right (hAavg_le_shift n) (hgavg_nonneg n)
      exact le_trans hleft (hreg_n n)
    exact (le_div_iff₀ hAavg_pos).2 (by simpa [M, mul_comm, mul_left_comm, mul_assoc] using hmul)
  have hg_int_bound :
      ∀ n, ∫ x, g n x ∂μ ≤ μ.real Set.univ * M := by
    intro n
    have havg := hgavg_bound n
    rw [MeasureTheory.setAverage_eq, smul_eq_mul] at havg
    rw [show ∫ x in Bhalf, g n x ∂volume = ∫ x, g n x ∂μ by
          simp [μ, Bhalf]] at havg
    rw [← hμ_eq] at havg
    rw [show (μ.real Set.univ)⁻¹ * ∫ x, g n x ∂μ =
        (∫ x, g n x ∂μ) / μ.real Set.univ by
          field_simp [hμ_pos.ne']] at havg
    have hmul : ∫ x, g n x ∂μ ≤ M * μ.real Set.univ := by
      exact (div_le_iff₀ hμ_pos).mp havg
    simpa [mul_comm, mul_left_comm, mul_assoc] using hmul
  have hright_eq :
      ∀ n,
        ∫⁻ x, ENNReal.ofReal (g n x) ∂μ =
          ENNReal.ofReal (∫ x, g n x ∂μ) := by
    intro n
    exact
      (MeasureTheory.ofReal_integral_eq_lintegral_ofReal
        (μ := μ) (f := g n) (hg_int n) (hg_nonneg n)).symm
  have hFatou :
      ∫⁻ x, ENNReal.ofReal (g0 x) ∂μ ≤
        atTop.liminf (fun n => ∫⁻ x, ENNReal.ofReal (g n x) ∂μ) := by
    have hmeas :
        ∀ n, AEMeasurable (fun x => ENNReal.ofReal (g n x)) μ := by
      intro n
      exact (hg_meas n).aemeasurable.ennreal_ofReal
    have hleft := MeasureTheory.lintegral_liminf_le' (μ := μ) hmeas
    have hlim :
        (fun x => Filter.liminf (fun n => ENNReal.ofReal (g n x)) atTop) =ᵐ[μ]
          fun x => ENNReal.ofReal (g0 x) := by
      filter_upwards [ae_restrict_mem measurableSet_ball] with x hx
      have hux_pos : 0 < u x := hu_pos_half x hx
      have hshift_tend :
          Filter.Tendsto (fun n => u x + εn n) atTop (nhds (u x)) := by
        simpa [εn, add_comm, one_div] using
          (tendsto_one_div_add_atTop_nhds_zero_nat.const_add (u x))
      have hinv_tend :
          Filter.Tendsto (fun n => (u x + εn n)⁻¹) atTop (nhds ((u x)⁻¹)) := by
        have hdiv :
            Filter.Tendsto (fun n : ℕ => (1 : ℝ) / (u x + εn n)) atTop
              (nhds ((1 : ℝ) / u x)) := by
          exact Tendsto.div tendsto_const_nhds hshift_tend hux_pos.ne'
        simpa [one_div] using hdiv
      have hreal_tend :
          Filter.Tendsto (fun n => g n x) atTop (nhds (g0 x)) := by
        dsimp [g, g0]
        exact ((Real.continuous_rpow_const hp_nonneg).continuousAt.tendsto.comp
          (continuous_abs.continuousAt.tendsto.comp hinv_tend))
      have henn_tend :
          Filter.Tendsto (fun n => ENNReal.ofReal (g n x)) atTop
            (nhds (ENNReal.ofReal (g0 x))) := by
        exact (ENNReal.continuous_ofReal.tendsto _).comp hreal_tend
      exact henn_tend.liminf_eq
    calc
      ∫⁻ x, ENNReal.ofReal (g0 x) ∂μ
          = ∫⁻ x, Filter.liminf (fun n => ENNReal.ofReal (g n x)) atTop ∂μ := by
              exact lintegral_congr_ae hlim.symm
      _ ≤ atTop.liminf (fun n => ∫⁻ x, ENNReal.ofReal (g n x) ∂μ) := hleft
  have hliminf_le :
      atTop.liminf (fun n => ∫⁻ x, ENNReal.ofReal (g n x) ∂μ) ≤
        ENNReal.ofReal (μ.real Set.univ * M) := by
    rw [Filter.liminf_congr (Eventually.of_forall hright_eq)]
    refine Filter.liminf_le_of_frequently_le' (Frequently.of_forall fun n => ?_)
    exact ENNReal.ofReal_le_ofReal (hg_int_bound n)
  have hmain_enn :
      ∫⁻ x, ENNReal.ofReal (g0 x) ∂μ ≤ ENNReal.ofReal (μ.real Set.univ * M) := by
    exact le_trans hFatou hliminf_le
  have hM_nonneg : 0 ≤ M := by
    dsimp [M, K]
    exact div_nonneg (by positivity) hAavg_pos.le
  have hg0_int : Integrable g0 μ := by
    have hlin_top : ∫⁻ x, ENNReal.ofReal (g0 x) ∂μ ≠ ⊤ := by
      exact ne_of_lt (lt_of_le_of_lt hmain_enn ENNReal.ofReal_lt_top)
    have hInt :=
      integrable_toReal_of_lintegral_ne_top
        (hg0_meas.aemeasurable.ennreal_ofReal) hlin_top
    refine hInt.congr ?_
    filter_upwards with x
    have hg0x_nonneg : 0 ≤ g0 x := by
      dsimp [g0]
      positivity
    rw [ENNReal.toReal_ofReal hg0x_nonneg]
  have hg0_int_bound :
      ∫ x, g0 x ∂μ ≤ μ.real Set.univ * M := by
    calc
      ∫ x, g0 x ∂μ = ENNReal.toReal (∫⁻ x, ENNReal.ofReal (g0 x) ∂μ) := by
            exact MeasureTheory.integral_eq_lintegral_of_nonneg_ae hg0_nonneg hg0_meas
      _ ≤ ENNReal.toReal (ENNReal.ofReal (μ.real Set.univ * M)) := by
            have hleft_ne_top : ∫⁻ x, ENNReal.ofReal (g0 x) ∂μ ≠ ⊤ := by
              exact ne_of_lt (lt_of_le_of_lt hmain_enn ENNReal.ofReal_lt_top)
            have hright_ne_top : ENNReal.ofReal (μ.real Set.univ * M) ≠ ⊤ := by
              exact ne_of_lt ENNReal.ofReal_lt_top
            exact (ENNReal.toReal_le_toReal hleft_ne_top hright_ne_top).2 hmain_enn
      _ = μ.real Set.univ * M := by
            rw [ENNReal.toReal_ofReal (mul_nonneg hμ_pos.le hM_nonneg)]
  have hg0_avg_bound :
      (⨍ x in Bhalf, g0 x ∂volume) ≤ M := by
    rw [MeasureTheory.setAverage_eq, smul_eq_mul, ← hμ_eq]
    rw [show (μ.real Set.univ)⁻¹ * ∫ x, g0 x ∂μ =
        (∫ x, g0 x ∂μ) / μ.real Set.univ by
          field_simp [hμ_pos.ne']]
    exact (div_le_iff₀ hμ_pos).2 (by simpa [mul_comm, mul_left_comm, mul_assoc] using hg0_int_bound)
  have hneg_target_eq :
      (⨍ x in Bhalf, |u x| ^ (-p) ∂volume) = ⨍ x in Bhalf, g0 x ∂volume := by
    exact MeasureTheory.setAverage_congr_fun measurableSet_ball <| by
      filter_upwards with x
      intro hx
      have hux_pos : 0 < u x := hu_pos_half x hx
      dsimp [g0]
      exact crossover_abs_rpow_neg_eq_abs_inv_rpow hux_pos
  rw [hneg_target_eq]
  calc
    (⨍ x in Bhalf, |u x| ^ p ∂volume) * (⨍ x in Bhalf, g0 x ∂volume)
        ≤ Aavg * M := by
          exact mul_le_mul_of_nonneg_left hg0_avg_bound hAavg_pos.le
    _ = K := by
          dsimp [M]
          field_simp [hAavg_pos.ne']
    _ ≤ C := by
          dsimp [C]
          exact le_max_right _ _

\end{lstlisting}
\begin{proof}
By a regularization step $u \mapsto u + \varepsilon$, it suffices to bound
\[
  J_\varepsilon := \Bigl(\fint_{\Bz{1/2}}(u+\varepsilon)^{p_0}\Bigr)
        \Bigl(\fint_{\Bz{1/2}}(u+\varepsilon)^{-p_0}\Bigr)
\]
uniformly in $\varepsilon > 0$ by $K_{1/2}(d)^2$. With $v_\varepsilon =
-\log(u+\varepsilon) + \log\varepsilon$, $a := \fint_{\Bz{1/48}} v_\varepsilon$,
$F_\pm(x) := \exp(\pm p_0(v_\varepsilon - a))$ and
$F_{\mathrm{abs}} := \exp(p_0|v_\varepsilon - a|)$, we have $F_\pm \leq
F_{\mathrm{abs}}$ pointwise. Theorem~\ref{thm:exp-half} together with
$\leanname{setAverage\_mono\_of\_nonneg\_local}$ yields
\[
  \fint_{\Bz{1/2}} F_+\,dx \leq K_{1/2}(d), \qquad
  \fint_{\Bz{1/2}} F_-\,dx \leq K_{1/2}(d).
\]
Now $(u+\varepsilon)^{p_0} = \exp(p_0(\log\varepsilon - a))\,F_+$ and
$(u+\varepsilon)^{-p_0} = \exp(p_0(a - \log\varepsilon))\,F_-$ pointwise on
$\Bz{1/2}$. The two scalar prefactors multiply to $1$ (cancellation
$\leanname{hcancel}$), so
\[
  J_\varepsilon \;=\; \Bigl(\fint F_+\Bigr)\Bigl(\fint F_-\Bigr)
    \;\leq\; K_{1/2}(d)^2.
\]
Letting $\varepsilon \downarrow 0$ along $\varepsilon_n = 1/(n+1)$, the
shifted integrals converge to the original ones by dominated/monotone
convergence (the integrability of $|u|^{p_0}$ is supplied by the
$W^{1,2}$ estimate using $p_0 < 1 \leq 2$, see
\leanname{hpow\_int}). Choose $C(d) := \max\{1, K_{1/2}(d)^2\}$.
\end{proof}

\begin{definition}[\leanname{C\_crossover'}]
\label{def:Ccrossover}
$C_{\mathrm{cross}}(d) := C(d)$ from Theorem~\ref{thm:crossover-exists}, so
$C_{\mathrm{cross}}(d) \geq 1$.
\end{definition}

\begin{lstlisting}[style=Matlab-editor,basicstyle=\mlttfamily,escapechar=",mlshowsectionrules=true,mathescape=true,morecomment={[s]{\%\{}{\%\}}},language=lean,numbers=none]
noncomputable def c_crossover' (d : ℕ) [NeZero d] : ℝ :=
  1 / (2 * C_JN d * c_crossover_bmo_scale d + 1)

\end{lstlisting}
\begin{theorem}[\leanname{crossover\_estimate}]
\label{thm:crossover-estimate}
Let $d \geq 3$, $A$ normalized elliptic on $\Bz 1$, and $u > 0$ a
supersolution. Then with $c := c'(d)/\Lambda_A^{1/2}$,
\[
  \Bigl(\fint_{\Bz{1/2}} u^c\,dx\Bigr)\Bigl(\fint_{\Bz{1/2}} u^{-c}\,dx\Bigr)
    \;\leq\; C_{\mathrm{cross}}(d).
\]
\end{theorem}

\begin{lstlisting}[style=Matlab-editor,basicstyle=\mlttfamily,escapechar=",mlshowsectionrules=true,mathescape=true,morecomment={[s]{\%\{}{\%\}}},language=lean,numbers=none]
/-! ### Crossover Estimate

Bridging positive and negative powers of a positive supersolution via
John-Nirenberg exponential integrability. -/

/-- Crossover estimate bridging positive and negative powers of a positive
supersolution. -/
theorem crossover_estimate
    (hd : 2 < (d : ℝ))
    (A : NormalizedEllipticCoeff d (Metric.ball (0 : E) 1))
    {u : E → ℝ}
    (hu_pos : ∀ x ∈ Metric.ball (0 : E) 1, 0 < u x)
    (hsuper : IsSupersolution A.1 u) :
    (⨍ x in Metric.ball (0 : E) (1 / 2 : ℝ),
        |u x| ^ (c_crossover' d / A.1.Λ ^ ((1 : ℝ) / 2)) ∂volume) *
      (⨍ x in Metric.ball (0 : E) (1 / 2 : ℝ),
        |u x| ^ (-(c_crossover' d / A.1.Λ ^ ((1 : ℝ) / 2))) ∂volume) ≤
        C_crossover' d := by
  -- Proof via local exponential integrability for `v = -log u`.
  -- Set v = -log u, c = c_crossover' d / Λ^{1/2}.
  -- Define w(x) = exp(c · (v_avg - v(x))) where v_avg = ⨍ v on B_{1/2}.
  -- Then log w = c(v_avg - v), so (log w)_{B_{1/2}} = 0.
  -- Also: w > 0 everywhere (exponential), w⁻¹(x) = exp(c(v(x) - v_avg)).
  -- The LHS equals ⨍ w · ⨍ w⁻¹ (the exp(±c·v_avg) factors cancel in the product).
  --
  --   Choose cutoff φ with φ = 1 on B_{3/4}, supp φ ⊆ B₁, |∇φ| ≤ C.
  --   ∫_{B_{3/4}} |∇v|² ≤ ∫_B φ²|∇v|² ≤ 4Λ ∫_B |∇φ|² ≤ C·Λ.
  --
  --   With c = c_crossover/Λ^{1/2}, this is ≤ c_crossover² · C.
  --   Choose c_crossover small enough → ≤ 1.
  --
  --   ∫_{B_{1/2}} |log w| ≤ C_P · ‖∇ log w‖_{L²} ≤ C_P.
  --
  --   dimension-only bounds for ⨍ w and ⨍ w⁻¹.
  --
  -- LHS = ⨍ |u|^c · ⨍ |u|^{-c}
  --      = ⨍ exp(c·log|u|) · ⨍ exp(-c·log|u|)
  --      = ⨍ exp(-c·v) · ⨍ exp(c·v)            (v = -log u)
  --      = exp(-c·v_avg) · ⨍ exp(c(v_avg-v)) · exp(c·v_avg) · ⨍ exp(c(v-v_avg))
  --        ^^^ the exp(±c·v_avg) factors cancel in the product ^^^
  --      = ⨍ w · ⨍ w⁻¹
  -- === Proof assembly ===
  set c := c_crossover' d / A.1.Λ ^ ((1 : ℝ) / 2)
  set v : E → ℝ := fun x => -Real.log (u x)
  set v_avg := ⨍ x in Metric.ball (0 : E) (1 / 2 : ℝ), v x ∂volume
  set w : E → ℝ := fun x => Real.exp (c * (v_avg - v x))
  set w_inv : E → ℝ := fun x => Real.exp (c * (v x - v_avg))
  have hw_pos : ∀ x, 0 < w x := fun x => Real.exp_pos _
  have hw_inv_pos : ∀ x, 0 < w_inv x := fun x => Real.exp_pos _
  have hw_mul : ∀ x, w x * w_inv x = 1 := by
    intro x; simp only [w, w_inv]
    rw [← Real.exp_add, show c * (v_avg - v x) + c * (v x - v_avg) = 0 by ring]
    exact Real.exp_zero
  have h_product_identity :
      (⨍ x in Metric.ball (0 : E) (1 / 2 : ℝ), |u x| ^ c ∂volume) *
        (⨍ x in Metric.ball (0 : E) (1 / 2 : ℝ), |u x| ^ (-c) ∂volume) =
      (⨍ x in Metric.ball (0 : E) (1 / 2 : ℝ), w x ∂volume) *
        (⨍ x in Metric.ball (0 : E) (1 / 2 : ℝ), w_inv x ∂volume) := by
    let B : Set E := Metric.ball (0 : E) (1 / 2 : ℝ)
    have hu_pos_half : ∀ x ∈ B, 0 < u x := by
      intro x hx
      exact hu_pos x (Metric.ball_subset_ball (by norm_num : (1 : ℝ) / 2 ≤ 1) hx)
    have hpow_eq :
        ∀ x ∈ B, |u x| ^ c = Real.exp (-c * v_avg) * w x := by
      intro x hx
      have hux : 0 < u x := hu_pos_half x hx
      have habs : |u x| = u x := abs_of_pos hux
      rw [habs, Real.rpow_def_of_pos hux]
      simp only [w, v]
      rw [← Real.exp_add]
      congr 1
      ring
    have hpow_inv_eq :
        ∀ x ∈ B, |u x| ^ (-c) = Real.exp (c * v_avg) * w_inv x := by
      intro x hx
      have hux : 0 < u x := hu_pos_half x hx
      have habs : |u x| = u x := abs_of_pos hux
      rw [habs, Real.rpow_def_of_pos hux]
      simp only [w_inv, v]
      rw [← Real.exp_add]
      congr 1
      ring
    have havg_pow :
        (⨍ x in B, |u x| ^ c ∂volume) =
          Real.exp (-c * v_avg) * (⨍ x in B, w x ∂volume) := by
      calc
        (⨍ x in B, |u x| ^ c ∂volume)
            = (volume.real B)⁻¹ * ∫ x in B, |u x| ^ c ∂volume := by
                rw [MeasureTheory.setAverage_eq, smul_eq_mul]
        _ = (volume.real B)⁻¹ * (Real.exp (-c * v_avg) * ∫ x in B, w x ∂volume) := by
              congr 1
              calc
                ∫ x in B, |u x| ^ c ∂volume
                    = ∫ x in B, Real.exp (-c * v_avg) * w x ∂volume := by
                        apply MeasureTheory.setIntegral_congr_fun measurableSet_ball
                        intro x hx
                        exact hpow_eq x hx
                _ = Real.exp (-c * v_avg) * ∫ x in B, w x ∂volume := by
                      rw [integral_const_mul]
        _ = Real.exp (-c * v_avg) * (⨍ x in B, w x ∂volume) := by
              rw [MeasureTheory.setAverage_eq, smul_eq_mul]
              ring
    have havg_pow_inv :
        (⨍ x in B, |u x| ^ (-c) ∂volume) =
          Real.exp (c * v_avg) * (⨍ x in B, w_inv x ∂volume) := by
      calc
        (⨍ x in B, |u x| ^ (-c) ∂volume)
            = (volume.real B)⁻¹ * ∫ x in B, |u x| ^ (-c) ∂volume := by
                rw [MeasureTheory.setAverage_eq, smul_eq_mul]
        _ = (volume.real B)⁻¹ * (Real.exp (c * v_avg) * ∫ x in B, w_inv x ∂volume) := by
              congr 1
              calc
                ∫ x in B, |u x| ^ (-c) ∂volume
                    = ∫ x in B, Real.exp (c * v_avg) * w_inv x ∂volume := by
                        apply MeasureTheory.setIntegral_congr_fun measurableSet_ball
                        intro x hx
                        exact hpow_inv_eq x hx
                _ = Real.exp (c * v_avg) * ∫ x in B, w_inv x ∂volume := by
                      rw [integral_const_mul]
        _ = Real.exp (c * v_avg) * (⨍ x in B, w_inv x ∂volume) := by
              rw [MeasureTheory.setAverage_eq, smul_eq_mul]
              ring
    rw [havg_pow, havg_pow_inv]
    calc
      (Real.exp (-c * v_avg) * (⨍ x in B, w x ∂volume)) *
          (Real.exp (c * v_avg) * (⨍ x in B, w_inv x ∂volume))
          = (Real.exp (-c * v_avg) * Real.exp (c * v_avg)) *
              ((⨍ x in B, w x ∂volume) * (⨍ x in B, w_inv x ∂volume)) := by
                ring
      _ = (⨍ x in B, w x ∂volume) * (⨍ x in B, w_inv x ∂volume) := by
            rw [show Real.exp (-c * v_avg) * Real.exp (c * v_avg) = 1 by
              rw [← Real.exp_add]
              simp]
            simp
  exact C_crossover'_spec.2 hd A hu_pos hsuper

\end{lstlisting}
\begin{proof}
Set $v(x) := -\log u(x)$, $v_{\mathrm{avg}} := \fint_{\Bz{1/2}} v\,dx$,
$w(x) := \exp(c(v_{\mathrm{avg}} - v(x)))$, and
$w^{-1}(x) := \exp(c(v(x) - v_{\mathrm{avg}}))$, so $w \cdot w^{-1} \equiv 1$.
On $\Bz{1/2}$ where $u > 0$, $u^c = e^{-c v_{\mathrm{avg}}}\,w$ and
$u^{-c} = e^{c v_{\mathrm{avg}}}\,w^{-1}$, hence
\[
  \Bigl(\fint u^c\Bigr)\Bigl(\fint u^{-c}\Bigr)
    \;=\; e^{-cv_{\mathrm{avg}}}e^{cv_{\mathrm{avg}}}
        \Bigl(\fint w\Bigr)\Bigl(\fint w^{-1}\Bigr)
    \;=\; \Bigl(\fint w\Bigr)\Bigl(\fint w^{-1}\Bigr).
\]
Apply Theorem~\ref{thm:crossover-exists}.
\end{proof}

\begin{corollary}[\leanname{crossover\_estimate\_unaveraged}]
\label{cor:crossover-unaveraged}
Multiplying through by $\vol(\Bz{1/2})^2 > 0$,
\[
  \Bigl(\int_{\Bz{1/2}} u^c\Bigr)\Bigl(\int_{\Bz{1/2}} u^{-c}\Bigr)
    \;\leq\; C_{\mathrm{cross}}(d)\,\vol(\Bz{1/2})^2.
\]
\end{corollary}

\begin{lstlisting}[style=Matlab-editor,basicstyle=\mlttfamily,escapechar=",mlshowsectionrules=true,mathescape=true,morecomment={[s]{\%\{}{\%\}}},language=lean,numbers=none]
/-- Un-averaged half-ball version of `crossover_estimate`.

This is the form needed by downstream weak-Harnack bookkeeping: multiply the
average-product estimate by `|B_{1/2}|^2`. -/
theorem crossover_estimate_unaveraged
    (hd : 2 < (d : ℝ))
    (A : NormalizedEllipticCoeff d (Metric.ball (0 : E) 1))
    {u : E → ℝ}
    (hu_pos : ∀ x ∈ Metric.ball (0 : E) 1, 0 < u x)
    (hsuper : IsSupersolution A.1 u) :
    (∫ x in Metric.ball (0 : E) (1 / 2 : ℝ),
        |u x| ^ (c_crossover' d / A.1.Λ ^ ((1 : ℝ) / 2)) ∂volume) *
      (∫ x in Metric.ball (0 : E) (1 / 2 : ℝ),
        |u x| ^ (-(c_crossover' d / A.1.Λ ^ ((1 : ℝ) / 2))) ∂volume) ≤
        C_crossover' d *
          (volume.real (Metric.ball (0 : E) (1 / 2 : ℝ))) ^ 2 := by
  let B : Set E := Metric.ball (0 : E) (1 / 2 : ℝ)
  set p₀ : ℝ := c_crossover' d / A.1.Λ ^ ((1 : ℝ) / 2)
  set Ipos : ℝ := ∫ x in B, |u x| ^ p₀ ∂volume
  set Ineg : ℝ := ∫ x in B, |u x| ^ (-p₀) ∂volume
  have havg := crossover_estimate (d := d) hd A hu_pos hsuper
  have hvol_pos : 0 < volume.real B := by
    exact ENNReal.toReal_pos
      (measure_ball_pos volume (0 : E) (by norm_num : (0 : ℝ) < 1 / 2)).ne'
      measure_ball_lt_top.ne
  have hvol_ne : volume.real B ≠ 0 := ne_of_gt hvol_pos
  have hscale_nonneg : 0 ≤ (volume.real B) ^ 2 := by positivity
  rw [MeasureTheory.setAverage_eq, MeasureTheory.setAverage_eq, smul_eq_mul, smul_eq_mul] at havg
  have hscaled :=
    mul_le_mul_of_nonneg_left havg hscale_nonneg
  calc
    Ipos * Ineg
        = (volume.real B) ^ 2 * ((volume.real B)⁻¹ * Ipos * ((volume.real B)⁻¹ * Ineg)) := by
            field_simp [hvol_ne]
    _ ≤ (volume.real B) ^ 2 * C_crossover' d := hscaled
    _ = C_crossover' d * (volume.real B) ^ 2 := by ring


\end{lstlisting}
%% --- Section 10 ---
\subsection{Weak Harnack Inequality}

\begin{definition}[\leanname{weakHarnackP0}, \leanname{weak\_harnack\_decay\_exp}]
\label{def:p0}
For a normalized elliptic $A$ on $\Bz 1$, set
\[
  p_0 := p_0(A) := \frac{c'(d)}{\Lambda_A^{1/2}}, \qquad
  \beta(d) := \frac{d\,\chi(d)}{c'(d)},
\]
where $\chi(d) := d/(d-2)$ is the Sobolev exponent (\leanname{moserChi}).
\end{definition}

\begin{lstlisting}[style=Matlab-editor,basicstyle=\mlttfamily,escapechar=",mlshowsectionrules=true,mathescape=true,morecomment={[s]{\%\{}{\%\}}},language=lean,numbers=none]
/-- The crossover exponent `p₀ = c_crossover'(d) / Λ^{1/2}`. -/
noncomputable def weakHarnackP0
    (A : NormalizedEllipticCoeff d (ball (0 : E) 1)) : ℝ :=
  c_crossover' d / A.1.Λ ^ ((1 : ℝ) / 2)

/-! ### Constants -/

/-- The exponent `dχ / c_crossover'(d)` on `(1-q)`. -/
noncomputable def weak_harnack_decay_exp (d : ℕ) [NeZero d] : ℝ :=
  ((d : ℝ) * moserChi d) / c_crossover' d

/-- The Sobolev gain factor `χ = d/(d-2)` in the elliptic Moser iteration. -/
noncomputable def moserChi (d : ℕ) [NeZero d] : ℝ :=
  (d : ℝ) / ((d : ℝ) - 2)

\end{lstlisting}
\begin{lemma}[\leanname{p\_0\_pos}, \leanname{p\_0\_lt\_one}, \leanname{Lambda\_mul\_p\_0\_sq}]
$0 < p_0(A) < 1$, $p_0(A) \leq c'(d)$, and $\Lambda_A\,p_0(A)^2 = c'(d)^2$.
\end{lemma}

\begin{proof}
Since $\Lambda_A \geq 1$, $\Lambda_A^{1/2} \geq 1$, so
$0 < p_0 = c'(d)/\Lambda_A^{1/2} \leq c'(d) < 1$. The identity follows from
$(c'(d)/\Lambda_A^{1/2})^2 \cdot \Lambda_A = c'(d)^2$.
\end{proof}

The proof of the weak Harnack inequality chains three pieces at half scale.
We rescale $v(z) := u(z/2)$ from $\Bz{1/2}$ to $\Bz 1$ via the rescaled
coefficient $A' := \leanname{rescaledCoeff}\,A$ which preserves $\Lambda$.
The three steps are:

\begin{itemize}
\item \emph{Forward Moser at level $p_0$}
  (\leanname{weak\_harnack\_stage\_one\_forward\_ball}):
\begin{equation*}
\Bigl(\int_{\Bz{1/4}} u^{q\,d/(d-2)}\Bigr)^{p_0(d-2)/(qd)}
\leq C_{\mathrm{wH},0\to}(d)\,(\Lambda_A p_0^2/(1-q)^2 + 1)^{d\chi/2}
\int_{\Bz{1/2}} u^{p_0}.
\end{equation*}

\item \emph{Crossover} (Corollary~\ref{cor:crossover-unaveraged}) at the level
$p_0$:
\begin{equation*}
\Bigl(\int_{\Bz{1/2}} u^{p_0}\Bigr)\Bigl(\int_{\Bz{1/2}} u^{-p_0}\Bigr)
\leq C_{\mathrm{cross}}(d)\,\vol(\Bz{1/2})^2.
\end{equation*}

\item \emph{Inverse Moser}
  (\leanname{weak\_harnack\_stage\_one\_inverse\_ball}):
\begin{equation*}
(\Lambda_A p_0^2 + 1)^{-d/2}\Bigl(\int_{\Bz{1/2}} u^{-p_0}\Bigr)^{-1}
\leq C_{\mathrm{wH},0}(d)\,\bigl(\essinf_{\Bz{1/4}} u\bigr)^{p_0}.
\end{equation*}
\end{itemize}

Multiplying these three estimates yields the chain.

\begin{theorem}[\leanname{weak\_harnack\_chain}]
\label{thm:wH-chain}
For $d \geq 3$, $A$ normalized on $\Bz 1$, $u > 0$ a supersolution, and
$0 < q < 1$:
\begin{equation*}
\Bigl(\int_{\Bz{1/4}} u^{q d/(d-2)}\,dx\Bigr)^{p_0(d-2)/(qd)}
\leq \frac{C_{\mathrm{chain}}(d)}{(1-q)^{d\chi(d)}}
\Bigl(\essinf_{\Bz{1/4}} u\Bigr)^{p_0}.
\end{equation*}
\end{theorem}

\begin{lstlisting}[style=Matlab-editor,basicstyle=\mlttfamily,escapechar=",mlshowsectionrules=true,mathescape=true,morecomment={[s]{\%\{}{\%\}}},language=lean,numbers=none]
/-- **Three-stage chain** at the `p₀`-power level, on `B_{1/4}`.

Uses the un-powered constant `C_chain` (not `C_weakHarnack`). The `1/c'`
power is introduced when raising to `1/p₀` in `weak_harnack`. -/
theorem weak_harnack_chain
    (hd : 2 < (d : ℝ))
    (A : NormalizedEllipticCoeff d (ball (0 : E) 1))
    {u : E → ℝ} {q : ℝ}
    (hq : 0 < q) (hq1 : q < 1)
    (hu_pos : ∀ x ∈ ball (0 : E) 1, 0 < u x)
    (hsuper : IsSupersolution A.1 u) :
    let p₀ := weakHarnackP0 A
    (∫ x in ball (0 : E) (1 / 4 : ℝ),
        |u x| ^ (q * (d : ℝ) / ((d : ℝ) - 2)) ∂volume) ^
          (p₀ * (((d : ℝ) - 2) / (q * (d : ℝ)))) ≤
      C_chain (d := d) hd / (1 - q) ^ ((d : ℝ) * moserChi d) *
        (essInf u (volume.restrict (ball (0 : E) (1 / 4 : ℝ)))) ^ p₀ := by
  intro p₀
  have hp₀_def : p₀ = weakHarnackP0 A := by rfl
  have hμquarter_ne_zero :
      volume.restrict (ball (0 : E) (1 / 4 : ℝ)) ≠ 0 := by
    intro hzero
    have hball_zero : volume (ball (0 : E) (1 / 4 : ℝ)) = 0 := by
      simpa [Measure.restrict_apply_univ] using
        congrArg (fun μ : Measure E => μ Set.univ) hzero
    exact (Metric.measure_ball_pos volume (0 : E) (by norm_num : (0 : ℝ) < 1 / 4)).ne'
      hball_zero
  have hquarter_nonneg :
      ∀ᵐ x ∂(volume.restrict (ball (0 : E) (1 / 4 : ℝ))), 0 ≤ u x := by
    filter_upwards [ae_restrict_mem measurableSet_ball] with x hx
    exact (hu_pos x ((Metric.ball_subset_ball (by norm_num : (1 / 4 : ℝ) ≤ 1)) hx)).le
  have hinf_nonneg :
      0 ≤ essInf u (volume.restrict (ball (0 : E) (1 / 4 : ℝ))) := by
    exact le_essInf_real_of_ae_le (d := d) hμquarter_ne_zero hquarter_nonneg
  by_cases hpq : p₀ < q
  · have hmain := weak_harnack_chain_main
        (d := d) hd A (u := u) (q := q) hq hq1 hu_pos hsuper hpq
    have hC_le :
        C_chainMain (d := d) hd / (1 - q) ^ ((d : ℝ) * moserChi d) ≤
          C_chain (d := d) hd / (1 - q) ^ ((d : ℝ) * moserChi d) := by
      let β : ℝ := (d : ℝ) * moserChi d
      have hden_pos : 0 < (1 - q) ^ β := by
        have h1q_pos : 0 < 1 - q := by linarith [hq1]
        exact Real.rpow_pos_of_pos h1q_pos β
      exact (div_le_div_iff_of_pos_right hden_pos).2 <| by
        simpa [β] using (C_chainMain_le_C_chain (d := d) hd)
    have hpow_nonneg :
        0 ≤ essInf u (volume.restrict (ball (0 : E) (1 / 4 : ℝ))) ^ p₀ := by
      exact Real.rpow_nonneg hinf_nonneg p₀
    have hfinal :
        C_chainMain (d := d) hd / (1 - q) ^ ((d : ℝ) * moserChi d) *
          (essInf u (volume.restrict (ball (0 : E) (1 / 4 : ℝ)))) ^ p₀ ≤
        C_chain (d := d) hd / (1 - q) ^ ((d : ℝ) * moserChi d) *
          (essInf u (volume.restrict (ball (0 : E) (1 / 4 : ℝ)))) ^ p₀ := by
      exact mul_le_mul_of_nonneg_right hC_le hpow_nonneg
    exact le_trans hmain hfinal
  · push_neg at hpq
    let q' : ℝ := (p₀ + 1) / 2
    let q_star : ℝ := q * (d : ℝ) / ((d : ℝ) - 2)
    let q'_star : ℝ := q' * (d : ℝ) / ((d : ℝ) - 2)
    have hp₀_pos : 0 < p₀ := by
      simpa using p₀_pos A
    have hp₀_lt : p₀ < 1 := by
      simpa using p₀_lt_one A
    have hq'_pos : 0 < q' := by
      dsimp [q']
      linarith
    have hq'_lt_one : q' < 1 := by
      dsimp [q']
      linarith
    have hp₀_lt_q' : p₀ < q' := by
      dsimp [q']
      linarith
    have hq_le_q' : q ≤ q' := by
      exact le_trans hpq (le_of_lt hp₀_lt_q')
    have hq_star_pos : 0 < q_star := by
      dsimp [q_star]
      have hdm2 : 0 < (d : ℝ) - 2 := by linarith
      exact div_pos (mul_pos hq (by positivity)) hdm2
    have hq'_star_pos : 0 < q'_star := by
      dsimp [q'_star]
      have hdm2 : 0 < (d : ℝ) - 2 := by linarith
      exact div_pos (mul_pos hq'_pos (by positivity)) hdm2
    have hqstar_le : q_star ≤ q'_star := by
      have hdm2_pos : 0 < (d : ℝ) - 2 := by linarith
      have hnum : q * (d : ℝ) ≤ q' * (d : ℝ) := by
        nlinarith [hq_le_q']
      simpa [q_star, q'_star] using
        (div_le_div_iff_of_pos_right hdm2_pos).2 hnum
    have hmain_q' := weak_harnack_chain_main
      (d := d) hd A (u := u) (q := q') hq'_pos hq'_lt_one hu_pos hsuper hp₀_lt_q'
    have hq'_le_two : q' ≤ 2 := by
      dsimp [q']
      linarith
    let huW : MemW1pWitness 2 u (ball (0 : E) 1) := hsuper.1.someWitness
    have hq'_int_ball1 :
        IntegrableOn (fun x => |u x| ^ q') (ball (0 : E) 1) volume := by
      haveI : IsFiniteMeasure (volume.restrict (ball (0 : E) 1)) := by
        rw [isFiniteMeasure_restrict]
        exact measure_ball_lt_top.ne
      have hmem_q' : MemLp u (ENNReal.ofReal q') (volume.restrict (ball (0 : E) 1)) := by
        exact huW.memLp.mono_exponent <| by
          simpa using (ENNReal.ofReal_le_ofReal hq'_le_two)
      have hint_q' :
          Integrable (fun x => ‖u x‖ ^ q') (volume.restrict (ball (0 : E) 1)) := by
        have hint' :=
          hmem_q'.integrable_norm_rpow (ENNReal.ofReal_pos.mpr hq'_pos).ne' ENNReal.ofReal_ne_top
        simpa [ENNReal.toReal_ofReal hq'_pos.le] using hint'
      simpa [IntegrableOn, Real.norm_eq_abs] using hint_q'
    have hq'_int_half :
        IntegrableOn (fun x => |u x| ^ q') (ball (0 : E) (1 / 2 : ℝ)) volume := by
      exact hq'_int_ball1.mono_set (Metric.ball_subset_ball (by norm_num : (1 / 2 : ℝ) ≤ 1))
    have hstep2 := supersolution_reverseHolder_forward
      (d := d) hd A (u := u) (p := q') (r := (1 / 4 : ℝ)) (s := (1 / 2 : ℝ))
      hq'_pos hq'_lt_one (by norm_num) (by norm_num) (by norm_num)
      hu_pos hsuper hq'_int_half
    have hq'_star_int :
        IntegrableOn (fun x => |u x| ^ q'_star) (ball (0 : E) (1 / 4 : ℝ)) volume := by
      have hq'_star_eq : moserChi d * q' = q'_star := by
        dsimp [q'_star, moserChi]
        rw [div_eq_mul_inv, div_eq_mul_inv]
        ring
      simpa [hq'_star_eq] using hstep2.1
    have hμquarter_le :
        volume.restrict (ball (0 : E) (1 / 4 : ℝ)) ≤
          volume.restrict (ball (0 : E) 1) := by
      have hsubset : ball (0 : E) (1 / 4 : ℝ) ⊆ ball (0 : E) 1 :=
        Metric.ball_subset_ball (by norm_num : (1 / 4 : ℝ) ≤ 1)
      exact Measure.restrict_mono_set volume hsubset
    have hu_aesm_quarter :
        AEStronglyMeasurable u (volume.restrict (ball (0 : E) (1 / 4 : ℝ))) := by
      let huW : MemW1pWitness 2 u (ball (0 : E) 1) := hsuper.1.someWitness
      exact huW.memLp.aestronglyMeasurable.mono_measure hμquarter_le
    have hnorm_mono :
        (∫ x in ball (0 : E) (1 / 4 : ℝ), |u x| ^ q_star ∂volume) ^ (1 / q_star) ≤
          (∫ x in ball (0 : E) (1 / 4 : ℝ), |u x| ^ q'_star ∂volume) ^ (1 / q'_star) := by
      exact quarterBall_lpNorm_mono
        (u := u) (r := q_star) (s := q'_star) hq_star_pos hqstar_le
        hu_aesm_quarter hq'_star_int
    have hIq_nonneg :
        0 ≤ ∫ x in ball (0 : E) (1 / 4 : ℝ), |u x| ^ q_star ∂volume := by
      exact integral_nonneg fun x => by positivity
    have hIq'_nonneg :
        0 ≤ ∫ x in ball (0 : E) (1 / 4 : ℝ), |u x| ^ q'_star ∂volume := by
      exact integral_nonneg fun x => by positivity
    have hpow_mono :
        (∫ x in ball (0 : E) (1 / 4 : ℝ), |u x| ^ q_star ∂volume) ^ (p₀ / q_star) ≤
          (∫ x in ball (0 : E) (1 / 4 : ℝ), |u x| ^ q'_star ∂volume) ^ (p₀ / q'_star) := by
      have hraw :
          ((∫ x in ball (0 : E) (1 / 4 : ℝ), |u x| ^ q_star ∂volume) ^ (1 / q_star)) ^ p₀ ≤
            ((∫ x in ball (0 : E) (1 / 4 : ℝ), |u x| ^ q'_star ∂volume) ^ (1 / q'_star)) ^ p₀ := by
        exact Real.rpow_le_rpow (by positivity) hnorm_mono hp₀_pos.le
      have hleft :
          ((∫ x in ball (0 : E) (1 / 4 : ℝ), |u x| ^ q_star ∂volume) ^ (1 / q_star)) ^ p₀ =
            (∫ x in ball (0 : E) (1 / 4 : ℝ), |u x| ^ q_star ∂volume) ^ (p₀ / q_star) := by
        rw [← Real.rpow_mul hIq_nonneg]
        rw [show (1 / q_star) * p₀ = p₀ / q_star by field_simp [hq_star_pos.ne']]
      have hright :
          ((∫ x in ball (0 : E) (1 / 4 : ℝ), |u x| ^ q'_star ∂volume) ^ (1 / q'_star)) ^ p₀ =
            (∫ x in ball (0 : E) (1 / 4 : ℝ), |u x| ^ q'_star ∂volume) ^ (p₀ / q'_star) := by
        rw [← Real.rpow_mul hIq'_nonneg]
        rw [show (1 / q'_star) * p₀ = p₀ / q'_star by field_simp [hq'_star_pos.ne']]
      rwa [hleft, hright] at hraw
    have hq_pow :
        p₀ / q_star = p₀ * (((d : ℝ) - 2) / (q * (d : ℝ))) := by
      dsimp [q_star]
      field_simp [hq.ne', show (d : ℝ) - 2 ≠ 0 by linarith]
    have hq'_pow :
        p₀ / q'_star = p₀ * (((d : ℝ) - 2) / (q' * (d : ℝ))) := by
      dsimp [q'_star]
      field_simp [hq'_pos.ne', show (d : ℝ) - 2 ≠ 0 by linarith]
    have hmain_to_public :
        C_chainMain (d := d) hd / (1 - q') ^ ((d : ℝ) * moserChi d) ≤
          C_chain (d := d) hd / (1 - q) ^ ((d : ℝ) * moserChi d) := by
      let β : ℝ := (d : ℝ) * moserChi d
      have h1q_pos : 0 < 1 - q := by linarith [hq1]
      have h1q_nonneg : 0 ≤ 1 - q := h1q_pos.le
      have hp₀_le_c : p₀ ≤ c_crossover' d := by
        simpa [hp₀_def] using (p₀_le_crossover' (d := d) A)
      have h1q'_eq : 1 - q' = (1 - p₀) / 2 := by
        dsimp [q']
        ring
      have hβ_pos : 0 < β := by
        dsimp [β]
        have hchi_pos : 0 < moserChi d := moserChi_pos (d := d) hd
        positivity
      have h1q'_pos : 0 < 1 - q' := by
        linarith [hq'_lt_one]
      have hsmall_pos : 0 < (1 - c_crossover' d) / 2 := by
        have hc_lt_one : c_crossover' d < 1 := c_crossover'_lt_one (d := d)
        linarith
      have hsmall_le : (1 - c_crossover' d) / 2 ≤ 1 - q' := by
        rw [h1q'_eq]
        linarith
      have hpow_small_le :
          ((1 - c_crossover' d) / 2) ^ β ≤ (1 - q') ^ β := by
        exact Real.rpow_le_rpow hsmall_pos.le hsmall_le hβ_pos.le
      have hpow_small_pos : 0 < ((1 - c_crossover' d) / 2) ^ β := by
        exact Real.rpow_pos_of_pos hsmall_pos _
      have hto_chain :
          C_chainMain (d := d) hd / (1 - q') ^ β ≤ C_chain (d := d) hd := by
        calc
          C_chainMain (d := d) hd / (1 - q') ^ β
              ≤ C_chainMain (d := d) hd / (((1 - c_crossover' d) / 2) ^ β) := by
                  have hmain_nonneg : 0 ≤ C_chainMain (d := d) hd := by
                    exact le_trans (by norm_num : (0 : ℝ) ≤ 1) (one_le_C_chainMain (d := d) hd)
                  exact div_le_div_of_nonneg_left hmain_nonneg hpow_small_pos hpow_small_le
          _ = C_chain (d := d) hd := by
                rw [C_chain, div_eq_mul_inv]
      have hpow_q_le_one : (1 - q) ^ β ≤ 1 := by
        have h1q_le_one : 1 - q ≤ 1 := by linarith [hq]
        exact Real.rpow_le_one h1q_nonneg h1q_le_one hβ_pos.le
      have hpow_q_pos : 0 < (1 - q) ^ β := by
        exact Real.rpow_pos_of_pos h1q_pos β
      have hC_nonneg : 0 ≤ C_chain (d := d) hd := by
        exact le_trans (by norm_num : (0 : ℝ) ≤ 1) (one_le_C_chain (d := d) hd)
      have hfrom_chain :
          C_chain (d := d) hd ≤ C_chain (d := d) hd / (1 - q) ^ β := by
        have hinv_ge_one : 1 ≤ ((1 - q) ^ β)⁻¹ := by
          exact (one_le_inv₀ hpow_q_pos).2 hpow_q_le_one
        calc
          C_chain (d := d) hd = C_chain (d := d) hd * 1 := by ring
          _ ≤ C_chain (d := d) hd * ((1 - q) ^ β)⁻¹ := by
                exact mul_le_mul_of_nonneg_left hinv_ge_one hC_nonneg
          _ = C_chain (d := d) hd / (1 - q) ^ β := by
                rw [div_eq_mul_inv]
      exact le_trans hto_chain hfrom_chain
    have hpow_nonneg :
        0 ≤ essInf u (volume.restrict (ball (0 : E) (1 / 4 : ℝ))) ^ p₀ := by
      exact Real.rpow_nonneg hinf_nonneg p₀
    calc
      (∫ x in ball (0 : E) (1 / 4 : ℝ),
        |u x| ^ q_star ∂volume) ^ (p₀ * (((d : ℝ) - 2) / (q * (d : ℝ))))
        = (∫ x in ball (0 : E) (1 / 4 : ℝ),
            |u x| ^ q_star ∂volume) ^ (p₀ / q_star) := by rw [hq_pow]
      _
        ≤ (∫ x in ball (0 : E) (1 / 4 : ℝ),
            |u x| ^ q'_star ∂volume) ^ (p₀ / q'_star) := hpow_mono
      _ ≤ C_chainMain (d := d) hd / (1 - q') ^ ((d : ℝ) * moserChi d) *
            (essInf u (volume.restrict (ball (0 : E) (1 / 4 : ℝ)))) ^ p₀ := by
              have hmain_q'' :
                  (∫ x in ball (0 : E) (1 / 4 : ℝ), |u x| ^ q'_star ∂volume) ^ (p₀ / q'_star) ≤
                    C_chainMain (d := d) hd / (1 - q') ^ ((d : ℝ) * moserChi d) *
                      (essInf u (volume.restrict (ball (0 : E) (1 / 4 : ℝ)))) ^ p₀ := by
                rw [hq'_pow]
                simpa [hp₀_def, q'_star] using hmain_q'
              exact hmain_q''
      _ ≤ C_chain (d := d) hd / (1 - q) ^ ((d : ℝ) * moserChi d) *
            (essInf u (volume.restrict (ball (0 : E) (1 / 4 : ℝ)))) ^ p₀ := by
              exact mul_le_mul_of_nonneg_right hmain_to_public hpow_nonneg

\end{lstlisting}
\begin{proof}
We split into the case $p_0 < q$ (the natural case) and the case $q \leq p_0$
(handled by interpolating with $q' := (p_0+1)/2$). Assume $p_0 < q$ first.

By regularization $u \mapsto u + \delta$ (which preserves the supersolution
property by \leanname{isSupersolution\_add\_const\_unitBall}), and using the
identity $\Lambda_A p_0^2 = c'(d)^2$, the product of the forward Moser
estimate, the crossover bound and the inverse Moser estimate gives
\begin{multline*}
  \Bigl(\int_{\Bz{1/4}}(u+\delta)^{q d/(d-2)}\Bigr)^{p_0(d-2)/(q d)}
  \times \bigl(c'(d)^2 + 1\bigr)^{d/2}
  \\[-2pt]
  \leq C_{\mathrm{wH},0\to}\,C_{\mathrm{wH},0}\,C_{\mathrm{cross}}\,
       \vol(\Bz{1/2})^2\,(c'(d)^2/(1-q)^2 + 1)^{d\chi/2}
       \bigl(\essinf_{\Bz{1/4}}(u+\delta)\bigr)^{p_0}.
\end{multline*}
By \leanname{weak\_harnack\_forward\_factor\_bound} the right factor is
absorbed into $(1-q)^{-d\chi}$ times a dimensional constant, so collecting
terms produces a constant $C_{\mathrm{chainMain}}(d)$ with
\[
  \Bigl(\int_{\Bz{1/4}} (u+\delta)^{qd/(d-2)}\Bigr)^{p_0(d-2)/(qd)}
  \;\leq\; \frac{C_{\mathrm{chainMain}}(d)}{(1-q)^{d\chi(d)}}
        \bigl(\essinf_{\Bz{1/4}}(u+\delta)\bigr)^{p_0}.
\]
Letting $\delta = 1/(n+1) \downarrow 0$ and using the dominated/monotone
convergence theorems on each side concludes the case $p_0 < q$.

For $q \leq p_0$ pick $q' := (p_0+1)/2 \in (p_0,1)$ and apply the previous
case with $q'$. Since $q d/(d-2) \leq q' d/(d-2)$, monotonicity of $L^p$ norms
on $\Bz{1/4}$ (\leanname{quarterBall\_lpNorm\_mono}) raises the resulting
$L^{q' d/(d-2)}$ norm to dominate the $L^{q d/(d-2)}$ norm, and a uniform
bound $C_{\mathrm{chainMain}}/(1-q')^{d\chi} \leq C_{\mathrm{chain}}/(1-q)^{d\chi}$
\leanname{(hmain\_to\_public)} then yields the claim.
\end{proof}

\begin{definition}[\leanname{C\_weakHarnack}]
\label{def:CwH}
$
  C_{wH}(d) \;:=\; C_{\mathrm{chain}}(d)^{1/c'(d)}.
$
By Lemma~\ref{lem:c-crossover-bounds} and $C_{\mathrm{chain}} \geq 1$, we have
$C_{wH}(d) \geq 1$ (\leanname{one\_le\_C\_weakHarnack}).
\end{definition}

\begin{lstlisting}[style=Matlab-editor,basicstyle=\mlttfamily,escapechar=",mlshowsectionrules=true,mathescape=true,morecomment={[s]{\%\{}{\%\}}},language=lean,numbers=none]
/-- The constant `C(d)` for the weak Harnack inequality.
The chain produces `(C_chain)^{1/p₀}` where the forward stage uses the honest
finite-measure-upgrade constant `C_weakHarnack0Forward`. Since
`1/p₀ = Λ^{1/2}/c'`, the stated bound has `(C_wH/(1-q)^{dχ/c'})^{Λ^{1/2}}`,
so `C_wH` must absorb `C_chain^{1/c'}`. -/
noncomputable def C_weakHarnack (d : ℕ) [NeZero d] (hd : 2 < (d : ℝ)) : ℝ :=
  (C_chain (d := d) hd) ^ (1 / c_crossover' d)

theorem one_le_C_weakHarnack (hd : 2 < (d : ℝ)) : 1 ≤ C_weakHarnack (d := d) hd := by
  unfold C_weakHarnack
  have hexp_nonneg : 0 ≤ 1 / c_crossover' d := by
    exact div_nonneg (by norm_num : (0 : ℝ) ≤ 1) (c_crossover'_pos (d := d)).le
  exact Real.one_le_rpow (one_le_C_chain (d := d) hd) hexp_nonneg

\end{lstlisting}
\begin{theorem}[\leanname{weak\_harnack}]
\label{thm:wH}
For $d \geq 3$, $A$ normalized elliptic on $\Bz 1$, $u > 0$ a supersolution,
and $0 < q < 1$:
\begin{equation*}
  \Bigl(\int_{\Bz{1/4}} u^{q d/(d-2)}\,dx\Bigr)^{(d-2)/(qd)}
    \leq \Bigl(\frac{C_{wH}(d)}{(1-q)^{\beta(d)}}\Bigr)^{\Lambda_A^{1/2}}
        \essinf_{\Bz{1/4}} u.
\end{equation*}
\end{theorem}

\begin{lstlisting}[style=Matlab-editor,basicstyle=\mlttfamily,escapechar=",mlshowsectionrules=true,mathescape=true,morecomment={[s]{\%\{}{\%\}}},language=lean,numbers=none]
/-! ### Main theorem -/

/-- **Weak Harnack inequality** for positive supersolutions on `B₁`.

For `u > 0` with `-∇·(A∇u) ≥ 0` on `B₁`, and `0 < q < 1`:
```
‖u‖_{L^{q*}(B_{1/4})} ≤ (C(d)/(1-q)^{d/c'})^{Λ^{1/2}} · essInf_{B_{1/4}} u
```

The estimate is stated on `B_{1/4}`, the ball naturally produced by the
forward and inverse Moser steps together with the crossover estimate. -/
theorem weak_harnack
    (hd : 2 < (d : ℝ))
    (A : NormalizedEllipticCoeff d (ball (0 : E) 1))
    {u : E → ℝ} {q : ℝ} (hq : 0 < q) (hq1 : q < 1)
    (hu_pos : ∀ x ∈ ball (0 : E) 1, 0 < u x)
    (hsuper : IsSupersolution A.1 u) :
    (∫ x in ball (0 : E) (1 / 4 : ℝ),
        |u x| ^ (q * (d : ℝ) / ((d : ℝ) - 2)) ∂volume) ^
          (((d : ℝ) - 2) / (q * (d : ℝ))) ≤
      (C_weakHarnack d hd / (1 - q) ^ (weak_harnack_decay_exp d)) ^
        (A.1.Λ ^ ((1 : ℝ) / 2)) *
        essInf u (volume.restrict (ball (0 : E) (1 / 4 : ℝ))) := by
  set p₀ := weakHarnackP0 A with hp₀_def
  have hp₀_pos : 0 < p₀ := p₀_pos A
  have hinv_p₀ : 1 / p₀ = A.1.Λ ^ ((1 : ℝ) / 2) / c_crossover' d := by
    rw [hp₀_def, weakHarnackP0]; field_simp
  -- The chain at the p₀ level, using C_chain (un-powered constant).
  have hchain := weak_harnack_chain hd A hq hq1 hu_pos hsuper
  -- Abbreviations.
  set I := ∫ x in ball (0 : E) (1 / 4 : ℝ),
      |u x| ^ (q * (d : ℝ) / ((d : ℝ) - 2)) ∂volume
  set α := p₀ * (((d : ℝ) - 2) / (q * (d : ℝ)))
  set inf_u := essInf u (volume.restrict (ball (0 : E) (1 / 4 : ℝ)))
  -- hchain: I^α ≤ (C_chain/(1-q)^d) · inf_u^{p₀}
  -- Raise to 1/p₀: I^{α/p₀} ≤ (C_chain/(1-q)^d)^{1/p₀} · inf_u
  have hI_nonneg : 0 ≤ I := integral_nonneg fun x => by positivity
  have hC_nonneg : 0 ≤ C_chain (d := d) hd / (1 - q) ^ ((d : ℝ) * moserChi d) := by
    exact div_nonneg (le_trans (by norm_num) (one_le_C_chain (d := d) hd))
      (Real.rpow_nonneg (by linarith) _)
  have hinf_nonneg : 0 ≤ inf_u := by
    have hμ_ne_zero :
        volume.restrict (ball (0 : E) (1 / 4 : ℝ)) ≠ 0 := by
      intro hzero
      have hball_zero : volume (ball (0 : E) (1 / 4 : ℝ)) = 0 := by
        simpa [Measure.restrict_apply_univ] using
          congrArg (fun μ : Measure E => μ Set.univ) hzero
      exact (Metric.measure_ball_pos volume (0 : E) (by norm_num : 0 < (1 / 4 : ℝ))).ne'
        hball_zero
    have hquarter_nonneg :
        ∀ᵐ x ∂(volume.restrict (ball (0 : E) (1 / 4 : ℝ))), (0 : ℝ) ≤ u x := by
      filter_upwards [ae_restrict_mem measurableSet_ball] with x hx
      exact (hu_pos x ((Metric.ball_subset_ball (by norm_num : (1 / 4 : ℝ) ≤ 1)) hx)).le
    exact le_essInf_real_of_ae_le (d := d) hμ_ne_zero hquarter_nonneg
  have hα_nonneg : 0 ≤ α := mul_nonneg hp₀_pos.le (div_nonneg (by linarith) (by positivity))
  have hstep := rpow_le_rpow_of_exponent hI_nonneg hC_nonneg hinf_nonneg hp₀_pos hα_nonneg hchain
  -- hstep: I^{α/p₀} ≤ (C_chain/(1-q)^d)^{1/p₀} · inf_u
  -- α/p₀ = (d-2)/(qd).
  have hα_div : α / p₀ = ((d : ℝ) - 2) / (q * (d : ℝ)) := by
    simp only [α]; field_simp
  rw [hα_div] at hstep
  -- Key identity: (C_chain/(1-q)^d)^{1/p₀} = (C_chain^{1/c'}/(1-q)^{d/c'})^{Λ^{1/2}}
  -- Since C_weakHarnack = C_chain^{1/c'}, this is (C_wH/(1-q)^{d/c'})^{Λ^{1/2}}.
  -- So hstep gives exactly the stated bound.
  have hconst_eq :
      (C_chain (d := d) hd / (1 - q) ^ ((d : ℝ) * moserChi d)) ^ (1 / p₀) =
        (C_weakHarnack d hd / (1 - q) ^ (weak_harnack_decay_exp d)) ^
          (A.1.Λ ^ ((1 : ℝ) / 2)) := by
    -- 1/p₀ = Λ^{1/2}/c'. So LHS = (C_chain/(1-q)^d)^{Λ^{1/2}/c'}.
    -- Distribute: = C_chain^{Λ^{1/2}/c'} · (1-q)^{-d·Λ^{1/2}/c'}.
    -- C_chain^{Λ^{1/2}/c'} = (C_chain^{1/c'})^{Λ^{1/2}} = C_wH^{Λ^{1/2}}.
    -- (1-q)^{-d·Λ^{1/2}/c'} = ((1-q)^{d/c'})^{-Λ^{1/2}}.
    -- Combined: (C_wH · (1-q)^{-d/c'})^{Λ^{1/2}} = (C_wH/(1-q)^{d/c'})^{Λ^{1/2}}.
    rw [hinv_p₀]
    -- Goal: (C_chain d / (1 - q) ^ (d : ℝ)) ^ (A.1.Λ ^ ((1:ℝ)/2) / c_crossover' d) =
    --       (C_weakHarnack d / (1 - q) ^ weak_harnack_decay_exp d) ^ (A.1.Λ ^ ((1:ℝ)/2))
    have hc'_pos := c_crossover'_pos (d := d)
    have hc'_ne : c_crossover' (d := d) ≠ 0 := ne_of_gt hc'_pos
    have h1q_pos : 0 < 1 - q := by linarith
    have h1q_nonneg : 0 ≤ 1 - q := h1q_pos.le
    have hbase_nonneg : 0 ≤ C_chain (d := d) hd / (1 - q) ^ ((d : ℝ) * moserChi d) :=
      div_nonneg (le_trans (by norm_num) (one_le_C_chain (d := d) hd)) (by positivity)
    have hexp_rw : A.1.Λ ^ ((1:ℝ)/2) / c_crossover' d =
        (1 / c_crossover' d) * A.1.Λ ^ ((1:ℝ)/2) := by
      rw [div_eq_inv_mul]; congr 1; exact (one_div _).symm
    rw [hexp_rw]
    rw [Real.rpow_mul hbase_nonneg]
    -- Goal: ((C_chain d / (1 - q) ^ (d : ℝ)) ^ (1 / c_crossover' d)) ^ (Λ^{1/2}) =
    --       (C_weakHarnack d / (1 - q) ^ weak_harnack_decay_exp d) ^ (Λ^{1/2})
    congr 1
    -- Goal: (C_chain d / (1 - q) ^ (d : ℝ)) ^ (1 / c_crossover' d) =
    --       C_weakHarnack d / (1 - q) ^ weak_harnack_decay_exp d
    rw [Real.div_rpow (le_trans (by norm_num) (one_le_C_chain (d := d) hd))
        (Real.rpow_nonneg h1q_nonneg _)]
    -- Goal: C_chain d ^ (1/c') / ((1-q)^d)^{1/c'} = C_wH / (1-q)^{d/c'}
    rw [← C_weakHarnack_eq_C_chain_rpow (d := d) hd]
    have : ((1 - q) ^ ((d : ℝ) * moserChi d)) ^ (1 / c_crossover' d) =
        (1 - q) ^ weak_harnack_decay_exp d := by
      rw [← Real.rpow_mul h1q_nonneg]
      unfold weak_harnack_decay_exp
      congr 1
      rw [mul_one_div]
    rw [this]
  rw [hconst_eq] at hstep
  exact hstep

\end{lstlisting}
\begin{proof}
By Theorem~\ref{thm:wH-chain},
$I^\alpha \leq C_{\mathrm{chain}}(d)/(1-q)^{d\chi}\cdot (\essinf u)^{p_0}$
where $I := \int_{\Bz{1/4}} u^{qd/(d-2)}$ and $\alpha := p_0(d-2)/(qd)$.
Raising both sides to the $1/p_0$ power (\leanname{rpow\_le\_rpow\_of\_exponent}):
\[
  I^{(d-2)/(qd)}
   \;\leq\; \Bigl(\frac{C_{\mathrm{chain}}(d)}{(1-q)^{d\chi(d)}}\Bigr)^{1/p_0}
       \essinf_{\Bz{1/4}} u.
\]
Now use $1/p_0 = \Lambda_A^{1/2}/c'(d)$ and the laws of real powers:
\begin{align*}
  \Bigl(\frac{C_{\mathrm{chain}}(d)}{(1-q)^{d\chi(d)}}\Bigr)^{1/p_0}
   &= \Bigl(\bigl(C_{\mathrm{chain}}(d)/(1-q)^{d\chi(d)}\bigr)^{1/c'(d)}\Bigr)^{\Lambda_A^{1/2}}\\
   &= \Bigl(C_{\mathrm{chain}}(d)^{1/c'(d)}\big/(1-q)^{d\chi(d)/c'(d)}\Bigr)^{\Lambda_A^{1/2}}\\
   &= \Bigl(C_{wH}(d)\big/(1-q)^{\beta(d)}\Bigr)^{\Lambda_A^{1/2}}.
\end{align*}
\end{proof}

\begin{theorem}[\leanname{weak\_harnack\_on\_ball}]
\label{thm:wH-on-ball}
On a general ball $\Ball R{x_0}$ with $R > 0$, the weak Harnack inequality
takes the form
\begin{multline*}
  \Bigl(R^{-d}\int_{\Ball{R/4}{x_0}} u^{qd/(d-2)}\,dx\Bigr)^{(d-2)/(qd)}\\
   \leq \Bigl(\frac{C_{wH}(d)}{(1-q)^{\beta(d)}}\Bigr)^{\Lambda_A^{1/2}}
        \essinf_{z \in \Bz{1/4}}\, u(x_0 + R z).
\end{multline*}
\end{theorem}

\begin{lstlisting}[style=Matlab-editor,basicstyle=\mlttfamily,escapechar=",mlshowsectionrules=true,mathescape=true,morecomment={[s]{\%\{}{\%\}}},language=lean,numbers=none]
/-- Weak Harnack on an arbitrary ball, obtained by rescaling the unit-ball
Chapter 06 statement. -/
theorem weak_harnack_on_ball
    (hd : 2 < (d : ℝ))
    {x₀ : E} {R : ℝ} (hR : 0 < R)
    (A : NormalizedEllipticCoeff d (Metric.ball x₀ R))
    {u : E → ℝ} {q : ℝ} (hq : 0 < q) (hq1 : q < 1)
    (hu_pos : ∀ x ∈ Metric.ball x₀ R, 0 < u x)
    (hsuper : IsSupersolution A.1 u) :
    (((R ^ Module.finrank ℝ E)⁻¹ *
        ∫ x in Metric.ball x₀ (R / 4 : ℝ),
          |u x| ^ (q * (d : ℝ) / ((d : ℝ) - 2)) ∂volume)) ^
          (((d : ℝ) - 2) / (q * (d : ℝ))) ≤
      (C_weakHarnack d hd / (1 - q) ^ (weak_harnack_decay_exp d)) ^
        (A.1.Λ ^ ((1 : ℝ) / 2)) *
        essInf (fun z : E => u (x₀ + R • z))
          (volume.restrict (Metric.ball (0 : E) (1 / 4 : ℝ))) := by
  let uR : E → ℝ := rescaleToUnitBall (d := d) (x₀ := x₀) (R := R) u
  let AR : NormalizedEllipticCoeff d (Metric.ball (0 : E) 1) :=
    rescaleNormalizedCoeffToUnitBall (d := d) (x₀ := x₀) (R := R) hR A
  have hposR : ∀ z ∈ Metric.ball (0 : E) 1, 0 < uR z := by
    intro z hz
    have hz_ball : x₀ + R • z ∈ Metric.ball x₀ R := by
      have hz' : z ∈ (fun y : E => x₀ + R • y) ⁻¹' Metric.ball x₀ (R * (1 : ℝ)) := by
        rw [affine_preimage_ball_mul (d := d) (x₀ := x₀) (R := R) (ρ := (1 : ℝ)) hR]
        exact hz
      simpa using hz'
    exact hu_pos (x₀ + R • z) hz_ball
  have hsuperR : IsSupersolution AR.1 uR := by
    change IsSupersolution
      (rescaleCoeffToUnitBall (d := d) (x₀ := x₀) (R := R) hR A.1)
      (rescaleToUnitBall (d := d) (x₀ := x₀) (R := R) u)
    exact rescaleToUnitBall_isSupersolution (d := d) (x₀ := x₀) (R := R) hR A.1 hsuper
  have hInt_eq :
      ∫ z in Metric.ball (0 : E) (1 / 4 : ℝ),
          |uR z| ^ (q * (d : ℝ) / ((d : ℝ) - 2)) ∂volume =
        (R ^ Module.finrank ℝ E)⁻¹ *
          ∫ x in Metric.ball x₀ (R / 4 : ℝ),
            |u x| ^ (q * (d : ℝ) / ((d : ℝ) - 2)) ∂volume := by
    dsimp [uR]
    simpa [show R * (1 / 4 : ℝ) = R / 4 by ring] using
      integral_comp_affine_ball
        (d := d) (x₀ := x₀) (R := R) (ρ := (1 / 4 : ℝ)) hR
        (fun x => |u x| ^ (q * (d : ℝ) / ((d : ℝ) - 2)))
  have hbase := weak_harnack (d := d) hd AR hq hq1 hposR hsuperR
  rw [hInt_eq] at hbase
  simpa [AR, uR, rescaleNormalizedCoeffToUnitBall, rescaleCoeffToUnitBall,
    rescaleToUnitBall] using hbase

\end{lstlisting}
\begin{proof}
Apply Theorem~\ref{thm:wH} to the rescaled function
$u_R(z) := u(x_0 + R z)$ on $\Bz 1$ with the rescaled coefficients
$A_R := \leanname{rescaleNormalizedCoeffToUnitBall}\,A$, which preserves
$\Lambda$ (so $p_0(A_R) = p_0(A)$). The change of variables
$\leanname{crossover\_integral\_comp\_affine\_ball}$ yields the
$R^{-d}$ Jacobian factor on the left, while the essential infimum of $u_R$
on $\Bz{1/4}$ equals that of $u$ on $\Ball{R/4}{x_0}$ pulled back through the
affine map.
\end{proof}

\subsubsection{Support constants}

\begin{definition}[\leanname{localMoserBase}, \leanname{localWeakHarnackBase}]
\label{def:local-bases}
\begin{align*}
  M_{\mathrm{loc}}(d) &:= C_{\mathrm{Moser}}(d)\,((d-1))^d\,4^d, \\
  W_{\mathrm{loc}}(d) &:= C_{wH,0}(d)\, d^{\beta(d)}.
\end{align*}
\end{definition}

\begin{lstlisting}[style=Matlab-editor,basicstyle=\mlttfamily,escapechar=",mlshowsectionrules=true,mathescape=true,morecomment={[s]{\%\{}{\%\}}},language=lean,numbers=none]
/-- The local Moser base constant appearing in the Chapter 07 endpoint bounds. -/
noncomputable def localMoserBase (d : ℕ) [NeZero d] : ℝ :=
  C_Moser d * (((d : ℝ) - 1) ^ (d : ℝ)) * ((4 : ℝ) ^ (d : ℝ))

/-- The local weak Harnack base constant appearing in the Chapter 07 endpoint
bounds. -/
noncomputable def localWeakHarnackBase (d : ℕ) [NeZero d] : ℝ :=
  C_weakHarnack_of d * ((d : ℝ) ^ (weak_harnack_decay_exp d))

\end{lstlisting}
\begin{lemma}[\leanname{one\_le\_localMoserBase}, \leanname{one\_le\_localWeakHarnackBase}]
For $d \geq 3$, $1 \leq M_{\mathrm{loc}}(d)$ and $1 \leq W_{\mathrm{loc}}(d)$.
Both are strictly positive.
\end{lemma}

\begin{lstlisting}[style=Matlab-editor,basicstyle=\mlttfamily,escapechar=",mlshowsectionrules=true,mathescape=true,morecomment={[s]{\%\{}{\%\}}},language=lean,numbers=none]
theorem one_le_localMoserBase (d : ℕ) [NeZero d] (hd : 2 < (d : ℝ)) :
    1 ≤ localMoserBase d := by
  unfold localMoserBase
  have hcm : 1 ≤ C_Moser d := one_le_C_Moser (d := d)
  have hpow1 : 1 ≤ ((d : ℝ) - 1) ^ (d : ℝ) := by
    have hbase : 1 ≤ (d : ℝ) - 1 := by linarith
    exact Real.one_le_rpow hbase (by positivity)
  have hpow2 : 1 ≤ (4 : ℝ) ^ (d : ℝ) := by
    exact Real.one_le_rpow (by norm_num : (1 : ℝ) ≤ 4) (by positivity)
  exact one_le_mul_of_one_le_of_one_le (one_le_mul_of_one_le_of_one_le hcm hpow1) hpow2

theorem one_le_localWeakHarnackBase (d : ℕ) [NeZero d] (hd : 2 < (d : ℝ)) :
    1 ≤ localWeakHarnackBase d := by
  unfold localWeakHarnackBase
  have hcm : 1 ≤ C_weakHarnack_of d := by
    exact one_le_C_weakHarnack_of (d := d)
  have hpow : 1 ≤ (d : ℝ) ^ (weak_harnack_decay_exp d) := by
    have hd_one : 1 ≤ (d : ℝ) := by
      exact_mod_cast Nat.succ_le_of_lt (Nat.pos_of_ne_zero (NeZero.ne d))
    have hdecay_nonneg : 0 ≤ weak_harnack_decay_exp d :=
      weak_harnack_decay_exp_nonneg (d := d) hd
    exact Real.one_le_rpow hd_one hdecay_nonneg
  exact one_le_mul_of_one_le_of_one_le hcm hpow

\end{lstlisting}
\begin{proof}
Each is a product of factors $\geq 1$: $C_{\mathrm{Moser}}(d) \geq 1$,
$(d-1)^d \geq 1$ (since $d-1 \geq 2$), $4^d \geq 1$, $C_{wH,0} \geq 1$,
$d^{\beta(d)} \geq 1$ (since $d \geq 1$ and $\beta(d) \geq 0$ by
\leanname{weak\_harnack\_decay\_exp\_nonneg}).
\end{proof}

\subsubsection{Measure-bound utilities}

The Harnack and Holder layers also use a small toolkit of $\esssup/\essinf$
arithmetic on restricted ball measures (\leanname{Support/MeasureBounds.lean}).

\begin{lemma}[\leanname{le\_essInf\_real\_of\_ae\_le}, \leanname{essSup\_le\_of\_ae\_bdd}, \leanname{essInf\_add\_const\_of\_bdd}, \leanname{essSup\_neg\_of\_bdd}]
\label{lem:meas-bd}
Let $\mu$ be a non-zero (Borel) measure on $E$ and $u : E \to \R$.
\begin{enumerate}
\item If $u \geq c$ $\mu$-a.e., then $c \leq \essinf_\mu u$.
\item If $K \leq u \leq C$ $\mu$-a.e., then $\essinf_\mu u, \esssup_\mu u
  \in [K,C]$.
\item If $u$ is bounded $\mu$-a.e., then
  $\essinf_\mu(u+c) = \essinf_\mu u + c$ and the symmetric statement for
  $\esssup$.
\item $\esssup_\mu(-u) = -\essinf_\mu u$ when $u$ is bounded.
\item For a measurable embedding $T : E\to E$ and measurable $g$,
  $\essinf_\mu(g\circ T) = \essinf_{T_*\mu} g$ \leanname{(essInf\_comp\_measurableEmbedding\_eq)}.
\item For $c \in \ENNReal\setminus\{0\}$, $\essinf_{c\cdot \mu}f =
  \essinf_\mu f$.
\end{enumerate}
Furthermore, $\vol(\Ball r c)\neq 0$ for $r > 0$, so the restriction
$\vol|_{\Ball r c}$ is non-zero \leanname{(restrict\_ball\_ne\_zero)}.
\end{lemma}

\begin{lstlisting}[style=Matlab-editor,basicstyle=\mlttfamily,escapechar=",mlshowsectionrules=true,mathescape=true,morecomment={[s]{\%\{}{\%\}}},language=lean,numbers=none]
omit [NeZero d] in
theorem le_essInf_real_of_ae_le
    {μ : Measure E} (hμ : μ ≠ 0) {u : E → ℝ} {c : ℝ}
    (hlow : ∀ᵐ x ∂μ, c ≤ u x) :
    c ≤ essInf u μ := by
  rw [essInf_eq_sSup]
  refine le_csSup ?_ ?_
  · obtain ⟨N, hN_nonzero⟩ : ∃ N : ℕ, μ {x | u x ≤ N} ≠ 0 := by
      by_contra hN
      have hnull : ∀ N : ℕ, μ {x | u x ≤ N} = 0 := by
        intro N
        exact Classical.byContradiction fun hNZ => hN ⟨N, hNZ⟩
      have huniv : (⋃ N : ℕ, {x | u x ≤ N}) = Set.univ := by
        ext x
        constructor
        · intro _
          simp
        · intro _
          obtain ⟨N, hNx⟩ := exists_nat_ge (u x)
          exact Set.mem_iUnion.2 ⟨N, hNx⟩
      have hzero : μ Set.univ = 0 := by
        calc
          μ Set.univ = μ (⋃ N : ℕ, {x | u x ≤ N}) := by rw [huniv]
          _ = 0 := measure_iUnion_null hnull
      exact hμ (Measure.measure_univ_eq_zero.mp hzero)
    refine ⟨(N : ℝ), ?_⟩
    intro a ha
    rw [Set.mem_setOf_eq] at ha
    by_contra hle
    push_neg at hle
    have hsubset : {x | u x ≤ (N : ℝ)} ⊆ {x | u x < a} := by
      intro x hx
      exact lt_of_le_of_lt hx hle
    exact hN_nonzero (measure_mono_null hsubset ha)
  · rw [Set.mem_setOf_eq]
    rw [ae_iff] at hlow
    simpa [not_le] using hlow

omit [NeZero d] in
theorem essSup_le_of_ae_bdd
    {μ : Measure E} (hμ : μ ≠ 0)
    {u : E → ℝ} {K C : ℝ}
    (hlow : ∀ᵐ x ∂ μ, K ≤ u x)
    (hupp : ∀ᵐ x ∂ μ, u x ≤ C) :
    essSup u μ ≤ C := by
  rw [essSup_eq_sInf]
  refine csInf_le ?_ ?_
  · refine ⟨K, ?_⟩
    intro b hb
    by_contra hKb
    have hb_lt : b < K := lt_of_not_ge hKb
    have hlt : ∀ᵐ x ∂ μ, b < u x := by
      filter_upwards [hlow] with x hx
      linarith
    have hb_ae : ∀ᵐ x ∂ μ, ¬ b < u x := by
      change μ {x | b < u x} = 0 at hb
      rw [ae_iff]
      simpa using hb
    have hfalse : ∀ᵐ x ∂ μ, False := by
      filter_upwards [hlt, hb_ae] with x hx hx'
      exact hx' hx
    rw [ae_iff] at hfalse
    have hzero : μ Set.univ = 0 := by
      simpa using hfalse
    exact hμ (Measure.measure_univ_eq_zero.mp hzero)
  · change μ {x | C < u x} = 0
    simpa [ae_iff, not_lt] using hupp

omit [NeZero d] in
theorem essInf_add_const_of_bdd
    {μ : Measure E} (hμ : μ ≠ 0) {u : E → ℝ} (c : ℝ)
    (hu_bdd_below : Filter.IsBoundedUnder (· ≥ ·) (ae μ) u)
    (hu_bdd_above : Filter.IsBoundedUnder (· ≤ ·) (ae μ) u) :
    essInf (fun x => u x + c) μ = essInf u μ + c := by
  haveI : NeZero μ := ⟨hμ⟩
  haveI : (ae μ).NeBot := Measure.ae.neBot
  dsimp [essInf]
  simpa using
    (liminf_add_const (ae μ) u c hu_bdd_above.isCoboundedUnder_ge hu_bdd_below)

omit [NeZero d] in
theorem essSup_neg_of_bdd
    {μ : Measure E} (hμ : μ ≠ 0) {u : E → ℝ}
    (hu_bdd_below : Filter.IsBoundedUnder (· ≥ ·) (ae μ) u)
    (hu_bdd_above : Filter.IsBoundedUnder (· ≤ ·) (ae μ) u)
    (hneg_bdd_below : Filter.IsBoundedUnder (· ≥ ·) (ae μ) (fun x => -u x))
    (hneg_bdd_above : Filter.IsBoundedUnder (· ≤ ·) (ae μ) (fun x => -u x)) :
    essSup (fun x => -u x) μ = - essInf u μ := by
  haveI : NeZero μ := ⟨hμ⟩
  haveI : (ae μ).NeBot := Measure.ae.neBot
  dsimp [essSup, essInf]
  simpa using
    (OrderIso.liminf_apply (β := ℝ) (γ := ℝᵒᵈ) (f := ae μ) (u := u) (g := OrderIso.neg ℝ)
      hu_bdd_below hu_bdd_above.isCoboundedUnder_ge
      hneg_bdd_above hneg_bdd_below.isCoboundedUnder_le).symm

omit [NeZero d] in
theorem essInf_comp_measurableEmbedding_eq
    {μ : Measure E} {f : E → E} (hf : MeasurableEmbedding f)
    {g : E → ℝ} (_hg : Measurable g) :
    essInf (g ∘ f) μ = essInf g (Measure.map f μ) := by
  rw [essInf_eq_sSup, essInf_eq_sSup]
  congr 1
  ext a
  change (μ {x | (g ∘ f) x < a} = 0) ↔ ((Measure.map f μ) {x | g x < a} = 0)
  rw [hf.map_apply μ {x | g x < a}]
  simp [Function.comp]

omit [NeZero d] in
theorem restrict_ball_ne_zero {c : E} {r : ℝ} (hr : 0 < r) :
    volume.restrict (Metric.ball c r) ≠ 0 := by
  intro hzero
  have hball_zero : volume (Metric.ball c r) = 0 := by
    simpa [Measure.restrict_apply_univ] using
      congrArg (fun μ : Measure E => μ Set.univ) hzero
  exact (Metric.measure_ball_pos volume c hr).ne' hball_zero

\end{lstlisting}
\begin{proof}
Each statement is a direct unfolding of $\essinf$/$\esssup$ as a $\sSup$/$\sInf$
over levels at which the corresponding sublevel set has positive/full measure.
For instance, in (1): assume $\mu \neq 0$ and $\mu(\{u < c\}) = 0$. Then
$\{u < a\}$ has measure zero for every $a < c$, so $a$ is dominated by every
upper bound used in the definition of $\essinf_\mu u$, hence $c \leq
\essinf_\mu u$. Statements (3),(4) follow from the corresponding
$\liminf/\limsup$ arithmetic on the filter $\ae \mu$ together with the fact
that the functions are eventually bounded. Statements (5),(6) follow from
the change-of-variables and scaling formulas for the underlying
$\sSup/\sInf$ over level sets.
\end{proof}

%% --- Section 11 ---
\subsection{Harnack Inequality and Holder Regularity}

\begin{definition}[\leanname{C\_harnack}]
\label{def:Charnack}
\[
  C_{\mathrm{Har}}(d) := 17\Bigl(\frac{d-2}{d-1}\,\log\bigl(C_{\mathrm{Moser}}(d)\,(d-1)^d\,4^d\bigr)
        + d + \log\bigl(C_{wH,0}(d)\,d^{\beta(d)}\bigr)\Bigr).
\]
\end{definition}

\begin{lstlisting}[style=Matlab-editor,basicstyle=\mlttfamily,escapechar=",mlshowsectionrules=true,mathescape=true,morecomment={[s]{\%\{}{\%\}}},language=lean,numbers=none]
/-- Quantitative constant used to absorb the local-ball Harnack chain into the
final exponential form. -/
noncomputable def C_harnack (d : ℕ) [NeZero d] : ℝ :=
  17 *
    (((d : ℝ) - 2) / ((d : ℝ) - 1) *
        Real.log
          (C_Moser d * (((d : ℝ) - 1) ^ (d : ℝ)) * ((4 : ℝ) ^ (d : ℝ))) +
      (d : ℝ) +
      Real.log (C_weakHarnack_of d * ((d : ℝ) ^ (weak_harnack_decay_exp d))))

\end{lstlisting}
\begin{lemma}[\leanname{C\_harnack\_nonneg}]
For $d \geq 3$, $0 \leq C_{\mathrm{Har}}(d)$.
\end{lemma}

\begin{lstlisting}[style=Matlab-editor,basicstyle=\mlttfamily,escapechar=",mlshowsectionrules=true,mathescape=true,morecomment={[s]{\%\{}{\%\}}},language=lean,numbers=none]
theorem C_harnack_nonneg
    (hd : 2 < (d : ℝ)) :
    0 ≤ C_harnack d := by
  unfold C_harnack
  have hcoeff_nonneg : 0 ≤ ((d : ℝ) - 2) / ((d : ℝ) - 1) := by
    have hnum_nonneg : 0 ≤ (d : ℝ) - 2 := by
      linarith
    have hden_nonneg : 0 ≤ (d : ℝ) - 1 := by
      linarith
    exact div_nonneg hnum_nonneg hden_nonneg
  have hlogM_nonneg :
      0 ≤ Real.log (C_Moser d * (((d : ℝ) - 1) ^ (d : ℝ)) * ((4 : ℝ) ^ (d : ℝ))) := by
    simpa [localMoserBase] using Real.log_nonneg (one_le_localMoserBase (d := d) hd)
  have hlogW_nonneg :
      0 ≤ Real.log (C_weakHarnack_of d * ((d : ℝ) ^ (weak_harnack_decay_exp d))) := by
    simpa [localWeakHarnackBase] using Real.log_nonneg (one_le_localWeakHarnackBase (d := d) hd)
  have hd_nonneg : 0 ≤ (d : ℝ) := by positivity
  refine mul_nonneg (by norm_num) ?_
  exact add_nonneg (add_nonneg (mul_nonneg hcoeff_nonneg hlogM_nonneg) hd_nonneg) hlogW_nonneg

\end{lstlisting}
\begin{proof}
$(d-2)/(d-1) \geq 0$, both logs are nonnegative since
$M_{\mathrm{loc}}(d) \geq 1$ and $W_{\mathrm{loc}}(d) \geq 1$, and $d \geq 0$.
\end{proof}

\begin{theorem}[\leanname{harnack}]
\label{thm:harnack}
Let $d \geq 3$, $A$ normalized elliptic on $\Bz 1$, and $u > 0$ a (local
weak) solution of $-\dv(A\nabla u) = 0$ on $\Bz 1$. Then
\begin{equation*}
  \esssup_{\Bz{1/2}} u \;\leq\; \exp\bigl(C_{\mathrm{Har}}(d)\,\Lambda_A^{1/2}\bigr)
        \essinf_{\Bz{1/2}} u.
\end{equation*}
\end{theorem}

\begin{lstlisting}[style=Matlab-editor,basicstyle=\mlttfamily,escapechar=",mlshowsectionrules=true,mathescape=true,morecomment={[s]{\%\{}{\%\}}},language=lean,numbers=none]
/-- Positive weak solutions satisfy Harnack on the unit ball.
-/
theorem harnack
    (hd : 2 < (d : ℝ))
    (A : NormalizedEllipticCoeff d (Metric.ball (0 : E) 1))
    {u : E → ℝ}
    (hu_pos : ∀ x ∈ Metric.ball (0 : E) 1, 0 < u x)
    (hsol : IsSolution A.1 u) :
    essSup u μhalf ≤
      Real.exp (C_harnack d * A.1.Λ ^ ((1 : ℝ) / 2)) *
        essInf u μhalf := by
  let K : ℝ := localHarnackConstant (d := d) A
  have hK_pos : 0 < K := localHarnackConstant_pos (d := d) hd A
  let centerOf : E → E := fun x => (15 / 16 : ℝ) • x
  have hcenter_mem :
      ∀ x ∈ Metric.closedBall (0 : E) (1 / 2 : ℝ),
        centerOf x ∈ Metric.ball (0 : E) (1 / 2 : ℝ) := by
    intro x hx
    rw [Metric.mem_ball, dist_eq_norm, sub_zero]
    have hxnorm : ‖x‖ ≤ 1 / 2 := by simpa [dist_eq_norm] using hx
    dsimp [centerOf]
    rw [norm_smul, Real.norm_of_nonneg (by positivity)]
    calc
      (15 / 16 : ℝ) * ‖x‖ ≤ (15 / 16 : ℝ) * (1 / 2 : ℝ) := by gcongr
      _ = 15 / 32 := by norm_num
      _ < 1 / 2 := by norm_num
  have hmem_cover_ball :
      ∀ x ∈ Metric.closedBall (0 : E) (1 / 2 : ℝ),
        x ∈ Metric.ball (centerOf x) (1 / 16 : ℝ) := by
    intro x hx
    rw [Metric.mem_ball, dist_eq_norm]
    have hxnorm : ‖x‖ ≤ 1 / 2 := by simpa [dist_eq_norm] using hx
    have hrepr : x - centerOf x = (1 / 16 : ℝ) • x := by
      dsimp [centerOf]
      calc
        x - (15 / 16 : ℝ) • x = (1 : ℝ) • x - (15 / 16 : ℝ) • x := by simp
        _ = ((1 : ℝ) - 15 / 16) • x := by rw [sub_smul]
        _ = (1 / 16 : ℝ) • x := by norm_num
    rw [hrepr, norm_smul, Real.norm_of_nonneg (by positivity)]
    calc
      (1 / 16 : ℝ) * ‖x‖ ≤ (1 / 16 : ℝ) * (1 / 2 : ℝ) := by gcongr
      _ = 1 / 32 := by norm_num
      _ < 1 / 16 := by norm_num
  obtain ⟨t, ht_mem, hcover⟩ :=
    (isCompact_closedBall (0 : E) (1 / 2 : ℝ)).elim_nhds_subcover
      (fun x : E => Metric.ball (centerOf x) (1 / 16 : ℝ))
      (fun x hx => IsOpen.mem_nhds Metric.isOpen_ball (hmem_cover_ball x hx))
  have hzero_mem : (0 : E) ∈ Metric.ball (0 : E) (1 / 2 : ℝ) := by
    exact Metric.mem_ball_self (by norm_num)
  have hsubset_cover :
      Metric.ball (0 : E) (1 / 2 : ℝ) ⊆ ⋃ x ∈ t, Metric.ball (centerOf x) (1 / 16 : ℝ) := by
    exact Metric.ball_subset_closedBall.trans hcover
  have hnonneg_half_ae : ∀ᵐ x ∂ μhalf, 0 ≤ u x := by
    filter_upwards [ae_restrict_mem measurableSet_ball] with x hx
    exact (hu_pos x (Metric.ball_subset_ball (by norm_num : (1 / 2 : ℝ) ≤ 1) hx)).le
  have hnonneg_on_eighthBall :
      ∀ {c z : E},
        c ∈ Metric.ball (0 : E) (1 / 2 : ℝ) →
        z ∈ Metric.ball c (1 / 8 : ℝ) →
        0 ≤ u z := by
    intro c z hc hz
    have hc_closed : c ∈ Metric.closedBall (0 : E) (1 / 2 : ℝ) := by
      exact Metric.mem_closedBall.2 (le_of_lt (Metric.mem_ball.1 hc))
    have hz_unit : z ∈ Metric.ball (0 : E) 1 :=
      ball_subset_unitBall_of_mem_closedBall_half (d := d) hc_closed (by positivity) (by norm_num) hz
    exact (hu_pos z hz_unit).le
  have hupper0 :
      ∀ᵐ x ∂ μhalf, u x ≤ K ^ 17 * quarterBallInf u 0 := by
    refine ae_on_set_of_ae_on_finite_cover (μ := volume)
      (s := Metric.ball (0 : E) (1 / 2 : ℝ))
      (U := fun x : E => Metric.ball (centerOf x) (1 / 16 : ℝ))
      hsubset_cover ?_
    intro x hx
    let c : E := centerOf x
    have hc_ball : c ∈ Metric.ball (0 : E) (1 / 2 : ℝ) := hcenter_mem x (ht_mem x hx)
    have hchain :
        quarterBallInf u c ≤ K ^ 16 * quarterBallInf u 0 := by
      simpa [K] using
        quarterBallInf_chain_le (d := d) hd A hu_pos hsol hc_ball hzero_mem
    filter_upwards [ae_le_localHarnack_on_eighthBall (d := d) hd A hu_pos hsol hc_ball] with x hx
    calc
      u x ≤ K * quarterBallInf u c := by simpa [K] using hx
      _ ≤ K * (K ^ 16 * quarterBallInf u 0) := by
            gcongr
      _ = K ^ 17 * quarterBallInf u 0 := by
            rw [pow_succ']
            ring
  have hcore :
      essSup u μhalf ≤ K ^ 17 * essInf u μhalf := by
    have hcenter_bound :
        ∀ {y : E}, y ∈ Metric.ball (0 : E) (1 / 2 : ℝ) →
          essSup u μhalf ≤ K ^ 17 * quarterBallInf u y := by
      intro y hy
      have hhalf :
          ∀ᵐ x ∂ μhalf, u x ≤ K ^ 17 * quarterBallInf u y := by
        refine ae_on_set_of_ae_on_finite_cover (μ := volume)
          (s := Metric.ball (0 : E) (1 / 2 : ℝ))
          (U := fun x : E => Metric.ball (centerOf x) (1 / 16 : ℝ))
          hsubset_cover ?_
        intro x hx
        let c : E := centerOf x
        have hc_ball : c ∈ Metric.ball (0 : E) (1 / 2 : ℝ) := hcenter_mem x (ht_mem x hx)
        have hchain :
            quarterBallInf u c ≤ K ^ 16 * quarterBallInf u y := by
          simpa [K] using
            quarterBallInf_chain_le (d := d) hd A hu_pos hsol hc_ball hy
        filter_upwards [ae_le_localHarnack_on_eighthBall (d := d) hd A hu_pos hsol hc_ball] with x hx
        calc
          u x ≤ K * quarterBallInf u c := by simpa [K] using hx
          _ ≤ K * (K ^ 16 * quarterBallInf u y) := by
                gcongr
          _ = K ^ 17 * quarterBallInf u y := by
                rw [pow_succ']
                ring
      exact essSup_halfBall_le_of_ae_bound hnonneg_half_ae hhalf
    have hlow_ae :
        ∀ᵐ y ∂ μhalf, essSup u μhalf / K ^ 17 ≤ u y := by
      refine ae_on_set_of_ae_on_finite_cover (μ := volume)
        (s := Metric.ball (0 : E) (1 / 2 : ℝ))
        (U := fun x : E => Metric.ball (centerOf x) (1 / 16 : ℝ))
        hsubset_cover ?_
      intro x hx
      let c : E := centerOf x
      have hc_ball : c ∈ Metric.ball (0 : E) (1 / 2 : ℝ) := hcenter_mem x (ht_mem x hx)
      have hcenter_div :
          essSup u μhalf / K ^ 17 ≤ quarterBallInf u c := by
        exact (div_le_iff₀ (pow_pos hK_pos 17)).2
          (by simpa [mul_comm, mul_left_comm, mul_assoc] using hcenter_bound hc_ball)
      have hu_ae_c :
          AEMeasurable u (volume.restrict (Metric.ball c (1 / 8 : ℝ))) := by
        have hsub_c : Metric.ball c (1 / 8 : ℝ) ⊆ Metric.ball (0 : E) 1 := by
          intro z hz
          have hc_closed : c ∈ Metric.closedBall (0 : E) (1 / 2 : ℝ) := by
            exact Metric.mem_closedBall.2 (le_of_lt (Metric.mem_ball.1 hc_ball))
          exact ball_subset_unitBall_of_mem_closedBall_half
            (d := d) hc_closed (by positivity) (by norm_num) hz
        exact aemeasurable_on_ball_of_isSolution (d := d) A hsol hsub_c
      have hquarter_small :
          ∀ᵐ z ∂ volume.restrict (Metric.ball c (1 / 16 : ℝ)),
            quarterBallInf u c ≤ u z := by
        refine ae_restrict_of_ae_restrict_of_subset
          (Metric.ball_subset_ball (by norm_num : (1 / 16 : ℝ) ≤ 1 / 8)) ?_
        exact quarterBallInf_le_ae (c := c) hu_ae_c (fun z hz => hnonneg_on_eighthBall hc_ball hz)
      filter_upwards [hquarter_small] with z hz
      exact le_trans hcenter_div hz
    have hlow :
        essSup u μhalf / K ^ 17 ≤ essInf u μhalf :=
      le_essInf_halfBall_of_ae_bound hlow_ae hupper0
    have hmul : essSup u μhalf ≤ essInf u μhalf * K ^ 17 :=
      (div_le_iff₀ (pow_pos hK_pos 17)).mp hlow
    simpa [mul_comm, mul_left_comm, mul_assoc, K] using hmul
  have hessInf_nonneg : 0 ≤ essInf u μhalf := by
    exact
      le_essInf_real_of_ae_le
        (μ := μhalf)
        (restrict_ball_ne_zero (c := (0 : E)) (r := (1 / 2 : ℝ)) (by norm_num))
        (by simpa using hnonneg_half_ae)
  have hKexp :
      K ^ 17 ≤ Real.exp (C_harnack d * A.1.Λ ^ ((1 : ℝ) / 2)) := by
    simpa [K] using localHarnackConstant_pow_seventeen_le_exp (d := d) hd A
  calc
    essSup u μhalf ≤ K ^ 17 * essInf u μhalf := hcore
    _ ≤ Real.exp (C_harnack d * A.1.Λ ^ ((1 : ℝ) / 2)) * essInf u μhalf := by
          exact mul_le_mul_of_nonneg_right hKexp hessInf_nonneg

\end{lstlisting}
\begin{proof}
Define
\[
K(A) := \bigl(M_{\mathrm{loc}}(d)\Lambda_A^{d/2}\bigr)^{1/p_{\mathrm{Har}}(d)}
        \cdot W_{\mathrm{loc}}(d)^{\Lambda_A^{1/2}},
\]
where $p_{\mathrm{Har}}(d) = (d-1)/(d-2) \cdot$ stuff (the auxiliary
exponent \leanname{harnack\_p}). The forward Moser estimate gives, on each
ball $\Ball{1/16}{c}$ centred at $c \in \Ball{1/2}{0}$,
\[
  \esssup_{\Ball{1/16}{c}} u
   \;\leq\; \bigl(M_{\mathrm{loc}}(d)\Lambda_A^{d/2}\bigr)^{1/p_{\mathrm{Har}}(d)}
        \cdot \Bigl(\fint_{\Ball{1/8}{c}} u^{p_{\mathrm{Har}}(d)}\Bigr)^{1/p_{\mathrm{Har}}(d)},
\]
and Theorem~\ref{thm:wH-on-ball} (weak Harnack on $\Ball{1/8}{c}$) gives
\[
  \Bigl(\fint_{\Ball{1/8}{c}} u^{p_{\mathrm{Har}}(d)}\Bigr)^{1/p_{\mathrm{Har}}(d)}
        \;\leq\; W_{\mathrm{loc}}(d)^{\Lambda_A^{1/2}}\,
        \essinf_{\Ball{1/16}{c}} u.
\]
Concatenating gives the local pointwise (a.e.) bound
$u(x) \leq K(A)\,\essinf_{\Ball{1/8}{c}} u$ for almost every $x \in \Ball{1/16}{c}$
(\leanname{ae\_le\_localHarnack\_on\_eighthBall}).

By compactness of $\overline{\Bz{1/2}}$, there is a finite cover of
$\overline{\Bz{1/2}}$ by balls $\Ball{1/16}{c_i}$ with $c_i = (15/16) x_i$,
$x_i \in \overline{\Bz{1/2}}$. Chaining at most $16$ such local Harnack
inequalities along a path connecting any two centres
(\leanname{quarterBallInf\_chain\_le}) yields, almost everywhere on
$\Bz{1/2}$,
\[
  u(x) \;\leq\; K(A)^{17}\,\essinf_{\Bz{1/2}} u.
\]
Hence $\esssup_{\Bz{1/2}} u \leq K(A)^{17}\,\essinf_{\Bz{1/2}} u$.

Finally, the elementary identity
\[
  17\Bigl[\frac{d-2}{d-1}\log M_{\mathrm{loc}}(d) + \frac{d}{2}\log\Lambda_A
    + \log W_{\mathrm{loc}}(d) \Lambda_A^{1/2}\Bigr]
  \;\leq\; C_{\mathrm{Har}}(d)\,\Lambda_A^{1/2}
\]
(using $\Lambda_A^{1/2} \geq 1$ and $\log\Lambda_A \leq 2\Lambda_A^{1/2}$,
established in \leanname{localHarnackConstant\_pow\_seventeen\_le\_exp})
gives $K(A)^{17} \leq \exp(C_{\mathrm{Har}}(d)\,\Lambda_A^{1/2})$, completing
the proof.
\end{proof}

\begin{corollary}[\leanname{harnack\_of\_homogeneousWeakSolution}]
The same conclusion holds for any positive homogeneous weak solution
$u$ on $\Bz 1$, since such a $u$ is in particular an
$\IsSolution$.
\end{corollary}

\begin{lstlisting}[style=Matlab-editor,basicstyle=\mlttfamily,escapechar=",mlshowsectionrules=true,mathescape=true,morecomment={[s]{\%\{}{\%\}}},language=lean,numbers=none]
/-- Harnack, packaged directly for the equality-form local weak-solution
interface from Chapter 04. -/
theorem harnack_of_homogeneousWeakSolution
    (hd : 2 < (d : ℝ))
    (A : NormalizedEllipticCoeff d (Metric.ball (0 : E) 1))
    {u : E → ℝ}
    (hu_pos : ∀ x ∈ Metric.ball (0 : E) 1, 0 < u x)
    (hweak : IsHomogeneousWeakSolution A.1 u) :
    essSup u μhalf ≤
      Real.exp (C_harnack d * A.1.Λ ^ ((1 : ℝ) / 2)) *
        essInf u μhalf := by
  exact harnack hd A hu_pos (isHomogeneousWeakSolution_isSolution hweak)

\end{lstlisting}
\subsubsection{Holder regularity}

\begin{definition}[\leanname{C\_holder\_Moser}, \leanname{moserDecayAlpha}]
\label{def:CholderMoser}
\begin{align*}
  C_{\mathrm{Hol}}(d) &:= 256\bigl(|C_{\mathrm{Har}}(d)| + M_{\mathrm{loc}}(d)
        + C_{\mathrm{Moser}}(d) 2^d + C_{wH,0}(d) d^d + 8\bigr), \\
  \alpha(A) &:= \exp\bigl(-(C_{\mathrm{Har}}(d)\,\Lambda_A^{1/2})\bigr).
\end{align*}
\end{definition}

\begin{lstlisting}[style=Matlab-editor,basicstyle=\mlttfamily,escapechar=",mlshowsectionrules=true,mathescape=true,morecomment={[s]{\%\{}{\%\}}},language=lean,numbers=none]
/-- Quantitative constant for the Moser Hölder theorem. -/
noncomputable def C_holder_Moser (d : ℕ) [NeZero d] : ℝ :=
  (256 : ℝ) *
    (|C_harnack d| +
      C_Moser d * (((d : ℝ) - 1) ^ (d : ℝ)) * ((4 : ℝ) ^ (d : ℝ)) +
      C_Moser d * ((2 : ℝ) ^ (d : ℝ)) +
      C_weakHarnack_of d * ((d : ℝ) ^ (d : ℝ)) + 8)

noncomputable def moserDecayAlpha
    (A : NormalizedEllipticCoeff d (Metric.ball (0 : E) 1)) : ℝ :=
  Real.exp (-(C_harnack d * A.1.Λ ^ ((1 : ℝ) / 2)))

\end{lstlisting}
\begin{lemma}[\leanname{moserDecayAlpha\_pos}, \leanname{moserDecayAlpha\_le\_one}, \leanname{moserLowerBound\_le\_decayAlpha}]
$0 < \alpha(A) \leq 1$, and
$\exp(-(C_{\mathrm{Hol}}(d)\,\Lambda_A^{1/2})) \leq \alpha(A)$.
\end{lemma}

\begin{lstlisting}[style=Matlab-editor,basicstyle=\mlttfamily,escapechar=",mlshowsectionrules=true,mathescape=true,morecomment={[s]{\%\{}{\%\}}},language=lean,numbers=none]
theorem moserDecayAlpha_pos
    (A : NormalizedEllipticCoeff d (Metric.ball (0 : E) 1)) :
    0 < moserDecayAlpha A := by
  unfold moserDecayAlpha
  positivity

theorem moserDecayAlpha_le_one
    (hd : 2 < (d : ℝ))
    (A : NormalizedEllipticCoeff d (Metric.ball (0 : E) 1)) :
    moserDecayAlpha A ≤ 1 := by
  unfold moserDecayAlpha
  refine Real.exp_le_one_iff.mpr ?_
  have hnonneg :
      0 ≤ C_harnack d * A.1.Λ ^ ((1 : ℝ) / 2) := by
    exact mul_nonneg (C_harnack_nonneg (d := d) hd) (Real.rpow_nonneg A.1.Λ_nonneg _)
  linarith

theorem moserLowerBound_le_decayAlpha
    (_hd : 2 < (d : ℝ))
    (A : NormalizedEllipticCoeff d (Metric.ball (0 : E) 1)) :
    Real.exp (-(C_holder_Moser d * A.1.Λ ^ ((1 : ℝ) / 2))) ≤ moserDecayAlpha A := by
  unfold moserDecayAlpha
  refine Real.exp_le_exp.mpr ?_
  have hs_nonneg : 0 ≤ A.1.Λ ^ ((1 : ℝ) / 2) := by
    exact Real.rpow_nonneg A.1.Λ_nonneg _
  have hmul :
      C_harnack d * A.1.Λ ^ ((1 : ℝ) / 2) ≤
        C_holder_Moser d * A.1.Λ ^ ((1 : ℝ) / 2) := by
    exact mul_le_mul_of_nonneg_right (C_harnack_le_C_holder_Moser (d := d)) hs_nonneg
  linarith

\end{lstlisting}
\begin{proof}
$\alpha(A) > 0$ as the exponential of a real, $\alpha(A) \leq 1$ since the
exponent is $\leq 0$ (because $C_{\mathrm{Har}} \geq 0$), and the lower
bound is the monotonicity of $\exp$ together with $C_{\mathrm{Har}}(d) \leq
C_{\mathrm{Hol}}(d)$ (\leanname{C\_harnack\_le\_C\_holder\_Moser}).
\end{proof}

\begin{lemma}[\leanname{one\_sub\_moserDecayAlpha\_le\_two\_rpow\_neg}]
$1 - \alpha(A) \leq (1/2)^{\alpha(A)}$.
\end{lemma}

\begin{lstlisting}[style=Matlab-editor,basicstyle=\mlttfamily,escapechar=",mlshowsectionrules=true,mathescape=true,morecomment={[s]{\%\{}{\%\}}},language=lean,numbers=none]
theorem one_sub_moserDecayAlpha_le_two_rpow_neg
    (hd : 2 < (d : ℝ))
    (A : NormalizedEllipticCoeff d (Metric.ball (0 : E) 1)) :
    1 - moserDecayAlpha A ≤ (1 / 2 : ℝ) ^ moserDecayAlpha A := by
  let α : ℝ := moserDecayAlpha A
  have hα_pos : 0 < α := by
    simpa [α] using moserDecayAlpha_pos (d := d) A
  have hα_le : α ≤ 1 := by
    simpa [α] using moserDecayAlpha_le_one (d := d) hd A
  have hlog_two_le_one : Real.log (2 : ℝ) ≤ 1 := by
    have := Real.log_le_sub_one_of_pos (by norm_num : 0 < (2 : ℝ))
    norm_num at this ⊢
    linarith
  have hneg_le : -α ≤ -α * Real.log 2 := by
    nlinarith
  have hlin : 1 - α ≤ Real.exp (-α) := by
    simpa [sub_eq_add_neg, add_comm, add_left_comm, add_assoc] using
      (Real.add_one_le_exp (-α))
  have hexp_mono : Real.exp (-α) ≤ Real.exp (-α * Real.log 2) := by
    exact Real.exp_le_exp.mpr hneg_le
  calc
    1 - moserDecayAlpha A = 1 - α := by rfl
    _ ≤ Real.exp (-α) := hlin
    _ ≤ Real.exp (-α * Real.log 2) := hexp_mono
    _ = (1 / 2 : ℝ) ^ α := by
          rw [Real.rpow_def_of_pos (by norm_num : 0 < (1 / 2 : ℝ))]
          congr 1
          have hlog_half : Real.log (1 / 2 : ℝ) = -Real.log 2 := by
            rw [show (1 / 2 : ℝ) = ((2 : ℝ)⁻¹) by norm_num, Real.log_inv]
          rw [hlog_half]
          ring
    _ = (1 / 2 : ℝ) ^ moserDecayAlpha A := by rfl

\end{lstlisting}
\begin{proof}
Set $\alpha = \alpha(A) \in (0,1]$. We have $1 - \alpha \leq e^{-\alpha}$
(by $1+x \leq e^x$ at $x = -\alpha$), and $-\alpha \leq -\alpha\log 2$ since
$\log 2 \leq 1$. Hence $e^{-\alpha} \leq e^{-\alpha\log 2} = (1/2)^\alpha$.
\end{proof}

\begin{theorem}[\leanname{holder\_Moser}]
\label{thm:holder-Moser}
Let $d \geq 3$, $A$ normalized elliptic on $\Bz 1$, $p_0 > 1$, and $u$ a
solution of $-\dv(A\nabla u) = 0$ with $|u|^{p_0} \in \Lp{1}(\Bz 1)$.
Then there exists a representative $v$ of $u$ on $\Bz 1$ (i.e.\ $v = u$
almost everywhere on $\Bz 1$) and an exponent
$\alpha = \alpha(A) > 0$ with
$\exp(-(C_{\mathrm{Hol}}(d)\,\Lambda_A^{1/2})) \leq \alpha$ such that for
every $x, y \in \Bz{1/2}$,
\begin{equation*}
  |v(x) - v(y)|
   \;\leq\; C_{\mathrm{Hol}}(d)\,\Lambda_A^{d/(2p_0)}
        \Bigl(\frac{p_0}{p_0 - 1}\Bigr)^{d/p_0}
        \Bigl(\int_{\Bz 1} |u|^{p_0}\,dz\Bigr)^{1/p_0}
        \,\|x - y\|^{\alpha}.
\end{equation*}
\end{theorem}

\begin{lstlisting}[style=Matlab-editor,basicstyle=\mlttfamily,escapechar=",mlshowsectionrules=true,mathescape=true,morecomment={[s]{\%\{}{\%\}}},language=lean,numbers=none]
theorem holder_Moser
    (hd : 2 < (d : ℝ))
    (A : NormalizedEllipticCoeff d (Metric.ball (0 : E) 1))
    {u : E → ℝ} {p₀ : ℝ} (hp₀ : 1 < p₀)
    (hsol : IsSolution A.1 u)
    (hInt : IntegrableOn (fun z => |u z| ^ p₀) (Metric.ball (0 : E) 1) volume) :
    ∃ v : E → ℝ,
      (∀ᵐ x ∂μ1, v x = u x) ∧
      ∃ α > 0,
        Real.exp (-(C_holder_Moser d * A.1.Λ ^ ((1 : ℝ) / 2))) ≤ α ∧
        ∀ x ∈ Metric.ball (0 : E) (1 / 2 : ℝ),
        ∀ y ∈ Metric.ball (0 : E) (1 / 2 : ℝ),
          |v x - v y| ≤
            C_holder_Moser d * A.1.Λ ^ ((d : ℝ) / (2 * p₀)) *
              (p₀ / (p₀ - 1)) ^ ((d : ℝ) / p₀) *
              (∫ z in Metric.ball (0 : E) 1, |u z| ^ p₀ ∂volume) ^ ((1 : ℝ) / p₀) *
              ‖x - y‖ ^ α := by
  let α : ℝ := moserDecayAlpha A
  have hα_pos : 0 < α := by
    simpa [α] using moserDecayAlpha_pos (d := d) A
  let v : E → ℝ := moserRepresentative u
  have hvu : v =ᵐ[μ1] u :=
    moserRepresentative_ae_eq (d := d) hd A hp₀ hsol hInt
  refine ⟨v, hvu, α, hα_pos, ?_, ?_⟩
  · simpa [α] using moserLowerBound_le_decayAlpha (d := d) hd A
  · intro x hx y hy
    by_cases hxy_eq : x = y
    · subst hxy_eq
      simp [hα_pos.ne', α]
    · by_cases hlarge : (1 / 8 : ℝ) ≤ ‖x - y‖
      · simpa [v, α, moserHolderNorm, mul_assoc, mul_left_comm, mul_comm] using
          (moserRepresentative_large_distance_bound
            (d := d) hd A hp₀ hsol hInt hx hy hxy_eq hlarge)
      · have hsmall : ‖x - y‖ < (1 / 8 : ℝ) := lt_of_not_ge hlarge
        simpa [v, α, moserHolderNorm, mul_assoc, mul_left_comm, mul_comm] using
          (moserRepresentative_small_distance_bound
            (d := d) hd A hp₀ hsol hInt hx hy hxy_eq hsmall)

\end{lstlisting}
\begin{proof}
Define $\alpha := \alpha(A)$ and let $v := \leanname{moserRepresentative}\,u$
be the dyadic-average representative supplied by
\leanname{Holder/Representative.lean}. By
\leanname{moserRepresentative\_ae\_eq}, $v = u$ almost everywhere on $\Bz 1$.

Fix $x, y \in \Bz{1/2}$. If $x = y$ the bound is trivial. Otherwise,
distinguish two cases.

\emph{Large distance.} If $\|x - y\| \geq 1/8$, by
\leanname{moserRepresentative\_large\_distance\_bound} there is a direct
control by the $L^{p_0}$ norm of $u$:
\[
  |v(x) - v(y)| \;\leq\; C_{\mathrm{Hol}}(d)\,\Lambda_A^{d/(2p_0)}
    \bigl(p_0/(p_0-1)\bigr)^{d/p_0}\,\|u\|_{L^{p_0}(\Bz 1)}.
\]
Since $\|x-y\|^\alpha \geq (1/8)^\alpha \geq 1/8$ and absorbing constants
into $C_{\mathrm{Hol}}$ via the factor $256$ in
Definition~\ref{def:CholderMoser} gives the claim.

\emph{Small distance.} If $\|x - y\| < 1/8$, choose $n$ with
$2^{-(n+1)} \leq \|x - y\| < 2^{-n}$ \leanname{(exists\_moserDyadicRadius\_near)},
and apply the dyadic oscillation decay
\leanname{moser\_oscillation\_decay\_on\_ball}, which iterates the Harnack
inequality on each ball $\Ball{2^{-k}}{c}$ centred near $x$:
\[
  \osc_{\Ball{2^{-k-1}}{c}} u
    \;\leq\; (1 - \alpha)\,\osc_{\Ball{2^{-k}}{c}} u.
\]
Iterating gives geometric decay
$\osc_{\Ball{2^{-n}}{c}} u \leq (1-\alpha)^n\,\osc_{\Ball 1 0} u$.
By the previous lemma $1 - \alpha \leq (1/2)^\alpha$, hence
$(1-\alpha)^n \leq 2^{-\alpha n} \leq C\,\|x - y\|^\alpha$.

The starting oscillation $\osc_{\Bz 1} u$ is controlled by twice the local
$L^\infty$ bound, which by the forward Moser estimate
\leanname{linfty\_subsolution\_Moser\_on\_ball} is bounded by
$M_{\mathrm{loc}}(d)\Lambda_A^{d/(2p_0)} (p_0/(p_0-1))^{d/p_0}
        \|u\|_{L^{p_0}(\Bz 1)}$.
This yields the desired pointwise Holder bound on the representative $v$.
\end{proof}

\begin{corollary}[\leanname{holder\_Moser\_of\_homogeneousWeakSolution}]
The same conclusion holds verbatim for homogeneous weak solutions
($\IsHomogeneousWeakSolution$), since these satisfy
$\IsSolution$.
\end{corollary}

\begin{lstlisting}[style=Matlab-editor,basicstyle=\mlttfamily,escapechar=",mlshowsectionrules=true,mathescape=true,morecomment={[s]{\%\{}{\%\}}},language=lean,numbers=none]
/-- Moser interior Hölder regularity for homogeneous local weak solutions. -/
theorem holder_Moser_of_homogeneousWeakSolution
    (hd : 2 < (d : ℝ))
    (A : NormalizedEllipticCoeff d (Metric.ball (0 : E) 1))
    {u : E → ℝ} {p₀ : ℝ} (hp₀ : 1 < p₀)
    (hweak : IsHomogeneousWeakSolution A.1 u)
    (hInt : IntegrableOn (fun z => |u z| ^ p₀) (Metric.ball (0 : E) 1) volume) :
    ∃ v : E → ℝ,
      (∀ᵐ x ∂μ1, v x = u x) ∧
      ∃ α > 0,
        Real.exp (-(C_holder_Moser d * A.1.Λ ^ ((1 : ℝ) / 2))) ≤ α ∧
        ∀ x ∈ Metric.ball (0 : E) (1 / 2 : ℝ),
        ∀ y ∈ Metric.ball (0 : E) (1 / 2 : ℝ),
          |v x - v y| ≤
            C_holder_Moser d * A.1.Λ ^ ((d : ℝ) / (2 * p₀)) *
              (p₀ / (p₀ - 1)) ^ ((d : ℝ) / p₀) *
              (∫ z in Metric.ball (0 : E) 1, |u z| ^ p₀ ∂volume) ^ ((1 : ℝ) / p₀) *
              ‖x - y‖ ^ α := by
  exact holder_Moser hd A hp₀ (isHomogeneousWeakSolution_isSolution hweak) hInt

\end{lstlisting}
\begin{remark}
The dependence on the ellipticity bound $\Lambda_A$ enters Holder regularity
in two ways: through the multiplicative prefactor $\Lambda_A^{d/(2p_0)}$
inherited from the forward Moser bound, and through the Holder exponent
$\alpha(A) = \exp(-C_{\mathrm{Har}}(d)\,\Lambda_A^{1/2})$ inherited from the
Harnack constant. As $\Lambda_A \to \infty$, $\alpha(A) \to 0$ exponentially,
quantifying the well-known degeneration of De Giorgi--Nash--Moser regularity
for highly anisotropic elliptic operators.
\end{remark}
\end{document}
%%% Local Variables:
%%% mode: LaTeX
%%% TeX-engine: xetex
%%% TeX-command-extra-options: "-shell-escape"
%%% End:
x
